% Preamble for the swarm-transcribed Weber volumes.
% Deliberately close to weber_v1_ge.tex so the verified p1-148 region and the swarm
% region can live in one document.
\documentclass[11pt]{article}
\usepackage[ngerman]{babel}
\usepackage[a4paper,margin=0.9in]{geometry}
\usepackage[T1]{fontenc}
\usepackage[utf8]{inputenc}
\usepackage{lmodern}
\usepackage[leqno]{amsmath}
\usepackage{amssymb,mathtools,mathrsfs}
\usepackage{microtype}
\usepackage{fancyhdr}
\usepackage{array,longtable,enumitem,multicol}
\usepackage{tikz}
\usepackage{graphicx}
\graphicspath{{../figures/b1/}{../figures/b2/}{../figures/b3/}}
\setlength{\parindent}{1.5em}
\setlength{\parskip}{0.18em}
\emergencystretch=6em
\setlength{\headheight}{14pt}
\pagestyle{fancy}
\fancyhf{}
\cfoot{\thepage}

% --- Weber conventions -------------------------------------------------------
% Sperrung (letter-spacing) is Weber's emphasis. There is no bold in his body text.
\providecommand{\gesperrt}[1]{\textls[130]{#1}}
% Printed-page marker: a no-op in the PDF. The machine-readable anchor is the
% "% !PRINTED n" comment in the source, which the collector and Codex both read.
\providecommand{\weberpage}[1]{}
% Displayed equation numbered in the LEFT margin, as Weber prints them.
\providecommand{\weq}[2]{\begin{equation}\tag{#1}#2\end{equation}}
% Figure placeholder. A later pass replaces it with TikZ or a cleaned scan crop.
% A marked gap is honest; a fabricated diagram is not.
\providecommand{\weberfigure}[2]{%
  \begin{center}\fbox{\parbox{0.78\linewidth}{\centering\small [Figur, S.~#1: #2]}}\end{center}}
% Source-faithful figure crop.  The asset search path works both from the live
% swarm_build directory and from a self-contained release bundle with assets/.
\providecommand{\weberfigureimage}[2]{%
  \begin{center}\includegraphics[width=#2\linewidth]{#1}\end{center}}

% Text symbols agents reach for that are not always defined. Providing them in the preamble
% fixes every occurrence at once and leaves the generated fragments untouched, which is the
% right trade whenever the fix is "define the macro" rather than "change the text".
\usepackage{textcomp}
% amsmath caps matrix environments at 10 columns. Weber prints permutation matrices far wider
% than that -- Bd III has an 11-column pmatrix of the w_nu -- and the failure is the opaque
% "Extra alignment tab has been changed to \cr".
\setcounter{MaxMatrixCols}{30}
\providecommand{\textquotedblbase}{\glqq}
\providecommand{\textquotedblleft}{``}
\providecommand{\textquotedblright}{''}
\providecommand{\textquotesinglbase}{\glq}

% --- Greek capitals are SLANTED in this fount -------------------------------------------
% RULING 2026-08-07, after the same observation was reported from five separate pages
% (printed 271, 275, 285, 333 and the nu/v cousin on 172).
%
% Verified at 8x zoom on printed 333 (b1 idx 358): every capital Delta on that page --
% "f(x) - f(alpha) = Delta", "Delta_alpha = f(alpha+delta) - f(alpha)" -- is the INCLINED
% form, which LaTeX calls \varDelta. \Delta, LaTeX's default, is upright and is not what
% Vieweg set. The same inclination is visible on capital Phi (printed 171, 282).
%
% This is a property of the TYPEFACE, not of any occurrence -- the same kind of fact as the
% varkappa ruling -- so it holds wherever the letter appears, and the right fix is one line
% here rather than several thousand edits across the fragments. \let leaves every fragment
% untouched, applies to the whole corpus at once, and is reversible by deleting these lines.
%
% Delta and Phi are the two verified by eye. The rest follow from its being a single fount;
% if any capital is later found upright on the page, delete its line and the fragments still
% say exactly what they said before.
\let\Gamma\varGamma
\let\Delta\varDelta
\let\Theta\varTheta
\let\Lambda\varLambda
\let\Xi\varXi
\let\uprightPi\Pi
\let\Pi\varPi
\let\Sigma\varSigma
\let\Upsilon\varUpsilon
\let\Phi\varPhi
\let\Psi\varPsi
\let\Omega\varOmega
% Lower-case rho is the OPEN-TAILED form in this fount. Verified at 10x on printed 419
% ("wenn also überhaupt ein \varrho vorhanden ist"): the tail curves down and left rather than
% descending straight, which is \varrho, not \rho. Same kind of fact as the capitals above, and
% the same one-line treatment. (kappa is already handled: the fragments write \varkappa
% explicitly, on the ruling recorded in resolve_class.py.)
\let\rho\varrho

% --- Fractions written in running text -------------------------------------------------
% Gauss's German prose says things like "von 6\tfrac{1}{2} Jahren" and "\tfrac{1}{2} Grad":
% a fraction as an ordinary word, not inside math. Bare \tfrac there is
%     ! Missing $ inserted.
% and the error points at the sentence, not at the macro.
%
% Wrapping the three fraction macros in \ensuremath fixes every occurrence at once and leaves
% the fragments untouched -- the right trade whenever the fix is "define the macro" rather than
% "change the text", as with the Greek capitals above. Inside math \ensuremath is a no-op, so
% the thousands of legitimate in-math uses are unaffected.
\let\swarmfrac\frac
\let\swarmtfrac\tfrac
\let\swarmdfrac\dfrac
\renewcommand{\frac}[2]{\ensuremath{\swarmfrac{#1}{#2}}}
\renewcommand{\tfrac}[2]{\ensuremath{\swarmtfrac{#1}{#2}}}
\renewcommand{\dfrac}[2]{\ensuremath{\swarmdfrac{#1}{#2}}}

% --- Unicode vulgar fractions -----------------------------------------------------------
% Agents sometimes type the single Unicode glyph the page shows -- "erstere ⅔" -- rather than a
% macro. utf8 inputenc has no default binding for U+2153..U+215E and stops the build with
%     ! LaTeX Error: Unicode character ⅔ (U+2154)
% Declaring them costs nothing, keeps the fragment reading exactly as the agent saw the page,
% and is the same "define it in the preamble" trade as the fraction macros above. Typing the
% glyph is not wrong; it is arguably closer to the page than \tfrac would be.
\DeclareUnicodeCharacter{2153}{\ensuremath{\swarmtfrac{1}{3}}}
\DeclareUnicodeCharacter{2154}{\ensuremath{\swarmtfrac{2}{3}}}
\DeclareUnicodeCharacter{2155}{\ensuremath{\swarmtfrac{1}{5}}}
\DeclareUnicodeCharacter{2156}{\ensuremath{\swarmtfrac{2}{5}}}
\DeclareUnicodeCharacter{2157}{\ensuremath{\swarmtfrac{3}{5}}}
\DeclareUnicodeCharacter{2158}{\ensuremath{\swarmtfrac{4}{5}}}
\DeclareUnicodeCharacter{2159}{\ensuremath{\swarmtfrac{1}{6}}}
\DeclareUnicodeCharacter{215A}{\ensuremath{\swarmtfrac{5}{6}}}
\DeclareUnicodeCharacter{215B}{\ensuremath{\swarmtfrac{1}{8}}}
\DeclareUnicodeCharacter{215C}{\ensuremath{\swarmtfrac{3}{8}}}
\DeclareUnicodeCharacter{215D}{\ensuremath{\swarmtfrac{5}{8}}}
\DeclareUnicodeCharacter{215E}{\ensuremath{\swarmtfrac{7}{8}}}
% Prime and double prime: arcminutes and arcseconds. Gauss's geodesy and astronomy volumes are
% full of "43° 37' 6\"741", and agents type the true Unicode prime characters rather than the
% ASCII apostrophe -- which is right, since that is what the page prints.
\DeclareUnicodeCharacter{2032}{\ensuremath{{}^{\prime}}}
\DeclareUnicodeCharacter{2033}{\ensuremath{{}^{\prime\prime}}}
\DeclareUnicodeCharacter{2034}{\ensuremath{{}^{\prime\prime\prime}}}
\DeclareUnicodeCharacter{00BD}{\ensuremath{\swarmtfrac{1}{2}}}
\DeclareUnicodeCharacter{00BC}{\ensuremath{\swarmtfrac{1}{4}}}
\DeclareUnicodeCharacter{00BE}{\ensuremath{\swarmtfrac{3}{4}}}

% Macros agents invent for a text fraction. \textfrac is not a real LaTeX command, but the
% intent is unambiguous and providing it is cheaper and safer than rewriting the fragment --
% the alternative is editing a table cell by hand and risking its alignment.
\providecommand{\textfrac}[2]{\ensuremath{\swarmtfrac{#1}{#2}}}

% Common shorthands inherited from the volume I file, so fragments that use them compile.
\providecommand{\eps}{\varepsilon}
\providecommand{\ii}{\mathrm{i}}
\providecommand{\tg}{\operatorname{tg}}
% Gauss writes the trigonometric functions in the period's Latin abbreviations -- "tang" for
% tangent, "cotang" for cotangent -- rather than Weber's "tg"/"cotg". Same function, different
% century's shorthand; both are reproduced as the author sets them.
\providecommand{\tang}{\operatorname{tang}}
\providecommand{\cotang}{\operatorname{cotang}}
\providecommand{\arctang}{\operatorname{arc\,tang}}
\providecommand{\ctg}{\operatorname{cotg}}
\providecommand{\cotg}{\operatorname{cotg}}
\providecommand{\sn}{\operatorname{sn}}
\providecommand{\cn}{\operatorname{cn}}
\providecommand{\dn}{\operatorname{dn}}
\providecommand{\am}{\operatorname{am}}
\providecommand{\dd}{\,d}
\providecommand{\tht}{\vartheta}
\providecommand{\vark}{\varkappa}
\providecommand{\Norm}{\operatorname{N}}
\providecommand{\ord}{\operatorname{ord}}
\providecommand{\ind}{\operatorname{ind}}
\providecommand{\abs}[1]{\left|#1\right|}
\providecommand{\sect}[2]{\section*{\S.\,#1. #2}}
\providecommand{\frp}{\mathfrak p}
\providecommand{\frP}{\mathfrak P}
\providecommand{\frG}{\mathfrak G}
\providecommand{\frA}{\mathfrak A}
\providecommand{\frB}{\mathfrak B}

% --- reader additions ---------------------------------------------------------
\usepackage{url}
\usepackage[unicode,bookmarksnumbered=false,bookmarksopen=true,bookmarksopenlevel=1,colorlinks=true,
  linkcolor=black,urlcolor=blue]{hyperref}
\usepackage{bookmark}
\renewcommand{\headrulewidth}{0pt}
\renewcommand{\weberpage}[1]{\phantomsection\label{wp:#1}}
\newcommand{\inhaltkopf}[1]{\begin{center}\large\textbf{#1}\end{center}\medskip}
\newcommand{\inhaltbuch}[2]{\par\medskip\begin{center}\if\relax\detokenize{#1}\relax\else#1\\[2pt]\fi{\large\textbf{#2}}\end{center}}
\newcommand{\inhaltabschnitt}[2]{\par\smallskip\begin{center}\small\if\relax\detokenize{#1}\relax\else#1\\\fi\textbf{#2}\end{center}}
\newcommand{\inhaltpar}[3]{\par\noindent\hangindent=3.4em\hbox to 3.4em{\S~#1.\hfil}#2\dotfill\,\hyperref[wp:#3]{#3}\par}
\newcommand{\inhaltzeile}[2]{\par\noindent #1\dotfill\,\hyperref[wp:#2]{#2}\par}

\begin{document}
\selectlanguage{ngerman}
\pagenumbering{roman}
\begin{center}
{\Large Heinrich Weber}\\[6pt]
{\LARGE\textbf{Lehrbuch der Algebra}}\\[8pt]
{\large Zweiter Band}\\[4pt]
Braunschweig: Friedrich Vieweg und Sohn, 1896\\[24pt]
{\large Deutscher Text}\\[4pt]
mit den Seitenzahlen des Originals
\end{center}
\clearpage
\noindent Transkribiert vom Scan \url{https://archive.org/details/lehrbuchderalgeb02weberich} (Internet Archive).
Der deutsche Text wurde von KI-Modellen (Anthropic Claude und OpenAI Codex) vom Scan transkribiert und am Scan
geprüft. Webers Text ist gemeinfrei; Transkription und Satz sind nach CC0~1.0 gemeinfrei gestellt
(\url{https://creativecommons.org/publicdomain/zero/1.0/}). Quelldateien: \url{https://github.com/KokunoYumeto/weber-lehrbuch-der-algebra}.
\clearpage

% ===== front matter, scan leaf 6 =====
\begin{center}
\vspace*{5cm}
{\Large LEHRBUCH}\\[1cm]
{\small DER}\\[1.2cm]
{\Huge \gesperrt{ALGEBRA.}}\\[1.5cm]
\rule{3cm}{0.4pt}
\end{center}

\clearpage
% ===== front matter, scan leaf 8 =====
\begin{center}
\vspace*{2cm}
{\LARGE LEHRBUCH}\\[1.2cm]
{\small DER}\\[1.2cm]
{\Huge \gesperrt{ALGEBRA}.}\\[2cm]
{\large \textsc{Von}}\\[1.2cm]
{\LARGE HEINRICH WEBER,}\\[0.2cm]
{\footnotesize PROFESSOR DER MATHEMATIK AN DER UNIVERSITÄT STRASSBURG.}\\[1.2cm]
\rule{2.5cm}{0.4pt}\\[1.2cm]
{\large \gesperrt{IN ZWEI BÄNDEN}.}\\[0.8cm]
\rule{2cm}{0.4pt}\\[0.6cm]
{\large \textbf{ZWEITER BAND.}}\\[3.5cm]
\rule{\textwidth}{1pt}\\[0.8cm]
{\Large BRAUNSCHWEIG,}\\[0.3cm]
{\small DRUCK UND VERLAG VON FRIEDRICH VIEWEG UND SOHN.}\\[0.3cm]
{\Large 1896.}
\end{center}

\clearpage
% ===== front matter, scan leaf 9 =====
% not printed: Bibliotheksstempel "MATH-STAT." mit handschriftlichem Kürzel; durchscheinende Titelseite der Rückseite
\begin{center}
\vspace*{0.35\textheight}
\rule{0.5\textwidth}{0.4pt}\\[1ex]
\gesperrt{Alle Rechte vorbehalten.}\\[0.5ex]
\rule{0.5\textwidth}{0.4pt}
\end{center}

\clearpage
% ===== front matter, scan leaf 10 =====
% not printed: handwritten call number QA153 W35 v.2; library stamps MATH.-STAT. LIBRARY and M777594
\begin{center}
{\large VORWORT ZUM ZWEITEN BANDE.}

\vspace{8pt}
\rule{2.5cm}{0.4pt}
\end{center}

\vspace{6pt}

Der in dem Vorworte zum ersten Bande angekündigten Absicht gemäss kann ich heute den zweiten Band meines Lehrbuches der Algebra der Oeffentlichkeit übergeben. Der dort aufgestellte Plan ist in den wesentlichen Punkten durchgeführt. Bei den Anwendungen bin ich bemüht gewesen, solche Probleme auszuwählen, die bereits in anderen Gebieten, der Geometrie oder Functionentheorie, ein selbständiges Interesse gewonnen haben, und die zugleich die Hauptpunkte der algebraischen Theorie möglichst vielseitig zur Anschauung bringen.

Die Anwendung der Theorie der algebraischen Zahlen ist bis zur Theorie der Kreistheilungszahlen durchgeführt. Wenn Leben und Arbeitskraft vorhalten, hoffe ich, in einer Fortsetzung meines Werkes die weiteren Anwendungen auf das Gebiet der elliptischen Functionen darzustellen, die nur zum Theil in meinem Buche „Elliptische Functionen und algebraische Zahlen“ enthalten sind.

Auch während der Ausarbeitung und des Druckes des zweiten Bandes hat mir die Hülfe und der Rath der Freunde zur Seite gestanden, die ich schon in der Vorrede zur ersten Auflage genannt habe. Aber auch manchen neuen Freund hat sich der erste Band bereits erworben, der meine Arbeit durch Winke und Rathschläge gefördert hat. Ihnen allen spreche ich an dieser Stelle meinen Dank aus, und füge die Bitte hinzu, dass sie dem Werke auch weiterhin ihr Interesse bewahren mögen.

\bigskip

\gesperrt{Strassburg}, im Juli 1896.

\begin{flushright}
\textbf{\gesperrt{Der Verfasser}.}\hspace*{2em}
\end{flushright}

\clearpage
% ===== front matter, scan leaf 12 =====
\inhaltkopf{INHALT DES ZWEITEN BANDES.}

\inhaltbuch{Erstes Buch.}{Gruppen.}

\inhaltabschnitt{Erster Abschnitt.}{Allgemeine Gruppentheorie.}
\inhaltpar{1}{Definition der Gruppen}{3}
\inhaltpar{2}{Die Divisoren endlicher Gruppen}{7}
\inhaltpar{3}{Normaltheiler einer Gruppe}{10}
\inhaltpar{4}{Composition der Theile}{12}
\inhaltpar{5}{Mehrstufiger Isomorphismus}{15}
\inhaltpar{6}{Die Compositionsreihe und der Satz von C. \gesperrt{Jordan}}{17}
\inhaltpar{7}{Weitere Sätze über die Compositionsreihen}{24}
\inhaltpar{8}{Metacyklische Gruppen}{27}

\inhaltabschnitt{Zweiter Abschnitt.}{Abel'sche Gruppen.}
\inhaltpar{9}{Darstellung \gesperrt{Abel}'scher Gruppen durch eine Basis}{32}
\inhaltpar{10}{Die Invarianten der \gesperrt{Abel}'schen Gruppen}{39}
\inhaltpar{11}{Gruppencharaktere}{43}
\inhaltpar{12}{Divisoren einer \gesperrt{Abel}'schen Gruppe. Reciproke Gruppen}{48}
\inhaltpar{13}{Die Geschlechter in einer \gesperrt{Abel}'schen Gruppe}{52}
\inhaltpar{14}{Indices nach einer ungeraden Primzahlpotenz als Modul}{54}
\inhaltpar{15}{Indices für eine Potenz von 2 als Modul}{58}
\inhaltpar{16}{Die Gruppen der Zahlclassen nach einem zusammengesetzten Modul}{60}

\inhaltabschnitt{Dritter Abschnitt.}{Die Gruppe der Kreistheilungskörper.}
\inhaltpar{17}{Die Resolventen der Kreistheilungstheorie}{63}
\inhaltpar{18}{Kreistheilungskörper}{67}
\inhaltpar{19}{Primäre und nicht primäre Theiler der Gruppe $\mathfrak{N}$}{72}
\inhaltpar{20}{Die Kreistheilungsperioden}{74}
\inhaltpar{21}{Kreistheilungskörper mit gegebener Gruppe}{79}
\inhaltpar{22}{Bestimmung der Gruppe $\mathfrak{A}$}{92}

\clearpage
% ===== front matter, scan leaf 13 =====
% !PRINTED VIII
\weberpage{VIII}
\inhaltabschnitt{Vierter Abschnitt.}{Cubische und biquadratische Abel'sche Körper.}
\inhaltpar{23}{Cubische Kreistheilungskörper}{94}
\inhaltpar{24}{Biquadratische Kreistheilungskörper}{101}
\inhaltpar{25}{Cubische \gesperrt{Abel}'sche Gleichungen}{107}
\inhaltpar{26}{Biquadratische \gesperrt{Abel}'sche Gleichungen}{110}
\inhaltabschnitt{Fünfter Abschnitt.}{Constitution der allgemeinen Gruppen.}
\inhaltpar{27}{Bildung von Gruppen nach \gesperrt{Cayley}}{114}
\inhaltpar{28}{Beziehung der allgemeinen Gruppen zu den Permutationsgruppen}{117}
\inhaltpar{29}{Der erste \gesperrt{Sylow}'sche Satz}{121}
\inhaltpar{30}{Der zweite \gesperrt{Sylow}'sche Satz}{125}
\inhaltpar{31}{Gruppen vom Grade $p^\alpha$}{127}
\inhaltpar{32}{Satz von \gesperrt{Frobenius}}{129}
\inhaltpar{33}{Gruppen vom Grade $p^\alpha q$}{134}
\inhaltpar{34}{Einfache Gruppen}{136}
\inhaltpar{35}{Gruppen vom Grade $p q$}{140}
\inhaltpar{36}{Grenzen des Index eines Theilers der symmetrischen Permutationsgruppe}{143}
\inhaltbuch{Zweites Buch.}{Lineare Gruppen.}
\inhaltabschnitt{Sechster Abschnitt.}{Gruppen linearer Substitutionen.}
\inhaltpar{37}{Lineare Substitutionen und ihre Zusammensetzung}{151}
\inhaltpar{38}{Substitution der Verhältnisse}{158}
\inhaltpar{39}{Permutationen als lineare Substitutionen}{159}
\inhaltpar{40}{Die Invarianten von endlichen Gruppen linearer Substitutionen}{161}
\inhaltpar{41}{Der Satz von \gesperrt{Hilbert}}{165}
\inhaltpar{42}{Endlichkeit des Invariantensystems einer endlichen linearen Substitutionsgruppe}{168}
\inhaltpar{43}{Das Formenproblem}{171}
\inhaltpar{44}{\gesperrt{Klein}'s Erweiterung des algebraischen Grundproblems}{176}
\inhaltpar{45}{Einfluss relativer Invarianten}{177}
\inhaltpar{46}{Der erweiterte Invariantenbegriff}{179}
\inhaltpar{47}{Normalformen}{181}
\inhaltabschnitt{Siebenter Abschnitt.}{Gruppen binärer linearer Substitutionen.}
\inhaltpar{48}{Ternäre orthogonale Substitutionen}{184}
\inhaltpar{49}{Lineare gebrochene Substitutionen}{189}

\clearpage
% ===== front matter, scan leaf 14 =====
% !PRINTED IX
\weberpage{IX}
\inhaltpar{50}{Realitätsbedingungen}{193}
\inhaltpar{51}{Endliche Gruppen linearer gebrochener Substitutionen. Pole der Gruppen}{195}
\inhaltpar{52}{Die verschiedenen Arten möglicher Gruppen}{199}
\inhaltpar{53}{Transformation der Substitutionen von $G$ auf einfache Formen}{203}
\inhaltpar{54}{Die Grundformen}{205}

\inhaltabschnitt{Achter Abschnitt.}{Die Polyëdergruppen.}
\inhaltpar{55}{Die cyklischen Gruppen und die Diëdergruppen}{209}
\inhaltpar{56}{Die Tetraëdergruppe}{212}
\inhaltpar{57}{Die Octaëdergruppe}{216}
\inhaltpar{58}{Die Ikosaëdergruppe}{220}
\inhaltpar{59}{Die Theiler der Ikosaëdergruppe}{227}
\inhaltpar{60}{Die Grundformen der Ikosaëdergruppe}{230}
\inhaltpar{61}{Die Invarianten des Ikosaëders}{232}
\inhaltpar{62}{Polyëdergruppen der zweiten Art. Krystallographische Gruppen}{234}

\inhaltabschnitt{Neunter Abschnitt.}{Congruenzgruppen.}
\inhaltpar{63}{Functionen-Congruenzen}{242}
\inhaltpar{64}{Congruenzkörper}{245}
\inhaltpar{65}{Congruenzgruppen im Körper $\mathfrak{E}$}{250}
\inhaltpar{66}{Einfachheit der Gruppe $E$}{254}
\inhaltpar{67}{Congruenzkörper zweiten Grades}{259}
\inhaltpar{68}{Die reelle lineare Congruenzgruppe $L_p$}{261}
\inhaltpar{69}{Imaginäre Form der Gruppe $L_p$}{266}
\inhaltpar{70}{Divisoren der Gruppe $L_p$, deren Grad durch $p$ theilbar ist}{271}
\inhaltpar{71}{Divisoren der Gruppe $L_p$, deren Grad nicht durch $p$ theilbar ist}{273}
\inhaltpar{72}{Constitution der Gruppe $L_7$ vom Grade 168}{282}

\inhaltbuch{Drittes Buch.}{Anwendungen der Gruppentheorie.}

\inhaltabschnitt{Zehnter Abschnitt.}{Allgemeine Theorie der metacyklischen Gleichungen.}
\inhaltpar{73}{Die Resolventen der Compositionsreihen}{289}
\inhaltpar{74}{Metacyklische Gleichungen}{292}
\inhaltpar{75}{Metacyklische Gleichungen, deren Grad eine Primzahlpotenz ist}{297}
\inhaltpar{76}{Darstellung der \gesperrt{Abel}'schen Gruppe $Q$}{298}
\inhaltpar{77}{Analytische Darstellung der Permutationen}{299}
\inhaltpar{78}{Darstellung der metacyklischen Gruppe $P$}{301}
\inhaltpar{79}{Ternäre lineare Congruenzgruppe für den Modul 2}{306}
\inhaltpar{80}{Resolventen der Gleichung 8\textsuperscript{ten} Grades}{308}

\clearpage
% ===== front matter, scan leaf 15 =====
% !PRINTED X
\weberpage{X}
\inhaltpar{81}{Die Tripelsysteme der Resolventen}{310}
\inhaltpar{82}{Anwendung auf Gleichungen 8\textsuperscript{ten} Grades}{313}
\inhaltpar{83}{Metacyklische Gleichungen 8\textsuperscript{ten} Grades}{314}
\inhaltpar{84}{Biquadratische Gleichungen}{319}
\inhaltabschnitt{Elfter Abschnitt.}{Die Wendepunkte einer Curve dritter Ordnung.}
\inhaltpar{85}{Ternäre Formen und algebraische Curven}{322}
\inhaltpar{86}{Singuläre Punkte. Wendepunkte. Doppeltangenten}{324}
\inhaltpar{87}{Fundamentale Covarianten einer ternären Form}{327}
\inhaltpar{88}{Die \gesperrt{Hesse}'sche Curve}{329}
\inhaltpar{89}{Inflexionspunkte einer Curve dritter Ordnung}{331}
\inhaltpar{90}{Transformation der cubischen Form auf die canonische Form}{333}
\inhaltpar{91}{Die Invarianten der Curve dritter Ordnung und die biquadratische Gleichung}{336}
\inhaltpar{92}{Tripelgleichungen}{342}
\inhaltpar{93}{Die Gruppe der Tripelgleichungen}{344}
\inhaltpar{94}{Realitätsverhältnisse der Tripelgleichungen}{348}
\inhaltabschnitt{Zwölfter Abschnitt.}{Die Doppeltangenten einer Curve vierter Ordnung.}
\inhaltpar{95}{Anzahl der Doppeltangenten einer Curve vierter Ordnung}{351}
\inhaltpar{96}{Die \gesperrt{Steiner}'schen Complexe}{357}
\inhaltpar{97}{Complexpaare und Complextripel}{363}
\inhaltpar{98}{Die \gesperrt{Aronhold}'schen Siebener-Systeme}{367}
\inhaltpar{99}{Der \gesperrt{Hesse-Cayley}'sche Algorithmus zur Bezeichnung der Doppeltangenten}{369}
\inhaltpar{100}{Rationale Bestimmung der Curve aus einem vollständigen Siebener-Systeme}{374}
\inhaltpar{101}{Die \gesperrt{Galois}'sche Gruppe des Doppeltangentenproblems}{380}
\inhaltpar{102}{Darstellung der Gruppe}{384}
\inhaltpar{103}{Einfachheit der Gruppe des Doppeltangentenproblems}{387}
\inhaltpar{104}{Realität der Doppeltangenten}{391}
\inhaltpar{105}{Beweis der Existenz der vier Fälle}{398}
\inhaltabschnitt{Dreizehnter Abschnitt.}{Allgemeine Theorie der Gleichung fünften Grades.}
\inhaltpar{106}{Fragestellung}{403}
\inhaltpar{107}{Satz von \gesperrt{Lüroth}}{404}
\inhaltpar{108}{Resolventen mit einem Parameter}{407}
\inhaltpar{109}{Gruppe der Resolventen mit einem Parameter}{410}
\inhaltpar{110}{Die Ikosaëdergleichung}{415}
\inhaltpar{111}{Die Resolventen der Ikosaëdergleichung}{419}
\inhaltpar{112}{Die Hauptresolvente fünften Grades}{422}

\clearpage
% ===== front matter, scan leaf 16 =====
% !PRINTED XI
\weberpage{XI}
\inhaltpar{113}{Resolventen sechsten Grades}{426}
\inhaltpar{114}{Ueber die Lösung der Ikosaëdergleichung durch transcendente Functionen}{429}
\inhaltabschnitt{Vierzehnter Abschnitt.}{Gruppen linearer ternärer Substitutionen.}
\inhaltpar{115}{Ternäre lineare Substitutionsgruppe vom 168\textsuperscript{sten} Grade}{433}
\inhaltpar{116}{Pole und Axen der ternären Gruppen}{438}
\inhaltpar{117}{Anwendung auf die Gruppe $G_{168}$. Siebenzählige Pole}{442}
\inhaltpar{118}{Die Hauptaxen}{444}
\inhaltpar{119}{Die drei- und sechszähligen Pole}{448}
\inhaltpar{120}{Die Configuration der Gruppe $G_{168}$}{451}
\inhaltpar{121}{Invariantencurven der Gruppe $G_{168}$}{453}
\inhaltpar{122}{Die erste Invariante der Gruppe $G_{168}$ und die Grundcurve}{454}
\inhaltpar{123}{Die höheren Invarianten}{459}
\inhaltpar{124}{Das volle Invariantensystem}{461}
\inhaltabschnitt{Fünfzehnter Abschnitt.}{Das Formenproblem der Gruppe $G_{168}$ und die Theorie der Gleichungen siebenten Grades.}
\inhaltpar{125}{Die Resolventen des Formenproblems}{466}
\inhaltpar{126}{Reduction der allgemeinen Resolvente siebenten Grades auf die specielle}{471}
\inhaltpar{127}{Permutationsgruppe von sieben Ziffern vom Grade 168}{473}
\inhaltpar{128}{Gleichungen siebenten Grades mit einer Gruppe 168\textsuperscript{sten} Grades}{476}
\inhaltpar{129}{Contragrediente Gruppen}{478}
\inhaltpar{130}{Lösung der Gleichung siebenten Grades mit der Gruppe $P_{168}$ durch das Formenproblem der Gruppe $G_{168}$}{481}
\inhaltpar{131}{Möglichkeit der Bestimmung der Functionen $X_1, X_2, X_3$}{484}
\inhaltbuch{Viertes Buch.}{Algebraische Zahlen.}
\inhaltabschnitt{Sechzehnter Abschnitt.}{Zahlen und Functionale eines algebraischen Körpers.}
\inhaltpar{132}{Definition der algebraischen Zahlen}{489}
\inhaltpar{133}{Ganze algebraische Zahlen}{490}
\inhaltpar{134}{Algebraische Körper}{493}
\inhaltpar{135}{Ganze Functionen in einem algebraischen Körper}{495}
\inhaltpar{136}{Die Functionale eines algebraischen Körpers $\Omega$ und der erweiterte Körper $\overline{\Omega}$}{499}
\inhaltpar{137}{Ganze Functionale}{504}
\inhaltpar{138}{Theilbarkeit. Associirte Functionale. Einheiten}{509}
\inhaltpar{139}{Grösster gemeinschaftlicher Theiler}{512}

\clearpage
% ===== front matter, scan leaf 17 =====
% !PRINTED XII
\weberpage{XII}

\inhaltpar{140}{Primfunctionale im Körper $\Omega$}{515}
\inhaltpar{141}{Zerlegung der ganzen und gebrochenen Functionale in Primfactoren}{516}
\inhaltpar{142}{Ganze Functionen im Körper $\overline{\Omega}$}{520}
\inhaltpar{143}{Die Primfactoren der Zahlen des Körpers $\Omega$}{523}

\inhaltabschnitt{Siebzehnter Abschnitt.}{Theorie der algebraischen Körper.}

\inhaltpar{144}{Basis eines algebraischen Zahlkörpers. Discriminanten}{527}
\inhaltpar{145}{Die Minimalbasis und die Körperdiscriminante}{529}
\inhaltpar{146}{Die Basen der Functionale}{532}
\inhaltpar{147}{Die absoluten Normen der Functionale}{536}
\inhaltpar{148}{Volles Restsystem nach einem Modul}{539}
\inhaltpar{149}{Congruenzen}{541}
\inhaltpar{150}{Der \gesperrt{Fermat}'sche Satz}{544}
\inhaltpar{151}{Die \gesperrt{Dedekind}'schen Ideale}{547}
\inhaltpar{152}{Aequivalenz}{552}
\inhaltpar{153}{Die Classenzahl des Körpers $\Omega$}{554}
\inhaltpar{154}{Die Gruppe der Idealclassen}{557}

\inhaltabschnitt{Achtzehnter Abschnitt.}{Discriminanten.}

\inhaltpar{155}{Minimum einer quadratischen Form}{559}
\inhaltpar{156}{Anwendung auf algebraische Körper}{566}
\inhaltpar{157}{Primfactoren der natürlichen Primzahlen}{569}
\inhaltpar{158}{\gesperrt{Dedekind}'s Satz über die Körperdiscriminante}{578}

\inhaltabschnitt{Neunzehnter Abschnitt.}{Beziehung eines Körpers auf seine Theiler.}

\inhaltpar{159}{Relativnormen}{583}
\inhaltpar{160}{Primideale im relativ normalen Körper}{588}
\inhaltpar{161}{Primitivwurzeln der Primideale}{591}
\inhaltpar{162}{Das Partial-Grundideal}{593}
\inhaltpar{163}{Die Ideale in den Theilern des Körpers $\Omega$}{597}

\inhaltabschnitt{Zwanzigster Abschnitt.}{Quadratische Körper.}

\inhaltpar{164}{Basis eines quadratischen Körpers}{601}
\inhaltpar{165}{Die Primideale in $\Omega$}{603}
\inhaltpar{166}{Functionale im quadratischen Körper und quadratische Irrationalzahlen}{607}
\inhaltpar{167}{Aequivalente Functionale und äquivalente Zahlen}{610}
\inhaltpar{168}{Composition der quadratischen Irrationalzahlen}{613}

\clearpage
% ===== front matter, scan leaf 18 =====
% !PRINTED XIII
\weberpage{XIII}

\inhaltabschnitt{Einundzwanzigster Abschnitt.}{Kreistheilungskörper.}
\inhaltpar{169}{Zerlegung der Primzahl $q$ im Kreistheilungskörper $\Omega_{q^\varkappa}$}{617}
\inhaltpar{170}{Minimalbasis des Körpers $\Omega_m$}{620}
\inhaltpar{171}{Die Primideale im Körper $\Omega_m$}{622}
\inhaltpar{172}{Die conjugirten Primideale}{623}
\inhaltpar{173}{Darstellung der Primfactoren von $p$}{628}
\inhaltpar{174}{Das \gesperrt{Kummer}'sche Theorem}{632}
\inhaltpar{175}{Die Einheitswurzeln im Körper $\Omega_m$}{638}
\inhaltpar{176}{Der in $\Omega_m$ enthaltene reelle Körper $H_m$}{641}
\inhaltpar{177}{Die Primideale des Körpers $H_m$}{643}
\inhaltpar{178}{Die Einheiten des Körpers $H_m$}{645}

\inhaltabschnitt{Zweiundzwanzigster Abschnitt.}{Abel'sche Körper und Kreistheilungskörper.}
\inhaltpar{179}{Einfache \gesperrt{Abel}'sche Körper}{648}
\inhaltpar{180}{Die Resolventen}{651}
\inhaltpar{181}{Vorbereitung zum Beweis}{654}
\inhaltpar{182}{Beweis des ersten Hülfssatzes für ein ungerades $m$}{659}
\inhaltpar{183}{Beweis des zweiten Hülfssatzes für eiu ungerades $m$}{666}
\inhaltpar{184}{Vorläufiges über den Fall eines geraden $m$}{667}

\inhaltabschnitt{Dreiundzwanzigster Abschnitt.}{Classenzahl.}
\inhaltpar{185}{Der \gesperrt{Dirichlet}'sche Satz über die Einheiten}{670}
\inhaltpar{186}{Systeme unabhängiger Einheiten und Exponentensysteme der Einheiten}{677}
\inhaltpar{187}{Fundamentalsysteme von Einheiten}{681}
\inhaltpar{188}{Reducirte Zahlen}{684}
\inhaltpar{189}{Grenzen der Anzahl der durch ein Ideal theilbaren ganzen Zahlen des Körpers $\Omega$}{686}
\inhaltpar{190}{Bestimmung des Volumens}{690}
\inhaltpar{191}{Sätze aus der Reihenlehre}{694}
\inhaltpar{192}{Anwendung auf die Bestimmung der Classenzahl}{701}

\inhaltabschnitt{Vierundzwanzigster Abschnitt.}{Classenzahl der Kreistheilungskörper.}
\inhaltpar{193}{Classenzahldarstellung im Kreistheilungskörper $\Omega_m$}{705}
\inhaltpar{194}{Bestimmung der Summen $X$}{708}
\inhaltpar{195}{Ueber die Classenzahl in dem in $\Omega_m$ enthaltenen reellen Körper}{711}
\inhaltpar{196}{Classenzahl im Körper der achten Einheitswurzeln}{717}
\inhaltpar{197}{Recurrente Berechnung der Classenzahl im Körpar $\Omega_m$, wenn $m$ eine Potenz von 2 ist}{719}
\inhaltpar{198}{Der Classenzahlfactor $A$}{722}
\inhaltpar{199}{Der Classenzahlfactor $B$}{726}

\clearpage
% ===== front matter, scan leaf 19 =====
% !PRINTED XIV
\weberpage{XIV}
\inhaltpar{200}{Normaleinheiten in $H_m$}{728}
\inhaltpar{201}{Fundamentalsystem von Einheiten des Körpers $H_m$}{735}
\inhaltpar{202}{Positive Einheiten}{742}
\inhaltabschnitt{Fünfundzwanzigster Abschnitt.}{Transcendente Zahlen.}
\inhaltpar{203}{Abzählbare Mengen}{745}
\inhaltpar{204}{Unzählbare Mengen}{748}
\inhaltpar{205}{Transcendenz der Zahl $e$}{751}
\inhaltpar{206}{Transcendenz der Zahl $\pi$}{756}
\inhaltpar{207}{Der allgemeine Satz von \gesperrt{Lindemann} über die Exponentialfunction}{760}
\inhaltbuch{}{Nachträge.}
\inhaltzeile{I. Irreducibilität der Kreistheilungsgleichung}{769}
\inhaltzeile{II. Die Irreducibilität der Kreistheilungsgleichung und der Satz über die in einer Linearform enthaltenen Primzahlen}{777}
\inhaltzeile{Register zu Band I und II}{785}
\inhaltzeile{Berichtigungen}{795}

\clearpage
\pagenumbering{arabic}

% ===== scan leaf 22, printed 1 =====
% !PRINTED 1
\weberpage{1}

\begin{center}
\gesperrt{ERSTES BUCH.}

\rule{2cm}{0.4pt}

\vspace{2em}
\gesperrt{GRUPPEN.}

\rule{4cm}{0.4pt}
\end{center}

\vspace{4em}
%% [running head, not body text -- commented out 2026-08-07] Page furniture (running
%% head, folio, signature) is excluded from this edition by standing policy.
%% \noindent{\small Weber, Algebra.\ II. \hfill 1}


% ===== scan leaf 23, printed 2 =====
\weberpage{2}% page number not printed on this page
% !PRINTED none
% Blank page (verso of the Book-I / "GRUPPEN." half-title, page 1). No printed
% text is visible in any of the three bands beyond scan noise/specks.


% ===== scan leaf 24, printed 3 =====
\weberpage{3}% page number not printed on this page
\pdfbookmark[0]{Erstes Buch. Gruppen.}{bm1}
\pdfbookmark[1]{Erster Abschnitt. Allgemeine Gruppentheorie.}{bm2}
\pdfbookmark[2]{§ 1. Definition der Gruppen}{bm3}
% !PRINTED none

\begin{center}
\gesperrt{Erster Abschnitt.}

\vspace{1em}
\gesperrt{Allgemeine Gruppentheorie.}

\rule{2cm}{0.4pt}

\vspace{1em}
\S.\,1.

\gesperrt{Definition der Gruppen.}
\end{center}

\noindent Wir haben im ersten Bande bei den Permutationen den Begriff einer Gruppe kennen gelernt und wichtige algebraische Anwendungen von ihm gemacht. Es muss nun unsere n\"achste Aufgabe sein, diesen in der ganzen neueren Mathematik so \"uberaus wichtigen Begriff allgemeiner zu fassen und die dabei herrschenden Gesetze kennen zu lernen. Wir stellen folgende Definition an die Spitze:

Ein System $P$ von Dingen (Elementen) irgend welcher Art wird zur Gruppe, wenn folgende Voraussetzungen erf\"ullt sind:

\noindent 1.\ \gesperrt{Es ist eine Vorschrift gegeben, nach der aus einem ersten und einem zweiten Elemente des Systems ein ganz bestimmtes drittes Element desselben Systems abgeleitet wird.}

Man schreibt symbolisch, wenn $a$ das erste, $b$ das zweite, $c$ das dritte Element ist:
\[
ab = c, \quad c = ab,
\]
und nennt $c$ aus $a$ und $b$ \gesperrt{componirt} und $a$ und $b$ die \gesperrt{Componenten} von $c$.

Bei dieser Composition wird im Allgemeinen nicht das commutative Gesetz vorausgesetzt, d.\,h. es kann $ab$ von $ba$ verschieden sein, dagegen wird

\noindent 2.\ \gesperrt{das associative Gesetz vorausgesetzt,} d.\,h. wenn $a, b, c$ irgend drei Elemente aus $P$ sind, so ist
\[
(ab)c = a(bc),
\]
%CONT>>


% ===== scan leaf 25, printed 4 =====
% !PRINTED 4
\weberpage{4}

%% [running head, not body text -- commented out 2026-08-07] Page furniture (running
%% head, folio, signature) is excluded from this edition by standing policy.
%% \noindent{\small 4 \hfill Erster Abschnitt. \hfill \S.\,1.}

%<<CONT
und hieraus folgt durch die Schlussweise der vollst\"andigen Induction, dass man immer zu demselben Resultate kommt, wenn man in einer beliebigen Reihe von Elementen aus $P$ in endlicher Anzahl, $a, b, c, d \ldots$ zuerst zwei benachbarte Elemente componirt, dann wieder zwei benachbarte u.\,s.\,w., bis die ganze Reihe auf ein Element reducirt ist, das mit $abcd \ldots$ bezeichnet wird. So ist z.\,B.:
\[
\begin{aligned}
abcd &= (ab)cd = [(ab)c]d = (ab)(cd)\\
&= a(bc)d = [a(bc)]d = a[(bc)d]\\
&= ab(cd) = (ab)(cd) = a[b(c.d)]^{1)}. % [sic: the source prints a dot between c and d]
\end{aligned}
\]

\noindent 3.\ \gesperrt{Es wird vorausgesetzt. dass, wenn $ab = ab'$ oder $ab = a'b$ ist, nothwendig im ersten Falle $b = b'$, im zweiten $a = a'$ sein muss.} % [sic: the source prints a full stop after vorausgesetzt]

Wenn $P$ eine endliche Anzahl von Elementen umfasst, so heisst die Gruppe eine \gesperrt{endliche} und die Anzahl ihrer Elemente ihr \gesperrt{Grad}.

F\"ur endliche Gruppen ergiebt sich aus 1., 2., 3. die \gesperrt{Folgerung:}

\noindent 4.\ \gesperrt{Wenn von den drei Elementen $a, b, c$ aus $P$ zwei beliebig gegeben sind, so kann man das dritte immer und nur auf eine Weise so bestimmen, dass}
\[
ab = c
\]
\gesperrt{ist.}

Sind n\"amlich $a, b$ die gegebenen Elemente, so f\"allt die Behauptung 4. mit 1. zusammen. Ist aber $a$ und $c$ gegeben, so lasse man in dem Compositum $ab$ die zweite Componente $b$ das ganze System $P$ durchlaufen, dessen Grad $= n$ sei. Dann erh\"alt man nach 1. und 3. in $ab$ lauter verschiedene Elemente von $P$, und da ihre Anzahl $= n$ ist, so m\"ussen alle Elemente von $P$, also auch $c$, darunter vorkommen. Ebenso schliesst man, wenn $b$ und $c$ gegeben sind, indem man $a$ das ganze System $P$ durchlaufen l\"asst.

F\"ur unendliche Gruppen kann nicht mehr so geschlossen werden. \gesperrt{F\"ur unendliche Gruppen wird also noch die Eigenschaft 4. als Forderung in die Begriffsbestimmung mit aufgenommen.}

\noindent\rule{3cm}{0.4pt}

\noindent{\footnotesize $^{1)}$ Vgl.\ \gesperrt{Dirichlet-Dedekind}, Vorlesungen \"uber Zahlentheorie, \S.\,2.}


% ===== scan leaf 26, printed 5 =====
% !PRINTED 5
\weberpage{5}

Bei den im ersten Bande betrachteten Permutationsgruppen wird man leicht die Merkmale des allgemeinen Gruppenbegriffes erkennen.

Wir ziehen nun aus dieser Definition zunächst einige ganz allgemeine Folgerungen.

Nach 4. giebt es für jedes gegebene $b$ \gesperrt{ein} Element $e$ in $P$, das der Bedingung
\weq{1}{eb = b}
genügt, und dies $e$ ist von $b$ unabhängig; denn aus (1) folgt für jedes $c$
\[
ebc = bc,
\]
und $bc$ kann nach 4. jedes Element in $P$ bedeuten. Ebenso giebt es ein Element $e'$, das für jedes $b$ der Bedingung
\weq{2}{be' = b}
genügt. Dies Element $e'$ ist aber von $e$ nicht verschieden; denn setzen wir $b = e'$ in (1) und $b = e$ in (2), so folgt
\[
e e' = e', \, . \, e e' = e, % [sic: the source prints a full stop between the two identities]
\]
also
\[
e = e'.
\]

Das Element $e$ ändert nichts, wenn es mit irgend welchen Elementen aus $P$ componirt wird, und wird die \gesperrt{Einheit der Gruppe} genannt.

In vielen Fällen kann es ohne Missverständniss geradezu mit „1“ bezeichnet werden.

Zu jedem Element $a$ giebt es nach 4. ein bestimmtes Element $a^{-1}$, das der Bedingung
\weq{3}{a^{-1} a = e}
genügt. Aus (3), (1) und (2) folgt
\[
a^{-1} a a^{-1} = e a^{-1} = a^{-1} = a^{-1} e,
\]
und folglich nach 3.
\weq{4}{a a^{-1} = e.}

Die beiden Elemente $a$, $a^{-1}$ heissen zu einander \gesperrt{entgegengesetzt} oder \gesperrt{reciprok}.

In besonderen Fällen kann bei der Composition der Elemente einer Gruppe $P$ auch das commutative Gesetz gelten, d.\,h. es kann für je zwei Elemente $a$, $b$ der Gruppe
\[
a b = b a
\]
sein.


% ===== scan leaf 27, printed 6 =====
% !PRINTED 6
\weberpage{6}

\gesperrt{5. Gruppen, die diese Eigenschaft haben, heissen commutative Gruppen oder auch Abel'sche Gruppen.}

Wenn sich die Elemente zweier Gruppen
\[
a, b, c, d \ldots
\]
und
\[
a', b', c', d' \ldots
\]
in der Weise gegenseitig eindeutig entsprechen, dass immer, wenn $ab = c$ ist, auch $a'b' = c'$ wird, so heissen die Gruppen \gesperrt{isomorph}, und es gilt der evidente Satz, dass zwei mit einer dritten isomorphe Gruppen unter einander isomorph sind. Man kann hiernach alle unter einander isomorphe Gruppen zu einer Classe von Gruppen zusammenfassen, die selbst wieder eine Gruppe ist, deren Elemente die Gattungsbegriffe sind, die man erhält, wenn man die entsprechenden Elemente der einzelnen isomorphen Gruppen zu einem Allgemeinbegriff zusammenfasst. Die einzelnen unter einander isomorphen Gruppen sind dann als verschiedene Repräsentanten eines Gattungsbegriffes aufzufassen.

Die Eigenschaften der Gruppen, die in Betracht kommen können, sind von verschiedener Art. Sie können nämlich entweder den besonderen Gruppen anhaften und aus der Natur der Elemente abgeleitet sein, aus denen die Gruppe besteht, oder auch aus der Natur des Compositionsgesetzes. Oder sie können den Gruppen als solchen anhaften und müssen sich dann lediglich aus der Definition des Gruppenbegriffes ableiten lassen. Die letzteren Eigenschaften kommen allen isomorphen Gruppen gemeinsam zu und können als invariante Eigenschaften der Gruppe bezeichnet werden. Wenn ein Vergleich gestattet ist, so könnte man an den Unterschied zwischen den metrischen und projectiven Eigenschaften in der Geometrie erinnern.

Zu der ersten Art der Eigenschaften, die aus der besonderen Natur der Elemente abgeleitet werden, gehören z.\,B. bei den Permutationsgruppen die Eigenschaften der Transitivität und Intransitivität, der Primitivität und Imprimitivität; zu den invarianten Eigenschaften gehören die Vertauschbarkeit oder Nichtvertauschbarkeit, der Grad, die Divisoren und ihr Index, die Normaltheiler. Mit diesen invarianten Eigenschaften, bei denen es gleichgültig ist, aus welchem Repräsentanten einer Classe isomorpher Gruppen sie abgeleitet sind, haben wir uns zunächst zu beschäftigen.


% ===== scan leaf 28, printed 7 =====
% !PRINTED 7
\weberpage{7}
\pdfbookmark[2]{§ 2. Die Divisoren endlicher Gruppen}{bm4}

\begin{center}
\S.\,2.

\gesperrt{Die Divisoren endlicher Gruppen.}
\end{center}

Es ist, wie wir schon im ersten Bande gesehen haben, eine besonders wichtige Frage, ob ein Theil der Elemente einer Gruppe selbst wieder eine Gruppe ist. Dazu ist nothwendig und hinreichend, dass je zwei Elemente dieses Theiles bei der Zusammensetzung wieder ein demselben Theile angehöriges Element ergeben. Diese kleinere Gruppe heisst dann ein \gesperrt{Theiler} oder \gesperrt{Divisor}\footnote{Auch Untergruppe genannt.} der ganzen Gruppe. Das Einheitselement „1“ bildet in jeder Gruppe für sich eine Gruppe, ist also in jeder Gruppe ein Divisor.

Wir beschränken uns jetzt, wie in der Folge fast immer, auf \gesperrt{endliche Gruppen}, und wollen annehmen, es sei
\[
P = 1, a_1, a_2 \ldots a_{n-1}
\]
eine Gruppe vom Grade $n$ mit den Elementen $a_0 = 1$, $a_1$, $a_2 \ldots a_{n-1}$ und
\[
Q = 1, a_1, a_2 \ldots a_{\nu-1}
\]
sei ein Theiler von $P$ vom Grade $\nu$. Jeder solche Theiler $Q$ muss die Grundeigenschaften einer Gruppe haben und enthält also sicher das Einheitselement „1“ und mit jedem seiner Elemente $a_1$ zugleich das entgegengesetzte $a_1^{-1}$. Ist $\nu < n$, so nennen wir $Q$ einen \gesperrt{echten Theiler} von $P$. Ist dann also $b$ ein nicht in $Q$ enthaltenes Element von $P$, so sind die Elemente
\[
Q_1 = b, a_1 b, a_2 b \ldots a_{\nu-1} b
\]
alle von einander verschieden (\S.\,1, 3.) und alle nicht in $Q$ enthalten; denn wenn etwa $a_1 b$ in $Q$ enthalten wäre, so müsste auch, da $Q$ eine Gruppe ist, $a_1^{-1} a_1 b = b$ in $Q$ enthalten sein, gegen die Voraussetzung. Ist mit $Q$ und $Q_1$ die ganze Gruppe $P$ noch nicht erschöpft, so nehmen wir eines der noch übrigen Elemente, das wir mit $c$ bezeichnen können, und bilden
\[
Q_2 = c, a_1 c, a_2 c \ldots a_{\nu-1} c,
\]
und überzeugen uns leicht, dass die Elemente von $Q_2$ alle nicht nur von einander, sondern auch von den Elementen von $Q$ und
%CONT>>


% ===== scan leaf 29, printed 8 =====
% !PRINTED 8
\weberpage{8}
%<<CONT
$Q_1$ verschieden sind. Denn wäre etwa $a_1 c = a_2 b$, so würde folgen, dass $c = a_1^{-1} a_2 b$ sein müsste; es ist aber $a_1^{-1} a_2$ in $Q$ enthalten, also mit einem der Elemente $1, a_1, a_2 \ldots a_{\nu-1}$ identisch, und es wäre also $c$ gegen die Voraussetzung in $Q_1$ enthalten. So fahren wir fort, die Systeme $Q, Q_1, Q_2 \ldots Q_{\mu-1}$ zu bilden, von denen das letzte
\[
Q_{\mu-1} = g, a_1 g, a_2 g \ldots a_{\nu-1} g
\]
sei. Dann muss aber $P$ durch die Systeme $Q, Q_1, Q_2 \ldots Q_{\mu-1}$ erschöpft sein, und es folgt, da diese Systeme lauter verschiedene Elemente enthalten,
\[
n = \nu \mu,
\]
oder der Satz:

\noindent 1.\ \gesperrt{Der Grad einer Gruppe ist durch den Grad jedes ihrer Divisoren theilbar. Der Quotient $\mu = n : \nu$ soll der Index des Theilers $Q$ heissen.}

Die Systeme $Q_1, Q_2 \ldots Q_{\mu-1}$ nennen wir, wie in dem speciellen Falle der Permutationsgruppen, die zu $Q$ gehörigen \gesperrt{Nebengruppen} und bezeichnen sie durch
\[
Q_1 = Qb, \quad Q_2 = Qc \ldots Q_{\mu-1} = Qg,
\]
so dass wir auch symbolisch
\weq{1}{P = Q + Qb + Qc + \cdots + Qg}
setzen können. Bisweilen werden wir auch das ganze System $Q, Q_1 \ldots Q_{\mu-1}$ ($Q$ selbst eingeschlossen) als ein System von Nebengruppen bezeichnen, und also die Darstellung (1) die \gesperrt{Zerlegung von $P$ in ein System von Nebengruppen} nennen.

Aus der Bildungsweise der Nebengruppen folgt noch, dass, wenn $b$, $c$ irgend zwei Elemente aus $P$ sind, die beiden Systeme $Qb$ und $Qc$ entweder ganz identisch sind oder kein gemeinsames Element enthalten. Denn ist $a_1 b = a_2 c$, worin $a_1$, $a_2$ zwei Elemente aus $Q$ sind, so ist auch für jedes andere Element $a$ aus $Q$:
\[
a a_1 b = a a_2 c,
\]
und wenn $a$ die ganze Gruppe $Q$ durchläuft, so durchläuft, wie schon in \S.\,1 bemerkt, auch jedes der beiden $a a_1$ und $a a_2$ die ganze Gruppe $Q$; folglich sind $Qb$ und $Qc$ identisch.

Wie in Bd.\,I, \S.\,154, 2. können wir $P$ auch in der folgenden Art in Nebengruppen zerlegen:
\[
P = Q + b^{-1} Q + c^{-1} Q + \cdots + g^{-1} Q.
\]


% ===== scan leaf 30, printed 9 =====
% !PRINTED 9
\weberpage{9}

Die Nebengruppen haben nicht die Merkmale einer Gruppe. Denn damit zwei
Elemente aus $Q_1$ bei der Zusammensetzung wieder ein Element von $Q_1$
ergeben, müsste etwa
\[
a_1 b a_2 b = a_3 b,
\]
also $a_1 b a_2 = a_3$, $b = a_1^{-1} a_3 a_2^{-1}$ sein. Es wäre also $b$ der
Voraussetzung entgegen in $Q$ enthalten.

Die Benennung Nebengruppe ist also nur uneigentlich zu verstehen.

Zwei Theiler $Q$, $Q'$ von $P$ haben immer das Element 1 mit einander
gemein. Sie können aber auch noch andere Elemente gemeinschaftlich haben,
und diese gemeinschaftlichen Elemente bilden eine Gruppe, die wir mit $R$
bezeichnen wollen. Denn gehören die Elemente $a$ und $b$ sowohl zu $Q$ als
zu $Q'$, so gilt wegen des Gruppencharakters von $Q$ und $Q'$ dasselbe von
dem Compositum $ab$; also ist $R$ eine Gruppe. Diese gemeinschaftliche
Gruppe nennen wir den \gesperrt{grössten gemeinschaftlichen Theiler} von $Q$ und $Q'$
oder auch, mit einem geometrischen Anklange, den \gesperrt{Durchschnitt von} $Q$ und
$Q'$. Ebenso folgt, dass die Elemente, die irgend einer Anzahl von
Theilern von $P$ gemeinsam sind, eine Gruppe bilden, die wir ebenso als
den grössten gemeinschaftlichen Theiler oder den Durchschnitt aller dieser
Gruppen bezeichnen.

In jeder Gruppe können wir nach folgendem Verfahren Theiler bilden.

Ein Theiler ist immer das Einheitselement für sich.

Bezeichnen wir die wiederholte Zusammensetzung eines Elementes mit sich
selbst durch Potenzen (mit $a^0$ das Einheitselement), so kann die Reihe
der Elemente
\[
1,\ a,\ a^2,\ a^3 \ldots,
\]
die alle der Gruppe $P$ angehören, nicht lauter verschiedene Elemente
enthalten. Ist also das erste Element, das zum zweiten Male wiederkehrt,
\[
a^\mu = a^{\mu+\alpha},
\]
so folgt, dass $a^\alpha = 1$ sein muss, und dass $\alpha$ die kleinste
positive Zahl ist, die dieser Bedingung genügt. Es ist dann, wenn
$m = q\alpha$ ein beliebiges Vielfaches von $\alpha$ ist, $a^m = 1$ und
umgekehrt: so oft $a^m = 1$ ist, muss $m$ durch $\alpha$ theilbar sein;
denn sonst könnte man $m = q\alpha + \alpha'$ setzen, worin $\alpha'$
positiv und kleiner als $\alpha$ ist, und es wäre $a^{\alpha'} = 1$, gegen
die Voraussetzung,
%CONT>>


% ===== scan leaf 31, printed 10 =====
% !PRINTED 10
\weberpage{10}
\pdfbookmark[2]{§ 3. Normaltheiler einer Gruppe}{bm5}

%<<CONT
dass $\alpha$ der kleinste positive Exponent sei, für den $a^\alpha = 1$
ist. Es ist immer und nur dann
\[
a^\mu = a^{\mu'},
\]
wenn $\mu \equiv \mu'\ (\mathrm{mod}\ \alpha)$, und hierin können die
Exponenten auch negativ sein, wenn $a^{-r}$ die $r^{\mathrm{te}}$ Potenz
des Elementes $a^{-1}$ oder das zu $a^r$ entgegengesetzte Element
bedeutet.

Die Reihe
\[
A = 1,\ a,\ a^2 \ldots a^{\alpha-1}
\]
enthält dann $\alpha$ von einander verschiedene Elemente von $P$, wobei
alle Potenzen von $a$ vertreten sind. Die Elemente $A$ bilden aber
offenbar eine Gruppe vom Grade $\alpha$, weil sich die Exponenten bei der
Zusammensetzung einfach addiren. Diese Gruppe ist ein Theiler von $P$ und
also ist $\alpha$ \gesperrt{ein Theiler von $n$}. Es ist also für jedes Element
$a^n = 1$.

\gesperrt{Die Gruppe $A$ heisst die Periode des Elementes $a$ und $\alpha$ wird auch
der Grad des Elementes $a$ genannt.}

\begin{center}
\S.\,3.

\gesperrt{Normaltheiler einer Gruppe.}
\end{center}

Ist $P$ wie oben eine Gruppe und $Q$ ein Divisor von $P$, ferner $b$ ein
nicht in $Q$ enthaltenes Element von $P$, so ist die Nebengruppe $Qb$
keine Gruppe. Dagegen bildet das System $b^{-1}Qb$, oder ausführlich
geschrieben
\[
1,\ b^{-1}a_1 b,\ b^{-1}a_2 b \ldots b^{-1}a_{\nu-1} b
\]
sicher eine Gruppe, weil
\[
b^{-1}a_1 b\ .\ b^{-1}a_2 b = b^{-1}a_1 a_2 b
\]
ist, und diese Gruppe ist mit $Q$ isomorph. Die Gruppe $b^{-1}Qb$ heisst
\gesperrt{die durch $b$ transformirte Gruppe}. Gehört $b$ selbst zu $Q$, so ist
$b^{-1}Qb$ mit $Q$ identisch. Nimmt man für $b$ die verschiedenen Elemente
von $P$, so erhält man eine ganze Schaar solcher Gruppen, die wir die
\gesperrt{zu $Q$ conjugirten Theiler von $P$} oder auch kurz \gesperrt{conjugirte Gruppen} nennen.

Es kann vorkommen, dass alle conjugirten Theiler mit einander identisch
sind. Wir haben schon bei den Permutationsgruppen gesehen, dass dieser
Fall von besonderer Wichtigkeit ist, und führen also nun allgemein
folgende Definition ein:
%CONT>>


% ===== scan leaf 32, printed 11 =====
% !PRINTED 11
\weberpage{11}

%<<CONT
\gesperrt{Wenn $Q$ ein Theiler von $P$ ist, der mit seinen sämmtlichen conjugirten
Theilern identisch ist, so heisst $Q$ ein Normaltheiler von
$P$}\footnote{Auch ausgezeichnete oder invariante Untergruppe genannt.}.

Die aus dem einzigen Element 1 gebildete Gruppe ist ein Normaltheiler von
jeder Gruppe. Wir erhalten ferner einen Normaltheiler in dem grössten
gemeinschaftlichen Theiler $R$ der sämmtlichen mit irgend einem Theiler
$Q$ von $P$ conjugirten Theiler.

Denn es ist schon oben bewiesen, dass $R$ als der Durchschnitt mehrerer
Theiler von $P$ eine Gruppe ist.

Sind nun
\weq{1}{Q,\ Q',\ Q'' \ldots}
die zu $Q$ conjugirten Theiler von $P$ und $b$ irgend ein Element von $P$,
dann ist das System der Gruppen
\weq{2}{b^{-1}Qb,\ b^{-1}Q'b,\ b^{-1}Q''b \ldots}
von dem System (1) nicht verschieden. Wenn aber $R$ der Durchschnitt der
Gruppen (1) ist, so ist $b^{-1}Rb$ der Durchschnitt von (2), und folglich
ist $R$ mit $b^{-1}Rb$ identisch, d.\,h. $R$ ist ein Normaltheiler von $P$.

Ist $N$ ein Normaltheiler irgend einer Gruppe $P$ und $b$ ein beliebiges
Element in $P$, so ist
\weq{3}{b^{-1}Nb = N \quad \text{oder} \quad Nb = bN.}

Ist $P$ vom Grade $n$, $N$ vom Grade $\nu$ und vom Index $\mu$, also
$n = \mu\nu$, so können wir die $\mu$ Elemente $1, b_1 \ldots b_{\mu-1}$ so
wählen, dass die Zerlegung von $P$ in die Nebengruppen
\[
\begin{aligned}
P &= N + Nb_1 + Nb_2 + \cdots + Nb_{\mu-1} \\
  &= N + b_1 N + b_2 N + \cdots + b_{\mu-1} N
\end{aligned}
\]
ergiebt. Es ist also
\weq{4}{N,\ N_1 = Nb_1 = b_1 N,\ N_2 = Nb_2 = b_2 N \ldots,\quad N_{\mu-1} = Nb_{\mu-1} = b_{\mu-1} N}
das System der Nebengruppen.

Ist $a$ irgend ein Element in $N$, so ist $Na$ mit $N$ identisch, also ist
auch $Nab$ mit $Nb$ identisch, und da jedes Element $c$ von $P$ in $N$
oder in einer seiner Nebengruppen vorkommt, also in der Form $ab_k$
enthalten ist, so muss $Nc = cN$ mit einem der Systeme (4)
übereinstimmen.


% ===== scan leaf 33, printed 12 =====
% !PRINTED 12
\weberpage{12}
\pdfbookmark[2]{§ 4. Composition der Theile}{bm6}

Wir erwähnen noch folgenden Satz, dessen Richtigkeit sich unmittelbar aus
dem Isomorphismus transformirter Gruppen ergiebt:

\gesperrt{Ist $Q$ ein Theiler von $P$, $R$ ein Theiler von $Q$, so ist, wenn $Q'$
und $R'$ aus $Q$ und $R$ durch dasselbe Element von $P$ transformirt
sind, auch $R'$ ein Theiler von $Q'$, und wenn $R$ Normaltheiler von $Q$
ist, so ist auch $R'$ Normaltheiler von $Q'$.}

\begin{center}
\S.\,4.

\gesperrt{Composition der Theile.}
\end{center}

Sind $A$ und $B$ irgend zwei Reihen von Elementen aus einer Gruppe $P$
(Gruppen oder nicht), so wollen wir unter dem symbolischen Producte $AB$
den Inbegriff aller Elemente verstehen, die man erhält, wenn man je ein
Element $a$ von $A$ mit je einem Element $b$ von $B$ nach der in $P$
geltenden Vorschrift zu $ab$ componirt. Diese Art der Zusammensetzung von
$A$ und $B$ zu $AB$ wollen wir die \gesperrt{Composition der Theile} (von $P$)
nennen. Wir unterscheiden hier zwischen einem \gesperrt{Theil} und einem \gesperrt{Theiler} von
$P$, so dass ein Theiler immer eine Gruppe sein soll, was bei einem Theil
nicht nothwendig ist.

Man kann ebenso drei und mehr Theile $A, B, C \ldots$ von $P$ componiren,
und es gilt bei der Composition der Theile das associative Gesetz, wie
unmittelbar daraus folgt, dass dieses Gesetz in $P$ gilt. Danach ist die
Bedeutung eines Compositums $ABC \ldots$ eindeutig bestimmt.

In der Zusammensetzung $AB$ kann einer der Theile $A$, $B$ auch aus einem
einzigen Elemente bestehen, und dann stimmt die Bezeichnung $AB$ mit der
oben für die Nebengruppen gebrauchten Bezeichnung überein.

Bei der Composition der Theile gelten folgende Sätze:

\noindent 1.\ \gesperrt{Die nothwendige und hinreichende Bedingung
dafür, dass ein Theil $A$ von $P$ eine Gruppe ist, besteht in der
Gleichung}
\weq{1}{AA = A.}
Denn diese Gleichung besagt nichts Anderes als dass, wenn $a$ und $a'$ in
$A$ enthalten sind, auch $aa'$ zu $A$ gehört, also dass $A$ eine Gruppe
ist.


% ===== scan leaf 34, printed 13 =====
% !PRINTED 13
\weberpage{13}

%% [running head, not body text -- commented out 2026-08-07] Page furniture (running
%% head, folio, signature) is excluded from this edition by standing policy.
%% \gesperrt{\S.\,4.} \hfill \gesperrt{Gruppe der Nebengruppen.} \hfill 13

\gesperrt{2. Durch die beiden Gleichungen}

\weq{2}{AA = A,}
\weq{3}{AB = BA,}

\gesperrt{worin $A$ ein bestimmter, $B$ jeder beliebige Theil von $P$ sein kann, wird ausgesagt, dass $A$ ein Normaltheiler von $P$ ist.}

Denn nach (2) ist $A$ ein Theiler von $P$ und nach (3) ist für jedes Element $b$ von $P$
\[
b^{-1}Ab = A,
\]
d.\,h. $A$ ist ein Normaltheiler.

\gesperrt{3. Sind $A$ und $B$ zwei Theiler von $P$ (Gruppen), so ist eine der beiden Gleichungen}

\weq{4}{AB = A, \quad BA = A,}

\gesperrt{die nothwendige und hinreichende Bedingung dafür, dass $B$ ein Theiler von $A$ ist.}

Denn wenn $B$ eine Gruppe und ein Theiler der Gruppe $A$ ist, und wenn $a$ in $A$, $b$ in $B$ und also auch in $A$ enthalten ist, so ist auch $ab$ in $A$ enthalten. $A$ enthält also jedes Element von $AB$.

Weil aber zweitens $B$ als Gruppe auch das Element 1 enthält, so enthält $AB$ auch jedes Element von $A$, und also ist $AB$ mit $A$ identisch.

Ebenso sieht man, dass $BA$ mit $A$ identisch ist. Andererseits folgt aus jeder der Gleichungen (4), da $A$ als Gruppe das Einheitselement enthält, dass jedes Element von $B$ in $A$ enthalten, also $B$ ein Theiler von $A$ ist.

\gesperrt{4. Bei der Composition der Theile bildet das System der Nebengruppen zu einem Normaltheiler von $P$ selbst eine Gruppe, in der der Normaltheiler $N$ die Einheit bildet, und}
\[
Nb, \quad Nb^{-1}
\]
\gesperrt{entgegengesetzte Elemente sind.}

Denn ist $N$ ein Normaltheiler von $P$ und
\weq{5}{N, \; N_1 = Nb_1, \; N_2 = Nb_2, \dots}
das System der Nebengruppen, so ist
\[
N_h N_k = N b_h N b_k.
\]


% ===== scan leaf 35, printed 14 =====
% !PRINTED 14
\weberpage{14}

%% [running head, not body text -- commented out 2026-08-07] Page furniture (running
%% head, folio, signature) is excluded from this edition by standing policy.
%% \gesperrt{Erster Abschnitt.} \hfill 14 \hfill \gesperrt{\S.\,4.}

Wegen der Eigenschaft der Normaltheiler ist aber $N$ mit jedem $b$ vertauschbar, also $b_h N = N b_h$, und folglich
\[
N_h N_k = N N b_h b_k.
\]

Weil $b_h b_k$ in $P$ enthalten ist, so ist $N b_h b_k$ mit einer der Nebengruppen (5) identisch, und da $NN = N$ gesetzt werden kann (nach 1.), so folgt, wenn wir $N b_h b_k = N_l$ setzen,
\weq{6}{N_h N_k = N_l.}

Damit ist die Eigenschaft \S.\,1, 1. der Gruppen für die Composition der Nebengruppen nachgewiesen. Dass auch \S.\,1, 2., nämlich das associative Gesetz gilt, haben wir schon hervorgehoben. Es ist also noch \S.\,1, 3. nachzuweisen, nämlich dass, wenn
\weq{7}{N_h N_i = N_k N_l}
gesetzt wird, und $N_h = N_k$ ist, auch $N_i = N_l$, und wenn $N_i = N_l$ ist, auch $N_h = N_k$ sein muss.

Nach der Darstellung (5) können wir aber (7) in der doppelten Weise darstellen:
\weq{8}{\begin{aligned}
N b_h b_i &= N b_k b_l \\
b_h b_i N &= b_k b_l N.
\end{aligned}}

Ist nun $N_i = N_l$, so können wir auch $b_i = b_l$ annehmen, und erhalten aus (8) durch Composition mit $b_i^{-1}$
\[
N b_h = N b_k;
\]
und wenn $N b_h = N b_k$ ist, so nehmen wir $b_h = b_k$ an und erhalten aus (8) durch Composition mit $b_h^{-1}$
\[
b_i N = b_l N,
\]
also auch $N_i = N_l$.

Damit ist also erwiesen, dass das System der Nebengruppen durch die Composition der Theile zu einer Gruppe wird; dass in dieser Gruppe der Normaltheiler $N$ das Einheitselement ist, folgt unmittelbar aus 1., weil danach
\[
N N_k = N N b_k = N_k
\]
ist. Und weil $N b_k N b_k^{-1} = N b_k b_k^{-1} = N$ ist, so sind $N b_k$ und $N b_k^{-1}$ entgegengesetzte Elemente.

Die Gruppe der Grössen $N_k$, deren Existenz also hiermit nachgewiesen ist, nennen wir die \gesperrt{Gruppe der Nebengruppen} oder die \gesperrt{zu $N$ complementäre Gruppe} in Bezug auf $P$. Wir
%CONT>>


% ===== scan leaf 36, printed 15 =====
% !PRINTED 15
\weberpage{15}
\pdfbookmark[2]{§ 5. Mehrstufiger Isomorphismus}{bm7}

%% [running head, not body text -- commented out 2026-08-07] Page furniture (running
%% head, folio, signature) is excluded from this edition by standing policy.
%% \gesperrt{\S.\,5.} \hfill \gesperrt{Mehrstufiger Isomorphismus.} \hfill 15

%<<CONT
bezeichnen sie nach dem Vorgange von Hölder\footnote{Mathematische Annalen, Bd.\,34. Das Zeichen erinnert an einen Quotienten, mit dem ja die complementäre Gruppe eine gewisse Analogie hat.} mit $P/N$. Der Grad der complementären Gruppe ist gleich dem Index der Gruppe $N$.

Wir erwähnen am Schluss dieser Betrachtungen noch einen Satz über Normaltheiler.

\noindent 5.\ \gesperrt{Ist $N$ ein Theiler von $P$, $N'$ ein Theiler von $N$ und zugleich Normaltheiler von $P$, so ist auch $N'$ Normaltheiler von $N$.}

Dies ist selbstverständlich, weil für jedes Element $b$ von $P$ $b^{-1}N'b = N'$ ist, und weil jedes Element von $N$ auch ein Element von $P$ ist.

\gesperrt{Man darf aber nicht umgekehrt schliessen, dass, wenn $N$ ein Normaltheiler von $P$, $N'$ ein Normaltheiler von $N$ ist, auch $N'$ ein Normaltheiler von $P$ sein müsse.}

\begin{center}
\S.\,5.

\gesperrt{Mehrstufiger Isomorphismus.}
\end{center}

Ist wie oben $N$ ein Normaltheiler von $P$ vom Grade $\nu$ und Index $\mu$, also $n = \mu \nu$, so ist die zu $P$ complementäre Gruppe $P/N$ vom Grade $\mu$. Wir wollen diese oder eine damit isomorphe Gruppe mit $Q$ bezeichnen und ihre Elemente, die den Nebengruppen $N, Nb_1, Nb_2 \dots$ entsprechen, mit $A, B, C \dots$ Jedem dieser Elemente entsprechen $\nu$ Elemente der Gruppe $P$, etwa

\noindent dem Elemente $A$ die Elemente $a, a_1 \dots a_{\nu-1}$

\noindent dem Elemente $B$ die Elemente $b, b_1 \dots b_{\nu-1}$

\noindent $. \; . \; . \; . \; . \; . \; . \; . \; . \; . \; . \; . \; . \; . \; . \; . \; . \; . \; . \; . \; .$

Dies Entsprechen ist dann derart, dass jedes zusammengesetzte Element $ab$ dem zusammengesetzten Elemente $AB$ entspricht. Denn die Elemente $A$ und $B$ entsprechen den Nebengruppen $Na$ und $Nb$, und es ist, weil $N$ ein Normaltheiler ist,
\[
Na \, Nb = Nab.
\]

Es gilt also hier bei der Zusammensetzung der $A, B \dots$ einerseits und der $a, b \dots$ andererseits dasselbe Gesetz, wie bei
%CONT>>


% ===== scan leaf 37, printed 16 =====
% !PRINTED 16
\weberpage{16}

%% [running head, not body text -- commented out 2026-08-07] Page furniture (running
%% head, folio, signature) is excluded from this edition by standing policy.
%% 16 \hfill \gesperrt{Erster Abschnitt.} \hfill \gesperrt{\S.\,5.}

%<<CONT
den isomorphen Gruppen (\S.\,1), nur mit dem Unterschiede, dass jedem Elemente $A$ nicht bloss ein Element $a$, sondern mehrere entsprechen. Diese Thatsache giebt Anlass, den Begriff des Isomorphismus zu erweitern, wie es durch folgende Definition geschieht:

\gesperrt{Man nennt eine Gruppe $P$ mit den Elementen $a, b, c \dots$ (mehrstufig) isomorph mit einer Gruppe $Q$ mit den Elementen $A, B, C \dots$, wenn beide Gruppen so auf einander bezogen werden können, dass jedem der Elemente $A, B, C \dots$ ein oder mehrere der Elemente $a, b, c \dots$ entsprechen, und zwar so, dass jedes der Elemente von $Q$ einem und nur einem der Elemente von $P$ entspricht, und dass, wenn $a$ und $A$, $b$ und $B$ einander entsprechen, $ab$ dem Elemente $AB$ entspricht.}
% [sic: Weber prints ``Elemente von Q''; the uniqueness condition requires P.]

So lässt sich zunächst beweisen, dass jedem der Elemente $A, B, C \dots$ eine gleiche Zahl von Elementen $a, b, c \dots$ entspricht und dass also \gesperrt{der Grad von $Q$ ein Theiler des Grades von $P$ ist}. Denn es sei $A$ das Einheitselement in $Q$, dem die Elemente $a, a_1 \dots a_{\nu-1}$ in $P$ entsprechen mögen. Diese letzteren Elemente müssen dann eine Gruppe vom Grade $\nu$ bilden, die wir mit $N$ bezeichnen wollen, weil nach der Definition des Isomorphismus das Element $a a_1$ dem Elemente $AA = A$ entsprechen muss. Ist dann $b$ ein dem Elemente $B$ entsprechendes Element von $P$, so müssen wiederum alle Elemente $Nb$ demselben Elemente $B$ aus $Q$ entsprechen. Denn jedes $ab$ muss dem Elemente $AB = B$ entsprechen. Sind $b_1, b$ zwei dem Elemente $B$ entsprechende Elemente, so entspricht $b_1 b^{-1} = a$ dem Elemente $A$, und ist also in $N$ enthalten, also ist $b_1 = ab$ in $Nb$ enthalten. Durch $Nb$ sind also alle Elemente, die dem Elemente $B$ entsprechen, erschöpft, und jedem Elemente von $Q$ entsprechen $\nu$ Elemente von $P$. Der Isomorphismus heisst \gesperrt{$\nu$-stufig}. Jedes Element $b^{-1}ab$ entspricht aber gleichfalls dem Einheitselemente $A$ und also ist $b^{-1}Nb = N$, d.\,h. $N$ ist Normaltheiler von $P$, und $Q$ ist einstufig isomorph mit $P/N$.

Wir sehen also, dass es einen anderen mehrstufigen Isomorphismus als den zwischen der complementären Gruppe eines Normaltheilers und der Gesammtgruppe nicht giebt. Man könnte aber den Begriff des Isomorphismus noch dahin erweitern, dass man zwei Gruppen, die zu einer dritten Gruppe $\mu$- und $\nu$-stufig
%CONT>>


% ===== scan leaf 38, printed 17 =====
% !PRINTED 17
\weberpage{17}
\pdfbookmark[2]{§ 6. Die Compositionsreihe und der Satz von C. Jordan}{bm8}
%<<CONT
isomorph sind, $\mu$-$\nu$-stufig isomorph zu einander nennt. Bei
Weitem der wichtigste ist der einstufige Isomorphismus, den wir
daher als Isomorphismus schlechtweg bezeichnen, während ein
mehrstufiger Isomorphismus immer durch einen Zusatz kenntlich
gemacht werden soll $^{1)}$.

\begin{center}
\S.\,6.

\gesperrt{Die Compositionsreihe und der Satz von C. Jordan.}
\end{center}

Eine Gruppe $P$ heisst \gesperrt{einfach}, wenn sie ausser sich selbst
und der Einheit keinen Normaltheiler hat, \gesperrt{zusammengesetzt},
wenn noch andere Normaltheiler vorhanden sind.

Ein Normaltheiler, der keinen anderen Normaltheiler \gesperrt{über
sich hat,} der also nicht Theil eines anderen echten Normaltheilers
 von $P$ ist, heisst \gesperrt{ein grösster Normaltheiler}
von $P$ $^{2)}$. Ist $P_1$ ein grösster Normaltheiler von $P$, $P_2$ ein
grösster Normaltheiler von $P_1$, $P_3$ ein grösster Normaltheiler
von $P_2$ u. s. f., so heisst die Reihe von Gruppen

\weq{1}{P,\, P_1,\, P_2,\, P_3 \ldots 1,}

die nothwendig mit der aus dem Einheitselemente allein gebildeten
 Gruppe 1 endigt, eine \gesperrt{Compositionsreihe} der
Gruppe $P$.

Wir wollen die Grade der Gruppen (1) mit $n, n_1, n_2, n_3 \ldots 1$
bezeichnen, und die Quotienten $n : n_1$, $n_1 : n_2$, $n_2 : n_3 \ldots$, also
die Indices von $P_1$ in Bezug auf $P$, $P_2$ in Bezug auf $P_1$ u. s. f.
mit $\nu_1, \nu_2, \nu_3 \ldots$ Die Reihe der ganzen Zahlen

\weq{2}{\nu_1,\, \nu_2,\, \nu_3 \ldots}

nennen wir die \gesperrt{Indexreihe} der Gruppe $P$.

Die grössten Normaltheiler $P_1, P_2, P_3 \ldots$ können für eine
gegebene Gruppe $P$ im Allgemeinen auf mehrere verschiedene
Arten ausgewählt werden, und danach können auch die Gradzahlen
 $n_1, n_2, n_3 \ldots$ und die Indices $\nu_1, \nu_2, \nu_3 \ldots$ verschieden
%CONT>>

$^{1)}$ Vgl. C.\,\gesperrt{Jordan}, Traité des substitutions. \gesperrt{Netto}, Substitutionentheorie
. Die Bezeichnung \glqq{}$\mu$-stufig\grqq{} rührt von \gesperrt{Netto} her. Nach \gesperrt{Jordan}
heisst der einstufige Isomorphismus \glqq{}holoëdrischer\grqq{}, der mehrstufige
\glqq{}meroëdrischer\grqq{} Isomorphismus. Der einstufige Isomorphismus wird bisweilen
 auch als Aequivalenz bezeichnet.

$^{2)}$ Ausgezeichnete Maximal-Untergruppe nach anderer Bezeichnung.

%% [page furniture, not body text -- commented out 2026-08-10] Page furniture (running
%% head, folio, signature) is excluded from this edition by standing policy.
%% Weber, Algebra. II. \hfill 2


% ===== scan leaf 39, printed 18 =====
% !PRINTED 18
\weberpage{18}
%<<CONT
ausfallen. Es gilt aber der folgende schöne und wichtige Satz
von C.\,\gesperrt{Jordan} $^{1)}$:

\gesperrt{I. Wie auch die Compositionsreihe einer Gruppe
$P$ gewählt sein mag, die Indexreihe ist, von
der Anordnung abgesehen, immer dieselbe.}

Der Beweis beruht auf der Betrachtung der im vorigen
Paragraphen eingeführten complementären Gruppen.

Es seien $Q$ und $Q_1$ Normaltheiler von $P$ und zugleich $Q_1$
ein Theiler (also nach \S.\,4, 5. Normaltheiler) von $Q$.

Dann gilt der Satz:

\gesperrt{1. Die Gruppe $Q/Q_1$ ist ein Normaltheiler der
Gruppe $P/Q_1$, und der Index von $Q/Q_1$ in Bezug
auf $P/Q_1$ ist gleich dem Index von $Q$ in Bezug
auf $P$.}

Um seine Richtigkeit einzusehen, braucht man nur die
Bildungsweise der einzelnen Gruppen genauer zu betrachten.

Es seien $n, q, q_1$ die Grade von $P, Q, Q_1$ und

\weq{3}{n = \nu q,\quad q = \nu_1 q_1,\quad n = \mu q_1,\quad \mu = \nu \nu_1.}

Man zerlegt nun $Q$ in die Nebengruppen zu $Q_1$, d. h. man
setzt, wenn $b_1, b_2 \ldots b_{\nu_1 - 1}$ passend bestimmte Elemente in $Q$ sind,

\weq{4}{Q = Q_1 + Q_1 b_1 + \cdots Q_1 b_{\nu_1 - 1}.}

Die Elemente der Gruppe $Q/Q_1$ sind

\weq{5}{Q_1,\, Q_1 b_1,\, \ldots\, Q_1 b_{\nu_1 - 1}.}

Wenn wir $P$ in die Nebengruppen zu $Q_1$ zerlegen, so kommen
 darunter sicher die Elemente (5) vor, und folglich ist $Q/Q_1$
ein Theiler von $P/Q_1$.

Es ist also noch nachzuweisen, dass dieser Theiler ein
Normaltheiler ist.

Bezeichnen wir aber mit $a$ irgend ein Element von $P$, so
sind alle Elemente von $P/Q_1$ in der Form $Q_1 a$ enthalten, und
wir haben also nur zu zeigen, dass, wenn $b$ irgend ein Element
in $Q$ ist, jedes Element

\weq{6}{Q_1 a^{-1} Q_1 b Q_1 a}

in $Q/Q_1$, d. h. in der Form $Q_1 b$ enthalten ist.

$^{1)}$ C.\,\gesperrt{Jordan}, Traité des substitutions, Paris 1870. \gesperrt{Netto}, Substitutionentheorie,
Leipzig 1882. Wir folgen hier dem Beweise, den
\gesperrt{Hölder} im Anschluss an \gesperrt{Netto} gegeben hat (Math. Annalen, Bd. 34).


% ===== scan leaf 40, printed 19 =====
% !PRINTED 19
\weberpage{19}
Dies folgt aber unmittelbar aus

\weq{7}{Q_1 a^{-1} Q_1 b Q_1 a = Q_1 a^{-1} b a,}

weil zugleich mit $b$ auch $a^{-1} b a$ in $Q$ vorkommt (weil $Q$ ein
Normaltheiler von $P$ ist). Und dass auch die Bestimmung der
Indices richtig ist, zeigt die Abzählung. Denn die Grade von
$P/Q_1$ und $Q/Q_1$ sind $\nu \nu_1$ und $\nu_1$, also der Index $\nu$.

Daneben besteht folgender Satz:

\gesperrt{2. Ist $Q_1$ ein Normaltheiler von $P$ vom Grade $q_1$,
und hat die Gruppe $P/Q_1$ selbst einen Theiler $M$
vom Grade $m$, so hat $P$ einen Theiler $Q$ vom
Grade $q = m q_1$. Zugleich ist $Q_1$ ein Normaltheiler
 von $Q$, und es ist $M = Q/Q_1$. Ist $M$ Normaltheiler
 von $P/Q_1$, so ist auch $Q$ Normaltheiler
von $P$.}

Es sei $\mu$ der Index von $Q_1$ in Bezug auf $P$, und $P$ sei in
die Nebengruppen zerlegt:

\weq{8}{P = Q_1 + Q_1 e_1 + Q_1 e_2 \cdots + Q_1 e_{\mu - 1},}

so dass $e_1, e_2 \ldots e_{\mu - 1}$ passend gewählte Elemente in $P$ und

\weq{9}{Q_1,\, Q_1 e_1,\, Q_1 e_2 \ldots Q_1 e_{\mu - 1}}

die Elemente von $P/Q_1$ sind. Wenn nun in der Gruppe $P/Q_1$
ein Theiler $M$ enthalten ist, so muss dieser aus einem Theile
der Elemente (9) bestehen, etwa aus

\weq{10}{Q_1,\, Q_1 b_1,\, Q_1 b_2 \ldots Q_1 b_{\nu_1 - 1},}

und wenn dies System eine Gruppe bilden soll, so muss für
irgend zwei Elemente $b_i, b_k$ das Product

\weq{11}{Q_1 b_i Q_1 b_k = Q_1 b_i b_k}

wieder in (10) enthalten, also etwa $= Q_1 b_h$ sein. Wenn also
das System $1, b_1, b_2 \ldots b_{\nu_1 - 1}$ mit $B$ bezeichnet wird, so lässt
sich nach der Methode der Composition der Theile die Relation
 (11) so schreiben:

\weq{12}{Q_1 B Q_1 B = Q_1 B,}

und so besagt sie nach \S.\,4, 1., dass die in

\weq{13}{Q = Q_1 B}

enthaltenen Elemente von $P$ eine Gruppe bilden, die die Gruppe
$Q_1$ als Theiler enthält und selbst ein Theiler von $P$ ist. Dass
$Q_1$ Normaltheiler von $Q$ ist, folgt aus \S.\,4, 5., weil $Q_1$ als Normal-
%CONT>>


% ===== scan leaf 41, printed 20 =====
% !PRINTED 20
\weberpage{20}
%<<CONT
theiler von $P$ vorausgesetzt ist, und aus der Darstellung (10)
der Gruppe $M$ ergiebt sich, dass $M = Q/Q_1$ ist.

Es ist noch zu zeigen, dass, wenn $M$ Normaltheiler von
$P/Q_1$ ist, auch $Q$ Normaltheiler von $P$ ist. Dies folgt so:

Ist $e$ irgend ein Element in $P$ und $Q_1 b_k$ ein Element in $M$,
so folgt aus der Annahme, dass $M$ Normaltheiler von $P/Q_1$ ist:

\[ Q_1 e^{-1} Q_1 b_k Q_1 e = Q_1 b_h, \]

was wir auch nach der Composition der Theile durch die Formel

\weq{14}{Q_1 e^{-1} Q_1 B Q_1 e = Q_1 B}

ausdrücken können. Da nun $Q_1$ Normaltheiler von $P$ ist, so
kann $Q_1$ mit jedem der anderen Factoren vertauscht werden, so
dass aus (14) folgt:

\[ e^{-1} Q_1 B e = Q_1 B, \]

oder

\weq{15}{e^{-1} Q e = Q,}

wodurch ausgedrückt ist, dass $Q$ Normaltheiler von $P$ ist.

Aus dem Theorem 1. und 2. ergiebt sich das Corollar:

\gesperrt{3. Ein Normaltheiler $Q$ von $P$ ist dann und nur
dann ein grösster Normaltheiler, wenn die complementäre
 Gruppe $P/Q$ einfach ist.}

Wenn $Q$ und $Q'$ zwei Normaltheiler von $P$ sind, so ist auch
das symbolische Product $Q Q'$ ein Normaltheiler.

Denn es ist, weil $Q$ ein Normaltheiler ist

\[ Q Q' = Q' Q, \]

also

\[ Q Q' Q Q' = Q Q Q' Q' = Q Q', \]

und dies ist nach \S.\,4, 1. das Kennzeichen, dass $Q Q'$ eine Gruppe
ist. Dass es ein Normaltheiler von $P$ ist, ergiebt sich dann
ebenfalls nach \S.\,4, 2. aus

\[ Q Q' B = Q B Q' = B Q Q', \]

wenn $B$ irgend ein Theil von $P$ ist.

Ist nun $Q$ ein \gesperrt{grösster Normaltheiler} von $P$, und $Q'$ von
$Q$ verschieden, so ist $Q Q'$ ein Normaltheiler von $P$, der sowohl
$Q$ als $Q'$ in sich enthält und also von $Q$ verschieden ist.
Folglich ist, da es über $Q$ keinen Normaltheiler von $P$ giebt,
ausser $P$ selbst,

\weq{16}{Q Q' = P,}

und ebenso muss auch

\weq{17}{Q' Q = P}
%CONT>>


% ===== scan leaf 42, printed 21 =====
% !PRINTED 21
\weberpage{21}
%<<CONT
sein. Diese Sätze gelten, wenn nur eine der beiden Gruppen $Q$, $Q'$
grösster Normaltheiler von $P$ ist. Wir nehmen jetzt an, dass sie
beide diese Eigenschaft haben, und verstehen unter $R$ den grössten
gemeinschaftlichen Theiler von $Q$ und $Q'$. Diese Gruppe $R$ ist
dann ein Normaltheiler von $P$ sowohl als von $Q$ und $Q'$. Denn ist
$a$ irgend ein Element in $P$, und $c$ in $Q$ sowohl als in $Q'$
enthalten, so ist auch $a^{-1}ca$ in $Q$ und in $Q'$, also auch in
$R$ enthalten. Also ist $a^{-1}Ra = R$, und $R$ ist Normaltheiler
von $P$ und mithin Normaltheiler von $Q$ und von $Q'$ (\S.\,4, 5.).
Um $Q$ in die Nebengruppen von $R$ zu zerlegen, wählen wir ein
passendes System von Elementen $1, b_1, b_2, b_3 \ldots$ in $Q$, so
dass $R, Rb_1, Rb_2, Rb_3 \ldots$ alle von einander verschieden
werden und setzen
\[
Q = R + Rb_1 + Rb_2 + Rb_3 + \ldots,
\]
was wir auch, wenn wir
\weq{18}{B = 1, b_1, b_2, b_3 \ldots}
setzen, kurz mit
\weq{19}{Q = RB}
bezeichnen können. Nun ist nach (13) $Q'Q = P$, also auch
\weq{20}{P = Q'RB = Q'B}
oder
\[
P = Q' + Q'b_1 + Q'b_2 + Q'b_3 \ldots
\]

Die Nebengruppen
\weq{21}{Q',\ Q'b_1,\ Q'b_2,\ Q'b_3 \ldots}
sind aber alle von einander verschieden. Denn wäre etwa
\[
Q'b_2 = Q'b_1,
\]
so würde folgen, da $Q'$ die Einheit enthält, dass es ein Element
$e$ in $Q'$ giebt, so dass
\[
b_2 = eb_1
\]
wäre. Dies Element $e$ ist aber gleich $b_2b_1^{-1}$, und also, da
die $b$ in $Q$ enthalten sind, gleichfalls in $Q$ und also auch in
$R$ enthalten. Dann aber wäre auch $Rb_2 = Rb_1$, was gegen die
Voraussetzung ist.

Hiernach können wir die Elemente der zu $R$ complementären Gruppe
in Bezug auf $Q$ und der zu $Q'$ complementären Gruppe in Bezug auf
$P$ bilden. Wir erhalten diese beiden Gruppen
\weq{22}{\begin{aligned}
Q/R &= R,\ Rb_1,\ Rb_2,\ Rb_3 \ldots \\
P/Q' &= Q',\ Q'b_1,\ Q'b_2,\ Q'b_3 \ldots
\end{aligned}}


% ===== scan leaf 43, printed 22 =====
% !PRINTED 22
\weberpage{22}
Daraus aber erkennt man leicht, dass diese Gruppen nicht nur von
gleichem Grade sind, sondern dass sie \gesperrt{isomorph} sind. Denn es ist
gleichzeitig
\[
Rb_1\,Rb_2 = Rb_1b_2
\]
und
\[
Q'b_1\,Q'b_2 = Q'b_1b_2.
\]
Ebenso kann man auch schliessen, dass die beiden Gruppen
\[
Q'/R,\quad P/Q
\]
isomorph sind.

Die Gruppen $P/Q$ und $P/Q'$ sind aber, da $Q$, $Q'$ grösste
Normaltheiler sind, nach 3) einfach, und folglich sind auch die
damit isomorphen Gruppen $Q'/R$ und $Q/R$ einfach, und folglich
haben wir den Satz:

\gesperrt{4. Sind $Q$ und $Q'$ zwei verschiedene grösste
Normaltheiler von $P$ und ist $R$ der Durchschnitt von $Q$ und
$Q'$, so ist $R$ grösster Normaltheiler sowohl von $Q$ als von
$Q'$, und der Index von $R$ in Bezug auf $Q$ ist gleich dem Index
von $Q'$ in Bezug auf $P$.}

Damit haben wir die Grundlage gewonnen, um sehr einfach den Satz
von der Constanz der Indexreihe nachzuweisen.

Wir schreiben zur besseren Uebersicht die verschiedenen in
Vergleich zu ziehenden Compositionsreihen so, dass wir den Index
eines jeden Gliedes in Bezug auf das vorangehende unter das Glied
setzen.

Danach mögen irgend zwei gegebene Compositionsreihen von $P$ mit
den zugehörigen Indices folgende sein:
\weq{23}{\begin{aligned}
&P,\ Q,\ Q_1,\ Q_2.\ Q_3 \ldots \\ % [sic: the source prints a full stop after Q_2.]
&\hphantom{P,\ }\nu,\ \nu_1,\ \nu_2,\ \nu_3 \ldots
\end{aligned}}
\weq{24}{\begin{aligned}
&P,\ Q',\ Q_1',\ Q_2',\ Q_3' \ldots \\
&\hphantom{P,\ }\nu',\ \nu_1',\ \nu_2',\ \nu_3' \ldots
\end{aligned}}
Es sei nun $R$ der Durchschnitt von $Q$ und $Q'$, und $\mu$ und
$\mu'$ seien die Indices von $R$ in Bezug auf $Q$ und $Q'$. Wegen
des Isomorphismus der Gruppen $P/Q$ und $Q'/R$ ist dann $\mu' = \nu$,
und aus dem entsprechenden Grunde $\mu = \nu'$. Wir bilden eine
Compositionsreihe von $R$ mit der zugehörigen Indexreihe
\weq{25}{\begin{aligned}
&R,\ R_1,\ R_2,\ R_3 \ldots \\
&\hphantom{R,\ }\mu_1,\ \mu_2,\ \mu_3 \ldots,
\end{aligned}}
%CONT>>


% ===== scan leaf 44, printed 23 =====
% !PRINTED 23
\weberpage{23}
%<<CONT
und da nun $R$ ein grösster Normaltheiler von $Q$ und von $Q'$ ist,
so können wir daraus zwei neue Compositionsreihen von $P$ bilden,
nämlich
\weq{26}{\begin{aligned}
&P,\ Q,\ R,\ R_1,\ R_2,\ R_3 \ldots \\
&\hphantom{P,\ }\nu,\ \nu',\ \mu_1,\ \mu_2,\ \mu_3 \ldots
\end{aligned}}
\weq{27}{\begin{aligned}
&P,\ Q',\ R,\ R_1,\ R_2,\ R_3 \ldots \\
&\hphantom{P,\ }\nu',\ \nu,\ \mu_1,\ \mu_2,\ \mu_3 \ldots
\end{aligned}}
Die beiden Compositionsreihen (26) und (27) von $P$ haben also
dieselbe Indexreihe.

Nehmen wir jetzt an, der zu beweisende Satz sei bereits als richtig
erwiesen für Gruppen, deren Grad $\leqq \frac{1}{2}n$ ist (wenn $n$
den Grad von $P$ bedeutet), so folgt, dass alle Indexreihen von
$Q$, dessen Grad ja ein echter Theiler von $n$ und also höchstens
$= \frac{1}{2}n$ ist, mit einander übereinstimmen, dass also die
Reihen
\[
\nu',\ \mu_1,\ \mu_2,\ \mu_3 \ldots
\]
\[
\nu_1,\ \nu_2,\ \nu_3,\ \nu_4 \ldots,
\]
von der Anordnung abgesehen, übereinstimmen. Ebenso stimmen die
Indexreihen von $Q'$
\[
\nu,\ \mu_1,\ \mu_2,\ \mu_3 \ldots
\]
\[
\nu_1',\ \nu_2',\ \nu_3',\ \nu_4' \ldots
\]
überein. Also stimmen auch die beiden Indexreihen von $P$
\[
\nu,\ \nu_1,\ \nu_2,\ \nu_3,\ \nu_4 \ldots
\]
\[
\nu',\ \nu_1',\ \nu_2',\ \nu_3',\ \nu_4' \ldots
\]
mit einander überein, da die erste mit $\nu, \nu', \mu_1, \mu_2, \mu_3
\ldots$, die zweite mit $\nu', \nu, \mu_1, \mu_2, \mu_3 \ldots$
übereinstimmt.

Für Gruppen vom Grade 2, die nur den einen Normaltheiler 1 haben,
ist aber der Satz evident, und also ist er durch vollständige
Induction allgemein bewiesen.

Wenn in zwei Compositionsreihen von $P$, die wir mit ihren
Indexreihen jetzt so darstellen wollen
\weq{28}{\begin{aligned}
&P,\ P_1,\ P_2,\ P_3 \ldots P_{\mu-1},\ 1 \\
&\hphantom{P,\ }j_1,\ j_2,\ j_3 \ldots j_{\mu-1},\ j_\mu
\end{aligned}}
\weq{29}{\begin{aligned}
&P,\ P_1',\ P_2',\ P_3' \ldots P_{\mu-1}',\ 1 \\
&\hphantom{P,\ }j_1',\ j_2',\ j_3' \ldots j_{\mu-1}',\ j_\mu'
\end{aligned}}
ein gemeinschaftliches Glied $P_r = P_s'$ vorkommt, so gilt der
Satz von der Constanz der Indexreihen auch für die Gruppe $P_r =
P_s'$, und daraus folgt zunächst, dass $s = r$ sein muss, und
%CONT>>


% ===== scan leaf 45, printed 24 =====
% !PRINTED 24
\weberpage{24}
\pdfbookmark[2]{§ 7. Weitere Sätze über die Compositionsreihen}{bm9}
%<<CONT
dass die Indexreihen $j_{r+1}, j_{r+2} \ldots j_\mu$ und $j_{r+1}',
j_{r+2}' \ldots j_\mu'$, von der Ordnung abgesehen, übereinstimmen.
Folglich müssen auch die vorangehenden Theile der Indexreihen
\[
j_1, j_2 \ldots j_r \quad \text{und} \quad j_1', j_2' \ldots j_r'
\]
übereinstimmen.

\begin{center}
\S.\,7.

\gesperrt{Weitere S\"atze \"uber die Compositionsreihen.}
\end{center}

Von den Compositionsreihen gelten noch mannigfaltige Sätze, von
denen wir hier noch die wichtigsten ableiten.

\gesperrt{II. Ist $P_\nu$ irgend ein Normaltheiler von $P$, so giebt
es eine Compositionsreihe von $P$, in der $P_\nu$ vorkommt.}

Ist $P_\nu$ ein grösster Normaltheiler von $P$, so können wir $P_\nu =
P_1$ setzen und die Compositionsreihe von da an beliebig
fortsetzen. Ist aber $P_\nu$ kein grösster Normaltheiler von $P$, so
suchen wir einen grössten Normaltheiler über $P_\nu$, den wir mit
$P_1$ bezeichnen. Dann ist $P_\nu$ ein Normaltheiler von $P_1$. Ist
es ein grösster, so setzen wir $P_\nu = P_2$ und setzen von da die
Compositionsreihe beliebig fort. Ist aber $P_\nu$ noch kein grösster
Normaltheiler von $P_1$, so suchen wir einen grössten
Normaltheiler von $P_1$ über $P_\nu$ und fahren so fort, bis wir
schliesslich zu $P_\nu$ gelangen, was nach einer endlichen Anzahl von
Schritten nothwendig eintreten muss. Wenn wir dann eine
Compositionsreihe von $P_\nu$ anhängen, so haben wir eine
Compositionsreihe von $P$, in der $P_\nu$ vorkommt, wie es der Satz
II. verlangt.

In einer Compositionsreihe einer Gruppe $P$ ist, vermöge der
Definition, jedes Glied ein Normaltheiler jedes vorangehenden. Das
erste Glied $P_1$ und das letzte Glied 1 sind zugleich
Normaltheiler von $P$ selbst. Bei den zwischenliegenden $P_2, P_3
\ldots$ wird dies im Allgemeinen nicht der Fall sein. In der
Compositionsreihe giebt es also gewisse ausgezeichnete Glieder, die
zugleich Normaltheiler von $P$ selbst sind, die eine besondere
Rolle spielen\footnote{Sie bilden die sogenannte Hauptreihe von $P$
(\gesperrt{Netto}).}. Hierüber leiten wir im Folgenden noch einen
wichtigen Satz ab.

Angenommen, es sei $Q$ ein Glied einer Compositionsreihe von $P$,
das zugleich Normaltheiler von $P$ ist, während das auf $Q$
%CONT>>


% ===== scan leaf 46, printed 25 =====
% !PRINTED 25
\weberpage{25}

%<<CONT
folgende Glied $Q_1$ nicht Normaltheiler von $P$ ist. Dann giebt
es unter den nach $P$ conjugirten Gruppen $a^{-1} Q_1 a$ mehrere verschiedene,
die wir mit

\weq{1}{Q_1,\ Q_1',\ Q_1'',\ Q_1''' \cdots}

bezeichnen wollen. \gesperrt{Der grösste gemeinschaftliche Theiler
$R$ aller dieser Gruppen ist wieder Normaltheiler von $P$.}

Nun sind aber alle die Gruppen (1) grösste Normaltheiler
von $Q$, wie sich nach dem Schlusssatze des \S.\,3 daraus ergiebt,
dass $Q, Q_1$ mit $a^{-1} Q a, a^{-1} Q_1 a$ isomorph sind, während $Q$ als
Normaltheiler von $P$ mit $a^{-1} Q a$ identisch ist. Wir können also
in der Compositionsreihe auf $Q$ jede der Gruppen (1) folgen
lassen, die alle denselben Index $\nu$ haben.

Ist nun $Q_2$ der grösste gemeinschaftliche Theiler von $Q_1, Q_1'$
(also von $Q_1$ verschieden), so ist $Q_2$ nach dem Satze \S.\,6, 4. ein
grösster Normaltheiler von $Q_1$ und $Q_1'$, dessen Index gleichfalls
$\nu$ ist. Also können wir die Compositionsreihe von $Q$ an auf die
folgenden beiden Arten fortsetzen:

\weq{2}{
\begin{array}{ccc@{\qquad}ccc}
Q, & Q_1, & Q_2, & Q & Q_1' & Q_2 \\
& \nu & \nu & & \nu & \nu
\end{array}
}

Das hierin ausgedrückte Gesetz wollen wir nun verallgemeinern
und durch vollständige Induction beweisen.

Wenn $Q_2$ noch nicht gleich $R$ ist, so fügen wir zu $Q_1, Q_1'$
eine dritte Gruppe $Q_1''$ aus der Reihe (1) hinzu, so dass der
Durchschnitt $Q_3$ von $Q_1, Q_1', Q_1''$ ein echter Theiler von $Q_2$ wird,
und fahren so fort, bis wir zum grössten gemeinschaftlichen
Theiler $R$ aller Gruppen (1) gelangt sind.

Es sei also:

\weq{3}{
\begin{array}{l}
Q_2 \text{ der Durchschnitt von } Q_1,\ Q_1' \\
Q_3 \text{ \ ,,\qquad ,,\qquad ,,\ } Q_1,\ Q_1',\ Q_1'' \\
\cdot\ \cdot\ \cdot\ \cdot\ \cdot\ \cdot\ \cdot\ \cdot\ \cdot\ \cdot\ \cdot\ \cdot\ \cdot \\
Q_r \text{ \ ,,\qquad ,,\qquad ,,\ } Q_1,\ Q_1',\ Q_1'' \cdots Q_1^{(r-1)}
\end{array}
}

und die Anordnung der Gruppen $Q_1, Q_1' \cdots Q_1^{(r-1)}$ sei so gewählt,
dass von den Gruppen $Q_1, Q_2 \cdots Q_r$ jede ein echter
Theiler der vorangehenden ist. Diese Anordnung lässt sich fortsetzen,
so lange $Q_r$ nicht gleich $R$ ist.

Wir wollen beweisen, dass wir die Compositionsreihe von $P$
so wählen können, dass darin

\weq{4}{
\begin{array}{cccc}
Q, & Q_1, & Q_2 & \cdots\ Q_r\\
& \nu' & \nu & \nu % [sic: the source prints \nu' although the text says all indices are \nu.]
\end{array}
}
%CONT>>


% ===== scan leaf 47, printed 26 =====
% !PRINTED 26
\weberpage{26}

%<<CONT
vorkommt, und dass die Indices in diesem Theile der Reihe alle
gleich $\nu$ sind.

Die Behauptung ist richtig für $r = 2$, und wir beweisen sie
allgemein durch vollständige Induction, unter der Voraussetzung,
dass sie für $r$ schon als richtig erwiesen sei.

Ist $Q_r$ nicht gleich $R$, so wählen wir $Q_1^{(r)}$ so, dass $Q_{r+1}$ der
Durchschnitt von $Q_1, Q_1' \cdots Q_1^{(r-1)}, Q_1^{(r)}$ ist, und bezeichnen überdies
 den Durchschnitt von $Q_1, Q_1' \cdots Q_1^{(r-2)}, Q_1^{(r)}$ mit $Q_r'$. Dann
ist $Q_r'$ von $Q_r$ verschieden und $Q_{r+1}$ ist der Durchschnitt von
$Q_r$ und $Q_r'$. Nun haben wir nach der Voraussetzung (4) zwei
entsprechende Stücke der Compositionsreihe mit der constanten
Indexreihe $\nu$

\weq{5}{
\begin{array}{l}
Q,\ Q_1 \cdots Q_{r-1}\ Q_r,\\
Q,\ Q_1 \cdots Q_{r-1}\ Q_r';
\end{array}
}

denn die zweite dieser Reihen ist genau wie die erste nach der
Vorschrift (3) gebildet, nur dass $Q_1^{(r)}$ an Stelle von $Q_1^{(r-1)}$ getreten
ist. Nach dem Satze \S.\,6, 2. ist daher $Q_{r+1}$ grösster Normaltheiler
 sowohl von $Q_r$ als von $Q_r'$, und zwar vom Index $\nu$, und
es folgt also, was bewiesen werden sollte, dass die Compositionsreihe
 so fortgesetzt werden kann:

\weq{6}{Q,\ Q_1 \cdots Q_{r-1},\ Q_r,\ Q_{r+1},}

und dass der neu hinzugekommene Index gleichfalls $\nu$ ist.

Dies können wir nun fortsetzen, bis wir zur Gruppe $R$ gelangt
 sind, und erhalten ein Stück der Compositionsreihe

\weq{7}{
\begin{array}{ccc}
Q,\ Q_1,\ Q_2 & \cdots & R,\\
\nu\quad\ \nu & & \nu
\end{array}
}

deren letztes Glied $R$ wieder Normaltheiler von $P$ ist und in
der alle Indices übereinstimmen.

Setzen wir von $R$ an die Compositionsreihe fort, so kann
sich der Index ändern; er kann aber von da an wieder gleich
gehalten werden, bis man abermals zu einem Normaltheiler von
$P$ gelangt.

Damit ist der folgende Satz bewiesen:

III.\ \gesperrt{Man kann die Compositionsreihe von $P$ mit ihrer
Indexreihe}

\weq{8}{
\begin{array}{c}
P,\ P_1,\ P_2,\ P_3 \cdots\\
j_1,\ j_2,\ j_3 \cdots
\end{array}
}

\gesperrt{so anordnen, dass, so oft eine Aenderung in der
Indexreihe eintritt, also $j_{r+1}$ von $j_r$ verschieden
ist, $P_r$ ein Normaltheiler von $P$ ist.}


% ===== scan leaf 48, printed 27 =====
% !PRINTED 27
\weberpage{27}
\pdfbookmark[2]{§ 8. Metacyklische Gruppen}{bm10}

\begin{center}
\S.\,8.

\gesperrt{Metacyklische Gruppen.}
\end{center}

Wir haben in \S.\,176 des ersten Bandes bei den Permutationsgruppen
 metacyklische Gruppen definirt und ihre Bedeutung
für die Auflösung der Gleichungen kennen gelernt. Der Begriff
ist aber von der besonderen Bedeutung der Gruppenelemente
unabhängig, und wir definiren daher jetzt allgemein:

\gesperrt{Eine Gruppe, deren Indexreihe aus lauter Primzahlen
 besteht, soll eine metacyklische Gruppe heissen.}

Für diese Gruppen gilt ein für die Anwendungen auf die
Algebra wichtiger Satz, der aber von diesen Anwendungen unabhängig
 ist und den wir daher hier noch ableiten. Der Satz
lautet:

IV.\ \gesperrt{Eine metacyklische Gruppe hat einen von der
Einheitsgruppe verschiedenen commutativen
Normaltheiler}\footnote{C. J o r d a n, Traité des substitutions, Livre IV.}.

Ist $P$ eine metacyklische Gruppe und

\weq{1}{
\begin{array}{c}
P,\ P_1,\ P_2 \cdots P_{\mu-1},\ 1\\
j_1,\ j_2 \cdots j_{\mu-1},\ j_\mu
\end{array}
}

die Compositionsreihe, deren Indexreihe $j_1, j_2 \cdots$ aus lauter
Primzahlen besteht, so ist $P_{\mu-1}$ jedenfalls eine commutative
Gruppe; denn wenn $\pi$ irgend ein von 1 verschiedenes Element
aus $P_{\mu-1}$ ist, so sind die Potenzen

\weq{2}{1,\ \pi,\ \pi^2 \cdots \pi^{j_\mu - 1}}

(weil $j_\mu$ eine Primzahl ist) alle von einander verschieden und
machen also die ganze Gruppe $P_{\mu-1}$ aus; die cyklische Gruppe (2)
ist aber gewiss commutativ, weil für je zwei Exponenten $h, k$
immer

\[
\pi^h \pi^k = \pi^k \pi^h = \pi^{h+k}
\]

ist. Zugleich ist, nach der Definition der Compositionsreihe. % [sic: the full stop is as printed; a comma is intended.]
$P_{\mu-1}$ Normaltheiler von $P_{\mu-2}$.

Wir nehmen demnach jetzt an, wir haben einen von der
Einheit verschiedenen commutativen Normaltheiler $Q_\nu$ irgend
einer Gruppe $P_\nu$ der Reihe (1) und setzen ausserdem noch voraus,
%CONT>>


% ===== scan leaf 49, printed 28 =====
% !PRINTED 28
\weberpage{28}

%<<CONT
wenn es zu $P_\nu$ mehrere Theiler von der Beschaffenheit wie $Q_\nu$
giebt, dass einer von möglichst niedrigem Grade für $Q_\nu$ genommen
 ist.

Wenn wir dann nachweisen, wie man aus $Q_\nu$ einen commutativen
 Normaltheiler $Q_{\nu-1}$ von $P_{\nu-1}$ ableiten kann, der zugleich
 $Q_\nu$ enthält, so können wir dies Verfahren fortsetzen, bis
wir zu einem commutativen Normaltheiler $Q$ von $P$ selbst gelangt
 sind, wie ihn unser Satz verlangt.

Wählen wir irgend ein Element $\gamma$ in der Gruppe $P_{\nu-1}$,
welches nicht in $P_\nu$ enthalten ist, so ist, da $P_\nu$ ein Normaltheiler
von $P_{\nu-1}$ ist,

\weq{3}{\gamma P_\nu = P_\nu \gamma,}

und wenn daher $h$ der niedrigste Exponent ist, für den $\gamma^h$ in
$P_\nu$ enthalten ist, so ist $h$ grösser als 1 und

\weq{4}{\gamma^h P_\nu = P_\nu,}

und die in dem System der Nebengruppen

\weq{5}{P_\nu,\ \gamma P_\nu,\ \gamma^2 P_\nu \ldots \gamma^{h-1} P_\nu}

enthaltenen Elemente, die alle von einander verschieden sind,
bilden eine Gruppe, deren Grad, wenn der Grad von $P_\nu$ mit $p_\nu$
bezeichnet wird, gleich $h p_\nu$ ist. Die Gruppe (5) ist aber auch
ein Theiler von $P_{\nu-1}$ und folglich muss $h p_\nu$ ein Theiler von
$p_{\nu-1} = j_\nu p_\nu$ sein, und $h$ ist ein Theiler von $j_\nu$. Da aber $j_\nu$
eine Primzahl ist, so muss $h = j_\nu$ sein, und durch (5) ist die
ganze Gruppe $P_{\nu-1}$ erschöpft:

\weq{6}{P_{\nu-1} = P_\nu + \gamma P_\nu + \gamma^2 P_\nu + \cdots \gamma^{j_\nu - 1} P_\nu.}

Wir betrachten nun das System der transformirten Gruppen

\weq{7}{\gamma^{-1} Q_\nu \gamma,}

worin $\gamma$ alle Elemente von $P_{\nu-1}$ durchläuft.

Diese Gruppen sind alle mit einander isomorph und sind
also ebenso wie $Q_\nu$ commutativ. Sie sind alle in $P_\nu$ enthalten,
und sind, da $\gamma^{-1} P_\nu \gamma = P_\nu$ ist, nach dem Schlusssatze von \S.\,3
Normaltheiler von $P_\nu$.

Wenn alle Gruppen (7) mit einander identisch sind, so ist
$Q_\nu$ Normaltheiler von $P_{\nu-1}$, und wir erreichen unseren Zweck
der Bestimmung von $Q_{\nu-1}$ schon dadurch, dass wir $Q_\nu$ selbst,
oder, falls noch ein commutativer Normaltheiler niedrigeren
Grades von $P_{\nu-1}$ vorhanden sein sollte, diesen letzteren, für $Q_{\nu-1}$
nehmen.


% ===== scan leaf 50, printed 29 =====
% !PRINTED 29
\weberpage{29}
Im anderen Falle wollen wir die von einander verschiedenen der Gruppen (7) mit
\weq{8}{Q_\nu, \; Q'_\nu, \; Q''_\nu \; \ldots}
bezeichnen, und daraus eine Gruppe
\weq{9}{R = (Q_\nu, \; Q'_\nu, \; Q''_\nu \; \ldots)}
ableiten, die dadurch definirt sein soll, dass es die kleinste Gruppe ist, die jede der Gruppen (8) als Theiler enthält, und die wir als das \gesperrt{kleinste gemeinschaftliche Vielfache} der Gruppen $Q_\nu, Q'_\nu, Q''_\nu \ldots$ bezeichnen können.

Dass es eine und nur eine solche Gruppe $R$ giebt, und dass jede Gruppe, die durch $Q_\nu, Q'_\nu, Q''_\nu \ldots$ theilbar ist, auch durch $R$ theilbar sein muss, ist leicht einzusehen. Denn erstens giebt es endliche Gruppen, z. B. die Gruppe $P_\nu$, die alle Gruppen (8) als Theiler enthalten, und zweitens, wenn zwei Gruppen $R, R'$ alle diese Gruppen enthalten, so gilt dasselbe auch von dem Durchschnitt von $R$ und $R'$. Ist daher $R$ von möglichst niedrigem Grade, so muss dieser Durchschnitt mit $R$ identisch sein, und $R$ ist also ein \gesperrt{Theiler von} $R'$. Wenn $R'$ auch von möglichst niedrigem Grade ist, so ist $R'$ von $R$ nicht verschieden.

Von dieser Gruppe $R$ weisen wir nun nach:

a) dass sie ein Normaltheiler von $P_{\nu-1}$ ist,

b) dass sie eine commutative Gruppe ist.

Wenn wir dies Beides nachgewiesen haben, so sind wir am Ziele; denn dann können wir, wenn es nicht noch einen commutativen Normaltheiler \gesperrt{niedrigeren Grades} von $P_{\nu-1}$ giebt, $R$ selbst für $Q_{\nu-1}$ wählen.

Das Erste ist aber ohne Weiteres klar. Denn bedeutet $\gamma$ irgend ein Element aus $P_{\nu-1}$, so ist
\[ \gamma^{-1} R \gamma = (\gamma^{-1} Q_\nu \gamma, \; \gamma^{-1} Q'_\nu \gamma, \; \gamma^{-1} Q''_\nu \gamma \; \ldots) \]
das kleinste gemeinschaftliche Vielfache der Gruppen
\[ \gamma^{-1} Q_\nu \gamma, \; \gamma^{-1} Q'_\nu \gamma, \; \gamma^{-1} Q''_\nu \gamma \; \ldots \]
Diese Gruppen aber sind nach der Definition (7) in ihrer Gesammtheit von den Gruppen
\[ Q_\nu, \; Q'_\nu, \; Q''_\nu \; \ldots \]
nicht verschieden, und folglich ist
\[ \gamma^{-1} R \gamma = R, \]
d. h. $R$ ist Normaltheiler von $P_{\nu-1}$.


% ===== scan leaf 51, printed 30 =====
% !PRINTED 30
\weberpage{30}
Um aber nachzuweisen, dass $R$ eine commutative Gruppe ist, müssen wir zwei Hülfsbetrachtungen voranschicken.

1) Es ist zu beweisen, dass irgend zwei der Gruppen $Q_\nu, Q'_\nu, Q''_\nu \ldots$ ausser der Einheit keinen gemeinschaftlichen Theiler haben.

Wir brauchen hierzu nur nachzuweisen, dass $Q_\nu, Q'_\nu$ diese Eigenschaft haben. Denn haben
\[ Q'_\nu = \gamma'^{-1} Q_\nu \gamma', \quad Q''_\nu = \gamma''^{-1} Q_\nu \gamma'' \]
den grössten gemeinschaftlichen Theiler $T$, so haben $Q_\nu$ und $Q'''_\nu = \gamma' Q''_\nu \gamma'^{-1}$ den grössten gemeinschaftlichen Theiler $\gamma' T \gamma'^{-1}$.

Nehmen wir also an, es sei $T$ der grösste gemeinschaftliche Theiler von $Q_\nu$ und $Q'_\nu$, so ist $T$ gewiss eine commutative Gruppe, weil $Q_\nu, Q'_\nu$ solche sind. Es ist aber $T$ auch Normaltheiler von $P_\nu$. Denn $Q_\nu$ und $Q'_\nu$ sind solche Theiler; und wenn also $\tau$ ein Element aus $T$ und $\pi$ ein Element aus $P_\nu$ ist, so ist $\pi^{-1}\tau\pi$ sowohl in $Q_\nu$ als in $Q'_\nu$, also auch in $T$ enthalten. Folglich ist $\pi^{-1} T \pi = T$.

Nun haben wir aber angenommen, dass $Q_\nu$ ein commutativer Normaltheiler von $P_\nu$ über der Einheit von möglichst niedrigem Grade sei; ferner waren $Q_\nu$ und $Q'_\nu$ von einander verschieden. Also kann $T$ nur die Einheitsgruppe selbst sein, w. z. b. w.

2) Um die Gruppe $R$ zu bilden, können wir so verfahren, dass wir die Elemente von $Q_\nu, Q'_\nu, Q''_\nu \ldots$ in beliebiger Reihenfolge und Wiederholung so lange mit einander componiren, als noch neue Elemente entstehen. Denn alle so gebildeten Zusammensetzungen bilden eine Gruppe $C$, die in $R$ enthalten ist, weil $R$ die einzelnen Elemente von $Q_\nu, Q'_\nu, Q''_\nu \ldots$ enthält. Andererseits sind auch die Gruppen $Q_\nu, Q'_\nu, Q''_\nu \ldots$ in $C$ enthalten, und folglich ist $R$ in $C$ enthalten und $C$ ist mit $R$ identisch.

Um also endlich zu beweisen, dass $R$ eine commutative Gruppe ist, bleibt uns nur zu zeigen, dass, wenn $\alpha, \beta$ irgend zwei Elemente aus $Q_\nu, Q'_\nu, Q''_\nu \ldots$ sind, die Vertauschbarkeit
\weq{10}{\alpha\beta = \beta\alpha}
besteht.

Wenn nun $\alpha, \beta$ beide in dieselbe Gruppe $Q_\nu$ gehören, so ist (10) Folge unserer Annahme, dass $Q_\nu$ commutativ sei; und ebenso ist es, wenn beide in einer der anderen Gruppen $Q'_\nu, Q''_\nu \ldots$ vorkommen.


% ===== scan leaf 52, printed 31 =====
% !PRINTED 31
\weberpage{31}
Sind aber $\alpha$ und $\beta$ in zwei verschiedenen Gruppen, etwa in $Q_\nu$ und $Q'_\nu$ enthalten, so schliessen wir so: Weil $\beta$ in $P_\nu$ enthalten und $Q_\nu$ ein Normaltheiler von $P_\nu$ ist, so ist auch
\weq{11}{\beta^{-1}\alpha\beta = \alpha'}
in $Q_\nu$ enthalten. Daraus folgt:
\weq{12}{\beta\alpha'\alpha^{-1} = \alpha\beta\alpha^{-1} = \beta',}
und weil $\alpha$ in $P_\nu$ enthalten und $Q'_\nu$ ein Normaltheiler von $P_\nu$ ist, so ist auch $\beta'$ in $Q'_\nu$ enthalten. Aus (12) aber folgt
\[ \alpha'\alpha^{-1} = \beta^{-1}\beta' = \delta, \]
und demnach ist das Element $\delta$ sowohl in $Q_\nu$ als in $Q'_\nu$ enthalten. Es kann also nach dem, was unter 1) bewiesen ist, nur das Einheitselement sein, d. h. es ist
\[ \alpha' = \alpha, \quad \beta' = \beta, \]
und demnach folgt aus (11), dass $\alpha\beta = \beta\alpha$ sein muss, wie bewiesen werden sollte.

\begin{center}
\rule{3cm}{0.4pt}
\end{center}
% [The source prints this plain horizontal rule below the paragraph, marking the close of \S.\,8; the remainder of the page is blank.]


% ===== scan leaf 53, printed 32 =====
\weberpage{32}% page number not printed on this page
\pdfbookmark[1]{Zweiter Abschnitt. Abel'sche Gruppen.}{bm11}
\pdfbookmark[2]{§ 9. Darstellung Abel'scher Gruppen durch eine Basis}{bm12}
% !PRINTED none
% [part-title / section-opening page: no running-head page number visible in any of the three bands; centred "Zweiter Abschnitt." / "Abel'sche Gruppen." title with rule, then \S.\,9 heading, then body text begins]
\begin{center}
\gesperrt{Zweiter Abschnitt.}

\gesperrt{Abel'sche Gruppen.}

\rule{0.3\linewidth}{0.4pt}

\S.\,9.

\gesperrt{Darstellung Abel'scher Gruppen durch eine Basis.}
\end{center}

Wir haben schon in \S.\,1 solche Gruppen, in denen bei der Composition das commutative Gesetz gilt, commutative oder \gesperrt{Abel'sche} Gruppen genannt. Bei diesen Gruppen, die in den Anwendungen von besonderer Wichtigkeit sind, herrschen viel einfachere Gesetze, als in den allgemeinen Gruppen, wie wir jetzt sehen werden. Es gelten für die Zusammensetzung der Elemente in diesen Gruppen dieselben Regeln, wie bei der Multiplication von Zahlen, und wir nennen demnach die Composita von Elementen einer solchen Gruppe, wie bei der Multiplication, Producte und Potenzen.

Auch hier beschränken wir uns auf die Betrachtung endlicher Gruppen.

Aus der Definition der Abel'schen Gruppen folgt, dass das commutative Gesetz auch bei der Composition der Theile einer solchen Gruppe gilt (\S.\,4). Jeder Theiler einer Abel'schen Gruppe ist selbst eine Abel'sche Gruppe, und ist zugleich Normaltheiler, so dass wir bei diesen Gruppen den Zusatz ,,Normal`` weglassen können.

Eine Gruppe, die nur aus den Wiederholungen eines Elementes besteht, wie
\[ 1, \; A, \; A^2 \; \ldots \; A^{a-1}, \]
haben wir schon früher (Bd. I, \S.\,148) eine cyklische Gruppe genannt, und wir behalten diese Bezeichnung hier bei. Wenn $a$ die kleinste positive Zahl ist, für die $A^a = 1$ ist, also $a$ der Grad der Gruppe, so heisst $a$ auch der Grad des Elementes $A$. Ist $m$ irgend ein Exponent, für den $A^m = 1$ ist, so ist $m$ ein
%CONT>>


% ===== scan leaf 54, printed 33 =====
% !PRINTED 33
\weberpage{33}
%<<CONT
Vielfaches von $a$. Das Element 1 hat den Grad 1. Alle cyklischen Gruppen sind commutativ.

Es gilt nun für jede Abel'sche Gruppe $S$ der Fundamentalsatz, dass sich immer ein System von Elementen $A, B, C \ldots$ von den Graden $a, b, c \ldots$ so auswählen lässt, dass sich in der Form
\[
\Theta = A^\alpha B^\beta C^\gamma \ldots
\]
jedes Element $\Theta$ von $S$, und jedes nur einmal, darstellen lässt, wenn $\alpha$ ein volles Restsystem nach dem Modul $a$, $\beta$ ein volles Restsystem nach dem Modul $b$, $\gamma$ ein volles Restsystem nach dem Modul $c$ etc. durchläuft.

Ein System von Elementen, das zu einer solchen Darstellung geeignet ist, heisst eine \gesperrt{Basis der Gruppe}, und wir können den zu beweisenden Satz demnach so aussprechen:

\gesperrt{I. Jede Abel'sche Gruppe lässt sich durch eine Basis darstellen.}

Der Beweis des Satzes gründet sich auf folgende Reihe von Schlüssen:

\gesperrt{1. Ist $n$ der Grad der Gruppe $S$, sind $A_1, A_2 \ldots A_n$ ihre Elemente und $a_1, a_2 \ldots a_n$ deren Grade; durchlaufen ferner $\alpha_1, \alpha_2 \ldots \alpha_n$ volle Restsysteme nach den Moduln $a_1, a_2 \ldots a_n$, so wird in der Form}
\weq{1}{\Theta = A_1^{\alpha_1} A_2^{\alpha_2} \ldots A_n^{\alpha_n}}
\gesperrt{jedes Element von $S$ und jedes gleich oft dargestellt.}

Dass man jedes Element in der Form (1) überhaupt erhält, sieht man unmittelbar; denn man erhält z.\,B. $A_1$, wenn man $\alpha_1 = 1, \alpha_2 = 0 \ldots \alpha_n = 0$ setzt.

Es ist noch zu beweisen, dass man jedes Element gleich oft erhält.

Wenn aber
\weq{2}{1 = A_1^{h_1} A_2^{h_2} \ldots A_n^{h_n}}
eine Darstellung des Elementes 1 ist, so ändert sich $\Theta$ nicht, wenn man $\alpha_1, \alpha_2 \ldots \alpha_n$ in (1) durch $\alpha_1 + h_1, \alpha_2 + h_2 \ldots \alpha_n + h_n$ ersetzt. Es wird also jedes Element $\Theta$ durch (1) mindestens so oft dargestellt, als das Element 1 durch (2).

Geben andererseits die beiden Exponentenreihen $\alpha_1, \alpha_2 \ldots \alpha_n$ und $\alpha_1', \alpha_2' \ldots \alpha_n'$ zwei Darstellungen des Elementes $\Theta$, so hat man
%CONT>>


% ===== scan leaf 55, printed 34 =====
% !PRINTED 34
\weberpage{34}
%<<CONT
\weq{3}{1 = A_1^{\alpha_1' - \alpha_1} A_2^{\alpha_2' - \alpha_2} \ldots A_n^{\alpha_n' - \alpha_n},}
also eine Darstellung des Elementes 1. Also können wir nach (2) setzen:
\[
\alpha_1' = \alpha_1 + h_1,\ \alpha_2' = \alpha_2 + h_2 \ldots \alpha_n' = \alpha_n + h_n.
\]
Es kann folglich auch nicht mehr verschiedene Darstellungen des Elementes $\Theta$ geben, als die Anzahl der Darstellungen des Elementes 1 beträgt. Wir bezeichnen die Anzahl dieser Darstellungen mit $h$. Im Ganzen ist aber die Anzahl der verschiedenen möglichen Bestimmungen der $\alpha$ in (1) gleich dem Producte $a_1 a_2 \ldots a_n$, und da $n$ die Anzahl der Elemente von $S$ ist, so folgt
\weq{4}{n h = a_1 a_2 \ldots a_n.}
Aus der Formel (4) ergiebt sich der Satz:

\gesperrt{2. Wenn $r$ eine im Grade $n$ der Gruppe aufgehende Primzahl ist, so giebt es in $S$ ein Element vom Grade $r$.}

Denn aus (4) sieht man zunächst, dass einer der Grade $a_1, a_2, \ldots a_n$ durch $r$ theilbar ist. Ist also etwa $a_1 = rk$, so ist $A_1^k$ ein Element vom Grade $r$.

\gesperrt{3. Sind $A, B \ldots$ irgend welche Elemente in $S$ von den Graden $a, b \ldots$, so ist auch $AB \ldots$ ein Element in $S$, und der Grad dieses Productes ist ein Theiler des kleinsten gemeinschaftlichen Vielfachen $m$ von $a, b \ldots$}

Denn es ist
\[
(AB \ldots)^m = A^m B^m \ldots = 1,
\]
und folglich $m$ ein Vielfaches des Grades von $AB \ldots$

\gesperrt{4. Ist der Grad $n$ der Gruppe $S$ in zwei Factoren $n = ab$ zerlegt, so dass $a$ und $b$ relativ prim sind, so giebt es in $S$ genau $a$ Elemente $A$, deren Grad ein Theiler von $a$ ist, und $b$ Elemente $B$, deren Grad ein Theiler von $b$ ist, und in der Form}
\weq{5}{\Theta = AB}
\gesperrt{sind sämmtliche Elemente von $S$ und jedes nur einmal enthalten.}

Um die Richtigkeit dieses Satzes einzusehen, stellen wir folgende Ueberlegung an. Der Inbegriff $\mathfrak{A}$ der Elemente $A$, deren
%CONT>>


% ===== scan leaf 56, printed 35 =====
% !PRINTED 35
\weberpage{35}
%<<CONT
Grad ein Theiler von $a$ ist, ist eine in $S$ enthaltene Gruppe; denn sind $A, A'$ zwei solche Elemente, so ist nach 3. der Grad von $AA'$ ein Theiler von $a$, und $AA'$ ist also auch in $\mathfrak{A}$ enthalten.

Ebenso ist der Inbegriff $\mathfrak{B}$ der Elemente $B$, deren Grad ein Theiler von $b$ ist, eine Gruppe. Das Element 1 kommt sowohl in $\mathfrak{A}$ als in $\mathfrak{B}$ vor, sonst enthalten beide kein gemeinschaftliches Element.

Der Grad $a'$ von $\mathfrak{A}$ ist relativ prim zu $b$. Denn ist $r$ eine in $a'$ aufgehende Primzahl, so giebt es nach 2. in $\mathfrak{A}$ ein Element vom Grade $r$. Da aber der Grad jedes Elementes von $\mathfrak{A}$ ein Theiler von $a$ ist, so ist auch $r$ ein Theiler von $a$ und nicht von $b$.

Ebenso beweist man, dass der Grad $b'$ von $\mathfrak{B}$ relativ prim zu $a$ ist.

Nun bestimme man (nach Bd.\,I, \S.\,118) zwei ganze Zahlen $x$ und $y$, so dass
\weq{6}{a x + b y = 1}
ist, und nehme irgend ein Element $\Theta$ von $S$. Dann ist
\weq{7}{\Theta = \Theta^{ax} \Theta^{by},}
und nun ist, da $(\Theta^{by})^a = 1$ ist, $\Theta^{by}$ in $\mathfrak{A}$ und aus dem gleichen Grunde $\Theta^{ax}$ in $\mathfrak{B}$ enthalten. Also folgt aus (7):
\weq{8}{\Theta = AB.}
Demnach ist jedes Element $\Theta$ in der Form $AB$ enthalten. Eine solche Darstellung ist aber auch nur auf eine Art möglich. Denn sind $A', B'$ zwei Elemente aus $\mathfrak{A}$ und $\mathfrak{B}$, und ist
\[
AB = A'B',
\]
so folgt, wenn man in die Potenz $by = 1 - ax$ erhebt, $A = A'$ und folglich auch $B = B'$. Die Anzahl der verschiedenen Producte der Form $AB$ ist aber $= a'b'$, und daher hat man
\[
n = ab = a'b',
\]
und da $a$ relativ prim zu $b'$ und $b$ relativ prim zu $a'$ ist:
\[
a' = a,\quad b' = b.
\]
Damit ist der Satz 4. in allen seinen Theilen bewiesen.

Wenn nun die beiden Gruppen $\mathfrak{A}, \mathfrak{B}$ durch Basen dargestellt sind, so folgt aus der Formel (8), dass auch $S$ durch eine Basis dargestellt ist, und die Basis von $S$ enthält die Elemente der Basen von $\mathfrak{A}$ und $\mathfrak{B}$ und keine anderen.


% ===== scan leaf 57, printed 36 =====
% !PRINTED 36
\weberpage{36}
Wenn $a$ und $b$ weiter in Factoren zerlegbar sind, die zu einander relativ prim sind, so können wir mit den Gruppen $\mathfrak{A}$ und $\mathfrak{B}$ wieder ebenso verfahren, und wir kommen also zu dem Resultate, dass unser Theorem I. allgemein bewiesen ist, wenn wir es noch für Gruppen nachweisen können, deren Grad eine Potenz einer Primzahl ist.

Es sei also jetzt der Grad $n$ der Gruppe $S$ eine Potenz einer Primzahl $p$
\weq{9}{n = p^k.}
Die Grade aller Elemente von $S$, die ja Divisoren von $n$ sind, müssen dann gleichfalls Potenzen von $p$ sein. Es ist nachzuweisen, dass eine solche Gruppe durch eine Basis darstellbar ist.

Man wähle in $S$ ein Element $A$ von möglichst hohem Grade $a$. Dann ist $a$ eine Potenz von $p$ und die Grade aller anderen Elemente sind Theiler von $a$, so dass für jedes Element $\Theta$ von $S$
\weq{10}{\Theta^a = 1}
ist. Die Elemente
\weq{11}{1, A, A^2 \ldots A^{a-1}}
sind alle von einander verschieden, und ihre Gesammtheit ist ein Theiler von $S$. Ist damit die Gruppe $S$ erschöpft, ist also jedes Element von der Form $A^\alpha$, so ist $S$ durch eine eingliedrige Basis dargestellt. Wenn aber $S$ mit (11) noch nicht erschöpft ist, so wird es doch für jedes Element $\Theta$ von $S$ einen gewissen Exponenten $h$ geben, so dass $\Theta^h$ in der Reihe (11) enthalten ist. Gewiss wird das eintreten, wenn $h$ der Grad von $\Theta$, also $\Theta^h = 1$ ist. Unter diesen Zahlen $h$ wird eine die kleinste positive sein, die wir mit $b$ bezeichnen wollen. Es giebt also für jedes Element $\Theta$ eine gewisse kleinste positive Zahl $b$, so dass
\weq{12}{\Theta^b = A^\lambda}
in der Reihe (11) enthalten ist.

Diese Zahl $b$ ist ein Theiler von $a$, also auch eine Potenz von $p$ und zugleich ein Theiler von $\lambda$. Denn setzen wir $a = qb + b'$, wo $0 \leqq b' < b$ ist, so folgt aus (10) und (12)
\[
\Theta^a = A^{\lambda q} \Theta^{b'} = 1,
\]
oder $\Theta^{b'} = A^{-\lambda q}$, und wenn also $b'$ nicht Null ist, so giebt es gegen die Voraussetzung eine noch kleinere Zahl als $b$, nämlich $b'$, die der Forderung (12) genügt.


% ===== scan leaf 58, printed 37 =====
% !PRINTED 37
\weberpage{37}
Aus (12) folgt ferner
\[
1 = \Theta^a = A^{\frac{\lambda a}{b}};
\]
also muss $\lambda a : b$ ein Vielfaches von $a$ und folglich $\lambda$ ein Vielfaches von $b$ sein. Wenn wir daher
\weq{13}{B = \Theta A^{-\frac{\lambda}{b}}}
setzen, so ist $B^b = 1$, und zugleich ist $B^b$ die niedrigste Potenz von $B$, die einer Potenz von $A$ gleich wird, weil, wenn $B^{b'}$ eine Potenz von $A$ ist, dasselbe nach (13) auch von $\Theta^{b'}$ gilt. Wir nehmen das Element $\Theta$ so gewählt an, dass $b$ so gross als möglich wird. Dann ist auch für jedes andere Element $\Theta_1$ aus $S$ immer $\Theta_1^b$ eine Potenz von $A$, deren Exponent durch $b$ theilbar ist.

Ist nun
\[
\begin{gathered}
\alpha = 0, 1, 2 \ldots a-1 \\
\beta = 0, 1, 2 \ldots b-1,
\end{gathered}
\]
so sind die Elemente
\weq{14}{A^\alpha B^\beta}
in der Anzahl $ab$ alle von einander verschieden. Sie bilden einen in $S$ enthaltenen Theiler $S'$, der durch die Basis $A, B$ dargestellt ist. Ist $S'$ mit $S$ identisch, so sind wir am Ziele.

Ist aber $S$ durch (14) noch nicht erschöpft, so fahren wir in derselben Weise fort.

Wir wollen durch Anwendung der vollständigen Induction gleich allgemein schliessen.

Es sei ein Theiler $S_{\nu-1}$ von $S$ ermittelt, der durch eine Basis in der Form dargestellt ist
\weq{15}{S_{\nu-1} = A_1^{\alpha_1} A_2^{\alpha_2} \ldots A_{\nu-1}^{\alpha_{\nu-1}},}
von dem wir folgende Voraussetzungen machen:

1) Die Grade von $A_1, A_2 \ldots A_{\nu-1}$ seien $a_1, a_2 \ldots a_{\nu-1}$, und es sei
\[
a_1 \geqq a_2 \geqq \ldots \geqq a_{\nu-1}.
\]

2) Für jedes in $S$ enthaltene Element $\Theta$ sei
\weq{16}{\Theta^{a_{\nu-1}} = A_1^{\lambda_1} A_2^{\lambda_2} \ldots A_{\nu-2}^{\lambda_{\nu-2}},}
d. h. in $S_{\nu-2}$ enthalten, und die Exponenten $\lambda_1, \lambda_2 \ldots \lambda_{\nu-2}$ seien durch $a_{\nu-1}$ theilbar.

Die oben abgeleitete Gruppe $S'$ genügt diesen Forderungen, wenn $\nu = 3$ und $a_{\nu-1} = b$ gesetzt wird, und es ist nun zu zeigen,
%CONT>>


% ===== scan leaf 59, printed 38 =====
% !PRINTED 38
\weberpage{38}
%<<CONT
wie man, wenn $S_{\nu-1}$ noch nicht die ganze Gruppe $S$ erschöpft, daraus eine eben solche umfassendere Gruppe $S_\nu$ ableiten kann.

Für jedes Element $\Theta$ von $S$ wird es einen gewissen niedrigsten positiven Exponenten $a_\nu$ geben, für den $\Theta^{a_\nu}$ in $S_{\nu-1}$ enthalten ist, und dies $a_\nu$ ist wegen (16) gleich oder kleiner als $a_{\nu-1}$, und ist ausserdem eine Potenz von $p$, da es ein Theiler des Grades von $\Theta$ sein muss. Wir wählen $\Theta$ so, dass der Exponent $a_\nu$ so gross als möglich wird. Ist $\Theta_1$ ein beliebiges anderes Element in $S$, und $\Theta_1^h$ die niedrigste Potenz von $\Theta_1$, die in $S_{\nu-1}$ enthalten ist, so ist auch $h$ eine Potenz von $p$ und gleich oder kleiner als $a_\nu$. Es ist daher $\Theta_1^{a_\nu}$ in $S_{\nu-1}$ enthalten.

Setzen wir aber
\weq{17}{\Theta_1^{a_\nu} = A_1^{\lambda_1} A_2^{\lambda_2} \ldots A_{\nu-1}^{\lambda_{\nu-1}},}
so sind die sämmtlichen Exponenten $\lambda$ durch $a_\nu$ theilbar. Denn es ist
\[
\Theta_1^{a_{\nu-1}} = A_1^{\frac{\lambda_1 a_{\nu-1}}{a_\nu}} A_2^{\frac{\lambda_2 a_{\nu-1}}{a_\nu}} \ldots A_{\nu-1}^{\frac{\lambda_{\nu-1} a_{\nu-1}}{a_\nu}},
\]
und nach der Voraussetzung 2) müssen die Exponenten von $A_1, A_2 \ldots A_{\nu-1}$ auf der rechten Seite dieser Formel durch $a_{\nu-1}$ theilbare ganze Zahlen sein. Folglich sind $\lambda_1, \lambda_2 \ldots \lambda_{\nu-1}$ durch $a_\nu$ theilbar. Wir können also, indem wir zu dem oben betrachteten speciellen $\Theta$ zurückkehren und unter $h_1, h_2 \ldots h_{\nu-1}$ ganze Zahlen verstehen,
\[
\Theta^{a_\nu} = A_1^{h_1 a_\nu} A_2^{h_2 a_\nu} \ldots A_{\nu-1}^{h_{\nu-1} a_\nu}
\]
setzen, und wenn wir dann
\[
A_\nu = \Theta A_1^{-h_1} A_2^{-h_2} \ldots A_{\nu-1}^{-h_{\nu-1}}
\]
annehmen, so ist
\weq{18}{A_\nu^{a_\nu} = 1}
die niedrigste Potenz von $A_\nu$, die in $S_{\nu-1}$ enthalten ist (weil sonst auch noch eine niedrigere Potenz von $\Theta$ in $S_{\nu-1}$ enthalten wäre).

Dann ist
\weq{19}{A_1^{\alpha_1} A_2^{\alpha_2} \ldots A_\nu^{\alpha_\nu}, \quad \alpha_i = 0, 1, 2 \ldots a_i - 1}
nur $= 1$, wenn $\alpha_1, \alpha_2 \ldots \alpha_\nu$ durch $a_1, a_2 \ldots a_\nu$ theilbar und folglich $= 0$ sind, und demnach sind die in (19) dargestellten Elemente alle von einander verschieden. Diese Elemente bilden aber eine Gruppe $S_\nu$ mit der Basis $A_1, A_2 \ldots A_\nu$, die der Forderung 1) genügt.


% ===== scan leaf 60, printed 39 =====
% !PRINTED 39
\weberpage{39}
\pdfbookmark[2]{§ 10. Die Invarianten der Abel'schen Gruppen}{bm13}
Zugleich ist, wie die Formel (17) zeigt, wenn $\Theta_1$ ein beliebiges Element in $S$ ist, $\Theta_1^{a_\nu}$ in $S_{\nu-1}$ enthalten, und die Exponenten $\lambda_1, \lambda_2 \ldots \lambda_{\nu-1}$ sind durch $a_\nu$ theilbar. Es ist daher auch die Forderung 2) befriedigt.

Damit ist also unser Satz I. bewiesen. Zugleich sehen wir aus dieser Ableitung, \gesperrt{dass man für eine beliebige Abel'sche Gruppe $S$ die Elemente der Basis immer so annehmen kann, dass ihre Grade Primzahlpotenzen sind.}

Aus der Darstellbarkeit durch eine Basis folgt auch, dass jede Abel'sche Gruppe metacyklisch ist; denn sind die Elemente einer solchen Gruppe $S$ in der Form dargestellt:
\[
A = A_1^{\alpha_1} A_2^{\alpha_2} \ldots A_\nu^{\alpha_\nu},
\]
und ist $p$ eine im Grade $a_1$ von $A_1$ aufgehende Primzahl, so bilden die Elemente
\[
A' = A_1^{\alpha_1 p} A_2^{\alpha_2} \ldots A_\nu^{\alpha_\nu},
\]
wenn $\alpha_1$ ein volles Restsystem nach dem Modul $a_1 : p$ durchläuft, einen Theiler $S'$ von $S$ vom Index $p$, der durch die Basis $A_1^p, A_2 \ldots A_\nu$ darstellbar ist. Von $S'$ kann man wieder auf die gleiche Weise einen Theiler von Primzahlindex finden u. s. f. Also ist $S$ metacyklisch (\S.\,8).

\begin{center}
\S.\,10.

\gesperrt{Die Invarianten der Abel'schen Gruppen.}
\end{center}

Der Beweis, den wir im vorigen Paragraphen für die Möglichkeit der Darstellung einer Abel'schen Gruppe durch eine Basis mitgetheilt haben, enthält zugleich einen Weg, eine solche Basis zu finden, und zwar eine, bei der die Grade der Basiselemente Primzahlpotenzen sind. Gleichwohl kann es vorkommen, dass eine und dieselbe Gruppe auf verschiedene Arten durch Basen dargestellt werden kann.

Betrachten wir z. B. zwei Elemente $A, B$, deren Grade $a, b$ relativ prim sind, so wird durch diese als Elemente einer zweigliedrigen Basis eine Gruppe
\weq{1}{A^\alpha B^\beta \qquad \begin{aligned}
\alpha &= 0, 1, 2 \ldots a-1 \\
\beta &= 0, 1, 2 \ldots b-1
\end{aligned}}
dargestellt. Dieselbe Gruppe kann aber auch durch die eingliedrige Basis $AB$ dargestellt werden in der Form
\weq{2}{(AB)^s \qquad s = 0, 1, 2 \ldots ab-1.}


% ===== scan leaf 61, printed 40 =====
% !PRINTED 40
\weberpage{40}
Denn sind $\alpha$ und $\beta$ beliebig gegeben, so kann man $s$ nach dem Modul $ab$ so bestimmen, dass
\[
s \equiv \alpha \ (\mathrm{mod}\ a), \quad s \equiv \beta \ (\mathrm{mod}\ b),
\]
wodurch (2) mit (1) identisch wird. Trotzdem ist in gewissem Sinne die Constitution der Basis durch die Natur der Gruppen völlig bestimmt, nach dem folgenden Satze:

\gesperrt{II. Sind
\[
A_1, A_2 \ldots A_\nu
\]
mit den Graden
\[
a_1, a_2 \ldots a_\nu
\]
und
\[
B_1, B_2 \ldots B_\mu
\]
mit den Graden
\[
b_1, b_2 \ldots b_\mu
\]
zwei Basen einer Abel'schen Gruppe $S$ vom Grade $n$, ist $p$ eine in $n$ aufgehende Primzahl und
\weq{1}{a_1 = p_1 a_1', \quad a_2 = p_2 a_2', \quad \ldots \quad a_\nu = p_\nu a_\nu',}
\weq{2}{b_1 = p_1' b_1', \quad b_2 = p_2' b_2', \quad \ldots \quad b_\mu = p_\mu' b_\mu',}
worin $p_1, p_2 \ldots p_\nu, p_1', p_2' \ldots p_\mu'$ die höchsten Potenzen von $p$ sind, die in $a_1, a_2 \ldots a_\nu, b_1, b_2 \ldots b_\mu$ aufgehen, so kommen alle Primzahlpotenzen $p_1, p_2 \ldots p_\nu$, die grösser als $1$ sind, auch unter den $p_1', p_2' \ldots p_\mu'$ vor, und umgekehrt.}

Um diesen Satz zu beweisen, nehmen wir die Elemente der beiden Basen $A$ und $B$ so geordnet an, dass
\weq{3}{\begin{aligned}
p_1 &\geqq p_2 \geqq \ldots \geqq p_\nu \\
p_1' &\geqq p_2' \geqq \ldots \geqq p_\mu'.
\end{aligned}}
Die in (1) und (2) vorkommenden ganzen Zahlen $a_1', a_2' \ldots a_\nu', b_1', b_2' \ldots b_\mu'$ sind nach ihrer Definition durch $p$ nicht theilbar, und wenn wir also mit $m$ das kleinste gemeinschaftliche Vielfache von $a_1', a_2' \ldots a_\nu'$ bezeichnen, so ist auch $m$ nicht durch $p$ theilbar. Ist aber
\weq{4}{\Theta = A_1^{\alpha_1} A_2^{\alpha_2} \ldots A_\nu^{\alpha_\nu}}
ein beliebiges Element von $S$, so folgt, dass
\[
\Theta^{p_1 m} = 1
\]
sein muss.


% ===== scan leaf 62, printed 41 =====
% !PRINTED 41
\weberpage{41}

Setzt man hierin $B_1, B_2 \ldots B_\mu$ für $\Theta$, so folgt, dass $p_1 m$
durch jede der Zahlen $b_1, b_2 \ldots b_\mu$ theilbar ist. Es ist also $p_1$
theilbar durch $p_1'$ und $m$ durch das kleinste gemeinschaftliche
Vielfache von $b_1', b_2' \ldots b_\mu'$. In diesem Schlusse können wir nun
durchweg $A$ mit $B$ vertauschen. Es muss also auch $p_1'$ durch $p_1$
theilbar sein, und daher
\weq{5}{p_1 = p_1',}
und ausserdem ergiebt sich, dass $m$ auch das kleinste gemeinschaftliche
 Vielfache von $b_1', b_2' \ldots b_\mu'$ ist.

Um daraus unseren Satz allgemein zu beweisen, nehmen wir an,
es sei für irgend ein $s$ bewiesen, dass
\weq{6}{p_1 = p_1', p_2 = p_2', \ldots p_{s-1} = p_{s-1}'}
sein müsse. Nach (4) ist für jedes Element $\Theta$:
\weq{7}{\Theta^{p_s m} = A_1^{\alpha_1 p_s m} A_2^{\alpha_2 p_s m} \ldots A_{s-1}^{\alpha_{s-1} p_s m},}
worin $m$ wie oben das kleinste gemeinschaftliche Vielfache von
$a_1', a_2' \ldots a_\nu'$ und zugleich von $b_1', b_2' \ldots b_\mu'$, also durch $p$ nicht
theilbar ist.

Wir bestimmen nun die Anzahl der in der Form (7) enthaltenen
 von einander verschiedenen Elemente.

Lassen wir $\alpha_1$ die Reihe der Zahlen $0, 1, 2 \ldots \dfrac{p_1}{p_s} - 1$
durchlaufen, so sind die Elemente
\weq{8}{1,\ A_1^{p_s m},\ A_1^{2 p_s m} \ldots A_1^{\left(\frac{p_1}{p_s}-1\right) p_s m}}
alle von einander verschieden, während $A_1^{\frac{p_1}{p_s} p_s m}$ wieder $= 1$
wird, so dass sich alle anderen Potenzen $A_1^{\alpha_1 p_s m}$ in der Reihe (8)
wiederfinden. Ebenso schliessen wir in Bezug auf die übrigen
Factoren von (7), woraus man die genaue Anzahl der von einander
 verschiedenen in $\Theta^{p_s m}$ enthaltenen Elemente
\weq{9}{z = \frac{p_1}{p_s} \frac{p_2}{p_s} \ldots \frac{p_{s-1}}{p_s}}
findet.

Von den beiden Zahlen $p_s$, $p_s'$ wird, wenn sie nicht gleich
sind, eine die grössere sein. Wir wollen also annehmen, es sei
\weq{10}{p_s' \geqq p_s.}

Nun drücken wir $\Theta$ durch die Basis $B$ aus, und setzen
\weq{11}{\Theta^{p_s m} = B_1^{\beta_1 p_s m} B_2^{\beta_2 p_s m} \ldots B_\mu^{\beta_\mu p_s m},}
%CONT>>


% ===== scan leaf 63, printed 42 =====
% !PRINTED 42
\weberpage{42}

%<<CONT
und wir zählen nun wieder ab, wie viel verschiedene Elemente
in dieser Form enthalten sind. Die beiden Werthe für $z$ müssen
dann übereinstimmen.

Zählen wir, wie oben in (7), die Anzahl der verschiedenen
in der Form
\weq{12}{B_1^{\beta_1 p_s m} B_2^{\beta_2 p_s m} \ldots B_{s-1}^{\beta_{s-1} p_s m}}
enthaltenen Elemente, so ergiebt sich auch hierfür nach der Annahme
 (6) der Werth $z$. Es kann daher in der Form
\weq{13}{B_s^{\beta_s p_s m} \ldots B_\mu^{\beta_\mu p_s m}}
nur das einzige Element $1$ enthalten sein, woraus zu schliessen
ist, dass $p_s$ durch $p_s'$ theilbar sein muss. Das ist aber mit (10)
nur unter der Voraussetzung vereinbar, dass
\weq{14}{p_s = p_s' \text{ ist.}}

Hiermit ist unser Theorem II. bewiesen.

Die in den Gradzahlen einer Basis von $S$ enthaltenen Primzahlpotenzen
 sind also von der besonderen Wahl der Basis ganz
unabhängig und wir nennen sie daher die \gesperrt{Invarianten der
Gruppe}. Das Product aller Invarianten ist gleich dem Grade $n$
der Gruppe.

Nach \S.\,9 können wir für $S$ eine Basis bestimmen, bei der
die Grade der Elemente lauter Primzahlpotenzen sind. Nach
dem Theorem II. sind diese Grade die Invarianten der Gruppe
und sind also durch die Gruppe vollständig bestimmt.

Dass die Invarianten auch die Natur der Gruppe vollständig
bestimmen, ergiebt sich aus dem Satze:

III. \gesperrt{Zwei Gruppen mit denselben Invarianten sind
isomorph, und isomorphe Gruppen haben dieselben
 Invarianten.}

Wenn nämlich zwei Gruppen $S$ und $S'$ der Anzahl und dem
Grade nach übereinstimmende Basiselemente haben, wenn etwa
\[ \Theta = A_1^{\alpha_1} A_2^{\alpha_2} \ldots A_\nu^{\alpha_\nu} \]
die Elemente von $S$ sind, und
\[ \Theta' = B_1^{\alpha_1} B_2^{\alpha_2} \ldots B_\nu^{\alpha_\nu} \]
die Elemente von $S'$, wo die $\alpha_1, \alpha_2 \ldots \alpha_\nu$ in beiden Fällen Restsysteme
 nach den Moduln $a_1, a_2 \ldots a_\nu$ durchlaufen, so brauchen
wir nur $\Theta$ und $\Theta'$ dann einander entsprechen zu lassen, wenn
%CONT>>


% ===== scan leaf 64, printed 43 =====
% !PRINTED 43
\weberpage{43}
\pdfbookmark[2]{§ 11. Gruppencharaktere}{bm14}

%<<CONT
die Exponenten $\alpha$ in beiden dieselben sind; dann sind beide
Gruppen isomorph auf einander bezogen.

Haben also zwei Gruppen dieselben Invarianten, so können
wir diese Invarianten in beiden Gruppen zu Graden der Basiselemente
 machen, und erhalten dann diesen Fall.

Wenn umgekehrt zwei Gruppen isomorph sind, so bilden
die den Basiselementen der einen Gruppe entsprechenden Elemente
 der anderen eine Basis der letzteren, und also müssen
auch die Invarianten in beiden dieselben sein\textsuperscript{1}).

\begin{center}
§. 11.

\gesperrt{Gruppencharaktere}.
\end{center}

Es ist ein Hauptproblem in der Theorie der \gesperrt{Abel}'schen
Gruppen, alle Theiler einer \gesperrt{Abel}'schen Gruppe zu finden. Die
Lösung dieser Aufgabe wird wesentlich erleichtert durch die
Einführung des Begriffes der Gruppencharaktere, der auch sonst
in mancher Beziehung von Wichtigkeit ist.

Es sei $S$ eine \gesperrt{Abel}'sche Gruppe $n^{\text{ten}}$ Grades, deren Elemente
 durch $A$ bezeichnet werden sollen, und es seien
\weq{1}{A_1, A_2 \ldots A_\nu}
die Elemente einer Basis von $S$ von den Graden $a_1, a_2 \ldots a_\nu$.
Jedes Element $A$ ist also, und zwar nur auf eine Weise, in
der Form
\weq{2}{A = A_1^{\alpha_1} A_2^{\alpha_2} \ldots A_\nu^{\alpha_\nu}}

\textsuperscript{1}) Ueber die Theorie der \gesperrt{Abel}'schen Gruppen ist zu vergleichen:

\gesperrt{Gauss}, D\'emonstration de quelques th\'eor\`emes concernant les
périodes des classes des formes binaires du second degré. Werke,
Bd. II, S. 266.

\gesperrt{Schering}, Die Fundamentalclassen der zusammensetzbaren
arithmetischen Formen, Göttinger Abhandlungen, Bd. 14.

\gesperrt{Frobenius} und \gesperrt{Stickelberger}, Ueber Gruppen vertauschbarer
 Elemente. Crelle's Journal, Bd. 86, S. 217.

\gesperrt{Weber}, ``Theorie der \gesperrt{Abel}'schen Zahlkörper'', Acta mathematica,
Bd. 8 u. 9.

Der Begriff der Invarianten ist zuerst eingeführt in der Abhandlung
von \gesperrt{Frobenius} und \gesperrt{Stickelberger}, aber etwas anders definirt als
hier. Dieser älteren Definition, die sich bei der Anwendung auf die Kreistheilungszahlen
 minder zweckmässig erweist, ist der Verfasser in der
citirten Abhandlung in den Acta mathematica gefolgt.
%CONT>>


% ===== scan leaf 65, printed 44 =====
% !PRINTED 44
\weberpage{44}

%<<CONT
darstellbar, so dass
\weq{3}{0 \leqq \alpha_1 < a_1,\ 0 \leqq \alpha_2 < a_2, \ldots 0 \leqq \alpha_\nu < a_\nu.}

Die Exponenten $\alpha_1, \alpha_2 \ldots \alpha_\nu$ oder irgend welche ihnen nach
den Moduln $a_1, a_2 \ldots a_\nu$ congruente Zahlen heissen die \gesperrt{Indices}
des Elementes $A$, und es gilt der Satz:

1. \gesperrt{Man erhält die Indices eines Compositums $A A'$,
wenn man die entsprechenden Indices der beiden
Factoren addirt.}

Wenn wir jedem Elemente $A$ irgendwie einen \gesperrt{Zahlenwerth}
 zuordnen, so können wir diese Zuordnung eine Function
von $A$ nennen. Eine solche Function $\chi(A)$ soll nun ein
\gesperrt{Gruppencharakter} genannt werden, wenn sie für je zwei Elemente
 $A$, $A'$ von $S$ der Bedingung genügt:
\weq{4}{\chi(AA') = \chi(A)\chi(A'),}
und folglich auch der allgemeineren
\[ \chi(AA'A''\ldots) = \chi(A)\chi(A')\chi(A'')\ldots \]
Zwei solche Functionen $\chi$ und $\chi_1$ werden als verschieden
angesehen, wenn es wenigstens ein Element $A$ giebt, wofür $\chi(A)$
von $\chi_1(A)$ verschieden ist. Diese Charaktere wollen wir nun
näher bestimmen.

Setzen wir $A' = 1$, so ergiebt sich aus (4):
\weq{5}{\chi(A) = \chi(A)\chi(1),}
und wenn wir ein für allemal den Fall ausschliessen, dass alle
$\chi(A) = 0$ sind:
\weq{6}{\chi(1) = 1,}
also der erste Satz:

2. \gesperrt{Jeder Charakter hat für das Einheitselement den
Werth 1.}

Setzen wir ferner $A' = A$, so folgt aus (4):
\[ \chi(A^2) = [\chi(A)]^2, \]
und daraus durch vollständige Induction für jeden Exponenten $h$:
\weq{7}{\chi(A^h) = \chi(A)^h.}
Bezeichnen wir also mit $a$ den Grad des Elementes $A$, so
folgt, wenn man in (7) $h = a$ setzt und (6) benutzt:
\weq{8}{\chi(A)^a = 1,}
also:
%CONT>>


% ===== scan leaf 66, printed 45 =====
% !PRINTED 45
\weberpage{45}
%<<CONT

3. \gesperrt{Die Werthe eines Charakters $\chi(A)$ sind Einheitswurzeln, deren Grad ein Theiler des Grades von $A$ ist.}

Danach können wir leicht alle Charaktere bestimmen. Setzen wir nämlich
\weq{9}{
\chi(A_1) = \omega_1,\ \chi(A_2) = \omega_2 \ldots \chi(A_\nu) = \omega_\nu,
}
so ist
\weq{10}{\omega_1^{a_1} = 1,\ \omega_2^{a_2} = 1, \ldots \omega_\nu^{a_\nu} = 1,}
und es ist nach (2) und (4)
\weq{11}{\chi(A) = \omega_1^{\alpha_1}\omega_2^{\alpha_2}\ldots\omega_\nu^{\alpha_\nu}.}

Umgekehrt ist, wenn $\omega_1, \omega_2 \ldots \omega_\nu$ irgend eine Lösung der Gleichungen (10) bedeutet, durch (11) eine Function von $A$ bestimmt, die nach 1. der Bedingung (4) genügt und also ein Charakter der Gruppe ist.

Jede der Gleichungen (10) hat $a_1, a_2 \ldots a_\nu$ Wurzeln, und wenn wir jede mit jeder combiniren, so ergeben sich
\[
a_1 a_2 \ldots a_\nu = n
\]
Combinationen. Alle diese Combinationen führen zu verschiedenen Charakteren $\chi(A)$. Denn bezeichnen wir mit $\omega_1', \omega_2' \ldots \omega_\nu'$ eine zweite Combination von Wurzeln der Gleichung (10), und ist für alle Elemente $A$
\[
\omega_1^{\alpha_1}\omega_2^{\alpha_2}\ldots\omega_\nu^{\alpha_\nu} = \omega_1'^{\alpha_1}\omega_2'^{\alpha_2}\ldots\omega_\nu'^{\alpha_\nu},
\]
so folgt, wenn man $\alpha_1 = 1, \alpha_2 = 0 \ldots \alpha_\nu = 0$ setzt, $\omega_1 = \omega_1'$ und ebenso $\omega_2 = \omega_2' \ldots \omega_\nu = \omega_\nu'$. Daraus folgt der Satz:

4. \gesperrt{Es giebt $n$ und nicht mehr verschiedene Charaktere einer Abel'schen Gruppe $n^{\text{ten}}$ Grades, die alle durch die Formel (11) dargestellt sind.}

Wenn wir unter $\varepsilon_1, \varepsilon_2 \ldots \varepsilon_\nu$ ein System primitiver Wurzeln der Gleichungen (10) verstehen, so können wir
\weq{12}{\omega_1 = \varepsilon_1^{\beta_1},\ \omega_2 = \varepsilon_2^{\beta_2}, \ldots \omega_\nu = \varepsilon_\nu^{\beta_\nu}}
setzen, und können die $\beta_1, \beta_2 \ldots \beta_\nu$ ebenso wie die $\alpha_1, \alpha_2 \ldots \alpha_\nu$ je einem vollen Restsysteme nach den Moduln $a_1, a_2 \ldots a_\nu$ entnehmen. Dann bekommen wir die sämmtlichen $n$ Gruppencharaktere bei feststehenden $\varepsilon_1, \varepsilon_2 \ldots \varepsilon_\nu$ in der Form
\weq{13}{\chi(A) = \varepsilon_1^{\alpha_1\beta_1}\varepsilon_2^{\alpha_2\beta_2}\ldots\varepsilon_\nu^{\alpha_\nu\beta_\nu}.}


% ===== scan leaf 67, printed 46 =====
% !PRINTED 46
\weberpage{46}

Unter diesen Charakteren kommt auch der vor, den man erhält, wenn man $\beta_1 = 0, \beta_2 = 0 \ldots \beta_\nu = 0$ setzt, der für alle Elemente $A$ den Werth 1 hat. Diesen wollen wir den \gesperrt{Einheitscharakter} nennen.

Die $n$ Charaktere von $S$ können nun selbst wieder zu einer \gesperrt{Abel}'schen Gruppe vereinigt werden, und zwar zu einer mit $S$ isomorphen Gruppe.

Nach (13) ist nämlich jeder der $n$ Charaktere durch ein nach dem Modulsysteme $a_1, a_2 \ldots a_\nu$ genommenes Zahlensystem $\beta_1, \beta_2 \ldots \beta_\nu$ bestimmt, und dies Zahlensystem der $\beta$ ist zugleich das System der Indices eines bestimmten Elementes $B$ von $S$, nämlich von
\weq{14}{B = A_1^{\beta_1}A_2^{\beta_2}\ldots A_\nu^{\beta_\nu},}
so dass jedem Elemente $B$ von $S$ ein bestimmter Charakter entspricht, den wir mit $\chi_B(A)$ oder, indem wir $A$ weglassen, mit $\chi_B$ bezeichnen. Diese Zuordnung der Charaktere $\chi$ zu den Elementen $B$ wird sich aber ändern, wenn eine andere Basis zu Grunde gelegt wird, oder wenn die Einheitswurzeln $\varepsilon_1, \varepsilon_2 \ldots \varepsilon_\nu$ anders angenommen werden.

Ist $B'$ ein zweites Element von $S$ mit den Indices $\beta_1', \beta_2' \ldots \beta_\nu'$, so erhalten wir in gleicher Weise den Charakter $\chi_{B'}$, und wenn wir nun unter $\chi_{BB'}$ den Charakter verstehen, der für jedes $A$ durch die Formel
\weq{15}{\chi_{BB'}(A) = \varepsilon_1^{\alpha_1(\beta_1+\beta_1')}\varepsilon_2^{\alpha_2(\beta_2+\beta_2')}\ldots\varepsilon_\nu^{\alpha_\nu(\beta_\nu+\beta_\nu')}}
bestimmt ist, so sind die Charaktere $\chi$ hierdurch zu einer mit $S$ isomorphen Gruppe verbunden. Das Einheitselement in dieser \gesperrt{Charakterengruppe} ist der Einheitscharakter.

Es lässt sich hiernach auch die Gruppe der Charaktere durch eine Basis darstellen, die man erhält, wenn man
\weq{16}{\chi_1(A) = \varepsilon_1^{\alpha_1},\ \chi_2(A) = \varepsilon_2^{\alpha_2} \ldots \chi_\nu(A) = \varepsilon_\nu^{\alpha_\nu}}
setzt. Dann ist jeder Charakter in der Form enthalten:
\weq{17}{\chi_B = \chi_1^{\beta_1}\chi_2^{\beta_2}\ldots\chi_\nu^{\beta_\nu},}
und es ist, wenn $B, B'$ zwei Elemente in $S$ sind:
\weq{18}{\chi_B\chi_{B'} = \chi_{BB'}.}

Sind $\chi_1, \chi_2$ irgend zwei der Charaktere und $\chi_1\chi_2$ der aus beiden zusammengesetzte, so ist nach (15) für jedes Element $A$
\weq{19}{\chi_1(A)\chi_2(A) = \chi_1\chi_2(A).}


% ===== scan leaf 68, printed 47 =====
% !PRINTED 47
\weberpage{47}

Wir wollen das gewonnene Resultat noch als Satz aussprechen:

5. \gesperrt{Die Charaktere einer Gruppe $S$ können zu einer mit $S$ isomorphen Gruppe verbunden werden.}

Es sei noch der Satz erwähnt, der sich aus der Definition von $\chi_B$ \gesperrt{durch} die Formel (13) unmittelbar ergiebt:
\weq{20}{\chi_B(A) = \chi_A(B).}

Endlich gilt auch noch der folgende Satz:

6. \gesperrt{Durchläuft $A$ alle Elemente der Gruppe $S$, so ist für einen feststehenden Charakter $\chi$:}
\weq{21}{\overset{A}{\Sigma}\,\chi(A) = n \text{ oder } = 0,}
\gesperrt{je nachdem $\chi$ der Einheitscharakter ist oder nicht. Ebenso ist, wenn das Element $A$ festgehalten wird und $\chi$ die Reihe der Charaktere durchläuft:}
\weq{22}{\overset{\chi}{\Sigma}\,\chi(A) = n \text{ oder } = 0,}
\gesperrt{je nachdem $A$ das Einheitselement ist oder nicht.}

Für den Fall, dass $\chi$ oder $A$ die Einheitselemente ihrer Gruppen sind, sind die Formeln (21) und (22) evident, da dann jedes der $n$ Glieder der Summe den Werth 1 hat. Ist aber $\chi$ nicht der Einheitscharakter, so giebt es ein Element $B$ in $S$, so dass $\chi(B)$ nicht $= 1$ ist. Setzen wir dann
\[
\overset{A}{\Sigma}\,\chi(A) = \sigma,
\]
so folgt durch Multiplication mit $\chi(B)$:
\[
\overset{A}{\Sigma}\,\chi(AB) = \sigma\chi(B).
\]
Da aber $AB$ zugleich mit $A$ die ganze Gruppe $S$ durchläuft, so ist auch $\overset{A}{\Sigma}\,\chi(AB) = \sigma$, also $\sigma[1 - \chi(B)] = 0$ oder $\sigma = 0$, w.\,z.\,b.\,w.

Ebenso beweist man die Formel (22), die übrigens auch nach (20) unmittelbar aus (21) folgt.


% ===== scan leaf 69, printed 48 =====
% !PRINTED 48
\weberpage{48}
\pdfbookmark[2]{§ 12. Divisoren einer Abel'schen Gruppe. Reciproke Gruppen}{bm15}

\begin{center}
\S.\,12.

\gesperrt{Divisoren einer Abel'schen Gruppe}.

\gesperrt{Reciproke Gruppen}.
\end{center}

Ein Theiler $T$ einer \gesperrt{Abel}'schen Gruppe $S$ ist, wie schon oben bemerkt, immer ein Normaltheiler, und daher giebt es nach \S.\,4 eine zu $T$ complementäre Gruppe
\[
S/T,
\]
deren Elemente die Nebengruppen von $T$ sind, nämlich
\weq{1}{T,\ TA',\ TA''\ldots,}
worin $A', A'' \ldots$ gewisse Elemente aus $S$ bedeuten. Die Anzahl der Elemente (1), also der Grad der Gruppe $S/T$ ist, wenn $t$ der Grad von $T$ ist, gleich dem Index des Theilers $T$ von $S$, also gleich
\[
\frac{n}{t} = j.
\]

Diese Gruppe $S/T$ ist aber selbst wieder eine \gesperrt{Abel}'sche; denn es ist
\[
TA'\,TA'' = TA'A'',\ \ TA''\,TA' = TA''A',
\]
also beides einander gleich.

Es sind nun die Charaktere der Gruppe $S/T$ zu bestimmen. Wir bezeichnen die Elemente dieser Gruppe mit
\weq{2}{T,\ T_1,\ T_2 \ldots T_{j-1},}
und einen ihrer Charaktere mit $\xi(T_i)$.

Aus dieser Function $\xi(T)$ können wir nun eine Function $\xi(A)$ der Elemente von $S$ ableiten, indem wir
\weq{3}{\xi(A_i) = \xi(T_i)}
setzen, wenn $A_i$ irgend ein in der Nebengruppe $T_i$ vorkommendes Element ist. Für die Elemente der Gruppe $T$ selbst ist dann $\xi(A) = 1$.

Diese Function $\xi(A_i)$ ist aber unter den Charakteren von $S$ enthalten. Denn wenn $A_i$ in $T_i$ und $A_k$ in $T_k$ vorkommt, so kommt $A_iA_k$ in $T_iT_k$ vor, und folglich ist
\weq{4}{\xi(A_i)\xi(A_k) = \xi(T_i)\xi(T_k) = \xi(T_iT_k) = \xi(A_iA_k).}
Dies aber ist nach \S.\,11,(4) die Definition für einen Charakter von $S$.


% ===== scan leaf 70, printed 49 =====
% !PRINTED 49
\weberpage{49}

Wenn umgekehrt einer der Charaktere $\chi = \xi$ der Gruppe $S$
für alle Elemente der Gruppe $T$ denselben Werth hat, so kann
dies nur der Werth 1 sein, da in $T$ sicher das Einheitselement
von $S$ vorkommt (\S.\,11, 2.). Wenn dann $A_i$ irgend ein Element
aus $T_i$ ist, und $A$ die ganze Gruppe $T$ durchläuft, so
durchläuft $AA_i$ die Nebengruppe $T_i$. Es ist dann aber
$\xi(A) = 1$ und folglich
\weq{5}{\xi(AA_i) = \xi(A_i),}
d.\,h. $\xi(A)$ hat für alle Elemente einer Nebengruppe $T_i$ einen
und denselben Werth, und $\xi(A)$ kann also als Function von $T_i$
aufgefasst und mit $\xi(T_i)$ bezeichnet werden.

Da nun, wenn $A_i$ in $T_i$ und $A_k$ in $T_k$ vorkommt, $A_iA_k$
in $T_iT_k$ enthalten ist, so folgt
\weq{6}{\xi(T_i)\xi(T_k) = \xi(T_iT_k),}
d.\,h. $\xi(T_i)$ ist unter den Charakteren der Gruppe $S/T$
enthalten. Es giebt also genau $j$ solche Charaktere $\xi(A)$, die
dadurch definirt sind, dass sie für jedes Element in $T$ den Werth
1 haben. Diese $j$ Charaktere $\xi$ bilden nach der Composition
der Charaktere eine Gruppe; denn ist $\xi_1(A) = 1$,
$\xi_2(A) = 1$, so ist auch $\xi_1\xi_2(A) = 1$ [\S.\,11, (18)].
Da die Functionen $\xi$ nach (4) auch als die Charaktere der
Gruppe $S/T$ aufgefasst werden können, so ist nach \S.\,11, 5. die
Gruppe der $\xi$ mit der Gruppe $S/T$ isomorph.

Damit ist also der folgende Satz bewiesen:

7.\ \gesperrt{Hat eine Abel'sche Gruppe $S$ einen Theiler $T$ vom
Grade $t$ und vom Index $j$, so giebt es unter den Charakteren von
$S$ genau $j$ und nicht mehr, die für alle Elemente von $T$ den
Werth 1 haben, während für jedes nicht in $T$ enthaltene Element
von $S$ wenigstens einer von ihnen von 1 verschieden ist. Diese
Charaktere bilden eine mit $S/T$ isomorphe Gruppe.}

Wir wollen diese Charaktere \gesperrt{zur Gruppe $T$ gehörig} nennen.

Da jeder Charakter $\chi_B$ einem bestimmten Elemente $B$ von $S$
entspricht, so wird auch, wenn $\chi_B$ die Gruppe der zu $T$
gehörigen Charaktere durchläuft, $B$ wegen der Formel \S.\,11, (18)
eine Gruppe durchlaufen, die vom Grade $j$ und mit der Gruppe der
$\chi_B$ und also auch mit der Gruppe $S/T$ isomorph ist. Diese
%CONT>>


% ===== scan leaf 71, printed 50 =====
% !PRINTED 50
\weberpage{50}
%<<CONT
Gruppe der $B$, die also auch ein Theiler von $S$ ist, wollen wir
mit $U$ bezeichnen und die zu $T$ \gesperrt{reciproke Gruppe} nennen. Auch
hier ist aber zu bemerken, dass der Begriff der reciproken Gruppe
im Allgemeinen von der Wahl der Basis und der Einheitswurzeln
$\varepsilon$ abhängt, also nicht zu den Gruppen $S$ und $T$ als
solchen gehört.

Die zu $T'$ reciproke Gruppe ist durch folgenden Satz charakterisirt:
% [sic: Weber prints T' here; the newly defined group is reciprocal to T.]

8.\ \gesperrt{Lässt man $A$ die Elemente eines Theilers $T$ von $S$
durchlaufen und sucht alle Elemente $B$ von $S$, die für jedes $A$
der Bedingung}
\weq{7}{\chi_B(A) = 1}
\gesperrt{genügen, so durchläuft $B$ die Elemente der zu $T$
reciproken Gruppe $U$, deren Grad gleich dem Index von $T$ ist.}

Nach dem Satze \S.\,11, (20) ist $T'$ die reciproke Gruppe zu $U$,
% [sic: Weber prints T' here; reciprocity requires T.]
die Beziehung dieser beiden Gruppen also eine gegenseitige.

Von diesen reciproken Gruppen gilt noch der Satz:

9.\ \gesperrt{Ist $T'$ ein Theiler von $T$, so ist umgekehrt die
reciproke Gruppe $U$ von $T'$ ein Theiler der zu $T''$ reciproken
Gruppe $U'$.}
% [sic: Weber prints T' and T'' here; the theorem requires T and T', respectively.]

Denn jedes Element $A'$ von $T''$ ist zugleich in $T'$ enthalten,
% [sic: Weber prints T'' and T' here; the proof requires T' and T, respectively.]
und folglich ist, wenn $B$ die Gruppe $U$ durchläuft,
$\chi_B(A') = 1$. Es muss also $B$ in der Gruppe $U'$ vorkommen.

Wählt man aus der Gesammtheit der Charaktere $\chi$ eine beliebige
Anzahl $\xi_1, \xi_2 \dots \xi_\mu$ aus, gleichviel ob sie eine
Gruppe bilden oder nicht, so bilden alle Elemente $A$, die den
$\mu$ Bedingungen
\weq{8}{\xi_1(A) = 1,\ \xi_2(A) = 1, \dots \xi_\mu(A) = 1}
genügen, eine Gruppe, weil aus $\xi_i(A) = 1$, $\xi_i(A') = 1$
folgt, dass auch $\xi_i(AA') = 1$ ist. Wenn die
$\xi_1, \xi_2 \dots \xi_\mu$ keine Gruppen sind, so folgen aus den
Gleichungen (8) noch so viele weitere $\xi_{\mu+1}(A) = 1 \dots$,
dass die $\xi_1, \xi_2 \dots \xi_\mu$ zu einer Gruppe ergänzt
werden.

10.\ \gesperrt{Aus dem Theorem 7. folgt, dass man so alle
möglichen Theiler von $S$ erhalten kann.}


% ===== scan leaf 72, printed 51 =====
% !PRINTED 51
\weberpage{51}

Es ist vielleicht erwünscht, diese Sätze an einem einfachen
Beispiele zu erläutern. Es sei der Grad der Gruppe $S$ das Quadrat
einer Primzahl $n = p^2$. Dann hat $S$ entweder nur die eine
Invariante $p^2$ und ist dann cyklisch:
\[
\text{a)}\quad S = 1, A, A^2 \dots A^{p^2-1},
\]
oder es hat $S$ zwei Invarianten, die beide gleich $p$ sind; dann
besteht $S$ aus den Elementen:
\[
\text{b)}\quad S = A_1^{\alpha_1}A_2^{\alpha_2}, \quad \alpha_1,
\alpha_2 = 0, 1 \dots p-1.
\]
Im Falle a) erhalten wir die $p^2$ Charaktere
\[
\chi_B(A^\alpha) = \varepsilon^{\beta\alpha},
\]
wenn $\varepsilon$ eine primitive Wurzel der Gleichung
$\varepsilon^{pp} = 1$ ist.
% [sic: the source prints \varepsilon^{pp}; the cyclic group order and surrounding p^2 formulas require \varepsilon^{p^2}.]
Nehmen wir, um nach dem Satze 10. die
Theiler von $S$ zu bilden, einen der Charaktere $\chi_B$, in dem
$\beta$ nicht durch $p$ theilbar ist, so wird $\chi_B(A^\alpha)$
nur dann $= 1$ sein können, wenn $\alpha$ durch $p^2$ theilbar
oder also $= 0$ ist, d.\,h. wir bekommen nur den aus dem
Einheitselemente bestehenden Theiler von $S$. Nehmen wir aber
$\beta = p\beta'$ durch $p$, aber nicht durch $p^2$ theilbar an, so
wird $\chi_{p\beta'}(A^\alpha)$ dann und nur dann $= 1$, wenn
$\alpha = p\alpha'$ durch $p$ theilbar ist. Wir erhalten also den
Theiler
\[
T = 1, A^p, A^{2p} \dots A^{(p-1)p},
\]
und der zu $T$ reciproke Theiler $U$ ist hier mit $T$ identisch.

Im Falle b) müssen wir, um die Charaktere zu bilden, zwei
$p^{\text{te}}$ Einheitswurzeln $\varepsilon^{\beta_1},
\varepsilon^{\beta_2}$ annehmen, und erhalten, wenn
\[
B = A_1^{\beta_1}A_2^{\beta_2}
\]
gesetzt ist,
\[
\chi_B(A_1^{\alpha_1}A_2^{\alpha_2}) = \varepsilon^{\alpha_1\beta_1 + \alpha_2\beta_2}.
\]
Setzen wir zwei dieser Charaktere $= 1$, also
\[
\alpha_1\beta_1 + \alpha_2\beta_2 \equiv 0, \quad
\alpha_1\beta_1' + \alpha_2\beta_2' \equiv 0,
\]
so folgt, wenn die Determinante $\beta_1\beta_2' - \beta_2\beta_1'$
nicht durch $p$ theilbar ist, dass $\alpha_1$ und $\alpha_2$ durch
$p$ theilbar sein müssen. Ist aber die Determinante durch $p$
theilbar, so ist die eine dieser beiden Congruenzen eine Folge der
anderen. Im ersten Falle bekommen wir also eine Gruppe, die nur
aus dem Einheitselemente besteht. Wir erhalten daher alle von 1
und $S$ verschiedenen Theiler $T_{\beta_1,\beta_2}$, wenn wir für
ein feststehendes $\beta_1, \beta_2$ die $\alpha_1, \alpha_2$ auf
alle möglichen Arten der Congruenz
\weq{9}{\alpha_1\beta_1 + \alpha_2\beta_2 \equiv 0\ (\mathrm{mod}\ p)}
%CONT>>


% ===== scan leaf 73, printed 52 =====
% !PRINTED 52
\weberpage{52}
\pdfbookmark[2]{§ 13. Die Geschlechter in einer Abel'schen Gruppe}{bm16}
%<<CONT
gemäss bestimmen. Ist $\alpha_1, \alpha_2$ eine Lösung dieser
Congruenz, in der nicht beide Zahlen Null sind, so erhalten wir
alle Lösungen, wenn wir in $h\alpha_1, h\alpha_2$ den Factor $h$
die Reihe der Zahlen $0, 1 \dots p-1$ durchlaufen lassen, und die
Gruppe $T$ wird also, wenn $\alpha_1, \alpha_2$ ein festes, der
Bedingung (9) genügendes Werthpaar ist,
\[
(A_1^{\alpha_1}A_2^{\alpha_2})^h, \quad h = 0, 1, \dots p-1.
\]

Die reciproke Gruppe $U$ erhält man, wenn man alle Werthe der
$\beta$ sucht, die der Bedingung (9) genügen, und man findet also
die Gruppe $U$ in der Form
\[
(A_1^{\beta_1}A_2^{\beta_2})^h, \quad h = 0, 1, \dots p-1,
\]
worin $\alpha_1, \alpha_2, \beta_1, \beta_2$ jetzt vier feste, der
Bedingung (9) genügende Zahlen sind. Setzt man $\beta_1 = 0$,
$\beta_2 = 1$, $\alpha_1 = 1$, $\alpha_2 = 0$, so erhält man den
besonderen Fall der beiden reciproken Gruppen
\[
A_1^h, \quad A_2^h, \quad h = 0, 1, \dots p-1.
\]

\begin{center}
\S.\,13.

\gesperrt{Die Geschlechter in einer Abel'schen Gruppe.}
\end{center}

Ein Element einer \gesperrt{Abel}'schen Gruppe $S$, das mit seinem
entgegengesetzten identisch ist, dessen zweite Potenz also gleich
dem Einheitselemente ist, wird ein \gesperrt{zweiseitiges
Element}\footnote{Vergl. Bd.~I, S.~390, Anmerkung.} genannt. Das
Einheitselement gehört also immer zu den zweiseitigen.

Stellen wir die Elemente von $S$ durch eine Basis
$A_1, A_2 \dots A_\nu$ dar, deren Elemente die Grade
$a_1, a_2 \dots a_\nu$ haben, so wird
\weq{1}{A = A_1^{\alpha_1}A_2^{\alpha_2} \dots A_\nu^{\alpha_\nu}}
dann und nur dann ein zweiseitiges Element sein, wenn
\weq{2}{2\alpha_1 \equiv 0\ (\mathrm{mod}\ a_1), \ 2\alpha_2 \equiv 0\ (\mathrm{mod}\ a_2), \dots 2\alpha_\nu \equiv 0\ (\mathrm{mod}\ a_\nu).}
Wenn nun unter den Invarianten der Gruppe $\lambda$ mal eine
Potenz von 2 vorkommt, so können wir die Basis von $S$ so
geordnet annehmen, dass $a_1, a_2 \dots a_\lambda$ gerade,
$a_{\lambda+1} \dots a_\nu$ ungerade Zahlen sind. Dann ergeben
sich die Lösungen von (2):
%CONT>>


% ===== scan leaf 74, printed 53 =====
% !PRINTED 53
\weberpage{53}
%<<CONT

\[
\alpha_1 = \frac{\eta_1 a_1}{2},\ \alpha_2 = \frac{\eta_2 a_2}{2} \cdots \alpha_\lambda = \frac{\eta_\lambda a_\lambda}{2},\ \alpha_{\lambda+1}=0 \ldots \alpha_\nu = 0,
\]

worin $\eta_1, \eta_2 \ldots \eta_\lambda$ gleich 0 oder gleich 1 sein können, und man erhält alle zweiseitigen Elemente in der Form

\[
A_1^{\frac{\eta_1 a_1}{2}} A_2^{\frac{\eta_2 a_2}{2}} \ldots A_\lambda^{\frac{\eta_\lambda a_\lambda}{2}},
\]

und ihre Anzahl ist, da alle Combinationen von 0 und 1 für die Exponenten $\eta$ zulässig sind, $= 2^\lambda$.

Ist $\chi$ ein beliebiger Charakter von $S$, so ist, wenn $A$ ein zweiseitiges Element ist, $\chi(A) = \pm 1$, da

\[
\chi(A)^2 = \chi(1) = 1
\]

ist.

Ebenso nennen wir einen zweiseitigen Charakter einen solchen, der in der Gruppe der Charaktere ein zweiseitiges Element ist. Da die Gruppe der Charaktere mit der Gruppe $S$ isomorph ist, so giebt es ebenso viele zweiseitige Charaktere als zweiseitige Elemente, nämlich $2^\lambda$. Sie werden nach \S.\,11, (17) dargestellt durch:

\[
\chi_1^{\frac{\eta_1 a_1}{2}} \chi_2^{\frac{\eta_2 a_2}{2}} \ldots \chi_\lambda^{\frac{\eta_\lambda a_\lambda}{2}}.
\]

\gesperrt{Die zweiseitigen Charaktere haben für jedes Element den Werth $\pm 1$.}

Die zweiseitigen Elemente bilden für sich eine Gruppe $T$ vom Grade $2^\lambda$. Ebenso bilden die zweiseitigen Charaktere eine damit isomorphe Gruppe, und man erhält die zu $T$ reciproke Gruppe $U$, wenn man alle Elemente $A$ aufsucht, für die alle zweiseitigen Charaktere den Werth $+1$ haben. Diese Gruppe, die nach \S.\,12, 7. ein Theiler von $S$ vom Index $2^\lambda$ ist, ist nach der letzten Definition weder von der Wahl der Basis noch von den die Charaktere darstellenden Einheitswurzeln abhängig, und ist durch die Natur der Gruppe $S$ vollständig bestimmt. Bezeichnen wir sie mit $G$, so ist das System der Nebengruppen

\weq{3}{G,\ G_1,\ G_2 \ldots G_{2^\lambda - 1}}

dadurch charakterisirt, dass für alle Elemente eines dieser Systeme die zweiseitigen Charaktere ein und dasselbe Werthsystem ergeben. Die Systeme $G, G_1, G_2 \ldots$ werden die in $S$ enthaltenen \gesperrt{Geschlechter} (Genera) genannt. Die Gruppe $G$ speciell heisst das \gesperrt{Hauptgeschlecht}.


% ===== scan leaf 75, printed 54 =====
% !PRINTED 54
\weberpage{54}
\pdfbookmark[2]{§ 14. Indices nach einer ungeraden Primzahlpotenz als Modul}{bm17}

\gesperrt{Die Anzahl der Geschlechter ist also so gross, wie die Anzahl der zweiseitigen Elemente}\footnote{\gesperrt{Gauss} hat in die Theorie der quadratischen Formen diese Begriffe zuerst eingeführt, und den Ausdruck \glqq Genera\grqq{} gebraucht. Disqu. arithm. art. 228 f.}.

\begin{center}
\S.\,14.
\end{center}

\begin{center}
\gesperrt{Indices nach einer ungeraden Primzahlpotenz}\\
\gesperrt{als Modul.}
\end{center}

Das wichtigste Beispiel einer commutativen Gruppe bieten die natürlichen Zahlen, wenn sie durch die gewöhnliche Multiplication mit einander verbunden werden.

Um daraus eine endliche Gruppe abzuleiten, nehmen wir eine beliebige ganze positive Zahl $m$ als Modul an und zerlegen $m$ in seine Primfactoren

\weq{1}{m = 2^\lambda q_1^{\varkappa_1} q_2^{\varkappa_2} \ldots,}

worin $q_1, q_2 \ldots$ verschiedene ungerade Primzahlen, $\lambda, \varkappa_1, \varkappa_2 \ldots$ positive Exponenten, $\lambda$ möglicherweise auch die Null, nämlich wenn $m$ ungerade ist, bedeuten.

Nun werden alle Zahlen, positive sowohl als negative, die nach dem Modul $m$ unter einander congruent sind, in eine \gesperrt{Zahlclasse} vereinigt, und jede dieser Zahlclassen wird durch einen Repräsentanten, etwa durch den Rest der Division, also durch eine der Zahlen $0, 1, 2 \ldots m-1$, dargestellt.

Sind $a$ und $a'$ zwei Zahlen einer solchen Classe und $b$ und $b'$ zwei Zahlen einer zweiten Classe, so gehören auch die Producte $ab$ und $a'b'$ in dieselbe Classe. Denn ist $a \equiv a'$, $b \equiv b'$, so ist auch $ab \equiv a'b'\ (\mathrm{mod}\ m)$. Durch die Multiplication werden also nicht bloss die Zahlen, sondern auch die Classen componirt.

Trotzdem bilden diese Zahlclassen in ihrer Gesammtheit noch keine Gruppe; denn aus einer Congruenz

\[
ab \equiv ac\ (\mathrm{mod}\ m)
\]

folgt nur dann nothwendig $b \equiv c$, wenn $a$ relativ prim zu $m$ ist. Es ist also die Forderung \S.\,1, 3. hier im Allgemeinen nicht erfüllt.

Der grösste gemeinschaftliche Theiler, den eine Zahl $a$ mit $m$ hat, ist bei der ganzen durch $a$ repräsentirten Classe der-
%CONT>>


% ===== scan leaf 76, printed 55 =====
% !PRINTED 55
\weberpage{55}
%<<CONT
selbe, und kann der Theiler der Classe genannt werden. Sind $d, d'$ die Theiler zweier Classen $a, a'$, so ist der grösste gemeinschaftliche Theiler von $m$ und $dd'$ der Theiler der Classe $aa'$. Daraus ergiebt sich, dass die Zahlclassen, deren Zahlen relativ prim zu $m$ sind, durch Composition immer wieder solche Zahlclassen ergeben, und für diese ist dann auch die Forderung \S.\,1, 3. erfüllt.

1.\ \gesperrt{Die Zahlclassen, deren Individuen relativ prim zum Modul sind, bilden also bei der Composition durch Multiplication eine Abel'sche Gruppe.}

Diese Gruppe ist der Gegenstand unserer Betrachtungen, wobei unser Hauptziel die Bestimmung einer Basis sein soll.

Wir wollen jede Zahlclasse, die nur zu $m$ theilerfremde Zahlen enthält, mit $N$ bezeichnen, und die Gruppe der Zahlclassen $N$, deren Existenz wir jetzt nachgewiesen haben, mit $\mathfrak{N}$. Mit $n$ wollen wir jede zu $m$ theilerfremde Zahl bezeichnen.

Der Grad dieser Gruppe ist so gross, wie die Anzahl der relativen Primzahlen zu $m$, die zugleich positiv und nicht grösser als $m$ sind, und diese Zahl haben wir in \S.\,132 des ersten Bandes bestimmt. Sie ist

\weq{2}{\mu = \varphi(m) = 2^{\lambda-1} q_1^{\varkappa_1 - 1}(q_1 - 1)\, q_2^{\varkappa_2 - 1}(q_2 - 1) \ldots}

oder

\[
= q_1^{\varkappa_1 - 1}(q_1 - 1)\, q_2^{\varkappa_2 - 1}(q_2 - 1) \ldots,
\]

wenn $\lambda = 0$ ist.

Der Grad eines Elementes $N$ dieser Gruppe ist der kleinste positive Exponent $e$, zu dem man einen Repräsentanten $n$ von $N$ erheben muss, damit $n^e$ der Einheit congruent wird nach dem Modul $m$. Da jedes $e$ ein Theiler des Grades der Gruppe $\varphi(m)$ sein muss, so folgt der \gesperrt{verallgemeinerte Fermat'sche Lehrsatz}:

\weq{3}{n^{\varphi(m)} \equiv 1\ (\mathrm{mod}\ m),}

eine Formel, die für jede zu $m$ theilerfremde Zahl $n$ gilt.

Ist nun $q$ einer der ungeraden Primfactoren von $m$, und $q^\varkappa$ die höchste in $m$ aufgehende Potenz von $q$, so nehmen wir eine primitive Wurzel $g$ von $q$ an, die wir, wenn $\varkappa$ grösser als 1 ist, so wählen, dass $g^q - g$ nicht durch $q^2$ theilbar ist (Bd. I, \S.\,184). Der Kürze wegen wollen wir eine dieser letzten Bedingung genügende Zahl $g$ eine \gesperrt{primitive Wurzel von $q^2$} nennen.


% ===== scan leaf 77, printed 56 =====
% !PRINTED 56
\weberpage{56}

Ist nun $\varkappa = 1$, so sind nach der Definition von $g$ die Zahlen

\[
1,\ g,\ g^2 \ldots g^{q-2}
\]

nach dem Modul $q$ alle von einander verschieden, während $g^{q-1}$ wieder congruent mit 1 ist.

Ist aber $\varkappa > 1$, so ist:

\weq{4}{g^{q-1} = 1 + hq,}

und nach unserer Voraussetzung über $g$ ist $h$ nicht durch $q$ theilbar. Wenn wir die Gleichung (4) in die Potenz $q$ erheben, und rechts den binomischen Lehrsatz anwenden, so ergiebt sich

\[
g^{q(q-1)} = 1 + hq^2 + h^2\, \frac{q^3(q-1)}{2} + \cdots,
\]

also

\weq{5}{g^{q(q-1)} = 1 + h_1 q^2,}

worin

\[
h_1 = h + h^2\, \frac{q(q-1)}{2} + \cdots
\]

nicht durch $q$ theilbar ist. (Das wäre für $q = 2$ nicht mehr richtig und darum verlangt die Primzahl 2 eine andere Behandlung.)

Erheben wir (5) nochmals in die $q^{\mathrm{te}}$ Potenz, so ergiebt sich

\[
g^{q^2(q-1)} = 1 + h_2 q^3,
\]

und so können wir fortfahren und erhalten für jeden beliebigen positiven Exponenten $\lambda$

\weq{6}{g^{q^{\lambda-1}(q-1)} = 1 + hq^\lambda,}

worin $h$ eine durch $q$ nicht theilbare ganze Zahl ist. Setzen wir zur Abkürzung:

\weq{7}{q^{\varkappa-1}(q-1) = \varphi(q^\varkappa) = c,}

so ist also

\weq{8}{g^c \equiv 1\ (\mathrm{mod}\ q^\varkappa),}

und es ist noch nachzuweisen, dass $c$ der kleinste positive Exponent ist, für den die Congruenz (8) erfüllt ist. Nehmen wir an, es sei $e$ dieser kleinste Exponent, also

\weq{9}{g^e \equiv 1\ (\mathrm{mod}\ q^\varkappa),}

so muss $e$ ein Theiler von $c$ sein, weil sonst die Congruenz (8) auch erfüllt wäre, wenn $c$ durch den Rest der Division von $c$ durch $e$, der kleiner als $e$ ist, ersetzt wird.

Andererseits muss $e$ durch $q-1$ theilbar sein, weil die in (9) enthaltene Congruenz $g^e \equiv 1\ (\mathrm{mod}\ q)$ nur für solche Exponenten, die durch $q-1$ theilbar sind, besteht.


% ===== scan leaf 78, printed 57 =====
% !PRINTED 57
\weberpage{57}

Es ist also $e$ von der Form $q^{\lambda-1}(q-1)$ mit einem positiven $\lambda$. Dass aber $\lambda$ nicht kleiner als $\varkappa$ sein kann, folgt aus (6), woraus zu sehen ist, dass $q^\lambda$ die höchste Potenz von $q$ ist, die in der Differenz

\[
g^{q^{\lambda-1}(q-1)} - 1
\]

aufgeht. Es ist also $c$ die kleinste positive Zahl, für die die Congruenz (8) erfüllt ist, und damit ist gleichbedeutend, dass die Zahlen

\weq{10}{1, g, g^2 \ldots g^{c-1}}

alle incongruent sind nach dem Modul $q^\varkappa$; $g$ kann also auch als \gesperrt{primitive Wurzel} von $q^\varkappa$ bezeichnet werden.

Die Anzahl der Glieder dieser Reihe (10) ist gleich $c$. Ebenso gross ist aber auch die Anzahl der nach dem Modul $q^\varkappa$ incongruenten, durch $q$ nicht theilbaren Zahlen $n$, und damit ist bewiesen:

\noindent\textbf{2.} \gesperrt{Wenn} $q$ \gesperrt{eine ungerade Primzahl,} $n$ \gesperrt{eine beliebige durch} $q$ \gesperrt{nicht theilbare Zahl,} $g$ \gesperrt{eine primitive Wurzel von} $q^2$ \gesperrt{und} $\varkappa$ \gesperrt{ein positiver Exponent ist, so lässt sich eine und nur eine Zahl} $\gamma$ \gesperrt{nach dem Modul} $c$ \gesperrt{bestimmen, die der Congruenz}

\[
n \equiv g^\gamma \pmod{q^\varkappa}
\]

\gesperrt{genügt.}

Die Zahl $c$ ist immer durch 2 theilbar, und wir heben also noch den besonderen Satz hervor, dass

\weq{11}{g^{\frac12 c} \equiv -1 \pmod{q^\varkappa}}

ist. Denn in dem durch $q^\varkappa$ theilbaren Producte

\[
g^c - 1 = (g^{\frac12 c} - 1)(g^{\frac12 c} + 1)
\]

ist der erste Factor nicht durch $q^\varkappa$ theilbar, und folglich muss der zweite Factor durch $q$ theilbar sein. Dann folgt aber, dass der erste Factor auch nicht durch $q$ theilbar sein kann, weil sonst auch die Summe der beiden Factoren, also auch $g$ selbst, durch $q$ theilbar sein müsste. Folglich muss der zweite Factor durch $q^\varkappa$ theilbar sein.

Hier besteht also vollständige Analogie mit den Sätzen über primitive Wurzeln und Indexsysteme, die wir im \S.\,136 des ersten Bandes kennen gelernt haben, und man kann also auch hier $\gamma$ als den Index von $n$ für den Modul $q^\varkappa$ bezeichnen.


% ===== scan leaf 79, printed 58 =====
% !PRINTED 58
\weberpage{58}
\pdfbookmark[2]{§ 15. Indices für eine Potenz von 2 als Modul}{bm18}

\begin{center}
\S.\,15.

\gesperrt{Indices für eine Potenz von 2 als Modul}.
\end{center}

Ein ähnlicher Satz muss nun für die Potenzen $2^\lambda$ der Primzahl 2 aufgestellt werden, die sich, wie schon vorhin bemerkt, anders verhält.

Ist $\lambda = 1$, so ist jede durch 2 nicht theilbare Zahl $n \equiv 1 \pmod{2}$. Ist $\lambda = 2$, so ist $-1$ als primitive Wurzel von 4 aufzufassen, denn jede ungerade Zahl $n$ genügt einer der Congruenzen

\[
n \equiv (-1)^\alpha \pmod{4},
\]

worin $\alpha = 0$ oder $= 1$ ist.

Aber schon für den Modul 8 existirt keine primitive Wurzel mehr, d.\,h. keine Zahl, durch deren Potenzen sich alle Zahlclassen ungerader Zahlen nach dem Modul 8 darstellen lassen. Denn ist $g$ irgend eine ungerade Zahl, so sind unter den Potenzen von $g$ höchstens die beiden $1, g$ nach dem Modul 8 verschieden, weil immer $g^2 \equiv 1 \pmod{8}$ ist. Nimmt man also $g$ nicht congruent 1 und nicht congruent $-1 \pmod{8}$, also $g = 3$ oder $= 5$, so ist jede ungerade Zahl einer der vier Zahlen

\[
(-1)^\alpha g^\beta \qquad \alpha = 0, 1, \quad \beta = 0, 1
\]

nach dem Modul 8 congruent.

Die Anzahl der Classen ungerader Zahlen nach dem Modul $2^\lambda$ ist $2^{\lambda-1}$. Nun ist aber für jede ungerade Zahl $g$, falls $\lambda > 2$ ist,

\weq{1}{g^{2^{\lambda-2}} \equiv 1 \pmod{2^\lambda}.}

Denn nehmen wir (1) als richtig an und setzen demgemäss

\[
g^{2^{\lambda-2}} = 1 + h\, 2^\lambda,
\]

und erheben ins Quadrat, so folgt:

\[
g^{2^{\lambda-1}} \equiv 1 \pmod{2^{\lambda+1}}.
\]

Ist also die Formel (1) für $\lambda$ richtig, so ist sie es auch für $\lambda + 1$, und da sie für $\lambda = 3$ gilt, so gilt sie allgemein. Daraus folgt, dass unter den Potenzen einer ungeraden Zahl $g$ höchstens $2^{\lambda-2}$ nach dem Modul $2^\lambda$ verschiedene vorkommen können.

Andererseits folgt aber leicht für $g = 5$, dass $5^{2^{\lambda-2}}$ die niedrigste Potenz ist, die nach dem Modul $2^\lambda$ mit der Einheit congruent wird. Denn ist $5^e$ die niedrigste Potenz, die der Ein-
%CONT>>


% ===== scan leaf 80, printed 59 =====
% !PRINTED 59
\weberpage{59}

%<<CONT
heit congruent ist, so ist $e$ nach (1) ein Theiler von $2^{\lambda-2}$, also eine Potenz von 2, und wenn $e < 2^{\lambda-2}$ wäre, so müsste

\[
5^{2^{\lambda-3}} \equiv 1 \pmod{2^\lambda}
\]

sein. Dies ist aber nicht möglich, denn es gilt für jedes $\lambda$, was grösser als 2 ist, die Formel

\[
5^{2^{\lambda-3}} = 1 + 2^{\lambda-1} h\ ^{1)}
\]

mit ungeradem $h$, eine Formel, die sich ebenso wie die Formel (1) durch vollständige Induction beweisen lässt. Es sind also, wenn $\lambda \geqq 3$ ist,
die $2^{\lambda-2}$ Potenzen

\weq{2}{1, 5, 5^2 \ldots 5^{2^{\lambda-2}-1}}

nach dem Modul $2^\lambda$ alle von einander verschieden. Nun ist eine Relation von der Form $5^h \equiv -5^\varkappa$ für den Modul 4, also um so mehr für jede höhere Potenz von 2, unmöglich, und folglich sind die Grössen

\weq{3}{-1, -5, -5^2 \ldots -5^{2^{\lambda-2}-1}}

nach dem Modul $2^\lambda$ sowohl unter einander als von den Grössen (2) verschieden, und da ihre gesammte Anzahl $2^{\lambda-1}$ beträgt, so ist jede ungerade Zahl einer und nur einer der Grössen (2), (3) nach dem Modul $2^\lambda$ congruent.

Setzen wir also

\weq{4}{a = 2, \quad b = 2^{\lambda-2},}

so erhalten wir folgenden Satz für $\lambda \geqq 3$:

\noindent\textbf{3.} \gesperrt{Für jede ungerade Zahl} $n$ \gesperrt{lässt sich ein nach dem Modulpaar} $a, b$ \gesperrt{völlig bestimmtes Zahlenpaar} $\alpha, \beta$ \gesperrt{angeben, so dass}

\[
n \equiv (-1)^\alpha 5^\beta \pmod{2^\lambda}
\]

\gesperrt{wird.}

Der Fall $\lambda = 2$ kann hierunter mit subsumirt werden, weil dann $b = 1$ wird und $\beta = 0$ gesetzt werden kann. Um auch den Fall $\lambda = 1$ mit darunter zu begreifen, der aber kein besonderes Interesse bietet, müsste man $a = 1, b = 1$ setzen.

\noindent\rule{4cm}{0.4pt}

\noindent $^{1)}$ Wollte man an Stelle der Zahl 5 die Zahl 3 als Basis nehmen, so würde diese Formel für $\lambda = 3$ noch nicht gültig sein, wohl aber für jedes grössere $\lambda$, und daher könnte 3 als Basis ebenso gut dienen, wie 5. Der Grund für die Bevorzugung der Basis 5 liegt darin, dass alle Potenzen dieser Basis $\equiv 1 \pmod{4}$ sind.


% ===== scan leaf 81, printed 60 =====
% !PRINTED 60
\weberpage{60}
\pdfbookmark[2]{§ 16. Die Gruppen der Zahlclassen nach einem zusammengesetzten Modul}{bm19}

\begin{center}
\S.\,16.

\gesperrt{Die Gruppen der Zahlclassen nach einem zusammengesetzten Modul}.
\end{center}

Aus der Verbindung der beiden Sätze 2. und 3. der beiden vorangegangenen Paragraphen ergiebt sich nun folgendes Resultat, durch welches die Aufgabe, die wir am Anfang des \S.\,14 gestellt haben, vollständig gelöst wird:

\noindent\textbf{4.} \gesperrt{Wenn der Modul}

\[
m = 2^\lambda q_1^{\varkappa_1} q_2^{\varkappa_2} \ldots
\]

\gesperrt{ist, und} $g_1, g_2 \ldots$ \gesperrt{primitive Wurzeln der Quadrate der Primzahlen} $q_1, q_2 \ldots$ \gesperrt{sind, wenn ferner}

\[
a = 2, \quad b = \tfrac12 \varphi(2^\lambda), \quad c_1 = \varphi(q_1^{\varkappa_1}), \quad c_2 = \varphi(q_2^{\varkappa_2}) \ldots
\]

\gesperrt{ist, so kann man für jede Zahl} $n$, \gesperrt{die zu} $m$ \gesperrt{relativ prim ist, ein System von Zahlen} $\alpha, \beta, \gamma_1, \gamma_2 \ldots$ \gesperrt{nach den Moduln} $a, b, c_1, c_2 \ldots$ \gesperrt{eindeutig bestimmen, die den Gleichungen}

\weq{1}{
\begin{aligned}
n &\equiv (-1)^\alpha 5^\beta \pmod{2^\lambda} \\
&\equiv g_1^{\gamma_1} \pmod{q_1^{\varkappa_1}} \\
&\equiv g_2^{\gamma_2} \pmod{q_2^{\varkappa_2}} \\
&\; \cdots\cdots\cdots\cdots\cdots\cdots
\end{aligned}
}

\gesperrt{genügen.}

\gesperrt{Die Zahlen $\alpha, \beta, \gamma_1, \gamma_2 \ldots$ heissen das System der Indices von $n$ für den Modul $m$.}

Hier ist zunächst $\lambda \geqq 3$ vorausgesetzt. Der Satz gilt aber auch für die übrigen Werthe von $\lambda$, wenn

\[
\text{für } \lambda = 0, 1, \quad a = 1, \quad b = 1
\]
\[
\text{\textquotedbl{}} \quad \lambda = 2, \qquad\quad a = 2, \quad b = 1
\]

gesetzt wird. Für $\lambda = 0, 1$ sind die Indices $\alpha$ und $\beta = 0$ zu setzen oder auch ganz wegzulassen, für $\lambda = 2$ fällt $\beta$ weg, während $\alpha$ gleich 0 oder gleich 1 sein kann.

Um nun also die Gruppe $\mathfrak{N}$ der Zahlclassen $N$ nach dem Modul $m$ durch eine Basis darzustellen, bestimme man die Zahlclassen

\weq{2}{A, B, C_1, C_2 \ldots,}

oder wenigstens Repräsentanten dieser Classen, was immer möglich ist, aus den Congruenzen
%CONT>>


% ===== scan leaf 82, printed 61 =====
% !PRINTED 61
\weberpage{61}
%<<CONT
\[
\begin{aligned}
A &\equiv -1 \pmod{2^{\lambda}} &&\equiv 1 \pmod{q_1^{\varkappa_1}} &&\equiv 1 \pmod{q_2^{\varkappa_2}} \ldots\\
B &\equiv \phantom{-}5\;'' &&\equiv 1\;'' &&\equiv 1\;'' \ldots\\
C_1 &\equiv \phantom{-}1\;'' &&\equiv g_1\;'' &&\equiv 1\;'' \ldots\\
C_2 &\equiv \phantom{-}1\;'' &&\equiv 1\;'' &&\equiv g_2\;'' \ldots
\end{aligned}
\]

\noindent
$\ldots\ldots\ldots\ldots\ldots\ldots\ldots\ldots\ldots\ldots\ldots\ldots\ldots\ldots\ldots\ldots\ldots\ldots\ldots\ldots\ldots\ldots$

und nach 4. erh\"alt man dann f\"ur jede Zahl $n$ unserer Gruppe

\weq{3}{n \equiv A^{\alpha} B^{\beta} C_1^{\gamma_1} C_2^{\gamma_2} \ldots \pmod{m}.}

Die $A, B, C_1, C_2 \ldots$ sind also die Elemente einer Basis der
Gruppe $\mathfrak{N}$ von den Graden $a, b, c_1, c_2 \ldots$

Wenn $m$ ungerade oder nur durch die erste Potenz von 2
theilbar ist, so fallen aus der Basis die beiden Elemente $A, B$
weg. Ist $m$ durch 4, aber nicht durch 8 theilbar, so f\"allt $B$ weg
und das Element $A$ vom zweiten Grade bleibt.

\gesperrt{5. Zerlegt man die Zahlen $a, b, c_1, c_2 \ldots$ in Potenzen
von einander verschiedener Primzahlen, so erh\"alt man die Invarianten der Gruppe $\mathfrak{N}$.}

Um die verschiedenen F\"alle zusammenzufassen, bezeichnen
wir die Elemente der Basis (2) mit

\weq{4}{C_{-1}, C_0, C_1 \ldots C_\mu.}

Hierin sollen $C_{-1}, C_0$ f\"ur $A$ und $B$ stehen und sind also
gleich 1 zu setzen, wenn $\lambda = 0$ oder $\lambda = 1$ ist. Wenn $\lambda = 2$
ist, so ist $B = C_{-1} = 1$, $A = C_0$, und wenn $\lambda > 1$ ist, $A = C_{-1}$,
$B = C_0$ zu setzen.
% [sic: Weber prints \lambda > 1; the preceding cases require \lambda > 2 here.]

Die Grade dieser Elemente seien mit

\weq{5}{c_{-1}, c_0, c_1 \ldots c_\mu}

bezeichnet, und die Indices einer Zahl aus $\mathfrak{N}$, die nach den
Gr\"ossen (5) als Moduln zu nehmen sind, mit

\weq{6}{\nu_{-1}, \nu_0, \nu_1 \ldots \nu_\mu.}

Ist $\lambda = 0$ oder 1, so haben $\nu_{-1}$ und $\nu_0$ nur den Werth 0,
ist $\lambda = 2$, so hat $\nu_{-1}$ den Werth 0 und $\nu_0$ den Werth 0 oder 1.
Ist $\lambda \geqq 3$, so hat $\nu_{-1}$ einen der beiden Werthe 0, 1, und $\nu_0$
einen der Werthe $0, 1, 2, \ldots c_0 - 1$; $\mu$ ist immer gleich der
Anzahl der von einander verschiedenen ungeraden Primzahlen,
die in $m$ aufgehen. Wir wollen, indem wir $\nu_{-1}, c_{-1}$ bei Seite
lassen, $\nu_0, \nu_1 \ldots \nu_\mu$ die den Primzahlen $2, q_1 \ldots q_\mu$ entsprechenden
\gesperrt{Indices} von $n$ und $c_0, c_1 \ldots c_\mu$ die denselben Primzahlen entsprechenden
\gesperrt{Indexmoduln} nennen. Diese Indexmoduln sind
\[
c_0 = \varphi(2^{\lambda-1}), \quad c_1 = \varphi(q_1^{\varkappa_1}), \ldots c_\mu = \varphi(q_\mu^{\varkappa_\mu}),
\]
%CONT>>


% ===== scan leaf 83, printed 62 =====
% !PRINTED 62
\weberpage{62}
%<<CONT
und nur wenn $\lambda = 2$ ist, ist $c_0 = 2$ und nicht $= 1$ zu setzen.

Wenn wir dann mit $\varepsilon_{-1}, \varepsilon_0, \varepsilon_1 \ldots \varepsilon_\mu$ ein System primitiver
Einheitswurzeln der Grade $c_{-1}, c_0, c_1 \ldots c_\mu$ bezeichnen, und mit
$\beta_{-1}, \beta_0, \beta_1 \ldots \beta_\mu$ die Indices einer Zahl $b$, so erhalten wir die
Charaktere der Gruppe $\mathfrak{N}$ in der Form

\weq{7}{\chi_b(n) = \varepsilon_{-1}^{\beta_{-1}\nu_{-1}} \varepsilon_0^{\beta_0\nu_0} \varepsilon_1^{\beta_1\nu_1} \ldots \varepsilon_\mu^{\beta_\mu\nu_\mu}.}

Darin ist $\varepsilon_{-1} = 1$, wenn $\lambda = 0, 1, 2$ ist; in den anderen
F\"allen ist $\varepsilon_{-1} = -1$, und $\varepsilon_1, \ldots \varepsilon_\mu$ sind primitive Einheitswurzeln
 der Grade
\[
c_1 = \varphi(q_1^{\varkappa_1}), \ldots c_\mu = \varphi(q_\mu^{\varkappa_\mu}).
\]

\begin{center}
\rule{3cm}{0.4pt}
\end{center}


% ===== scan leaf 84, printed 63 =====
\weberpage{63}% page number not printed on this page
\pdfbookmark[1]{Dritter Abschnitt. Die Gruppe der Kreistheilungskörper.}{bm20}
\pdfbookmark[2]{§ 17. Die Resolventen der Kreistheilungstheorie}{bm21}
% !PRINTED none
\begin{center}
\gesperrt{Dritter Abschnitt.}

\gesperrt{Die Gruppe der Kreistheilungsk\"orper.}

\rule{3cm}{0.4pt}

\S.\,17.

\gesperrt{Die Resolventen der Kreistheilungstheorie.}
\end{center}

Von den S\"atzen \"uber \gesperrt{Abel'sche} Gruppen machen wir eine
Anwendung auf die Kreistheilungstheorie f\"ur den Fall, dass der
Grad der Einheitswurzeln nicht eine Primzahl, sondern eine
h\"ohere Potenz einer Primzahl ist. Ist $q$ eine ungerade Primzahl
und $m = q^{\varkappa}$, $\varkappa > 1$, so nehmen wir eine primitive Congruenzwurzel
 $g$ von $m$ und setzen f\"ur jede durch $q$ nicht theilbare
 Zahl $n$
\weq{1}{n \equiv g^{\nu} \pmod{m}.}
Dann ist $\nu$ der \gesperrt{Index} von $n$. Durchl\"auft $n$ die Gruppe $\mathfrak{N}$
der durch $q$ nicht theilbaren Zahlclassen nach dem Modul $m$
vom Grade
\weq{2}{c = \varphi(q^{\varkappa}) = q^{\varkappa - 1}(q - 1),}
so durchl\"auft $\nu$ ein volles Restsystem nach dem Modul $c$. Wir
nehmen nun eine primitive $m^{\text{te}}$ Einheitswurzel $r$ und eine primitive
 $c^{\text{te}}$ Einheitswurzel $\varepsilon$ und bilden die \gesperrt{Lagrange'schen}
\gesperrt{Resolventen}
\weq{3}{(\varepsilon^{\beta}, r) = \sum^{n} \varepsilon^{\beta \nu} r^n,}
worin $\beta$ ein beliebiger Exponent ist. Es handelt sich um die
Frage, wann eine solche Resolvente verschwinden kann. Ist
$\varkappa = 1$, also $m$ eine Primzahl, so verschwindet sie, wie wir im
Bd.\ I, \S.\,169 gesehen haben, f\"ur keinen Werth von $\beta$. Im allgemeinen
 Falle eines beliebigen $\varkappa$ bedeute $n'$ eine beliebige Zahl
aus $\mathfrak{N}$. Dann durchl\"auft $n n'$ zugleich mit $n$ die ganze Gruppe $\mathfrak{N}$.
Man kann also in (3) $n$ durch $n n'$ ersetzen, wenn man zugleich
%CONT>>


% ===== scan leaf 85, printed 64 =====
% !PRINTED 64
\weberpage{64}
%<<CONT
$\nu$ durch $\nu + \nu'$ ersetzt, wenn $\nu'$ den Index von $n'$ bedeutet.
Dadurch folgt aus (3):
\[
\varepsilon^{-\beta \nu'} (\varepsilon^{\beta}, r) = \sum^{n} \varepsilon^{\beta \nu} r^{n n'},
\]
und wenn man also mit $r^{n'}$ multiplicirt und die Summe \"uber $n'$
nimmt:
\weq{4}{(\varepsilon^{-\beta}, r)(\varepsilon^{\beta}, r) = \sum^{n, n'} \varepsilon^{\beta \nu} r^{n'(n+1)}.}
Es ist also zun\"achst die Summe
\weq{5}{\sigma = \sum^{n'} r^{n'(n+1)}}
zu bestimmen. Die Summe $\sigma$ verschwindet aber immer, wenn
$n + 1$ nicht durch $q^{\varkappa - 1}$ theilbar ist (nach Bd.\ I, \S.\,133); denn
dann ist $r^{n+1}$ eine Einheitswurzel, deren Grad eine h\"ohere als
die erste Potenz von $q$ ist. $\sigma$ kann also nur dann von Null verschieden
 sein, wenn $n$ die Form hat
\weq{6}{n \equiv -1 + t q^{\varkappa - 1} \pmod{m},}
und dann wollen wir seinen Werth mit $\sigma_t$ bezeichnen. Wir erhalten
 alle Werthe von $n$ nach dem Modul $m$, die in der Form
(6) enthalten sind, wenn wir $t$ ein volles Restsystem nach dem
Modul $q$ durchlaufen lassen, also etwa
\weq{7}{t = 0, 1, 2 \ldots q - 1}
setzen. Die Summe $\sigma$ besteht aus $\varphi(m)$ Gliedern. Ist $t = 0$, so
wird jedes dieser Glieder $= 1$ und es folgt
\weq{8}{\sigma_0 = \varphi(m) = q^{\varkappa} - q^{\varkappa - 1};}
f\"ur jeden anderen Werth von $t$ ist $r^{n+1}$ eine primitive $q^{\text{te}}$ Einheitswurzel
, und in $\sigma_t$ kommt dann jede solche Einheitswurzel
$\varphi(m) : (q - 1) = q^{\varkappa - 1}$ mal vor. Die Summe aller primitiven
$q^{\text{ten}}$ Einheitswurzeln ist aber gleich $-1$ und also folgt f\"ur
$t = 1, 2 \ldots q - 1$
\weq{9}{\sigma_t = -q^{\varkappa - 1}.}
Damit ist die Summe $\sigma$ bestimmt.

Um aber den Werth der Summe in (4) daraus abzuleiten,
ist es noch n\"othig, den Index $\nu$ der Zahlen $n$ von der Form (6)
zu ermitteln.

F\"ur diese Zahlen ist aber
\[
g^{\nu} \equiv -1 \pmod{q^{\varkappa - 1}},
\]
und da nach dem \gesperrt{Fermat'schen} Satze [\S.\,14, (11)]
\weq{10}{g^{1/2 \varphi(q^{\varkappa})} \equiv -1 \pmod{q^{\varkappa}}}
%CONT>>


% ===== scan leaf 86, printed 65 =====
% !PRINTED 65
\weberpage{65}
%<<CONT
ist, so folgt
\[
g^\nu \equiv g^{\frac{1}{2}\varphi(q^\varkappa)} \ (\mathrm{mod}\ q^{\varkappa-1});
\]
also $\nu \equiv \tfrac{1}{2}\varphi(q^\varkappa)\ \bigl[\mathrm{mod}\ \varphi(q^{\varkappa-1})\bigr]$, da der Index einer Zahl für den Modul $q^{\varkappa-1}$ völlig bestimmt ist nach dem Modul $\varphi(q^{\varkappa-1})$. Es wird also
\weq{11}{\nu \equiv \tfrac{1}{2}\varphi(q^\varkappa) + \tau\varphi(q^{\varkappa-1})\ \bigl[\mathrm{mod}\ \varphi(q^\varkappa)\bigr],}
und hierin durchläuft nun, da $\varphi(q^\varkappa) = q\,\varphi(q^{\varkappa-1})$ ist, $\tau$ zugleich mit $t$ ein volles Restsystem nach dem Modul $q$. Dem Werthe $t=0$ entspricht der Werth $\tau=0$ wegen (10).

Es ist aber
\[
\varepsilon^{\frac{1}{2}\varphi(q^\varkappa)} = -1, \qquad \varepsilon^{\varphi(q^{\varkappa-1})} = \varrho,
\]
wenn $\varrho$ eine primitive $q^{\text{te}}$ Einheitswurzel ist, also nach (11)
\[
\varepsilon^\nu = -\varrho^\tau,
\]
und danach ergiebt sich aus (8) und (9):
\[
(\varepsilon^{-\beta}, r)(\varepsilon^{\beta}, r) = \sum \varepsilon^{\beta\nu}\sigma_t
\]
\[
= (-1)^\beta (q^\varkappa - q^{\varkappa-1}) - (-1)^\beta q^{\varkappa-1} \sum_{1,\,q-1}^{\tau} \varrho^{\tau\beta}.
\]

Nun hat die Summe $\sum\limits_{1,\,q-1}^{\tau} \varrho^{\tau\beta}$, wenn $\beta$ nicht durch $q$ theilbar ist, den Werth $-1$, und wenn $\beta$ durch $q$ theilbar ist, den Werth $q-1$, so dass man folgendes Resultat erhält:
\weq{12}{\begin{aligned}
(\varepsilon^{-\beta}, r)(\varepsilon^{\beta}, r) &= 0, & \beta &\equiv 0\ (\mathrm{mod}\ q) \\
&= (-1)^\beta q^\varkappa, & \beta\ \text{nicht} &\equiv 0\ (\mathrm{mod}\ q).
\end{aligned}}

Wenn also $\beta$ nicht durch $q$ theilbar ist, so kann von den Factoren $(\varepsilon^{-\beta}, r), (\varepsilon^{\beta}, r)$ keiner verschwinden; wenn aber $\beta$ durch $q$ theilbar ist, so verschwindet wenigstens einer der beiden Factoren. Dass sie dann beide verschwinden, zeigt die folgende directe Betrachtung der Summe.

Wenn $\beta$ durch $q$ theilbar ist, so ist für jedes ganzzahlige $t$
\[
\varepsilon^{\beta[\nu + t\varphi(q^{\varkappa-1})]} = \varepsilon^{\beta\nu}
\]
\[
g^{t\varphi(q^{\varkappa-1})} \equiv 1 + \tau q^{\varkappa-1}\ (\mathrm{mod}\ q^\varkappa),
\]
worin nun $\tau$ zugleich mit $t$ ein volles Restsystem nach dem Modul $q$ durchläuft. Wenn man also in (3) unter der Voraussetzung, dass $\beta$ durch $q$ theilbar sei, $\nu$ durch $\nu + t\varphi(q^{\varkappa-1})$, d.\,h. $n$ durch $n(1+\tau q^{\varkappa-1})$ ersetzt, so folgt
\[
(\varepsilon^\beta, r) = \sum^{n} \varepsilon^{\beta\nu} r^n r^{n q^{\varkappa-1}\tau}.
\]


% ===== scan leaf 87, printed 66 =====
% !PRINTED 66
\weberpage{66}
Diese Summe ist also von dem willkürlich anzunehmenden $\tau$ unabhängig. Summirt man hier von $\tau=0$ bis $\tau=q-1$, so ergiebt sich, da $\sum r^{n q^{\varkappa-1}\tau}=0$ ist:
\weq{13}{(\varepsilon^\beta, r) = 0.}
Wir haben also den Satz:

1.~\gesperrt{Ist der Grad der Einheitswurzel $r$ eine Potenz einer ungeraden Primzahl, so verschwindet die Resolvente $(\varepsilon^\beta, r)$ dann und nur dann, wenn $\beta$ durch $q$ theilbar ist.}

Wir haben noch den Fall zu betrachten, dass der Grad der Einheitswurzel $r$ eine Potenz von 2 ist. Wir setzen also
\[
m = 2^\lambda
\]
und nehmen zunächst $\lambda \geqq 3$ an. Dann können wir für jede ungerade Zahl $n$ ein Indexpaar $\nu, \nu_1$ aus den Congruenzen
\[
n \equiv (-1)^{\nu_1} 5^\nu\ (\mathrm{mod}\ m)
\]
bestimmen, und zwar ist $\nu_1$ nach dem Modul 2, $\nu$ nach dem Modul $2^{\lambda-2}$ bestimmt. Unter den Resolventen verstehen wir in diesem Falle die Summen
\weq{14}{\bigl((-1)^{\beta_1}, \Theta^\beta, r\bigr) = \sum^{n} (-1)^{\beta_1 \nu_1} \Theta^{\beta \nu} r^n,}
worin $\Theta$ eine primitive Einheitswurzel vom Grade $2^{\lambda-2}$ bedeutet. Nun verfahren wir ganz ähnlich wie vorher. Wenn wir in (14) $n$ durch $nn'$ ersetzen und mit $\nu_1', \nu'$ die Indices von $n'$ bezeichnen, so ergiebt sich
\weq{15}{(-1)^{-\beta_1 \nu_1'} \Theta^{-\beta \nu'} \bigl((-1)^{\beta_1}, \Theta^\beta, r\bigr) = \sum^{n} (-1)^{\beta_1 \nu_1} \Theta^{\beta \nu} r^{n n'},}
und daraus durch Multiplication mit $r^{n'}$ und Summation nach $n'$
\weq{16}{\bigl((-1)^{-\beta_1}, \Theta^{-\beta}, r\bigr)\bigl((-1)^{\beta_1}, \Theta^\beta, r\bigr) = \sum^{n n'} (-1)^{\beta_1 \nu_1} \Theta^{\beta \nu} r^{n'(n+1)}.}
Nun ist aber aus den oben angeführten Gründen
\[
\begin{aligned}
\sum^{n'} r^{n'(n+1)} &= 0, & &\text{wenn } n+1 \text{ nicht durch } 2^{\lambda-1} \text{ theilbar ist,} \\
&= 2^{\lambda-1} & &\text{für } n+1 \equiv 0\ (\mathrm{mod}\ 2^\lambda),\ \nu=0,\ \nu_1=1, \\
&= -2^{\lambda-1} & &\text{für } n+1 \equiv 2^{\lambda-1}\ (\mathrm{mod}\ 2^\lambda),\ \nu=2^{\lambda-3},\ \nu_1=1,
\end{aligned}
\]
und danach giebt (16)
\weq{17}{\begin{aligned}
\bigl((-1)^{-\beta_1}, \Theta^{-\beta}, r\bigr)\bigl((-1)^{\beta_1}, \Theta^\beta, r\bigr) &= (-1)^{\beta_1} 2^\lambda, & \beta &\equiv 1\ (\mathrm{mod}\ 2) \\
&= 0, & \beta &\equiv 0\ (\mathrm{mod}\ 2).
\end{aligned}}


% ===== scan leaf 88, printed 67 =====
% !PRINTED 67
\weberpage{67}
\pdfbookmark[2]{§ 18. Kreistheilungskörper}{bm22}
Dass im Falle eines geraden $\beta$ jeder der beiden Factoren verschwindet, sieht man, wenn man in (15)
\[
n' = 1 + 2^{\lambda-1},
\]
also
\[
r^{n'} = -r, \qquad \nu_1' = 0, \qquad \nu' = 2^{\lambda-3}
\]
setzt. Dann giebt diese Formel, da $n$ ungerade ist:
\[
\bigl((-1)^{\beta_1}, \Theta^\beta, r\bigr) = -\Theta^{\beta 2^{\lambda-3}} \bigl((-1)^{\beta_1}, \Theta^\beta, r\bigr).
\]
Bei geradem $\beta$ ist aber $\Theta^{\beta 2^{\lambda-3}} = +1$ und folglich:
\weq{18}{\bigl((-1)^{\beta_1}, \Theta^\beta, r\bigr) = 0;}
also:

2.~\gesperrt{Ist der Grad der Einheitswurzel $r$ eine Potenz von 2 und grösser als 4, so verschwindet die Resolvente $\bigl((-1)^{\beta_1}, \Theta^\beta, r\bigr)$ dann und nur dann, wenn $\beta$ gerade ist.}

Im Falle $m=4$, also $r=i$, hat man nur die zwei Resolventen $i+i^3, i-i^3$, die wir unter der Bezeichnung $((-1)^\beta, i)$, $\beta=0,1$ zusammenfassen können, und es ist $((-1)^\beta, i)$ dann und nur dann $=0$, wenn $\beta=0$ ist.

\begin{center}
\S.\,18.\\
\gesperrt{Kreistheilungsk\"orper.}
\end{center}

Die Theorie der \gesperrt{Abel'schen} Gruppen eröffnet uns einen tieferen Einblick in die Theorie der Einheitswurzeln und der daraus entspringenden algebraischen Zahlen.

Es möge jetzt $m$ irgend eine ganze positive Zahl sein, die in ihre Primfactoren zerlegt sei:
\weq{1}{m = 2^\lambda q_1^{\varkappa_1} q_2^{\varkappa_2} \cdots,}
und es sei $r$ eine primitive $m^{\text{te}}$ Einheitswurzel. Den Fall $\lambda=1$ können wir ein- für allemal von unserer Betrachtung ausschliessen; denn wenn $m$ ungerade ist und $r$ die primitiven $m^{\text{ten}}$ Einheitswurzeln durchläuft, so kommen darunter keine zwei entgegengesetzte vor, und $-r$ durchläuft die primitiven $2m^{\text{ten}}$ Einheitswurzeln.

$r$ ist die Wurzel einer ganzzahligen irreduciblen \gesperrt{Abel'schen} Gleichung vom Grade
\weq{2}{\nu = \varphi(m),}
wie wir im \S.\,134 des ersten Bandes nachgewiesen haben.


% ===== scan leaf 89, printed 68 =====
% !PRINTED 68
\weberpage{68}
Der Inbegriff aller rationalen Functionen von $r$ mit rationalen Zahlen als Co\"efficienten ist also ein Zahlkörper $\Omega(r)$ vom Grade $\nu$, den wir einen \gesperrt{Kreistheilungskörper} nennen. Wir wollen aber den Begriff des Kreistheilungskörpers noch etwas allgemeiner fassen und darunter jeden Körper verstehen, dessen Zahlen lauter rationale Functionen irgend welcher Einheitswurzeln sind.

Den Körper $\nu^{\text{ten}}$ Grades $\Omega(r)$, der aus \gesperrt{allen} rationalen Functionen einer $m^{\text{ten}}$ Einheitswurzel $r$ besteht, nennen wir zur genaueren Unterscheidung den \gesperrt{vollen Kreistheilungskörper} der Ordnung $m$ und bezeichnen ihn mit $\Omega_m$.

Beliebige Einheitswurzeln $r, r', r'' \ldots$ beliebiger Grade $m, m', m'' \ldots$ kann man immer auffassen als Potenzen einer und derselben Einheitswurzel $\varrho$, deren Grad das kleinste gemeinschaftliche Vielfache von $m, m', m'' \ldots$ ist. Demnach ist ein Kreistheilungskörper, der nur rationale Functionen von $r, r', r'' \ldots$ enthält, ein Theiler des vollen Kreistheilungskörpers $\Omega(\varrho)$, und wir bekommen also alle überhaupt existirenden Kreistheilungskörper, wenn wir die sämmtlichen Divisoren aller vollen Kreistheilungskörper aufsuchen.

Die \gesperrt{Galois'sche} Gruppe des Körpers $\Omega_m$ besteht aus den sämmtlichen Substitutionen
\weq{3}{(r, r^n),}
wenn $n$ jede nach dem Modul $m$ genommene relative Primzahl zu $m$ bedeutet. Denn die Kreistheilungsgleichung $\nu^{\text{ten}}$ Grades, deren Wurzeln die $r^n$ sind, ist eine Normalgleichung und also ihre eigene \gesperrt{Galois'sche} Resolvente. Da, wenn $a, b$ zwei dieser Zahlen $n$ sind,
\[
(r, r^a)(r, r^b) = (r, r^{ab})
\]
ist, so ist diese Gruppe isomorph mit der Gruppe $\mathfrak{N}$ aller Zahlclassen $N$ der zu $m$ theilerfremden Zahlen, die wir im vorigen Paragraphen betrachtet haben.

Ist $\mathfrak{A}$ ein Theiler von $\mathfrak{N}$, und $\varrho$ eine zu $\mathfrak{A}$ gehörige Function aus $\Omega_m$, so ist der Inbegriff der rationalen Functionen von $\varrho$, $\Omega(\varrho)$, ein in $\Omega_m$ enthaltener Körper.

Ist umgekehrt $\Omega$ ein Theiler von $\Omega_m$ und $\varrho$ eine primitive Zahl des Körpers $\Omega$, also auch eine Zahl in $\Omega_m$, so kann $\Omega$ als der Inbegriff der rationalen Functionen von $\varrho$ dargestellt und mit $\Omega(\varrho)$ bezeichnet werden.


% ===== scan leaf 90, printed 69 =====
% !PRINTED 69
\weberpage{69}
Diese Function $\varrho$ gehört dann zu einer gewissen Gruppe $\mathfrak{A}$,
die ein Theiler von $\mathfrak{N}$ ist. Wir nennen auch die Gruppe $\mathfrak{A}$ und
den Körper $\Omega(\varrho)$ \gesperrt{zusammengehörig} und ziehen daraus den
Satz:

\gesperrt{1. Zu jedem Theiler $\mathfrak{A}$ von $\mathfrak{N}$ gehört ein gewisser
in $\Omega_m$ enthaltener Kreistheilungskörper $\Omega(\varrho)$, und
umgekehrt gehört zu jedem Theiler $\Omega(\varrho)$ von $\Omega_m$
ein gewisser Theiler $\mathfrak{A}$ von $\mathfrak{N}$ als Gruppe, in dem
Sinne, dass, wenn $a$ eine Zahl aus $\mathfrak{A}$ ist, alle
Zahlen des Körpers $\Omega(\varrho)$ die Permutationen $(r,r^a)$
gestatten, und dass umgekehrt jede Zahl in $\Omega_m$,
die diese Permutationen gestattet, in $\Omega(\varrho)$ enthalten ist.}

Die Galois'sche Gruppe eines solchen Körpers $\Omega(\varrho)$ erhalten wir nach Bd. I, §. 156, wenn wir $\mathfrak{N}$ in das System der
Nebengruppen zerlegen:
\[
\mathfrak{N} = \mathfrak{A} + \mathfrak{A}_1 + \mathfrak{A}_2 + \cdots,
\]
und die unter den Nebengruppen $\mathfrak{A}, \mathfrak{A}_1, \mathfrak{A}_2 \ldots$ durch Composition
mit den Elementen von $\mathfrak{N}$ hervorgerufenen Permutationen aufsuchen. Diese Gruppe ist aber nach §. 4 isomorph mit der
Gruppe $\mathfrak{N}/\mathfrak{A}$, also auch isomorph mit der zu $\mathfrak{A}$ reciproken Gruppe
(§. 12), die wir mit $\mathfrak{B}$ bezeichnen wollen. Ist $a$ der Grad von $\mathfrak{A}$
und $b$ der von $\mathfrak{B}$, so ist $ab = \nu$, und $\varrho$ genügt einer irreduciblen
Abel'schen Gleichung vom Grade $b$.

\gesperrt{2. Wir bekommen also alle Kreistheilungskörper,
wenn wir zu jedem Modul $m$ die sämmtlichen
Divisoren $\mathfrak{A}$ der Gruppe $\mathfrak{N}$ bilden, zu jeder dieser
Gruppen $\mathfrak{A}$ eine zugehörige Function $\varrho$ suchen
und daraus die Körper $\Omega(\varrho)$ ableiten.}

Es ist aber noch die Frage, ob bei diesem Processe ein und
derselbe Körper $\Omega$ mehrmals auftreten kann, wodurch wir auf
die Untersuchung der gemeinschaftlichen Theiler zweier Kreistheilungskörper geführt werden.

Nach der oben gegebenen Definition ist wohl zu unterscheiden zwischen der Gruppe, zu der ein Körper $\Omega$ \gesperrt{gehört},
und der Galois'schen Gruppe des Körpers; beide Gruppen sind
zu einander reciprok. So gehört der Körper $\Omega_m$ selbst zur
Einheitsgruppe, während seine Galois'sche Gruppe $\mathfrak{N}$ ist,


% ===== scan leaf 91, printed 70 =====
% !PRINTED 70
\weberpage{70}
Es gilt nun der Satz:

\gesperrt{3. Sind $\Omega'$ und $\Omega''$ zwei Theiler von $\Omega_m$, die zu den
Gruppen $\mathfrak{A}'$ und $\mathfrak{A}''$ gehören, so ist $\Omega''$ dann und
nur dann ein Theiler von $\Omega'$, wenn $\mathfrak{A}'$ ein Theiler
von $\mathfrak{A}''$ ist.}

Denn wenn $\Omega''$ ein Theiler von $\Omega'$ ist, so müssen alle Zahlen
von $\Omega''$ die Substitutionen der Gruppe $\mathfrak{A}'$, zu der $\Omega'$ gehört, gestatten, also muss $\mathfrak{A}'$ in $\mathfrak{A}''$ enthalten sein. Und wenn umgekehrt $\mathfrak{A}'$ in $\mathfrak{A}''$
enthalten ist, so gestatten alle Zahlen von $\Omega''$
die sämmtlichen Substitutionen von $\mathfrak{A}'$ und sind also rational
durch eine primitive Zahl des Körpers $\Omega'$ darstellbar; d. h. $\Omega''$
ist in $\Omega'$ enthalten (§. 144 des ersten Bandes).

Der Körper $\Omega_m$ enthält als Theiler alle Körper $\Omega_{m_1}$, bei
denen $m_1$ ein Theiler von $m$ ist; denn $\Omega_{m_1}$ besteht aus allen
rationalen Functionen von $r^{\frac{m}{m_1}}$. Nun ist dann und nur dann
\[
r^{a\frac{m}{m_1}} = r^{\frac{m}{m_1}},
\]
wenn
\weq{4}{a \equiv 1\,(\mathrm{mod}\ m_1).}

Die Zahlen $a$ aus $\mathfrak{N}$, die der Congruenz (4) genügen, bilden
also die Gruppe, zu der der Körper $\Omega_{m_1}$ gehört. Wir wollen
diese Gruppe mit $\mathfrak{A}_{m_1}$ bezeichnen und symbolisch
\weq{5}{\mathfrak{A}_{m_1} \equiv 1\,(\mathrm{mod}\ m_1).}

setzen. Da der Grad des Körpers $\Omega_{m_1}$ gleich $\varphi(m_1)$ ist, was
zugleich der Index des Theilers $\mathfrak{A}_{m_1}$ von $\mathfrak{N}$ ist, so ist der Grad
von $\mathfrak{A}_{m_1}$ gleich $\varphi(m):\varphi(m_1)$.

Ist $m = m_1 m_2$ und $m_1$ relativ prim zu $m_2$, so erhält man
alle Zahlen von $\mathfrak{A}_{m_1}$ aus $a = 1 + x m_1$, wenn man $x$ so bestimmt,
dass $a$ nach dem Modul $m_2$ mit allen relativen Primzahlen zu
$m_2$ congruent wird. Daraus ergiebt sich in Uebereinstimmung
mit der obigen allgemeinen Bestimmung der Grad von $\mathfrak{A}_{m_1}$ für
diesen Fall $=\varphi(m_2)$.

Sind nun aber $m_1$ und $m_2$ irgend zwei Theiler von $m$, und
$\mu$ der grösste gemeinschaftliche Theiler von $m_1$ und $m_2$, so wird
der grösste gemeinschaftliche Theiler der beiden Körper $\Omega_{m_1}, \Omega_{m_2}$
zu der Gruppe gehören, die man als kleinstes gemeinschaftliches
Vielfaches der beiden Gruppen
\weq{6}{\mathfrak{A}_{m_1} \equiv 1\,(\mathrm{mod}\ m_1),\quad \mathfrak{A}_{m_2} \equiv 1\,(\mathrm{mod}\ m_2)}
%CONT>>


% ===== scan leaf 92, printed 71 =====
% !PRINTED 71
\weberpage{71}
%<<CONT
erhält. Das ist aber die Gruppe
\weq{7}{\mathfrak{A}_\mu \equiv 1\,(\mathrm{mod}\ \mu),}

d. h. die Gruppe, die aus allen den Zahlen von $\mathfrak{N}$ besteht, die
der Congruenz $a \equiv 1\,(\mathrm{mod}\ \mu)$ genügen. Denn zunächst ist klar,
dass $\mathfrak{A}_{m_1}$ und $\mathfrak{A}_{m_2}$ Divisoren von $\mathfrak{A}_\mu$ sind. Also ist auch das
kleinste gemeinschaftliche Vielfache von $\mathfrak{A}_{m_1}$ und $\mathfrak{A}_{m_2}$ in $\mathfrak{A}_\mu$
enthalten. Ist aber andererseits $\alpha = 1 + \xi\mu$ irgend eine Zahl in $\mathfrak{A}_\mu$, so
kann man die Zahlen $a_1 = 1 + x_1 m_1$ in $\mathfrak{A}_{m_1}$ und $a_2 = 1 + x_2 m_2$
in $\mathfrak{A}_{m_2}$ so bestimmen, dass $\alpha \equiv a_1 a_2\,(\mathrm{mod}\ m)$ wird. Man hat
nur $\xi\mu = x_1 m_1 + x_2 m_2$ zu setzen, was nach Bd. I, §. 118 immer
möglich ist. Es ist also auch $\mathfrak{A}_\mu$ in dem kleinsten gemeinschaftlichen Vielfachen von $\mathfrak{A}_{m_1}$ und $\mathfrak{A}_{m_2}$ enthalten, welches demnach
$\mathfrak{A}_\mu$ selbst ist. Daraus also der Satz:

\gesperrt{4. Sind $m_1$ und $m_2$ irgend zwei natürliche Zahlen
und $\mu$ der grösste gemeinschaftliche Theiler von
$m_1$ und $m_2$, so ist der Körper $\Omega_\mu$ der grösste
gemeinschaftliche Theiler von $\Omega_{m_1}$ und $\Omega_{m_2}$ und
$\Omega_{\frac{m_1 m_2}{\mu}}$ ihr kleinstes gemeinsames Multiplum.}

Daraus folgt nun: wenn zwei volle Kreistheilungskörper $\Omega_m$
und $\Omega_{m'}$ einen gemeinsamen Theiler $\Omega$ haben, und $m'$ ist kleiner
als $m$, so muss es einen echten Theiler $d$ von $m$ geben, so dass
$\Omega$ auch ein Theiler von $\Omega_d$ ist.

Wir wollen einen Theiler von $\Omega_m$ \gesperrt{primär} nennen, wenn er
nicht zugleich in einem vollen Kreistheilungskörper $\Omega_{m'}$ von
niedrigerem $m'$ enthalten ist. Dann folgt also, dass man alle
nicht primären Theiler von $\Omega_m$ erhält, wenn man in $\Omega_d$ den
Index $d$ alle echten Theiler von $m$ durchlaufen lässt und die
Theiler von $\Omega_d$ aufsucht.

Bezeichnen wir nun mit $q_1, q_2, \ldots q_\tau$ die sämmtlichen in $m$
aufgehenden verschiedenen Primzahlen (2 eingeschlossen) und
setzen
\weq{8}{m = q_1 m_1 = q_2 m_2 = \cdots = q_\tau m_\tau,}

so ist jedes $d$ Theiler von einem der $m_1, m_2, \ldots m_\tau$, und wir erhalten die nicht primären Theiler von $\Omega_m$, wenn wir alle Theiler
der Körper
\weq{9}{\Omega_{m_1},\ \Omega_{m_2},\ \ldots\ \Omega_{m_\tau}}

aufsuchen. Diese Körper gehören aber zu der Gruppe
\weq{10}{Q_1 \equiv 1\,(\mathrm{mod}\ m_1),\ \ Q_2 \equiv 1\,(\mathrm{mod}\ m_2),\ \ldots\ Q_\tau \equiv 1\,(\mathrm{mod}\ m_\tau),}
%CONT>>


% ===== scan leaf 93, printed 72 =====
% !PRINTED 72
\weberpage{72}
\pdfbookmark[2]{§ 19. Primäre und nicht primäre Theiler der Gruppe N}{bm23}
%<<CONT
und also wird nach 3. ein Theiler $\Omega$ von $\Omega_m$ dann und nur
dann nicht primär sein, wenn in der zu $\Omega$ gehörigen Gruppe $\mathfrak{A}$
eine der Gruppen $Q_1, Q_2, \ldots Q_\tau$ enthalten ist. Wenn wir also
solche Theiler der Gruppe $\mathfrak{N}$, die keine der Gruppen $Q_1, Q_2 \ldots Q_\tau$
als Theiler enthalten, primäre Theiler nennen, so können wir
den Satz aussprechen:

\gesperrt{5. Um alle primären Theiler von $\Omega_m$ zu erhalten,
hat man alle primären Theiler der Gruppe $\mathfrak{N}$
aufzusuchen und die zugehörigen Körper zu
bilden.}

\gesperrt{6. Wenn man aber alle primären Theiler aller
vollen Kreistheilungskörper aufstellt, so erhält
man jeden Kreistheilungskörper, und jeden
nur einmal, und zwar jeden dargestellt durch
Einheitswurzeln möglichst niedrigen Grades.}

Der Körper $\Omega_{2m}$ hat, wenn $m$ ungerade ist, gar keinen primären Theiler und ist mit $\Omega_m$ identisch. In allen anderen Fällen
ist $\Omega_m$ wenigstens sein eigener primärer Theiler.

Es ist für das Folgende noch erforderlich, dass wir uns
über den Umfang und die Constitution der Gruppen $Q$ klar
werden.

\begin{center}
\S.\,19.

\gesperrt{Primäre und nicht primäre Theiler der Gruppe $\mathfrak{N}$.}
\end{center}

Wir bezeichnen mit $q$ irgend eine der in $m$ aufgehenden
Primzahlen und setzen $m = q m'$. Dann betrachten wir die
Gruppe
\[
Q \equiv 1\,(\mathrm{mod}\ m').
\]

Um die Bedingungen für eine Zahl in $Q$ zu ermitteln, wenden
wir die Darstellung der Gruppe $\mathfrak{N}$ durch eine Basis und die
Bezeichnungsweise der Indices an, wie wir sie im §. 16 eingeführt
und erklärt haben, und bezeichnen danach die zu $Q$ reciproke
Gruppe mit $P$.

Die Indices einer Zahl in $Q$ bezeichnen wir mit
\[
\gamma_{-1}, \gamma_0, \gamma_1, \ldots \gamma_\mu,
\]
und einer Zahl in $P$ mit
\[
\delta_{-1}, \delta_0, \delta_1, \ldots \delta_\mu.
\]


% ===== scan leaf 94, printed 73 =====
% !PRINTED 73
\weberpage{73}

Dann müssen die $\gamma$ der Bedingung genügen:

\weq{1}{C_{-1}^{\gamma_{-1}} C_0^{\gamma_0} C_1^{\gamma_1} \dots C_\mu^{\gamma_\mu} \equiv 1 \, (\mathrm{mod}\ m').}

Diese Bedingung fordert, dass alle Indices mit Ausnahme des der Primzahl $q$
entsprechenden, den wir mit $\gamma$ bezeichnen, Null sein müssen. In Bezug auf
$\gamma$ ist aber zu unterscheiden, ob $q$ noch in $m'$ aufgeht oder nicht, d.\,h.
ob $q$ ein mehrfacher oder nur ein einfacher Factor von $m$ ist. Ist $q$ nicht
mehr in $m'$ enthalten, so enthält die Congruenz (1) gar keine Beschränkung für
$\gamma$, und $\gamma$ kann jeden Werth nach dem Modul $c$, d.\,h. jeden Werth
$0, 1, \dots q - 2$ annehmen.

Ist aber $q$ noch in $m'$ enthalten, also $m$ durch $q^\varkappa$ theilbar, und $\varkappa > 1$,
so fordert die Bedingung (1):
\[
C^\gamma \equiv 1 \, (\mathrm{mod}\ q^{\varkappa-1}),
\]
also
\weq{2}{\gamma \equiv 0 \, \left(\mathrm{mod}\ \dfrac{c}{q}\right),}
und $\gamma$ kann also jeden der Werthe
\[
0, \dfrac{c}{q}, \dfrac{2c}{q}, \dots \dfrac{(q-1)c}{q}
\]
erhalten. Im ersten Falle ist der Grad der Gruppe $Q$ gleich $q - 1$, im
zweiten gleich $q$.

Die Indices $\delta$ einer Zahl aus der Gruppe $P$ erhält man nach \S.\,12, 7.
und \S.\,16, (7) aus der Bedingung, dass für alle zulässigen $\gamma$
\weq{3}{\varepsilon_{-1}^{\gamma_{-1}\delta_{-1}} \varepsilon_0^{\gamma_0\delta_0} \varepsilon_1^{\gamma_1\delta_1} \dots \varepsilon_\mu^{\gamma_\mu\delta_\mu} = 1}
sein soll. Nach dem, was eben über die Indices $\gamma$ bewiesen ist, fordert
aber (3) nur das eine, dass der der Primzahl $q$ entsprechende Index $\delta$
durch $q$ oder durch $q - 1$ theilbar sein soll, je nachdem $q$ mehrmals oder
nur einmal in $m$ aufgeht. Wir heben also den Satz hervor:

7. \gesperrt{Die zu $Q$ reciproke Gruppe $P$ ist dadurch charakterisirt, dass
der $q$ entsprechende Index $\delta$ aller seiner Zahlen durch $q$ oder durch
$q - 1$ theilbar ist, je nachdem $q$ mehrmals oder nur einmal in $m$ aufgeht.}

Und daraus:

8. \gesperrt{Ein nicht primärer Theiler $\mathfrak{A}$ von $\mathfrak{N}$ ist
dadurch charakterisirt, dass in der reciproken Gruppe $\mathfrak{B}$}
%CONT>>


% ===== scan leaf 95, printed 74 =====
% !PRINTED 74
\weberpage{74}
\pdfbookmark[2]{§ 20. Die Kreistheilungsperioden}{bm24}
%<<CONT
\gesperrt{der einer Primzahl $q$ entsprechende Index $\beta$ aller Zahlen $b$ durch $q$ oder durch $q - 1$ theilbar ist, je nachdem $q$ mehrmals oder nur einmal unter den Primfactoren von $m$ vorkommt.}

\begin{center}
\S.\,20.

\gesperrt{Die Kreistheilungsperioden.}
\end{center}

Als die einfachsten Functionen, durch die man die Kreistheilungskörper darzustellen versuchen kann, bieten sich die Kreistheilungsperioden dar, die eine unmittelbare Verallgemeinerung der im \S.\,167 des ersten Bandes betrachteten Gauss'schen Perioden sind. Wir verstehen darunter Folgendes.

Es bedeute $\mathfrak{A}$ einen Theiler der Gruppe $\mathfrak{N}$ vom Index $e$ und $a$ durchlaufe die Zahlen von $\mathfrak{A}$. Ist $r$ eine primitive $m^{\text{te}}$ Einheitswurzel, so heisst die Summe
\weq{1}{\eta = \overset{a}{\Sigma}\, r^a}
\gesperrt{eine zu der Gruppe $\mathfrak{A}$ gehörige Kreistheilungsperiode} vom Index $e$.

Machen wir in (1) eine der Substitutionen $(r, r^a)$, so bleibt $\eta$ ungeändert, wie unmittelbar aus der Gruppeneigenschaft der $a$ folgt.

Um $\mathfrak{N}$ in die zu $\mathfrak{A}$ gehörigen Nebengruppen zu zerlegen, müssen wir die Zahlen $1, n_1, n_2, \ldots n_{e-1}$ aus $\mathfrak{N}$ passend auswählen, dass man
\weq{2}{\mathfrak{N} = \mathfrak{A} + \mathfrak{A} n_1 + \mathfrak{A} n_2 + \cdots + \mathfrak{A} n_{e-1}}
erhält. Machen wir dann in $\eta$ die Substitutionen
\weq{3}{(r, r), (r, r^{n_1}), \ldots (r, r^{n_{e-1}}),}
so geht $\eta$ in die conjugirten Perioden
\weq{4}{\eta, \eta_1, \ldots \eta_{e-1}}
über, und wenn diese alle von einander verschieden sind, so gehört $\eta$ zur Gruppe $\mathfrak{A}$. Der zu $\mathfrak{A}$ gehörige Körper $\Omega(\eta)$ besteht aus allen rationalen Functionen von $\eta$, und die Zahlen $\eta_1 \ldots \eta_{e-1}$ gehören alle zu derselben Gruppe $\mathfrak{A}$.

Wenn aber die Grössen (4) nicht alle von einander verschieden sind, so gehört die Zahl $\eta$ nicht zu der Gruppe $\mathfrak{A}$, sondern zu einer umfassenderen Gruppe $\mathfrak{A}'$, von der $\mathfrak{A}$ ein Theiler
%CONT>>


% ===== scan leaf 96, printed 75 =====
% !PRINTED 75
\weberpage{75}
%<<CONT
ist. Um die Bedingungen f\"ur diesen Fall zu ermitteln, erweitern
wir den Begriff der \gesperrt{Resolventen}, wie wir ihn schon im \S.\,17
betrachtet haben, noch etwas.

Wir lassen $n$ die Reihe der Zahlen der Gruppe $\mathfrak{N}$ durchlaufen
 und bezeichnen mit $\chi(n)$ einen der Charaktere dieser
Gruppe. Die erweiterten Resolventen sind dann, wenn $r$ eine
$m^{\text{te}}$ Einheitswurzel ist:
\weq{5}{(\chi, r) = \sum^{n} \chi(n) r^n.}
Verstehen wir unter $\mathfrak{B}$ die zu $\mathfrak{A}$ reciproke Gruppe und lassen
$b$ die Zahlen von $\mathfrak{B}$ durchlaufen, so giebt es nach \S.\,12, 7. eine
dem Grade von $\mathfrak{B}$ gleiche Anzahl von Charakteren $\chi_b$, die dadurch
 ausgezeichnet sind, dass
\[
\chi_b(a) = 1
\]
ist, f\"ur jede Zahl $a$ aus der Gruppe $\mathfrak{A}$. Daraus folgt dann nach
der Definition der Charaktere \S.\,11, (4) f\"ur jedes $n$:
\weq{6}{\chi_b(a n) = \chi_b(n),}
d. h. der Charakter $\chi_b$ hat f\"ur alle Zahlen einer jeden Nebengruppe
 $\mathfrak{A} n_1, \mathfrak{A} n_2 \ldots$ einen und denselben Werth. Demnach wird
die Resolvente $(\chi_b, r)$ nur von den Perioden $\eta$ abh\"angen, und
den Ausdruck erhalten:
\weq{7}{(\chi_b, r) = \eta + \chi_b(n_1) \eta_1 + \cdots + \chi_b(n_{e-1}) \eta_{e-1}.}
Wenn $\eta$ zu einer Gruppe $\mathfrak{A}'$ geh\"ort, so geh\"oren die conjugirten
 Zahlen $\eta_1, \eta_2, \ldots \eta_{e-1}$ zu derselben Gruppe (weil $\mathfrak{A}'$ ein
Normaltheiler von $\mathfrak{N}$ ist, Bd.\ I, \S.\,154). Wenn also $a'$ irgend
eine Zahl aus $\mathfrak{A}'$ ist, so bleiben die Gr\"ossen $\eta, \eta_1, \eta_2, \ldots \eta_{e-1}$
durch die Substitution $(r, r^{a'})$ unge\"andert und nach (7) ist auch
\weq{8}{(\chi_b, r) = (\chi_b, r^{a'}).}
Andererseits erh\"alt man aus (5), wenn man bedenkt, dass
$a' n$ zugleich mit $n$ die ganze Gruppe $\mathfrak{N}$ durchl\"auft,
\[
\begin{aligned}
(\chi_b, r) &= \sum^{n} \chi_b(n a') r^{n a'} = \chi_b(a') \sum^{n} \chi_b(n) r^{n a'}\\
&= \chi_b(a') (\chi_b, r^{a'}),
\end{aligned}
\]
also nach (8):
\weq{9}{(\chi_b, r) = \chi_b(a') (\chi_b, r).}
Ist $\mathfrak{A}$ nicht mit $\mathfrak{A}'$ identisch, sondern ein echter Theiler
von $\mathfrak{A}'$, so ist die zu $\mathfrak{A}'$ reciproke Gruppe $\mathfrak{B}'$ ein echter Theiler
von $\mathfrak{B}$ (nach \S.\,12, 8.); wenn also $b$ eine Zahl in $\mathfrak{B}$ ist, die nicht
%CONT>>


% ===== scan leaf 97, printed 76 =====
% !PRINTED 76
\weberpage{76}
%<<CONT
zugleich in $\mathfrak{B}'$ enthalten ist, so kann man $a'$ so wählen, dass $\chi_b(a')$ nicht $= 1$ ist. Dann folgt aber aus (9)
\weq{10}{(\chi_b, r) = 0;}
also der Satz:

\gesperrt{1. Wenn die Periode $\eta$ nicht zu $\mathfrak{A}$, sondern zu einer umfassenderen Gruppe $\mathfrak{A}'$ gehört, dann verschwindet jede Resolvente $(\chi_b, r)$, wenn $b$ eine Zahl aus $\mathfrak{B}$ ist, die nicht in $\mathfrak{B}'$ vorkommt.}

Um nun die Bedingungen für diesen Satz weiter zu verfolgen, müssen wir die Bildungsweise der Charaktere berücksichtigen.

Wir wählen die Bezeichnung so, wie wir sie am Schluss des \S.\,16 eingeführt haben. Dann ist, wenn $\beta_{-1}, \beta_0, \beta_1, \ldots \beta_\mu$ die Indices von $b$ und $\nu_{-1}, \nu_0, \nu_1, \ldots \nu_\mu$ die von $n$ sind,
\weq{11}{\chi_b(n) = \varepsilon_{-1}^{\beta_{-1}\,\nu_{-1}}\, \varepsilon_0^{\beta_0\nu_0}\, \varepsilon_1^{\beta_1\nu_1} \ldots \varepsilon_\mu^{\beta_\mu\nu_\mu}.}

Bezeichnen wir mit $r_0, r_1, \ldots r_\mu$ primitive Einheitswurzeln der Grade $2^\lambda, q_1^{\varkappa_1}, q_2^{\varkappa_2}, \ldots q_\mu^{\varkappa_\mu}$, so können wir jede primitive $m^{\text{te}}$ Einheitswurzel in der Form darstellen (Bd.\,I, \S.\,132)
\weq{12}{r = r_0 r_1 \ldots r_\mu,}
und wenn wir $n_0, n_1, \ldots n_\mu$ aus den Congruenzen bestimmen
\weq{13}{n \equiv n_0 \ (\mathrm{mod}\ 2^\lambda), \quad n \equiv n_1 \ (\mathrm{mod}\ q_1^{\varkappa_1}), \ldots n \equiv n_\mu \ (\mathrm{mod}\ q_\mu^{\varkappa_\mu}),}
so zerfällt $(\chi_b, r)$ nach (5) in das Product der folgenden Summen:
\weq{14}{\overset{n_0}{\Sigma}\, \varepsilon_{-1}^{\beta_{-1}\,\nu_{-1}}\, \varepsilon_0^{\beta_0\nu_0}\, r_0^{n_0}, \quad \overset{n_1}{\Sigma}\, \varepsilon_1^{\beta_1\nu_1}\, r_1^{n_1}, \ldots \overset{n_\mu}{\Sigma}\, \varepsilon_\mu^{\beta_\mu\nu_\mu}\, r_\mu^{n_\mu}.}

Wenn nun $b$ eine Zahl ist, die in $\mathfrak{B}$, aber nicht in $\mathfrak{B}'$ vorkommt, so muss nach dem Satze 1. eine von diesen Summen verschwinden.

Dafür ist aber nach \S.\,17 die nothwendige und hinreichende Bedingung die, dass eine der Congruenzen
\weq{15}{\beta_0 \equiv 0\ (\mathrm{mod}\ 2), \quad \beta_1 \equiv 0\ (\mathrm{mod}\ q_1), \ldots \beta_\mu \equiv 0\ (\mathrm{mod}\ q_\mu)}
befriedigt ist, und zwar eine solche, deren Modul mehrfach in $m$ aufgeht.

Diese Bedingung muss zunächst, wenn die Voraussetzung des Satzes 1. zutrifft, für jede Zahl $b$, die in $\mathfrak{B}$, aber nicht in $\mathfrak{B}'$ vorkommt, erfüllt sein. Ist aber $b'$ eine Zahl in $\mathfrak{B}'$, so ist für jeden ganzzahligen Exponenten $x$ das Product $b'^{x} b$ in $\mathfrak{B}$, aber
%CONT>>


% ===== scan leaf 98, printed 77 =====
% !PRINTED 77
\weberpage{77}

%<<CONT
nicht in $\mathfrak{B}'$ enthalten, und also muss, wenn $\beta_0', \beta_1', \dots \beta_\mu'$ die
Indices von $b'$ sind, für jedes $x$ eine der Congruenzen

\weq{16}{
\begin{aligned}
x\beta_0' + \beta_0 &\equiv 0 \pmod{2} \\
x\beta_1' + \beta_1 &\equiv 0 \pmod{q_1} \\
&\cdots\cdots\cdots\cdots \\
x\beta_\mu' + \beta_\mu &\equiv 0 \pmod{q_\mu}
\end{aligned}
}

befriedigt sein.

Sind nun zwei Zahlen $\beta', \beta$ nicht beide durch eine Primzahl $q$ theilbar, so kann man immer über $x$ so verfügen, dass $x\beta' + \beta$ nicht durch $q$ theilbar wird; man braucht nur, wenn $\beta$ durch $q$ theilbar ist, $x$ durch $q$ nicht theilbar, und wenn $\beta$ durch $q$ nicht theilbar ist, $x$ durch $q$ theilbar anzunehmen. Man kann also $x$ auch so bestimmen, dass es gleichzeitig mehreren solchen Forderungen genügt. Es folgt also aus (16), dass entweder $\beta_0', \beta_0$ durch $2$ oder $\beta_1', \beta_1$ durch $q_1 \dots$ oder $\beta_\mu', \beta_\mu$ durch $q_\mu$ theilbar sein müssen; d.\,h. eine der Congruenzen (15) muss auch erfüllt sein, wenn $b$ irgend eine Zahl aus $\mathfrak{B}$ ist, gleichviel ob sie in $\mathfrak{B}'$ vorkommt oder nicht. Endlich folgt wieder aus (16), indem wir unter $b, b'$ zwei beliebige Zahlen aus $\mathfrak{B}$ verstehen, dass von den Congruenzen (15) eine und dieselbe für alle Zahlen aus $\mathfrak{B}$ bestehen muss. Dies können wir nun so zusammenfassen:

\gesperrt{2. Wenn die Periode $\eta$ nicht zu $\mathfrak{A}$ gehört, so muss es unter den Primzahlen (einschliesslich 2), die mehr als einmal in $m$ aufgehen, eine geben, $q$, so dass für jede Zahl $b$ aus $\mathfrak{B}$ der dem $q$ entsprechende Index $\beta$ durch $q$ theilbar ist.}

Wir betrachten jetzt die Gruppe $Q$, die nach der Bezeichnungsweise (5) des \S.\,18 durch die Congruenz

\weq{17}{Q \equiv 1 \pmod{\dfrac{m}{q}}}

definirt ist, und bezeichnen mit $P$ die dazu reciproke Gruppe.

Nach dem Satze 7., \S.\,19 ist jede Gruppe, die wie $\mathfrak{B}$ die Eigenschaft hat, dass der dem $q$ entsprechende Index $\beta$ durch $q$ theilbar ist, ein Theiler der Gruppe $P$, und folglich ist nach \S.\,12, 8. $Q$ ein Theiler von $\mathfrak{A}$, also $\mathfrak{A}$ nach \S.\,19, 8. ein nicht primärer Theiler von $\mathfrak{N}$, und damit sind wir zu dem Satze gelangt:

\gesperrt{3. Ist $\mathfrak{A}$ ein primärer Theiler von $\mathfrak{N}$, so gehört die Periode $\eta$ zu der Gruppe $\mathfrak{A}$.}


% ===== scan leaf 99, printed 78 =====
% !PRINTED 78
\weberpage{78}

Alle primären Theiler des vollen Kreistheilungskörpers $\Omega_m$ sind demnach in der Form $\Omega(\eta)$ darstellbar, d.\,h. alle Zahlen eines solchen Theilers sind rational durch die Kreistheilungsperioden $\eta$ darstellbar, und da man die nicht primären Theiler von $\Omega_m$ als primäre Theiler von niedrigeren Kreistheilungskörpern wiederfindet, so folgt:

\gesperrt{4. Alle Kreistheilungskörper sind in der Form $\Omega(\eta)$ darstellbar.}

Oder auch:

\gesperrt{5. Jede rationale Function von Einheitswurzeln kann als rationale Function einer Kreistheilungsperiode dargestellt werden.}

Und dieser Satz lässt sich auch in der Form aussprechen:

\gesperrt{6. Ist $\mathfrak{A}$ ein beliebiger primärer oder nicht primärer Theiler von $\mathfrak{N}$ vom Index $e$, so giebt es einen Theiler $m_1$ von $m$ und eine in $\Omega_{m_1}$ gelegene Kreistheilungsperiode $\eta$ vom Index $e$, so dass $\eta$, als Function von $r$ aufgefasst, eine zur Gruppe $\mathfrak{A}$ gehörige Function ist, also, wenn $a$ zu $\mathfrak{A}$ gehört, durch die Substitution $(r, r^a)$ ungeändert bleibt und für alle Substitutionen von $\mathfrak{N}$ $e$ verschiedene Werthe erhält.}

Durchläuft $a$ eine Gruppe $\mathfrak{A}$, die ein primärer Theiler von $\mathfrak{N}$ ist, und setzen wir

\[
\eta_h = \sum^a r^{ah},
\]

auch wenn $h$ nicht relativ prim zu $m$ ist, so können wir diese Grössen $\eta_h$, die alle die Permutationen der Gruppe $\mathfrak{A}$ gestatten, rational durch $\eta = \sum r^a$ darstellen. Andererseits lässt sich aber auch jede rationale Function von $\eta$ linear durch die $\eta_h$ darstellen. Dies wird bewiesen sein, wenn gezeigt ist, dass man das Product zweier $\eta_h$ linear durch die $\eta_h$ ausdrücken kann. Es ist aber:

\[
\eta_h \eta_k = \sum^{a, a'} r^{ha + ka'};
\]

nimmt man zuerst die Summe nach $a'$ bei feststehendem $a$, so kann man $a'$ durch $aa'$ ersetzen, da $aa'$ zugleich mit $a'$ die Gruppe $\mathfrak{A}$ durchläuft. Man erhält so

\[
\eta_h \eta_k = \sum^{a, a'} r^{a(h + ka')} = \sum^{a'} \eta_{h + ka'}.
\]


% ===== scan leaf 100, printed 79 =====
% !PRINTED 79
\weberpage{79}
\pdfbookmark[2]{§ 21. Kreistheilungskörper mit gegebener Gruppe}{bm25}

Wenn die Primzahl $q$ nur einfach in $m$ aufgeht und $\mathfrak{A}$ durch $Q$ theilbar ist, so ist zwar $\mathfrak{A}$ nicht primär; es gehört aber trotzdem die Periode $\eta$ zu der Gruppe $\mathfrak{A}$.

Es kann aber in diesem Falle $\eta$ durch Einheitswurzeln von niedrigerem Grade dargestellt werden. In folgender Weise lässt sich diese Darstellung finden. Alle Zahlen von $Q$ sind von der Form $1 + hm'$, worin $h$ die Reihe der Zahlen $0, 1, \dots q - 1$ durchläuft, mit Ausnahme des einzigen Werthes, für den

\weq{18}{1 + hm' = kq}

durch $q$ theilbar wird. Setzt man also unter der Voraussetzung, dass $\mathfrak{A}$ durch $Q$ theilbar ist,

\[
\mathfrak{A} = Q + Qa_1 + Qa_2 \dots,
\]

so folgt, dass jede Zahl $a$ in $\mathfrak{A}$ von der Form ist:

\weq{19}{a = a'(1 + hm'),}

worin $a'$ die Reihe der Zahlen $1, a_1, a_2 \dots$ durchläuft. Daraus ist noch zu schliessen, dass die Reste der Zahlen $1, a', a'' \dots$ nach dem Modul $m'$ eine Gruppe bilden. Nimmt man nun die Summe über alle Werthe $h = 0, 1, \dots q - 1$ und nimmt dann den einen durch (18) bestimmten Werth wieder weg, so folgt:

\[
\eta = \sum^a r^a = \sum^{a'} \sum^h r^{a'(1 + hm')} - \sum^{a'} r^{a'kq},
\]

oder, da $\sum^h r^{a'm'h} = 0$ ist:

\[
\eta = -\sum^{a'} r^{a'kq}.
\]

Nun ist $r^q$ eine $m'^{\text{te}}$ Einheitswurzel, und $\eta$ also, vom Vorzeichen abgesehen, gleich einer aus $m'^{\text{ten}}$ Einheitswurzeln gebildeten Periode\footnote{Hier ist die Abhandlung von \gesperrt{Fuchs} zu vergleichen: Ueber die Perioden, welche aus den Wurzeln der Gleichung $\omega^n = 1$ gebildet sind, wenn $n$ eine zusammengesetzte Zahl ist. Crelle's Journal, Bd. 61 (1863).}.

\begin{center}
\S.\,21.

\gesperrt{Kreistheilungskörper mit gegebener Gruppe}.
\end{center}

Im Vorhergehenden hat sich also ergeben, dass jede Kreistheilungsperiode

\[
\eta = \sum^a r^a,
\]

in der $r$ eine primitive $m^{\text{te}}$ Einheitswurzel ist und $a$ die Zahlen einer Gruppe $\mathfrak{A}$ durchläuft, wenn $\mathfrak{A}$ ein primärer Theiler von $\mathfrak{N}$
%CONT>>


% ===== scan leaf 101, printed 80 =====
% !PRINTED 80
\weberpage{80}

%<<CONT
ist, die Wurzel einer Abel'schen Gleichung ist, deren Gruppe mit der zu $\mathfrak{A}$ reciproken Gruppe $\mathfrak{B}$ isomorph ist.

Die Aufgabe, die wir jetzt stellen und lösen wollen, ist folgende:

\gesperrt{I. Es soll der Modul $m$ und die Gruppe $\mathfrak{A}$ auf alle mögliche Arten so bestimmt werden, dass $\mathfrak{B}$ ein beliebig gegebenes System von Invarianten hat.}

Oder was dasselbe ist:

\gesperrt{Es sollen alle Kreistheilungskörper von gegebener Gruppe bestimmt werden.}

Wir machen dabei immer die Voraussetzung, dass $\mathfrak{A}$ ein primärer Theiler von $\mathfrak{N}$ sein soll, weil wir sonst einen und denselben Körper mehrmals erhalten würden. Die Aufgabe zerfällt in zwei Theile, nämlich:

\gesperrt{II. Welche Moduln $m$ sind geeignet, eine Gruppe $\mathfrak{B}$ von gegebenen Invarianten zu erzeugen?}

\gesperrt{III. Wie findet man diese Gruppe $\mathfrak{B}$ und die zugehörige Gruppe $\mathfrak{A}$?}

Wir müssen bei dieser Untersuchung die Bezeichnung gegen das Frühere etwas ändern, damit die Uebersichtlichkeit nicht verloren geht. Die exceptionelle Stellung der Primzahl 2, die bei den bisherigen Betrachtungen immer berücksichtigt werden musste, machte eine etwas umständliche Bezeichnung nothwendig. Diese Unterscheidung ist im Folgenden nicht mehr in dem Maasse nothwendig, und darum lässt sich die Bezeichnung jetzt vereinfachen.

Wir bezeichnen die Indices einer Zahl $a$ aus der Gruppe $\mathfrak{A}$ mit

\[
\alpha_1, \alpha_2, \dots \alpha_\mu \quad \text{(Indices von $a$)}
\]

und die entsprechenden Indexmoduln mit

\[
c_1, c_2, \dots c_\mu \quad \text{(Indexmoduln)},
\]

so dass $\mu$ gleich der Anzahl der in $m$ aufgehenden Primzahlen, oder wenn $m$ durch 8 theilbar ist, um 1 grösser ist.

Es seien ferner

\[
\omega_1, \omega_2, \dots \omega_\mu
\]

primitive Einheitswurzeln der Grade $c_1, c_2, \dots c_\mu$. Sind nun die Indices einer Zahl $b$ in der zu $\mathfrak{A}$ reciproken Gruppe $\mathfrak{B}$

\[
\beta_1, \beta_2, \dots \beta_\mu \quad \text{(Indices von $b$),}
\]
%CONT>>


% ===== scan leaf 102, printed 81 =====
% !PRINTED 81
\weberpage{81}
%<<CONT
so ist die Beziehung zwischen den $\alpha$ und den $\beta$ durch die
Gleichung

\weq{1}{\omega_1^{\alpha_1\beta_1}\,\omega_2^{\alpha_2\beta_2}\,\dots\,\omega_\mu^{\alpha_\mu\beta_\mu}=1}
ausgedrückt.

Die Invarianten der Gruppe $\mathfrak{B}$, die nach \S.\,10 lauter Primzahlpotenzen sind, wollen wir mit
\[
i_1, i_2, \dots i_\nu \quad (\text{Invarianten von } \mathfrak{B})
\]
bezeichnen. Wir nehmen diese Invarianten als beliebig gegebene
Primzahlpotenzen an und bezeichnen mit $J$ ihr kleinstes gemeinschaftliches Vielfache, so dass
\weq{2}{J=i_1i_1'=i_2i_2'=\dots=i_\nu i_\nu'}
gesetzt werden kann.

Ausserdem soll $\mathfrak{A}$ als primärer Theiler von $\mathfrak{N}$ vorausgesetzt werden, was nach \S.\,19, 8. mit der Bedingung gleichbedeutend ist:

\gesperrt{IV. Keiner der Indices $\beta$ soll für alle Zahlen $b$, wenn die entsprechende Primzahl $q$ mehrmals in $m$ aufgeht, durch $q$, oder wenn $q$ nur einmal in $m$ aufgeht, durch $q-1$ theilbar sein.}

Nach der Definition der Invarianten muss die Gruppe $\mathfrak{B}$ eine Basis haben, deren Elemente
\[
\begin{aligned}
&\phantom{\text{von den Graden}} \quad g_1, g_2, \dots g_\nu \quad (\text{Basis von } \mathfrak{B})\\
&\text{von den Graden} \quad i_1, i_2, \dots i_\nu
\end{aligned}
\]
sind, so dass jede Zahl $b$ einmal und nur einmal in der Form
\weq{3}{b\equiv g_1^{x_1}g_2^{x_2}\dots g_\nu^{x_\nu}\ (\mathrm{mod}\ m)}
enthalten ist, wenn die Exponenten $x_1, x_2, \dots x_\nu$ je ein volles Restsystem nach den Moduln $i_1, i_2, \dots i_\nu$ durchlaufen.

Alle Grössen von der Form (3) bilden, wenn $g_1, g_2, \dots g_\nu$ beliebige Zahlen sind, bei unbeschränkter Veränderlichkeit der ganzzahligen Exponenten $x$ gewiss eine Gruppe. Sollen aber die $g_h$ die Elemente einer Basis dieser Gruppe, und ihre Grade $i_h$ sein, so darf $1\equiv g_1^{x_1}g_2^{x_2}\dots g_\nu^{x_\nu}\ (\mathrm{mod}\ m)$ nur dann erfüllt sein, wenn $x_h$ durch $i_h$ theilbar ist, also:

\gesperrt{1. Die nothwendige und hinreichende Bedingung dafür, dass $g_1, g_2, \dots g_\nu$ die Elemente einer Basis einer Gruppe $\mathfrak{B}$ von den Graden $i_1, i_2, \dots i_\nu$ seien, ist die, dass die Congruenz}
\weq{4}{g_1^{x_1}g_2^{x_2}\dots g_\nu^{x_\nu}\equiv 1\ (\mathrm{mod}\ m)}
%CONT>>


% ===== scan leaf 103, printed 82 =====
% !PRINTED 82
\weberpage{82}
%<<CONT
\gesperrt{dann und nur dann besteht, wenn zugleich die
Congruenzen}
\weq{5}{x_1\equiv 0\ (\mathrm{mod}\ i_1),\quad x_2\equiv 0\ (\mathrm{mod}\ i_2),\quad \dots\ x_\nu\equiv 0\ (\mathrm{mod}\ i_\nu)}
\gesperrt{erfüllt sind.}

Hiernach ist also $J$ die kleinste positive Zahl, die für alle $b$ der Congruenz
\[
b^J\equiv 1\ (\mathrm{mod}\ m)
\]
genügt, oder die kleinste positive Zahl, für die für alle Systeme der $\beta$
\weq{6}{J\beta_1\equiv 0\ (\mathrm{mod}\ c_1),\quad J\beta_2\equiv 0\ (\mathrm{mod}\ c_2),\quad \dots\ J\beta_\mu\equiv 0\ (\mathrm{mod}\ c_\mu).}

Daraus ergeben sich die ersten Schlüsse über die Zusammensetzung der Zahl $m$.

a) \gesperrt{Eine ungerade Primzahl $q$, die nicht in $J$ enthalten ist, kann nur einfach in $m$ aufgehen.}

Denn ist $m$ durch $q^\varkappa$ theilbar, so ist eine der Grössen $c$, etwa $c_1=\varphi(q^\varkappa)=q^{\varkappa-1}(q-1)$, also wenn $\varkappa>1$ ist, so ist $c_1$ noch durch $q$ theilbar, und aus (6) folgt dann, dass alle $\beta_1$ durch $q$ theilbar sein müssten, was der Voraussetzung IV. widerspricht.

b) \gesperrt{Eine ungerade Primzahl $q$, die in $J$ aufgeht, kann höchstens einmal mehr in $m$, als in $J$ enthalten sein.}

Denn man kann ebenso schliessen, dass, wenn $c_1=q^{\varkappa-1}(q-1)$, und $J$ nicht durch $q^{\varkappa-1}$ theilbar ist, alle $\beta_1$ durch $q$ theilbar sein müssten.

c) \gesperrt{Ist $J$ ungerade, so muss auch $m$ ungerade sein.}

Denn ist $m$ durch eine Potenz von 2, also mindestens durch 4 theilbar, so ist der der Zahl 2 entsprechende Indexmodul gerade, und die entsprechenden $\beta$ müssten nach (6) alle gerade sein, entgegen der Forderung IV.

d) \gesperrt{Ist $J$ durch eine Potenz von 2 theilbar, so kann $m$ den Factor 2 höchstens zweimal öfter enthalten, als $J$.}

Denn ist $m$ durch $2^\lambda$ theilbar und $\lambda>2$, so sind zwei der Indexmoduln
\[
c_1=2,\quad c_2=2^{\lambda-2}.
\]
%CONT>>


% ===== scan leaf 104, printed 83 =====
% !PRINTED 83
\weberpage{83}
%<<CONT
Wäre also $J$ nicht durch $2^{\lambda-2}$ theilbar, so müssten alle $\beta_2$ durch 2 theilbar sein, was wieder gegen die Forderung IV ist.

e) \gesperrt{Wenn $q$ eine einfach in $m$ aufgehende Primzahl ist, so muss $q-1$ wenigstens durch eine der in $J$ aufgehenden Primzahlen theilbar sein.}

Denn wäre $c_1=q-1$ relativ prim zu $J$, so müssten alle $\beta_1$ durch $q-1$ theilbar sein, im Widerspruch mit IV. Die weiteren Bedingungen für $m$ ergeben sich später.

\gesperrt{Bestimmung der Gruppe $\mathfrak{B}$:}

Die Gruppe $\mathfrak{B}$ ist bestimmt, wenn ihre Basis $g_1, g_2, \dots g_\nu$ gegeben ist. Die Elemente der Basis wollen wir durch ihre Indices bestimmen, und führen also folgende Bezeichnung ein. Es seien
\weq{7}{
\begin{aligned}
&\gamma_{1,1}, \gamma_{1,2}, \dots \gamma_{1,\mu} \quad \text{die Indices von } g_1\\
&\gamma_{2,1}, \gamma_{2,2}, \dots \gamma_{2,\mu} \quad \text{\textquotedbl{}} \quad \text{\textquotedbl{}} \quad \text{\textquotedbl{}} \quad g_2\\
&\dots\dots\dots\dots\dots\dots\dots\dots\dots\dots\dots\dots\dots\dots\\
&\gamma_{\nu,1}, \gamma_{\nu,2}, \dots \gamma_{\nu,\mu} \quad \text{\textquotedbl{}} \quad \text{\textquotedbl{}} \quad \text{\textquotedbl{}} \quad g_\nu.
\end{aligned}
}
Die Indices $\beta_1, \beta_2, \dots \beta_\mu$ eines beliebigen Elementes $b$ erhält man dann nach (3) in der Form:
\weq{8}{
\begin{aligned}
\beta_1 &\equiv \gamma_{1,1}x_1+\gamma_{2,1}x_2+\dots+\gamma_{\nu,1}x_\nu\ (\mathrm{mod}\ c_1)\\
\beta_2 &\equiv \gamma_{1,2}x_1+\gamma_{2,2}x_2+\dots+\gamma_{\nu,2}x_\nu\ (\mathrm{mod}\ c_2)\\
&\dots\dots\dots\dots\dots\dots\dots\dots\dots\dots\dots\dots\dots\dots\\
\beta_\mu &\equiv \gamma_{1,\mu}x_1+\gamma_{2,\mu}x_2+\dots+\gamma_{\nu,\mu}x_\nu\ (\mathrm{mod}\ c_\mu),
\end{aligned}
}
und wir fragen nun nach den nothwendigen und hinreichenden Bedingungen für die Zahlen $\gamma_{h,k}$, damit $g_1, g_2 \dots g_\nu$ die Basis einer Gruppe $\mathfrak{B}$ von den verlangten Eigenschaften ist. Nach 1. ist hierfür zuerst nothwendig:

\gesperrt{2. Die Zahlen $\gamma_{h,k}$ müssen so beschaffen sein, dass die Congruenzen}
\weq{9}{
\begin{aligned}
\gamma_{1,1}x_1+\gamma_{2,1}x_2+\dots+\gamma_{\nu,1}x_\nu &\equiv 0\ (\mathrm{mod}\ c_1)\\
\gamma_{1,2}x_1+\gamma_{2,2}x_2+\dots+\gamma_{\nu,2}x_\nu &\equiv 0\ (\mathrm{mod}\ c_2)\\
&\dots\dots\dots\dots\dots\dots\dots\dots\dots\dots\dots\dots\\
\gamma_{1,\mu}x_1+\gamma_{2,\mu}x_2+\dots+\gamma_{\nu,\mu}x_\nu &\equiv 0\ (\mathrm{mod}\ c_\mu)
\end{aligned}
}
\gesperrt{dann und nur dann erfüllt sind. wenn zugleich die Congruenzen}
% [sic: Weber prints a full stop before the lower-case ``wenn''; a comma is grammatically required.]
\weq{10}{x_1\equiv 0\ (\mathrm{mod}\ i_1),\quad x_2\equiv 0\ (\mathrm{mod}\ i_2),\quad \dots\ x_\nu\equiv 0\ (\mathrm{mod}\ i_\nu)}
\gesperrt{bestehen.}


% ===== scan leaf 105, printed 84 =====
% !PRINTED 84
\weberpage{84}
Ist diese Bedingung erfüllt, so ist gewiss $g_1, g_2, \dots g_\nu$ die Basis einer Gruppe mit den Invarianten $i_1, i_2, \dots i_\nu$. Wenn aber auch noch die in IV. aufgestellte Forderung erfüllt sein soll, dann folgt noch weiter:

\gesperrt{3. Ist $q_h$ eine in $m$ aufgehende Primzahl, der der Index $\beta_h$ entspricht, so dürfen die $\nu$ Grössen}
\[
\gamma_{1,h}, \gamma_{2,h}, \dots \gamma_{\nu,h},
\]
\gesperrt{wenn $q_h$ mehrfach in $m$ aufgeht, nicht alle durch $q_h$, und wenn $q_h$ nur einmal in $m$ aufgeht, nicht alle durch $q_h-1$ theilbar sein.}

Die Bedingungen 2., 3. sind dann also die nothwendigen und hinreichenden, und es kommt jetzt nur noch darauf an, sie in eine Form zu bringen, dass die Möglichkeit ihrer Erfüllung beurtheilt werden kann.

Wir bezeichnen mit $\delta_{k,h}$ \gesperrt{den grössten gemeinschaftlichen Theiler von $i_k$ und $c_h$ und setzen}
\weq{11}{i_k=\delta_{k,h}\,i_{k,h}, \quad c_h=\delta_{k,h}\,c_{k,h}, \quad \begin{aligned}k&=1,2,\dots\nu,\\ h&=1,2,\dots\mu,\end{aligned}}
so dass $i_{k,h}$ und $c_{k,h}$ \gesperrt{relativ prim} sind. Es ist dann zu bemerken, dass die Grössen
\[
\delta_{k,h}, \quad i_{k,h}, \quad c_{k,h}
\]
durch $m$ und die Invarianten $i_k$ völlig bestimmt sind. Nach der Definition der Invarianten müssen die $\delta_{k,h}$ und $i_{k,h}$ \gesperrt{Primzahlpotenzen} sein.

Aus den Voraussetzungen, die wir in a) bis d) über die Zahl $m$ gemacht haben, ergiebt sich eine Folgerung für die $c_{k,h}$, die wir hervorheben müssen.

\gesperrt{4. Ist $q_h^\varkappa$ einer der Factoren von $m$, der dem Index $\beta_h$ entspricht, ist also $c_h=q_h^{\varkappa-1}(q_h-1)$, oder wenn $q=2$ und $\varkappa>2$ ist, $c_h=2^{\varkappa-2}$, und wenn $\varkappa=2$ ist, $c_h=2$, so können, wenn $\varkappa>1$ ist, die Zahlen}
\weq{12}{c_{1,h}, c_{2,h}, \dots c_{\nu,h}}
\gesperrt{nicht alle durch $q_h$, und wenn $\varkappa=1$ ist, nicht alle durch $q-1$ theilbar sein.}
% [sic: Weber switches from q_h to q in the exceptional 2-adic clause and in the final q-1 clause; q_h is intended in both places.]

Wenn nämlich die Grössen (12) alle durch $q_h$ theilbar sind, so sind nach (11) die Zahlen
\weq{13}{\delta_{1,h}.\ \delta_{2,h}, \dots \delta_{\nu,h}}
% [sic: Weber prints a full stop, rather than a comma, after the first term.]
%CONT>>


% ===== scan leaf 106, printed 85 =====
% !PRINTED 85
\weberpage{85}
%<<CONT
alle nicht durch $q_h^{\varkappa-1}$ theilbar, und die Zahlen
% [class ruling, 2026-08-06, 6x zoom on b2 idx 106 = printed 85]
% The glyph is VARKAPPA, the hooked form. On that page it is legible in
% q_h^{varkappa-1} and again in running text: "nur varkappa - 1 oder 2
% (bei varkappa = 2) an Stelle von varkappa". It is the curved form with no
% straight vertical stem, which is the standard German mathematical kappa of the
% period -- not x and not u. This is a property of the typeface, so it holds
% wherever the letter appears in these volumes. \varkappa is correct.
\weq{14}{i_{1,h},\ i_{2,h},\ \dots\ i_{\nu,h}}
sind alle nicht durch $q_h$ theilbar, folglich ist nach (11) keine der
Zahlen $i_1, i_2, \dots i_\nu$, und also auch $J$ nicht durch $q_h^{\varkappa-1}$ theilbar,
was den Annahmen a) und b) widerspricht. Für $q_h = 2$ ist in
diesem Schlusse nur $\varkappa - 1$ oder $2$ (bei $\varkappa = 2$) an Stelle von $\varkappa$
zu setzen, um denselben Widerspruch gegen c) und d) zu erhalten.

    Ist aber $\varkappa = 1$ und sind die Zahlen (12) alle durch $q_h - 1$
theilbar, so müssen sie gleich $q_h - 1$ sein; die Zahlen (13) sind
alle $= 1$ und es würde also folgen, dass die Zahlen $i_k$ und
also auch $J$ relativ prim zu $q_h - 1$ wären, was der Voraussetzung e) widerspricht.

    Wenn wir jetzt zu der in 2. ausgesprochenen Forderung für
die Grössen $\gamma_{k,h}$ zurückkehren, so ergiebt sich zunächst, dass
die Congruenzen (9) erfüllt sein müssen, wenn $x_k = i_k$ und die
übrigen $x = 0$ gesetzt werden, also:
\[
\gamma_{k,h}\, i_k \equiv 0\ (\mathrm{mod}\ c_h),
\]
und daraus mit Berücksichtigung von (11):
\[
i_{k,h}\, \gamma_{k,h} \equiv 0\ (\mathrm{mod}\ c_{k,h}).
\]

    Weil aber $i_{k,h}$ und $c_{k,h}$ relativ prim sind, so folgt, dass $\gamma_{k,h}$
durch $c_{k,h}$ theilbar sein muss. Wir führen also ein neues System
von ganzen Zahlen $e_{k,h}$ ein durch die Gleichungen
\weq{15}{\gamma_{k,h} = c_{k,h}\, e_{k,h}}
und suchen nun für diese Zahlen die aus 2. folgenden Bedingungen. Da die Zahlen $\gamma_{k,h}$, wie aus ihrer Definition hervorgeht, nur nach dem Modul $c_h$ bestimmt sind, so ist $e_{k,h}$ nur nach
dem Modul $\delta_{k,h}$ zu bestimmen. Nehmen wir eine von den Congruenzen (9):
\[
\gamma_{1,h}x_1 + \gamma_{2,h}x_2 + \cdots + \gamma_{\nu,h}x_\nu \equiv 0\ (\mathrm{mod}\ c_h),
\]
und führen darin (15) ein, so erhält sie die Form
\weq{16}{c_{1,h}e_{1,h}x_1 + c_{2,h}e_{2,h}x_2 + \cdots + c_{\nu,h}e_{\nu,h}x_\nu \equiv 0\ (\mathrm{mod}\ c_h).}
Hierin setzen wir nun nach (11):
\[
c_{k,h} = \frac{c_h}{\delta_{k,h}} = \frac{c_h\, i_{k,h}}{i_k}\,.
\]
Demnach sind die Congruenzen (16) gleichbedeutend mit
der Forderung, dass
%CONT>>


% ===== scan leaf 107, printed 86 =====
% !PRINTED 86
\weberpage{86}
%<<CONT
\[
\frac{e_{1,h}i_{1,h}x_1}{i_1} + \frac{e_{2,h}i_{2,h}x_2}{i_2} + \cdots \frac{e_{\nu,h}i_{\nu,h}x_\nu}{i_\nu}
\]
ganze Zahlen sein müssen, oder, mit Rücksicht auf die Bezeichnung (2), mit den Congruenzen:
\weq{17}{e_{1,h}i_{1,h}i'_1x_1 + e_{2,h}i_{2,h}i'_2x_2 + \cdots e_{\nu,h}i_{\nu,h}i'_\nu x_\nu \equiv 0\ (\mathrm{mod}\ J).}
Diese Congruenzen sind, wie man sieht, immer erfüllt, wenn
$x_1$ durch $i_1$, $x_2$ durch $i_2 \dots$, $x_\nu$ durch $i_\nu$ theilbar ist, und die
Forderung 2. reducirt sich also jetzt darauf:

    5. \gesperrt{dass die $e_{k,h}$ so zu bestimmen sind, dass die Congruenzen (17) für keine anderen als den Congruenzen $x_k \equiv 0\ (\mathrm{mod}\ i_k)$ genügende Werthe der $x$ befriedigt sein sollen.}

    Dazu kommt noch als Umformung der Forderung 3.:

    6. Die Zahlen
\[
c_{1,h}e_{1,h},\ c_{2,h}e_{2,h},\ \dots\ c_{\nu,h}e_{\nu,h}
\]
\gesperrt{dürfen nicht alle durch $q_h$, wenn $\varkappa > 1$, und
durch $(q_h - 1)$, wenn $\varkappa = 1$ ist, theilbar sein.}

    Die Forderung 6. enthält wegen 4. keine Unmöglichkeit.

    \gesperrt{Durch 5. und 6. sind die nothwendigen und hinreichenden Bedingungen für die Zahlen $e_{h,k}$ ausgedrückt.}

    Es handelt sich also jetzt noch um die Frage, ob und unter welchen Voraussetzungen diese Forderungen durch geeignete Wahl der $e_{h,k}$ befriedigt werden können.
% [sic: Weber switches in these two sentences from the defined e_{k,h} to e_{h,k}; e_{k,h} is intended.]

    Diese Frage wird dadurch sehr vereinfacht, dass sich die
Congruenz (17) in mehrere spalten lässt. Die Invarianten $i_k$
sind, wie wir angenommen haben, Primzahlpotenzen. Es möge
eine von diesen Primzahlen mit $p$ bezeichnet sein, und es sei
\weq{18}{i_1 = p^{\pi_1},\ i_2 = p^{\pi_2},\ \dots\ i_\varrho = p^{\pi_\varrho},}
während die übrigen $i_k$, wenn noch solche vorhanden sind, nicht
mehr durch $p$ theilbar sein sollen. Den grössten der Exponenten $\pi_1, \pi_2, \dots \pi_\varrho$ wollen wir mit $\pi$ bezeichnen.

    Dann ist
\[
J = p^\pi J',
\]
und $J'$ ist nicht mehr durch $p$ theilbar. Es ist ferner
\[
i'_1 = p^{\pi-\pi_1}J',\ i'_2 = p^{\pi-\pi_2}J',\ \dots\ i'_\varrho = p^{\pi-\pi_\varrho}J',
\]
während $i'_{\varrho+1}, \dots i'_\nu$ durch $p^\pi$ theilbar sind.


% ===== scan leaf 108, printed 87 =====
% !PRINTED 87
\weberpage{87}
Demnach hat die Congruenz (17) zur Folge:
\begin{gather}
e_{1,h}i_{1,h}p^{\pi-\pi_1}x_1 + e_{2,h}i_{2,h}p^{\pi-\pi_2}x_2 + \cdots e_{\varrho,h}i_{\varrho,h}p^{\pi-\pi_\varrho}x_\varrho \notag\\
\equiv 0\ (\mathrm{mod}\ p^\pi) \qquad h = 1, 2, \dots \mu, \tag{19}
\end{gather}
und die übrigen in $J$ aufgehenden Primzahlen ergeben jede ein
eben solches System von Congruenzen, die in ihrer Gesammtheit
vollständig gleichbedeutend mit der Congruenz (17) sind. Diese
Congruenzen (19) dürfen nicht anders befriedigt werden können,
als dadurch, dass die Zahlen
\[
p^{\pi-\pi_1}x_1,\ p^{\pi-\pi_2}x_2,\ \dots\ p^{\pi-\pi_\varrho}x_\varrho
\]
durch $p^\pi$ theilbar sind.

    Setzen wir also zur Abkürzung:
\[
p^{\pi-\pi_1}x_1 = y_1,\ p^{\pi-\pi_2}x_2 = y_2,\ \dots\ p^{\pi-\pi_\varrho}x_\varrho = y_\varrho,
\]
so erhält (19) die Form:
\weq{20}{e_{1,h}i_{1,h}y_1 + e_{2,h}i_{2,h}y_2 + \cdots e_{\varrho,h}i_{\varrho,h}y_\varrho \equiv 0\ (\mathrm{mod}\ p^\pi),}
und diese Congruenzen sollen
\[
y_1 \equiv 0,\ y_2 \equiv 0,\ \dots\ y_\varrho \equiv 0\ (\mathrm{mod}\ p^\pi)
\]
zur nothwendigen Folge haben.

    Das System von linearen Congruenzen (20) kann nun mit
denselben Mitteln discutirt werden, wie ein System linearer Gleichungen, nämlich durch die Determinanten (vgl. den Abschnitt II
des ersten Bandes).

    Wir wollen zur besseren Uebersicht die Congruenzen (20)
für den Augenblick so schreiben:
\begin{gather}
a_{1,1}y_1 + a_{2,1}y_2 + \cdots + a_{\varrho,1}y_\varrho \equiv 0 \notag\\
a_{1,2}y_1 + a_{2,2}y_2 + \cdots + a_{\varrho,2}y_\varrho \equiv 0\ (\mathrm{mod}\ p^\pi). \tag{21}\\
\dots\dots\dots\dots\dots\dots\dots\dots\dots\dots \notag\\
a_{1,\mu}y_1 + a_{2,\mu}y_2 + \cdots + a_{\varrho,\mu}y_\varrho \equiv 0 \notag
\end{gather}
Wenn die Coëfficienten $a_{k,h}$ alle durch $p$ theilbar sind, dann
sind diese Congruenzen schon befriedigt, wenn die $y$ alle nur
durch $p^{\pi-1}$ theilbar sind, und dieser Fall darf also nicht eintreten.

    Wir wollen nun annehmen, es lasse sich aus der Matrix
\weq{22}{
\begin{matrix}
a_{1,1}, & a_{2,1}, & \dots & a_{\varrho,1} \\
a_{1,2}, & a_{2,2}, & \dots & a_{\varrho,2} \\
\dots & \dots & \dots & \dots \\
a_{1,\mu}, & a_{2,\mu}, & \dots & a_{\varrho,\mu}
\end{matrix}
}
%CONT>>


% ===== scan leaf 109, printed 88 =====
% !PRINTED 88
\weberpage{88}
%<<CONT
eine Determinante von $\tau$ Reihen bilden, die nicht durch $p$ theilbar ist, während alle Determinanten von mehr als $\tau$ Reihen,
wenn solche vorhanden sind, durch $p$ theilbar sind. Wir können
voraussetzen, dass die Indices so angeordnet sind, dass
\[
\Delta = \sum \pm a_{1,1}a_{2,2} \dots a_{\tau,\tau}
\]
nicht durch $p$ theilbar ist. Bezeichnen wir mit $\Delta_{k,h}$ die Unterdeterminanten von $\Delta$ und setzen
\[
\Delta_{k,1}a_{s,1} + \Delta_{k,2}a_{s,2} + \cdots \Delta_{k,\tau}a_{s,\tau} = D_{k,s},
\]
so ist $D_{k,s} = \Delta$, wenn $k = s$, und $= 0$, wenn $s \leqq \tau$ und von
$k$ verschieden ist, und ist, wenn $s > \tau$ ist, eine $\tau$-reihige Determinante der Matrix (22). Mit Hülfe dieser Bezeichnung lassen
sich nun aus den $\tau$ ersten Congruenzen (21) die folgenden herleiten:
\begin{gather}
\Delta y_1 + D_{1,\tau+1}y_{\tau+1} + \cdots D_{1,\varrho}y_\varrho \equiv 0 \notag\\
\Delta y_2 + D_{2,\tau+1}y_{\tau+1} + \cdots D_{2,\varrho}y_\varrho \equiv 0\ (\mathrm{mod}\ p^\pi) \tag{23}\\
\dots\dots\dots\dots\dots\dots\dots\dots\dots \notag\\
\Delta y_\tau + D_{\tau,\tau+1}y_{\tau+1} + \cdots D_{\tau,\varrho}y_\varrho \equiv 0. \notag
\end{gather}
Setzen wir zur Abkürzung
\[
A_{s,r} = \Delta a_{s,r} - \sum_{1,\tau}^{k} a_{k,r}D_{k,s} = \Delta a_{s,r} - \sum_{1,\tau}^{h}\sum_{1,\tau}^{k} \Delta_{k,h}a_{s,h}a_{k,r},
\]
so ist $A_{s,r}$ nach Bd. I, \S.\,23 eine $(\tau+1)$-reihige Determinante der
Matrix (22) und also nach Voraussetzung durch $p$ theilbar.

    Nach dieser Bezeichnungsweise erhalten wir aus (23):
\[
\Delta (a_{1,r}y_1 + a_{2,r}y_2 + \cdots + a_{\varrho,r}y_\varrho)
\equiv A_{\tau+1,r}y_{\tau+1} + \cdots + A_{\varrho,r}y_\varrho\ (\mathrm{mod}\ p^\pi),
\]
und es zeigt sich also, dass, wenn $\tau < \varrho$ ist, die Congruenzen (20)
schon befriedigt sind, wenn $y_{\tau+1}, \dots y_\varrho$ durch $p^{\pi-1}$ theilbar angenommen werden, weil alle $A_{s,r}$ durch $p$ theilbar sind. Dies
ist also nicht zulässig. Es muss also $\tau = \varrho$ sein, und dies genügt
auch, weil dann (23) sich auf
\[
\Delta y_1 \equiv 0,\ \Delta y_2 \equiv 0 \dots \Delta y_\varrho \equiv 0\ (\mathrm{mod}\ p^\pi)
\]
reducirt, und folglich, da $\Delta$ durch $p$ nicht theilbar ist, $y_1, y_2, \dots y_\varrho$
durch $p^\pi$ theilbar sein müssen.

    Hierdurch ist die Forderung 6. in eine andere Form gebracht, die wir, indem wir zu unserer ursprünglichen Bezeichnung
in (20) zurückkehren, so aussprechen können:
% [sic: Weber prints Forderung 6; the determinant argument reformulates Forderung 5, while Forderung 6 is the separate primarity condition.]
%CONT>>


% ===== scan leaf 110, printed 89 =====
% !PRINTED 89
\weberpage{89}
%<<CONT

7. \gesperrt{Es muss die Anzahl} $\varrho$ \gesperrt{der durch eine Primzahl} $p$
\gesperrt{theilbaren Invarianten gleich oder kleiner als die Anzahl} $\mu$
\gesperrt{der Indexmoduln sein, und es muss sich aus der Matrix}

\weq{24}{
\begin{gathered}
e_{1,1}\,i_{1,1}, \quad e_{2,1}\,i_{2,1}, \; \dots \; e_{\varrho,1}\,i_{\varrho,1} \\
e_{1,2}\,i_{1,2}, \quad e_{2,2}\,i_{2,2}, \; \dots \; e_{\varrho,2}\,i_{\varrho,2} \\
\cdot\ \cdot\ \cdot\ \cdot\ \cdot\ \cdot\ \cdot\ \cdot\ \cdot\ \cdot \\
e_{1,\mu}\,i_{1,\mu}, \quad e_{2,\mu}\,i_{2,\mu}, \; \dots \; e_{\varrho,\mu}\,i_{\varrho,\mu}
\end{gathered}
}

\gesperrt{wenigstens eine durch} $p$ \gesperrt{untheilbare Determinante von}
$\varrho$ \gesperrt{Reihen bilden lassen.}

Dies involvirt zun\"achst eine Forderung f\"ur die $i_{k,h}$, d.\,h. also
in letzter Instanz eine Bedingung f\"ur die Zahl $m$.

Alle Glieder irgend einer $\varrho$-reihigen Determinante der Matrix (24)
enthalten einen Factor der Form

\[
i_{1,h_1}\,i_{2,h_2}, \; \dots \; i_{\varrho,h_\varrho},
\]

worin $h_1, h_2, \dots h_\varrho$ irgend eine Combination von $\varrho$
verschiedenen Zahlen der Reihe $1, 2, \dots \mu$ bedeuten. W\"aren nun alle
diese Producte durch $p$ theilbar, so w\"aren sicher alle Determinanten der
Matrix (24) von $\varrho$ Reihen auch durch $p$ theilbar, und es muss also
wenigstens eine solche Combination geben, die nicht durch $p$ theilbar ist.
Weil aber die $i_{k,h}$ nur Potenzen von $p$ sein k\"onnen, so muss

\[
i_{1,h_1} = 1, \quad i_{2,h_2} = 1, \; \dots \; i_{\varrho,h_\varrho} = 1
\]

sein. Geht man nun auf die Bedeutung der $i_{k,h}$ in (11) zur\"uck, so
k\"onnen wir die letzte Forderung f\"ur den Grad $m$ folgendermaassen
ausdr\"ucken:

f) \gesperrt{Wenn eine Primzahl} $p$ \gesperrt{in} $\varrho$
\gesperrt{Invarianten} $i_1, i_2, \dots i_\varrho$ \gesperrt{aufgeht, so muss
die Anzahl} $\mu$ \gesperrt{der Indexmoduln gleich oder gr\"osser als}
$\varrho$ \gesperrt{sein, und es m\"ussen sich} $\varrho$ \gesperrt{dieser
Indexmoduln} $c_{h_1}, c_{h_2}, \dots c_{h_\varrho}$ \gesperrt{so
ausw\"ahlen lassen, dass} $c_{h_1}$ \gesperrt{durch} $i_1$, $c_{h_2}$
\gesperrt{durch} $i_2 \dots$, $c_{h_\varrho}$ \gesperrt{durch} $i_\varrho$
\gesperrt{theilbar ist} $^1$).

\noindent\rule{4cm}{0.4pt}

\footnotesize
$^1$) Der Nachweis der Thatsache, dass dieser Forderung immer auf
unendlich viele Arten entsprochen werden kann, st\"utzt sich auf den Satz
der Zahlentheorie, dass unter den Zahlen einer arithmetischen Progression,
deren Anfangsglied und Differenz ohne gemeinsamen Theiler sind, unendlich
viele Primzahlen vorkommen, ein Satz, der bis jetzt nur von
\normalsize
%CONT>>


% ===== scan leaf 111, printed 90 =====
% !PRINTED 90
\weberpage{90}
%<<CONT

Hiermit aber ist alles ersch\"opft, was wir von $m$ fordern m\"ussen. Denn
wenn diese Bedingung f) befriedigt ist, so brauchen wir z.\,B. nur die
$e_{1,h_1}, e_{2,h_2}, \dots e_{\varrho,h_\varrho}$ durch $p$ untheilbar,
die \"ubrigen $e$ durch $p$ theilbar anzunehmen; dann ist die Determinante

\[
\begin{vmatrix}
e_{1,h_1}\,i_{1,h_1}, & \dots & e_{1,h_\varrho}\,i_{1,h_\varrho} \\
\cdot\ \cdot\ \cdot & \cdots & \cdot\ \cdot\ \cdot \\
e_{\varrho,h_1}\,i_{\varrho,h_1}, & \dots & e_{\varrho,h_\varrho}\,i_{\varrho,h_\varrho}
\end{vmatrix}
\equiv e_{1,h_1}\,e_{2,h_2} \dots e_{\varrho,h_\varrho} \pmod{p}
\]

also sicher durch $p$ untheilbar.

Und bei dieser Annahme kann auch noch die Bedingung 6. befriedigt werden.
Um dies einzusehen, bezeichne man mit $p^\pi$ irgend eine der Invarianten,
etwa $i_1$, und mit $c_h$ den Indexmodul, der nach f) durch $p^\pi$
theilbar ist, so dass

\weq{25}{c_h = p^\pi c_{1,h}}

ist. Ist dann $q_h$ der zum Indexmodul $c_h$ geh\"orige Primtheiler von
$m$, so ist $p$ entweder gleich $q_h$ oder ein Theiler von $q_h - 1$. Nach
der zuletzt gemachten Annahme ist $e_{1,h}$ nicht durch $p$ theilbar,
w\"ahrend $e_{2,h}, e_{3,h}, \dots e_{\nu,h}$ durch $p$ theilbar sind. Es
sollen dann die Producte

\[
c_{1,h}\,e_{1,h}, \quad c_{2,h}\,e_{2,h}, \; \dots \; c_{\nu,h}\,e_{\nu,h}
\]

nicht alle durch $q_h$ oder durch $q_h - 1$ theilbar sein.

Ist zun\"achst $q_h$ ein einfacher Theiler von $m$, so ist $c_h = q_h -
1$, und wegen (25) ist $c_{1,h}$ nicht durch $q_h - 1$ theilbar. Man kann
also $e_{1,h}$ noch so annehmen, dass $c_{1,h}\,e_{1,h}$ nicht durch
$q_h - 1$ theilbar ist. Ist aber $q_h$ ein $\varkappa$-facher Theiler von
$m$ und $\varkappa > 1$, und ist $q_h = p$, so ist, wenn $p$ ungerade
ist, $c_h = p^{\varkappa - 1}(p-1)$, und nach der Voraussetzung b)
$\pi \leqq \varkappa - 1$. Wenn also $c_h$ durch $p^\pi$ theilbar sein
soll, so muss $\pi = \varkappa - 1$ sein, und (25) ergiebt $c_{1,h} =
p - 1$. Ist aber $p = 2$, so ist $c_h = 2$ oder $= 2^{\varkappa - 1}$,
und es muss also nach der Voraussetzung d) $\pi = 1$ oder $= \varkappa -
2$ sein, und es ist $c_{1,h} = 1$. Man kann also
% [sic: Weber prints pi <= kappa-1; hypothesis b) gives pi >= kappa-1, which is what makes divisibility force equality.]
% [sic: Weber prints c_h=2^{kappa-1} in the second 2-adic case; the index modulus is 2^{kappa-2}, as stated in Theorem 4, and that value is required for c_{1,h}=1.]

\noindent\rule{4cm}{0.4pt}

\footnotesize
\gesperrt{Dirichlet} mit Anwendung der Integralrechnung bewiesen ist. (Abhandlungen
der Berliner Akademie 1837; \gesperrt{Dirichlet's Werke}, Bd.\,I, Nr.\,21;
\gesperrt{Dirichlet-Dedekind}, Vorlesungen \"uber Zahlentheorie, Supplement VI;
\gesperrt{Bachmann}, Analytische Zahlentheorie, Abschn.\,IV). Behalten wir die oben
benutzte Bezeichnung bei, so kann man, um f\"ur ein gegebenes $p$ der
Forderung f) zu gen\"ugen, $p^{\pi+1}$ als Factor in $m$ aufnehmen. Dann
aber m\"ussen, wenn $\varrho > 1$ ist, noch $\varrho - 1$ Primfactoren
$q$ in $m$ vorkommen, die an Congruenzen von der Form $q \equiv 1
\pmod{i_1} \dots$ gebunden sind.
\normalsize
%CONT>>


% ===== scan leaf 112, printed 91 =====
% !PRINTED 91
\weberpage{91}

%<<CONT
auch in diesen F\"allen $e_{1,h}$ so annehmen, dass $c_{1,h}\,e_{1,h}$
durch $q_h$ nicht theilbar wird.

Die Forderungen, die hiermit f\"ur die $e_{k,h}$ gestellt sind, bestehen
also nur darin, dass sie durch gewisse Primzahlen theilbar, durch andere
nicht theilbar sein sollen, und sind also noch auf unendlich viele Arten
mit einander vertr\"aglich. Aber diese letzte Annahme \"uber die $e_{k,h}$
hatte nur den Zweck, zu zeigen, dass die fr\"uheren allgemeineren
Forderungen immer befriedigt werden k\"onnen. Es kann noch viele andere
Arten geben, ihnen zu gen\"ugen.

\noindent\rule{4cm}{0.4pt}

\begin{center}
\gesperrt{Uebersicht der Resultate.}
\end{center}

Wenn es sich darum handelt, alle Kreistheilungsk\"orper, und jeden nur
einmal, und zwar auf m\"oglichst einfache Art, durch
Kreistheilungsperioden darzustellen, so verfahre man so:

Man nehme ein beliebiges System von Primzahlpotenzen als Invarianten

\[
i_1, i_2, \dots i_v
\]

an, und w\"ahle dann eine Zahl $m$ nach folgenden Bedingungen.

Wenn $p^\pi$ die h\"ochste Potenz einer Primzahl $p$ ist, die unter den
Invarianten vorkommt, so nehme man $p$ h\"ochstens $(\pi+1)$-, oder,
wenn $p = 2$ ist, $(\pi+2)$ mal in $m$ auf.

Eine Primzahl $q$, die gar nicht in den Invarianten aufgeht, darf nur
einfach in $m$ aufgenommen werden; aber nur solche einfache Primfactoren
$q$ darf $m$ haben, bei denen $q - 1$ durch einen der Primfactoren $p$
der Invarianten theilbar ist. Also kann auch $p$ selbst nur dann einfach
in $m$ vorkommen, wenn $p - 1$ durch eines der anderen $p$ theilbar ist.

Ferner muss $m$ noch der Bedingung gen\"ugen, dass unter den
Indexmoduln

\[
c_1, c_2, \dots c_\mu
\]

$\varrho$ verschiedene, wie $c_1, c_2, \dots c_\varrho$, durch $i_1, i_2,
\dots i_\varrho$ theilbar sind, wenn $i_1, i_2, \dots i_\varrho$ Potenzen
von einer und derselben Primzahl sind. Dann bestimme man die Gr\"ossen
$i_{k,h}$ und $c_{k,h}$ nach den Formeln (11) und theile die gegebenen
Invarianten in Systeme ein, so dass alle Potenzen einer und derselben
Primzahl in einem Systeme vereinigt sind.

F\"ur jedes dieser Systeme bilde man die Matrix (24) und w\"ahle die
Zahlen $c_{k,h}$ so, dass aus jeder solchen Matrix eine
% [sic: Weber prints c_{k,h}; these quantities were fixed by (11), whereas the matrix entries and the ensuing construction require choosing e_{k,h}.]
%CONT>>


% ===== scan leaf 113, printed 92 =====
% !PRINTED 92
\weberpage{92}
\pdfbookmark[2]{§ 22. Bestimmung der Gruppe A}{bm26}

%<<CONT
Determinante von $\varrho$ Reihen gebildet werden kann, die durch die
betreffende Primzahl nicht theilbar ist, und dass nicht alle Producte

\[
c_{1,h}\,e_{1,h}, \quad c_{2,h}\,e_{2,h}, \; \dots \; c_{\nu,h}\,e_{\nu,h}
\]

durch $q_h$ oder durch $q_h - 1$ theilbar werden. Dann setze man

\[
\gamma_{k,h} = c_{k,h}\,e_{k,h},
\]

und hat so nach (7) die Indices einer Basis einer Abel'schen Gruppe
$\mathfrak{B}$, deren reciproke Gruppe $\mathfrak{A}$ ein prim\"arer
Theiler der ganzen Gruppe $\mathfrak{N}$ ist.

\noindent\rule{4cm}{0.4pt}

\begin{center}
\S.\,22.

\gesperrt{Bestimmung der Gruppe $\mathfrak{A}$.}
\end{center}

Das letzte Ziel dieser Untersuchung ist nicht sowohl die Bestimmung der
Gruppe $\mathfrak{B}$, als die der Gruppe $\mathfrak{A}$, aus der man
direct die Kreistheilungsperioden $\eta = \Sigma r^a$ bilden kann. Diese
Aufgabe kann man auf folgende Art l\"osen: Nach (1) sind die Indices
$\alpha$ der Zahlen in $\mathfrak{A}$ von den $\beta$ abh\"angig durch
die Gleichung

\weq{1}{\omega_1^{\alpha_1\beta_1}\,\omega_2^{\alpha_2\beta_2} \dots
\omega_\mu^{\alpha_\mu\beta_\mu} = 1,}

worin $\omega_1, \omega_2, \dots \omega_\mu$ feste primitive
Einheitswurzeln der Grade $c_1, c_2, \dots c_\mu$ bedeuten. Versteht man
aber unter $C$ das kleinste gemeinschaftliche Vielfache von $c_1, c_2,
\dots c_\mu$ und setzt

\weq{2}{C = c_1 c_1' = c_2 c_2' = \dots = c_\mu c_\mu',}

und bezeichnet mit $\omega$ eine primitive $C^{\mathrm{te}}$
Einheitswurzel, so kann man statt (1) auch setzen

\[
\omega^{c_1'\alpha_1\beta_1 + c_2'\alpha_2\beta_2 + \dots +
c_\mu'\alpha_\mu\beta_\mu} = 1,
\]

und bekommt daher f\"ur die $\alpha$ die Congruenz

\[
c_1'\alpha_1\beta_1 + c_2'\alpha_2\beta_2 + \dots +
c_\mu'\alpha_\mu\beta_\mu \equiv 0 \pmod{C},
\]

die gleichbedeutend ist mit der Bedingung, dass die Summen

\[
\frac{\alpha_1\beta_1}{c_1} + \frac{\alpha_2\beta_2}{c_2} + \dots +
\frac{\alpha_\mu\beta_\mu}{c_\mu} = \sum_{1,\mu}^{h} \frac{\alpha_h
\beta_h}{c_h}
\]

ganze Zahlen sein m\"ussen. Wenn man darin f\"ur $\beta_k$ die
Ausdr\"ucke \S.\,21, (8) substituirt, so folgt, dass auch

\[
\sum^{k} x_k \sum^{h} \frac{\alpha_h\,\gamma_{k,h}}{c_h},
\]
%CONT>>


% ===== scan leaf 114, printed 93 =====
% !PRINTED 93
\weberpage{93}
%<<CONT
oder, da die ganzen Zahlen $x_k$ beliebig sind, die Summen
\[
\sum_{1,\mu}^{h}\frac{\alpha_h\,\gamma_{k,h}}{c_h}\qquad k=1,2,\ldots \nu
\]
ganze Zahlen sein müssen. Setzt man dann nach (15) und (11) \S.\,21
\[
\gamma_{k,h}=c_{k,h}\,e_{k,h},\quad c_h=\delta_{k,h}\,c_{k,h},\quad i_k=\delta_{k,h}\,i_{k,h},
\]
so folgt, dass
\[
\sum^{h}\frac{\alpha_h\,e_{k,h}}{\delta_{k,h}}=\sum^{h}\frac{\alpha_h\,i_{k,h}\,e_{k,h}}{i_k}
\]
ganze Zahlen sein müssen, und diese Forderung ist gleichbedeutend mit dem System der Congruenzen
\[
\sum^h \alpha_h\,i_{k,h}\,e_{k,h}\equiv 0\ (\mathrm{mod}\ i_k),
\]
oder, ausführlicher geschrieben
\weq{3}{\begin{aligned}
&e_{1,1}i_{1,1}\alpha_1+e_{1,2}i_{1,2}\alpha_2+\cdots+e_{1,\mu}i_{1,\mu}\alpha_\mu\equiv 0\ (\mathrm{mod}\ i_1)\\
&e_{2,1}i_{2,1}\alpha_1+e_{2,2}i_{2,2}\alpha_2+\cdots+e_{2,\mu}i_{2,\mu}\alpha_\mu\equiv 0\ (\mathrm{mod}\ i_2)\\
&\cdots\cdots\cdots\cdots\cdots\cdots\cdots\cdots\cdots\cdots\cdots\cdots\cdots\cdots\cdots\cdots\cdots\cdots\cdots\cdots\\
&e_{\nu,1}i_{\nu,1}\alpha_1+e_{\nu,2}i_{\nu,2}\alpha_2+\cdots+e_{\nu,\mu}i_{\nu,\mu}\alpha_\mu\equiv 0\ (\mathrm{mod}\ i_\mu). % [sic: Weber prints i_\mu; the row pattern requires i_\nu.]
\end{aligned}}
Durch diese Congruenzen sind, wie es sein muss, die Zahlen $\alpha_1, \alpha_2, \ldots \alpha_\mu$ nur nach den Moduln $c_1, c_2, \ldots c_\mu$ bestimmt, weil $i_{k,h}c_h=i_k c_{k,h}$ durch $i_k$ theilbar ist, und \gesperrt{um die Indices aller Zahlen der Gruppe $\mathfrak{A}$ zu finden, hat man einfach alle Lösungen der Congruenzen (3) aufzusuchen.}

\begin{center}
\rule{4cm}{0.4pt}
\end{center}


% ===== scan leaf 115, printed 94 =====
\weberpage{94}% page number not printed on this page
\pdfbookmark[1]{Vierter Abschnitt. Cubische und biquadratische Abel'sche Körper.}{bm27}
\pdfbookmark[2]{§ 23. Cubische Kreistheilungskörper}{bm28}
% !PRINTED none
\begin{center}
\gesperrt{Vierter Abschnitt.}
\end{center}

\begin{center}
\gesperrt{Cubische und biquadratische Abel'sche K\"orper.}
\end{center}

\begin{center}
\rule{3cm}{0.4pt}
\end{center}

\begin{center}
\S.\,23.
\end{center}

\begin{center}
\gesperrt{Cubische Kreistheilungsk\"orper.}
\end{center}

Die allgemeinen Untersuchungen, die in den letzten Paragraphen dargestellt sind, gestatten mannigfache Anwendungen auf specielle Fälle, von denen die einfachsten hier betrachtet werden sollen. Es ist dabei bemerkenswerth, dass schon die einfachsten Fälle, wenn man sie direct angreift, fast in dieselben Schwierigkeiten hineinführen, die wir durch die allgemeine Untersuchung überwunden haben.

Wir wollen zuerst die Aufgabe behandeln, alle \gesperrt{cubischen Kreistheilungskörper, also alle Kreistheilungsperioden, die einer rationalen Gleichung $3^{\mathrm{ten}}$ Grades genügen,} aufzufinden, und zwar so, dass von mehreren Ausdrücken, deren dieselbe Grösse fähig ist, der einfachste gesetzt wird.

Wir haben hier nur eine Invariante $i_1=3$, und nach den Vorschriften in der Zusammenstellung des \S.\,21 dürfen in $m$ aufgenommen werden der Factor 9, nicht aber 3, ferner Primzahlen in beliebiger Menge von der Form $6N+1$, also die Primzahlen 7, 13, 19, 31, 37, 43, 61, 67, 73, 79, 97 . . ., und so bekommen wir also als geeignete Werthe von $m$ unter 200:
\[
\begin{aligned}
m &= 7,\ 9,\ 13,\ 19,\ 31,\ 37,\ 43,\ 61,\ \mathbf{63},\ 67,\ 73,\ 79,\ \mathbf{91},\ 97,\ 103,\ 109,\\
&\quad \mathbf{117},\ 127,\ \mathbf{133},\ 139,\ 151,\ 157,\ 163,\ \mathbf{171},\ 181,\ 193,\ 199,
\end{aligned}
\]
unter denen die Zahlen mit mehreren Primtheilern durch fetten Druck ausgezeichnet sind. Die kleinste Zahl $m$, in der mehr als zwei Primzahlen aufgehen, ist $7.9.13=819$.


% ===== scan leaf 116, printed 95 =====
% !PRINTED 95
\weberpage{95}
Verstehen wir also unter $q_1, q_2, q_3 \ldots$ Zahlen, die aus der Reihe der Zahlen 9 und der Primzahlen 7, 13, 19, 31 . . . genommen sind, so ist der allgemeine Ausdruck für $m$:
\weq{1}{m = q_1\,q_2\,q_3\cdots q_\mu.}
Die Indexmoduln sind, wenn $q_1=9$ ist, $c_1=6$, $c_2=q_2-1$, $c_3=q_3-1 \ldots$ Wenn 9 und 7 unter den Factoren von $m$ nicht vorkommen, so fällt der Indexmodul 6 weg.

Die Grössen $\delta_{1,1}, \delta_{1,2}, \ldots \delta_{1,\mu}$ sind die grössten gemeinschaftlichen Theiler von $i_1=3$ mit den $c_1, c_2, \ldots c_\mu$; sie sind also alle $=3$, und es folgt nach \S.\,21, (11):
\[
i_{1,h}=1,\quad c_{1,h}=\tfrac{1}{3}\,c_1.
\]
% [sic: Weber prints c_1 on the right; because h is arbitrary and c_{1,h}=c_h/delta_{1,h}, the intended subscript is h.]
Die Matrix \S.\,21, (24) besteht hier nur aus der einen Verticalreihe $e_{1,1}, e_{1,2}, \ldots e_{1,\mu}$, und diese Grössen kann man beliebig wählen, wenn nur wenigstens eine darunter ist, die durch 3 untheilbar ist. Diese Zahlen $e_{1,1}, e_{1,2}, \ldots e_{1,\mu}$ sind übrigens nur nach dem Modul 3 bestimmt, und können also gleich 0, 1 oder 2 angenommen werden. Die Bedingung 6. des \S.\,21 verlangt aber auch, dass, wenn $q_1=9$ ist, $c_{1,1}e_{1,1}=2e_{1,1}$ nicht durch 3 theilbar sein soll, und wenn $q_2$ eine Primzahl ist, dass
\[
c_{1,2}\,e_{1,2}=\tfrac{1}{3}(q_2-1)\,e_{1,2}
\]
nicht durch $q_2-1$ theilbar sei, d. h. also, es darf \gesperrt{keine der Zahlen} $e_{1,1}, e_{1,2}, \ldots e_{1,\mu}$ durch 3 theilbar sein, und wir können für jede von ihnen $\pm 1$ wählen.

Um also die Indices $\alpha_1, \alpha_2, \ldots \alpha_\mu$ aller Zahlen $a$ zu finden, die in der Periode
\weq{2}{\eta = \sum r^a}
vorkommen können, hat man nach \S.\,22, (3) alle Zahlen $\alpha$ (nach dem Modul 3) zu nehmen, die einer Congruenz
\weq{3}{\pm\alpha_1\pm\alpha_2\pm\alpha_3\cdots\pm\alpha_\mu\equiv 0\ (\mathrm{mod}\ 3)}
genügen, wobei für \gesperrt{eine bestimmte Gruppe eine Vorzeichen-Combination festgehalten werden muss}.

Die Anzahl der möglichen cubischen Körper hängt bei gegebenem $m$ nur von der Zahl $\mu$ der Indexmoduln ab und beträgt, da eines der Vorzeichen in (3) willkürlich angenommen werden kann, $2^{\mu-1}$.

Ganz ähnlich würde sich das Resultat gestalten, wenn man statt 3 eine andere Primzahl $p$ für $i_1$ setzen würde. Es würde dann nur in die Congruenz (3) an Stelle der doppelten Vor-
%CONT>>


% ===== scan leaf 117, printed 96 =====
% !PRINTED 96
\weberpage{96}
%<<CONT
zeichen je ein volles Restsystem von $p$ (mit Ausschluss der 0) treten.

Wir wollen aber bei der Annahme $i_1=3$ stehen bleiben, und unter $\varepsilon_1, \varepsilon_2, \ldots \varepsilon_\mu$ irgend eine Combination der Zahlen $\pm 1$ verstehen. Dann können wir die Congruenz (3) auch so schreiben
\weq{4}{\varepsilon_1\alpha_1+\varepsilon_2\alpha_2+\cdots+\varepsilon_\mu\alpha_\mu\equiv 0\ (\mathrm{mod}\ 3).}

Die Zeichencombination $\varepsilon_1, \varepsilon_2, \ldots \varepsilon_\mu$, in der ein Zeichen willkürlich angenommen werden kann, da die gleichzeitige Aenderung aller Vorzeichen in (4) nichts ändert, bestimmt die einzelne Gruppe $\mathfrak{A}$, und offenbar sind diese Gruppen $\mathfrak{A}$ auch alle von einander verschieden. Denn ist z. B. in einer Gruppe $\varepsilon_1=\varepsilon_2=+1$, in einer anderen $\varepsilon_1=+1$, $\varepsilon_2=-1$, so enthält die erste die Zahlencombination:
\[
\alpha_1\equiv 1,\quad \alpha_2\equiv -1,\quad \alpha_3\equiv 0,\quad \ldots\ \alpha_\mu\equiv 0\ (\mathrm{mod}\ 3),
\]
die in der zweiten nicht vorkommt.

Um die Gruppen $\mathfrak{A}$ definitiv zu finden, legt man ein System primitiver Wurzeln $g_1, g_2, \ldots g_\mu$ von $q_1, q_2, \ldots q_\mu$ zu Grunde und bestimmt $a$ nach dem Modul $m$ aus den Congruenzen:
\weq{5}{\begin{aligned}
a &\equiv g_1^{\alpha_1}\ (\mathrm{mod}\ q_1)\\
&\equiv g_2^{\alpha_2}\ (\mathrm{mod}\ q_2)\\
&\cdots\cdots\cdots\cdots\cdots\\
&\equiv g_\mu^{\alpha_m}\ (\mathrm{mod}\ q_\mu). % [sic: Weber prints \alpha_m; the pattern requires \alpha_\mu.]
\end{aligned}}
Setzen wir, indem wir unter $\sigma$ einen der Reste 0, $\pm 1$ verstehen,
\weq{6}{\sigma \equiv \varepsilon_1\nu_1+\varepsilon_2\nu_2+\cdots+\varepsilon_\mu\nu_\mu\ (\mathrm{mod}\ 3),}
so besteht $\mathfrak{A}$ aus allen Zahlen $a$, deren Indices $\nu_1=\alpha_1, \nu_2=\alpha_2, \ldots \nu_\mu=\alpha_\mu$ der Bedingung $\sigma=0$ genügen. Die beiden anderen Werthe $\sigma=1$, $\sigma=-1$ bestimmen die beiden Nebengruppen
$\mathfrak{A}'$, $\mathfrak{A}''$ von $\mathfrak{A}$, so dass
\[
\mathfrak{N} = \mathfrak{A} + \mathfrak{A}' + \mathfrak{A}''
\]
ist. Ist $a'$ und $a''$ je eine Zahl aus $\mathfrak{A}'$ und $\mathfrak{A}''$, so ist $\mathfrak{A}'=\mathfrak{A}a'$, $\mathfrak{A}''=\mathfrak{A}a''$.

Wir bezeichnen nun mit $\eta$ die Kreistheilungsperiode, die zu $\mathfrak{A}$ gehört, und mit $\eta'$, $\eta''$ die conjugirten Perioden, also
\weq{7}{\eta = \sum r^a,\quad \eta' = \sum r^{a'},\quad \eta'' = \sum r^{a''}.}

Wenn dann $\varrho$ eine imaginäre dritte Einheitswurzel ist, so erhalten wir die Resolvente
\[
(\varrho, \eta) = \eta + \varrho\,\eta' + \varrho^2\,\eta'',
\]
%CONT>>


% ===== scan leaf 118, printed 97 =====
% !PRINTED 97
\weberpage{97}
%<<CONT
und dafür können wir mit Benutzung der Bezeichnung (6) setzen

\weq{8}{(\varrho, \eta) = \overset{n}{\Sigma}\, \varrho^{\sigma} r^n,}

worin $n$ alle Zahlen der Gruppe $\mathfrak{N}$ durchläuft, und in $\sigma$ für
$\nu_1, \nu_2 \ldots \nu_\mu$ jedesmal die Indices von $n$ zu setzen sind.

Es mögen nun mit $r_1, r_2, \ldots r_\mu$ primitive Einheitswurzeln der Grade
$q_1, q_2, \ldots q_\mu$ bezeichnet sein, und

\[ r = r_1 r_2, \ldots r_\mu \]

gesetzt werden. Dann lässt sich der Ausdruck (8) in ein Product zerlegen,
nämlich, wenn wir

\[ n \equiv n_h \ (\mathrm{mod}\ q_h) \]

setzen, in:

\weq{9}{(\varrho, \eta) = \Sigma\, \varrho^{\varepsilon_1 \nu_1} r_1^{n_1}\, \Sigma\, \varrho^{\varepsilon_2 \nu_2} r_2^{n_2} \ldots \Sigma\, r_\mu^{\varepsilon_\mu \nu_\mu} r_\mu^{n_\mu}.} % [sic: Weber prints r_\mu as the final first base; the parallel factors require \varrho.]

Die Factoren dieses Productes sind nun genau die Resolventen, die wir in
\S.\,172 des ersten Bandes untersucht haben, und wir erhalten, wenn wir

\weq{10}{\overset{n_1}{\Sigma}\, \varrho^{\nu_h} r^{n_h} = (\varrho, \eta_h), \quad h = 1, 2, \ldots \mu} % [sic: Weber prints the upper bound n_1 and the base r^{n_h}; the general h factor requires n_h and r_h^{n_h}.]

setzen:

\weq{11}{(\varrho, \eta) = (\varrho^{\varepsilon_1}, \eta_1)(\varrho^{\varepsilon_2}, \eta_2) \ldots (\varrho^{\varepsilon_\mu}, \eta_\mu).}

Wenn wir hierauf die Formeln (5), (6), (25), (26) des eben angeführten
Paragraphen anwenden, so erhalten wir:

\weq{12}{(\varrho, \eta)(\varrho^{-1}, \eta) = m,}

\weq{13}{(\varrho, \eta)^3 = m(a + b\varrho), \quad (\varrho^2, \eta)^3 = m(a + b\varrho^2),}

worin $a$ und $b$ gewisse ganze Zahlen sind, die sich aus den genannten
Formeln berechnen lassen, die aber nach der Wahl der $\varepsilon$
verschieden ausfallen.

Aus (12) und (13) folgt dann noch:

\weq{14}{m = a^2 - ab + b^2.}

Hiernach können wir die cubische Gleichung aufstellen, deren Wurzeln die
Grössen $\eta, \eta', \eta''$ sind. Sie möge lauten:

\weq{15}{\eta^3 - \alpha \eta^2 + \beta \eta - \gamma = 0,}

worin $\alpha, \beta, \gamma$ ganze Zahlen sind. Was zunächst $\alpha$
betrifft, so ist es gleich der Summe aller primitiven $m^{\mathrm{ten}}$
Einheitswurzeln, also

\[ \alpha = \Sigma\, r_1\, \Sigma\, r_2 \ldots \Sigma\, r_\mu, \]

und dieser Werth ist $= 0$ oder $= (-1)^{\mu}$,
je nachdem 9 unter den Factoren von $m$ vorkommt oder nicht. Die beiden
anderen Co\"efficienten lassen sich ganz so berechnen, wie an der er-
%CONT>>


% ===== scan leaf 119, printed 98 =====
% !PRINTED 98
\weberpage{98}
%<<CONT
wähnten Stelle gezeigt ist. Wir wollen hier die Rechnung in etwas
veränderter Form anordnen. Wir setzen:

\[ \xi = \eta - \frac{\alpha}{3}, \]

was zur Folge hat, dass $(\varrho, \eta) = (\varrho, \xi)$ ist, wegen
$1 + \varrho + \varrho^2 = 0$. Dann ergiebt sich für $\xi$ die cubische
Gleichung:

\weq{16}{\xi^3 + P\xi - Q = 0,}

worin

\weq{17}{P = \beta - \frac{\alpha^2}{3}, \quad Q = \frac{2\alpha^3}{27} - \frac{\alpha\beta}{3} + \gamma.}

Nun ergiebt sich aus (12):

\[ \begin{aligned}
(\varrho, \xi)(\varrho^{-1}, \xi) = m &= \xi^2 + \xi'^2 + \xi''^2 - \xi\xi' - \xi\xi'' - \xi'\xi'' \\
&= -3P
\end{aligned} \]

oder

\weq{18}{P = -\frac{m}{3},}

und aus (13), nach einigen einfachen Rechnungen, wenn

\[ S = \xi^2\xi' + \xi'^2\xi'' + \xi''^2\xi, \quad S' = \xi\xi'^2 + \xi'\xi''^2 + \xi''\xi^2 \]

gesetzt wird:

\[ \begin{aligned}
m(a + b\varrho) &= 9Q + 3\varrho S + 3\varrho^2 S' \\
m(a + b\varrho^2) &= 9Q + 3\varrho^2 S + 3\varrho S' \\
0 &= 9Q + 3S + 3S',
\end{aligned} \]

woraus durch Addition:

\weq{19}{27Q = m(2a - b),}

und durch Subtraction der beiden ersten:

\weq{20}{mb = 3(S - S') = 3(\xi - \xi')(\xi'' - \xi)(\xi'' - \xi').}

Wenn $m$ durch 9 theilbar ist, so ist $\alpha = 0$, und $P$ und $Q$ sind
\gesperrt{ganze} Zahlen und $P$ durch 3 theilbar. Ist aber $m$ nicht durch 9 theilbar,
so sind erst $3P$ und $27Q$ ganze, nicht durch 3 theilbare Zahlen.

Ist $D$ die Discriminante der Gleichung (15), also eine ganze Zahl, so ist
auch

\weq{21}{\begin{aligned}
\sqrt{D} &= (\eta - \eta')(\eta'' - \eta)(\eta'' - \eta') = (\xi - \xi')(\xi'' - \xi)(\xi'' - \xi') \\
&= \frac{mb}{3}
\end{aligned}}

eine ganze Zahl, und also ist, wenn $m$ nicht durch 9 theilbar ist, $b$
durch 3 theilbar. Nun ist nach (14)
%CONT>>


% ===== scan leaf 120, printed 99 =====
% !PRINTED 99
\weberpage{99}
%<<CONT
\weq{22}{4m = (2a - b)^2 + 3b^2,}

und wenn $m$ durch 9 theilbar ist, so ist $2a - b$ nach (19) durch 3
theilbar, und aus (22) folgt dann, dass auch in diesem Falle $b$ durch 3
theilbar ist. Wir können also in allen Fällen setzen

\[ 2a - b = A, \quad b = 3B, \]

so dass $A, B$ ganze Zahlen sind, die der Bedingung

\weq{23}{4m = A^2 + 27B^2}

genügen, und die Gleichung für $\xi$ wird dann

\weq{24}{\xi^3 - \frac{m}{3}\xi - \frac{mA}{27} = 0,}

und für die Discriminante dieser Gleichung folgt

\weq{25}{\sqrt{D} = mB.}

Um für die einfachsten Beispiele die Rechnung durchzuführen, nehmen wir für
$q_1 = 9, q_2 = 7, q_3 = 13$ aus \S.\,172 des ersten Bandes die Werthe:

\[ \begin{aligned}
q_1 &= 9, & a + b\varrho &= 3\varrho, \\
q_2 &= 7, & a + b\varrho &= -(1 + 3\varrho), \\
q_3 &= 13, & a + b\varrho &= -(4 + 3\varrho).
\end{aligned} \]

Daraus ergeben sich für $m = 63$ die beiden Werthe:

\[ \begin{aligned}
m = 63:\quad a + b\varrho &= -3\varrho(1 + 3\varrho) = 9 + 6\varrho, \\
a + b\varrho &= -3\varrho(1 + 3\varrho^2) = -9 - 3\varrho,
\end{aligned} \]

für

\[ \begin{aligned}
m = 91:\quad a + b\varrho &= (1 + 3\varrho)(4 + 3\varrho) = -5 + 6\varrho, \\
a + b\varrho &= (1 + 3\varrho)(4 + 3\varrho^2) = 10 + 9\varrho,
\end{aligned} \]

und für $m = 819$ die Werthe:

\[ \begin{aligned}
a + b\varrho &= 3\varrho(1 + 3\varrho)(4 + 3\varrho) = -3(6 + 11\varrho), \\
&= 3\varrho(1 + 3\varrho)(4 + 3\varrho^2) = -3(9 - \varrho), \\
&= 3\varrho(1 + 3\varrho^2)(4 + 3\varrho) = 3(9 + 10\varrho), \\
&= 3\varrho(1 + 3\varrho^2)(4 + 3\varrho^2) = 3(6 - 5\varrho).
\end{aligned} \]

Daraus finden sich die cubischen Gleichungen für $\eta$ in diesen drei
Fällen:

\[ \begin{aligned}
m = 63:\quad &\eta^3 - 21\eta - 28 = 0, & (\eta = \xi) \\
&\eta^3 - 21\eta + 35 = 0, \\
m = 91:\quad &\eta^3 - \eta^2 - 30\eta + 64 = 0, & (\eta = \xi + \tfrac{1}{3}) \\
&\eta^3 - \eta^2 - 30\eta - 27 = 0,
\end{aligned} \]
%CONT>>


% ===== scan leaf 121, printed 100 =====
% !PRINTED 100
\weberpage{100}
%<<CONT
\[ \begin{aligned}
m = 819:\quad &\eta^3 - 273\eta + 91 = 0, & (\eta = \xi) \\
&\eta^3 - 273\eta + 91 . 19 = 0, \\
&\eta^3 - 273\eta + 91 . 8 = 0, \\
&\eta^3 - 273\eta + 91 . 17 = 0.
\end{aligned} \]

Die Wurzeln dieser einzelnen Gleichungen findet man am besten aus der
Gleichung (9) oder (11), indem man darin nach Potenzen von $\varrho$
ordnet, dabei aber nur $\varrho^3 = 1$, nicht $\varrho^2 + \varrho + 1 = 0$
benutzt.

So bekommt man z.\,B. für den Fall $m = 63$:

\[ \begin{aligned}
(\varrho, \eta) &= (\eta_1 + \varrho\eta_1' + \varrho^2\eta_1'')(\eta_2 + \varrho\eta_2' + \varrho^2\eta_2'') \\
\text{oder} \\
&= (\eta_1 + \varrho\eta_1' + \varrho^2\eta_1'')(\eta_2 + \varrho\eta_2'' + \varrho^2\eta_2'),
\end{aligned} \]

worin

\[ \begin{aligned}
\eta_1 &= r_1 + r_1^{-1}, & \eta_1' &= r_1^2 + r_1^{-2}, & \eta_1'' &= r_1^4 + r_1^{-4} \\
\eta_2 &= r_2 + r_2^{-1}, & \eta_2' &= r_2^3 + r_2^{-3}, & \eta_2'' &= r_2^2 + r_2^{-2}
\end{aligned} \]

zu setzen ist (Bd.\,I, \S.\,172), und man findet so:

\[ \begin{aligned}
\eta &= \eta_1\eta_2 + \eta_1'\eta_2'' + \eta_1''\eta_2' \\
\text{oder} \\
\eta &= \eta_1\eta_2 + \eta_1'\eta_2' + \eta_1''\eta_2''.
\end{aligned} \]

Um die Gruppen $\mathfrak{A}$ daraus zu finden, hat man zwei Zahlen $x, y$
so zu bestimmen, dass $7x + 9y = 1$ wird, und dann

\[ r = r_1 r_2, \quad \text{also} \quad r_1 = r^{7x}, \ r_2 = r^{9y} \]

zu setzen. Nimmt man $x = 4, y = -3$ an, so erhält man für $x$ die beiden % [sic: Weber prints ``für x'', although x has just been fixed and values of \eta follow.]
Werthe von $\eta$:

\[ \begin{aligned}
\eta &= (r^{28} + r^{-28})(r^{27} + r^{-27}) + (r^7 + r^{-7})(r^9 + r^{-9}) \\
&\qquad + (r^{14} + r^{-14})(r^{18} + r^{-18}) \\
\eta &= (r^{28} + r^{-28})(r^{27} + r^{-27}) + (r^7 + r^{-7})(r^{18} + r^{-18}) \\
&\qquad + (r^{14} + r^{-14})(r^9 + r^{-9}),
\end{aligned} \]

also

\[ \mathfrak{A} = \pm 1, \ \pm 8, \ \pm 2, \ \pm 16, \ \pm 4, \ \pm 32, \]

oder

\[ \mathfrak{A} = \pm 1, \ \pm 8, \ \pm 11, \ \pm 25, \ \pm 5, \ \pm 23. \]

Die beiden Gruppen können durch die Potenzen von 2 und von 11 dargestellt
werden.

Die cubischen Gleichungen, deren Bildungsgesetze wir jetzt kennen gelernt
haben, enthalten mit ihren \gesperrt{Tschirnhausen-Transformationen alle
cubischen Kreistheilungsgleichungen.} Aber diese Gleichungen sind auch alle
von einander
%CONT>>


% ===== scan leaf 122, printed 101 =====
% !PRINTED 101
\weberpage{101}
\pdfbookmark[2]{§ 24. Biquadratische Kreistheilungskörper}{bm29}
%<<CONT
verschieden und können nicht durch \gesperrt{Tschirnhausen-Transformation} in einander übergeführt werden, weil sie zu primären Theilern von verschiedenen vollen Kreistheilungskörpern gehören, oder, wenn sie in demselben vollen Kreistheilungskörper enthalten sind, verschiedene Gruppen haben.

Man kann freilich noch aus anderen Einheitswurzeln Kreistheilungsperioden bilden, die zu cubischen Gleichungen führen. Diese Perioden können aber durch niedrigere Einheitswurzeln ausgedrückt werden, und sind nicht primär. So erhält man z.\,B., wenn $r$ eine 35\textsuperscript{ste} Einheitswurzel ist, eine Periode von 8 Gliedern, deren Exponenten aus den Potenzen von $-8$ gebildet sind:

\[
\eta = r + r^{-1} + r^6 + r^{-6} + r^8 + r^{-8} + r^{13} + r^{-13},
\]
die also auch einer cubischen Gleichung genügt.

Es ist aber
\[
1 + r^7 + r^{-7} + r^{14} + r^{-14} = 0,
\]
und wenn man diese Gleichung mit $(r^{15} + r^{-15})$ multiplicirt:
\[
r^{15} + r^{-15} + r + r^{-1} + r^6 + r^{-6} + r^8 + r^{-8} + r^{13} + r^{-13} = 0,
\]
also
\[
\eta = -(r^{15} + r^{-15}),
\]
was unter den Perioden der 7\textsuperscript{ten} Einheitswurzeln schon vorkommt. Diese Erscheinung erklärt sich daraus, dass $m = 35$ nicht die in \S.\,21 verlangte Eigenschaft hat, dass $q - 1$ für alle in $m$ aufgehenden Primzahlen $q$ durch 3 theilbar ist.

\begin{center}
\S.\,24.

\gesperrt{Biquadratische Kreistheilungskörper.}
\end{center}

Bei den biquadratischen Kreistheilungsgleichungen hat man zwei Arten zu unterscheiden. Wir können zwei Invarianten $i_1 = i_2 = 2$ oder nur eine Invariante $i_1 = 4$ annehmen. Wir wollen den ersten Fall voranschicken.

Für $m$ haben wir in diesem Falle nach der Zusammenstellung \S.\,21 folgende Bedingungen:

$m$ kann ungerade sein, durch 4 oder durch 8, aber durch keine höhere Potenz von 2 theilbar sein. Keine ungerade Primzahl kann mehr als einmal in $m$ aufgehen, und es müssen mindestens zwei Indexmoduln vorhanden sein; also müssen ausser in
%CONT>>


% ===== scan leaf 123, printed 102 =====
% !PRINTED 102
\weberpage{102}
%<<CONT
dem Falle $m = 8$ mindestens zwei verschiedene Primzahlen in $m$ enthalten sein.

Setzen wir:
\[
m = q_1\, q_2 \cdots q_\mu,
\]
wenn $m$ ungerade oder nur durch 4 theilbar ist,
\[
m = q_2\, q_3 \cdots q_\mu,
\]
wenn $m$ durch 8 theilbar ist, so dass $q_1 = 4$ sein kann und in der zweiten Formel $q_2 = 8$ ist, während die übrigen $q$ ungerade Primzahlen sind, so sind die Indexmoduln in den drei Fällen:
\[
\begin{aligned}
& q_1 - 1,\quad q_2 - 1,\ \ldots,\ q_\mu - 1\\
& 2\phantom{{}_1 -{}1,}\ ,\quad q_2 - 1,\ \ldots,\ q_\mu - 1\\
& 2,\quad 2\phantom{{}_1 -{}1}\ ,\quad q_3 - 1,\ \ldots,\ q_\mu - 1.
\end{aligned}
\]
Weil hiernach alle Indexmoduln durch 2 theilbar sind, so sind die Bedingungen für $m$ erschöpft.

Es ist [\S.\,21, (11)]:
\weq{1}{\delta_{k,h} = 2,\ i_{k,h} = 1,\ c_{k,h} = \tfrac{1}{2} c_h,\ k = 1, 2,\ h = 1, 2, \ldots \mu.}
Nun hat man die Grössen $e_{k,h}$ so zu bestimmen, dass aus der Matrix
\weq{2}{\begin{array}{cccc}
e_{1,1}, & e_{1,2}, & \ldots & e_{1,\mu}\\
e_{2,1}, & e_{2,2}, & \ldots & e_{2,\mu}
\end{array}}
wenigstens eine ungerade zweireihige Determinante gebildet werden kann.

Damit $\mathfrak{A}$ primärer Theiler von $\mathfrak{N}$ sei, kommt noch die Bedingung \S.\,21, 6. hinzu, wo nun zu unterscheiden ist, ob $m$ durch 8 theilbar ist oder nicht. Ist $m$ durch 8 theilbar, so ist die Anzahl der Indexmoduln um eins grösser, als die Anzahl der Primfactoren von $m$, und es giebt einen Indexmodul, dem kein Primfactor von $m$ entspricht; im anderen Falle stimmen beide Zahlen überein. Nach \S.\,21, 6. hat man die $e_{k,h}$ so zu wählen, dass in keinem der Paare
\[
e_{1,1}, e_{2,1};\ e_{1,2}, e_{2,2};\ \ldots\ e_{1,\mu}, e_{2,\mu},
\]
denen ein Primfactor von $m$ entspricht, beide Zahlen gerade sind. Wenn $m$ durch 8 theilbar ist, so ist ein dem Indexmodul $c_k = 2$ entsprechendes Paar darunter, dem keine solche Beschränkung auferlegt ist. Sind die Zahlen $e_{k,h}$ (nach dem Modul 2) so bestimmt, so
%CONT>>


% ===== scan leaf 124, printed 103 =====
% !PRINTED 103
\weberpage{103}
%<<CONT
erhält man nach \S.\,22 für die Indices der Zahlen in $\mathfrak{A}$ die beiden Congruenzen:
\weq{3}{\begin{aligned}
e_{1,1}\alpha_1 + e_{1,2}\alpha_2 + \cdots + e_{1,\mu}\alpha_\mu &\equiv 0\\
e_{2,1}\alpha_1 + e_{2,2}\alpha_2 + \cdots + e_{2,\mu}\alpha_\mu &\equiv 0
\end{aligned}\ \ (\mathrm{mod}\ 2).}
Um die Zahl der möglichen Gruppen $\mathfrak{A}$ zu ermitteln, bedenke man, dass man die Relationen (3), ohne ihren Inhalt zu ändern, durch zwei von einander unabhängige lineare Combinationen ersetzen kann. Danach kann man eines der Zahlenpaare in der Matrix (2), z.\,B. $e_{1,1}, e_{2,1} = 1, 0$ annehmen und so für (2) setzen:
\[
\begin{array}{cccc}
1, & e_{1,2}, & \ldots & e_{1,\mu}\\
0, & e_{2,2}, & \ldots & e_{2,\mu}.
\end{array}
\]
Jetzt kann man für die $e_{2,2} \ldots e_{2,\mu}$ irgend eine Combination der Zahlen 0, 1 setzen, mit Ausnahme der einen, bei der alle $e_{2,h} = 0$ werden. Zu jedem $e_{2,h} = 0$ muss das zugehörige $e_{1,h} = 1$ sein, während einem $e_{2,h} = 1$ ein $e_{1,h} = 0$ oder $= 1$ entsprechen kann. Die Anzahl der auf diese Weise entstandenen Combinationen ist:
\[
1 + (\mu - 1)2 + \frac{(\mu - 1)(\mu - 2)}{1.2} 2^2 + \cdots + 2^{\mu - 1} - 1 = 3^{\mu - 1} - 1.
\]
Nun kann man aber die zweite der Gleichungen (3), ohne ihre Bedeutung zu ändern, zu der ersten noch addiren, so dass also je zwei der gewonnenen Combinationen dieselbe Gruppe ergeben. Es bleiben also
\weq{4}{\frac{3^{\mu - 1} - 1}{2}}
verschiedene Combinationen der $e_{h,k}$, die auch, wie leicht einzusehen ist, zu verschiedenen Gruppen $\mathfrak{A}$ führen. Denn angenommen, man habe das eine Mal $e_{1,2}, e_{2,2} = 1, 0$, das andere Mal $= 0, 1$ genommen, so genügt der ersten Annahme $\alpha_1 = 0, \alpha_2 = 1, \alpha_3 = 0, \ldots \alpha_\mu = 0$, was der zweiten nicht genügt. % [sic: Weber prints e_{h,k}; the defining matrix and congruences consistently use e_{k,h}.]

Die Zahl, die wir eben bestimmt haben, ist nur dann erschöpfend, wenn $m$ nicht durch 8 theilbar ist, weil dabei angenommen ist, dass dabei niemals $e_{1,h}, e_{2,h} = 0, 0$ sei. Wenn aber $m$ durch 8 theilbar ist, dann kann bei einem Paare diese Werthcombination vorkommen, und wir müssen also noch die Fälle hinzufügen, die durch
\[
\begin{array}{ccccc}
1 & 0 & e_{1,3} & \ldots & e_{1,\mu}\\
0 & 0 & e_{2,3} & \ldots & e_{2,\mu}
\end{array}
\]
%CONT>>


% ===== scan leaf 125, printed 104 =====
% !PRINTED 104
\weberpage{104}
%<<CONT
angedeutet sind, und deren Zahl man ebenso wie oben
\[
= \frac{3^{\mu - 2} - 1}{2}
\]
findet. Wenn also $m$ durch 8 theilbar ist, so ist die Gesammtzahl der Gruppen $\mathfrak{A}$
\weq{5}{\frac{3^{\mu - 1} - 1}{2} + \frac{3^{\mu - 2} - 1}{2} = 2 . 3^{\mu - 2} - 1.}
Die kleinsten Werthe, die $m$ in diesem Falle annehmen kann, sind
\[
m = 8, 12, 15, 20, 21, 24, 28, 33, 35, 39, 40, 44 \ldots
\]
Die kleinste Zahl, die zu mehr als einer solchen Gruppe Anlass giebt, ist $m = 24$, für die man $\mu = 3$ hat, und die also nach (5) fünf verschiedene Gruppen giebt. Man erhält nämlich folgende zulässige und von einander verschiedene Combinationen der $e_{h,k}$: % [sic: Weber again prints e_{h,k}; the displayed combinations and defining congruences use e_{k,h}.]
\[
\begin{array}{ccccc}
1\ 0\ 1 & 1\ 1\ 0 & 1\ 1\ 0 & 1\ 0\ 0 & 1\ 0\ 0\\
0\ 1\ 0, & 0\ 0\ 1, & 0\ 1\ 1, & 0\ 1\ 1, & 0\ 0\ 1,
\end{array}
\]
was zu folgenden Bestimmungen über die Indices führt:
\[
\begin{aligned}
& \alpha_1 \equiv \alpha_3,\quad \alpha_2 \equiv 0,\\
& \alpha_1 \equiv \alpha_2,\quad \alpha_3 \equiv 0,\\
& \alpha_1 \equiv \alpha_2 \equiv \alpha_3,\qquad (\mathrm{mod}\ 2)\\
& \alpha_1 \equiv 0,\quad \alpha_2 \equiv \alpha_3\\
& \alpha_1 \equiv 0,\quad \alpha_3 \equiv 0,
\end{aligned}
\]
und wenn man also die Zahlen $a$ durch die Congruenzen
\[
a \equiv (-1)^{\alpha_2} 5^{\alpha_1} (\mathrm{mod}\ 8), \quad a \equiv 2^{\alpha_3} (\mathrm{mod}\ 3)
\]
bestimmt, so ergeben sich folgende fünf zweigliedrige Gruppen:
\[
\begin{aligned}
\mathfrak{A}_1 &= 1,\ 5\\
\mathfrak{A}_2 &= 1,\ 19\\
\mathfrak{A}_3 &= 1,\ 11\\
\mathfrak{A}_4 &= 1,\ 23\\
\mathfrak{A}_5 &= 1,\ 7.
\end{aligned}
\]
Um die Perioden $\eta$ zu berechnen, können wir ähnlich verfahren, wie im vorigen Paragraphen. Wir bezeichnen mit $\nu_1, \nu_2, \ldots \nu_\mu$ die Indices einer Zahl $n$ und setzen
\[
\begin{aligned}
e_{1,1} \nu_1 + e_{1,2} \nu_2 + \cdots + e_{1,\mu} \nu_\mu &= \sigma_1,\\
e_{2,1} \nu_1 + e_{2,2} \nu_2 + \cdots + e_{2,\mu} \nu_\mu &= \sigma_2,
\end{aligned}
\]
%CONT>>


% ===== scan leaf 126, printed 105 =====
% !PRINTED 105
\weberpage{105}
%<<CONT
und erhalten, wenn $\varepsilon_1$, $\varepsilon_2$ den Werth 0 oder 1 haben, die Resolventen

\weq{6}{\begin{aligned}
\Omega_{\varepsilon_1,\varepsilon_2} &= \sum^n (-1)^{\varepsilon_1\sigma_1+\varepsilon_2\sigma_2}\, r^n\\
&= \eta + (-1)^{\varepsilon_1}\eta' + (-1)^{\varepsilon_2}\eta'' + (-1)^{\varepsilon_1+\varepsilon_2}\eta'''.
\end{aligned}}

Ist nun $m = q_1\, q_2 \ldots q_\mu$ nicht durch 8 theilbar, so können
wir, wenn $r_1, r_2, \ldots r_\mu$ primitive Wurzeln der Grade $q_1, q_2, \ldots q_\mu$
bedeuten,
\[ r = r_1\, r_2 \ldots r_\mu \]
setzen, ferner
\[ \begin{aligned}
n &\equiv n_1\ (\mathrm{mod}\ q_1)\\
&\equiv n_2\ (\mathrm{mod}\ q_2)\\
&\cdot\ \cdot\ \cdot\ \cdot\ \cdot\ \cdot\ \cdot\\
&\equiv n_\mu\ (\mathrm{mod}\ q_\mu),
\end{aligned} \]
wodurch sich ergiebt:
\weq{7}{\Omega_{\varepsilon_1,\varepsilon_2} = \prod_{1,\mu}^{h} \sum^{n_h} (-1)^{\nu_h(\varepsilon_1 e_{1,h}+\varepsilon_2 e_{2,h})}\, r_h^{n_h}.}

Die einzelnen Factoren dieses Productes sind die in \S.\,171
(Bd.\ I) bestimmten \gesperrt{Gauss'schen Summen}, die dort durch
Quadratwurzeln ausgedrückt sind. Nur wenn $q_1 = 4$ ist und
$\varepsilon_1\, e_{1,1} + \varepsilon_2\, e_{1,2} = 0$, ist $\Omega_{\varepsilon_1,\varepsilon_2} = 0$. % [sic: Weber prints e_{1,2}; specializing the preceding general factor to h=1 requires e_{2,1}.]

Ist aber $m$ durch 8 theilbar und
\[ m = q_2\, q_3 \ldots q_\mu, \quad q_2 = 8, \]
so erhält man, wenn man $r = r_2\, r_3 \ldots r_\mu$ setzt:
\[ \begin{aligned}
\Omega_{\varepsilon_1,\varepsilon_2} &= \sum^{n_2} (-1)^{\varepsilon_1(e_{1,1}\nu_1+e_{1,2}\nu_2)}(-1)^{\varepsilon_2(e_{2,1}\nu_1+e_{2,2}\nu_2)}\, r_2^{n_2}\\
&\times \prod_{3,\mu}^{h} \sum^{n_h} (-1)^{\nu_h(\varepsilon_1 e_{1,h}+\varepsilon_2 e_{2,h})}\, r_h^{n_h}.
\end{aligned} \]

In den fünf oben aufgestellten Gruppen, die zu $m = 24$
gehören, kann man die Perioden aus dieser Formel oder auch
leicht direct berechnen, wenn man
\[ r = e^{\frac{\pi i}{12}} = e^{\frac{\pi i}{3}}\, e^{-\frac{\pi i}{4}} \]
setzt. Man erhält dann für die fünf Grössen:
\[ \mathfrak{A}_1)\quad \eta = \frac{(1+i)\sqrt{3}}{\sqrt{2}} \]
\[ \mathfrak{A}_2)\quad \eta = \frac{\sqrt{3}-i}{\sqrt{2}} \]
%CONT>>


% ===== scan leaf 127, printed 106 =====
% !PRINTED 106
\weberpage{106}
%<<CONT
\[ \mathfrak{A}_3)\quad \eta = i\,\frac{\sqrt{3}-1}{\sqrt{2}} \]
\[ \mathfrak{A}_4)\quad \eta = \frac{1+\sqrt{3}}{\sqrt{2}} \]
\[ \mathfrak{A}_5)\quad \eta = \frac{1+i\sqrt{3}}{\sqrt{2}}. \]

Es soll nun zuletzt noch der Fall einer einzigen Invariante
$i_1 = 4$ betrachtet werden. Die Bedingungen für $m$ bestehen
dann einfach darin, dass $m$ durch 4, 8, 16, aber keine höhere
Potenz von 2 theilbar sein kann, dass ungerade Primzahlen nur
einfach in $m$ aufgehen können, und dass mindestens einer der
Indexmoduln durch 4 theilbar sein muss. Die ersten Werthe
sind also
\[ m = 5,\ 13,\ 15,\ 16,\ 17,\ 20,\ 29,\ 33,\ 37,\ 39,\ 40\ \ldots \]

Wir wollen annehmen, es seien von den Indexmoduln
$c_1, c_2, \ldots c_\varrho$ durch 4 theilbar, $c_{\varrho+1}, c_{\varrho+2}, \ldots c_\mu$ durch 2, aber
nicht durch 4 theilbar. Es muss dann $\varrho$ mindestens gleich 1
sein, und folglich ist
\[ \begin{aligned}
&\delta_{1,1} = \delta_{1,2} = \cdots = \delta_{1,\varrho} = 4, \quad \delta_{1,\varrho+1} = \cdots = \delta_{1,\mu} = 2\\
&i_{1,1} = i_{1,2} = \cdots = i_{1,\varrho} = 1, \quad i_{1,\varrho+1} = \cdots = i_{1,\mu} = 2\\
&c_{1,1} = \frac{c_1}{4}, \cdots\cdots, c_{1,\varrho} = \frac{c_\varrho}{4}, \quad c_{1,\varrho+1} = \frac{c_{\varrho+1}}{2}, \cdots, c_{1,\mu} = \frac{c_\mu}{2}.
\end{aligned} \]

Nun sind die $e_{1,h}$ so zu bestimmen, dass von den Zahlen
\[ e_{1,1}, e_{1,2}, \ldots e_{1,\varrho} \]
wenigstens eine ungerade ist, und dass (wegen \S.\,21, 6.) von den
Zahlen
\[ e_{1,1}, e_{1,2}, \ldots e_{1,\varrho} \]
keine durch 4, und von den
\[ e_{1,\varrho+1}, e_{1,\varrho+2}, \ldots e_{1,\mu} \]
keine durch 2 theilbar sei. Eine Ausnahme bildet hierbei wieder
im Falle, wo $m$ durch 8 oder durch 16 theilbar ist, die Zahl $e_{1,h}$,
der keiner der Primfactoren von $m$ entspricht, die auch gerade
sein kann.

Dann erhalten wir für die Indices $\alpha$ der Zahlen $a$ die Congruenz:
\[ e_{1,1}\alpha_1 + \cdots + e_{1,\varrho}\alpha_\varrho + 2\,(e_{1,\varrho+1}\alpha_{\varrho+1} + \cdots + e_{1,\mu}\alpha_\mu) \equiv 0\ (\mathrm{mod}\ 4). \]


% ===== scan leaf 128, printed 107 =====
% !PRINTED 107
\weberpage{107}
\pdfbookmark[2]{§ 25. Cubische Abel'sche Gleichungen}{bm30}
Nehmen wir als Beispiel $m = 48$, so haben wir zwei Möglichkeiten, nämlich:
\[ \begin{aligned}
e_{1,1} &= 1, & e_{1,2} &= 0, & e_{1,3} &= 1\\
e_{1,1} &= 1, & e_{1,2} &= 1, & e_{1,3} &= 1,
\end{aligned} \]
und man findet aus ihnen die beiden Gruppen:
\[ \begin{aligned}
\mathfrak{A}_1 &= 1,\ 17,\ 31,\ 47\\
\mathfrak{A}_2 &= 1,\ \ 7,\ 41,\ 47.
\end{aligned} \]

\begin{center}
\S.\,25.
\end{center}

\begin{center}
\gesperrt{Cubische Abel'sche Gleichungen.}
\end{center}

Wir wollen diese Betrachtungen über die cubischen und
biquadratischen Kreistheilungsgleichungen mit dem Beweise eines
merkwürdigen Satzes abschliessen, der geeignet ist, diesen Gleichungen
 ein weit höheres und allgemeineres Interesse zu verleihen
 und der ein specieller Fall eines ganz allgemeinen Satzes
ist, den wir erst später kennen lernen werden.

Der Satz lautet so:

\gesperrt{Alle cubischen und biquadratischen Abel'schen
Gleichungen im Körper der rationalen Zahlen sind
Kreistheilungsgleichungen.}

Dadurch gewinnen die Resultate der beiden letzten Paragraphen
 einen höheren Werth. Es folgt dann nämlich, dass
durch die dort gebildeten Kreistheilungsperioden $\eta$ \gesperrt{alle Abel'schen
Zahlkörper} dritten und vierten Grades dargestellt sind,
dass andere als die dort besprochenen Körper dieser Art nicht
existiren.

Die Möglichkeit, diesen Satz für die speciellen Fälle der
cubischen und biquadratischen Gleichungen einfach und ohne
weitere Vorbereitung zu beweisen, beruht darauf, dass in den
Körpern der dritten und vierten Einheitswurzeln $R(\varrho)$, $R(i)$ dieselben
 Gesetze der Zerlegung der Zahlen in ihre Primfactoren
gelten, wie in den Körpern der rationalen Zahlen, wie in den
\S\S.\,173, 174 des ersten Bandes nachgewiesen ist. Es ist darum
auch von Interesse, diese besonderen Fälle genauer zu betrachten,
weil dadurch der Gang und das Ziel des später beizubringenden
 allgemeinen Beweises deutlicher erkannt wird.


% ===== scan leaf 129, printed 108 =====
% !PRINTED 108
\weberpage{108}
Beginnen wir also mit den cubischen Gleichungen, und verstehen
 unter $x_0, x_1, x_2$ die Wurzeln irgend einer \gesperrt{Abel'schen}
Gleichung dritten Grades. Da 3 eine Primzahl ist, so muss die
Gleichung cyklisch sein, d.\,h. \gesperrt{die cyklischen Functionen
der Grössen} $x_0, x_1, x_2$ \gesperrt{sind rationale Zahlen.}

Bezeichnen wir mit $\varrho$ eine imaginäre dritte Einheitswurzel,
so haben wir die beiden Resolventen
\weq{1}{\begin{aligned}
(\varrho, x_0) &= x_0 + \varrho\, x_1 + \varrho^2 x_2\\
(\varrho^2, x_0) &= x_0 + \varrho^2 x_1 + \varrho\, x_2,
\end{aligned}}
deren Cuben also. Zahlen des Körpers $R(\varrho)$ und deren Product % [sic: Weber prints a full stop after ``also'', although the sentence continues.]
eine rationale Zahl sein muss.

Wir setzen demnach, indem wir mit $a, b, c$ rationale (ganze
oder gebrochene) Zahlen bezeichnen,
\weq{2}{\begin{aligned}
(\varrho, x_0)^3 &= a + b\varrho\\
(\varrho^2, x_0)^3 &= a + b\varrho^2\\
(\varrho, x_0)(\varrho^2, x_0) &= c,
\end{aligned}}
also auch
\weq{3}{(a+b\varrho)(a+b\varrho^2) = c^3.}

Nun haben wir im Körper $R(\varrho)$ ausser den sechs Einheiten
\[ \pm 1,\ \pm\varrho,\ \pm\varrho^2, \]
die alle als Potenzen von $-\varrho$ dargestellt werden können, die
Primzahlen $\sqrt{-3} = \varrho - \varrho^2$, die reellen Primzahlen $q$ von der
Form $3N+2$ und die beiden complexen Factoren der reellen
Primzahlen $p$ von der Form $3N+1$. Die Zerlegung dieser
Primzahlen in die beiden complexen Factoren $\pi, \pi'$:
\weq{4}{p = \pi\pi'}
haben wir im \S.\,172 des ersten Bandes dargestellt in der Form
\weq{5}{p = \psi_1(\varrho)\,\psi_1(\varrho^2).}
Wir können also
\weq{6}{\pi = \psi_1(\varrho), \quad \pi' = \psi_1(\varrho^2),}
setzen und haben dadurch unter den verschiedenen zu $\pi$ associirten
 Zahlen eine bestimmte ausgewählt, und diese Zahlen
$\pi, \pi'$ stehen in einer bestimmten Beziehung zu den Kreis\-theilungsperioden
 der $p^{\text{ten}}$ Einheitswurzeln $\eta$:
\weq{7}{(\varrho, \eta)^3 = p\pi, \quad (\varrho^2, \eta)^3 = p\pi'}
[Bd.\ I, \S.\,172, (5)].


% ===== scan leaf 130, printed 109 =====
% !PRINTED 109
\weberpage{109}

Nun ist $a + b\varrho$ eine ganze oder gebrochene Zahl des Körpers $R(\varrho)$.
Wenn wir Zähler und Nenner dieser Zahl in ihre Primfactoren zerlegen, gemeinsame
Factoren wegheben, und wenn wir unter $q_1, q_2, \ldots$ reelle Primzahlen der
Form $3N+2$, unter $\pi_1, \pi_1', \pi_2, \pi_2', \ldots$ conjugirte Paare
imaginärer Primzahlen der Form (6) verstehen, so erhalten wir

\weq{8}{\begin{aligned}
a + b\varrho &= (-\varrho)^\lambda (\sqrt{-3})^n q_1^{s_1} q_2^{s_2} \ldots \pi_1^{t_1} \pi_1'^{t_1'} \pi_2^{t_2} \pi_2'^{t_2'} \ldots \\
a + b\varrho^2 &= (-\varrho^2)^\lambda (-\sqrt{-3})^n q_1^{s_1} q_2^{s_2} \ldots \pi_1^{t_1'} \pi_1'^{t_1} \pi_2^{t_2'} \pi_2'^{t_2} \ldots
\end{aligned}}

worin

\[
\lambda, \ n, \ s_1, \ s_2 \ldots t_1, \ t_1', \ t_2, \ t_2' \ldots
\]

ganze positive oder negative Zahlen sind.

Nun ist aber nach (3) das Product der beiden Zahlen (8) der Cubus einer
rationalen Zahl, und dies giebt, da eine rationale Zahl nur auf eine Art in
Primfactoren zerlegbar ist, die folgenden Bedingungen:

\weq{8}{\begin{aligned}
n &\equiv 0, \quad s_1 \equiv 0, \quad s_2 \equiv 0 \ldots \\
t_1+t_1' &\equiv 0, \quad t_2+t_2' \equiv 0 \ldots \quad \text{(mod 3),}
\end{aligned}} % [sic: Weber repeats equation number (8); (9) follows.]

und wir setzen also

\weq{9}{\begin{aligned}
n &= 3\mu, \quad s_1 = 3\sigma_1, \quad s_2 = 3\sigma_2 \ldots \\
t_1+t_1' &= 3\tau_1, \quad t_2+t_2' = 3\tau_2 \ldots,
\end{aligned}}

so dass $n, \sigma_1, \sigma_2 \ldots \tau_1, \tau_2 \ldots$ ganze Zahlen sind. % [sic: Weber prints n, although n=3\mu has just been set and \mu is evidently intended.]
Nun ist nach den Formeln (7):

\weq{10}{\begin{aligned}
\pi^t \pi'^{t'} &= p^{-(t+t')} (\varrho,\eta)^{3t} (\varrho^2,\eta)^{3t'} \\
\pi^{t'} \pi'^{t} &= p^{-(t+t')} (\varrho,\eta)^{3t'} (\varrho^2,\eta)^{3t}.
\end{aligned}}

Bezeichnen wir endlich noch mit $\varepsilon$ eine $9^{\text{te}}$
Einheitswurzel, so können wir

\[
-\varrho = (-\varepsilon)^3
\]

setzen, und stellen dadurch $a+b\varrho$ und $a+b\varrho^2$ als Cuben dar. Wir
erhalten also nach (2):

\weq{11}{(\varrho, x_0)^3 = (-\varepsilon)^{3\lambda} (\sqrt{-3})^{3\mu} q_1^{3\sigma_1} q_2^{3\sigma_2} \ldots p_1^{-3\tau_1} p_2^{-3\tau_2} \ldots H_1^3,}

worin

\weq{12}{H_1 = (\varrho,\eta_1)^{t_1} (\varrho^2,\eta_1)^{t_1'} (\varrho,\eta_2)^{t_2} (\varrho^2,\eta_2)^{t_2'} \ldots}

eine Kreistheilungszahl ist. Bezeichnen wir mit $H_2$ die aus $H_1$ durch
Vertauschung von $\varrho$ mit $\varrho^2$ hervorgehende Zahl, so ist $H_1 H_2$
eine rationale Zahl.

Wenn wir aus (11) die dritte Wurzel ziehen, so folgt, da hierbei noch eine
dritte Einheitswurzel $\varrho_1$ als Factor auftreten kann,

\weq{13}{(\varrho, x_0) = (-\varepsilon)^\lambda \varrho_1 (\sqrt{-3})^\mu q_1^{\sigma_1} q_2^{\sigma_2} \ldots p_1^{-\tau_1} p_2^{-\tau_2} \ldots H_1,}
%CONT>>


% ===== scan leaf 131, printed 110 =====
% !PRINTED 110
\weberpage{110}
\pdfbookmark[2]{§ 26. Biquadratische Abel'sche Gleichungen}{bm31}
%<<CONT
und ebenso erhalten wir

\weq{14}{(\varrho^2, x_0) = (-\varepsilon^2)^\lambda \varrho_2 (-\sqrt{-3})^u q_1^{\sigma_1} q_2^{\sigma_2} \ldots p_1^{-\tau_1} p_2^{-\tau_2} \ldots H_2.} % [sic: Weber prints exponent u; the conjugate of (13) requires \mu.]

Weil das Product dieser beiden Ausdrücke rational sein muss, so ergiebt sich
noch zwischen den Einheitswurzeln $\varrho, \varrho_1, \varrho_2$ die Relation

\[
\varrho^\lambda \varrho_1 \varrho_2 = 1.
\]

Setzen wir

\weq{15}{x_0 + x_1 + x_2 = A,}

so ist auch $A$ eine rationale Zahl, und aus (13), (14), (15) folgt, dass

\weq{16}{\begin{aligned}
x_0 &= 1/3 \left[A + (\varrho, x_0) + (\varrho^2, x_0)\right] \\
x_1 &= 1/3 \left[A + \varrho^2 (\varrho, x_0) + \varrho (\varrho^2, x_0)\right] \\
x_2 &= 1/3 \left[A + \varrho (\varrho, x_0) + \varrho^2 (\varrho^2, x_0)\right]
\end{aligned}}

Kreistheilungzahlen sind. % [sic] printed \textquotedbl{}Kreistheilungzahlen\textquotedbl{}, elsewhere on the page \textquotedbl{}Kreistheilungszahlen\textquotedbl{}
Aus den Formeln (13), (14) können wir die Zusammensetzung dieser Zahlen aus den
Perioden $\eta_1, \eta_2, \ldots$ ersehen. Ein näheres Eingehen auf diesen
Gegenstand ist aber nicht mehr erforderlich, weil wir schon im \S.\,23 alle
Kreistheilungszahlen, die cubischen Gleichungen genügen, vollständig gebildet
haben.

\noindent\hrulefill

\begin{center}
\S.\,26.

\gesperrt{Biquadratische Abel'sche Gleichungen.}
\end{center}

Ganz ähnlich kann der Beweis des entsprechenden Satzes für die biquadratischen
\gesperrt{Abel}'schen Gleichungen geführt werden. Nur sind hier zwei Fälle zu
unterscheiden, nämlich Gleichungen mit nicht cyklischer Gruppe, und
Gleichungen mit cyklischer Gruppe.

Für die nicht cyklischen \gesperrt{Abel}'schen Gleichungen mit den Wurzeln $x_0, x_1, x_2,
x_3$ ist die Gruppe:

\weq{1}{1, \ \ (0,1)(2,3), \ \ (0,2)(1,3), \ \ (0,3)(1,2),}

und es sind also die drei Quadrate

\weq{2}{\begin{aligned}
(x_0+x_1-x_2-x_3)^2 &= a \\
(x_0-x_1+x_2-x_3)^2 &= b \\
(x_0-x_1-x_2+x_3)^2 &= c
\end{aligned}}

und das Product

\weq{3}{(x_0+x_1-x_2-x_3)(x_0-x_1+x_2-x_3)(x_0-x_1-x_2+x_3) = c,} % [sic: Weber prints RHS c, reusing c from the third square in (2); this rational product requires a distinct symbol.]
%CONT>>


% ===== scan leaf 132, printed 111 =====
% !PRINTED 111
\weberpage{111}
%<<CONT
und ferner die Summe

\weq{4}{x_0 + x_1 + x_2 + x_3 = 4A}

rationale Zahlen.

Wenn wir aus (2) die Quadratwurzeln ausziehen und berücksichtigen, dass sich
$\sqrt{a}$ nach (3) von $\sqrt{bc}$ nur durch einen rationalen Factor
unterscheidet, so folgt durch Addition:

\weq{5}{x_0 = A + B\sqrt{b} + C\sqrt{c} + D\sqrt{bc},}

worin $A, B, C, D$ rationale Zahlen sind; und daraus erhält man $x_1, x_2,
x_3$, wenn man die Vorzeichen von $\sqrt{b}, \sqrt{c}$ ändert.

Nach Band I, \S.\,171 können aber alle Quadratwurzeln, nöthigenfalls unter
Zuziehung von $i$, was ja selbst eine Kreistheilungszahl ist, rational durch
die Kreistheilungsperioden (die \gesperrt{Gauss}'schen Summen) ausgedrückt
werden, so dass also in (5) schon der Beweis unseres Satzes liegt.

\noindent\hrulefill

Ist die Gruppe der biquadratischen Gleichung cyklisch, so können wir die
Wurzeln so anordnen, dass die Gruppe aus den Permutationen

\weq{6}{1, \ \ (0,1,2,3), \ \ (0,2)(1,3), \ \ (0,3,2,1)}

besteht, und dann ist zu setzen:

\weq{7}{\begin{aligned}
4A &= x_0+x_1+x_2+x_3 \\
(i,x_0) &= x_0+ix_1-x_2-ix_3 \\
(-1,x_0) &= x_0-x_1+x_2-x_3 \\
(-i,x_0) &= x_0-ix_1-x_2+ix_3;
\end{aligned}}

darin ist $A$ und $(-1,x_0)^2=m$ rational, und daher kann $(-1,x_0)$ durch
Kreistheilungszahlen ausgedrückt werden. Die vierten Potenzen

\weq{8}{(i,x_0)^4 = a+bi, \quad (-i,x_0)^4 = a-bi}

sind Zahlen des Körpers $R(i)$, und um dieselben Schlüsse wie bei den
cubischen Gleichungen ziehen zu können, kommt es also nur noch darauf an,
$a+bi$ und $a-bi$ als vierte Potenzen von Kreistheilungszahlen darzustellen.
Dazu dient noch der Satz, dass

\weq{9}{(i,x_0)(-i,x_0) = c}

eine rationale Zahl ist.


% ===== scan leaf 133, printed 112 =====
% !PRINTED 112
\weberpage{112}
Im Körper $R(i)$ haben wir (Bd. I, \S.\,173) die Einheiten $\pm 1, \pm i$, ferner
als Primzahlen die associirten Factoren $1\pm i$ von 2, die reellen Primzahlen
$q$ der Form $4N+3$ und die beiden complexen Factoren $\pi, \pi'$ der reellen
Primzahlen $p$ von der Form $4N+1$.

In Band I, \S.\,172 haben wir die Zerlegung

\[
p = \psi_1(i)\,\psi_1(-i)
\]

gefunden, die uns erlaubt,

\weq{10}{\pi=\psi_1(i), \quad \pi'=\psi_1(-i)}

zu setzen, und, wenn $\eta$ eine $N$-gliedrige Periode von $p^{\text{ten}}$
Einheitswurzeln ist,

\[
(i,\eta)^4 = p\,\psi_1(i)^2, \quad (-i,\eta)^4 = p\,\psi_1(-i)^2,
\]

oder

\weq{11}{(i,\eta)^4=\pi^3\pi', \quad (-i,\eta)^4=\pi\pi'^3, \quad (i,\eta)(-i,\eta)=p.}

Zerlegen wir nun, wie bei den cubischen Gleichungen, die Zahlen $a+bi$,
$a-bi$ in ihre Primfactoren in $R(i)$, so ergiebt sich unter Anwendung einer
entsprechenden Bezeichnung

\weq{12}{\begin{aligned}
a+bi &= i^\lambda (1+i)^n q_1^{s_1} q_2^{s_2} \ldots \pi_1^{t_1} \pi_1'^{t_1'} \pi_2^{t_2} \pi_2'^{t_2'} \ldots \\
a-bi &= i^{-\lambda}(1-i)^n q_1^{s_1} q_2^{s_2} \ldots \pi_1^{t_1'} \pi_1'^{t_1} \pi_2^{t_2'} \pi_2'^{t_2} \ldots
\end{aligned}}

Das Product dieser beiden Zahlen muss aber die vierte Potenz einer rationalen
Zahl sein [nach (8) und (9)], und daraus folgt:

\weq{13}{\begin{aligned}
n &\equiv 0 & &\text{(mod 4)} \\
s_1 &\equiv 0, \quad s_2 \equiv 0 \ldots & &\text{(mod 2)} \\
t_1+t_1' &\equiv 0, \quad t_2+t_2' \equiv 0 \ldots & &\text{(mod 4)} \\
t_1-t_1' &\equiv 0, \quad t_2-t_2' \equiv 0 \ldots & &\text{(mod 2).}
\end{aligned}}

Es folgt aber jetzt aus (11):

\weq{14}{\pi^t\pi'^{t'} = (\pi^3\pi')^{\frac{t-t'}{2}}(\pi\pi')^{\frac{-t+3t'}{2}} = (i,\eta)^{2(t-t')}\,p^{\frac{-t+3t'}{2}},}

und ferner

\weq{15}{(1+i)^n = (-1)^{\frac{n}{4}}\,2^{\frac{n}{2}},}

und hiernach können wir also, mit Rücksicht auf (13), wenn wir mit $H_1, H_2$
zwei Kreistheilungszahlen, nämlich das Product aller in der Zerlegung von
$a+bi$ vorkommenden Zahlen

\[
(i,\eta)^{\frac{t-t'}{2}},
\]
%CONT>>


% ===== scan leaf 134, printed 113 =====
% !PRINTED 113
\weberpage{113}
%<<CONT
und die durch die Vertauschung von $i$ mit $-i$ daraus hervorgehende
 Zahl, ferner mit $c$ eine rationale Zahl bezeichnen, setzen
\[
\begin{aligned}
a + b i &= i^{\lambda}\, c^2 H_1^4\\
a - b i &= i^{-\lambda} c^2 H_2^4.
\end{aligned}
\]
Aus (8) folgt durch Ausziehen der $4^{\text{ten}}$ Wurzel, wodurch
eine Potenz von $i$ als Factor hinzukommen kann,
\weq{15}{\begin{aligned}
(i, x_0) &= i^h \sqrt[4]{i}^{\lambda} \sqrt{c}\, H_1\\
(-i, x_0) &= i^{-h} \sqrt[4]{-i}^{\lambda} \sqrt{c}\, H_2\\
(-1, x_0) &= \sqrt{m},
\end{aligned}}
und daraus ergiebt sich, da $\sqrt[4]{i}$, $\sqrt{c}$ und $\sqrt{m}$ Kreistheilungszahlen
 sind, die Richtigkeit unseres Satzes auch in diesem Falle.

Auch hier kann die weitere Betrachtung des Baues der in
(15) vorkommenden Ausdr\"ucke dazu dienen, die Zusammensetzung
 der $x_0, x_1, x_2, x_3$ durch die Perioden $\eta$ n\"aher zu erforschen,
was aber wieder zu keinen anderen Resultaten f\"uhren
kann, als zu den schon in \S.\,24 abgeleiteten.

Wir haben also hiermit den Satz allgemein bewiesen, \gesperrt{dass
alle Abel'schen K\"orper dritten und vierten Grades
Kreistheilungsk\"orper sind, und dass alle diese K\"orper
rational durch die Kreistheilungsperioden dargestellt
werden k\"onnen.}

\rule{3cm}{0.4pt}


% ===== scan leaf 135, printed 114 =====
\weberpage{114}% page number not printed on this page
\pdfbookmark[1]{Fünfter Abschnitt. Constitution der allgemeinen Gruppen.}{bm32}
\pdfbookmark[2]{§ 27. Bildung von Gruppen nach Cayley}{bm33}
% !PRINTED none
\weberpage{none}
\begin{center}
\gesperrt{Fünfter Abschnitt.}

\gesperrt{Constitution der allgemeinen Gruppen.}
\end{center}

\rule{4cm}{0.4pt}

\begin{center}
\S.\,27.

\gesperrt{Bildung von Gruppen nach Cayley.}
\end{center}

Die allgemeine Definition der Gruppe, die im \S.\,1 gegeben ist, lässt über die Natur dieses Begriffes noch manches im Dunkel, und auch die verschiedenen einzelnen Gruppen, die wir im Verlauf unserer Betrachtungen kennen gelernt haben, geben nur Hinweise auf allgemeine Gesetze, und zeigen, dass der Gruppenbegriff an sich nichts Widersprechendes hat. In der Definition der Gruppe ist mehr enthalten, als es auf den ersten Blick den Anschein hat, und die Zahl der möglichen Gruppen, die aus einer gegebenen Anzahl von Elementen zusammengesetzt werden können, ist eine sehr beschränkte. Die allgemeinen Gesetze, die hier herrschen, sind erst zum kleinsten Theile erkannt, so dass jede neue specielle Gruppe, namentlich bei kleinerer Gliederzahl, ein neues Interesse bietet und zu eingehendem Studium auffordert.

\gesperrt{Welche Gruppen sind zwischen einer gegebenen Zahl von Elementen, d.\,h. bei gegebenem Grade überhaupt möglich?} Das ist die allgemeine Frage, um die es sich handelt, von deren vollständiger Lösung wir aber noch weit entfernt sind. \gesperrt{Cayley} hat diese Aufgabe zuerst für die niedrigsten Gradzahlen in Angriff genommen\footnote{On the theory of groups, as depending on the symbolic equation $\Theta^n = 1$. Philosophical Magazine, Vol.\ II, 1854. (\gesperrt{Cayley's} mathematical papers, Vol.\ II, 125.) American Journal of mathematics, Vol.\ I.}.


% ===== scan leaf 136, printed 115 =====
% !PRINTED 115
\weberpage{115}
Für jede beliebige Gradzahl $n$ haben wir immer eine, nämlich die cyklische Gruppe, die wir erhalten, wenn wir $a^n = 1$ und die $n$ Elemente
\[
1, a, a^2 \ldots a^{n-1}
\]
als verschieden annehmen.

Wenn $n$ eine Primzahl ist, so ist diese die einzige Gruppe vom Grade $n$ (wenn isomorphe Gruppen als identisch betrachtet werden). Denn der Grad eines jeden Elementes einer Gruppe ist ein Theiler des Grades der Gruppe, und also hat in einer Gruppe von Primzahlgrad jedes Element, mit Ausnahme des Einheitselementes, den Grad $n$.

Um eine Gruppe vom Grade $n$ vollständig darzustellen, müsste man eine quadratische Tafel construiren mit $n^2$ Feldern; die in $n$ Zeilen und $n$ Colonnen angeordnet sind. Man bezeichnet jede Zeile und jede Colonne durch eines der gegebenen Elemente, und setzt in das Feld, in dem diese beiden sich schneiden, das zusammengesetzte Element, wobei etwa der Zeilenzeiger die erste, der Colonnenzeiger die zweite Componente bedeutet:
\[
\begin{array}{c!{\vrule width 0.8pt}c|c|c|c}
 & 1 & \alpha & \beta & \gamma \ldots \\
\noalign{\hrule height 0.8pt}
1 & 1 & \alpha & \beta & \gamma \ldots \\ \hline
\alpha & \alpha & \alpha^2 & \alpha\beta & \alpha\gamma \ldots \\ \hline
\beta & \beta & \beta\alpha & \beta^2 & \beta\gamma \ldots \\ \hline
\gamma & \gamma & \gamma\alpha & \gamma\beta & \gamma^2 \ldots \\
\multicolumn{5}{c}{\cdots\cdots\cdots\cdots\cdots\cdots\cdots\cdots\cdots\cdots}
\end{array}
\]

Man kann aber die Felder einer solchen Tafel nicht ganz beliebig mit den Elementen ausfüllen, sondern man muss sich dabei an das associative Gesetz halten, so dass man, wenn man $\alpha\beta\gamma$ aufsucht, indem man zuerst in der Zeile $\alpha$ und in der Colonne $\beta$ das Element $(\alpha\beta)$, dann in der Zeile $(\alpha\beta)$ und der Colonne $\gamma$ das Element $(\alpha\beta)\gamma$ aufsucht, dasselbe Resultat findet, wie wenn man zuerst $(\beta\gamma)$ in der Zeile $\beta$ und der Colonne $\gamma$ und dann $\alpha(\beta\gamma)$ in der Zeile $\alpha$ und in der Colonne $(\beta\gamma)$ aufsucht.

Für $n = 4$ haben wir, wenn ein Element vom $4^{\mathrm{ten}}$ Grade existirt, die cyklische Gruppe
\[
1, \alpha, \alpha^2, \alpha^3,
\]
%CONT>>


% ===== scan leaf 137, printed 116 =====
% !PRINTED 116
\weberpage{116}
%<<CONT
und wenn alle Elemente vom $1^{\mathrm{sten}}$ oder $2^{\mathrm{ten}}$ Grade sind, die Gruppen
\[
1, \alpha, \beta, \alpha\beta
\]
mit der Bedingung $\alpha\beta = \beta\alpha$. Es giebt also nur diese zwei Gruppen vom $4^{\mathrm{ten}}$ Grade. Wenn die Elemente mit $1, \alpha, \beta, \gamma$ bezeichnet werden, so haben wir die beiden Tabellen:
\[
\begin{array}{c|cccc}
 & 1 & \alpha & \beta & \gamma \\
\hline
1 & 1 & \alpha & \beta & \gamma \\
\alpha & \alpha & \beta & \gamma & 1 \\
\beta & \beta & \gamma & 1 & \alpha \\
\gamma & \gamma & 1 & \alpha & \beta
\end{array}
\qquad
\begin{array}{c|cccc}
 & 1 & \alpha & \beta & \gamma \\
\hline
1 & 1 & \alpha & \beta & \gamma \\
\alpha & \alpha & 1 & \gamma & \beta \\
\beta & \beta & \gamma & 1 & \alpha \\
\gamma & \gamma & \beta & \alpha & 1
\end{array}
\]

Der Fall $n = 6$ lässt sich in folgender Weise vollständig erledigen: Wir bezeichnen die sechs Elemente mit $1, \alpha, \beta, \gamma, \delta, \varepsilon$, so dass 1 die Einheit der Gruppe ist. Wenn darunter ein Element vom Grade 6 vorkommt, so ist die Gruppe cyklisch.

Wenn wir also von diesem Falle absehen, so können die Grade der Elemente $\alpha, \beta, \gamma, \delta, \varepsilon$ nur 2 oder 3 sein. Sind $\alpha$ und $\beta$ vom $2^{\mathrm{ten}}$ Grade, so kann $\alpha\beta$ nicht vom $2^{\mathrm{ten}}$ Grade sein; denn sonst wäre
\[
\alpha = \alpha^{-1}, \quad \beta = \beta^{-1}, \quad \alpha\beta = \beta^{-1}\alpha^{-1} = \beta\alpha,
\]
und es wäre also $1, \alpha, \beta, \alpha\beta$ eine Gruppe $4^{\mathrm{ten}}$ Grades; eine solche kann aber nicht Theiler einer Gruppe $6^{\mathrm{ten}}$ Grades sein.

Es muss also mindestens ein Element $3^{\mathrm{ten}}$ Grades vorkommen, und wenn wir ein solches mit $\alpha$ bezeichnen und $\gamma$ nicht in $1, \alpha, \alpha^2$ enthalten ist, so ordnet sich die Gruppe so:
\weq{1}{S = 1, \alpha, \alpha^2, \gamma, \gamma\alpha, \gamma\alpha^2.}
Um die Bedingung zu ermitteln, dass dies eine Gruppe sei, bilden wir
\[
\gamma S = \gamma, \gamma\alpha, \gamma\alpha^2, \gamma^2, \gamma^2\alpha, \gamma^2\alpha^2,
\]
was mit $S$ identisch sein muss; und es muss also
\[
\gamma^2 = 1 \text{ oder } = \alpha \text{ oder } = \alpha^2
\]
sein. Ist aber $\gamma^2 = \alpha$ oder $= \alpha^2$, so ist $\gamma^3 = \alpha\gamma$ oder $= \alpha^2\gamma$, und $\gamma^6$, aber keine niedrigere Potenz von $\gamma = 1$, d.\,h. $S$ ist
%CONT>>


% ===== scan leaf 138, printed 117 =====
% !PRINTED 117
\weberpage{117}
\pdfbookmark[2]{§ 28. Beziehung der allgemeinen Gruppen zu den Permutationsgruppen}{bm34}
%<<CONT
cyklisch ($S = 1, \gamma, \gamma^2, \gamma^3, \gamma^4, \gamma^5$). Es bleibt also nur die Annahme, dass $\gamma$ vom $2^{\mathrm{ten}}$ Grade, also
\weq{2}{\gamma^2 = 1}
ist. Da wir aber $\gamma$ durch $\gamma\alpha$ oder $\gamma\alpha^2$ ersetzen können, so müssen auch diese vom $2^{\mathrm{ten}}$ Grade sein und es folgt:
\weq{3}{\gamma\alpha = \alpha^2\gamma, \quad \gamma\alpha^2 = \alpha\gamma,}
mit deren Hülfe man jedes Compositum aus beliebig vielen Potenzen von $\alpha$ und $\gamma$ immer auf eines und nur eines der Elemente $S$ zurückführen kann. Und wenn man jetzt
\[
1, \alpha, \alpha^2, \gamma, \gamma\alpha, \gamma\alpha^2
\]
mit
\[
1, \alpha, \beta, \gamma, \delta, \varepsilon
\]
bezeichnet, so erhält man folgende Tabelle:
\[
\begin{array}{c!{\vrule width 0.8pt}c|c|c|c|c|c}
 & 1 & \alpha & \beta & \gamma & \delta & \varepsilon \\
\noalign{\hrule height 0.8pt}
1 & 1 & \alpha & \beta & \gamma & \delta & \varepsilon \\ \hline
\alpha & \alpha & \beta & 1 & \varepsilon & \gamma & \delta \\ \hline
\beta & \beta & 1 & \alpha & \delta & \varepsilon & \gamma \\ \hline
\gamma & \gamma & \delta & \varepsilon & 1 & \alpha & \beta \\ \hline
\delta & \delta & \varepsilon & \gamma & \beta & 1 & \alpha \\ \hline
\varepsilon & \varepsilon & \gamma & \delta & \alpha & \beta & 1
\end{array}
\]

Es kommen, wie es sein muss, in jeder Zeile und in jeder Colonne alle Elemente vor.

\begin{center}
\S.\,28.

\gesperrt{Beziehung der allgemeinen Gruppen zu den Permutationsgruppen.}
\end{center}

Die Gruppen, die wir bisher am meisten angewandt haben, sind die Permutationsgruppen; diese Art von Gruppen gewinnen eine erhöhte Bedeutung auch für die allgemeine Gruppentheorie durch die folgenden Betrachtungen.

Es sei
\weq{1}{P = a_0, a_1, a_2, \ldots, a_{n-1}}
eine beliebige Gruppe vom Grade $n$. Greifen wir aus $P$ irgend
%CONT>>


% ===== scan leaf 139, printed 118 =====
% !PRINTED 118
\weberpage{118}
%<<CONT
ein Element $b$ heraus, so ist der Complex $Pb$ mit $P$ völlig
identisch. Es können sich also die beiden Reihen
\[
A = a_0,\ a_1,\ a_2,\ \ldots,\ a_{n-1}
\]
\[
Ab = a_0b,\ a_1b,\ a_2b,\ \ldots,\ a_{n-1}b
\]
nur durch die Anordnung von einander unterscheiden. Der
Uebergang von $A$ zu $Ab$ ist also eine Permutation von $n$ Ziffern
$0, 1, 2 \ldots n - 1$, und jedem Elemente $b$ von $P$ entspricht eine
solche Permutation, die wir mit $\pi_b$ bezeichnen wollen. Zwei verschiedenen
Elementen $b, c$ entsprechen immer zwei verschiedene
Permutationen $\pi_b, \pi_c$, und dem Einheitselemente entspricht die
identische Permutation. Setzen wir
\[
\pi_b = (A,\ Ab), \quad \pi_c = (A,\ Ac),
\]
so können wir $\pi_c$ auch mit $(Ab, Abc)$ bezeichnen, und daraus
ergiebt sich
\weq{2}{\pi_b \pi_c = \pi_{bc},}
d.\,h. die Permutationen $\pi$ bilden eine mit $P$ isomorphe Permutationsgruppe. Diese Permutationsgruppe ist transitiv, da z.\,B.
das Element $a_0$ in jedes beliebige andere Element von $P$ übergehen kann. Also haben wir den Satz:

1.\ \gesperrt{Jede Gruppe vom Grade $n$ ist isomorph mit einer
transitiven Permutationsgruppe von $n$ Ziffern.}

Es giebt natürlich noch andere mit einer gegebenen Gruppe
isomorphe Permutationsgruppen, und es wäre von besonderem
Interesse, eine solche Gruppe mit möglichst geringer Ziffernzahl
zu bilden. Diese Aufgabe kann bis jetzt nicht allgemein gelöst
werden. Wir müssen uns hier mit wenigen Sätzen begnügen.

Wir nehmen an, es sei $R$ irgend ein Divisor von $P$ vom
Index $j$, und bezeichnen die Elemente von $R$ mit $c$, setzen ferner,
indem wir $P$ in ein System von Nebengruppen zerlegen,
\weq{3}{P = R + Rb_1 + Rb_2 + \cdots + Rb_{j-1}.}
Ist dann $a$ irgend ein Element von $P$, so werden die beiden
Systeme
\weq{4}{\begin{aligned}
B &= R,\ Rb_1,\ Rb_2,\ \ldots,\ Rb_{j-1}\\
Ba &= Ra,\ Rb_1a,\ Rb_2a,\ \ldots,\ Rb_{j-1}a,
\end{aligned}}
von der Reihenfolge abgesehen, übereinstimmen, und es entspricht also jedem Element $a$ eine bestimmte Permutation $\pi_a$ der
$j$ Ziffern $0, 1 \ldots j - 1$, die wir mit
\weq{5}{\pi_a = (B,\ Ba)}
%CONT>>


% ===== scan leaf 140, printed 119 =====
% !PRINTED 119
\weberpage{119}
%<<CONT
bezeichnen können. Auch hier gilt das Gesetz
\weq{6}{\pi_a \pi_{a'} = \pi_{aa'},}
wenn $a'$ gleichfalls irgend ein Element aus $P$ ist. Es ist noch
die Frage, ob verschiedene Elemente $a$ dieselbe Permutation $\pi$
hervorrufen können. Wenn $\pi_a = \pi_{a'}$ ist, so folgt aus (6):
\[
\pi_{a'a^{-1}} = 1,
\]
und wenn wir also mit $a_0$ alle der Bedingung
\weq{7}{\pi_{a_0} = 1}
genügenden Elemente bezeichnen, so ist $a' = a_0 a$, und es kommt
darauf an, die Gesammtheit der Elemente $a_0$ zu ermitteln. Wenn
aber
\weq{8}{Rba_0 = Rb}
sein soll, so muss
\[
ba_0 = cb
\]
oder
\[
a_0 = b^{-1}cb
\]
sein, d.\,h. $a_0$ muss der Gruppe $b^{-1}Rb$ angehören, und dies ist
auch hinreichend für das Bestehen von (8).

Daraus ergiebt sich, dass die der Bedingung (7) genügenden
Elemente $a_0$ die ganze Gruppe $R_0$ erfüllen, die der Durchschnitt
aller mit $R$ conjugirten Theiler von $P$:
\[
R,\ b_1^{-1}Rb_1,\ b_2^{-1}Rb_2,\ \ldots,\ b_{j-1}^{-1}Rb_{j-1},
\]
und, wie wir früher nachgewiesen haben, ein Normaltheiler
von $P$ ist.

Daraus ergiebt sich für den Fall, dass $R_0 = 1$ ist, der Satz:

2.\ \gesperrt{Hat eine Gruppe $P$ einen Theiler vom Index $j$,
der mit seinen conjugirten Theilern ausser
dem Einheitselemente kein Element gemein
hat, so ist $P$ (einstufig) isomorph mit einer
transitiven Permutationsgruppe von $j$ Ziffern.}

Dass die Permutationsgruppe transitiv ist, sieht man aus (4),
wo für $a$ jedes der Elemente $b_1, b_2, \ldots, b_{j-1}$ gesetzt werden kann.

Der Satz 2. findet immer statt, wenn $P$ eine \gesperrt{einfache Gruppe}
und $R$ ein echter Theiler von $P$ ist, weil dann $R_0$, als Normaltheiler von $P$, sich auf die Einheitsgruppe reduciren muss\footnotemark.

Es mögen jetzt $R$ und $Q$ irgend zwei Theiler der Gruppe $P$
bedeuten, $R$ vom Index $j$, $Q$ vom Grade $q$.
\footnotetext{\gesperrt{Hölder}, Mathematische Annalen, Bd.\ 40, S.\ 57.}


% ===== scan leaf 141, printed 120 =====
% !PRINTED 120
\weberpage{120}
Wir zerlegen $P$ wie in (4) in die $j$ Nebengruppen $B$, bezeichnen
mit $e$ alle Elemente von $Q$, und untersuchen nun die Permutation
\[
\pi_e = (B,\ Be),
\]
die durch die Elemente $e$ unter den $j$ Elementen (4) hervorgerufen
wird. Diese Permutationen bilden nach (5) und (6) gewiss eine
mit $Q$ ein- oder mehrstufig isomorphe Gruppe $\Pi$.

Ist $a_1$ irgend ein Element von $P$, so werden in dem nach
der Composition der Theile gebildeten Complex
\[
P_1 = Ra_1Q
\]
eine gewisse Anzahl der Nebengruppen $B$ enthalten sein. Diese
Anzahl finden wir, wenn wir die Bedingung aufsuchen, dass
$Ra_1 = Ra_1e$ sei, oder dass $a_1e$ in $Ra_1$ oder $e$ in $a_1^{-1}Ra_1$ enthalten
sei. Ist also
\[
Q_1 \text{ der Durchschnitt von } Q \text{ und } a_1^{-1}Ra_1
\]
vom Index $h_1$ in Bezug auf $Q$, so enthält $P_1$ genau $h_1$ und nicht
mehr von den Complexen $B$.

Ist $a_2$ ein zweites Element von $P$, so ist der Complex
\[
P_2 = Ra_2Q
\]
entweder völlig mit $P_1$ übereinstimmend, oder beide enthalten
gar kein gemeinsames Element. Denn wenn $P_1$ und $P_2$ ein
gemeinsames Element enthalten, so haben sie einen der Complexe $B$ gemein. Das ist aber nur möglich, wenn $a_2e$ in $Ra_1e'$,
also $a_2$ in $P_1$ enthalten ist, und dann ist $P_1$ und $P_2$ identisch.
Wir können also mit der Bildung der Complexe $P_1, P_2 \ldots$ fortfahren, bis die Gruppe $P$ erschöpft ist, und erhalten die Zerlegung der Gruppe $P$ nach den zwei Gruppen $Q$ und $R$:
\weq{9}{\begin{aligned}
P &= Ra_1Q + Ra_2Q + \cdots\\
&= P_1 \phantom{{}={}} + P_2 \phantom{{}={}} + \cdots
\end{aligned}}

Aus der Bemerkung, dass, wenn $a_2$ nicht in $Ra_1Q$ vorkommt,
$a_2^{-1}$ nicht in $Qa_1^{-1}R$ enthalten ist, ergiebt sich die zweite
Zerlegung:
\weq{10}{P = Qa_1^{-1}R + Qa_2^{-1}R + \cdots\footnotemark.}

Die Bestandtheile $P_1, P_2, \ldots$ von (9) enthalten je $h_1, h_2, \ldots$
Elemente der Reihe $B$, wenn $h_1, h_2, \ldots$ die Indices der Durchschnitte von $Q$ mit $a_1^{-1}Ra_1, a_2^{-1}Ra_1, \ldots$ in Bezug auf $Q$, also % [sic: Weber prints a_2^{-1}Ra_1; the conjugate pattern requires a_2^{-1}Ra_2.]
Theiler des Grades $q$ von $Q$ bedeuten.
\footnotetext{Diese Zerlegungen verdanke ich einer Mittheilung von \gesperrt{Dedekind}.}


% ===== scan leaf 142, printed 121 =====
% !PRINTED 121
\weberpage{121}
\pdfbookmark[2]{§ 29. Der erste Sylow'sche Satz}{bm35}
Wenn nun $e$ ein beliebiges Element aus $Q$ ist, so ist $Qe = Q$,
und folglich auch $Ra_iQe = Ra_iQ$. Daraus folgt, dass durch
Composition mit $e$ die Gesammtheit der $h_i$ Elemente, aus denen
$P_i$ zusammengesetzt ist, nicht geändert wird, und dass sie also
durch die Permutationsgruppe $\Pi$ nur unter einander vertauscht
werden. Man sieht aber auch unmittelbar, dass diese $h_i$ Elemente $B$ durch $\Pi$ \gesperrt{transitiv} verbunden sind; denn $P_1$ besteht eben
aus allen den Complexen $B$, die aus einem unter ihnen durch
Anwendung von Elementen aus $Q$ abgeleitet werden können, und
daher kann jeder dieser Complexe $B$ aus $P_i$ durch Permutationen
aus $\Pi$ in jeden anderen übergeführt werden.

Wir wollen das gewonnene Resultat in folgender Weise als
Theorem formuliren:

3.\ \gesperrt{Ist $P$ eine Gruppe vom Grade $n$, $R$ ein Theiler
von $P$ vom Index $j$, $Q$ ein Theiler von $P$ vom
Grade $q$, so ist mit $Q$ eine Permutationsgruppe $\Pi$
von $j$ Ziffern ein- oder mehrstufig isomorph. Die
$j$ Ziffern zerfallen in Reihen von je $h_1, h_2, h_3, \ldots$
durch $\Pi$ transitiv verbundenen Ziffern, so dass}
\weq{11}{j = h_1 + h_2 + h_3 + \cdots}
\gesperrt{ist, und $h_1, h_2, h_3, \ldots$ Theiler von $q$ sind. Die
Zahlen $h_1, h_2, h_3, \ldots$ sind die Indices von Theilern
von $Q$, die sich als Durchschnitte von $Q$ mit
den conjugirten Gruppen $a^{-1}Ra$ ergeben.}

Haben die verschiedenen Gruppen $Q_1, Q_2, \ldots$ ausser dem
Einheitselemente keinen gemeinschaftlichen Theiler, was bei einfachen Gruppen $P$ immer eintritt, so ist der Isomorphismus einstufig.

\begin{center}
\S.\,29.

\gesperrt{Der erste \gesperrt{Sylow}'sche Satz.}
\end{center}

Es ist nicht leicht, auf dem von \gesperrt{Cayley} eingeschlagenen
directen Wege bei der Bildung von Gruppen über die niedrigsten
Gradzahlen hinauszugehen. Für weitergehende Untersuchungen
in dieser Richtung ist ein Satz von grossem Nutzen, der in beschränktem Umfange von \gesperrt{Cauchy} herrührt, allgemein aber
von \gesperrt{Sylow} bewiesen ist. Wir wollen diesen wichtigen Satz
%CONT>>


% ===== scan leaf 143, printed 122 =====
% !PRINTED 122
\weberpage{122}
%<<CONT
hier nach einem Verfahren von \gesperrt{Frobenius} beweisen, das sich
nur auf den allgemeinen Gruppenbegriff stützt\footnote{\gesperrt{Cauchy}, Exerc. d'analyse, Tom. III. \gesperrt{Sylow}, Mathem. Annalen, Bd. 5. \gesperrt{Frobenius}, Journ. für Mathematik, Bd. 100. \gesperrt{Netto}, Mathem. Annalen, Bd. 13; Substitutionentheorie. \S.\,48.}.

Der Satz lässt sich einfach so aussprechen:

\gesperrt{I. Ist $P$ eine Gruppe von irgend welchen Elementen vom Grade $n$ und $p^{\varkappa}$ eine in $n$ aufgehende Primzahlpotenz, so hat die Gruppe $P$ einen Theiler vom Grade $p^{\varkappa}$.}

\gesperrt{Cauchy} hat diesen Satz für $\varkappa = 1$ bewiesen. Für den Fall, dass $n$ eine Primzahl ist, ist er evident; für $n = 4$ und $n = 6$ kann man ihn aus den Zusammenstellungen in \S.\,27 leicht ablesen, so dass er also für die Fälle $n = 2, 3, 4, 5, 6, 7$ als erwiesen betrachtet werden kann. Auf Grund dieser Wahrnehmung lässt sich die vollständige Induction anwenden, und wir setzen also voraus, \gesperrt{der Satz sei bewiesen für Gruppen, deren Grad niedriger ist als} $n$.

Wir wollen die Elemente der Gruppe $P$ durch die Buchstaben $a, b, c, \ldots$ bezeichnen und unter $x$ ein Zeichen für ein veränderliches Element in $P$ verstehen, das die ganze Gruppe $P$ durchläuft.

Wir fassen zunächst alle die Elemente $a$ zusammen, die der Bedingung
\weq{1}{x^{-1}ax = a \quad \text{oder} \quad ax = xa}
genügen, so dass $a$ jedes Element bedeutet, das mit allen Elementen von $P$ \gesperrt{vertauschbar} ist. Dazu gehört jedenfalls das Einheitselement, und die Gesammtheit dieser Elemente $a$ bildet eine in $P$ enthaltene Gruppe $A$. Denn aus
\[ x^{-1}ax = a,\quad x^{-1}a'x = a' \]
folgt
\[ x^{-1}aa'x = aa'. \]

\gesperrt{$A$ ist eine Abel'sche Gruppe und ein Normaltheiler von $P$}; denn aus der Definition (1) folgt $x^{-1}Ax = A$, und für irgend zwei Elemente $a, a'$ von $A$:
\[ aa' = a'a. \]
Wir bezeichnen mit $\nu$ und $j$ Grad und Index dieser Gruppe, so dass
\weq{2}{n = \nu j}
%CONT>>


% ===== scan leaf 144, printed 123 =====
% !PRINTED 123
\weberpage{123}
%<<CONT
ist. Wenn nun $c$ ein nicht zu $A$ gehöriges Element von $P$ ist, so wird $x^{-1}cx$ nicht für alle Elemente $x$ gleich $c$ sein. Es werden sich aber gewisse Elemente $b, b', \ldots$ unter den $x$ finden, darunter gewiss das Einheitselement, die der Bedingung
\[ x^{-1}cx = c \]
genügen, und diese Elemente bilden eine Gruppe, weil aus
\weq{3}{b^{-1}cb = c,\quad b'^{-1}cb' = c}
folgt, dass auch
\[ (b'b)^{-1}c(b'b) = b^{-1}b'^{-1}cb'b = c \]
sein muss. Diese Gruppe bezeichnen wir mit $B$ und ihren Grad und Index mit $\mu, \varepsilon$, so dass
\weq{4}{n = \mu\varepsilon}
wird; $B$ ist ein echter Theiler von $P$, und daher ist $\varepsilon$ jedenfalls grösser als 1, weil sonst $c$ in $A$ enthalten wäre, gegen die Annahme.

Ist $g$ ein Element von $P$, das nicht in $B$ vorkommt, und $b$ ein beliebiges Element von $B$, so ist nach (3):
\weq{5}{g^{-1}b^{-1}cbg = g^{-1}cg,}
was von $c$ verschieden ist, weil $g$ nicht in $B$ vorkommt. Wenn man also $P$ in die Nebengruppen zerlegt:
\weq{6}{P = B + Bg_1 + Bg_2 + \cdots + Bg_{\varepsilon-1},}
so erhält man, wenn man in $x^{-1}cx$ für $x$ alle Elemente einer dieser Nebengruppen setzt, immer dasselbe Element $g^{-1}cg$, und man erhält also $\varepsilon$ und nicht mehr Elemente, und jedes gleich oft:
\weq{7}{c,\ g_1^{-1}cg_1,\ g_2^{-1}cg_2,\ \ldots,\ g_{\varepsilon-1}^{-1}cg_{\varepsilon-1};}
alle diese Elemente, deren Gesammtheit wir mit $C$ bezeichnen wollen, sind unter einander und von den Elementen $a$ verschieden. $C$ ist aber keine Gruppe, da das Einheitselement darunter nicht vorkommt.

Ist nun mit $A$ und $C$ die Gesammtheit der Elemente von $P$ noch nicht erschöpft, so nehme man ein in $A$ und $C$ nicht vorkommendes Element $c'$ und bilde nach derselben Vorschrift die Gruppe $B'$ vom Grade $\mu'$ aller Elemente, die der Bedingung
\[ b'^{-1}c'b' = c' \]
genügen, und das System $C'$:
\weq{8}{c',\ g_1'^{-1}c'g_1',\ g_2'^{-1}c'g_2',\ \ldots,\ g_{\varepsilon'-1}'^{-1}c'g_{\varepsilon'-1}',}
%CONT>>


% ===== scan leaf 145, printed 124 =====
% !PRINTED 124
\weberpage{124}
%<<CONT
und man sieht dann, dass diese Elemente alle von den Elementen $C$ verschieden sind, da, wenn $g'^{-1}c'g' = g^{-1}cg$ wäre, auch $c' = g'g^{-1}cgg'^{-1}$ wäre, und es würde also gegen die Voraussetzung $c'$ schon in $C$ vorkommen. Auf diese Weise fährt man mit der Bildung der Gruppen $B, B', B'', \ldots$ und der Systeme $C, C', C'', \ldots$ fort, bis alle Elemente von $P$ in den Systemen $A, C, C', C'', \ldots$ untergebracht sind. Man hat dann die Zahlengleichungen:
\weq{9}{n = \nu j = \mu\varepsilon = \mu'\varepsilon' = \mu''\varepsilon''\ldots}
\weq{10}{n = \nu + \varepsilon + \varepsilon' + \varepsilon'' + \cdots}

Auf Grund dieser Formeln lässt sich nun der \gesperrt{Sylow'sche} Satz durch vollständige Induction beweisen, wobei zwei Fälle zu unterscheiden sind:

1) \gesperrt{Der Grad} $\nu$ \gesperrt{von} $A$ \gesperrt{ist durch} $p$ \gesperrt{theilbar.} In diesem Falle giebt es, weil $A$ eine Abel'sche Gruppe ist, nach \S.\,9, 2. ein Element $a$ in $A$ vom Grade $p$, und die cyklische Gruppe $p^{\mathrm{ten}}$ Grades
\[ 1,\ a,\ a^2,\ \ldots,\ a^{p-1}, \]
die wir mit $Q$ bezeichnen, ist ein Normaltheiler von $P$, weil ja für ein Element $a$ von $A$ immer $x^{-1}ax = a$ sein sollte. Die complementäre Gruppe zu $Q$, $P/Q$ ist vom Grade $n:p$, und ihr Grad ist also kleiner als $n$ und durch $p^{\varkappa-1}$ theilbar. Nach unserer Voraussetzung giebt es also einen Theiler von $P/Q$ vom Grade $p^{\varkappa-1}$, und folglich giebt es nach \S.\,6, 2. einen Theiler von $P$ vom Grade $p^{\varkappa}$.

2) \gesperrt{Der Grad} $\nu$ \gesperrt{von} $A$ \gesperrt{ist nicht durch} $p$ \gesperrt{theilbar.} Da $n$ durch $p$ theilbar ist, so können in diesem Falle nach (10) nicht alle Zahlen $\varepsilon, \varepsilon', \varepsilon'', \ldots$, die alle grösser als 1 sind, durch $p$ theilbar sein. Wenn nun $\varepsilon$ nicht durch $p$ theilbar ist, so ist $\mu$ durch $p^{\varkappa}$ theilbar. Der Grad der Gruppe $B$ ist kleiner als $n$ und durch $p^{\varkappa}$ theilbar, und also giebt es nach Voraussetzung einen Theiler von $B$, der also auch Theiler von $P$ ist, vom Grade $p^{\varkappa}$.

Hiermit ist der Satz I bewiesen.

Zu diesem Satze kommen aber noch wesentliche Ergänzungen.
%CONT>>


% ===== scan leaf 146, printed 125 =====
% !PRINTED 125
\weberpage{125}
\pdfbookmark[2]{§ 30. Der zweite Sylow'sche Satz}{bm36}
%<<CONT
\begin{center}
\S.\,30.\\
\gesperrt{Der zweite Sylow'sche Satz.}
\end{center}

Es sei nun, wie oben, $P$ eine Gruppe vom Grade $n$ und $p^{\varkappa}$ die höchste Potenz der Primzahl $p$, die in $n$ aufgeht; $Q$ sei ein Theiler von $P$ vom Grade $p^{\varkappa}$, der nach dem Satze I, \S.\,29 existirt.

Der Inbegriff aller Elemente $c$ von $P$, die der Bedingung genügen:
\weq{1}{c^{-1}Qc = Q,}
unter denen gewiss alle Elemente von $Q$ enthalten sind, bildet eine Gruppe; denn aus
\[ Qc = cQ,\quad Qc' = c'Q \]
folgt, dass auch $Qcc' = cc'Q$ ist. Diese Elemente $c$ können wir \gesperrt{die mit der Gruppe $Q$ vertauschbaren Elemente von $P$} nennen. Bezeichnen wir die Gruppe der $c$ mit $R$, so ist $R$ ein Theiler von $P$, und $Q$ ein Theiler von $R$, und zwar ein Normaltheiler. In besonderen Fällen kann $R$ sowohl mit $Q$, als mit $P$ identisch sein. Ist aber $b$ irgend ein nicht in $R$ enthaltenes Element von $P$, so giebt es unter den Elementen der Gruppe $b^{-1}Qb$ gewiss wenigstens eines, das nicht in $Q$ enthalten ist.

Bezeichnen wir mit $r$ den Index von $Q$ in Bezug auf $R$, mit $j$ den Index von $R$ in Bezug auf $P$, so ist
\weq{2}{n = p^{\varkappa}rj,}
und $r$ und $j$ sind durch $p$ nicht theilbar.

Der Satz, den wir zu beweisen haben, lautet:

\gesperrt{II. Der Index $j$ von $R$ ist von der Form $pk + 1$, worin $k$ irgend eine ganze Zahl ist $[j \equiv 1\ (\mathrm{mod}\ p)]$; es giebt $j$ und nicht mehr verschiedene Theiler von $P$ vom Grade $p^{\varkappa}$, die alle mit einander conjugirt sind. Jeder Theiler von $P$, dessen Grad eine Potenz von $p$ ist, ist in einer dieser conjugirten Gruppen enthalten.}

Ist $c$ ein Element aus $R$ vom Grade $\gamma$, und $s$ der kleinste positive Exponent, für den $c^s$ in $Q$ enthalten ist, so muss jeder andere Exponent $t$, für den $c^t$ zu $Q$ gehört, durch $s$ theilbar
%CONT>>


% ===== scan leaf 147, printed 126 =====
% !PRINTED 126
\weberpage{126}

%<<CONT
sein. Denn setzen wir $t = m s + s_1$, worin $s_1 < s$ ist, so ist $c^{s_1}$
in $Q$ enthalten, und $s_1$ muss also $= 0$ sein. Demnach ist $s$ ein
Theiler von $\gamma$. Nun bilden aber die Elemente

\[
Q + c\,Q + c^2\,Q + \cdots c^{s-1}\,Q
\]

wegen (1) eine in $P$ enthaltene Gruppe, und zwar vom Grade $p^{\varkappa} s$,
und da also $p^{\varkappa} s$ ein Theiler von $n$ sein muss, so folgt, dass $s$
nicht durch $p$ theilbar sein kann. Wäre nun der Grad $\gamma$ von $c$
eine Potenz von $p$, so müsste auch $s$ eine Potenz von $p$ sein, was
hiernach nicht möglich ist.

Es kann also der Grad von $c$ gewiss nicht eine Potenz von
$p$ sein; d. h. \gesperrt{ausser den Elementen $Q$ giebt es in $R$ keine
Elemente, deren Grad eine Potenz von $p$ ist.}

Bedeutet nun $b$ ein Element von $P$, das nicht zu $R$ gehört,
so giebt es nach der Definition von $R$ in der Gruppe

\[
Q' = b^{-1}\,Q\,b
\]

mindestens ein Element, das nicht in $Q$ vorkommt, und der
Grad dieses Elementes muss eine Potenz von $p$ sein, weil der
Grad von $Q'$ gleich $p^{\varkappa}$, also eine Potenz von $p$ ist. Dies Element
kann also nicht in $R$ vorkommen, \gesperrt{und $Q'$ ist daher kein
Theiler von $R$.}

Wählen wir die Elemente $b_1, b_2, \ldots, b_{j-1}$ aus $P$ so aus,
dass wir

\weq{10}{P = R + R\,b_1 + R\,b_2 + \cdots + R\,b_{j-1}} % [sic: Weber jumps from equation (2) to (10) in this section and subsequently cites (10).]

setzen können, so sind also die $j - 1$ Gruppen:

\[
b_1^{-1}\,Q\,b_1,\ \ b_2^{-1}\,Q\,b_2,\ \ \ldots,\ \ b_{j-1}^{-1}\,Q\,b_{j-1},
\]

die alle vom Grade $p^{\varkappa}$ sind, von $Q$ verschieden. Sie sind aber
auch von einander verschieden; denn wäre z. B.

\[
b_1^{-1}\,Q\,b_1 = b_2^{-1}\,Q\,b_2,
\]

so würde folgen:

\[
b_2\,b_1^{-1}\,Q\,b_1\,b_2^{-1} = Q,
\]

d. h. $b_2\,b_1^{-1} = c$ wäre in $R$ enthalten, also $b_2 = c\,b_1$ in $R\,b_1$, was
wegen (10) nicht möglich ist. Wir haben also gewiss $j$ verschiedene
conjugirte Gruppen $p^{\varkappa\,\mathrm{ten}}$ Grades:

\weq{11}{Q,\ \ b_1^{-1}\,Q\,b_1,\ \ b_2^{-1}\,Q\,b_2,\ \ \ldots,\ \ b_{j-1}^{-1}\,Q\,b_{j-1}.}

Ist $Q' = b^{-1}\,Q\,b$ eine von ihnen, so enthält $R' = b^{-1}\,R\,b$
den Theiler $Q'$, aber ausser diesem kein Element, dessen Grad
%CONT>>


% ===== scan leaf 148, printed 127 =====
% !PRINTED 127
\weberpage{127}
\pdfbookmark[2]{§ 31. Gruppen vom Grade pα}{bm37}
%<<CONT
eine Potenz von $p$ ist. $R'$ ist der Inbegriff aller mit $Q'$ vertauschbaren Elemente von $P$.

Wir wenden nun weiter den Satz 3., \S.\,28 an, indem wir unter $P$, $R$, $Q$ die Gruppen verstehen, die hier mit denselben Buchstaben bezeichnet sind. Da hier $Q$ ein Theiler von $R$ ist, so ist $h_1 = 1$. Von den conjugirten Gruppen $b^{-1}Rb$ ist aber keine ausser $R$ durch $Q$ theilbar, und also sind die übrigen $h_2, h_3, \ldots$ alle grösser als $1$. Sie sind aber, da sie Theiler des Grades von $Q$ sein müssen, Potenzen von $p$ und aus der Gleichung

\[
j = h_1 + h_2 + h_3 + \cdots
\]

folgt also $j \equiv 1\ (\mathrm{mod}\ p)$.

Ist aber $Q_1$ irgend ein Divisor von $P$, dessen Grad eine Potenz von $p$ ist, so wenden wir wieder die Zerlegung von $j$ nach dem Theorem 3., \S.\,28 an, indem wir für die dort vorkommenden Gruppen $P$, $R$, $Q$ die Gruppen $P$, $R$, $Q_1$ setzen. Dann ist wieder

\[
j = h_1 + h_2 + h_3 + \cdots,
\]

worin jetzt $h_1, h_2, h_3, \ldots$ die Indices (in Bezug auf $Q$) der Durchschnitte von $Q_1$ mit der Gruppe $R$ und seinen conjugirten bedeuten, die also auch Potenzen von $p$ sind. % [sic: Weber prints "in Bezug auf Q"; after substituting Q_1 for theorem 3's Q, these indices are with respect to Q_1.]
Da aber $j \equiv 1\ (\mathrm{mod}\ p)$ ist, so muss mindestens eine von den Zahlen $h = 1$ sein, d. h. eine der Gruppen $b_\lambda^{-1}Rb_\lambda$ und folglich auch eine der Gruppen (11) ist durch $Q_1$ theilbar. Ist $Q_1$ vom Grade $p^\varkappa$, so ist es mit einer der Gruppen (11) identisch.

Damit ist also das Theorem II vollständig bewiesen.

\begin{center}
\S.\,31.

\gesperrt{Gruppen vom Grade $p^\alpha$.}
\end{center}

\gesperrt{Sylow} hat aus seinen Sätzen eine merkwürdige Eigenschaft der Gruppen abgeleitet, deren Grad eine Potenz einer Primzahl $p$, $p^\alpha$, ist. Es sei also $P$ eine solche Gruppe und $\beta$ eine positive Zahl, kleiner als $\alpha$. Dann enthält nach dem Satze I, \S.\,29 die Gruppe $P$ einen Theiler $Q$ vom Grade $p^\beta$, deren Index $p^{\alpha-\beta}$ ist. Wir wenden den Satz 3., \S.\,28 an, indem wir für $R$ und $Q$ diese Gruppe $Q$ vom Grade $p^\beta$ setzen, dann sind $h_1, h_2, h_3, \ldots$
%CONT>>


% ===== scan leaf 149, printed 128 =====
% !PRINTED 128
\weberpage{128}
%<<CONT
die Indices von $Q$ und den mit $Q$ conjugirten Theilern von $P$ in Bezug auf $Q$. Daher ist $h_1 = 1$ und

\[
p^{\alpha-\beta} = h_1 + h_2 + h_3 \ldots % [sic: Weber prints no plus sign before the ellipsis; the preceding decomposition requires one.]
\]

Die Summanden $h_1, h_2, h_3, \ldots$ können hier als Theiler des Grades von $Q$ nur Potenzen von $p$ sein, und da $h_1 = 1$ ist; so folgt, dass mindestens $p$ dieser Summanden $= 1$ sein müssen, da sonst ihre Summe nicht durch $p$ theilbar sein könnte. Das heisst aber nichts Anderes, als dass wenigstens $p - 1$ der conjugirten Gruppen $b^{-1}Qb$, wo $b$ nicht in $Q$ enthalten ist, durch $Q$ theilbar und also gleich $Q$ sein müssen. Wenn aber

\[
b^{-1}Qb = Q
\]

ist, so ist auch für jeden Exponenten $s$

\[
b^{-s}Qb^s = Q.
\]

Ist $b^r$ die niedrigste Potenz von $b$, die in $Q$ enthalten ist, so muss, da der Grad von $b$ eine Potenz von $p$ ist, auch $r$ eine Potenz von $p$ sein, und

\[
b^{\frac{r}{p}} = c
\]

ist ein Element von $P$, das selbst nicht in $Q$ vorkommt, dessen $p^{\text{te}}$ Potenz aber in $Q$ enthalten ist, und das zugleich der Bedingung

\[
cQc^{-1} = Q
\]

genügt. Daraus folgt, dass die Elemente

\[
R = Q + Qc + Qc^2 + \cdots + Qc^{p-1}
\]

eine Gruppe vom Grade $p^{\beta+1}$ bilden, von der $Q$ ein Normaltheiler ist. Wir haben also den Satz:

III. \gesperrt{Ist $P$ eine Gruppe vom Grade $p^\alpha$ und $Q$ ein Theiler von $P$ vom Grade $p^\beta$ ($\beta < \alpha$), so giebt es einen Theiler $R$ von $P$ vom Grade $p^{\beta+1}$, von dem $Q$ Normaltheiler ist.}

Nimmt man in diesem Satze $\beta = \alpha - 1$ an, so fällt $R$ mit $P$ zusammen, und es ergiebt sich, dass $Q$ ein Normaltheiler von $P$ ist. Wendet man denselben Satz auf die Gruppe $Q$ an u. s. f., so erhält man:

IV. \gesperrt{Jede Gruppe, deren Grad eine Primzahlpotenz ist, ist metacyklisch (\S.\,8).}


% ===== scan leaf 150, printed 129 =====
% !PRINTED 129
\weberpage{129}
\pdfbookmark[2]{§ 32. Satz von Frobenius}{bm38}
\begin{center}
\S.\,32.

\gesperrt{Satz von Frobenius.}
\end{center}

\gesperrt{Frobenius} hat einen Satz aufgestellt, der als ein Gegenstück zum vierten Sylow'schen betrachtet werden kann, der sich so aussprechen lässt\footnote{\gesperrt{Frobenius}, Sitzungsberichte der Berl. Akademie, 4. Mai 1893.}:

V. \gesperrt{Jede Gruppe, deren Grad ein Product von lauter verschiedenen Primzahlen ist, ist metacyklisch.}

Wir setzen also voraus, dass der Grad $n$ einer Gruppe $P$ \gesperrt{durch keine Quadratzahl theilbar} sei. Ist $t$ der grösste Primtheiler von $n$, so giebt es nach dem \gesperrt{Cauchy-Sylow}'schen Satze (\S.\,29) einen Theiler $T$ vom Grade $t$ von $P$, der, weil $t$ eine Primzahl ist, eine cyklische Gruppe ist, also aus den Elementen

\[
T = 1, \; a, \; a^2, \; \ldots \; a^{t-1}
\]

besteht. Wenn sich nun nachweisen liesse, dass unter den gemachten Voraussetzungen $T$ ein Normaltheiler von $P$ ist, so könnte man dieselbe Schlussweise auf die Gruppe $P/T$, deren Grad $n:t$ ist, anwenden, und würde, wenn $t_1$ die zweitgrösste in $n$ aufgehende Primzahl ist, auf einen Normaltheiler $N$ der Gruppe $P/T$ vom Grade $t_1$ schliessen. Daraus würde dann aber folgen (\S.\,6, 2.), dass $P$ einen Normaltheiler $T_1$ vom Grade $t\,t_1$ hat, von dem $T$ selbst wieder Normaltheiler ist. Indem man in gleicher Weise die Gruppe $P/T_1$ betrachtet u. s. f., ergiebt sich eine Compositionsreihe von $P$:

\[
P, \; T_k, \; T_{k-1}, \; \ldots \; T_1, \; T, \; 1,
\]

in der die Indexreihe aus den der Grösse nach aufsteigend geordneten Primfactoren von $n$ besteht, und $P$ ergiebt sich als metacyklisch. Alles kommt also darauf an, nachzuweisen, dass $T$ ein \gesperrt{Normaltheiler} von $P$ ist.

Dieser Satz ist leichter zu beweisen, wenn man ihn als speciellen Fall eines allgemeineren betrachtet, als wenn man ihn direct angreift. Dieser allgemeinere Satz lautet unter den über $P$ und $n$ gemachten Voraussetzungen:

$\alpha$) \gesperrt{Ist $n = \mu\nu$ in zwei Factoren zerlegt und ist jeder Primtheiler von $\nu$ grösser als jeder Prim-}
%CONT>>


% ===== scan leaf 151, printed 130 =====
% !PRINTED 130
\weberpage{130}
%<<CONT
\gesperrt{theiler von $\mu$, so giebt es $\nu$ und nicht mehr Elemente in $P$, deren Grad in $\nu$ aufgeht.}

Wenden wir diesen Satz auf $\nu = t$ an, so folgt, dass es ausser $T$ keine Elemente in $P$ geben kann, deren Grad $= t$ ist. Wenn aber $c$ irgend ein Element von $P$ ist, so ist die mit $T$ conjugirte Gruppe $c^{-1}Tc$ gleichfalls vom Grade $t$, und muss also mit $T$ identisch sein. Damit ist dann bewiesen, dass $T$ ein Normaltheiler von $P$ ist.

Zum Beweis des Satzes $\alpha$) wendet man die vollständige Induction an. Derselbe ist offenbar richtig, wenn $\mu = 1$, $\nu = n$ ist, und wir setzen also als bewiesen voraus:

$\alpha'$) \gesperrt{Ist $\mu'\nu'$ durch kein Quadrat theilbar, jeder Primtheiler von $\nu'$ grösser als jeder Primtheiler von $\mu'$, und die Anzahl der Primtheiler von $\mu'$ kleiner als die Anzahl der Primtheiler von $\mu$, so giebt es in jeder Gruppe vom Grade $\mu'\nu'$ genau $\nu'$ Elemente, deren Grad ein Theiler von $\nu'$ ist.}

Um daraus den Beweis des Satzes $\alpha$) abzuleiten, bezeichnen wir mit $p$ den grössten Primtheiler von $\mu$ und mit $Q$ eine in $P$ enthaltene Gruppe vom Grade $p$, deren Existenz nach dem \gesperrt{Cauchy-Sylow'schen} Theorem feststeht. Dann sind alle Primfactoren von $\nu$ grösser als $p$. Wie in \S.\,30 bezeichnen wir mit $R$ die Gruppe, die aus allen der Bedingung
\[
c^{-1}Qc = Q
\]
genügenden Elementen $c$ von $P$ besteht, deren Grad wir, wie oben, mit $pr$ bezeichnen; $j$ bedeutet, wie früher, den Index des Theilers $R$ von $P$. Wir setzen $r = r_1 r_2$, und verstehen unter $r_1$ den grössten gemeinschaftlichen Theiler von $\mu$ und $r$, so dass $r_2$ der grösste gemeinschaftliche Theiler von $\nu$ und $r$ ist. Es ist dann
\[
n = \mu\nu = r_1 p r_2 j.
\]
Nach der Hypothese $\alpha'$) enthält nun $P$ genau $p\nu$ Elemente, deren Grad ein Theiler von $p\nu$ ist; es möge mit $U$ die Gesammtheit dieser Elemente bezeichnet sein. Die Gruppe $R$ enthält aber, gleichfalls nach $\alpha'$), genau $pr_2$ Elemente, deren Grad ein Theiler von $pr_2$, also ein Theiler von $p\nu$ ist. Dieses System wollen wir mit $V$ bezeichnen. Offenbar sind alle Elemente von $V$ zugleich in $U$ enthalten.


% ===== scan leaf 152, printed 131 =====
% !PRINTED 131
\weberpage{131}
Wenn wir nun noch beweisen können, dass es in $U$ \gesperrt{genau $(p-1)\nu$ Elemente giebt, deren Grad durch $p$ theilbar ist,} so sind wir am Ziele; denn dann folgt, dass es unter den Elementen $U$ und also auch unter den Elementen von $P$ genau
\[
p\nu - (p-1)\nu = \nu
\]
giebt, deren Grad ein Theiler von $\nu$ ist.

Dies ergiebt sich aber aus Folgendem:

$\beta$) \gesperrt{Ist $v$ ein Element aus $V$, so gehören alle Elemente $vQ$ (deren Anzahl $p$ beträgt) zu $V$. Es giebt darunter $p-1$, deren Grad durch $p$ theilbar ist, und eines, dessen Grad nicht durch $p$ theilbar ist. In $V$ giebt es $(p-1)r_2$ Elemente, deren Grad durch $p$ theilbar ist, und $r_2$ Elemente, deren Grad nicht durch $p$ theilbar ist.}

Denn ist $a$ ein von $1$ verschiedenes Element von $Q$, so ist, da $v$ zu $R$ gehört, also $v^{-1}av$ in $Q$ enthalten ist, und da $Q$ aus den Potenzen von $a$ besteht,
\weq{1}{v^{-1}av = a^s.}
Durch wiederholte Anwendung ergiebt sich daraus für jeden Exponenten $h$
\weq{2}{v^{-h}av^h = a^{s^h}.}
Nehmen wir für $h$ den Grad des Elementes $v$, der nach der Voraussetzung ein Theiler von $p\nu$ ist, setzen also $v^h = 1$, so folgt
\weq{3}{a = a^{s^h}, \quad s^h \equiv 1 \ (\mathrm{mod}\ p).}
In $h$ gehen aber keine anderen Primfactoren auf als solche, die gleich oder grösser als $p$ sind, und also ist $p-1$ relativ prim zu $h$ und man kann der Congruenz
\weq{4}{xh \equiv 1 \ (\mathrm{mod}\ p-1)}
genügen; demnach ergiebt sich aus dem \gesperrt{Fermat'schen} Lehrsatze:
\weq{5}{s^{xh} \equiv s \equiv 1 \ (\mathrm{mod}\ p),}
also nach (1):
\weq{6}{av = va,}
d. h. $a$ und $v$ sind vertauschbar. Daraus folgt für jeden Exponenten $\lambda$
\weq{7}{(va)^\lambda = v^\lambda a^\lambda.}


% ===== scan leaf 153, printed 132 =====
% !PRINTED 132
\weberpage{132}
Ist, wie vorhin, $h$ der Grad von $v$, und setzen wir, wenn $h$ durch $p$ theilbar ist, $\lambda = h$, und wenn $h$ nicht durch $p$ theilbar ist, $\lambda = hp$, so folgt aus (7) $(va)^\lambda = 1$, d. h. der Grad von $va$ ist ein Theiler von $h$ oder von $hp$, also jedenfalls ein Theiler von $p\nu$; d. h. $va$ \gesperrt{ist in} $V$ \gesperrt{enthalten}.

Ersetzen wir nun in (7) $a$ durch $a^x$ und lassen $x$ die Reihe der Zahlen $0, 1, \ldots p-1$ durchlaufen, so durchläuft $a^x$ die Gruppe $Q$ und $va^x$ das System $vQ$, und es ist für jedes $\lambda$
\weq{8}{(va^x)^\lambda = v^\lambda a^{x\lambda}.}
Wenn nun der Grad $h$ von $v$ nicht durch $p$ theilbar ist, so ist
\[
(va^x)^h = a^{xh},
\]
also nur dann $= 1$, wenn $x = 0$ ist; dagegen ist
\[
(va^x)^{ph} = 1,
\]
d. h. der Grad von $va^x$ ist, wenn $x$ nicht $= 0$ ist, ein Theiler von $ph$, aber nicht von $h$, also durch $p$ theilbar; dagegen ist er $= h$, wenn $x = 0$ ist.

Ist aber der Grad $h$ von $v$ durch $p$ theilbar, so ist er von der Form $pg$, und $g$ ist durch $p$ nicht theilbar, weil $p$ nur einfach in $n$, also auch in $h$ aufgeht. Es ist also $v^g$ ein Element aus $R$ vom Grade $p$, und muss also nach \S.\,30 in $Q$ enthalten sein. Setzen wir also hiernach
\[
v^g = a^y,
\]
so wird
\[
(va^x)^g = a^{y+gx},
\]
und dies ist dann und nur dann $= 1$, wenn $y + gx \equiv 0 \ (\mathrm{mod}\ p)$ wird, was nur für \gesperrt{einen} Werth von $x$ eintritt. Für diesen Werth von $x$ ist der Grad von $va^x$ ein Theiler von $g$; für die anderen Werthe von $x$ ist er Theiler von $pg$, aber nicht von $g$, also durch $p$ theilbar. Damit sind die beiden ersten Behauptungen des Satzes $\beta$) erwiesen.

Um auch den letzten Theil einzusehen, bemerke man, dass, wenn $v_1, v_2$ zwei Elemente aus $V$ bedeuten, die Systeme $v_1 Q, v_2 Q$ entweder ganz identisch sind, oder kein einziges gemeinschaftliches Element haben. Daraus folgt, dass man $V$ in der Weise darstellen kann:
\[
V = v_1 Q + v_2 Q + \cdots + v_{r_2} Q
\]
(obwohl $V$ im Allgemeinen keine Gruppe ist), und in jedem dieser $r_2$ Theilsysteme $vQ$ giebt es $p-1$ Elemente, deren Grad
%CONT>>


% ===== scan leaf 154, printed 133 =====
% !PRINTED 133
\weberpage{133}
%<<CONT
durch $p$ theilbar ist, und ein Element, dessen Grad nicht durch $p$ theilbar ist.

Nach dem Satze II, \S.\,30 giebt es $j$ und nicht mehr von einander verschiedene conjugirte Gruppen $Q, Q', Q'', \ldots Q^{(j-1)}$ in $P$, und $j$ von einander verschiedene conjugirte Gruppen $R, R', R'', \ldots R^{(j-1)}$.

Es lässt sich hiernach der Satz beweisen:

$\gamma$) \gesperrt{Jedes Element von $P$, dessen Grad durch $p$ theilbar ist, kommt in einer und nur in einer der Gruppen $R, R', R'' \ldots R^{(j-1)}$ vor.}

Die Gruppen $Q, Q', Q'' \ldots$ sind, da ihr Grad eine Primzahl ist, cyklisch, und je zwei enthalten, weil sie von einander verschieden sind, ausser dem Einheitselemente kein gemeinsames Element. Es sei
\[
Q = 1, a, a^2, \ldots a^{p-1}.
\]
Die Gruppe $R$, die aus allen der Bedingung
\[
c^{-1}Qc = Q
\]
genügenden Elementen besteht, enthält ausser den Potenzen von $a$ kein Element vom Grade $p$; und Entsprechendes gilt für die anderen Gruppen $Q', R', Q'', R'' \ldots$

Ist nun $b$ ein Element von $P$, dessen Grad durch $p$ theilbar ist und gleich $ps$ sein mag, so dass $s$ nicht durch $p$ theilbar ist, so ist $b^s$ ein Element, dessen Grad $p$ ist, und das also in einer und nur in einer der cyklischen Gruppen $Q, Q', Q'' \ldots$ vorkommt. Ist aber $b^s = a$ ein Element von $Q$, so ist
\[
b^{-1}ab = a,
\]
d. h. $b$ kommt in der Gruppe $R$ vor, und kann in keiner der anderen Gruppen, z. B. in $R'$, vorkommen, weil sonst $b^s$ auch in $Q'$ vorkäme, was nicht möglich ist, da $Q$ und $Q'$ nur das Einheitselement gemein haben. Damit ist $\gamma$) bewiesen.

$\delta$) \gesperrt{In $U$ giebt es genau $jr_2(p-1)$ Elemente, deren Grad durch $p$ theilbar ist.}

Unter $U$ haben wir die Gesammtheit der Elemente von $P$ verstanden, deren Grad ein Theiler von $pr$ ist. Nun giebt es nach $\beta$) in $R$ genau $r_2(p-1)$ Elemente aus $U$, deren Grad durch $p$ theilbar ist, und da jede der $j$ Gruppen $R, R', R'', \ldots$ an Stelle von $R$ treten kann, so folgt $\delta$) aus $\gamma$). Um also noch zu zeigen, worauf nach dem Obigen alles ankommt, dass in $U$
%CONT>>


% ===== scan leaf 155, printed 134 =====
% !PRINTED 134
\weberpage{134}
\pdfbookmark[2]{§ 33. Gruppen vom Grade pα q}{bm39}
%<<CONT
genau $(p-1)\nu$ Elemente vorkommen, deren Grad durch $p$ theilbar ist, ist also nur noch die Formel

\[ \varepsilon)\quad \nu = j\,r_2 \]

zu beweisen. Diese ergiebt sich aus folgender Ueberlegung. Nach $\alpha')$ ist die Anzahl der Elemente $U$ gleich $p\,\nu$. Unter diesen Elementen $U$ ist aber gewiss eines, das Einheitselement, dessen Grad nicht durch $p$ theilbar ist, und folglich ist nach $\delta)$

\[ j\,r_2\,(p-1) < p\,\nu. \]

Andererseits ist $n = \mu\,\nu = p\,r_1 r_2 j$, und da $\nu$ relativ prim zu $p\,r_1$ ist, so ist $j\,r_2$ durch $\nu$ theilbar. Also ist $j\,r_2 : \nu$ eine ganze positive Zahl, die kleiner als $p : p-1$ und folglich auch kleiner als 2 ist, und die daher nur gleich 1 sein kann. Dadurch ist $\varepsilon)$ bewiesen und zugleich das ganze Theorem.

\begin{center}
\S.\,33.

\gesperrt{Gruppen vom Grade $p^{\alpha} q$.}
\end{center}

\gesperrt{Frobenius} hat in der erwähnten Abhandlung den Satz ausgesprochen:

VI. \gesperrt{Jede Gruppe vom Grade $p^{\alpha} q$ ist, wenn $p$ und $q$ von einander verschiedene Primzahlen sind und $\alpha$ einen beliebigen positiven Exponenten bedeutet, metacyklisch.}

Der folgende Beweis dieses Satzes entstammt einer brieflichen Mittheilung von \gesperrt{Frobenius} an den Verfasser.

Zunächst ist klar, dass es genügt, zu beweisen, dass keine Gruppe $P$ vom Grade $p^{\alpha} q$ einfach ist. Unser Satz ist nämlich bewiesen für $\alpha = 1$. Nehmen wir ihn also für jedes kleinere $\alpha$ schon als bewiesen an, und setzen voraus, dass $P$ einen echten Normaltheiler $Q$ habe, der mehr als das Einheitselement umfasst, so sind die beiden Gruppen $Q$, $P/Q$ nach der Voraussetzung oder nach dem Satze IV metacyklisch, und also ist auch $P$ metacyklisch.

Wir nehmen also jetzt an, es sei die Gruppe $P$ vom Grade $p^{\alpha} q$ einfach, und wir haben aus dieser Annahme einen Widerspruch abzuleiten, um ihre Unmöglichkeit nachzuweisen.

$\alpha)$ Zunächst haben wir nach dem Satze I., \S.\,29 einen Theiler von $P$ vom Grade $q$. Wir wollen annehmen, es gebe einen
%CONT>>


% ===== scan leaf 156, printed 135 =====
% !PRINTED 135
\weberpage{135}
%<<CONT
Theiler von $P$ vom Grade $p^{\beta} q$, wobei $0 \leqq \beta < \alpha$ vorausgesetzt ist, und es sei $\beta$ so gross als möglich angenommen. Dann ist wegen der vorausgesetzten Einfachheit der Gruppe $P$ der Satz \S.\,28, 2. anwendbar, wonach $P$ isomorph ist mit einer Permutationsgruppe von $p^{\alpha-\beta}$ Ziffern. Unter den Permutationen dieser Gruppe giebt es auch solche, deren Grad durch $q$ theilbar ist, und die also, in ihre Cyklen zerlegt, einen Cyklus von $q$ Ziffern enthalten müssen, und daraus folgt:

\weq{1}{p^{\alpha-\beta} > q.}

$\beta)$ Es sei nun zweitens $Q$ ein Theiler von $P$ vom Grade $p^{\alpha}$ und Index $q$. Da nach der Voraussetzung $P$ einfach ist, so kann die in \S.\,30 mit $R$ bezeichnete Gruppe, die aus den mit $Q$ vertauschbaren Elementen von $P$ besteht, nicht mit $P$ identisch sein, da sonst $Q$ ein Normaltheiler von $P$ wäre, was doch, da $P$ als einfache Gruppe vorausgesetzt ist, unmöglich ist. Da aber $Q$ ein Theiler von $R$ ist, und der Index von $Q$ in Bezug auf $P$ eine Primzahl, so muss $R = Q$ sein. Nach II, \S.\,30 giebt es also $q$ und nicht mehr von einander verschiedene conjugirte Gruppen

\[ Q,\ Q_1,\ Q_2,\ \dots\ Q_{q-1}, \]

die alle vom Grade $p^{\alpha}$ sind. Unter diesen Gruppen nehmen wir zwei, etwa $Q$, $Q_1$, die einen grössten gemeinschaftlichen Theiler $Q_0$ vom Grade $p^{\gamma}$ haben, und wir nehmen diese beiden Gruppen so gewählt an, dass $\gamma$ \gesperrt{so gross als möglich} wird. Es ist dann $0 \leqq \gamma < \alpha$. Wenden wir nun auf die Gruppe $P$ das Theorem \S.\,28, 3. an, indem wir für die beiden Gruppen $Q$, $R$ jenes Theorems die Gruppe $Q$ setzen, so ist

\[ q = h_1 + h_2 + h_3 + \cdots, \]

und darin ist $h_1 = 1$, und die übrigen Summanden $h_2, h_3, \dots$ sind die Indices von grössten gemeinschaftlichen Theilern von $Q$ mit je einer der Gruppen $Q_1, Q_2, \dots$. Also sind $h_1, h_2, \dots$ durch $p^{\alpha-\gamma}$ theilbar, % [sic: Weber includes h_1 among the divisible terms although h_1 = 1.]
und es folgt

\[ q \equiv 1 \pmod{p^{\alpha-\gamma}}, \]

also

\weq{2}{q > p^{\alpha-\gamma},}

und folglich wegen (1):

\weq{3}{\gamma > \beta,\quad \gamma > 0.}

$\gamma)$ Ist nun also $Q_0$ der Durchschnitt von $Q$ und $Q_1$ vom Grade $\gamma$, % [sic: Weber prints gamma although Q_0 was defined to have degree p^gamma.]
so können wir nach \S.\,31, III. einen Theiler $R$ von $Q$ vom Grade $p^{\gamma+1}$ bestimmen, von dem $Q_0$ ein Normaltheiler ist, und $R$
%CONT>>


% ===== scan leaf 157, printed 136 =====
% !PRINTED 136
\weberpage{136}
\pdfbookmark[2]{§ 34. Einfache Gruppen}{bm40}
%<<CONT
ist dann in keiner der Gruppen $Q_1, Q_2, \dots, Q_{q-1}$ enthalten, weil angenommen war, dass $Q$ mit keiner dieser Gruppen einen Theiler gemein hat, dessen Grad grösser als $p^{\gamma}$ ist. Ebenso können wir einen Theiler $R_1$ von $Q_1$ vom Grade $p^{\gamma+1}$ finden, von dem $Q_0$ Normaltheiler ist, und der in keiner der Gruppen $Q, Q_2, \dots, Q_{q-1}$ enthalten ist. Nun ist $R$ sowohl als $R_1$ in $P$ enthalten. Wir betrachten die kleinste Gruppe $S$, die $R$ und $R_1$ zugleich enthält, die jedenfalls in $P$ enthalten ist (das kleinste gemeinschaftliche Vielfache von $R$ und $R_1$). Der Grad dieser Gruppe $S$ kann nicht eine Potenz von $p$ sein, da sonst nach dem Satze II., \S.\,30 $S$ und mithin $R$ und $R_1$ in einer der Gruppen $Q, Q_1, \dots, Q_{q-1}$ enthalten sein müssten. Es ist aber $R$ nur in $Q$, $R_1$ nur in $Q_1$ enthalten; also ist diese Annahme unzulässig. Der Grad von $S$ muss also von der Form $p^{\lambda} q$ sein.

Es kann aber $\lambda$ nicht kleiner als $\alpha$ sein; denn es ist, da eine Gruppe vom Grade $p^{\gamma+1}$ in $S$ enthalten ist, nämlich $R$,

\weq{4}{\lambda \geqq \gamma+1 > \beta.}

Nach der in $\alpha)$ gemachten Voraussetzung ist aber in $P$ kein Theiler vom Grade $p^{\lambda} q$ enthalten, in dem $\lambda$ zugleich grösser als $\beta$ und kleiner als $\alpha$ ist. Nach (4) ist also nothwendig $\lambda = \alpha$, d.\,h. $S$ ist mit $P$ identisch.

Nun ist $Q_0$ ein Normaltheiler von $R$ und $R_1$. Fassen wir also alle Elemente $c$ von $P$, die der Bedingung

\[ c^{-1}\,Q_0\,c = Q_0 \]

genügen, zu einer Gruppe zusammen, so enthält diese Gruppe sowohl $R$ als $R_1$, und ist also die ganze Gruppe $P$. Es ist also $Q_0$ auch Normaltheiler von $P$. Da nach (3) $\gamma$ grösser als Null ist, so ist der Grad von $Q_0$ grösser als 1, und wir stossen auf einen Widerspruch mit der Annahme, dass $P$ einfach sei.

\begin{center}
\S.\,34.

\gesperrt{Einfache Gruppen.}
\end{center}

Die Sätze, die wir in den vorhergehenden Paragraphen kennen gelernt haben, gestatten ziemlich weitgehende Schlüsse über die Natur der Gruppen. Sie zeigen, dass wenigstens bei den niedrigeren Gradzahlen die metacyklischen (und cyklischen) Gruppen entschieden überwiegen. Um so höheres Interesse beanspruchen
%CONT>>


% ===== scan leaf 158, printed 137 =====
% !PRINTED 137
\weberpage{137}
%<<CONT
die wenigen darunter enthaltenen nicht metacyklischen und besonders die einfachen nicht cyklischen Gruppen.

Es ist bis jetzt nicht gelungen, die Gradzahlen der einfachen Gruppen in einem bestimmten Gesetze zusammenzufassen, und man hat sich begnügt, einerseits mit Hülfe der oben entwickelten Sätze, andererseits durch besondere Methoden alle einfachen Gruppen bis zu gewissen Grenzen der Gradzahlen zu ermitteln. \gesperrt{Hölder} hat diese Untersuchungen für die Gradzahlen bis 200 und \gesperrt{Cole} für die Gradzahlen bis 500 durchgeführt\footnote{\gesperrt{Hölder}, Mathematische Annalen, Bd.\,40. \gesperrt{Cole}, American Journal, Vol.\,14.}. Von einfachen nicht cyklischen Gruppen haben sich dabei nur die auch schon aus anderen Untersuchungen bekannten von den Graden 60, 168, 360 ergeben. Es ist ferner noch eine einfache Gruppe vom Grade 660 bekannt, und neuerdings haben \gesperrt{Cole} und \gesperrt{Moore} noch auf eine einfache Gruppe vom Grade 504 aufmerksam gemacht\footnote{Bulletin of the New York mathem. society, Oct.\,1893.}. Diese Untersuchungen werden mit dem Wachsen der Gradzahlen sehr mühsam. Wir wollen uns hier auf die Betrachtung der Gradzahlen des ersten Hunderts beschränken.

Hier erweist sich nach den drei allgemeinen Sätzen von \gesperrt{Sylow} und \gesperrt{Frobenius} die Mehrzahl der Gruppen als metacyklisch, und es bleiben nur die Gradzahlen

\[ 36,\ 60,\ 72,\ 84,\ 90,\ 100 \]

übrig. Dass aber eine Gruppe $36^{\mathrm{sten}}$ Grades nicht einfach sein kann, ergiebt sich nach dem \gesperrt{Sylow}'schen Satze. Denn eine solche Gruppe müsste einen Theiler $9^{\mathrm{ten}}$ Grades haben und müsste also durch die Permutationen von vier Ziffern darstellbar sein. Das ist aber unmöglich, weil es nur 24 Permutationen von vier Ziffern giebt. Eine Gruppe vom Grade $72 = 8 . 9$ muss nach \S.\,30, II., da 8 nicht $\equiv 1 \pmod{3}$ ist, einen Theiler vom Grade 18 enthalten, und wäre also, wenn sie einfach wäre, ebenfalls durch die Permutationen von vier Ziffern darstellbar, was noch weniger möglich ist. Dass einfache Gruppen von den Graden 84 und 100 nicht existiren können, ergiebt sich gleichfalls sofort aus den \gesperrt{Sylow}'schen Sätzen, da, wenn wir die Gruppen $7^{\mathrm{ten}}$ oder $25^{\mathrm{sten}}$ Grades aussondern, kein Theiler übrig bleibt, der nach dem Modul 7 oder 5 mit 1 congruent ist. Dass alle diese Gradzahlen nur metacyklischen Gruppen angehören können, folgt dann nach \S.\,6, 2. ohne Weiteres daraus, dass die Gruppen, deren Grade
%CONT>>


% ===== scan leaf 159, printed 138 =====
% !PRINTED 138
\weberpage{138}
%<<CONT
echte Theiler dieser Zahlen sind, schon als metacyklisch erwiesen sind.

Es bleibt nur noch die Untersuchung der Zahlen 60 und 90 übrig. Dass eine einfache Gruppe vom 60\textsuperscript{sten} Grade existirt, wissen wir schon; nämlich die alternirende Gruppe der Permutationen von fünf Buchstaben (Bd.\,I, \S.\,177). Es ist aber noch fraglich, ob dies die einzige ist.

Eine einfache Gruppe $P$ vom Grade 60 enthält nach \S.\,29, I. als Theiler eine Gruppe 5\textsuperscript{ten} Grades und eine Gruppe 3\textsuperscript{ten} Grades. Nimmt man in dem zweiten \gesperrt{Sylow'schen} Satze (\S.\,30) für $Q$ die Gruppe 5\textsuperscript{ten} Grades, so muss $R$ vom 10\textsuperscript{ten} Grade sein, weil der Index von $R$ congruent mit 1 nach dem Modul 5, also nur $= 6$ sein kann. Also kann die Gruppe $P$ vom 60\textsuperscript{sten} Grade nach \S.\,28, 2. als transitive Permutationsgruppe von sechs Ziffern dargestellt werden; sie kann aber keine cyklische Permutation von nur drei Ziffern enthalten, weil sie als einfache Gruppe (nach Bd.\,I, \S.\,158, 2.) primitiv sein muss, und daher, wenn sie einen dreigliedrigen Cyklus enthielte, nach Bd.\,I, \S.\,153, 10. die ganze alternirende Gruppe 360\textsuperscript{sten} Grades enthalten müsste. Wir können also eine der Permutationen 3\textsuperscript{ten} Grades in der Form annehmen:
\[
a = (1, 2, 3) (4, 5, 6).
\]
Hierzu wollen wir eine Permutation 5\textsuperscript{ten} Grades, $b$, fügen, die nur eine cyklische sein kann. Lassen wir darin die Ziffer 6 fehlen, so können wir statt $b$ eine solche Potenz von $b$ nehmen (die ja auch in $P$ vorkommen muss), dass etwa 1, 2 die beiden ersten Ziffern werden, und dann bleiben noch sechs Möglichkeiten für $b$ übrig:
\[
(1, 2, 3, 4, 5), \quad (1, 2, 4, 3, 5), \quad (1, 2, 4, 5, 3),
\]
\[
(1, 2, 3, 5, 4), \quad (1, 2, 5, 3, 4), \quad (1, 2, 5, 4, 3);
\]
von diesen sechs Annahmen sind aber wieder die drei in einer Reihe stehenden nicht wesentlich (d.\,h. nur durch die Bezeichnung) verschieden. Denn ersetzt man z.\,B. bei der zweiten Annahme $b = (1, 2, 4, 3, 5)$ das Element $a$ durch $a^2$ und $b$ durch $b^3$, so erhält man
\[
(1, 3, 2) (4, 6, 5); \quad (1, 3, 2, 5, 4),
\]
was durch Vertauschung der Ziffern 2 mit 3 und 4 mit 5 in die erste Annahme übergeht. Die dritte Annahme $b = (1, 2, 4, 5, 3)$ geht, wenn man $b$ durch $b^4$ und $a$ durch $a^2$ ersetzt und dann
%CONT>>


% ===== scan leaf 160, printed 139 =====
% !PRINTED 139
\weberpage{139}
%<<CONT
1 mit 2 und 4 mit 5 vertauscht, in die erste über, und ebenso lassen sich die drei anderen Annahmen über $b$ auf einander zurückführen. Es bleiben also nur zwei Möglichkeiten zu untersuchen:
\[
\begin{gathered}
a = (1, 2, 3) (4, 5, 6), \quad b = (1, 2, 3, 4, 5) \\
a = (1, 2, 3) (4, 5, 6), \quad b = (1, 2, 3, 5, 4).
\end{gathered}
\]
Davon ist aber die erste zu verwerfen, weil sie in
\[
a^2 b = (1, 4, 6)
\]
einen dreigliedrigen Cyklus ergeben würde, der nicht vorkommen kann. Es bleibt also nur die einzige Möglichkeit:
\[
a = (1, 2, 3) (4, 5, 6), \quad b = (1, 2, 3, 5, 4),
\]
und es giebt also, wenn man äquivalente Gruppen als nicht verschieden betrachtet, nur \gesperrt{eine} einfache Gruppe 60\textsuperscript{sten} Grades.

Diese beiden Elemente $a, b$ kann man als erzeugende Elemente der ganzen Gruppe betrachten und man kann durch ihre wiederholte Zusammensetzung auf sehr mannigfaltige Art die ganze Gruppe vom 60\textsuperscript{sten} Grade herleiten. Wir bekommen ausser dem Einheitselemente die Permutationen 5\textsuperscript{ten} Grades:
\[
\begin{gathered}
(2, 3, 6, 5, 4), \quad (1, 3, 5, 6, 4), \quad (1, 2, 5, 4, 6), \\
(1, 2, 6, 3, 5), \quad (1, 2, 4, 6, 3), \quad (1, 2, 3, 5, 4),
\end{gathered}
\]
die mit ihren Potenzen 24 ergeben; ferner die Permutationen 3\textsuperscript{ten} Grades:
\[
\begin{array}{ll}
(1, 2, 3) (4, 5, 6), & (1, 3, 5) (2, 4, 6), \\
(1, 2, 4) (3, 6, 5), & (1, 3, 6) (2, 4, 5), \\
(1, 2, 5) (3, 6, 4), & (1, 4, 5) (2, 3, 6), \\
(1, 2, 6) (3, 4, 5), & (1, 4, 6) (2, 3, 5), \\
(1, 3, 4) (2, 6, 5), & (1, 5, 6) (2, 3, 4),
\end{array}
\]
die mit ihren Quadraten zusammen 20 ergeben. Endlich erhält man noch 15 Permutationen 2\textsuperscript{ten} Grades:
\[
\begin{array}{lll}
(1, 2) (3, 4), & (1, 2) (5, 6), & (3, 4) (5, 6), \\
(1, 3) (4, 5), & (1, 3) (2, 6), & (2, 6) (4, 5), \\
(1, 4) (2, 5), & (1, 4) (3, 6), & (2, 5) (3, 6), \\
(1, 5) (2, 3), & (1, 5) (4, 6), & (2, 3) (4, 6), \\
(1, 6) (2, 4), & (1, 6) (3, 5), & (2, 4) (3, 5),
\end{array}
\]
von denen je drei in einer Reihe stehenden mit dem Einheitselemente zusammen eine Gruppe 4\textsuperscript{ten} Grades bilden.

Dass diese Gruppe auch dargestellt werden kann durch die Permutationen von fünf Ziffern, ergiebt sich nun schon daraus,
%CONT>>


% ===== scan leaf 161, printed 140 =====
% !PRINTED 140
\weberpage{140}
\pdfbookmark[2]{§ 35. Gruppen vom Grade p q}{bm41}
%<<CONT
dass eine einfache Permutationsgruppe 60\textsuperscript{sten} Grades bei fünf Ziffern wirklich existirt, nämlich die alternirende Gruppe.

Man kann aber auch die Thatsache dadurch verificiren, dass man einen Theiler 12\textsuperscript{ten} Grades von $P$ nachweist (\S.\,28, 3.). Man erhält diesen Theiler 12\textsuperscript{ten} Grades z. B. daraus, dass die Gruppe 4\textsuperscript{ten} Grades:
\[
Q = 1, (1, 2) (3, 4), (1, 2) (5, 6), (3, 4) (5, 6)
\]
mit den Permutationen 3\textsuperscript{ten} Grades:
\[
c = (1, 3, 5) (2, 4, 6)
\]
vertauschbar ist, d. h. der Bedingung $c^{-1} Q c = Q$ genügt, und erhält also dann eine Gruppe 12\textsuperscript{ten} Grades:
\[
Q, \; Q c, \; Q c^2.
\]
Die einfache Gruppe 60\textsuperscript{sten} Grades wird auch die \gesperrt{Ikosaëdergruppe} genannt aus einem Grunde, den wir später kennen lernen werden.

Dieselbe Betrachtung zeigt auch, dass eine einfache Gruppe 90\textsuperscript{sten} Grades nicht existirt. Denn diese Gruppe müsste ein Element 5\textsuperscript{ten} und ein Element 3\textsuperscript{ten} Grades enthalten. Da $90 = 5 . 18$ ist und 18 keinen anderen Theiler als 6 hat, der nach dem Modul 5 mit 1 congruent ist, so muss, wenn wir für $Q$ eine Gruppe 5\textsuperscript{ten} Grades wählen, $R$ vom Index 6 sein und die Gruppe $P$ wäre also als transitive Permutationsgruppe von sechs Ziffern darstellbar. Ganz wie oben schliesst man, dass sie die beiden Elemente $a, b$ enthalten müsste, was, wie wir gesehen haben, zur Ikosaëdergruppe führt, die nicht Theiler einer Gruppe 90\textsuperscript{sten} Grades sein kann.

\begin{center}
\S.\,35.

\gesperrt{Gruppen vom Grade $pq$.}
\end{center}

Es ist nun noch von Interesse, zu untersuchen, wie bei einer gegebenen Gradzahl die Gruppen constituirt sind. Dafür sind die Sätze, die in den vorangegangenen Paragraphen abgeleitet sind, die Grundlagen. Freilich sind wir bis jetzt nur bei den einfacheren Gradzahlen im Stande, die vorhandenen Gruppen vollständig zu übersehen. Am weitesten geht hierin eine Arbeit von \gesperrt{Hölder}\footnote{\gesperrt{Hölder}, Die Gruppen der Ordnungen $p^3, pq^2, pqr, p^4$. Mathematische Annalen, Bd. 43.}. Wir betrachten hier nur den einfachen Fall,
%CONT>>


% ===== scan leaf 162, printed 141 =====
% !PRINTED 141
\weberpage{141}
%<<CONT
dass der Grad $pq$ das Product zweier Primzahlen ist, die auch einander gleich sein können.

Eine solche Gruppe $P$ ist metacyklisch und hat einen Normaltheiler $Q$ vom Grade $p$, der also eine cyklische Gruppe ist, die wir so darstellen:
\weq{1}{Q = 1, a, a^2, \dots a^{p-1} \quad (a^p = 1).}
Ist nun $b$ ein nicht in $Q$ enthaltenes Element von $P$, so ist
\[
b^{-1} Q b = Q.
\]
Wenn $b^h$ die niedrigste Potenz von $b$ ist, die in $Q$ vorkommt, so ist $Q + Qb + \cdots Qb^{h-1}$ eine Gruppe, die mit $P$ identisch sein muss, und folglich muss $h = q$ sein. Wir können also $P$ so darstellen:
\weq{2}{P = Q + Qb + Qb^2 + \cdots + Qb^{q-1}.}
Ist nun $g$ eine primitive Wurzel der Primzahl $p$ und $\nu$ ein vorläufig noch unbestimmter Exponent, der nach dem Modul $p - 1$ zu nehmen ist, so folgt, da $b^{-1} a b$ in $Q$ enthalten ist,
\weq{3}{b^{-1} a b = a^{g^\nu},}
woraus sich für jeden Exponenten $\alpha$ ergiebt:
\weq{4}{b^{-1} a^\alpha b = a^{\alpha g^\nu},}
und durch wiederholte Zusammensetzung mit $b$:
\weq{5}{b^{-\beta} a^\alpha b^\beta = a^{\alpha g^{\nu\beta}}.}
Nun ist jedenfalls $b^{pq} = 1$, und wenn wir also in (5) $\beta = pq$ setzen, so folgt, dass für jedes $\alpha$
\[
a^{\alpha g^{\nu pq}} = a^\alpha,
\]
also $g^{\nu pq} \equiv 1 \pmod{p}$ oder
\weq{6}{\nu pq \equiv 0 \pmod{p - 1}}
sein muss. Hieraus schliesst man weiter, dass, wenn $p - 1$ nicht durch $q$ theilbar ist, was immer eintritt, wenn $p = q$ ist, $\nu$ durch $p - 1$ theilbar, also nach (3)
\[
ab = ba,
\]
und in Folge dessen auch für jedes $\alpha, \beta, \alpha', \beta'$:
\[
a^\alpha b^\beta = b^\beta a^\alpha, \quad a^\alpha b^\beta a^{\alpha'} b^{\beta'} = a^{\alpha'} b^{\beta'} a^\alpha b^\beta,
\]
d. h. dass die Gruppe eine \gesperrt{Abel}'sche sein muss. Diesen Fall haben wir aber schon im zweiten Abschnitte dieses Bandes untersucht und gefunden, dass es, wenn $p = q$ ist, zwei Gruppen (mit
%CONT>>


% ===== scan leaf 163, printed 142 =====
% !PRINTED 142
\weberpage{142}
%<<CONT
den beiden Invarianten $p, p$ oder mit einer Invariante $p^2$), und
wenn $p$ von $q$ verschieden ist, eine Gruppe (mit den beiden Invarianten
 $p, q$) giebt.

Es bleibt uns also noch der Fall
\[
p \equiv 1\ (\mathrm{mod}\ q)
\]
zu untersuchen, in dem $\nu$ jeden der Bedingung
\[
\nu \equiv 0\ \left(\mathrm{mod}\ \frac{p-1}{q}\right)
\]
genügenden Werth haben kann, deren es $q$ nach dem Modul
$p - 1$ verschiedene giebt, nämlich:
\weq{7}{\nu = 0,\quad \frac{p-1}{q},\quad 2\,\frac{p-1}{q},\quad \dots\ \frac{(q-1)(p-1)}{q}.}
Der Fall $\nu = 0$ führt auch hier auf eine \gesperrt{Abel'sche} Gruppe.
Es sind aber auch die anderen Werthe von $\nu$ zulässig, die zu
ebenso vielen nicht commutativen Gruppen führen. Diese nicht
commutativen Gruppen sind unter einander isomorph und werden
auf einander zurückgeführt, wenn man $b$ durch ein $b^\beta$ ersetzt,
wie die Formel (5) zeigt.

Wir können $b^q = 1$ annehmen, denn es ist gewiss $b^q$ in $Q$
enthalten und $b^{pq} = 1$. Wenn also $b^q$ nicht $= 1$ ist, so ersetzen
wir $b$ durch $b^p$.

Man erhält eine der nicht commutativen Gruppen, wenn
man in
\weq{8}{\Theta = a^\alpha b^\beta}
$\alpha$ die Reihe der Zahlen $0, 1, \dots p - 1$ und $\beta$ die Zahlen
$0, 1, \dots q - 1$ durchlaufen lässt. Die Zusammensetzung zweier
Elemente dieser Gruppe ergiebt sich nach (5) aus
\weq{9}{b^\beta a^\alpha = a^{\alpha g^{-\nu\beta}} b^\beta, \quad a^p = 1,\ b^q = 1,}
wodurch man jedes Compositum aus Elementen $\Theta$ auf die Form
$\Theta$ zurückführen kann.

Um zu zeigen, dass die Elemente (8) nach den Compositionsregeln
 (9) wirklich eine Gruppe bilden, hat man die Eigenschaften
der Gruppe, nämlich, dass aus $\Theta\Theta' = \Theta\Theta''$ und aus $\Theta'\Theta = \Theta''\Theta$
folgt, dass $\Theta' = \Theta''$ ist, und ferner das associative Gesetz
\[
\Theta(\Theta'\Theta'') = (\Theta\Theta')\Theta''
\]
nachzuweisen. Beides aber ergiebt sich sehr leicht aus der Zusammensetzung,
die aus (9) folgt:
\weq{10}{a^\alpha b^\beta a^{\alpha'} b^{\beta'} = a^{\alpha+\alpha' g^{-\nu\beta}} b^{\beta+\beta'}.}


% ===== scan leaf 164, printed 143 =====
% !PRINTED 143
\weberpage{143}
\pdfbookmark[2]{§ 36. Grenzen des Index eines Theilers der symmetrischen Permutationsgruppe}{bm42}
\begin{center}
\S.\,36.

\gesperrt{Grenzen des Index eines Theilers der symmetrischen
Permutationsgruppe.}
\end{center}

Wir beschliessen diese Betrachtungen mit dem Beweise eines
Satzes über Permutationsgruppen, der durch die Schwierigkeit,
die sein Beweis anfangs bot, eine gewisse Berühmtheit erlangt
hat, und der für die Beurtheilung algebraischer Fragen von
Wichtigkeit ist\,$^{1}$).

Es handelt sich dabei um die symmetrische Permutationsgruppe
 $P$ von $n$ Ziffern und um ihre Theiler von möglichst
kleinem Index. Wir wissen, dass die Gruppe $P$ immer einen
Theiler vom Index 2 hat, nämlich die alternirende Gruppe.
Ausserdem ist noch ein Theiler vom Index $n$ bekannt, der alle
Permutationen von $P$ umfasst, die eine Ziffer ungeändert lassen,
der also intransitiv ist.

Zunächst gilt der folgende Satz:

I. \gesperrt{Der Index eines imprimitiven Theilers von $P$ ist
immer grösser als $n$, und der Index eines intransitiven
 Theilers ist gleich oder grösser
als $n$, und nur dann gleich $n$, wenn der Theiler
eine Ziffer in Ruhe lässt, und die übrigen
$n - 1$ Ziffern auf alle mögliche Arten permutirt.}

Nehmen wir an, es sei $Q$ ein imprimitiver Theiler von $P$
vom Index $j$, und es bestehen $r$ Systeme der Imprimitivität von
je $s$ Ziffern, so dass $n = rs$ ist. Eine Zahl, die der Grad der
Gruppe $Q$ sicher nicht übersteigen kann, erhalten wir, wenn wir
alle Permutationen in jedem einzelnen der $r$ Systeme und dann
noch sämmtliche Permutationen der Systeme abzählen. Der Grad
von $Q$ ist also kleiner oder gleich dem Producte $[\Pi(s)]^r\,\Pi(r)$,
wenn für jede ganze Zahl $n$
\[
\Pi(n) = 1.2.3\dots n
\]
%CONT>>

$^{1}$) \gesperrt{Bertrand}, Journal de Mathématiques, Tome XV (1845). \gesperrt{Serret},
Algèbre supér., Section IV, Chapitre III. C. \gesperrt{Jordan}, Traité des substitutions,
p.\,67. \gesperrt{Netto}, Substitutionentheorie, Capitel VI. Crelle's Journ.,
Bd. 100.


% ===== scan leaf 165, printed 144 =====
% !PRINTED 144
\weberpage{144}
%<<CONT
ist, und folglich ist
\weq{1}{j \geqq \frac{\Pi(n)}{[\Pi(s)]^r\,\Pi(r)}.}
Es ist leicht einzusehen, dass diese Zahl, wenn keiner der
Factoren $r, s$ gleich 1 ist, grösser als $n$ ist. Dies ergiebt sich,
wenn wir den Quotienten so schreiben:
\[
\frac{\Pi(n)}{[\Pi(s)]^r\,\Pi(r)} = \frac{(r+1)(r+2)\dots n}{(2.3\dots s)^r} =
\]
\[
\frac{n}{2}\cdot\left(\frac{r+1}{2}\cdot\frac{r+2}{2}\dots\frac{2r-1}{2}\right)\left(\frac{2r}{3}\cdot\frac{2r+1}{3}\dots\frac{3r-1}{3}\right)\dots
\]
\[
\left(\frac{(s-1)r}{s}\cdot\frac{(s-1)r+1}{s}\dots\frac{n-1}{s}\right).
\]
Denn hiernach ist
\weq{2}{j > \frac{n}{2}\left(\frac{r+1}{2}\right)^{r-1}\left(\frac{2r}{3}\right)^r \dots \left(\frac{(s-1)r}{s}\right)^r.}
Die Factoren $\dfrac{r+1}{2}, \dfrac{2r}{3}, \dots \dfrac{(s-1)r}{s}$ sind alle gleich oder
grösser als 1, denn es ist
\[
\frac{(h-1)r}{h} - 1 = \frac{(h-1)(r-1)-1}{h},
\]
also immer positiv, und der erste Factor
\[
\left(\frac{r+1}{2}\right)^{r-1}
\]
ist grösser als 2, wenn $r > 2$ ist. Ist aber $r = 2$, so ist das
Product der beiden Factoren:
\[
\left(\frac{r+1}{2}\right)^{r-1}\left(\frac{2r}{3}\right)^r = \frac{3}{2}\left(\frac{4}{3}\right)^2 = \frac{8}{3} > 2,
\]
und folglich ist unter allen Umständen $j > n$.

Ist zweitens die Gruppe $Q$ intransitiv, und zerfällt das
System der $n$ Ziffern in zwei Systeme von je $a$ und $b$ Ziffern, so
dass die Ziffern dieser beiden Systeme durch $Q$ nur unter einander
 vertauscht werden, und $n = a + b$ ist, so sind alle Permutationen
 von $Q$ in der Gruppe enthalten, die aus allen Permutationen
 der $a$ und der $b$ Ziffern besteht; d. h. der Grad
von $Q$ ist gleich oder kleiner als $\Pi(a)\Pi(b)$, und folglich der
Index $j$ von $Q$
\weq{3}{j \geqq \frac{\Pi(n)}{\Pi(a)\,\Pi(b)} = \frac{n}{1}\,\frac{n-1}{2}\dots\frac{b+1}{a}.}


% ===== scan leaf 166, printed 145 =====
% !PRINTED 145
\weberpage{145}
Nehmen wir, was freisteht, an, dass $b \geqq a$ sei, so ist der
Ausdruck auf der rechten Seite dieser Ungleichung grösser
als $n$, und nur dann gleich $n$, wenn $b = n - 1, a = 1$ ist. In
diesem speciellen Falle kann $j = n$ werden, aber nur dann,
wenn die $n - 1$ Elemente durch $Q$ auf alle möglichen Arten
permutirt werden, was in dem Satze I. ausgesprochen ist.

Die Ausdrücke (3) für die untere Grenze von $j$ sind nichts
Anderes, als die Binomialco\"efficienten $B_a^{(n)}$, deren Bildung
sofort zeigt, dass sie bis zur Mitte hin, d. h. so lange $2a \leqq n$,
eine \gesperrt{wachsende Zahlenreihe} bilden.

Wir wollen für den weiteren Gebrauch hieraus den Schluss
ziehen:

$\alpha$) \gesperrt{Ist $a = 1$, lässt also die Gruppe $Q$ eine Ziffer
ungeändert, so ist ihr Grad ein Theiler von
$\Pi(n-1)$, ist aber $a > 1$, so ist der Grad von $Q$
$\leqq \Pi(n-2)\,\Pi(2)$.}

Der Satz, den wir ferner noch beweisen wollen, lautet nun:

\gesperrt{II. Ausser der alternirenden Gruppe giebt es keinen
transitiven und primitiven Theiler $Q$ von $P$,
dessen Index $\leqq n$ ist, ausgenommen in den
beiden Fällen $n = 4, n = 6$.}

Beim Beweise machen wir Gebrauch von den Sätzen (\S.\,153,
9., 10. des ersten Bandes), dass ein transitiver und primitiver
Theiler von $P$, der nicht die ganze alternirende Gruppe enthält,
und daher nicht mit der alternirenden oder der symmetrischen
Gruppe selbst identisch ist, keine Transposition und keine
cyklische Permutation von nur drei Ziffern enthalten kann.

Es sei also $P$ die symmetrische Permutationsgruppe der
$n$ Ziffern $0, 1, 2 \dots n - 1$ und $Q$ ein primitiver und transitiver
Theiler von $P$ vom Index $j$, der nicht die alternirende Gruppe
enthält. Es sei ferner $P$ in die Nebengruppen zerlegt:
\weq{4}{P = Q + Q\pi_1 + Q\pi_2 + \dots + Q\pi_{j-1}.}
Wir betrachten das System der Transpositionen:
\weq{5}{(0, 1), (0, 2), \dots (0, n-1),}
deren keine in $Q$ vorkommen kann. Ist nun $j < n$, so müssen
wenigstens zwei dieser Permutationen, etwa $(0,1), (0,2)$, in der-
%CONT>>


% ===== scan leaf 167, printed 146 =====
% !PRINTED 146
\weberpage{146}
%<<CONT
selben Nebengruppe, etwa in $Q\pi_1$, vorkommen, und folglich giebt es zwei
Permutationen $\varkappa_1, \varkappa_2$ in $Q$, die der Bedingung

\[ \varkappa_1 \pi_1 = (0, 1), \quad \varkappa_2 \pi_1 = (0, 2) \]

genügen. Es ist daher

\[ \varkappa_1 \pi_1 \pi_1^{-1} \varkappa_2^{-1} = \varkappa_1 \varkappa_2^{-1} = (0,1)(0,2) = (0,1,2) \]

in $Q$ enthalten. Dies aber widerspricht unserer Voraussetzung, dass $Q$ keinen
dreigliedrigen Cyklus enthalten soll. Demnach kann $j$ nicht $< n$ sein. Ist aber
$j = n$, also der Grad von $Q$ gleich $\Pi(n-1)$, so müssen die $n-1$ Permutationen
(5) in $n-1$ verschiedenen Nebengruppen vorkommen, und $P$ lässt sich so darstellen:

\weq{6}{P = Q + Q(0,1) + Q(0,2) + \cdots + Q(0, n-1).}

Betrachten wir irgend eine andere Transposition, z. B. $(2,3)$, so kann diese
nicht in $Q$ und nicht in $Q(0,2)$ oder in $Q(0,3)$ vorkommen, weil sonst
$(2,3)(0,2) = (0,2,3)$ oder $(2,3)(0,3) = (0,3,2)$ in $Q$ vorkäme. Wir können also,
ohne die Allgemeinheit zu beeinträchtigen, annehmen, dass $(2,3)$ in $Q(0,1)$
vorkommt, und daraus ergiebt sich:

\gesperrt{$\beta$) Die Gruppe $Q$ enthält das Transpositionspaar}

\weq{7}{(0,1)(2,3).}

Die Gruppe $Q$ hat einen Theiler $Q_0$, der aus allen den Permutationen besteht,
die die Ziffer 0 an ihrer Stelle lassen. Da $Q$ transitiv ist, so ist, wenn
$\varkappa_1, \varkappa_2, \ldots, \varkappa_{n-1}$ Elemente aus $Q$ sind, die die 0
in 1, in 2, \ldots, in $n-1$ überführen,

\weq{8}{Q = Q_0 + Q_0 \varkappa_1 + \cdots + Q_0 \varkappa_{n-1},}

und da $Q$ vom Grade $\Pi(n-1)$ ist, so folgt hieraus, dass $Q_0$ vom Grade

\weq{9}{g = \frac{\Pi(n-1)}{n}}

ist. Da $g$ eine ganze Zahl, also $\Pi(n-1)$ durch $n$ theilbar sein muss, so
schliessen wir zunächst, dass der Fall $j = n$ niemals eintreten kann, wenn $n$
eine Primzahl ist, und für diesen Fall ist also unser Theorem II. bewiesen.

Im Allgemeinen können wir aber schliessen:

\gesperrt{$\gamma$) Ist $n$ nicht $= 4$, so sind die $n-1$ Ziffern
$1, 2, \ldots, n-1$ durch $Q_0$ noch transitiv verbunden.}


% ===== scan leaf 168, printed 147 =====
% !PRINTED 147
\weberpage{147}
Wäre nämlich $Q_0$ in diesen $n-1$ Ziffern intransitiv, so müsste nach $\alpha$)

\[
\begin{aligned}
&\text{entweder } \Pi(n-2) \text{ theilbar durch } g,\\
&\text{oder } \Pi(n-3)\,\Pi(2) \geqq g
\end{aligned}
\]

sein; also nach (9):

\[
\begin{aligned}
&\text{entweder } \frac{n}{n-1} \text{ eine ganze Zahl, also } n \geqq 2(n-1)\\
&\text{oder } n^2 - 5n + 2 \leqq 0.
\end{aligned}
\]

Das Eine ist unmöglich, wenn $n > 2$ ist, das Andere, wenn $n > 4$ ist.

Wir sehen von dem Falle $n = 4$ ab und betrachten jetzt den Theiler $Q_{0,1}$ von
$Q$, der die beiden Ziffern 0, 1 ungeändert lässt.

Da $Q_{0,1}$, als Permutationsgruppe der $n-1$ Ziffern betrachtet, ein transitiver
Theiler von $Q_0$ ist, so ist, wie man durch nochmalige Anwendung der Zerlegung (8)
schliesst, der Grad $g_1$ von $Q_{0,1}$ gleich $g : n-1$, also

\weq{10}{g_1 = \frac{\Pi(n-1)}{n(n-1)} = \frac{\Pi(n-2)}{n}.}

Daraus schliessen wir ähnlich wie oben:

\gesperrt{$\delta$) Ist $n$ nicht $= 6$, so sind die $n-2$ Ziffern
$2, 3, \ldots, n-1$ durch $Q_{0,1}$ transitiv verbunden.}

Denn wären sie es nicht, so müsste nach $\alpha$):

\[
\begin{aligned}
&\text{entweder } \Pi(n-3) \text{ theilbar durch } g_1,\\
&\text{oder } \Pi(n-4)\,\Pi(2) \geqq g_1
\end{aligned}
\]

sein; also nach (10):

\[
\begin{aligned}
&\text{entweder } \frac{n}{n-2} \text{ eine ganze Zahl, also } n \geqq 2n - 4,\\
&\text{oder } n^2 - 7n + 6 \leqq 0.
\end{aligned}
\]

Das Erste ist nicht möglich, wenn $n > 4$, das Zweite, wenn $n > 6$ ist.

Ist nun $n > 6$, so gilt noch Folgendes:

\gesperrt{$\varepsilon$) Die Gruppe $Q_{0,1,2}$, die die drei Ziffern 0, 1, 2
ungeändert lässt, hat den Grad}

\[ g_2 = \frac{g_1}{n-2} = \frac{\Pi(n-3)}{n}, \]

\gesperrt{und es ist nicht möglich, dass durch die ganze Gruppe $Q_{0,1,2}$ noch eine
vierte Ziffer 3 ungeändert bleibt.}


% ===== scan leaf 169, printed 148 =====
% !PRINTED 148
\weberpage{148}
Der Grad $g_2$ ergiebt sich genau wie der von $Q_0$ und $Q_{0,1}$. Wenn aber
durch $Q_{0,1,2}$ noch eine vierte Ziffer 3 ungeändert bliebe, so wäre $g_2$ ein
Theiler des Grades der symmetrischen Gruppe von $n-4$ Ziffern, also ein Theiler
von $\Pi(n-4)$. Es müsste also

\[ \frac{n\,\Pi(n-4)}{\Pi(n-3)} = \frac{n}{n-3} \]

eine ganze Zahl sein, also $n \geqq 2n - 6$ oder $n \leqq 6$.

Aus $\varepsilon$) schliessen wir, dass es in $Q$ eine Permutation $\varkappa$
giebt, durch die $0, 1, 2, 3$ in $0, 1, 2, 4$ übergeht, worin 4 eine von 3
verschiedene Ziffer ist. Nach $\beta$) enthält also $Q$ auch die Permutation:

\[ \varkappa^{-1}(0,1)(2,3)\varkappa = (0,1)(2,4), \]

folglich auch:

\[ (0,1)(2,3)(0,1)(2,4) = (2,3,4), \]

also einen dreigliedrigen Cyklus, was der Voraussetzung widerspricht. Hiernach
ist das Theorem II. vollständig bewiesen.

Dass die Fälle $n = 4$ und $n = 6$ wirklich Ausnahmen bilden, geht aus Bd. I,
\S.\,160 und \S.\,182 hervor, wo wir gesehen haben, \gesperrt{dass die symmetrische
Permutationsgruppe von vier Ziffern einen transitiven Theiler vom Index 3 und
die von sechs Ziffern einen transitiven Theiler vom Index 6 besitzt.}

\begin{center}
\rule{3cm}{0.4pt}
\end{center}


% ===== scan leaf 170, printed 149 =====
\weberpage{149}% page number not printed on this page
% !PRINTED none
% Part-title (Buch-Titel) page: mostly blank, with two centred rules framing
% the book/part heading. No running head, no printed page number.
\begin{center}
\gesperrt{ZWEITES BUCH.}

\rule{3cm}{0.4pt}

\vspace{2cm}

\gesperrt{LINEARE GRUPPEN.}

\vspace{2cm}

\rule{5cm}{0.4pt}
\end{center}


% ===== scan leaf 171, printed 150 =====
\weberpage{150}% page number not printed on this page
% !PRINTED none
% Blank page (verso). All three bands show only faint bleed-through ghosting
% from the facing/adjacent printed page; no legible running head, page
% number, or body text is present anywhere on this page image.


% ===== scan leaf 172, printed 151 =====
\weberpage{151}% page number not printed on this page
\pdfbookmark[0]{Zweites Buch. Lineare Gruppen.}{bm43}
\pdfbookmark[1]{Sechster Abschnitt. Gruppen linearer Substitutionen.}{bm44}
\pdfbookmark[2]{§ 37. Lineare Substitutionen und ihre Zusammensetzung}{bm45}
% !PRINTED none
% Chapter-opening page: no running head/folio visible in any of the three
% bands (consistent with an Abschnitt-opening page).

\begin{center}
\gesperrt{Sechster Abschnitt.}

\gesperrt{Gruppen linearer Substitutionen.}
\end{center}

\begin{center}
\rule{3cm}{0.4pt}
\end{center}

\begin{center}
§. 37.

\gesperrt{Lineare Substitutionen und ihre Zusammensetzung.}
\end{center}

Eines der wirksamsten Mittel zur Bildung von Gruppen, auf welches zugleich viele Anwendungen führen, sind die linearen Substitutionen und ihre Zusammensetzung. Wir sind schon mehrfach solchen linearen Substitutionen begegnet und haben sie z. B. im zweiten Abschnitte des ersten Bandes bei Gelegenheit des Multiplicationsgesetzes der Determinanten, und sodann im elften Abschnitte bei den äquivalenten Zahlen betrachtet.

Unter einer \gesperrt{linearen Substitution} von $n$ Variablen verstehen wir ein System von Gleichungen, durch das ein System von $n$ Veränderlichen $y_1, y_2, \dots y_n$ linear durch ein anderes System $x_1, x_2, \dots x_n$ ausgedrückt wird. Wir unterscheiden nach der Anzahl der Variablen \gesperrt{unäre, binäre, ternäre, quaternäre} Substitutionen, und wollen im Allgemeinen die Zahl der Variablen die \gesperrt{Dimension} der Substitution nennen. Wir beschränken uns fürs erste auf \gesperrt{homogene Substitutionen}, so dass, wenn mit $a_i^{(\varkappa)}$ die Coëfficienten bezeichnet werden, die Substitution durch das Gleichungssystem

\weq{1}{y_\varkappa = \sum_{1,n}^{i} a_i^{(\varkappa)} x_i \qquad \varkappa = 1, 2, \dots n}

dargestellt ist. Oft kommt es auf die Variablen selbst nicht an, so dass eine solche Substitution durch ihre Coëfficienten $a_i^{(\varkappa)}$ hinlänglich gekennzeichnet ist. Man benutzt zuweilen auch zur Be-
%CONT>>


% ===== scan leaf 173, printed 152 =====
% !PRINTED 152
\weberpage{152}
%<<CONT
zeichnung einer Substitution einen einfachen Buchstaben, etwa $A$, und setzt dann

\weq{2}{A = \begin{pmatrix} a_1^{(1)}, & a_2^{(1)}, & \dots & a_n^{(1)} \\ a_1^{(2)}, & a_2^{(2)}, & \dots & a_n^{(2)} \\ \multicolumn{4}{c}{\cdots\cdots\cdots\cdots} \\ a_1^{(n)}, & a_2^{(n)}, & \dots & a_n^{(n)} \end{pmatrix}.}

Wir werden auch abkürzend $A = (a_i^{(\varkappa)})$ setzen und die Determinante der Substitutionscoëfficienten mit $|A|$ bezeichnen. Diese Determinante wird immer von Null verschieden vorausgesetzt. Das Gleichungssystem (1) stellen wir auch symbolisch so dar:

\weq{3}{(y_1, y_2, \dots y_n) = A (x_1, x_2, \dots x_n),}

oder noch kürzer:

\weq{4}{(y) = A (x).}

Die Substitution:

\weq{5}{y_1 = x_1, \quad y_2 = x_2, \quad \dots \quad y_n = x_n}

oder

\weq{6}{J = \begin{pmatrix} 1, & 0, & \dots & 0 \\ 0, & 1, & \dots & 0 \\ \multicolumn{3}{c}{\cdots\cdots\cdots} \\ 0, & 0, & \dots & 1 \end{pmatrix}}

heisst die \gesperrt{identische Substitution}. Eine Substitution der Form

\weq{7}{y_1 = \mu_1 x_1, \quad y_2 = \mu_2 x_2, \quad \dots, \quad y_n = \mu_n x_n}

oder

\weq{8}{M = \begin{pmatrix} \mu_1, & 0, & \dots & 0 \\ 0, & \mu_2, & \dots & 0 \\ \multicolumn{3}{c}{\cdots\cdots\cdots} \\ 0, & 0, & \dots & \mu_n \end{pmatrix}}

soll eine \gesperrt{multiplicative Substitution} oder kurz eine \gesperrt{Multiplication} und die Elemente $\mu_1, \mu_2, \dots, \mu_n$ die \gesperrt{Multiplicatoren} genannt werden, ein Ausdruck, der sich durch die Gleichungen (7) rechtfertigt.

Nehmen wir eine zweite lineare Substitution der Dimension $n$ an, durch die ein neues System von Variablen $z$ eingeführt wird:

\weq{9}{z_i = \sum_{1,n}^{h} b_h^{(i)} y_h, \quad (z) = B (y),}

und führen für $y$ die Werthe aus (1) ein, so erhalten wir die $z$
%CONT>>


% ===== scan leaf 174, printed 153 =====
% !PRINTED 153
\weberpage{153}
%<<CONT
durch die $x$ ausgedrückt mittelst einer neuen linearen Substitution $E$, die wir so bezeichnen können:

\weq{10}{z = E (x) = B A (x),}

und die Substitution $E = BA$ heisst aus $B$ und $A$ \gesperrt{zusammengesetzt} oder \gesperrt{componirt}. Es ist dabei aber zwischen $AB$ und $BA$ zu unterscheiden. Zwei Substitutionen $A, B$ von der besonderen Eigenschaft, dass $AB = BA$ ist, heissen mit einander \gesperrt{vertauschbar} oder \gesperrt{commutativ}. Die Substitutionscoëfficienten von $E$ ergeben sich durch Einsetzen der Ausdrücke (1) in (9):

% [sic: Weber prints i below the first sum although h is its summation index; 1,n is expected.]
\weq{11}{\begin{gathered} z_k = \sum_{i}^{h} e_h^{(k)} x_h, \\ e_h^{(k)} = \sum_{1,n}^{i} b_i^{(k)} a_h^{(i)}. \end{gathered}}

Diese Formeln sind ganz dieselben, die wir im §. 27 des ersten Bandes benutzt haben\footnote{Nur war dort aus einem leicht ersichtlichen Grunde die Bezeichnung etwas anders.}, und wir können demnach das in (11) ausgedrückte Gesetz der Composition der Substitutionen so ausdrücken:

\gesperrt{1. Um die aus zwei Substitutionen $B, A$ zusammengesetzte Substitution $BA$ zu bilden, verfährt man ganz so, als ob die beiden Determinanten $|B|, |A|$ nach der Multiplicationsregel mit einander multiplicirt werden sollten. Es sind dabei, um die Elemente einer Zeile zu bilden, die Elemente einer Zeile der ersten Componente mit den entsprechenden Elementen der Colonnen der zweiten Componente zu multipliciren und dann zu addiren. Man drückt dies auch kurz so aus, dass in der ersten Componente nach Zeilen, in der zweiten nach Colonnen summirt wird.}

Aus dieser Regel ergiebt sich die Folgerung:

\gesperrt{2. Die Determinante einer zusammengesetzten Substitution ist gleich dem Producte aus den Determinanten der Componenten.}


% ===== scan leaf 175, printed 154 =====
% !PRINTED 154
\weberpage{154}
Um drei Substitutionen derselben Dimension, $C$, $B$, $A$, zusammenzusetzen, muss man die Ausdrücke (10) für $z$ in eine neue lineare Substitution

\weq{12}{(u) = C(z)}

einführen und die Variablen $u$ durch die $x$ ausdrücken:

\weq{13}{(u) = CBA(x).}

Da es offenbar gleichgültig ist, ob man die Ausdrücke (10) in (12) einführt, oder ob man zuerst $z$ nach (9) durch $y$ und dann $y$ nach (4) durch $x$ ausdrückt, so gilt für diese Composition das \gesperrt{associative Gesetz}:

\weq{14}{C(BA) = (CB)A = CBA,}

was sich auch leicht durch Rechnung bestätigen lässt, wenn man nach (11) die Elemente von $C(BA)$ und $(CB)A$ bildet. Man findet für beide den Ausdruck:

\[
\sum^i \sum^h c_i^{(k)}\, b_h^{(i)}\, a_l^{(h)}.
\]

Wir sprechen also den Satz aus:

\par
3.\enspace\gesperrt{Bei der Zusammensetzung der linearen Substitutionen gilt das associative, aber nicht immer das commutative Gesetz.}

Eine Multiplication, bei der alle Multiplicatoren einander gleich sind, also

\[
N = \begin{pmatrix} \nu, & 0, & \ldots & 0 \\ 0, & \nu, & \ldots & 0 \\ \cdot & \cdot & \cdot & \cdot \\ 0, & 0, & \ldots & \nu \end{pmatrix}
\]

soll eine \gesperrt{Aehnlichkeitssubstitution} genannt werden, und zwei Substitutionen, die, wie $A$ und $AN$ oder $A$ und $NA$, durch Zusammensetzung mit einer Aehnlichkeitssubstitution aus einander abgeleitet werden können, heissen \gesperrt{ähnliche Substitutionen}. Aus dem Gesetze der Composition ergeben sich sofort die Sätze:

\par
4.\enspace\gesperrt{Eine Aehnlichkeitssubstitution ist mit jeder Substitution $A$ derselben Dimension vertauschbar; zwei Aehnlichkeitssubstitutionen, mit einander componirt, geben wieder eine Aehnlichkeitssubstitution, und zwei mit einer dritten ähnlichen Substitutionen sind auch unter einander ähnlich.}


% ===== scan leaf 176, printed 155 =====
% !PRINTED 155
\weberpage{155}
Ebenso leicht erhalten wir:

\par
5.\enspace\gesperrt{Durch Zusammensetzung mit der identischen Substitution $J$ wird keine Substitution geändert.}

Die Substitution $J$ wird also bei der Composition als \gesperrt{Einheit} betrachtet, und kann, wo sie mit anderen Substitutionen componirt auftritt, weggelassen werden.

Nun gilt weiter der Satz:

\par
6.\enspace\gesperrt{Zu jeder Substitution $A$ giebt es eine und nur eine inverse Substitution $A^{-1}$, die der Bedingung}

\weq{15}{AA^{-1} = A^{-1}A = J}

\gesperrt{genügt.}

Dieser letzte Satz ergiebt sich aus den Grundformeln der Determinantentheorie, wenn man die Elemente $\alpha_h^{(k)}$ von $A^{-1}$ aus den linearen Gleichungen

\weq{16}{\begin{aligned} \sum^i a_i^{(h)}\alpha_k^{(i)} &{}= 0, & h \gtrless k\\ &{}= 1, & h = k \end{aligned}}

bestimmt, aus denen sich, wenn wir mit $A_h^{(k)}$ die Unterdeterminanten von $|A|$ bezeichnen,

\weq{17}{|A|\,\alpha_h^{(k)} = A_k^{(h)}}

ergiebt, und die das andere System

\weq{18}{\begin{aligned} \sum^i \alpha_i^{(h)} a_k^{(i)} &{}= 0, & h \gtrless k\\ &{}= 1, & h = k \end{aligned}}

zur Folge haben. Die inverse Substitution zu $A$ ist nichts Anderes, als die Auflösung des Gleichungssystems (1) der directen Substitution $A$, so dass aus $(y) = A(x)$ folgt:

\weq{19}{(x) = A^{-1}(y).}

Für die Composition der inversen Substitutionen ergiebt sich aus $ABB^{-1}A^{-1} = J$ der Satz:

\weq{20}{(AB)^{-1} = B^{-1}A^{-1}.}

Sind $A$, $B$, $C$ drei Substitutionen, so ergiebt sich durch Zusammensetzung mit $A^{-1}$ aus jeder der beiden Gleichungen

\[
AB = AC, \quad BA = CA,
\]

dass $B = C$ sein muss. Demnach sind für die Zusammensetzung der Substitutionen die charakteristischen Merkmale für eine Gruppe erfüllt, und es ist also der Inbegriff aller Substitutionen
%CONT>>


% ===== scan leaf 177, printed 156 =====
% !PRINTED 156
\weberpage{156}
%<<CONT
von bestimmter Dimension eine (unendliche) \gesperrt{Gruppe} (\S.\,1). Wenn wir aus dieser Gesammtheit irgend eine Menge herausheben, die so in sich abgeschlossen ist, dass irgend zwei ihrer Elemente durch Composition ein Element derselben Menge ergeben, so bildet diese Menge gleichfalls eine \gesperrt{Gruppe}, die endlich oder unendlich sein kann.

Bedeutet $L$ eine feste Substitution, so kann man aus jeder Substitution $A$ von derselben Dimension eine Substitution

\weq{21}{A' = L^{-1}AL}

ableiten, die die \gesperrt{Transformirte von $A$ durch $L$} heisst.

Setzen wir

\weq{22}{(x) = L(x'), \quad (y) = L(y'),}

so folgt aus (4):

\weq{23}{(y') = A'(x'),}

so dass der Uebergang zu der transformirten Substitution gleichbedeutend ist mit der gleichzeitigen Transformation beider Reihen von Veränderlichen durch $L^{-1}$.

Ist

\[
A' = L^{-1}AL, \quad B' = L^{-1}BL,
\]

so folgt aus den Gesetzen der Composition:

\[
A'B' = L^{-1}ABL,
\]

und daraus also der Satz:

\par
7.\enspace\gesperrt{Durchläuft $A$ die Substitutionen einer Gruppe, so durchläuft bei feststehendem $L$ die Transformirte $L^{-1}AL$ die Substitutionen einer isomorphen Gruppe.}

Von den inversen Substitutionen sind wohl zu unterscheiden die \gesperrt{transponirten Substitutionen}.

Man erhält nämlich aus jeder Substitution $A$ eine bestimmte andere, die die transponirte Substitution zu $A$ heisst, und die wir für den Augenblick mit $A_1$ bezeichnen wollen, wenn man in $A$ die Zeilen zu Colonnen macht, und umgekehrt, wenn man also in (2) die oberen mit den unteren Indices der $a$ vertauscht.

Wenn man in der Summe (11), $\sum^i b_i^{(k)} a_h^{(i)}$, die unteren mit den oberen Indices und gleichzeitig $a$ mit $b$ vertauscht, so erhält man einen Ausdruck, der nach (11) gleich $c_k^{(h)}$ ist.


% ===== scan leaf 178, printed 157 =====
% !PRINTED 157
\weberpage{157}
Diese Bemerkung giebt die Vorschrift, nach der die transponirten Substitutionen zusammengesetzt werden, die sich in dem Satze ausspricht:

\par
8.\enspace\gesperrt{Sind $A_1$, $B_1$ die Transponirten zu $A$, $B$, so ist $A_1 B_1$ die Transponirte zu $BA$. In Zeichen:}

\weq{24}{(BA)_1 = A_1 B_1.}

Wenn wir die Substitutionsformeln (1) mit einem unbestimmten Factor $\eta_k$ multipliciren und dann die Summe in Bezug auf $k$ nehmen, so folgt:

\weq{25}{\sum y_k \eta_k = \sum^i \sum^k x_i\, a_i^{(k)}\, \eta_k,}

oder wenn wir

\weq{26}{\sum^k a_i^{(k)} \eta_k = \xi_i}

setzen,

\weq{27}{\sum y_k \eta_k = \sum x_i \xi_i.}

Wenn wir also die Substitutionen (1) durch (4) darstellen, so sind die Formeln (26) der Ausdruck für die Substitution

\[
(\xi) = A_1(\eta),
\]

und wir können also noch den Satz aussprechen:

\par
9.\enspace\gesperrt{Sind $A$, $A_1$ transponirte Substitutionen von einander von der Dimension $n$, und sind $x$, $y$, $\xi$, $\eta$ vier Systeme von Variablen, die mit einander durch die Substitutionen}

\weq{28}{(y) = A(x), \quad (\xi) = A_1(\eta)}

\gesperrt{zusammenhängen, so besteht die Identität}

\weq{29}{y_1\eta_1 + y_2\eta_2 + \cdots y_n\eta_n = x_1\xi_1 + x_2\xi_2 + \cdots x_n\xi_n.}

Wenn zwei Reihen von Variablen mit

\[
\begin{gathered} x_1, x_2, \ldots x_n \\ \xi_1, \xi_2, \ldots \xi_n \end{gathered}
\]

gleichzeitig durch die Substitutionen (28) in zwei neue Reihen

\[
\begin{gathered} y_1, y_2, \ldots y_n \\ \eta_1, \eta_2, \ldots \eta_n \end{gathered}
\]

transformirt werden, so heissen die beiden Variablenreihen mit einander \gesperrt{contragradient}, und die Formeln (28) stellen zwei mit einander \gesperrt{contragradiente Transformationen} dar.


% ===== scan leaf 179, printed 158 =====
% !PRINTED 158
\weberpage{158}
\pdfbookmark[2]{§ 38. Substitution der Verhältnisse}{bm46}

\begin{center}
\S.\,38.

\gesperrt{Substitution der Verhältnisse.}
\end{center}

Wir haben schon oben gesehen, dass zwei mit einer dritten ähnliche Substitutionen gleicher Dimension unter einander ähnlich sind. Demnach können wir alle mit einander ähnlichen Substitutionen $n^{\mathrm{ter}}$ Dimension in eine Classe vereinigen, und jede Substitution kann in einer und nur in einer solchen Classe untergebracht werden. Die einzelnen Substitutionen einer Classe heissen die \gesperrt{Repräsentanten} der Classe, und jede Classe ist durch irgend einen ihrer Repräsentanten völlig bestimmt.

Ist $A$ ähnlich mit $A'$, $B$ ähnlich mit $B'$, so ist auch $AB$ ähnlich mit $A'B'$.

Bezeichnen wir also mit $\mathfrak{A}, \mathfrak{B}$ die Classen, in die $A, A'$ und $B, B'$ gehören, so gelangt man immer in dieselbe Classe, welchen Repräsentanten aus $\mathfrak{A}$ und aus $\mathfrak{B}$ man auch zusammensetzen mag. Diese Classe, die durch $AB$ oder $A'B'$ repräsentirt wird, nennen wir daher aus $\mathfrak{A}$ und $\mathfrak{B}$ zusammengesetzt und bezeichnen sie mit $\mathfrak{A}\mathfrak{B}$. Bei dieser Zusammensetzung gelten dieselben Regeln, wie bei der Zusammensetzung der Substitutionen selbst, und die Gesammtheit der Classen bildet also auch eine Gruppe.

Diese Auffassung der Substitutionen und ihrer Zusammensetzung ist immer dann zweckmässig, wenn es, wie z.\,B. in der projectiven Geometrie, nur auf die Verhältnisse der Variablen ankommt, und wir bezeichnen daher die Substitutionsclassen als \gesperrt{Substitutionen der Verhältnisse.}

Wir werden aber diese Substitutionen der Verhältnisse ebenso bezeichnen, wie die anderen Substitutionen, nämlich jede Classe durch einen Repräsentanten, wodurch nicht leicht ein Missverständniss entstehen wird. Den Repräsentanten kann man nach sehr verschiedenen Gesichtspunkten auswählen.

Oft empfiehlt es sich, ihn so anzunehmen, dass die Determinante $= 1$ ist. Dadurch ist aber bei einer Substitution $n^{\mathrm{ter}}$ Dimension die multiplicative Substitution $(\mu)$ nur bis auf eine $n^{\mathrm{te}}$ Einheitswurzel bestimmt. Sind die Repräsentanten so gewählt, so bleibt die Eigenschaft bei der Composition erhalten, und diese Repräsentanten der Classen bilden also auch unter sich eine Substitutionsgruppe.


% ===== scan leaf 180, printed 159 =====
% !PRINTED 159
\weberpage{159}
\pdfbookmark[2]{§ 39. Permutationen als lineare Substitutionen}{bm47}

Im Gebiete der binären Substitutionen ist die Substitution $A$:
\[
(y_1, y_2) = \begin{pmatrix} a, & b \\ c, & d \end{pmatrix} (x_1, x_2)
\]
als Substitution der Verhältnisse aufgefasst, gleichbedeutend mit der \gesperrt{linearen gebrochenen Substitution}:
\[
\eta = \frac{a\,\xi + b}{c\,\xi + d},
\]
wenn $\xi = x_1 : x_2$ und $\eta = y_1 : y_2$ gesetzt wird; kommt eine zweite Substitution $A'$:
\[
(z_1, z_2) = \begin{pmatrix} a', & b' \\ c', & d' \end{pmatrix} (y_1, y_2),
\]
oder
\[
\zeta = \frac{a'\eta + b'}{c'\eta + d'}
\]
hinzu, so erhält man die zusammengesetzte Substitution
\[
A'' = A'A,
\]
durch welche $\zeta$ durch $\xi$ ausgedrückt wird:
\[
\zeta = \frac{a''\xi + b''}{c''\xi + d''},
\]
nach den Regeln der Composition in der Form
\[
\begin{pmatrix} a'', & b'' \\ c'', & d'' \end{pmatrix} = \begin{pmatrix} a', & b' \\ c', & d' \end{pmatrix} \begin{pmatrix} a, & b \\ c, & d \end{pmatrix} = \begin{pmatrix} a'a + b'c, & a'b + b'd \\ c'a + d'c, & c'b + d'd \end{pmatrix}.
\]
Ist eine solche Substitution multiplicativ, hat sie also die Form
\[
\begin{pmatrix} a, & 0 \\ 0, & d \end{pmatrix},
\]
so werden wir auch das Verhältniss $a : d$ den \gesperrt{Multiplicator} nennen.

\begin{center}
\S.\,39.

\gesperrt{Permutationen als lineare Substitutionen.}
\end{center}

Die Wichtigkeit der linearen Substitutionen und besonders der daraus gebildeten \gesperrt{endlichen} Gruppen für die Algebra ergiebt sich daraus, dass die Permutationsgruppen von $n$ Elementen als specielle Fälle solcher Substitutionsgruppen aufgefasst werden können.


% ===== scan leaf 181, printed 160 =====
% !PRINTED 160
\weberpage{160}

Bezeichnen wir nämlich mit $x_1, x_2, \dots, x_n$ ein System von $n$ Veränderlichen, und mit $\alpha_1, \alpha_2, \dots \alpha_n$ irgend eine Anordnung der $n$ Ziffern $1, 2, \dots n$, so bestimmen die Gleichungen:
\weq{1}{x_1 = x'_{\alpha_1}, \quad x_2 = x'_{\alpha_2}, \quad \dots \quad x_n = x'_{\alpha_n}}
eine lineare Substitution $n^{\mathrm{ter}}$ Dimension:
\weq{2}{A = \begin{pmatrix} a_1^{(1)}, & \dots & a_n^{(1)} \\ \cdot & \cdot & \cdot \\ a_1^{(n)}, & \dots & a_n^{(n)} \end{pmatrix}, \quad (x) = A(x'),}
bei der in jeder Zeile und in jeder Colonne nur ein Coëfficient von Null verschieden ist, und dieser eine den Werth 1 hat.

Die Determinante $|A|$ der Substitution ist also $= \pm 1$. Setzt man aber nach Ausführung der Substitution für $x'_i$ wieder $x_i$, so ist das Ergebniss nichts Anderes, als die Permutation
\[
\begin{pmatrix} 1, & 2, & \dots & n \\ \alpha_1, & \alpha_2, & \dots & \alpha_n \end{pmatrix}
\]
der Indices von $x$.

Ist $B$ eine zweite ebenso gebildete Substitution
\weq{3}{(x') = B(x''),}
oder ausführlicher:
\weq{4}{x'_1 = x''_{\beta_1}, \quad x'_2 = x''_{\beta_2}, \quad \dots, \quad x'_n = x''_{\beta_n},}
so ergiebt die Zusammensetzung nach den Regeln des \S.\,37:
\weq{5}{x_1 = x''_{\beta_{\alpha_1}}, \quad x_2 = x''_{\beta_{\alpha_2}}, \quad \dots \quad x_n = x''_{\beta_{\alpha_n}},}
was abgekürzt durch
\weq{6}{(x) = AB(x'')}
zu bezeichnen ist.

Nach Bd.\,I, \S.\,148 ist aber (nach der Zusammensetzung der Permutationen):
\weq{7}{\begin{pmatrix} 1, & 2, & \dots & n \\ \alpha_1, & \alpha_2 & \dots & \alpha_n \end{pmatrix} \begin{pmatrix} 1, & 2, & \dots & n \\ \beta_1, & \beta_2 & \dots & \beta_n \end{pmatrix} = \begin{pmatrix} 1, & 2, & \dots & n \\ \beta_{\alpha_1}, & \beta_{\alpha_2} & \dots & \beta_{\alpha_n} \end{pmatrix}.}
Fassen wir also die Substitutionen $A, B$ als Permutationen der Indices auf, so ist die Substitution $AB$ gleichbedeutend mit der zusammengesetzten Permutation $AB$.

Die Permutationsgruppen von $n$ Ziffern sind hiernach nichts Anderes, als ein specieller Fall \gesperrt{endlicher Gruppen} linearer Substitutionen $n^{\mathrm{ter}}$ Dimension.

Die Permutationen der ersten Art entsprechen Substitutionen mit der Determinante $+1$, und die Permutationen der zweiten Art Substitutionen mit der Determinante $-1$.


% ===== scan leaf 182, printed 161 =====
% !PRINTED 161
\weberpage{161}
\pdfbookmark[2]{§ 40. Die Invarianten von endlichen Gruppen linearer Substitutionen}{bm48}

Dies ergiebt sich einfach daraus, dass eine Transposition, z.\,B. $(1, 2)$, der Substitution
\[
\begin{pmatrix} 0, & 1, & 0, & \dots & 0 \\ 1, & 0, & 0, & \dots & 0 \\ 0, & 0, & 1, & \dots & 0 \\ \cdot & \cdot & \cdot & \cdot & \cdot \\ 0, & 0, & 0, & \dots & 1 \end{pmatrix},
\]
deren Determinante $-1$ ist, entspricht, und dass man alle Permutationen der ersten Art aus einer geraden, und alle Permutationen der zweiten Art aus einer ungeraden Anzahl von Transpositionen zusammensetzen kann.

\begin{center}
\S.\,40.

\gesperrt{Die Invarianten von endlichen Gruppen linearer Substitutionen.}
\end{center}

Wir beschränken jetzt unsere Betrachtungen auf \gesperrt{endliche Gruppen} linearer homogener Substitutionen. Das Beispiel der Permutationsgruppen zeigt, dass es solche Gruppen giebt. Andere werden wir später noch kennen lernen.

Wir bezeichnen eine solche Gruppe mit $S$, und ihre Substitutionen mit $A, B, C, \dots$ Werden in irgend einer Function der Variablen $x_1, x_2, \dots, x_n$ die Variablen $(x)$ durch $A(x)$ ersetzt, so sagen wir, \gesperrt{die Substitution $A$ werde auf die Variablen $(x)$ angewandt}. Wir definiren nun zunächst folgenden Begriff:

1. \gesperrt{Eine Form $\Phi(x_1, x_2, \dots, x_n)$ von $n$ Variablen heisst eine Invariante oder invariante Form der Gruppe $S$, wenn sie ungeändert bleibt, wenn auf die Variablen $(x)$ die sämmtlichen Substitutionen $A, B, C, \dots$ der Gruppe $S$ angewandt werden.}

Wie früher (Bd.\,I, \S.\,15) ist hier unter einer Form eine ganze homogene Function der Variablen $x$ zu verstehen.

Wir können diese Definition auch durch die Formeln:
\weq{1}{\Phi[A(x)] = \Phi[B(x)] = \Phi[C(x)] \dots}
ausdrücken.

Dass es für jede endliche Gruppe $S$ invariante Formen in beliebiger Menge giebt, ist leicht einzusehen. Man braucht nur
%CONT>>


% ===== scan leaf 183, printed 162 =====
% !PRINTED 162
\weberpage{162}
%<<CONT
eine beliebige Form $\varphi(x)$ der $n$ Veränderlichen $x$ zu nehmen,
die Functionen:
\weq{2}{\varphi[A(x)],\ \varphi[B(x)],\ \varphi[C(x)],\ \ldots}
für alle Substitutionen der Gruppe zu bilden, und irgend eine
symmetrische Function der Formen (2) für $\Phi$ zu nehmen.

Denn wendet man auf $(x)$ eine Substitution der Gruppe $S$
an, so ändert sich die Gesammtheit der Functionen (2) nicht; es
wird nur ihre Reihenfolge eine andere.

Man kann den Begriff der Invarianten noch allgemeiner
fassen, wie folgt:

2. \gesperrt{Eine Form $\mathit\Psi(x_1,x_2,\ldots,x_n)$ heisst auch dann eine
Invariante der Gruppe $S$, wenn sie constante
Factoren annimmt, wenn auf die Variablen $(x)$
die Substitutionen $A,B,C \ldots$ der Gruppe $S$
angewandt werden.}

Durch Formeln wird diese Eigenschaft so ausgedrückt:
\weq{3}{\mathit\Psi[A(x)]=\alpha\,\mathit\Psi(x),\ \mathit\Psi[B(x)]=\beta\,\mathit\Psi(x),\ \mathit\Psi[C(x)]=\gamma\,\mathit\Psi(x)\ldots}
worin die Co\"efficienten $\alpha,\beta,\gamma\ldots$ von den $x$ unabhängig sind.

Wenn eine Unterscheidung n\"othig ist, wollen wir die in 1.
definirten Formen \gesperrt{absolute Invarianten}, und die in 2. \gesperrt{relative
Invarianten} der Gruppe $S$ nennen.

Aus den Formeln (3) ergiebt sich:
\weq{4}{\mathit\Psi[AB(x)] = \alpha\beta\,\mathit\Psi(x),}
und daraus folgt, dass die Factoren $\alpha,\beta,\gamma,\ldots$ in ihrer Gesammtheit,
bei der Zusammensetzung durch Multiplication, eine Gruppe
bilden müssen.

Zwischen dieser (commutativen) Gruppe und der Gruppe $S$
besteht ein im Allgemeinen mehrstufiger Isomorphismus, da
verschiedene Substitutionen aus $S$ zu demselben Factor $\alpha$ führen
können.

Ist $\mu$ der Grad der Gruppe $S$, so ist der Grad eines jeden
Elementes $A, B \ldots$ von $S$ ein Theiler von $\mu$, und folglich ist
$A^\mu = B^\mu \cdots$ gleich der identischen Substitution.

Hieraus folgt, dass die Factoren $\alpha,\beta,\gamma\ldots$ \gesperrt{$\mu^{\mathrm{te}}$ Einheitswurzeln} sind.

Ist $e$ die kleinste positive, der Bedingung
\weq{5}{\alpha^e = \beta^e = \gamma^e = \cdots = 1}
genügende Zahl, so sind die $\alpha, \beta, \gamma, \ldots$ zugleich $e^{\mathrm{te}}$ Einheits
%CONT>>


% ===== scan leaf 184, printed 163 =====
% !PRINTED 163
\weberpage{163}
%<<CONT
wurzeln, und $e$ soll der \gesperrt{Index der Invariante} $\mathit\Psi(x)$ heissen.
Ist $\varepsilon$ eine primitive $e^{\mathrm{te}}$ Einheitswurzel, so können wir
\[
\alpha=\varepsilon^a,\quad \beta=\varepsilon^b,\quad \gamma=\varepsilon^c \ldots
\]
setzen, und die Exponenten $a,b,c\ldots$ können keinen gemeinschaftlichen
 Theiler mit $e$ haben. Daraus folgt, dass man die
ganzen Zahlen $x,y,z,\ldots$ so bestimmen kann, dass
\[
ax+by+cz+\cdots \equiv 1\ (\mathrm{mod}\ e)
\]
wird (Bd. I, \S.\,118). Da nun wegen der Gruppennatur unter
den Factoren $\alpha,\beta,\gamma,\ldots$ auch die Zahl
\[
\alpha^x\,\beta^y\,\gamma^z \cdots = \varepsilon
\]
vorkommt, so folgt, dass die Gesammtheit der Factoren $\alpha,\beta,\gamma\ldots$
mit den Potenzen von $\varepsilon$:
\[
1,\ \varepsilon,\ \varepsilon^2,\ \ldots\ \varepsilon^{e-1}
\]
zusammenfallen muss.

Suchen wir in $S$ alle Substitutionen $A$, die der Bedingung
\weq{4}{\mathit\Psi[A(x)] = \mathit\Psi(x)}
genügen, zu denen gewiss die identische Substitution gehört, so
erhalten wir eine neue Gruppe $T$, die ein Theiler von $S$ ist.
Bedeutet dann $E$ eine der Bedingung
\weq{5}{\mathit\Psi[E(x)] = \varepsilon\,\mathit\Psi(x)}
genügende Substitution aus $S$, so haben alle Substitutionen $AE$
und $EA$ die gleiche Eigenschaft, und wir erhalten die Zerlegung
von $S$ in die Nebengruppen:
\weq{6}{S = T + TE + TE^2 + \cdots + TE^{e-1};}
der Index des Theilers $T$ ist also $=\varepsilon$. % [sic]
Zugleich ergiebt sich
noch, da jede Substitution $E^{-1}AE$ der Bedingung (4) genügt,
\weq{7}{E^{-1}TE = T,}
woraus hervorgeht, dass $T$ ein Normaltheiler von $S$ ist. Wir
haben damit den Satz bewiesen:

3. \gesperrt{Ist $\mathit\Psi(x)$ eine Invariante der Gruppe $S$ vom
Index $e$, so bilden alle Substitutionen von $S$,
durch die $\mathit\Psi(x)$ ungeändert bleibt, einen Normaltheiler
 $T$ von $S$ vom Index $e$.}

Für die Gruppe $T$ ist $\mathit\Psi(x)$ absolute Invariante. Die Invarianten
 vom Index 1 sind die absoluten Invarianten von $S$.
Die Gruppe $T$ heisst die \gesperrt{zur Invariante $\mathit\Psi(x)$ geh\"orige
Gruppe.}


% ===== scan leaf 185, printed 164 =====
% !PRINTED 164
\weberpage{164}
Hat die Gruppe $S$ vom Grade $\mu$ eine Invariante vom Index
 $\mu$, so wird $T$ die Einheitsgruppe und $S$ ist eine cyklische
Gruppe, die aus den Elementen $1, E, E^2 \ldots E^{\mu-1}$ besteht.

Betrachten wir als Beispiel die symmetrische Permutationsgruppe
 $P$ von $n$ Elementen $x_1,x_2,\ldots,x_n$, die ja nach \S.\,39 unter
den Gruppen linearer Substitutionen enthalten ist, so haben wir
als absolute Invarianten die \gesperrt{symmetrischen Functionen} der
$n$ Variablen $x$. Das Differenzenproduct
\[
\sqrt{\varDelta} = (x_1-x_2)(x_1-x_3)\ldots(x_{n-1}-x_n)
\]
ist eine relative Invariante vom Index 2, zu der die alternirende
Gruppe gehört. Da die Gruppe $P$ ausser der alternirenden
Gruppe keinen Normaltheiler hat und auch nicht cyklisch ist, so
giebt es keine relativen Invarianten von höherem Index als 2.

Die absoluten Invarianten, d. h. die symmetrischen Functionen
, sind hier durch eine endliche Anzahl solcher Formen
rational darstellbar, nämlich durch die symmetrischen Grundfunctionen
, und die Invarianten vom Index 2 sind das Product
von $\sqrt{\varDelta}$ mit absoluten Invarianten. Dass analoge S\"atze auch
im allgemeinen Falle gelten, werden wir in der Folge beweisen.

Wir schliessen hier mit dem Beweise eines allgemeinen
Satzes, der bei allen Anwendungen f\"ur die Bildung der Invarianten
einer Gruppe $S$ von grossem Nutzen ist.

Im \S.\,60 des ersten Bandes haben wir f\"ur irgend eine Form
der $n$ Variablen $F(x_1,x_2,\ldots x_n) = F(x)$ gewisse Formen derselben
 Variablen $C(x_1,x_2,\ldots x_n) = C(x)$ als Covarianten
definirt, die dadurch charakterisirt waren, dass, wenn $F(x)$ durch
irgend eine lineare Substitution in eine neue Form $F'(y)$ transformirt
 wird, die Form $C$ der Bedingung genügt
\[
C'(y) = r^\lambda\,C(x),
\]
wo $C'$ ebenso von den Co\"efficienten von $F'$, wie $C$ von den
Co\"efficienten von $F$ abhängt. Darin bedeutet $r$ die Substitutionsdeterminante
, ist also eine Constante. Die Co\"efficienten der
Form $F$ kommen in $C$ nur homogen vor. Als Beispiel einer Covariante
 f\"uhren wir die \gesperrt{Hesse'sche Determinante} an.

Wenn nun $F(x)$ eine Invariante der Gruppe $S$ ist, und die
$y$ mit den $x$ gleichfalls durch eine Substitution aus $S$ zusammen-
h\"angen, so unterscheiden sich die Co\"efficienten von $F$ nur durch
einen gemeinschaftlichen constanten Factor von den Co\"efficienten
von $F'$, und Gleiches gilt also auch von den beiden Formen $C$
%CONT>>


% ===== scan leaf 186, printed 165 =====
% !PRINTED 165
\weberpage{165}
\pdfbookmark[2]{§ 41. Der Satz von Hilbert}{bm49}
%<<CONT
und $C'$. Daraus schliessen wir, dass auch $C$ zu den Invarianten
der Gruppe $S$ gehört, und sprechen dies als Satz aus:

4. \gesperrt{Bildet man aus einer invarianten Form der
Gruppe $S$ beliebige Covarianten, so erhält man
neue invariante Formen der Gruppe.}

\begin{center}
\S.\ 41.

\gesperrt{Der Satz von Hilbert.}
\end{center}

Der Beweis des Satzes von der Endlichkeit des Invariantensystems
 einer linearen Substitutionsgruppe beruht auf einem sehr
allgemeinen Satze über Formensysteme irgend welcher Art, den
\gesperrt{Hilbert} entdeckt und in mannigfachen Untersuchungen über
die Endlichkeit von Invariantensystemen mit ausgezeichnetem
Erfolge angewandt hat, zu dessen Ableitung wir jetzt übergehen
wollen $^{1}$).

I. \gesperrt{Bedeutet $\mathfrak{S}$ irgend ein System von Formen der
$n$ Ver\"anderlichen $x_1,x_2,\ldots x_n$ in endlicher oder
unendlicher Anzahl, so l\"asst sich aus $\mathfrak{S}$ eine
endliche Anzahl von Formen $F_1,F_2,\ldots,F_\mu$ so
ausw\"ahlen, dass jede Form $F$ von $\mathfrak{S}$ durch einen
Ausdruck}
\weq{1}{F = A_1F_1 + A_2F_2 + \cdots + A_\mu F_\mu}
\gesperrt{dargestellt werden kann, worin $A_1,A_2,\ldots A_\mu$
Formen der Variablen $x_1,x_2,\ldots x_n$ sind.}

Die Definition des Formensystemes $\mathfrak{S}$ muss so vollst\"andig
sein, dass von jeder einzelnen Form der Variablen $x$ entschieden
ist, ob sie zu $\mathfrak{S}$ gehört oder nicht, ist aber übrigens an keine
Voraussetzung gebunden.

Besteht $\mathfrak{S}$ nur aus einer endlichen Zahl von Formen, so ist
unser Satz selbstverst\"andlich, denn man kann ja in diesem Falle
die s\"ammtlichen Formen von $\mathfrak{S}$ f\"ur $F_1,F_2,\ldots F_\mu$ nehmen. Der
Beweis wird sich also nur noch mit dem Falle eines unendlichen
Systemes $\mathfrak{S}$ zu befassen haben.

Zur Vereinfachung des Ausdruckes wollen wir das Formensystem
 $F_1,F_2,\ldots F_\mu$ eine \gesperrt{Basis des Systemes} $\mathfrak{S}$ nennen.
Die Functionen $A_1,A_2,\ldots A_\mu$ müssen so beschaffen sein, dass

\bigskip
\noindent\rule{4cm}{0.4pt}

\noindent $^{1}$) \gesperrt{Hilbert}, \glqq Ueber die Theorie der algebraischen Formen\grqq{}. Mathematische
 Annalen, Bd. 36 (1890).
%CONT>>


% ===== scan leaf 187, printed 166 =====
% !PRINTED 166
\weberpage{166}
%<<CONT
die $\mu$ Producte $A_1F_1, A_2F_2, \ldots, A_\mu F_\mu$ alle von gleichem Grade, dem Grade von $F$ sind. Nat\"urlich aber wird im Allgemeinen nicht gefordert, dass umgekehrt alle Functionen von der Form $A_1F_1+A_2F_2+\cdots+A_\mu F_\mu$ bei beliebigen $A_i$ zu dem Systeme $\mathfrak S$ geh\"oren\footnote{Dies findet nur bei besonderen Systemen $\mathfrak S$ statt, die man Moduln nennt.}.

Der Satz, den wir zu beweisen haben, ist evident, wenn es sich um Functionen einer einzigen Ver\"anderlichen $x_1$ handelt. Denn dann sind alle Formen eines Systemes $\mathfrak S$ Potenzen der Variablen $x_1$ mit nicht negativen Exponenten und mit irgend welchen constanten Co\"efficienten multiplicirt. Identisch verschwindende Functionen brauchen wir nicht zu ber\"ucksichtigen. Nehmen wir dann f\"ur $F_1$ eine dieser Functionen von \gesperrt{m\"oglichst niedrigem Grade}, so kann jede andere Function von $\mathfrak S$ in der Form eines Productes $A_1F_1$ dargestellt werden, worin $A_1$ ebenfalls eine Potenz von $x_1$ mit nicht negativem Exponenten und constantem Co\"efficienten ist.

Um also durch Anwendung der vollst\"andigen Induction zum allgemeinen Beweise zu gelangen, nehmen wir zun\"achst an, der Satz I. sei als richtig erwiesen f\"ur jedes System $\mathfrak S_0$ von Formen von $n$ Variablen $x$ und betrachten zun\"achst ein System $\mathfrak S_r$ von Formen $F$, die ausser den $x$ noch eine $(n+1)^{\text{te}}$ Variable $y$, aber nicht in h\"oherer als der $r^{\text{ten}}$ Potenz enthalten, wenn $r$ irgend eine positive ganze Zahl ist. Jede Function $F$ l\"asst sich dann auf eine einzige Art in die Form setzen:
\weq{2}{F = y^r\varphi + \psi,}
worin die Variable $y$ in $\varphi$ gar nicht mehr und in $\psi$ h\"ochstens bis zur $(r-1)^{\text{ten}}$ Potenz vorkommt.

Wenn $F$ das System $\mathfrak S_r$ durchl\"auft, so durchl\"auft $\varphi$ ein gewisses System $\mathfrak S_0$, das sich nach unserer Voraussetzung durch eine Basis darstellen l\"asst, nehmen wir an in der Form:
\weq{3}{\varphi = a_1\varphi_1 + a_2\varphi_2 + \cdots + a_\mu\varphi_\mu.}
Da nun $\varphi_1,\varphi_2,\ldots,\varphi_\mu$ zu dem Systeme $\mathfrak S_0$ geh\"oren, so giebt es Functionen $F_1, F_2,\ldots,F_\mu$ in $\mathfrak S_r$, so dass
\weq{4}{F_1=y^r\varphi_1+\psi_1,\ F_2=y^r\varphi_2+\psi_2,\ \ldots,\ F_\mu=y^r\varphi_\mu+\psi_\mu,}
und dass $y$ in $\psi_1,\psi_2,\ldots,\psi_\mu$ h\"ochstens bis zur $(r-1)^{\text{ten}}$ Potenz vorkommt.


% ===== scan leaf 188, printed 167 =====
% !PRINTED 167
\weberpage{167}
Aus (2), (3), (4) ergiebt sich aber, wenn wir noch
\[ \Psi = \psi - a_1\psi_1 - a_2\psi_2 \cdots - a_\mu\psi_\mu \]
setzen:
\weq{5}{F = a_1F_1+a_2F_2+\cdots+a_\mu F_\mu+\Psi,}
worin $\Psi$ die Variable $y$ h\"ochstens bis zum Grade $r-1$ enth\"alt. Alle Functionen $\Psi$, die der Bedingung (5) gen\"ugen, bilden ein System $\mathfrak S_{r-1}$.

Wir machen nun weiter die Annahme, das Theorem I. sei bewiesen f\"ur Functionen eines jeden Systemes $\mathfrak S_{r-1}$. Dann k\"onnen wir also $\Psi$ durch eine Basis darstellen in der Form:
\weq{6}{\Psi = b_1\Psi_1+b_2\Psi_2+\cdots+b_\nu\Psi_\nu,}
worin $\Psi_1,\Psi_2,\ldots,\Psi_\nu$ specielle Functionen unseres Systemes $\mathfrak S_{r-1}$ sind. Man kann also \"uber die Co\"efficienten $a_{i,k}$ so verf\"ugen, dass die Functionen
\weq{7}{\begin{aligned}
F_1' &= a_{1,1}F_1+a_{2,1}F_2+\cdots+a_{\mu,1}F_\mu+\Psi_1\\
F_2' &= a_{1,2}F_1+a_{2,2}F_2+\cdots+a_{\mu,2}F_\mu+\Psi_2\\
&\cdots\cdots\cdots\cdots\cdots\cdots\cdots\cdots\cdots\cdots\cdots\cdots\cdots\cdots\\
F_\nu' &= a_{1,\nu}F_1+a_{2,\nu}F_2+\cdots+a_{\mu,\nu}F_\mu+\Psi_\nu
\end{aligned}}
in $\mathfrak S_r$ enthalten sind.

Multiplicirt man mit $b_1,b_2,\ldots b_\nu$ und setzt in der Summe nach (5) und (6):
\[ \begin{aligned}
b_1\Psi_1+b_2\Psi_2+\cdots+b_\nu\Psi_\nu &= \Psi\\
&= F - a_1F_1-a_2F_2-\cdots-a_\mu F_\mu,
\end{aligned} \]
so folgt
\weq{8}{\begin{aligned}
F &= b_1F_1'+b_2F_2'+\cdots+b_\nu F_\nu'\\
&\quad+A_1F_1+A_2F_2+\cdots+A_\mu F_\mu,
\end{aligned}}
worin
\weq{9}{\begin{aligned}
A_1 &= a_1 - a_{1,1}b_1 - \cdots - a_{1,\nu}b_\nu\\
&\cdots\cdots\cdots\cdots\cdots\cdots\cdots\cdots\cdots\cdots\cdots\\
A_\mu &= a_\mu - a_{\mu,1}b_1 - \cdots - a_{\mu,\nu}b_\nu.
\end{aligned}}
Die Formen $F_1',\ldots F_\nu', F_1,\ldots F_\mu$ geh\"oren dem Systeme $\mathfrak S_r$ an und bilden eine Basis dieses Systemes, und das Theorem I. ist also unter der Voraussetzung, dass es f\"ur $\mathfrak S_0$ gilt, f\"ur jedes System $\mathfrak S_r$ bewiesen.

Es sei nun $\mathfrak S$ ein beliebiges System von Formen von $n+1$ Variablen. Wir greifen irgend eine Form $F_0$ vom Grade $r$ aus diesem Systeme heraus, und setzen f\"ur die Variablen, von denen das System abh\"angt,
\weq{10}{y,\ x_1+\lambda_1y,\ \ldots,\ x_n+\lambda_ny,}
%CONT>>


% ===== scan leaf 189, printed 168 =====
% !PRINTED 168
\weberpage{168}
\pdfbookmark[2]{§ 42. Endlichkeit des Invariantensystems einer endlichen linearen Substitutionsgruppe}{bm50}
%<<CONT
worin $\lambda_1,\ldots\lambda_n$ Constanten sind, \"uber die wir so verf\"ugen, dass der Co\"efficient von $y^r$ in $F_0$ nicht verschwindet, d. h. dass $F_0(1,\lambda_1,\ldots\lambda_n)$ von Null verschieden wird. Dann geht das System $\mathfrak S$, wie \"uberhaupt jede Form der Variablen (10), in ein System von Formen der Variablen $y,x_1,\ldots x_n$ \"uber, und umgekehrt kann jede Form, die von diesen Variablen abh\"angt, auch als Form von den linearen Verbindungen (10) dargestellt werden.

Irgend eine Form $F$ des Systemes $\mathfrak S$ wird nun nach Potenzen von $y$ geordnet und dann in Bezug auf $y$ die Division mit $F_0$ ausgef\"uhrt, wobei sich
\weq{11}{F = a_0F_0 + \Phi}
ergeben mag, so dass $a_0$ der Quotient und $\Phi$ der Rest der Division ist. $a_0$ und $\Phi$ sind \gesperrt{ganze Functionen} der Variablen $x, y$, weil der Co\"efficient der h\"ochsten Potenz von $y$ im Divisor $F_0$ constant ist. $\Phi$ \"ubersteigt in Bezug auf $y$ nicht den Grad $r-1$. Durchl\"auft nun $F$ das System $\mathfrak S$, so bildet die Gesammtheit der durch (11) definirten Functionen $\Phi$ ein System $\mathfrak S_{r-1}$, von dem wir die Darstellbarkeit durch eine Basis als schon erwiesen annehmen. Wir k\"onnen also setzen:
\weq{12}{\Phi = A_1\Phi_1 + \cdots + A_\mu\Phi_\mu,}
so dass $\Phi_1,\ldots\Phi_\mu$ dem Systeme $\mathfrak S_{r-1}$ angeh\"oren, d. h. so, dass sich in dem Systeme $\mathfrak S$ die Formen
\weq{13}{F_1 = a_1F_0+\Phi_1,\ \ldots,\ F_\mu = a_\mu F_0+\Phi_\mu}
bestimmen lassen. Setzen wir also
\weq{14}{A_0 = a_0 - a_1A_1 - \cdots - a_\mu A_\mu,}
so folgt aus (11), (12) und (13):
\weq{15}{F = A_0F_0+A_1F_1+\cdots+A_\mu F_\mu,}
wodurch das Theorem I. allgemein bewiesen ist.

\begin{center}
\S. 42.\\
\gesperrt{Endlichkeit des Invariantensystemes einer endlichen linearen Substitutionsgruppe.}
\end{center}

Der im vorigen Paragraphen gegebene Beweis des Satzes I. ist an sich keinerlei Ausnahmen unterworfen. Wir verlieren aber nichts Wesentliches an seiner Allgemeinheit, wenn wir ein- f\"ur allemal identisch verschwindende Formen ausschliessen. Wenn
%CONT>>


% ===== scan leaf 190, printed 169 =====
% !PRINTED 169
\weberpage{169}
%<<CONT
ferner das System $\mathfrak{S}$ Formen 0\textsuperscript{ten} Grades, d.\,h. von Null verschiedene Constanten enthält, so ist, da wir für $F_1$ eine solche Constante nehmen können, unser Satz selbstverständlich, da, wenn $A_1 = F : F_1$ gesetzt wird, $F = A_1F_1$ ist. In dieser Form ist aber der Satz inhaltlos. Wenn wir aber von $\mathfrak{S}$ alle constanten Formen ausschliessen, so bleibt ein System $\mathfrak{S}'$, das keine Formen 0\textsuperscript{ten} Grades mehr enthält, für das unser Satz gleichfalls gilt. Die Basis $F_1, F_2, \ldots F_\mu$ enthält dann gleichfalls keine Formen 0\textsuperscript{ten} Grades, und wir können daher den Satz I. auch so ausdrücken:

\gesperrt{II. Alle Formen des Systemes $\mathfrak{S}$ von positivem Grade lassen sich durch eine Basis $F_1, F_2, \ldots F_\mu$, deren Elemente von positivem Grade sind, in der Form ausdrücken:}
\weq{1}{F = \Phi_1F_1 + \Phi_2F_2 + \cdots + \Phi_\mu F_\mu.}
\gesperrt{Die Grade von $\Phi_1, \Phi_2, \ldots \Phi_\mu$ sind dann niedriger als der Grad von $F$.}

Die wichtigsten Anwendungen findet dieses Theorem bei Untersuchungen über die Möglichkeit, alle Formen eines gewissen Systemes $\mathfrak{S}$ als \gesperrt{ganze rationale Functionen einer endlichen Anzahl} unter ihnen darzustellen. Man nennt ein solches Formensystem ein \gesperrt{endliches} (nicht in dem Sinne, dass es nur aus einer endlichen Zahl von Formen besteht). Es gilt der folgende Satz:

\gesperrt{III. Wenn sich die Co\"efficienten $\Phi_1, \Phi_2, \ldots \Phi_\mu$ in der Darstellung (1) des Theorems II. für jede Form $F$ in $\mathfrak{S}$ so wählen lassen, dass sie, wenn sie nicht constant sind, selbst dem Systeme $\mathfrak{S}$ angehören, so ist das System $\mathfrak{S}$ endlich.}

Dies ergiebt sich unmittelbar daraus, dass die Grade der Formen $\Phi_1, \Phi_2, \ldots \Phi_\mu$ niedriger sind, als der Grad von $F$. Wendet man also die Darstellung (1) auf die nicht constanten unter den Functionen $\Phi$ an, so gelangt man zu Co\"efficienten von noch niedrigerem Grade und muss also schliesslich bei wiederholter Anwendung dieses Verfahrens auf Constanten kommen.

Daraus folgt nun durch eine sehr einfache Schlussweise, die ich einer mündlichen Mittheilung von \gesperrt{Hurwitz} verdanke, die Endlichkeit des Invariantensystemes $\mathfrak{J}$ einer endlichen Gruppe linearer Substitutionen $S$.


% ===== scan leaf 191, printed 170 =====
% !PRINTED 170
\weberpage{170}
Es sei $F_1, F_2, \ldots, F_\mu$ eine nach II. bestimmte Basis des Systemes $\mathfrak{J}$, und $F$ irgend eine andere nicht constante Invariante von $S$. Dann lassen sich die Formen $A_1, A_2, \ldots, A_\mu$ so bebestimmen, dass % [sic]
\weq{2}{F = A_1F_1 + A_2F_2 + \cdots + A_\mu F_\mu}
wird. Wendet man auf diese identische Gleichung sämmtliche Substitutionen der Gruppe $S$ an, so bleiben nach Voraussetzung $F, F_1, F_2, \ldots F_\mu$ ungeändert, während $A_\varkappa$ in $A_\varkappa, A_\varkappa', A_\varkappa'' \ldots$ übergehen mag. Bildet man die Summe der so aus (2) abgeleiteten Gleichungen und setzt, wenn $m$ den Grad der Gruppe $S$ bedeutet,
\weq{3}{m\Phi_\varkappa = A_\varkappa + A_\varkappa' + A_\varkappa'' + \cdots,}
so folgt:
\weq{4}{F = \Phi_1F_1 + \Phi_2F_2 + \cdots + \Phi_\mu F_\mu.}
Die Functionen $\Phi_1, \Phi_2, \ldots \Phi_\mu$ sind aber nach \S.\,40 Invarianten von $S$, und damit ist nach III. die Endlichkeit des Systemes $\mathfrak{J}$ bewiesen.

Derselbe Schluss lässt sich auch auf die relativen Invarianten anwenden, wie folgt:

Wir bezeichnen mit $F(x)$ das ganze System der Functionen, die den Bedingungen \S.\,40, (3):
\[
F[A(x)] = \alpha F(x), \quad F[B(x)] = \beta F(x) \ldots
\]
bei feststehenden Factoren $\alpha, \beta, \ldots$, die, wie wir gesehen haben, Einheitswurzeln sind, genügen.

Nach dem \gesperrt{Hilbert'schen} Satze lässt sich ein speci\"elles System solcher Functionen $F_1, F_2, \ldots, F_m$ derart auswählen, dass man
\weq{5}{F = A_1F_1 + A_2F_2 + \cdots + A_mF_m}
setzen kann, worin $A_1, A_2, \ldots, A_m$ Formen der Variablen $(x)$ sind. Behandelt man diese Formel so wie die Formel (2), indem man die Substitutionen der Gruppe $S$ darauf anwendet und dann die Summe bildet, so erhält man, entsprechend der Formel (4):
\[
F = \Phi_1F_1 + \Phi_2F_2 + \cdots + \Phi_mF_m,
\]
worin die Co\"efficienten $\Phi_1, \Phi_2, \ldots, \Phi_m$ \gesperrt{absolute Invarianten} sind.

Zur Vervollständigung ist noch hinzuzufügen, dass inhomogene Functionen der Variablen nur dann Invarianten sein können, wenn ihre einzelnen homogenen Bestandtheile Invarianten sind.


% ===== scan leaf 192, printed 171 =====
% !PRINTED 171
\weberpage{171}
\pdfbookmark[2]{§ 43. Das Formenproblem}{bm51}
Endlich können wir auch noch nach gebrochenen Invarianten fragen. Ist
\[
\frac{F(x)}{F_1(x)}
\]
eine solche gebrochene Invariante, so nehmen wir zunächst an, die beiden ganzen rationalen Functionen $F(x), F_1(x)$ von den $n$ Variablen $x$ seien von gemeinschaftlichen Theilern befreit (Bd.\,I, \S.\,51).

Beschränken wir uns fürs Erste auf absolute Invarianten, und ist demnach
\[
\frac{F(x)}{F_1(x)} = \frac{F[A(x)]}{F_1[A(x)]},
\]
so haben, da die $x$ hier als unabhängige Variable angesehen werden, weder rechts noch links Zähler und Nenner einen gemeinschaftlichen Theiler, und es folgt:
\[
F[A(x)] = \alpha F(x), \quad F_1[A(x)] = \alpha F_1(x),
\]
worin $\alpha$ ein constanter Factor ist, der, wie wir früher gesehen haben, eine Einheitswurzel ist. Zähler und Nenner einer gebrochenen Invariante müssen daher selbst, wenn auch nur relative, Invarianten sein.

Wenn wir aber nicht gerade die einfachste Darstellung suchen, so können wir absolute gebrochene Invarianten auch als Quotienten von absoluten ganzen Invarianten darstellen. Wir brauchen den Bruch nur durch eine geeignete Potenz des Nenners zu erweitern, also wenn die $\alpha$ $e^{\text{te}}$ Einheitswurzeln sind,
\[
\frac{F(x)}{F_1(x)} = \frac{F(x) F_1(x)^{e-1}}{[F_1(x)]^e}
\]
zu setzen. Hiernach können wir auch alle relativen Invarianten mit einem bestimmten Factorensysteme $\alpha, \beta, \ldots$ darstellen als Product von einer von ihnen mit absoluten Invarianten, die aber gebrochen sein können.

\begin{center}
\S.\,43.

\gesperrt{Das Formenproblem.}
\end{center}

Im vorigen Paragraphen ist nachgewiesen, dass es zu einer endlichen Gruppe linearer Substitutionen eine endliche Anzahl unabhängiger Invarianten giebt. Ist $n$ die Dimension der linearen
%CONT>>


% ===== scan leaf 193, printed 172 =====
% !PRINTED 172
\weberpage{172}
%<<CONT
Substitution, so sind diese Formen homogene Functionen von
$n$ Variablen, und es können also nicht mehr als $n$ von einander
unabhängige existiren. Damit ist nicht gesagt, dass sich alle
diese Formen \gesperrt{rational} durch $n$ unter ihnen ausdrücken lassen,
aber zwischen $n+1$ Invarianten muss immer eine rationale
Gleichung bestehen, die sich durch Elimination der $n$ Variablen
ergiebt.

Es ist aber noch die umgekehrte Frage zu untersuchen, ob
es wirklich für eine lineare Substitutionsgruppe von der Dimension
$n$ immer $n$ unabhängige Invarianten giebt.

Wenn wir für die Variablen feste Werthe setzen, so erhalten
dadurch die sämmtlichen Invarianten der Gruppe gleichfalls
bestimmte Werthe. Die Frage, die wir noch zu beantworten
haben, ist nun die, ob zu einem Werthsysteme der Invarianten
bestimmte Werthsysteme der Variablen, und zwar in endlicher
Anzahl, existiren. Wenn dies bewiesen ist, so können wir die
Variablen als (mehrwerthige) algebraische Functionen der Invarianten
auffassen, und die Anzahl der von einander unabhängigen
Invarianten kann nicht kleiner sein, als die Anzahl
der Variablen, weil sonst ein Theil der Variablen, wenn die
Werthe der Invarianten gegeben sind, noch willkürlich bleiben
würde.

Wenn $(x)$ ein einem bestimmten Werthsysteme der Invarianten
entsprechendes Werthsystem der Variablen ist, und $A$ eine Substitution
der Gruppe $S$, so entspricht nach der Natur der Invarianten
das System $A(x)$ demselben Werthsysteme der Invarianten,
und unser Problem hat also mindestens so viele
Lösungen, als der Grad der Gruppe beträgt. Dass dies aber die
genaue Anzahl der Lösungen ist, geht aus den folgenden Betrachtungen
hervor.

In besonderen Fällen, d. h. für besondere Werthe der Invarianten,
können von diesen Werthsystemen der Variablen
mehrere zusammenfallen. Dies kann aber nur für solche Werthsysteme
der $(x)$ geschehen, für die eine Relation von der Form
$(x) = A(x)$ besteht, denn aus $A(x) = B(x)$ würde $(x) = A^{-1}B(x)$
folgen, was in der Form $(x) = A(x)$ enthalten ist.

Wir nehmen eine lineare homogene Function $\Theta$ der Variablen
$x_1, x_2, \dots x_n$ an, die wir so bezeichnen:
\weq{1}{\Theta = \Theta(x) = c_1 x_1 + c_2 x_2 + \cdots + c_n x_n.}


% ===== scan leaf 194, printed 173 =====
% !PRINTED 173
\weberpage{173}
Bedeutet $A$ eine Substitution der endlichen Gruppe $S$, so
können wir aus $\Theta(x)$ eine neue Function:
\weq{2}{\Theta[A(x)] = c_1' x_1 + c_2' x_2 + \cdots + c_n x_n} % [sic]
ableiten, deren Co\"efficienten $c_i'$ nach \S.\,37, 9. mit den ursprünglichen
Co\"efficienten $c_i$ durch die zu $A$ \gesperrt{transponirte Substitution}
\[ (c') = A_1(c) \]
zusammenhängen.

Wenn nun $A, B, C, \dots$ die sämmtlichen Substitutionen der
Gruppe $S$ sind, so kann man in gleicher Weise die Functionen
\weq{3}{\Theta[A(x)],\ \Theta[B(x)],\ \Theta[C(x)],\ \dots}
bilden, die wir auch kürzer durch
\weq{4}{\Theta,\ \Theta_1,\ \Theta_2,\ \dots,\ \Theta_{\mu-1}}
bezeichnen, wenn $\mu$ der Grad der Gruppe $S$ ist.

Nun kann man über die Co\"efficienten $c_i$, die bis jetzt noch
ganz willkürlich sind, so verfügen, dass keine zwei der Functionen
(4) mit einander identisch werden (Bd.\ I, \S.\,143, 1.).
Aus (3) aber ergiebt sich:

\gesperrt{1. Wenn man in den $\mu$ Functionen (4) gleichzeitig
irgend eine Substitution aus $S$ anwendet, so
ändert sich die Gesammtheit dieser Functionen
nicht, sondern sie erleiden nur eine Permutation.}

Daraus ergiebt sich ferner:

\gesperrt{2. Jede symmetrische Function der Grössen (4) ist
eine absolute Invariante der Gruppe $S$.}

Wir bemerken noch, dass man an Stelle der Function $\Theta$
irgend eine andere auch nicht lineare, selbst eine gebrochene
oder inhomogene Function setzen könnte, wenn nur die $\mu$ Functionen
(4) von einander verschieden sind.

Bilden wir nun das Product:
\weq{5}{\begin{aligned}
\Phi(t) &= (t-\Theta)(t-\Theta_1)\dots(t-\Theta_{\mu-1}) \\
&= t^\mu + A_1 t^{\mu-1} + A_2 t^{\mu-2} + \cdots + A_\mu,
\end{aligned}}
welches eine ganze rationale Function $\mu^{\text{ten}}$ Grades von $t$ ist, so
sind die Co\"efficienten $A_1, A_2, \dots, A_\mu$ nach 2. Invarianten der
Gruppe $S$, und die Grössen $\Theta, \Theta_1, \dots, \Theta_{\mu-1}$ sind die Wurzeln
der Gleichung
\weq{6}{\Phi(t) = 0.}


% ===== scan leaf 195, printed 174 =====
% !PRINTED 174
\weberpage{174}
Wir betrachten nun als Rationalitätsbereich $\Omega$ den Körper,
der aus den absoluten Invarianten der Gruppe $S$ und allen
Zahlen\footnote{Man kann sich auch auf einen besonderen Zahlkörper beschränken, wenn nur darin die Substitutionsco\"efficienten der Gruppe $S$ enthalten sind.} besteht. Diesem Rationalitätsbereich gehören die Co\"efficienten
von $\Phi(t)$ an, und wir beweisen zunächst den Satz:

\gesperrt{3. Die Function $\Phi(t)$ ist in $\Omega$ irreducibel.}

Ist nämlich $\Psi(t)$ irgend eine Function in $\Omega$, die für $t = \Theta$
verschwindet, so können wir in der Gleichung $\Psi(\Theta) = 0$, da
die $x_1, x_2, \dots, x_n$ unabhängige Variable sind, das System $(x)$
dieser Variablen durch $A(x)$ ersetzen, wenn $A$ irgend eine Substitution
aus $S$ ist. Dadurch kann $\Theta$ in jede der Functionen
$\Theta_1, \Theta_2, \dots, \Theta_{\mu-1}$ übergeführt werden, während die Co\"efficienten
von $\Psi$ ungeändert bleiben, und folglich ist $\Psi(\Theta_1) = 0$,
$\Psi(\Theta_2) = 0, \dots \Psi(\Theta_{\mu-1}) = 0$. Es muss also $\Psi(t)$ durch $\Phi(t)$
theilbar sein, wodurch die Irreducibilität erwiesen ist.

Es bedeute nun $\omega$ irgend eine Function der Variablen $x$ und
\[ \omega,\ \omega_1,\ \omega_2,\ \dots,\ \omega_{\mu-1} \]
mögen die Functionen sein, die aus $\omega$ durch Anwendung der
Substitutionen von $S$ entstehen, von denen nun nicht vorausgesetzt
zu werden braucht, dass sie alle von einander verschieden
sind. Ist $t$ eine Variable, so ist
\weq{7}{\Phi(t)\left(\frac{\omega}{t-\Theta} + \frac{\omega_1}{t-\Theta_1} + \cdots + \frac{\omega_{\mu-1}}{t-\Theta_{\mu-1}}\right) = \Psi(t)}
eine ganze rationale Function $(\mu-1)^{\text{ten}}$ Grades von $t$, und zugleich
ist es eine Invariante von $S$, also eine Function in $\Omega$.
Setzen wir darin $t = \Theta$, so folgt durch einen schon früher oft
angewandten Schluss (Bd.\ I, \S.\,143, 155):
\weq{8}{\omega = \frac{\Psi(\Theta)}{\Phi'(\Theta)},}
worin der Satz enthalten ist:

\gesperrt{4. Jede rationale Function der Variablen $(x)$ kann
rational durch $\Theta$ ausgedrückt werden, gehört
also dem Körper $\Omega(\Theta)$ an.}

Aus diesem Satze können wir einen zweiten Beweis dafür
ableiten, dass alle Invarianten der Gruppe $S$ rational durch eine
endliche Anzahl von ihnen, nämlich die Co\"efficienten der Function
$\Phi(t)$, darstellbar sind. Denn wenn wir eine absolute In-
%CONT>>


% ===== scan leaf 196, printed 175 =====
% !PRINTED 175
\weberpage{175}
%<<CONT
variante $J$ nach dem Satze 4. als rationale Function von $\Theta$
darstellen, so kann sich diese nicht ändern, wenn eine der Substitutionen
$(\Theta, \Theta_1), (\Theta, \Theta_2), \dots$ ausgeführt wird, und diese Function
ist also rational durch die Co\"efficienten von $\Phi(t)$ ausdrückbar.
Ob diese Darstellung freilich durch ganze Functionen
möglich ist, würde bei diesem Beweise unentschieden bleiben.

Unter den Functionen $\omega$ der Variablen $x$ sind auch die
Variablen $x$ selbst enthalten, und die Frage, von der wir ausgegangen
sind, ob die Variablen $x$ als algebraische Functionen
der Invarianten angesehen werden können, ist damit bejahend
entschieden.

Die Aufgabe, die Variablen $x$ als algebraische Functionen
der Invarianten der Gruppe $S$ darzustellen, also die Bestimmung
des Körpers $\Omega(\Theta)$ heisst nach \gesperrt{F. Klein} das \gesperrt{Formenproblem
der Gruppe $S$}\footnote{Vorlesungen über das Ikosaëder, S.\,123. (Leipzig 1884.)}. Ist die Anzahl der Variablen $n$, so nennen
wir das Formenproblem von der $n^{\text{ten}}$ Dimension.

Der Körper $\Omega(\Theta)$ ist ein durch die Gruppe $S$ völlig bestimmter
algebraischer Körper über $\Omega$. Er ist ein Normalkörper,
denn nach 4. sind die conjugirten Grössen $\Theta, \Theta_1, \Theta_2, \dots \Theta_{\mu-1}$
alle im Körper $\Omega(\Theta)$ selbst enthalten. Die Gleichung $\Phi(t) = 0$
ist eine Normalgleichung und ist die \gesperrt{Galois'sche Resolvente
des Formenproblems} (Bd.\ I, \S.\,145).

\gesperrt{5. Die Galois'sche Gruppe des Formenproblems,
d. h. die Galois'sche Gruppe der Gleichung
$\Phi(t) = 0$, ist mit der Gruppe $S$ isomorph.}

Um dies nachzuweisen, bezeichnen wir mit $A, B$ zwei Substitutionen
aus $S$ und mit $AB = C$ die daraus zusammengesetzte
Substitution. Ist nun
\[ \Theta_1 = \Theta[A(x)],\ \Theta_2 = \Theta[B(x)],\ \Theta_3 = \Theta[C(x)], \]
so ist die Substitution
\[ (\Theta, \Theta_2) = (\Theta_1, \Theta_3), \]
und folglich
\[ (\Theta, \Theta_1)(\Theta, \Theta_2) = (\Theta, \Theta_1)(\Theta_1, \Theta_3) = (\Theta, \Theta_3), \]
d. h. die Gruppe der Substitutionen $(\Theta, \Theta_1), (\Theta, \Theta_2), \dots$ ist mit
der Gruppe der $A, B, \dots$ isomorph.

Nehmen wir statt der Function $\Theta$ eine Function $\eta$, die nicht
lauter verschiedene Werthe hat, sondern die Substitutionen eines
%CONT>>


% ===== scan leaf 197, printed 176 =====
% !PRINTED 176
\weberpage{176}
\pdfbookmark[2]{§ 44. Klein's Erweiterung des algebraischen Grundproblems}{bm52}
%<<CONT
Theilers $S'$ von $S$ vom Index $j$, aber keine anderen gestattet, so
genügt $\eta$ einer Gleichung $j^{\mathrm{ten}}$ Grades, die als eine Resolvente
des Formenproblems zu betrachten ist. Jede Function, die die
Permutationen von $S'$ gleichfalls gestattet, ist dann eine rationale
Function von $\eta$ und von den Invarianten der Gruppe $S$. Die
Resolvente der $\eta$ ist eine Partial- oder Totalresolvente, je
nachdem die Gruppe $S'$ mit den zu ihr conjugirten Theilern
von $S$ einen gemeinschaftlichen Theiler hat oder relativ prim
ist (Bd. I, \S.\,156).

\begin{center}
\S.\ 44.

\gesperrt{Klein's Erweiterung des algebraischen Grundproblems.}
\end{center}

Die Betrachtungen, die im vorigen Paragraphen durchgeführt
sind, bilden eine directe Verallgemeinerung der \gesperrt{Galois}'schen
Theorie für eine allgemeine Gleichung $n^{\mathrm{ten}}$ Grades; und diese
Theorie ist als Specialfall in der Theorie der linearen Substitutionsgruppen enthalten.

Es sind nämlich nach \S.\ 39 die symmetrischen Functionen
von $n$ unabhängigen Variablen $x_1, x_2, \ldots, x_n$ die Invarianten der
Gruppe $S$, die aus den Permutationen dieser $n$ Variablen besteht,
und die Gleichung $\Phi(t) = 0$ ist also für diesen Fall nach Bd. I,
\S.\ 145 die \gesperrt{Galois}'sche Resolvente der Gleichung $n^{\mathrm{ten}}$ Grades,
deren Wurzeln die Grössen $x_i$ sind. Wir können also die allgemeine Aufgabe der Algebra, eine Gleichung $n^{\mathrm{ten}}$ Grades aufzulösen, als ein Formenproblem einer linearen Substitutionsgruppe
$n^{\mathrm{ter}}$ Dimension auffassen. Nun giebt es aber specielle Gleichungen, die durch Formenprobleme von niedrigerer Dimension
gelöst werden können, so insbesondere die \gesperrt{reinen Gleichungen}
die durch ein Formenproblem der ersten Dimension lösbar sind.

Lineare homogene Substitutionen von einer Dimension sind
nämlich nur von der Form
\weq{1}{x' = \alpha x,}
und wenn diese eine endliche Gruppe vom Grade $\mu$ bilden sollen,
so müssen die Co\"efficienten $\alpha$ Einheitswurzeln vom Grade $\mu$ sein.
Lassen wir umgekehrt $\alpha$ in (1) sämmtliche $\mu^{\mathrm{te}}$ Einheitswurzeln
durchlaufen, so haben wir eine Gruppe vom Grade $\mu$. Diese
Gruppe hat eine absolute Invariante $x^{\mu}$, und wenn diese gegeben
ist, so erhält man $x$ als $\mu^{\mathrm{te}}$ Wurzel daraus.


% ===== scan leaf 198, printed 177 =====
% !PRINTED 177
\weberpage{177}
\pdfbookmark[2]{§ 45. Einfluss relativer Invarianten}{bm53}
Die Auflösung der allgemeinen Gleichungen $2^{\mathrm{ten}}$, $3^{\mathrm{ten}}$ und
$4^{\mathrm{ten}}$ Grades sind also auf Formenprobleme von nur einer Dimension zurückführbar. Wir werden in einem späteren Abschnitte
sehen, dass die allgemeine Gleichung $5^{\mathrm{ten}}$ Grades auf ein binäres
Formenproblem zurückführbar ist, und man kann sich nun als
eine unmittelbare Erweiterung der Aufgabe, die Lösung einer
Gleichung auf reine Gleichungen zurückzuführen, die Frage
stellen: welches ist die geringste Dimensionenzahl eines Formenproblems, durch das sich eine gegebene Gleichung lösen lässt?
Die Aufgabe würde dann so formulirt werden müssen:

\gesperrt{Es sollen aus den Wurzeln einer gegebenen Gleichung rationale Functionen in möglichst kleiner Zahl so gebildet werden, dass sie in homogene lineare Functionen ihrer selbst übergehen, wenn die Wurzeln den Permutationen der Galois'schen Gruppe der gegebenen Gleichung unterworfen werden.}

Auf diese Weise hat F. \gesperrt{Klein} die Aufgabe der algebraischen
Auflösung einer Gleichung erweitert. Er hat, um die Beantwortung der Frage anzubahnen, für die allgemeine Gleichung
$6^{\mathrm{ten}}$ und $7^{\mathrm{ten}}$ Grades bewiesen, dass sie auf quaternäre Formenprobleme zurückgeführt werden können.

Die allgemeine Gleichung $n^{\mathrm{ten}}$ Grades ist, wie wir gesehen
haben, unmittelbar einem Formenproblem von $n$ Dimensionen
äquivalent. Die Frage aber, ob die allgemeine Gleichung $n^{\mathrm{ten}}$
Grades, wenn $n$ grösser als 7 ist, einem Formenprobleme von
weniger als $n$ Dimensionen entspricht, ist noch nicht beantwortet; sie muss wahrscheinlich verneint werden $^{1})$.

\begingroup
\renewcommand{\thefootnote}{1)}
\footnotetext[1]{F. \gesperrt{Klein}, Zur Theorie der allgemeinen Gleichungen $6^{\mathrm{ten}}$ und
$7^{\mathrm{ten}}$ Grades; Mathematische Annalen, Bd. XXVIII, S. 18. Vergl. auch \glqq The
Evanston Colloquium, Lectures on Mathematics by \gesperrt{Felix Klein}\grqq{}, Lecture IX, London und New York, Macmillan and Co. (1894).}
\endgroup

\begin{center}
\S.\ 45.

\gesperrt{Einfluss relativer Invarianten.}
\end{center}

Bei der Definition des Formenproblems im \S.\ 43 haben wir
nur die absoluten Invarianten der Gruppe $S$ benutzt. Nun giebt
es, wie wir gesehen haben, auch Fälle, in denen ausser den ab-
%CONT>>


% ===== scan leaf 199, printed 178 =====
% !PRINTED 178
\weberpage{178}
%<<CONT
soluten auch relative Invarianten existiren, und diese können zur
Vereinfachung des Formenproblems benutzt werden.

Ist $r$ eine solche relative Invariante, so wird eine gewisse
Potenz von $r$, deren Grad $e$ ein Theiler des Grades $\mu$ von $S$
ist, eine absolute Invariante, und $r$ wird also durch eine \gesperrt{reine
Gleichung} in $\Omega$ bestimmt.

Alle Substitutionen von $S$, durch die $r$ ungeändert bleibt,
bilden für sich eine Gruppe $T_r$, und die Gruppe $T_r$ ist, wie wir
in \S.\,40 gesehen haben, ein Normaltheiler von $S$.

Ferner aber sehen wir, dass jede Function der Variablen $x$,
die durch die Substitutionen der Gruppe $T_r$ ungeändert bleibt,
rational durch $r$ ausgedrückt werden kann, d.\,h. in dem durch
Adjunction von $r$ aus $\Omega$ abgeleiteten Körper $\Omega_r$ enthalten ist.

Nach \S.\,40 nämlich giebt es eine Substitution $E$ in $S$, und
eine primitive $e^{\text{te}}$ Einheitswurzel $\varepsilon$, so dass $r$ durch $E$ in $\varepsilon r$
übergeht, und alle Werthe, die $r$ annehmen kann:
\[
r, \varepsilon r, \varepsilon^2 r, \ldots \varepsilon^{e-1} r,
\]
erhält man durch Wiederholung der Substitution $E$. Bedeutet
ferner $\tau$ eine Function der $x$, die die Substitutionen der Gruppe
$T_r$ gestattet, und durch $E$ und seine Potenzen in
\[
\tau, \tau_1, \tau_2, \ldots \tau_{e-1}
\]
übergeht, so gestatten alle diese Functionen gleichfalls die
Substitutionen von $T_r$, weil $T_r$ ein Normaltheiler von $S$ ist. Die
Function
\weq{1}{(\varepsilon^{-h}, \tau) = \tau + \varepsilon^{-h}\tau_1 + \cdots + \varepsilon^{-h(e-1)}\tau_{e-1}, \qquad (h = 0, 1, 2, \ldots e-1),}
die nach Analogie der \gesperrt{Lagrange}'schen Resolventen (Bd.\,I,
\S.\,164) gebildet ist, nimmt durch Anwendung der Substitution $E$
den Factor $\varepsilon^h$ an, und wenn wir also
\weq{2}{(\varepsilon^{-h}, \tau) = r^h \Psi_h
}
setzen, so ist $\Psi_h$ eine absolute Invariante, also im Körper $\Omega$
enthalten.

Aus (1) und (2) ergiebt sich aber
\weq{3}{e\tau = \Psi + r\Psi_1 + r^2\Psi_2 + \cdots + r^{e-1}\Psi_{e-1},}
wodurch die Behauptung erwiesen ist.

Setzen wir
\weq{4}{\mu = e\nu,}
%CONT>>


% ===== scan leaf 200, printed 179 =====
% !PRINTED 179
\weberpage{179}
\pdfbookmark[2]{§ 46. Der erweiterte Invariantenbegriff}{bm54}
%<<CONT
so ist $\nu$ der Grad der Gruppe $T_r$, und wenn die Function $\Theta$, die
wir in \S.\ 43 zur Lösung des Formenproblems angewandt haben,
durch die Substitutionen von $T_r$ in
\[ \Theta,\ \Theta_1,\ \ldots \Theta_{\nu-1} \]
übergeht, so ist
\[ \Phi_r(t) = (t - \Theta)(t - \Theta_1) \ldots (t - \Theta_{\nu-1}) \]
eine Function von $t$ in dem erweiterten Rationalitätsbereiche
$\Omega_r$; der Körper $\Omega(\Theta)$ ist identisch mit $\Omega_r(\Theta)$. Er ist also ein algebraischer Körper $\mu^{\mathrm{ten}}$ Grades über $\Omega$ und $\nu^{\mathrm{ten}}$ Grades über $\Omega_r$.
Das Formenproblem $\mu^{\mathrm{ten}}$ Grades wird demnach durch Adjunction
eines Radicals auf ein Formenproblem $\nu^{\mathrm{ten}}$ Grades zurückgeführt.

\begin{center}
\S.\ 46.

\gesperrt{Der erweiterte Invariantenbegriff.}
\end{center}

Die relativen Invarianten sind im Grunde ein specieller Fall
eines allgemeineren Begriffes, den wir in ähnlicher Weise, wie
die relativen Invarianten, zur Reduction des Formenproblems
anwenden können.

Wir suchen nach Systemen von homogenen Formen gleichen
Grades der Variablen $x_1, x_2 \ldots x_n$:
\weq{1}{X_1, X_2, \ldots X_m,}
von der Beschaffenheit, dass durch die Anwendungen der Substitutionen der Gruppe $S$ die Functionen $X_1, X_2, \ldots X_m$ ein
System $\Sigma$ von \gesperrt{linearen Substitutionen} erleiden; die Substitutionen $\Sigma$ bilden dann gleichfalls eine Gruppe, und zwar
eine Gruppe, die mit $S$ ein- oder mehrstufig isomorph ist. Ist
$m = 1$, so erhalten wir, wie man sieht, den Begriff der relativen
Invarianten. Wir suchen nun alle Substitutionen von $S$, die
jede einzelne der Functionen (1) ungeändert lassen. Diese Substitutionen bilden eine Gruppe, die wir mit $T$ bezeichnen, und
von der wir nachweisen wollen, dass sie ein Normaltheiler von
$S$ ist. Bezeichnen wir mit $s$ eine Substitution aus $S$, durch die
die $X$ die Substitution $\sigma$ erleiden, so erleiden die $X$ durch $s^{-1}$
die Substitution $\sigma^{-1}$. Ist dann ferner $\tau$ eine Substitution aus $T$,
so bleiben durch $\tau$ die $X$ ungeändert. Demnach erleiden durch
$s^{-1}\tau s$ die $X$ die Substitution $\sigma^{-1}\sigma = 1$, d.\,h. sie bleiben un-
%CONT>>


% ===== scan leaf 201, printed 180 =====
% !PRINTED 180
\weberpage{180}

%% [running head, not body text -- commented out 2026-08-07] Page furniture (running
%% head, folio, signature) is excluded from this edition by standing policy.
%% \noindent 180 \hfill \textit{Sechster Abschnitt.} \hfill \S.\,46.

%<<CONT
geändert. $s^{-1}\tau s$ ist also auch in $T$ enthalten, und $T$ ist folglich ein Normaltheiler von $S$.

Wir bezeichnen mit $\mu, \nu$ die Grade von $S$ und $T$ und setzen
\[
\mu = e\nu;
\]
dann ist $e$ der Index von $T$ und zugleich der Grad der Gruppe $\Sigma$.

Jede absolute Invariante der Gruppe $T$, d.\,h. jede Function von $x$, die die Substitutionen von $T$ gestattet, kann rational in $\Omega$ durch die $X$ ausgedrückt werden.

Dies folgt durch das schon oft angewandte Schlussverfahren: wenn wir mit $\varrho$ eine Function der $X$ bezeichnen, die $e$ verschiedene Werthe $\varrho, \varrho_1, \ldots \varrho_{e-1}$ annimmt, und
\weq{2}{\Phi(t) = (t-\varrho)(t-\varrho_1)\ldots(t-\varrho_{e-1})}
setzen, so ist $\Phi(t)$ eine rationale Function der Variablen $t$ in $\Omega$.

Eine absolute Invariante $r$ von $T$ nimmt höchstens $e$ verschiedene Werthe an, die den Werthen $\varrho, \varrho_1, \ldots \varrho_{e-1}$ entsprechen, nämlich $r, r_1, \ldots r_{e-1}$, wobei unter Umständen auch gleiche Werthe vorkommen können. Die Function von $t$:
\[
\frac{\Phi(t)\,r}{t-\varrho} + \frac{\Phi(t)\,r_1}{t-\varrho} + \cdots + \frac{\Phi(t)\,r_{e-1}}{t-\varrho_{e-1}} = \Psi(t) % [sic] second denominator printed t-\varrho, not t-\varrho_1
\]
ist dann in $\Omega$ enthalten, und für $t=\varrho$ ergiebt sich:
\weq{3}{r = \frac{\Psi(\varrho)}{\Phi'(\varrho)},}
wodurch, da $\Phi'(\varrho)$ von Null verschieden ist, $r$ rational durch $\varrho$ ausgedrückt ist. Es ist also $r$ im Körper $\Omega(X_1, X_2, \ldots X_m)$ enthalten.

Die Invarianten der Gruppe $T$ sind zugleich Invarianten der Gruppe $S$; durch die Lösung des Formenproblems für die Gruppe $\Sigma$ sind dann die $X_1, \ldots X_m$ bekannt, also auch die Invarianten der Gruppe $T$, und das Formenproblem der Gruppe $S$ ist also zurückgeführt auf die successive Lösung der beiden Formenprobleme der Gruppen $\Sigma$ und $T$, die von niedrigerem Grade als das Formenproblem für $S$ sind. Ein Formenproblem, was in dieser Weise in zwei Formenprobleme niedrigeren Grades zerlegbar ist, können wir ein \gesperrt{imprimitives Formenproblem} nennen. Da es für eine solche Reduction nöthig ist, dass $T$ ein von $S$ und von der Einheit verschiedener Normaltheiler von $S$ sei, so ist, wenn $S$ eine einfache Gruppe ist, das entsprechende Formenproblem stets primitiv.


% ===== scan leaf 202, printed 181 =====
% !PRINTED 181
\weberpage{181}
\pdfbookmark[2]{§ 47. Normalformen}{bm55}

%% [running head, not body text -- commented out 2026-08-07] Page furniture (running
%% head, folio, signature) is excluded from this edition by standing policy.
%% \noindent \S.\,47. \hfill \textit{Normalformen.} \hfill 181

Wir können aber auch umgekehrt schliessen, dass, wenn $S$ einen Normaltheiler vom Index $e$ besitzt, das Formenproblem imprimitiv ist. Wir brauchen nämlich nur, um ein System der Functionen $X_1, \ldots X_m$ zu erhalten, eine Invariante $\varrho$ der Gruppe $T$ zu nehmen, die $e$ verschiedene Werthe in $S$ enthält, und können für $X_1, \ldots X_m$ geradezu die Werthe $\varrho, \varrho_1, \ldots \varrho_{e-1}$ nehmen. Die Gruppe $\Sigma$ ist dann die durch $S$ unter den $\varrho$ hervorgerufene Permutationsgruppe.

Im Sinne des \S.\,44 wird es aber immer darauf ankommen, $m$ so klein als möglich zu machen, und eine Reduction des Problems in diesem Sinne wird nur dann erzielt sein, wenn $m < n$ ist.

\begin{center}
\S.\,47.

\gesperrt{Normalformen.}
\end{center}

Noch eine allgemeine Betrachtung müssen wir anstellen, ehe wir zu speciellen Anwendungen übergehen. Wir haben schon in \S.\,37 gesehen, dass wir aus jeder Gruppe $S$ von linearen Substitutionen unendlich viele isomorphe Gruppen ableiten können, indem wir mit einer willkürlichen Substitution $L$ von derselben Dimension transformiren, also die transformirte Gruppe
\weq{1}{L^{-1} S L}
bilden; und dies ist gleichbedeutend mit der Einführung anderer Variablen $y$ an Stelle von $x$ durch die Substitution
\weq{2}{(x) = L(y).}

Diese Transformation können wir dazu verwenden, um die Gruppe $S$ in einer einfachen Normalform darzustellen, und dann erhalten auch die Invarianten gewisse feste Normalformen, die in manchen Fällen sehr einfache Gestalten annehmen können. Bilden wir für die Normalform die Resolvente des Formenproblems [\S.\,43, (8)], die wir jetzt in der Form schreiben wollen:
\weq{3}{\Phi(\Theta, A_1, A_2, \ldots) = 0,}
worin $A_1, A_2, \ldots$ Invarianten der Gruppe $S$ bedeuten, so wird auch diese, wenn wir die Normalform benutzen, eine einfache Gestalt erhalten. Die Variablen $x_i$ lassen sich, wie wir gesehen
%CONT>>


% ===== scan leaf 203, printed 182 =====
% !PRINTED 182
\weberpage{182}

%% [running head, not body text -- commented out 2026-08-07] Page furniture (running
%% head, folio, signature) is excluded from this edition by standing policy.
%% \noindent 182 \hfill \textit{Sechster Abschnitt.} \hfill \S.\,47.

%<<CONT
haben, rational durch $\Theta$ und die $A_k$ ausdrücken, und wir setzen
\weq{4}{x_i = \varphi_i(\Theta, A_1, A_2, \ldots).}

Auch diese Ausdrücke werden, wenn über die Function $\Theta$ verfügt ist, für die Normalform feste Gestalten annehmen.

Nun kann aber auch der Fall vorkommen, dass die Invarianten nicht in der Normalform, sondern in einer beliebigen anderen Form, die wir die allgemeine Form nennen wollen, gegeben sind. Dann wird es sich darum handeln, die Substitution $L$ zu finden, durch die die allgemeine Form in eine von vornherein als möglich erkannte Normalform transformirt wird. Diese Aufgabe ist, wie wir nun sehen wollen, keine andere, als das für die Normalform gestellte Formenproblem selbst.

Nehmen wir an, die Invarianten in der allgemeinen Form, als Functionen von $y$, seien $B_1, B_2, \ldots$, so dass durch die Substitution (2) die Identitäten
\weq{5}{A_1 = B_1, \ A_2 = B_2, \ldots}
hergestellt werden. Es ist die Aufgabe, wenn die Functionen $A_i, B_i$ der Form nach gegeben sind, die Substitution $L$ zu bestimmen, die die Gleichungen (5) zu identischen macht. Diese Aufgabe ist gelöst, wenn wir die Gleichung (3) für ein passend gewähltes specielles Werthsystem der $A_i$ als gelöst voraussetzen. Um dies nachzuweisen, bilden wir die vollständigen Differentiale der Gleichungen (3) und (4):
\[
0 = \Phi'(\Theta)\,d\Theta + \sum^s \Phi'(A_s)\,dA_s
\]
\[
dx_i = \varphi_i'(\Theta)\,d\Theta + \sum^s \varphi_i'(A_s)\,dA_s,
\]
worin $\Phi'(\Theta)$, $\Phi'(A_s)$, $\varphi_i'(\Theta)$, $\varphi_i'(A_s)$ die partiellen Ableitungen bedeuten. Eliminiren wir $d\Theta$, so folgt:
\weq{6}{dx_i = \sum^s \frac{\varphi_i'(A_s)\,\Phi'(\Theta) - \Phi'(A_s)\,\varphi_i'(\Theta)}{\Phi'(\Theta)}\, dA_s.}

Nun ist in Folge der Gleichungen (5):
\weq{7}{dA_s = \sum^k B_s'(y_k)\, dy_k,}
und wenn wir also
\weq{8}{x_i = \alpha_{1,i} y_1 + \cdots + \alpha_{n,i} y_n}
\weq{9}{dx_i = \alpha_{1,i}\, dy_1 + \cdots + \alpha_{n,i}\, dy_n}
%CONT>>


% ===== scan leaf 204, printed 183 =====
% !PRINTED 183
\weberpage{183}

%% [running head, not body text -- commented out 2026-08-07] Page furniture (running
%% head, folio, signature) is excluded from this edition by standing policy.
%% \noindent \S.\,47. \hfill \textit{Normalformen.} \hfill 183

%<<CONT
setzen, so ergiebt die Vergleichung von (6) mit (9):
\weq{10}{\alpha_{k,i} = \sum^s \frac{\varphi_i'(A_s)\,\Phi'(\Theta) - \Phi'(A_s)\,\varphi_i'(\Theta)}{\Phi'(\Theta)}\, B_s'(y_k).}

Die rechte Seite dieser Gleichungen muss sich also auf eine Constante reduciren, und wir können ihren Werth finden, wenn wir für die $y$ irgend ein specielles Werthsystem setzen, das nur an die eine Bedingung gebunden ist, dass $\Phi'(\Theta)$ nicht verschwindet. Für dieses specielle Werthsystem sind die Werthe der $A_i$ durch die Gleichungen (5) bestimmt, und $\Theta$ ist bekannt, wenn wir für dies specielle Werthsystem der $A_i$ das Formenproblem (3) als gelöst voraussetzen. Dann sind durch (10) die Co\"efficienten $\alpha_{k,i}$ und damit die Substitution $L$ vollständig bestimmt.

\noindent\hrulefill


% ===== scan leaf 205, printed 184 =====
\weberpage{184}% page number not printed on this page
\pdfbookmark[1]{Siebenter Abschnitt. Gruppen binärer linearer Substitutionen.}{bm56}
\pdfbookmark[2]{§ 48. Ternäre orthogonale Substitutionen}{bm57}
% !PRINTED none

\begin{center}
\gesperrt{Siebenter Abschnitt.}

Gruppen bin\"arer linearer Substitutionen.

\rule{4cm}{0.4pt}

\S.\,48.

\gesperrt{Tern\"are orthogonale Substitutionen.}
\end{center}

Es ist nun unsere Aufgabe, aus der Gesammtheit der linearen
Substitutionen engere Gruppen auszusondern und schliesslich
zu endlichen Gruppen zu gelangen. Solche engere Gruppen, die
immer noch unendlich sein k\"onnen, erh\"alt man, wenn man die
Forderung stellt, dass gegebene homogene Functionen der
Variablen invariant bleiben sollen. Wir wollen aber die Aufgabe
nicht in dieser Allgemeinheit weiter verfolgen, sondern gleich
zur Betrachtung des wichtigsten speciellen Falles \"ubergehen. Wir
wollen uns auf tern\"are Substitutionen beschr\"anken und fordern,
dass eine quadratische Form von nicht verschwindender Determinante invariant bleiben soll. Da, wie wir fr\"uher gesehen
haben (Bd. I, \S.\,57), jede solche quadratische Form durch lineare
Transformation in eine Summe von Quadraten verwandelt werden
kann, so beschr\"anken wir das Problem nicht weiter, wenn wir
f\"ur diese quadratische Form die Summe der Quadrate annehmen.
Solche Substitutionen heissen \gesperrt{orthogonal}.

Die Substitution

\weq{1}{(y_1, y_2, y_3) = A (x_1, x_2, x_3)}

ist also orthogonal, wenn die Substitutionscoëfficienten so bestimmt sind, dass die Identit\"at besteht:

\weq{2}{y_1^2 + y_2^2 + y_3^2 = x_1^2 + x_2^2 + x_3^2.}

Wenn man die Ausdr\"ucke (1) in (2) substituirt und die
Coëfficienten entsprechender Glieder einander gleich setzt, so
%CONT>>


% ===== scan leaf 206, printed 185 =====
% !PRINTED 185
\weberpage{185}
%<<CONT
erh\"alt man sechs Relationen zwischen den neun Coëfficienten
von $A$.

Diese Relationen lauten, wenn

\weq{3}{A = \begin{pmatrix} a_1, & a_2, & a_3\\ b_1, & b_2, & b_3\\ c_1, & c_2, & c_3 \end{pmatrix}}

angenommen wird:

\weq{4}{\begin{array}{r@{}c@{}l@{\qquad}r@{}c@{}l}
a_1^2+b_1^2+c_1^2 &{}={}& 1, & a_2a_3+b_2b_3+c_2c_3 &{}={}& 0\\
a_2^2+b_2^2+c_2^2 &{}={}& 1, & a_3a_1+b_3b_1+c_3c_1 &{}={}& 0\\
a_3^2+b_3^2+c_3^2 &{}={}& 1, & a_1a_2+b_1b_2+c_1c_2 &{}={}& 0,
\end{array}}

und sind aus der analytischen Geometrie wohl bekannt. F\"ur das
Quadrat der Substitutionsdeterminante $|A|$ ergiebt sich aus diesen
Relationen nach der Multiplicationsregel der Determinanten der
Werth 1, und folglich hat $|A|$ den Werth $\pm 1$.

Aus den Formeln (4) ergiebt sich, dass die \gesperrt{inverse Substitution} zu $A$

\[
A^{-1} = \begin{pmatrix} a_1, & b_1, & c_1\\ a_2, & b_2, & c_2\\ a_3, & b_3, & c_3 \end{pmatrix}
\]

mit der transponirten identisch ist, und diese Eigenschaft k\"onnte
auch als Definition der orthogonalen Substitutionen dienen.

Die Gesammtheit der orthogonalen Substitutionen bildet
eine Gruppe. Darunter ist eine engere Gruppe enthalten, die
durch den Werth

\weq{5}{|A| = + 1}

ausgezeichnet ist, die wir als die Gruppe der \gesperrt{eigentlichen
orthogonalen Substitutionen} bezeichnen wollen.

Man kann die lineare Substitution (1) als den Uebergang
von einem rechtwinkeligen Coordinatensysteme zu einem zweiten
mit demselben Anfangspunkte deuten, wenn man $x_1, x_2, x_3$ und
$y_1, y_2, y_3$ als rechtwinkelige Coordinaten eines und desselben
(ver\"anderlichen) Punktes in dem ersten und dem zweiten Coordinatensysteme ansieht. Durch $A$ ist die gegenseitige Lage der
beiden Coordinatensysteme bestimmt und umgekehrt. Besteht
die Bedingung (5), wie wir jetzt voraussetzen wollen, so kann
das erste Coordinatensystem mit dem zweiten zur Deckung gebracht werden durch Drehung um eine feste Axe mit einem
bestimmten Drehungswinkel.


% ===== scan leaf 207, printed 186 =====
% !PRINTED 186
\weberpage{186}
Denn setzen wir

\[
D = \begin{vmatrix} a_1 - 1, & a_2, & a_3\\ b_1, & b_2 - 1, & b_3\\ c_1, & c_2, & c_3 - 1 \end{vmatrix},
\]

so erhalten wir aus den Relationen (4):

\[
|A|\, D = - D,
\]

also, wenn (5) besteht, $D = 0$. Demnach lassen sich drei
Gr\"ossen $\lambda, \mu, \nu$ aus den Gleichungen:

\weq{6}{\begin{array}{l}
\lambda = a_1 \lambda + a_2 \mu + a_3 \nu\\
\mu = b_1 \lambda + b_2 \mu + b_3 \nu\\
\nu = c_1 \lambda + c_2 \mu + c_3 \nu\\
\lambda^2 + \mu^2 + \nu^2 = 1
\end{array}}

bestimmen, und diese drei Gr\"ossen $\lambda, \mu, \nu$ bestimmen die Richtung einer geraden Linie, die mit den Axen $y_1, y_2, y_3$ dieselben
Winkel einschliesst, wie mit den Axen $x_1, x_2, x_3$. Eine in dieser
Richtung durch den Coordinatenanfangspunkt gelegte Linie ist
die Drehungsaxe.

In dem Falle $|A| = - 1$ trifft dieser Schluss nicht mehr zu.

Daneben besteht noch eine andere Deutung der orthogonalen
Substitution $(y) = A (x)$, wonach $(x)$ und $(y)$ die Coordinaten
zweier verschiedener Punkte $x$ und $y$ sind, bezogen auf ein und
dasselbe Coordinatensystem. Wenn $x$ einen Raumtheil (Linie,
Fl\"ache oder K\"orper) \"uberstreicht, so erf\"ullt der entsprechende
Punkt $y$ einen congruenten Raumtheil (wenn $|A| = - 1$ ist, einen
spiegelbildlich gleichen). Dies wird durch die folgende Betrachtung dargethan.

Die Gruppe der \gesperrt{eigentlichen} orthogonalen Substitutionen
ist \"aquivalent mit einer Gruppe, die man aus den verschiedenen
Stellungen eines um einen festen Punkt drehbaren K\"orpers bilden
kann. Diese Gruppe erh\"alt man, wenn man eine beliebige Stellung $E$ als Einheit annimmt, aus der man in irgend eine andere
Stellung $A$ gelangt durch Drehung um eine bestimmte Axe mit
einem bestimmten Winkel. Ist $B$ eine dritte Stellung, so hat
man unter der zusammengesetzten Stellung $AB$ die Stellung zu
verstehen, die man erh\"alt, wenn man die Drehung, die zu $A$
gef\"uhrt hat, nicht von der Einheitsstellung, sondern von der
Stellung $B$ aus vollf\"uhrt. Denn nimmt man ein mit dem K\"orper
in starrer Verbindung stehendes Axensystem an, z. B. das System
der Haupttr\"agheitsaxen, und bezeichnet mit $(x)$ die Coordinaten
%CONT>>


% ===== scan leaf 208, printed 187 =====
% !PRINTED 187
\weberpage{187}
%<<CONT
eines beliebigen, im Raume festen Punktes, bezogen auf dies
Axensystem in der Einheitsstellung $E$, so erh\"alt man die Coordinaten desselben Punktes, bezogen auf das Axensystem in der
Stellung $A$ oder $B$ durch zwei orthogonale Substitutionen $A(x)$
und $B(x)$.

Demnach sind $AB(x)$ die auf das System $A$ bezogenen
Coordinaten eines Punktes $y$ mit den Coordinaten $(y) = B(x)$
im Systeme $E$. Der Punkt $y$ hat also im Systeme $E$ dieselben
Coordinaten, wie der Punkt $x$ im Systeme $B$, d. h. $y$ liegt zur
Einheitsstellung so wie $x$ zur Stellung $B$. F\"uhren wir nun die
Drehung, die von $E$ zu $B$ f\"uhrt, so aus, dass wir den Punkt $x$
festhalten, aber den Punkt $y$ und das Axensystem $A$ die Drehung
begleiten lassen, so gelangt der Punkt $y$ nach $x$, und die Coordinaten von $x$, bezogen auf die neue Lage des Systemes $A$, sind
dieselben, wie die des Punktes $y$ in Bezug auf die urspr\"ungliche
Lage von $A$, d. h. $AB(x)$. Diese neue Lage des Systemes $A$
kann man aber offenbar auch so erreichen, dass man die Drehung,
die zu der Stellung $A$ f\"uhrt, nicht von der Einheitsstellung,
sondern von der Stellung $B$ aus vollzieht.

Aus der Gruppe der eigentlichen erh\"alt man die Gesammtheit
aller orthogonalen Substitutionen durch Zusammensetzung mit
\gesperrt{einer} uneigentlichen, etwa mit $(x_1, x_2, x_3) = (y_1, y_2, - y_3)$,
die, nach der zweiten geometrischen Auffassung, eine Spiegelung
an der Ebene $x_1, x_2$ ist, bei der der Punkt $y$ das Spiegelbild
des Punktes $x$ ist.

Von besonderem Interesse sind nun die in der Gesammtheit
der eigentlichen orthogonalen Substitutionen enthaltenen \gesperrt{endlichen Gruppen}, auf die wir sp\"ater noch n\"aher eingehen werden.
Wir wollen hier nur noch \"uber die geometrische Seite dieser
Frage Folgendes bemerken. Einer solchen endlichen Gruppe $G$
vom Grade $g$ von orthogonalen Substitutionen entspricht eine
endliche Gruppe von Stellungen eines K\"orpers. Denkt man sich
den K\"orper in diesen verschiedenen Stellungen gleichzeitig fixirt
und das Ganze zu einem neuen starren K\"orper vereinigt, so
erh\"alt man ein Gebilde, das in einer endlichen Anzahl $g$ verschiedener Stellungen sich selbst decken kann. Man kann f\"ur
jede Gruppe K\"orper von unendlich vielen verschiedenen Gestalten
finden. Die anschaulichsten und bekanntesten Verh\"altnisse ergeben sich, wenn man den K\"orper von ebenen Fl\"achen begrenzt
annimmt.


% ===== scan leaf 209, printed 188 =====
% !PRINTED 188
\weberpage{188}

Solche Gebilde sind die regulären Pyramiden, die Doppelpyramiden
und die regulären Körper (Tetra\"eder, Octa\"eder, Würfel,
Dodeka\"eder und Ikosa\"eder).

Die reguläre Pyramide von $g$ Seiten gelangt durch Drehung um die
Hauptaxe mit einem Winkel $2\pi:g$ und durch Wiederholung dieser
Drehung auf $g$ verschiedene Arten mit sich zur Deckung.

Ist $g$ gerade, so kann man eine reguläre Doppelpyramide von $g$
Seitenflächen ausser durch Drehung um ihre Hauptaxe mit dem Winkel
$4\pi:g$ noch (auf $\tfrac{1}{2}g$ verschiedene Arten) durch Drehung
um eine in der Aequatorialebene liegende Axe mit einem
Drehungswinkel von 180 Grad mit sich zur Deckung bringen.

Das Tetra\"eder gestattet 12 verschiedene Stellungen, in denen es
denselben Raum einnimmt.

Um sie zu erhalten, bezeichne man den Ort einer Ecke in der
Einheitsstellung mit 1. Dann erhält man drei Stellungen des
Tetra\"eders, bei denen die Ecke 1 fest bleibt. Man kann aber jede
Ecke an die Stelle von 1 bringen, wodurch die Zahl sich
vervierfacht.

Dieselbe Zahl findet man auch, wenn man die Ebene einer
Seitenfläche festhält, wobei man noch drei Stellungen des
Tetra\"eders findet, und jede Seitenfläche in die Ausgangsebene
bringt.

Endlich kann man auch so zählen, dass man jede Kante auf zwei
Arten mit einer festen Strecke zur Deckung bringt. Ebenso
verfährt man bei den übrigen regulären Körpern. Man findet so den
Grad der Gruppe

\begin{quote}
gleich dem Producte aus der Anzahl der Ecken mit der Anzahl der in
einer Ecke zusammenstossenden Kanten oder Seitenflächen, oder

gleich dem Producte aus der Anzahl der Seitenflächen mit der
Anzahl der Seiten oder Ecken einer Grenzfläche, oder

gleich der doppelten Anzahl der Kanten.
\end{quote}

Jede dieser Zählungen ergiebt

\begin{quote}
für das Tetra\"eder 12 Stellungen,\\
für das Octa\"eder und den Würfel 24 Stellungen,\\
für das Dodeka\"eder und Ikosa\"eder 60 Stellungen.
\end{quote}


% ===== scan leaf 210, printed 189 =====
% !PRINTED 189
\weberpage{189}
\pdfbookmark[2]{§ 49. Lineare gebrochene Substitutionen}{bm58}

Es sei schliesslich noch bemerkt, dass, anstatt der \gesperrt{Stellungen} des
Körpers, auch die Drehungen selbst als Elemente der Gruppe
aufgefasst werden können.

\begin{center}
\S.\,49.

\gesperrt{Lineare gebrochene Substitutionen.}
\end{center}

Die Gruppe der orthogonalen ternären Substitutionen ist, wie
jetzt nachgewiesen werden soll, isomorph mit der Gruppe der
linearen gebrochenen Substitutionen einer Veränderlichen, oder der
binären Substitutionen der Verhältnisse (\S.\,38).

Wir bezeichnen
eine ternäre orthogonale Substitution mit
\weq{1}{(y_1,\ y_2,\ y_3) = A\,(x_1,\ x_2,\ x_3),}
\weq{2}{y_1^2 + y_2^2 + y_3^2 = x_1^2 + x_2^2 + x_3^2.}
Hierauf wenden wir zunächst nach \S.\,37, (21) die
Transformation durch die feste (nicht orthogonale) Substitution
\weq{3}{L = \left\{\begin{array}{ccc}
i, & 0, & 0 \\[4pt]
0, & \dfrac{1}{\sqrt{2}}, & \dfrac{1}{\sqrt{2}} \\[6pt]
0, & \dfrac{i}{\sqrt{2}}, & \dfrac{-i}{\sqrt{2}}
\end{array}\right\},\quad
L^{-1} = \left\{\begin{array}{ccc}
-i, & 0, & 0 \\[4pt]
0, & \dfrac{1}{\sqrt{2}}, & \dfrac{-i}{\sqrt{2}} \\[6pt]
0, & \dfrac{1}{\sqrt{2}}, & \dfrac{i}{\sqrt{2}}
\end{array}\right\}}
an, worin $i = \sqrt{-1}$ ist, deren Determinante den Werth 1 hat,
und setzen
\weq{4}{\begin{aligned}
(y_1,\ y_2,\ y_3) &= L\,(y_1',\ y_2',\ y_3')\\
(x_1,\ x_2,\ x_3) &= L\,(x_1',\ x_2',\ x_3')' % [sic] source prints a stray prime after the closing parenthesis
\end{aligned}}
\weq{5}{(y_1',\ y_2',\ y_3') = A'\,(x_1',\ x_2',\ x_3'),}
worin
\weq{6}{A' = L^{-1} A L}
gleichfalls eine lineare Substitution bedeutet, deren
Determinante $|A'| = |A|$, also gleich $\pm 1$ ist. Die Relation
(2) geht durch die Substitutionen (4) in
\weq{7}{-y_1'^2 + 2\,y_2'\,y_3' = -x_1'^2 + 2\,x_2'\,x_3'}
über, woraus sich sechs Relationen zwischen den Co\"efficienten
von $A'$ ergeben, die wir aber hier nicht aufstellen wollen, da
wir die Substitution $A'$ auf andere Weise einfacher finden
können.


% ===== scan leaf 211, printed 190 =====
% !PRINTED 190
\weberpage{190}

Wenn wir nämlich unter $\xi_1, \xi_2$ neue Variable verstehen und
\weq{8}{x_1' = \sqrt{2}\,\xi_1\,\xi_2,\quad x_2' = \xi_1^2,\quad x_3' = \xi_2^2}
setzen, so verschwindet $x_1'^2 - 2\,x_2'\,x_3'$ identisch und
$y_1', y_2', y_3'$ gehen durch (5) und (8) in binäre quadratische
Formen der Variablen $\xi_1, \xi_2$ über, die nach (7) der
Gleichung
\weq{9}{y_1'^2 - 2\,y_2'\,y_3' = 0}
identisch genügen müssen, d. h. $y_2' y_3'$ muss ein
vollständiges Quadrat werden. Die quadratischen Formen $y_1', y_2',
y_3'$ zerlegen wir nun in je zwei lineare Factoren. Dabei können
$y_2'$ und $y_3'$ keinen gemeinsamen Factor haben, da sonst auch
der andere Factor in beiden Functionen übereinstimmen und also
die drei Co\"efficienten von $y_2'$ und $y_3'$ mit einander
proportional sein müssten. Dann würde aber $|A'|$ verschwinden,
was nicht möglich ist. Demnach ergiebt sich aus (9), dass $y_2'$
und $y_3'$ Quadrate linearer Functionen sein müssen. Setzen wir
\weq{10}{\eta_1 = \alpha\,\xi_1 + \beta\,\xi_2,\quad \eta_2 = \gamma\,\xi_1 + \delta\,\xi_2,}
so können wir also hiernach mit Rücksicht auf (9)
\weq{11}{y_1' = \sqrt{2}\,\eta_1\,\eta_2,\quad y_2' = \eta_1^2,\quad y_3' = \eta_2^2}
setzen; wenn wir die Multiplicationen ausführen und dann die
Substitution (8) im umgekehrten Sinne ausführen, so erhalten wir
aus (10) und (11):
\[
\begin{aligned}
y_1' &= (\alpha\delta + \beta\gamma)\,x_1' + \sqrt{2}\,\alpha\gamma\,x_2' + \sqrt{2}\,\beta\delta\,x_3', \\
y_2' &= \sqrt{2}\,\alpha\beta\,x_1' + \alpha^2\,x_2' + \beta^2\,x_3', \\
y_3' &= \sqrt{2}\,\gamma\delta\,x_1' + \gamma^2\,x_2' + \delta^2\,x_3'.
\end{aligned}
\]
Also lautet die Substitution $A'$:
\weq{12}{A' = \begin{pmatrix}
\alpha\delta + \beta\gamma, & \sqrt{2}\,\alpha\gamma, & \sqrt{2}\,\beta\delta \\
\sqrt{2}\,\alpha\beta, & \alpha^2, & \beta^2 \\
\sqrt{2}\,\gamma\delta, & \gamma^2, & \delta^2
\end{pmatrix}.}
Durch diese Substitution ist die Bedingung (7) noch nicht völlig
befriedigt, sondern es folgt nur, dass die rechte und linke Seite
sich durch einen constanten Factor unterscheiden, der von den
Co\"efficienten $\alpha, \beta, \gamma, \delta$ abhängt. Nun ist
aber der Co\"efficient von $x_1'^2$ in der Verbindung
$y_1'^2 - 2\,y_2'\,y_3'$:
\[
(\alpha\delta + \beta\gamma)^2 - 4\,\alpha\delta\beta\gamma = (\alpha\delta - \beta\gamma)^2,
\]
%CONT>>


% ===== scan leaf 212, printed 191 =====
% !PRINTED 191
\weberpage{191}
%<<CONT
und wir müssen also die Determinante $\alpha\delta - \beta\gamma = \pm 1$
setzen. Berechnet man die Determinante $|A'|$, so ergiebt sich
dafür $(\alpha\delta - \beta\gamma)^3$, so dass also, wenn wir nur
die eigentlichen orthogonalen Substitutionen berücksichtigen,
\weq{13}{\alpha\delta - \beta\gamma = 1}
setzen müssen, während
\weq{14}{\alpha\delta - \beta\gamma = -1}
den uneigentlichen orthogonalen Substitutionen entspricht.
Dadurch ist die Substitution $A'$ völlig bestimmt, und ist
zurückgeführt auf die binäre Substitution
\weq{15}{(\eta_1, \eta_2) = \begin{pmatrix} \alpha, & \beta \\ \gamma, & \delta \end{pmatrix} (\xi_1, \xi_2),\quad \alpha\delta - \beta\gamma = \pm 1,}
oder, wenn $\eta_1 : \eta_2 = \eta$, $\xi_1 : \xi_2 = \xi$ gesetzt
wird, auf die lineare gebrochene Substitution
\weq{16}{\eta = \dfrac{\alpha\xi + \beta}{\gamma\xi + \delta}.}
Hierzu ist aber nun noch Folgendes zu bemerken. Die beiden
Substitutionen $A, A'$ sind Transformationen von einander, und
entsprechen sich also gegenseitig eindeutig.

Durch $A'$ sind aber die Zahlen $\alpha, \beta, \gamma, \delta$
nur bis auf das gemeinsame Vorzeichen bestimmt. Wenn wir also die
lineare Substitution
\weq{17}{\begin{pmatrix} \alpha, & \beta \\ \gamma, & \delta \end{pmatrix}}
mit einer der beiden Bedingungen (13) und (14) betrachten, so ist
zwar hierdurch die Substitution $A'$ und daher auch $A$ eindeutig
bestimmt; aber umgekehrt entsprechen jedem $A'$ zwei
Substitutionen der Form (17):
\[
\begin{pmatrix} \alpha, & \beta \\ \gamma, & \delta \end{pmatrix}
\quad\text{und}\quad
\begin{pmatrix} -\alpha, & -\beta \\ -\gamma, & -\delta \end{pmatrix}.
\]
Wir erhalten aber wieder ein eindeutiges Entsprechen, wenn wir
diese beiden Substitutionen zu einem gemeinsamen Begriffe, den
wir mit $A''$ bezeichnen, zusammenfassen.

Was nun endlich die Substitution (16) betrifft, die wir mit
\weq{18}{A''' = \begin{pmatrix} \alpha, & \beta \\ \gamma, & \delta \end{pmatrix}}
bezeichnen, und nach \S.\,38 als Substitution der Verhältnisse
auffassen, so können wir einen gemeinsamen Factor von
$\alpha, \beta, \gamma, \delta$
%CONT>>


% ===== scan leaf 213, printed 192 =====
% !PRINTED 192
\weberpage{192}
%<<CONT
so bestimmen, dass sowohl die Bedingung (13) als (14) befriedigt
wird. % [sic/wording] the common factor has two choices, satisfying (13) or (14), respectively
Demnach bekommen wir aus jeder Substitution $A'''$ zwei
orthogonale Substitutionen $A$, von denen die eine eigentlich, die
andere uneigentlich ist.

Die drei Substitutionen
\weq{19}{A,\ A',\ A''}
entsprechen sich also hiernach gegenseitig eindeutig, die Substitutionen

\weq{20}{A,\ A',\ A'''}
aber nur dann, wenn wir noch die Forderung stellen, dass $A$
eine eigentliche orthogonale Substitution sein soll.

Ist nun
\weq{21}{B,\ B',\ B''}
ein zweites System von Substitutionen der Form (19), so sind
auch $AB$, $A'B'$ zwei einander entsprechende Substitutionen, wie
unmittelbar aus der Bedeutung der $A'$ als transformirte Substitution
 der $A$ hervorgeht. Es ist aber noch nachzuweisen,
dass auch
\weq{22}{AB,\ A'B',\ A''B''}
ein zusammengehöriges System von der Form (19) ist, dass also
die Gruppen der $A, A', A''$ isomorph sind. Daraus ergiebt sich
dann von selbst aus der Beziehung, in der $A''$ und $A'''$ zu einander
 stehen, dass (bei Beschränkung auf eigentlich orthogonale
Substitutionen $A$) auch die Gruppe der $A'''$ damit isomorph ist.

Setzen wir, um dies zu beweisen:
\weq{23}{(\eta_1,\eta_2) = A''(\xi_1,\xi_2),\quad (\xi_1,\xi_2) = B''(\zeta_1,\zeta_2),}
so folgt
\weq{24}{(\eta_1,\eta_2) = A''B''(\zeta_1,\zeta_2).}

Aus $A''$, $B''$, $A''B''$ leiten wir nun nach (12) drei ternäre
Substitutionen $A'$, $B'$, $C'$ her. So erhalten wir nach (8):
\weq{25}{\begin{aligned}
(\sqrt{2}\,\eta_1\eta_2,\ \eta_1^2,\ \eta_2^2) &= A'(\sqrt{2}\,\xi_1\xi_2,\ \xi_1^2,\ \xi_2^2)\\
(\sqrt{2}\,\xi_1\xi_2,\ \xi_1^2,\ \xi_2^2) &= B'(\sqrt{2}\,\zeta_1\zeta_2,\ \zeta_1^2,\ \zeta_2^2)\\
(\sqrt{2}\,\eta_1\eta_2,\ \eta_1^2,\ \eta_2^2) &= C'(\sqrt{2}\,\zeta_1\zeta_2,\ \zeta_1^2,\ \zeta_2^2).
\end{aligned}}

Diese Formeln sind in Bezug auf $\zeta_1, \zeta_2$ identisch, und sie
müssen also richtig bleiben, wenn $\sqrt{2}\,\zeta_1\zeta_2, \zeta_1^2, \zeta_2^2$ durch drei
%CONT>>


% ===== scan leaf 214, printed 193 =====
% !PRINTED 193
\weberpage{193}
\pdfbookmark[2]{§ 50. Realitätsbedingungen}{bm59}
%<<CONT
unabhängige Variable $z_1', z_2', z_3'$ ersetzt werden. Bezeichnen wir
die Ausdrücke, die sich dadurch für
\[ \sqrt{2}\,\xi_1\xi_2,\ \xi_1^2,\ \xi_2^2;\quad \sqrt{2}\,\eta_1\eta_2,\ \eta_1^2,\ \eta_2^2 \]
ergeben, mit
\[ x_1', x_2', x_3';\quad y_1', y_2', y_3', \]
so ergeben die Formeln (25):
\weq{26}{\begin{aligned}
(y_1', y_2', y_3') &= A'(x_1', x_2', x_3')\\
(x_1', x_2', x_3') &= B'(z_1', z_2', z_3')\\
(y_1', y_2', y_3') &= C'(z_1', z_2', z_3'),
\end{aligned}}
d. h. es ist
\[ C' = A'B', \]
w. z. b. w.

\begin{center}
\S.\,50.

\gesperrt{Realitätsbedingungen.}
\end{center}

Es bleibt uns noch eine Frage zu beantworten. Bisher haben
wir nirgends auf die Realität der Coëfficienten Rücksicht genommen.
Wenn aber irgend welche geometrische Anwendung
gemacht werden soll, so ist es nöthig, dass die orthogonale Substitution
 $A$ reell sei. Es ist also noch zu untersuchen, welchen
Bedingungen die Substitutionscoëfficienten $\alpha, \beta, \gamma, \delta$ in $A''$ zu
unterwerfen sind, damit die Coëfficienten von $A$ reell werden.

Um diese Frage zu entscheiden, bilden wir nach \S.\,49, (3), (12)
die Zusammensetzung
\[ A = LA'L^{-1}, \]
und erhalten:
\weq{1}{
A = \left\{
\begin{array}{ccc}
\alpha\delta+\beta\gamma, & i(\alpha\gamma+\beta\delta), & \alpha\gamma-\beta\delta\\[4pt]
-i(\alpha\beta+\gamma\delta), & \dfrac{\alpha^2+\beta^2+\gamma^2+\delta^2}{2}, & i\,\dfrac{-\alpha^2+\beta^2-\gamma^2+\delta^2}{2}\\[8pt]
\alpha\beta-\gamma\delta, & i\,\dfrac{\alpha^2+\beta^2-\gamma^2-\delta^2}{2}, & \dfrac{\alpha^2-\beta^2-\gamma^2+\delta^2}{2}
\end{array}
\right\}.
}

Die Coëfficienten dieser Substitution sollen also reell sein,
und ausserdem
\weq{2}{\alpha\delta - \beta\gamma = \pm 1.}

Wenn von den vier Coëfficienten $\alpha, \beta, \gamma, \delta$ einer verschwindet,
so muss auch noch ein zweiter verschwinden. Denn wenn z. B.
$\beta = 0$ ist, so muss $i\alpha\gamma$ und $\alpha\gamma$ reell sein. Dies ist aber mit (2)
%CONT>>


% ===== scan leaf 215, printed 194 =====
% !PRINTED 194
\weberpage{194}
%<<CONT
nur verträglich, wenn $\gamma = 0$ ist. Ebenso können $\alpha$ und $\delta$ nur
zugleich verschwinden.

Ist nun $\beta$ und $\gamma = 0$, so müssen $\alpha^2 + \delta^2$ und $i(\alpha^2 - \delta^2)$
reell sein, also $\alpha^2$ und $\delta^2$ conjugirt imaginär, und wegen (2) ihr
Product $= 1$. Je nachdem also in (2) das obere oder untere
Zeichen gilt, müssen $\alpha, \pm\delta$ conjugirt imaginär sein. Ebenso ergiebt
 sich, dass, wenn $\alpha$ und $\delta = 0$ sind, $\beta$ und $\mp\gamma$ conjugirt
imaginär sein müssen.

Nehmen wir also jetzt an, dass keine der vier Zahlen $\alpha, \beta, \gamma, \delta$
verschwinde. Dann müssen, damit die Grössen
\[
\begin{aligned}
&\alpha\delta+\beta\gamma,\quad i(\alpha\gamma+\beta\delta),\quad -i(\alpha\beta+\gamma\delta)\\
&\alpha\delta-\beta\gamma,\quad \alpha\gamma-\beta\delta,\quad \alpha\beta-\gamma\delta
\end{aligned}
\]
reell sein können,
\[
\begin{aligned}
&\alpha\delta \text{ und } \beta\gamma \text{ reell,}\\
&\alpha\gamma \text{ conjugirt imaginär mit } -\beta\delta,\\
&\alpha\beta \quad \text{\textquotedbl{}}\quad \text{\textquotedbl{}}\quad \text{\textquotedbl{}}\quad -\gamma\delta
\end{aligned}
\]
sein, und danach ist auch
\[
\alpha^2 = \frac{\alpha\gamma \cdot \alpha\beta}{\beta\gamma} \text{ conjugirt imaginär mit } \delta^2 = \frac{\beta\delta \cdot \gamma\delta}{\beta\gamma}
\]
und
\[
\beta^2 = \frac{\alpha\beta \cdot \delta\beta}{\alpha\delta} \quad \text{\textquotedbl{}}\quad \text{\textquotedbl{}}\quad \text{\textquotedbl{}}\quad \gamma^2 = \frac{\gamma\delta \cdot \alpha\gamma}{\alpha\delta};
\]
also muss $\alpha$ conjugirt mit $\pm\delta$, $\beta$ conjugirt mit $\mp\gamma$ sein. Dass
beide Male das gleiche Zeichen stehe, widerspricht der Forderung,
dass $\alpha\gamma$ und $-\beta\delta$ conjugirt sein sollen. Da das Product
zweier conjugirt imaginärer Grössen stets positiv ist, so gelten
die oberen oder die unteren Zeichen, je nachdem in (2) das
obere oder das untere Zeichen steht, d. h. je nachdem $A$ eigentlich
 oder uneigentlich orthogonal ist. Wir kommen also zu
folgendem allgemeinen Resultate, in dem auch die speciellen
Fälle von verschwindenden $\alpha, \beta, \gamma, \delta$ mit enthalten sind:

\gesperrt{Die nothwendige und hinreichende Bedingung
dafür, dass die orthogonale Substitution (1), (2) reell
sei, besteht darin, dass $\alpha, \pm\delta$ und $\beta, \mp\gamma$ zwei conjugirt
imaginäre Paare bilden, worin die oberen Zeichen bei
den eigentlich, die unteren bei den uneigentlich
orthogonalen Substitutionen gelten.}


% ===== scan leaf 216, printed 195 =====
% !PRINTED 195
\weberpage{195}
\pdfbookmark[2]{§ 51. Endliche Gruppen linearer gebrochener Substitutionen. Pole der Gruppen}{bm60}
\begin{center}
\S.\,51.

\gesperrt{Endliche Gruppen linearer gebrochener Substitutionen.}

\gesperrt{Pole der Gruppen.}
\end{center}

Aus den bisherigen Entwickelungen ergiebt sich, dass, wenn
alle endlichen Gruppen linearer gebrochener Substitutionen gefunden
 sind, daraus ohne Weiteres alle endlichen Gruppen
eigentlich orthogonaler ternärer Substitutionen und alle endlichen
 Gruppen binärer linearer Substitutionen mit der Determinante
 1 gefunden werden können.

Ausserdem giebt es noch endliche Gruppen, die neben den
eigentlichen auch uneigentlich orthogonale Substitutionen enthalten
.

Diese Gruppen, auf die wir später zurückkommen, können
nicht aus den Gruppen linearer gebrochener Substitutionen allein
abgeleitet werden, weil sich bei den gebrochenen Substitutionen
der Unterschied zwischen eigentlich und uneigentlich orthogonalen
 Substitutionen verwischt.

Wir suchen also jetzt zunächst \gesperrt{alle endlichen Gruppen
linearer gebrochener Substitutionen} zu ermitteln $^{1)}$.

Wir bezeichnen mit $G$ eine solche Gruppe vom Grade $n$,
und mit
\weq{1}{x,\ \Theta_1(x),\ \Theta_2(x),\ \ldots,\ \Theta_{n-1}(x)}
ihre Elemente, worin die $\Theta$ Symbole für lineare Functionen
\weq{2}{\Theta(x) = \frac{\alpha x + \beta}{\gamma x + \delta}}
sind, die auch mit
\weq{3}{\Theta = \begin{pmatrix}\alpha, & \beta\\ \gamma, & \delta\end{pmatrix}}
bezeichnet werden.

Aus jeder Gruppe von der Form (1) können wir nach \S.\,37
eine ganze Schaar isomorpher Gruppen ableiten, wenn wir mit
$L$ eine willkürliche lineare Substitution bezeichnen:
\weq{4}{1,\ L\Theta_1L^{-1},\ L\Theta_2L^{-1},\ \ldots,\ L\Theta_{n-1}L^{-1},}
%CONT>>

\par\noindent\rule{4cm}{0.4pt}\par
$^{1)}$ Ueber die Theorie dieser Gruppen ist besonders zu vergleichen:
\gesperrt{Gordan}, Ueber endliche Gruppen linearer Transformationen einer
Veränderlichen. Mathematische Annalen, Bd. XII, S. 23 (1877);
\gesperrt{Klein}, Vorlesungen über das Ikosaëder (Leipzig 1884).


% ===== scan leaf 217, printed 196 =====
% !PRINTED 196
\weberpage{196}
%<<CONT
und wir betrachten unsere Aufgabe als gelöst, wenn von jeder
dieser Schaaren ein Repräsentant bestimmt ist.

Wir werden diese Freiheit später benutzen, um die gefundenen
Gruppen möglichst einfach darzustellen.

Die Determinante
\[
\alpha\delta - \beta\gamma = \Delta
\]
muss von Null verschieden sein, und wenn es die Einfachheit
verlangt, können wir sie immerhin $= 1$ annehmen, was wir bisweilen
thun werden.

Wir bezeichnen im Sinne der Gruppentheorie die Wiederholung
einer Substitution durch Exponenten: $\Theta, \Theta^2, \Theta^3, \ldots$,
worunter also nicht Potenzen, sondern immer wieder lineare
Substitutionen der Gruppe (1) zu verstehen sind. Die identische
Substitution $x = x$, die als die Einheit der Gruppe anzusehen
ist, wird auch mit $\Theta^0$ oder mit 1 bezeichnet.

Zwei inverse Substitutionen $\Theta, \Theta^{-1}$ können in der Form
dargestellt werden:
\[
\Theta = \begin{pmatrix} \alpha, & \beta \\ \gamma, & \delta \end{pmatrix}, \quad
\Theta^{-1} = \begin{pmatrix} \delta, & -\beta \\ -\gamma, & \alpha \end{pmatrix}.
\]

Da die Gruppe nach der Voraussetzung endlich ist, so hat
jedes ihrer Elemente einen bestimmten \gesperrt{Grad}, d. h. es giebt für
jedes $\Theta$ eine bestimmte kleinste positive Zahl $t$, für die $\Theta^t = 1$
ist. Diese Zahl $t$ muss ein Theiler von $n$ sein.

Für die Folge ist es von Wichtigkeit, für jede der Substitutionen
der Gruppe (ausgenommen die identische Substitution)
die Werthe der Variablen zu betrachten, die ihren Transformirten
gleich werden, also die Wurzeln der Gleichungen
\weq{5}{x = \Theta(x).}
Diese Werthe wollen wir die \gesperrt{Pole} der Substitution $\Theta$
nennen.

Die Gleichung (5) ist quadratisch und nimmt, wenn man
für $\Theta$ den Ausdruck (2) einsetzt, die Form an:
\weq{6}{\gamma x^2 + (\delta - \alpha) x - \beta = 0.}
Diese Gleichung hat zwei Wurzeln, die nur dann einander
gleich sind, wenn
\weq{7}{(\delta - \alpha)^2 + 4\beta\gamma = 0}
oder
\weq{8}{(\alpha + \delta)^2 = 4\Delta}
%CONT>>


% ===== scan leaf 218, printed 197 =====
% !PRINTED 197
\weberpage{197}
%<<CONT
ist. Es ist leicht einzusehen, dass dieser Fall, wenn $\Theta$ einen
endlichen Grad hat, nicht vorkommen kann.

Denn ist zunächst $\beta$ oder $\gamma = 0$, so folgt aus (7), dass
$\alpha = \delta$ sein muss, und beide können $= 1$ angenommen werden.

Wenn nun $\Theta$ eine der beiden Substitutionen
\[
\begin{pmatrix} 1, & 0 \\ \gamma, & 1 \end{pmatrix}, \quad
\begin{pmatrix} 1, & \beta \\ 0, & 1 \end{pmatrix}
\]
ist, so ist für jedes $\lambda$:
\[
\Theta^\lambda = \begin{pmatrix} 1, & 0 \\ \lambda\gamma, & 1 \end{pmatrix} \text{ oder } \begin{pmatrix} 1, & \lambda\beta \\ 0, & 1 \end{pmatrix},
\]
wie man leicht durch vollständige Induction findet. Es kann
also, da $\beta$ und $\gamma$ nicht zugleich Null sein können, wenn $\Theta$ nicht
die identische Substitution ist, $\Theta^\lambda$ für kein positives $\lambda$ gleich 1
werden, wie es doch sein müsste, wenn $\lambda$ gleich dem Grade von
$\Theta$ wäre.

Ist aber $\beta$ von Null verschieden, so setzen wir, indem wir
jetzt $\Delta = 1$ und nach (8) $\alpha + \delta = 2$ annehmen,
\[
L = \begin{pmatrix} \beta, & 0 \\ 1-\alpha, & \beta^{-1} \end{pmatrix},
\]
und erhalten
\[
\begin{aligned}
L^{-1}\Theta L &= \begin{pmatrix} \beta^{-1}, & 0 \\ \alpha-1, & \beta \end{pmatrix} \begin{pmatrix} \alpha, & \beta \\ \gamma, & \delta \end{pmatrix} \begin{pmatrix} \beta, & 0 \\ 1-\alpha, & \beta^{-1} \end{pmatrix} \\
&= \begin{pmatrix} 1, & \beta^{-1} \\ 0, & 1 \end{pmatrix},
\end{aligned}
\]
also
\[
\Theta = L \begin{pmatrix} 1, & \beta^{-1} \\ 0, & 1 \end{pmatrix} L^{-1},
\]
und daraus
\[
\Theta^\lambda = L \begin{pmatrix} 1, & \lambda\beta^{-1} \\ 0, & 1 \end{pmatrix} L^{-1},
\]
was wieder für kein positives $\lambda$ gleich 1 werden kann. Wir
haben daher den Satz:

\par\noindent 1.\ \gesperrt{Jede der $n-1$ nicht identischen Substitutionen
einer Gruppe $n^{\text{ten}}$ Grades hat zwei von einander
verschiedene Pole.}

Jede nicht identische Substitution der Gruppe $G$ giebt uns
also zwei Pole. Es kann aber ein und derselbe Werth bei
mehreren verschiedenen Substitutionen als Pol auftreten. Zählen
wir einen dieser Werthe $h$ mal, wenn er in $h$ Substitutionen der
Gruppe als Pol vorkommt, so ergiebt sich die Anzahl der Pole
%CONT>>


% ===== scan leaf 219, printed 198 =====
% !PRINTED 198
\weberpage{198}
%<<CONT
gleich $2n-2$. Diese Werthe sollen die \gesperrt{Pole der Gruppe}
heissen. Die genaue Abzählung dieser Pole giebt uns die
wichtigsten Aufschlüsse über Zahl und Beschaffenheit der
möglichen endlichen Gruppen.

Wenn die Substitutionsco\"efficienten $\beta, \gamma$ gleich Null sind, so
ist $\Theta$ eine multiplicative Substitution
\[
\begin{pmatrix} \alpha, & 0 \\ 0, & \delta \end{pmatrix}.
\]

Damit diese Substitution von endlichem Grade sein kann,
muss der Quotient $\alpha : \delta = \varepsilon$ eine \gesperrt{Einheitswurzel sein,
deren Grad ein Theiler von $n$ ist}. Als Pole dieser Substitution
hat man $x = 0$ und $x = \infty$ anzusehen, die der Bedingung
\[
x = \varepsilon x
\]
genügen.

Nach (4) kann man aus $G$ eine isomorphe Gruppe
\[
L G L^{-1} = G'
\]
ableiten, wenn für $L$ eine beliebige lineare Substitution
\[
L = \begin{pmatrix} A, & B \\ C, & D \end{pmatrix}
\]
genommen wird. Die Pole dieser Gruppe erhält man, wenn
man auf die Pole von $G$ die Substitution $L$ anwendet. Wenn also
$a$ einer der Pole von $G$ ist, so ist
\[
a' = \frac{Aa + B}{Ca + D}
\]
der entsprechende Pol von $G'$.

Man kann die Substitutionsco\"efficienten $A, B, C, D$ so bestimmen,
dass drei der Pole von $G'$ vorgeschriebene Werthe erhalten.
Um z. B. den Polen $a, b, c$ von $G$ die Pole $0, \infty, 1$ von $G'$
entsprechen zu lassen, setze man
\[
L(x) = \frac{c-b}{c-a} \, \frac{x-a}{x-b}.
\]
Wenn $a$ und $b$ die Pole einer und derselben Substitution $\Theta$
von $G$ sind, so ist also die entsprechende Substitution von $G'$
multiplicativ.

Wir erhalten hieraus den Satz:

\par\noindent 2.\ \gesperrt{Zu jeder Gruppe linearer Substitutionen kann
man eine transformirte Gruppe finden, in der einer
beliebigen der gegebenen Substitutionen,}
%CONT>>


% ===== scan leaf 220, printed 199 =====
% !PRINTED 199
\weberpage{199}
\pdfbookmark[2]{§ 52. Die verschiedenen Arten möglicher Gruppen}{bm61}
%<<CONT
\gesperrt{nur nicht gerade der identischen, eine Multiplication
entspricht. Die transformirende Substitution $L$ ist dadurch
selbst bis auf eine Multiplication bestimmt.}

\begin{center}
\S.\,52.

\gesperrt{Die verschiedenen Arten möglicher Gruppen.}
\end{center}

Es sei $a$ einer der Pole der Gruppe $G$, und wir wollen
annehmen, es gebe $\nu - 1$ und nicht mehr Elemente in $G$,
$\Theta_1, \Theta_2, \ldots, \Theta_{\nu-1}$, so dass
\weq{1}{a = \Theta_1(a) = \Theta_2(a) \cdots = \Theta_{\nu-1}(a)}
ist. Ein solcher Pol soll ein $\nu$-zähliger Pol heissen. Es ist
nun zunächst klar, dass die Elemente
\weq{2}{1, \Theta_1, \Theta_2, \ldots, \Theta_{\nu-1}}
für sich eine Gruppe, und zwar einen Theiler von $G$ bilden;
denn aus
\[
a = \Theta_1(a), \quad a = \Theta_2(a)
\]
folgt, wenn man auf der rechten Seite der zweiten Gleichung
$a$ durch das ihm gleiche $\Theta_1(a)$ ersetzt:
\[
a = \Theta_2 \Theta_1(a).
\]
\gesperrt{Es muss also $\nu$ ein Theiler von $n$ sein:}
\weq{3}{n = \nu\mu.}
Wir bezeichnen die Gruppe (2) mit $Q$.

Es lässt sich beweisen, dass diese Gruppe cyklisch sein
muss, also aus den Wiederholungen eines ihrer Elemente
besteht.

Denn wenn wir durch Transformation nach dem Satze 2.
$a$ nach unendlich werfen, so erhalten die Substitutionen (2) die
Form:
\[
x, \; \varepsilon_1 x + c_1, \; \varepsilon_2 x + c_2, \ldots, \varepsilon_{\nu-1} x + c_{\nu-1},
\]
und die Composition von zweien unter ihnen giebt:
\[
\begin{aligned}
\Theta_1 \Theta_2 &= \varepsilon_1 \varepsilon_2 x + (c_2 \varepsilon_1 + c_1) \\
\Theta^\lambda &= \varepsilon^\lambda x + c(\varepsilon^{\lambda-1} + \varepsilon^{\lambda-2} + \cdots + 1) \\
\Theta^{-1} &= \varepsilon^{-1} x - c \varepsilon^{-1}.
\end{aligned}
\]
Daraus schliesst man, dass alle $\varepsilon_1, \varepsilon_2, \ldots, \varepsilon_{\nu-1}$ Einheitswurzeln
vom Grade $\nu$ sein müssen. Ausserdem können nicht zwei
der $\varepsilon$ einander gleich sein. Denn wäre z. B. $\varepsilon_1 = \varepsilon_2$, so wäre
\[
\Theta_1 \Theta_2^{-1} = x + (c_1 - c_2).
\]


% ===== scan leaf 221, printed 200 =====
% !PRINTED 200
\weberpage{200}
Diese Substitution muss in $G$ vorkommen und nach dem Satze 1. muss $c_1 = c_2$, also $\Theta_1$ mit $\Theta_2$ identisch sein. Die Zahlen $1, \varepsilon_1, \varepsilon_2, \ldots, \varepsilon_{\nu-1}$ müssen also zusammen alle $\nu^{\text{ten}}$ Einheitswurzeln enthalten, und also mit den Potenzen einer primitiven unter ihnen übereinstimmen. Da andererseits alle Potenzen von einer der Grössen $\Theta$ in der Gruppe (1) vorkommen müssen, so wird, wenn das in $\Theta$ vorkommende $\varepsilon$ eine primitive $\nu^{\text{te}}$ Einheitswurzel ist, die ganze Gruppe (1) durch die Potenzen von $\Theta$ erschöpft. % [sic/reference: both occurrences of ``Gruppe (1)'' are printed; the stabilizer group Q in (2), named in the next sentence, is required by the argument.]
Also besteht diese Gruppe $Q$ aus den Elementen

\weq{4}{1, \Theta, \Theta^2, \ldots, \Theta^{\nu-1}.}

Nehmen wir (nach dem Satze 2.) $\Theta$ als multiplicative Substitution an, so ist 0 der zweite Pol von $\Theta$ und zugleich der zweite Pol der sämmtlichen Substitutionen (4). Es ist daher dieser Pol gleichfalls ein mindestens \gesperrt{$\nu$-zähliger}. Da aber beide Pole von $\Theta$ in dieser Betrachtung vertauscht werden können, so erhalten wir die Sätze:

3. \gesperrt{Ein $\nu$-zähliger Pol bestimmt eine in $G$ enthaltene cyklische Gruppe $Q$ vom $\nu^{\text{ten}}$ Grade.}

4. \gesperrt{Die beiden Pole irgend einer Substitution $\Theta$ der Gruppe $G$ sind gleichzählige Pole und sind die beiden einzigen Pole der Gruppe $Q$.}

Nach \S.\,2 lassen sich $\mu - 1$ Elemente

\weq{5}{\psi_1, \psi_2, \ldots, \psi_{\mu-1}}

in $G$ so auswählen, dass $\psi_1 Q, \psi_2 Q, \ldots, \psi_{\mu-1} Q$ die Nebengruppen zu $Q$ sind, und dass also

\weq{6}{G = Q + \psi_1 Q + \psi_2 Q + \cdots + \psi_{\mu-1} Q}

wird. Wir setzen nun

\weq{7}{a_1 = \psi_1(a), a_2 = \psi_2(a), \ldots, a_{\mu-1} = \psi_{\mu-1}(a).}

Die Grössen $a_1, a_2, \ldots, a_{\mu-1}$ sind nicht nur alle von $a$, sondern auch unter einander verschieden. Denn wäre etwa

\[ \psi_1(a) = \psi_2(a), \]

so würde folgen:

\[ a = \psi_1^{-1}\psi_1(a) = \psi_1^{-1}\psi_2(a), \]

und $\psi_1^{-1}\psi_2$ wäre in $Q$ enthalten, also $\psi_2$ in $\psi_1 Q$, was gegen die Voraussetzung ist. Es ist nun leicht nachzuweisen, dass diese Werthe $a_1, a_2, \ldots, a_{\mu-1}$ sämmtlich $\nu$-zählige Pole der Gruppe sind. Es genügt, dies für $a_1$ zu zeigen.


% ===== scan leaf 222, printed 201 =====
% !PRINTED 201
\weberpage{201}
Wir nehmen irgend eine der Gleichungen (1)

\weq{8}{a = \Theta(a)}

und setzen in der Gleichung

\weq{9}{a_1 = \psi_1(a)}

für $a$ den ihm gleichen Werth $\Theta(a)$. So erhalten wir

\weq{10}{a_1 = \psi_1 \Theta(a).}

Nach (9) ist aber $a = \psi_1^{-1}(a_1)$, und folglich ergiebt sich aus (10):

\weq{11}{a_1 = \psi_1 \Theta \psi_1^{-1}(a_1).}

Daraus ersieht man, da wir $\nu - 1$ verschiedene Substitutionen $\Theta$ haben, dass $a_1$ ein $\nu$-zähliger Pol ist und zwar zu der mit $Q$ conjugirten Gruppe $\psi_1 Q \psi_1^{-1}$ gehörig. Aus diesem Grunde nennen wir

\weq{12}{a, a_1, a_2, \ldots, a_{\mu-1}}

\gesperrt{ein System conjugirter $\nu$-zähliger Pole von $G$.}

Eine beliebige Substitution $\chi$ von $G$ kann nach (6) immer in die Form $\psi_i \Theta$ gebracht werden, so dass $\Theta$ zu $Q$ gehört, und daraus folgt, dass $\chi(a) = \psi_i \Theta(a) = a_i$ ist, und ebenso, wenn man $\psi_i \Theta \psi_k = \psi_h \Theta'$ setzt, so dass auch $\Theta'$ zu $Q$ gehört:

\[ \chi(a_k) = \psi_i \Theta \psi_k(a) = \psi_h \Theta'(a) = a_h. \]

Da ferner zwei Grössen $\chi(a_h), \chi(a_k)$ nur dann einander gleich sein können, wenn $a_h = a_k$ ist, so folgt der Satz:

5. \gesperrt{Die beiden Reihen}

\[ a, \quad a_1, \quad a_2, \quad \ldots, \quad a_{\mu-1} \]
\[ \chi(a), \quad \chi(a_1), \quad \chi(a_2), \quad \ldots, \quad \chi(a_{\mu-1}) \]

\gesperrt{stimmen, welche Substitution aus $G$ auch für $\chi$ genommen werden mag, abgesehen von der Reihenfolge, mit einander überein.}

Ist dann $b$ ein in dem Systeme (12) nicht enthaltener Pol, so kann auch $\chi(b)$ nicht in (12) vorkommen, und daraus ergiebt sich, dass zwei Systeme conjugirter Pole, wenn sie nicht ganz identisch sind, keinen Pol gemein haben.

Hiernach können wir die sämmtlichen Pole der Gruppe $G$ in Systeme conjugirter Pole anordnen, und wir erhalten nach \S.\,51 ihre Gesammtzahl $2n - 2$, wenn wir jeden $\nu$-zähligen Pol $(\nu - 1)$mal mitrechnen.


% ===== scan leaf 223, printed 202 =====
% !PRINTED 202
\weberpage{202}
Diese Bemerkung giebt uns eine wesentliche Begrenzung der Zahlen $\nu$. Es ist nämlich danach:

\weq{13}{2n - 2 = \mu(\nu - 1) + \mu'(\nu' - 1) + \mu''(\nu'' - 1) + \cdots}

oder, wenn wir mit $h$ die Anzahl der Systeme conjugirter Pole bezeichnen, und $n = \mu \nu = \mu' \nu' = \cdots$ setzen:

\weq{14}{2n - 2 = nh - \mu - \mu' - \mu'' - \cdots}

Die Zahlen $\nu, \nu', \nu'' \ldots$ sind mindestens gleich 2, also ist

\[ 1 \leqq \mu \leqq \frac{n}{2}, \quad 1 \leqq \mu' \leqq \frac{n}{2}, \cdots \]

und daher nach (14)

\[ \frac{nh}{2} \leqq 2n - 2 \leqq (n-1)h, \]

oder

\[ 2 \leqq h \leqq 4 - \frac{4}{n}. \]

Es kann also $h$ nur einen der beiden Werthe 2 oder 3 haben, und wir finden fünf Arten, die Gleichung (13) zu befriedigen.

Wenn zunächst $h = 2$ ist, so folgt aus (14)

\[ \mu + \mu' = 2, \quad \mu = \mu' = 1, \quad \nu = \nu' = n. \]

Wir haben also:

I. \gesperrt{Kreistheilungsgruppe} oder \gesperrt{cyklische Gruppe:}

\[ \nu = \nu' = n, \quad n \text{ beliebig.} \]
\[ \mu = \mu' = 1, \]

Ist ferner $h = 3$, so folgt aus (14):

\weq{15}{\mu + \mu' + \mu'' = n + 2.}

Daraus ist zu schliessen, dass mindestens eine der Zahlen $\nu, \nu', \nu''$ gleich 2 sein muss. Denn wären sie alle $\geqq 3$, so wäre

\[ \mu \leqq \frac{n}{3}, \quad \mu' \leqq \frac{n}{3}, \quad \mu'' \leqq \frac{n}{3}, \]

also $\mu + \mu' + \mu'' \leqq n$, was mit (15) im Widerspruch steht.

Ist also $\nu = 2$, $\mu = \dfrac{n}{2}$, so folgt aus (15):

\weq{16}{\mu' + \mu'' = \frac{n}{2} + 2.}

Wir nehmen zunächst an, dass auch noch $\nu' = 2$ sei. Setzen wir dann $\nu'' = m$, so folgt aus (16):

\[ \mu'' = 2, \quad n = 2m, \]
%CONT>>


% ===== scan leaf 224, printed 203 =====
% !PRINTED 203
\weberpage{203}
\pdfbookmark[2]{§ 53. Transformation der Substitutionen von G auf einfache Formen}{bm62}
%<<CONT
und wir erhalten eine zweite Möglichkeit:

\gesperrt{II. Diëdergruppe,} $\nu = \nu' = 2, \nu'' = m, n = 2m, m \geqq 2$

\[ \mu = \mu' = m, \quad \mu'' = 2. \]

Ist ferner keine der Zahlen $\nu', \nu''$ gleich 2, so muss eine von ihnen gleich 3 sein. Denn sind sie beide $\geqq 4$, so ist

\[ \mu' + \mu'' \leqq \frac{n}{2}, \]

was mit (16) im Widerspruche steht. Ist also $\nu' = 3$, $\mu' = \dfrac{n}{3}$, so folgt aus (16):

\weq{17}{\mu'' = \frac{n}{6} + 2, \quad \nu'' = \frac{6n}{n+12},}

woraus folgt, dass $\nu'' < 6$ sein muss, also nur einen der Werthe 3, 4, 5 haben kann.

Danach bekommen wir noch drei mögliche Fälle:

\gesperrt{III. Tetraëdergruppe:} $\nu = 2, \nu' = 3, \nu'' = 3, n = 12$

\[ \mu = 6, \quad \mu' = 4, \quad \mu'' = 4. \]

\gesperrt{IV. Octaëdergruppe:} $\nu = 2, \nu' = 3, \nu'' = 4, n = 24$

\[ \mu = 12, \quad \mu' = 8, \quad \mu'' = 6. \]

\gesperrt{V. Ikosaëdergruppe:} $\nu = 2, \nu' = 3, \nu'' = 5, n = 60$

\[ \mu = 30, \quad \mu' = 20, \quad \mu'' = 12. \]

Hiermit sind alle Möglichkeiten erschöpft.

Es ist freilich noch nicht bewiesen, dass diese Gruppen, die wir einstweilen mit den gebräuchlichen Namen aufgeführt haben, und die wir unter dem gemeinsamen Namen der \gesperrt{Polyëdergruppen} zusammenfassen, wirklich existiren, noch wie gross ihre Mannigfaltigkeit ist. Zu diesem Beweise führen erst die folgenden Betrachtungen.

\begin{center}
\S.\,53.

\gesperrt{Transformation der Substitutionen von $G$ auf einfache Formen.}
\end{center}

Ehe wir zur definitiven Aufstellung dieser Gruppen gehen, leiten wir einen Satz ab, der die Möglichkeit der vorkommenden Substitutionen noch weiter beschränkt.


% ===== scan leaf 225, printed 204 =====
% !PRINTED 204
\weberpage{204}

Ist $a$ ein $\nu$-zähliger Pol der Gruppe $G$ und

\weq{1}{a = \Theta(a) = \Theta^2(a)\cdots = \Theta^{\nu-1}(a),}

so ist der zweite Pol $a'$ von $\Theta$ nach \S.\,52, 4. gleichfalls $\nu$-zählig,
und es ist

\weq{2}{a' = \Theta(a') = \Theta^2(a')\cdots = \Theta^{\nu-1}(a').}

Giebt es nun in der Gruppe $G$ mehr als ein System conjugirter
$\nu$-zähliger Pole, so können $a, a'$ entweder in demselben
oder auch in verschiedenen dieser Systeme vorkommen.

Giebt es aber nur \gesperrt{ein} System $\nu$-zähliger Pole, so muss $a$
und $a'$ in demselben Systeme vorkommen, und es muss also in $G$
eine nicht unter den Potenzen von $\Theta$ enthaltene Substitution $\psi$
existiren, so dass

\weq{3}{a' = \psi(a)}

ist. Daraus folgt mit Anwendung von (1) und (2):
\[
\psi(a) = \Theta\psi(a), \quad a = \psi^{-1}\Theta\psi(a),
\]
woraus zu schliessen ist, dass $\psi^{-1}\Theta\psi$ unter den Potenzen von
$\Theta$ vorkommt, und dass daher nach (2) auch
\[
a' = \psi^{-1}\Theta\psi(a'), \quad \psi(a') = \Theta\psi(a')
\]
sein muss. Es ist also auch $\psi(a')$ ein Pol von $\Theta$, und weil $\psi(a')$
nicht $= a'$ sein kann, weil sonst $a'$ ein Pol von $\psi$ wäre, was
nicht sein kann, da $\psi$ nicht unter den Potenzen von $\Theta$ vorkommt,
so ist

\weq{4}{\psi(a') = a.}

Wenn man nun durch Transformation der Gruppe die Pole
$a$ und $a'$ nach $0$ und $\infty$ bringt, so erhält $\Theta$ die Form
\[
\Theta(x) = \varepsilon x,
\]
worin $\varepsilon$ eine primitive $\nu^{\text{te}}$ Einheitswurzel ist, und $\psi$ muss wegen
(3) und (4) die Form haben:
\[
\psi(x) = \frac{c}{x},
\]
worin $c$ eine Constante ist. Durch eine abermalige Transformation
der Gruppe kann man der Constanten $c$ jeden beliebigen Werth,
z. B. durch Transformation mit der Substitution $x\sqrt{c}$ für $x$, den
Werth $1$ geben. Wir erhalten also folgenden Satz:

1.~\gesperrt{Wenn in der Gruppe $G$ nur ein System conjugirter
 $\nu$-zähliger Pole vorkommt, so kann man}
%CONT>>


% ===== scan leaf 226, printed 205 =====
% !PRINTED 205
\weberpage{205}
\pdfbookmark[2]{§ 54. Die Grundformen}{bm63}
%<<CONT
\gesperrt{die Gruppe so transformiren, dass sie die beiden
Substitutionen}
\[
\Theta = \varepsilon x, \quad \psi = \frac{c}{x}
\]
\gesperrt{enthält, worin $\varepsilon$ eine primitive $\nu^{\text{te}}$ Einheitswurzel
bedeutet, und $c$ ein beliebig vorgeschriebener
Werth, z. B. auch $1$, sein kann.}

Die Voraussetzung dieses Satzes ist bei den in \S.\,52 aufgezählten
 Fällen immer für einen der verschiedenen Werthe $\nu$
erfüllt, ausgenommen bei der cyklischen Gruppe und bei der
Diëdergruppe mit $m = 2$.

Endlich beweisen wir noch den folgenden Satz:

2.~\gesperrt{Sind $a, a'$ die Pole einer Substitution $\Theta$ von $G$,
so sind die mit $a$ und $a'$ conjugirten Pole}
\[
b = \chi(a), \quad b' = \chi(a'),
\]
\gesperrt{worin $\chi$ eine beliebige Substitution aus $G$ ist, die
beiden Pole einer und derselben Substitution,
nämlich der Substitution $\chi\Theta\chi^{-1}$.}

Die Richtigkeit ergiebt sich unmittelbar aus dem Anblick
der Gleichungen:
\[
\begin{aligned}
a &= \Theta(a), & \chi\Theta\chi^{-1}\chi(a) &= \chi\Theta(a) = \chi(a) \\
a' &= \Theta(a'), & \chi\Theta\chi^{-1}\chi(a') &= \chi\Theta(a') = \chi(a').
\end{aligned}
\]

\begin{center}
\S.\,54.

\gesperrt{Die Grundformen.}
\end{center}

Um zu der endgültigen Aufstellung aller endlichen Gruppen
$G$ zu gelangen, ist es nothwendig, auf die Invarianten der entsprechenden
 binären Substitutionsgruppen näher einzugehen.

Wir müssen daher neben den linearen gebrochenen Substitutionen


\weq{1}{y = \frac{\alpha x + \beta}{\gamma x + \delta},}

die wir im \S.\,49 mit $A'''$ bezeichnet haben, noch die binären
Substitutionen

\weq{2}{(y_1, y_2) = \begin{pmatrix}\alpha, & \beta \\ \gamma, & \delta\end{pmatrix}(x_1, x_2)}
%CONT>>


% ===== scan leaf 227, printed 206 =====
% !PRINTED 206
\weberpage{206}
%<<CONT
betrachten, die dort mit $A''$ bezeichnet waren, und die, wenn man
$y_1 : y_2 = y$, $x_1 : x_2 = x$ setzt, wieder auf die Substitution (1)
führen. In (2) ist immer

\weq{3}{\alpha\delta - \beta\gamma = 1}

vorausgesetzt, aber die beiden Substitutionen
\[
\begin{pmatrix}\alpha, & \beta \\ \gamma, & \delta\end{pmatrix}, \quad
\begin{pmatrix}-\alpha, & -\beta \\ -\gamma, & -\delta\end{pmatrix}
\]
werden als identisch angesehen.

Es kann nicht zu einem Missverständniss führen, wenn wir
beide Arten von Substitutionen übereinstimmend durch ein
Symbol wie
\[
\Theta = \begin{pmatrix}\alpha, & \beta \\ \gamma, & \delta\end{pmatrix}
\]
bezeichnen, da ja dann in beiden die Compositionsregel genau
dieselbe ist.

Wenn nun

\weq{4}{a, a_1, a_2, \ldots, a_{\mu-1}}

ein System conjugirter Pole der Gruppe $G$ ist, so ist

\weq{5}{f(x) = (x-a)(x-a_1)\cdots(x-a_{\mu-1})}

eine ganze Function $\mu^{\text{ten}}$ Grades, deren Wurzeln jene conjugirten
Pole sind. Diese Function $f(x)$ hat folgende Eigenschaft:

Nehmen wir irgend eine Substitution der Gruppe $G$
\[
\psi(x) = \frac{\alpha x + \beta}{\gamma x + \delta},
\]
so sind die Wurzeln der Function

\weq{6}{(\gamma x + \delta)^{\mu} f[\psi(x)]}

die Grössen $\psi^{-1}(a), \psi^{-1}(a_1), \ldots, \psi^{-1}(a_{\mu-1})$. Diese aber stimmen,
von der Reihenfolge abgesehen, nach \S.\,52, 5. mit den
Grössen $a, a_1, a_2, \ldots, a_{\mu-1}$ überein, und folglich können sich die
Gleichungen (5) und (6) nur durch einen constanten Factor
unterscheiden. Wenn wir also

\weq{7}{x_2^{\mu} f\left(\frac{x_1}{x_2}\right) = f(x_1, x_2)}

setzen, so bleibt diese Form $f$, von einem constanten Factor abgesehen,
ungeändert, wenn auf $(x_1, x_2)$ irgend eine Substitution
der Gruppe $G$ angewandt wird.

Nach der Definition \S.\,40, 2. ist also $f(x_1, x_2)$ eine
\gesperrt{Invariante der Gruppe} $G$, und wir haben den Satz:


% ===== scan leaf 228, printed 207 =====
% !PRINTED 207
\weberpage{207}

1.~\gesperrt{Jedem Systeme conjugirter Pole der Gruppe $G$
entspricht eine invariante Form, deren Grad
gleich der Anzahl der Pole des Systemes ist, und
deren Wurzeln eben diese Pole sind.}

Wir bezeichnen mit $f_1, f_2, \ldots$ die zu den verschiedenen
Systemen conjugirter Pole gehörigen invarianten Formen, die
von den Graden $\mu, \mu', \ldots$ sind und die wir die \gesperrt{Grundformen}
der Gruppe nennen wollen.

Ist dann $F(x_1, x_2)$ irgend eine invariante Form der
Gruppe $G$ und $f(x, y)$ eine Grundform, und hat $F(x, 1) = 0$
mit $f(x, 1) = 0$ eine gemeinsame Wurzel, so müssen alle Wurzeln
 von $f = 0$ zugleich Wurzeln von $F = 0$ sein. Denn
nach Voraussetzung haben die beiden Functionen $F(x, 1)$ und
$F[\psi(x), 1]$ dieselben Wurzeln. Wenn also $F(a, 1) = 0$ ist,
so ist auch $F[\psi(a), 1] = F(a_1, 1) = 0$. Wir haben also den
folgenden Satz:

2.~\gesperrt{Hat eine zu $G$ gehörige invariante Form $F(x_1, x_2)$
mit einer der Grundformen $f(x_1, x_2)$ einen Theiler
gemein, so ist $F(x_1, x_2)$ durch $f(x_1, x_2)$ theilbar.}

Ist $F(x_1, x_2)$ eine invariante Form der Gruppe $G$, und $\xi$
eine Wurzel der Gleichung $F(x, 1) = 0$, so sind auch, wenn
$\psi$ die Elemente der Gruppe $G$ durchläuft, die sämmtlichen
$n$ Grössen $\psi(\xi)$ Wurzeln derselben Gleichung. Wenn also der
Grad von $F$ niedriger ist, als der Grad der Gruppe, so können
diese Grössen nicht alle von einander verschieden sein, und es
folgt für irgend zwei von einander verschiedene Substitutionen
$\chi, \psi$ von $G$
\[
\chi(\xi) = \psi(\xi)
\]
oder
\[
\xi = \chi^{-1}\psi(\xi),
\]
d. h. $\xi$ muss unter den Polen der Gruppe $G$ vorkommen, und
es ist also $F$ nach dem Satze 2. durch eine der Grundformen
theilbar. Da wir auf den Quotienten der Division dieselbe Schlussweise
 anwenden können, so folgt:

3.~\gesperrt{Eine invariante Form $F$ der Gruppe $G$, deren
Grad niedriger ist, als der Grad der Gruppe,
ist, von einem constanten Factor abgesehen, ein
Product von Potenzen der Grundformen.}
%CONT>>


% ===== scan leaf 229, printed 208 =====
% !PRINTED 208
\weberpage{208}
%<<CONT
Es ist hierbei immer angenommen, dass die Co\"efficienten
von $F(x_1, x_2)$ nicht alle gleich Null sind. Wir k\"onnen also,
wenn wir von dieser Voraussetzung absehen, den Satz 3. auch
so aussprechen:

4.\ \gesperrt{Ist $F(x_1, x_2)$ eine invariante Form der Gruppe $G$
von niedrigerem Grade als $G$, die sich nicht als
Product aus den Grundformen darstellen l\"asst,
so muss $F(x_1, x_2)$ identisch verschwinden.}

\rule{4cm}{0.4pt}


% ===== scan leaf 230, printed 209 =====
\weberpage{209}% page number not printed on this page
\pdfbookmark[1]{Achter Abschnitt. Die Polyëdergruppen.}{bm64}
\pdfbookmark[2]{§ 55. Die cyklischen Gruppen und die Diëdergruppen}{bm65}
% !PRINTED none

\begin{center}
\gesperrt{Achter Abschnitt.}

\textbf{\gesperrt{Die Poly\"edergruppen.}}

\rule{3cm}{0.4pt}

\S.\ 55.
\end{center}

\begin{center}
\gesperrt{Die cyklischen Gruppen und die Di\"edergruppen.}
\end{center}

Wir gehen nun dazu über, die allgemeinen Principien zur
wirklichen Bildung der verschiedenen Poly\"edergruppen anzuwenden,
zunächst also zu zeigen, dass die in \S.\,52 als möglich
erkannten Arten dieser Gruppen alle existiren.

Bei den cyklischen Gruppen $n^{\text{ten}}$ Grades haben wir nur zwei
Systeme conjugirter Pole, deren jedes nur einen $(n-1)$fachen
Pol enthält. Diese beiden Pole müssen also die gemeinsamen
Pole aller Substitutionen der Gruppe sein und können nach
\S.\,51, 2. als $0$ und $\infty$ angenommen werden. Wir haben dann
nur die beiden linearen Grundformen
\weq{1}{f_1 = x_1, \quad f_2 = x_2.}
Alle Substitutionen der Gruppe sind multiplicativ, und der
Multiplicator muss eine $n^{\text{te}}$ Einheitswurzel sein. Sie müssen also
die Form haben:
\weq{2}{x, \varepsilon x, \varepsilon^2 x, \ldots, \varepsilon^{n-1} x,}
wenn $\varepsilon$ eine primitive $n^{\text{te}}$ Einheitswurzel ist.

Wir erhalten also in der That für jedes $n$ eine cyklische
Gruppe, die wir mit $C_n$ bezeichnen wollen:
\weq{3}{C_n = \begin{pmatrix} \varepsilon^r, & 0 \\ 0, & 1 \end{pmatrix}, \quad r = 0, 1, \ldots, n-1,}
oder mit der Determinante $1$ geschrieben:
\weq{4}{C_n = \begin{pmatrix} \varepsilon^{1/2\,r}, & 0 \\ 0, & \varepsilon^{-1/2\,r} \end{pmatrix}.}


% ===== scan leaf 231, printed 210 =====
% !PRINTED 210
\weberpage{210}
Bei der Di\"edergruppe vom Grade $n = 2m$, die wir mit $D_m$
bezeichnen wollen, haben wir nach dem vorigen Paragraphen
zwei Systeme von je $m$ conjugirten zweizähligen Polen und ein
System von zwei conjugirten $m$-zähligen Polen. Die letzteren
müssen also nach dem Satze \S.\,53 die Pole einer Substitution
sein. Wenn wir sie mit $0$ und $\infty$ zusammenfallen lassen, so erhalten
 wir die Grundform zweiten Grades:
\[
f_1 = x_1 x_2.
\]
Dies kann aber nur dann eine invariante Form der Gruppe
sein, wenn alle Substitutionen in einer der beiden Formen
\[
\begin{pmatrix} \alpha, & 0 \\ 0, & \delta \end{pmatrix}, \quad \begin{pmatrix} 0, & \beta \\ \gamma, & 0 \end{pmatrix}
\]
enthalten sind, und hierin müssen $\alpha:\delta$ und $-\beta:\gamma$ $m^{\text{te}}$ Einheitswurzeln
 sein. Wir bekommen also die Substitutionen der Gruppe,
wenn $\varepsilon$ eine primitive $m^{\text{te}}$ Einheitswurzel ist:
\[
x, \varepsilon x, \varepsilon^2 x, \ldots, \varepsilon^{m-1} x
\]
\[
\frac{1}{x}, \frac{\varepsilon}{x}, \frac{\varepsilon^2}{x}, \cdots, \frac{\varepsilon^{m-1}}{x},
\]
d.\,h.
\weq{5}{D_m = \begin{pmatrix} \varepsilon^r, & 0 \\ 0, & 1 \end{pmatrix}, \begin{pmatrix} 0, & \varepsilon^r \\ 1, & 0 \end{pmatrix},}
oder, mit der Determinante $+1$ dargestellt:
\weq{6}{D_m = \begin{pmatrix} \varepsilon^{1/2\,r}, & 0 \\ 0, & \varepsilon^{-1/2\,r} \end{pmatrix}, \begin{pmatrix} 0, & i\varepsilon^{1/2\,r} \\ i\varepsilon^{-1/2\,r}, & 0 \end{pmatrix}, \quad r = 0, 1, \ldots, m-1.}
Man sieht sofort, dass diese Substitutionen wirklich eine
Gruppe bilden.

Ausser $0$ und $\infty$ haben wir hier noch die aus den Gleichungen
\[
x = \frac{\varepsilon^h}{x}
\]
hervorgehenden Pole $\pm \varepsilon^{1/2\,h}$, aus denen man noch die beiden
Grundformen $m^{\text{ten}}$ Grades
\weq{7}{f_2 = x_1^m + x_2^m, \quad f_3 = x_1^m - x_2^m}
erhält.

Hieraus ersieht man, dass die Di\"edergruppen und die cyklischen
 Gruppen ganz von einander verschieden sind, da ihre
Grundformen verschiedene Grade haben.


% ===== scan leaf 232, printed 211 =====
% !PRINTED 211
\weberpage{211}
Durch Transformation mittelst der Substitution

\[
\begin{pmatrix} e^{\frac{\pi i}{4}}, & 0 \\ 0, & e^{-\frac{\pi i}{4}} \end{pmatrix}
\]

geht die Di\"edergruppe $D_m$ in

\[
\begin{pmatrix} \varepsilon^{1/2\,\lambda}, & 0 \\ 0, & \varepsilon^{-1/2\,\lambda} \end{pmatrix}, \quad
\begin{pmatrix} 0, & \varepsilon^{1/2\,\lambda} \\ -\varepsilon^{-1/2\,\lambda}, & 0 \end{pmatrix}
\]

oder in

\[
\begin{gathered}
x, \; \varepsilon x, \; \varepsilon^2 x, \; \ldots, \; \varepsilon^{m-1} x\\[4pt]
\dfrac{-1}{x}, \; \dfrac{-\varepsilon}{x}, \; \dfrac{-\varepsilon^2}{x}, \; \cdots, \; \dfrac{-\varepsilon^{m-1}}{x}
\end{gathered}
\]

\"uber, und die Grundformen f\"ur diese Darstellung sind:

\[
f_1 = x_1 x_2, \quad f_2 = x_1^m + i^m x_2^m, \quad f_3 = x_1^m - i^m x_2^m.
\]

Die Di\"edergruppe $D_m$ enth\"alt als Theiler die cyklische
Gruppe $C_m$, und zwar ist $C_m$ ein \gesperrt{Normaltheiler} von $D_m$.

Durch Transformation mit irgend einer multiplicativen
Substitution

\[
\begin{pmatrix} \alpha, & 0 \\ 0, & \delta \end{pmatrix}
\]

\"andert sich $C_m$ nicht, w\"ahrend die nicht in $C_m$ enthaltenen Substitutionen
 von $D_m$ ihre Form \"andern; denn es ist:

\[
\begin{gathered}
\begin{pmatrix} \delta, & 0 \\ 0, & \alpha \end{pmatrix}
\begin{pmatrix} \varepsilon, & 0 \\ 0, & 1 \end{pmatrix}
\begin{pmatrix} \alpha, & 0 \\ 0, & \delta \end{pmatrix}
=
\begin{pmatrix} \varepsilon, & 0 \\ 0, & 1 \end{pmatrix},\\[6pt]
\begin{pmatrix} \delta, & 0 \\ 0, & \alpha \end{pmatrix}
\begin{pmatrix} 0, & \varepsilon \\ 1, & 0 \end{pmatrix}
\begin{pmatrix} \alpha, & 0 \\ 0, & \delta \end{pmatrix}
=
\begin{pmatrix} 0, & \varepsilon \delta^2 \\ \alpha^2, & 0 \end{pmatrix}.
\end{gathered}
\]

Abgesehen von einer solchen multiplicativen Substitution ist
aber die Gruppe $D_m$ durch die darin enthaltene Gruppe $C_m$
vollkommen bestimmt. Denn setzen wir

\[
D_m = C_m + \varphi C_m,
\]
worin

\[
\varphi = \begin{pmatrix} \alpha, & \beta \\ \gamma, & \delta \end{pmatrix}
\]

gesetzt ist, so muss, da $C_m$ Normaltheiler von $D_m$ sein muss,
$\varphi C_m = C_m \varphi$ sein, d.\,h. es muss sich zu jedem Exponenten $r$
ein Exponent $s$, und umgekehrt, bestimmen lassen, so dass

\[
\begin{pmatrix} \alpha, & \beta \\ \gamma, & \delta \end{pmatrix}
\begin{pmatrix} \varepsilon^r, & 0 \\ 0, & 1 \end{pmatrix}
=
\begin{pmatrix} \varepsilon^s, & 0 \\ 0, & 1 \end{pmatrix}
\begin{pmatrix} \alpha, & \beta \\ \gamma, & \delta \end{pmatrix}
\]
%CONT>>


% ===== scan leaf 233, printed 212 =====
% !PRINTED 212
\weberpage{212}
\pdfbookmark[2]{§ 56. Die Tetraëdergruppe}{bm66}
%<<CONT
oder
\[
\begin{pmatrix} \alpha\,\varepsilon^r, & \beta \\ \gamma\,\varepsilon^r, & \delta \end{pmatrix} = \begin{pmatrix} \varepsilon^s\alpha, & \varepsilon^s\beta \\ \gamma, & \delta \end{pmatrix}.
\]

Wären $\beta$ und $\gamma = 0$, so wäre $\varphi$ selbst von der Form $\varepsilon x$,
und $\varepsilon$ eine $2m^{\text{te}}$ Einheitswurzel, und $D_m$ wäre also eine cyklische
Gruppe und keine Diëdergruppe. Ist aber $\beta$ oder $\gamma$ von Null
verschieden, so folgt $\varepsilon^s = \varepsilon^{-r}$ und $\alpha = \delta = 0$, also
\[
\varphi = \begin{pmatrix} 0, & \beta \\ \gamma, & 0 \end{pmatrix},
\]
was durch eine multiplicative Transformation auf die Form
$\begin{pmatrix} 0, & 1 \\ 1, & 0 \end{pmatrix}$ gebracht werden kann.

Unter den Diëdergruppen ist die Gruppe $D_2$ besonders
hervorzuheben, in der drei Systeme von je zwei zweizähligen
Polen vorhanden sind. Sie besteht aus den vier Substitutionen
\[
\begin{pmatrix} 1, & 0 \\ 0, & 1 \end{pmatrix},\quad
\begin{pmatrix} -1, & 0 \\ 0, & 1 \end{pmatrix},\quad
\begin{pmatrix} 0, & 1 \\ 1, & 0 \end{pmatrix},\quad
\begin{pmatrix} 0, & 1 \\ -1, & 0 \end{pmatrix}
\]
und wird die \gesperrt{Vierergruppe} genannt.

\begin{center}
\S.\,56. \\[6pt]
\gesperrt{Die Tetraëdergruppe.}
\end{center}

Bei der Tetraëdergruppe haben wir nach \S.\,52, III. ein
System von sechs conjugirten zweizähligen Polen. Wir können
also nach \S.\,53 annehmen, dass in der Gruppe die beiden Substitutionen

\[
\Theta(x) = -x, \quad \psi(x) = \frac{1}{x}
\]
vorkommen, und folglich auch
\[
\psi_1(x) = \psi\,\Theta(x) = -\frac{1}{x}.
\]

Die beiden Substitutionen $\psi$ und $\psi_1$ geben die Pole $\pm 1$
und $\pm i$, und diese müssen, da $\psi$ und $\psi_1$ vom $2^{\text{ten}}$ Grade sind
und in $G$ überhaupt nur zwei- oder dreizählige Pole vorkommen,
zu den zweizähligen gehören (\S.\,52, 3.). Die Substitution $\Theta(x)$
giebt die zweizähligen Pole $0, \infty$. Demnach lautet die Gleichung,
von der die zweizähligen Pole abhängen, $x(x^4-1) = 0$ und wir
erhalten die erste Grundform $6^{\text{ten}}$ Grades:
\weq{1}{f(x_1, x_2) = x_1 x_2 (x_1^4 - x_2^4).}


% ===== scan leaf 234, printed 213 =====
% !PRINTED 213
\weberpage{213}
Um die übrigen Substitutionen der Gruppe zu finden, bemerken
 wir, dass wir nach dem Theorem \S.\,53, 2., wenn wir $\pm 1$
für $a, a'$ und $0, \infty$ für $b, b'$ nehmen, auf die Existenz einer Substitution
 $\chi$ in $G$ schliessen können, die den Bedingungen genügt
\[
\chi(-1) = 0, \quad \chi(+1) = \infty,
\]
und dass also $\chi$ die Form haben muss:
\weq{2}{\chi = \lambda\,\frac{x+1}{x-1},}
wo $\lambda$ ein constanter Factor ist.

Wenn wir in $\chi$ die sechs conjugirten Pole $0, \infty, +1, -1,$
$+i, -i$ einsetzen, so müssen wir dieselben Werthe in einer
anderen Reihenfolge erhalten (\S.\,52, 5.). Diese Werthe sind aber
\[
-\lambda,\ \lambda,\ \infty,\ 0,\ -\lambda i,\ \lambda i.
\]

Es muss also $\lambda = \pm 1$ oder $= \pm i$ sein. Der Werth $\pm 1$
ist nicht zulässig, weil sonst
\[
\chi(i) = \mp i, \quad \Theta\chi(i) = \pm i,
\]
und also $\pm i$ die Pole von $\chi$ oder von $\Theta\chi$ wären, während sie
doch nur zweizählig und die Pole von $\Theta\psi$ sind. Es muss also
$\lambda = \pm i$ sein, und wir können ohne Beschränkung der Allgemeinheit
 das obere Zeichen nehmen, denn das untere Zeichen
entspricht dann der Substitution $\Theta\chi$. Es ist daher
\[
\chi(x) = i\,\frac{x+1}{x-1}.
\]

Daraus erhält man
\[
\chi^2(x) = \frac{x+i}{x-i}, \quad \chi^3(x) = x,
\]
und es ergeben sich die 12 Substitutionen:
\weq{3}{\begin{array}{cccc}
1, & \Theta, & \psi, & \Theta\psi \\
\chi, & \Theta\chi, & \psi\chi, & \Theta\psi\chi \\
\chi^2, & \Theta\chi^2, & \psi\chi^2, & \Theta\psi\chi^2,
\end{array}}
oder
\weq{4}{\pm x,\ \pm\frac{1}{x},\ \pm i\,\frac{x+1}{x-1},\ \pm i\,\frac{x-1}{x+1},\ \pm\frac{x+i}{x-i},\ \pm\frac{x-i}{x+i}.}

Dass diese Substitutionen wirklich eine Gruppe, die \gesperrt{Tetraëdergruppe},
bilden, von der $\Theta, \psi, \chi$ die erzeugenden Elemente
sind, ergiebt sich aus
\weq{5}{\chi\Theta = \Theta\psi\chi,\ \chi^2\Theta = \psi\chi^2,\ \psi\Theta = \Theta\psi,\ \chi\psi = \Theta\chi,\ \chi^2\psi = \Theta\psi\chi^2,}
%CONT>>


% ===== scan leaf 235, printed 214 =====
% !PRINTED 214
\weberpage{214}
%<<CONT
die sich mit Rücksicht auf $\Theta^2 = \psi^2 = \chi^3 = 1$ aus den drei
Relationen
\[
\psi\Theta = \Theta\psi, \quad \chi\Theta = \Theta\psi\chi, \quad \chi\psi = \Theta\chi
\]
ableiten lassen.

Aus (5) geht noch hervor, dass die Vierergruppe
\weq{6}{1,\ \Theta,\ \psi,\ \Theta\psi}
ein \gesperrt{Normaltheiler} der Tetraëdergruppe ist.

Die Substitutionen der Tetraëdergruppe können in jeder der
beiden Formen
\[
\chi^\lambda \psi^\mu \Theta^\nu, \quad \Theta^\nu \psi^\mu \chi^\lambda
\]
\[
\lambda = 0, 1, 2, \quad \mu = 0, 1, \quad \nu = 0, 1
\]
dargestellt werden.

Um die zu den dreizähligen Polen gehörigen beiden Grundformen
 vierter Ordnung $\Phi_1, \Phi_2$ zu erhalten, kann man entweder
aus (5) die noch fehlenden Pole berechnen, oder man kann so
verfahren:

Da die Grundformen $\Phi_1, \Phi_2$ die Substitutionen $\Theta$ und $\psi$
gestatten müssen und nicht durch $x_1$ und $x_2$ theilbar sein können,
so können sie nur von folgender Form sein:
\[
\begin{aligned}
\Phi_1 &= x_1^4 + m_1 x_1^2 x_2^2 + x_2^4 \\
\Phi_2 &= x_1^4 + m_2 x_1^2 x_2^2 + x_2^4,
\end{aligned}
\]
worin $m_1, m_2$ noch zu bestimmende Constanten sind. Bildet man
aber die \gesperrt{Hesse'sche} Covariante von $\Phi_1$, so erhält man nach
\S.\,40, 4. eine Invariante der Gruppe:
\[
\frac{1}{12}\left[\frac{\partial^2 \Phi_1}{\partial x_1^2}\,\frac{\partial^2 \Phi_1}{\partial x_2^2} - \left(\frac{\partial^2 \Phi_1}{\partial x_1\,\partial x_2}\right)^2\right] = 2m_1(x_1^4+x_2^4) + (12-m_1^2)\,x_1^2 x_2^2.
\]

Diese Form kann, da keine anderen Grundformen vierter
Ordnung existiren, nur entweder mit $\Phi_1$ oder mit $\Phi_2$ übereinstimmen
 (von einem constanten Factor abgesehen). Es muss also
\[
\frac{12-m_1^2}{2m_1} = m_1 \ \text{oder}\ m_2
\]
sein. Die erste Annahme giebt aber $m_1 = \pm 2$. Dann wäre $\Phi_1$
ein Quadrat, was nicht sein kann. Es bleibt also nur:
\[
m_1^2 + 2m_1 m_2 = 12,
\]
und ebenso aus $\Phi_2$:
\[
m_2^2 + 2m_1 m_2 = 12,
\]
also
\[
m_1 = -m_2 = \pm 2\sqrt{-3},
\]
%CONT>>


% ===== scan leaf 236, printed 215 =====
% !PRINTED 215
\weberpage{215}
%<<CONT
und die beiden Grundformen werden:
\weq{7}{\begin{aligned}
\Phi_1 &= x_1^4 + 2\sqrt{-3}\,x_1^2 x_2^2 + x_2^4 \\
\Phi_2 &= x_1^4 - 2\sqrt{-3}\,x_1^2 x_2^2 + x_2^4.
\end{aligned}}

Zwischen den drei Grundformen $f, \Phi_1, \Phi_2$ kann man eine
Relation herleiten, wenn man $x_1, x_2$ eliminirt. Man findet aus (7)
\[
2(x_1^4+x_2^4) = \Phi_1 + \Phi_2
\]
\[
4\sqrt{-3}\,x_1^2 x_2^2 = \Phi_1 - \Phi_2,
\]
und aus (1)
\[
f^2 = x_1^2 x_2^2 (x_1^4+x_2^4)^2 - 4x_1^6 x_2^6,
\]
woraus man leicht berechnet
\weq{8}{12\sqrt{-3}\,f^2 = \Phi_1^3 - \Phi_2^3.}

Die Substitutionen $\Theta, \psi, \chi$ erhalten, wenn sie mit der Determinante
 1 dargestellt werden, den Ausdruck:
\weq{9}{\begin{pmatrix} i, & 0 \\ 0, & -i \end{pmatrix},\quad
\begin{pmatrix} 0, & i \\ i, & 0 \end{pmatrix},\quad
\begin{pmatrix} \dfrac{1-i}{2}, & \dfrac{1-i}{2} \\[6pt] -\dfrac{1+i}{2}, & \dfrac{1+i}{2} \end{pmatrix},}
woraus nach den im \S.\,50 aufgestellten Bedingungen folgt, dass
die entsprechende Gruppe orthogonaler ternärer Substitutionen
reell ist.

Wendet man die Substitutionen (9) mit der Determinante 1
auf die Grundform $f(x_1, x_2)$ an, so ergiebt eine sehr einfache
\gesperrt{Rechnung, dass die Grundform $f$ eine absolute Invariante}
der binären Gruppe $G$ ist.

Dieselbe Eigenschaft hat auch die \gesperrt{Hesse'sche} Determinante
von $f$:
\[
\begin{aligned}
H &= -\frac{1}{25}\left[f''(x_1,x_1)\,f''(x_2,x_2) - f''(x_1,x_2)^2\right] \\
&= (x_1^4+x_2^4)^2 + 12\,x_1^4 x_2^4,
\end{aligned}
\]
die nichts Anderes ist, als das Product der beiden Functionen
$\Phi_1, \Phi_2$; während die Functionen $\Phi_1$ und $\Phi_2$ selbst bei der dritten
der Substitutionen (9) eine dritte Einheitswurzel als Factor
annehmen.


% ===== scan leaf 237, printed 216 =====
% !PRINTED 216
\weberpage{216}
\pdfbookmark[2]{§ 57. Die Octaëdergruppe}{bm67}

\begin{center}
\S.\,57.

\gesperrt{Die Octaëdergruppe.}
\end{center}

Bei der Octaëdergruppe kommt ein System von sechs conjugirten
vierzähligen Polen vor. Wir haben demnach eine Substitution
4\textsuperscript{ten} Grades, deren Periode
\[
1, \ \Theta, \ \Theta^2, \ \Theta^3, \ \Theta^4 = 1
\]
ein Paar dieser conjugirten vierzähligen Pole giebt. Nach \S.\,53, 1.
können wir in der Gruppe die Substitutionen annehmen:
\weq{1}{\Theta(x) = ix, \quad \psi(x) = \frac{1}{x},}
und es ist
\weq{2}{\psi^2 = 1, \quad \Theta\psi = \psi\Theta^3, \quad \Theta^2\psi = \psi\Theta^2, \quad \Theta^3\psi = \psi\Theta.}
Die zu dem Systeme der vierzähligen Pole gehörige Grundform
6\textsuperscript{ten} Grades muss bis auf einen constanten Factor
ungeändert bleiben durch die Substitutionen $\Theta$ und $\psi$, und
sie muss ausserdem den Factor $x_1x_2$ enthalten, also von einer der
beiden Formen sein
\[
x_1x_2(x_1^4 - x_2^4), \quad x_1x_2(x_1^4 + x_2^4).
\]
Es ist gleichgültig, welche der beiden Annahmen wir verfolgen, da
die eine durch die Substitution $ix_2$ für $x_2$, was dem Uebergange
zu einer transformirten Gruppe entspricht, in die andere übergeht.
Nehmen wir
\weq{3}{f(x_1, x_2) = x_1x_2(x_1^4 - x_2^4)}
als Grundform 6\textsuperscript{ten} Grades an, so sind die sechs
vierzähligen Pole
\weq{4}{0, \ \infty, \ 1, \ -1, \ i, \ -i.}
Es muss nun nach \S.\,53, 2. in $G$ eine Substitution $\chi$ geben,
so dass
\[
\chi(-1) = 0, \quad \chi(+1) = \infty
\]
ist, und wenn man den constanten Factor wie bei der Tetraëdergruppe
bestimmt [\S.\,56, (2)], so findet sich
\weq{5}{\chi(x) = i\,\frac{x+1}{x-1}.}
Hieraus erhalten wir
\[
\chi^2(x) = \frac{x+i}{x-i}, \quad \chi^3(x) = x,
\]
%CONT>>


% ===== scan leaf 238, printed 217 =====
% !PRINTED 217
\weberpage{217}

%<<CONT
und die Relationen
\weq{6}{\Theta\chi = \chi^2\psi\Theta^3, \quad \Theta^2\chi = \chi\psi, \quad \Theta^3\chi = \chi^2\Theta \\
\psi\chi = \chi\psi\Theta^2, \quad \psi\chi^2 = \chi^2\Theta^2.}
Diese Relationen in Verbindung mit (2) zeigen, dass die
Substitutionen
\weq{7}{\chi^{\lambda}\psi^{\mu}\Theta^{\nu}, \quad \lambda = 0, 1, 2; \quad \mu = 0, 1; \quad \nu = 0, 1, 2, 3}
in der That eine Gruppe bilden, weil man mit ihrer Hülfe jede
Substitution der Form
\[
\chi^{\lambda}\psi^{\mu}\Theta^{\nu}\chi, \quad \chi^{\lambda}\psi^{\mu}\Theta^{\nu}\psi,
\]
und folglich durch Wiederholung auch jede Substitution
\[
\chi^{\lambda}\psi^{\mu}\Theta^{\nu}\,\chi^{\lambda'}\psi^{\mu'}\Theta^{\nu'}
\]
in die Form (7) bringen kann. Die Gruppe kann explicite in
der Form
\weq{8}{i^{\nu}x, \quad \frac{i^{\nu}}{x}, \quad i^{\nu}\,\frac{x - i^{\nu'}}{x + i^{\nu'}}, \quad \nu, \nu' = 0, 1, 2, 3}
dargestellt werden und ist vom 24\textsuperscript{sten} Grade. Es ist die
\gesperrt{Octaëdergruppe}.

Die Gruppe hat einen Normaltheiler 12\textsuperscript{ten} Grades, den man
in der Form $\chi^{\lambda}\psi^{\mu}\Theta^{2\nu}$ darstellen kann und der eine
Tetraëdergruppe ist.

Die Darstellung wird in gewisser Beziehung einfacher, wenn
man an Stelle von $\psi$ ein neues Element
\weq{9}{\omega = \psi\Theta = \frac{-i}{x}}
einführt, was ebenso wie $\psi$ von der zweiten Ordnung ist. Dann
kann die Gruppe dargestellt werden durch
\[
\chi^{\lambda}\omega^{\mu}\Theta^{\nu},
\]
und an Stelle der Relationen (2) und (6) treten die folgenden:
\[
\omega\chi = \chi^2\omega, \quad \omega\chi^2 = \chi\omega
\]
\weq{10}{\Theta\chi = \chi^2\omega\Theta^2, \; \Theta^2\chi = \chi\omega\Theta^3, \; \Theta^3\chi = \chi^2\Theta, \; \Theta^2\chi^2 = \chi^2\omega\Theta \\
\Theta\omega = \omega\Theta^3, \quad \Theta^2\omega = \omega\Theta^2, \quad \Theta^3\omega = \omega\Theta.}
Alle diese Relationen aber ergeben sich als Folgerungen
aus den vieren:
\weq{11}{\omega\chi = \chi^2\omega, \quad \Theta\omega = \omega\Theta^3, \quad \Theta\chi = \chi^2\omega\Theta^2, \quad \Theta^2\chi = \chi\omega\Theta^3,}
in Verbindung mit den die Grade ausdrückenden Formeln:
\weq{12}{\chi^3 = 1, \quad \omega^2 = 1, \quad \Theta^4 = 1.}


% ===== scan leaf 239, printed 218 =====
% !PRINTED 218
\weberpage{218}

Es folgt nämlich zunächst aus der zweiten der Relationen (11):
\[
\Theta^2\omega = \Theta\omega\Theta^3 = \omega\Theta^2, \quad \Theta^3\omega = \Theta\omega\Theta^2 = \omega\Theta,
\]
ferner aus der letzten (11):
\[
\Theta^3\chi = \Theta\chi\omega\Theta^3 = \chi^2\omega\Theta^2\omega\Theta^3 = \chi^2\Theta,
\]
und weiter:
\[
\omega\chi^2 = \chi^2\omega\chi = \chi\omega,
\]
\[
\Theta^2\chi^2 = \chi\omega\Theta^3\chi = \chi\omega\chi^2\Theta = \chi^2\omega\Theta.
\]
Wenn wir nun irgend drei Elemente $\chi$, $\omega$, $\Theta$ haben, die
sich nach irgend einer Regel componiren lassen, wenn dabei $\chi$ vom
dritten, $\omega$ vom zweiten, $\Theta$ vom vierten Grade ist, so folgt aus
dem Bestehen der Relationen (11), dass die Elemente
$\chi^{\lambda}\omega^{\mu}\Theta^{\nu}$ eine Gruppe 24\textsuperscript{sten} Grades bilden, die mit
der Octaëdergruppe isomorph ist. Dazu ist nur noch nachzuweisen,
dass aus diesen Voraussetzungen folgt, dass die 24 Elemente
$\chi^{\lambda}\omega^{\mu}\Theta^{\nu}$ alle von einander verschieden sind. Nehmen wir
an, es sei
\[
\chi^{\lambda}\omega^{\mu}\Theta^{\nu} = \chi^{\lambda'}\omega^{\mu'}\Theta^{\nu'},
\]
so würde folgen:
\[
\chi^{\lambda - \lambda'} = \omega^{\mu'}\Theta^{\nu' - \nu}\omega^{-\mu},
\]
und nach (11):
\[
\chi^{\lambda - \lambda'} = \omega^{\mu' - \mu}\Theta^{\pm(\nu - \nu')}.
\]
Nun folgt aber aus (11), dass $\omega\Theta$, $\omega\Theta^2$, $\omega\Theta^3$ vom
zweiten Grade sind, und da $\chi$ vom dritten Grade ist, so kann diese
Beziehung nur stattfinden, wenn $\lambda - \lambda'$ durch 3, $\mu - \mu'$ durch 2 und
$\nu - \nu'$ durch 4 theilbar, also $\chi^{\lambda} = \chi^{\lambda'}$, $\omega^{\mu} = \omega^{\mu'}$, $\Theta^{\nu} = \Theta^{\nu'}$ ist.

Nach der aus (10) folgenden Relation
\[
\omega = \chi\Theta\chi\Theta^2
\]
können wir $\omega$ aus $\chi$ und $\Theta$ zusammensetzen und daher $\chi$ und
$\Theta$ als \gesperrt{erzeugende Substitutionen} der Gruppe ansehen.

Die noch fehlenden beiden Grundformen achten und zwölften
Grades findet man sehr leicht, wenn man zunächst die Hesse'sche
Covariante von $f$ bildet.

Man erhält so die Grundform achten Grades:
\weq{13}{W = x_1^8 + 14\,x_1^4x_2^4 + x_2^8,}
und wenn man aus $f$ und $W$ die Functionaldeterminante bildet,
die ja wieder eine Covariante von $f$ ist (Bd. I, \S.\,59, 60), so
ergiebt sich die Grundform zwölften Grades:
\weq{14}{K = x_1^{12} - 33\,x_1^8x_2^4 - 33\,x_1^4x_2^8 + x_2^{12}.}


% ===== scan leaf 240, printed 219 =====
% !PRINTED 219
\weberpage{219}

Die Wurzeln von $f$ entsprechen den sechs Octaëderecken oder
den sechs Würfelflächen; die Wurzeln von $W$ den acht Octaëderflächen
oder Würfelecken, und die Wurzeln von $K$ sowohl beim Octaëder
als beim Würfel den 12 Kanten (vgl. \S.\,48).

Wenn man die Formen $W$, $K$ so darstellt:
\[
W = (x_1^4 + x_2^4)^2 + 12\,x_1^4x_2^4, \quad K = (x_1^4 + x_2^4)^3 - 36\,x_1^4x_2^4(x_1^4 + x_2^4),
\]
so lässt sich leicht die zwischen den drei Functionen $f$, $W$, $K$
bestehende identische Relation ableiten:
\weq{15}{W^3 - K^2 = 108\,f^4.}
Die Octaëdergruppe enthält, wie man sieht, die Tetraëdergruppe
und entsteht aus ihr durch Hinzunahme der einen Substitution
 $\Theta(x) = ix$. Die Grundformen des Octaëders sind also
zugleich invariante Formen des Tetraëders, und in der That
stimmen die Formen $f$ des Tetraëders und Octaëders mit einander
 genau überein, und es ist $W = \Phi_1\Phi_2$,
\[
K = \tfrac{1}{2}(\Phi_1^3 + \Phi_2^3), \quad (\S.\,56).
\]
Mit der Determinante 1 geschrieben, lauten die beiden erzeugenden
 Substitutionen der Octaëdergruppe so:
\[
\Theta = \begin{pmatrix} \sqrt{i}, & 0 \\ 0, & \dfrac{1}{\sqrt{i}} \end{pmatrix}, \quad
\chi = \begin{pmatrix} \dfrac{1}{\sqrt{2i}}, & \dfrac{1}{\sqrt{2i}} \\[6pt] \dfrac{1}{i\sqrt{2i}}, & \dfrac{-1}{i\sqrt{2i}} \end{pmatrix},
\]
worin $\sqrt{2i} = 1 + i$ zu setzen ist.

Durch die Substitution $\Theta$ \gesperrt{ändert} nun die Form $f(x_1, x_2)$,
wie die Formel (3) unmittelbar zeigt, ihr \gesperrt{Vorzeichen}. Durch
Anwendung von $\chi$ werden die linearen Factoren von $f$ folgendermaassen
 verändert:
\[
\begin{gathered}
x_1, \quad x_2, \quad x_1 + x_2, \quad x_1 - x_2, \quad x_1 + ix_2, \quad x_1 - ix_2 \\
\frac{x_1 + x_2}{\sqrt{2i}}, \quad \frac{x_1 - x_2}{i\sqrt{2i}}, \quad -i(x_1 + ix_2), \quad (x_1 - ix_2), \quad \frac{2x_1}{\sqrt{2i}}, \quad \frac{2x_2}{\sqrt{2i}},
\end{gathered}
\]
und daraus geht hervor, dass $f(x_1, x_2)$ durch Anwendung der
Substitution $\chi$ \gesperrt{ungeändert bleibt}, und dadurch sind die
Aenderungen der Form $f$ durch alle anderen Octaëdersubstitutionen
 zugleich mit bestimmt.

Die \gesperrt{Hesse}'sche Covariante $W$ von $f$ bleibt ungeändert,
wenn $f$ in $-f$ verwandelt wird, und folglich bleibt $W$ bei den
%CONT>>


% ===== scan leaf 241, printed 220 =====
% !PRINTED 220
\weberpage{220}
\pdfbookmark[2]{§ 58. Die Ikosaëdergruppe}{bm68}
%<<CONT
Octa\"edersubstitutionen absolut unge\"andert, w\"ahrend $K$ wieder
die gleichen Vorzeichen\"anderungen wie $f$ erleidet.

\begin{center}
\S.\,58.

\gesperrt{Die Ikosa\"edergruppe.}
\end{center}

Da wir bei der Ikosa\"edergruppe nur ein System conjugirter
f\"unfz\"ahliger Pole haben, so k\"onnen wir (\S.\,53, 1.) in dieser Gruppe
die beiden Substitutionen

\weq{1}{\Theta(x) = \varepsilon x, \quad \psi(x) = \frac{-1}{x}}

annehmen, worin $\varepsilon$ eine primitive f\"unfte Einheitswurzel bedeutet.
Wir nehmen hier, was freisteht, die Substitutionen $\psi$ in der
Form $-1:x$ an, weil dadurch die Formeln einfacher werden.
Die zu dem Systeme der f\"unfz\"ahligen Pole geh\"orige Grundform
12\textsuperscript{ten} Grades muss, da sie die Substitutionen (1) gestattet, von
der Form sein:

\weq{2}{f(x_1,x_2) = x_1x_2(x_1^{10} + m\,x_1^5x_2^5 - x_2^{10}),}

und es handelt sich noch um die Bestimmung des constanten
Factors $m$. Die beiden anderen Grundformen sind vom 20\textsuperscript{sten}
und 30\textsuperscript{sten} Grade. Wir k\"onnen nun $m$ nach dem Satze \S.\,54, 4.
bestimmen, wenn wir eine Covariante von $f(x)$ bilden k\"onnen,
deren Grad niedriger als 60 ist, und die sich nicht als ein Product
 aus Potenzen von drei Functionen 12\textsuperscript{ten}, 20\textsuperscript{sten}, 30\textsuperscript{sten}
Grades darstellen l\"asst, die nach dem erw\"ahnten Satze identisch
Null sein muss.

Nun l\"asst sich leicht eine Covariante 16\textsuperscript{ten} Grades von der
Form $f$ bilden, wenn wir nach Bd.\ I, \S.\,60 die vierte Polare der
Form $f(x_1,x_2)$ nehmen:

\[
12\,.\,11\,.\,10\,.\,9\;P_4(x,\xi) = u_0\xi_1^4 + u_1\xi_1^3\xi_2 + u_2\xi_1^2\xi_2^2 + u_3\xi_1\xi_2^3 + u_4\xi_2^4,
\]

worin

\[
u_0 = \frac{\partial^4 f}{\partial x_1^4}, \quad u_1 = 4\,\frac{\partial^4 f}{\partial x_1^3\,\partial x_2}, \quad u_2 = 6\,\frac{\partial^4 f}{\partial x_1^2\,\partial x_2^2},
\]
\[
u_3 = 4\,\frac{\partial^4 f}{\partial x_1\,\partial x_2^3}, \quad u_4 = \frac{\partial^4 f}{\partial x_2^4}.
\]

Daraus erhalten wir eine Covariante 16\textsuperscript{ten} Grades von $f$ als
%CONT>>


% ===== scan leaf 242, printed 221 =====
% !PRINTED 221
\weberpage{221}
%<<CONT
erste Invariante der in Bezug auf die Variablen $\xi_1, \xi_2$ biquadratischen
 Form, n\"amlich (Bd.\ I, \S.\,64):

\weq{3}{u_2^2 - 3\,u_1u_3 + 12\,u_0u_4.}

Da aber eine Form 16\,ten Grades sich nicht als Product von
Formen 12\,ten, 20\,sten und 30\,sten Grades darstellen lassen kann, so
muss diese Covariante identisch verschwinden. Nun ist hier

\[
u_0 = 11\,.\,10\,.\,9\,.\,8\,x_1^7x_2 + 6\,.\,5\,.\,4\,.\,3\,m\,x_1^2x_2^6,
\]
\[
u_1 = 4\,.\,11\,.\,10\,.\,9\,x_1^8 + 4\,.\,6\,.\,5\,.\,4\,.\,6\,m\,x_1^3x_2^5,
\]
\[
u_2 = 6^3\,.\,5^2\,m\,x_1^4x_2^4,
\]
\[
u_3 = -4\,.\,11\,.\,10\,.\,9\,x_2^8 + 4\,.\,6\,.\,5\,.\,4\,.\,6\,m\,x_2^3x_1^5,
\]
\[
u_4 = -11\,.\,10\,.\,9\,.\,8\,x_2^7x_1 + 6\,.\,5\,.\,4\,.\,3\,m\,x_2^2x_1^6,
\]

und wenn wir daraus die Covariante (3) bilden und den Co\"efficienten
 von $x_1^8x_2^8$ gleich 0 setzen, so ergiebt sich $m = \pm 11$.
Beide Zahlen sind hier zul\"assig. Die eine Annahme wird auf
die andere zur\"uckgef\"uhrt durch die Vertauschung von $x_1$ mit
$-x_1$, also durch eine Transformation der Gruppe. Wir setzen
demnach:

\weq{4}{f(x_1,x_2) = x_1x_2(x_1^{10} + 11\,x_1^5x_2^5 - x_2^{10}).}

Die zehn noch fehlenden f\"unfz\"ahligen Pole erh\"alt man also
durch Aufl\"osung der Gleichung

\[
x^{10} + 11\,x^5 - 1 = 0,
\]

aus der man

\[
x^5 = \frac{-11 \pm 5\sqrt5}{2} = \left(\frac{-1 \pm \sqrt5}{2}\right)^5
\]

findet. Setzen wir demnach

\weq{5}{\begin{aligned}
\omega &= \varepsilon + \varepsilon^4 = 2\cos\frac{2\pi}{5} = \frac{-1+\sqrt5}{2} \\
\omega' &= \varepsilon^2 + \varepsilon^3 = 2\cos\frac{4\pi}{5} = \frac{-1-\sqrt5}{2},
\end{aligned}}

so dass

\weq{6}{\omega^2 + \omega = 1, \quad \omega + \omega' = -1, \quad \omega\omega' = -1}

ist, so sind die f\"unfz\"ahligen Pole ausser 0 und $\infty$:

\weq{7}{\xi = \varepsilon^\lambda\omega, \quad \varepsilon^\lambda\omega' = \frac{-\varepsilon^\lambda}{\omega}, \quad \lambda = 0, 1, 2, 3, 4.}

Nun wenden wir den Satz \S.\,53, 2.\ an, nach dem, wenn $a, a'$
zwei Pole einer Substitution $\Theta$ sind, und $b$ ein zu $a$ conjugirter
Pol ist, eine Substitution $\chi$ in der Gruppe existiren muss, so
%CONT>>


% ===== scan leaf 243, printed 222 =====
% !PRINTED 222
\weberpage{222}
%<<CONT
dass $b = \chi(a)$, $b' = \chi(a')$ die beiden Pole der Substitution
$\chi^{-1}\Theta\chi$ sind. % [sic/mathematics: with b=chi(a), theorem 53.2 requires chi Theta chi^{-1}; the source prints the inverse conjugation order and the ensuing calculation follows that printed order.]
Darin k\"onnen wir $\infty$ und 0 f\"ur $a$ und $a'$ nehmen. % [sic: Weber prints a full stop before the lower-case ``und''; a comma is intended.]
und $\omega$ f\"ur $b$. Dann wird $b'$ ein noch n\"aher zu bestimmender
Pol $\xi$ aus der Reihe (7), und f\"ur $\chi$ erhalten wir die Form

\[
\chi(x) = \frac{\alpha x + \beta}{\gamma x + \delta} = \begin{pmatrix} \alpha, & \beta \\ \gamma, & \delta \end{pmatrix},
\]

worin

\weq{8}{\omega = \frac{\alpha}{\gamma}, \quad \xi = \frac{\beta}{\delta}}

zu setzen ist.
Nun k\"onnen wir f\"ur $\Theta$ jede der Substitutionen $\varepsilon^h x$ nehmen,
wenn nur $h$ nicht durch 5 theilbar ist, und erhalten so die Substitution


\[
\begin{aligned}
\chi^{-1}\Theta\chi &= \begin{pmatrix} \delta, & -\beta \\ -\gamma, & \alpha \end{pmatrix}\begin{pmatrix} \varepsilon^h, & 0 \\ 0, & 1 \end{pmatrix}\begin{pmatrix} \alpha, & \beta \\ \gamma, & \delta \end{pmatrix} \\
&= \begin{pmatrix} \varepsilon^h\alpha\delta - \beta\gamma, & (\varepsilon^h-1)\beta\delta \\ -(\varepsilon^h-1)\alpha\gamma, & \alpha\delta - \varepsilon^h\beta\gamma \end{pmatrix},
\end{aligned}
\]

deren Pole $\omega$ und $\xi$ sein m\"ussen. Nach \S.\,51, (6) m\"ussen also
$\omega$ und $\xi$ die Wurzeln der quadratischen Gleichung

\[
\alpha\gamma x^2 + (\alpha\delta + \beta\gamma)x + \beta\delta = 0
\]

sein, woraus

\[
\frac{\beta\delta}{\alpha\gamma} = \omega\xi, \quad \frac{\delta}{\gamma} + \frac{\beta}{\alpha} = -\omega - \xi
\]

folgt. Aus der zweiten dieser Gleichungen ergiebt sich mittelst (8):

\[
(\alpha + \delta)(\alpha\delta + \beta\gamma) = 0,
\]

und da $\alpha\delta + \beta\gamma$ nicht verschwinden kann, weil sonst $\omega + \xi = 0$
sein m\"usste, was nach (7) nicht m\"oglich ist, so ist $\alpha + \delta = 0$.
Danach erhalten wir f\"ur $\chi$, wenn wir der Einfachheit halber
$\gamma = 1$ annehmen:

\weq{9}{\chi = \begin{pmatrix} \omega, & \beta \\ 1, & -\omega \end{pmatrix}.}

Hierin ist $\beta$, was nach (8) den Werth $-\omega\xi$ hat, noch zu
bestimmen.

Wenn wir die Substitution $\chi$ als bekannt annehmen, so
k\"onnen wir die ganze Ikosa\"edergruppe bilden.

Man hat n\"amlich, wenn man $r$ und $s$ die Zahlen 0, 1, 2, 3, 4
durchlaufen l\"asst, in dieser Gruppe die Substitutionen

\weq{10}{\Theta^r, \quad \psi\Theta^r, \quad \Theta^r\chi\Theta^s, \quad \Theta^r\chi\psi\Theta^s,}
%CONT>>


% ===== scan leaf 244, printed 223 =====
% !PRINTED 223
\weberpage{223}
%<<CONT
deren Zahl gerade 60 betr\"agt. Dass sie alle von einander verschieden
 sind, sieht man, wenn man sie in der Form darstellt:

\weq{11}{\Theta^r = \begin{pmatrix} \varepsilon^r, & 0 \\ 0, & 1 \end{pmatrix},}

\weq{12}{\psi\Theta^r = \begin{pmatrix} 0, & 1 \\ -\varepsilon^r, & 0 \end{pmatrix},}

\weq{13}{\Theta^r\chi\Theta^s = \begin{pmatrix} \varepsilon^r\omega, & \varepsilon^{r-s}\beta \\ 1, & -\varepsilon^{-s}\omega \end{pmatrix},}

\weq{14}{\begin{aligned}
\Theta^r\chi\psi\Theta^s &= \begin{pmatrix} -\varepsilon^{r+s}\beta, & \varepsilon^r\omega \\ \varepsilon^s\omega, & 1 \end{pmatrix} \\
&= \begin{pmatrix} -\varepsilon^r\beta\omega^{-1}, & \varepsilon^{r-s} \\ 1, & \varepsilon^{-s}\omega^{-1} \end{pmatrix},
\end{aligned}}

von denen, da $-\omega^2$ keine Potenz von $\varepsilon$ ist, ersichtlich keine
zwei einander gleich sind.

Die Ikosa\"edergruppe hat, wie wir gesehen haben, 12 f\"unf-
z\"ahlige Pole. Je zwei dieser Pole sind die Pole von vier Substitutionen
 5\textsuperscript{ten} Grades, die mit der Identit\"at zusammen einen
f\"unfgliedrigen Cyklus bilden. Folglich giebt es in der Gruppe
24 Elemente 5\textsuperscript{ten} Grades. Ebenso giebt es 20 dreiz\"ahlige Pole,
die zu 20 Substitutionen dritter Ordnung f\"uhren, und 30 zwei-
z\"ahlige Pole, die zu zweien die Pole von je einer Substitution
zweiter Ordnung sind, so dass es 15 Substitutionen zweiter Ordnung
 giebt. Dies giebt mit der Identit\"at zusammen

\[
24 + 20 + 15 + 1 = 60.
\]

Die Bestimmung von $\beta$ in der Substitution $\chi$ ergiebt sich
nun durch Betrachtung der Grade von (13) und (14). Wir bilden
dazu zun\"achst f\"ur eine beliebige lineare Substitution $S = \begin{pmatrix} a, & b \\ c, & d \end{pmatrix}$
die zweite und dritte Wiederholung:

\[
\begin{aligned}
S^2 &= \begin{pmatrix} a^2 + bc, & b(a+d) \\ c(a+d), & bc + d^2 \end{pmatrix} \\
S^3 &= \begin{pmatrix} a^3 + 2abc + bcd, & b(a^2+bc+ad+d^2) \\ c(a^2+bc+ad+d^2), & d^3+abc+2bcd \end{pmatrix},
\end{aligned}
\]

und wenn wir von dem Falle absehen, dass $b$ oder $c$ verschwindet
, der hier nicht in Betracht kommt, so erhalten wir die
nothwendige und hinreichende Bedingung:
%CONT>>


% ===== scan leaf 245, printed 224 =====
% !PRINTED 224
\weberpage{224}

%<<CONT
für eine Substitution zweiter Ordnung
\weq{15}{a + d = 0,}
und für eine Substitution dritter Ordnung
\weq{16}{a^2 + bc + ad + d^2 = 0.}

Nun sind die Substitutionen (11), abgesehen von der darunter enthaltenen identischen, von der fünften Ordnung. Die fünf Substitutionen (12) sind von der zweiten Ordnung. Von den Substitutionen (13) sind nach (15) die fünf in der Form $\Theta^r\chi\Theta^{-r}$ enthaltenen (und nur diese) von der zweiten Ordnung, und folglich müssen noch fünf von den Substitutionen (14) von der zweiten Ordnung sein. Dies ist aber nach (15) nur möglich, \gesperrt{wenn $\beta$ eine Potenz von $\varepsilon$ ist}.

Die Substitutionen dritter Ordnung müssen sich nun auf die Formen (13) und (14) vertheilen. Für diese ergeben sich nach (16) die Bedingungen:

dass (13) von der dritten Ordnung sei
\weq{17}{\omega^2(\varepsilon^{r+s} + \varepsilon^{-(r+s)} - 1) = -\beta,}
und dass (14) von der dritten Ordnung sei
\weq{18}{\omega^2 = \beta - \varepsilon^{-(r+s)} - \beta^2\varepsilon^{r+s}.}

Soll aber auch nur eine der beiden Bedingungen (17) oder (18) befriedigt werden können, so muss
\[
\beta = 1
\]
sein. Denn $\omega^2$ ist reell und die ganze linke Seite von (17) ist also gleichfalls reell; folglich muss $\beta$ reell, und da es eine Potenz von $\varepsilon$ ist, gleich 1 sein.

Soll aber (18) befriedigt werden, so darf sich die linke Seite nicht ändern, wenn man zu den conjugirt imaginären Grössen übergeht; d. h. es muss
\[
\beta - \varepsilon^{-(r+s)} - \beta^2\varepsilon^{r+s} = \beta^{-1} - \varepsilon^{r+s} - \beta^{-2}\varepsilon^{-(r+s)}
\]
oder
\[
(\beta - \beta^{-1})(1 - \beta^{-1}\varepsilon^{-(r+s)} - \beta\varepsilon^{r+s}) = 0
\]
sein, und dies ist, weil $\beta$ eine Potenz von $\varepsilon$ ist, nur möglich, wenn $\beta = \beta^{-1}$, also $\beta = 1$ ist.

Hiernach sind fünf von den Substitutionen (14) von der zweiten Ordnung, nämlich $\Theta^r\chi\psi\Theta^{-r}$, und es ergiebt sich, dass (13) von der dritten Ordnung ist, wenn
\[
\omega^2(\varepsilon^{r+s} + \varepsilon^{-(r+s)} - 1) = -1,
\]
%CONT>>


% ===== scan leaf 246, printed 225 =====
% !PRINTED 225
\weberpage{225}
%<<CONT
oder, wenn man mit $\omega'^2$ multiplicirt, nach (5) und (6)
\[
\begin{aligned}
\varepsilon^{r+s} + \varepsilon^{-(r+s)} - 1 &= -\varepsilon - \varepsilon^{-1} - 2 \\
\varepsilon + \varepsilon^{-1} + \varepsilon^{r+s} + \varepsilon^{-(r+s)} + 1 &= 0,
\end{aligned}
\]
eine Bedingung, die dann und nur dann erfüllt ist, wenn
$r + s \equiv \pm 2 \pmod{5}$ ist.

Wenn (14) von der dritten Ordnung sein soll, so muss
nach (18) und (5)
\[
\begin{aligned}
\varepsilon^2 + \varepsilon^{-2} + 2 &= 1 - \varepsilon^{-(r+s)} - \varepsilon^{r+s} \\
\varepsilon^2 + \varepsilon^{-2} + \varepsilon^{r+s} + \varepsilon^{-(r+s)} + 1 &= 0
\end{aligned}
\]
sein, und diese Bedingung ist dann und nur dann befriedigt,
wenn $r + s \equiv \pm 1 \pmod{5}$ ist.

Hiernach erhalten wir, wie es sein muss, in (13) und (14)
gerade 10 Substitutionen zweiter und 20 Substitutionen dritter
Ordnung, und die anderen Fälle bleiben also für die fünfte Ordnung
übrig.

Fassen wir das Resultat dieser Betrachtung zusammen, so
haben wir:

Die Substitution $\chi$ muss den Ausdruck haben:
\weq{19}{\chi = \begin{pmatrix} \omega, & 1 \\ 1, & -\omega \end{pmatrix},}
und unter den 60 Substitutionen (10):
\weq{20}{\Theta^r, \; \psi\Theta^r, \; \Theta^r\chi\Theta^s, \; \Theta^r\chi\psi\Theta^s}
kommen ausser der Identität vor:

15 \gesperrt{Substitutionen} $2^{\text{ten}}$ \gesperrt{Grades}:

$\psi\Theta^r, \; \Theta^r\chi\Theta^{-r}, \; \Theta^r\chi\psi\Theta^{-r}, \; r = 0, 1, 2, 3, 4,$
\weq{21}{
\begin{aligned}
&\text{20 }\text{\gesperrt{Substitutionen}}\ 3^{\text{ten}}\ \text{\gesperrt{Grades}:} \\
&\Theta^r\chi\Theta^s, \quad r+s \equiv \pm 2 \pmod{5} \\
&\Theta^r\chi\psi\Theta^s, \quad r+s \equiv \pm 1 \pmod{5}, \\
&\text{24 }\text{\gesperrt{Substitutionen}}\ 5^{\text{ten}}\ \text{\gesperrt{Grades}:} \\
&\Theta^r, \; \Theta^r\chi\Theta^s, \quad r+s \equiv \pm 1 \pmod{5} \\
&\Theta^r\chi\psi\Theta^s, \quad r+s \equiv \pm 2 \pmod{5}.
\end{aligned}
}
In expliciter Form erhält man für die Substitutionen (20)
den Ausdruck:
\weq{22}{
\begin{pmatrix} \varepsilon^r, & 0 \\ 0, & 1 \end{pmatrix}, \;
\begin{pmatrix} 0, & \varepsilon^r \\ -1, & 0 \end{pmatrix}, \;
\begin{pmatrix} \varepsilon^r\omega, & \varepsilon^{r-s} \\ 1, & -\varepsilon^{-s}\omega \end{pmatrix}, \;
\begin{pmatrix} -\varepsilon^{r+s}, & \varepsilon^r\omega \\ \varepsilon^s\omega, & 1 \end{pmatrix}
}
oder
\[
\varepsilon^r x, \quad \frac{-\varepsilon^r}{x}, \quad \frac{\varepsilon^r\omega x + \varepsilon^{r-s}}{x - \varepsilon^{-s}\omega}, \quad \frac{-\varepsilon^{r+s} x + \varepsilon^r\omega}{\varepsilon^s\omega x + 1}.
\]


% ===== scan leaf 247, printed 226 =====
% !PRINTED 226
\weberpage{226}
Um also endlich die Existenz der Ikosa\"edergruppe festzustellen, ist noch nachzuweisen, dass die Gesammtheit der Substitutionen (20) eine Gruppe bildet. Dies folgt aber aus den nun abzuleitenden Compositionsgesetzen.

Zunächst ergiebt sich sehr einfach aus der Bedeutung von $\Theta$, $\psi$, $\chi$:
\[
\Theta = \begin{pmatrix} \varepsilon, & 0 \\ 0, & 1 \end{pmatrix}, \quad
\psi = \begin{pmatrix} 0, & 1 \\ -1, & 0 \end{pmatrix}, \quad
\chi = \begin{pmatrix} \omega, & 1 \\ 1, & -\omega \end{pmatrix}
\]
\weq{23}{\psi\Theta = \Theta^{-1}\psi, \quad \chi\psi = \psi\chi.}

Ferner erhält man für jeden beliebigen Exponenten $\lambda$:
\[
\chi\Theta^\lambda\chi = \begin{pmatrix} \varepsilon^\lambda\omega^2+1, & (\varepsilon^\lambda-1)\omega \\ (\varepsilon^\lambda-1)\omega, & \varepsilon^\lambda+\omega^2 \end{pmatrix},
\]
oder, indem man nach (5) und (6)
\weq{24}{\omega = \varepsilon+\varepsilon^{-1}, \quad \omega' = \varepsilon^2+\varepsilon^{-2}, \quad \omega\omega' = -1}
setzt:
\[
\chi\Theta^\lambda\chi = \begin{pmatrix} \varepsilon^\lambda\omega-\omega', & \varepsilon^\lambda-1 \\ \varepsilon^\lambda-1, & -\varepsilon^\lambda\omega'+\omega \end{pmatrix}.
\]

Setzen wir darin zunächst $\lambda=1$, so folgt nach (24), (13), (14):
\[
\chi\Theta\chi = \begin{pmatrix} 1-\varepsilon^3, & \varepsilon-1 \\ \varepsilon-1, & \varepsilon-\varepsilon^3 \end{pmatrix}
= \begin{pmatrix} \varepsilon^3(\varepsilon+1), & 1 \\ 1, & -\varepsilon(\varepsilon+1) \end{pmatrix}
\]
\[
= \begin{pmatrix} \varepsilon\omega', & 1 \\ 1, & -\varepsilon^{-1}\omega' \end{pmatrix}
= \begin{pmatrix} -\varepsilon^2, & \varepsilon\omega \\ \varepsilon\omega, & 1 \end{pmatrix}
= \Theta\chi\psi\Theta,
\]
woraus, wenn man beiderseits zur entgegengesetzten Substitution übergeht,
\[
\chi\Theta^{-1}\chi = \Theta^{-1}\chi\psi\Theta^{-1}
\]
folgt. Setzt man andererseits $\lambda=2$, so folgt:
\[
\chi\Theta^2\chi = \begin{pmatrix} \varepsilon-\varepsilon^2, & \varepsilon^2-1 \\ \varepsilon^2-1, & \varepsilon-1 \end{pmatrix}
= \begin{pmatrix} \varepsilon^2(\varepsilon^2+1), & 1 \\ 1, & -\varepsilon(\varepsilon^2+1) \end{pmatrix}
\]
\[
= \begin{pmatrix} \varepsilon^3\omega, & 1 \\ 1, & -\varepsilon^2\omega \end{pmatrix}
= \Theta^{-2}\chi\Theta^{-2},
\]
worin wieder $\Theta$ durch $\Theta^{-1}$ ersetzt werden kann. Man hat daher die Compositionsformeln:

\weq{25}{\chi\Theta^{\pm1}\chi = \Theta^{\pm1}\chi\psi\Theta^{\pm1}, \quad \chi\Theta^{\pm2}\chi = \Theta^{\mp2}\chi\Theta^{\mp2}.}

Durch Anwendung der Formeln (23) und (25) kann man nun je zwei der Substitutionen (20) componiren und gelangt immer wieder auf eine Substitution von der Form (20), wodurch
%CONT>>


% ===== scan leaf 248, printed 227 =====
% !PRINTED 227
\weberpage{227}
\pdfbookmark[2]{§ 59. Die Theiler der Ikosaëdergruppe}{bm69}
%<<CONT
die Gruppennatur nachgewiesen und die Ikosa\"edergruppe gebildet ist.

Aus der ersten Formel (25) ergiebt sich noch
\[
\psi = \chi\Theta^{\mp1}\chi\Theta^{\pm1}\chi\Theta^{\mp1},
\]
woraus man schliesst, dass $\psi$ aus $\chi$ und $\Theta$ abgeleitet werden kann, dass also $\chi$ und $\Theta$ allein schon als \gesperrt{erzeugende Substitutionen der Ikosa\"edergruppe} betrachtet werden können.

Wollen wir die Substitution $\chi$ mit der Determinante 1 darstellen, so beachten wir die Relation
\[
-\omega^2-1 = \omega-2 = (\varepsilon^2-\varepsilon^{-2})^2,
\]
und erhalten, wenn wir noch $\varepsilon+\varepsilon^{-1}$ für $\omega$ einsetzen:
\weq{26}{\chi = \begin{pmatrix} \dfrac{1}{\varepsilon-\varepsilon^{-1}}, & \dfrac{1}{\varepsilon^2-\varepsilon^{-2}} \\[2mm] \dfrac{1}{\varepsilon^2-\varepsilon^{-2}}, & \dfrac{-1}{\varepsilon-\varepsilon^{-1}} \end{pmatrix}.}

\S.\,59.

\gesperrt{Die Theiler der Ikosa\"edergruppe.}

Die Theiler der Ikosa\"edergruppe müssen unter den niedrigeren Poly\"edergruppen gesucht werden. Darunter finden sich zunächst die cyklischen Gruppen $C_n$, deren jede aus der Periode einer der Ikosa\"edersubstitutionen besteht. Die Abzählung der Anzahl dieser Gruppen ergiebt sich sofort aus der Zusammenstellung der Substitutionen nach ihren Graden, wie wir sie im vorigen Paragraphen gegeben haben; nämlich:
\[
\begin{aligned}
&15 \text{ Gruppen } C_2, \\
&10 \quad \text{„} \quad C_3, \\
&\phantom{0}6 \quad \text{„} \quad C_5.
\end{aligned}
\]
Es sind ferner in der Ikosa\"edergruppe Di\"edergruppen $D_2$, $D_3$, $D_5$ enthalten, als deren Repräsentanten wir folgende aufstellen:
\[
\begin{aligned}
D_2 &: 1, \psi, \chi, \psi\chi, \\
D_3 &: 1, \chi\Theta^2, \Theta^{-2}\chi, \psi\chi, \psi\Theta^2, \Theta^{-2}\psi\chi\Theta^2, \\
D_5 &: \Theta^r, \psi\Theta^r, \quad r=0,1,2,3,4,5.
\end{aligned}
\]
% [sic: Weber prints r=0,1,2,3,4,5; since \Theta^5=1, the last value duplicates r=0.]


% ===== scan leaf 249, printed 228 =====
% !PRINTED 228
\weberpage{228}
Die Anzahl der Di\"edergruppen erhalten wir daraus, dass die Di\"edergruppe durch die in ihr enthaltene cyklische Gruppe vollständig bestimmt ist (\S.\,55). Bei der Vierergruppe ist aber noch zu beachten, dass man dieselbe Gruppe erhält, wenn man von $\psi$, von $\chi$ oder von $\psi\chi$ ausgeht, und dass also die direct erhaltene Zahl durch 3 zu dividiren ist. Die Anzahl der $D_2$ ist also 5, die der $D_3$ ist 10 und die der $D_5$ ist 6.

Von besonderer Wichtigkeit sind aber die in der Ikosa\"edergruppe enthaltenen Tetra\"edergruppen. Wir wollen eine von ihnen, die wir mit $Q$ bezeichnen, bestimmen.

Die Tetra\"edergruppe muss (nach \S.\,56) eine Vierergruppe $D_2$ enthalten. Wir gehen von einer solchen Gruppe $D_2$ aus und wählen dazu
\[
1, \psi, \chi, \psi\chi.
\]

Es muss nun weiter in $Q$ eine Substitution dritter Ordnung vorkommen. Diese können wir in einer der beiden Formen $\Theta^r\chi\Theta^s$ oder $\Theta^r\chi\psi\Theta^s$ annehmen. Da beide Annahmen zu demselben Resultate führen, wählen wir als Substitution dritter Ordnung
\[
\varphi = \Theta^r\chi\Theta^s, \quad r+s\equiv\pm2 \ (\mathrm{mod}\ 5) \quad [\S.\,58, (21)].
\]

Von den Zahlen $r, s$ kann aber keine $\equiv0$ sein, weil sonst entweder $\chi\varphi$ oder $\varphi\chi$ eine Potenz von $\Theta$, also vom $5^{\mathrm{ten}}$ Grade wäre, während doch in $Q$ kein Element $5^{\mathrm{ten}}$ Grades vorkommen kann.

Wir bilden ferner nach \S.\,58, (23), (25):
\[
\begin{aligned}
\chi\varphi &= \Theta^r\chi\psi\Theta^{r+s}, \quad r\equiv\pm1 \\
\chi\varphi &= \Theta^{-r}\chi\Theta^{-r+s}, \quad r\equiv\pm2,
\end{aligned}
\]
wodurch aus dem oben angeführten Grunde im ersten Falle $r\equiv-s$, im zweiten $r\equiv s$ ausgeschlossen ist, und im ersten Falle $2r+s\equiv\pm1$, im zweiten $-2r+s\equiv\pm2$ gefordert wird. Danach bleiben die vier möglichen Fälle
\[
r\equiv\pm1, \quad s\equiv\pm2, \quad r\equiv\pm2, \quad s\equiv\pm1 \ (\mathrm{mod}\ 5)
\]
übrig.

Wenn von den so bestimmten vier Substitutionen dritter Ordnung eine in $Q$ vorkommt, so kommen auch die drei anderen vor. Denn setzen wir
\weq{1}{\varphi = \Theta\chi\Theta^2,}
so ergiebt sich nach \S.\,58, (23), (25):
\weq{2}{\varphi^{-1} = \Theta^{-2}\chi\Theta^{-1}, \quad \varphi\chi = \Theta^{-1}\chi\Theta^{-2}, \quad \chi\varphi^{-1} = \Theta^2\chi\Theta.}


% ===== scan leaf 250, printed 229 =====
% !PRINTED 229
\weberpage{229}
Bildet man ausserdem noch $\varphi\psi$, $\varphi^{-1}\psi$, $\varphi\chi\psi$, $\chi\varphi^{-1}\psi$, so folgt:
\[
\begin{aligned}
\varphi\psi &= \Theta\chi\psi\Theta^{-2}, \quad \varphi^{-1}\psi = \Theta^{-2}\chi\psi\Theta, \quad \varphi\chi\psi = \Theta^{-1}\chi\psi\Theta^2, \\
&\chi\varphi^{-1}\psi = \Theta^2\chi\psi\Theta^{-1},
\end{aligned}
\]
woraus man sieht, dass man zu keinem anderen Resultate kommen würde, wenn man die Substitution dritter Ordnung, von der man ausgeht, in der zweiten Form $\Theta^r\chi\psi\Theta^s$ annehmen wollte. Die ganze Gruppe $Q$ ist also durch die angenommene Vierergruppe $D_2$ völlig bestimmt, und man erhält sie in der Form
\weq{3}{
\begin{array}{ll}
1 & \psi \\
\varphi=\Theta\chi\Theta^2, & \varphi\psi=\Theta\chi\psi\Theta^{-2}, \\
\varphi^{-1}=\Theta^{-2}\chi\Theta^{-1}, & \varphi^{-1}\psi=\Theta^{-2}\chi\psi\Theta, \\
\chi & \chi\psi \\
\varphi\chi=\Theta^{-1}\chi\Theta^{-2}, & \varphi\chi\psi=\Theta^{-1}\chi\psi\Theta^2, \\
\varphi^{-1}\chi=\Theta^2\chi\psi\Theta^{-1}, & \varphi^{-1}\chi\psi=\Theta^2\chi\Theta.
\end{array}
}

Man kann diese Gruppe aus den Substitutionen $\varphi, \psi, \chi$ als den Erzeugenden ableiten und sie in die Form setzen:
\weq{4}{Q = \varphi^r\chi^s\psi^t, \qquad \begin{gathered}r=0,1,2,\\ s=0,1; \quad t=0,1.\end{gathered}}

Dass dadurch wirklich eine Tetra\"edergruppe dargestellt ist, ergiebt sich aus den Zusammensetzungen, die man mittelst \S.\,58, (23), (25) leicht aus (1) findet:
\weq{5}{
\begin{aligned}
\psi\varphi &= \varphi\chi\psi, \quad \chi\varphi = \varphi\psi, \\
\psi\varphi^2 &= \chi\varphi^2\psi, \quad \chi\varphi^2 = \varphi^2\chi\psi. \quad \psi\chi = \chi\psi. % [sic: Weber prints a full stop between the two relations, unlike the commas separating the parallel relations; the printed punctuation is preserved.]
\end{aligned}
}

Da in der Tetra\"edergruppe fünf Vierergruppen $D_2$ enthalten sind, so hat die Ikosa\"edergruppe fünf Tetra\"edergruppen zu Theilern. % [sic: Weber prints Tetraëdergruppe; mathematically Ikosaëdergruppe is intended.]
Diese können wir aus $Q$ durch Transformation mittelst der Potenzen von $\Theta$ ableiten und erhalten sie in der Form
\[
\Theta^{-r}Q\Theta^r \quad r=0,1,2,3,4.
\]

Diese Gruppen haben, ausser der Identität, keine Substitution mit einander gemein. Denn die beiden Gruppen $Q$ und $\Theta^{-1}Q\Theta$ haben ausser der Identität nur die beiden Elemente
\[
\Theta^{-2}\chi\Theta^{-1}, \quad \Theta\chi\Theta^2
\]
mit einander gemein, und von diesen kommt keines in $\Theta Q\Theta^{-1}$ vor.

Daraus ergiebt sich nun nach dem Satze 2. in \S.\,28, dass die Ikosa\"edergruppe isomorph ist mit einer Permutationsgruppe
%CONT>>


% ===== scan leaf 251, printed 230 =====
% !PRINTED 230
\weberpage{230}
\pdfbookmark[2]{§ 60. Die Grundformen der Ikosaëdergruppe}{bm70}
%<<CONT
60\textsuperscript{sten} Grades von fünf Ziffern. Diese Permutationsgruppe ergiebt sich, wenn wir mit den Nebengruppen
\[
Q, \quad Q\Theta, \quad Q\Theta^2, \quad Q\Theta^3, \quad Q\Theta^4
\]
die sämmtlichen Elemente $\sigma$ der Ikosa\"edergruppe verbinden, also
\[
Q\sigma, \quad Q\Theta\sigma, \quad Q\Theta^2\sigma, \quad Q\Theta^3\sigma, \quad Q\Theta^4\sigma
\]
bilden, und die dadurch bewirkte Permutation dieser Nebengruppen untersuchen. Für die erzeugenden Substitutionen der Ikosa\"edergruppe $\sigma = \Theta, \psi, \chi$ erhalten wir so die Permutationen
\[
(1, \ \Theta, \ \Theta^2, \ \Theta^3, \ \Theta^4) \qquad \sigma = \Theta
\]
\[
\begin{pmatrix} Q, & Q\Theta, & Q\Theta^2, & Q\Theta^3, & Q\Theta^4 \\ Q, & Q\Theta^4, & Q\Theta^3, & Q\Theta^2, & Q\Theta \end{pmatrix} \qquad \sigma = \psi
\]
\[
\begin{pmatrix} Q, & Q\Theta, & Q\Theta^2, & Q\Theta^3, & Q\Theta^4 \\ Q, & Q\Theta^3, & Q\Theta^4, & Q\Theta, & Q\Theta^2 \end{pmatrix} \qquad \sigma = \chi.
\]

Letzteres findet man aus den Formeln (1) und (2), wonach
\[
\Theta\chi = \varphi\Theta^3, \quad \Theta^2\chi = \chi\varphi^2\Theta^4, \quad \Theta^3\chi = \varphi^2\Theta, \quad \Theta^4\chi = \varphi\chi\Theta^2.
\]

Man sieht, dass diese Permutationen alle zur ersten Art gehören, und dass also die Permutationsgruppe, um die es sich handelt, keine andere als die alternirende Gruppe von fünf Ziffern (Bd. I, \S.\,177) sein kann. Daraus ergiebt sich auch, dass die Ikosa\"edergruppe einfach ist, und folglich mit der Gruppe übereinstimmt, die wir schon in \S.\,34 vorläufig als Ikosa\"edergruppe bezeichnet haben.

\begin{center}
\S.\,60.\\[2pt]
\gesperrt{Die Grundformen der Ikosa\"edergruppe.}
\end{center}

Wir haben oben die eine der drei Grundformen der Ikosa\"edergruppe $f$ gefunden:
\weq{1}{f = x_1 x_2 (x_1^{10} + 11 x_1^5 x_2^5 - x_2^{10}).}
Es fehlen uns also noch zwei dieser Formen, eine vom 20\textsuperscript{sten} und eine vom 30\textsuperscript{sten} Grade.

Die Grundform 20\textsuperscript{sten} Grades ergiebt sich als die \gesperrt{Hesse}'sche Covariante von $f$. Wir setzen sie
\weq{2}{H = \frac{1}{121}\,[f''(x_1,x_1)\, f''(x_2,x_2) - f''(x_1,x_2)^2],}
%CONT>>


% ===== scan leaf 252, printed 231 =====
% !PRINTED 231
\weberpage{231}
%<<CONT
und finden durch einfache Rechnung:
\weq{3}{H = -(x_1^{20} + x_2^{20}) + 228\,(x_1^{15}x_2^5 - x_2^{15}x_1^5) - 494\,x_1^{10}x_2^{10}.}
Die Grundform 30\textsuperscript{sten} Grades können wir als die Functionaldeterminante von $H$ und $f$ definiren. Setzen wir
\weq{4}{T = \frac{1}{20}\,[f'(x_1)\,H'(x_2) - f'(x_2)\,H'(x_1)],}
so findet sich
\weq{5}{T = (x_1^{30} + x_2^{30}) + 522\,(x_1^{25}x_2^5 - x_1^5x_2^{25}) - 10005\,(x_1^{20}x_2^{10} + x_2^{20}x_1^{10}).}
Auch zwischen diesen drei Formen besteht eine identische Relation. Um sie zu finden, setzen wir
\[
\lambda = x_1^{10} - x_2^{10}, \quad \mu = x_1^5 x_2^5,
\]
und erhalten, wenn wir bei $T$ den Factor $x_1^{10} + x_2^{10}$ herausnehmen und dann $T^2$ bilden:
\[
f^5 = \mu\,(\lambda + 11\mu)^5,
\]
\[
H = -\lambda^2 + 228\,\lambda\mu - 496\,\mu^2,
\]
\[
T^2 = (\lambda^2 + 4\mu^2)\,(\lambda^2 + 522\,\lambda\mu - 10004\,\mu^2)^2.
\]
Daraus erhält man durch die numerische Berechnung
\weq{6}{T^2 + H^3 = 1728\,f^5,}
was die gesuchte Relation ist.

Da die Ikosa\"edergruppe, wie wir gesehen haben, einfach ist, so kann sie nach \S.\,40 keine relativen Invarianten haben. Daraus ergiebt sich, dass die Ikosa\"ederformen $f$, $H$, $T$ absolut ungeändert bleiben, wenn man sie den mit der Determinante $1$ dargestellten binären Ikosa\"edersubstitutionen unterwirft. Dasselbe lässt sich aber auch, ohne jene allgemeinen Sätze zu benutzen, in folgender Weise direct nachweisen.

Die Relation (6) lässt sich in der Form darstellen:
\weq{7}{\frac{T^2}{f^5} + \frac{H^3}{f^5} = 1728,}
und die beiden Quotienten $T^2 : f^5$, $H^3 : f^5$ können wir als Functionen der einen Veränderlichen $x = x_1 : x_2$ auffassen. Wenden wir auf diese Variable eine Substitution der Ikosa\"edergruppe an, so ändern sich diese beiden Quotienten nur um constante Factoren, d. h. wenn wir
\weq{8}{\frac{T^2}{f^5} = \Phi(x), \quad \frac{H^3}{f^5} = \Psi(x), \quad y = \frac{\alpha x + \beta}{\gamma x + \delta}}
%CONT>>


% ===== scan leaf 253, printed 232 =====
% !PRINTED 232
\weberpage{232}
\pdfbookmark[2]{§ 61. Die Invarianten des Ikosaëders}{bm71}
%<<CONT
setzen und mit $h$, $k$ zwei Constanten bezeichnen, so ist
\[
\Phi(y) = h\,\Phi(x), \quad \Psi(y) = k\,\Psi(x),
\]
vorausgesetzt, dass $\begin{pmatrix} \alpha, & \beta \\ \gamma, & \delta \end{pmatrix}$ eine der Ikosa\"edersubstitutionen ist. Da nun aber $y$ sowohl als $x$ eine unabhängige Variable ist, so besteht nach (7) die Identität
\[
\Phi(x) + \Psi(x) = h\,\Phi(x) + k\,\Psi(x) = 1728,
\]
und da $\Phi$ und $\Psi$ nicht in constantem Verhältnisse stehen, so muss $h = k = 1$ sein, d. h. die Quotienten $\Phi$ und $\Psi$ bleiben bei den Ikosa\"edersubstitutionen absolut ungeändert.

Daraus ergiebt sich weiter, \gesperrt{dass $f$, $H$, $T$ bei den homogenen Ikosa\"edersubstitutionen mit der Determinante $1$ absolut ungeändert bleiben.}

Denn wenn $f$ durch eine solche Substitution in $cf$ übergeht, so gehen $H$ und $T$ in $c^2H$ und $c^3T$ über [nach (2) und (4)], und die Quotienten $\Phi$ und $\Psi$ in $c\Phi$ und $c\Psi$. Da andererseits $\Phi$ und $\Psi$ ungeändert bleiben, so ist $c = 1$.

Setzen wir $\Psi(x) = z$, so wird $\Phi(x) = 1728 - z$, und wir erhalten aus (8) die beiden Gleichungen:
\weq{9}{H^3 - z f^5 = 0,}
\weq{10}{T^2 - (1728 - z)\,f^5 = 0,}
von denen wegen der Identität (6) die eine aus der anderen folgt, so dass es eigentlich nur zwei verschiedene Formen für eine und dieselbe Gleichung sind. Betrachten wir nun darin $z$ als gegeben, so haben wir eine Gleichung 60\textsuperscript{sten} Grades für $x$, die die \gesperrt{Ikosa\"edergleichung} heisst.

Auf die Eigenschaften dieser Gleichung und ihre Beziehung zu der allgemeinen Gleichung 5\textsuperscript{ten} Grades kommen wir in einem späteren Abschnitte zurück.

\begin{center}
\S.\,61.\\[2pt]
\gesperrt{Die Invarianten des Ikosa\"eders.}
\end{center}

Durch die Grundformen der Poly\"edergruppen, die wir bisher kennen gelernt haben, ist, wie wir jetzt beweisen können, das Gebiet der Invarianten der betreffenden Gruppen erschöpft. Wir
%CONT>>


% ===== scan leaf 254, printed 233 =====
% !PRINTED 233
\weberpage{233}
%<<CONT
beschränken uns bei diesem Beweise auf die Betrachtung der Ikosa\"edergruppe, da für die anderen Gruppen ganz ähnliche Schlüsse zu machen sind, die wir dem Leser überlassen können. Der Umstand, dass bei der Tetra\"eder- und Octa\"edergruppe neben den absoluten auch relative Invarianten vorkommen, während die Ikosa\"edergruppe als einfache Gruppe nur absolute Invarianten hat, ist hierbei von keinem wesentlichen Einflusse.

Wir beweisen also, dass alle Invarianten der Ikosa\"edergruppe sich als ganze rationale Functionen der drei Formen $f$, $H$, $T$ darstellen lassen.

Dazu führt uns folgende Schlusskette:

1. \gesperrt{Keine zwei der Ikosa\"ederformen $f$, $H$, $T$ haben einen gemeinschaftlichen Theiler.}

Dies folgt unmittelbar aus der Definition dieser Functionen, wonach die Gleichungen $f = 0$, $H = 0$, $T = 0$ die fünfzähligen, dreizähligen und zweizähligen Pole liefern, die alle von einander verschieden sind.

2. \gesperrt{Eine Doppelwurzel der Ikosa\"edergleichung}
\weq{1}{H^3 - z f^5 = 0}
\gesperrt{ist nothwendig eine Wurzel von $f$, $H$ oder $T$, und kann also nur für einen der Werthe $z = 0, \infty, 1728$ eintreten.}

Wenn nämlich (1) eine Doppelwurzel hat, so müssen mit der Function zugleich die erste Ableitung, oder, wenn man die homogene Form anwendet, nach der \gesperrt{Euler}'schen Theorie [Bd. I, \S.\,17, (5)] die beiden Ableitungen nach $x_1$ und $x_2$ zugleich verschwinden, also
\[
3H^2 H'(x_1) - 5z f^4 f'(x_1) = 0,
\]
\[
3H^2 H'(x_2) - 5z f^4 f'(x_2) = 0;
\]
folglich bleiben, da $H$ und $f$ nach 1. nicht zugleich verschwinden, nur drei Möglichkeiten: entweder $H = 0$, $z = 0$, oder $f = 0$, $z = \infty$, oder endlich
\[
H'(x_1)\,f'(x_2) - H'(x_2)\,f'(x_1) = 0,
\]
d. h. $T = 0$, $z = 1728$ [\S.\,60, (4)]. Hieran schliesst sich auch der evidente Satz:
%CONT>>


% ===== scan leaf 255, printed 234 =====
% !PRINTED 234
\weberpage{234}
\pdfbookmark[2]{§ 62. Polyëdergruppen der zweiten Art. Krystallographische Gruppen}{bm72}
%<<CONT

\gesperrt{3. Wenn} $J(x_1, x_2)$ \gesperrt{irgend eine invariante Form der Ikosa\"edergruppe ist, und} $\xi$ \gesperrt{eine ihrer Wurzeln, so dass} $J(\xi, 1) = 0$ \gesperrt{ist, so sind auch alle die Gr\"ossen Wurzeln von} $J$, \gesperrt{die aus} $\xi$ \gesperrt{durch die gebrochenen Ikosa\"edersubstitutionen hervorgehen.}

Wenn daher $J$ mit einer der Functionen $f$, $H$, $T$ einen Theiler gemein hat, so ist $J$ durch die betreffende Form theilbar.

Wir denken uns nun zun\"achst aus $J$ m\"oglichst hohe Potenzen der drei Grundformen $f$, $H$, $T$ weggehoben, so dass $J$ zu diesen Functionen theilerfremd ist. Ist dann $\xi$ eine Wurzel von $J$, so k\"onnen wir $\eta$ so bestimmen, dass $\xi$ eine Wurzel der Form
\[
\Theta = H^3 - \eta f^5
\]
ist, und $\xi$ ist dann nach 2. eine einfache Wurzel von $\Theta$. Es m\"ussen dann auch nach 3. alle \"ubrigen Wurzeln von $\Theta$ zugleich Wurzeln von $J$ sein, d. h. $J$ ist durch $\Theta$ theilbar. Der Quotient ist wieder eine invariante Form des Ikosa\"eders, und der Schluss l\"asst sich wiederholen. Wir kommen so zu dem Satze:

\gesperrt{4. Jede Invariante des Ikosa\"eders l\"asst sich in der Form darstellen:}
\weq{2}{J(x_1,x_2) = C f^\alpha H^\beta T^\gamma F(H^3, f^5),}
\gesperrt{worin $\alpha$, $\beta$, $\gamma$ nicht negative ganze Zahlen, $C$ eine Constante und $F$ eine ganze homogene Function bedeuten.}

\begin{center}
\S.\,62.\\
\gesperrt{Poly\"edergruppen der zweiten Art. Krystallographische Gruppen.}
\end{center}

Wir haben schon im \S.\,51 auf Gruppen linearer tern\"arer Substitutionen aufmerksam gemacht, die ausser den eigentlichen auch uneigentlich orthogonale Substitutionen enthalten, und die wir \gesperrt{Gruppen der zweiten Art} nennen.

Wir wollen die endlichen unter ihnen jetzt als \gesperrt{Poly\"edergruppen} der zweiten Art oder auch als \gesperrt{erweiterte Poly\"edergruppen}, und die darin enthaltenen Substitutionen mit der
%CONT>>


% ===== scan leaf 256, printed 235 =====
% !PRINTED 235
\weberpage{235}

%<<CONT
\gesperrt{Determinante} $-1$ \gesperrt{als Substitutionen der zweiten Art} bezeichnen.

Nach den Resultaten des \S.\,50 sind diese Gruppen isomorph mit den Gruppen bin\"arer linearer Substitutionen mit der Determinante $+1$ und $-1$, w\"ahrend bei den gebrochenen Substitutionen, die uns auf die Poly\"edergruppen gef\"uhrt haben, die beiden Arten nicht unterscheidbar sind.

Die endlichen Gruppen der zweiten Art sind, abgesehen von ihrem allgemeinen gruppentheoretischen Interesse, wichtig wegen ihrer Anwendung in der Krystallographie. Ihre vollst\"andige Bestimmung bietet jetzt keine wesentlichen Schwierigkeiten mehr. Wir verstehen unter Substitutionen schlechtweg bin\"are lineare Substitutionen mit der Determinante $\pm 1$, und erinnern daran, dass zwei solche Substitutionen, deren Coëfficienten sich nur durch das gemeinschaftliche Vorzeichen unterscheiden, wie
\[
\begin{pmatrix} \alpha, & \beta \\ \gamma, & \delta \end{pmatrix}, \quad
\begin{pmatrix} -\alpha, & -\beta \\ -\gamma, & -\delta \end{pmatrix},
\]
als nicht verschieden gelten (\S.\,49).

Da eine Gruppe linearer Substitutionen die identische Substitution enth\"alt, die von der Determinante $+1$ ist, so m\"ussen in jeder Gruppe $Q$ der zweiten Art neben den uneigentlichen auch eigentliche Substitutionen vorkommen, und diese eigentlichen Substitutionen bilden f\"ur sich eine Gruppe $P$. Ist $\varphi$ irgend eine in $Q$ vorkommende Substitution der zweiten Art, so kommt jede Composition von $\varphi$ mit einer Substitution aus $Q$ der zweiten Art in $P$ vor, und daher k\"onnen wir $Q$ immer so zerlegen:
\weq{1}{Q = P + P\varphi.}
Die Gruppe $P$ muss eine der fr\"uher betrachteten Poly\"edergruppen sein. Es muss $\varphi^{-1}P\varphi = P$ sein, und $P$ ist also ein \gesperrt{Normaltheiler} von $Q$. Ist umgekehrt $P$ eine Poly\"edergruppe und $\varphi$ eine Substitution zweiter Art, die der Bedingung $\varphi^{-1}P\varphi = P$ gen\"ugt, so ist $Q = P + P\varphi$ eine \gesperrt{Gruppe der zweiten Art}.

Auch hier betrachten wir zwei Gruppen, die durch Transformation aus einander hervorgehen, als nicht wesentlich verschieden. Wir haben, um alle $Q$ zu bilden, die verschiedenen F\"alle durchzugehen.

I. Ist $P$ die Einheitsgruppe, die wir als cyklische Gruppe ersten Grades mit $C_1$ bezeichnen, besteht also $P$ nur aus der
%CONT>>


% ===== scan leaf 257, printed 236 =====
% !PRINTED 236
\weberpage{236}

%<<CONT
identischen Substitution, so muss $\varphi^2 = 1$ sein, und dies ist auch ausreichend. Wir k\"onnen hier durch Transformation $\varphi$ auf die Form $\begin{pmatrix} \alpha, & 0 \\ 0, & -\alpha^{-1} \end{pmatrix}$ bringen und erhalten als Bedingung $\alpha^2 = \pm 1$, und demnach zwei Formen der Substitution $\varphi$:
\[
\varphi' = \begin{pmatrix} i, & 0 \\ 0, & i \end{pmatrix}, \quad
\varphi'' = \begin{pmatrix} 1, & 0 \\ 0, & -1 \end{pmatrix},
\]
von denen der erste die \gesperrt{Inversion}, der zweite die \gesperrt{Spiegelung} genannt wird. Uebertr\"agt man sie nach \S.\,50, (1) auf eine rechtwinkelige Coordinatentransformation, so bedeutet $\varphi'$ die gleichzeitige Vorzeichen\"anderung aller drei Coordinaten, $\varphi''$ die Vertauschung von $x$ mit $-x$. Es ergeben sich hieraus die beiden Gruppen der zweiten Art:
\weq{2}{C_1' = \begin{pmatrix} i^h, & 0 \\ 0, & i^h \end{pmatrix}, \quad C_1'' = \begin{pmatrix} 1, & 0 \\ 0, & (-1)^h \end{pmatrix}, \quad h = 0, 1.}

II. Es sei $P$ eine cyklische Gruppe $C_n$, $n > 1$.
Setzen wir
\weq{3}{c = \begin{pmatrix} \varepsilon, & 0 \\ 0, & \varepsilon^{-1} \end{pmatrix}, \quad \varepsilon = e^{\frac{\pi i}{n}},}
so ist nach \S.\,55:
\weq{4}{C_n = 1, c, c^2, \ldots, c^{n-1},}
worin
\[
c^s = \begin{pmatrix} \varepsilon^s, & 0 \\ 0, & \varepsilon^{-s} \end{pmatrix};
\]
wenn nun
\weq{5}{\varphi = \begin{pmatrix} \alpha, & \beta \\ \gamma, & \delta \end{pmatrix}, \quad \alpha\delta - \beta\gamma = -1}
ist, so bilden wir zun\"achst
\weq{6}{\varphi^{-1} c \varphi = \begin{pmatrix} \alpha\delta\varepsilon - \beta\gamma\varepsilon^{-1}, & \beta\delta(\varepsilon - \varepsilon^{-1}) \\ -\alpha\gamma(\varepsilon - \varepsilon^{-1}), & \alpha\delta\varepsilon^{-1} - \beta\gamma\varepsilon \end{pmatrix},}
\weq{7}{\varphi^2 = \begin{pmatrix} \alpha^2 + \beta\gamma, & (\alpha + \delta)\beta \\ (\alpha + \delta)\gamma, & \beta\gamma + \delta^2 \end{pmatrix}.}
Diese beiden Substitutionen m\"ussen in der Reihe (4) vorkommen, und es muss also $\beta\delta = 0$, $\alpha\gamma = 0$ sein; also m\"ussen entweder $\beta$ und $\gamma$ oder $\alpha$ und $\delta = 0$ sein. Ist $\beta = \gamma = 0$, so folgt aus (7), dass $\alpha$ von der Form $\varepsilon^{1/2\,r}$, also
\[
\varphi = \begin{pmatrix} e^{\frac{\pi i}{2n}\,r}, & 0 \\ 0, & -e^{-\frac{\pi i}{2n}\,r} \end{pmatrix}
\]
%CONT>>


% ===== scan leaf 258, printed 237 =====
% !PRINTED 237
\weberpage{237}

%<<CONT
sein muss, worin $r$ eine ganze Zahl bedeutet. Je nachdem $r$ ungerade oder gerade angenommen wird, erhalten wir zwei verschiedene Gruppen, die sich nicht durch Transformation auf einander zur\"uckf\"uhren lassen:
\weq{1}{C_n' = \begin{pmatrix} e^{\frac{\pi i h}{2n}}, & 0 \\ 0, & (-1)^h e^{-\frac{\pi i h}{2n}} \end{pmatrix} \quad h = 0, 1, \ldots, 2n-1,}
\weq{2}{C_n'' = \begin{pmatrix} e^{\frac{\pi i h}{n}}, & 0 \\ 0, & \pm e^{-\frac{\pi i h}{n}} \end{pmatrix} \quad h = 0, 1, \ldots, n-1. \quad \text{(beide Vorzeichen)}}
Wenn $n$ gerade ist, so kommen $C_1'$ und $C_1''$ beide unter $C_n''$ vor; ist aber $n$ ungerade, so kommt $C_1'$ in $C_n'$, $C_1''$ in $C_n''$ vor.

Ist sodann $\alpha = \delta = 0$, $\beta\gamma = 1$, so k\"onnen wir $\varphi$ durch die Transformation
\[
\begin{pmatrix} \beta^{-1/2}, & 0 \\ 0, & \beta^{1/2} \end{pmatrix}
\begin{pmatrix} 0, & \beta \\ \gamma, & 0 \end{pmatrix}
\begin{pmatrix} \beta^{1/2}, & 0 \\ 0, & \beta^{-1/2} \end{pmatrix}
= \begin{pmatrix} 0, & 1 \\ 1, & 0 \end{pmatrix},
\]
durch die die Gruppe $P$ unge\"andert bleibt, umformen, und erhalten noch eine dritte Gruppe:
\weq{3}{C_n''' = \begin{pmatrix} e^{\frac{\pi i h}{n}}, & 0 \\ 0, & e^{-\frac{\pi i h}{n}} \end{pmatrix}, \quad \begin{pmatrix} 0, & e^{\frac{\pi i h}{n}} \\ e^{-\frac{\pi i h}{n}}, & 0 \end{pmatrix} \qquad h = 0, 1, \ldots, n-1.}

III. Es sei $P$ die Di\"edergruppe $D_n$, die aus den Substitutionen
\weq{8}{\begin{aligned} D_n = {}& 1, c, c^2, \ldots, c^{n-1} \\ & d, cd, c^2d, \ldots, c^{n-1}d, \\ = {}& C_n + C_n d \end{aligned} \qquad d = \begin{pmatrix} 0, & i \\ i, & 0 \end{pmatrix}}
besteht. Die Substitution $c$ ist vom $n^{\text{ten}}$, die $c^\lambda d$ sind alle vom $2^{\text{ten}}$ Grade. Es ist aber $\varphi^{-1}c\varphi$ vom $n^{\text{ten}}$ Grade, und muss daher, wenn $n > 2$ ist, unter den Potenzen von $c$ enthalten sein. Folglich sind zur Bestimmung von $\varphi$ die vorigen Betrachtungen anwendbar, und in der Gruppe $Q$ muss eine der Gruppen $C_n'$, $C_n''$, $C_n'''$ enthalten sein. Da $d$ nicht in diesen Gruppen vorkommt, so ergeben sich f\"ur die Di\"edergruppen zweiter Art die Formen
\weq{9}{C_n' + C_n'd, \quad C_n'' + C_n''d, \quad C_n''' + C_n'''d,}
%CONT>>


% ===== scan leaf 259, printed 238 =====
% !PRINTED 238
\weberpage{238}
%<<CONT
von denen die beiden ersten
\[
1)\quad D_n' = \begin{pmatrix} e^{\frac{\pi i h}{2n}}, & 0 \\ 0, & (-1)^h e^{-\frac{\pi i h}{2n}} \end{pmatrix},\quad
\begin{pmatrix} 0, & -e^{\frac{\pi i h}{2n}} \\ (-1)^{h+n} e^{-\frac{\pi i h}{2n}}, & 0 \end{pmatrix}
\]
\[
h = 0, 1, \ldots, 2n-1,
\]
\[
2)\quad D_n'' = \begin{pmatrix} e^{\frac{\pi i h}{n}}, & 0 \\ 0, & \pm\, e^{-\frac{\pi i h}{n}} \end{pmatrix},\quad
\begin{pmatrix} 0, & e^{\frac{\pi i (2h+n)}{2n}} \\ \pm\, e^{-\frac{\pi i (2h+n)}{2n}}, & 0 \end{pmatrix}
\]
\[
h = 0, 1, \ldots, n-1
\]
sind, während die dritte Gruppe (9)
\[
\begin{pmatrix} e^{\frac{\pi i h}{n}}, & 0 \\ 0, & e^{-\frac{\pi i h}{n}} \end{pmatrix},\quad
\begin{pmatrix} 0, & e^{\frac{\pi i h}{n}} \\ e^{-\frac{\pi i h}{n}}, & 0 \end{pmatrix},
\]
\[
\begin{pmatrix} 0, & e^{\frac{\pi i}{2n}(2h+n)} \\ -e^{-\frac{\pi i}{2n}(2h+n)}, & 0 \end{pmatrix},\quad
\begin{pmatrix} e^{\frac{\pi i}{2n}(2h+n)}, & 0 \\ 0, & -e^{-\frac{\pi i}{2n}(2h+n)} \end{pmatrix}
\]
\[
h = 0, 1, \ldots, n-1
\]
ergiebt, was bei geradem $n$ mit $D_n''$, bei ungeradem $n$ mit $D_n'$ übereinstimmt.

Der Fall $n=2$, wo $c = \begin{pmatrix} i, & 0 \\ 0, & -i \end{pmatrix}$ ist, bildet nur eine scheinbare Ausnahme, weil in diesem Falle
\weq{10}{\varphi^{-1} c \varphi = d \quad \text{oder} \quad \varphi^{-1} c \varphi = c d}
sein kann. Verfolgt man diese Annahmen durch einfache Rechnung, so ergeben sich noch zwei Gruppen:
\weq{11}{
\begin{aligned}
&1,\ c,\ d,\ cd,\ \varphi_1,\ \varphi_1 c,\ \varphi_1 d,\ \varphi_1 c d\\
&1,\ c,\ d,\ cd,\ \varphi_2,\ \varphi_2 c,\ \varphi_2 d,\ \varphi_2 c d,
\end{aligned}
}
worin
\[
\varphi_1 = \begin{pmatrix} \dfrac{1}{\sqrt2}, & \dfrac{1}{\sqrt2} \\ \dfrac{1}{\sqrt2}, & \dfrac{-1}{\sqrt2} \end{pmatrix},\quad
\varphi_2 = \begin{pmatrix} \dfrac{1}{\sqrt2}, & -\dfrac{i}{\sqrt2} \\ \dfrac{i}{\sqrt2}, & -\dfrac{1}{\sqrt2} \end{pmatrix}.
\]


% ===== scan leaf 260, printed 239 =====
% !PRINTED 239
\weberpage{239}
Es ist aber
\[
\varphi_2^{-1} \varphi_1 \varphi_2 = \begin{pmatrix} 0, & e^{\frac{\pi i}{4}} \\ e^{-\frac{\pi i}{4}}, & 0 \end{pmatrix}
\]
\[
\varphi_1^{-1} \varphi_2 \varphi_1 = \begin{pmatrix} 0, & e^{\frac{\pi i}{4}} \\ e^{-\frac{\pi i}{4}}, & 0 \end{pmatrix},
\]
\[
\begin{aligned}
&c\varphi_1 = \varphi_1 d,\quad d\varphi_1 = \varphi_1 c,\quad cd\varphi_1 = \varphi_1 dc\\
&c\varphi_2 = \varphi_2 dc,\quad d\varphi_2 = \varphi_2 d,\quad cd\varphi_2 = \varphi_2 c,
\end{aligned}
\]
und aus diesen Formeln folgt, dass die Gruppen (11) durch Transformation mit $\varphi_1$ und $\varphi_2$ in $D_2'$ transformirt werden.

Bei der Bestimmung der übrigen Poly\"edergruppen zweiter Art sind die folgenden allgemeinen Bemerkungen von Nutzen.
Die Inversion $j = \begin{pmatrix} i, & 0 \\ 0, & i \end{pmatrix}$ ist eine Aehnlichkeitssubstitution (\S.\,37) und daher mit jeder anderen Substitution vertauschbar, und folglich können wir aus jeder Poly\"edergruppe erster Art wenigstens eine Gruppe zweiter Art herleiten:
\weq{12}{P + Pj.}
Um die anderen etwa noch vorhandenen Gruppen dieser Art zu finden, erinnern wir uns, dass nach dem \gesperrt{Sylow}'schen Satze (\S.\,29, I.) in jeder Gruppe $G$ eine Gruppe enthalten sein muss, deren Grad die höchste Potenz von 2 ist, die im Grade von $G$ aufgeht. Die erweiterten Tetra\"eder-, Octa\"eder- und Ikosa\"edergruppen haben aber die Grade 24, 48, 120, müssen also einen Theiler vom Grade 8, 16, 8 enthalten.

Es sei nun $T$ die Tetra\"eder-, $O$ die Octa\"eder- und $J$ die Ikosa\"edergruppe. In $T$ und $J$ ist eine Vierergruppe $D_2$ enthalten, aber keine Substitution vierter Ordnung, in $O$ eine Di\"odergruppe $D_4$ und keine Substitution achter Ordnung, und es muss daher unter den erweiterten Poly\"edergruppen eine der unter III. betrachteten erweiterten Di\"edergruppen enthalten sein.
% [sic] printed with \"o here (Di\"odergruppe) vs \"e elsewhere on this page (Di\"edergruppen) — reproduced as read.

Von den erweiterten Di\"edergruppen entsteht aber $D_2''$ aus $D_2$ durch Inversion, während
\[
\begin{aligned}
D_2' = &\begin{pmatrix}1, & 0\\ 0, & 1\end{pmatrix},\ \begin{pmatrix}\sqrt{i}, & 0\\ 0, & +i\sqrt{i}\end{pmatrix},\ \begin{pmatrix}i, & 0\\ 0, & -i\end{pmatrix},\ \begin{pmatrix}i\sqrt{i}, & 0\\ 0, & +\sqrt{i}\end{pmatrix},\\
&\begin{pmatrix}0, & -1\\ 1, & 0\end{pmatrix},\ \begin{pmatrix}0, & -i\sqrt{i}\\ \sqrt{i}, & 0\end{pmatrix},\ \begin{pmatrix}0, & i\\ i, & 0\end{pmatrix},\ \begin{pmatrix}0, & -\sqrt{i}\\ i\sqrt{i}, & 0\end{pmatrix}
\end{aligned}
\]
%CONT>>


% ===== scan leaf 261, printed 240 =====
% !PRINTED 240
\weberpage{240}
%<<CONT
durch Composition von $D_2$ mit
\[
j_1 = \begin{pmatrix} 0, & -\sqrt{i} \\ i\sqrt{i}, & 0 \end{pmatrix}
\]
entsteht, also
\[
D_2' = D_2 + D_2 j_1
\]
ergiebt. Nun ist allgemein:
\[
j_1^{-1} \begin{pmatrix} \alpha, & \beta \\ \gamma, & \delta \end{pmatrix} j_1 = \begin{pmatrix} \delta, & \gamma i \\ -\beta i, & \alpha \end{pmatrix},
\]
und daraus schliesst man nach \S.\,56, (4), dass $j_1^{-1} T j_1 = T$ ist, dass also aus der Tetra\"edergruppe zwei erweiterte Gruppen $T' = T+Tj$ und $T'' = T+Tj_1$ abgeleitet werden können; aber auch nur zwei, weil nach \S.\,56 die Tetra\"edergruppe durch eine darin enthaltene Vierergruppe völlig bestimmt ist. Denn wenn wir statt der Erweiterung $D_2'$ die Erweiterung der Vierergruppe zu einer der Gruppen (11) wählen, so erhalten wir eine erweiterte Tetra\"edergruppe $T'''$, die durch Transformation mit $\varphi_1$ oder $\varphi_2$ in $T''$ übergeht.

Da die Substitution
\[
\begin{pmatrix} 0, & e^{\frac{\pi i}{8}} \\ e^{-\frac{\pi i}{8}}, & 0 \end{pmatrix}
\]
mit der Octa\"edergruppe \S.\,57, (8) nicht vertauschbar ist, so lässt sich aus der Octa\"edergruppe ausser durch Inversion keine weitere Gruppe zweiter Art ableiten, und dasselbe gilt von der Ikosa\"edergruppe. Man muss, um diese Schlussweise anzuwenden, die Ikosa\"edergruppe so transformiren, dass eine der darin enthaltenen Vierergruppen die einfachste Gestalt annimmt, was etwas weitläufig, aber durchaus nicht schwierig ist, und hier nicht weiter ausgeführt werden soll.

Wir erhalten also noch folgende Gruppen zweiter Art:
\[
\begin{aligned}
&\text{IV.}\quad 1)\ T' = T+Tj, \qquad 2)\ T'' = T+Tj_1.\\
&\text{V.}\qquad\qquad O' = O+Oj.\\
&\text{IV.}\qquad\qquad J' = J+Jj.
\end{aligned}
\]
% [sic] the third line is printed "IV." again (as the first); expected "VI." given the sequence — reproduced as printed.


% ===== scan leaf 262, printed 241 =====
% !PRINTED 241
\weberpage{241}
Hierin sind die 32 Symmetriesysteme der Krystallographen enthalten. Es sind die Gruppen:
\[
\begin{array}{ll}
C_1,\ C_1',\ C_1'' & D_2,\ D_2',\ D_2''\\
C_2,\ C_2',\ C_2'',\ C_3''' & D_3,\ D_3',\ D_3''\\ % [sic: Weber prints C_3''' in the C_2 row; C_2''' is mathematically expected.]
C_3,\ C_3',\ C_3'',\ C_3''' & D_4,\ D_4'\\
C_4,\ C_4',\ C_4'' & D_6,\ D_6'\\
C_6,\ C_6',\ C_6'' & T,\ T',\ T''\\
 & O,\ O'.
\end{array}
\]
Die übrigen Poly\"edergruppen sind in der Krystallographie durch das krystallographische Gesetz der rationalen Indices ausgeschlossen\footnote{Vgl. \gesperrt{Sch\"onfliess}, Krystallsysteme u. Krystallstructur. Leipzig 1891.}.


% ===== scan leaf 263, printed 242 =====
\weberpage{242}% page number not printed on this page
\pdfbookmark[1]{Neunter Abschnitt. Congruenzgruppen.}{bm73}
\pdfbookmark[2]{§ 63. Functionen-Congruenzen}{bm74}
% !PRINTED none
% Section-opening page: no page number is visible in the running head of any of
% the three bands. Page carries the Abschnitt title, the \S.\,63 sub-heading, the
% opening paragraphs of the section, and a footnote citing the literature.

\gesperrt{Neunter Abschnitt.}

\gesperrt{Congruenzgruppen.}

\S.\,63.

\gesperrt{Functionen-Congruenzen.}

Aus den linearen Substitutionen, deren Bildung und Zusammensetzung wir in den früheren Abschnitten kennen gelernt haben, lassen sich noch eine ganz andere Art endlicher Gruppen ableiten, die \gesperrt{Congruenzgruppen}.
% [sic] Weber prints the plural verb "lassen" with the singular head "eine ... Art"; the source wording is preserved.

Diese Congruenzgruppen haben ein mannigfaches Interesse. Sie stammen zunächst aus der Theorie der elliptischen Functionen, wo sie sich als \gesperrt{Galois}'sche Gruppen der Modulargleichungen einstellen. Sie sind aber auch für die allgemeine Gruppentheorie von Wichtigkeit, weil sie uns ein Mittel geben, ganze Reihen von \gesperrt{einfachen Gruppen} zu bilden. Dies ist um so bemerkenswerther, als, wie wir im vierten Abschnitte gesehen haben, die einfachen Gruppen, wenigstens unter den niedrigeren Gradzahlen, sehr selten sind.

Die Theorie dieser Congruenzgruppen wird nicht nur ausserordentlich verallgemeinert, sondern die herrschenden Gesetze treten weit schärfer hervor, wenn man sich auf eine von \gesperrt{Gauss} herrührende Erweiterung des Congruenzbegriffes stützt, die seitdem mehrfach für die Probleme der Gruppentheorie, besonders von \gesperrt{Galois}, ausgebildet und angewandt worden ist\footnote{\gesperrt{Galois}, \glqq Sur la théorie des nombres\grqq{}. Bulletin des sciences math. de Ferussac. 1830. (Vergl. die Note zu \S.\,151 des ersten Bandes.) --- \gesperrt{Schönemann}, Grundzüge einer allgemeinen Theorie der höheren Congruenzen, deren Modulus eine reelle Primzahl ist. (Crelle's Journ. f. Mathematik, Bd. 31, 1846.) --- \gesperrt{Dedekind}, Abriss einer Theorie der höheren Congruenzen in Bezug auf einen reellen Primzahlmodulus. (Crelle's Journ. f. Mathematik, Bd. 54, 1856.)}.


% ===== scan leaf 264, printed 243 =====
% !PRINTED 243
\weberpage{243}
Um eine Grundlage für diese Theorie zu gewinnen, betrachten wir eine ganze Function einer veränderlichen Grösse $t$

\weq{1}{f(t) = a_0 t^n + a_1 t^{n-1} + \cdots + a_{n-1} t + a_n,}

deren Co\"efficienten ganze Zahlen sein sollen. Wir wählen ausserdem eine Primzahl $p$ als Modulus, und nehmen an, $a_0$ sei durch $p$ nicht theilbar.

Zwei ganze Functionen von $t$, in denen entsprechende Co\"efficienten nach dem Modul $p$ congruent sind, sollen selbst \gesperrt{nach dem Modul $p$ congruent} genannt werden.

Die Function $f(t)$ heisst \gesperrt{nach dem Modul $p$ reducibel}, wenn es zwei ganze ganzzahlige Functionen $\varphi(t)$, $\psi(t)$ giebt, in deren jeder wenigstens ein von $t$ abhängiges Glied einen durch $p$ nicht theilbaren Co\"efficienten hat, so dass

\weq{2}{f(t) \equiv \varphi(t)\,\psi(t)\ (\mathrm{mod}\ p)}

ist. In $\varphi(t)$ und $\psi(t)$ kann man alle Glieder, deren Co\"efficienten durch $p$ theilbar sind, weglassen, und dann muss der Grad von $\varphi(t)\,\psi(t)$ mit dem Grade von $f(t)$ übereinstimmen. Es müssen also die Grade von $\varphi(t)$ und $\psi(t)$ niedriger als $n$ sein. Giebt es keine solche Functionen $\varphi(t)$, $\psi(t)$, so heisst $f(t)$ \gesperrt{nach dem Modul $p$ irreducibel}.

Eine nach dem Modul $p$ irreducible Function ist gewiss immer im Körper der rationalen Zahlen absolut irreducibel. Das Umgekehrte ist aber nicht nothwendig.

So ist z.\,B. die Function $2^{\mathrm{ten}}$ Grades $t^2 - R$ reducibel nach $p$, wenn $R$ quadratischer Rest von $p$ ist; denn ist $a^2 \equiv R\ (\mathrm{mod}\ p)$, so ist

\[
t^2 - R \equiv (t-a)(t+a)\ (\mathrm{mod}\ p).
\]

Dagegen ist $t^2 - N$, wenn $N$ Nichtrest von $p$ ist, irreducibel, weil sonst $t^2 - N$ für ein rationales $t$ durch $p$ theilbar werden müsste.

Ein anderes Beispiel bieten die Functionen

\[
t^2 + t + 1, \quad t^3 + t + 1,
\]

die für den Modul $2$ irreducibel sind. Denn wären sie reducibel, so müsste einer der Factoren linear sein, und es müsste eine ganze Zahl geben, die, für $t$ eingesetzt, diese Functionen zu geraden Zahlen macht. Das aber ist offenbar unmöglich, da beide Functionen für $t = 0$ und $t = 1$ ungerade Zahlen sind.


% ===== scan leaf 265, printed 244 =====
% !PRINTED 244
\weberpage{244}
Bedeutet

\weq{3}{P = t^n + p_1 t^{n-1} + p_2 t^{n-2} + \cdots + p_{n-1} t + p_n}

eine nach dem Modul $p$ irreducible Function $n^{\mathrm{ten}}$ Grades, in der wir der Einfachheit halber den Co\"efficienten der höchsten Potenz $t^n$ gleich $1$ annehmen, so lässt sich aus jeder anderen ganzen Function $F(t)$ mit ganzzahligen Co\"efficienten ein Rest $\Phi(t)$ ableiten, dessen Grad niedriger als $n$ ist, indem man (nach \S.\,3 des ersten Bandes)

\weq{4}{F(t) = QP + \Phi(t)}

setzt. $\Phi(t)$ hat hierin ganzzahlige Co\"efficienten, und wir bezeichnen ihre Beziehung zur Function $F(t)$ als eine Congruenz nach dem Modul $P$.

\gesperrt{Es heissen also zwei Functionen $F(t)$, $F_1(t)$, die denselben Rest $\Phi(t)$ haben, congruent nach dem Modul $P$}, was durch die Formel

\[
F(t) \equiv F_1(t)\ (\mathrm{mod}\ P)
\]

ausgedrückt wird. Damit ist gleichbedeutend, dass $F(t) - F_1(t)$ durch $P$ theilbar ist.

Wenn zwei Functionen $F(t)$ und $F_1(t)$ nicht gleichen Rest geben, sondern zwei Reste, die nach dem Modul $p$ congruent sind, so heissen die Functionen $F$, $F_1$ congruent nach dem Doppelmodul $P$, $p$, und man drückt dies durch eine Formel so aus:

\weq{5}{F(t) \equiv F_1(t)\ (\mathrm{mod}\ P, p).}

Für diese Art der Congruenz gilt der Satz:

\gesperrt{1. Das Product zweier Functionen $F(t)$, $F_1(t)$ kann nicht mit Null congruent sein, wenn nicht der eine Factor mit Null congruent ist.}

Der Beweis ergiebt sich aus dem Algorithmus des grössten gemeinschaftlichen Theilers. Ist nämlich $P_1$ eine ganze Function von niedrigerem Grade als $P$, die nicht nach dem Modul $p$ mit Null congruent ist, so kann man eine Reihe von eben solchen Functionen $P_2$, $P_3$, $\ldots$ von abnehmendem Grade, und die Quotienten $Q_1$, $Q_2$, $\ldots$ gleichfalls als ganze Functionen so bestimmen, dass

\weq{6}{P \equiv Q_1 P_1 + P_2,\quad P_1 \equiv Q_2 P_2 + P_3, \ldots,\quad P_{\nu-2} \equiv Q_{\nu-1} P_{\nu-1} + P_\nu\ (\mathrm{mod}\ p)}

wird, und die Reihe dieser Gleichungen bricht ab, wenn $P_\nu$ eine Constante (ganze Zahl) geworden ist. Diese Constante $P_\nu$ kann
%CONT>>


% ===== scan leaf 266, printed 245 =====
% !PRINTED 245
\weberpage{245}
\pdfbookmark[2]{§ 64. Congruenzkörper}{bm75}
%<<CONT
aber nicht $\equiv 0\ (\mathrm{mod}\ p)$ sein; denn sonst liesse sich aus (6) eine Congruenz ableiten:

\[
P \equiv T P_{\nu-1}\ (\mathrm{mod}\ p),
\]

in der $T$ eine nicht constante Function von $t$ ist, und $P$ wäre nicht irreducibel nach dem Modul $p$.

Ist nun $Q$ irgend eine ganze Function von $t$, die der Bedingung

\[
Q P_1 \equiv 0\ (\mathrm{mod}\ P, p)
\]

genügt, so folgt aus (6) durch Multiplication mit $Q$:

\[
Q P_2 \equiv 0,\ Q P_3 \equiv 0,\ \ldots,\ Q P_\nu \equiv 0\ (\mathrm{mod}\ P, p),
\]

also $Q \equiv 0$. Und wenn wir also für $P_1$, $Q$ die Reste der Functionen $F(t)$, $F_1(t)$ setzen, so ist hiermit der Satz 1. bewiesen.

Die Reste aller Functionen $F(t)$ nach dem Modul $P$ sind von der Form

\weq{7}{X = X(t) = x_0 + x_1 t + \cdots + x_{n-1} t^{n-1},}

und wenn man nur die nach dem Modul $p$ incongruenten unter ihnen haben will, so genügt es, die Co\"efficienten $x_0, x_1, \ldots x_{n-1}$ die Reihe der Zahlen $0, 1, 2, \ldots, p-1$ durchlaufen zu lassen.

\gesperrt{Es giebt also $p^n$ und nicht mehr nach den Moduln $P$, $p$ incongruente Functionen.}

\S.\,64.

\gesperrt{Congruenzkörper.}

Es ist nur eine andere Ausdrucksweise für die im vorigen Paragraphen durchgeführten Betrachtungen, wenn man mit \gesperrt{Galois} eine imaginäre Zahl $\varepsilon$ einführt, die der Congruenz

\weq{1}{P(\varepsilon) \equiv 0\ (\mathrm{mod}\ p)}

genügt. Eine ganze rationale Zahl dieser Art giebt es nicht, sobald $n > 1$ ist, weil sonst $P(t) - P(\varepsilon) = (t-\varepsilon)Q$ durch $t-\varepsilon$ theilbar, mithin $P(t) \equiv (t-\varepsilon)Q$, und $P(t)$ nicht irreducibel wäre. Gleichwohl kann man, wie man in der gewöhnlichen Zahlenlehre die imaginäre Einheit $i = \sqrt{-1}$ einführt, ein solches Symbol $\varepsilon$ in der Rechnung benutzen. Man rechnet damit, sofern es sich um Addition, Subtraction und Multiplication handelt, nach den Regeln der Buchstabenrechnung, und kann dabei nach Belieben die Gleichung $P(\varepsilon) = 0$ benutzen. Man kann sogar unter $\varepsilon$ geradezu eine Wurzel der Gleichung $P(x) = 0$
%CONT>>


% ===== scan leaf 267, printed 246 =====
% !PRINTED 246
\weberpage{246}
%<<CONT
im gewöhnlichen Sinne verstehen. Alle Zahlen, auf die man hierbei kommt, lassen sich auf die Form

\weq{2}{\alpha = a_0 + a_1\varepsilon + \cdots + a_{n-1}\varepsilon^{n-1}}

bringen. Da bei der Rechnung mit solchen Zahlen der Modul $p$ immer festgehalten wird, so wollen wir zwei Zahlen, die nach dem Modul $p$ congruent sind, geradezu als \gesperrt{gleich} bezeichnen. Objecte der Rechnung sind dann eigentlich nicht die einzelnen Zahlen selbst, sondern die aus allen unter einander congruenten Zahlen bestehenden \gesperrt{Zahlclassen}. Diese Zahlen $\alpha$ werden die \gesperrt{Galois'schen Imaginären} genannt. Das zu einem bestimmten $\varepsilon$ gehörige System solcher imaginärer Zahlen bezeichnen wir der Abkürzung wegen mit $\mathfrak{C}$.

Wir haben dann zunächst den Satz:

\gesperrt{2. Es giebt $p^n$ und nicht mehr verschiedene Zahlen in $\mathfrak{C}$.}

Aus dem Satze 1. aber folgt noch:

\gesperrt{3. Das Product von zwei oder mehr Zahlen aus $\mathfrak{C}$ ist dann und nur dann gleich Null, wenn wenigstens einer der Factoren gleich Null ist.}

Wir erhalten das vollständige Zahlensystem $\mathfrak{C}$, wenn man die rationalen Zahlen $a_0, a_1, \ldots, a_{n-1}$ in (2) je ein volles Restsystem nach dem Modul $p$ durchlaufen lässt. Jedes System $\mathfrak{C}$ enthält die Reste der natürlichen Zahlen $0, 1, \ldots, p-1$, und für $n = 1$ ist $\mathfrak{C}$ mit diesem Systeme identisch.

Ist $\alpha$ eine feste von Null verschiedene Zahl in $\mathfrak{C}$, so durchläuft $\alpha\xi$ zugleich mit $\xi$ das volle System $\mathfrak{C}$. Denn aus 3. folgt, dass $\alpha\xi$ nur dann gleich $\alpha\xi'$ sein kann, wenn $\xi = \xi'$ ist. Daraus folgt:

\gesperrt{4. Sind $\alpha, \beta$ zwei gegebene Zahlen in $\mathfrak{C}$ und $\alpha$ von Null verschieden, so giebt es eine und nur eine Zahl $\gamma$ in $\mathfrak{C}$, die der Bedingung}

\[
\alpha\gamma = \beta
\]

\gesperrt{genügt.}

Damit ist auch die Operation der Division in dem Systeme $\mathfrak{C}$ als erlaubt nachgewiesen. Wir bezeichnen die Zahl $\gamma$, deren Existenz in 4. ausgesprochen ist, mit

\[
\beta : \alpha \quad \text{oder} \quad \frac{\beta}{\alpha}.
\]


% ===== scan leaf 268, printed 247 =====
% !PRINTED 247
\weberpage{247}
Man kann das System $\mathfrak{C}$ einen \gesperrt{endlichen Körper} nennen, da es nur eine endliche Anzahl von Zahlen enthält, die das charakteristische Merkmal eines Körpers, nämlich die unbeschränkte Ausführbarkeit der Rechenoperationen, ausgenommen die Division durch Null, aufweisen (Bd.\,I, \S.\,139)\textsuperscript{1)}. Wir wollen daher $\mathfrak{C}$ einen \gesperrt{Congruenzkörper} nennen.

Es soll $n$ der \gesperrt{Grad} und $p$ der \gesperrt{Modul} des Körpers $\mathfrak{C}$ heissen.

Ist $\alpha$ eine von Null verschiedene Zahl in $\mathfrak{C}$, so durchläuft das Product $\alpha\xi = \eta$ zugleich mit $\xi$ das volle Zahlensystem $\mathfrak{C}$. Schliessen wir $\xi = 0$ aus, so kommt auch $\eta = 0$ nicht vor, und das Product $\Pi$ aller $\xi$ stimmt mit dem Producte aller $\eta$ überein und ist von Null verschieden. Durch Multiplication aller Gleichungen $\alpha\xi = \eta$ folgt aber

\[
\alpha^{p^n-1}\Pi = \Pi,
\]

und nach Abwerfung des gemeinsamen Factors $\Pi$ ergiebt sich

\gesperrt{5. der Fermat'sche Satz:}

\weq{3}{\alpha^{p^n-1} = 1.}

Multiplicirt man mit $\alpha$, so erhält man diesen Satz in der Form

\weq{4}{\alpha^{p^n} = \alpha,}

in der er auch noch für $\alpha = 0$ besteht. Ist $\alpha$ von Null verschieden, so giebt es wegen (3) einen \gesperrt{kleinsten} positiven Exponenten $e$, für den

\weq{5}{\alpha^e = 1}

ist, und für den folglich die Potenzen $1, \alpha, \alpha^2, \ldots, \alpha^{e-1}$ von einander verschieden sind. Man sagt dann, $\alpha$ \gesperrt{gehört zum Exponenten} $e$. Dieser Exponent $e$ muss ein Theiler von $p^n - 1$ sein. Denn ist $\alpha^m = 1$ für irgend einen Exponenten $m$, so setzen wir $m = qe + e'$, worin $q$ eine ganze Zahl und $e' < e$ ist. Dann ist auch $\alpha^{e'} = 1$, und folglich muss $e' = 0$, d.\,h. $m$ durch $e$ theilbar sein. Unter diesen Exponenten $m$ findet sich auch $p^n - 1$. Setzen wir also

\[
p^n - 1 = ef,
\]

so wird $\alpha^h$ immer dann zum Exponenten $e$ gehören, wenn $h$ relativ prim zu $e$ ist. Denn dann und nur dann ist $he$ das kleinste positive, durch $e$ theilbare Vielfache von $h$.

\textsuperscript{1)} Man sehe des Verfassers Abhandlung: Die allgemeinen Grundlagen der \gesperrt{Galois'schen} Gleichungstheorie. Mathematische Annalen, Bd.\,43.


% ===== scan leaf 269, printed 248 =====
% !PRINTED 248
\weberpage{248}
Giebt es also überhaupt eine zum Exponenten $e$ gehörige Zahl $\alpha$, so giebt es so viele wie relative Primzahlen zu $e$, die positiv und kleiner als $e$ sind, eine Zahl, die wir schon früher mit $\varphi(e)$ bezeichnet haben, und die, wenn $e$ alle Divisoren von $p^n - 1$ durchläuft, der Relation

\weq{6}{\textstyle\sum \varphi(e) = p^n - 1}

genügt (Bd.\,I, \S.\,132, 133).

Ist $\psi(e)$ die Anzahl der Zahlen $\alpha$, die zum Exponenten $e$ gehören, so ist $\psi(e) = \varphi(e)$ oder $= 0$, und da die Anzahl aller Zahlen $\alpha$ gleich $p^n - 1$ ist und jede Zahl $\alpha$ zu einem und nur zu einem Exponenten gehört, so ist auch

\weq{7}{\textstyle\sum \psi(e) = p^n - 1,}

und aus (6) und (7) folgt, dass $\psi(e)$ \gesperrt{immer gleich} $\varphi(e)$ ist.

Es giebt also $\varphi(p^n - 1)$ Zahlen $\gamma$ in $\mathfrak{C}$, die zum Exponenten $p^n - 1$ gehören, und die folglich die Eigenschaft haben, dass die Potenzen

\weq{8}{1, \gamma, \gamma^2, \ldots, \gamma^{p^n-2}}

alle von einander verschieden sind. Durch diese Reihe ist daher die Gesammtheit der von Null verschiedenen Zahlen in $\mathfrak{C}$ erschöpft.

Solche Zahlen $\gamma$ heissen \gesperrt{primitive Wurzeln von $\mathfrak{C}$}.

Unter den Zahlen in $\mathfrak{C}$ giebt es solche, die als Quadrat einer anderen Zahl in $\mathfrak{C}$ darstellbar sind, die wir \gesperrt{Quadrate} nennen, und andere, bei denen dies nicht zutrifft, die wir \gesperrt{Nichtquadrate} nennen.

Wenn $p = 2$ ist, so ist jede Zahl in $\mathfrak{C}$ ein Quadrat, wie die Formel (4) zeigt.

Wenn aber $p$ ungerade ist, so besteht die eine Hälfte der von Null verschiedenen Zahlen in $\mathfrak{C}$ aus Quadraten, die andere aus Nichtquadraten.

Denn zwei entgegengesetzte Zahlen, wie $+\alpha$ und $-\alpha$, sind bei ungeradem $p$ von einander verschieden, und geben trotzdem dasselbe Quadrat. Wenn man also die sämmtlichen Zahlen von $\mathfrak{C}$ zum Quadrat erhebt, so erhält man höchstens $\frac{1}{2}(p^n - 1)$ verschiedene Quadrate. Nun kann andererseits $\beta^2$ nur dann gleich $\alpha^2$ sein, wenn $\beta = \pm\alpha$ ist; denn aus $\beta^2 - \alpha^2 = (\beta - \alpha)(\beta + \alpha) = 0$ folgt, dass $\beta - \alpha$ oder $\beta + \alpha$ verschwinden muss. Es giebt also wirklich $\frac{1}{2}(p^n - 1)$ Quadrate und ebenso viele Nichtquadrate.


% ===== scan leaf 270, printed 249 =====
% !PRINTED 249
\weberpage{249}
Wenn man die Null mit zu den Quadraten zählt, so erhält man den Satz:

\gesperrt{6. Wenn $p = 2$ ist, so sind alle Zahlen in $\mathfrak{C}$ Quadrate. Ist $p$ ungerade, so giebt es $\frac{1}{2}(p^n + 1)$ Quadrate, $\frac{1}{2}(p^n - 1)$ Nichtquadrate in $\mathfrak{C}$.}

Bei ungeradem $p$ muss eine primitive Wurzel immer ein Nichtquadrat sein und in der Reihe (8) sind die Zahlen mit geraden Exponenten die Quadrate, die mit ungeraden Exponenten die Nichtquadrate.

Das Product und der Quotient zweier Quadrate ist offenbar wieder ein Quadrat.

Daraus folgt, dass das Product aus einem Nichtquadrat und einem von Null verschiedenen Quadrate ein Nichtquadrat ist. Denn ist das Product $\alpha\beta$ ein Quadrat und der eine Factor $\alpha$ ein Quadrat, so muss auch $\beta = \alpha\beta : \alpha$ ein Quadrat sein.

Weiter schliesst man daraus, dass das Product von zwei Nichtquadraten ein Quadrat ist. Denn bedeutet $\beta$ irgend ein Nichtquadrat, so lasse man in dem Producte $\beta\xi$ den Factor $\xi$ die sämmtlichen von Null verschiedenen Zahlen von $\mathfrak{C}$ durchlaufen. Das Product durchläuft dann dieselbe Zahlenreihe, und da für alle quadratischen $\xi$ das Product $\beta\xi$ Nichtquadrat ist, so muss es für die nichtquadratischen $\xi$ Quadrat sein.

Alles dies ergiebt sich auch sehr einfach aus der Darstellung (8) der Zahlen von $\mathfrak{C}$ durch eine primitive Wurzel.

Wir schliessen diese allgemeinen Betrachtungen mit dem Satze, den wir später brauchen werden:

\gesperrt{7. Jede nicht quadratische Zahl ist die Summe von zwei Quadraten.}

Hierbei ist $p$ als ungerade vorauszusetzen, da es nur dann nichtquadratische Zahlen giebt.

Es lässt sich dann zunächst jede Zahl $\beta$ als Differenz zweier Quadrate darstellen, wie die Identität

\[
\left(\beta + \tfrac{1}{4}\right)^2 - \left(\beta - \tfrac{1}{4}\right)^2 = \beta
\]

zeigt. Ist nun $-1$ ein Quadrat in $\mathfrak{C}$, so setze man

\[
\left(\beta + \tfrac{1}{4}\right)^2 = \xi^2, \quad -\left(\beta - \tfrac{1}{4}\right)^2 = \eta^2,
\]

und erhält

\[
\beta = \xi^2 + \eta^2.
\]

Ist aber $-1$ und also auch $p - 1$ Nichtquadrat, so suche man in der Reihe der natürlichen Zahlen $1, 2, \ldots, p - 1$, deren
%CONT>>


% ===== scan leaf 271, printed 250 =====
% !PRINTED 250
\weberpage{250}
\pdfbookmark[2]{§ 65. Congruenzgruppen im Körper E}{bm76}
%<<CONT
erstes Glied ein Quadrat und deren letztes ein Nichtquadrat ist, ein quadratisches Glied auf, auf welches ein Nichtquadrat folgt.

Ist dann also $a = \xi^2$ ein Quadrat und $\xi^2 + 1$ ein Nichtquadrat, so ist, wenn $\beta$ ein Nichtquadrat ist, $\beta : \xi^2 + 1 = \eta^2$ ein Quadrat, und daraus folgt:
\[ \beta = \xi^2 \eta^2 + \eta^2 \]
w. z. b. w.

\begin{center}
\S.\,65.

\gesperrt{Congruenzgruppen im K\"orper $\mathfrak{C}$.}
\end{center}

In jedem Congruenzk\"orper $\mathfrak{C}$ kann man eine endliche Gruppe von linearen Substitutionen bilden, wenn man in
\weq{1}{A = \begin{pmatrix} \alpha, & \beta \\ \gamma, & \delta \end{pmatrix}}
die Elemente $\alpha, \beta, \gamma, \delta$ alle Zahlen von $\mathfrak{C}$ durchlaufen l\"asst, und wenn man je zwei dieser Substitutionen nach der Vorschrift
\weq{2}{\begin{pmatrix} \alpha, & \beta \\ \gamma, & \delta \end{pmatrix} \begin{pmatrix} \alpha', & \beta' \\ \gamma', & \delta' \end{pmatrix} = \begin{pmatrix} \alpha\alpha' + \beta\gamma', & \alpha\beta' + \beta\delta' \\ \gamma\alpha' + \delta\gamma', & \gamma\beta' + \delta\delta' \end{pmatrix}}
componirt. Der Grad dieser Gruppe ist $p^{4n}$, wenn $p$ der Modul und $n$ der Grad von $\mathfrak{C}$ ist.

Darin ist aber eine Gruppe niedrigeren Grades enthalten. Nennen wir $\alpha\delta - \beta\gamma$ die Determinante der Substitution (1), so folgt aus (2), dass die Determinante der aus $A$ und $A'$ componirten Substitution $AA'$ gleich dem Producte der Determinanten von $A$ und $A'$ ist. Wenn wir daher den Zahlen $\alpha, \beta, \gamma, \delta$ die Bedingung auferlegen, dass
\weq{3}{\alpha\delta - \beta\gamma = 1}
sein soll, so bilden auch diese Elemente $A$ eine Gruppe. Der Grad dieser Gruppe ist gleich der Anzahl der L\"osungen von (3). Um diese Anzahl zu finden, bemerken wir zun\"achst, dass wir f\"ur $\alpha, \beta$ irgend zwei Zahlen aus $\mathfrak{C}$ setzen k\"onnen, die nicht beide gleich Null sind. Denn ist etwa $\alpha$ von Null verschieden, so kann man $\gamma = 0, \delta = 1 : \alpha$ setzen und hat so die Bedingung (3) erf\"ullt. Die Anzahl der brauchbaren Zahlenpaare $\alpha, \beta$ ist also $p^{2n} - 1$. Hat man aber zu einem Zahlenpaare $\alpha, \beta$ eine L\"osung $\gamma_0, \delta_0$ von (3) gefunden, so m\"ussen alle \"ubrigen der Bedingung
\[ \alpha (\delta - \delta_0) = \beta (\gamma - \gamma_0) \]
%CONT>>


% ===== scan leaf 272, printed 251 =====
% !PRINTED 251
\weberpage{251}
%<<CONT
gen\"ugen, und wenn man also $\gamma - \gamma_0 = \alpha\xi$ setzt (falls $\alpha$ von Null verschieden ist), so folgt $\delta - \delta_0 = \beta\xi$, also
\[ \gamma = \gamma_0 + \alpha\xi, \quad \delta = \delta_0 + \beta\xi, \]
und umgekehrt gen\"ugt auch jedes in dieser Form enthaltene Zahlenpaar $\gamma, \delta$. Dies gilt offenbar auch noch in dem ausgenommenen Falle $\alpha = 0$. Da wir nun $p^n$ verschiedene Zahlen f\"ur $\xi$ setzen k\"onnen, so ergiebt sich der Grad unserer Gruppe gleich
\weq{4}{p^n (p^{2n} - 1).}

Ist $p$ ungerade, so k\"onnen wir eine noch kleinere Gruppe ableiten:

Ist n\"amlich $p = 2$, so ist eine Zahl $\alpha$ von $-\alpha$ nicht verschieden, es sind also auch die Elemente
\weq{5}{\begin{pmatrix} \alpha, & \beta \\ \gamma, & \delta \end{pmatrix}, \quad \begin{pmatrix} -\alpha, & -\beta \\ -\gamma, & -\delta \end{pmatrix}}
mit einander identisch. Wenn aber $p$ ungerade ist, so sind die beiden Elemente (5) von einander verschieden, und das ganze System dieser Substitutionen zerf\"allt in Paare von der Form (5).

Wenn wir zwei solche Paare nehmen und componiren je ein Element des einen Paares mit je einem Element des anderen nach der Regel (2), so erhalten wir nur zwei verschiedene Elemente, die wieder ein solches Paar bilden. Diese Paare k\"onnen also selbst als Elemente einer neuen Gruppe vom Grade
\weq{6}{\dfrac{p^n (p^{2n} - 1)}{2}}
aufgefasst werden. Wir k\"onnen uns auch so ausdr\"ucken, dass zwei Elemente $A$, deren entsprechende Zahlen nur durch das Vorzeichen unterschieden sind, als nicht verschieden anzusehen sind.

Die so definirte Gruppe, deren Grad f\"ur $p = 2$ durch (4) und f\"ur ein ungerades $p$ durch (6) ausgedr\"uckt ist, wollen wir \gesperrt{die zum K\"orper $\mathfrak{C}$ geh\"orige Congruenzgruppe nennen und mit $E$ bezeichnen}.

F\"ur die genauere Untersuchung dieser Gruppe ist es von Wichtigkeit, ein System erzeugender Substitutionen zu kennen. Stellen wir nach \S.\,64, (2) die Zahlen von $\mathfrak{C}$ in der Form dar:
\weq{7}{\xi = x_0 + x_1 \varepsilon_1 + \cdots + x_{n-1} \varepsilon^{n-1},} % [sic] first power printed as subscript \varepsilon_1, not \varepsilon or \varepsilon^1
worin die $x_i$ rationale, nach dem Modul $p$ genommene Zahlen
%CONT>>


% ===== scan leaf 273, printed 252 =====
% !PRINTED 252
\weberpage{252}
%<<CONT
bedeuten, so erhalten wir ein System erzeugender Substitutionen in der Form:
\weq{8}{A_0 = \begin{pmatrix} 0, & 1 \\ -1, & 0 \end{pmatrix}, \quad B_h = \begin{pmatrix} 1, & 0 \\ \varepsilon^h, & 1 \end{pmatrix}, \quad h = 0, 1, \ldots, n-1.}

Um dies nachzuweisen, bemerken wir, dass
\[ \begin{pmatrix} 1, & 0 \\ \gamma, & 1 \end{pmatrix} \begin{pmatrix} 1, & 0 \\ \gamma', & 1 \end{pmatrix} = \begin{pmatrix} 1, & 0 \\ \gamma + \gamma', & 1 \end{pmatrix} \]
ist; und durch wiederholte Anwendung hiervon folgt, wenn
\[ \gamma = c_0 + c_1 \varepsilon + \cdots + c_{n-1} \varepsilon^{n-1} \]
mit ganzen rationalen Co\"efficienten $c$ eine beliebige Zahl in $\mathfrak{C}$ ist,
\weq{9}{\begin{pmatrix} 1, & 0 \\ \gamma, & 1 \end{pmatrix} = B_0^{c_0} B_1^{c_1} \cdots B_{n-1}^{c_{n-1}}.}

Ferner folgt durch Zusammensetzung mit $A_0$:
\weq{10}{\begin{gathered} \begin{pmatrix} 0, & 1 \\ -1, & 0 \end{pmatrix}\begin{pmatrix} 1, & 0 \\ \alpha, & 1 \end{pmatrix} = \begin{pmatrix} \alpha, & 1 \\ -1, & 0 \end{pmatrix}, \quad \begin{pmatrix} 1, & 0 \\ \delta, & 1 \end{pmatrix}\begin{pmatrix} 0, & 1 \\ -1, & 0 \end{pmatrix} = \begin{pmatrix} 0, & 1 \\ -1, & \delta \end{pmatrix} \\ \begin{pmatrix} 0, & 1 \\ -1, & 0 \end{pmatrix}\begin{pmatrix} 1, & 0 \\ -\beta, & 1 \end{pmatrix}\begin{pmatrix} 0, & -1 \\ 1, & 0 \end{pmatrix} = \begin{pmatrix} 1, & \beta \\ 0, & 1 \end{pmatrix}. \end{gathered}}

Daraus geht hervor, dass man aus den Elementen (8) alle Substitutionen von der Form
\[ \begin{pmatrix} \alpha, & 1 \\ -1, & 0 \end{pmatrix}, \quad \begin{pmatrix} 1, & \beta \\ 0, & 1 \end{pmatrix}, \quad \begin{pmatrix} 1, & 0 \\ \gamma', & 1 \end{pmatrix}, \quad \begin{pmatrix} 0, & 1 \\ -1, & \delta \end{pmatrix} \]
zusammensetzen kann, worin $\alpha, \beta, \gamma, \delta$ beliebige Zahlen in $\mathfrak{C}$ sind. Ferner ist
\weq{11}{\begin{pmatrix} \alpha, & 1 \\ -1, & 0 \end{pmatrix}\begin{pmatrix} 1, & \beta \\ 0, & 1 \end{pmatrix}\begin{pmatrix} 1, & 0 \\ \gamma, & 1 \end{pmatrix} = \begin{pmatrix} \alpha + \gamma + \alpha\beta\gamma, & \alpha\beta + 1 \\ -\beta\gamma - 1, & -\beta \end{pmatrix}.}

Nimmt man $\beta$ von Null verschieden an, so kann man $\alpha$ und $\gamma$ so w\"ahlen, dass $\alpha\beta + 1$, $-\beta\gamma - 1$ beliebige Zahlen werden, und nach (11) kann man also alle Substitutionen $\begin{pmatrix} \alpha, & \beta \\ \gamma, & \delta \end{pmatrix}$, worin $\delta$ von Null verschieden ist, aus $A_0, B_h$ zusammensetzen. Die Beschr\"ankung, dass $\delta$ von Null verschieden sei, kann aber nach der Formel
\[ \begin{pmatrix} \alpha, & \beta \\ \gamma, & 0 \end{pmatrix} = \begin{pmatrix} -\beta, & \alpha \\ 0, & \gamma \end{pmatrix}\begin{pmatrix} 0, & -1 \\ 1, & 0 \end{pmatrix} \]
fallen gelassen werden.


% ===== scan leaf 274, printed 253 =====
% !PRINTED 253
\weberpage{253}
Die Gruppe $E$ enth\"alt einen Theiler vom Index $p^n + 1$, der aus allen Substitutionen der Form
\[ \begin{pmatrix} \alpha, & \beta \\ 0, & \delta \end{pmatrix} \]
besteht, worin $\beta$ jede beliebige, $\alpha$ jede von Null verschiedene Zahl in $\mathfrak{C}$ bedeutet, und $\delta = \alpha^{-1}$ ist.

Diesem Theiler entspricht eine der Gruppe $E$ isomorphe Permutationsgruppe von $p^n + 1$ Ziffern, die man so bilden kann:

Der lineare Ausdruck
\weq{12}{\eta = \dfrac{\alpha\xi + \beta}{\gamma\xi + \delta} \quad \text{oder} \quad \xi = \dfrac{\delta\eta - \beta}{-\gamma\eta + \alpha}}
giebt f\"ur jede Zahl $\xi$ aus $\mathfrak{C}$ eine entsprechende Zahl $\eta$, ausgenommen, wenn $\gamma\xi + \delta = 0$ ist. Ebenso giebt es zu jedem $\eta$ ein bestimmtes $\xi$, ausgenommen f\"ur $\gamma\eta - \alpha = 0$. Um diese Ausnahme zu beseitigen, gen\"ugt es, das System $\mathfrak{C}$ durch Hinzuf\"ugung \gesperrt{eines} Elementes, das wir \glqq Unendlich\grqq{} nennen und mit $\infty$ bezeichnen, zu erweitern. Mit diesem Zeichen wird nach den folgenden Regeln gerechnet, worin $\alpha$ irgend ein Element des erweiterten Systemes $E$ bedeutet:
% [sic] Weber druckt gewoehnliches E sowohl hier als auch unten in "Systemes E", obwohl beide Stellen den erweiterten Koerper \mathfrak{C} meinen; originalgetreu bewahrt. Gewoehnliches E bezeichnet sonst die Congruenzgruppe.
\[
\begin{array}{l@{\quad}l}
\alpha + \infty = \alpha - \infty = \infty, & \text{ausser f\"ur } \alpha = \infty \\
\alpha \cdot \infty = \infty & \text{\textquotedbl}\ \ \text{\textquotedbl}\ \ \alpha = 0 \\
\alpha : 0 = \infty & \text{\textquotedbl}\ \ \text{\textquotedbl}\ \ \alpha = 0 \\
\alpha : \infty = 0 & \text{\textquotedbl}\ \ \text{\textquotedbl}\ \ \alpha = \infty.
\end{array}
\]
Dann f\"uhrt das Ergebniss jeder Rechnung mit den vier Species immer zu einer bestimmten Zahl des Systems $E$ und nur die Verbindungen $\infty \pm \infty, 0 . \infty, 0 : 0, \infty : \infty$ bleiben ohne Bedeutung.

Nach (12) entspricht dann der Zahl $\xi = -\delta : \gamma$ die Zahl $\eta = \infty$, und der Zahl $\xi = \infty$ die Zahl $\eta = \alpha : \gamma$, und jeder Substitution $A$ in $E$ entspricht eine bestimmte Permutation der $p^n + 1$ Elemente des erweiterten Systemes $\mathfrak{C}$. Man sieht leicht, dass nur die identische Substitution alle diese Elemente unge\"andert l\"asst, und dass folglich auch zwei verschiedene Substitutionen aus $E$ immer verschiedene Permutationen hervorrufen. Der Isomorphismus ist also einstufig.


% ===== scan leaf 275, printed 254 =====
% !PRINTED 254
\weberpage{254}
\pdfbookmark[2]{§ 66. Einfachheit der Gruppe E}{bm77}

\begin{center}
\S.\,66.\\
\gesperrt{Einfachheit der Gruppe $E$.}
\end{center}

Die erste Frage bei einer eingehenderen Untersuchung der Gruppe $E$ ist die nach einem etwa vorhandenen Normaltheiler. Es l\"asst sich beweisen, dass ein solcher Normaltheiler, von zwei ganz einfachen Ausnahmen abgesehen, nicht vorhanden ist.

Nehmen wir an, es sei $G$ ein Normaltheiler von $E$, der wenigstens eine von der Identit\"at verschiedene Substitution
\weq{1}{A = \begin{pmatrix} \alpha, & \beta \\ \gamma, & \delta \end{pmatrix}}
enth\"alt. Nach dem Begriffe des Normaltheilers ist dann, wenn $T$ ein beliebiges Element in $E$ ist,
\weq{2}{TAT^{-1}}
auch in $G$ enthalten, und es ist zu zeigen, dass man auf diese Weise und durch Zusammensetzung solcher Substitutionen alle Elemente von $E$ ableiten kann, woraus dann folgt, dass $G$ mit $E$ identisch sein muss. Es gen\"ugt aber dazu, das Vorkommen der erzeugenden Substitutionen $A_0$, $B_h$ [\S.\,65, (8)] in $G$ nachzuweisen.

Wir gehen dazu schrittweise vor.

1) In $G$ kommt eine Substitution von der Form
\weq{3}{\begin{pmatrix} 0, & \beta \\ \gamma, & \delta \end{pmatrix}}
vor. Um dies zu zeigen, bilden wir nach (2):
\weq{4}{\begin{pmatrix} \xi, & 1 \\ -1, & 0 \end{pmatrix}\begin{pmatrix} \alpha, & \beta \\ \gamma, & \delta \end{pmatrix}\begin{pmatrix} 0, & -1 \\ 1, & \xi \end{pmatrix} = \begin{pmatrix} \beta\xi+\delta, & \beta\xi^2+(\delta-\alpha)\xi-\gamma \\ -\beta, & -\beta\xi+\alpha \end{pmatrix}.}
Wenn nun $\beta$ von Null verschieden ist, so kann man $\xi$ aus $\beta\xi+\delta=0$ bestimmen und hat das Ziel erreicht. Ist aber $\beta=0$, so ist nicht gleichzeitig $\gamma=0$, $\delta-\alpha=0$, weil sonst $A$ die identische Substitution w\"are; man kann daher $\xi$ so annehmen, dass $(\delta-\alpha)\xi-\gamma$ von Null verschieden ist. Dann hat man durch (4) eine Substitution $A$ in $G$ gebildet, in der $\beta$ nicht Null ist, und auf die man die Formel (4) nun nochmals anwenden kann, um $\alpha$ zu Null zu machen.

2) In $G$ kommt eine Substitution von der Form
\[
\begin{pmatrix} 0, & 1 \\ -1, & \xi \end{pmatrix}
\]
%CONT>>


% ===== scan leaf 276, printed 255 =====
% !PRINTED 255
\weberpage{255}
%<<CONT
vor. Um dies nachzuweisen, bilden wir nach (2):
\weq{5}{\begin{pmatrix} \lambda, & \mu \\ \nu, & \varrho \end{pmatrix}\begin{pmatrix} 0, & \beta \\ \gamma, & \delta \end{pmatrix}\begin{pmatrix} \varrho, & -\mu \\ -\nu, & \lambda \end{pmatrix} = \begin{pmatrix} \mu\varrho\gamma-\lambda\nu\beta-\mu\nu\delta, & -\mu^2\gamma+\lambda^2\beta+\lambda\mu\delta \\ \varrho^2\gamma-\nu^2\beta-\varrho\nu\delta, & -\varrho\mu\gamma+\lambda\nu\beta+\lambda\varrho\delta \end{pmatrix},}
und es ist also nur zu zeigen, dass man $\lambda,\mu,\nu,\varrho$ so bestimmen kann, dass die drei Gleichungen:
\weq{6}{
\begin{aligned}
\lambda\varrho-\mu\nu &= 1\\
\mu\varrho\gamma-\lambda\nu\beta-\mu\nu\delta &= 0\\
-\mu^2\gamma+\lambda^2\beta+\lambda\mu\delta &= 1
\end{aligned}
}
befriedigt sind. Aus den beiden letzten Gleichungen folgt aber, wenn man sie zuerst mit $\mu,\varrho$, dann mit $\lambda,\nu$ multiplicirt und jedesmal addirt, mit Benutzung der ersten Gleichung:
\weq{7}{\lambda\beta+\mu\delta = \varrho, \quad \mu\gamma = \nu,}
und wenn man diese Werthe von $\varrho,\nu$ in die erste Gleichung (6) einsetzt:
\weq{8}{\lambda^2\beta+\lambda\mu\delta-\mu^2\gamma = 1.}
Diese Gleichung kann aber immer befriedigt werden. Denn wegen $\beta\gamma=-1$ kann keine der beiden Zahlen $\beta,\gamma$ verschwinden. Ist nun $\beta$ ein Quadrat, so k\"onnen wir $\mu=0$ setzen und $\lambda$ aus $\lambda^2\beta=1$ bestimmen. Dies ist \gesperrt{immer} m\"oglich bei $p=2$. Ist aber $p$ ungerade, so erhalten wir, wenn $\beta$ ein Nichtquadrat ist, aus (8):
\[
(2\lambda\beta+\delta\mu)^2+\mu^2(4-\delta^2) = 4\beta.
\]
Ist nun zun\"achst $4-\delta^2$ Nichtquadrat, so kann man $\mu$ aus $\mu^2(4-\delta^2)=4\beta$ und dann $\lambda$ aus $2\lambda\beta+\delta\mu=0$ bestimmen. Ist aber $4-\delta^2$ Quadrat, so zerlege man $4\beta$ nach \S.\,64, 7. in die Summe aus zwei Quadraten $4\beta=\sigma^2+\tau^2$ und bestimme $\mu$ und $\lambda$ aus
\[
\mu^2(4-\delta^2) = \sigma^2, \quad 2\lambda\beta+\delta\mu = \tau.
\]
Hat man so $\lambda$ und $\mu$ ermittelt, so erh\"alt man $\varrho$ und $\nu$ aus (7), und dadurch sind die Gleichungen (6) thats\"achlich befriedigt. Damit ist das Ziel erreicht.
% [sic] Der gedruckte Beweis uebergeht den Fall 4-\delta^2=0. Ueber F_5 gilt fuer \beta=\gamma=\delta=2 zwar \beta\gamma=-1, doch hat (8), 2\lambda^2+2\lambda\mu-2\mu^2=1, keine Loesung. Damit beweist das gedruckte Argument die behauptete allgemeine Loesbarkeit nicht; der Quelltext bleibt originalgetreu bewahrt.

3) Die Gruppe $G$ enth\"alt die Substitution
\[
A_0 = \begin{pmatrix} 0, & 1 \\ -1, & 0 \end{pmatrix},
\]
%CONT>>


% ===== scan leaf 277, printed 256 =====
% !PRINTED 256
\weberpage{256}
%<<CONT
mit alleiniger Ausnahme des Falles $n=1$, $p=2$. Zun\"achst haben wir in $G$ die Substitutionen
\[
\begin{pmatrix} 0, & -1 \\ 1, & 0 \end{pmatrix}\begin{pmatrix} 0, & 1 \\ -1, & \xi \end{pmatrix}\begin{pmatrix} 0, & 1 \\ -1, & 0 \end{pmatrix} = \begin{pmatrix} \xi, & 1 \\ -1, & 0 \end{pmatrix},
\]
also auch
\[
\begin{pmatrix} 0, & 1 \\ -1, & \xi \end{pmatrix}\begin{pmatrix} \xi, & 1 \\ -1, & 0 \end{pmatrix} = \begin{pmatrix} 1, & 0 \\ 2\xi, & 1 \end{pmatrix},
\]
und mithin
\[
\begin{pmatrix} 1, & 0 \\ 2\xi, & 1 \end{pmatrix}\begin{pmatrix} 0, & 1 \\ -1, & \xi \end{pmatrix} = \begin{pmatrix} 0, & 1 \\ -1, & 3\xi \end{pmatrix},
\]
\[
\begin{pmatrix} 1, & 0 \\ 2\xi, & 1 \end{pmatrix}\begin{pmatrix} 0, & 1 \\ -1, & 3\xi \end{pmatrix} = \begin{pmatrix} 0, & 1 \\ -1, & 5\xi \end{pmatrix} \ \text{u. s. f.,}
\]
\[
\begin{pmatrix} 0, & 1 \\ -1, & m\xi \end{pmatrix}
\]
f\"ur jede ungerade ganze Zahl $m$; wenn also $p$ ungerade ist, so k\"onnen wir $m=p$ annehmen und erhalten $A_0$. Ist aber $p=2$, so haben wir, wenn wir mit $\varrho$ eine noch zu bestimmende Zahl in $\mathfrak{C}$ bezeichnen, folgende Substitutionen in $G$:
\[
\begin{pmatrix} 0, & \varrho \\ -\varrho^{-1}, & 0 \end{pmatrix}\begin{pmatrix} 0, & 1 \\ -1, & \xi \end{pmatrix}\begin{pmatrix} 0, & -\varrho \\ \varrho^{-1}, & 0 \end{pmatrix} = \begin{pmatrix} \xi, & \varrho^2 \\ -\varrho^{-2}, & 0 \end{pmatrix}.
\]
Da hier nun jede Zahl ein Quadrat ist, so kann man $\varrho^2$ auch durch $\varrho$ ersetzen, und findet also, dass in $G$ die Substitution
\[
\begin{pmatrix} \xi, & \varrho \\ -\varrho^{-1}, & 0 \end{pmatrix}
\]
f\"ur jedes beliebige von Null verschiedene $\varrho$, und folglich auch die entgegengesetzte Substitution
\[
\begin{pmatrix} 0, & -\varrho \\ \varrho^{-1}, & \xi \end{pmatrix}
\]
vorkommen muss. Bedeuten also $\varrho,\sigma$ zwei von Null verschiedene Zahlen, so haben wir in $G$ auch
\weq{9}{\begin{pmatrix} 0, & \varrho^{-1}\sigma^{-1} \\ -\varrho\sigma, & \xi \end{pmatrix}\begin{pmatrix} \xi, & \sigma^{-1} \\ -\sigma, & 0 \end{pmatrix}\begin{pmatrix} 0, & \varrho \\ -\varrho^{-1}, & \xi \end{pmatrix} = \begin{pmatrix} 0, & 1 \\ -1, & \xi\varrho(1+\sigma+\sigma\varrho) \end{pmatrix}.}


% ===== scan leaf 278, printed 257 =====
% !PRINTED 257
\weberpage{257}

Wenn nun $n$ gr\"osser als 1 ist, so kann man $\varrho$ und $1+\varrho$ von Null verschieden annehmen und dann $\sigma=-1:(1+\varrho)$ setzen, wodurch die Substitution (9) in $A_0$ \"ubergeht.

Ist aber $n=1$ und $p=2$, so tritt der \gesperrt{erste Ausnahmefall} ein; denn dann ist immer eine der beiden Zahlen $\varrho$ und $1+\varrho$ gleich Null. Dieser Fall f\"uhrt auf eine Gruppe $6^{\text{ten}}$ Grades
\[
E = \begin{pmatrix}1,&0\\0,&1\end{pmatrix}, \begin{pmatrix}1,&0\\1,&1\end{pmatrix}, \begin{pmatrix}0,&1\\1,&1\end{pmatrix}, \begin{pmatrix}1,&1\\1,&0\end{pmatrix}, \begin{pmatrix}1,&1\\0,&1\end{pmatrix}, \begin{pmatrix}0,&1\\1,&0\end{pmatrix},
\]
in der ein Normaltheiler $3^{\text{ten}}$ Grades enthalten ist:
\[
G = \begin{pmatrix}1,&0\\0,&1\end{pmatrix}, \begin{pmatrix}0,&1\\1,&1\end{pmatrix}, \begin{pmatrix}1,&1\\1,&0\end{pmatrix}.
\]

Von diesem Ausnahmefalle sehen wir jetzt ab. Wir haben dann weiter:

4) Die Gruppe $G$ enth\"alt jede Substitution von der Form
\weq{10}{\begin{pmatrix} 1, & 0 \\ \xi, & 1 \end{pmatrix},}
worin $\xi$ eine beliebige Zahl in $\mathfrak{E}$ ist, ausgenommen in dem Falle $n=1$, $p=3$.

Denn nach 3) enth\"alt $G$ die Substitution
\[
\begin{pmatrix} 0, & \varrho \\ -\varrho^{-1}, & 0 \end{pmatrix}\begin{pmatrix} 0, & 1 \\ -1, & 0 \end{pmatrix}\begin{pmatrix} 0, & -\varrho \\ \varrho^{-1}, & 0 \end{pmatrix}\begin{pmatrix} 0, & 1 \\ -1, & 0 \end{pmatrix} = \begin{pmatrix} \varrho^2, & 0 \\ 0, & \varrho^{-2} \end{pmatrix},
\]
wenn $\varrho$ eine beliebige von Null verschiedene Zahl in $\mathfrak{E}$ ist. Daraus folgt, dass auch
\weq{11}{\begin{pmatrix} 1, & 0 \\ -\lambda, & 1 \end{pmatrix}\begin{pmatrix} \varrho^{-2}, & 0 \\ 0, & \varrho^2 \end{pmatrix}\begin{pmatrix} 1, & 0 \\ \lambda, & 1 \end{pmatrix}\begin{pmatrix} \varrho^2, & 0 \\ 0, & \varrho^{-2} \end{pmatrix} = \begin{pmatrix} 1, & 0 \\ \lambda(\varrho^4-1), & 1 \end{pmatrix}}
f\"ur jedes beliebige $\lambda$ in $G$ vorkommt. Kann man nun $\varrho$ von Null verschieden so annehmen, dass $\varrho^4-1$ nicht verschwindet, so kann man f\"ur jedes $\xi$ die Zahl $\lambda$ aus $\lambda(\varrho^4-1)=\xi$ bestimmen, und erh\"alt also aus (11) jede Substitution der Form (10). Um die M\"oglichkeit hiervon zu beurtheilen, nehmen wir eine primitive Wurzel $\gamma$ von $\mathfrak{E}$ an und setzen $\varrho=\gamma^h$.

Kann man den Exponenten $h$ so annehmen, dass $4h$ nicht durch $p^n-1$ theilbar ist, so gen\"ugt $\varrho$ unserer Forderung. Dies ist immer m\"oglich, wenn $p=2$, $n>1$ ist, und wenn $p^n-1>4$ ist, und es bleiben also nur die beiden F\"alle $n=1$, $p=3$ und $n=1$, $p=5$ noch zweifelhaft.
% [signature line "Weber, Algebra. II.  17" omitted -- gathering signature, not body text]


% ===== scan leaf 279, printed 258 =====
% !PRINTED 258
\weberpage{258}

In dem Falle $n=1$, $p_{,}=5$ haben wir aber in $G$ nach (2) % [sic] odd subscript comma printed after p
und 3) die Substitution

\[
\begin{pmatrix} 1, & 1 \\ 2, & -2 \end{pmatrix}
\begin{pmatrix} 0, & 1 \\ -1, & 0 \end{pmatrix}
\begin{pmatrix} -2, & -1 \\ -2, & 1 \end{pmatrix}
\begin{pmatrix} 0, & 1 \\ -1, & 0 \end{pmatrix}
=
\begin{pmatrix} -2, & 0 \\ 0, & 2 \end{pmatrix},
\]
und folglich auch
\weq{12}{
\begin{pmatrix} 1, & 0 \\ -\xi, & 1 \end{pmatrix}
\begin{pmatrix} -2, & 0 \\ 0, & 2 \end{pmatrix}
\begin{pmatrix} 1, & 0 \\ \xi, & 1 \end{pmatrix}
\begin{pmatrix} 2, & 0 \\ 0, & -2 \end{pmatrix}
=
\begin{pmatrix} 1, & 0 \\ -2\xi, & 1 \end{pmatrix},
}
also, da $\xi$ und folglich auch $-2\xi$ beliebig ist, wieder die Substitutionen (10) und damit also alle Erzeugenden der Gruppe $E$,
\gesperrt{und also ist in allen diesen Fällen $G$ mit $E$ identisch und
folglich die Gruppe $E$ einfach}.

Der Fall $n=1$, $p=3$ bietet aber die zweite wirkliche
Ausnahme. In diesem Falle ist die Gruppe $E$ vom $12^{\mathrm{ten}}$ Grade
und sie hat den Normaltheiler $4^{\mathrm{ten}}$ Grades

\[
G =
\begin{pmatrix} 1, & 0 \\ 0, & 1 \end{pmatrix},
\begin{pmatrix} 0, & 1 \\ -1, & 0 \end{pmatrix},
\begin{pmatrix} -1, & 1 \\ 1, & 1 \end{pmatrix},
\begin{pmatrix} 1, & 1 \\ 1, & -1 \end{pmatrix}.
\]

Um sich davon zu überzeugen, braucht man nur die beiden
Nebengruppen:
\[
\begin{aligned}
\begin{pmatrix} 1, & 0 \\ 1, & 1 \end{pmatrix} G &=
\begin{pmatrix} 1, & 0 \\ 1, & 1 \end{pmatrix},
\begin{pmatrix} 0, & 1 \\ -1, & 1 \end{pmatrix},
\begin{pmatrix} -1, & 1 \\ 0, & -1 \end{pmatrix},
\begin{pmatrix} 1, & 1 \\ -1, & 0 \end{pmatrix} \\
\begin{pmatrix} 1, & 0 \\ -1, & 1 \end{pmatrix} G &=
\begin{pmatrix} 1, & 0 \\ -1, & 1 \end{pmatrix},
\begin{pmatrix} 0, & 1 \\ -1, & -1 \end{pmatrix},
\begin{pmatrix} -1, & 1 \\ -1, & 0 \end{pmatrix},
\begin{pmatrix} 1, & 1 \\ 0, & 1 \end{pmatrix}
\end{aligned}
\]
mit $G\begin{pmatrix} 1, & 0 \\ 1, & 1 \end{pmatrix}$, $G\begin{pmatrix} 1, & 0 \\ -1, & 1 \end{pmatrix}$
zu vergleichen.

Damit ist also allgemein bewiesen, dass, von den beiden
Fällen $n=1$, $p=2$ und $n=1$, $p=3$ abgesehen, die Gruppe $E$
einfach ist. Für die ersten Fälle erhält man für die Grade $g$
dieser einfachen Gruppen:

\[
\begin{array}{lll@{\qquad}lll}
n=1, & p=\ 5, & g=\ \ 60 & n=2, & p=2, & g=\ \ 60 \\
     & p=\ 7, & g=\ 168  &      & p=3, & g=\ 360 \\
     & p=11,  & g=\ 660  &      & p=5, & g=7800 \\
     & p=13,  & g=1092   & n=3, & p=2, & g=\ 504 \\
     &        &          &      & p=3, & g=9828 \\
     &        &          & n=4, & p=2, & g=4080^{1)}.
\end{array}
\]

\rule{4cm}{0.4pt}

$^{1)}$ Siehe Bulletin of the New York math. Society. October 1893. In
dem Berichte über den mathematischen Congress in Chicago finden sich
zwei Notizen über diese Gruppen von \gesperrt{Cole} und \gesperrt{Moore}.


% ===== scan leaf 280, printed 259 =====
% !PRINTED 259
\weberpage{259}
\pdfbookmark[2]{§ 67. Congruenzkörper zweiten Grades}{bm78}

\begin{center}
\S.\,67.

\gesperrt{Congruenzkörper zweiten Grades.}
\end{center}

In jedem Congruenzkörper $\mathfrak{E}_{n,p}$ für den Modul $p$ vom $n^{\mathrm{ten}}$
Grade ist der Congruenzkörper ersten Grades $\mathfrak{E}_{1,p}$, der aus den
nach dem Modul $p$ genommenen rationalen Zahlen besteht, als
Theiler enthalten. Daher ist auch in der Gruppe der linearen
Substitutionen $E_{n,p}$ eine Gruppe als Theiler enthalten, die aus
allen Substitutionen der Form

\weq{1}{A = \begin{pmatrix} a, & b \\ c, & d \end{pmatrix}, \quad ab-cd=1 % [sic] printed ab-cd, not ad-bc
}
besteht, worin $a, b, c, d$ nach dem Modul $p$ genommene ganze
rationale Zahlen sind.

Diese Gruppe, die als die Gruppe der Modulargleichungen
in der Theorie der elliptischen Functionen auftritt$^{1)}$, ist das
eigentliche Ziel unserer Betrachtungen. Es bietet aber für die
Untersuchung dieser Gruppe, und besonders für die Aufsuchung
ihrer Theiler, den grössten Vortheil, diese Gruppe nicht selbstständig
für sich zu betrachten, sondern als Theiler einer
Gruppe $E_{2,p}$. Bezeichnen wir nämlich mit $N$ irgend einen feststehenden
quadratischen Nichtrest von $p$, so ist, wie schon oben
bemerkt, $t^2-N$ nach dem Modul $p$ irreducibel, und wir erhalten
einen Congruenzkörper $2^{\mathrm{ten}}$ Grades, wenn wir eine
\gesperrt{Galois'sche} imaginäre Zahl durch die Gleichung

\weq{2}{\varepsilon^2 = N}
einführen.

Es soll der Deutlichkeit wegen für die nächsten Betrachtungen
festgesetzt sein, dass die kleinen lateinischen Buchstaben
ganze rationale Zahlen (nach dem Modul $p$ genommen), die wir
hier auch \gesperrt{reelle Zahlen} nennen, die griechischen Buchstaben
Zahlen des Körpers $\mathfrak{E}_{2,p}$ bedeuten sollen.

Dieser Körper $\mathfrak{E}_{2,p}$ hat die für uns sehr wichtige Eigenschaft,
dass \gesperrt{in ihm alle reellen Zahlen Quadrate sind}.

Wenn nämlich $a$ quadratischer Rest von $p$ ist, so können
wir $a=x^2$ setzen, und wenn $b$ quadratischer Nichtrest ist, so
kann $x^2=bN$ befriedigt werden, also $b=(x\varepsilon^{-1})^2$.

\rule{4cm}{0.4pt}

$^{1)}$ Man vergleiche des Verfassers Buch: \glqq Elliptische Functionen und
algebraische Zahlen\grqq{}. Braunschweig 1891.


% ===== scan leaf 281, printed 260 =====
% !PRINTED 260
\weberpage{260}

Eine Zahl von der Form $a\varepsilon$ soll eine \gesperrt{rein imaginäre}
\gesperrt{Zahl} heissen, zwei Zahlen der Form $a+b\varepsilon$, $a-b\varepsilon$ \gesperrt{con-}
\gesperrt{jugirt imaginäre Zahlen}. Die Uebertragung dieser Bezeichnungen
aus der Theorie der gewöhnlichen complexen Zahlen
auf die Zahlen des Körpers $\mathfrak{E}_{2,p}$ rechtfertigt sich durch die
Uebereinstimmung der Rechenregeln und kann zu keinem Missverständniss
führen, da in diesen Betrachtungen von den gewöhnlichen
complexen Zahlen nicht die Rede ist.

Weil immer, wenn $N$ ein quadratischer Nichtrest ist,
$N^{1/2(p-1)}=-1$ ist, so folgt aus (2):

\weq{3}{\varepsilon^p = -\varepsilon.}

Wendet man den binomischen Satz auf die $p^{\mathrm{te}}$ Potenz des
Binoms $\alpha=a+b\varepsilon$ an, und beachtet, dass alle Binomialco\"efficienten,
mit Ausnahme des ersten und des letzten, durch $p$
theilbar, also hier $=0$ sind, dass ferner nach dem \gesperrt{Fermat'schen}
Satze $a^p=a$, $b^p=b$ ist, so ergiebt sich nach (3):

\weq{4}{\alpha = a+b\varepsilon, \quad \alpha^p = a-b\varepsilon,}
und folglich durch Multiplication:

\weq{5}{\alpha^{p+1} = a^2 - Nb^2,}
woraus folgt, dass $\alpha^{p+1}$ immer eine reelle Zahl ist.

\gesperrt{Ist $\gamma$ eine primitive Wurzel des Körpers $\mathfrak{E}_{2,p}$ [\S.\,64, (8)],
so ist $\gamma^r$ dann und nur dann reell, wenn $r$ durch $p+1$
theilbar ist, und $\gamma^{p+1}$ ist eine primitive Wurzel der
Primzahl $p$ im gewöhnlichen Sinne}.

\gesperrt{Die nothwendige und hinreichende Bedingung
dafür, dass eine Zahl $\alpha$ in $\mathfrak{E}_{2,p}$ reell ist, besteht in der
Gleichung $\alpha^p=\alpha$}.

\gesperrt{Jede reelle Zahl ist als $(p+1)^{\mathrm{te}}$ Potenz einer Zahl
in $\mathfrak{E}_{2,p}$ darstellbar}.

Denn dass $\gamma^r$ reell ist, wenn $r$ durch $p+1$ theilbar ist,
folgt aus (5). Dass es aber nicht anders reell sein kann, ergiebt
sich daraus, dass nach dem Fermat'schen Satze jede reelle von
Null verschiedene Zahl der Bedingung $a^{p-1}=1$ genügt, dass
also, wenn $\gamma^r$ reell ist, $\gamma^{r(p-1)}=1$ sein muss, und es muss
daher $r(p-1)$ durch $p^2-1$ und folglich $r$ durch $p+1$ theilbar
sein. Da ferner $\gamma^{r(p+1)}$ nur dann $=1$ ist, wenn $r$ durch $p-1$
theilbar ist, so ist $\gamma^{p+1}$ primitive Wurzel von $p$.

\gesperrt{Auch der Begriff der Einheitswurzeln lässt sich auf
die Zahlen des Körpers $\mathfrak{E}_{2,p}$ übertragen}.


% ===== scan leaf 282, printed 261 =====
% !PRINTED 261
\weberpage{261}
\pdfbookmark[2]{§ 68. Die reelle lineare Congruenzgruppe Lp}{bm79}

Ist $m$ ein Theiler von $p^2-1$, und

\[
\varrho = \gamma^{\frac{p^2-1}{m}},
\]
so ist $\varrho^m=1$ und es giebt keine niedrigere als die $m^{\mathrm{te}}$
Potenz von $\varrho$, die gleich 1 wird. Daher nennen wir $\varrho$ eine
primitive $m^{\mathrm{te}}$ Einheitswurzel (in $\mathfrak{E}_{2,p}$). Alle anderen $m^{\mathrm{ten}}$
Einheitswurzeln sind dann Potenzen $\varrho^s$ von $\varrho$ und unter diesen
sind die und nur die primitiv, bei denen $s$ relativ prim zu $m$ ist.

Dies folgt daraus, dass jede von Null verschiedene Zahl $\xi$
in der Form $\xi=\gamma^r$ darstellbar ist, und dass $\xi^m=\gamma^{rm}$ dann
und nur dann $=1$ ist, wenn $rm$ durch $p^2-1$ theilbar ist.

Jede von Null verschiedene Zahl in $\mathfrak{E}_{2,p}$ ist eine Einheitswurzel
 vom Grade $p^2-1$.

\begin{center}
\S.\,68.

\gesperrt{Die reelle lineare Congruenzgruppe $L_p$.}
\end{center}

Wir gehen nun, mit diesen Hülfsmitteln ausgestattet, an die
Untersuchung der reellen Congruenzgruppe $L_p$, die aus allen
Substitutionen der Form

\[
A = \begin{pmatrix} a, & b \\ c, & d \end{pmatrix}
\]
besteht, worin $a, b, c, d$ reelle, nach dem Modul $p$ genommene
ganze Zahlen sind, die der Bedingung $ad-bc\equiv 1$ genügen.
Die Gruppe $L_p$ ist nach der früheren Bezeichnung mit $E_{1,p}$ zu
bezeichnen, und ihr Grad ist, wenn wir den Fall $p=2$ ausschliessen:

\[
\frac{p(p^2-1)}{2}.
\]

Sind $A, U$ irgend zwei Elemente aus einer Gruppe, so haben
wir schon früher

\[
U^{-1}AU = A'
\]
das durch $U$ aus $A$ transformirte Element genannt. $A'^2, A'^3, \ldots$
sind transformirt aus $A^2, A^3, \ldots$, woraus folgt, dass alle aus
einander durch Transformation gewonnenen Elemente denselben
Grad haben. Entnehmen wir die Elemente $A$ aus irgend einer
Gruppe $G$, so bilden die transformirten Elemente eine mit $G$
isomorphe Gruppe

\[
U^{-1}GU = G',
\]
die wir die transformirte Gruppe von $G$ nennen.


% ===== scan leaf 283, printed 262 =====
% !PRINTED 262
\weberpage{262}
Es möge nun $A$ irgend ein Element der Gruppe $L_p$ sein,
und $U = \begin{pmatrix} \lambda, & \mu \\ \nu, & \varrho \end{pmatrix}$ ein Element aus $E_{2,p}$. Wir wollen aus $A$ durch
Transformation mit $U$ eine gewisse Normalform ableiten.

Es soll zunächst versucht werden, $A$ in die Normalform
\weq{1}{S = \begin{pmatrix} \sigma, & 0 \\ 0, & \sigma^{-1} \end{pmatrix}}
zu transformiren. Aus der Gleichung
\[
\begin{pmatrix} a, & b \\ c, & d \end{pmatrix} \begin{pmatrix} \lambda, & \mu \\ \nu, & \varrho \end{pmatrix} = \begin{pmatrix} \lambda, & \mu \\ \nu, & \varrho \end{pmatrix} \begin{pmatrix} \sigma, & 0 \\ 0, & \sigma^{-1} \end{pmatrix}
\]
erhält man
\weq{2}{\begin{aligned}
&(a-\sigma)\,\lambda + b\,\nu = 0, \quad (a-\sigma^{-1})\,\mu + b\,\varrho = 0,\\
&c\,\lambda + (d-\sigma)\,\nu = 0, \quad\ \ c\,\mu + (d-\sigma^{-1})\,\varrho = 0,
\end{aligned}}
und daraus ergiebt sich, dass $\sigma$ und $\sigma^{-1}$ die Wurzeln der quadratischen
Gleichung
\[
(a-\sigma)\,(d-\sigma) - b\,c = 0
\]
oder
\weq{3}{\sigma^2 - \sigma\,(a+d) + 1 = 0}
sein müssen. Hat man $\sigma$ und $\sigma^{-1}$ hieraus bestimmt, so findet
man aus (2) die Verhältnisse $\lambda : \nu$ und $\mu : \varrho$. Die Grössen $\lambda, \mu, \nu, \varrho$
selbst sind dann noch so zu bestimmen, dass $\lambda\varrho - \mu\nu = 1$ wird.
Dies ist immer möglich, wenn die aus (2) bestimmten Werthe
nicht die Gleichung $\lambda\varrho = \mu\nu$ erfüllen. Sind $b$ und $c$ beide $= 0$,
so hat $A$ schon die Form (1). Ist aber eine der beiden Zahlen
$b, c$, etwa $b$, von Null verschieden, so ergiebt sich aus (2):
\weq{4}{\frac{\nu}{\lambda} = -\frac{(a-\sigma)}{b}, \quad \frac{\varrho}{\mu} = -\frac{(a-\sigma^{-1})}{b},}
und es kann nur dann $\lambda\varrho = \mu\nu$ oder $\nu : \lambda = \varrho : \mu$ sein, wenn
$\sigma = \sigma^{-1}$, also $\sigma = \pm 1$ ist, d.\,h. nach (3), wenn $a + d = \pm 2$ ist.

Die Gleichung (3) ist aber im Körper $\mathfrak{E}_{2,p}$ immer lösbar, weil
jede \gesperrt{reelle} Zahl in diesem Körper ein Quadrat ist, und ergiebt
\weq{5}{\begin{aligned}
\sigma &= \frac{a+d}{2} + \sqrt{\left(\frac{a+d}{2}\right)^2 - 1}\\
\sigma^{-1} &= \frac{a+d}{2} - \sqrt{\left(\frac{a+d}{2}\right)^2 - 1}
\end{aligned}}


% ===== scan leaf 284, printed 263 =====
% !PRINTED 263
\weberpage{263}
Wir kommen also zu dem ersten Resultate:

1.~\gesperrt{Eine Substitution $A$ ist immer auf die Normalform $S$ transformirbar, wenn $a+d$ nicht $=\pm 2$
ist. $\sigma$ ist reell oder imaginär, je nachdem $(a+d)^2 - 4$ quadratischer Rest oder Nichtrest
von $p$ ist.}

In dem noch übrigen Falle, wo $a+d = \pm 2$ ist, lässt sich
die Normalform $S$ nicht herstellen; dagegen können wir $A$ in
diesem Falle auf eine andere Normalform, nämlich
\weq{6}{T = \begin{pmatrix} 1, & t \\ 0, & 1 \end{pmatrix}}
bringen, in der $t$ reell ist.

Um dies zu beweisen, können wir zunächst
\weq{7}{a + d = +2}
annehmen, weil der andere Fall $a+d = -2$ durch Aenderung
aller Vorzeichen von $a, b, c, d$ auf diesen zurückgeführt wird.

Aus (7) aber folgt noch
\weq{8}{a - 1 = -(d-1), \quad (a-1)(d-1) = b\,c.}
Wenn also zunächst $b$ oder $c = 0$ ist, so ist $a = d = 1$
und wir haben entweder
\[
A = \begin{pmatrix} 1, & b \\ 0, & 1 \end{pmatrix},
\]
was schon die Normalform $T$ hat, oder
\[
A = \begin{pmatrix} 1, & 0 \\ c, & 1 \end{pmatrix} = \begin{pmatrix} 0, & 1 \\ -1, & 0 \end{pmatrix} \begin{pmatrix} 1, & -c \\ 0, & 1 \end{pmatrix} \begin{pmatrix} 0, & -1 \\ 1, & 0 \end{pmatrix},
\]
so dass $\begin{pmatrix} 0, & -1 \\ 1, & 0 \end{pmatrix} A \begin{pmatrix} 0, & 1 \\ -1, & 0 \end{pmatrix}$ die Normalform $T$ hat. Sind aber $b$ und $c$ von Null verschieden,
so erhalten wir aus (8) die Transformation
\[
\begin{pmatrix} (d-1)\,b^{-1}, & 0, \\ -1, & (a-1)\,c^{-1} \end{pmatrix} \begin{pmatrix} a, & b \\ c, & d \end{pmatrix} \begin{pmatrix} (a-1)\,c^{-1}, & 0, \\ 1, & (d-1)\,b^{-1} \end{pmatrix} = \begin{pmatrix} 1, & -c \\ 0, & 1 \end{pmatrix},
\]
was die Form $T$ hat. Damit ist also bewiesen:

2.~\gesperrt{Eine Substitution $A$, in der $a+d = \pm 2$ ist, kann immer durch
Transformation auf die Normalform $T$ gebracht werden.}


% ===== scan leaf 285, printed 264 =====
% !PRINTED 264
\weberpage{264}
Hiernach lassen sich leicht die Grade der Elemente unserer
Gruppe bestimmen. Wir haben dabei drei Fälle zu unterscheiden.

I.~\gesperrt{$a + d = \pm 2$.} \gesperrt{Jede Substitution dieser Art lässt sich
in die Normalform $T$ transformiren. Es ist aber für jeden
Exponenten $r$}
\[
T^r = \begin{pmatrix} 1, & r\,t \\ 0, & 1 \end{pmatrix},
\]
\gesperrt{und folglich ist der Grad dieser Substitution $p$, ausser
wenn $t = 0$, also $T$ die identische Substitution ist.}

Die Anzahl der Elemente dieser Art ist leicht zu bestimmen.
Nehmen wir zunächst $c = 0$, so können wir $b$ beliebig wählen
und für $a, d$ erhalten wir aus $a + d = \pm 2$, $a\,d = 1$ die beiden
Bestimmungen $a = d = \pm 1$. Dies giebt $2p$ Formen, die identische
 Substitution eingeschlossen. Nehmen wir sodann $a$ und $c$
beliebig, jedoch $c$ von Null verschieden, was $p\,(p-1)$ Möglichkeiten
 giebt, so folgt aus $a\,d - b\,c = 1$
\[
b = -c^{-1} + a\,d\,c^{-1},
\]
und zu jedem $a$ kann $d$ auf zwei Arten bestimmt werden. Die
Gesammtzahl $2\,p^2$, die wir so erhalten, ist noch zu halbiren, weil
hier die Substitutionen als zwei verschiedene auftreten, in denen
alle Zeichen entgegengesetzt sind; und wir erhalten also $p^2$ Substitutionen
 dieser Art (die identische eingeschlossen).

II.~\gesperrt{$(a+d)^2 - 4$ quadratischer Rest von $p$. Eine solche
Substitution lässt sich in die Normalform $S$ transformiren
mit reellem $\sigma$.}

Verstehen wir unter $\gamma$ eine primitive Wurzel des Körpers
$\mathfrak{E}_{2,p}$, so ist
\[
\gamma^{\frac{p^2-1}{2}} = -1,
\]
und wir können $r$ so bestimmen, dass
\[
\sigma = \gamma^{r(p+1)}, \quad \sigma^{\frac{p-1}{2}} = (-1)^r
\]
wird (\S.\,67). Nach (5) erhalten wir dann
\weq{9}{a + d = \gamma^{r(p+1)} + \gamma^{-r(p+1)},}
und weil nach (5) die beiden Werthe $\sigma, \sigma^{-1}$, von der Reihenfolge
abgesehen, durch $a+d$ bestimmt sind, so werden zwei Werthe
des Ausdrucks (9) nur dann einander gleich oder entgegengesetzt,
wenn die entsprechenden Werthe, mit positivem oder
%CONT>>


% ===== scan leaf 286, printed 265 =====
% !PRINTED 265
\weberpage{265}
%<<CONT
negativem Zeichen genommen, von $r$ nach dem Modul $\frac{1}{2}(p-1)$
congruent sind. Man erhält daher alle von einander und von
$\pm 2$ verschiedenen Werthe dieses Ausdruckes, wenn man
\weq{10}{r = 1, 2, \cdots \frac{p-3}{2}}
setzt. Da hier für jeden Exponenten $k$
\[
S^k = \begin{pmatrix} \sigma^k, & 0 \\ 0, & \sigma^{-k} \end{pmatrix}
\]
ist, so ist der Grad von $S$, und also auch von jeder Substitution
der II$^{\mathrm{ten}}$ Art entweder $\frac{1}{2}(p-1)$ oder ein Theiler davon, je
nachdem $r$ relativ prim zu $\frac{1}{2}(p-1)$ ist oder nicht.

Um die Anzahl dieser Substitutionen zu bestimmen, verfahren
 wir wie oben. Ist zunächst $c = 0$, so ergiebt sich
\[
a = \gamma^{\pm r(p+1)}, \quad d = \gamma^{\mp r(p+1)},
\]
und $b$ kann $p$ Werthe haben. Die Zahl $r$ hat die Werthe (10),
und also ergeben sich hiernach $p\,(p-3)$ Substitutionen. Ist
dann $c$ von Null verschieden, so ist $p\,(p-1)$ die Anzahl der
verschiedenen Annahmen über $a, c$. Ist $a$ angenommen, so können
wir $d$ aus (9) bestimmen und
\[
b = -c^{-1} + a\,d\,c^{-1},
\]
setzen, woraus wir $\frac{1}{2}\,p\,(p-1)(p-3)$ Bestimmungen erhalten,
also mit den für $c = 0$ gezählten zusammen $\frac{1}{2}\,p\,(p-3)(p+1)$.
Auch diese Zahl ist noch zu halbiren, so dass wir
\[
\frac{1}{4}\,p\,(p-3)(p+1)
\]
Substitutionen der II$^{\mathrm{ten}}$ Art erhalten.

III.~$(a+d)^2 - 4$ quadratischer Nichtrest von $p$. In diesem
Falle ist $\sigma$ nicht reell, aber mit seinem reciproken Werth
conjugirt. Also ist nach \S.\,67, (4) $\sigma^p = \sigma^{-1}$, und folglich
\[
\sigma^{\frac{p+1}{2}} = \pm 1.
\]
Daraus ergiebt sich, dass der Grad einer solchen Substitution
 $\frac{1}{2}(p+1)$ oder ein Theiler dieser Zahl ist.
In diesem dritten Falle setzen wir
\weq{11}{\begin{aligned}
&\sigma = \gamma^{r(p-1)}, \quad a+d = \gamma^{r(p-1)} + \gamma^{-r(p-1)}\\
&r = 1, 2, \ldots, \tfrac{1}{2}(p-1).
\end{aligned}}
Hier kann der Fall $c = 0$ oder $b = 0$ nicht vorkommen,
weil sich aus $a\,d = 1$ ergeben würde, dass $a = \sigma^{\pm 1}$, $d = \sigma^{\mp 1}$
%CONT>>


% ===== scan leaf 287, printed 266 =====
% !PRINTED 266
\weberpage{266}
\pdfbookmark[2]{§ 69. Imaginäre Form der Gruppe Lp}{bm80}
%<<CONT
sein müsste, und es würden $a$ und $d$ nicht reell ausfallen. Also
nehmen wir für $a$ und $c$, wie oben, die $p\,(p-1)$ verschiedenen
Werthsysteme, und erhalten aus (11) für jedes von ihnen
$1/2\,(p-1)$ Bestimmungen von $b$ und $d$; auch diese Zahl ist zu
halbiren und es ergeben sich
\[
1/4\;p\,(p-1)\,(p-1)
\]
Substitutionen der III\textsuperscript{ten} Art. Die Gesammtzahl aller Substitutionen
 der Gruppe $L_p$ ist hiernach, wie es sein muss,
\[
p^2 + 1/4\,p\,(p-3)\,(p+1) + 1/4\,p\,(p-1)^2 = 1/2\,p\,(p^2-1).
\]

\begin{center}
§. 69.

\gesperrt{Imaginäre Form der Gruppe $L_p$.}
\end{center}

Die Substitutionen der beiden ersten Arten der Gruppe $L_p$
haben die Eigenschaft, dass die ihnen entsprechende Normalform
 $T$ oder $S$ gleichfalls in $L_p$ enthalten ist, und dass man sie
in diese Normalformen transformiren kann durch \gesperrt{reelle} Substitutionen
, d. h. durch Substitutionen, die selbst in $L_p$ vorkommen
. Beides trifft bei den Substitutionen der dritten Art
nicht zu.

Man kann aber, wie wir jetzt beweisen werden, die ganze
Gruppe $L_p$ in eine andere Gruppe $\Lambda_p$ so transformiren, dass $\Lambda_p$ ein
mit $L_p$ \gesperrt{isomorpher Theiler der Gruppe} $E_{2,p}$ ist, dass die
Normalformen der Transformationen dritter Art in $\Lambda_p$ enthalten
sind und dass jede Substitution dritter Art in $\Lambda_p$ in die
Normalform transformirt werden kann durch Substitutionen von
$\Lambda_p$ selbst.

Um eine solche Transformation von $L_p$ zu finden, betrachten
wir die Substitution
\weq{1}{S = \begin{pmatrix} \sigma, & 0 \\ 0, & \sigma^{-1} \end{pmatrix}}
unter der Voraussetzung, dass $\sigma$ und $\sigma^{-1}$ conjugirt imaginär
sind, und suchen die Substitution
\weq{2}{R = \begin{pmatrix} \lambda, & \mu \\ \nu, & \varrho \end{pmatrix}}
so zu bestimmen, dass
\weq{3}{R\,S\,R^{-1}}
%CONT>>


% ===== scan leaf 288, printed 267 =====
% !PRINTED 267
\weberpage{267}
%<<CONT
reell wird. Es ergiebt sich aber nach (1) und (2)
\[
R\,S\,R^{-1} = \begin{pmatrix} \lambda\,\varrho\,\sigma - \mu\,\nu\,\sigma^{-1}, & -\lambda\,\mu\,(\sigma-\sigma^{-1}) \\ \varrho\,\nu\,(\sigma-\sigma^{-1}), & -\mu\,\nu\,\sigma+\varrho\,\lambda\,\sigma^{-1} \end{pmatrix},
\]
und dies wird reell, wenn $\lambda\mu$ und $\varrho\nu$ rein imaginär, $\lambda\varrho$ und
$-\mu\nu$ conjugirt imaginär sind. Diesen Bedingungen kann man
auf mehrfache Art genügen: am einfachsten wohl, und zwar
zugleich so, dass die Determinante von $R = 1$ wird, wenn man
\weq{4}{R = \begin{pmatrix} \frac{1}{2}, & -\frac{1}{2}\,\varepsilon^{-1} \\ \varepsilon, & 1 \end{pmatrix}}
setzt. Ist andererseits
\[
A = \begin{pmatrix} a, & b \\ c, & d \end{pmatrix}
\]
eine beliebige reelle Substitution aus $L_p$, so ergiebt sich nach (4)
wegen $\varepsilon^2 = N$:
\[
R^{-1}\,A\,R =
\]
\weq{5}{\begin{pmatrix} \dfrac{a+d}{2}+b\,\varepsilon+\dfrac{c}{4}\,\varepsilon^{-1}, & -\dfrac{a-d}{2}\,\varepsilon^{-1}+b-\dfrac{c}{4}\,N^{-1} \\[1.2ex] -N\left(\dfrac{a-d}{2}\,\varepsilon^{-1}+b-\dfrac{c}{4}\,N^{-1}\right), & \dfrac{a+d}{2}-b\,\varepsilon-\dfrac{c}{4}\,\varepsilon^{-1} \end{pmatrix},}
wofür wir auch
\weq{6}{A = \begin{pmatrix} \alpha, & \beta \\ -N\beta', & \alpha' \end{pmatrix} = \begin{pmatrix} \alpha, & \beta \\ -N\beta^p, & \alpha^p \end{pmatrix}, \quad \alpha\alpha'+N\beta\beta' = 1}
setzen, worin $\alpha,\alpha'$ und $\beta,\beta'$ zwei conjugirt imaginäre Paare
sind, so dass nach §.\,65, (4) $\alpha' = \alpha^p$, $\beta' = \beta^p$ gesetzt werden
kann.

Wenn wir umgekehrt irgend eine Substitution $A$ von der
Form (6) nehmen, so ergiebt sich
\[
A = R\,A\,R^{-1} =
\]
\weq{7}{\begin{pmatrix} \dfrac{\alpha+\alpha'}{2}+\varepsilon\,\dfrac{\beta'-\beta}{2}, & \dfrac{\beta+\beta'}{4}+\varepsilon^{-1}\,\dfrac{\alpha-\alpha'}{4} \\[1.2ex] -N(\beta+\beta')+\varepsilon(\alpha-\alpha'), & \dfrac{\alpha+\alpha'}{2}-\varepsilon\,\dfrac{\beta'-\beta}{2} \end{pmatrix}}
als eine \gesperrt{reelle} Substitution aus $L_p$.

\gesperrt{Daraus geht hervor, dass, wenn $A$ die ganze Gruppe $L_p$
durchläuft, die durch (5) und (6) bestimmte Substitution $A$ eine
mit $L_p$ isomorphe Gruppe durchläuft, die wir mit $\Lambda_p$ bezeichnen}.


% ===== scan leaf 289, printed 268 =====
% !PRINTED 268
\weberpage{268}
In der Gruppe $\Lambda_p$ ist die Normalform $S$ für die Substitutionen
 dritter Art enthalten.

Nehmen wir nun eine Substitution $A$ von der dritten Art in
der Gruppe $\Lambda_p$ an, so können wir sie, wie jetzt noch nachgewiesen
 werden soll, durch Transformation mit einer Substitution
 $U$, die selbst der Gruppe $\Lambda_p$ angehört, in die Normalform
 transformiren. Denn setzen wir
\weq{8}{A = \begin{pmatrix} \alpha, & \beta \\ -N\beta', & \alpha' \end{pmatrix}, \quad U = \begin{pmatrix} \lambda, & \mu \\ -N\mu', & \lambda' \end{pmatrix},}
\weq{9}{\alpha\alpha'+N\beta\beta' = 1, \qquad \lambda\lambda'+N\mu\mu' = 1,}
worin $\lambda,\lambda'$ und $\mu,\mu'$ conjugirte Paare sind, so erhalten die
Gleichungen (3), (4), §.\,68 die Form:
\weq{10}{(\sigma-\alpha)(\sigma-\alpha')+N\beta\beta' = \sigma^2-\sigma(\alpha+\alpha')+1 = 0,}
\weq{11}{N\,\dfrac{\mu'}{\lambda} = \dfrac{\alpha-\sigma}{\beta}, \qquad \dfrac{\lambda'}{\mu} = -\dfrac{\alpha-\sigma^{-1}}{\beta},}
und die Gleichung (10) muss, da $A$ von der dritten Art sein
soll, zwei zu einander reciproke, conjugirt imaginäre Wurzeln
haben.

Aus der zweiten Gleichung (11) folgt aber durch Uebergang
zu den conjugirt imaginären Zahlen
\[
\dfrac{\lambda}{\mu'} = -\dfrac{\alpha'-\sigma}{\beta'},
\]
und dies ist nach (10) eine Folge der ersten Gleichung (11).
Setzen wir also
\[
N\mu' = \varkappa\,(\alpha-\sigma), \qquad N\mu = \varkappa'\,(\alpha'-\sigma^{-1})
\]
\[
\lambda = \varkappa\beta, \qquad \lambda' = \varkappa'\beta',
\]
worin $\varkappa,\varkappa'$ zwei conjugirt imaginäre Zahlen sind, so sind die
Gleichungen (11) befriedigt, und man kann dann $\varkappa\varkappa'$ noch so
bestimmen, dass $\lambda\lambda'+N\mu\mu' = 1$ wird.

Daraus ergiebt sich noch ein wichtiges Resultat. Bezeichnen
wir mit $\gamma$ eine primitive Wurzel des Körpers $\mathfrak{E}_{2,p}$, so können
wir für die Substitutionen der Normalform für die zweite und
dritte Art die Ausdrücke annehmen (§.\,68):
\[
S_2^r = \begin{pmatrix} \gamma^{(p+1)r}, & 0 \\ 0, & \gamma^{-(p+1)r} \end{pmatrix}, \qquad S_3^r = \begin{pmatrix} \gamma^{(p-1)r}, & 0 \\ 0, & \gamma^{-(p-1)r} \end{pmatrix}.
\]

Ist nun $A$ eine Substitution zweiter Art aus $L_p$ und $A$ eine
%CONT>>


% ===== scan leaf 290, printed 269 =====
% !PRINTED 269
\weberpage{269}
%<<CONT
Substitution dritter Art aus $\Lambda_p$, so können wir die Substitutionen
 $U$, $V$ aus $L_p$ oder aus $\Lambda_p$ so bestimmen, dass
\[
A = U\,S_2^r\,U^{-1}, \qquad A = V\,S_3^r\,V^{-1}.
\]
Setzen wir also, dem Werth $r=1$ entsprechend,
\[
A_1 = U\,S_2\,U^{-1}, \qquad A_1 = V\,S_3\,V^{-1},
\]
so ergiebt sich
\[
A = A_1^r, \qquad A = A_1^r,
\]
und $A_1$ kommt in $L_p$, $A_1$ in $\Lambda_p$ vor.

Ist ferner $B$ eine Substitution erster Art aus $L_p$, so kann
man wieder, wenn man
\[
T_1 = \begin{pmatrix} 1, & 1 \\ 0, & 1 \end{pmatrix}
\]
setzt, $U$ aus $L_p$ so bestimmen, dass
\[
B = U\,T_1^r\,U^{-1} = B_1^r
\]
wird. Hieraus ergiebt sich der Satz:

\gesperrt{Man kann in der Gruppe $L_p$ (oder $\Lambda_p$) die Substitutionen
 der ersten, zweiten und dritten Art mit Hinzuziehung
der Identität in Cyklen von $p$, $\dfrac{p-1}{2}$, $\dfrac{p+1}{2}$
Gliedern anordnen.}

Zwei solche Cyklen können ausser dem Einheitselemente
kein Glied gemein haben.

Dies ist zunächst evident bei den Substitutionen der ersten
Art, deren Grad gleich $p$ ist; denn bei diesen lässt sich der
ganze Cyklus durch Potenzirung aus einem beliebigen von der
Einheit verschiedenen Elemente ableiten.

Wenn aber in zwei Cyklen der zweiten oder dritten Art ein
gemeinsames Element vorkommt, so kann die Gruppe so transformirt
 werden, dass die beiden Cyklen die Gestalt bekommen:
\weq{12}{\begin{aligned} &1, \quad U\,S\,U^{-1}, \quad U\,S^2\,U^{-1}, \ \ldots \\ &1, \quad S, \quad S^2, \ \ldots, \end{aligned}}
worin $S = \begin{pmatrix} \sigma, & 0 \\ 0, & \sigma^{-1} \end{pmatrix}$ die Normalform hat.

Ist nun für irgend zwei von Null verschiedene Exponenten
$r,t$
\[
S^r = U\,S^t\,U^{-1},
\]
%CONT>>


% ===== scan leaf 291, printed 270 =====
% !PRINTED 270
\weberpage{270}
%<<CONT
so ergiebt sich, wenn $U = \begin{pmatrix}\lambda, & \mu\\ \nu, & \varrho\end{pmatrix}$ gesetzt wird, als Bedingung
\[
\begin{aligned}
\sigma^r\lambda &= \pm\sigma^t\lambda, & \sigma^r\mu &= \pm\sigma^{-t}\mu\\
\sigma^{-r}\nu &= \pm\sigma^t\nu, & \sigma^{-r}\varrho &= \pm\sigma^{-t}\varrho,
\end{aligned}
\]
wo überall dasselbe Zeichen gelten muss. Daraus folgt, dass entweder
\[
\sigma^r = \pm\sigma^t, \quad \nu = 0, \quad \mu = 0, \quad \lambda\varrho = 1
\]
oder
\[
\sigma^r = \pm\sigma^{-t}, \quad \lambda = 0, \quad \varrho = 0, \quad \mu\nu = -1
\]
sein muss, und dann ergiebt sich:
\[
USU^{-1} = S \quad \text{oder} \quad = S^{-1},
\]
und in beiden Fällen stimmen also die beiden Cyklen (12) vollständig mit einander überein.

Hiernach ergiebt sich aus den Zahlen am Schluss des \S.\,68 für die Anzahl der Cyklen erster, zweiter und dritter Art:
\[
p+1, \quad \frac{p(p+1)}{2}, \quad \frac{p(p-1)}{2}.
\]

Es ist zweckmässig, neben den Gruppen $L_p$ und $\Lambda_p$ noch eine dritte Gruppe in $E_{2,p}$ zu betrachten, die mit diesen isomorph ist, wenn sie sich auch nicht durch Transformation darauf zurückführen lässt.

Da $N$ reell ist, so können wir nach \S.\,67 eine Zahl $\nu$ in $\mathfrak{E}_{2,p}$ so bestimmen, dass $N = \nu^{p+1}$ ist. Wir lassen nun der Substitution
\[
A = \begin{pmatrix}\alpha, & \beta\\ -N\beta^p, & \alpha^p\end{pmatrix}
\]
aus $\Lambda_p$ eine Substitution $B$ entsprechen, die so gebildet ist:
\weq{13}{B = \begin{pmatrix}\alpha, & \nu\beta\\ -\nu^p\beta^p, & \alpha^p\end{pmatrix}, \quad \alpha^{p+1} + N\beta^{p+1} = 1.}
Setzen wir zwei Substitutionen $B$, $B_1$ von der Form (13) zusammen, so erhalten wir
\[
\begin{aligned}
BB_1 &= \begin{pmatrix}\alpha, & \nu\beta\\ -\nu^p\beta^p, & \alpha^p\end{pmatrix}\begin{pmatrix}\alpha_1, & \nu\beta_1\\ -\nu^p\beta_1^p, & \alpha_1^p\end{pmatrix}\\
&= \begin{pmatrix}\alpha\alpha_1 - N\beta\beta_1^p, & \nu(\alpha\beta_1+\beta\alpha_1^p)\\ -\nu^p(\beta^p\alpha_1+\alpha^p\beta_1^p), & \alpha^p\alpha_1^p - N\beta^p\beta_1\end{pmatrix};
\end{aligned}
\]
darin sind

$\alpha\alpha_1 - N\beta\beta_1^p, \ \alpha^p\alpha_1^p - N\beta^p\beta_1$ und $\alpha\beta_1+\beta\alpha_1^p, \ \beta^p\alpha_1+\alpha^p\beta_1^p$
%CONT>>


% ===== scan leaf 292, printed 271 =====
% !PRINTED 271
\weberpage{271}
\pdfbookmark[2]{§ 70. Divisoren der Gruppe Lp, deren Grad durch p theilbar ist}{bm81}
%<<CONT
zwei conjugirte Paare, und daher hat $BB_1$ auch die Form (13) und steht in derselben Beziehung zu $AA_1$, wie $B$ zu $A$ und $B_1$ zu $A_1$.

\gesperrt{Daraus geht hervor, dass die Gesammtheit der Substitutionen $B$ eine mit $\Lambda_p$ isomorphe Gruppe bildet, und diese Gruppe wollen wir mit $\Gamma_p$ bezeichnen.}

Setzen wir $\beta$ an Stelle von $\nu\beta$ und bezeichnen mit $\alpha', \beta'$ die zu $\alpha, \beta$ conjugirten Grössen, so können wir die Substitutionen der Gruppe $\Gamma_p$ auch in der einfacheren Form annehmen:
\weq{14}{B = \begin{pmatrix}\alpha, & \beta\\ -\beta', & \alpha'\end{pmatrix}, \quad \alpha\alpha' + \beta\beta' = 1.}
Um von einer dieser Substitutionen $B$ zu der entsprechenden reellen Substitution $A$ überzugehen, leitet man zunächst aus (14) die entsprechende Substitution $A$ her:
\weq{15}{A = \begin{pmatrix}\alpha, & \nu^{-1}\beta\\ -\nu\beta', & \alpha'\end{pmatrix},}
und bildet daraus nach (3) und (7):
\weq{16}{A = RAR^{-1}.}
Trotz des Isomorphismus lässt sich die Gruppe $\Gamma_p$ nicht in $L_p$ oder $\Lambda_p$ transformiren, wenigstens nicht durch Substitutionen, deren Zahlen dem Körper $\mathfrak{E}_{2,p}$ angehören. Man müsste, um die Transformation auszuführen, in einen höheren Körper gehen, in dem \gesperrt{alle} Zahlen von $\mathfrak{E}_{2,p}$ Quadrate sind.

\begin{center}
\S.\,70.\\
\gesperrt{Divisoren der Gruppe $L_p$, deren Grad durch $p$}\\
\gesperrt{theilbar ist.}
\end{center}

Es ist ein Problem von grösstem Interesse, alle Divisoren der Gruppe $L_p$ zu bestimmen. Von der Lösung dieses Problems hängt es ab, in welcher Weise man die Gruppe $L_p$ durch Permutationen von Ziffern darstellen kann, bei welchen algebraischen Gleichungen also diese Gruppen auftreten können (\S.\,28, 2.).

Die cyklischen Gruppen, die in $L_p$ enthalten sind, haben wir in den vorangehenden Paragraphen schon betrachtet. Wir haben gesehen, dass der Grad einer cyklischen Gruppe entweder gleich $p$ oder ein Theiler von $1/2\,(p-1)$ oder von $1/2\,(p+1)$
%CONT>>


% ===== scan leaf 293, printed 272 =====
% !PRINTED 272
\weberpage{272}
%<<CONT
ist, und jeder Theiler dieser Zahlen tritt auch unter den Graden der cyklischen Gruppen auf. Der Index eines cyklischen Theilers kann niemals kleiner sein als $1/2\,(p^2-1)$, $p(p+1)$, $p(p-1)$.

Wenn der Grad eines Theilers $G$ von $L_p$ durch $p$ theilbar ist, so enthält er eine cyklische Gruppe vom Grade $p$, und wir können nach \S.\,68 die Gruppe $G$ so transformiren, dass diese cyklische Gruppe aus den Substitutionen
\[
T = \begin{pmatrix}1, & t\\ 0, & 1\end{pmatrix}, \quad t = 0, 1, 2, \ldots, p-1
\]
besteht. Wenn nun $G$ noch eine Substitution
\[
A = \begin{pmatrix}a, & b\\ c, & d\end{pmatrix}
\]
enthält, in der $c$ von Null verschieden ist, so kommt darin auch
\weq{1}{\begin{pmatrix}1, & -ac^{-1}\\ 0, & 1\end{pmatrix}\begin{pmatrix}a, & b\\ c, & d\end{pmatrix}\begin{pmatrix}1, & -dc^{-1}\\ 0, & 1\end{pmatrix} = \begin{pmatrix}0, & -c^{-1}\\ c, & 0\end{pmatrix}}
vor. Daraus folgt weiter, dass auch
\[
\begin{pmatrix}0, & -c^{-1}\\ c, & 0\end{pmatrix}\begin{pmatrix}1, & t\\ 0, & 1\end{pmatrix}\begin{pmatrix}0, & -c^{-1}\\ c, & 0\end{pmatrix} = \begin{pmatrix}1, & 0\\ -c^2 t, & 1\end{pmatrix},
\]
und mithin jede Substitution
\[
U = \begin{pmatrix}1, & 0\\ u, & 1\end{pmatrix}, \quad u = 0, 1, 2, \ldots, p-1
\]
vorkommt. Setzt man $U$ für $u = -1$ an Stelle von $A$ in die Formel (1), so folgt, dass $G$ auch die Substitution $\begin{pmatrix}0, & 1\\ -1, & 0\end{pmatrix}$ enthält, die mit $U$ zusammen ein System erzeugender Elemente von $L_p$ giebt (\S.\,65), so dass also $G$ mit $L_p$ identisch ist. Es folgt hieraus, dass $G$, wenn es nicht mit $L_p$ identisch sein soll, nur Substitutionen von der Form
\[
A = \begin{pmatrix}a, & b\\ 0, & a^{-1}\end{pmatrix}
\]
enthalten kann. Umgekehrt bilden alle diese Substitutionen eine Gruppe vom Grade $1/2\,p(p-1)$, also einen Theiler vom Index $p+1$.

Es sind darunter noch gewisse Theiler von grösserem Index enthalten, die man erhält, wenn man für $a$ nicht alle Werthe nimmt, sondern alle Werthe von der Form $a^s$, wenn $s$ ein Theiler von $1/2\,(p-1)$ ist. Der Grad einer solchen Gruppe ist $1/2\,p(p-1) : s$.


% ===== scan leaf 294, printed 273 =====
% !PRINTED 273
\weberpage{273}
\pdfbookmark[2]{§ 71. Divisoren der Gruppe Lp, deren Grad nicht durch p theilbar ist}{bm82}
Fassen wir die Gruppe $L_p$ als Gruppe der linearen Substitutionen
\[
\eta = \frac{a\xi+b}{c\xi+d}
\]
auf, und lassen darin $\xi, \eta$ nach \S.\,65 die Zahlen $\infty, 0, 1, \ldots, p-1$ durchlaufen, so giebt uns $G$ die Gruppe der ganzen linearen Substitutionen
\[
\eta = a^2\xi + ab.
\]
Das sind die Substitutionen, die die Ziffer $\infty$ ungeändert lassen, und also eine Permutationsgruppe von nur $p$ Ziffern liefern. Es sind dies dieselben speciellen metacyklischen Gruppen, die wir im siebzehnten Abschnitte des ersten Bandes betrachtet haben.

Die verschiedenen aus $G$ transformirten Gruppen lassen irgend einen der übrigen $p+1$ Indices ungeändert, und die Gesammtzahl dieser Gruppen ist $p+1$.

\begin{center}
\S.\,71.\\
\gesperrt{Divisoren der Gruppe $L_p$, deren Grad nicht durch $p$}\\
\gesperrt{theilbar ist.}
\end{center}

Um die übrigen Theiler von $L_p$ zu finden, deren Grad nicht durch $p$ theilbar ist, können wir Schritt für Schritt denselben Weg gehen, der uns im siebenten und achten Abschnitte zu der Bestimmung der Polyëdergruppen geführt hat. Wir können uns hier damit begnügen, die Hauptmomente der Ableitung hervorzuheben und wegen der Beweise auf die genannte Stelle zu verweisen, wo ganz dieselben Schlüsse zu machen waren.

Es ist hierzu erforderlich, die Gruppe $L_p$, wie im \S.\,65 auseinander gesetzt, als Gruppe linearer gebrochener Substitutionen
\weq{1}{\Theta(\xi) = \frac{a\xi+b}{c\xi+d}}
aufzufassen, und darin der Veränderlichen $\xi$ alle $p^2+1$ Zahlenwerthe des Congruenzkörpers $\mathfrak{E}_{2,p}$, einschliesslich $\infty$, beizulegen. Dadurch gewinnen wir den Vortheil, dass die Gleichung
\weq{2}{\xi = \frac{a\xi+b}{c\xi+d}}
%CONT>>


% ===== scan leaf 295, printed 274 =====
% !PRINTED 274
\weberpage{274}
%<<CONT
immer zwei Wurzeln hat. Diese beiden Wurzeln sind von einander
verschieden, wenn wir Substitutionen $p^{\mathrm{ten}}$ Grades ausschliessen,
die ja in einer Gruppe, deren Grad nicht durch $p$ theilbar ist, nicht
vorkommen können, und die nach \S.\,68, 1.\ die einzigen sind, für welche
die beiden Wurzeln von (2) zusammenfallen.

Diese beiden Wurzeln nennen wir die Pole von $\Theta$.

Wir untersuchen eine Gruppe $G$, deren Grad $n$ nicht durch $p$
theilbar ist, die ein Theiler von $L_p$ sein soll. Wenn in $G$ die
Substitutionen
\[
\Theta_1,\ \Theta_2,\ \ldots,\ \Theta_{\nu-1},
\]
aber keine anderen vorkommen, die denselben Pol $\alpha$ haben, so
nennen wir $\alpha$ \gesperrt{einen $\nu$-zähligen Pol} (\S.\,52).

Ist nun $S$ irgend eine Substitution der Gruppe $\mathfrak{E}_{2,p}$, so
erhalten wir, wenn wir $L_p$ durch $S^{-1}$ transformiren, eine mit $L_p$
isomorphe Gruppe $S L_p S^{-1}$, und $G$ geht durch dieselbe
Transformation in eine isomorphe Gruppe $S G S^{-1}$ über.

Obwohl die Substitutionen dieser letzteren Gruppe keine reellen
Co\"efficienten haben, so hat doch jede von ihnen zwei Pole, die man
erhält, wenn man die Substitution $S$ auf die Pole von $\Theta$
anwendet. Denn es ist, wenn $\alpha$ ein Pol von $\Theta$ ist:
\[
S\,\Theta\,S^{-1}\,S(\alpha) = S\,\Theta(\alpha) = S(\alpha).
\]
Wir haben also, wie in \S.\,51, 2.:

\gesperrt{1. Die Gruppe $G$ lässt sich so transformiren, dass eine
beliebige ihrer Substitutionen die Pole 0 und $\infty$ erhält, dass
also diese Substitution multiplicativ wird.}

Wir führen nun der Reihe nach die Sätze des \S.\,52 auf, soweit sie
hier in Frage kommen, und werden nur da, wo die veränderten
Voraussetzungen es erfordern, eine Ausführung hinzufügen:

\gesperrt{2. Ist $\alpha$ ein $\nu$-zähliger Pol, so bilden die
Substitutionen 1, $\Theta_1$, $\Theta_2$, . . ., $\Theta_{\nu-1}$ eine
cyklische Gruppe. Ihre Substitutionen können alle (durch
Transformation) in die Form gesetzt werden:}
\weq{3}{\Theta(\xi) = \gamma^{h\frac{p^2-1}{\nu}}\,\xi, \quad h = 0, 1, \ldots, \nu-1 \quad (\S.\,52, 3.).}


% ===== scan leaf 296, printed 275 =====
% !PRINTED 275
\weberpage{275}
Hierin bedeutet $\gamma$ eine primitive Wurzel, also $\gamma^{\frac{p^2-1}{\nu}}$
eine primitive $\nu^{\mathrm{te}}$ Einheitswurzel des Congruenzkörpers
$\mathfrak{E}_{2,p}$ (\S.\,67).

3. \gesperrt{Beide Pole einer Substitution $\Theta$ sind gleichzählig}
(\S.\,52, 4.).

Ist $Q$ die cyklische Gruppe $\nu^{\mathrm{ten}}$ Grades der Potenzen
von $\Theta$ und $n = \nu\mu$, so erhält man
\[
G = Q + \psi_1 Q + \cdots + \psi_{\mu-1} Q,
\]
worin $\psi_1, \ldots, \psi_{\mu-1}$ Substitutionen aus $G$ sind, und

4. \gesperrt{die Zahlen}
\weq{4}{\alpha,\ \psi_1(\alpha) = \alpha_1,\ \psi_2(\alpha) = \alpha_2,\ \ldots,\ \psi_{\mu-1}(\alpha) = \alpha_{\mu-1}}
\gesperrt{bilden ein System gleichzähliger conjugirter Pole der Gruppe $G$.}

Die sämmtlichen Pole der Gruppe $G$ lassen sich also in Systeme
conjugirter Pole zusammenfassen, und wir bekommen so die unbestimmte
Gleichung
\weq{5}{\begin{aligned}
2n - 2 &= \mu(\nu-1) + \mu'(\nu'-1) + \mu''(\nu''-1) \ldots \\
&= nh - \mu - \mu' - \mu'' \ldots,
\end{aligned}}
wenn $h$ die Anzahl der Systeme conjugirter Pole ist, und von dieser
Gleichung haben wir in \S.\,52 gesehen, dass sie nur eine beschränkte
Anzahl von Lösungen zulässt.

Um die diesen Lösungen entsprechenden Theiler der Congruenzgruppen zu
finden, können wir geradezu in den Formeln der Polyëdergruppen, die im
achten Abschnitte aufgestellt sind, für die darin vorkommenden
Einheitswurzeln die in \S.\,67 definirten Einheitswurzeln des Körpers
$\mathfrak{E}_{2,p}$ setzen, mit denen ja nach denselben Regeln
gerechnet wird. Dabei können natürlich nur solche Einheitswurzeln
vorkommen, deren Grad ein Theiler von $p^2 - 1$ ist.

Es ist schliesslich bei jeder solchen Gruppe noch zu untersuchen, ob
ihre Substitutionen in $L_p$ oder in einer damit isomorphen Gruppe
enthalten sind. Wir haben dann folgende Fälle von Lösungen der
Gleichung (5), wobei $\gamma$ stets eine primitive Wurzel des Körpers
$\mathfrak{E}_{2,p}$ bedeutet (\S.\,55).

\gesperrt{I. Cyklische Gruppen vom Grade $n$,}\quad
$\begin{aligned} h &= 2,\ \nu = \nu' = n \\ \mu &= \mu' = 1 \end{aligned}$
\[
C_n = \begin{pmatrix} \gamma^{\frac{p^2-1}{2n}\lambda}, & 0 \\[2pt] 0, & \gamma^{-\frac{(p^2-1)}{2n}\lambda} \end{pmatrix},
\quad \lambda = 0, 1, \ldots, n-1.
\]


% ===== scan leaf 297, printed 276 =====
% !PRINTED 276
\weberpage{276}
Diese Gruppe ist reell, wenn $n$ ein Theiler von $1/2\,(p-1)$ ist, weil
dann und nur dann die Exponenten von $\gamma$ durch $p+1$ theilbar
sind. Die Gruppe ist in $\Gamma_p$ und zugleich in $\Lambda_p$ (\S.\,69)
enthalten, wenn $\gamma^{\frac{p^2-1}{2n}}$, $\gamma^{-\frac{p^2-1}{2n}}$
conjugirt sind, wenn also
\[
\gamma^{-\frac{p^2-1}{2n}} = \gamma^{p\frac{p^2-1}{2n}}
\quad \text{oder} \quad
\gamma^{(p+1)\frac{p^2-1}{2n}} = 1
\]
ist, d. h. wenn $n$ ein Theiler von $1/2\,(p+1)$ ist. Dies stimmt mit
\S.\,68 überein, wonach andere cyklische Gruppen, als solche, deren
Grad gleich $p$ oder ein Theiler von $1/2\,(p-1)$ oder $1/2\,(p+1)$
ist, nicht vorkommen. In den übrigen Polyëdergruppen ist $h = 3$, und
wir haben:

II. \gesperrt{Diëdergruppen} vom Grade $2m$,\quad
$\begin{aligned} \nu &= \nu' = 2,\ \nu'' = m \\ \mu &= \mu' = m,\ \mu'' = 2 \end{aligned}$
\[
D_m = \begin{pmatrix} \gamma^{\frac{p^2-1}{2m}\lambda}, & 0 \\[2pt] 0, & \gamma^{-\frac{p^2-1}{2m}\lambda} \end{pmatrix},
\begin{pmatrix} 0, & \gamma^{\frac{p^2-1}{2m}\lambda} \\[2pt] -\gamma^{-\frac{p^2-1}{2m}\lambda}, & 0 \end{pmatrix},
\quad \lambda = 0, 1, \ldots, m-1.
\]
Diese Gruppe ist in $L_p$ enthalten, wenn $p \equiv 1 \pmod{2m}$, und
in $\Gamma_p$, wenn $p \equiv -1 \pmod{2m}$, wie im Falle I.

Hier ist die in \S.\,55 gegebene zweite Darstellung der Diëdergruppe
gewählt, in der die Unterscheidung der Fälle etwas einfacher wird, als
bei der ersten.

III. \gesperrt{Tetraëdergruppe,}\quad
$\begin{aligned} \nu &= 2,\ \nu' = 3,\ \nu'' = 3,\ n = 12 \\ \mu &= 6,\ \mu' = 4,\ \mu'' = 4. \end{aligned}$

Um diese Gruppe darzustellen, bemerken wir, dass im Körper
$\mathfrak{E}_{2,p}$ immer eine $8^{\mathrm{te}}$ Einheitswurzel
$\gamma^{\frac{p^2-1}{8}}$ existirt. Wir erhalten also nach \S.\,56 die
drei erzeugenden Substitutionen
\[
\Theta = \begin{pmatrix} \gamma^{\frac{p^2-1}{4}}, & 0 \\[2pt] 0, & \gamma^{-\frac{p^2-1}{4}} \end{pmatrix},
\quad
\psi = \begin{pmatrix} 0, & \gamma^{\frac{p^2-1}{4}} \\[2pt] -\gamma^{-\frac{p^2-1}{4}}, & 0 \end{pmatrix},
\quad
\psi_1 = \psi\,\Theta = \begin{pmatrix} 0, & 1 \\ -1, & 0 \end{pmatrix}
\]
\[
\chi = \begin{pmatrix} \dfrac{1}{\sqrt{-2}}\,\gamma^{\frac{p^2-1}{8}}, & \dfrac{1}{\sqrt{-2}}\,\gamma^{\frac{p^2-1}{8}} \\[8pt] \dfrac{1}{\sqrt{-2}}\,\gamma^{-\frac{p^2-1}{8}}, & -\dfrac{1}{\sqrt{-2}}\,\gamma^{-\frac{p^2-1}{8}} \end{pmatrix}.
\]
1) Ist $p \equiv 1 \pmod 4$, so ist $\gamma^{\frac{p^2-1}{4}}$ reell,
und $\gamma^{\frac{p^2-1}{8}}$ ist reell oder rein imaginär, je
nachdem $p \equiv 1$ oder $\equiv 5 \pmod 8$ ist.


% ===== scan leaf 298, printed 277 =====
% !PRINTED 277
\weberpage{277}
Uebereinstimmend damit ist auch $\sqrt{-2}$ reell oder rein imaginär
(Bd.\,I, \S.\,138, 4., 6.), und folglich ist in diesen Fällen $\Theta$,
$\psi$, $\chi$ und mithin die ganze Tetraëdergruppe reell.

2) Ist aber $p \equiv 3 \pmod 4$, so sind $\gamma^{\frac{p^2-1}{4}}$,
$\gamma^{-\frac{p^2-1}{4}}$ conjugirt imaginär, weil
$\gamma^{(p+1)\frac{p^2-1}{4}} = 1$ ist.

Ist $p \equiv 3 \pmod 8$, so ist $\sqrt{-2}$ reell und
$\gamma^{\frac{p^2-1}{8}(p+1)} = -1$, also $\gamma^{\frac{p^2-1}{8}}$
und $-\gamma^{\frac{p^2-1}{8}}$ conjugirt imaginär, und ist $p \equiv 7
\pmod 8$, so ist $\sqrt{-2}$ rein imaginär, $\gamma^{\frac{p^2-1}{8}}$
und $\gamma^{-\frac{p^2-1}{8}}$ conjugirt imaginär, und folglich gehört
$\Theta$, $\psi_1$, $\chi$ in diesen Fällen zur Gruppe $\Gamma_p$. Es
giebt also in \gesperrt{allen Fällen in der Gruppe} $L_p$ eine Tetraëdergruppe.

Für $p = 5$ z. B. kann man, da $-2$ quadratischer Nichtrest von 5 ist,
$\varepsilon = \sqrt{-2}$, und, da $\pm 2$ die beiden primitiven
Wurzeln von 5 sind, etwa
\[
\gamma^6 = -2, \quad \gamma^3 = \sqrt{-2}, \quad \gamma = 1 + \sqrt{-2}
\]
setzen, dann folgt
\weq{6}{\begin{aligned}
&\Theta = \begin{pmatrix} 2, & 0 \\ 0, & 3 \end{pmatrix},\quad
\psi = \begin{pmatrix} 0, & 2 \\ 2, & 0 \end{pmatrix},\quad
\psi_1 = \begin{pmatrix} 0, & 1 \\ -1, & 0 \end{pmatrix} \\
&\chi = \begin{pmatrix} 1, & 1 \\ 2, & -2 \end{pmatrix},\quad
\chi^2 = \begin{pmatrix} 2, & 1 \\ 2, & -1 \end{pmatrix},
\end{aligned}}
und daraus ergiebt sich die gesuchte Tetraëdergruppe für $p = 5$
[\S.\,56, (3)]:
\weq{7}{\begin{aligned}
&\begin{pmatrix} 1, & 0 \\ 0, & 1 \end{pmatrix},
\begin{pmatrix} 2, & 0 \\ 0, & -2 \end{pmatrix},
\begin{pmatrix} 0, & 2 \\ 2, & 0 \end{pmatrix},
\begin{pmatrix} 0, & 1 \\ -1, & 0 \end{pmatrix} \\
&\begin{pmatrix} 1, & 1 \\ 2, & -2 \end{pmatrix},
\begin{pmatrix} 2, & 2 \\ 1, & -1 \end{pmatrix},
\begin{pmatrix} -1, & 1 \\ 2, & 2 \end{pmatrix},
\begin{pmatrix} -2, & 2 \\ 1, & 1 \end{pmatrix} \\
&\begin{pmatrix} 2, & 1 \\ 2, & -1 \end{pmatrix},
\begin{pmatrix} -1, & 2 \\ 1, & 2 \end{pmatrix},
\begin{pmatrix} 1, & 2 \\ 1, & -2 \end{pmatrix},
\begin{pmatrix} -2, & 1 \\ 2, & 1 \end{pmatrix}.
\end{aligned}}
Diese Gruppe ist ein Theiler von $L_5$ vom Index 5.

IV. \gesperrt{Octaëdergruppe.}

In der Octaëdergruppe ist eine cyklische Gruppe vom Grade 4 enthalten,
und eine solche Gruppe kann also nur existiren, wenn $1/2\,(p-1)$ oder
$1/2\,(p+1)$ durch 4 theilbar, d. h. wenn $p \equiv \pm 1 \pmod 8$ ist.


% ===== scan leaf 299, printed 278 =====
% !PRINTED 278
\weberpage{278}
Unter dieser Voraussetzung führen die Formeln §. 57 zur
Aufstellung einer Octaëdergruppe.

Die erzeugenden Substitutionen sind:
\[
\Theta = \begin{pmatrix} \gamma^{\frac{p^2-1}{8}}, & 0 \\ 0, & \gamma^{-\frac{p^2-1}{8}} \end{pmatrix}, \quad
\psi = \begin{pmatrix} 0, & \gamma^{\frac{p^2-1}{4}} \\ \gamma^{\frac{p^2-1}{4}}, & 0 \end{pmatrix}
\]
\weq{8}{
\chi = \begin{pmatrix} \dfrac{1}{\sqrt{2}}\,\gamma^{-\frac{p^2-1}{8}}, & \dfrac{1}{\sqrt{2}}\,\gamma^{-\frac{p^2-1}{8}} \\[2mm] \dfrac{-1}{\sqrt{2}}\,\gamma^{\frac{p^2-1}{8}}, & \dfrac{1}{\sqrt{2}}\,\gamma^{\frac{p^2-1}{8}} \end{pmatrix}.
}
Diese Gruppe ist reell, wenn $p\equiv 1\ (\mathrm{mod}\ 8)$ ist, und imaginär, aber in $\Gamma_p$ enthalten, wenn $p\equiv -1\ (\mathrm{mod}\ 8)$ ist.

Das erste Beispiel ist $p=7$. Hier ist $-1$ quadratischer Nichtrest, und man kann also $\varepsilon=\sqrt{-1}=i$ setzen. Nehmen wir $\gamma=2+i$, $\gamma^7=2-i$, $\gamma^6=2\,(1+i)$, $\gamma^8=5$, so ist $\gamma$ eine primitive Wurzel, da $\gamma,\gamma^2,\dots,\gamma^7$ von einander verschieden sind und $5$ reelle primitive Wurzel von $7$ ist, so dass jede nicht verschwindende Zahl des Körpers $\mathfrak{E}_{2,7}$ in der Form
\[
5^{\mu}\gamma^{\nu} = \gamma^{8\mu+\nu}
\]
dargestellt werden kann, wenn
\[
\begin{gathered}
\mu = 0,\ 1,\ 2,\ 3,\ 4,\ 5,\\
\nu = 0,\ 1,\ 2,\ 3,\ 4,\ 5,\ 6,\ 7
\end{gathered}
\]
gesetzt wird.

Es ist dann ferner $\sqrt{2}=3$, und folglich sind die Substitutionen $\Theta,\psi,\chi$:
\weq{9}{
\begin{pmatrix} 2\,(1+i), & 0 \\ 0, & 2\,(1-i) \end{pmatrix}, \
\begin{pmatrix} 0, & i \\ i, & 0 \end{pmatrix}, \
\begin{pmatrix} 4\,(1-i), & 4\,(1-i) \\ -4\,(1+i), & 4\,(1+i) \end{pmatrix}.
}
Will man zur reellen Form übergehen, so muss man nach §. 69 verfahren. Man geht von den Substitutionen $B$ zu den $A$ über, indem man $\nu=\gamma^3=2-3i$ setzt, und erhält aus (9) für $\Theta,\psi,\chi$ in der Gruppe $\Lambda_p$:
\weq{10}{
\begin{pmatrix} 2+2i, & 0 \\ 0, & 2-2i \end{pmatrix}, \
\begin{pmatrix} 0, & 3-2i \\ 3+2i, & 0 \end{pmatrix}, \
\begin{pmatrix} 4-4i, & 1-4i \\ 1+4i, & 4+4i \end{pmatrix}.
}
Hier ist ferner
\[
R = \begin{pmatrix} 4, & 4i \\ i, & 1 \end{pmatrix}
\]
%CONT>>


% ===== scan leaf 300, printed 279 =====
% !PRINTED 279
\weberpage{279}
%<<CONT
zu setzen, woraus sich durch Bildung von $R\,A\,R^{-1}$ die erzeugenden Substitutionen der Octaëdergruppe in reeller Form ergeben:
\[
\Theta = \begin{pmatrix} 2, & 1 \\ 3, & 2 \end{pmatrix}, \
\psi = \begin{pmatrix} 2, & 2 \\ 1, & -2 \end{pmatrix}, \
\omega = \psi\,\Theta = \begin{pmatrix} 3, & -1 \\ 3, & -3 \end{pmatrix}
\]
\weq{11}{
\chi = \begin{pmatrix} 0, & 2 \\ 3, & 1 \end{pmatrix}, \
\chi^2 = \begin{pmatrix} -1, & 2 \\ 3, & 0 \end{pmatrix}.
}
Nun ist es leicht, nach §. 57 die vollständige Gruppe zu bilden, was nicht nöthig ist, weiter auszuführen.

\gesperrt{Die Octaëdergruppe für $p=7$ ist ein Theiler von $L_7$ vom Index 7.}

V. \gesperrt{Ikosaëdergruppe.}

In dem Falle $p=5$ ist $L_p$ selbst eine Ikosaëdergruppe. Für andere Werthe von $p$ kann nur dann eine Ikosaëdergruppe in $L_p$ enthalten sein, wenn $p^2-1$ durch 5 theilbar, also $p\equiv\pm 1\ (\mathrm{mod}\ 5)$ ist. Dann erhalten wir aus §. 58 die erzeugenden Substitutionen einer solchen Gruppe.

Setzen wir zur Abkürzung
\[
\gamma^{\frac{p^2-1}{10}} = \varrho,
\]
so haben wir $\varrho^2$ für $\varepsilon$ in die Formeln des §. 58 einzusetzen, und erhalten
\weq{12}{
\Theta = \begin{pmatrix} \varrho, & 0 \\ 0, & \varrho^{-1} \end{pmatrix}, \
\psi = \begin{pmatrix} 0, & 1 \\ -1, & 0 \end{pmatrix}, \quad
\chi = \begin{pmatrix} \dfrac{-1}{\varrho^2-\varrho^{-2}}, & \dfrac{1}{\varrho-\varrho^{-1}} \\[2mm] \dfrac{1}{\varrho-\varrho^{-1}}, & \dfrac{1}{\varrho^2-\varrho^{-2}} \end{pmatrix}\ [\S.\,58,\,(26)].
}
Ist $p\equiv 1\ (\mathrm{mod}\ 5)$, so ist $\varrho$, und damit die ganze Gruppe reell. Ist aber $p\equiv -1\ (\mathrm{mod}\ 5)$, so sind $\varrho$ und $\varrho^{-1}$ conjugirt, und folglich ist die gefundene Gruppe in $\Gamma_p$ enthalten, und um die Ikosaëdergruppe in $L_p$ zu erhalten, müssen wir erst nach §. 69 transformiren.

\gesperrt{Für $p=11$ erhalten wir unmittelbar eine reelle Ikosaëdergruppe vom Index 11.}


% ===== scan leaf 301, printed 280 =====
% !PRINTED 280
\weberpage{280}
Wir können für $\varrho=\gamma^{12}$ eine beliebige primitive Wurzel von 11, z. B. 8, wählen. Dann wird $\varrho^{-1}=7$ und
\[
\varrho-\varrho^{-1}=1,\quad \varrho^2-\varrho^{-2}=4,
\]
also
\weq{13}{
\Theta = \begin{pmatrix} 3, & 0 \\ 0, & 4 \end{pmatrix}, \
\psi = \begin{pmatrix} 0, & 1 \\ -1, & 0 \end{pmatrix}, \
\chi = \begin{pmatrix} -3, & 1 \\ 1, & 3 \end{pmatrix}.
}
Damit ist also festgestellt, welche Arten von Theilern in der Gruppe $L_p$ vorkommen können, und es ist auch gezeigt, dass alle diese Theiler wirklich vorhanden sind. Die Frage, wie man für jeden Typus alle überhaupt möglichen Theiler erhält, haben wir hier bei Seite gelassen. Wir wollen darüber nur noch folgende Bemerkungen machen.

Wenn man eine gefundene Gruppe $G$ durch irgend eine Substitution der Gruppe $L_p$ transformirt, so erhält man einen conjugirten Theiler. Wenn man aber durch eine imaginäre Substitution der Gruppe $E_p$ transformirt, so kann es vorkommen, dass die transformirte Gruppe von $G$ trotzdem reell wird und nicht mit $G$ innerhalb $L_p$ conjugirt (wiewohl beide conjugirte Theiler von $E_p$ sind). So erhält man allgemein aus der Octaëder- und der Ikosaëdergruppe eine zweite nicht conjugirte, wenn man
mit $\begin{pmatrix} \varepsilon, & 0 \\ 0, & \varepsilon^{-1} \end{pmatrix}$ transformirt. Bei der Tetraëdergruppe giebt dies nur dann eine neue Gruppe, wenn $p\equiv\pm 1\ (\mathrm{mod}\ 8)$ ist. Die zweite Gruppe ergiebt sich nach der Formel:
\weq{14}{
\begin{pmatrix} \varepsilon, & 0 \\ 0, & \varepsilon^{-1} \end{pmatrix}
\begin{pmatrix} a, & b \\ c, & d \end{pmatrix}
\begin{pmatrix} \varepsilon^{-1}, & 0 \\ 0, & \varepsilon \end{pmatrix}
=
\begin{pmatrix} a, & bN \\ cN^{-1}, & d \end{pmatrix}
\quad (N=\varepsilon^2).
}
Um diese Verhältnisse an den Beispielen $p=5,7,11$ aufzuweisen, ist es zweckmässig, die oben gefundene Gruppe (6) für $p=7$ so zu transformiren, dass die Substitution $\chi$, die von der dritten Ordnung ist, die Normalform $S$ erhält. Man findet diese Transformation leicht, und erhält für diesen Fall:
\[
\begin{pmatrix} 1, & 1 \\ -2, & -1 \end{pmatrix}
\begin{pmatrix} 0, & 2 \\ 3, & 1 \end{pmatrix}
\begin{pmatrix} -1, & -1 \\ 2, & 1 \end{pmatrix}
=
\begin{pmatrix} 3, & 0 \\ 0, & 5 \end{pmatrix}
= \chi'
\]
\weq{15}{
\begin{pmatrix} 1, & 1 \\ -2, & -1 \end{pmatrix}
\begin{pmatrix} 2, & 1 \\ 3, & 2 \end{pmatrix}
\begin{pmatrix} -1, & -1 \\ 2, & 1 \end{pmatrix}
=
\begin{pmatrix} 1, & -2 \\ -1, & 3 \end{pmatrix}
= \Theta'
}
\[
\begin{pmatrix} 1, & 1 \\ -2, & -1 \end{pmatrix}
\begin{pmatrix} 2, & 2 \\ 1, & -2 \end{pmatrix}
\begin{pmatrix} -1, & -1 \\ 2, & 1 \end{pmatrix}
=
\begin{pmatrix} -3, & -3 \\ 1, & 3 \end{pmatrix}
= \psi',
\]
%CONT>>


% ===== scan leaf 302, printed 281 =====
% !PRINTED 281
\weberpage{281}

%<<CONT
woraus man durch Zusammensetzung

\weq{16}{\psi' \Theta' \chi' = \begin{pmatrix} 0, & 1 \\ -1, & 0 \end{pmatrix}}

erhält. Wir können demnach in den drei Fällen folgende erzeugende Substitutionen der Tetra\"eder-, Octa\"eder- und Ikosa\"edergruppe [(6), (13), (15)] annehmen:

\[
p = 5. \quad \begin{pmatrix} 2, & 0 \\ 0, & 3 \end{pmatrix}, \ \begin{pmatrix} 0, & 1 \\ -1, & 0 \end{pmatrix}, \ \begin{pmatrix} 2, & 2 \\ 1, & -1 \end{pmatrix}
\]
\[
p = 7. \quad \begin{pmatrix} 3, & 0 \\ 0, & 5 \end{pmatrix}, \ \begin{pmatrix} 0, & 1 \\ -1, & 0 \end{pmatrix}, \ \begin{pmatrix} 1, & -2 \\ -1, & 3 \end{pmatrix}
\]
\[
p = 11. \quad \begin{pmatrix} 3, & 0 \\ 0, & 4 \end{pmatrix}, \ \begin{pmatrix} 0, & 1 \\ -1, & 0 \end{pmatrix}, \ \begin{pmatrix} -3, & 1 \\ 1, & 3 \end{pmatrix}.
\]

Durch wiederholte Zusammensetzung, die sich auf sehr verschiedene Arten anordnen lässt, ergeben sich hieraus leicht die vollständigen Gruppen:

\[
\begin{gathered}
p = 5.\\[2pt]
\begin{pmatrix} x, & 0 \\ 0, & x^{-1} \end{pmatrix}, \ \begin{pmatrix} 0, & x \\ -x^{-1}, & 0 \end{pmatrix}, \ \begin{pmatrix} x, & y \\ 2y^{-1}, & 3x^{-1} \end{pmatrix} \quad \begin{aligned} x &= 1, 2\\ y &= \pm 1, \pm 2 \end{aligned}
\end{gathered}
\]
\[
\begin{gathered}
p = 7.\\[2pt]
\begin{pmatrix} x, & 0 \\ 0, & x^{-1} \end{pmatrix}, \ \begin{pmatrix} 0, & x \\ -x^{-1}, & 0 \end{pmatrix}, \ \begin{pmatrix} x, & -xz \\ -2x^{-1}z^{-1}, & 3x^{-1} \end{pmatrix}, \ \begin{pmatrix} xz, & x \\ -3x^{-1}, & -2x^{-1}z^{-1} \end{pmatrix}\\[2pt]
x = 1, 2, 3; \quad z = 1, 2, 4\\
(z \text{ quadratischer Rest von } 7).
\end{gathered}
\]
\[
\begin{gathered}
p = 11.\\[2pt]
\begin{pmatrix} x, & 0 \\ 0, & x^{-1} \end{pmatrix}, \ \begin{pmatrix} 0, & x \\ -x^{-1}, & 0 \end{pmatrix}, \ \begin{pmatrix} x, & -xz \\ -x^{-1}z^{-1}, & 2x^{-1} \end{pmatrix}, \ \begin{pmatrix} xz, & x \\ -2x^{-1}, & -x^{-1}z^{-1} \end{pmatrix}\\[2pt]
x = 1, 2, 3, 4, 5; \quad z = 1, 3, 4, 5, 9\\
(z \text{ quadratischer Rest von } 11).
\end{gathered}
\]

Die Anwendung der Transformation (14) führt bei $p = 5$ zu keiner neuen Gruppe; bei $p = 7, 11$ erhält man je eine andere Gruppe, die aus diesen hervorgeht, wenn man für $z$ die Reihe der quadratischen Nichtreste statt der Reste setzt\textsuperscript{1}).

\bigskip
\hrule
\bigskip

\textsuperscript{1}) Die Eigenschaft der Congruenzgruppen, einfach zu sein, hat schon \gesperrt{Galois} gekannt. Ebenso waren ihm die Theiler vom Index $p$ für $p = 5, 7, 11$ bekannt. Eingehender untersucht sind diese Gruppen von
%CONT>>


% ===== scan leaf 303, printed 282 =====
% !PRINTED 282
\weberpage{282}
\pdfbookmark[2]{§ 72. Constitution der Gruppe L7 vom Grade 168}{bm83}

\begin{center}
\S.\,72.

\gesperrt{Constitution der Gruppe $L_7$ vom Grade 168.}
\end{center}

Unter den hier gefundenen einfachen Gruppen ist die nächste nach der Ikosa\"edergruppe, die hier als $L_5$ wiederkehrt, die Gruppe $L_7$ vom Grade 168, die, wie wir gesehen haben, eine Octa\"edergruppe als Theiler enthält. Da uns diese merkwürdige Gruppe später noch mehrfach begegnen wird, so wollen wir hier noch etwas näher auf ihren Bau eingehen.

Im \S.\,57 ist die Octa\"edergruppe durch drei Elemente $\chi, \omega, \Theta$ in der Form dargestellt:

\weq{1}{\chi^\lambda \omega^u \Theta^v, \quad \lambda = 0, 1, 2; \ \mu = 0, 1; \ \nu = 0, 1, 2, 3.} % [sic] formula prints Latin u, v; the parameter list prints Greek mu, nu

Zwischen diesen Elementen bestehen die Relationen

\weq{2}{\omega\chi = \chi^2\omega, \quad \Theta\omega = \omega\Theta^3, \quad \Theta\chi = \chi^2\omega\Theta^2, \quad \Theta^2\chi = \chi\omega\Theta^3,}

und diese Bedingungen haben wir, wenn noch die Grade 3, 2, 4 der Elemente $\chi, \omega, \Theta$ hinzukommen, als ausreichend nachgewiesen, um das System (1) als Octa\"edergruppe zu charakterisiren.

Aus (2) haben wir im \S.\,57 als Folgerungen die Formeln abgeleitet:

\weq{3}{\begin{gathered}
\omega\chi = \chi^2\omega, \quad \omega\chi^2 = \chi\omega, \quad \Theta^2\chi^2 = \chi^2\omega\Theta,\\
\Theta\chi = \chi^2\omega\Theta^2, \quad \Theta^2\chi = \chi\omega\Theta^3, \quad \Theta^3\chi = \chi^2\Theta,\\
\Theta\omega = \omega\Theta^3, \quad \Theta^2\omega = \omega\Theta^2, \quad \Theta^3\omega = \omega\Theta, \quad \Theta\chi^2 = \chi\Theta^3,
\end{gathered}}

die wir später mehrfach benutzen werden.

Im \S.\,71 haben wir die Octa\"edergruppe auch als Congruenzgruppe nach dem Modul 7 dargestellt und haben in den dortigen Formeln (11)

\weq{4}{\chi = \begin{pmatrix} 0, & 2 \\ 3, & 1 \end{pmatrix}, \ \omega = \begin{pmatrix} 3, & -1 \\ 3, & -3 \end{pmatrix}, \ \Theta = \begin{pmatrix} 2, & 1 \\ 3, & 2 \end{pmatrix}}

gefunden. Für das Folgende ist es aber, im Interesse einer einfacheren Rechnung, zweckmässig, diese Gruppe noch zu trans-
%CONT>>

\bigskip
\hrule
\bigskip

%<<CONT
\gesperrt{Serret} (Cours d'alg\`ebre supérieure) und von C. \gesperrt{Jordan} (Traité des Substitutions). Vergl. \gesperrt{Weber}, Elliptische Functionen etc., \S.\,84 f. Die vollständige Aufstellung aller Theiler rührt von \gesperrt{Gierster} her (Math. Ann., Bd. XVIII). Ausführliche Behandlung der Congruenzgruppen in \glqq \gesperrt{Klein-Fricke}, Vorlesungen über die Theorie der elliptischen Modulfunctionen\grqq{}, Bd. I, Leipzig 1890.


% ===== scan leaf 304, printed 283 =====
% !PRINTED 283
\weberpage{283}

%<<CONT
formiren. Es hat sich nämlich früher schon gezeigt, dass in der Gruppe $L_7$ Theiler vom Grade 21 enthalten sind, deren Index $= 8$ ist, und die gerade für unseren Zweck von besonderer Wichtigkeit sind. Unter den conjugirten Theilern des Index 8 ist aber der einfachste und für die Rechnung bequemste der aus allen Substitutionen

\weq{5}{\begin{pmatrix} a, & b \\ 0, & a^{-1} \end{pmatrix} \quad (\mathrm{mod}\ 7)}

bestehende, und wir wollen die Transformation also so einrichten, dass $\chi$ in dieser Form auftritt. Dies erreichen wir durch Transformation der Substitutionen (4) mittelst $\begin{pmatrix} 1, & 1 \\ -1, & 0 \end{pmatrix}$, und dadurch ergiebt sich aus (4)

\weq{6}{\chi = \begin{pmatrix} 2, & 3 \\ 0, & -3 \end{pmatrix}, \ \omega = \begin{pmatrix} 1, & -3 \\ 3, & -1 \end{pmatrix}, \ \Theta = \begin{pmatrix} 1, & 3 \\ -2, & 2 \end{pmatrix}.}

Wir stellen hiernach noch zur besseren Uebersicht die ganze Octa\"edergruppe in eine Tafel zusammen, deren sechs Zeilen die Elemente $\Theta^\lambda, \omega\Theta^\lambda, \chi\Theta^\lambda, \chi\omega\Theta^\lambda, \chi^2\Theta^\lambda, \chi^2\omega\Theta^\lambda$ für $\lambda = 0, 1, 2, 3$ enthalten:

\weq{7}{\begin{gathered}
\begin{pmatrix} 1, & 0 \\ 0, & 1 \end{pmatrix}, \ \begin{pmatrix} 1, & 3 \\ -2, & 2 \end{pmatrix}, \ \begin{pmatrix} 2, & 2 \\ 1, & -2 \end{pmatrix}, \ \begin{pmatrix} 2, & -3 \\ 2, & 1 \end{pmatrix};\\
\begin{pmatrix} 1, & -3 \\ 3, & -1 \end{pmatrix}, \ \begin{pmatrix} 0, & 3 \\ 2, & 0 \end{pmatrix}, \ \begin{pmatrix} -1, & 1 \\ -2, & 1 \end{pmatrix}, \ \begin{pmatrix} 3, & 1 \\ -3, & -3 \end{pmatrix};\\
\begin{pmatrix} 2, & 3 \\ 0, & -3 \end{pmatrix}, \ \begin{pmatrix} 3, & -2 \\ -1, & 1 \end{pmatrix}, \ \begin{pmatrix} 0, & 2 \\ 3, & 1 \end{pmatrix}, \ \begin{pmatrix} 3, & -3 \\ 1, & -3 \end{pmatrix};\\
\begin{pmatrix} 3, & 2 \\ 2, & -3 \end{pmatrix}, \ \begin{pmatrix} 1, & 1 \\ -1, & 0 \end{pmatrix}, \ \begin{pmatrix} 1, & 2 \\ 1, & 3 \end{pmatrix}, \ \begin{pmatrix} -3, & 0 \\ 2, & 2 \end{pmatrix};\\
\begin{pmatrix} 3, & 3 \\ 0, & -2 \end{pmatrix}, \ \begin{pmatrix} -3, & 1 \\ -3, & 3 \end{pmatrix}, \ \begin{pmatrix} 2, & 0 \\ -2, & -3 \end{pmatrix}, \ \begin{pmatrix} -2, & 1 \\ 3, & -2 \end{pmatrix};\\
\begin{pmatrix} -2, & 2 \\ 1, & 2 \end{pmatrix}, \ \begin{pmatrix} -1, & 2 \\ 3, & 0 \end{pmatrix}, \ \begin{pmatrix} 2, & 1 \\ 3, & 2 \end{pmatrix}, \ \begin{pmatrix} 0, & 1 \\ -1, & -1 \end{pmatrix}.
\end{gathered}}
Um die ganze Gruppe $L_7$ zu erhalten, müssen wir noch ein Element $7^{\text{ten}}$ Grades hinzufügen, und dafür wählen wir

\weq{8}{\tau = \begin{pmatrix} 1, & 1 \\ 0, & 1 \end{pmatrix}.}


% ===== scan leaf 305, printed 284 =====
% !PRINTED 284
\weberpage{284}

Dann ist die gesammte Gruppe $L_7$ so dargestellt:

\weq{9}{\tau^\varrho \chi^\lambda \omega^\mu \Theta^\nu, \quad \begin{aligned} \varrho &= 0, 1, 2, 3, 4, 5. 6\\ \lambda &= 0, 1, 2; \ \mu = 0, 1\\ \nu &= 0, 1, 2, 3. \end{aligned}} % [sic] full stop printed between 5 and 6

Dass alle diese Elemente (9) von einander verschieden sind, ergiebt sich einfach daraus, dass keine Potenz von $\tau$, deren Exponent nicht durch 7 theilbar ist, in der Octa\"edergruppe enthalten sein kann, weil die Octa\"edergruppe kein Element vom $7^{\text{ten}}$ Grade hat.

Die in der Gruppe $L_7$ enthaltene Gruppe (5) vom $21^{\text{sten}}$ Grade ist dann in der Form $\tau^\varrho \chi^\lambda$ dargestellt.

Um die Composition der Elemente (9) zu charakterisiren, genügt es offenbar, wenn für jedes $\varrho$ die Elemente

\weq{10}{\Theta\tau^\varrho, \ \omega\tau^\varrho, \ \chi\tau^\varrho}

in der Form (9) dargestellt sind, weil dadurch, zusammen mit den Compositionen in der Octa\"edergruppe, das symbolische Product je zweier Elemente der Form (9) wieder in der Form (9) dargestellt werden kann.

Da nun

\weq{11}{\tau^\varrho = \begin{pmatrix} 1, & \varrho \\ 0, & 1 \end{pmatrix}.}

ist, so erhält man zunächst sehr einfach aus (6)

\weq{12}{\chi\tau^\varrho = \tau^{4\varrho}\chi,}

und daraus

\weq{13}{\chi^2\tau^\varrho = \tau^{2\varrho}\chi^2,}

und man sieht leicht, dass diese sechs Relationen eine Folge von der einen sind:

\weq{14}{\chi\tau = \tau^4\chi.}

Die übrigen Compositionen (10) erhält man aber nur durch wirkliche Ausrechnung in den einzelnen Fällen, wobei die Tabelle (7) gute Dienste leistet. Man wendet sie in der Weise an, dass man für jeden Werth $\varrho$ die Elemente

\[
\tau^{\varrho'} \omega \tau^\varrho, \quad \tau^{\varrho'} \Theta \tau^\varrho
\]

bildet, und $\varrho'$ so bestimmt, dass sich diese Elemente in der Tabelle finden, was immer nur auf eine Art möglich ist.


% ===== scan leaf 306, printed 285 =====
% !PRINTED 285
\weberpage{285}
Man findet so durch leichte, wenn auch etwas umständliche
Rechnung
\weq{15}{\begin{array}{rl@{\qquad}rl}
\omega\tau &{}= \tau^2\chi^2\omega\Theta^2, & \Theta\tau &{}= \tau^3\omega\Theta\\
\omega\tau^2 &{}= \tau\chi^2\Theta^3, & \Theta\tau^2 &{}= \tau\chi\omega\Theta^3\\
\omega\tau^3 &{}= \tau^5\chi\Theta^2, & \Theta\tau^3 &{}= \tau^5\chi\omega\\
\omega\tau^4 &{}= \tau^4\chi^2\Theta, & \Theta\tau^4 &{}= \tau^6\omega\Theta^2\\
\omega\tau^5 &{}= \tau^3\chi^2\omega\Theta, & \Theta\tau^5 &{}= \tau^2\Theta^3\\
\omega\tau^6 &{}= \tau^6\omega\Theta^3, & \Theta\tau^6 &{}= \tau^4\chi^2\Theta^2.
\end{array}}
Diese zwölf Relationen lassen sich aber alle aus vieren von
ihnen, die man auf mannigfaltige Art auswählen kann, als Folgerungen
 ableiten. Man kann z.\,B. für diese vier fundamentalen
Relationen die folgenden wählen:
\weq{16}{\begin{array}{c}
\omega\tau = \tau^2\chi^2\omega\Theta^2, \quad \omega\tau^3 = \tau^5\chi\Theta^2, \quad \omega\tau^4 = \tau^4\chi^2\Theta,\\
\Theta\tau = \tau^3\omega\Theta,
\end{array}}
aus denen mit Benutzung der Formeln (3), (12), (13) alle anderen
leicht folgen, z.\,B.:
\[
\omega\tau^5 = \tau^4\chi^2\Theta\tau = \tau^4\chi^2\tau^3\omega\Theta = \tau^3\chi^2\omega\Theta,
\]
und so die übrigen.

Wie wir früher gesehen haben, dass wir als erzeugende
Elemente der Octa\"edergruppe $\chi$ und $\Theta$ betrachten können, so
können wir jetzt als Erzeugende der Gruppe $L_7$ die zwei Elemente
 $\omega, \tau$ ansehen, denn es ergiebt sich aus (15):
\weq{17}{\Theta^3 = \omega\tau\omega\tau^6, \quad \chi = \tau^2\Theta\tau^3\omega.}
Fassen wir zusammen, so ergiebt sich folgendes Resultat:

\gesperrt{I. Sind vier Elemente $\tau, \chi, \omega, \Theta$ der Grade 7, 3, 2, 4
gegeben, bei deren Zusammensetzung die Relationen
 bestehen:}
\[
\begin{array}{ccc}
\omega\chi = \chi^2\omega, & \Theta\omega = \omega\Theta^3, & \Theta\chi = \chi^2\omega\Theta^2,\\
\Theta^2\chi = \chi\omega\Theta^3, & \omega\tau = \tau^2\chi^2\omega\Theta^2, & \omega\tau^3 = \tau^5\chi\Theta^2,\\
\omega\tau^4 = \tau^4\chi^2\Theta, & \Theta\tau = \tau^3\omega\Theta, & \chi\tau = \tau^4\chi,
\end{array}
\]
\gesperrt{so bilden die 168 Elemente}
\[
\sigma = \tau^\varrho \chi^\lambda \omega^u \Theta^v,
\]
\gesperrt{wenn $\varrho, \lambda, \mu, \nu$ volle Restsysteme nach den Moduln
 7, 3, 2, 4 durchlaufen, eine einfache Gruppe
168\textsuperscript{sten} Grades, und $\omega$ und $\tau$ können als erzeugende
 Elemente dieser Gruppe angesehen
werden.} % [sic] exponents in the enunciation text (mu, nu) do not match the exponents u, v used in the preceding formula for sigma


% ===== scan leaf 307, printed 286 =====
% !PRINTED 286
\weberpage{286}
Unter den Theilern dieser Gruppe sind hervorzuheben die
Octa\"edergruppe
\[
\chi^\lambda \omega^u \Theta^v
\]
vom Index 7, sodann die Elemente
\[
\tau^\varrho \chi^\lambda,
\]
die nach den Relationen (12) und (13) eine Gruppe 21\textsuperscript{sten} Grades,
also einen Theiler der Gesammtgruppe vom Index 8 bilden;
ferner die Gruppe 8\textsuperscript{ten} Grades $\omega^u \Theta^v$. Die gesammte Gruppe
lässt sich auch in anderer Reihenfolge so darstellen:
\[
\chi^\lambda \omega^u \Theta^v \tau^\varrho, \quad \omega^u \Theta^v \chi^\lambda \tau^\varrho, \quad \tau^\varrho \omega^u \Theta^v \chi^\lambda.
\]

\bigskip
\begin{center}
\rule{4cm}{0.4pt}
\end{center}


% ===== scan leaf 308, printed 287 =====
\weberpage{287}% page number not printed on this page
% !PRINTED none
% Part-title page (no running head, no page number).
\begin{center}
\gesperrt{DRITTES BUCH.}

\rule{2.5cm}{0.4pt}

\vspace{1.5cm}

\gesperrt{ANWENDUNGEN}

\vspace{0.5cm}

\gesperrt{DER}

\vspace{0.5cm}

\gesperrt{GRUPPENTHEORIE.}

\vspace{1.5cm}

\rule{5cm}{0.4pt}
\end{center}


% ===== scan leaf 309, printed 288 =====
\weberpage{288}% page number not printed on this page
% !PRINTED none
% Blank page (verso of the part-title leaf). No text or page number is printed.


% ===== scan leaf 310, printed 289 =====
\weberpage{289}% page number not printed on this page
\pdfbookmark[0]{Drittes Buch. Anwendungen der Gruppentheorie.}{bm84}
\pdfbookmark[1]{Zehnter Abschnitt. Allgemeine Theorie der metacyklischen Gleichungen.}{bm85}
\pdfbookmark[2]{§ 73. Die Resolventen der Compositionsreihen}{bm86}
% !PRINTED none
\weberpage{none}

\begin{center}
\gesperrt{Zehnter Abschnitt.}

\bigskip

\gesperrt{Allgemeine Theorie der metacyklischen Gleichungen.}
\end{center}

\bigskip
\begin{center}
\rule{4cm}{0.4pt}
\end{center}
\bigskip

\begin{center}
\S.\,73.

\medskip

\gesperrt{Die Resolventen der Compositionsreihe.}
\end{center}

\bigskip

Aus den S\"atzen der allgemeinen Gruppentheorie, die im ersten Buche dieses Bandes
behandelt sind, ergeben sich wichtige algebraische Folgerungen, wenn man sie auf die
\gesperrt{Galois}'sche Gruppe einer Gleichung anwendet.

Es wird jetzt ein beliebiger K\"orper $\Omega$ als Rationalit\"atsbereich angenommen,
$f(x) = 0$ sei eine Gleichung in diesem K\"orper ohne mehrfache Wurzeln und $P$ ihre
\gesperrt{Galois}'sche Gruppe. Diese Gruppe $P$ habe einen Normaltheiler $Q$. Wir bezeichnen mit
$n$ den Grad von $P$ und mit $j$ den Index des Theilers $Q$, so dass $n : j$ der Grad
von $Q$ ist.

Wir haben im ersten Bande (\S.\,156) gesehen, dass die Gruppe von $f(x)$ durch
Adjunction einer Wurzel einer irreduciblen Normalgleichung $j^{\mathrm{ten}}$ Grades
auf $Q$ reducirt wird und diese H\"ulfsgleichung $j^{\mathrm{ten}}$ Grades haben wir
als \gesperrt{Partialresolvente} bezeichnet.

Die \gesperrt{Galois}'sche Gruppe dieser Partialresolvente haben wir so erhalten: Bedeutet
$\psi$ eine zu der Gruppe $Q$ geh\"orige Function der Wurzeln von $f(x)$ und ist $P$
in die Nebengruppen
\[
P = Q + Qa + Qb + Qc + \cdots
\]
zerlegt, wo also $a, b, c, \ldots$ gewisse Permutationen aus $P$ sind, so geht $\psi$
durch die ganze Nebengruppe $Qa$ in ein und dieselbe Function $\psi_a$ \"uber, und die
\gesperrt{Galois}'sche Gruppe der Resolvente besteht aus den Substitutionen
\[
(\psi, \psi), (\psi, \psi_a), (\psi, \psi_b), (\psi, \psi_c), \ldots
\]


% ===== scan leaf 311, printed 290 =====
% !PRINTED 290
\weberpage{290}

Setzt man zwei dieser Substitutionen zusammen, so hat man zu beachten, dass
\[
(\psi, \psi_b) = (\psi_a, \psi_{ab})
\]
ist, so dass man
\[
(\psi, \psi_a)(\psi, \psi_b) = (\psi, \psi_{ab})
\]
hat. Es setzen sich also diese Substitutionen ganz in derselben Weise zusammen, wie
nach \S.\,4 dieses Bandes die Nebengruppen, und es ergiebt sich daraus:

\noindent 1. \gesperrt{Die Galois'sche Gruppe der zu $Q$ geh\"origen Partialresolvente
ist isomorph mit der zu $Q$ complement\"aren Gruppe $P/Q$.}

Nehmen wir jetzt irgend eine Compositionsreihe von $P$ mit der zugeh\"origen
Indexreihe (\S.\,6):
\weq{1}{
\begin{gathered}
P,\ P_1,\ P_2,\ \ldots,\ P_{\mu-1},\ 1 \\
j_1,\ j_2,\ \ldots,\ j_{\mu-1},\ j_\mu,
\end{gathered}
}
so wird nach dem soeben Gesagten die Gruppe der Gleichung $f = 0$ von $P$ auf $P_1$
reducirt durch Adjunction \gesperrt{einer} Wurzel einer Normalgleichung vom Grade $j_1$. Dann
wird sie auf $P_2$ reducirt durch eine Wurzel einer Normalgleichung vom Grade $j_2$
u.\,s.\,f. und endlich wird die Gleichung vollst\"andig gel\"ost durch eine Wurzel
einer Normalgleichung vom Grade $j_\mu$.

Auf dies Aufl\"osungsverfahren f\"allt nun von dem Satze \"uber die
Unver\"anderlichkeit der Indexreihe (\S.\,6, I.) ein neues Licht.

\noindent 2. \gesperrt{Um die Gleichung $f = 0$ zu l\"osen, hat man nach einander je
eine Wurzel einer Normalgleichung der Grade $j_1, j_2, \ldots, j_\mu$ zu adjungiren.
Die Grade dieser Resolventen k\"onnen zwar in der Reihenfolge, nicht aber in der
Gesammtheit abge\"andert werden.}

Die Zahlen $j_1, j_2, \ldots, j_\mu$ h\"angen also weit tiefer mit der Natur einer
Gleichung oder allgemeiner mit den durch die Gleichung definirten K\"orpern zusammen,
als etwa der Grad der Gleichung; denn w\"ahrend der Grad durch Transformation auf
mannigfache Weise ver\"andert werden kann, wenn man Functionen der Wurzeln als neue
Unbekannte einf\"uhrt, bleiben die Zahlen $j_1, j_2, \ldots, j_\mu$ immer erhalten und
sind als wahre \gesperrt{Invarianten} des durch die Gleichung definirten
Normalk\"orpers zu betrachten.


% ===== scan leaf 312, printed 291 =====
% !PRINTED 291
\weberpage{291}

Zu diesen Invarianten geh\"ort auch der Grad $n$ der Gruppe $P$ selbst, der durch
die $j$ so bestimmt ist:
\weq{2}{n = j_1 j_2 \cdots j_{\mu-1} j_\mu.}
Denn nach der Bedeutung der Indices ist $n : j_1$ der Grad von $P_1$, $n : j_1 j_2$
der Grad von $P_2$ u.\,s.\,f., und da der letzte der Grade gleich $1$ ist, so ergiebt
sich die Formel (2).

Im Allgemeinen ist in einer Compositionsreihe von $P$ jedes Glied Normaltheiler nur
des n\"achst vorangehenden. Es k\"onnen aber auch einzelne Glieder vorkommen, die
auch noch von weiter vorangehenden Gliedern Normaltheiler sind. Besonders wichtig
sind solche Glieder, die Normaltheiler von $P$ selbst und also auch von allen ihnen
vorangehenden Gliedern der Compositionsreihe sind. Wir wollen sehen, welche
algebraische Consequenzen aus diesem Umstande zu ziehen sind.

Es sei $P_\nu$ ein Glied der Reihe (1), welches zugleich Normaltheiler von $P$ ist.
Der Index von $P_\nu$ in Bezug auf $P$ ist $j_1 j_2 \ldots j_\nu$, und dies Product
ist der Grad der Partialresolvente, durch die die Gruppe $P$ auf $P_\nu$ reducirt
wird, die wir mit $\chi(y) = 0$ bezeichnen wollen. Die Gruppe dieser Resolvente, die
eine Normalgleichung ist, erhalten wir nach dem Satze 1. in der Form $P/P_\nu$.

Wir k\"onnen nun leicht eine Compositionsreihe f\"ur diese Gruppe nebst der
zugeh\"origen Indexreihe finden, n\"amlich:
\weq{3}{
\begin{gathered}
P/P_\nu,\ P_1/P_\nu,\ P_2/P_\nu,\ \ldots,\ P_{\nu-1}/P_\nu,\ 1 \\
j_1,\quad j_2,\quad \ldots,\quad j_{\nu-1},\quad j_\nu.
\end{gathered}
}
Denn nach dem Satze 1., \S.\,6 ist $P_1/P_\nu$ ein Normaltheiler von $P/P_\nu$ vom
Index $j_1$, und es ist ein gr\"osster Normaltheiler,
weil nach 2., \S.\,6 \"uber $P_1/P_\nu$ kein Normaltheiler von $P/P_\nu$ stehen
kann, wenn \"uber $P_1$ kein Normaltheiler von $P$ steht. Und ebenso kann man in
Bezug auf die folgenden Glieder von (3) schliessen.

Daraus ziehen wir noch eine wichtige Folgerung. Wir haben schon im \S.\,7 gezeigt,
dass sich, wenn $Q$ irgend ein Normaltheiler von $P$ ist, eine Compositionsreihe von
$P$ finden l\"asst, in der $Q$ vorkommt. Ist nun $R$ ein anderer Normaltheiler von
$P$ und zugleich ein Theiler von $Q$, so kann man die Compositionsreihe von $P$ so
einrichten, dass $Q$ und $R$ darin vorkommen.

Wir denken uns eine solche Compositionsreihe bestimmt:
\weq{4}{P,\ Q',\ Q'',\ \ldots,\ Q,\ \ldots,\ R.}


% ===== scan leaf 313, printed 292 =====
% !PRINTED 292
\weberpage{292}
\pdfbookmark[2]{§ 74. Metacyklische Gleichungen}{bm87}

Nach (3) k\"onnen wir die Compositionsreihen der beiden Gruppen $P/Q$ und $P/R$
daraus herleiten, die wir so andeuten wollen:
\weq{5}{P/Q,\ Q'/Q,\ Q''/Q,\ \ldots,\ 1}
\weq{6}{P/R,\ Q'/R,\ Q''/R,\ \ldots,\ Q/R \ldots}
Nun ist der Index des Theilers $Q'/Q$ von $P/Q$ gleich dem Index des Theilers $Q'$
von $P$, also auch gleich dem Index des Theilers $Q'/R$ von $P/R$ (\S.\,6, 1.), und
Gleiches gilt von den folgenden Gliedern. Wir sprechen also den Satz aus:

\noindent 3. \gesperrt{Sind $Q$ und $R$ Normaltheiler von $P$, und ist $R$ ein
Theiler von $Q$, so ist die Indexreihe von $P/Q$ ein Theil der Indexreihe von
$P/R$..} % [sic] double period as printed

Sind die Indices $j_1, j_2, \ldots, j_\nu$ lauter Primzahlen, so ist die L\"osung
der Resolvente $\chi(y) = 0$ auf die L\"osung einer Kette von cyklischen
Gleichungen der Grade $j_1, j_2, \ldots, j_\nu$ reducirbar; diese Resolvente ist
also metacyklisch in dem Sinne, wie wir diesen Begriff im siebzehnten Abschnitte
des ersten Bandes festgestellt haben, wonach die metacyklischen Gleichungen mit den
sonst algebraisch l\"osbar genannten identisch sind.

Eine irreducible Gleichung $f(x) = 0$ ist metacyklisch, wenn die Indexreihe ihrer
Gruppe aus lauter Primzahlen besteht. Wir haben an der erw\"ahnten Stelle die
Bedingungen f\"ur metacyklische Gleichungen von Primzahlgrad untersucht, und
m\"ussen jetzt diese Betrachtungen f\"ur den allgemeinen Fall durchf\"uhren, dass
der Grad der Gleichung beliebig zusammengesetzt ist.

\begin{center}
\S.\,74.

\medskip

\gesperrt{Metacyklische Gleichungen.}
\end{center}

\bigskip

Wenn $f(x) = 0$ eine irreducible Gleichung $m^{\mathrm{ten}}$ Grades und $P$ ihre
\gesperrt{Galois}'sche Gruppe ist, so ist $P$ transitiv. Eine Compositions- und Indexreihe
f\"ur $P$ sei
\weq{1}{
\begin{gathered}
P,\ P_1,\ P_2,\ \ldots,\ P_{\mu-1},\ 1 \\
j_1,\ j_2,\ \ldots,\ j_{\mu-1},\ j_\mu.
\end{gathered}
}
Es wird, da die letzte Gruppe $1$ intransitiv ist, in der Compositionsreihe einmal
eine intransitive Gruppe auftreten, und es sei also $P_\lambda$ die erste
intransitive Gruppe der Reihe (1). Dann
%CONT>>


% ===== scan leaf 314, printed 293 =====
% !PRINTED 293
\weberpage{293}
%<<CONT
sind auch alle folgenden Gruppen $P_{\lambda+1}, P_{\lambda+2}, \ldots$, die ja alle
Theiler von $P_\lambda$ sind, intransitiv.

Es ist nun an die Sätze 1., 2., 3. im §.\,158 des ersten Bandes zu erinnern. Da
$P_\lambda$ ein Normaltheiler von $P_{\lambda-1}$ ist, so folgt aus jenen Sätzen,
dass $P_{\lambda-1}$ imprimitiv ist, und wenn durch die nöthigen Adjunctionen die
Gruppe von $f(x) = 0$ auf $P_{\lambda-1}$ reducirt ist, so wird durch weitere
Adjunction einer Wurzel einer Normalgleichung vom Grade $j_\lambda$ die Function
$f(x)$ in mehrere irreducible Factoren
\weq{2}{f(x) = f_1(x)\, f_2(x) \ldots f_h(x)}
zerfallen, die alle von gleichem Grade sind und deren Anzahl $h$ ein Theiler von
$j_\lambda$ ist.

Bezeichnen wir mit $m$ den Grad von $f(x)$, mit $m_1$ den Grad von $f_1(x)$, so ist
\weq{3}{m = h\, m_1.}
Nach dem angeführten Satze 1. im §.\,158, Bd. I haben die Gleichungen
$f_1 = 0, f_2 = 0, \ldots, f_h = 0$ alle dieselbe Gruppe, die wir mit $Q$ bezeichnen
wollen. Sind $\alpha, \alpha_1, \alpha_2, \ldots, \alpha_{m_1-1}$ die Wurzeln von
$f_1 = 0$, so erhalten wir die Gruppe $Q$, wenn wir die Permutationen der $\alpha$
sammeln, die durch $P_\lambda$ hervorgerufen werden. Es handelt sich zunächst um
das Verhältniss der Grade $p_\lambda$ und $q$ der beiden Gruppen $P_\lambda$ und $Q$.

Es kann mehrere verschiedene Permutationen in $P_\lambda$ geben, die dasselbe
Element in $Q$ erzeugen, die also dieselbe Permutation der $\alpha$ enthalten.
Sind $\pi_1, \pi_2$ zwei verschiedene Permutationen aus $P_\lambda$, die unter den
$\alpha$ dieselbe Permutation hervorrufen, so wird $\pi_2 \pi_1^{-1} = \pi_0$ eine
Permutation sein, die die $\alpha$ in Ruhe lässt, und es ist also
$\pi_2 = \pi_0 \pi_1$. Wenn umgekehrt $\pi_0$ irgend eine Permutation aus
$P_\lambda$ ist, die die Wurzeln $\alpha$ nicht permutirt, so wird die Permutation
$\pi_2 = \pi_0 \pi_1$ dieselbe Aenderung unter den $\alpha$ bedingen, wie $\pi_1$.

Es folgt hieraus, dass jede Permutation der $\alpha$ gleich oft durch die
Permutationen von $P_\lambda$ erzeugt wird, nämlich ebenso oft, als die $\alpha$
durch Permutationen aus $P_\lambda$ ungeändert bleiben. Bezeichnen wir diese Zahl
mit $p_0$, so ist also
\weq{4}{p_\lambda = p_0\, q,}
\gesperrt{oder der Grad $p_\lambda$ der Gruppe $P_\lambda$ ist ein Vielfaches des Grades $q$
von $Q$.}


% ===== scan leaf 315, printed 294 =====
% !PRINTED 294
\weberpage{294}
Wir haben nun im §.\,154, 7. des ersten Bandes den Satz bewiesen, dass der Grad
einer transitiven Permutationsgruppe immer durch die Anzahl der permutirten
Ziffern theilbar ist. Hier ist aber $f_1(x)$ irreducibel und demnach die Gruppe
$Q$ transitiv. Es ist also $q$ durch den Grad von $f_1(x)$, d.\,h. durch $m_1$
theilbar, und folglich ist nach (4) auch $p_\lambda$ durch $m_1$ theilbar.

Den Grad $p_{\lambda-1}$ von $P_{\lambda-1}$ erhalten wir nach der Bedeutung von
$j_\lambda$ in der Form:
\weq{5}{p_{\lambda-1} = j_\lambda\, p_\lambda.}
Nun ist $p_\lambda$ ein Vielfaches von $m_1$, $j_\lambda$ ein Vielfaches von $h$,
und $h\, m_1 = m$ gleich dem Grade von $f(x)$. Demnach ist
\weq{6}{p_{\lambda-1} = k\, m}
ein Vielfaches von $m$.

Aus der Bedeutung der $j$ ergiebt sich aber [§.\,73, (2)]
$p_{\lambda-1} = j_\lambda j_{\lambda-1} \ldots j_\mu$, % [sic] second factor $j_{\lambda-1}$ as printed; equation (7) gives $j_{\lambda+1}$.
und folglich haben wir die
Relation
\weq{7}{j_\lambda j_{\lambda+1} j_{\lambda+2} \ldots j_\mu = k\, m,}
woraus sich eine sehr merkwürdige Folgerung ziehen lässt.

Angenommen, es werde die irreducible Function $f(x)$ durch successive Adjunction
von Wurzeln cyklischer Gleichungen reducibel, dann lässt sich nach §.\,176, Bd. I
die Compositionsreihe (1) so geordnet annehmen, dass die Indices
$j_1, j_2, \ldots, j_\lambda$ Primzahlen sind. Es ist also, da $h$ ein Theiler von
$j_\lambda$ ist, $j_\lambda = h$, und $j_\lambda$ ist nach (3) eine in $m$
aufgehende Primzahl.

Wenn nun in $m$ ausser $j_\lambda$ noch eine zweite Primzahl $p$ aufgeht, so muss
nach (7) einer der Factoren $j_\lambda, j_{\lambda+1}, j_{\lambda+2}, \ldots, j_\mu$
durch $p$ theilbar sein, und sie können also gewiss nicht alle gleich $j_\lambda$
sein.

Daraus aber folgt nach dem Satze III, §.\,7, dass man eine Compositionsreihe von
$P_\lambda$:
\weq{8}{P_\lambda,\ P_{\lambda+1},\ \ldots,\ P_{\mu-1},\ 1}
so finden kann, dass unter den Gruppen $P_\lambda, P_{\lambda+1}, \ldots,
P_{\mu-1}$ eine, etwa $P_\nu$, ein \gesperrt{Normaltheiler von} $P$ ist.

Dieses $P_\nu$ ist aber als Theiler der intransitiven Gruppe $P_\lambda$ selbst
intransitiv, und daher muss $P$ selbst \gesperrt{imprimitiv} sein (Bd. I, §.\,158, 2).

Wir sprechen dies in folgender Form als Satz aus:

\gesperrt{1. Wenn im Grade einer irreduciblen Gleichung mehrere verschiedene
Primzahlen aufgehen, so}
%CONT>>


% ===== scan leaf 316, printed 295 =====
% !PRINTED 295
\weberpage{295}
%<<CONT
\gesperrt{kann diese Gleichung nur dann durch successive Adjunction von Wurzeln
cyklischer Gleichungen reducibel werden, wenn sie imprimitiv ist.}

Wir können über die Art und Weise der Reduction, ihre Möglichkeit vorausgesetzt,
noch einiges Nähere anführen. Ist $P_\nu$ die erste Gruppe der Reihe (8), die
Normaltheiler von $P$ ist, so ist nach dem eben angeführten Satze III, §.\,7 bei
richtiger Anordnung der Compositionsreihe:
\weq{9}{j_\lambda = j_{\lambda+1} \cdots = j_{\nu+1},}
also alle gleich derselben Primzahl, und die Resolvente $\chi = 0$, durch die die
Gruppe $P$ auf $P_\nu$ reducirt wird, hat folglich eine metacyklische Gruppe.

Bezeichnen wir mit $A, B, \ldots, S$ die Systeme der Intransitivität der Gruppe
$P_\nu$, deren Anzahl $s$ sei, und setzen
\weq{10}{m = r\, s,}
so wird durch Adjunction einer Wurzel von $\chi = 0$ die Function $f(x)$ in $s$
Factoren $r^{\mathrm{ten}}$ Grades zerfallen.

Die $A, B, \ldots, S$ sind Systeme der Imprimitivität von $P$ (Bd. I, §.\,158, 2.).

Bezeichnen wir mit $Q$ die Gesammtheit aller Permutationen von $P$, die die
einzelnen Systeme $A, B, \ldots, S$ an ihrer Stelle lassen und nur die Elemente
der Systeme unter sich vertauschen, so ist $Q$, wie wir im §.\,158 des ersten
Bandes gesehen haben, ein Normaltheiler von $P$, und durch Adjunction der Wurzeln
einer Hülfsgleichung $s^{\mathrm{ten}}$ Grades $\varphi(y) = 0$ zerfällt $f(x)$ in
$s$ Factoren $r^{\mathrm{ten}}$ Grades, denen die einzelnen Systeme
$A, B, \ldots, S$ als Wurzeln angehören:
\[
f(x) = f(x, y_\alpha)\, f(x, y_\beta)\, f(x, y_\gamma) \ldots
\]
Da durch die Hülfsgleichung $\varphi(y) = 0$ die Gruppe $P$ auf $Q$ reducirt wird,
so ist die Gruppe von $\varphi(y) = 0$ isomorph mit $P/Q$. Andererseits ist aber
auch die intransitive Gruppe $P_\nu$ ein Theiler von $Q$, da die Systeme
$A, B, \ldots, S$ durch $P_\nu$ nicht verschoben werden, und wir können also den
Satz §.\,73, 3. auf die drei Gruppen $P, Q, P_\nu$ anwenden. Da die Indexreihe
von $P/P_\nu$ nach Voraussetzung aus lauter Primzahlen besteht, so gilt nach
jenem Satze dasselbe von $P/Q$. Diese Gruppe ist metacyklisch und also ist auch
die Hülfsgleichung $\varphi(y) = 0$ metacyklisch. Es folgt also der Satz:
%CONT>>


% ===== scan leaf 317, printed 296 =====
% !PRINTED 296
\weberpage{296}
%<<CONT
\gesperrt{2. Wenn eine irreducible Gleichung, in deren Grad mehr als eine
Primzahl aufgeht, durch successive Adjunction von Radicalen in Factoren
zerfällt, so wird eine Zerfällung in $s$ Factoren $r^{\mathrm{ten}}$ Grades
herbeigeführt durch Adjunction der Wurzeln einer metacyklischen Gleichung
$r^{\mathrm{ten}}$ Grades.}

Dieser Satz enthält als speciellen Fall den von \gesperrt{Abel} herrührenden
Satz, dass eine irreducible Gleichung, deren Grad $m$ nicht die Potenz einer
Primzahl ist, nur dann durch Radicale lösbar sein kann, wenn sie durch
Adjunction der Wurzeln einer lösbaren Gleichung niedrigeren Grades, deren Grad
ein Theiler von $m$ ist, in Factoren zerfällt\textsuperscript{1)}. Damit ist die
Frage nach der Auflösung einer Gleichung durch Radicale oder auch nur der
Reduction einer Gleichung durch Radicale ausserordentlich vereinfacht. Man kann
in der That die fernere Untersuchung auf Gleichungen beschränken, deren Grad
eine Primzahl oder eine Primzahlpotenz ist, weil darauf alle anderen Fälle durch
wiederholte Anwendung des Satzes 2. zurückgeführt sind.

So gestattet z.\,B. die Frage nach allen durch Radicale lösbaren irreduciblen
Gleichungen $6^{\mathrm{ten}}$ Grades eine geradezu triviale Antwort:

\gesperrt{Um alle metacyklischen Gleichungen $6^{\mathrm{ten}}$ Grades in irgend einem
Körper $\Omega$ zu erhalten, adjungire man dem Körper $\Omega$ eine
Quadratwurzel und bilde in dem erweiterten Körper alle cubischen Gleichungen,
oder man adjungire die Wurzel einer cubischen Gleichung und bilde in dem
erweiterten Körper alle quadratischen Gleichungen.}

Als Beispiel einer solchen Gleichung führen wir die von \gesperrt{Hesse}
behandelte an, von der die Kreisschnitte einer nicht auf die Hauptaxen
bezogenen Fläche $2^{\mathrm{ten}}$ Grades abhängen. Dies Problem wird durch
Quadratwurzeln gelöst, wenn vorher durch die Lösung einer cubischen Gleichung
die Hauptaxen bestimmt sind\textsuperscript{2)}.

\noindent\rule{4cm}{0.4pt}\par
\textsuperscript{1)} \gesperrt{Abel} giebt den Satz ohne Beweis in der
Abhandlung \glqq Sur la rés. algébr. des équations\grqq{}. Oeuvres complètes
1881, Bd. II, S. 217. Vgl. auch die Abhandlung von \gesperrt{Galois} in
\gesperrt{Liouville's} Journal, Bd. 11.\par
\textsuperscript{2)} \gesperrt{Hesse}, \glqq Ueber die Auflösung derjenigen
Gleichungen $6^{\mathrm{ten}}$ Grades etc.\grqq{} Crelle's Journal, Bd. 41
(1851).


% ===== scan leaf 318, printed 297 =====
% !PRINTED 297
\weberpage{297}
\pdfbookmark[2]{§ 75. Metacyklische Gleichungen, deren Grad eine Primzahlpotenz ist}{bm88}
\begin{center}
\S.\,75.\\
\gesperrt{Metacyklische Gleichungen, deren Grad eine Primzahlpotenz ist.}
\end{center}

Nach den letzten S\"atzen concentrirt sich das Interesse weiterer Untersuchungen \"uber metacyklische Gleichungen haupts\"achlich auf den Fall, dass der Grad der Gleichung eine Potenz $p^k$ einer Primzahl $p$ ist. Wir setzen die Gleichung als irreducibel voraus und beschr\"anken uns auf die Betrachtung \gesperrt{primitiver Gleichungen}; denn die Aufl\"osung der imprimitiven reducirt sich auf die successive L\"osung zweier (oder mehrerer) primitiver Gleichungen, deren Grade gleichfalls Potenzen von $p$ sind, und die, wenn die urspr\"ungliche Gleichung metacyklisch ist, auch metacyklisch sein m\"ussen.

Wir nehmen also jetzt an, es sei $P$ die Gruppe einer irreduciblen primitiven metacyklischen Gleichung $f(x) = 0$ vom Grade $p^k$. Nach Bd. I, \S.\,158, 2. muss nicht nur $P$ selbst, sondern alle seine Normaltheiler (mit Ausnahme der Einheitsgruppe) noch \gesperrt{transitiv} sein. Nun wenden wir den in \S.\,8 allgemein bewiesenen Satz IV. an, nach dem $P$ einen Normaltheiler $Q$ mit lauter commutativen Elementen besitzt, und wenn mehrere solche Theiler vorhanden sind, so verstehen wir unter $Q$ einen von ihnen, dessen Grad m\"oglichst niedrig, aber noch gr\"osser als 1 ist. Diese Gruppe $Q$ ist also gleichfalls noch transitiv. $Q$ kann aber keinen von der Einheit verschiedenen echten Theiler mehr haben, der zugleich Normaltheiler von $P$ ist, weil sonst dieser an die Stelle von $Q$ treten w\"urde.

Wenn durch die geh\"origen Adjunctionen die Gruppe unserer Gleichung von $P$ auf $Q$ reducirt ist, so ist $f(x) = 0$ in dem erweiterten Rationalit\"atsbereiche zu einer \gesperrt{Abel'schen Gleichung} geworden, und da sie noch irreducibel geblieben ist, so ist nach Bd. I, \S.\,162 der Grad der Gleichung gleich dem Grade der Gruppe; d.\,h. $Q$ ist eine \gesperrt{Abel'sche Gruppe} vom Grade $p^k$. Die Grade aller Elemente von $Q$ sind also Potenzen von $p$. Wir wollen nachweisen, dass ausser dem Einheitselemente in $Q$ \gesperrt{nur Elemente vom Grade $p$ selbst vorkommen.}

Nehmen wir an, es sei $p^\lambda$ der h\"ochste Grad, der unter den Elementen von $Q$ vorkommt, und es sei $\lambda > 1$. Dann bilden alle Elemente von $Q$, deren Grad ein Theiler von $p^{\lambda-1}$ ist, einen
%CONT>>


% ===== scan leaf 319, printed 298 =====
% !PRINTED 298
\weberpage{298}
\pdfbookmark[2]{§ 76. Darstellung der Abel'schen Gruppe Q}{bm89}
%<<CONT
Theiler $Q'$ von $Q$, der sowohl von $Q$ selbst als von der Einheitsgruppe verschieden ist. Dieser Theiler $Q'$ muss aber ein Normaltheiler von $P$ sein, weil, wenn $\pi$ in $Q$, $\gamma$ in $P$ enthalten ist, $\gamma^{-1}\pi\gamma$, was vom selben Grade wie $\pi$ ist, gleichfalls in $Q$ enthalten sein muss, und also zu $Q'$ geh\"ort, wenn $\pi$ zu $Q'$ geh\"ort. Da nun aber, wie oben bemerkt, $Q$ keinen echten Theiler haben kann, der gr\"osser als die Einheitsgruppe und zugleich Normaltheiler von $P$ ist, so muss $\lambda = 1$ sein.

\begin{center}
\S.\,76.\\
\gesperrt{Darstellung der Abel'schen Gruppe $Q$.}
\end{center}

Die Betrachtungen des vorigen Paragraphen haben gezeigt, dass in einer metacyklischen Gruppe $P$, die als \gesperrt{Galois}'sche Gruppe einer irreduciblen primitiven Gleichung $f(x) = 0$ des Grades $n = p^k$ auftreten kann, eine \gesperrt{Abel'sche Gruppe} $Q$ als Normaltheiler enthalten sein muss, deren Grad $p^k$ ist, und die ausser dem Einheitselemente nur Elemente vom Grade $p$ enth\"alt.

Nach \S.\,9 l\"asst sich diese Gruppe $Q$ durch eine Basis darstellen, die aus $k$ Elementen $A_1, A_2, \ldots, A_k$ vom Grade $p$ besteht, so dass jedes Element von $Q$ die Form erh\"alt:
\weq{1}{A_1^{z_1} A_2^{z_2} \ldots A_k^{z_k},}
und hierin durchlaufen $z_1, z_2, \ldots, z_k$ von einander unabh\"angig je ein volles Restsystem nach dem Modul $p$. Wir gehen nun, um die entsprechende Permutationsgruppe zu finden, von einer beliebigen Wurzel $x$ von $f(x)$ aus, bezeichnen die Wurzel, in die $x$ durch die Permutation (1) \"ubergeht, mit
\weq{2}{[z_1, z_2, \ldots, z_k],}
und setzen fest, dass dies Zeichen seine Bedeutung nicht \"andern soll, wenn die Zahlen $z_1, z_2, \ldots, z_k$ beliebig um Vielfache von $p$ ver\"andert werden. Der Wurzel $x$, von der wir ausgingen, kommt dann das Zeichen $[0, 0, \ldots, 0]$ zu, und durch das Zeichen (2) ist jede Wurzel von $f(x)$ ein und nur einmal dargestellt.

Wenn $x$ durch eine Permutation $A'$ in $x'$ und $x'$ durch $A''$ in $x''$ \"ubergeht, so geht $x$ durch die Permutation $A' A''$ in $x''$ \"uber. Daraus folgt aber, dass $[z_1, z_2, \ldots, z_k]$ durch die Permutation
\weq{3}{A = A_1^{\alpha_1} A_2^{\alpha_2} \ldots A_k^{\alpha_k}}
%CONT>>


% ===== scan leaf 320, printed 299 =====
% !PRINTED 299
\weberpage{299}
\pdfbookmark[2]{§ 77. Analytische Darstellung der Permutationen}{bm90}
%<<CONT
in $[z_1+\alpha_1, z_2+\alpha_2, \ldots, z_k+\alpha_k]$ \"ubergeht. Demnach k\"onnen wir die Permutation $A$ auch so bezeichnen:
\weq{4}{\begin{pmatrix} z_1, & z_2, & \ldots, & z_k \\ z_1+\alpha_1, & z_2+\alpha_2, & \ldots, & z_k+\alpha_k \end{pmatrix},}
und man erh\"alt die ganze Gruppe $Q$, wenn man $\alpha_1, \alpha_2, \ldots, \alpha_k$ je ein volles Restsystem nach dem Modul $p$ durchlaufen l\"asst.

Die Gruppe $Q$, deren Bildungsweise jetzt festgestellt ist, muss ein Normaltheiler von $P$ sein, und daraus l\"asst sich eine allgemeine Form herleiten, in der alle Permutationen von $P$ enthalten sein m\"ussen.

Dazu bed\"urfen wir eines H\"ulfssatzes \"uber die analytische Darstellung der Permutationen, den wir zun\"achst ableiten.

\begin{center}
\S.\,77.\\
\gesperrt{Analytische Darstellung der Permutationen.}
\end{center}

\gesperrt{H\"ulfssatz. Wenn bei irgend einer Permutation der Gr\"ossen (2), \S.\,76, $z_1$ in $z_1'$, $z_2$ in $z_2'$, \ldots, $z_k$ in $z_k'$ \"ubergeht, so kann man immer
\weq{1}{
\begin{aligned}
z_1' &\equiv \varphi_1(z_1, z_2, \ldots, z_k)\\
z_2' &\equiv \varphi_2(z_1, z_2, \ldots, z_k)\\
&\;\cdots\cdots\cdots\cdots\\
z_k' &\equiv \varphi_k(z_1, z_2, \ldots; z_k)
\end{aligned}
\quad (\mathrm{mod}\ p).
}
setzen, worin $\varphi_1, \varphi_2, \ldots, \varphi_k$ ganze rationale Functionen der Variablen $z_1, z_2, \ldots, z_k$ mit ganzzahligen Co\"efficienten bedeuten, die in Bezug auf keine der Variablen den Grad $p - 1$ \"ubersteigen.}

Der Satz ist eine Verallgemeinerung des entsprechenden Satzes im \S.\,180 des ersten Bandes. Dort ist der Satz f\"ur $k = 1$ bewiesen. Dem allgemeinen Beweise, der auf der vollst\"andigen Induction beruht, schicken wir folgende Erw\"agungen voraus.

Da wir in (1) f\"ur die Argumente $z_i$ alle Combinationen von $k$ Reihen nach dem Modul $p$ zu setzen haben, so ergeben sich $p^k$ solcher Systeme. Jede der Functionen $\varphi$ hat $p^k$ unbestimmte Co\"efficienten, und wenn also die $z_k'$ gegeben sind, so erhalten wir $kp^k$ lineare Congruenzen f\"ur ebenso viele unbekannte Zahlen, n\"amlich die Co\"efficienten der $\varphi_k$. Es ist nur noch nachzuweisen,
%CONT>>


% ===== scan leaf 321, printed 300 =====
% !PRINTED 300
\weberpage{300}
%<<CONT
dass diese Congruenzen von einander unabh\"angig sind. Hierbei k\"onnen wir jede der Functionen $\varphi$ f\"ur sich betrachten. Setzen wir also
\weq{2}{z' \equiv \varphi(z_1, z_2, \ldots, z_k) \quad (\mathrm{mod}\ p),}
so haben wir, wenn wir f\"ur jede Combination der $z_i$ den zugeh\"origen Werth von $z'$ kennen, $p^k$ lineare Congruenzen zur Bestimmung der Co\"efficienten von $\varphi$. Es ist zu beweisen, dass diese Congruenzen immer l\"osbar sind, wobei wir voraussetzen k\"onnen, dass diese M\"oglichkeit f\"ur Functionen von weniger Variablen schon erwiesen sei.

Wenn wir nun $\varphi$ nach Potenzen der Variablen $z_1$ ordnen, so erhalten wir
\weq{3}{z' \equiv \chi_0 z_1^{p-1} + \chi_1 z_1^{p-2} + \cdots + \chi_{p-1} \quad (\mathrm{mod}\ p),}
worin die Co\"efficienten $\chi_0, \chi_1, \ldots, \chi_{p-1}$ ebensolche Functionen sind wie $\varphi$, nur dass sie von einer Variablen weniger abh\"angen.

Halten wir irgend eine der Combinationen der $z_2, \ldots, z_k$ fest, und setzen $z_1 = 0, 1, \ldots, p - 1$, so erhalten wir aus (3) ein System von $p$ linearen Congruenzen, dessen Determinante nicht durch $p$ theilbar ist, und wir k\"onnen daraus, wie im \S.\,180 des ersten Bandes, die Werthe von $\chi_0, \chi_1, \ldots, \chi_{p-1}$ f\"ur diese Combination der $z_2, \ldots, z_k$ bestimmen, und daher sind f\"ur jede Combination der Variablen die Werthe der Function $\chi_i$ (nach dem Modul $p$) bekannt. Der Voraussetzung nach k\"onnen wir aber die Co\"efficienten der Functionen $\chi_i$, die ja nur von $k - 1$ Variablen abh\"angen, daraus bestimmen, und damit ist der H\"ulfssatz bewiesen.

Hiernach k\"onnen wir jede Permutation der Wurzeln $[z_1, \ldots, z_k]$ so darstellen:
\weq{4}{\begin{pmatrix} z_1, & z_2, & \ldots \\ \varphi_1(z_1, z_2, \ldots), & \varphi_2(z_1, z_2, \ldots), & \ldots \end{pmatrix},}
oder in noch abgek\"urzterer Schreibweise:
\weq{5}{\begin{pmatrix} z \\ \varphi(z) \end{pmatrix}.}

In dem Symbol (4) k\"onnen wir, ohne seine Bedeutung zu \"andern, $z_1, z_2, \ldots$ durch irgend eine andere Combination $z_1', z_2', \ldots$ ersetzen, und wenn also nach dem H\"ulfssatze
\[
z_i' = \psi_i(z_1, z_2, \ldots)
\]
%CONT>>


% ===== scan leaf 322, printed 301 =====
% !PRINTED 301
\weberpage{301}
\pdfbookmark[2]{§ 78. Darstellung der metacyklischen Gruppe P}{bm91}
%<<CONT
ist, so können wir (4) oder (5) auch so darstellen:

\weq{6}{\begin{pmatrix} \psi_1(z_1, z_2, \dots), & \psi_2(z_1, z_2, \dots), \dots \\ \varphi_1(\psi_1, \psi_2, \dots), & \varphi_2(\psi_1, \psi_2, \dots), \dots \end{pmatrix} = \begin{pmatrix} \psi(z) \\ \varphi[\psi(z)] \end{pmatrix}.}

Dies führt zur Zusammensetzung der Permutationen:

\weq{7}{\begin{pmatrix} z \\ \psi(z) \end{pmatrix}\begin{pmatrix} z \\ \varphi(z) \end{pmatrix} = \begin{pmatrix} z \\ \varphi[\psi(z)] \end{pmatrix}.}

\begin{center}
§.\,78.

\gesperrt{Darstellung der metacyklischen Gruppe $P$.}
\end{center}

Nun soll die Gruppe $Q$ ein Normaltheiler der Gruppe $P$ sein, d. h. wenn $A$ eine Permutation aus $Q$, $B$ eine Permutation aus $P$ ist, so soll sich eine zweite Permutation $A'$ aus $Q$ so
bestimmen lassen, dass

\weq{1}{AB = BA'}

ist. Setzen wir in der abgekürzten Bezeichnung des vorigen Paragraphen

\[
A = \begin{pmatrix} z \\ z+\alpha \end{pmatrix}, \quad A' = \begin{pmatrix} z \\ z+\alpha' \end{pmatrix}, \quad B = \begin{pmatrix} z \\ \varphi(z) \end{pmatrix},
\]

so wird

\[
AB = \begin{pmatrix} z \\ \varphi(z+\alpha) \end{pmatrix}, \quad BA' = \begin{pmatrix} z \\ \varphi(z)+\alpha' \end{pmatrix},
\]

oder

\[
\varphi(z+\alpha) \equiv \varphi(z) + \alpha'.
\]

Diese Congruenz aber ist nur ein abgekürztes Symbol für das System der Congruenzen nach dem Modul $p$:

\weq{2}{\begin{aligned}
\varphi_1(z_1+\alpha_1, z_2+\alpha_2, \dots) &\equiv \varphi_1(z_1, z_2, \dots) + \alpha_1' \\
\varphi_2(z_1+\alpha_1, z_2+\alpha_2, \dots) &\equiv \varphi_2(z_1, z_2, \dots) + \alpha_2' \\
&\cdots\cdots\cdots\cdots\cdots\cdots\cdots\cdots\cdots\cdots
\end{aligned}}

Hierin kann das System der Zahlen $\alpha_1, \alpha_2, \dots$ nach dem Modul $p$ beliebig gegeben sein, $\alpha_1', \alpha_2', \dots$ sind dadurch bestimmt.

Setzen wir in der ersten der Formeln (2) je eine der Zahlen $\alpha_1, \alpha_2, \dots$ gleich 1, die übrigen gleich 0, so mag sich ergeben:

\weq{3}{\begin{aligned}
\varphi_1(z_1+1, z_2 \dots) &\equiv \varphi_1(z_1, z_2, \dots) + \alpha_{1,1} \\
\varphi_1(z_1, z_2+1, \dots) &\equiv \varphi_1(z_1, z_2, \dots) + \alpha_{1,2} \\
&\cdots\cdots\cdots\cdots\cdots\cdots\cdots\cdots\cdots\cdots
\end{aligned}}


% ===== scan leaf 323, printed 302 =====
% !PRINTED 302
\weberpage{302}
Wenn man die erste dieser Formeln $\alpha_1$ mal, die zweite $\alpha_2$ mal
nach einander anwendet u. s. f., so folgt:

\[
\begin{aligned}
\varphi_1(z_1+\alpha_1, z_2, \dots) &\equiv \varphi_1(z_1, z_2, \dots) + \alpha_1\alpha_{1,1} \\
\varphi_1(z_1, z_2+\alpha_2, \dots) &\equiv \varphi_1(z_1, z_2, \dots) + \alpha_2\alpha_{1,2} \\
&\cdots\cdots\cdots\cdots\cdots\cdots\cdots\cdots\cdots,
\end{aligned}
\]

und wenn man in der ersten $z_2$ in $z_2+\alpha_2$ verwandelt, und die
zweite anwendet und so fortfährt, so ergiebt sich schliesslich:

\weq{4}{\begin{aligned}
\varphi_1(z_1+\alpha_1, z_2+\alpha_2, \dots) &\equiv \\
\varphi_1(z_1, z_2, \dots) + \alpha_1\alpha_{1,1} + \alpha_2\alpha_{1,2} + \cdots
\end{aligned}}

Hierin setzen wir nun $z_1, z_2, \cdots = 0$ und schreiben dann an Stelle von
$\alpha_1, \alpha_2, \dots$ wieder $z_1, z_2, \dots$. Ferner bezeichnen
wir $\varphi_1(0, 0, \dots)$ durch $\alpha_1$ und erhalten so aus (4):

\weq{5}{\varphi_1(z_1, z_2, \dots) \equiv \alpha_1 + z_1\alpha_{1,1} + z_2\alpha_{1,2} + \cdots,}

d. h. $\varphi_1$ muss eine \gesperrt{lineare Function} sein.

Genau auf dieselbe Weise verfahren wir mit sämmtlichen Congruenzen (2) und gelangen so zu folgendem Endresultate:

\gesperrt{I. Alle Permutationen der Galois'schen Gruppe $P$ einer primitiven irreduciblen metacyklischen Gleichung vom Grade $p^k$}

\[
\begin{pmatrix} z_1, z_2, \dots, z_k \\ z_1', z_2', \dots, z_k' \end{pmatrix}
\]

\gesperrt{sind von der Form}

\weq{6}{\begin{aligned}
z_1' &\equiv \alpha_{1,1}z_1 + \alpha_{1,2}z_2 + \cdots + \alpha_{1,k}z_k + \alpha_1 \\
z_2' &\equiv \alpha_{2,1}z_1 + \alpha_{2,2}z_2 + \cdots + \alpha_{2,k}z_k + \alpha_2 \quad (\mathrm{mod}\ p). \\
&\cdots\cdots\cdots\cdots\cdots\cdots\cdots\cdots\cdots\cdots
\end{aligned}}

\[
z_k' \equiv \alpha_{k,1}z_1 + \alpha_{k,2}z_2 + \cdots + \alpha_{k,k}z_k + \alpha_k
\]

Das System der Congruenzen (6) stellt immer dann und nur dann eine Permutation dar, wenn die Determinante

\[
\sum \pm \alpha_{1,1}\alpha_{2,2}\cdots\alpha_{k,k}
\]

nicht durch $p$ theilbar ist, weil dann und nur dann auch umgekehrt zu jedem Systeme der $z'$ ein bestimmtes System der $z$ (nach dem Modul $p$) gehört.

Der Inbegriff aller Permutationen (6) ist in der That eine Gruppe, die die \gesperrt{allgemeine lineare Congruenzgruppe} heisst, und in ihr müssen die metacyklischen Gruppen $P$ ent-
%CONT>>


% ===== scan leaf 324, printed 303 =====
% !PRINTED 303
\weberpage{303}
%<<CONT
halten sein. Es sind aber nicht alle linearen Congruenzgruppen auch umgekehrt metacyklisch, und diese daraus auszusondern, ist noch ein weiteres Problem, das bis jetzt nur in besonderen Fällen gelöst ist $^{1)}$.

Um aber das Problem wenigstens genauer zu formuliren, fügen wir noch folgende Betrachtungen bei.

Wir bezeichnen die vollständige lineare Congruenzgruppe (6) mit $R$, und sondern daraus die darin enthaltene Gruppe der linearen homogenen Substitutionen

\weq{7}{\begin{aligned}
z_1' &\equiv \alpha_{1,1}z_1 + \alpha_{1,2}z_2 + \cdots + \alpha_{1,k}z_k \\
z_2' &\equiv \alpha_{2,1}z_1 + \alpha_{2,2}z_2 + \cdots + \alpha_{2,k}z_k \quad (\mathrm{mod}\ p) \\
&\cdots\cdots\cdots\cdots\cdots\cdots\cdots\cdots\cdots\cdots \\
z_k' &\equiv \alpha_{k,1}z_1 + \alpha_{k,2}z_2 + \cdots + \alpha_{k,k}z_k
\end{aligned}}

mit nicht verschwindender Determinante ab, die \gesperrt{lineare homogene Congruenzgruppe} genannt, die wir mit $S$ bezeichnen wollen. Ferner ist in (6) die $k$ fach cyklische Gruppe $Q$ enthalten, die aus den Substitutionen

\weq{8}{z_1' = z_1 + \alpha_1, \quad z_2' = z_2 + \alpha_2, \dots, \quad z_k' = z_k + \alpha_k}

besteht.

Bedeutet nun $\sigma$ ein Element aus $S$ und $\gamma$ ein Element aus $Q$, so ergiebt sich aus der Compositionsregel §.\,77, (7) ohne Schwierigkeit, dass sich jede Substitution $\pi$ aus der Gruppe $R$ \gesperrt{auf eine und nur auf eine Weise} in jeder der beiden Formen

\weq{9}{\pi = \sigma\gamma = \gamma'\sigma}

darstellen lässt, worin, wenn $\sigma, \gamma$ durch (7), (8) bestimmt sind, $\gamma'$ die Substitution

\weq{10}{z_1' = z_1 + \alpha_1', \quad z_2' = z_2 + \alpha_2', \dots, \quad z_k' = z_k + \alpha_k'}

bedeutet, wenn die $\alpha_i'$ aus den $\alpha_i$ durch die Substitution $\sigma$ abgeleitet sind.

$Q$ ist ein Normaltheiler von $R$; denn ist $\gamma_1$ ein beliebiges Element aus $Q$, so ist bei mehrmaliger Anwendung von (9):

\[
\sigma\gamma\gamma_1 = \gamma'\gamma_1'\sigma = \gamma'\gamma_1'\sigma\gamma^{-1}\gamma = \gamma'\gamma_1'\gamma'^{-1}\sigma\gamma,
\]

also, wenn

\[
\sigma\gamma = \pi, \quad \gamma'\gamma_1'\gamma'^{-1} = \gamma_1''
\]

gesetzt wird,

\[
\pi\gamma_1 = \gamma_1''\pi,
\]
%CONT>>

\begin{center}
\rule{4cm}{0.4pt}
\end{center}

$^{1)}$ \gesperrt{C. Jordan}, Liouville's Journ. de Math\'em. S\'er. (2), Bd. XIII.


% ===== scan leaf 325, printed 304 =====
% !PRINTED 304
\weberpage{304}
%<<CONT
und folglich

\weq{11}{\pi^{-1}Q\pi = Q.}

Wenn nun $\pi$ irgend einen Theiler $R_1$ von $R$ durchläuft, der seinerseits die Gruppe $Q$ enthält, so ist $Q$ Normaltheiler von $R_1$ und die in (9) vorkommende Substitution $\sigma$ muss einen Theiler $S_1$ von $S$ durchlaufen. Man kann setzen:

\[
R_1 = S_1 Q = Q S_1,
\]

und da die in der Form (9) enthaltenen Substitutionen alle von einander verschieden sind, so ist der Grad von $R_1$ gleich dem Producte der Grade von $S_1$ und $Q$. Ist $R_2$ ein Normaltheiler von $R_1$, der gleichfalls noch $Q$ (als Normaltheiler) enthält, so ist ebenso

\[
R_2 = S_2 Q,
\]

und $S_2$ ist ein Normaltheiler von $S_1$.

Denn nach Voraussetzung ist für jedes Element $\pi_1$ aus $R_1$

\[
\pi_1^{-1}S_2 Q \pi_1 = S_2 Q,
\]

also nach (11):

\[
\pi_1^{-1}S_2 Q \pi_1 = \pi_1^{-1}S_2\pi_1 Q = S_2 Q,
\]

und folglich

\[
\pi_1^{-1}S_2\pi_1 = S_2,
\]

und da dies auch für $\pi_1 = \sigma_1$ gilt, wo $\sigma_1$ in $S_1$ enthalten ist, so ist $S_2$ Normaltheiler von $S_1$. Umgekehrt ist, wenn $S_2$ ein Normaltheiler von $S_1$ ist, $S_2 Q$ Normaltheiler von $S_1 Q$. Denn wenn $S_2\sigma_1 = \sigma_1 S_2$ ist, so folgt

\[
\sigma_1 S_2 Q = S_2\sigma_1 Q = S_2 Q\sigma_1,
\]

und daraus

\[
\gamma\sigma_1 S_2 Q = \gamma S_2 Q\sigma_1 = S_2\gamma' Q\sigma_1 = S_2 Q\sigma_1 = S_2 Q\gamma^{-1}\gamma\sigma_1 = S_2 Q\gamma\sigma_1
\]

(weil $\gamma' Q = Q\gamma^{-1} = Q$ ist); d. h.

\[
\pi_1 S_2 Q = S_2 Q\pi_1,
\]

w. z. b. w. Hieraus aber ergiebt sich Folgendes:

\gesperrt{II. Ist $P$ die Galois'sche Gruppe einer primitiven irreduciblen metacyklischen Gleichung vom Grade $p^k$, so ist $P$ in der Form darstellbar:}

\[
P = TQ,
\]

\gesperrt{worin $T$ eine in der Congruenzgruppe $S$ enthaltene metacyklische Gruppe ist.}


% ===== scan leaf 326, printed 305 =====
% !PRINTED 305
\weberpage{305}

Die Aufgabe der Auffindung aller dieser metacyklischen
Gleichungen ist also darauf zur\"uckgef\"uhrt, alle metacyklischen
Theiler der homogenen Congruenzgruppe $S$ zu finden.

Nehmen wir als Beispiel den Fall $p = 3$, $k = 2$, fragen
also nach den metacyklischen Gleichungen 9\textsuperscript{ten} Grades, so handelt
es sich nach diesem Satze um die Gruppe $S$ der nach dem
Modul 3 genommenen linearen Substitutionen

\weq{12}{\sigma = \begin{pmatrix} a, & b \\ c, & d \end{pmatrix},}

mit denen wir uns, unter etwas anderem Gesichtspunkte, im \S.\,66
besch\"aftigt haben.

Wir haben dort zun\"achst einen Theiler $E$ von $S$ betrachtet,
der aus allen Substitutionen $\sigma$ mit der Bedingung $a\,d - b\,c$
$\equiv 1$ (mod 3) besteht, und haben ausserdem zwei Substitutionen

\weq{13}{\begin{pmatrix} a, & b \\ c, & d \end{pmatrix}, \quad \begin{pmatrix} -a, & -b \\ -c, & -d \end{pmatrix}}

zu einem einzigen Elemente zusammengefasst. Die so specialisirte
Gruppe war vom Grade 12 und hatte einen Normaltheiler
vom Index 3:

\[ G = \begin{pmatrix} 1, & 0 \\ 0, & 1 \end{pmatrix}, \begin{pmatrix} 0, & 1 \\ -1, & 0 \end{pmatrix}, \begin{pmatrix} -1, & 1 \\ 1, & 1 \end{pmatrix}, \begin{pmatrix} 1, & 1 \\ 1, & -1 \end{pmatrix}, \]

aus dem $E$ so zusammengesetzt war:

\[ E = G + \begin{pmatrix} 1, & 0 \\ 1, & 1 \end{pmatrix} G + \begin{pmatrix} 1, & 0 \\ -1, & 1 \end{pmatrix} G. \]

In der Gruppe $S$ gelten aber die beiden Substitutionen (13)
als verschieden, und ausserdem m\"ussen noch die Substitutionen
$\sigma$, deren Determinante $a\,d - b\,c \equiv -1$ (mod 3) ist, dazu
genommen werden, wodurch sich der Grad der Gruppe auf 48
erh\"oht. $G$ erweitert sich durch die Aenderung der Vorzeichen
aller Elemente zu einer Gruppe 8\textsuperscript{ten} Grades und $E$ zu einer
Gruppe 24\textsuperscript{sten} Grades, von der das erweiterte $G$ Normaltheiler
ist. Hier ist nun die ganze Gruppe $S$ \gesperrt{metacyklisch}, wie man
aus folgender Zusammensetzung sieht, in der die erweiterte
Gruppe $G$ mit $S_2$ und die erweiterte Gruppe $E$ mit $S_1$ bezeichnet
 ist:

\[ S_4 = \begin{pmatrix} 1, & 0 \\ 0, & 1 \end{pmatrix}, \begin{pmatrix} -1, & 0 \\ 0, & -1 \end{pmatrix} \]
%CONT>>


% ===== scan leaf 327, printed 306 =====
% !PRINTED 306
\weberpage{306}
\pdfbookmark[2]{§ 79. Ternäre lineare Congruenzgruppe für den Modul 2}{bm92}
%<<CONT
\[
\begin{aligned}
S_3 &= S_4 + \begin{pmatrix} -1, & 1 \\ 1, & 1 \end{pmatrix} S_4 \\
S_2 &= S_3 + \begin{pmatrix} 0, & 1 \\ -1, & 0 \end{pmatrix} S_3 \\
S_1 &= S_2 + \begin{pmatrix} 1, & 0 \\ 1, & 1 \end{pmatrix} S_2 + \begin{pmatrix} 1, & 0 \\ -1, & 1 \end{pmatrix} S_2 \\
S &= S_1 + \begin{pmatrix} -1, & 0 \\ 0, & 1 \end{pmatrix} S_1.
\end{aligned}
\]

Wir haben hiernach die Compositions- und Indexreihe von $S$

\[
\begin{gathered}
S,\ S_1,\ S_2,\ S_3,\ S_4, \\
2 \quad 3 \quad 2 \quad 2
\end{gathered}
\]

\gesperrt{und wir kommen also hier zu dem Ergebniss, dass alle Gleichungen
 9\textsuperscript{ten} Grades mit linearer Congruenzgruppe
metacyklisch sind.}

\begin{center}
\S.\,79.

\gesperrt{Tern\"are lineare Congruenzgruppe f\"ur den Modul 2.}
\end{center}

Mannigfache interessante Beziehungen ergeben sich, wenn
wir die tern\"are lineare Congruenzgruppe $R$ f\"ur den Modul 2
betrachten, die als transitive Permutationsgruppe von acht Elementen
 aufgefasst werden kann, und die Gruppe der metacyklischen
 Gleichungen 8\textsuperscript{ten} Grades enthalten muss. Nach dem
oben Bewiesenen kommt es vor Allem darauf an, die homogene
Congruenzgruppe $S$ zu betrachten, die aus den Elementen

\weq{1}{\sigma = \begin{pmatrix} a, & b, & c \\ a_1, & b_1, & c_1 \\ a_2, & b_2, & c_2 \end{pmatrix}}

besteht, worin die Ziffern $a, b, c, \ldots$ nach dem Modul 2, also $= 0$
oder $= 1$, anzunehmen sind, und die nach der im \S.\,37 gegebenen
Vorschrift componirt werden. Es kommen dabei nur solche
Systeme der Zahlen $a, b, c, \ldots$ in Betracht, deren Determinante
$= 1$, d.\,h. ungerade ist.

Um den Grad der Gruppe $S$ zu ermitteln, beachte man, dass
die erste Zeile von $\sigma$ sieben verschiedene Formen haben kann:

\[
(1,0,0),\ (0,1,0),\ (0,0,1),\ (0,1,1),\ (1,0,1),\ (1,1,0),\ (1,1,1).
\]


% ===== scan leaf 328, printed 307 =====
% !PRINTED 307
\weberpage{307}

F\"ur die erste Annahme $(a, b, c) = (1, 0, 0)$ wird die Determinante
 von $\sigma$ gleich $b_1 c_2 - c_1 b_2$, was auf sechs verschiedene
Arten $= 1$ werden kann, n\"amlich:

\[
\begin{pmatrix} b_1, & c_1 \\ b_2, & c_2 \end{pmatrix} = \begin{pmatrix} 1, & 0 \\ 0, & 1 \end{pmatrix}, \begin{pmatrix} 1, & 0 \\ 1, & 1 \end{pmatrix}, \begin{pmatrix} 0, & 1 \\ 1, & 0 \end{pmatrix}, \begin{pmatrix} 0, & 1 \\ 1, & 1 \end{pmatrix}, \begin{pmatrix} 1, & 1 \\ 1, & 0 \end{pmatrix}, \begin{pmatrix} 1, & 1 \\ 0, & 1 \end{pmatrix}.
\]

Dann k\"onnen noch $a_1, a_2$ auf vier verschiedene Arten angenommen
 werden. Also hat man, wenn die erste Zeile $(1, 0, 0)$
ist, 24 verschiedene Formen von $\sigma$. Dieselbe Zahl ergiebt sich
f\"ur die Annahme der ersten Zeile in den Formen $(0, 1, 0), (0, 0, 1)$.
Das Gleiche erh\"alt man aber auch f\"ur die anderen Annahmen
\"uber die erste Zeile. Denn ist diese z.\,B. $(0, 1, 1)$, so muss

\[
a_1 (b_2 - c_2) - a_2 (b_1 - c_1) \equiv 1 \text{ (mod 2)}
\]

sein, was wieder sechs M\"oglichkeiten f\"ur $a_1, b_1 - c_1, a_2, b_2 - c_2$
ergiebt, und da man $c_1, c_2$ auf vier Arten annehmen kann, so
erh\"alt man wieder 24 M\"oglichkeiten. Endlich muss, wenn die
erste Zeile $(1, 1, 1)$ ist,

\[
(a_1 - c_1)(b_2 - c_2) - (a_2 - c_2)(b_1 - c_1) \equiv 1 \text{ (mod 2)}
\]

sein, was zu derselben Zahl f\"uhrt. Die Gesammtzahl aller verschiedenen
 Substitutionen $\sigma$, d.\,h. der Grad von $S$ ist also 168.
Die Zahl 168 als Gradzahl einer Gruppe ist uns schon einmal
begegnet, n\"amlich bei der Gruppe $L_7$ (\S.\,66, 72), und es liegt
daher die Vermuthung nahe, dass die Gruppen $S$ und $L_7$
isomorph sein m\"ochten. Wenn sich diese Vermuthung best\"atigt,
so w\"urde die Untersuchung der Gruppe $S$ dadurch wesentlich
 erleichtert sein, dass wir die Divisoren der
Gruppe $L_7$ schon kennen.

Diese Vermuthung zu pr\"ufen, dient aber das Theorem I.
in \S.\,72.

Wir suchen zu diesem Zweck erzeugende Elemente $\chi, \omega, \Theta, \tau$
der Gruppe $S$ so auszuw\"ahlen, dass sie den charakteristischen
Bedingungen des Theorems I., \S.\,72 gen\"ugen. Dies kann auf
mehrfache Weise geschehen. Es fehlt freilich an einem allgemeinen
 Verfahren dazu, gelingt aber leicht durch einige Versuche.

Man w\"ahlt zun\"achst ein Element 7\textsuperscript{ten} Grades beliebig aus
und bildet daraus die Periode, etwa
%CONT>>


% ===== scan leaf 329, printed 308 =====
% !PRINTED 308
\weberpage{308}
\pdfbookmark[2]{§ 80. Resolventen der Gleichung 8ten Grades}{bm93}
%<<CONT
\weq{2}{
\begin{aligned}
\tau &= \begin{pmatrix} 1, & 0, & 1 \\ 1, & 0, & 0 \\ 0, & 1, & 0 \end{pmatrix}, &
\tau^2 &= \begin{pmatrix} 1, & 1, & 1 \\ 1, & 0, & 1 \\ 1, & 0, & 0 \end{pmatrix}, &
\tau^3 &= \begin{pmatrix} 0, & 1, & 1 \\ 1, & 1, & 1 \\ 1, & 0, & 1 \end{pmatrix}. \\
\tau^4 &= \begin{pmatrix} 1, & 1, & 0 \\ 0, & 1, & 1 \\ 1, & 1, & 1 \end{pmatrix}, &
\tau^5 &= \begin{pmatrix} 0, & 0, & 1 \\ 1, & 1, & 0 \\ 0, & 1, & 1 \end{pmatrix}, &
\tau^6 &= \begin{pmatrix} 0, & 1, & 0 \\ 0, & 0, & 1 \\ 1, & 1, & 0 \end{pmatrix}
\end{aligned}
}

Hierauf sucht man unter den Elementen 2\textsuperscript{ten} Grades, deren
es 21 giebt, ein passendes f\"ur $\omega$ aus; es eignet sich z.\,B. dieses

\weq{3}{\omega = \begin{pmatrix} 1, & 0, & 0 \\ 0, & 1, & 0 \\ 0, & 1, & 1 \end{pmatrix};}

und dann sind nach den Formeln \S.\,72, (17) die Elemente $\Theta$
und $\chi$ bestimmt:

\weq{4}{\Theta = \begin{pmatrix} 1, & 1, & 1 \\ 0, & 1, & 0 \\ 0, & 1, & 1 \end{pmatrix}, \quad \Theta^2 = \begin{pmatrix} 1, & 1, & 0 \\ 0, & 1, & 0 \\ 0, & 0, & 1 \end{pmatrix}, \quad \Theta^3 = \begin{pmatrix} 1, & 0, & 1 \\ 0, & 1, & 0 \\ 0, & 1, & 1 \end{pmatrix}}

\weq{5}{\chi = \begin{pmatrix} 1, & 0, & 0 \\ 0, & 0, & 1 \\ 0, & 1, & 1 \end{pmatrix}, \quad \chi^2 = \begin{pmatrix} 1, & 0, & 0 \\ 0, & 1, & 1 \\ 0, & 1, & 0 \end{pmatrix}.}

Hierauf ist es Sache einer einfachen Rechnung, die Formeln
des Theorems I., \S.\,72 zu best\"atigen, wodurch die Uebereinstimmung
 der Gruppen $S$ und $L_7$ nachgewiesen ist.

Es folgt daraus, dass die Gruppe $S$ ebenso wie $L_7$ \gesperrt{einfach}
 ist.

\begin{center}
\S.\,80.

\gesperrt{Resolventen der Gleichung 8\textsuperscript{ten} Grades.}
\end{center}

Um nun die Anwendung auf die Theorie der Gleichungen
8\textsuperscript{ten} Grades zu machen, bezeichnen wir wie fr\"uher mit $Q$ die
cyklische Gruppe 8\textsuperscript{ten} Grades: % [sic] the displayed translation group is elementary abelian of order 8, not cyclic.

\[
\begin{pmatrix} z_1, & z_2, & z_3 \\ z_1 + h_1, & z_2 + h_2, & z_3 + h_3 \end{pmatrix} \text{ (mod 2)},
\]

oder k\"urzer

\[
(z_i,\ z_i + h_i),
\]

und mit $S$ die im \S.\,79 betrachtete tern\"are Congruenzgruppe,
und setzen

\[
R = S\,Q,
\]
%CONT>>


% ===== scan leaf 330, printed 309 =====
% !PRINTED 309
\weberpage{309}
%% [page furniture, not body text -- commented out 2026-08-10] Page furniture (running
%% head, folio, signature) is excluded from this edition by standing policy.
%% \noindent \S.\,80. \hfill Gleichungen 8ten Grades.

%<<CONT
so dass $R$ die gesammte ternäre lineare Congruenzgruppe für den Modul 2 ist.

Wir betrachten ein System von acht unabhängigen Veränderlichen $X_{z_1,z_2,z_3}$, worin die Indices $z_1, z_2, z_3$ nach dem Modul 2 zu nehmen sind.

Hieraus bilden wir die sieben Functionen (\gesperrt{Lagrange'sche} Resolventen, Bd.\,I, \S.\,164).
\weq{1}{\overset{z}{\Sigma} (-1)^{\Sigma z\xi} X_{z_1,z_2,z_3} = \Psi_{\xi_1,\xi_2,\xi_3},}
worin zur Abkürzung
\[
\Sigma z\xi = z_1\xi_1 + z_2\xi_2 + z_3\xi_3
\]
gesetzt ist, und $\xi_1, \xi_2, \xi_3$ je ein volles Restsystem nach dem Modul 2, jedoch mit Ausschluss der Combination $0, 0, 0$ durchlaufen.

Wenden wir auf die Indices $z_i$ der Grössen $X$ die Substitutionen der Gruppe $Q$ an, so erleiden die $X$ die Permutationen einer Gruppe, die wir gleichfalls mit $Q$ bezeichnen. Die Aenderungen der $\Psi$, die diesen Permutationen entsprechen, ergeben sich aus (1), wenn wir auf die Indices der $X$ die Substitution $(z_i, z_i+h)$ anwenden:
\weq{2}{\Sigma (-1)^{\Sigma z\xi} X_{z_1+h_1, z_2+h_2, z_3+h_3} = \overset{z}{\Sigma} (-1)^{\Sigma(z+h)\xi} X_{z_1,z_2,z_3} = (-1)^{\Sigma h\xi} \Psi_{\xi_1,\xi_2,\xi_3}.}

\gesperrt{1. Die Functionen $\Psi$ ändern also durch die Permutationen von $Q$ nur ihr Zeichen, und die Quadrate der $\Psi$ bleiben durch $Q$ ungeändert.}

Es ist ferner der Einfluss einer Substitution $\sigma$ aus $S$ auf die Functionen $\Psi$ festzustellen.

Bezeichnen wir zu diesem Zwecke mit $\varrho$ die \gesperrt{transponirte Substitution} zu $\sigma$ und setzen
\weq{3}{z' = \sigma(z), \quad \xi = \varrho(\xi'), \quad \xi' = \varrho^{-1}(\xi),}
so ist (\S.\,37, 9.):
\weq{4}{\Sigma \xi z = \Sigma \xi' z'.}

Machen wir die Substitutionen $\sigma$ der Gruppe $S$ in den Indices von $X$, so erhalten wir eine mit $S$ zu bezeichnende Permutationsgruppe der $X$, und aus (1) ergiebt sich
\[
\overset{z}{\Sigma} (-1)^{\Sigma z\xi} X_{z_1', z_2', z_3'} = \overset{z}{\Sigma} (-1)^{\Sigma \xi' z'} X_{z_1', z_2', z_3'},
\]
%CONT>>


% ===== scan leaf 331, printed 310 =====
% !PRINTED 310
\weberpage{310}
\pdfbookmark[2]{§ 81. Die Tripelsysteme der Resolventen}{bm94}
%% [page furniture, not body text -- commented out 2026-08-10] Page furniture (running
%% head, folio, signature) is excluded from this edition by standing policy.
%% \noindent Zehnter Abschnitt. \hfill \S.\,81.

%<<CONT
und da das System der $z_i'$ dieselbe Werthreihe durchläuft, wie das System der $z_i$, so ist diese Summe gleich $\Psi_{\xi_1',\xi_2',\xi_3'}$, und wir erhalten:

\gesperrt{2. Die Anwendung der Permutation $\sigma$ auf die $X$ in den Functionen $\Psi$ ist gleichbedeutend mit der Anwendung von $\varrho^{-1}$ auf die Indices $\xi_1, \xi_2, \xi_3$.}

\begin{center}
\S.\,81.\\
\gesperrt{Die Tripelsysteme der Resolventen.}
\end{center}

Wir bezeichnen jetzt die Resolvente $\Psi_{\xi_1,\xi_2,\xi_3}$ kürzer durch $\Psi^\xi$ und bemerken, dass nach der Formel (2), \S.\,80 ausser den Quadraten dieser Functionen auch noch gewisse Producte die Permutationen der Gruppe $Q$ gestatten. Zu diesen gehört zunächst das Product aller Grössen $\Psi$, das wir mit $A$ bezeichnen wollen:
\weq{1}{A = \overset{\xi}{\prod} \Psi_\xi,}
sodann aber die Producte von dreien $\Psi_\xi \Psi_\eta \Psi_\zeta$, und folglich auch ihre reciproken Werthe
\weq{2}{v = \frac{1}{\Psi_\xi \Psi_\eta \Psi_\zeta},}
wenn die Indices den Bedingungen genügen:
\weq{3}{\begin{aligned}
\xi_1 + \eta_1 + \zeta_1 &\equiv 0\\
\xi_2 + \eta_2 + \zeta_2 &\equiv 0 \quad (\mathrm{mod}\ 2).\\
\xi_3 + \eta_3 + \zeta_3 &\equiv 0
\end{aligned}}
Ein solches System von drei Functionen $\Psi$ nennen wir ein \gesperrt{Tripel}. Ebenso wollen wir drei Indexsysteme $(\xi), (\eta), (\zeta)$, die den Bedingungen (3) genügen, ein Tripel nennen.

Es giebt im Ganzen sieben und nicht mehr solcher Tripel; denn unter den drei Systemen $(\xi), (\eta), (\zeta)$ können nicht zwei einander gleich sein, und durch zwei ist das dritte eindeutig bestimmt. Man kann also für $(\xi), (\eta)$ jedes Paar aus den sieben möglichen Systemen $(\xi)$ nehmen, erhält aber dann jedes Tripel sechsmal. Die Gesammtanzahl ist also $7 . 6 : 6 = 7$.

Um die einzelnen Tripelsysteme zu charakterisiren, führen wir drei nach dem Modul 2 zu nehmende Zahlen $\alpha_1, \alpha_2, \alpha_3$ ein, die nicht alle drei verschwinden, und bemerken, dass die Congruenz
\weq{4}{\alpha_1 \xi_1 + \alpha_2 \xi_2 + \alpha_3 \xi_3 \equiv 0 \ (\mathrm{mod}\ 2)}
%CONT>>


% ===== scan leaf 332, printed 311 =====
% !PRINTED 311
\weberpage{311}
%% [page furniture, not body text -- commented out 2026-08-10] Page furniture (running
%% head, folio, signature) is excluded from this edition by standing policy.
%% \noindent \S.\,81. \hfill Tripelsysteme.

%<<CONT
für drei und nur für drei Systeme der Unbekannten $\xi_1, \xi_2, \xi_3$ befriedigt werden, und dass diese drei ein Tripel bilden, so dass das Tripel durch die drei Congruenzen
\weq{5}{\Sigma \alpha\xi \equiv 0, \quad \Sigma \alpha\eta \equiv 0, \quad \Sigma \alpha\zeta \equiv 0 \quad (\mathrm{mod}\ 2)}
bestimmt ist. Man erkennt dies ohne weitläufige allgemeine Betrachtungen, wenn man die Fälle, in denen nur ein $\alpha$ oder zwei oder alle drei $\alpha \equiv 1$ (mod 2) sind, einzeln betrachtet.

Da es sieben verschiedene Zahlensysteme $\alpha_1, \alpha_2, \alpha_3$ giebt, so ist durch die Relationen (5) aus jedem solchen Systeme der $\alpha$ ein Tripel völlig bestimmt, und wir können die Function $v$ eindeutig durch $v_{\alpha_1,\alpha_2,\alpha_3}$ oder kürzer durch $v_\alpha$ bezeichnen:
\weq{6}{v_\alpha = \frac{1}{\Psi_\xi \Psi_\eta \Psi_\zeta}.}

\gesperrt{3. Die Functionen $v_\alpha$, als Functionen der $X$ aufgefasst, gestatten die Permutationen der Gruppe $Q$.}

Wenn wir auf die $z$ eine Substitution $\sigma$ aus $S$ anwenden, so erleiden nach 2.\ (\S.\,80) die $\xi, \eta, \zeta$ gleichzeitig die Substitution $\varrho^{-1}$; daraus aber ergiebt sich nach (5), dass die $\alpha$ die Substitution $\sigma$ erleiden, und wir haben also:

\gesperrt{4. Die Anwendung der Substitution $\sigma$ aus $S$ auf die Indices $z$ von $X$ hat die Anwendung derselben Substitution auf die Indices $\alpha$ von $v_\alpha$ zur Folge.}

Setzen wir in den Summen (5), in denen $(\xi), (\eta), (\zeta)$ das zu $(\alpha)$ gehörige Tripel ist, für $(\alpha)$ ein anderes System $(\beta)$, so ist wegen (3)
\[
\Sigma \beta\xi + \Sigma \beta\eta + \Sigma \beta\zeta \equiv 0.
\]

Ist nun $\beta$ von $\alpha$ verschieden, so können diese drei Summen nicht alle drei congruent mit 0 sein, weil sonst zwei verschiedene Systeme $(\alpha), (\beta)$ zu demselben Tripel führen würden; \gesperrt{folglich müssen von diesen Summen zwei ungerade und eine gerade sein.} Wir haben daher den Satz:

\gesperrt{5. Durchläuft $(\xi)$ das zu $(\alpha)$ gehörige Tripel, und ist $(\beta)$ ein von $(\alpha)$ verschiedenes System, so sind unter den Summen $\Sigma \beta\xi$ zwei ungerade und eine gerade.}


% ===== scan leaf 333, printed 312 =====
% !PRINTED 312
\weberpage{312}
%% [page furniture, not body text -- commented out 2026-08-10] Page furniture (running
%% head, folio, signature) is excluded from this edition by standing policy.
%% \noindent Zehnter Abschnitt. \hfill \S.\,81.

Wenn wir in der Summe $\Sigma \alpha\xi$ für das System $(\xi)$ alle sieben möglichen Annahmen machen, so erhält man dreimal den Werth 0 und folglich viermal den Werth 1.

Die Congruenz
\weq{7}{\Sigma \alpha\xi \equiv 1 \ (\mathrm{mod}\ 2)}
hat also vier und nur vier verschiedene Lösungen für $\xi$. Die Gesammtheit dieser vier Lösungen wollen wir ein \gesperrt{Quadrupel} nennen. Jedes Tripel bestimmt eindeutig ein Quadrupel und umgekehrt, und es giebt also auch sieben verschiedene Quadrupel, die nach (7) durch das Zahlensystem $(\alpha)$ vollständig bestimmt sind. Wir beweisen nun den Satz:

\gesperrt{6. Durchläuft $(\xi)$ das zu $(\alpha)$ gehörige Quadrupel und ist $(\beta)$ ein von $(\alpha)$ verschiedenes System, so sind unter den vier Summen $\Sigma \beta\xi$ zwei gerade und zwei ungerade.}

Denn wenn $(\xi)$ die Gesammtheit aller sieben Systeme durchläuft, so finden sich unter den Summen $\Sigma \beta\xi$ drei gerade und vier ungerade; von diesen liefert das zu $(\alpha)$ gehörige Tripel zwei ungerade und eine gerade; es bleiben also für das Quadrupel zwei gerade und zwei ungerade.

Wir ändern nun die Bezeichnung bei den Quadrupeln und nehmen für das System $(\alpha)$ ein anderes Zeichen $(h_1, h_2, h_3) = (h)$. Das zu $(h)$ gehörige Quadrupel bezeichnen wir mit $(\alpha), (\beta), (\gamma), (\delta)$, so dass die vier Congruenzen:
\weq{8}{\Sigma h\alpha \equiv 1, \quad \Sigma h\beta \equiv 1, \quad \Sigma h\gamma \equiv 1, \quad \Sigma h\delta \equiv 1 \quad (\mathrm{mod}\ 2)}
bestehen.

Im Nenner des Productes der vier Functionen
\[
v_\alpha, \quad v_\beta, \quad v_\gamma, \quad v_\delta
\]
kommen nach (6) im Ganzen 12 Factoren $\Psi_\xi$ vor, und zwar kommt jeder Factor $\Psi_\xi$ so oft darunter vor, als die Anzahl der geraden Zahlen unter den Summen
\[
\Sigma \xi\alpha, \quad \Sigma \xi\beta, \quad \Sigma \xi\gamma, \quad \Sigma \xi\delta
\]
beträgt, d.\,h.\ $\Psi_h$ kommt gar nicht vor, während jedes andere $\Psi_\xi$ zweimal vorkommt (nach 6.).

Wenn wir also noch die Relation (1) benutzen, so ergiebt sich:
\weq{9}{\Psi_h^2 = A^2 v_\alpha v_\beta v_\gamma v_\delta.}


% ===== scan leaf 334, printed 313 =====
% !PRINTED 313
\weberpage{313}
\pdfbookmark[2]{§ 82. Anwendung auf Gleichungen 8ten Grades}{bm95}

Wir stellen für die Anwendung die sieben Tripel zusammen, aus denen man die Quadrupel als die Ergänzung zu dem vollen Systeme ablesen kann. In der ersten Columne steht das Zeichen ($\alpha$), zu dem das Tripel gehört:

\[
\begin{array}{c||ccc|}
1\,0\,0 & 0\,1\,0 & 0\,0\,1 & 0\,1\,1 \\
0\,1\,0 & 1\,0\,0 & 0\,0\,1 & 1\,0\,1 \\
0\,0\,1 & 1\,0\,0 & 0\,1\,0 & 1\,1\,0 \\
0\,1\,1 & 1\,0\,0 & 0\,1\,1 & 1\,1\,1 \\
1\,0\,1 & 0\,1\,0 & 1\,0\,1 & 1\,1\,1 \\
1\,1\,0 & 0\,0\,1 & 1\,1\,0 & 1\,1\,1 \\
1\,1\,1 & 0\,1\,1 & 1\,0\,1 & 1\,1\,0 \\
\end{array}
\]

Bezeichnen wir die Symbole ($\alpha$) in der Reihenfolge, in der sie in der ersten Columne dieser Tabelle stehen, mit $1, 2, 3, 4, 5, 6, 7$, so können wir die Tabelle für die Tripel und Quadrupel einfacher so darstellen:

\[
\begin{array}{c|ccc|cccc}
1 & 2 & 3 & 4 & 1 & 5 & 6 & 7 \\
2 & 1 & 3 & 5 & 2 & 4 & 6 & 7 \\
3 & 1 & 2 & 6 & 3 & 4 & 5 & 7 \\
4 & 1 & 4 & 7 & 2 & 3 & 5 & 6 \\
5 & 2 & 5 & 7 & 1 & 3 & 4 & 6 \\
6 & 3 & 6 & 7 & 1 & 2 & 4 & 5 \\
7 & 4 & 5 & 6 & 1 & 2 & 3 & 7 \\
\end{array}
\]

\begin{center}
\S.\,82.

\gesperrt{Anwendung auf Gleichungen $8^{\mathrm{ten}}$ Grades.}
\end{center}

In den abgeleiteten Formeln setzen wir nun für die Variablen $X_{z_1,z_2,z_3}$ die Wurzeln einer Gleichung $8^{\mathrm{ten}}$ Grades, von der wir annehmen, dass ihre \gesperrt{Galois}'sche Gruppe $P = TQ$ in dem festgesetzten Rationalitätsbereiche $\Omega$ in der linearen Congruenzgruppe $R = SQ$ enthalten sei, so dass $T$ ein Theiler von $S$ ist. Alle Functionen, die die Permutationen dieser Gruppe gestatten, sind rational.

Dazu gehört zunächst die Summe der $X$:

\weq{1}{\sum^{z} X_{z_1, z_2, z_3} = B,}
%CONT>>


% ===== scan leaf 335, printed 314 =====
% !PRINTED 314
\weberpage{314}
\pdfbookmark[2]{§ 83. Metacyklische Gleichungen 8ten Grades}{bm96}
%<<CONT
ferner die durch (1) des vorigen Paragraphen definirte Grösse $A$; sodann aber auch die symmetrischen Functionen der sieben Grössen $\Psi^2$, und ferner, wenn wir nun der Einfachheit halber die Annahme hinzufügen, auf die wir noch zurückkommen, dass von den Grössen $\Psi$ keine verschwinde, die symmetrischen Functionen der Grössen $v$; aber nicht nur die symmetrischen Functionen der $v$, sondern alle Functionen der $v$, die die Permutationen der Gruppe $T$ gestatten.

Die sieben Grössen

\weq{2}{\begin{aligned}
v_{1,0,0} &= v_1, & v_{0,1,0} &= v_2, & v_{0,0,1} &= v_3, \\
v_{0,1,1} &= v_4, & v_{1,0,1} &= v_5, & v_{1,1,0} &= v_6, & v_{1,1,1} &= v_7
\end{aligned}}

sind alsdann die Wurzeln einer rationalen Gleichung $7^{\mathrm{ten}}$ Grades mit der Gruppe $T$, und durch Addition der Gleichung (1) und der Gleichungen [\S.\,80, (1)]:

\[
\Sigma (-1)^{\Sigma z\xi} X_{z_1, z_2, z_3} = \Psi_{\xi_1, \xi_2, \xi_3}
\]

ergiebt sich mit Rücksicht auf (9) des vorigen Paragraphen:

\weq{3}{8\,X_{0,0,0} = A\Bigl\{\begin{aligned}[t]
&\sqrt{v_1}\sqrt{v_5}\sqrt{v_6}\sqrt{v_7} + \sqrt{v_2}\sqrt{v_4}\sqrt{v_6}\sqrt{v_7} \\
&+ \sqrt{v_3}\sqrt{v_4}\sqrt{v_5}\sqrt{v_7} + \sqrt{v_2}\sqrt{v_3}\sqrt{v_5}\sqrt{v_6} \\
&+ \sqrt{v_1}\sqrt{v_3}\sqrt{v_4}\sqrt{v_6} + \sqrt{v_1}\sqrt{v_2}\sqrt{v_4}\sqrt{v_5} \\
&+ \sqrt{v_1}\sqrt{v_2}\sqrt{v_3}\sqrt{v_7}
\end{aligned}\Bigr\} + B.}

Es giebt 15 Vorzeichenänderungen der sieben Quadratwurzeln $\sqrt{v}$, bei denen dieser Ausdruck ungeändert bleibt, nämlich, wenn alle sieben Vorzeichen gleichzeitig geändert werden oder wenn die Vorzeichen eines Tripels oder eines Quadrupels geändert werden. \gesperrt{Folglich erhält der Ausdruck (3) nur acht verschiedene Werthe, wie man auch die sieben darin vorkommenden Quadratwurzeln bestimmen mag.}

\begin{center}
\S.\,83.

\gesperrt{Metacyklische Gleichungen $8^{\mathrm{ten}}$ Grades.}
\end{center}

Um nun nach diesen Vorbereitungen die primitiven metacyklischen Gleichungen $8^{\mathrm{ten}}$ Grades zu finden, haben wir zunächst die metacyklischen Theiler der Gruppe $S$ zu ermitteln. Die Gruppe $S$ selbst ist, wie schon erwähnt, nicht metacyklisch. Ihre
%CONT>>


% ===== scan leaf 336, printed 315 =====
% !PRINTED 315
\weberpage{315}
%<<CONT
Theiler haben wir schon betrachtet (\S.\,70 bis 72), und alle echten Theiler von $S$ sind metacyklisch. Darunter sind Theiler vom $21^{\mathrm{sten}}$ und $7^{\mathrm{ten}}$ Grade, die wir mit $T_{21}$, $T_7$ bezeichnen wollen, ferner Theiler vom $24^{\mathrm{sten}}$ Grade, und noch andere Theiler, deren Gradzahlen in 24 aufgehen, die wir hier alle in dem Zeichen $T_{24}$ zusammenfassen wollen.

Diese Gruppen $T_{24}$ können nicht die sieben Grössen $(\xi) = (\xi_1, \xi_2, \xi_3)$ transitiv mit einander verbinden, weil der Grad einer transitiven Permutationsgruppe von sieben Elementen durch 7 theilbar sein muss (Bd.\,I, \S.\,154, 7.). Wir wollen nachweisen, dass die den Gruppen $T_{24}\,Q$ entsprechenden Permutationsgruppen der $X$ imprimitiv sind. % [sic] Full stop before the syntactic continuation with \glqq und dass\grqq{}.
und dass diese Gruppen daher von unserer Betrachtung ausscheiden\footnote{Aus der Literatur über die Gleichungen $8^{\mathrm{ten}}$ Grades ist zu erwähnen: \gesperrt{Nöther}, Mathem. Annalen, Bd. XV. \gesperrt{Wilshaus}, Ueber die algebraische Auflösbarkeit der Gleichungen $8^{\mathrm{ten}}$ Grades. Dissertation. Marburg 1888.}.

Da alle Permutationen $\sigma$, die zwei der sieben Elemente $(\xi)$ ungeändert lassen oder nur unter einander permutiren, ein drittes Element, nämlich die Ergänzung zum Tripel, ungeändert lassen, so sind zwei Fälle möglich:

\begin{itemize}
\item[1)] die Gruppe $T_{24}$ lässt ein Element in Ruhe, oder
\item[2)] die sieben Elemente werden in zwei Theile von drei und vier Elementen zerlegt, von denen jeder Theil nur unter sich permutirt wird.
\end{itemize}

Im letzten Falle können wir die drei Elemente als einem Tripel, und demnach die vier Elemente als einem Quadrupel angehörig betrachten. Denn werden drei Elemente, die nicht einem Tripel angehören, unter sich permutirt, so werden auch die drei Ergänzungen je zweier Elemente zum Tripel unter sich permutirt, und das eine übrig bleibende Element bleibt ungeändert. Wir kommen also auf den Fall 1) zurück.

Wir fügen nun zu den sieben Systemen $(\xi)$ noch das achte $(0,0,0) = (0)$ hinzu und bezeichnen die acht Grössen $X_{\xi_1, \xi_2, \xi_3}$ mit

\weq{1}{0, 1, 2, 3, 4, 5, 6, 7.}

Im Falle 1) bleiben dann durch die Substitutionen von $T_{24}$ zwei dieser acht Elemente, etwa 0, 1, ungeändert. Es giebt ferner in $Q$ eine und nur eine Substitution $\gamma_0$, durch die 0 mit 1 vertauscht wird.


% ===== scan leaf 337, printed 316 =====
% !PRINTED 316
\weberpage{316}
Die beiden Substitutionen $1, \gamma_0$ bilden eine Gruppe $Q_0$. Irgend eine Substitution $\gamma_1$ von $Q$, die nicht in $Q_0$ vorkommt, führt 0, 1 in zwei davon verschiedene Elemente, etwa 2, 3, über, und durch die zwei Substitutionen $Q_0\gamma_1$ geht 0, 1 in 2, 3 oder in 3, 2 über. Wenn ferner durch $\gamma_2$ das Paar 0, 1 in ein drittes Paar 4, 5 übergeht, so ist 4, 5 sowohl von 0, 1 als von 2, 3 verschieden, und so können wir $Q$ in die Nebengruppen zerlegen:

\[
Q = Q_0 + Q_0\gamma_1 + Q_0\gamma_2 + Q_0\gamma_3,
\]

wodurch die acht Elemente (1) in vier Paare

\weq{2}{0, 1;\quad 2, 3;\quad 4, 5;\quad 6, 7}

zerlegt sind, so dass durch die beiden Substitutionen einer der Nebengruppen $Q_0\gamma_i$ das Paar 0, 1 in eines dieser vier Paare übergeht.

Ist nun $k$ irgend eine Substitution aus der Gruppe

\[
P = T_{24}\,Q,
\]

so lassen sich die Elemente $\tau_0, \tau_1, \tau_2, \tau_3$ aus $T_{24}$ und $\gamma_0', \gamma_1', \gamma_2', \gamma_3'$ aus $Q$ so bestimmen, dass

\[
k = \tau_0\gamma_0', \quad \gamma_1 k = \tau_1\gamma_1', \quad \gamma_2 k = \tau_2\gamma_2', \quad \gamma_3 k = \tau_3\gamma_3'
\]

oder

\[
k = \tau_0\gamma_0' = \gamma_1^{-1}\tau_1\gamma_1' = \gamma_2^{-1}\tau_2\gamma_2' = \gamma_3^{-1}\tau_3\gamma_3',
\]

und diese Darstellung zeigt, dass durch $k$ irgend eines der Paare (2) in ein Paar desselben Systems übergeführt wird; denn z.\,B. durch $\gamma_1^{-1}$ wird 2, 3 in 0, 1 übergeführt, durch $\tau_1$ bleibt 0, 1 ungeändert und durch $\gamma_1'$, was zu $Q$ gehört, wird 0, 1 in eines der Paare (2) übergeführt.

Die Gruppe $T_{24}\,Q$ ist also imprimitiv, und zwar so, dass wir vier Systeme der Imprimitivität haben. Die Wurzeln 0, 1 genügen einer quadratischen Gleichung, deren Coëfficienten von den Wurzeln einer biquadratischen Gleichung abhängen.

Gehen wir nun zu dem Falle 2) über, in dem die Elemente je eines Tripels 1, 2, 3 und eines Quadrupels 4, 5, 6, 7 durch $T_{24}$ transitiv unter einander verbunden sind. Die acht Grössen (1) zerfallen hier in zwei Quadrupel:

\weq{3}{0, 1, 2, 3;\quad 4, 5, 6, 7,}

und durch die Substitutionen von $Q$ werden die Grössen eines dieser Quadrupel entweder nur unter sich vertauscht, oder sie werden in die Grössen des anderen Quadrupels übergeführt.


% ===== scan leaf 338, printed 317 =====
% !PRINTED 317
\weberpage{317}

Um dies einzusehen, bezeichnen wir in der früheren Bezeichnungsweise 1, 2, 3 mit $(\xi)$, $(\eta)$, $(\zeta)$; da diese drei Grössen ein Tripel bilden, so giebt es [nach \S.\,81, (5)] ein System $(\alpha)$, so dass
\weq{4}{\Sigma\alpha\xi \equiv 0, \quad \Sigma\alpha\eta \equiv 0, \quad \Sigma\alpha\zeta \equiv 0 \ (\mathrm{mod}\ 2)}
ist. Wenden wir nun eine Substitution aus $Q$ an:
\[
\begin{pmatrix} \xi_1, & \xi_2, & \xi_3 \\ \xi_1', & \xi_2', & \xi_3' \end{pmatrix},
\]
worin $\xi_i' = \xi_i + h_i$ sein mag, so folgt aus (4):
\[
\Sigma\alpha\xi' \equiv \Sigma\alpha\eta' \equiv \Sigma\alpha\zeta' \equiv \Sigma\alpha h,
\]
d.\,h. die $(\xi')$, $(\eta')$, $(\zeta')$ bilden das zu $(\alpha)$ gehörige Tripel, wenn $(h)$ selbst zu diesem Tripel gehört, oder im anderen Falle ist $(h)$, $(\xi')$, $(\eta')$, $(\zeta')$, das zu $(\alpha)$ gehörige Quadrupel.

Bezeichnet also nun wieder $k$ irgend eine Substitution aus $P$ und $\gamma$ eine Substitution aus $Q$, durch die das Quadrupel 0, 1, 2, 3 in 4, 5, 6, 7 übergeht, so können wir die Elemente $\tau'$, $\tau''$ aus $T_{24}$ und $\gamma'$, $\gamma''$ aus $Q$ so bestimmen, dass
\[
k = \tau'\gamma' = \gamma^{-1}\tau''\gamma'',
\]
wodurch, wie oben, bewiesen wird, dass die beiden Systeme (3) durch $k$ imprimitiv verbunden sind.

Nach Adjunction einer Quadratwurzel zum Rationalitätsbereiche sind dann die Grössen 0, 1, 2, 3 die Wurzeln einer biquadratischen Gleichung.

Durch diese Betrachtungen wird für primitive Gleichungen 8\textsuperscript{ten} Grades auch die Möglichkeit ausgeschlossen, dass eine der sieben Grössen $\Psi_{\xi_1,\xi_2,\xi_3}$ [\S.\,80, (1)] verschwinde. Diese Functionen sind nämlich rationale Functionen der Wurzeln $X$; wenn also eine von ihnen verschwindet, so kann auf die rationale Gleichung $\Psi = 0$ jede Permutation der Gruppe der Gleichung angewandt werden. Daraus ergiebt sich nach dem Theorem \S.\,80, 2., dass entweder alle sieben Functionen $\Psi$ verschwinden müssen, in welchem Falle die Wurzeln rational wären, oder dass die Substitutionen $\varrho^{-1}$, auf die $(\xi)$ angewandt, eine intransitive Gruppe bilden müssen. Der Grad der Gruppe könnte dann nicht durch 7 theilbar sein. Nun ist die Gruppe der $\varrho^{-1}$ isomorph mit der Gruppe $T$ der $\sigma$ (\S.\,37, 8.), und daher ist auch der Grad von $T$ nicht durch 7 theilbar, und folglich ist
%CONT>>


% ===== scan leaf 339, printed 318 =====
% !PRINTED 318
\weberpage{318}
%<<CONT
auch $T$ intransitiv. Daraus folgt, wie oben gezeigt, dass $TQ$
imprimitiv ist.

Als Gruppe für primitive metacyklische Gleichungen 8\textsuperscript{ten}
Grades bleiben also nur die beiden Typen
\[
T_7 Q, \quad T_{21} Q
\]
übrig. Im ersten Falle sind die in der Formel [\S.\,82, (3)] vorkommenden
Grössen $v$ die Wurzeln einer cyklischen Gleichung 7\textsuperscript{ten}
Grades mit der Gruppe $(z, z+b)$, im zweiten einer metacyklischen mit der
Gruppe $(z, az+b)$ (Bd. I, \S.\,180), worin $z, a, b$ nach dem Modul 7
genommen sind und $a$ die Werthe 1, 2, 4 durchläuft.

Um eine Zuordnung der $v_i$, wie sie in der Formel (3), \S.\,82 vorkommen,
zu den $v_z$, worin $z$ der in den Gruppen $(z, z+b)$ oder $(z, az+b)$
vorkommende Index ist, zu erhalten, können wir von einer beliebigen
Substitution siebenter Ordnung ausgehen, etwa von
\[
\tau = \begin{pmatrix} 1, & 0, & 1 \\ 1, & 0, & 0 \\ 0, & 1, & 0 \end{pmatrix}
\quad [\S.\,79, (2)].
\]

Nehmen wir dann $v_{1,0,0}$ für $v_0$, so findet man durch wiederholte
Anwendung von $\tau$ auf $1, 0, 0$ folgende Anordnung:
\[
(1,0,0)=0, \quad (1,1,0)=1, \quad (1,1,1)=2, \quad (0,1,1)=3, \quad (1,0,1)=4,
\]
\[
(0,1,0)=5, \quad (0,0,1)=6,
\]
und nach \S.\,82, (2) hat man daher
\[
v_1, v_2, v_3, v_4, v_5, v_6, v_7
\]
durch
\[
v_0, v_5, v_6, v_3, v_4, v_1, v_2
\]
zu ersetzen. Dadurch ergiebt sich aus \S.\,82, (3):
\[
8X = A \sum_{0,\,6}^{z} \sqrt{v_z}\,\sqrt{v_{z+1}}\,\sqrt{v_{z+2}}\,\sqrt{v_{z+4}} + B,
\]
worin $A, B$ rationale Grössen, und $v_z$ die Wurzeln einer cyklischen oder
einer metacyklischen Gleichung von der Gruppe $(z, az+b)$ vom Grade 7 sind.

Umgekehrt ist es auch nicht schwer, nachzuweisen, dass die acht durch
Wechsel der Vorzeichen hieraus abgeleiteten Grössen immer die Wurzeln
einer primitiven metacyklischen Gleichung
%CONT>>


% ===== scan leaf 340, printed 319 =====
% !PRINTED 319
\weberpage{319}
\pdfbookmark[2]{§ 84. Biquadratische Gleichungen}{bm97}
%<<CONT
8\textsuperscript{ten} Grades sind, worauf hier nicht näher eingegangen
werden soll (vgl. Bd. I, \S.\,185).

\begin{center}
\S.\ 84.

\gesperrt{Biquadratische Gleichungen.}
\end{center}

Wir haben uns auf die Betrachtung der primitiven metacyklischen
Gleichungen 8\textsuperscript{ten} Grades beschränkt, weil die imprimitiven
auf Gleichungen 4\textsuperscript{ten} Grades zurückkommen. Um also auch die
Wurzeln der letzteren in der Weise des vorigen Paragraphen darstellen zu
können, haben wir noch eine analoge Darstellung der Wurzeln biquadratischer
Gleichungen zu suchen. Hierbei sind dieselben Betrachtungen, wie im
\S.\,82, nur in wesentlich einfacherer Gestalt anwendbar.

Wir bezeichnen die Wurzeln der biquadratischen Gleichung mit
\weq{1}{X_{0,0}, X_{1,0}, X_{0,1}, X_{1,1},}

und können dann die ganze Permutationsgruppe 24\textsuperscript{sten}
Grades als binäre lineare Congruenzgruppe für den Modul 2 darstellen:
\weq{2}{P = SQ.}

Die darin enthaltene homogene Gruppe $S$ ist vom 6\textsuperscript{ten}
Grade und besteht aus den Substitutionen
\weq{3}{\sigma = \begin{pmatrix}1, & 0\\0, & 1\end{pmatrix},
\begin{pmatrix}1, & 0\\1, & 1\end{pmatrix},
\begin{pmatrix}0, & 1\\1, & 0\end{pmatrix},
\begin{pmatrix}0, & 1\\1, & 1\end{pmatrix},
\begin{pmatrix}1, & 1\\1, & 0\end{pmatrix},
\begin{pmatrix}1, & 1\\0, & 1\end{pmatrix},}

und man kann
\weq{4}{\sigma_1 = \begin{pmatrix}0, & 1\\1, & 0\end{pmatrix}, \quad
\sigma_2 = \begin{pmatrix}1, & 1\\0, & 1\end{pmatrix}}

als die erzeugenden Substitutionen der Gruppe $S$ ansehen.

Wir definiren nun wie im \S.\,80 die drei Grössen:
\weq{5}{\begin{aligned}
\Psi_{1,0} &= X_{0,0} - X_{1,0} + X_{0,1} - X_{1,1}\\
\Psi_{0,1} &= X_{0,0} + X_{1,0} - X_{0,1} - X_{1,1}\\
\Psi_{1,1} &= X_{0,0} - X_{1,0} - X_{0.1} + X_{1,1}, % [sic] printed with a period after 0 instead of a comma
\end{aligned}}

so dass $\Psi^2_{1,0}, \Psi^2_{0,1}, \Psi^2_{1,1}$, aber auch das Product
$\Psi_{1,0}\Psi_{0,1}\Psi_{1,1}$ durch die Substitutionen der cyklischen
Gruppe $Q$ ungeändert bleiben.


% ===== scan leaf 341, printed 320 =====
% !PRINTED 320
\weberpage{320}
Wendet man auf die Indices $\xi, \eta$ der $X$ die Substitution $\sigma$
an, so erleiden die Indices der $\Psi$ [wie im \S.\,80, (2)] die
Substitution $\varrho^{-1}$, wenn $\varrho$ die transponirte Substitution
von $\sigma$ ist.

Es ist aber, den Substitutionen (4) entsprechend,
\weq{6}{\varrho_1^{-1} = \begin{pmatrix}0, & 1\\1, & 0\end{pmatrix}, \quad
\varrho_2^{-1} = \begin{pmatrix}1, & 0\\1, & 1\end{pmatrix}.}

Indem wir nun der Einfachheit halber voraussetzen, dass die drei Grössen
(5) von Null verschieden sind, setzen wir
\weq{7}{\begin{aligned}
v_{1,0} &= \frac{\Psi_{0,1}\Psi_{1,1}}{\Psi_{1,0}}
= \frac{\Psi_{0,1}\Psi_{1,0}\Psi_{1,1}}{\Psi_{1,0}^2};\\
v_{0,1} &= \frac{\Psi_{1,0}\Psi_{1,1}}{\Psi_{0,1}}
= \frac{\Psi_{0,1}\Psi_{1,0}\Psi_{1,1}}{\Psi_{0,1}^2};\\
v_{1,1} &= \frac{\Psi_{1,0}\Psi_{0,1}}{\Psi_{1,1}}
= \frac{\Psi_{1,0}\Psi_{0,1}\Psi_{1,1}}{\Psi_{1,1}^2}.
\end{aligned}}

Die Grössen $v_{\xi,\eta}$ bleiben durch die Substitutionen von $Q$
ungeändert und erleiden durch $\sigma_1, \sigma_2$ die durch (6)
bestimmten Permutationen
\[
\begin{array}{l l l l}
 & v_{1,0}, & v_{0,1}, & v_{1,1}\\
\sigma_1 & v_{0,1}, & v_{1,0}, & v_{1,1}\\
\sigma_2 & v_{1,1}, & v_{0,1}, & v_{1,0}.
\end{array}
\]

Daraus ergiebt sich, dass die symmetrischen Functionen der drei Grössen
$v_{1,0}, v_{0,1}, v_{1,1}$ durch die Permutationen $P$ ungeändert
bleiben, und folglich in dem Rationalitätsbereiche der symmetrischen
Functionen der $X$ enthalten sind. Die drei Grössen $v$ sind also die
Wurzeln einer cubischen Gleichung, und es würde keine Schwierigkeit
haben, diese cubische Resolvente der allgemeinen biquadratischen
Gleichung zu bilden.

Aus (7) folgt aber ferner
\[
v_{1,0}\,v_{0,1} = \psi_{1,1}^2, \quad v_{1,0}\,v_{1,1} = \psi_{0,1}^2,
\quad v_{0,1}\,v_{1,1} = \psi_{1,0}^2, % [sic] printed with lowercase psi, elsewhere capital Psi is used for these quantities
\]
und demnach ergiebt sich aus (5), wenn wir noch die rationale Grösse
\[
X_{0,0} + X_{1,0} + X_{0,1} + X_{1,1} = a
\]
setzen,
\weq{8}{4X_{0,0} = a + \sqrt{v_{1,0}}\,\sqrt{v_{0,1}}
+ \sqrt{v_{1,0}}\,\sqrt{v_{1,1}} + \sqrt{v_{0,1}}\,\sqrt{v_{1,1}}.}

Der Ausdruck (8) giebt nur vier verschiedene Werthe, wenn wir die
Vorzeichen der drei darin vorkommenden Quadratwurzeln auf alle mögliche
Arten bestimmen, und man kann auch um-
%CONT>>


% ===== scan leaf 342, printed 321 =====
% !PRINTED 321
\weberpage{321}
%<<CONT
gekehrt leicht zeigen, dass die vier in der Form (8) enthaltenen
Grössen $X$, wenn die $v_{1,0}, v_{0,1}, v_{1,1}$ die Wurzeln irgend einer
cubischen Gleichung sind, einer biquadratischen Gleichung mit
rationalen Co\"efficienten genügen\textsuperscript{1)}.

\noindent\rule{0.35\linewidth}{0.4pt}

\footnotesize
\noindent\textsuperscript{1)} Diese Form der Lösung der biquadratischen Gleichung ist das
Analogon zu der \gesperrt{Cayley'schen} Form der \gesperrt{Cardanischen Formel} (Bd. I,
\S.\,36). Ich weiss nicht, ob diese Form der Auflösung biquadratischer
Gleichungen in der elementaren Algebra sonst schon bekannt ist. Ich
habe sie zuerst in einer englischen Aufgaben-Sammlung \glqq Mathematical
Tripos, Part I, June 1894\grqq{} gefunden, wo auch die cubische Resolvente
der $v$ gebildet ist.
\normalsize


% ===== scan leaf 343, printed 322 =====
\weberpage{322}% page number not printed on this page
\pdfbookmark[1]{Elfter Abschnitt. Die Wendepunkte einer Curve dritter Ordnung.}{bm98}
\pdfbookmark[2]{§ 85. Ternäre Formen und algebraische Curven}{bm99}
% !PRINTED none
% Part-title (Abschnitt-opening) page: no running head / page number printed on this leaf.

\begin{center}
\gesperrt{Elfter Abschnitt.}

\textbf{\gesperrt{Die Wendepunkte einer Curve dritter Ordnung.}}

\noindent\rule{0.25\linewidth}{0.4pt}

\S.\,85.

\gesperrt{Ternäre Formen und algebraische Curven.}
\end{center}

Um im weiteren Verlauf unserer Darstellung dem Leser
einige von den mannigfaltigen und interessanten Anwendungen
der Algebra auf geometrische Probleme vorführen zu können,
ohne ihn auf andere Hülfsmittel zu verweisen, sollen hier die
nothwendigen geometrischen Grundlagen für diese Anwendungen
in der Kürze abgeleitet werden. Wir beschränken uns dabei
auf die Geometrie der Ebene und lassen auch da Alles bei Seite,
was für unsere Anwendungen nicht direct erforderlich ist\textsuperscript{1)}.

Wir wenden die homogenen Dreieckscoordinaten an, die dadurch erklärt sind, dass die Lage eines Punktes $x$ durch seine
Abstände von den drei Seiten des Coordinatendreiecks, in drei
beliebigen Maasseinheiten gemessen (oder mit drei beliebigen
constanten Factoren multiplicirt), bestimmt sind. Diese drei
Grössen heissen die \gesperrt{Coordinaten} des Punktes $x$, und werden
meist mit $x_1, x_2, x_3$ oder mit $y_1, y_2, y_3$ etc. bezeichnet.

Zwischen diesen drei Coordinaten des Punktes $x$ besteht
eine lineare, nicht homogene Relation, die man erhält, wenn man
den Flächeninhalt des Coordinatendreiecks $(1, 2, 3)$ gleich der
Summe der Inhalte der drei Dreiecke $(x, 2, 3)$, $(x, 3, 1)$, $(x, 1, 2)$
setzt. Eine Coordinate und also auch der Ausdruck für den
Flächeninhalt des entsprechenden der drei letztgenannten Drei-
%CONT>>

\noindent\rule{0.25\linewidth}{0.4pt}

\footnotesize
\noindent\textsuperscript{1)} Ausführlicheres findet man in den Lehrbüchern der analytischen
Geometrie, unter denen das Werk von \gesperrt{George Salmon}, \glqq Higher plane
curves\grqq{}, deutsch von \gesperrt{Fiedler}, hervorzuheben ist.
\normalsize


% ===== scan leaf 344, printed 323 =====
% !PRINTED 323
\weberpage{323}
%<<CONT
ecke ändern ihr Vorzeichen, so oft der Punkt $x$ durch eine Seite
des Coordinatendreiecks hindurchgeht.

Mit Hülfe dieser nicht homogenen Relation kann die Einheit
linear und homogen durch $x_1, x_2, x_3$ dargestellt und dadurch
jede nicht homogene Function in eine homogene verwandelt
werden.

Die Coordinaten $x_1, x_2, x_3$ sind nach diesen Bestimmungen
zunächst nicht unabhängige Variable, sondern an die erwähnte
Relation gebunden. Wir erweitern ihre Bedeutung aber jetzt so,
dass wir sie als unabhängige Variable betrachten können, wenn
wir festsetzen, dass wir sie nur in homogenen Formen oder
Gleichungen benutzen wollen;
und dass nicht $x_1, x_2, x_3$ selbst,
sondern nur die Verhältnisse $x_1 : x_2 : x_3$ die Lage des Punktes $x$
bestimmen sollen. Dann können wir von der erwähnten Relation
zwischen den $x_1, x_2, x_3$ ganz absehen, und die $x_1, x_2, x_3$ als unabhängige Variable betrachten, für die nur die einzige Werthcombination $x_1 = 0, x_2 = 0, x_3 = 0$ ausgeschlossen ist. Die
Werthe $h x_1, h x_2, h x_3$ bestimmen dann für jedes von Null verschiedene $h$ einen und denselben Punkt $x$. Genau gesagt, sind
dann die Coordinaten $x_1, x_2, x_3$ des Punktes $x$ nicht die oben
beschriebenen Abstände von den Coordinatenaxen selbst, sondern irgend welche Grössen, die mit diesen Abständen proportional sind.

Jede lineare Substitution

\weq{1}{\begin{aligned}
x_1 &= \alpha_1^{(1)} y_1 + \alpha_1^{(2)} y_2 + \alpha_1^{(3)} y_3 \\
x_2 &= \alpha_2^{(1)} y_1 + \alpha_2^{(2)} y_2 + \alpha_2^{(3)} y_3 \\
x_3 &= \alpha_3^{(1)} y_1 + \alpha_3^{(2)} y_2 + \alpha_3^{(3)} y_3,
\end{aligned}}

\noindent deren Determinante

\weq{2}{r = \Sigma \pm \alpha_1^{(1)} \alpha_2^{(2)} \alpha_3^{(3)}}

\noindent von Null verschieden ist, kann in doppelter Weise geometrisch
gedeutet werden, einmal als eine \gesperrt{Abbildung}, indem wir $x_1, x_2, x_3$
und $y_1, y_2, y_3$ als Coordinaten zweier Punkte $x, y$ in demselben
Coordinatensysteme ansehen, so dass jedem Punkte $y$ ein Punkt $x$
entspricht und umgekehrt, sodann auch als \gesperrt{Coordinatentransformation}, wenn wir $x_1, x_2, x_3$ und $y_1, y_2, y_3$ als Coordinaten eines und desselben Punktes, bezogen auf zwei verschiedene
Coordinatensysteme, betrachten. Für die nächsten Betrachtungen
werden wir an der letzteren Interpretation festhalten.


% ===== scan leaf 345, printed 324 =====
% !PRINTED 324
\weberpage{324}
\pdfbookmark[2]{§ 86. Singuläre Punkte. Wendepunkte. Doppeltangenten}{bm100}
Irgend eine ternäre Form $n^{\text{ten}}$ Grades $f(x_1, x_2, x_3)$ stellt,
gleich Null gesetzt, eine Curve $n^{\text{ter}}$ Ordnung dar, die mit einer
geraden Linie $n$ Schnittpunkte hat. Diese Curve bezeichnen wir
kurz als die Curve $f$.

Durch die Substitution (1) geht $f$ in eine andere Form desselben Grades in den Variablen $y_1, y_2, y_3$ über:

\weq{3}{f(x_1, x_2, x_3) = \varphi(y_1, y_2, y_3),}

\noindent und $\varphi = 0$ stellt dieselbe Curve dar, wie $f = 0$, bezogen auf
das Coordinatensystem $y_1, y_2, y_3$.

Bilden wir die Ableitungen der Identität (3) nach den
Variablen $y_1, y_2, y_3$, so folgt mit Rücksicht auf (1):

\weq{4}{\begin{aligned}
\varphi'(y_1) &= \alpha_1^{(1)} f'(x_1) + \alpha_2^{(1)} f'(x_2) + \alpha_3^{(1)} f'(x_3) \\
\varphi'(y_2) &= \alpha_1^{(2)} f'(x_1) + \alpha_2^{(2)} f'(x_2) + \alpha_3^{(2)} f'(x_3) \\
\varphi'(y_3) &= \alpha_1^{(3)} f'(x_1) + \alpha_2^{(3)} f'(x_2) + \alpha_3^{(3)} f'(x_3).
\end{aligned}}

\begin{center}
\S.\,86.

\gesperrt{Singuläre Punkte. Wendepunkte. Doppeltangenten.}
\end{center}

Wenn für einen Punkt $x$ die drei Gleichungen

\weq{1}{f'(x_1) = 0, \quad f'(x_2) = 0, \quad f'(x_3) = 0}

\noindent zugleich erfüllt sind, so liegt nach dem \gesperrt{Euler'schen Theorem}

\weq{2}{n f = x_1 f'(x_1) + x_2 f'(x_2) + x_3 f'(x_3)}

\noindent (Bd. I, \S.\,17) dieser Punkt auf der Curve $f$. Die Gleichungen \S.\,85,(4)
zeigen, dass in demselben Punkte auch $\varphi'(y_1), \varphi'(y_2), \varphi'(y_3)$
verschwinden müssen, dass also die durch die Gleichungen (1)
charakterisirte Eigenschaft eines Punktes bei linearer Transformation erhalten bleibt, also, wie wir uns auch ausdrücken, zu
den \gesperrt{invarianten} Eigenschaften gehört.

Nicht auf jeder Curve $n^{\text{ter}}$ Ordnung kommen Punkte vor, die
den Bedingungen (1) genügen, wie z.\,B. die Annahme $f(x_1, x_2, x_3)$
$= x_1^n + x_2^n + x_3^n$ zeigt. Es wird vielmehr, wenn auch nur ein
solcher Punkt existiren soll, eine Gleichung zwischen den Co\"efficienten von $f$ erfüllt sein müssen, die man durch Elimination
von $x_1, x_2, x_3$ aus den drei Gleichungen (1) erhält (Bd. I, \S.\,50).
Man kann diese Relation in der Form darstellen, dass eine
gewisse ganze rationale und homogene Function der sämmtlichen
Co\"efficienten von $f$ gleich Null sein muss, und diese Function
%CONT>>


% ===== scan leaf 346, printed 325 =====
% !PRINTED 325
\weberpage{325}
%<<CONT
heisst die Discriminante von $f$. Sie ist durch diese Definition nur bis auf einen numerischen Factor bestimmt, und kann bis jetzt nur in den einfachsten Fällen berechnet werden.

Die Punkte der Curve $f$, deren Coordinaten den Bedingungen (1) genügen, heissen \gesperrt{singuläre Punkte}.

Man kann auch die Frage aufwerfen, unter welcher Bedingung es möglich ist, dass auf einer Curve unendlich viele singuläre Punkte vorkommen. Damit dies eintritt, muss zunächst das Eliminationsresultat von einer der Unbekannten, etwa von $x_3$, aus den zwei Gleichungen $f'(x_1) = 0$, $f'(x_2) = 0$ identisch verschwinden, d.\,h. die beiden Functionen $f'(x_1)$ und $f'(x_2)$ müssen, als Functionen von $x_3$ betrachtet, einen gemeinschaftlichen Theiler haben, und nach Bd. I, \S.\,51 müssen sie daher auch einen (nicht constanten) gemeinsamen Theiler im Gebiete der Formen $x_1, x_2, x_3$ haben. Dieser Theiler muss dann wieder, wie aus denselben Erwägungen hervorgeht, einen gemeinschaftlichen Theiler mit $f$, oder, was dasselbe ist, mit $f'(x_3)$ haben, und es ergiebt sich also, dass nur dann unendlich viele singuläre Punkte vorhanden sein können, wenn die vier Formen $f, f'(x_1), f'(x_2), f'(x_3)$ einen gemeinschaftlichen Theiler haben.

Es sei nun $u$ ein solcher gemeinschaftlicher Theiler, den wir als \gesperrt{irreducibel} voraussetzen.

Wir setzen
\[
f(x_1, x_2, x_3) = uv
\]
\[
f'(x_1) = u\,\frac{\partial v}{\partial x_1} + v\,\frac{\partial u}{\partial x_1}
\]
\[
f'(x_2) = u\,\frac{\partial v}{\partial x_2} + v\,\frac{\partial u}{\partial x_2}
\]
\[
f'(x_3) = u\,\frac{\partial v}{\partial x_3} + v\,\frac{\partial u}{\partial x_3},
\]
und diese Gleichungen zeigen, da $u$ irreducibel ist und daher
\[
u, \quad \frac{\partial u}{\partial x_1}, \quad \frac{\partial u}{\partial x_2}, \quad \frac{\partial u}{\partial x_3}
\]
keinen gemeinschaftlichen Theiler haben können, dass $v$ durch $u$ theilbar sein muss; d.\,h. es können nur dann unendlich viele singuläre Punkte vorkommen, wenn $f(x_1, x_2, x_3)$ einen \gesperrt{quadratischen Theiler} hat. In diesem Falle sagen wir, dass die Curve $f$ einen doppelt zählenden Bestandtheil enthält.


% ===== scan leaf 347, printed 326 =====
% !PRINTED 326
\weberpage{326}
Um die geometrische Bedeutung der singulären Punkte zu erkennen, denken wir uns die Function $f$ nach Potenzen der einen Variablen $x_3$ geordnet:

\weq{2}{f(x_1, x_2, x_3) = x_3^n f_0 + x_3^{n-1} f_1 + x_3^{n-2} f_2 + \cdots,}
worin $f_0, f_1, f_2, \dots$ binäre Formen der Variablen $x_1, x_2$ sind, deren Grad durch den Index angegeben ist, so dass $f_0$ eine Constante, $f_1$ eine lineare, $f_2$ eine quadratische Form ist u.\,s.\,f.

Wenn nun ein singulärer Punkt vorhanden ist, so können wir wegen der Invarianz dieser Eigenschaft annehmen, dass dieser Punkt in die Ecke $x_1 = 0$, $x_2 = 0$ des Coordinatendreiecks, die wir mit $\xi_3$ bezeichnen wollen, falle. Dann müssen $f_0$ und $f_1$ verschwinden, und wenn wir dann $x_1 = 0$ setzen, so erhält die Gleichung $f = 0$ den Factor $x_2^2$. Dies besagt, dass zwei der Schnittpunkte der Linie $x_1 = 0$ in dem singulären Punkte zusammenfallen. Die Linie $x_2 = 0$ kann aber jede durch den singulären Punkt hindurchgehende Gerade sein, und so folgt, dass der singuläre Punkt die Eigenschaft hat, dass \gesperrt{jede} durch ihn hindurchgehende Gerade die Curve dort in zwei zusammenfallenden Punkten schneidet. % [sic: the source prints $x_2=0$ here; the preceding derivation uses the cutting line $x_1=0$]

Aus diesem Grunde nennt man die singulären Punkte auch \gesperrt{Doppelpunkte}; womit aber nicht ausgeschlossen ist, dass auch mehr als zwei Schnittpunkte zusammenfallen können, wodurch die \gesperrt{höheren Singularitäten} entstehen.

An die Form der Gleichung (2) knüpfen wir noch einige in der Folge wichtige einfache Betrachtungen.

Wenn der Punkt $\xi_3$ auf der Curve $f$ liegt, aber kein singulärer Punkt ist, so muss die Constante $f_0$ verschwinden und $f_1 = 0$ ist die Gleichung der \gesperrt{Tangente} der Curve in diesem Punkte. Nehmen wir diese Tangente für die Linie $x_1 = 0$, so lautet die Gleichung der Curve, wenn wir einen constanten Factor gleich 1 annehmen:

\weq{3}{f = x_3^{n-1} x_1 + x_3^{n-2} f_2 + \cdots = 0.}
Wenn die quadratische Form $f_2$ durch $x_1$ theilbar ist, so erhält die Gleichung (3) für $x_1 = 0$ den Factor $x_2^3$ und der Punkt $\xi_3$ zählt als dreifacher Schnittpunkt von $x_1 = 0$ mit der Curve. Ein solcher Punkt heisst ein \gesperrt{Wendepunkt} oder \gesperrt{Inflexionspunkt} der Curve $f$, und $x_1 = 0$ heisst die \gesperrt{Wendetangente}. Die Wendepunkte gehören nicht zu den singulären Punkten.


% ===== scan leaf 348, printed 327 =====
% !PRINTED 327
\weberpage{327}
\pdfbookmark[2]{§ 87. Fundamentale Covarianten einer ternären Form}{bm101}
Ist der Punkt $\xi_3$ ein Wendepunkt und $x_1 = 0$ die Wendetangente, so hat $f$ die Gestalt

\weq{4}{f = x_1 \varphi + x_2^3 \psi,}
worin $\varphi$ eine Form $(n-1)^{\text{ten}}$, $\psi$ eine Form $(n-3)^{\text{ten}}$ Grades ist.

Umgekehrt ist, wenn $f$ diese Form hat, der Punkt $\xi_3$ ein Wendepunkt, und $x_1 = 0$ die Wendetangente.

Eine gerade Linie, die die Curve $f$ in zwei getrennten Punkten berührt, heisst eine \gesperrt{Doppeltangente}. Nehmen wir eine solche Doppeltangente zur Linie $x_3 = 0$ und die Berührungspunkte als die auf $x_3 = 0$ gelegenen Ecken des Coordinatendreiecks, so muss, wenn $x_3 = 0$ in $f = 0$ eingesetzt wird, eine Gleichung entstehen, die sowohl $x_1$ als $x_2$ zum quadratischen Factor hat. Die Function $f$ muss also die Gestalt haben:

\weq{4}{f = x_3 \varphi + x_1^2 x_2^2 \psi,} % [sic] number (4) reprinted; page also numbers preceding display (4)
worin $\varphi$ eine Form $(n-1)^{\text{ten}}$ Grades, $\psi$ eine Form $(n-4)^{\text{ten}}$ Grades ist.

Etwas allgemeiner können wir auch sagen: Wenn $x_3 = 0$ eine Doppeltangente ist, so kann $f$ so dargestellt werden:

\weq{5}{f = x_3 \varphi + \chi^2 \psi,}
worin $\chi = 0$ die Gleichung eines Kegelschnittes ist. Die Berührungspunkte der Doppeltangente sind die Schnittpunkte von $x_3 = 0$ und $\chi = 0$.

\begin{center}
\S.\,87.

\gesperrt{Fundamentale Covarianten einer ternären Form.}
\end{center}

Im \S.\,60 des ersten Bandes haben wir den Begriff der Invarianten und Covarianten einer Form kennen gelernt. Auch in der Geometrie spielen diese Formen eine grosse Rolle.

Wenn durch die lineare Substitution \S.\,85, (1) die ternäre Form $f(x_1, x_2, x_3)$, die wir abgekürzt auch mit $f(x)$ bezeichnen, in $\varphi(y)$ übergeht, und wenn wir mit $a$ die Coëfficienten von $f$, mit $b$ die Coëfficienten von $\varphi$ und mit $r$ die Substitutionsdeterminante bezeichnen, so heisst eine Form $\Phi(x, a)$, die sowohl in Bezug auf die $x$ wie in Bezug auf die $a$ homogen ist, eine \gesperrt{Covariante} von $f(x)$, wenn sie der Bedingung

\weq{1}{\Phi(y, b) = r^\lambda \Phi(x, a)}
%CONT>>


% ===== scan leaf 349, printed 328 =====
% !PRINTED 328
\weberpage{328}
%<<CONT
genügt. Wenn insbesondere die Form von den Variablen $x$ unabhängig und nur von den Coëfficienten $a$ abhängig ist, so wird sie zur \gesperrt{Invariante}. Der Exponent $\lambda$, der das \gesperrt{Gewicht} der Covariante heisst, ist eine ganze Zahl.

Wir haben im Bd. I, \S.\,59 eine bei allen Formen, deren Grad grösser als 1 ist, vorhandene Covariante vom Gewicht 2 kennen gelernt, nämlich die \gesperrt{Hesse'sche Determinante}:

\weq{2}{H(x_1, x_2, x_3) = \begin{vmatrix} f''(x_1, x_1), & f''(x_1, x_2), & f''(x_1, x_3) \\ f''(x_2, x_1), & f''(x_2, x_2), & f''(x_2, x_3) \\ f''(x_3, x_1), & f''(x_3, x_2), & f''(x_3, x_3) \end{vmatrix}.}

Ausser dieser Covariante $H$ wollen wir noch zwei andere allgemeine Covarianten betrachten, die bei allen ternären Formen von höherem als dem $2^{\text{ten}}$ Grade existiren, und die sich leicht auch für Formen von beliebig vielen Variablen verallgemeinern lassen. Die erste dieser Formen ist, wenn $H$ wie oben die \gesperrt{Hesse}'sche Covariante bedeutet:

\weq{3}{C_1(x, a) = \begin{vmatrix} f''(x_1, x_1), & f''(x_1, x_2), & f''(x_1, x_3), & H'(x_1) \\ f''(x_2, x_1), & f''(x_2, x_2), & f''(x_2, x_3), & H'(x_2) \\ f''(x_3, x_1), & f''(x_3, x_2), & f''(x_3, x_3), & H'(x_3) \\ H'(x_1), & H'(x_2), & H'(x_3) & 0 \end{vmatrix}.}

Um die Invarianten-Eigenschaft für diese Form nachzuweisen, machen wir darin die Substitution \S.\,85, (1). Dadurch geht $f(x, a)$ in $\varphi(y, b)$ und $r^2 H(x, a)$ in $H(y, b)$ über. Wenn wir dann unter $H'(y_1), H'(y_2), H'(y_3)$ die Derivirten von $H(y, b)$ verstehen, so haben wir die Formeln

\weq{4}{
\begin{aligned}
\varphi''(y_\mu, y_\nu) &= \sum\limits^{i,k}_{i} f''(x_i, x_k)\, \alpha_i^{(\mu)} \alpha_k^{(\nu)} \\
H'(y_\nu) &= r^2 \sum H'(x_i)\, \alpha_i^{(\nu)},
\end{aligned}
}
in denen $\mu, \nu, i, k$ von 1 bis 3 laufen.

Multiplicirt man nun die Determinante $C_1(x, a)$ zweimal nach einander mit der in der Form

\[
r = \begin{vmatrix} \alpha_1^{(1)}, & \alpha_1^{(2)}, & \alpha_1^{(3)}, & 0 \\ \alpha_2^{(1)}, & \alpha_2^{(2)}, & \alpha_2^{(3)}, & 0 \\ \alpha_3^{(1)}, & \alpha_3^{(2)}, & \alpha_3^{(3)}, & 0 \\ 0, & 0, & 0, & 1 \end{vmatrix}
\]
geschriebenen Substitutionsdeterminante, indem man das eine
%CONT>>


% ===== scan leaf 350, printed 329 =====
% !PRINTED 329
\weberpage{329}
\pdfbookmark[2]{§ 88. Die Hesse'sche Curve}{bm102}
%<<CONT
Mal nach Zeilen, das andere Mal nach Columnen multiplicirt, so
folgt nach den Formeln (4):

\weq{5}{C_1\,(y,b) = r^6\,C_1\,(x,a),}

wodurch die Invarianten-Eigenschaft ausgedrückt und das Gewicht
$= 6$ bestimmt ist.

Die dritte Covariante $C_2$, die wir betrachten, ist die Functionaldeterminante
 aus den drei Formen $f$, $H$, $C_1$, also

\weq{6}{C_2\,(x, a) = \begin{vmatrix} f'\,(x_1), & H'\,(x_1), & C_1'\,(x_1) \\ f'\,(x_2), & H'\,(x_2), & C_1'\,(x_2) \\ f'\,(x_3), & H'\,(x_3), & C_1'\,(x_3) \end{vmatrix},}

die, wie schon aus \S.\,59 des ersten Bandes hervorgeht, die
Invarianten-Eigenschaft hat:

\weq{7}{C_2\,(y,b) = r^9\,C_2\,(x,a),}

und die also das Gewicht 9 hat.

Bezeichnen wir mit $\lambda$ das Gewicht, mit $\nu$ den Grad in den
Variablen, mit $\mu$ den Grad in den Coëfficienten, so ergiebt sich
folgende Zusammenstellung:

\begin{center}
\begin{tabular}{lccc}
 & $\nu$, & $\mu$, & $\lambda$ \\
\hline
$f$ & $n$, & 1, & 0 \\
$H$ & $3n-6$, & 3, & 2 \\
$C_1$ & $8n-18$, & 8, & 6 \\
$C_2$ & $12n-27$, & 12, & 9.
\end{tabular}
\end{center}

\begin{center}
\S.\,88.

\gesperrt{Die Hesse'sche Curve.}
\end{center}

Die erste der eingeführten Covarianten $H$ hat eine einfache
geometrische Bedeutung, die wir jetzt kennen lernen wollen.
Diese Untersuchung wird dadurch wesentlich erleichtert, dass
wir, wegen der Invarianten-Eigenschaft der Form $H$, nichts an
Allgemeinheit verlieren, wenn wir dem Coordinatensysteme irgend
eine specielle Lage gegen die Curve $f$ geben.

Legen wir die Ecke $\xi_3$ des Coordinatendreiecks in einen
Punkt der Curve $f$, der nicht zu den singulären gehört, und
wählen die Tangente in diesem Punkte zur $x_1$-Axe, so erhält
nach \S.\,86, (3) die Form $f$ den Ausdruck:

\weq{1}{f = h\,x_3^{n-1}x_1 + x_3^{n-2}(a\,x_1^2+2\,b\,x_1x_2+c\,x_2^2)+x_3^{n-3}f_3+\cdots,}
%CONT>>


% ===== scan leaf 351, printed 330 =====
% !PRINTED 330
\weberpage{330}
%<<CONT
worin $h$, $a$, $b$, $c$ Constanten sind, von denen $h$ von Null verschieden
 ist, und $f_3$ eine binäre cubische Form von $x_1, x_2$.

Wenn wir hiernach die nach absteigenden Potenzen von $x_3$
geordnete Function $H$ berechnen:

\weq{2}{H = x_3^{3n-6}\,H_0 + x_3^{3n-7}\,H_1 + \cdots,}

so findet sich

\weq{3}{\begin{aligned} H_0 &= -\,2\,(n-1)^2\,h^2\,c \\ H_1 &= -(n-1)^2\,h^2\,f_3''\,(x_2,x_2) + 4\,(n-1)(n-2)\,h\,(b^2-a\,c)\,x_1. \end{aligned}}

Die durch die Gleichung $H=0$ dargestellte Curve heisst
die \gesperrt{Hesse'sche Curve} der Curve $f$. Die Formeln (2) und (3)
zeigen, dass die Hesse'sche Curve dann und nur dann durch den
Punkt $\xi_3$ hindurchgeht, wenn $H_0=0$, also $c=0$, d.\,h. nach \S.\,86,
wenn dieser Punkt ein Wendepunkt ist. Aus dem Ausdrucke
für $H_1$ kann man noch weiter schliessen, dass die Wendetangente
$x_1=0$ dann und nur dann zugleich Tangente der Curve $H=0$
ist, wenn $f_3''\,(x_2,x_2)$ und folglich auch $f_3$ selbst durch $x_1$ theilbar
ist. In diesem Falle wäre, wie aus (3) hervorgeht, $x_1=0$ eine
im Punkte $\xi_3$ \gesperrt{vierpunktig berührende Tangente}. Dieser
Fall kann bei den Curven dritter Ordnung nur dann vorkommen,
wenn $f$ durch $x_1$ theilbar ist, also niemals bei einer irreduciblen
Curve dritter Ordnung.

Die \gesperrt{Hesse'sche} Curve geht ausser durch die Wendepunkte
noch durch die singulären Punkte der Curve $f$, was hier nicht
weiter untersucht werden soll.

Wenn wir also unsere Betrachtungen auf Curven $f$ ohne
singulären Punkt beschränken, so können wir sagen, dass jeder
Schnittpunkt der Curve $f$ mit ihrer \gesperrt{Hesse'schen} Curve einen
Wendepunkt der Curve $f$ giebt, und dass diese Curve $f$ keine
anderen Wendepunkte hat. Wenn Punkte mit vierpunktig berührender
 Tangente vorkommen, so berühren sich die Curven $f$
und $H$, und wir können solche Punkte auffassen als durch das
Zusammenfallen von zwei Wendepunkten entstanden. Wenden
wir noch das Theorem von Bezout (Bd.\,I, \S.\,49) an, so erhalten
wir das Ergebniss:

\begin{quote}
\gesperrt{Eine Curve $n^{\mathrm{ter}}$ Ordnung ohne singulären Punkt
hat $3\,n\,(n-2)$ Inflexionspunkte;}
\end{quote}

und speciell:

\begin{quote}
\gesperrt{Eine Curve dritter Ordnung ohne singulären Punkt
hat neun getrennt liegende Inflexionspunkte.}
\end{quote}


% ===== scan leaf 352, printed 331 =====
% !PRINTED 331
\weberpage{331}
\pdfbookmark[2]{§ 89. Inflexionspunkte einer Curve dritter Ordnung}{bm103}
\begin{center}
\S.\,89.

\gesperrt{Inflexionspunkte einer Curve dritter Ordnung.}
\end{center}

Es sei jetzt $f=0$ die Gleichung einer Curve dritter Ordnung
 ohne singulären Punkt. Wir legen zwei der Ecken des
Coordinatendreiecks $\xi_1, \xi_2$ in zwei von den neun Inflexionspunkten
und nehmen die entsprechenden Inflexionstangenten für die
Seiten $x_2, x_1$. Dann muss sich $f$, wenn $x_1=0$ oder $x_2=0$ gesetzt
wird, auf das Glied $h\,x_3^3$ reduciren, worin $h$ eine Constante ist,
und folglich müssen alle Glieder, die in der cubischen Form
einen von Null verschiedenen Coëfficienten haben, ausgenommen
$h\,x_3^3$, durch $x_1\,x_2$ theilbar sein. Demnach erhält $f$ die Form

\weq{1}{f\,(x_1,x_2,x_3) = x_1\,x_2\,(a_1\,x_1+a_2\,x_2+a_3\,x_3) + h\,x_3^3,}

worin $a_1, a_2, a_3$ Constanten sind. Daraus aber geht hervor, dass
die gerade Linie

\[
a_1\,x_1 + a_1\,x_2 + a_3\,x_3 = 0
\] % [sic] printed coefficient of x_2 is a_1; expected a_2

gleichfalls Inflexionstangente der Curve $f$ ist, und dass ihr Schnittpunkt
 mit der Linie $x_3=0$ ein dritter Inflexionspunkt ist.
Dieser Punkt ist von den beiden ersten $\xi_1, \xi_2$ verschieden, weil
weder $a_1$ noch $a_2=0$ sein kann, wenn die Curve keinen singulären
 Punkt hat. Wir wollen dieses wichtige Resultat in Form
eines Satzes aussprechen:

1. \gesperrt{Wenn man irgend zwei Inflexionspunkte einer
Curve dritter Ordnung ohne singulären Punkt
durch eine gerade Linie verbindet, so schneidet
diese gerade Linie die Curve in einem dritten
Inflexionspunkte.}

Da es neun Inflexionspunkte giebt, so gehen von jedem
dieser Punkte vier gerade Linien aus, deren jede noch zwei
weitere Inflexionspunkte enthält. Es giebt neun solcher Büschel
von vier Geraden, und weil jede dieser Geraden in drei Büscheln
vorkommt, so ist die Anzahl der Geraden $9\,.\,4 : 3 = 12$. Wir
können demnach den Satz 1. so ergänzen:

2. \gesperrt{Die neun Inflexionspunkte einer Curve dritter
Ordnung liegen zu je dreien auf zwölf geraden
Linien, und durch jeden Inflexionspunkt gehen
vier von diesen Geraden.}


% ===== scan leaf 353, printed 332 =====
% !PRINTED 332
\weberpage{332}
Aus dem Ausdruck (1) können wir eine \gesperrt{canonische Darstellung}
 für die cubische Form $f$ herleiten:
Wir bezeichnen mit $\varepsilon$ eine imaginäre dritte Einheitswurzel
und führen zwei Variable $z_1, z_2$ ein, die durch die Gleichungen

\weq{2}{\begin{aligned} a_1\,x_1 &= \varepsilon\,z_1 + \varepsilon^2\,z_2 - \tfrac{1}{3}\,a_3\,x_3 \\ a_2\,x_2 &= \varepsilon^2\,z_1 + \varepsilon\,z_2 - \tfrac{1}{3}\,a_3\,x_3 \end{aligned}}

definirt sind, woraus noch folgt:

\weq{3}{-\,(a_1\,x_1 + a_2\,x_2 + a_3\,x_3) = z_1 + z_2 - \tfrac{1}{3}\,a_3\,x_3.}

Die Variablen $z_1, z_2$ sind durch (2) als lineare Functionen
von $x_1, x_2, x_3$ bestimmt, und die linearen Functionen $z_1, z_2, x_3$
sind als Functionen von $x_1, x_2, x_3$ linear unabhängig. Wenn wir
(2) und (3) in (1) einführen, so ergiebt sich

\[
-\,a_1\,a_2\,f = z_1^3 + z_2^3 - \left(a_1\,a_2\,h + \frac{a_3^3}{27}\right)x_3^3 + a_3\,z_1\,z_2\,x_3.
\]

Diesem Ausdrucke giebt man eine mehr symmetrische Gestalt
, indem man für $z_1, z_2, x_3$ drei neue lineare Functionen
$y_1, y_2, y_3$ einführt, die sich von diesen nur um constante Factoren
unterscheiden:

\[
f = \varphi\,(y_1,y_2,y_3) = \alpha_1\,y_1^3 + \alpha_2\,y_2^3 + \alpha_3\,y_3^3 + 6\,m\,y_1\,y_2\,y_3,
\]

worin $\alpha_1, \alpha_2, \alpha_3, m$ Constanten bedeuten.

Wenn kein singulärer Punkt vorhanden ist, so können die
Coëfficienten $\alpha_1, \alpha_2, \alpha_3$ nicht verschwinden, und man kann sie
ohne Beschränkung der Allgemeinheit einander gleich, etwa $=h$,
annehmen. Man könnte sie sogar $=1$ setzen, was aber, um die
Homogenëität aufrecht zu erhalten, zunächst besser nicht geschieht.
Wir haben also die \gesperrt{canonische Form} für die ternäre
cubische Form:

\weq{4}{\varphi\,(y_1,y_2,y_3) = h\,(y_1^3+y_2^3+y_3^3) + 6\,m\,y_1\,y_2\,y_3.}

Nach unserem Ausgangspunkte war die Linie $x_3$, oder was
dasselbe ist, $y_3$ die Verbindungslinie dreier Inflexionspunkte.
Solcher Linien giebt es aber zwölf, und wir können daher die
Linie $y_3$ auf zwölf Arten wählen. Setzen wir aber $y_1$ oder $y_2$ an
Stelle von $y_3$, so bekommen wir dieselbe Darstellung in der
canonischen Form, und es giebt also nur vier wesentlich verschiedene
 Arten dieser Darstellung (abgesehen von dem willkürlichen
 $h$), d.\,h. vier Arten, die Linien $y_1, y_2, y_3$ zu bestimmen. Wir
haben daher den Satz:
%CONT>>


% ===== scan leaf 354, printed 333 =====
% !PRINTED 333
\weberpage{333}
\pdfbookmark[2]{§ 90. Transformation der cubischen Form auf die canonische Form}{bm104}
%<<CONT

\noindent 3.\ \gesperrt{Eine ternäre cubische Form $f(x_1,x_2,x_3)$, deren Discriminante von Null verschieden ist, lässt sich auf vier verschiedene Arten in die canonische Form}
\[
\varphi(y_1,y_2,y_3) = h(y_1^3+y_2^3+y_3^3) + 6\,m\,y_1y_2y_3
\]
\gesperrt{transformiren.}

Da auf jeder der Linien $y_1=0$, $y_2=0$, $y_3=0$ drei Inflexionspunkte liegen, so ist das Problem der Inflexionspunkte wesentlich identisch mit dem der Transformation der cubischen Form $f$ auf die canonische Form, mit dem wir uns zunächst beschäftigen wollen $^1$).

\begin{center}
\S.\,90.

\gesperrt{Transformation der cubischen Form auf die canonische Form.}
\end{center}

Um die Transformation der cubischen ternären Form auf die canonische Form durchzuführen, d. h. auf die Lösung gewisser algebraischer Gleichungen zurückzuführen, müssen wir zunächst die Covarianten für die canonische Form bilden.

Es sei also
\[
\varphi(y) = h(y_1^3+y_2^3+y_3^3) + 6\,m\,y_1y_2y_3,
\]
wofür wir zur Abkürzung auch setzen:
\weq{1}{\varphi(y) = h\,\Sigma y_1^3 + 6\,m\,y_1y_2y_3,}
indem wir unter dem Zeichen $\Sigma$ eine Summe über drei durch cyklische Vertauschung der Variablen $y_1,y_2,y_3$ aus dem ersten gebildete Glieder verstehen.

Die \gesperrt{Hesse'sche} Covariante wollen wir zur Vereinfachung der Formeln mit einem geeigneten numerischen Factor multipliciren, und demnach
\weq{2}{\Delta(x) = \frac{1}{6^3}\left|\begin{array}{ccc} f_{11}, & f_{12}, & f_{13} \\ f_{21}, & f_{22}, & f_{23} \\ f_{31}, & f_{32}, & f_{33} \end{array}\right|}

\medskip
\noindent\rule{4cm}{0.4pt}

\smallskip
\noindent $^1$) Aus der Literatur über die Wendepunkte der Curven dritter Ordnung und über die Theorie der ternären cubischen Formen erwähnen wir ausser den älteren Arbeiten von \gesperrt{Newton} und \gesperrt{Mac Laurin}: \gesperrt{Hesse}, \glqq Ueber die Elimination etc.\grqq{} und \glqq Ueber die Wendepunkte der Curven dritter Ordnung\grqq{}, \gesperrt{Crelle's} Journal, Bd. 28 (1844). \gesperrt{Aronhold}, Theorie der homogenen Functionen 3\textsuperscript{ten} Grades von drei Veränderlichen, \gesperrt{Crelle's} Journal, Bd. 55 (1858). \gesperrt{Salmon}, Higher plane curves (deutsch von \gesperrt{Fiedler}).
%CONT>>


% ===== scan leaf 355, printed 334 =====
% !PRINTED 334
\weberpage{334}
%<<CONT
setzen, worin $f_{ik}$ für $f''(x_i,x_k)$ steht. Für die canonische Form erhält man hieraus durch einfache Rechnung
\weq{3}{\Delta(y) = -h\,m^2\,\Sigma y_1^3 + (h^3+2m^3)\,y_1y_2y_3.}
Daraus leiten wir die zweite Covariante $C_1$ (\S.\,87) her, in der wir gleichfalls einen numerischen Factor einführen, und demnach
\weq{4}{J(x) = \frac{1}{36}\left|\begin{array}{cccc} f_{11}, & f_{12}, & f_{13}, & \Delta_1 \\ f_{21}, & f_{22}, & f_{23}, & \Delta_2 \\ f_{31}, & f_{32}, & f_{33}, & \Delta_3 \\ \Delta_1, & \Delta_2, & \Delta_3, & 0 \end{array}\right|}
setzen, worin $\Delta_i$ für $\Delta'(x_i)$ steht. Auch diese Covariante muss in der canonischen Form dargestellt werden, was keine Schwierigkeit hat, wenn man die Determinante nach der Formel Bd. I, \S.\,23, (12) bildet. Man findet so:
\weq{5}{\begin{aligned}
J(y) &= -h^2(h^3+8m^3)^2\,\Sigma y_2^3y_3^3\\
&\quad + 3m^2(5h^6+26h^3m^3-4m^6)\,y_1^2y_2^2y_3^2\\
&\quad + hm(2h^6+5h^3m^3+20m^6)\,y_1y_2y_3\,\Sigma y_1^3\\
&\quad + 9h^2m^6(\Sigma y_1^3)^2.
\end{aligned}}
Danach lassen sich die drei symmetrischen Functionen $y_1y_2y_3$, $\Sigma y_1^3$, $\Sigma y_2^3y_3^3$ durch die drei Formen $\varphi(y)$, $\Delta(y)$, $J(y)$, und also auch durch die entsprechenden Formen in den ursprünglichen Variablen $x$, die wir kurz durch $f,\Delta,J$ bezeichnen, nach den Invarianten-Relationen
\weq{6}{\varphi(y) = f,\quad \Delta(y) = r^2\Delta,\quad J(y) = r^6 J}
ausdrücken. Man erhält zunächst aus (1) und (3):
\weq{7}{(h^3+8m^3)\,y_1y_2y_3 = m^2 f + r^2\Delta,}
\weq{8}{h(h^3+8m^3)\,\Sigma y_1^3 = (h^3+2m^3)\,f - 6\,m\,r^2\Delta,}
und sodann aus (5)
\weq{9}{\begin{aligned}
h^2(h^3+8m^3)^2\,\Sigma y_2^3y_3^3 &= (2h^3+m^3)\,m^3 f^2 + m(2h^3-5m^3)\,r^2 f\Delta\\
&\quad + 3m^2 r^4 \Delta^2 - r^6 J.
\end{aligned}}
Sehr einfach ist nun die Berechnung der dritten Covariante $C_2$ (\S.\,87) für die canonische Form.

Wir führen auch hier einen vereinfachenden numerischen Factor ein, und setzen
\weq{10}{K = K(x) = \frac{1}{9}\left|\begin{array}{ccc} f_1, & \Delta_1, & J_1 \\ f_2, & \Delta_2, & J_2 \\ f_3, & \Delta_3, & J_3 \end{array}\right|,}
%CONT>>


% ===== scan leaf 356, printed 335 =====
% !PRINTED 335
\weberpage{335}
%<<CONT
worin $f_k,\Delta_k,J_k$ die Ableitungen nach $x_k$ bedeuten. Für die canonische Form ist dann
\weq{11}{K(y) = r^9 K.}
Um $K(y)$ zu finden, brauchen wir nur die Functionaldeterminante aus den linken Theilen der Formeln (7), (8), (9) zu bilden, woraus sich mit Anwendung einfacher Determinantensätze ergiebt (Bd. I, \S.\,22, \S.\,27):
\[
K(y) = -h^3(h^3+8m^3)^3 \left|\begin{array}{ccc} 1, & y_1^3, & y_1^6 \\ 1, & y_2^3, & y_2^6 \\ 1, & y_3^3, & y_3^6 \end{array}\right|,
\]
und mithin ist
\weq{12}{r^9 K = h^3(h^3+8m^3)^3\,(y_1^3-y_2^3)(y_1^3-y_3^3)(y_2^3-y_3^3).}
Zu bemerken ist noch zu diesen Formeln, dass der in ihnen auftretende Factor $h^3+8m^3$ nicht verschwinden kann, wenn die Curve $f$ keinen singulären Punkt hat. Denn ein singulärer Punkt wird bestimmt durch die drei Gleichungen:
\[
h\,y_1^2 + 2m\,y_2y_3 = 0,\quad h\,y_2^2 + 2m\,y_3y_1 = 0,\quad h\,y_3^2 + 2m\,y_1y_2 = 0,
\]
und diese drei Gleichungen sind mit einander verträglich, wenn $h^3+8m^3=0$. (Ist z.\,B. $h=-2m$, so sind sie befriedigt, wenn $y_1=y_2=y_3$ ist.)

Der Co\"efficient $h$ kann für eine gegebene Form $f$ jeden vorgeschriebenen von Null verschiedenen Werth haben, und erst wenn dieser Werth gegeben ist, ist das Problem bestimmt. Wir wollen der Einfachheit halber jetzt $h=1$ setzen, und erhalten aus (7) bis (9) das Resultat:

Setzt man
\weq{13}{\begin{aligned}
P &= \frac{(2+m^3)m^3f^2+(2-5m^3)mr^2f\Delta+3m^2r^4\Delta^2-r^6J}{(1+8m^3)^2}, \\
Q &= \frac{(1+2m^3)f-6mr^2\Delta}{1+8m^3}, \\
R &= \frac{m^2f+r^2\Delta}{1+8m^3},
\end{aligned}}
so sind $y_1^3,y_2^3,y_3^3$ die Wurzeln der cubischen Gleichung
\weq{14}{u^3 - Qu^2 + Pu - R^3 = 0,}
und die Discriminante dieser cubischen Gleichung ist:
\weq{15}{D_1 = \frac{r^{18}K^2}{(1+8m^3)^6}.}


% ===== scan leaf 357, printed 336 =====
% !PRINTED 336
\weberpage{336}
\pdfbookmark[2]{§ 91. Die Invarianten der Curve dritter Ordnung und die biquadratische Gleichung}{bm105}

Die Co\"efficienten der cubischen Gleichung sind bekannt, sobald $m^3, mr^2, r^6$, oder auch, wenn $m$ und $r^2$ bekannt sind.

Ist dann die Gleichung (14) gelöst, so sind $y_1,y_2,y_3$ bis auf dritte Einheitswurzeln, die aber der Relation $y_1y_2y_3=R$ genügen müssen, bestimmt, und weiter können auch diese Grössen der Natur der Sache nach nicht bestimmt werden.

Wenn die Curve $f$ eine reelle Curve ist, so ergiebt sich aus den Formeln (13) und (15) noch der Satz, dass, \gesperrt{wenn $m$ und $r$ reell sind, $D_1$ positiv ist, und folglich die Wurzeln der Gleichung (14) alle drei reell sind.}

Drückt man die Discriminante $D_1$ durch die Coëfficienten der Gleichung $P,Q,R$ aus, so fliesst aus der Formel (15) noch das Resultat, dass man das Quadrat der Covariante $K$ als ganze rationale Function der Covarianten $f,\Delta,J$ ausdrücken kann. Der Ausdruck wird ziemlich complicirt und soll hier nicht weiter verfolgt werden.

Um aber die biquadratische Gleichung zu bilden, von der $m$ und $r$ abhängen, ist es n\"othig, auf die Invarianten dritter Ordnung näher einzugehen.

\begin{center}
\S.\,91.

\gesperrt{Die Invarianten der Curve dritter Ordnung und die biquadratische Gleichung.}
\end{center}

Alle Curven dritter Ordnung, deren Gleichung in der canonischen Form durch dieselben Grössen $\Sigma y_1^3$, $y_1y_2y_3$ linear dargestellt werden kann, haben dieselben Punkte zu Wendepunkten, beispielsweise also die Curven $f=0$ und $\Delta=0$, und allgemeiner alle Curven der Schaar
\weq{1}{F = \lambda f + 6\mu\Delta = 0,}
wo $\lambda,\mu$ beliebige Parameter sind. Alle Curven dieser Schaar schneiden sich in neun festen Punkten, und eine solche Schaar wird ein \gesperrt{Curvenbüschel} genannt. Die neun festen Punkte, die für alle Curven des Büschels die Wendepunkte sind, heissen die \gesperrt{Grundpunkte}. Bilden wir von der Form $F$ die \gesperrt{Hesse'sche} Determinante, so wird auch diese Form linear durch $\Sigma y_1^3$ und $y_1y_2y_3$ ausgedrückt und kann daher nach \S.\,90, (7) und (8) auch linear durch $f$ und $\Delta$ ausgedrückt werden, d. h. wir erhalten eine Form desselben Büschels.


% ===== scan leaf 358, printed 337 =====
% !PRINTED 337
\weberpage{337}
Wir setzen demnach diese Determinante
\weq{2}{\Delta_{\lambda,\mu}(x) = Lf + M\Delta,}

worin $L$ und $M$ binäre Formen 3\textsuperscript{ten} Grades von
$\lambda, \mu$ sind.

Wenn wir eine beliebige lineare Transformation auf die Variablen $x$
anwenden, so ist, wegen der Invarianten-Eigenschaft der Function $\Delta$
\weq{3}{\Delta_{\lambda,\mu}(y,b) = r^2\Delta_{\lambda,\mu\,r^2}(x,a),}

und daraus folgt, dass die Coëfficienten der einzelnen Potenzen und
Producte von $\lambda, \mu$ Covarianten der ursprünglichen Form sind, und
die Coëfficienten $L, M$ von $f$ und $\Delta$ in diesen Ausdrücken, die
rational von den Coëfficienten $a$ der Form $f$ abhängen, sind daher
\gesperrt{Invarianten} dieser Form. Wir haben darin ein Mittel, diese Invarianten
thatsächlich zu berechnen.

Bezeichnen wir nämlich mit $a_i, A_i, B_i$ die Coëfficienten der
einzelnen Potenzen und Producte der Variablen $x$ in $f, \Delta$ und
$\Delta_{\lambda,\mu}$, so ergiebt sich aus (2) ein System von zehn
Gleichungen:
\[
L a_i + M A_i = B_i,
\]
woraus $L$ und $M$ bestimmt werden können. Man erhält, wenn $i, k$ zwei
verschiedene Indices sind,
\weq{4}{\begin{aligned}
L\,(a_i A_k - a_k A_i) &= B_i A_k - B_k A_i\\
M\,(a_i A_k - a_k A_i) &= -B_i a_k + B_k a_i.
\end{aligned}}

Daraus schliesst man noch, dass $L, M$ \gesperrt{ganze} Functionen der
Coëfficienten $a_i$ sind. Denn in (4) sind die rechten Seiten ganze
Functionen, und wenn also $L$ oder $M$, in der einfachsten Gestalt
dargestellt, noch einen Nenner hätten, so müssten die sämmtlichen
Determinanten $a_iA_k - a_kA_i$, welches ganze Functionen
4\textsuperscript{ten} Grades der $a$ sind, einen gemeinschaftlichen
Theiler haben. Dass dies aber nicht möglich ist, kann man an irgend
einem speciellen Falle nachweisen, z.\,B. an der Form
\[
f = 3x_3(\alpha x_1^2 + \beta x_2^2) + a x_1^3 + b x_2^3,
\]
deren Hesse'sche Form:
\[
-\beta^2(\alpha x_3 + a x_1)x_2^2 - \alpha^2(\beta x_3 + b x_2)x_1^2
\]
ist. Hier kommen unter den $a_iA_k - a_kA_i$ z.\,B. die beiden
Functionen $a^2\beta^2, b^2\alpha^2$ vor, die keinen gemeinschaftlichen
Theiler haben.


% ===== scan leaf 359, printed 338 =====
% !PRINTED 338
\weberpage{338}
% running head: 338 | Elfter Abschnitt. | §. 91.

Um $\varDelta_{\lambda,\mu}$ für die canonische Form zu berechnen, hat man
in \S.\,90, (3):
\[
h, \qquad m
\]
durch
\[
h(\lambda-6m^2\mu), \qquad m\lambda+(h^3+2m^3)\mu
\]
zu ersetzen, und man findet so:
\[
\varDelta_{\lambda,\mu} = -\,h(\lambda-6m^2\mu)\,[m\lambda+(h^3+2m^3)\mu]^2\,\Sigma y_1^3
+ \{h^3(\lambda-6m^2\mu)^3+2[m\lambda+(h^3+2m^3)\mu]^3\}\,y_1y_2y_3,
\]
woraus nach \S.\,90, (7), (8) mit Rücksicht auf (2) nach einigen
einfachen Rechnungen folgt:
\weq{5}{\begin{gathered}
L = -2\lambda^2\mu(h^3-m^3)m - \lambda\mu^2(h^6-20h^3m^3-8m^6) \\
\qquad{}+ 8\mu^3m^2(h^3-m^3)^2, \\
M = \lambda^3+12\lambda\mu^2m(h^3-m^3)+2\mu^3(h^6-20h^3m^3-8m^6).
\end{gathered}}

Man bekommt also auf diese Weise nur zwei Invarianten vom
$4^{\mathrm{ten}}$ und $6^{\mathrm{ten}}$ Grade, nämlich in der canonischen Form:
\weq{6}{m(h^3-m^3), \quad h^6-20h^3m^3-8m^6.}

Das Gewicht dieser Invarianten ist [nach Bd. I, \S.\,60, (3)]
gleichfalls 4 und 6, und wenn wir sie also in der allgemeinen
Form mit $S$ und $T$ bezeichnen und $h=1$ setzen, so findet sich
\weq{7}{\begin{gathered}
m(1-m^3) = r^4S, \\
1-20m^3-8m^6 = r^6T.
\end{gathered}}

Für $L$ und $M$ erhält man aus (5) die Darstellung:
\weq{8}{\begin{gathered}
L = -2S\lambda^2\mu - T\lambda\mu^2 + 8S^2\mu^3 \\
M = \lambda^3+12S\lambda\mu^2+2T\mu^3.
\end{gathered}}

Die Functionen $S$ und $T$ vom $4^{\mathrm{ten}}$ und $6^{\mathrm{ten}}$ Grade sind
zuerst von \gesperrt{Aronhold} berechnet worden. Wir wollen die langen
Ausdrücke nicht hierher setzen, betrachten aber diese beiden
Grössen jetzt als bekannt und gegeben.

Aus (7) lässt sich eine biquadratische Gleichung ableiten
für das Verhältniss $m^2:r^2$. Man findet nämlich daraus:
\weq{9}{\begin{gathered}
m^5-m^8 = r^4m^4S \\
m^2-2m^5+m^8 = r^8S^2 \\
m^2-20m^5-8m^8 = r^6m^2T.
\end{gathered}}

Wenn man daraus $m^2$ und $m^5$ eliminirt, so findet man
\weq{10}{27m^8+18Sm^4r^4+Tm^2r^6-S^2r^8=0.}


% ===== scan leaf 360, printed 339 =====
% !PRINTED 339
\weberpage{339}
% running head: §. 91. | Invarianten der Curve dritter Ordnung. | 339

Diese Gleichung erhält eine einfachere Gestalt durch die
Substitution
\weq{11}{z = \frac{3m^2}{r^2}.}
Sie wird dann
\weq{12}{z^4+6Sz^2+Tz-3S^2=0.}
Diese biquadratische Gleichung hat die Eigenschaft, dass
ihre erste Invariante $A=0$ ist, und dass ihre Discriminante
(Bd. I, \S.\,47):
\weq{13}{D_2 = -27(T^2+64S^3)^2 = -27D^2,}
also stets negativ ist. Die Grösse $D=T^2+64S^3$ erhält für
die canonische Form (5) den Ausdruck:
\weq{14}{h^3(h^3+8m^3)^3,}
und kann also nicht verschwinden, so lange die Curve $f$ keinen
singulären Punkt hat. (Sie ist in der That nichts Anderes, als
die Discriminante der Form $f$.)

Aus (13) ergiebt sich nun, dass die biquadratische Gleichung
 (12) bei reellen Curven $f$ eine negative Discriminante, und
folglich (Bd. I, \S.\,78) zwei reelle und zwei conjugirt imaginäre
Wurzeln hat, und weil das Product aller vier Wurzeln $-3S^2$
negativ ist, \gesperrt{so ist von den reellen Wurzeln eine positiv},
eine negativ.

Aus $z$ kann man mit Hülfe von (11) und der ersten Gleichung
 (7) die beiden Grössen $m$ und $r$ berechnen. Man erhält so
\weq{15}{m^3 = \frac{z^2}{9S+z^2}, \quad mr^2 = \frac{3z}{9S+z^2}, \quad r^6 = \frac{27z}{(9S+z^2)^2},}
und hier könnte, wie man aus (12) schliesst, wenn man
$z^2=-9S$ setzt, $9S+z^2$ nur dann verschwinden, wenn entweder
 $S$ oder $D$ Null ist. Damit sind die Co\"efficienten $P$, $Q$, $R^3$
der cubischen Gleichung \S.\,90, (14) eindeutig durch $z$ und durch
bekannte Grössen bestimmt, und für die reelle positive Wurzel $z$
wird $m$ und $r$ reell, und folglich auch $y_1$, $y_2$, $y_3$ reell.

Eine leichte Ausnahme bildet der Fall $S=0$, in dem eine
Wurzel der Gleichung (12) verschwindet, und die zweite reelle
Wurzel $-\sqrt[3]{T}$ wird, und für die zugehörigen Werthe von $m$
und $r^6$ erhält man aus (7) und (11)
\[
m = 0, 1; \quad r^6 = \frac{1}{T}, \ \frac{-27}{T};
\]
%CONT>>


% ===== scan leaf 361, printed 340 =====
% !PRINTED 340
\weberpage{340}
% running head: 340 | Elfter Abschnitt. | §. 91.
%<<CONT
je nachdem also $T$ positiv oder negativ ist, fällt der erste oder
der zweite Werth von $r$ reell aus, und im Wesentlichen bleibt
Alles wie vorher.

Es ergiebt sich hieraus der Satz:

\gesperrt{Jede reelle ternäre cubische Form mit nicht verschwindender
 Discriminante kann durch eine reelle
Transformation auf die reelle canonische Form gebracht
 werden.}

Man kann jetzt leicht die Coordinaten der Wendepunkte in
dem canonischen Coordinatensysteme angeben. Man erhält sie,
wenn man in der Gleichung
\[
y_1^3+y_2^3+y_3^3+6m\,y_1y_2y_3=0
\]
je einmal $y_1$, $y_2$, $y_3=0$ setzt. Da eine der nicht verschwindenden
Coordinaten $=1$ gesetzt werden kann, so findet man, wenn $\varepsilon$
eine imaginäre dritte Einheitswurzel ist, die folgenden Werthe
der Coordinaten $(y_1, y_2, y_3)$:
\weq{16}{\begin{array}{ccc}
(0, 1, -1), & (0, 1, -\varepsilon), & (0, 1, -\varepsilon^2), \\
(-1, 0, 1), & (-\varepsilon, 0, 1), & (-\varepsilon^2, 0, 1), \\
(1, -1, 0), & (1, -\varepsilon, 0), & (1, -\varepsilon^2, 0),
\end{array}}
und es folgt daraus, dass \gesperrt{immer drei der Inflexionspunkte
reell und sechs imaginär sind.}

In der Tabelle (16) liegen je drei in einer Zeile stehende
Punkte auf einer Geraden, nämlich auf den Geraden
\[
y_1=0, \quad y_2=0, \quad y_3=0.
\]

Es ist danach leicht, auch die übrigen Tripel von Punkten
zusammenzustellen, die auf einer geraden Linie liegen, und damit
diese zwölf Linien zu bestimmen. Man hat dabei nur immer
solche Punkte zusammenzustellen, deren Coordinaten eine verschwindende
 Determinante geben. Man kann diese zwölf Linien
auf vier Arten in Tripel anordnen, so dass in jedem Tripel alle
neun Inflexionspunkte vorkommen. Das erste dieser Tripel wird
von den Zeilen von (16) gebildet. Die drei anderen lauten so:
\weq{17}{\begin{array}{ccc}
(0, 1, -1), & (-1, 0, 1), & (1, -1, 0), \\
(0, 1, -\varepsilon), & (-\varepsilon, 0, 1), & (1, -\varepsilon, 0), \\
(0, 1, -\varepsilon^2), & (-\varepsilon^2, 0, 1), & (1, -\varepsilon^2, 0),
\end{array}}
%CONT>>


% ===== scan leaf 362, printed 341 =====
% !PRINTED 341
\weberpage{341}
% running head: §. 91. | Realität. | 341
%<<CONT

\weq{18}{\begin{array}{ccc}
(0, 1, -1), & (-\varepsilon, 0, 1), & (1, -\varepsilon^2, 0) \\
(-1, 0, 1), & (1, -\varepsilon, 0), & (0, 1, -\varepsilon^2) \\
(1, -1, 0), & (0, 1, -\varepsilon), & (-\varepsilon^2, 0, 1),
\end{array}}
\weq{19}{\begin{array}{ccc}
(0, 1, -1), & (-\varepsilon^2, 0, 1), & (1, -\varepsilon, 0) \\
(-1, 0, 1), & (1, -\varepsilon^2, 0), & (0, 1, -\varepsilon) \\
(1, -1, 0), & (0, 1, -\varepsilon^2), & (-\varepsilon, 0, 1).
\end{array}}

Von den geraden Linien (16) enthält jede einen reellen
und zwei conjugirt imaginäre Punkte, und alle drei Linien sind
reell. Von den Geraden (17) enthält die erste die drei reellen
Punkte und ist reell, während die beiden anderen conjugirt
imaginär sind. In (18) und (19) sind alle drei Linien imaginär,
und zwar sind die Linien von (18) conjugirt zu den Linien
von (19).

Eine noch übersichtlichere Bezeichnung ist folgende. Wenn
$\xi$ und $\eta$ je ein volles Restsystem nach dem Modul 3 durchlaufen
, etwa $0, 1, 2$, und zwei nach dem Modul 3 congruente
Zahlen als nicht verschieden angesehen werden, so giebt es neun
Combinationen $(\xi, \eta)$, die wir als Bezeichnung für die neun
Wendepunkte benutzen können. Setzen wir fest, dass $(\xi, \eta)$ und
$(\eta, \xi)$ conjugirt imaginäre Wendepunkte, also $(0, 0)$, $(1, 1)$, $(2, 2)$
die drei reellen bedeuten, so erhalten wir aus (16) bis (19) folgende
 Tabelle, in der wieder die in einer Zeile stehenden Wendepunkte
 auf einer geraden Linie liegen:
\weq{20}{\begin{array}{ccc@{\qquad}ccc}
(0,0) & (1,2) & (2,1) & (0,0) & (1,1) & (2,2) \\
(1,1) & (2,0) & (0,2) & (1,2) & (2,0) & (0,1) \\
(2,2) & (0,1) & (1,0) & (2,1) & (0,2) & (1,0) \\[4pt]
(0,0) & (2,0) & (1,0) & (0,0) & (0,2) & (0,1) \\
(1,1) & (0,1) & (2,1) & (1,1) & (1,0) & (1,2) \\
(2,2) & (1,2) & (0,2) & (2,2) & (2,1) & (2,0).
\end{array}}

Man sieht hieraus, dass drei Wendepunkte $(\xi_1, \eta_1)$, $(\xi_2, \eta_2)$,
$(\xi_3, \eta_3)$ dann und nur dann auf einer geraden Linie liegen, wenn
die beiden Congruenzen
\weq{21}{\begin{aligned}
\xi_1+\xi_2+\xi_3 &\equiv 0 \\
\eta_1+\eta_2+\eta_3 &\equiv 0
\end{aligned} \quad (\mathrm{mod}\ 3)}
erfüllt sind.


% ===== scan leaf 363, printed 342 =====
% !PRINTED 342
\weberpage{342}
\pdfbookmark[2]{§ 92. Tripelgleichungen}{bm106}

Man kann die Bezeichnung der vier Systeme von geraden
Linien auch aus dem einen Schema:
\[
\begin{array}{ccc}
(0,\,0) & (1,\,2) & (2,\,1)\\
(1,\,1) & (2,\,0) & (0,\,2)\\
(2,\,2) & (0,\,1) & (1,\,0)
\end{array}
\]
ableiten, wenn man die Zeilen, die Colonnen und die je drei
den positiven und negativen Gliedern der Determinantenbildung
entsprechenden Combinationen aufstellt.

\S.\,92.

\gesperrt{Tripelgleichungen.}

Das System der neun Wendepunkte hat die Eigenschaft,
dass man aus zwei beliebigen von ihnen einen ganz bestimmten
dritten auf rationalem Wege ableiten kann. Dies drückt sich
als eine Eigenschaft der Gleichung $9^{\text{ten}}$ Grades aus, von der die
Bestimmung der Wendepunkte abhängt, die z.\,B. die Abscissen
dieser Punkte zu Wurzeln hat. Wir stellen folgende Definition
auf:

\gesperrt{Eine Gleichung ohne gleiche Wurzeln heisst eine
Tripelgleichung, wenn je zwei ihrer Wurzeln eine
dritte bestimmen, so dass jede Wurzel eines solchen
Tripels rational durch die beiden anderen ausgedrückt
werden kann, und zwar so, dass man in einer Gleichung}
\weq{1}{x_3 = \Theta\,(x_1, x_2),}
\gesperrt{in der $\Theta$ eine festgehaltene rationale Function bedeutet, für
$x_1, x_2, x_3$ die Wurzeln irgend eines Tripels
in beliebiger Reihenfolge setzen kann}\footnote{Die Tripelgleichungen $9^{\text{ten}}$ Grades sind zuerst behandelt von
\gesperrt{Hesse} in der Abhandlung: \glqq Algebraische Auflösung derjenigen Gleichung
$9^{\text{ten}}$ Grades etc.\grqq{} in Crelle's Journal, Bd.\,34 (1847).}.

Die Gleichung, von der die Wendepunkte abhängen, ist eine
Tripelgleichung $9^{\text{ten}}$ Grades. Es giebt aber auch Tripelgleichungen
von anderem, z.\,B. vom $7^{\text{ten}}$ Grade, zu denen die im vorigen
Abschnitte betrachteten gehören\footnote{Vgl. \gesperrt{N\"other}, Mathem. Annalen, Bd.\,15 (1879): \glqq Ueber die Gleichungen
$8^{\text{ten}}$ Grades und ihr Auftreten in der Theorie der Curven vierter
Ordnung.\grqq{}}.


% ===== scan leaf 364, printed 343 =====
% !PRINTED 343
\weberpage{343}

Wir wollen hier also nur auf die Tripelgleichungen $9^{\text{ten}}$ Grades
eingehen, und zunächst zeigen, wie sich aus der Tripeleigenschaft
die Bezeichnung und Anordnung der Wurzeln ableiten lässt, die
wir im vorigen Paragraphen für die Wendepunkte kennen gelernt
haben.

Jede der neun Wurzeln kommt in vier Tripeln vor, und im
Ganzen existiren daher $4.9:3=12$ verschiedene Tripel. Nehmen
wir ein beliebiges von diesen Tripeln und bezeichnen dessen
Wurzeln mit
\[
(0,\,0)\ (1,\,2)\ (2,\,1).
\]
Dann sei $(1,\,1)$ eine beliebige vierte Wurzel. Diese bestimmt
mit den drei ersten zunächst drei Tripel, die wir mit
\[
(1,\,1)\ (0,\,0)\ (2,\,2), \quad (1,\,1)\ (1,\,2)\ (1,\,0), \quad (1,\,1)\ (2,\,1)\ (0,\,1)
\]
bezeichnen, und es bleiben noch zwei Wurzeln übrig, die das
vierte Tripel mit $(1,\,1)$ bilden, und das wir nun mit
\[
(1,\,1)\ (2,\,0)\ (0,\,2)
\]
bezeichnen.

Wir suchen nun das durch $(0,\,0), (2,\,0)$ bestimmte Tripel.
In diesem können keine Wurzeln vorkommen, die mit einer der
beiden Wurzeln $(0,\,0), (2,\,0)$ in einer der schon bestimmten Tripel % [sic] page prints feminine "einer" with neuter "Tripel"; expected "einem"
vorkommen, also nicht
\[
(1,\,2)\ (2,\,1)\ (1,\,1)\ (2,\,2)\ (0,\,2),
\]
und es bleibt für die dritte Wurzel dieses Tripels nur $(0,\,1)$
oder $(1,\,0)$ übrig, und dieselben beiden Wurzeln bleiben auch nur
übrig für das durch $(0,\,0)\,(0,\,2)$ bestimmte Tripel. Da wir aber
noch $(0,\,2)$ mit $(2,\,0)$ vertauschen können, so beschränken wir
die Allgemeinheit nicht, wenn wir die zwei weiteren Tripel
\[
(0,\,0)\ (2,\,0)\ (1,\,0), \quad (0,\,0)\ (0,\,2)\ (0,\,1)
\]
annehmen. Nun bestimmen wir die Tripel, die $(1,\,0)$ enthalten,
von denen zwei schon bekannt sind. Die beiden anderen
müssen die Wurzeln $(2,\,2), (0,\,1), (0,\,2), (2,\,1)$ enthalten. Da aber
$(0,\,1)\,(1,\,0)\,(0,\,2)$ und $(0,\,1)\,(1,\,0)\,(2,\,1)$ keine Tripel sind [weil
$(0,\,1)\,(0,\,2)$ das Tripel $(0,\,0)\,(0,\,2)\,(0,\,1)$ und $(0,\,1)\,(2,\,1)$ das
Tripel $(1,\,1)\,(2,\,1)\,(0,\,1)$ bestimmt], so haben wir die weiteren
Tripel
\[
(2,\,2)\ (0,\,1)\ (1,\,0), \quad (1,\,0)\ (0,\,2)\ (2,\,1).
\]
Ebenso ergiebt sich
\[
(0,\,1)\ (2,\,0)\ (1,\,2).
\]


% ===== scan leaf 365, printed 344 =====
% !PRINTED 344
\weberpage{344}
\pdfbookmark[2]{§ 93. Die Gruppe der Tripelgleichungen}{bm107}

Endlich bilden wir noch je ein $(0,\,2)$ und $(2,\,0)$ enthaltendes
Tripel
\[
(0,\,2)\ (1,\,2)\ (2,\,2), \quad (2,\,0)\ (2,\,1)\ (2,\,2),
\]
womit alle zwölf Tripel, und damit die Anordnung \S.\,91, (20)
bestimmt sind.

\S.\,93.

\gesperrt{Die Gruppe der Tripelgleichungen.}

Es soll nun die \gesperrt{Galois}'sche Gruppe $P$ der Tripelgleichungen
bestimmt werden. Zunächst ist klar, dass in dieser Gruppe nur
solche Permutationen vorkommen, bei denen jedes Tripel wieder
in ein Tripel übergeht. Denn nehmen wir an, die Wurzeln eines
Tripels $x_1, x_2, x_3$ gehen durch eine Permutation in $y_1, y_2, y_3$
über, so folgt, da diese Permutation auf die Gleichung (1) (\S.\,92)
anwendbar ist,
\[
y_1 = \Theta\,(y_2, y_3),
\]
und folglich ist $y_1$ die dritte Wurzel des durch $y_2, y_3$ bestimmten
Tripels.

Gehen wir nun zu der Bezeichnung für die Wurzeln:
\weq{1}{x = (\xi, \eta)}
über, so ergiebt sich aus der Zusammenstellung der Tripel in
den §§.\,91 und 92, dass drei Wurzeln
\weq{2}{x_1 = (\xi_1, \eta_1), \quad x_2 = (\xi_2, \eta_2), \quad x_3 = (\xi_3, \eta_3)}
immer dann und nur dann ein Tripel bilden, wenn
\weq{3}{\begin{array}{l}
\xi_1 + \xi_2 + \xi_3 \equiv 0\\
\eta_1 + \eta_2 + \eta_3 \equiv 0
\end{array} \quad (\mathrm{mod}\ 3).}
Nennen wir eine Permutationsgruppe der $(\xi, \eta)$, deren Substitutionen
 alle der vollständigen linearen Congruenzgruppe für
den Modul 3 angehören (\S.\,78), eine \gesperrt{lineare Gruppe}, so gilt
der Satz:

\gesperrt{1. Die Gruppe $P$ einer Tripelgleichung ist immer
linear.}

Nach \S.\,77 lässt sich jede Permutation der Grössen $(\xi, \eta)$
überhaupt in der Form darstellen:
\weq{4}{\xi' \equiv \varphi\,(\xi, \eta), \quad \eta' \equiv \psi\,(\xi, \eta) \quad (\mathrm{mod}\ 3),}
%CONT>>


% ===== scan leaf 366, printed 345 =====
% !PRINTED 345
\weberpage{345}

%<<CONT
worin $\varphi, \psi$ ganze ganzzahlige Functionen der Variablen $\xi, \eta$
sind, deren Grad in Bezug auf keine der Variablen den Werth 2
übersteigt. Ordnen wir diese Functionen nach $\eta$, so können
wir setzen:
\weq{5}{\xi' = A\eta^2 + B\eta + C, \quad \eta' = A'\eta^2 + B'\eta + C',}
worin $A, B, C, A', B', C'$ ganze Functionen von $\xi$, höchstens
vom $2^{\text{ten}}$ Grade sind. Wir wollen in den Gleichungen (4) die
Variable $\xi$ festhalten und $\eta = 0, 1, 2$ setzen, d.\,h. wir wenden
(5) auf das Tripel $(\xi, 0)\,(\xi, 1)\,(\xi, 2)$ an. Die entsprechenden
$(\xi_1', \eta_1'), (\xi_2', \eta_2'), (\xi_3', \eta_3')$ müssen dann den Relationen (3)
genügen, also:
\[
\begin{array}{ll}
\xi_1' \equiv C, & \eta_1' \equiv C',\\
\xi_2' \equiv A + B + C, & \eta_2' \equiv A' + B' + C',\\
\xi_3' \equiv A - B + C, & \eta_3' \equiv A' - B' + C',
\end{array}
\]
und daraus folgt, dass $A$ und $A'$ für jedes $\xi$ congruent mit
Null sein müssen.

Ebenso kann man schliessen, dass in den Substitutionen (4)
kein Glied mit $\xi^2$ vorkommen kann, und dass daher diese Substitutionen
 die Form haben müssen:
\[
\begin{array}{l}
\xi' \equiv m\,\xi\,\eta + a\,\xi + b\,\eta + c\\
\eta' \equiv m'\,\xi\,\eta + a'\,\xi + b'\,\eta + c'
\end{array} \quad (\mathrm{mod}\ 3).
\]
Wenn man aber diese Substitutionen auf das Tripel $(0,\,0)\,(1,\,1)\,(2,\,2)$
anwendet, und die Summe der entsprechenden $\xi'$ und $\eta'$ gleich
Null setzt, so folgt, dass $m = m' = 0$ sein muss, und dass also
jede Permutation der Gruppe $P$ in der Form
\weq{6}{\begin{array}{l}
\xi' \equiv a\,\xi + b\,\eta + c\\
\eta' \equiv a'\,\xi + b'\,\eta + c'
\end{array} \quad (\mathrm{mod}\ 3)}
enthalten sein muss, wie behauptet war.

\gesperrt{2. Umgekehrt gehen durch jede lineare Substitution
von der Form (6) drei Grössen $(\xi, \eta)$ eines Tripels
in drei $(\xi', \eta')$ über, die gleichfalls ein Tripel
bilden.}

Denn aus (6) folgt für irgend drei Zahlenpaare $(\xi_1, \eta_1)$,
$(\xi_2, \eta_2), (\xi_3, \eta_3)$:
\weq{7}{\begin{array}{l}
\xi_1' + \xi_2' + \xi_3' \equiv a\,(\xi_1 + \xi_2 + \xi_3) + b\,(\eta_1 + \eta_2 + \eta_3)\\
\eta_1' + \eta_2' + \eta_3' \equiv a'\,(\xi_1 + \xi_2 + \xi_3) + b'\,(\eta_1 + \eta_2 + \eta_3)
\end{array} \quad (\mathrm{mod}\ 3);}
%CONT>>


% ===== scan leaf 367, printed 346 =====
% !PRINTED 346
\weberpage{346}
%<<CONT
also wenn $\xi_1 + \xi_2 + \xi_3$ und $\eta_1 + \eta_2 + \eta_3$ mit Null congruent
ist, so gilt dasselbe von den linken Seiten von (7).

Nun gehen auch umgekehrt durch jede lineare Substitution der Form (6) immer drei
Grössen ($\xi, \eta$) eines Tripels in ein Tripel über, und daraus lässt sich beweisen,
dass jede Gleichung 9$^{\mathrm{ten}}$ Grades, deren Galois'sche Gruppe $P$ in der
vollständigen linearen Congruenzgruppe enthalten ist, eine Tripelgleichung ist, wenn
noch die weitere Bedingung hinzukommt, dass durch die Gruppe $P$ irgend zwei gegebene
Wurzeln in zwei andere gleichfalls beliebig gegebene Wurzeln übergehen, oder, was damit
gleichbedeutend ist, wenn $P$ \gesperrt{zweifach transitiv} ist.

Denn wenn $P$ aus lauter linearen Substitutionen besteht, so reducirt sich $P$ durch
Adjunction zweier beliebiger Wurzeln $x_1 = (\xi_1, \eta_1)$, $x_2 = (\xi_2, \eta_2)$
auf eine Gruppe $P_1$, die nicht nur $x_1, x_2$, sondern auch die dritte Wurzel
$x_3 = (\xi_3, \eta_3)$ des durch $x_1, x_2$ bestimmten Tripels ungeändert lässt, und
folglich gehört $x_3$ dem neuen Rationalitätsbereiche an, und mithin ist $x_3$ rational
durch $x_1, x_2$ ausdrückbar in der Form

\weq{8}{x_3 = \Theta (x_1, x_2).}

Wenn nun die Gruppe $P$ zweifach transitiv ist, so giebt es eine Permutation in $P$,
durch die $x_1, x_2$ in zwei beliebige andere Wurzeln $y_1, y_2$ übergehen, und durch
die folglich auch $x_3$ in die dritte Wurzel des Tripels $y_1, y_2, y_3$ übergeht, und
diese Permutation ist auf (8) anwendbar. Also ist auch

\[
y_3 = \Theta (y_1, y_2).
\]

Wir sprechen dies als Satz aus:

3. \gesperrt{Jede Gleichung 9$^{\mathrm{ten}}$ Grades, deren Gruppe linear
und zweifach transitiv ist, ist eine Tripelgleichung.}

Durch die Permutationen einer zweifach transitiven linearen Gruppe kann jedes Tripel
in jedes andere, und zwar in beliebiger Anordnung, übergeführt werden.

Die Voraussetzung der zweifachen Transitivität kann auch so ausgedrückt werden, dass
alle die Permutationen aus $P$, die eine Wurzel ungeändert lassen, die übrigen noch
transitiv mit einander verbinden. Bezeichnen wir mit $P_0$ die in $P$ enthaltene
Gruppe, durch deren Permutationen eine der Wurzeln $x_0$ ungeändert bleibt, so muss
$P_0$ in Bezug auf die anderen Wurzeln noch transitiv sein. Ist aber die Gruppe $P$
nur einfach transitiv,
%CONT>>


% ===== scan leaf 368, printed 347 =====
% !PRINTED 347
\weberpage{347}
%<<CONT
so ist $P_0$ nicht mehr transitiv, d.\,h. nach Adjunction von $x_0$
zerfällt die Gleichung für die übrigen acht Wurzeln in mehrere Factoren. Wenn die
Gruppe $P$ überhaupt transitiv ist, so kann hierbei $x_0$ jede beliebige der Wurzeln
sein; denn wenn $\pi$ eine Permutation von $P$ ist, durch die $x_0$ in $x_1$ übergeht,
so bleibt durch $P_1 = \pi^{-1} P_0 \pi$ die Wurzel $x_1$ ungeändert, und wenn $P_0$
für acht Wurzeln intransitiv ist, so muss auch $P_1$ für acht Wurzeln intransitiv sein,
da sich die Permutationen von $P_1$ von denen von $P_0$ nur durch die Bezeichnung der
Wurzeln unterscheiden (Bd. I, §.\,154, 6.).

Die einfach transitiven unter den linearen Gruppen haben einen specielleren Charakter
als die Gruppe der Tripelgleichungen.

Wir haben im §.\,78 eine Zerlegung der allgemeinen linearen Congruenzgruppe für den
Modul 3 kennen gelernt, an der wir diese Verhältnisse übersehen können.

Die \gesperrt{Abel'sche} Gruppe $Q$, die aus den Substitutionen

\weq{9}{\xi' \equiv \xi + \alpha, \quad \eta' \equiv \eta + \beta \quad (\mathrm{mod}\ 3)}

besteht, ist gewiss nur einfach transitiv; denn wenn eine Wurzel durch (9) ungeändert
bleiben soll, so muss $\alpha \equiv 0$, $\beta \equiv 0$ sein und alle Wurzeln bleiben
ungeändert. Jede Wurzel ist rational durch jede andere ausdrückbar.

Suchen wir nun in der Gruppe $P$ die Gruppe $P_0$ der Substitutionen auf, durch die
$(0,0)$ ungeändert bleibt, so enthält $P_0$ nur homogene Substitutionen und ist also
der grösste gemeinschaftliche Theiler von $P$ und der homogenen Congruenzgruppe $S$,
und es ist $P = P_0 Q$.

Bilden wir für $S$ die Compositionsreihe (§.\,78):

\[
S,\ S_1,\ S_2,\ S_3,\ S_4,
\]

so ist $S_4 = \begin{pmatrix} 1, & 0 \\ 0, & 1 \end{pmatrix}, \begin{pmatrix} -1, & 0 \\ 0, & -1 \end{pmatrix}$,
und durch $S_4$ gehen die Wurzeln $(1,1)$, $(2,2)$ nur in einander über. Durch $S_4 Q$
kann das Tripel $(0,0)$ $(1,1)$ $(2,2)$ zwar jede Anordnung seiner Wurzeln erfahren.
Es kann aber nur noch in die beiden anderen Tripel $(1,0)$ $(2,1)$ $(0,2)$;
$(0,1)$ $(1,2)$ $(2,0)$ übergehen. $S_4 Q$ ist also nur \gesperrt{einfach transitiv} und
ist überdies \gesperrt{imprimitiv}.

Durch die Gruppe $S_3$, die aus $S_4$ durch Hinzunahme der
Substitution

\[
\begin{pmatrix} -1, & 1 \\ 1, & 1 \end{pmatrix}
\]

%CONT>>


% ===== scan leaf 369, printed 348 =====
% !PRINTED 348
\weberpage{348}
\pdfbookmark[2]{§ 94. Realitätsverhältnisse der Tripelgleichungen}{bm108}
%<<CONT
entsteht, kann $(1,1)$, $(2,2)$ noch in $(0,2)$, $(0,1)$, aber nicht in
$(2,0)$, $(1,0)$ oder in $(1,2)$, $(2,1)$ übergehen. Also ist auch die Gruppe $S_3 Q$ noch nicht
zweifach transitiv, wohl aber ist sie primitiv. Nehmen wir aber, um $S_2$
zu bilden, noch die Substitution

\[
\begin{pmatrix} 0, & 1 \\ -1, & 0 \end{pmatrix}
\]

hinzu, so geht $(1,1)$, $(2,2)$ in $(0,2)$, $(0,1)$ und in $(1,2)$, $(2,1)$ und
$(2,0)$, $(1,0)$ über. Also ist die Gruppe $S_2 Q$ zweifach transitiv und folglich
auch $S_1 Q$ und $S Q$. Es fangen also erst bei $S_2 Q$ die Tripelgleichungen an.
Die Gleichungen mit engeren Gruppen sind noch keine eigentlichen Tripelgleichungen.

Wir können noch andere in $S Q$ enthaltene lineare Gruppen bilden, z.\,B. wenn wir
mit $S'$ die cyklische Gruppe 3$^{\mathrm{ten}}$ Grades:

\[
S' = \begin{pmatrix} 1, & 0 \\ 0, & 1 \end{pmatrix}, \begin{pmatrix} 1, & 0 \\ 1, & 1 \end{pmatrix}, \begin{pmatrix} 1, & 0 \\ -1, & 1 \end{pmatrix}
\]

bezeichnen, die Gruppe 27$^{\mathrm{sten}}$ Grades $S' Q = Q S'$. Durch die Gruppe $S'$
geht $(1,1)$ $(2,2)$ in $(1,1)$ $(2,2)$, $(1,2)$ $(2,1)$, $(1,0)$ $(2,0)$ über, während
$(0,1)$ $(0,2)$ nicht daraus entsteht. Also ist auch dies nicht die Gruppe einer
Tripelgleichung.

Hier wäre noch ein Feld weitergehender Untersuchungen über diese speciellen
metacyklischen Gleichungen 9$^{\mathrm{ten}}$ Grades und die Darstellung ihrer
Wurzeln nach den Principien von \gesperrt{Kronecker} (Bd. I, achtzehnter Abschnitt).

\begin{center}
\S.\,94.\\[2pt]
\gesperrt{Realitätsverhältnisse der Tripelgleichungen.}
\end{center}

Wenn der Rationalitätsbereich reell ist, so können wir über die Realität der Wurzeln
einer Tripelgleichung 9$^{\mathrm{ten}}$ Grades einige allgemeine Schlüsse machen.
Wenn in diesem Falle nämlich in einem Tripel zwei reelle Wurzeln vorkommen, so muss,
wie aus der Gleichung $x_3 = \Theta (x_1, x_2)$ folgt, auch die dritte Wurzel reell
sein. Wenn also überhaupt eine imaginäre Wurzel vorhanden ist, so muss in jedem der
vier Tripel, dem diese Wurzel angehört, mindestens noch eine zweite imaginäre Wurzel
vorkommen, und folglich ist die Anzahl der imaginären Wurzeln mindestens gleich 5,
oder, da die Zahl der imaginären Wurzeln
%CONT>>


% ===== scan leaf 370, printed 349 =====
% !PRINTED 349
\weberpage{349}
%<<CONT
gerade sein muss, gleich 6. Zwei reelle oder zwei conjugirt imaginäre Wurzeln
$x_1, x_2$ gehören immer einem Tripel an, dessen dritte Wurzel reell ist; dies
ergiebt sich aus der Gleichung:

\[
x_3 = \Theta (x_1, x_2) = \Theta (x_2, x_1).
\]

Wir wollen demnach ein Tripel, was zwei reelle oder zwei conjugirt imaginäre Wurzeln
enthält, ein \gesperrt{reelles Tripel} nennen. Wenn in einem Tripel irgend imaginäre
Wurzeln vorkommen, so bilden auch die conjugirt imaginären Wurzeln ein Tripel, und
zwei solche Tripel sollen \gesperrt{conjugirt imaginäre Tripel} heissen.

Es giebt dann drei Möglichkeiten:

1. \gesperrt{Lauter reelle Wurzeln.}

2. \gesperrt{Eine reelle und acht imaginäre Wurzeln.}

In diesem Falle müssen die vier Paare conjugirt imaginärer Wurzeln vier Tripel
bestimmen, die alle dieselbe reelle Wurzel als dritte enthalten. Bezeichnen wir
diese reelle Wurzel mit $(0,0)$, so müssen

\[
(1,2)\ (2,1),\quad (1,1)\ (2,2),\quad (1,0)\ (2,0),\quad (0,1)\ (0,2)
\]

die vier Paare conjugirter Wurzeln sein.

3. \gesperrt{Drei reelle und sechs imaginäre Wurzeln.}

Die drei reellen Wurzeln müssen hier ein Tripel bilden. Nehmen wir für die reellen
Wurzeln

\weq{1}{(0,0)\ (1,1)\ (2,2),}

so können in den beiden Reihen

\weq{2}{\begin{aligned} &(1,2)\ (0,2)\ (1,0) \\ &(2,1)\ (2,0)\ (0,1) \end{aligned}}

conjugirt imaginäre Wurzeln nicht in derselben Reihe vorkommen, weil sonst die
dritte Wurzel der betreffenden Reihe reell wäre.

Es ist zu zeigen, dass die durch die conjugirten Paare bestimmten drei reellen
Tripel zusammen alle drei reellen Punkte enthalten müssen.

Nehmen wir nämlich an, unter den drei Tripeln

\[
(0,0)\ (1,2)\ (2,1),\quad (0,0)\ (2,0)\ (1,0),\quad (0,0)\ (0,1)\ (0,2)
\]

kommen zwei reelle vor, so muss auch das dritte reell sein, d.\,h. es müssen

\[
(1,2)\ (2,1),\quad (2,0)\ (1,0),\quad (0,1)\ (0,2)
\]
%CONT>>


% ===== scan leaf 371, printed 350 =====
% !PRINTED 350
\weberpage{350}
%<<CONT
conjugirte Paare sein. Dann aber bilden die drei Wurzeln $(1,1)$ $(2,0)$ $(0,2)$ ein Tripel, deren conjugirte $(1,1)$ $(1,0)$ $(0,1)$ kein Tripel bilden, was unmöglich ist. Jede der drei reellen Wurzeln $(0,0)$, $(1,1)$, $(2,2)$ kommt also ausser in (1) noch in einem und nur in einem reellen Tripel vor, dessen beide anderen Wurzeln conjugirt imaginär sind, und der dritte Fall führt also zu der Anordnung der reellen und imaginären Wurzeln, die wir bei den Wendepunkten der Curve dritter Ordnung kennen gelernt haben.

Dass aber auch die Fälle 1. und 2. vorkommen können, lehrt folgende Betrachtung: Aus einer allgemeinen Gleichung $9^{\text{ten}}$ Grades erhält man eine Tripelgleichung durch Adjunction einer zu der linearen Gruppe $SQ$ gehörigen Function. Eine solche Function ist z.\,B.:

\[
\begin{aligned}
v &= (0,0)(1,1)(2,2) + (0,0)(1,2)(2,1) + (0,0)(1,0)(2,0) \\
&\quad + (0,0)(0,1)(0,2) + (1,1)(1,2)(1,0) + (1,1)(2,1)(0,1) \\
&\quad + (1,1)(0,2)(2,0) + (1,2)(0,1)(2,0) + (2,2)(2,1)(2,0) \\
&\quad + (2,2)(1,2)(0,2) + (2,2)(0,1)(1,0) + (2,1)(0,2)(1,0),
\end{aligned}
\]

und diese Function ist reell in den Fällen 1., 2., 3. Also ist auch der erweiterte Rationalitätsbereich, in dem unsere Gleichung eine Tripelgleichung ist, reell.

\begin{center}
\rule{0.22\textwidth}{0.4pt}
\end{center}


% ===== scan leaf 372, printed 351 =====
\weberpage{351}% page number not printed on this page
\pdfbookmark[1]{Zwölfter Abschnitt. Die Doppeltangenten einer Curve vierter Ordnung.}{bm109}
\pdfbookmark[2]{§ 95. Anzahl der Doppeltangenten einer Curve vierter Ordnung}{bm110}
% !PRINTED none
% Abschnitt-opening page: no running head/page number printed on this recto.
\begin{center}
\gesperrt{Zwölfter Abschnitt.}

\textbf{Die Doppeltangenten einer Curve vierter Ordnung.}

\rule{0.3\linewidth}{0.4pt}

§. 95.

\gesperrt{Anzahl der Doppeltangenten einer Curve\\
vierter Ordnung.}
\end{center}

Ein geometrisches Problem, das wegen seiner mannigfachen Beziehungen zu anderen Gebieten von besonderer Bedeutung ist, was zugleich in ähnlicher Weise, wie das Problem der Wendepunkte der Curven dritter Ordnung, merkwürdige algebraische Verhältnisse bietet, ist das der Doppeltangenten der Curven vierter Ordnung.

Unter einer \gesperrt{Doppeltangente} einer Curve versteht man, wie schon im \S.\,86 bemerkt ist, eine gerade Linie, die die Curve in zwei verschiedenen Punkten berührt. Bei Curven von niedrigerer als der vierten Ordnung können Doppeltangenten nicht auftreten.

Bei den Curven vierter Ordnung hat eine Doppeltangente ausser den Berührungspunkten keinen Punkt mit der Curve gemein. In besonderen Fällen können die beiden Berührungspunkte auch zusammenfallen. Dann haben wir Linien mit vierpunktiger Berührung.

Die Bestimmung der Doppeltangenten wird als algebraisches Problem von einer gewissen algebraischen Gleichung abhängen, deren Grad gleich der Anzahl der Doppeltangenten ist, und die erste Frage ist die nach dem Grade dieser Gleichung, also nach der Anzahl der Doppeltangenten. Wir beschränken uns hier auf die Betrachtung von Curven vierter Ordnung, obwohl der Weg, den wir gehen, auch auf Curven höherer Ordnung anwendbar
%CONT>>


% ===== scan leaf 373, printed 352 =====
% !PRINTED 352
\weberpage{352}
%<<CONT
ist. Wir setzen auch voraus, dass die Curve vierter Ordnung frei von singulären Punkten sei\footnote{\gesperrt{Jacobi}, \glqq Ueber die Anzahl der Doppeltangenten ebener algebraischer Curven\grqq{}, \gesperrt{Crelle's} Journal, Bd.\,40 (1850). \gesperrt{Clebsch}, \glqq Bemerkung zu \gesperrt{Jacobi's} Beweis für die Anzahl der Doppeltangenten\grqq{}, ebend., Bd.\,63 (1864). Aus der ziemlich umfassenden Literatur über die Theorie der Doppeltangenten einer Curve vierter Ordnung sind die folgenden Schriften hervorzuheben. Zunächst mehrere Arbeiten von \gesperrt{Hesse} in \gesperrt{Crelle's} Journal, Bd.\,41, 49, 52. Ferner \gesperrt{Steiner}, Eigenschaften der Curven vierter Ordnung rücksichtlich ihrer Doppeltangenten, \gesperrt{Crelle's} Journal, Bd.\,49. \gesperrt{Aronhold}, Monatsber. d. Berl. Akademie 1864. \gesperrt{Geiser}, \glqq Ueber die Doppeltangenten einer ebenen Curve 4\,ten Grades\grqq{}, Math. Annalen, Bd.\,1 (1868). \gesperrt{Cayley}, \gesperrt{Crelle's} Journal, Bd.\,68 (1868). Auch das oben citirte Werk von \gesperrt{Salmon} ist hier hervorzuheben. In neuerer Zeit wurde die Theorie der Doppeltangenten im Zusammenhange mit der Theorie der \gesperrt{Abel'schen} Functionen weiter ausgebildet. \gesperrt{Riemann}, Gesammelte mathematische Werke (2.\,Aufl., 1892) XXXI, aus dem Nachlass. \gesperrt{Weber}, Theorie der \gesperrt{Abel'schen} Functionen vom Geschlecht 3, Berlin 1876. \gesperrt{Frobenius}, \gesperrt{Crelle's} Journal, Bd.\,99 u.\,103. Die algebraische Seite der Frage ist in zwei Abhandlungen: \gesperrt{Nöther}, Math. Annalen, Bd.\,15 und \gesperrt{Weber}, ebend., Bd.\,23 behandelt.}. Es möge jetzt $f(x_1, x_2, x_3)$ oder kürzer $f(x)$ eine ternäre Form 4\,ten Grades sein, die, gleich Null gesetzt, eine Curve vierter Ordnung ohne singulären Punkt darstellt, die wir die Curve $f$ nennen.

Setzen wir in dieser Gleichung
\weq{1}{f(x_1, x_2, x_3) = 0}
an Stelle der Variablen $x_1, x_2, x_3$ Ausdrücke von der Form

\[
x_1+t y_1, \quad x_2+t y_2, \quad x_3+t y_3,
\]

so ergiebt sich eine Gleichung 4\,ten Grades in Bezug auf $t$:
\weq{2}{f(x+ty) = 0,}
deren Wurzeln, wenn $(x)$ und $(y)$ zwei feste Punkte sind, in $x+ty$ eingesetzt, die Coordinaten der Schnittpunkte der Verbindungslinie der Punkte $(x), (y)$, die wir als die Linie $(x,y)$ bezeichnen wollen, mit der Curve $f$ geben.

Es werde nun die Function $f(x+ty)$ nach Potenzen von $t$ geordnet. Dies giebt (nach Bd.\ I, \S.\,60):
\weq{3}{f(x+yt) = P_0(x,y)+4tP_1(x,y)+6t^2P_2(x,y)+4t^3P_3(x,y)+t^4P_4(x,y),}
worin die $P_\nu(x,y)$ die Polaren von $f(x)$ sind, also homogene
%CONT>>


% ===== scan leaf 374, printed 353 =====
% !PRINTED 353
\weberpage{353}
%<<CONT
Functionen $\nu^{\text{ter}}$ Ordnung in Bezug auf $y$, und $(4-\nu)^{\text{ter}}$ Ordnung in Bezug auf $x$. Es ist nämlich

\weq{4}{
\begin{aligned}
P_0(x,y) &= f(x) \\
P_1(x,y) &= \tfrac{1}{4}\sum_i y_i f'(x_i) \\
P_2(x,y) &= \tfrac{1}{12}\sum_i y_i y_k f''(x_i x_k) \\
P_3(x,y) &= \tfrac{1}{4}\sum x_i f'(y_i) \\
P_4(x,y) &= f(y).
\end{aligned}
}

Wenn $x$ auf der Curve $f$ liegt, so wird $P_0=0$, und eine Wurzel der Gleichung (2) verschwindet. Nehmen wir ausserdem noch $(y)$ auf der Tangente im Punkte $(x)$ an, so ist auch $P_1(x,y)=0$, und es verschwinden zwei Wurzeln von (2). Die beiden übrigen werden nach (3) durch die quadratische Gleichung
\weq{5}{6P_2+4tP_3+t^2P_4=0}
bestimmt. Wenn diese Gleichung zwei gleiche Wurzeln hat, so wird die Linie $(x,y)$ noch einen zweiten Berührungspunkt mit der Curve $f$ haben, d.\,h. sie wird Doppeltangente sein. Die Bedingung hierfür ist aber die, dass die Discriminante der Gleichung (5) verschwindet, oder dass
\weq{6}{R(x,y) = 2P_3^2 - 3P_4P_2}
gleich Null sei. Hierin ist $R(x,y)$ eine in Bezug auf jedes der beiden Systeme $(x)$ und $(y)$ homogene Function, und zwar für die $x$ vom $2^{\text{ten}}$, für die $y$ vom $6^{\text{ten}}$ Grade. Es ist aber noch zu bemerken, dass die Gleichung $R=0$ auch dann erfüllt ist, wenn $(x)$ ein beliebiger Punkt der Curve ist, und $(y)$ mit $(x)$ zusammenfällt, weil dann

\[
P_2(x,x)=P_3(x,x)=P_4(x,x)=f(x)=0
\]

wird. Sonst aber wird $R$ nur dann verschwinden, wenn die Linie $(x,y)$ eine Doppeltangente ist. Wir stellen also den Satz auf:

1. \gesperrt{Die Gleichung $R(x,y)=0$ ist, wenn $(x)$ ein Punkt der Curve $f$ und $(y)$ ein von $(x)$ verschiedener Punkt der Tangente in $(x)$ ist, die nothwendige und hinreichende Bedingung dafür, dass $(x)$ ein Berührungspunkt einer Doppeltangente sei.}

Es kommt nun darauf an, aus dieser Bedingung eine andere herzuleiten, die nur die $x$ allein enthält, und die für keine
%CONT>>


% ===== scan leaf 375, printed 354 =====
% !PRINTED 354
\weberpage{354}
%<<CONT
anderen Punkte als die Berührungspunkte der Doppeltangenten
befriedigt ist.

Die Variablen $y$ genügen der Tangentengleichung

\weq{7}{y_1 f'(x_1) + y_2 f'(x_2) + y_3 f'(x_3) = 0.}

Bezeichnen wir mit $b_1, b_2, b_3$ irgend drei willkürliche Constanten, und setzen

\weq{8}{b_x = b_1 x_1 + b_2 x_2 + b_3 x_3,}

so dass $b_x = 0$ die Gleichung einer willkürlichen geraden Linie ist, so können wir zu (7) noch die Gleichung

\weq{9}{b_1 y_1 + b_2 y_2 + b_3 y_3 = 0}

hinzunehmen, also $(y)$ als den Schnittpunkt der Tangente in $(x)$ mit der beliebigen geraden Linie $b_x$ auffassen. Dann erhalten wir

\weq{10}{\begin{aligned}
y_1 &= b_2 f'(x_3) - b_3 f'(x_2),\\
y_2 &= b_3 f'(x_1) - b_1 f'(x_3),\\
y_3 &= b_1 f'(x_2) - b_2 f'(x_1),
\end{aligned}}

und hierdurch ist die Bedingung (7) identisch befriedigt. Durch diese Substitution geht $R(x,y)$ in eine homogene Function $20^{\text{sten}}$ Grades der $x$ und $6^{\text{ten}}$ Grades der $b$ über, die wir mit $D(x,b)$ bezeichnen wollen. Diese Function $D(x,b)$ verschwindet aber ausser in den Berührungspunkten der Doppeltangenten noch in den vier Schnittpunkten der Geraden $b_x$ mit der Curve $f$. Diese Bedingung muss nun noch so umgeformt werden, dass sie von den $b$ unabhängig wird.

Gehen wir von dem Punkte $(y)$ zu einem beliebigen anderen Punkte der Tangente über, so können wir dies dadurch erreichen, dass wir $y$ durch $\mu x + \lambda y$ ersetzen, worin $\lambda, \mu$ zwei willkürliche Parameter bedeuten. Dadurch geht $f(x+ty)$ in

\[
f[(1+\mu t)x + \lambda t y] = (1+\mu t)^4 f\left(x + \frac{\lambda t}{1+\mu t}\, y\right)
\]

über, und die Entwickelung (3) ergiebt, wenn wir

\[
P_0(x,y),\ P_1(x,y),\ P_1(x,\mu x + \lambda y)
\]

gleich Null setzen:

\[
\begin{aligned}
&6t^2(1+\mu t)^2\lambda^2 P_2(x,y) + 4t^3(1+\mu t)\lambda^3 P_3(x,y) + t^4\lambda^4 P_4(x,y)\\
&= 6t^2 P_2(x,\mu x + \lambda y) + 4t^3 P_3(x,\mu x + \lambda y) + t^4 P_4(x,\mu x + \lambda y),
\end{aligned}
\]
%CONT>>


% ===== scan leaf 376, printed 355 =====
% !PRINTED 355
\weberpage{355}
%<<CONT
also, wenn wir beiderseits nach Potenzen von $t$ ordnen:

\[
\begin{aligned}
P_2(x,\mu x + \lambda y) &= \lambda^2 P_2(x,y)\\
P_3(x,\mu x + \lambda y) &= 3\lambda^2\mu\, P_2(x,y) + \lambda^3 P_3(x,y)\\
P_4(x,\mu x + \lambda y) &= 6\lambda^2\mu^2 P_2(x,y) + 4\lambda^3\mu\, P_3(x,y) + \lambda^4 P_4(x,y),
\end{aligned}
\]

und hieraus erhält man nach (6)

\weq{11}{R(x,\mu x+\lambda y) = \lambda^6 R(x,y).}

Diese Formel gilt aber nicht identisch, sondern nur unter der Voraussetzung, dass $(x)$ ein Punkt der Curve sei, also dass $f(x) = 0$ ist.

Der Punkt $(\mu x+\lambda y)$ kann ein ganz beliebiger Punkt der Tangente in $(x)$ sein, und daher können wir, wenn $a_1, a_2, a_3$ drei willkürliche Grössen sind, über $\mu, \lambda$ so verfügen, dass

\weq{12}{\begin{aligned}
\mu x_1 + \lambda y_1 &= a_2 f'(x_3) - a_3 f'(x_2)\\
\mu x_2 + \lambda y_2 &= a_3 f'(x_1) - a_1 f'(x_3)\\
\mu x_3 + \lambda y_3 &= a_1 f'(x_2) - a_2 f'(x_1)
\end{aligned}}

wird. Dann ist $(\mu x+\lambda y)$ der Schnittpunkt der Tangente mit der Geraden $a_x = 0$, wenn

\weq{13}{a_x = a_1 x_1 + a_2 x_2 + a_3 x_3}

gesetzt ist. Dadurch geht $R(x,\mu x+\lambda y)$ in $D(x,a)$ über.

Um nun $\lambda$ und $\mu$ zu bestimmen, multipliciren wir die Gleichungen (10) und (12) mit $a_1, a_2, a_3$, sodann (12) mit $b_1, b_2, b_3$ und addiren jedesmal. Dadurch erhalten wir mit Rücksicht auf $b_y = 0$:

\[
\begin{aligned}
a_y &= -\mu b_x = \Sigma \pm a_1 b_2 f'(x_3)\\
\mu a_x + \lambda a_y &= 0,
\end{aligned}
\]

und daraus

\[
\lambda b_x = a_x.
\]

Hiernach lässt sich die Gleichung (11) so darstellen:

\weq{14}{\frac{D(x,a)}{a_x^6} = \frac{D(x,b)}{b_x^6},}

und zeigt in dieser Form, dass der Quotient $D(x,a) : a_x^6$ von den willkürlichen Grössen $a$ unabhängig ist. Die Gleichung (14) ist aber keine Identität, sondern sie ist nur befriedigt, so lange $f(x) = 0$ ist. Wir können aber eine Identität daraus herleiten, wenn wir annehmen, dass $f(x)$ irreducibel sei (oder wenigstens keinen mehrfach zählenden Factor enthält). Es ist nämlich die Form $26^{\text{sten}}$ Grades

\[
b_x^6\, D(x,a) - a_x^6\, D(x,b)
\]
%CONT>>


% ===== scan leaf 377, printed 356 =====
% !PRINTED 356
\weberpage{356}
%<<CONT
gleich Null für alle der Gleichung $f(x)=0$ genügenden Werthe von $(x)$, und folglich muss sie durch $f(x)$ theilbar sein. Wir können also, wenn wir mit $\Phi(x,a,b)$ eine Form $22^{\text{sten}}$ Grades bezeichnen, setzen:

\weq{15}{b_x^6\, D(x,a) - a_x^6\, D(x,b) = f(x)\, \Phi(x,a,b).}

Denken wir uns in (15) für einen Augenblick ein neues Coordinatensystem eingeführt, in dem die Linien $a_x, b_x$ zwei Axen sind, von denen wir voraussetzen, dass sie sich nicht auf der Curve $f$ schneiden, so folgt aus der Vergleichung der rechten und linken Seite der identischen Gleichung (15), dass in $\Phi$ kein Glied vorkommen kann, was nicht entweder mit $a_x^6$ oder mit $b_x^6$ multiplicirt ist, und folglich hat $\Phi(x)$ die Form

\[
\Phi(x,a,b) = a_x^6\, \Phi_2(x) - b_x^6\, \Phi_1(x),
\]

worin $\Phi_1, \Phi_2$ Formen $16^{\text{ter}}$ Ordnung sind. Demnach erhält die identische Gleichung (15) die Form

\[
b_x^6\,[D(x,a) + f(x)\,\Phi_1(x)] = a_x^6\,[D(x,b) + f(x)\,\Phi_2(x)],
\]

und sie zeigt, dass $D(x,a) + f(x)\,\Phi_1(x)$ durch $a_x^6$ theilbar ist. Es existirt also eine Form $14^{\text{ten}}$ Grades $\chi(x,a,b)$, so dass

\weq{16}{\frac{D(x,a)}{a_x^6} = \chi(x,a,b) - f(x)\,\frac{\Phi_1(x,a,b)}{a_x^6}.}

Setzen wir hierin für die $a$ und $b$ irgend specielle, z. B. rationale Werthe $c, d$, und bezeichnen $\chi(x,c,d)$, was dann nur noch von $x$ und von den Coëfficienten der Form $f(x)$ abhängig ist, mit $\chi(x)$, so folgt:

\[
\frac{D(x,c)}{c_x^6} = \chi(x) - f(x)\,\frac{\Phi_1(x,c,d)}{c_x^6},
\]

und wenn wir dies von (16) subtrahiren und die Relation (15) für $b=c$ benutzen:

\weq{17}{\chi(x,a,b) - \chi(x) = f(x)\,\Psi(x),}

worin $\Psi(x)$ eine Function ist, die jedenfalls keinen anderen Nenner enthalten kann, als ein Product von Potenzen von $a_x$ und $c_x$. Da aber $f(x)$ nicht durch $a_x$ und $c_x$ theilbar ist, so muss $\Psi(x)$ eine ganze Function sein, und (16) ergiebt, wenn wir mit $\Theta(x)$ eine neue Form von $x$ (vom Grade 16) bezeichnen:

\weq{18}{\frac{D(x,a) - f(x)\,\Theta(x)}{a_x^6} = \chi(x).}


% ===== scan leaf 378, printed 357 =====
% !PRINTED 357
\weberpage{357}
\pdfbookmark[2]{§ 96. Die Steiner'schen Complexe}{bm111}
Diese Gleichung zeigt aber, dass die Curve $14^{\text{ten}}$ Grades:

\weq{19}{\chi(x) = 0,}

die von den $a$ gänzlich unabhängig ist, durch die Berührungspunkte der Doppeltangenten, aber durch keinen anderen Punkt der Curve $f(x)$ hindurchgeht.

Die Function $\chi(x)$ ist rational von den Coëfficienten der Gleichung $f(x)$ abhängig. Es giebt unendlich viele verschiedene solcher Curven, unter denen sich auch Covarianten von $f(x)$ finden. Man kann sie, wie \gesperrt{Hesse} nachgewiesen hat, in einfacher Weise durch Determinanten ausdrücken.

Hier ziehen wir daraus den Schluss:

2.\ \gesperrt{Eine Curve vierter Ordnung ohne singulären Punkt hat 28 Doppeltangenten.}

Einem Bedenken gegen diesen Schluss ist aber noch zu begegnen. Es wäre denkbar, dass die Curve $\chi$ die Curve $f$ berührt, oder dass die Curve $f$ durch singuläre Punkte der Curve $\chi$ hindurchgeht. Dann würde sich die Anzahl der Schnittpunkte und möglicherweise auch die Anzahl der Doppeltangenten vermindern, so dass unsere Schlussweise eigentlich nur lehrt, dass eine Curve vierter Ordnung nicht mehr als 28 Doppeltangenten haben kann. Dies Bedenken wird sich aber durch die folgenden Betrachtungen von selbst dadurch erledigen, dass wir Formen der Curvengleichung kennen lernen werden, bei denen die Existenz von 28 verschiedenen Doppeltangenten ersichtlich ist.

\begin{center}
\S.\,96.

\gesperrt{Die Steiner'schen Complexe.}
\end{center}

Wenn $x_1 = 0$ die Gleichung irgend einer Doppeltangente der Curve vierter Ordnung $f = 0$ ist, die wir kurz die Doppeltangente $x_1$ nennen, so muss, wenn man $x_1 = 0$ setzt, $f$ in ein Quadrat übergehen. Daraus ergiebt sich, dass $f$ von der Form sein muss:

\weq{1}{f = x_1 V - u^2,}

worin $V$ eine cubische, $u$ eine quadratische Form ist. Der Function $f$ kann aber auf dreifach unendlich viele Arten die Gestalt (1) gegeben werden. Denn bedeutet $p$ eine beliebige
%CONT>>


% ===== scan leaf 379, printed 358 =====
% !PRINTED 358
\weberpage{358}
%<<CONT
lineare Function, die drei willkürliche Constanten enthält, so
folgt aus (1):
\[
f = x_1 (V + 2pu + x_1p^2) - (u + px_1)^2,
\]
was wieder von der Form (1) ist.

Ist nun $y_1 = 0$ die Gleichung einer zweiten von $x_1$ verschiedenen Doppeltangente, so können wir die Constanten in $p$ (und zwar noch auf unendlich viele verschiedene Arten) so wählen, dass der Kegelschnitt $u + px_1 = 0$ durch die beiden Berührungspunkte von $y_1$ geht, und wir können daher annehmen, dass schon in der Form (1) die Function $u$ so gewählt sei. Dann muss in denselben Punkten auch die cubische Form $V$ verschwinden.

Aber noch mehr: Da die Linie $y_1 = 0$ Doppeltangente sein soll, so muss, wenn wir $x_1, y_1$ und irgend eine dritte davon unabhängige lineare Function $z$ als Variable einführen, in den Berührungspunkten
\[
f'(x_1) = 0, \quad f'(z) = 0
\]
sein, und daraus folgt nach (1):
\[
V'(x_1) = 0, \quad V'(z) = 0,
\]
d. h. die Curve dritter Ordnung $V = 0$ wird von der Linie $y_1 = 0$ in den Berührungspunkten mit $f$ berührt, und dies ist nur möglich, wenn $V$ zerfällt und die Function $y_1$ als Theiler enthält. Hiernach erhalten wir eine neue Gestalt der Function $f$:
\weq{2}{f = x_1 y_1 v - u^2,}
worin $u, v$ Functionen zweiten Grades sind.

Sind $x_1, y_1$ irgend beliebige lineare, $u, v$ quadratische Formen von drei Variablen, so stellt die durch (2) bestimmte Function $f$, gleich Null gesetzt, eine Curve vierter Ordnung dar, von der $x_1$ und $y_1$ zwei Doppeltangenten sind.

Sind umgekehrt $x_1$ und $y_1$ zwei beliebige Doppeltangenten einer Curve vierter Ordnung $f = 0$, so kann $f$ auf die Form (2) gebracht werden.

Denken wir uns die Co\"efficienten von $x_1$ und $y_1$ als variabel, und lassen diese Grössen so variiren, dass $x_1$ und $y_1$ sich auf der Curve $f$ schneiden, so wird ein Doppelpunkt eintreten.

Die Doppeltangenten arten aus in zwei von dem Doppelpunkte auslaufende Tangenten.


% ===== scan leaf 380, printed 359 =====
% !PRINTED 359
\weberpage{359}
Fallen die Doppeltangenten in eine Linie zusammen, so treten zwei Doppelpunkte auf.

Dagegen können sehr wohl, ohne dass singuläre Punkte entstehen, die beiden Berührungspunkte einer Doppeltangente zusammenfallen, und so die Doppeltangente in eine vierpunktig berührende Tangente übergehen. Dies tritt ein, wenn die Linie $x_1$ den Kegelschnitt $u$ berührt.

Auch die Form (2) ist noch mit Beibehaltung von $x_1, y_1$ auf unendlich viele Arten herzustellen.

Nehmen wir, um dies einzusehen, an, es sei
\[
f = x_1 y_1 v - u^2 = x_1 y_1 v_1 - u_1^2,
\]
so folgt die identische Gleichung:
\weq{3}{x_1 y_1 (v - v_1) = (u - u_1)(u + u_1).}
Wenn nun $u - u_1$ durch $x_1$, $u + u_1$ durch $y_1$ theilbar wäre, so würde im Schnittpunkte von $x_1$ und $y_1$ auch $u$ und folglich $f$ verschwinden. Dieser Schnittpunkt würde auf der Curve $f$ liegen, was, so lange die Curve keinen singulären Punkt hat, nicht möglich ist. Es muss also einer der beiden Factoren $u - u_1$, $u + u_1$ in (3) durch $x_1 y_1$ theilbar sein. Da das Vorzeichen von $u_1$ noch nicht bestimmt ist, so können wir annehmen, es sei $u - u_1$ durch $x_1 y_1$ theilbar, und also bis auf einen constanten Factor mit $x_1 y_1$ identisch. Es wird dann, wenn $\lambda$ ein solcher Factor ist,
\[
u_1 = u + \lambda x_1 y_1,
\]
und die Function $f$ erhält die Form:
\weq{4}{f = x_1 y_1 (v + 2\lambda u + \lambda^2 x_1 y_1) - (u + \lambda x_1 y_1)^2.}
Umgekehrt sind für jedes beliebige $\lambda$ die Formeln (2) und (4) mit einander identisch.

Wenn wir nun $\lambda$ so bestimmen können, dass die quadratische Function $v + 2\lambda u + \lambda^2 x_1 y_1$ in zwei lineare Factoren zerfällt, so erhält, wenn wir $u + \lambda x_1 y_1 = u_1$ setzen, $f$ nach (4) die Gestalt
\weq{5}{f = x_1 y_1 x_2 y_2 - u_1^2,}
und diese Form zeigt, dass $x_2, y_2$ \gesperrt{zwei weitere Doppeltangenten sind}, und dass die acht Berührungspunkte der Doppeltangenten $x_1, y_1, x_2, y_2$ auf einem Kegelschnitte $u_1 = 0$ liegen.


% ===== scan leaf 381, printed 360 =====
% !PRINTED 360
\weberpage{360}
Nun ist die nothwendige und hinreichende Bedingung dafür, dass eine ternäre quadratische Form in zwei lineare Factoren zerfalle, die, dass die Determinante der quadratischen Form verschwinde (Bd. I, \S.\,56).

Drücken wir die Formen $u, v$ durch $x_1, y_1$ und irgend eine dritte Variable $y$ aus, und bezeichnen die Co\"efficienten in $v$ und $u$ mit $a_{ik}, b_{ik}$, so dass $2a_{23}, 2b_{23}$ die Co\"efficienten von $x_1 y_1$ werden, so erhalten wir die Bedingung für das Zerfallen von $v + 2\lambda u + \lambda^2 x_1 y_1$ in der Form:
\weq{6}{\begin{vmatrix} a_{11}+2\lambda b_{11}, & a_{12}+2\lambda b_{12}, & a_{13}+2\lambda b_{13} \\ a_{21}+2\lambda b_{21}, & a_{22}+2\lambda b_{22}, & a_{23}+2\lambda b_{23}+\tfrac{1}{2}\lambda^2 \\ a_{31}+2\lambda b_{31}, & a_{32}+2\lambda b_{32}+\tfrac{1}{2}\lambda^2, & a_{33}+2\lambda b_{33} \end{vmatrix} = 0,}
und dies ist eine Gleichung $5^{\mathrm{ten}}$ Grades in Bezug auf $\lambda$, die uns also fünf Zerlegungen giebt:
\[
x_2 y_2,\ x_3 y_3,\ x_4 y_4,\ x_5 y_5,\ x_6 y_6.
\]
Ueber die Co\"efficienten in $f$ lässt sich, wie (6) zeigt, so verfügen, dass alle die so bestimmten Functionen $x_i y_i$ von einander verschieden sind, und daraus folgt, wie oben gezeigt, dass sie verschieden bleiben, so lange die Curve $f$ keine singulären Punkte hat. Es gilt also der folgende Satz:

1. \gesperrt{Zu jedem Paare von Doppeltangenten $x_1 y_1$ gehören fünf weitere Paare $x_i y_i$ von der Art, dass die acht Berührungspunkte von $x_1 y_1 x_i y_i$ auf einem Kegelschnitte liegen.}

Die sechs Paare von Doppeltangenten, die auf diese Weise bestimmt sind:
\[
x_1 y_1,\ x_2 y_2,\ x_3 y_3,\ x_4 y_4,\ x_5 y_5,\ x_6 y_6,
\]
wollen wir einen \gesperrt{Steiner'schen Complex} nennen\footnote{Es ist dafür bisher gewöhnlich der Ausdruck \gesperrt{Steiner'sche Gruppe} gebraucht. Da aber das Wort \glqq Gruppe\grqq{} in der Algebra eine ganz bestimmte andere Bedeutung hat, so ziehen wir es vor, diesen Ausdruck hier zu vermeiden.}.

Betrachten wir die Paare $x_1 y_1, x_2 y_2, x_3 y_3$ eines solchen Complexes, und setzen
\[
f = x_1 y_1 x_2 y_2 - u_1^2 = x_1 y_1 x_3 y_3 - u_2^2,
\]
%CONT>>


% ===== scan leaf 382, printed 361 =====
% !PRINTED 361
\weberpage{361}
%<<CONT
so folgt daraus die Identität
\[
x_1 y_1 (x_2 y_2 - x_3 y_3) = (u_1 - u_2)(u_1 + u_2),
\]
und daraus wie oben
\[
\begin{aligned}
u_1 - u_2 &= h x_1 y_1\\
u_1 + u_2 &= \frac{x_2 y_2 - x_3 y_3}{h},
\end{aligned}
\]
worin $h$ eine Constante bedeutet, also
\[
2u_1 = h x_1 y_1 + \frac{x_2 y_2 - x_3 y_3}{h},
\]
und danach wird
\weq{7}{4f = 4 x_1 y_1 x_2 y_2 - \left(h x_1 y_1 + \frac{x_2 y_2 - x_3 y_3}{h}\right)^2.}
Dieser Gleichungsform können wir eine elegantere Gestalt geben, wenn wir
\[
h x_1,\quad \frac{x_2}{h},\quad \frac{x_3}{h}
\]
durch
\[
x_1,\quad x_2,\quad x_3
\]
ersetzen. Dann bedeuten die neuen $x_1, x_2, x_3$, gleich Null gesetzt, dieselben Linien wie die ursprünglichen, da sich ja beide nur durch einen constanten Factor unterscheiden, und es ergiebt sich aus (7):
\weq{8}{\begin{aligned}
-4f &= (x_1 y_1 + x_2 y_2 - x_3 y_3)^2 - 4 x_1 y_1 x_2 y_2\\
&= x_1^2 y_1^2 + x_2^2 y_2^2 + x_3^2 y_3^2\\
&\quad - 2 x_2 y_2 x_3 y_3 - 2 x_1 y_1 x_2 y_2 - 2 x_1 y_1 x_3 y_3,
\end{aligned}}
und die Gleichung $f = 0$ kann auch in der eleganten irrationalen Form
\[
x_1 y_1 + x_2 y_2 - x_3 y_3 = -2\sqrt{x_1 y_1 x_2 y_2}
\]
oder
\weq{9}{\sqrt{x_1 y_1} + \sqrt{x_2 y_2} + \sqrt{x_3 y_3} = 0}
dargestellt werden. Aus (9) können wir wieder rückwärts die Form (8) herleiten. Weil aber (9) ganz symmetrisch ist, so können wir die drei Paare vertauschen und erhalten z. B. auch
\[
4 x_2 y_2 x_3 y_3 = (-x_1 y_1 + x_2 y_2 + x_3 y_3)^2,
\]
woraus zu ersehen ist, dass auch die Berührungspunkte von $x_2 y_2 x_3 y_3$ auf einem Kegelschnitte liegen, und dass in dem Complex, den man aus $x_2 y_2$ erhält, nicht nur das Paar $x_1 y_1$, son-
%CONT>>


% ===== scan leaf 383, printed 362 =====
% !PRINTED 362
\weberpage{362}
%<<CONT
dern auch das Paar $x_3 y_3$ und folglich alle Paare $x_i y_i$ vorkommen, dass also dieser Complex von dem aus $x_1 y_1$ abgeleiteten überhaupt nicht verschieden ist. Wir haben also den Satz:

\gesperrt{2. Die Paare eines \gesperrt{Steiner}'schen Complexes haben die Eigenschaft, dass die acht Berührungspunkte von irgend zweien dieser Paare auf einem Kegelschnitte liegen, und dass man immer denselben Complex erhält, von welchem der sechs Paare man ausgehen mag.}

Hieraus ergiebt sich, dass drei Doppeltangenten eines \gesperrt{Steiner}'schen Complexes, wie $x_1, y_1, x_2$, von denen zwei ein Paar des Complexes bilden, ihre sechs Berührungspunkte auf einem Kegelschnitte haben. Dagegen giebt es wieder Systeme von drei Doppeltangenten (Tripel), deren sechs Berührungspunkte nicht auf einem Kegelschnitte liegen.

Nach einer von \gesperrt{Frobenius} eingeführten Bezeichnung bilden drei Doppeltangenten, wie $x_1, y_1, x_2$, deren Berührungspunkte auf einem Kegelschnitte liegen, ein \gesperrt{syzygetisches Tripel}. Drei Doppeltangenten, deren sechs Berührungspunkte nicht auf einem Kegelschnitte liegen, bilden ein \gesperrt{azygetisches Tripel}. Entsprechend wollen wir vier Doppeltangenten, deren acht Berührungspunkte auf einem Kegelschnitte liegen, ein \gesperrt{syzygetisches Quadrupel}, und irgend ein System von Doppeltangenten, von denen je drei azygetisch sind, ein \gesperrt{azygetisches System} nennen.

Hier gilt nun der folgende wichtige Satz:

\gesperrt{3. Drei Doppeltangenten, die in einem \gesperrt{Steiner}'schen Complexe vorkommen, so dass keine zwei von ihnen ein Paar bilden, wie $x_1, x_2, x_3$, sind immer azygetisch.}

Der Beweis dieses Satzes ergiebt sich einfach aus der Gleichungsform (9). Aus ihr ersieht man zunächst, dass die drei Linien $x_1, x_2, x_3$ sich nicht in einem Punkte schneiden; denn ein solcher Schnittpunkt würde auf der Curve $f$ liegen und wäre daher ein singulärer Punkt. Wir können also $x_1, x_2, x_3$ als Coordinaten einführen und demnach
\[
\begin{gathered}
y_1 = \alpha_1 x_1 + \alpha_2 x_2 + \alpha_3 x_3\\
y_2 = \beta_1 x_1 + \beta_2 x_2 + \beta_3 x_3\\
y_3 = \gamma_1 x_1 + \gamma_2 x_2 + \gamma_3 x_3
\end{gathered}
\]
%CONT>>


% ===== scan leaf 384, printed 363 =====
% !PRINTED 363
\weberpage{363}
\pdfbookmark[2]{§ 97. Complexpaare und Complextripel}{bm112}
%<<CONT
setzen. Hierin kann keine der drei Constanten $\alpha_1, \beta_2, \gamma_3$ verschwinden. Denn wenn z. B. $\alpha_1 = 0$ ist, so schneiden sich nach (9) die Linien $x_2, x_3, y_1$ auf der Curve $f$, und dieser Schnittpunkt müsste ein singulärer Punkt sein.

Nach der Gleichung (9) ergeben sich aber die Coordinaten der Berührungspunkte von $x_1, x_2, x_3$ aus den folgenden drei Paaren von Gleichungen, von denen jedes Paar zwei Berührungspunkte giebt:
\[
\begin{gathered}
x_1 = 0,\quad x_2(\beta_2 x_2 + \beta_3 x_3) - x_3(\gamma_2 x_2 + \gamma_3 x_3) = 0\\
x_2 = 0,\quad x_3(\gamma_3 x_3 + \gamma_1 x_1) - x_1(\alpha_3 x_3 + \alpha_1 x_1) = 0\\
x_3 = 0,\quad x_1(\alpha_1 x_1 + \alpha_2 x_2) - x_2(\beta_1 x_1 + \beta_2 x_2) = 0.
\end{gathered}
\]
Sollen nun diese sechs Punkte auf einem Kegelschnitte $\varphi = 0$ liegen, so muss $\varphi$, von einem constanten Factor $h$ abgesehen, für $x_1 = 0$ in den linken Theil der zweiten Gleichung des ersten Paares $x_2(\beta_2 x_2 + \beta_3 x_3) - x_3(\gamma_2 x_2 + \gamma_3 x_3)$ übergehen. Bezeichnen wir also mit $a_1, a_2, a_3$ die Co\"efficienten von $x_1^2, x_2^2, x_3^2$ in $\varphi$, so muss
\weq{10}{
\begin{gathered}
a_2 = h_1 \beta_2,\quad a_3 = -h_1 \gamma_3\\
a_3 = h_2 \gamma_3,\quad a_1 = -h_2 \alpha_1\\
a_1 = h_3 \alpha_1,\quad a_2 = -h_3 \beta_2
\end{gathered}
}
sein, und hierin sind $h_1, h_2, h_3$ drei Constanten.

Aus (10) folgt aber:
\[
h_2 = -h_3,\quad h_3 = -h_1,\quad h_1 = -h_2,
\]
und dies wäre nur möglich, wenn $h_1 = h_2 = h_3 = 0$ wäre.
Dies ist aber nur dann der Fall, wenn $a_1, a_2, a_3$ verschwinden, wenn also der Kegelschnitt $\varphi$ durch die Schnittpunkte der drei Geraden $x_1, x_2, x_3$ geht. Er soll aber durch die Berührungspunkte dieser drei Doppeltangenten gehen, die sicher von ihren drei Schnittpunkten verschieden sind. Also ist unsere Annahme als unmöglich nachgewiesen und der Satz 3. bewiesen.

\begin{center}
§. 97.

\gesperrt{Complexpaare und Complextripel.}
\end{center}

Die Sätze des vorigen Paragraphen geben ein vorzügliches Hülfsmittel, um die mannigfachen geometrischen und algebraischen Beziehungen der Doppeltangenten zu erforschen und dar-
%CONT>>


% ===== scan leaf 385, printed 364 =====
% !PRINTED 364
\weberpage{364}
%<<CONT
zustellen. Wir beschränken uns hier auf das, was für die Erreichung unseres Hauptzieles, nämlich der Bestimmung der \gesperrt{Galois}'schen Gruppe des Problems, erforderlich ist.

Wir gehen von einem beliebig herausgegriffenen \gesperrt{Steiner}'schen Complexe aus:
\weq{1}{x_1 y_1,\ x_2 y_2,\ x_3 y_3,\ x_4 y_4,\ x_5 y_5,\ x_6 y_6.}
Die Form der Gleichung \S.\,96, (5) zeigt dann, dass in dem durch das Paar $x_1 x_2$ bestimmten Complexe das Paar $y_1 y_2$, und in dem durch $x_1 y_2$ bestimmten Complexe das Paar $x_2 y_1$ vorkommen muss.

Wir betrachten also neben (1) die zwei weiteren Complexe:
\weq{2}{x_1 x_2,\quad y_1 y_2,\quad \ldots}
\weq{3}{x_1 y_2,\quad x_2 y_1,\quad \ldots}
Jeder der beiden Complexe (2), (3) enthält ausser den schon bekannten noch acht weitere Doppeltangenten, und diese müssen alle von den $x_i, y_i$ verschieden sein, weil, wenn z. B. $x_3$ in (2) vorkäme, $x_1 x_2 x_3$ syzygetisch wäre, was dem Satze 3., (\S.\,96) widerspricht. Ebenso können die beiden Complexe (2) und (3) ausser $x_1, x_2, y_1, y_2$ keine gemeinsame Doppeltangente enthalten. Denn wenn etwa $z$ in beiden vorkäme, so wären $x_1 x_2 z$ nach (2) syzygetisch, nach (3) azygetisch, was ein Widerspruch ist. Daraus folgt:

\gesperrt{4. In den drei Complexen (1), (2), (3) zusammengenommen kommen alle 28 Doppeltangenten vor.}

Hieraus folgt, dass jede Doppeltangente, die mit irgend einem Paare $x_1 y_1$ ein syzygetisches Tripel bildet, in dem Complexe $x_1 y_1$ vorkommen muss, dass also jedes syzygetische Tripel durch eine bestimmte weitere Doppeltangente zu einem syzygetischen Quadrupel ergänzt wird, und dass folglich ein Kegelschnitt, der durch die Berührungspunkte von drei Doppeltangenten hindurchgeht, die Curve $f$ in den Berührungspunkten einer vierten Doppeltangente schneidet.

Zwei Paare eines Complexes bilden immer ein syzygetisches Quadrupel, und wie man ein solches Quadrupel auch in zwei Paare theilen mag, beide Paare gehören immer demselben Complexe an.

Da man aus 28 Dingen 14.27 Paare bilden kann, da jedes Paar von Doppeltangenten einen Complex bestimmt und in jedem
%CONT>>


% ===== scan leaf 386, printed 365 =====
% !PRINTED 365
\weberpage{365}
%<<CONT
Complexe sechs Paare vorkommen, so ist die Gesammtzahl der Complexe $14.27 : 6 = 63$.

\gesperrt{5. Es giebt 63 \gesperrt{Steiner}'sche Complexe.}

Wenn wir aus den Paaren des Complexes (1) statt $x_1 y_1, x_2 y_2$ irgend ein anderes Paar von Paaren herausgreifen, so können wir daraus nach dem Schema von (2) und (3) jedesmal zwei neue Complexe bilden. Da es 15 solcher Paare von Paaren giebt, so erhalten wir 30 neue Complexe vom Typus (2), (3), die alle in gleicher Weise aus dem Complexe (1) abgeleitet sind, und die alle unter einander verschieden sind.

Nehmen wir nun irgend eine von den in (2) und (3) neu hinzutretenden Doppeltangenten $z_1$, so erhalten wir zwei neue Complexe, wenn wir von den beiden Paaren $x_1 z_1, y_1 z_1$ ausgehen, und da wir $z_1$ auf 16 verschiedene Arten wählen können, so ergeben sich so 32 neue Complexe, womit die Gesammtzahl aller Complexe erschöpft ist. Es muss aber noch die Vertheilung der Doppeltangenten auf die Complexe $x_1 z_1, y_1 z_1$ genauer untersucht werden.

Wir können, ohne die Allgemeinheit zu beschränken, da wir n\"othigenfalls $x_1$ mit $y_1$ vertauschen können, die Annahme machen, dass $z_1$ in dem Complexe (2) und darin in dem Paare $z_1 z_2$ vorkomme.

Dann erhalten wir die beiden Complexe:
\weq{4}{x_1 z_1,\quad x_2 z_2,\quad \ldots}
\weq{5}{y_1 z_1,\quad y_2 z_2,\quad \ldots}
Wir betrachten jetzt die drei Complexe:
\weq{4}{x_1 z_1,\quad x_2 z_2,\quad \ldots}
\weq{2}{x_1 x_2,\quad z_1 z_2,\quad \ldots}
\weq{4a}{x_1 z_2,\quad x_2 z_1,\quad \ldots,}
die nach dem Satze 4. alle Doppeltangenten enthalten müssen, darunter also auch $x_3, y_3$. Diese kommen aber nicht in (2) vor, und ebenso können nicht beide in (4) oder beide in (4a) vorkommen, da sonst $x_1, x_3, y_3$ azygetisch sein müssten, während sie doch [nach (1)] syzygetisch sind. Da wir eventuell $x_3$ und $y_3$ in der Bezeichnung vertauschen dürfen, so können wir annehmen, dass $x_3$ in (4) vorkomme, und zwar in einem Paare $x_3 z_3$.

Dann kann, da $z_1 z_2 z_3$ azygetisch sind, $z_3$ nicht in dem Complexe (2) vorkommen und muss folglich in (3) enthalten sein.


% ===== scan leaf 387, printed 366 =====
% !PRINTED 366
\weberpage{366}
Nun haben wir die beiden Complexe:

\weq{4}{x_2 z_2,\ x_3 z_3,\ \dots}
\weq{4\text{b}}{x_2 x_3,\ z_2 z_3,\ \dots{},}

und da $x_2 x_3 y_2 y_3$ ein syzygetisches Quadrupel sind, so enthält die Gruppe (4 b) auch das Paar $y_2 y_3$, und folglich sind auch $y_2 y_3 z_2 z_3$ ein syzygetisches Quadrupel. Daraus folgt weiter, dass $y_2 z_2$ und $y_3 z_3$ in denselben Complex gehören. Da man dieselbe Betrachtung wie für $x_3 y_3$ auch für die Paare $x_4 y_4, x_5 y_5, x_6 y_6$ durchführen kann, so lassen sich hiernach die Complexe (4), (5) vollständig bilden, und sie erhalten den Ausdruck:

\weq{4}{x_1 z_1,\ x_2 z_2,\ x_3 z_3,\ x_4 z_4,\ x_5 z_5,\ x_6 z_6}
\weq{5}{y_1 z_1,\ y_2 z_2,\ y_3 z_3,\ y_4 z_4,\ y_5 z_5,\ y_6 z_6.}

Hierin bilden $z_1 z_2$ ein Paar des Complexes (2) und $z_3, z_4, z_5, z_6$, die ein azygetisches System bilden, kommen alle in dem Complex (3) vor, in dem keine zwei gepaart sind.

Da wir, wie vorhin gezeigt, nach dem Typus (2), (3), (4), (5) aus dem willkürlich angenommenen Complex (1) alle überhaupt existirenden Complexe ableiten können, so ergiebt sich der Satz:

6. \gesperrt{Irgend zwei Complexe haben entweder ein syzygetisches Quadrupel oder ein azygetisches System von sechs Doppeltangenten gemein.}

Zwei Complexe, die ein syzygetisches Quadrupel gemein haben, wollen wir ein \gesperrt{syzygetisches Complexpaar} nennen. Zu jedem solchen Paare giebt es einen dritten Complex, der dasselbe Quadrupel enthält. Drei solche Complexe nennen wir ein \gesperrt{syzygetisches Complextripel} [z.\,B. (1), (2), (3)].

Ebenso nennen wir zwei Complexe der zweiten Art, d.\,h. solche, die sechs azygetische Elemente gemein haben, ein \gesperrt{azygetisches Complexpaar}.

Jedes azygetische Complexpaar wird gleichfalls durch einen bestimmten Complex, der mit jedem der beiden Complexe ein azygetisches Paar bildet, zu einem Tripel ergänzt [wie (1), (4), (5)]. Ein solches nennen wir ein \gesperrt{azygetisches Complextripel}.

In einem syzygetischen Tripel kommen alle 28 Doppeltangenten vor, in einem azygetischen nur 18.


% ===== scan leaf 388, printed 367 =====
% !PRINTED 367
\weberpage{367}
\pdfbookmark[2]{§ 98. Die Aronhold'schen Siebener-Systeme}{bm113}
\S.\,98.

\gesperrt{Die Aronhold'schen Siebener-Systeme.}

Von besonderer Wichtigkeit sind die azygetischen Systeme von sieben Doppeltangenten, die zuerst von \gesperrt{Aronhold} betrachtet sind, und die daher \gesperrt{Aronhold'sche Siebener-Systeme} heissen. Wir nennen sie auch \gesperrt{vollständige Siebener-Systeme} oder kurz \gesperrt{vollständige Systeme}.

Dass es solche Systeme giebt, zeigen die Zusammenstellungen des vorigen Paragraphen. Denn wenn

\weq{1}{\begin{aligned}
&x_1 y_1,\ x_2 y_2,\ x_3 y_3,\ x_4 y_4,\ x_5 y_5,\ x_6 y_6\\
&x_1 z_1,\ x_2 z_2,\ x_3 z_3,\ x_4 z_4,\ x_5 z_5,\ x_6 z_6
\end{aligned}}

ein azygetisches Complexpaar ist, so ist

\[ x_1,\ x_2,\ x_3,\ x_4,\ x_5,\ y_6,\ z_6 \]

ein vollständiges Siebener-System (weil in dem Complexe $y_6 z_6$ keines der $x$ vorkommt). Es wird sich zeigen, dass keine azygetischen Systeme von mehr als sieben Elementen bestehen. Es gilt zunächst der Satz:

7. \gesperrt{Irgend sechs Elemente eines vollständigen Systemes kommen in einem und nur in einem Complexe vor.}

Wir beweisen zunächst den zweiten Theil der Behauptung, d.\,h. wir beweisen, dass, wenn

\weq{2}{x_1,\ x_2,\ x_3,\ x_4,\ x_5,\ x_6,\ x_7}

ein vollständiges System ist, $x_1, x_2, x_3, x_4, x_5, x_6$ nicht in zwei verschiedenen Complexen vorkommen können. Nehmen wir an, es sei dies möglich, so müssen die beiden Complexe ein azygetisches Paar wie (1) bilden. In dem syzygetischen Tripel

\[\begin{aligned}
&y_1 z_1,\ y_2 z_2,\ y_3 z_3,\ y_4 z_4,\ y_5 z_5,\ y_6 z_6\\
&y_1 y_2,\ z_1 z_2,\ x_1 x_2\\
&y_1 z_2,\ y_2 z_1
\end{aligned}\]

müssen die Doppeltangenten $x_3, x_4, x_5, x_6, x_7$ vorkommen, und da sie weder im ersten noch im zweiten dieser Complexe vorkommen, so müssen sie im dritten vorkommen. Da in diesem Complexe aber nur noch vier Paare übrig sind, so müssen minde-
%CONT>>


% ===== scan leaf 389, printed 368 =====
% !PRINTED 368
\weberpage{368}
%<<CONT
stens zwei von den $x_3, \dots x_7$ ein Paar darin bilden. Das ist aber nicht möglich, da dieses Paar sonst mit einem der übrigen $x$ ein syzygetisches Tripel bilden würde.

Um nachzuweisen, dass die Doppeltangenten $x_1, x_2, x_3, x_4, x_5, x_6$ immer in einem Complexe vorkommen, nehmen wir zwei beliebige von ihnen, $x_1 x_2$, heraus und wählen ein in dem Complexe $x_1 x_2$ vorkommendes anderes Paar $y_1 y_2$, was auf fünf Arten möglich ist. Daraus bilden wir das syzygetische Complextripel

\weq{2}{\begin{aligned}
\alpha)\ & x_1 x_2,\ y_1 y_2\\
\beta)\ & x_1 y_1,\ x_2 y_2\\
\gamma)\ & x_1 y_2,\ x_2 y_1,
\end{aligned}}

welches fünf verschiedene solcher Tripel repräsentirt. In $\beta$) und $\gamma$) müssen nun $x_3, x_4, x_5, x_6, x_7$ vorkommen, und zwar keine zwei von ihnen gepaart. Wir zeigen nun zunächst, dass sich diese fünf $x$ nicht zu zwei und drei auf die beiden Complexe $\beta$), $\gamma$) vertheilen können, sondern nur zu eins und vier. Nehmen wir nämlich an, die Complexe $\beta$), $\gamma$) seien so zusammengesetzt:

\weq{2}{\begin{aligned}
\beta)\ & x_1 y_1,\ x_2 y_2,\ x_3 y_3,\ x_4 y_4,\ x_5 y_5,\ \cdot\\
\gamma)\ & x_1 y_2,\ x_2 y_1,\ x_6 y_6,\ x_7 y_7,\ \cdot\quad\cdot\text{,}
\end{aligned}}

so können wir noch den Complex

\weq{2}{\delta)\quad x_6 x_7,\ y_6 y_7\ \cdots{},}

betrachten, der mit $\beta$) zusammen ein syzygetisches Paar bildet, weil weder $x_1$ noch $y_1$ in $\delta$) vorkommt, das Paar also nicht azygetisch sein kann. Es müssen also $\beta$) und $\delta$) ein syzygetisches Quadrupel gemein haben, und da in zwei Paaren von $\beta$) mindestens ein von $x_6, x_7$ verschiedenes $x$ vorkommt, so muss dieses $x$ auch in $\delta$) vorkommen, was unmöglich ist, weil kein $x$ mit $x_6, x_7$ syzygetisch ist. Es bleibt also für $\beta$), $\gamma$) nur eine Zusammensetzung übrig, wie die folgende:

\weq{3}{\begin{aligned}
\beta)\ & x_1 y_1,\ x_2 y_2,\ x_3 y_3,\ x_4 y_4,\ x_5 y_5,\ x_6 y_6\\
\gamma)\ & x_1 y_2,\ x_2 y_1,\ x_7 y_7,\ \cdot\quad\cdot\quad\cdot
\end{aligned}}

Dass wir gerade diese Annahme machen, und nicht die sechs $x$ in $\gamma$) aufnehmen, ist keine Beschränkung, da wir, wenn n\"othig, $y_1$ mit $y_2$ vertauschen können.

Wenn wir nun unter Festhaltung von $x_1 x_2$ an Stelle von $y_1 y_2$ die anderen Paare von (2) $\alpha$) treten lassen, so bekommen
%CONT>>


% ===== scan leaf 390, printed 369 =====
% !PRINTED 369
\weberpage{369}
\pdfbookmark[2]{§ 99. Der Hesse-Cayley'sche Algorithmus zur Bezeichnung der Doppeltangenten}{bm114}
%<<CONT
wir fünf Complexbildungen vom Typus (3). Zwei solche Complexbildungen können sich aber nur dadurch unterscheiden, dass an Stelle von $x_7$ jedesmal ein anderes der Elemente $x_3, x_4, x_5, x_6, x_7$ tritt, und alle diese Möglichkeiten müssen auch vorkommen, weil wir sonst zwei verschiedene Complexe erhalten würden, die dieselben sechs Elemente des vollständigen Systems der $x$ enthalten, was nicht möglich ist, wie wir bewiesen haben. Demnach können wir annehmen, dass die in dem Complexe (3) $\beta$) vorkommenden $x_1, x_2, x_3, x_4, x_5, x_6$ \gesperrt{irgend welche sechs Elemente} des vollständigen Systemes seien, was bewiesen werden sollte.

\S.\,99.

\gesperrt{Der Hesse-Cayley'sche Algorithmus zur Bezeichnung der Doppeltangenten.}

Der Satz 7. des vorigen Paragraphen führt zu einer Bezeichnungsweise für die Doppeltangenten, durch die eine sehr übersichtliche Darstellung aller dieser Verhältnisse möglich ist, die von \gesperrt{Cayley} (im Anschluss an \gesperrt{Hesse}) ausgebildet ist.

Wir legen ein vollständiges Siebener-System

\[ x_1,\ x_2,\ x_3,\ x_4,\ x_5,\ x_6,\ x_7 \]

zu Grunde, dessen Elemente wir einfach durch die Ziffern 1, 2, 3, 4, 5, 6, 7 bezeichnen. Wir sondern eine beliebige Doppeltangente, etwa $x_1$, dieses Systemes aus, und bilden den Complex, der die übrigen sechs enthält:

\weq{1}{x_2 y_2,\ x_3 y_3,\ x_4 y_4,\ x_5 y_5,\ x_6 y_6,\ x_7 y_7.}

Die Doppeltangente $y_2$ ist dann durch das gewählte $x_1$ und durch $x_2$ völlig bestimmt und kann daher durch $[1\,2]$ bezeichnet werden. Ebenso bezeichnen wir $y_3$ durch $[1\,3]$ u.\,s.\,f. Die Doppeltangenten $[1\,2], [1\,3], \dots, [1\,7]$ sind hierdurch vollständig bestimmt und von einander verschieden. Aus dieser Bestimmung geht auch hervor, was man allgemein unter $[\mu\,\nu]$ zu verstehen hat, wenn $\mu, \nu$ zwei verschiedene Ziffern aus der Reihe 1 bis 7 bedeuten.

Es ist nun zunächst zu zeigen, dass $[\mu\,\nu] = [\nu\,\mu]$ ist. Es genügt, wenn wir nachweisen, dass $[1\,2] = [2\,1]$ ist. Dazu brauchen wir nur den Complex zu bilden, der $x_1, x_3, x_4, x_5, x_6, x_7$
%CONT>>


% ===== scan leaf 391, printed 370 =====
% !PRINTED 370
\weberpage{370}
% Source-faithful three-stroke Cayley configuration sign.
\providecommand{\weberthreeup}{\mathord{\tikz[baseline=-0.2ex,x=0.23em,y=0.23em,line width=0.45pt]{\draw (0,0)--(-1,1) (0,0)--(0,1.08) (0,0)--(1,1);}}}
%<<CONT
enthält, und der mit (1) ein azygetisches Paar bildet. Er muss
also von der Gestalt sein:
\weq{2}{x_1 y_2,\ x_3 z_3,\ x_4 z_4,\ x_5 z_5,\ x_6 z_6,\ x_7 z_7,}
und demnach ist $y_2$ auch durch $[2\,1]$ zu bezeichnen, w. z. b. w.

Die $z$ in (2) sind von den $y$ in (1) verschieden, und da z. B.
$z_3 = [2\,3]$ ist, so folgt, dass allgemein zwei $[\mu\,\nu]$, die nicht in
beiden Ziffern $\mu, \nu$ übereinstimmen, von einander verschieden
sind. Da man aus sieben Ziffern einundzwanzig Paare bilden
kann, so erhält man auf diese Weise alle Doppeltangenten und
jede nur einmal.

Aus dieser Darstellungsweise ergiebt sich auch, dass keine
azygetischen Systeme von mehr als sieben Doppeltangenten existiren
; denn fügen wir zu den sieben $x$ noch eine beliebige
weitere Doppeltangente, für die wir bei der Gleichberechtigung
der Ziffern 1 bis 7 etwa $y_2 = [1\,2]$ wählen können, so ist $y_2 x_1 x_2$
ein syzygetisches Tripel.

Es ist zweckmässig, eine achte Ziffer 8 einzuführen, und die
Elemente des ursprünglichen Siebener-Systems nicht durch die
einfachen Ziffern, sondern durch die Paare
\[
[1\,8],\ [2\,8],\ [3\,8],\ [4\,8],\ [5\,8],\ [6\,8],\ [7\,8]
\]
zu bezeichnen, wobei dann auch gelten soll, dass $[1\,8] = [8\,1]$ u. s. f.
ist. Dann werden alle 28 Doppeltangenten übereinstimmend
durch die Paare $[\mu\,\nu]$ bezeichnet, in denen $\mu$ und $\nu$ zwei verschiedene
 Ziffern der Reihe 1 bis 8 bedeuten.

Dabei ist es für die Uebersichtlichkeit sehr förderlich, wenn
man eine anschauliche Bezeichnung anwendet\footnote{Nach \gesperrt{Cayley}; vergl. \gesperrt{Salmon}, \glqq Higher plane curves\grqq{}.}. Man deutet
eine Doppeltangente $[\mu\,\nu]$ durch einen einfachen Strich $\mid$ an, an
dessen Enden man sich die beiden Ziffern $\mu, \nu$ gesetzt denkt.
Dann bedeuten zwei Striche ohne gemeinsamen Punkt $\mid\mid$ zwei
Doppeltangenten, in deren Bezeichnung $[\mu\,\nu]$ keine gemeinschaftliche
 Ziffer vorkommt, und zwei von einem Punkte auslaufende
Striche, $\vee$, zwei Doppeltangenten, die in ihren Symbolen $[\mu\,\nu]$
eine gemeinschaftliche Ziffer haben. Hiernach sind die complicirteren
 Zeichen, die wir nachher anwenden, von selbst verständlich.
Für ein Tripel von Doppeltangenten haben wir z. B. folgende
fünf Zeichen $\mid\mid\mid$, $\sqcap$, $\triangle$, $\weberthreeup$, $\vee\mid$.


% ===== scan leaf 392, printed 371 =====
% !PRINTED 371
\weberpage{371}
% Source-faithful three-stroke Cayley configuration sign.
\providecommand{\weberthreeup}{\mathord{\tikz[baseline=-0.2ex,x=0.23em,y=0.23em,line width=0.45pt]{\draw (0,0)--(-1,1) (0,0)--(0,1.08) (0,0)--(1,1);}}}
Es soll jetzt zunächst untersucht werden, wie sich in dieser
Bezeichnungsweise die azygetischen und die syzygetischen Tripel
unterscheiden.

Dabei ist zu beachten, dass nach der bis jetzt gegebenen
Erklärung die Ziffer 8 eine besondere Stelle einnimmt, während
die Ziffern 1 bis 7 vollständig gleichartig auftreten und beliebig
permutirt werden können.

Wir leiten aus den beiden Complexen (1), (2) noch die Ergänzung
 zu einem azygetischen Complextripel her, nämlich:
\weq{3}{
\begin{aligned}
&x_2 y_2,\ x_3 y_3,\ x_4 y_4,\ x_5 y_5,\ x_6 y_6,\ x_7 y_7,\\
&x_1 y_2,\ x_3 z_3,\ x_4 z_4,\ x_5 z_5,\ x_6 z_6,\ x_7 z_7,\\
&x_1 x_2,\ y_3 z_3,\ y_4 z_4,\ y_5 z_5,\ y_6 z_6,\ y_7 z_7,
\end{aligned}
}
und gehen nun die einzelnen Zeichen für die Doppeltangententripel
 durch.

Wir beginnen mit dem Zeichen $\weberthreeup$, für welches wir mit
Rücksicht auf die Ausnahmestellung der Ziffer 8 drei Typen zu
betrachten haben:
\weq{4}{
\begin{aligned}
&[1\,8]\,[2\,8]\,[3\,8] = x_1 x_2 x_3\\
&[1\,5]\,[1\,6]\,[1\,7] = y_5 y_6 y_7\\
&[1\,6]\,[1\,7]\,[1\,8] = y_6 y_7 x_1,
\end{aligned}
}
und der Anblick der drei Complexe (3) zeigt (nach dem Satze
\S.\,96, 3.), dass alle diese Tripel azygetisch sind.

Zweitens betrachten wir das Zeichen $\vee\mid$, für welches vier
Typen zu berücksichtigen sind:
\weq{5}{
\begin{aligned}
&[8\,3]\,[8\,4]\,[1\,2] = x_3 x_4 y_2\\
&[8\,4]\,[1\,4]\,[2\,3] = x_4 y_4 z_3\\
&[1\,2]\,[1\,3]\,[4\,8] = y_2 y_3 x_4\\
&[1\,4]\,[1\,5]\,[2\,3] = y_4 y_5 z_3,
\end{aligned}
}
und auch diese Tripel sind azygetisch.

Für das Zeichen $\triangle$ ist zu betrachten:
\weq{6}{
\begin{aligned}
&[1\,8]\,[2\,8]\,[1\,2] = x_1 x_2 y_2\\
&[1\,2]\,[1\,3]\,[2\,3] = y_2 y_3 z_3,
\end{aligned}
}
die sich gleichfalls als azygetisch erweisen, weil $y_2$ in dem Complexe
 (3), der die Paare $x_1 x_2$, $y_3 z_3$ enthält, nicht vorkommt.

Für das Zeichen $\mid\mid\mid$ giebt es zwei Typen:
\weq{7}{
\begin{aligned}
&[1\,3]\,[2\,4]\,[5\,8] = y_3 z_4 x_5\\
&[1\,3]\,[2\,4]\,[5\,6] = y_3 z_4 [5\,6].
\end{aligned}
}


% ===== scan leaf 393, printed 372 =====
% !PRINTED 372
\weberpage{372}
% Source-faithful three-stroke Cayley configuration sign.
\providecommand{\weberthreeup}{\mathord{\tikz[baseline=-0.2ex,x=0.23em,y=0.23em,line width=0.45pt]{\draw (0,0)--(-1,1) (0,0)--(0,1.08) (0,0)--(1,1);}}}
Betrachten wir das syzygetische Complextripel
\weq{8}{
\begin{aligned}
&y_3 z_4,\ y_4 z_3,\ \cdot\ \cdot\ \cdot\ \cdot\\
&y_3 y_4,\ z_3 z_4,\ x_3 x_4\ \cdot\ \cdot\\
&y_3 z_3,\ y_4 z_4,\ \cdot\ \cdot\ \cdot\ \cdot,
\end{aligned}
}
in dem alle Doppeltangenten vorkommen müssen, so finden wir
$[5\,8]$ und $[5\,6]$ nicht in den beiden letzten Complexen von (8),
weil
\weq{9}{
\begin{aligned}
&y_3 z_3 [5\,8] = [1\,3]\,[2\,3]\,[5\,8],\quad y_3 y_4 [5\,8] = [1\,3]\,[1\,4]\,[5\,8]\\
&y_3 z_3 [5\,6] = [1\,3]\,[2\,3]\,[5\,6],\quad y_3 y_4 [5\,6] = [1\,3]\,[1\,4]\,[5\,6]
\end{aligned}
}
das Zeichen $\vee\mid$ haben und daher azygetisch sind. Folglich kommen
 $[5\,8]$ und $[5\,6]$ im ersten der Complexe (8) vor, und die
beiden Tripel (7) sind syzygetisch.

Endlich haben wir auch das Zeichen $\sqcap$ zu betrachten, das
wieder drei Typen giebt:
\weq{10}{
\begin{aligned}
&[1\,2]\,[2\,3]\,[3\,4] = y_2 z_3 [3\,4]\\
&[1\,2]\,[2\,3]\,[3\,8] = y_2 z_3 x_3\\
&[1\,8]\,[8\,2]\,[2\,3] = x_1 x_2 z_3.
\end{aligned}
}
Auch diese Tripel sind syzygetisch, was für die beiden letzten
unmittelbar aus dem Complextripel (3) zu ersehen ist, und für
das erste aus dem syzygetischen Complextripel
\[
\begin{aligned}
&y_2 z_3,\ x_1 x_3,\ \cdot\ \cdot\ \cdot\ \cdot\\
&y_2 x_1,\ z_3 x_3,\ \cdot\ \cdot\ \cdot\ \cdot\\
&y_2 x_3,\ z_3 x_1,\ \cdot\ \cdot\ \cdot\ \cdot,
\end{aligned}
\]
folgt, von denen die beiden letzten $[3\,4]$ nicht enthalten, weil
$y_2 x_1 [3\,4]$ und $y_2 x_3 [3\,4]$ beide das Zeichen $\vee\mid$ haben.

Hier ist aber die Ausnahmestellung der Ziffer 8 gänzlich
verschwunden, und wir kommen zu dem Resultate:

\gesperrt{Unter den Tripeln von Doppeltangenten sind die mit den
Zeichen}
\[
\mid\mid\mid,\ \sqcap
\]
\gesperrt{syzygetisch, und die mit den Zeichen}
\[
\weberthreeup,\ \triangle,\ \vee\mid
\]
\gesperrt{azygetisch.}

Hieraus erhält man sehr leicht die Zeichen für sämmtliche
syzygetische und azygetische Quadrupel:
%CONT>>


% ===== scan leaf 394, printed 373 =====
% !PRINTED 373
\weberpage{373}
%<<CONT
% Source-faithful Cayley configuration signs: three/four concurrent strokes and a seven-rayed star.
\providecommand{\weberthreeup}{\mathord{\tikz[baseline=-0.2ex,x=0.23em,y=0.23em,line width=0.45pt]{\draw (0,0)--(-1,1) (0,0)--(0,1.08) (0,0)--(1,1);}}}
\providecommand{\weberfourup}{\mathord{\tikz[baseline=-0.2ex,x=0.23em,y=0.23em,line width=0.45pt]{\draw (0,0)--(-1.05,1) (0,0)--(-0.35,1.08) (0,0)--(0.35,1.08) (0,0)--(1.05,1);}}}
\providecommand{\webersevenstar}{\mathord{\tikz[baseline=-0.35ex,x=0.21em,y=0.21em,line width=0.45pt]{\draw (0,0)--(90:1) (0,0)--(38.57:1) (0,0)--(-12.86:1) (0,0)--(-64.29:1) (0,0)--(-115.71:1) (0,0)--(-167.14:1) (0,0)--(141.43:1);}}}
\gesperrt{Die Quadrupel von Doppeltangenten mit den
Zeichen}
\[
\mid\mid\mid\mid,\ \square
\]
\gesperrt{sind syzygetisch, und die mit den Zeichen}
\[
\weberfourup,\ \triangle\mid,\ \weberthreeup\mid,\ \vee\vee
\]
\gesperrt{azygetisch.}

Alle übrigen Quadrupel sind weder syzygetisch noch azygetisch
.

Hiernach ist es leicht, die Zeichen für sämmtliche vollständige
 Siebener-Systeme zu bilden.

Ein solches Zeichen muss aus sieben Strichen bestehen, die
nicht mehr als acht Eckpunkte haben können und die eine Figur
bilden, aus der sich keine der beiden Figuren $\mid\mid\mid$, $\sqcap$ ablösen
lässt. Daraus folgt zunächst, dass diese Figur aus nicht mehr als
zwei getrennten Theilen bestehen kann, weil sonst die Figur $\mid\mid\mid$
darin enthalten wäre, und dass kein Theil mehr als ein Centrum
haben kann, von dem mehrere Striche auslaufen, ausser wenn
dieser Theil das Dreieck $\triangle$ ist, weil sonst die Figur $\sqcap$ vorkommen
 würde.

Wenn nun die Figur eintheilig ist, so muss sie ein siebenstrahliger
 Stern $\webersevenstar$ sein, und da man jede der acht Ziffern als
Mittelpunkt wählen kann, so sind dies acht Möglichkeiten, von
denen eine die oben betrachtete Annahme ist:
\[
[1\,8]\,[2\,8]\,[3\,8]\,[4\,8]\,[5\,8]\,[6\,8]\,[7\,8].
\]
Ist aber die Figur zweitheilig, so ist zunächst auszuschliessen,
dass der eine Theil aus einem oder aus zwei Strichen oder
einem dreistrahligen Sterne besteht; denn in diesen Fällen müsste
der andere Theil ein Stern mit sechs, fünf oder vier Strahlen
sein. Dazu aber bleiben von den acht Ziffern nicht mehr genug
übrig. Es bleibt also nur noch übrig, dass der eine Theil ein
Dreieck, der andere ein vierstrahliger Stern ist, $\triangle\,\weberfourup$, und diese
Annahme ist auch in der That immer zulässig. Ein Repräsentant
eines solchen Systems ist:
\[
[1\,2]\,[1\,3]\,[2\,3]\,[4\,5]\,[4\,6]\,[4\,7]\,[4\,8].
\]
Das Dreieck kann man auf $8.7.6 : 2.3 = 56$ verschiedene
Arten wählen, und zu jedem Dreieck giebt es noch fünf Möglichkeiten
, den vierstrahligen Stern anzunehmen, da man jeden der
übrigen fünf Punkte zum Centrum machen kann. Die Anzahl
%CONT>>


% ===== scan leaf 395, printed 374 =====
% !PRINTED 374
\weberpage{374}
\pdfbookmark[2]{§ 100. Rationale Bestimmung der Curve aus einem vollständigen Siebener-Systeme}{bm115}
% Source-faithful Cayley configuration signs: a four-stroke star and a seven-rayed star.
\providecommand{\weberfourup}{\mathord{\tikz[baseline=-0.2ex,x=0.23em,y=0.23em,line width=0.45pt]{\draw (0,0)--(-1.05,1) (0,0)--(-0.35,1.08) (0,0)--(0.35,1.08) (0,0)--(1.05,1);}}}
\providecommand{\webersevenstar}{\mathord{\tikz[baseline=-0.35ex,x=0.21em,y=0.21em,line width=0.45pt]{\draw (0,0)--(90:1) (0,0)--(38.57:1) (0,0)--(-12.86:1) (0,0)--(-64.29:1) (0,0)--(-115.71:1) (0,0)--(-167.14:1) (0,0)--(141.43:1);}}}
%<<CONT
dieser Bestimmungen ist daher 280, und wir kommen zu folgendem Resultate:

\gesperrt{Es giebt im Ganzen 288 Aronhold'sche Systeme, deren Zeichen eine der beiden Gestalten hat}

\[
\webersevenstar;\ \Delta\,\weberfourup.
\]

Aus diesen Zeichen darf man aber nicht etwa schliessen, dass diese Siebener-Systeme in zwei verschiedene Arten zerfallen. Der Unterschied der beiden Figuren liegt nur in der Bezeichnung. In der That können wir ja von einem ganz beliebigen der vollständigen Siebener-Systeme ausgehen, um die Bezeichnung abzuleiten.

Hiernach findet man leicht die Bezeichnung für die \gesperrt{Steiner'schen} Complexe, die nach einem der beiden folgenden Typen zu bilden sind:

\[
[1\,2][3\,4],\ [1\,3][2\,4],\ [1\,4][2\,3],\ [5\,6][7\,8],\ [5\,7][6\,8],\ [5\,8][6\,7],
\]
\[
[1\,7][1\,8],\ [2\,7][2\,8],\ [3\,7][3\,8],\ [4\,7][4\,8],\ [5\,7][5\,8],\ [6\,7][6\,8].
\]

\begin{center}
\S.\,100.

\gesperrt{Rationale Bestimmung der Curve aus einem vollständigen Siebener-Systeme.}
\end{center}

Die grosse Bedeutung der \gesperrt{Aronhold'schen} Systeme für das Problem der Doppeltangenten spricht sich in folgenden beiden Sätzen aus:

I. \gesperrt{Ist bei einer Curve vierter Ordnung ein vollständiges Siebener-System gegeben, so können daraus alle übrigen Doppeltangenten rational bestimmt werden.}

II. \gesperrt{Sind sieben beliebige gerade Linien in einer Ebene gegeben, so kann man im Allgemeinen, d.\,h. wenn gewisse rationale Functionen von den Co\"efficienten in den Gleichungen dieser Geraden nicht verschwinden, auf rationalem Wege die Gleichung einer Curve vierter Ordnung ohne singulären Punkt bestimmen, für die die gegebenen sieben Linien ein Aronhold'sches System bilden.}

Um den ersten Satz zu beweisen, würde es bei der vollkommenen Gleichberechtigung der Ziffern 1 bis 7 genügen, die
%CONT>>


% ===== scan leaf 396, printed 375 =====
% !PRINTED 375
\weberpage{375}
%<<CONT
Bestimmung von einer achten Doppeltangente aus einem gegebenen vollständigen Systeme durchzuführen. Wir bestimmen aber besser gleichzeitig drei. Es sei also jetzt

\weq{1}{x_1,\ x_2,\ x_3,\ x_4,\ x_5,\ x_6,\ x_7}

ein als bekannt vorausgesetztes vollständiges Siebener-System. Wir denken uns die drei \gesperrt{Steiner'schen} Complexe gebildet, in denen je sechs der Doppeltangenten (1), mit Ausschluss zuerst von $x_1$, dann von $x_2$, zuletzt von $x_3$ enthalten sind.

Die darin neu hinzutretenden 15 Doppeltangenten bezeichnen wir mit dem Buchstaben $\xi$ und erhalten diese drei Complexe [nach \S.\,99, (1), (2)] in der Gestalt:

\weq{2}{\begin{aligned}
x_2\xi_3 &,\ x_3\xi_2,\ x_4\xi_{41},\ x_5\xi_{51},\ x_6\xi_{61},\ x_7\xi_{71}\\
x_3\xi_1 &,\ x_1\xi_3,\ x_4\xi_{42},\ x_5\xi_{52},\ x_6\xi_{62},\ x_7\xi_{72}\\
x_1\xi_2 &,\ x_2\xi_1,\ x_4\xi_{43},\ x_5\xi_{53},\ x_6\xi_{63},\ x_7\xi_{73},
\end{aligned}}

und nach der Bezeichnung des \S.\,99 ist

\[
\xi_1 = [2\,3],\ \xi_2 = [3\,1],\ \xi_3 = [1\,2],\ \xi_{41} = [1\,4],\ \ldots
\]

Aus (2) ergiebt sich noch ein \gesperrt{Steiner'scher} Complex, der mit jedem der Complexe (2) ein syzygetisches Paar bildet:

\weq{3}{x_1\xi_1,\ x_2\xi_2,\ x_3\xi_3,\ \ldots}

und dieser Complex enthält keine der Doppeltangenten $x_4, x_5, x_6, x_7$.

Indem wir nun die gesuchten Functionen $\xi_1, \xi_2, \xi_3$ mit den geeigneten constanten Factoren multiplicirt annehmen, können wir die Gleichung der Curve vierter Ordnung nach \S.\,96, (9) in die Form

\weq{4}{\sqrt{x_1\xi_1} + \sqrt{x_2\xi_2} + \sqrt{x_3\xi_3} = 0}

setzen, und also die rationale Function $f$, die, gleich Null gesetzt, die Gleichung der Curve giebt, in jeder der drei mit einander identischen Formen

\weq{5}{f = 4x_2\xi_2x_3\xi_3 - u_1^2 = 4x_3\xi_3x_1\xi_1 - u_2^2 = 4x_1\xi_1x_2\xi_2 - u_3^2,}

annehmen, worin

\weq{6}{\begin{aligned}
u_1 &= -x_1\xi_1 + x_2\xi_2 + x_3\xi_3\\
u_2 &= x_1\xi_1 - x_2\xi_2 + x_3\xi_3\\
u_3 &= x_1\xi_1 + x_2\xi_2 - x_3\xi_3
\end{aligned}}

gesetzt ist. Nun bilden [nach (2)] auch $x_2\xi_3\, x_4\xi_{41}$ ein syzygetisches Quadrupel, und folglich können wir, wenn über einen
%CONT>>


% ===== scan leaf 397, printed 376 =====
% !PRINTED 376
\weberpage{376}
%<<CONT
constanten Factor in $\xi_{41}$ verfügt wird, $f$ auch in der Form darstellen:

\weq{7}{f = 4x_2\xi_3\,x_4\xi_{41} - v^2,}

worin $v$ eine quadratische Form ist. Hieraus und aus der ersten Darstellung in (5) ergiebt sich aber die Identität

\[
4x_2\xi_3(x_3\xi_2 - x_4\xi_{41}) = (u_1 - v)(u_1 + v),
\]

und daraus schliessen wir, dass einer der beiden Factoren rechts durch $x_2\xi_3$ theilbar sein muss [\S.\,96, (3)]. Nehmen wir an, es sei dies $u_1 - v$, was durch Verfügung über das Vorzeichen von $v$ erreicht werden kann, so folgt, dass ein von Null verschiedener constanter Factor $\lambda$ existiren muss, so dass
% [sic] The source prints \lambda here; the following display and later prose use \lambda_1.

\[
\begin{aligned}
u_1 - v &= 2\lambda_1 x_2\xi_3,\\
u_1 + v &= \frac{2(x_3\xi_2 - x_4\xi_{41})}{\lambda_1},
\end{aligned}
\]

woraus

\[
u_1 = \lambda_1 x_2\xi_3 + \frac{x_3\xi_2 - x_4\xi_{41}}{\lambda_1}
\]

folgt. Ersetzen wir hierin $u_1$ durch seinen Ausdruck aus (6), so ergiebt sich

\[
x_4\xi_{41} = x_3\xi_2 - \lambda_1(-x_1\xi_1 + x_2\xi_2 + x_3\xi_3) + \lambda_1^2 x_2\xi_3.
\]

Solcher Gleichungen erhalten wir zunächst drei, wenn wir die Indices 1, 2, 3 cyklisch vertauschen und an Stelle der unbekannten Constanten $\lambda_1$ drei Constanten $\lambda_1, \lambda_2, \lambda_3$ setzen:

\weq{8}{\begin{aligned}
x_4\xi_{41} &= x_3\xi_2 - \lambda_1(-x_1\xi_1 + x_2\xi_2 + x_3\xi_3) + \lambda_1^2 x_2\xi_3\\
x_4\xi_{42} &= x_1\xi_3 - \lambda_2(x_1\xi_1 - x_2\xi_2 + x_3\xi_3) + \lambda_2^2 x_3\xi_1\\
x_4\xi_{43} &= x_2\xi_1 - \lambda_3(x_1\xi_1 + x_2\xi_2 - x_3\xi_3) + \lambda_3^2 x_1\xi_2.
\end{aligned}}

Um die Constanten $\lambda_1, \lambda_2, \lambda_3$ zu bestimmen, dividiren wir die beiden letzten dieser Gleichungen mit $\lambda_2$ und $\lambda_3$ und addiren. Dadurch folgt die identische Gleichung

\weq{9}{\begin{aligned}
x_4\left(\frac{\xi_{42}}{\lambda_2} + \frac{\xi_{43}}{\lambda_3}\right) &=\\
x_1\left(-2\xi_1 + \lambda_3\xi_2 + \frac{\xi_3}{\lambda_2}\right) &+ \xi_1\left(\lambda_2 x_3 + \frac{x_2}{\lambda_3}\right).
\end{aligned}}

Da nun $x_4, x_1, \xi_1$ azygetisch sind, und folglich nicht durch einen Punkt gehen können, so muss die Linie

\[
\lambda_2 x_3 + \frac{x_2}{\lambda_3} = 0
\]
%CONT>>


% ===== scan leaf 398, printed 377 =====
% !PRINTED 377
\weberpage{377}
%<<CONT
durch den Schnitt von $x_1$ und $x_4$ gehen, und folglich ist $x_4$ aus $x_1$ und $\lambda_2 x_3 + \dfrac{x_2}{\lambda_3}$ linear zusammengesetzt. Nun aber können wir $x_4$ linear durch $x_1, x_2, x_3$ ausdrücken, etwa in der Form

\weq{10}{x_4 = a_1 x_1 + a_2 x_2 + a_3 x_3,}

und darin sind $a_1, a_2, a_3$ als bekannt zu betrachten, und keine dieser Constanten kann verschwinden. Es ist also $\lambda_2\lambda_3 = a_3 : a_2$, und wenn wir eine neue Constante $h_1$ einführen, so ist

\weq{11}{\lambda_2 = h_1 a_3,\quad \frac{1}{\lambda_3} = h_1 a_2,}

und die Gleichung (9) ergiebt die Identität:

\[
h_1 x_4\left(\frac{\xi_{42}}{\lambda_2} + \frac{\xi_{43}}{\lambda_3} - h_1\xi_1\right) = x_1\left(-h_1(2+a_1h_1)\xi_1 + \frac{\xi_2}{a_2} + \frac{\xi_3}{a_3}\right).
\]

Daraus schliesst man weiter, wenn $k_1$ eine neue Constante ist,

\[
k_1 x_4 = -h_1(2+a_1h_1)\xi_1 + \frac{\xi_2}{a_2} + \frac{\xi_3}{a_3},
\]

\weq{12}{\frac{\xi_{42}}{\lambda_2} + \frac{\xi_{43}}{\lambda_3} = h_1\xi_1 + \frac{k_1}{h_1} x_1.}

Wenn man hierin die Ziffern 1, 2, 3 cyklisch vertauscht, so ergeben sich aus (8) drei solche Systeme; zunächst:

\weq{13}{\begin{aligned}
k_1 x_4 &= -h_1(2+a_1h_1)\xi_1 + \frac{\xi_2}{a_2} + \frac{\xi_3}{a_3}\\
k_2 x_4 &= \frac{\xi_1}{a_1} - h_2(2+a_2h_2)\xi_2 + \frac{\xi_3}{a_3}\\
k_3 x_4 &= \frac{\xi_1}{a_1} + \frac{\xi_2}{a_2} - h_3(2+a_3h_3)\xi_3.
\end{aligned}}

Da diese drei Ausdrücke für $x_4$ mit einander identisch sein müssen, so folgt:

\[
-\frac{h_1(2+a_1h_1)}{k_1} = \frac{1}{k_2 a_1} = \frac{1}{k_3 a_1},
\]

also $k_2 = k_3$ und ebenso $k_2 = k_1$. Es sind also $k_1, k_2, k_3$ einander gleich und wir setzen dafür $k$. Dann folgt weiter

\[
a_1 h_1(2+a_1h_1) + 1 = (a_1h_1+1)^2 = 0,
\]

also $h_1 a_1 = -1$, und ebenso

\[
h_1 = \frac{-1}{a_1},\quad h_2 = \frac{-1}{a_2},\quad h_3 = \frac{-1}{a_3}.
\]


% ===== scan leaf 399, printed 378 =====
% !PRINTED 378
\weberpage{378}
Demnach liefern die Formeln (13) übereinstimmend

\[
k\,x_4 = \frac{\xi_1}{a_1} + \frac{\xi_2}{a_2} + \frac{\xi_3}{a_3}
\]
oder
\weq{14}{k\,(a_1\,x_1 + a_2\,x_2 + a_3\,x_3) = \frac{\xi_1}{a_1} + \frac{\xi_2}{a_2} + \frac{\xi_3}{a_3}.}
Aus (11) aber folgt noch
\weq{15}{\lambda_1 = -\frac{a_2}{a_3}, \quad \lambda_2 = -\frac{a_3}{a_1}, \quad \lambda_3 = -\frac{a_1}{a_2}.}
Aus (12) ergeben sich dann ferner die drei Relationen
\[
\begin{aligned}
-\frac{\xi_{42}}{\lambda_2} - \frac{\xi_{43}}{\lambda_3} &= \frac{\xi_1}{a_1} + k\,a_1\,x_1 \\
-\frac{\xi_{43}}{\lambda_3} - \frac{\xi_{41}}{\lambda_1} &= \frac{\xi_2}{a_2} + k\,a_2\,x_2 \\
-\frac{\xi_{41}}{\lambda_1} - \frac{\xi_{42}}{\lambda_2} &= \frac{\xi_3}{a_3} + k\,a_3\,x_3.
\end{aligned}
\]
Daraus durch Addition mit Rücksicht auf (14)
\[
\frac{\xi_{41}}{\lambda_1} + \frac{\xi_{42}}{\lambda_2} + \frac{\xi_{43}}{\lambda_3} = -k\,(a_1\,x_1 + a_2\,x_2 + a_3\,x_3),
\]
und folglich
\weq{16}{
\begin{aligned}
\frac{\xi_{41}}{\lambda_1} &= \frac{\xi_1}{a_1} - k\,(a_2\,x_2 + a_3\,x_3) \\
\frac{\xi_{42}}{\lambda_2} &= \frac{\xi_2}{a_2} - k\,(a_3\,x_3 + a_1\,x_1) \\
\frac{\xi_{43}}{\lambda_3} &= \frac{\xi_3}{a_3} - k\,(a_1\,x_1 + a_2\,x_2).
\end{aligned}
}
Unbekannt ist in diesen Formeln noch die Constante $k$. Diese bestimmen wir aus der Bemerkung, dass wir das ganze bisher betrachtete Formelsystem vervierfachen können, indem wir an Stelle von $x_4$ treten lassen $x_4$, $x_5$, $x_6$, $x_7$.

Der Formel (10) entsprechend wollen wir diese vier Functionen linear durch $x_1$, $x_2$, $x_3$ ausdrücken in der Weise:
\weq{17}{
\begin{aligned}
x_4 &= a_{4,1}\,x_1 + a_{4,2}\,x_2 + a_{4,3}\,x_3 \\
x_5 &= a_{5,1}\,x_1 + a_{5,2}\,x_2 + a_{5,3}\,x_3 \\
x_6 &= a_{6,1}\,x_1 + a_{6,2}\,x_2 + a_{6,3}\,x_3 \\
x_7 &= a_{7,1}\,x_1 + a_{7,2}\,x_2 + a_{7,3}\,x_3,
\end{aligned}
}
%CONT>>


% ===== scan leaf 400, printed 379 =====
% !PRINTED 379
\weberpage{379}
%<<CONT
worin die Co\"efficienten $a$ als bekannt anzusehen sind. Dann bekommen wir aus (14) vier Relationen, wenn wir an Stelle von $k$ vier verschiedene Constanten $k_4$, $k_5$, $k_6$, $k_7$ treten lassen:
\weq{18}{
\begin{aligned}
k_4\,(a_{4,1}\,x_1+a_{4,2}\,x_2+a_{4,3}\,x_3) &= \frac{\xi_1}{a_{4,1}}+\frac{\xi_2}{a_{4,2}}+\frac{\xi_3}{a_{4,3}} \\
k_5\,(a_{5,1}\,x_1+a_{5,2}\,x_2+a_{5,3}\,x_3) &= \frac{\xi_1}{a_{5,1}}+\frac{\xi_2}{a_{5,2}}+\frac{\xi_3}{a_{5,3}} \\
k_6\,(a_{6,1}\,x_1+a_{6,2}\,x_2+a_{6,3}\,x_3) &= \frac{\xi_1}{a_{6,1}}+\frac{\xi_2}{a_{6,2}}+\frac{\xi_3}{a_{6,3}} \\
k_7\,(a_{7,1}\,x_1+a_{7,2}\,x_2+a_{7,3}\,x_3) &= \frac{\xi_1}{a_{7,1}}+\frac{\xi_2}{a_{7,2}}+\frac{\xi_3}{a_{7,3}},
\end{aligned}
}
und nun sind die Constanten $k$ so zu bestimmen, dass von diesen vier Gleichungen die eine aus den drei anderen folgt. Die Bedingungen dafür kann man in symmetrischer Weise dadurch bilden, dass man vier Factoren $l_4$, $l_5$, $l_6$, $l_7$ einführt, deren Verhältnisse man aus den Gleichungen bestimmt:
\weq{19}{\frac{l_4}{a_{4,i}} + \frac{l_5}{a_{5,i}} + \frac{l_6}{a_{6,i}} + \frac{l_7}{a_{7,i}} = 0, \quad i = 1, 2, 3,}
und die dann auch den drei Gleichungen
\weq{20}{k_4\,l_4\,a_{4,i} + k_5\,l_5\,a_{5,i} + k_6\,l_6\,a_{6,i} + k_7\,l_7\,a_{7,i} = 0, \quad i = 1, 2, 3}
genügen müssen, woraus die Verhältnisse der $k$ bestimmt sind. Ein gemeinschaftlicher Factor der vier Grössen $k$ bleibt der Natur der Sache nach unbestimmt und kann beliebig angenommen werden. Hiernach können aus den Gleichungen (18) die Functionen $\xi_1$, $\xi_2$, $\xi_3$ rational bestimmt werden, und durch (16) sind dann auch die Functionen $\xi_{4i}$, $\xi_{5i}$, $\xi_{6i}$, $\xi_{7i}$ bestimmt.

Es fehlen noch sechs Doppeltangenten, die man durch geeignete Permutationen unter den Functionen des Systemes (1) erhalten kann. Damit ist der an die Spitze gestellte Satz I. bewiesen.

\par\noindent\rule{0.3\textwidth}{0.4pt}\par

Um auch die Richtigkeit des Satzes II. einzusehen, braucht man nur unsere Analyse rückwärts zu verfolgen, indem man die Co\"efficienten $a_{k,i}$ als unabhängige Variable ansieht. Dann sind durch die Gleichungen (18), (19), (20) die Functionen $\xi_1$, $\xi_2$, $\xi_3$
%CONT>>


% ===== scan leaf 401, printed 380 =====
% !PRINTED 380
\weberpage{380}
\pdfbookmark[2]{§ 101. Die Galois'sche Gruppe des Doppeltangentenproblems}{bm116}
%<<CONT
rational durch diese $a_{k,i}$ bestimmt, und aus (16) und (17) erhält man sodann $\xi_{4i}$, $\xi_{5i}$, $\xi_{6i}$, $\xi_{7i}$, $x_4$, $x_5$, $x_6$, $x_7$.

Durch Substitution der $\xi_1$, $\xi_2$, $\xi_3$ in die Gleichung (4) erhält man die Gleichung einer Curve vierter Ordnung, deren Co\"efficienten rationale Functionen der $a_{k,i}$ sind, und die Discriminante dieser Gleichung kann nicht identisch verschwinden, weil man umgekehrt, wie wir gesehen haben, aus der Gleichung einer Curve vierter Ordnung ohne singulären Punkt ein Gleichungssystem (18), (19), (20) ableiten kann.

Aus den Gleichungen (18), (16) folgen dann auch die Formeln (7), und die daraus durch Vertauschung von 4 mit 5, 6, 7 hervorgehenden, woraus zu schliessen ist, dass $x_4$, $x_5$, $x_6$, $x_7$ zusammen mit $x_1$, $x_2$, $x_3$ ein vollständiges Siebener-System bilden.

\S.\,101.

\gesperrt{Die Galois'sche Gruppe des Doppeltangentenproblems.}

Die Sätze, die wir abgeleitet haben, sind ausreichend, um die \gesperrt{Galois'sche} Gruppe der algebraischen Gleichung zu bestimmen, von der die Doppeltangenten abhängen. Wir betrachten hierbei als \gesperrt{Rationalitätsbereich den Körper, der aus allen rationalen Functionen der 14 Verhältnisse der Co\"efficienten einer allgemeinen ternären Form $4^{\text{ten}}$ Grades} besteht, worin diese Co\"efficienten als \gesperrt{unabhängige Variable} gelten. Die Gleichung $28^{\text{sten}}$ Grades können wir uns dann etwa in der Weise gebildet denken, dass wir als Unbekannte die Abscissen der Schnittpunkte der Doppeltangenten mit einer beliebigen festen geraden Linie $L$ betrachten. Durch die Wurzeln $\xi$ dieser Gleichung, die wir die \gesperrt{Doppeltangentengleichung} nennen, sind dann die Doppeltangenten rational darstellbar.

Benutzt man ein \gesperrt{Cartesisches} Coordinatensystem $x$, $y$, dessen $x$-Axe die Linie $L$ ist, so erhält die Gleichung einer Doppeltangente die Gestalt
\weq{1}{y = \Theta\,(x - \xi),}
und die Doppeltangentengleichung kann gebildet werden, wenn $F(x,y) = 0$ die Gleichung der Curve vierter Ordnung ist, indem man die Bedingungen aufsucht, dass die Function von $x$
%CONT>>


% ===== scan leaf 402, printed 381 =====
% !PRINTED 381
\weberpage{381}
%<<CONT
$4^{\text{ten}}$ Grades, $F[x, \Theta(x-\xi)]$ ein vollständiges Quadrat sei. Dies giebt zwei Gleichungen zwischen $\xi$ und $\Theta$, aus denen man durch Elimination von $\Theta$ die Doppeltangentengleichung erhält. Da zu jedem $\xi$ nur ein Werth von $\Theta$ gehört (so lange sich nicht zwei Doppeltangenten auf der Linie $L$ schneiden), so kann $\Theta$ rational durch $\xi$ ausgedrückt werden, und zwar in einer Form
\weq{2}{\Theta = \varphi\,(\xi),}
die für jedes zusammengehörige Paar $\xi$, $\Theta$ gilt.

Die Wurzeln der Doppeltangentengleichung ordnen sich ebenso, wie die entsprechenden Doppeltangenten in Complexe, Siebener-Systeme u.\,s.\,w. Wir bezeichnen diese Wurzeln ebenso wie die Doppeltangenten in \S.\,99 durch das Symbol $[ik]$, worin $i$, $k$ die Paare der Ziffern 1 bis 8 durchlaufen und $[ik] = [ki]$ ist.

Betrachten wir irgend zwei von den Doppeltangenten, $p$, $q$, als bekannt, so können wir auf rationalem Wege die Gleichung $u = 0$ eines Kegelschnittes daraus ableiten, der durch die vier Berührungspunkte dieser Doppeltangenten geht, und wir können also die Gleichung \S.\,96, (2)
\[
f = p\,q\,v - u^2
\]
in rationaler Form aufstellen. Dann giebt es nach \S.\,96, (4) unter der Kegelschnittschaar
\[
v + 2\,\lambda\,u + \lambda^2\,p\,q = 0
\]
fünf, die in ein Linienpaar zerfallen, und wenn wir das Product
\[
\prod_{\lambda} (v + 2\,\lambda\,u + \lambda^2\,p\,q) = \Phi
\]
bilden, über die fünf Wurzeln der Gleichung $5^{\text{ten}}$ Grades, von der $\lambda$ abhängt [\S.\,96, (6)], so ist dies Product gleichfalls rational durch die Co\"efficienten von $f$ und von $p$, $q$ ausdrückbar. Dies Product ist aber eine Form $10^{\text{ten}}$ Grades, die in zehn lineare Factoren zerfällt, die mit $p\,q$ zusammen einen \gesperrt{Steiner'schen} Complex bilden.

Die Co\"efficienten in $\Phi$ können nun auch rational durch die beiden den $p$, $q$ entsprechenden Wurzeln $\xi_1$, $\xi_2$ der Doppeltangentengleichung ausgedrückt werden, und wenn wir dann die Abscissen der Schnittpunkte der Linie $L$ mit $\Phi = 0$ aufsuchen, so erhalten wir eine Gleichung $10^{\text{ten}}$ Grades für $\xi$
\[
X(\xi_1, \xi_2, \xi) = 0,
\]
%CONT>>


% ===== scan leaf 403, printed 382 =====
% !PRINTED 382
\weberpage{382}
%<<CONT
deren Wurzeln die zehn mit $\xi_1\xi_2$ syzygetischen Wurzeln sind. Bedeutet $\xi_3$ eine von diesen, so ist also

\weq{3}{X(\xi_1,\xi_2,\xi_3) = 0,}

und es folgt daraus der Satz

\gesperrt{1. Es giebt eine rationale Function von drei Variablen $X(\xi_1,\xi_2,\xi_3)$, die verschwindet, wenn $\xi_1\xi_2\xi_3$ irgend ein syzygetisches Tripel von Wurzeln der Doppeltangentengleichung ist, und die nicht verschwindet, wenn $\xi_1\xi_2\xi_3$ ein azygetisches Tripel ist.}

Nach \S.\,100 können durch ein vollständiges Siebener-System dieser Wurzeln

\weq{4}{\xi_1,\ \xi_2,\ \xi_3,\ \xi_4,\ \xi_5,\ \xi_6,\ \xi_7}

alle Wurzeln rational ausgedrückt werden, und zwar in der Weise, dass z.\,B.

\weq{5}{\xi_{12} = \Psi(\xi_1,\xi_2 \mid \xi_3,\xi_4,\xi_5,\xi_6,\xi_7)}

eine Wurzel wird, wo $\Psi$ eine rationale Function bedeutet, die sich nicht ändert, wenn $\xi_1$ und $\xi_2$ vertauscht oder wenn $\xi_3,\xi_4,\xi_5,\xi_6,\xi_7$ beliebig permutirt werden. Wird aber in (5) bei festgehaltener Function $\Psi$ an Stelle von $\xi_1,\xi_2$ ein anderes Paar $\xi_i,\xi_k$ gesetzt, so erhält man eine andere Wurzel $\xi_{ik}$. Die Wurzeln (4) sind auch mit $[1\,8],\dots[7\,8]$ und $\xi_{ik}$ mit $[i\,k]$ zu bezeichnen.

Durch jede Permutation der Wurzeln (1) wird nach (2) das ganze System der Wurzeln $[i\,k]$ eine gewisse Permutation erfahren.

Ersetzt man aber das vollständige Siebener-System (4) durch ein anderes, so ergiebt die Formel (5) eine bestimmte andere Wurzel, und das ganze System der Wurzeln $[i\,k]$ wird einer zweiten Permutation unterworfen. Es ist dann zunächst leicht zu beweisen:

\gesperrt{2. Die Permutationsgruppe $P$ der Wurzeln der Doppeltangentengleichung, die man erhält, wenn man in (4) und (5) an Stelle von $\xi_1,\xi_2,\dots\xi_7$ alle vollständigen Systeme, jedes in jeder beliebigen Ordnung, setzt, ist die Galois'sche Gruppe der Doppeltangentengleichung.}

Um dies nachzuweisen, haben wir Zweierlei zu zeigen:
%CONT>>


% ===== scan leaf 404, printed 383 =====
% !PRINTED 383
\weberpage{383}
%<<CONT
\gesperrt{a) Jede Permutation $\pi$ der Wurzeln der Doppeltangentengleichung, die auf alle rationalen Gleichungen zwischen diesen Wurzeln anwendbar ist, gehört zu $P$.}

Dies ergiebt sich so: Wenn $\pi$ auf alle rationalen Gleichungen zwischen den Wurzeln anwendbar ist, so gilt dasselbe von den Potenzen von $\pi$. Nach 1. kann niemals durch $\pi$ ein syzygetisches Tripel in ein azygetisches oder umgekehrt übergeführt werden; denn $\pi$ ist, wenn $\xi_1\xi_2\xi_3$ ein syzygetisches Tripel ist, auf die Gleichung (3) anwendbar; also kann $\xi_1\xi_2\xi_3$ nicht in ein azygetisches Tripel übergehen. Und auch das Umgekehrte ist nicht möglich, weil sonst durch $\pi^{-1}$ ein syzygetisches in ein azygetisches Tripel übergeführt würde. Daher geht auch durch $\pi$ irgend ein vollständiges Siebener-System wieder in ein solches System in irgend welcher Anordnung über, und wenn man dann $\pi$ auf alle Gleichungen von der Form (5) anwendet, so ergiebt sich eine Permutation der $\xi_i, \xi_{ik}$, die zu $P$ gehört.

\gesperrt{b) Jede rationale Gleichung zwischen den Wurzeln der Doppeltangentengleichung gestattet alle Permutationen der Gruppe $P$.}

Eine rationale Relation zwischen den Wurzeln hat die Form

\weq{6}{\Phi(\xi_1,\dots,\xi_7,\xi_{12},\dots,a\dots) = 0,}

worin $\Phi$ eine rationale Function ist, und $a\dots$ die Verhältnisse der Coëfficienten von $f$ bedeuten. Hierin kann man durch (5) die $\xi_{12},\xi_{13},\dots$ rational durch $\xi_1,\xi_2,\dots,\xi_7$ ausdrücken, und nach \S.\,100, II., nachdem die Curve $f$ durch ein vollständiges Siebener-System ihrer Doppeltangenten rational bestimmt ist, lassen sich dann die $a$ rational (mit nur numerischen Coëfficienten) durch die sieben Grössenpaare (2)

\[
\xi_1, \Theta_1;\ \xi_2, \Theta_2;\ \dots;\ \xi_7, \Theta_7
\]

ausdrücken. Da diese aber ganz beliebig gegeben sein können, so muss die Gleichung (6) durch diese Substitutionen in eine Identität übergehen. Die Gleichung (6) muss also auch richtig bleiben, wenn man darin die $\xi_1,\xi_2,\dots,\xi_7$ durch ein anderes Siebener-System oder auch durch dasselbe in anderer Ordnung ersetzt, und gleichzeitig unter den $\xi_{12},\dots$ die durch die Formel (5) vorgeschriebene Permutation vornimmt, d.\,h. wenn man unter den Wurzeln der Doppeltangentengleichung eine Permutation der Gruppe $P$ ausführt.


% ===== scan leaf 405, printed 384 =====
% !PRINTED 384
\weberpage{384}
\pdfbookmark[2]{§ 102. Darstellung der Gruppe}{bm117}
\providecommand{\weberfourup}{\mathord{\tikz[baseline=-0.2ex,x=0.23em,y=0.23em,line width=0.45pt]{\draw (0,0)--(-1.05,1) (0,0)--(-0.35,1.08) (0,0)--(0.35,1.08) (0,0)--(1.05,1);}}}
\begin{center}
\S.\,102.

\gesperrt{Darstellung der Gruppe.}
\end{center}

Der Grad der Gruppe $P$ ist sofort anzugeben. Da es 288 vollständige Siebener-Systeme giebt, und da die Elemente eines solchen Systemes auf $\Pi(7)$ Arten permutirt werden können, so ist der Grad der Gruppe:

\[
288\,\Pi(7) = 36\,\Pi(8) = 1\,451\,520.
\]

Bei der Bildung der Permutationen der Wurzeln benutzen wir zweierlei Bezeichnung. Zunächst

\weq{1}{\xi_1,\ \xi_2,\ \dots,\ \xi_7,\ \xi_{ik},}

worin $i,k$ von 1 bis 7 geht, und die einheitliche Bezeichnung $[i\,k]$, wobei $i,k$ von 1 bis 8 geht, und wobei $\xi_1,\xi_2\dots$ durch $[1\,8],[2\,8],\dots$ zu bezeichnen sind.

Wenn wir nun zwei Ziffern aus der Reihe $1,2,\dots,7$ permutiren, z.\,B. 1 mit 2, so geht $\xi_1$ in $\xi_2$ über, und nach der Definition \S.\,99 bleibt $\xi_{12}$ ungeändert, $\xi_{13}$ geht in $\xi_{23}$ über u.\,s.\,f.

Wenn wir also in der Reihe der Wurzeln $[i\,k]$ die Ziffern 1 bis 7 beliebig permutiren, so erhalten wir lauter Permutationen der Gruppe $P$.

Nun haben wir im \S.\,100 gesehen, dass bei der Anordnung in syzygetische und azygetische Tripel und in Folge dessen auch in Complexe und vollständige Systeme die Ziffer 8 mit den übrigen Ziffern 1 bis 7 vollständig gleichberechtigt auftritt, und da die Zuordnung der Wurzel $\xi_{ik}$ zu dem Paare $\xi_i\xi_k$ [durch die Formel \S.\,101, (5)] nach \S.\,99 nur von dieser Anordnung abhängt, so können wir auch die Ziffern $1,2,\dots,7,8$ permutiren, ohne dass wir aus der Gruppe $P$ herauskommen. Es ist also in $P$ ein Theiler enthalten, der mit der symmetrischen Gruppe der Permutationen von acht Elementen isomorph ist, der also den Index 36 hat und den wir mit $S$ bezeichnen wollen.

Um die noch fehlenden Permutationen von $P$ zu bestimmen, lassen wir an Stelle des Siebener-Systemes $\xi_1,\dots,\xi_7$ ein neues vom Typus $\Delta\,\weberfourup$ treten, etwa so:

% [sic] The source prints a full stop after [6 8] in the first row.
\weq{2}{\begin{pmatrix} [1\,8], & [2\,8], & [3\,8], & [4\,8], & [5\,8], & [6\,8]. & [7\,8] \\ {}[2\,3], & [3\,1], & [1\,2], & [4\,8], & [5\,8], & [6\,8], & [7\,8] \end{pmatrix},}
%CONT>>


% ===== scan leaf 406, printed 385 =====
% !PRINTED 385
\weberpage{385}
%<<CONT
und bezeichnen die hierdurch bedingte Permutation in doppelter Weise mit

\weq{3}{\Pi_{1,2,3,8} = \Pi_{4,5,6,7},}

indem wir festsetzen, dass $\Pi_{\alpha_1\alpha_2\alpha_3\alpha_4}$ und $\Pi_{\beta_1\beta_2\beta_3\beta_4}$ dasselbe bedeuten sollen, wenn $\beta_1,\beta_2,\beta_3,\beta_4$ irgend eine Permutation von $\alpha_1,\alpha_2,\alpha_3,\alpha_4$ ist, oder wenn $\alpha_1,\alpha_2,\alpha_3,\alpha_4,\beta_1,\beta_2,\beta_3,\beta_4$ zusammen alle acht Ziffern umfassen.

Durch dies Zeichen ist die Vertauschung (2) eindeutig bezeichnet, und die Anzahl der verschiedenen Permutationen dieser Art beträgt genau 35, so dass wir die ganze Gruppe $P$ durch die Nebengruppen so darstellen können:

\weq{4}{P = S + \Sigma\, S\, \Pi_{\alpha_1\alpha_2\alpha_3\alpha_4},}

worin $\alpha_1,\alpha_2,\alpha_3,\alpha_4$ alle Systeme von vier Ziffern aus der Reihe $1,2,\dots,8$ durchlaufen, wobei eine beliebige Ziffer, z.\,B. $\alpha_4 = 8$, festgehalten werden kann.

Um den Einfluss von $\Pi_{1,2,3,8}$ auf irgend eine Wurzel $[i\,k]$ zu erkennen, genügt es, die drei Wurzeln $[1\,2], [1\,4], [4\,5]$ zu betrachten, weil $1,2,3$ einerseits, $4,5,6,7$ andererseits ganz gleichartig in $\Pi_{1,2,3,8}$ vorkommen.

Diese Vertauschungen erhält man einfach aus der Bemerkung, dass $[1\,2], [1\,4]$ nach \S.\,99 die in dem Complex

\weq{5}{[2\,8][1\,2],\ [3\,8][1\,3],\ [4\,8][1\,4],\ [5\,8][1\,5],\ [6\,8][1\,6],\ [7\,8][1\,7]}

mit $[2\,8]$ und $[4\,8]$ verbundenen Wurzeln sind, und dass ebenso $[4\,5]$ die mit $[5\,8]$ verbundene Wurzel in dem Complexe

\weq{6}{[1\,8][1\,4],\ [2\,8][2\,4],\ [3\,8][3\,4],\ [5\,8][5\,4],\ [6\,8][6\,4],\ [7\,8][7\,4]}

ist. Durch die Vertauschung (2) gehen aber die Complexe (5) und (6) in folgende über, wie man leicht aus der Darstellung der Complexe im \S.\,99 findet [der Complex (5) bleibt als Ganzes ungeändert]:

\[
[3\,1][3\,8],\ [2\,1][2\,8],\ [4\,8][1\,4],\ [5\,8][5\,1],\ [6\,8][6\,1],\ [7\,8][7\,1]
\]

\[
[2\,3][1\,4],\ [3\,1][2\,4],\ [1\,2][3\,4],\ [5\,8][6\,7],\ [6\,8][5\,7],\ [7\,8][5\,6],
\]

woraus man folgende durch $\Pi_{1,2,3,8}$ bewirkte Vertauschungen erhält:

\weq{7}{\begin{array}{l} [1\,2],\ [1\,4],\ [4\,5] \\ {}[3\,8],\ [1\,4],\ [6\,7]. \end{array}}

Aus (2) und (7) lässt sich nun folgende allgemeine und einfache Regel ableiten.


% ===== scan leaf 407, printed 386 =====
% !PRINTED 386
\weberpage{386}

3.\ \gesperrt{Bedeuten $\alpha_1, \alpha_2, \alpha_3, \alpha_4, \beta_1, \beta_2, \beta_3, \beta_4$ zusammen alle acht Ziffern, so hat man, um den Einfluss von}
\[
\Pi_{\alpha_1\alpha_2\alpha_3\alpha_4} = \Pi_{\beta_1\beta_2\beta_3\beta_4}
\]
\gesperrt{auf irgend eine Wurzel $[\mu\,\nu]$ zu bestimmen, zu unterscheiden, ob $\mu, \nu$ beide unter den $\alpha$ oder beide unter den $\beta$, oder ob die Ziffer $\mu$ unter den $\alpha$, $\nu$ unter den $\beta$ vorkommt. In den ersten Fällen geht $[\mu\,\nu]$ in $[\mu'\,\nu']$ über, wenn $\mu, \nu, \mu', \nu'$ entweder alle $\alpha$ oder alle $\beta$ bedeuten; im dritten Falle bleibt $[\mu\,\nu]$ ungeändert.}

In allen Fällen sind die Permutationen $\Pi_{\alpha_1\alpha_2\alpha_3\alpha_4}$ nur vom 2\textsuperscript{ten} Grade, führen also bei einmaliger Wiederholung zur Identität zurück.

Damit ist die Gruppe $P$ vollständig bestimmt und dargestellt, und um die Gesetze der Composition in $P$ festzustellen, sind nur noch wenige Formeln nöthig, die sich aus der oben aufgestellten Regel leicht ergeben. Dabei ist zu bemerken, dass man zwei von einander verschiedene der Permutationen $\Pi_{\alpha_1\alpha_2\alpha_3\alpha_4}$ immer so annehmen kann, dass sie im Index zwei oder drei Ziffern gemein haben, da man, wenn sie nur eine Ziffer gemein haben, für den Index der einen seine Ergänzung nehmen kann.

Es möge nun
\[
\sigma = \begin{pmatrix} 1, & 2, & 3, & 4, & 5, & 6, & 7, & 8 \\ \alpha_1, & \alpha_2, & \alpha_3, & \alpha_4, & \alpha_5, & \alpha_6, & \alpha_7, & \alpha_8 \end{pmatrix}
\]
irgend eine Permutation der acht Ziffern sein und $(\alpha, \beta)$ die Transposition der beiden Ziffern $\alpha$ und $\beta$, also ein Element aus $S$, bedeuten. Dann ist
\weq{8}{\Pi_{1,2,3,4}\,\sigma = \sigma\,\Pi_{\alpha_1\alpha_2\alpha_3\alpha_4},}
\weq{9}{\Pi_{1,2,3,4}\,\Pi_{1,2,3,4} = 1,}
\weq{10}{\Pi_{1,2,3,4}\,\Pi_{1,2,3,5} = (4,5)\,\Pi_{1,2,3,4},}
\weq{11}{\Pi_{1,2,3,4}\,\Pi_{1,2,5,6} = (1,2)(3,4)(5,6)(7,8)\,\Pi_{1,2,7,8}.}

Man beweist diese Formeln leicht nach der Regel 3., wenn man die einzelnen Fälle durchgeht, wobei natürlich nur eine ganz kleine Zahl von Typen zu betrachten sind; so geht z.\,B. $[1\,4]$ durch $\Pi_{1,2,3,4}$ in $[2\,3]$, dies durch $\Pi_{1,2,3,5}$ in $[1\,5]$ über, und $[1\,5]$ wird durch $\Pi_{1,2,3,4}$ nicht mehr geändert, folglich bewirkt
\weq{12}{\Pi_{1,2,3,4}\,\Pi_{1,2,3,5}\,\Pi_{1,2,3,4}}
%CONT>>


% ===== scan leaf 408, printed 387 =====
% !PRINTED 387
\weberpage{387}
\pdfbookmark[2]{§ 103. Einfachheit der Gruppe des Doppeltangentenproblems}{bm118}

%<<CONT
die Vertauschung von $[1\,4]$ mit $[1\,5]$ in Uebereinstimmung mit der Formel (10). Ebenso leicht erkennt man, dass z.\,B. $[1\,2]$ durch (12) nicht geändert wird.

\begin{center}
\S.\,103.

\gesperrt{Einfachheit der Gruppe des Doppeltangentenproblems.}
\end{center}

Die Darstellung der Gruppe $P$, die wir im vorigen Paragraphen entwickelt haben, liefert uns nun einen ganz einfachen Beweis dafür, dass diese Gruppe \gesperrt{keinen Normaltheiler} hat, also nach unserer früher gebrauchten Ausdrucksweise \gesperrt{einfach} ist, woraus dann folgt, dass die Gruppe nicht durch Adjunction von Irrationalitäten mit kleinerer Gruppe, also beispielsweise nicht durch cyklische Gleichungen erniedrigt werden kann.

Wir stellen $P$ in der Form (4), \S.\,102, dar:
\weq{1}{P = S + \Sigma\, S\, \Pi_{\alpha_1, \alpha_2, \alpha_3, \alpha_4},}
worin $S$ die ganze Gruppe aller Permutationen von acht Ziffern ist. Die Gruppe $S$ hat einen Normaltheiler $S'$ vom Index 2, nämlich die alternirende Gruppe der acht Ziffern, die ihrerseits einfach ist, und jeder Normaltheiler von $S$, der nicht aus der einzigen identischen Permutation besteht, muss die ganze Gruppe $S'$ enthalten. Wir setzen
\weq{2}{S = S' + S'',}
worin $S'' = S'\sigma$ ist, wenn $\sigma$ irgend eine Permutation der zweiten Art, z.\,B. eine Transposition bedeutet.

Wir nehmen jetzt an, es sei $Q$ ein Normaltheiler von $P$, der nicht aus der einzigen identischen Substitution besteht.

Der grösste gemeinschaftliche Theiler von $Q$ und $S$ ist dann ein Normaltheiler von $S$, und muss daher, wenn er nicht die identische Gruppe ist, die Gruppe $S'$ enthalten. Daraus folgt:

1.\ \gesperrt{Wenn $Q$ eine nicht identische Permutation aus $S$ enthält, so enthält $Q$ die ganze Gruppe $S'$.}

Wir beweisen sodann, dass $Q$, wenn es die Gruppe $S'$ enthält, mit $P$ identisch sein muss.

Wenn nämlich $Q$ die ganze Gruppe $S'$ enthält, so enthält
%CONT>>


% ===== scan leaf 409, printed 388 =====
% !PRINTED 388
\weberpage{388}

%<<CONT
es als Normaltheiler von $P$ auch, wenn $(4,5,6)$ ein dreigliedriger Cyklus aus $S'$ ist [nach \S.\,102, (8), (10)]:
\[
\Pi_{1,2,3,5}\,(5,4,6)\,\Pi_{1,2,3,5} = (5,4,6)\,\Pi_{1,2,3,4}\,\Pi_{1,2,3,5}
\]
\[
= (5,4,6)(4,5)\,\Pi_{1,2,3,4},
\]
und folglich auch
\[
S'(4,5)\,\Pi_{1,2,3,4} = S''\,\Pi_{1,2,3,4}.
\]
Ist nun
\[
\sigma = \begin{pmatrix} 1, & 2, & 3, & 4, & 5, & 6, & 7, & 8 \\ \alpha_1, & \alpha_2, & \alpha_3, & \alpha_4, & \alpha_5, & \alpha_6, & \alpha_7, & \alpha_8 \end{pmatrix}
\]
eine Permutation aus $S$, in der $\alpha_1, \alpha_2, \alpha_3, \alpha_4$ beliebig gegeben sind, so kann man die Anordnung der $\alpha_5, \alpha_6, \alpha_7, \alpha_8$ noch so wählen, dass $\sigma$ nach Belieben zu $S'$ oder zu $S''$ gehört. Wählt man $\sigma$ in $S'$ und beachtet, dass dann $S'' = S''\sigma^{-1}$ ist, so ergiebt sich, dass auch
\[
S''\,\sigma^{-1}\,\Pi_{1,2,3,4}\,\sigma = S''\,\Pi_{\alpha_1,\alpha_2,\alpha_3,\alpha_4}
\]
in $Q$ enthalten sein muss. Nun ist aber auch
\[
\Pi_{1,2,3,4}\,(3,5)(4,6)\,\Pi_{1,2,3,4} = (3,5)(4,6)\,\Pi_{1,2,5,6}\,\Pi_{1,2,3,4}
\]
\[
= (3,5)(4,6)(1,2)(3,4)(5,6)(7,8)\,\Pi_{1,2,7,8}
\]
\[
= (1,2)(3,6)(4,5)(7,8)\,\Pi_{1,2,7,8},
\]
und da $(1,2)(3,6)(4,5)(7,8)$ zu $S'$ gehört, so enthält $Q$ auch $S'\,\Pi_{1,2,7,8}$, und folglich, wie oben, alle $S'\,\Pi_{\alpha_1,\alpha_2,\alpha_3,\alpha_4}$, also auch alle $S\,\Pi_{\alpha_1\alpha_2\alpha_3\alpha_4}$, und mithin auch
\[
S\,\Pi_{\alpha_1,\alpha_2,\alpha_3,\alpha_4}\,\Pi_{\alpha_1,\alpha_2,\alpha_3,\alpha_4} = S,
\]
d.\,h. $Q$ umfasst die ganze Gruppe $P$. Daraus folgt in Verbindung mit 1.:

2.\ \gesperrt{Wenn ein Normaltheiler $Q$ von $P$ ausser der identischen Permutation noch irgend eine Permutation mit $S$ gemein hat, so ist $Q$ mit $P$ identisch.}

Es kann nun ferner die Frage sein, ob $Q$ eine Permutation aus einer der Nebengruppen von (1), also ein Element von der Form $\sigma\,\Pi_{1,2,3,4}$ enthalten kann, ohne mit $P$ identisch zu sein. Ist zunächst $\sigma = 1$, enthält also $Q$ das Element $\Pi_{1,2,3,4}$, so enthält es als Normaltheiler von $P$ auch alle anderen $\Pi_{\alpha_1\alpha_2\alpha_3\alpha_4}$, wie aus der Formel \S.\,102, (8) zu ersehen ist, und damit auch $\Pi_{1,2,3,4}\,\Pi_{1,2,3,5}\,\Pi_{1,2,3,4} = (4,5)$, und ist also nach 2. mit $P$ identisch.


% ===== scan leaf 410, printed 389 =====
% !PRINTED 389
\weberpage{389}

Ist aber $\sigma$ von 1 verschieden, so kann man ein Ziffernpaar $\alpha, \beta$, beide unter den $1,2,3,4$ oder beide unter den $5,6,7,8$, so wählen, dass die Transposition
\[
\sigma\,(\alpha,\beta)\,\sigma^{-1} = (\alpha', \beta')
\]
von $(\alpha, \beta)$ verschieden ist. Denn nimmt man zunächst für $\alpha$ eine durch $\sigma^{-1}$ veränderte Ziffer, so ist $\alpha'$ von $\alpha$ verschieden (Bd.\,I, \S.\,154, 6.), wählt man dann für $\beta$ eine Ziffer, die durch $\sigma^{-1}$ nicht in $\alpha$ übergeht und die mit $\alpha$ zugleich in der ersten oder in der zweiten Hälfte der acht Ziffern vorkommt, die gewiss immer existirt, weil nur eine Ziffer durch $\sigma^{-1}$ in $\alpha$ übergeht, so ist $(\alpha', \beta')$ von $(\alpha, \beta)$ verschieden. Dann folgt aber, dass in der Gruppe $Q$ das Element $(\alpha,\beta)\,\sigma\,\Pi_{1,2,3,4}\,(\alpha,\beta)$, und folglich auch [nach \S.\,102, (8), (9)]
\[
(\alpha,\beta)\,\sigma\,\Pi_{1,2,3,4}\,(\alpha,\beta)\,\Pi_{1,2,3,4}\,\sigma^{-1} = (\alpha,\beta)\,\sigma\,(\alpha,\beta)\,\sigma^{-1} = (\alpha,\beta)(\alpha',\beta')
\]
vorkommt, und dies ist eine von der Identität verschiedene Permutation aus $S$. Damit ist also bewiesen:

3.\ \gesperrt{Die Gruppe $P$ der Doppeltangentengleichung ist einfach.}

Als specielle Anwendung können wir hervorheben, dass die Gruppe $P$ unter den 28 Wurzeln der Doppeltangentengleichung nur Permutationen der ersten Art bewirken kann, und dass \gesperrt{folglich die Discriminante dieser Gleichung ein Quadrat} ist.

Bezeichnen wir nämlich für den Augenblick mit $G$ die Gruppe aller Permutationen der 28 Wurzeln und mit $A$ die darin als Normaltheiler enthaltene alternirende Gruppe, so ist der grösste gemeinschaftliche Theiler von $A$ und $P$ ein Normaltheiler von $P$ und muss also mit $P$ identisch sein; d.\,h. $P$ ist in $A$ enthalten.

Wenn man die Gruppe der Permutationen aufsucht, die eine der Wurzeln der Doppeltangentengleichung, etwa $[1\,2]$, ungeändert lassen, so findet man eine Gruppe, die für die übrigen 27 Wurzeln noch transitiv ist\footnote{Diese Gruppe ist nach einem geometrischen Satze von \gesperrt{Geiser} isomorph mit der Gruppe der Gleichung 27\textsuperscript{sten} Grades, von der die Lösung des Problems der 27 Geraden auf einer Fläche dritter Ordnung abhängt. (Mathem. Annalen, Bd.\,I.).}.


% ===== scan leaf 411, printed 390 =====
% !PRINTED 390
\weberpage{390}

Um dies nachzuweisen, genügt es zu zeigen, dass durch die
Permutationen dieser Gruppe irgend eine Wurzel, etwa [1\ 3], in
jede andere (mit Ausnahme von [1\ 2]) übergehen kann. Nun
geht aber [1,\,3] % [sic] comma printed here, unlike the bracket pairs elsewhere on this page which are space-separated without a comma
durch Permutationen der Gruppe $S$, durch die
1, 2 ungeändert bleiben, in [1\ 4], \dots, [1\ 8] über; ebenso [2\ 3] in
[2\ 4], \dots, [2\ 8]. Es bleibt also noch zu zeigen, dass man [1\ 3]
auch in [2\ 3] und in [4\ 5] überführen kann. Dies zeigt aber
der Anblick der drei folgenden vollständigen Siebener-Systeme:

\[
\begin{aligned}
&[1\ 2]\ [1\ 3]\ [1\ 4]\ [1\ 5]\ [1\ 6]\ [1\ 7]\ [1\ 8]\\
&[1\ 2]\ [2\ 3]\ [2\ 4]\ [2\ 5]\ [2\ 6]\ [2\ 7]\ [2\ 8]\\
&[1\ 2]\ [4\ 5]\ [3\ 4]\ [3\ 5]\ [1\ 6]\ [1\ 7]\ [1\ 8].
\end{aligned}
\]

Die Gruppe $P$ ist hiernach zweifach transitiv. Sie kann
aber nicht mehr als zweifach transitiv sein. Denn lässt man
zwei Wurzeln ungeändert, so kann eine mit diesen beiden syzygetische
Wurzel nicht in eine azygetische übergeführt werden.
Durch Adjunction von zwei Wurzeln wird die Doppeltangentengleichung
reducibel. Es löst sich ein Factor $10^{\text{ten}}$ Grades ab,
dessen Wurzeln mit den beiden gegebenen einen \gesperrt{Steiner}'schen
Complex bilden (die Function $X$ im Satze 1., \S.\,101).

Unter den Divisoren der Gruppe $P$ ist besonders die
Gruppe $S$ vom Index 36 bemerkenswerth. Diese Gruppe ist noch
transitiv, weil durch sie [1\ 2] in jede beliebige Wurzel $[\varkappa\ \lambda]$
übergeführt werden kann. Sie ist aber nur noch einfach transitiv,
weil bei festgehaltenem [1\ 2] die Wurzel [1\ 3] nicht mehr
in [4\ 5] übergehen kann.

Durch Adjunction \gesperrt{einer} Wurzel einer Gleichung $36^{\text{sten}}$ Grades
wird also die Gruppe der Doppeltangentengleichung auf die
Gruppe einer allgemeinen Gleichung $8^{\text{ten}}$ Grades reducirt.

Setzt man z.\,B.

\[
v_1 = [1\ 2]\ [1\ 3]\ [1\ 4]\ [1\ 5]\ [1\ 6]\ [1\ 7]\ [1\ 8],
\]

und definirt entsprechend $v_2, v_3, \dots, v_8$, so ist jede symmetrische
Function dieser acht Grössen, z.\,B.

\[
u = v_1 + v_2 + v_3 + v_4 + v_5 + v_6 + v_7 + v_8,
\]

Wurzel einer Gleichung $36^{\text{sten}}$ Grades. Adjungirt man dem
Problem die Grösse $u$, so werden die $v_1, v_2, \dots, v_8$ die Wurzeln
einer Gleichung $8^{\text{ten}}$ Grades, die keinen Affect hat.


% ===== scan leaf 412, printed 391 =====
% !PRINTED 391
\weberpage{391}
\pdfbookmark[2]{§ 104. Realität der Doppeltangenten}{bm119}

\S.\,104.

\gesperrt{Realität der Doppeltangenten.}

Wir wollen noch die Frage erörtern, wie sich bei einer
reellen Curve vierter Ordnung ohne singulären Punkt die Doppeltangenten
in Bezug auf ihre Realität verhalten können. Wir
nehmen also jetzt die Co\"efficienten der Gleichung der Curve
vierter Ordnung reell an, d.\,h. wir setzen einen reellen Rationalitätsbereich
voraus.

Wenn dann eine Doppeltangente $\xi$ imaginär ist, so muss
auch die conjugirte Gerade $\xi'$ Doppeltangente sein.

Wenn wir in einer rationalen Gleichung zwischen den Wurzeln
der Doppeltangentengleichung jede imaginäre Wurzel durch
die conjugirte ersetzen, so entsteht wieder eine richtige Gleichung.
Es folgt daraus, dass zu syzygetischen oder azygetischen
Systemen vom % [sic] \textquotedbl{}vom\textquotedbl{} printed here where \textquotedbl{}von\textquotedbl{} would be expected
Doppeltangenten conjugirte Systeme desselben
Charakters gehören, die man erhält, wenn man überall $i$ durch
$-i$ ersetzt.

Ein \gesperrt{Steiner}'scher Complex $C$ geht daher durch Uebergang
von $i$ zu $-i$ in einen \gesperrt{Steiner}'schen Complex $C'$ über und wir
unterscheiden zwei Fälle:

Ist $C$ mit $C'$ identisch, so nennen wir $C = C'$ einen \gesperrt{reellen
Complex}.

Ist aber $C$ von $C'$ verschieden, so bilden sie ein \gesperrt{conjugirtes
Complexpaar}.

Die Paare eines reellen Complexes bestehen entweder aus
zwei reellen Doppeltangenten oder es sind conjugirte Paare $\xi, \xi'$,
oder sie enthalten zwei nicht conjugirte imaginäre Doppeltangenten
$\xi, \eta$. Im letzteren Falle muss dann im Complex auch
das aus den conjugirten Doppeltangenten $\xi', \eta'$ gebildete Paar
vorkommen. Zwei solche Paare $(\xi, \eta)$, $(\xi', \eta')$ wollen wir ein
\gesperrt{conjugirtes Doppelpaar} nennen.

Niemals kommt in einem reellen Complexe eine reelle
Wurzel $x$ mit einer imaginären $\xi$ gepaart vor, weil sonst neben
dem Paare $(x, \xi)$ auch das conjugirte Paar $(x, \xi')$ in demselben
Complexe vorkommen müsste, was unmöglich ist, da diese beiden
Paare ein gemeinschaftliches Element $x$ enthalten würden.

Ein \gesperrt{Aronhold}'sches Siebener-System $S$ geht durch Vertauschung
aller seiner Elemente mit den conjugirten in ein eben-
%CONT>>


% ===== scan leaf 413, printed 392 =====
% !PRINTED 392
\weberpage{392}

%<<CONT
solches System $S'$ über, und wir nennen ein solches System
\gesperrt{reell}, wenn $S$ mit $S'$ identisch ist. Ein reelles Siebener-System
enthält daher zu jeder in ihm vorkommenden Wurzel $\xi$ die conjugirte
$\xi'$, und muss folglich wenigstens eine, immer aber eine
ungerade Anzahl von reellen Doppeltangenten enthalten.

Wir leiten jetzt der Reihe nach die Hauptsätze ab:

\gesperrt{1. Es können nicht alle Doppeltangenten imaginär sein.}

Es sei nämlich $\xi, \xi'$ ein conjugirtes Paar und $\eta$ eine dritte
imaginäre Doppeltangente, so dass die drei $\xi, \xi', \eta$ \gesperrt{azygetisch}
sind. (Dies wäre sicher möglich, wenn alle Doppeltangenten
imaginär wären.) Ist $\eta'$ zu $\eta$ conjugirt, so haben wir die beiden
azygetischen Tripel:

\[
\xi,\ \xi',\ \eta;\quad \xi,\ \xi',\ \eta'.
\]

Da hiernach in dem durch das Paar $\xi\eta$ bestimmten \gesperrt{Steiner}'schen
Complexe $\xi'$ nicht vorkommt, so ist dieser Complex
imaginär.

Wir stellen ihn mit seinem conjugirten Complexe zusammen:

\weq{1}{
\begin{aligned}
&1)\quad \xi\eta,\ \xi_1\eta_1,\ \xi_2\eta_2,\ \xi_3\eta_3,\ \xi_4\eta_4,\ \xi_5\eta_5\\
&2)\quad \xi'\eta',\ \xi_1'\eta_1',\ \xi_2'\eta_2',\ \xi_3'\eta_3',\ \xi_4'\eta_4',\ \xi_5'\eta_5',
\end{aligned}
}

indem wir unter $\xi_i', \eta_i'$ die zu $\xi_i, \eta_i$ conjugirten Elemente verstehen,
so dass, falls $\xi_i$ reell ist, $\xi_i = \xi_i'$ zu setzen ist, und umgekehrt
auch aus $\xi_i = \xi_i'$ folgt, dass $\xi_i$ reell ist.

Es sind nun zwei Möglichkeiten zu unterscheiden. Wenn
\gesperrt{erstens} die Complexe 1), 2) ein syzygetisches Paar bilden (\S.\,97),
so haben sie vier syzygetische Elemente gemein. Darunter
können $\xi, \eta, \xi', \eta'$ nicht vorkommen, und wir beschränken daher
die Allgemeinheit nicht, wenn wir annehmen, es seien $\xi_1, \eta_1, \xi_2, \eta_2$
die gemeinschaftlichen Elemente. Diese können in ihrer Gesammtheit
aber nicht verschieden sein von $\xi_1', \eta_1', \xi_2', \eta_2'$, da der
Uebergang zu den accentuirten Buchstaben, d.\,h. zu den conjugirt
imaginären Grössen, wodurch 1) in 2) übergeht, überall
gestattet ist.

Nun ist die Annahme $\xi_1 = \eta_1'$ ausgeschlossen, weil daraus
$\xi_1' = \eta_1$ folgen würde und 1) von 2) nicht verschieden wäre.

Ist $\xi_1 = \xi_1'$, so ist $\xi_1$ \gesperrt{reell}.

Ist aber $\xi_1 = \xi_2'$, $\xi_2 = \xi_1'$,
so ist $\eta_1 = \eta_1'$, also $\eta_1$ \gesperrt{reell},
und davon ist die Annahme $\xi_1 = \eta_2'$, $\xi_1' = \eta_2$, $\xi_2 = \xi_2'$ nicht
wesentlich verschieden.


% ===== scan leaf 414, printed 393 =====
% !PRINTED 393
\weberpage{393}
Dieser Fall führt also immer auf eine reelle Doppeltangente.

Wenn aber \gesperrt{zweitens} das Complexpaar 1), 2) azygetisch ist, so
enthält jedes Paar des einen Complexes ein Element, was auch im
anderen Complexe vorkommt, und ein Element, was im anderen nicht
vorkommt. Kommt also etwa $\eta$ im Complexe 2) vor, so können wir,
ohne Beeinträchtigung der Allgemeinheit $\eta = \eta_1'$, $\eta' = \eta_1$
annehmen, und erhalten folgende Paare in 1) und 2)
\weq{2}{
\begin{array}{rl}
1) & \xi\eta,\ \xi_1\eta' \\
2) & \xi_1'\eta,\ \xi'\eta'.
\end{array}
}
Wenn ferner $\eta_2$ in 1) und 2) vorkommt, so kann $\eta_2$ nicht gleich
$\xi_2'$ sein, weil sonst 1) das Paar $\xi_2\xi_2'$ enthielte und folglich
reell wäre.

Wenn $\eta_2 = \eta_2'$ ist, so ist $\eta_2$ \gesperrt{reell}. Ist aber $\eta_2$ gleich
einem Element der drei letzten Paare von 2), so können wir es ohne
Beschränkung $= \eta_3'$ annehmen, also $\eta_2 = \eta_3'$, $\eta_3 = \eta_2'$.

Dann sind aber nach dem Satze \S.\,96, 3. sowohl $\eta, \eta_2, \eta_2'$ als
$\eta, \eta', \eta_2$ azygetisch. Die beiden Paare $\eta\eta_2$, $\eta'\eta_2'$
bestimmen zwei Complexe, deren zweiter weder $\eta$ noch $\eta_2$ enthält,
und die daher ein \gesperrt{syzygetisches Paar} bilden. Damit sind wir auf
die erste Annahme zurückgeführt und unser Satz 1. ist bewiesen.

Es giebt also immer mindestens \gesperrt{zwei reelle} Doppeltangenten.

\gesperrt{2. Es giebt immer mindestens ein System von vier reellen
syzygetischen Doppeltangenten.}

Beim Beweise dieses Satzes gehen wir aus von einem nach 1. immer
existirenden reellen Paare $x\,y$, und betrachten den reellen Complex,
der durch dies Paar bestimmt ist. Da die conjugirten Doppelpaare je
zwei Paare sind, so muss unter den sechs Paaren dieses Complexes
entweder ein zweites reelles Paar vorkommen, und dann trifft der Satz
2. zu, oder er muss ein conjugirtes Paar $\xi\xi'$ enthalten.

Im letzteren Falle betrachten wir das durch diese beiden Paare
bestimmte syzygetische Complextripel, in dem alle 28 Doppeltangenten
vorkommen (\S.\,97, 4.):
\[
\begin{array}{rl}
1) & x\,y,\ \xi\xi' \\
2) & x\xi,\ y\xi',\ \xi_1\eta_1 \\
3) & x\xi',\ y\xi,\ \xi_1'\eta_1',
\end{array}
\]
%CONT>>


% ===== scan leaf 415, printed 394 =====
% !PRINTED 394
\weberpage{394}
%<<CONT
von denen die beiden letzten conjugirt imaginär sind, und folglich, da
sie ausser $x, y, \xi, \xi'$ kein gemeinschaftliches Element haben, nur
noch imaginäre Paare enthalten, von denen wir eins, $\xi_1\eta_1$, nebst
dem dazu conjugirten $\xi_1'\eta_1'$ mit aufgeführt haben.

Es genügt demnach, wenn wir die Existenz von einer weiteren reellen
Doppeltangente beweisen können. Denn diese muss dann in 1) vorkommen
und muss in diesem Complexe zu einem reellen Paare gehören.

Jetzt bilden wir noch die beiden conjugirten Complexe:
\[
\begin{array}{rl}
4) & x\xi_1,\ \xi\eta_1 \\
5) & x\xi_1',\ \xi'\eta_1',
\end{array}
\]
in denen $y$ gewiss nicht vorkommt, da sonst $x, \xi, y$ azygetisch wären
(\S.\,96, 3.). Aus demselben Grunde kommt $\xi$ nicht in 5) und $\xi'$
nicht in 4) vor.

Wenn nun \gesperrt{zunächst} die Complexe 4) und 5) syzygetisch sind, so muss
$\xi_1$ in 5) und $\xi_1'$ in 4) vorkommen, und wenn 4) das Paar $p\xi_1'$
enthält, so muss in 5) das Paar $p\xi_1$ vorkommen. Es ist also $p$ mit
seinem conjugirten Elemente identisch, d. h. \gesperrt{reell}.

Wenn \gesperrt{zweitens} die Complexe 4) und 5) azygetisch sind, so kommen
$\xi_1, \xi$ nicht in 5) vor, und es muss $\eta_1$ in 5) enthalten sein.
Wenn $\eta_1$ in 5) mit einem imaginären Elemente $\xi_2'$ gepaart ist, und
wenn in 4) noch ein weiteres imaginäres Paar $\xi_3\eta_3$ vorkommt, so
setzen sich diese Complexe in folgender Weise fort:
\[
\begin{array}{rl}
4) & x\xi_1,\ \xi\eta_1,\ \xi_2\eta_1',\ \xi_3\eta_3,\ \xi_3'\eta_4 \\
5) & x\xi_1',\ \xi'\eta_1',\ \xi_2'\eta_1,\ \xi_3'\eta_3',\ \xi_3\eta_4',
\end{array}
\]
worin auch $\eta_4$ imaginär ist, und durch die accentuirten Buchstaben
immer die conjugirten Elemente zu den unaccentuirten verstanden sind.
Die beiden Complexe enthalten dann nur noch je ein Paar $\xi_5\eta_5$,
$\xi_5'\eta_5'$, die ein gemeinschaftliches Element enthalten müssen. Es
kann aber nicht $\xi_5 = \eta_5'$ sein, weil sonst auch $\xi_5' = \eta_5$, und
mithin beide Paare identisch wären. Also muss $\xi_5 = \xi_5'$ (oder
$\eta_5 = \eta_5'$) sein; d. h. eine dieser beiden Doppeltangenten
ist \gesperrt{reell}, und damit ist unser Satz 2. bewiesen.

3. \gesperrt{Wenn ausser den vier reellen Doppeltangenten $x, y, p, q$,
deren Existenz der Satz 2. behauptet, noch eine weitere vorhanden ist,
so existiren acht reelle Doppeltangenten, die in einem
Steiner'schen Complexe vier Paare bilden.}


% ===== scan leaf 416, printed 395 =====
% !PRINTED 395
\weberpage{395}
Wir bilden das reelle syzygetische Complextripel, in dem alle
Doppeltangenten vorkommen müssen:
\[
\begin{array}{rl}
1) & x\,y,\ p\,q \\
2) & x\,p,\ y\,q \\
3) & x\,q,\ y\,p.
\end{array}
\]
Ist eine fünfte reelle Doppeltangente $r$ vorhanden, so können wir
annehmen, sie komme im Complexe 1) vor. Dann muss aber dieser Complex
ein drittes reelles Paar $r\,s$ enthalten, und entweder ein viertes
reelles Paar, in welchem Falle der Satz 3. schon zutrifft, oder ein
conjugirtes Paar $\varrho\varrho'$.

Enthält der Complex 2) nicht lauter reelle Doppeltangenten, ein Fall,
in dem gleichfalls der Satz 3. zutreffen würde, so muss in 2) ein
imaginäres Paar $\xi\eta$ vorkommen, und danach betrachten wir also die
folgenden Complexe:
\[
\begin{array}{rl}
1) & x\,y,\ p\,q,\ r\,s,\ \varrho\varrho' \\
2) & x\,p,\ y\,q,\ \xi\eta \\
4) & x\,r,\ y\,s.
\end{array}
\]
Die beiden Complexe 2), 4) sind aber \gesperrt{azygetisch}, da z. B. $r$ in 2)
nicht vorkommt [weil es in 1) vorkommt]. Folglich kann 4) auch nur
eine der beiden $\xi, \eta$ enthalten, und folglich können, da der Complex
4) reell ist, $\xi, \eta$ nicht conjugirt sein. Es ergiebt sich daraus für
2) und 4) je ein conjugirtes Doppelpaar, und wir haben:
\[
\begin{array}{rl}
2) & x\,p,\ y\,q,\ \xi\eta,\ \xi'\eta' \\
4) & x\,r,\ y\,s,\ \xi\zeta,\ \xi'\zeta'.
\end{array}
\]
worin $\zeta, \zeta'$ wieder ein Paar conjugirter Doppeltangenten bedeutet.

Nun bilden wir den sowohl mit 2) als mit 4) syzygetischen reellen
Complex
\[
5)\quad \xi\xi',\ \eta\eta',\ \zeta\zeta',
\]
der keines der Elemente $x, y, p, q, r, s$ enthalten kann, und ferner die
beiden conjugirt imaginären Complexe
\[
\begin{array}{rl}
6) & x\xi,\ p\eta,\ r\zeta \\
7) & x\xi',\ p\eta',\ r\zeta',
\end{array}
\]
und 5), 6), 7) bilden ein \gesperrt{azygetisches Tripel}. $y, q, s$ kommen in 6)
und 7) nicht vor, denn sonst müssten sie mit je zweien der $x, p, r$
azygetisch sein, was nach 2) und 4) unmöglich ist.


% ===== scan leaf 417, printed 396 =====
% !PRINTED 396
\weberpage{396}
Es ist also jeder der Complexe 6), 7) azygetisch mit dem Complexe 1),
und es muss also eine und nur eine der beiden Doppeltangenten
$\varrho, \varrho'$ in 6) vorkommen; ist dies $\varrho$, so ist $\varrho'$ in
7) enthalten. Ist $\varrho\sigma$ ein Paar von 6), so ist $\sigma$ von seinem
conjugirten $\sigma'$ verschieden, weil sonst $\sigma\varrho'$ in 7) und folglich
$\varrho\varrho'$ in 5) vorkommen müsste, was nicht der Fall ist. Also haben
wir, wenn $\tau, \tau'$ zwei weitere conjugirt imaginäre Doppeltangenten
sind, die Complexe:
\[
\begin{array}{rl}
6) & x\xi,\ p\eta,\ r\zeta,\ \varrho\sigma,\ \sigma'\tau \\
7) & x\xi',\ p\eta',\ r\zeta',\ \varrho'\sigma',\ \sigma\tau'.
\end{array}
\]
Ist $t\lambda$ das letzte Paar des Complexes 6), so kommt das Paar
$t'\lambda'$ in 7) vor, und diese beiden Paare müssen ein gemeinsames
Element enthalten. Da aber $\lambda$ nicht gleich $t'$ sein kann, weil sonst
beide Paare identisch wären, so muss $t = t'$ (oder $\lambda = \lambda'$, was
nicht wesentlich verschieden ist) sein; es ist also $t$ reell, und es
wird:
\[
\begin{array}{rl}
6) & x\xi,\ p\eta,\ r\zeta,\ \varrho\sigma,\ \sigma'\tau,\ t\lambda \\
7) & x\xi',\ p\eta',\ r\zeta',\ \varrho'\sigma',\ \sigma\tau',\ t\lambda'.
\end{array}
\]
Jetzt kehren wir zu dem syzygetischen Complextripel 1), 2), 3) zurück.
Da wir aus 6) schliessen, dass $x\,p\,t$ und $x\,r\,t$ azygetisch sind, so
kann $t$ weder in 2) noch in 3) vorkommen, und muss also in 1)
enthalten sein. Da aber 1) ein reeller Complex ist, so muss darin $t$
mit einer reellen Doppeltangente $u$ gepaart erscheinen, und der
Complex 1) wird
\[
x\,y,\ p\,q,\ r\,s,\ t\,u,
\]
wodurch unser Satz 3. bewiesen ist.

\gesperrt{4. Wenn in einem Complexe fünf reelle Paare vorkommen, so
sind alle Doppeltangenten reell.}

Gehen wir aus von einem Complexe mit fünf reellen Paaren
\[
1)\quad x_1y_1,\ x_2y_2,\ x_3y_3,\ x_4y_4,\ x_5y_5,
\]
und nehmen zunächst an, dass ausser dem letzten Paare dieses
Complexes noch eine imaginäre Doppeltangente $\xi$ existire, dann ist
auch eine conjugirte Doppeltangente $\xi'$ vorhanden, und der durch das
Paar $\xi\xi'$ bestimmte reelle Complex 2) ist mit 1) syzygetisch [weil
$\xi$ und $\xi'$ nicht in 1) vorkommen]. Die vier gemeinschaftlichen
Elemente von 1) und 2) müssen aber reell
%CONT>>


% ===== scan leaf 418, printed 397 =====
% !PRINTED 397
\weberpage{397}
%<<CONT
sein, weil 2) als reeller Complex keine reelle Doppeltangente
mit einer imaginären gepaart enthalten kann.

Es sei also
\[
2)\quad \xi\xi',\ x_1x_2,\ y_1y_2,
\]
und daraus leiten wir die zwei Complexe her:
\[
\begin{aligned}
3)&\quad x_1\xi,\ x_2\xi'\\
4)&\quad x_1\xi',\ x_2\xi,
\end{aligned}
\]
die mit einander syzygetisch sind, und mithin ausser den vier
angegebenen kein Element gemein haben. Sie sind aber zugleich
conjugirt, und daher kann in 3) ausser $x_1, x_2$ keine reelle Doppeltangente
 vorkommen. Nun sind aber 1) und 3) azygetisch, weil
$\xi$ in 1) nicht vorkommt; mithin muss aus jedem Paare von 1)
ein Element in 3) vorkommen. Also enthält 3) ausser $x_1, x_2$ noch
reelle Elemente, wodurch sich der Widerspruch ergiebt.

Wir schliessen daraus, dass die beiden Complexe, die durch
die Paare $x_1x_2$, $y_1y_2$ bestimmt sind und mit 1) ein syzygetisches
Complextripel ausmachen, \gesperrt{nur reelle Elemente} enthalten.

Lassen wir einen dieser Complexe, etwa $x_1x_2$, an Stelle
von 1) treten und wiederholen dann unseren Schluss, so ergiebt
sich, dass überhaupt alle Doppeltangenten, also auch die des
letzten Paares von 1), reell sein müssen.

5. \gesperrt{Wenn mehr als acht reelle Doppeltangenten
vorhanden sind, so giebt es sechzehn, die in
einem syzygetischen Complextripel je vier reelle
Paare bilden.}

Nehmen wir mehr als acht reelle Doppeltangenten an, so
können wir nach den Sätzen 3., 4. folgendes syzygetische Complextripel
 zusammenstellen:
\[
\begin{aligned}
1)&\quad x_1y_1,\ x_2y_2,\ x_3y_3,\ x_4y_4\\
2)&\quad x_1x_2,\ y_1y_2,\ z_1z_2\\
3)&\quad x_1y_2,\ x_2y_1,
\end{aligned}
\]
worin die $x_i, y_i, z_i$ reell sind.

Damit verbinden wir den Complex
\[
4)\quad x_1z_1,\ x_2z_2,
\]
der mit 1) azygetisch ist, und daher aus den Paaren $x_3y_3, x_4y_4$
je ein und nur ein Element enthalten kann, etwa $x_3$ und $x_4$.
Da aber 4) ein reeller Complex ist, so muss es zwei weitere
%CONT>>


% ===== scan leaf 419, printed 398 =====
% !PRINTED 398
\weberpage{398}
\pdfbookmark[2]{§ 105. Beweis der Existenz der vier Fälle}{bm120}
%<<CONT
reelle Doppeltangenten $z_3, z_4$ geben, so dass der Complex 4) sich
so fortsetzt:
\[
4)\quad x_1z_1,\ x_2z_2,\ x_3z_3,\ x_4z_4;
\]
$z_3$ und $z_4$ kommen nicht in 2) und auch nicht in 1) vor, und
müssen also in 3) enthalten sein.

Sie können auch in 3) nicht gepaart vorkommen, weil
$x_1, z_3, z_4$ azygetisch sind. Hieraus schliessen wir, wenn $t_3, t_4$ zwei
weitere reelle Doppeltangenten sind, auf folgende Zusammensetzung
 des Complexes 3):
\[
3)\quad x_1y_2,\ x_2y_1,\ z_3t_3,\ z_4t_4.
\]

Betrachtet man nun an Stelle des Complexes 4) den Complex
\[
5)\quad x_1z_3,\ y_2t_3,
\]
so kann man genau ebenso auf ein viertes reelles Paar $u_1u_2$ im
Complex 2) schliessen, und damit ist 5. bewiesen.

Fassen wir das Ergebniss zusammen, so erkennen wir, dass
in Bezug auf die Realität der Doppeltangenten einer reellen
Curve vierter Ordnung nur vier Fälle möglich sind:

1) \gesperrt{Vier reelle syzygetische Doppeltangenten.}

2) \gesperrt{Acht reelle Doppeltangenten, und zwar vier Paare
eines Steiner'schen Complexes.}

3) \gesperrt{Sechzehn reelle Doppeltangenten, die in einem
syzygetischen Complextripel je vier reelle Paare bilden.}

4) \gesperrt{Achtundzwanzig reelle Doppeltangenten.}

\begin{center}
\S.\,105.\\
\gesperrt{Beweis der Existenz der vier Fälle.}
\end{center}

Wir haben im vorigen Paragraphen zunächst nur bewiesen,
dass es keine anderen als die Fälle 1), 2), 3), 4) geben kann.
Dass diese vier Fälle aber wirklich alle möglich sind, ist jetzt
auch leicht zu zeigen, auf Grund des Satzes \S.\,100, nach dem
man aus sieben beliebig gegebenen geraden Linien auf \gesperrt{rationalem
 Wege} eine Curve vierter Ordnung ableiten kann, für
die die gegebenen Linien ein \gesperrt{Aronhold}'sches System bilden,
aus dem sich alle Doppeltangenten \gesperrt{rational} ableiten lassen.

Unter den sieben gegebenen geraden Linien können auch
imaginäre vorkommen, und als Bedingung der Realität der Curve
%CONT>>


% ===== scan leaf 420, printed 399 =====
% !PRINTED 399
\weberpage{399}
%<<CONT
(d. h. der Coëfficienten in ihrer Gleichung) ergiebt sich die, dass
das System, was man erhält, wenn man jede imaginäre Gerade
durch ihre conjugirte ersetzt, wieder ein Aronhold'sches System
derselben Curve ist. Denn dann ändern sich die Coëfficienten
nicht, wenn $i$ durch $-i$ ersetzt wird.

Die Curve wird also gewiss reell, wenn in dem gegebenen
Siebener-Systeme zu jeder imaginären Geraden die conjugirte
vorkommt; dann bilden sie ein \gesperrt{reelles Siebener-System}.

Ein solches reelles System kann nun enthalten:

1) eine reelle, drei Paar conjugirt imaginäre Geraden;

2) drei reelle und zwei Paar conjugirt imaginäre Geraden;

3) fünf reelle und ein Paar conjugirt imaginäre Geraden;

4) sieben reelle Geraden.

Wir werden sehen, dass diese vier Annahmen in derselben
Ordnung zu den im vorigen Paragraphen aufgezählten vier
Fällen führen.

Dieser Nachweis wird sehr einfach, wenn man sich der
Bezeichnungsweise der Doppeltangenten, die wir im \S.\,99 dargelegt
 haben, bedient, nach der wir die Steiner'schen Complexe
 unmittelbar aus der Bezeichnung bilden können.

1) Im ersten Falle bezeichnen wir die gegebenen sieben
Geraden mit
\[
0,\ 1,\ 1',\ 2,\ 2',\ 3,\ 3',
\]
setzen $0$ als reell voraus, $1$ mit $1'$, $2$ mit $2'$, $3$ mit $3'$ conjugirt
imaginär. Die übrigen 21 Doppeltangenten werden dann durch
die Zeichen $[01], [01'], [11'], \dots$ bezeichnet.

Nach der Vorschrift des \S.\,99 erhalten wir die beiden
folgenden Steiner'schen Complexe:
\[
\begin{aligned}
(1)&\quad 0[01],\ 1'[11'],\ 2[12],\ 2'[12'],\ 3[13],\ 3'[13']\\
(1')&\quad 0[01'],\ 1[11'],\ 2[1'2],\ 2'[1'2'],\ 3[1'3],\ 3'[1'3'].
\end{aligned}
\]

Der zu (1) conjugirte Complex muss die Elemente $0,1,2,2',3,3'$
enthalten, und ist also nach \S.\,98, 7. mit (1') identisch, der
diese sechs Elemente gleichfalls enthält.

Daraus folgt, dass $[11']$ reell ist, und dass
\[
\begin{aligned}
&{}[01]\quad [01']\\
&{}[12]\quad [1'2']\\
&{}[12']\quad [1'2]
\end{aligned}
\]
%CONT>>


% ===== scan leaf 421, printed 400 =====
% !PRINTED 400
\weberpage{400}
%<<CONT
conjugirte Paare sind. Da man in dieser Betrachtung 1 mit 2
und 3 vertauschen kann, so folgt, dass $0, [11'], [22'], [33']$
reell und alle übrigen Doppeltangenten imaginär sind.

Nach \S.\,99 können wir leicht einen Steiner'schen Complex
bilden, der die vier reellen Doppeltangenten enthält:
\[
0[11'],\ [22'][33'],\ 1[01'],\ 1'[01],\ [23][2'3'],\ [23'][2'3],
\]
woraus noch zu sehen ist, dass dieser Complex ausser den reellen
Paaren zwei conjugirte Paare und ein conjugirtes Doppelpaar
enthält.

Dies ist also der Fall 1) des vorigen Paragraphen.

2) Es seien die sieben gegebenen Geraden:
\[
0,\ 1,\ 2,\ 3,\ 3',\ 4,\ 4',
\]
und darin $0, 1, 2$ reell, $3$ zu $3'$ und $4$ zu $4'$ conjugirt. Wir
bilden drei Complexe:
\[
\begin{aligned}
(0)&\quad 1[01],\ 2[02],\ 3[03],\ 3'[03'],\ 4[04],\ 4'[04']\\
(3)&\quad 0[03],\ 1[13],\ 2[23],\ 3'[33'],\ 4[34],\ 4'[34']\\
(3')&\quad 0[03'],\ 1[13'],\ 2[23'],\ 3[33'],\ 4[3'4],\ 4'[3'4'].
\end{aligned}
\]

Der zu (0) conjugirte Complex enthält $1, 2, 3, 3', 4, 4'$ und
ist folglich mit (0) identisch, d. h. der Complex (0) ist reell.
Daraus folgt, dass $[01], [02]$ reell, $[03]$ mit $[03']$ und $[04]$
mit $[04']$ conjugirt ist. Ferner ergiebt sich auf die gleiche
Weise, dass (3), (3') conjugirte Complexe sind, und dass also
$[33']$ reell, $[34]$ mit $[3'4']$, $[34']$ mit $[3'4]$ conjugirt imaginär
sind. Da man 0, 1, 2 beliebig vertauschen darf, und ebenso 3
mit 4, so erhalten wir folgende Zusammenstellung:

\begin{center}
\gesperrt{Reelle Doppeltangenten:}
\end{center}
\[
0,\ 1,\ 2,\ [12],\ [20],\ [01],\ [33'],\ [44'].
\]

\begin{center}
\gesperrt{Conjugirte Paare:}
\end{center}
\[
\begin{aligned}
3\ \ 4 &\quad [03][04][13][14][23][24][34][34']\\
3'\ 4' &\quad [03'][04'][13'][14'][23'][24'][3'4'][3'4].
\end{aligned}
\]

Der Steiner'sche Complex, der vier reelle Paare enthält,
ist hier
\[
0[12],\ 1[20],\ 2[01],\ [33'][44'],\ [34][3'4'],\ [34'][3'4],
\]
und er enthält noch zwei conjugirte Paare. Dies ist also der
Fall 2) des \S.\,104.


% ===== scan leaf 422, printed 401 =====
% !PRINTED 401
\weberpage{401}

3) Es seien die gegebenen Linien:
\[
0,\ 1,\ 2,\ 3,\ 4,\ 5,\ 5',
\]
und davon 0, 1, 2, 3, 4 reell, 5 und 5' conjugirt imagin\"ar. Die
Betrachtung der drei Complexe:
\[
\begin{aligned}
(0)\quad & 1\,[1\,0],\ 2\,[2\,0],\ 3\,[3\,0],\ 4\,[4\,0],\ 5\,[5\,0],\ 5'\,[5'\,0] \\
(5)\quad & 0\,[0\,5],\ 1\,[1\,5],\ 2\,[2\,5],\ 3\,[3\,5],\ 4\,[4\,5],\ 5'\,[5'\,5] \\
(5')\quad & 0\,[0\,5'],\ 1\,[1\,5'],\ 2\,[2\,5'],\ 3\,[3\,5'],\ 4\,[4\,5'],\ 5\ [5\,5'],
\end{aligned}
\]
von denen der erste reell, die beiden anderen conjugirt imagin\"ar
sind, und der durch Vertauschung von 0, 1, 2, 3, 4 aus (0) abgeleiteten
 giebt, genau wie oben, folgendes Resultat:

\begin{center}
\gesperrt{Reelle Doppeltangenten:}
\end{center}
\[
\begin{aligned}
&0,\quad 1,\quad 2,\quad 3,\quad 4,\quad [0\,1],\ [0\,2],\ [0\,3] \\
&[0\,4]\ [1\,2],\ [1\,3]\ [1\,4],\ [2\,3],\ [2\,4],\ [3\,4],\ [5\,5'].
\end{aligned}
\]

\begin{center}
\gesperrt{Conjugirte Paare:}
\end{center}
\[
\begin{aligned}
&5,\ [0\,5],\ [1\,5],\ [2\,5],\ [3\,5],\ [4\,5] \\
&5',\ [0\,5'],\ [1\,5'],\ [2\,5'];\ [3\,5'],\ [4\,5'], % [sic] semicolon printed where a comma is expected
\end{aligned}
\]
und wir finden ein syzygetisches Complextripel:
\[
\begin{aligned}
&0,\quad 1,\quad [0\,2][1\,2],\ [0\,3][1\,3],\ [0\,4][1\,4\,],\ [0\,5][1\,5\,],\ [0\,5'][1\,5'] \\
&2,\quad 3,\quad [0\,2][0\,3],\ [1\,2][1\,3],\ [4\,2][4\,3\,],\ [2\,5][3\,5\,],\ [2\,5'][3\,5'] \\
&[0\,1],\ [2\,3],\ [0\,2][1\,3],\ [0\,3][1\,2],\quad 4\ \ [5\,5'],\quad 5\ \ [4\,5'],\quad 5'\ \ [4\,5],
\end{aligned}
\]
von denen jeder Complex noch vier reelle Paare und ein imagin
\"ares Doppelpaar enth\"alt. Solcher Tripel lassen sich aber noch
mehrere bilden, da man 0, 1, 2, 3, 4 beliebig vertauschen darf.
Dies ist der dritte Fall von \S.\,104.

4) Dass endlich, wenn wir alle Elemente des gegebenen
Siebener-Systemes reell voraussetzen, auch alle \"ubrigen Doppeltangenten
 reell ausfallen, ist eine unmittelbare Folge der rationalen
 Darstellung (\S.\,100).

Hiermit ist gezeigt, dass die vier F\"alle des vorigen Paragraphen
 wirklich alle vorkommen k\"onnen. Ob die hier besprochene
 Erzeugungsweise die einzig m\"ogliche ist, mit anderen
Worten, ob bei jeder reellen Curve vierter Ordnung ein reelles
\gesperrt{Aronhold'sches} System existirt, diese Frage m\"ussen wir unentschieden
 lassen.


% ===== scan leaf 423, printed 402 =====
% !PRINTED 402
\weberpage{402}

Auch ist hier noch darauf hinzuweisen, dass aus der Realit\"at
einer Doppeltangente noch keineswegs die Realit\"at der Ber\"uhrungspunkte
 folgt, weil diese Ber\"uhrungspunkte wieder von einer
quadratischen Gleichung abh\"angen. Die reellen Doppeltangenten
zerfallen also wieder in zwei Arten, solche mit reellen und
solche mit imagin\"aren Ber\"uhrungspunkten. Was hier f\"ur m\"ogliche
 F\"alle zu unterscheiden sind, diese Frage er\"ortern wir nicht
weiter \textsuperscript{1}).

\begin{center}
\rule{0.3\linewidth}{0.4pt}
\end{center}

\textsuperscript{1}) Vergl. \"uber die ganze Frage von der Realit\"at der Doppeltangenten
vom geometrischen Gesichtspunkte: \gesperrt{Zeuthen}, \glqq Sur les diff\'erentes formes
des courbes planes du quatri\`eme ordre\grqq{}. Mathem. Ann., Bd. VII (1873).


% ===== scan leaf 424, printed 403 =====
\weberpage{403}% page number not printed on this page
\pdfbookmark[1]{Dreizehnter Abschnitt. Allgemeine Theorie der Gleichung fünften Grades.}{bm121}
\pdfbookmark[2]{§ 106. Fragestellung}{bm122}
% !PRINTED none

\begin{center}
\gesperrt{Dreizehnter Abschnitt.}

\medskip
\gesperrt{Allgemeine Theorie der Gleichung f\"unften Grades.}

\medskip
\rule{0.3\linewidth}{0.4pt}

\medskip
\S.\,106.

\gesperrt{Fragestellung.}
\end{center}

Wir haben im \S.\,44 gesehen, dass es \"uber die Frage nach
der \gesperrt{Galois'schen} Gruppe einer algebraischen Gleichung noch
ein weiter gehendes Problem giebt, n\"amlich die Frage nach der
linearen Substitutionsgruppe von m\"oglichst geringer Dimensionenzahl
, auf deren Formenproblem die gegebene Gleichung zur\"uck-
f\"uhrbar ist. Bei den metacyklischen Gleichungen, insbesondere
also auch bei den allgemeinen Gleichungen 3\textsuperscript{ten} und 4\textsuperscript{ten} Grades,
ist diese Dimensionenzahl gleich 1, wodurch eben ausgedr\"uckt
ist, dass diese Gleichungen durch Radicale l\"osbar sind. Die
zun\"achst zu untersuchenden Gleichungen sind dann die vom
5\textsuperscript{ten} Grade, und es wird sich zeigen, dass die L\"osung der allgemeinen
 Gleichung 5\textsuperscript{ten} Grades auf eine bin\"are lineare Substitutionsgruppe
, n\"amlich auf das Ikosa\"ederproblem f\"uhrt.

Damit im Zusammenhange steht aber noch eine andere Frage.
Die allgemeine Gleichung 5\textsuperscript{ten} Grades enth\"alt f\"unf Co\"efficienten
, die als unabh\"angige Variable betrachtet werden k\"onnen,
und folglich ist die Wurzel einer solchen Gleichung eine algebraische
 Function von f\"unf Variablen. Nun kann man aber
schon durch die einfachen linearen Substitutionen die Zahl dieser
Variablen vermindern. Durch \gesperrt{Tschirnhausen}-Transformation,
z.\,B. auf die \gesperrt{Jerrard'sche} Form, kann man die Gleichung sogar
nur von einem variablen Co\"efficienten (einem \gesperrt{Parameter}) ab-
h\"angig machen, und man kann also, durch Vermittelung von
Gleichungen, die den 5\textsuperscript{ten} Grad nicht erreichen, die Wurzel einer
allgemeinen Gleichung 5\textsuperscript{ten} Grades von einer algebraischen
Function von einer Variablen abh\"angig machen.


% ===== scan leaf 425, printed 404 =====
% !PRINTED 404
\weberpage{404}
\pdfbookmark[2]{§ 107. Satz von Lüroth}{bm123}

Die Frage, auf die wir hier gef\"uhrt werden, ist also die, wie
man die L\"osung einer algebraischen Gleichung, deren Co\"efficienten
 von einer gewissen Anzahl von Variablen abh\"angen, auf
eine Gleichung mit einer m\"oglichst geringen Anzahl von Parametern
, und insbesondere, unter welchen Umst\"anden und mit
welchem H\"ulfsmittel man sie auf Gleichungen mit nur einem
Parameter zur\"uckf\"uhren kann.

Wir stellen uns demnach jetzt die Frage, unter welchen
Voraussetzungen bei einer \gesperrt{allgemeinen Gleichung $n^{\text{ten}}$ Grades
Resolventen mit nur einem Parameter existiren.}

Es sei $x_0, x_1, \ldots, x_{n-1}$ ein System von $n$ unabh\"angigen Variablen,
\weq{1}{u = \Phi(x_0, x_1, \ldots, x_{n-1})}
eine rationale Function dieser Variablen, die durch die Permutationen
 der \gesperrt{alternirenden Gruppe} der $n$ Buchstaben $x$ in die
von einander verschiedenen Functionen
\weq{2}{u,\ u_1, u_2, \ldots u_{\nu-1}}
\"ubergeht, worin $\nu$ gleich oder kleiner als der Grad der alternirenden
 Gruppe sein kann, und $z$ eine Function der Variablen $x$
ist, die durch die alternirende Gruppe unge\"andert bleibt. Es
fragt sich, wann eine rationale Gleichung $\nu^{\text{ten}}$ Grades in Bezug
auf $u$ besteht,
\weq{3}{F(u, z) = 0,}
die durch die $\nu$ Functionen (2) identisch befriedigt wird? Die
Co\"efficienten in (3) m\"ussen von den $x$ unabh\"angig und also reine
Zahlen sein. Eine Beschr\"ankung des Rationalit\"atsbereiches im
Gebiete der Zahlen lassen wir einstweilen nicht eintreten.

Die Gleichung (3) ist, wenn $\nu > 1$ ist, eine Resolvente der
Gleichung $n^{\text{ten}}$ Grades, deren Wurzeln die $x$ sind, die nur von dem
\gesperrt{einen Parameter} $z$ abh\"angt. Den Fall $\nu = 1$ schliessen wir aus.

\begin{center}
\S.\,107.

\gesperrt{Satz von L\"uroth.}
\end{center}

Es ist zun\"achst folgender H\"ulfssatz zu beweisen \textsuperscript{1}):

1. \gesperrt{Ist $u, u_1, u_2, \ldots, u_{\nu-1}$ ein System rationaler Functionen
 einer Variablen $t$, so giebt es eine rationale
Function der $u, u_1, \ldots, u_{\nu-1}$:}
\[
\vartheta = \chi(u, u_1, \ldots, u_{\nu-1})
\]

\begin{center}
\rule{0.3\linewidth}{0.4pt}
\end{center}

\textsuperscript{1}) \gesperrt{L\"uroth}, Mathematische Annalen, Bd. IX.
%CONT>>


% ===== scan leaf 426, printed 405 =====
% !PRINTED 405
\weberpage{405}
%<<CONT
\gesperrt{von der Art, dass $u, u_1, \ldots, u_{\nu-1}$ rational durch $\vartheta$
allein ausgedrückt werden können.}

Um ihn zu beweisen, setzen wir

\weq{1}{u = \dfrac{\varphi(t)}{\psi(t)}, \quad u_1 = \dfrac{\varphi_1(t)}{\psi_1(t)}, \quad \cdots, \quad u_{\nu-1} = \dfrac{\varphi_{\nu-1}(t)}{\psi_{\nu-1}(t)},}

und verstehen unter $\varphi(t)$, $\psi(t)$ ganze Functionen von $t$
ohne gemeinschaftlichen Theiler; ebenso unter $\varphi_1(t)$,
$\psi_1(t)$ u.\,s.\,f. Ausserdem dürfen wir noch voraussetzen, dass
von den Functionen $u, u_1, \ldots, u_{\nu-1}$ keine eine Constante sei.
Wir führen neben der Variablen $t$ eine zweite Variable $\tau$ ein und
setzen:

\weq{2}{v = \dfrac{\varphi(\tau)}{\psi(\tau)}, \quad v_1 = \dfrac{\varphi_1(\tau)}{\psi_1(\tau)}, \quad \cdots, \quad v_{\nu-1} = \dfrac{\varphi_{\nu-1}(\tau)}{\psi_{\nu-1}(\tau)}.}

Nun bilden wir die ganzen Functionen der beiden Variablen $t, \tau$,
deren keine identisch verschwindet:

\weq{3}{\begin{aligned}
\varphi(t)\,\psi(\tau) - \varphi(\tau)\,\psi(t) &= \Phi(t,\tau) = -\Phi(\tau,t)\\
\varphi_1(t)\,\psi_1(\tau) - \varphi_1(\tau)\,\psi_1(t) &= \Phi_1(t,\tau) = -\Phi_1(\tau,t)\\
&\cdots\cdots\cdots\cdots\cdots\cdots\cdots\cdots\cdots\cdots\cdots\cdots\cdots\cdots\\
\varphi_{\nu-1}(t)\,\psi_{\nu-1}(\tau) - \varphi_{\nu-1}(\tau)\,\psi_{\nu-1}(t) &= \Phi_{\nu-1}(t,\tau) = -\Phi_{\nu-1}(\tau,t).
\end{aligned}}

Betrachten wir sie als Functionen von $t$, so verschwinden sie alle,
wenn $t = \tau$ wird, und sie haben also einen grössten
gemeinschaftlichen Theiler, der in Bezug auf $t$ mindestens vom
ersten Grade ist, und vom $\mu^{\text{ten}}$ Grade sein möge:

\weq{4}{R = R(t, \tau).}

Diese Function $R$ ist auch in Bezug auf $\tau$ rational, und wenn
man sie so einrichtet, dass $\tau$ nicht im Nenner und nicht in einem
überflüssigen Factor vorkommt, so kann $R(t, \tau)$ bei der
Vertauschung von $t$ und $\tau$ höchstens sein Vorzeichen ändern.
Denn die Functionen (3) sind, auch als Functionen der beiden
Variablen $t, \tau$ betrachtet, durch $R(t, \tau)$, und da sie
alternirend sind, durch $R(\tau, t)$ theilbar (Bd.\,I, \S.\,51). Es
ist also $R(t, \tau)$ durch $R(\tau, t)$ theilbar, und umgekehrt, und
folglich unterscheiden sich beide nur durch einen constanten Factor,
der $= \pm 1$ sein muss, weil die nochmalige Vertauschung von $t$ mit
$\tau$ die ursprüngliche Function wieder herstellt.

Es lässt sich noch beweisen, dass keine der Functionen (3) (bei
unbestimmtem $\tau$) als Function von $t$ betrachtet, einen Factor
%CONT>>


% ===== scan leaf 427, printed 406 =====
% !PRINTED 406
\weberpage{406}
%<<CONT

mehrfach enth\"alt, und dass in Folge dessen auch $R\,(t, \tau)$ durch
kein Quadrat theilbar ist. Angenommen n\"amlich, es h\"atten
\[
\begin{aligned}
\varPhi\,(t,\tau) &= \varphi\,(t)\,\psi\,(\tau) - \varphi\,(\tau)\,\psi\,(t) \\
\varPhi'(t,\tau) &= \varphi'(t)\,\psi\,(\tau) - \varphi\,(\tau)\,\psi'(t)
\end{aligned}
\]
als Functionen von $t$ einen gemeinsamen Theiler $P$, so m\"usste $P$
auch Theiler von
\[
\psi'(t)\,\varPhi - \psi(t)\,\varPhi' = [\varphi\,(t)\,\psi'(t) - \psi\,(t)\,\varphi'(t)]\,\psi\,(\tau)
\]
sein. Es w\"are daher $P$ Theiler von $\varphi\,(t)\,\psi'(t) - \psi\,(t)\,\varphi'(t)$, und
k\"onnte also von $\tau$ unabh\"angig angenommen werden.

Nehmen wir nun zwei Werthe $\tau_1, \tau_2$ von $\tau$ so an, dass
$\varPhi\,(\tau_2, \tau_1)$ von Null verschieden wird, so ist nach (3)
\[
\varPhi\,(t,\tau_1)\,\varphi\,(\tau_2) - \varPhi\,(t,\tau_2)\,\varphi\,(\tau_1) = \varphi\,(t)\,\varPhi\,(\tau_2,\tau_1).
\]
Darin ist die linke Seite durch $P$ theilbar, und also ist
auch $\varphi\,(t)$ und folglich $\psi\,(t)$ durch $P$ theilbar, was der Voraussetzung
 widerspricht, dass diese beiden Functionen ohne gemeinschaftlichen
 Theiler sein sollen.

Daraus folgt beil\"aufig, dass $R\,(t,\tau) = -R\,(\tau,t)$ sein muss,
da $R$ durch $t - \tau$, aber nicht durch $(t-\tau)^2$ theilbar ist.

Um die Function $R\,(t,\tau)$ zu bilden, kann man den gr\"ossten
gemeinschaftlichen Theiler der Functionen $\varphi_h(t) - v_h\,\psi_h(t)$ aufsuchen
, woraus folgt, dass $R\,(t,\tau)$, abgesehen von einem von $t$
unabh\"angigen Factor, rational durch $v, v_1, \ldots, v_{\nu-1}$ dargestellt
werden kann. Sind also $a$ und $b$ irgend zwei feste numerische
Werthe, so ist
\weq{5}{\frac{R\,(a,\tau)}{R\,(b,\tau)} = \vartheta = \chi\,(v, v_1, \ldots, v_{\nu-1})}
eine rationale Function von $v, v_1, \ldots, v_{\nu-1}$. Dabei ist, wenn $\xi$
eine neue Variable bedeutet:
\weq{6}{R\,(a,\xi) - \vartheta\,R\,(b,\xi) = X}
in Bezug auf $\xi$ vom $\mu^{\mathrm{ten}}$ Grade.

Ist $\tau$ ein willk\"urlicher Werth, so m\"ogen die Wurzeln der
Gleichung $R\,(t,\tau) = 0$
\weq{7}{t = \tau, \tau', \tau'', \ldots}
sein, so dass f\"ur jeden Index $h$
\[
\varPhi_h\,(\tau, \tau), \quad \varPhi_h\,(\tau', \tau), \quad \varPhi_h\,(\tau'', \tau), \ldots
\]
verschwinden.


% ===== scan leaf 428, printed 407 =====
% !PRINTED 407
\weberpage{407}
\pdfbookmark[2]{§ 108. Resolventen mit einem Parameter}{bm124}
Sind $\tau', \tau''$ irgend zwei dieser Wurzeln, so folgt aus (3):
\[
\begin{aligned}
\varphi_h\,(\tau')\,\psi_h\,(\tau) - \varphi_h\,(\tau)\,\psi_h\,(\tau') &= 0\\
\varphi_h\,(\tau'')\,\psi_h\,(\tau) - \varphi_h\,(\tau)\,\psi_h\,(\tau'') &= 0,
\end{aligned}
\]
und daraus, da $\psi_h\,(\tau)$ und $\varphi_h\,(\tau)$ nicht zugleich
verschwinden können,
\[
\varphi_h\,(\tau')\,\psi_h\,(\tau'') - \varphi_h\,(\tau'')\,\psi_h\,(\tau') = 0,
\]
d.\,h. es ist, wenn $\tau', \tau''$ irgend zwei der Grössen (7) sind,
\[
\Phi_h\,(\tau', \tau'') = 0.
\]

Daraus folgt, dass die $\nu$ Functionen $\Phi_h\,(t,\tau')$, von einem von
$t$ unabhängigen Factor abgesehen, den nämlichen grössten
gemeinschaftlichen Theiler haben, wie die Functionen $\Phi_h\,(t,\tau)$.
Das Gleiche gilt für die Functionen $\Phi_h\,(t,\tau'')$ u.\,s.\,f., oder
die Gleichungen
\[
R\,(t,\tau) = 0, \quad R\,(t,\tau') = 0, \quad R\,(t,\tau'') = 0, \ldots
\]
haben alle dieselben Wurzeln. Daraus folgt nach (5):
\[
\vartheta = \dfrac{R\,(a,\tau)}{R\,(b,\tau)} = \dfrac{R\,(a,\tau')}{R\,(b,\tau')} = \dfrac{R\,(a,\tau'')}{R\,(b,\tau'')} = \cdots,
\]
und mithin sind nach (6) die $\mu$ Werthe
\[
\xi = \tau, \tau', \tau'', \ldots
\]
die Wurzeln der Gleichung $X = 0$. Andererseits folgt aus
$\Phi_h\,(\tau',\tau) = 0$, $\Phi_h\,(\tau'',\tau) = 0, \ldots$
\[
v_h = \dfrac{\varphi_h\,(\tau)}{\psi_h\,(\tau)} = \dfrac{\varphi_h\,(\tau')}{\psi_h\,(\tau')} = \dfrac{\varphi_h\,(\tau'')}{\psi_h\,(\tau'')} \cdots = \dfrac{1}{\mu}\,\Sigma\,\dfrac{\varphi_h\,(\tau)}{\psi_h\,(\tau)},
\]
und es kann folglich $v_h$ als symmetrische Function der Wurzeln von (6)
rational durch die Coëfficienten dieser Gleichung, d.\,h. rational durch
$\vartheta$ dargestellt werden. Vertauscht man wieder $t$ mit $\tau$, so
erhält man nach (1) und (2) Ausdrücke von der Form

\weq{10}{\begin{aligned}
u &= f\,(\vartheta), \quad u_1 = f_1\,(\vartheta), \ldots, u_{\nu-1} = f_{\nu-1}(\vartheta)\\
\vartheta &= \chi\,(u, u_1, \ldots, u_{\nu-1}),
\end{aligned}}

worin $f, f_1, \ldots, f_{\nu-1}, \chi$ rationale Functionen bedeuten.

\begin{center}
\S.\ 108.\\
\gesperrt{Resolventen mit einem Parameter.}
\end{center}

Wir kehren jetzt zu unseren anfänglichen Voraussetzungen (\S.\,106)
zurück, und bezeichnen mit $u, u_1, \ldots, u_{\nu-1}$ ein System
%CONT>>


% ===== scan leaf 429, printed 408 =====
% !PRINTED 408
\weberpage{408}
%<<CONT
rationaler Functionen von $x_0, x_1, \ldots, x_{n-1}$, die durch die
Permutationen der alternirenden Gruppe aus einer von ihnen hervorgehen und
Wurzeln der Gleichung

\weq{1}{F\,(u, z) = 0}

sind, worin $z$ eine rationale Function der $x$ ist, die durch die
Permutationen der alternirenden Gruppe ungeändert bleibt.

Wir beweisen folgenden zweiten Hülfssatz:

\gesperrt{2. Wenn die rationale Function $\Psi\,(u, u_1, \ldots, u_{\nu-1})$
identisch verschwindet, wenn für die
\[
x_0, x_1, \ldots, x_{n-1}
\]
gewisse rationale Functionen einer Variablen $t$}

\weq{2}{x_h = \varphi_h\,(t)}

\gesperrt{gesetzt werden, und wenn durch diese Substitution $z$ nicht von
$t$ unabhängig wird, so verschwindet $\Psi$ identisch auch als Function der
unabhängigen Variablen $x_0, x_1, \ldots, x_{\nu-1}$.}
% [sic] The source prints x_{\nu-1} here, although the independent variables were introduced through x_{n-1}.

Die Galois'sche Gruppe der Gleichung (1) in dem Körper der rationalen
Functionen von $z$ wird aus gewissen Permutationen der Wurzeln
$u, u_1, \ldots, u_{\nu-1}$ bestehen. Führen wir diese Permutationen in
der Function $\Psi$ aus, und bilden das Product

\weq{3}{\Pi\,\Psi\,(u, u_1, \ldots, u_{\nu-1}) = f\,(z)}

aller so erhaltenen Functionen, so ergiebt sich eine rationale Function
$f\,(z)$ von $z$. Wenn nun einer der Factoren des Productes (3) nach der
Substitution (2) identisch verschwindet, und $z$ ist nicht von $t$
unabhängig, so muss $f\,(z)$ identisch gleich Null sein, und folglich muss
einer der Factoren des Productes (3) auch als Function der $x$ identisch
verschwinden. Wenn aber einer dieser Factoren identisch gleich Null ist,
so verschwinden auch alle anderen Factoren, weil man in jeder rationalen
Gleichung zwischen den $u$ alle Permutationen der Galois'schen Gruppe
ausführen kann. Damit ist der Satz 2. bewiesen.

Dieser Satz giebt nun mit dem im vorigen Paragraphen bewiesenen Satze 1.
zusammengenommen das folgende Resultat:

\gesperrt{3. Sind $u, u_1, \ldots, u_{\nu-1}$ die Wurzeln einer von einem
Parameter $z$ abhängigen Gleichung (1), und sind $u, u_1, \ldots,
u_{\nu-1}, z$ rationale Functionen der}
%CONT>>


% ===== scan leaf 430, printed 409 =====
% !PRINTED 409
\weberpage{409}
%<<CONT
\gesperrt{unabhängigen Variablen $x_0, x_1, \ldots, x_{n-1}$, so kann
man}
\weq{4}{u = f(\vartheta),\ u_1 = f_1(\vartheta), \ldots, u_{\nu-1} = f_{\nu-1}(\vartheta)}
\gesperrt{setzen, worin $f, f_1, \ldots, f_{\nu-1}$ rationale Functionen
von $\vartheta$ sind und}
\weq{5}{\vartheta = \chi(u, u_1, \ldots, u_{\nu-1}) = \psi(x_0, x_1, \ldots, x_{n-1})}
\gesperrt{eine rationale Function der $u$, oder auch der
$x$ ist.}

Nach dem Satze 1. nämlich können wir zunächst, wenn wir
$x_0, x_1, \ldots, x_{n-1}$ durch rationale Functionen einer Variablen $t$
ersetzen, so dass $z$ nicht von $t$ unabhängig wird, die Functionen
$u_h$ durch die Formeln (4), (5) darstellen. Die Relationen
\[
u_h = f_h(\vartheta) = f_h[\chi(u, u_1, \ldots, u_{\nu-1})]
\]
sind dann in Bezug auf $t$ identisch befriedigt, und müssen also
nach dem Satze 2. auch in den Variablen $x$ identisch sein,
w. z. b. w.

4.~\gesperrt{Wenn von zwei Variablen $\vartheta, \vartheta_1$ jede eine ratio\-
nale Function der anderen ist, so sind sie lineare
Functionen von einander.}

Zum Beweise nehmen wir an, es sei
\weq{6}{\vartheta = \frac{\varphi(\vartheta_1)}{\psi(\vartheta_1)},\ \vartheta_1 = \frac{\varphi_1(\vartheta)}{\psi_1(\vartheta)},}
und verstehen unter $\varphi, \psi$ zwei ganze Functionen ohne gemeinsamen
 Theiler, ebenso unter $\varphi_1, \psi_1$. Ordnen wir die zweite der
Gleichungen (6) in der Form $\varphi_1(\vartheta) - \vartheta_1\psi_1(\vartheta) = 0$ nach Potenzen
 von $\vartheta$, so mag sie die Gestalt annehmen:
\[
a_0\vartheta^m + a_1\vartheta^{m-1} + \cdots + a_{m-1}\vartheta + a_m = 0,
\]
worin die Coëfficienten $a_0, a_1, \ldots, a_m$ \gesperrt{ganze lineare Functionen}
 von $\vartheta_1$ sind. Substituiren wir darin für $\vartheta$ den Werth
aus der ersten Gleichung (6), so folgt die in Bezug auf $\vartheta_1$ identische
 Gleichung
\[
a_0\varphi^m + a_1\varphi^{m-1}\psi + \cdots + a_{m-1}\varphi\psi^{m-1} + a_m\psi^m = 0.
\]
Hiernach muss $a_0\varphi^m$ durch $\psi$ theilbar sein, und weil $\varphi$
und $\psi$ relativ prim vorausgesetzt sind, so muss $a_0$ durch $\psi$
theilbar, also $\psi$ constant oder linear sein. Ebenso schliessen
wir, dass $\varphi$ constant oder linear sein muss, und da nicht beide
%CONT>>


% ===== scan leaf 431, printed 410 =====
% !PRINTED 410
\weberpage{410}
\pdfbookmark[2]{§ 109. Gruppe der Resolventen mit einem Parameter}{bm125}
%<<CONT
Functionen constant sein können, so ist nach (6) der Satz 4.
bewiesen.

Diesen Satz wenden wir auf die in 3. vorkommende Function
 $\vartheta$ an. Wenn wir mit den Variablen $x$ irgend eine Permutation
 der alternirenden Gruppe vornehmen, so erfahren die
Functionen $u, u_1, \ldots, u_{\nu-1}$ gleichfalls eine Permutation. Nach
(5) geht $\vartheta$ durch diese Permutation in eine andere Function $\vartheta_1$
über, die nach (4) rational durch $\vartheta$ darstellbar ist. Ebenso ist
aber auch $\vartheta$ rational durch $\vartheta_1$ darstellbar, da man in dem
Satze 3. $\vartheta$ durch $\vartheta_1$ ersetzen kann, und auch $\vartheta$ aus $\vartheta_1$ durch
eine Permutation der alternirenden Gruppe entsteht. Es sind
also nach 4. $\vartheta$ und $\vartheta_1$ \gesperrt{lineare Functionen} von einander.
Daraus ergiebt sich das folgende Resultat:

5.~\gesperrt{Wenn bei einer allgemeinen Gleichung $n^{\text{ten}}$
Grades eine Resolvente mit einem Parameter
besteht, so giebt es eine rationale Function $\vartheta$
von $n$ Variablen $x$, die durch die Permutationen
der alternirenden Gruppe in eine lineare Function
 von sich selbst}
\weq{7}{\chi(\vartheta) = \frac{a\vartheta + b}{c\vartheta + d}}
\gesperrt{übergeht, worin $a, b, c, d$ Constante, d. h. Zahlen
sind.}

\begin{center}
\S.\,109.

\gesperrt{Gruppe der Resolventen mit einem Parameter.}
\end{center}

Die Function $\vartheta$, die im Satze 5. des vorigen Paragraphen
vorkommt, ist nach dem Satze 3. selbst die Wurzel einer Resolvente
 mit \gesperrt{einem Parameter} $\Phi(\vartheta, z) = 0$. Diese Gleichung
ist, da jede Wurzel rational durch jede andere ausdrückbar ist,
eine Normalgleichung, und ihre \gesperrt{Galois}'sche Gruppe ist isomorph
 mit der Gruppe der linearen Substitutionen $\chi(\vartheta)$ des
Satzes 5. Da man nun in der identischen Gleichung $\Phi(\vartheta, z) = 0$
alle Permutationen der alternirenden Gruppe $A$ der Variablen $x$
ausführen kann, wodurch sich $z$ nicht ändert, während $\vartheta$ in jede
andere Wurzel von $\Phi$ übergeht, so ist die Gruppe der linearen
Substitutionen $\chi(\vartheta)$, d. h. die \gesperrt{Galois}'sche Gruppe der Glei-
%CONT>>


% ===== scan leaf 432, printed 411 =====
% !PRINTED 411
\weberpage{411}
%<<CONT
chung $\Phi$, mit der alternirenden Permutationsgruppe von $n$ Ziffern
(ein- oder mehrstufig) isomorph. Die Gruppe der linearen Substitutionen
 $\chi(\vartheta)$ möge mit $L$ bezeichnet sein. Sie muss, da sie
endlich ist, mit einer der im siebenten Abschnitte betrachteten
Polyëdergruppen identisch sein.

Zur Vereinfachung machen wir nun von dem in \S.\,51, 2.
bewiesenen Satze Gebrauch, nach dem sich die Gruppe $L$ so
transformiren lässt, dass eine beliebige der nicht identischen
Substitutionen von $L$ eine Multiplication wird. Eine Transformation
 der Gruppe $L$ ist aber gleichbedeutend damit, dass für $\vartheta$
eine lineare Function von $\vartheta$ gesetzt wird, der dieselbe Eigenschaft
 wie der ursprünglichen Function $\vartheta$ zukommt.

Bezeichnen wir also mit $\vartheta_\pi$ die Function der Variablen
$x_0, x_1, \ldots, x_{n-1}$, die aus $\vartheta$ durch Anwendung der Permutation $\pi$
hervorgeht, so können wir $\vartheta$ so wählen, dass für eine bestimmte,
aber beliebige Permutation $\pi$
\[
\vartheta_\pi = \varepsilon\vartheta
\]
wird. Wendet man $\pi$ wiederholt an, so folgt
\[
\vartheta_{\pi^2} = \varepsilon^2\vartheta,\ \vartheta_{\pi^3} = \varepsilon^3\vartheta, \ldots,
\]
und wenn also $p$ der Grad der Permutation $\pi$ ist, so ist $\varepsilon$ eine
$p^{\text{te}}$ Einheitswurzel.

Ist zunächst $n = 3$, so besteht die alternirende Gruppe aus
den Potenzen der cyklischen Permutation $\gamma = (0, 1, 2)$. Wir
können annehmen, dass die der Permutation entsprechende Substitution
 $\chi(\vartheta)$ multiplicativ sei, dass also $\vartheta_\gamma = \varepsilon\vartheta$ sei, worin $\varepsilon$
eine dritte Einheitswurzel ist. $\vartheta^3$ ist daher eine alternirende
(oder symmetrische) Function. Als Resolvente mit einem Parameter
 erhalten wir also die reine Gleichung
\[
\vartheta^3 = z,
\]
wo $z$ eine alternirende Function ist. Wir können etwa für $\vartheta$
die \gesperrt{Lagrange}'sche Resolvente $x_0 + \varepsilon x_1 + \varepsilon^2 x_2$ nehmen, und
erhalten die bekannte Reduction der cubischen Gleichung auf
eine reine Gleichung (Bd. I, \S.\,159).

Wir gehen zu dem Falle $n = 4$ über, in dem die alternirende
 Gruppe $A$ die Permutationen
\[
1,\ \alpha_1 = (0, 1)(2, 3),\ \alpha_2 = (0, 2)(1, 3),\ \alpha_3 = (0, 3)(1, 2)
\]
enthält, die eine Vierergruppe $B$ bilden; ausserdem kommen
%CONT>>


% ===== scan leaf 433, printed 412 =====
% !PRINTED 412
\weberpage{412}
%<<CONT
in $A$ noch acht cyklische Permutationen von je drei Ziffern
$\gamma = (0, 1, 2) \ldots$ vor, und $A$ kann so dargestellt werden:
\[
A = B + B\gamma + B\gamma^2.
\]
Je eine Permutation $\alpha$ und eine Permutation $\gamma$ können als
Erzeugende der Gruppe betrachtet werden. Denn ist etwa
\[
\gamma = (0, 1, 2),\ \alpha_1 = (0, 1)(2, 3),
\]
so ist
\weq{1}{\alpha_1\gamma\alpha_1 = \gamma' = (0, 3, 1),}
und aus $(0, 1, 2), (0, 3, 1)$ kann die ganze Gruppe $A$ abgeleitet
werden (Bd. I, \S.\,153, 7.).

Ebenso kann man $\gamma$ und $\gamma'$ als erzeugende Permutationen
von $A$ auffassen. Insbesondere ist
\weq{2}{\gamma'\gamma\gamma' = \alpha_1.}
Wir wählen $\vartheta$ so, dass der Permutation $\gamma$ eine Multiplication
$\varepsilon\vartheta$ entspricht, in der, da $\gamma$ vom Grade 3 ist, $\varepsilon$ eine dritte Einheitswurzel
 bedeutet. Wir setzen
\weq{3}{\vartheta = \frac{\varphi(x_0, x_1, x_2, x_3)}{\psi(x_0, x_1, x_2, x_3)},}
und verstehen unter $\varphi, \psi$ zwei ganze Functionen der vier
Variablen $x$ ohne gemeinschaftlichen Theiler.

Aus der identischen Gleichung
\[
\frac{\varphi_\gamma}{\psi_\gamma} = \varepsilon\,\frac{\varphi}{\psi}
\]
folgt dann, dass
\weq{4}{\varphi_\gamma = \varepsilon_1\varphi,\ \psi_\gamma = \varepsilon_2\varphi} % [sic]
sein muss, worin $\varepsilon_1, \varepsilon_2$ gleichfalls dritte Einheitswurzeln sind.

Wir wenden eine der Permutationen $\alpha$ auf $\vartheta$ an, und erhalten
 aus $\vartheta_\alpha = \chi(\vartheta)$:
\weq{5}{\frac{\varphi_\alpha}{\psi_\alpha} = \frac{a\varphi + b\psi}{c\varphi + d\psi}.}
Da nun $\varphi, \psi$, also auch $a\varphi + b\psi$ und $c\varphi + d\psi$, und
ebenso $\varphi_\alpha, \psi_\alpha$ ohne gemeinsamen Theiler sind, so muss in der
Identität (5) der Zähler dem Zähler und der Nenner dem
Nenner, wenigstens bis auf einen constanten Factor, gleich sein.
Diesen constanten Factor können wir in die Constanten $a, b, c, d$
einrechnen und erhalten
\[
\text{6)}\quad \varphi_\alpha = a\varphi + b\psi,\ \psi_\alpha = c\varphi + d\psi.
\]


% ===== scan leaf 434, printed 413 =====
% !PRINTED 413
\weberpage{413}
Nehmen wir hierin $\alpha=\alpha_1$ und $\alpha=\alpha_2$ an, und setzen zur
Abkürzung $\varphi_{\alpha_1}=\varphi_1$, $\varphi_{\alpha_2}=\varphi_2$, so
können wir aus den beiden Gleichungen
\[
\varphi_1 = a_1\varphi + b_1\psi, \quad \varphi_2 = a_2\varphi + b_2\psi
\]
$\psi$ eliminiren und erhalten eine Gleichung von der Form:
\weq{7}{h\varphi + h_1\varphi_1 + h_2\varphi_2 = 0,}
worin $h$, $h_1$, $h_2$ Constanten sind, unter denen wenigstens zwei von
Null verschieden sind. Auf die identische Gleichung (7) können wir nun die
Permutationen $\alpha_1$, $\alpha_2$ anwenden, und erhalten, da
$\alpha_1\alpha_2 = \alpha_2\alpha_1 = \alpha_3$, $\alpha_1\alpha_1 = 1$,
$\alpha_2\alpha_2 = 1$ ist,
\[
\begin{array}{l|l}
h\varphi + h_1\varphi_1 + h_2\varphi_2 = 0 & h,\\
h\varphi_1 + h_1\varphi + h_2\varphi_3 = 0 & h_1,\\
h\varphi_2 + h_1\varphi_3 + h_2\varphi = 0 & -h_2,
\end{array}
\]
woraus durch Multiplication mit den daneben stehenden Factoren und
Addition:
\weq{8}{(h^2 + h_1^2 - h_2^2)\varphi + 2hh_1\varphi_1 = 0.}
Wenn also $h$ und $h_1$ von Null verschieden sind, so unterscheiden sich
$\varphi$ und $\varphi_1$ nur durch einen constanten Factor (der $=\pm 1$
sein muss, da $\alpha_1$ vom 2\textsuperscript{ten} Grade ist). Ist aber
$h$ oder $h_1=0$, so folgt aus (7) $\varphi_1=\pm\varphi_2$ oder
$\varphi=\pm\varphi_2$, und die Anwendung von $\alpha_1$ auf die erste
Gleichung giebt $\varphi=\pm\varphi_3$. Es findet also jedenfalls \gesperrt{eine der
Relationen} statt:
\weq{9}{\varphi = \pm\varphi_1, \quad \varphi = \pm\varphi_2, \quad \varphi = \pm\varphi_3.}
Da man nun ebensowohl $\gamma,\alpha_1$ als $\gamma,\alpha_2$, als auch
$\gamma,\alpha_3$ als erzeugende Elemente der Gruppe $A$ betrachten kann,
so ergiebt dies Resultat, zusammengenommen mit (4), den Satz:

1. \gesperrt{Die Function $\varphi$ hat die Eigenschaft, sich durch alle
Permutationen der alternirenden Gruppe nur um einen constanten Factor zu
ändern.}

Ganz derselbe Schluss ist aber auch auf $\psi$ anwendbar, und also auch
auf den Quotienten $\vartheta$ beider Functionen. Es ist sonach für jede
Permutation $\pi$ der alternirenden Gruppe
\weq{10}{\vartheta_\pi = \varepsilon\vartheta.}
Hierin ist, wenn $\pi$ zu den cyklischen Permutationen $\gamma$ gehört,
$\varepsilon$ eine dritte Einheitswurzel. Daraus ist aber ferner zu
schliessen, weil alle Permutationen $\pi$ aus $\gamma$, $\gamma'$
zusammengesetzt
%CONT>>


% ===== scan leaf 435, printed 414 =====
% !PRINTED 414
\weberpage{414}
%<<CONT
werden können, dass die in (10) vorkommende Constante $\varepsilon$ \gesperrt{immer}
eine dritte Einheitswurzel ist. Wenn $\pi$ zu den $\alpha$ gehört, so ist
$\varepsilon$ zugleich eine zweite Einheitswurzel und muss also $=1$ sein.
Daraus folgern wir den Satz:

2. \gesperrt{Die Function $\vartheta$ bleibt bei den Permutationen der
Vierergruppe $A$ ungeändert.}

Die dritte Potenz von $\vartheta$ gestattet die Permutationen der
alternirenden Gruppe und $\vartheta$ ist daher die Wurzel einer reinen
cubischen Gleichung $\vartheta^3 = z$. Dies ist also die Resolvente
$\Phi(\vartheta,z) = 0$, die in diesem Falle eine \gesperrt{Partialresolvente} ist
(Bd. I, \S.\,161).

Aus diesen Betrachtungen ergiebt sich nun ohne
Schwierigkeit der folgende Satz, dessen Beweis das Ziel dieser ganzen
Betrachtungen ist:

3. \gesperrt{Eine allgemeine Gleichung von höherem als
4\textsuperscript{ten} Grade hat keine Resolventen mit einem Parameter.}

Denn ist $n$ grösser als 4, so können wir, wenn wir $\vartheta$ als
Function von je vieren der Variablen $x$ betrachten, den Satz 2. anwenden.
Es folgt dann, dass $\vartheta$ ungeändert bleibt, wenn unter den
Variablen irgend ein Paar von Transpositionen vorgenommen wird. Da man
nun aus Transpositionspaaren die ganze alternirende Gruppe zusammensetzen
kann (Bd. I, \S.\,153, 8.), so folgt, dass $\vartheta$ überhaupt durch
die alternirende Gruppe ungeändert bleibt, also eine alternirende oder
eine symmetrische Function ist. $\Phi(\vartheta,z) = 0$ reducirt sich auf
die Identität $\vartheta = z$ und ist keine Resolvente.

Es sei noch bemerkt, dass, wenn man an Stelle der alternirenden Gruppe
die symmetrische treten lässt, dieser Satz a fortiori gilt.

Den Satz, dessen Beweis wir hier entwickelt haben, hat zuerst
\gesperrt{Kronecker} ausgesprochen, ohne einen Beweis dafür zu
veröffentlichen. Der erste Beweis ist von \gesperrt{F. Klein} gegeben, der
dann von \gesperrt{Gordan} noch vereinfacht
ist\footnote{\gesperrt{F. Klein}, Vorlesungen über das Ikosaëder. Der
Beweis von \gesperrt{Gordan}, dem unsere Darstellung in den Grundzügen
folgt, findet sich in Bd. 29 der mathematischen Annalen (1887).}.


% ===== scan leaf 436, printed 415 =====
% !PRINTED 415
\weberpage{415}
\pdfbookmark[2]{§ 110. Die Ikosaëdergleichung}{bm126}
\begin{center}
\S.\,110.

\gesperrt{Die Ikosaëdergleichung.}
\end{center}

Das Resultat dieser Betrachtungen ist zunächst, sofern es die allgemeinen
Gleichungen von höherem als dem 4\textsuperscript{ten} Grade betrifft,
ein negatives. Es gelingt nicht, durch rationale Resolventenbildung die
allgemeine Gleichung 5\textsuperscript{ten} Grades auf eine Gleichung mit
\gesperrt{einem} Parameter zu reduciren. Wollen wir gleichwohl dies Ziel noch nicht
aufgeben, so bleibt nichts übrig, als dass wir die Variablen
$x_0, x_1, \ldots, x_{n-1}$ nicht mehr als völlig unabhängige Variable
auffassen, dass wir zwischen ihnen gewisse Gleichungen bestehen lassen.
Wir haben es freilich dann nicht mehr mit der allgemeinen Gleichung des
betreffenden Grades zu thun, sondern mit einer specielleren, bei der die
Co\"efficienten nicht unabhängig von einander sind, und es stellt sich
die zweite Frage ein, wie und durch welche irrationale Functionen sich
die allgemeine Gleichung auf diese specielle zurückführen lässt. Gelingt
dies durch Gleichungen niedrigeren Grades, z. B. durch Quadratwurzeln, so
wird immerhin eine Reduction des Problems erreicht sein.

Wir nehmen also jetzt an, dass die Variablen $x_0, x_1, \ldots, x_{n-1}$
durch eine beliebige Anzahl von Relationen
\weq{1}{A(x_0, x_1, \ldots, x_{n-1}) = 0, \quad B(x_0, x_1, \ldots, x_{n-1}) = 0, \ldots}
mit einander verknüpft seien, worin die $A, B, \ldots$ rationale
Functionen der $x$ sind, die an Anzahl und gegenseitiger Abhängigkeit so
beschaffen sind, dass die Werthe der $x$ nicht numerisch festgelegt,
sondern noch variabel sind. Ausserdem legen wir eine Permutationsgruppe
$\mathfrak{A}$ der $x_0, x_1, \ldots, x_{n-1}$ zu Grunde, und betrachten
die Functionen der $x$, die die Permutationen von $\mathfrak{A}$
gestatten, und nur diese als rational, so dass $\mathfrak{A}$ die
Galois'sche Gruppe der Gleichung ist, deren Wurzeln die
$x_0, x_1, \ldots, x_{n-1}$ sind. \gesperrt{Die Relationen (1) müssen dann so
beschaffen sein, dass auch sie die Permutationen $\mathfrak{A}$
gestatten.}

Wir fragen unter diesen Voraussetzungen nach der Möglichkeit, zwei
rationale Functionen $u, z$ von $x_0, x_1, \ldots, x_{n-1}$ so zu
bestimmen, dass zwischen ihnen eine Gleichung
\weq{2}{F(u, z) = 0}
%CONT>>


% ===== scan leaf 437, printed 416 =====
% !PRINTED 416
\weberpage{416}
%<<CONT
besteht, die für alle den Relationen (1) genügenden Werthe der $x$
befriedigt ist und worin $z$, immer mit Rücksicht auf die Relationen (1),
ungeändert bleibt, wenn die $x$ den Permutationen der Gruppe
$\mathfrak{A}$ unterworfen werden. Die Function $u$ soll dadurch in
$\nu$ verschiedene Functionen
\weq{3}{u, u_1, \ldots, u_{\nu-1}}
übergehen, so dass die Grössen (3) die Wurzeln der Gleichung (2) sind, und
wenn $\nu > 1$ ist, (2) eine Resolvente der Gleichung für die $x$ ist.

In dieser Allgemeinheit haben wir für die Lösung der Frage bis jetzt noch
keinen Angriffspunkt. Wir fügen daher eine weitere beschränkende
Voraussetzung hinzu, nämlich: \gesperrt{Es soll möglich sein, für die $x_h$
rationale Functionen einer Variablen $t$}
\weq{4}{x_h = \varphi_h(t)}
\gesperrt{zu setzen, so dass die Relationen (1) identisch befriedigt sind, und dass
zugleich $z$ nicht von $t$ unabhängig
wird}\footnote{Betrachtet man die $x$ und die $u$ als Coordinaten je eines
Punktes $X$ und $U$ in einem Raume von $n$ und $\nu$ Dimensionen, so
bestimmen die Relationen (1) für $x$ eine Mannigfaltigkeit von weniger
als $n$ Dimensionen. Jedem Punkte $X$ entspricht ein Punkt $U$, und der
Gesammtheit der $X$ nach (2) eine Mannigfaltigkeit der $U$ von einer
Dimension (einer Curve). Unsere Voraussetzung ist die, dass dies eine
sogenannte rationale oder unicursale Curve oder (in der Sprache der
Functionentheorie) eine Curve vom Geschlecht Null sei.}.

Um gleich hier ein Beispiel anzuführen, mit dem wir uns später noch
eingehend beschäftigen werden, erwähnen wir den Fall von fünf Variablen
$x_0, x_1, x_2, x_3, x_4$, die den Bedingungen
\weq{5}{\Sigma x = 0, \quad \Sigma x^2 = 0}
genügen, und die also die Wurzeln einer \gesperrt{Hauptgleichung}
5\textsuperscript{ten} Grades sind (Bd. I, \S.\,54). Dass die
Voraussetzung, die wir gemacht haben, in diesem Falle zutrifft, zeigt die
Darstellung in Bd. I, \S.\,74, (4):
\[
y = t_3 F_0 + t_2(\alpha_2 F_1 + \alpha_1 F_2 + \alpha_0 F_3),
\]
was, wenn die dort angeführten Bestimmungen getroffen sind, für die $x$
gesetzt, den Bedingungen (5) für alle Werthe des Verhältnisses
$t = t_2 : t_3$ genügt.


% ===== scan leaf 438, printed 417 =====
% !PRINTED 417
\weberpage{417}
Auch wie die allgemeine Gleichung $5^{\mathrm{ten}}$ Grades auf eine diesen Voraussetzungen entsprechende transformirt werden kann, ist dort gezeigt.

Unter der Voraussetzung (4) hat das Bestehen der Relationen (1) keinen Einfluss auf die Schlüsse des \S.\,108, und wir haben dann die folgenden Sätze.

Es sei $x_0, x_1, \ldots, x_{n-1}$ ein System von Variablen, das den Relationen $A=0$, $B=0$, \ldots\ unterworfen, aber sonst unabhängig ist, und es sei $\mathfrak{A}$ eine Permutationsgruppe der $x$, deren Permutationen die Functionen $A, B, \ldots$ ungeändert lassen.

\gesperrt{Soll die Gleichung, deren Wurzeln die Grössen $x$ sind, eine Resolvente}
\weq{6}{F(u,z) = 0}
\gesperrt{besitzen, die ausser numerischen Coëfficienten nur einen Parameter enthält, der eine durch die Gruppe $\mathfrak{A}$ ungeänderte Function der $x$ ist, so muss eine Function $\vartheta$ der Variablen $x_0, \ldots, x_{n-1}$ existiren, die durch die Permutationen $\alpha$ von $\mathfrak{A}$ (mit Rücksicht auf die Relationen $A=0$, $B=0$, \ldots) lineare Substitutionen erleidet}
\weq{7}{\vartheta_\alpha = \dfrac{a\vartheta+b}{c\vartheta+d},}
\gesperrt{worin $a, b, c, d$ numerische Coëfficienten sind. Die Wurzeln der Gleichung (6) sind rational durch $\vartheta$ ausdrückbar, und folglich ist auch $\vartheta$ Wurzel einer Gleichung mit einem Parameter}
\weq{8}{\Phi(\vartheta, z) = 0.}

Ist (6) eine Totalresolvente, was immer eintritt, wenn die Gruppe $\mathfrak{A}$ einfach ist, so können auch die $x$ rational durch $\vartheta$ und die zur Gruppe $\mathfrak{A}$ gehörigen Functionen ausgedrückt werden.

Bis hierher findet also vollständige Uebereinstimmung mit den Resultaten des \S.\,108 statt.

Die Beschränkung aber, auf die wir im \S.\,109 gestossen sind, wird hier nicht mehr nothwendig eintreten, weil jetzt nicht mehr, wie im Falle völlig unabhängiger Variablen, aus der Gleichheit zweier gebrochener Functionen auf die Uebereinstimmung von Zähler und Nenner geschlossen werden kann.

Die linearen Substitutionen (7) bilden eine mit $\mathfrak{A}$ (ein- oder mehrstufig) isomorphe Gruppe $P$, die mit einer unserer
%CONT>>


% ===== scan leaf 439, printed 418 =====
% !PRINTED 418
\weberpage{418}
%<<CONT
\gesperrt{Polyëdergruppen} identisch sein muss. Wir wollen hier nur den interessantesten Fall eingehender behandeln, dass $P$ die \gesperrt{Ikosaëdergruppe} ist, die wir überdies ohne Beschränkung der Allgemeinheit in der Normalform (\S.\,58) annehmen können.

Für diese Gruppe haben wir in \S.\,60 die drei Grundformen abgeleitet. Bezeichnen wir die Variablen dieser Grundformen mit $y_1, y_2$, so sind es die drei homogenen Functionen $12^{\mathrm{ten}}$, $20^{\mathrm{sten}}$ und $30^{\mathrm{sten}}$ Grades:
\weq{9}{\begin{aligned}
f(y) &= y_1 y_2 (y_1^{10} + 11 y_1^5 y_2^5 - y_2^{10}) \\
H(y) &= -y_1^{20} - y_2^{20} + 228 (y_1^{15} y_2^5 - y_2^{15} y_1^5) - 494 y_1^{10} y_2^{10} \\
T(y) &= y_1^{30} + y_2^{30} + 522 (y_1^{25} y_2^5 - y_1^5 y_2^{25}) - 10005 (y_1^{20} y_2^{10} + y_2^{20} y_1^{10}),
\end{aligned}}
zwischen denen noch die Relation
\weq{10}{T^2 + H^3 = 1728 f^5}
besteht.

Wenn wir in dem Quotienten $T^2 : f^5$ für das Verhältniss $y_1 : y_2$ einen der 60 Werthe $\vartheta_\alpha$ setzen, die aus $\vartheta$ durch die Ikosaëdersubstitutionen entstehen, so erhält dieser Quotient immer denselben Werth, der sich also rational durch die Coëfficienten der Gleichung (8), d.\,h. rational durch $z$ ausdrücken lässt. Wir können ihn geradezu gleich $z$ setzen, und erhalten für die Resolvente (8) die Form
\weq{11}{T^2 - zf^5 = 0,}
was wegen (10) mit
\[
H^3 - (1728 - z) f^5 = 0
\]
gleichbedeutend ist. Diese Gleichung ist in Bezug auf $y_1, y_2$ homogen und vom $60^{\mathrm{sten}}$ Grade, und wenn wir $y_2 = 1$ setzen, so sind ihre 60 Wurzeln $y_1 = \vartheta_\alpha$. Es ist die \gesperrt{Ikosaëdergleichung}, die schon im \S.\,60 definirt war, und die also ausführlich, in der Form (11) geschrieben, so lautet:
\[
(\vartheta^{30} + 522 \vartheta^{25} - 10005 \vartheta^{20} - 10005 \vartheta^{10} - 522 \vartheta^5 + 1)^2 = z \vartheta^5 (\vartheta^{10} + 11 \vartheta^5 - 1)^5.
\]

Das Problem, wie wir es also jetzt gefasst haben, kommt auf die Frage hinaus:

\gesperrt{Welche Gleichungen lassen sich durch die Ikosaëdergleichung lösen?}


% ===== scan leaf 440, printed 419 =====
% !PRINTED 419
\weberpage{419}
\pdfbookmark[2]{§ 111. Die Resolventen der Ikosaëdergleichung}{bm127}
\begin{center}
\S.\,111.

\gesperrt{Die Resolventen der Ikosaëdergleichung.}
\end{center}

Die Formulirung des Problems, die wir am Schluss des vorigen Paragraphen gegeben haben, führt uns auf die Frage nach den verschiedenen Resolventen der Ikosaëdergleichung. Jedem Theiler der Ikosaëdergruppe entspricht eine solche Resolvente, deren Grad gleich dem Index des Theilers, also immer ein Theiler von 60 ist. Nun enthält die Ikosaëdergruppe (\S.\,59):

\begin{tabular}{@{}ll c l@{}}
1) Tetraëdergruppen (\S.\,56): & Resolventen & $5^{\mathrm{ten}}$ & Grades \\
2) Diëdergruppen $D_5$ (\S.\,55): & \textquotedbl & $6^{\mathrm{ten}}$ & \textquotedbl \\
3) Diëdergruppen $D_3$: & \textquotedbl & $10^{\mathrm{ten}}$ & \textquotedbl \\
4) Cyklische Gruppen $C_5$: & \textquotedbl & $12^{\mathrm{ten}}$ & \textquotedbl \\
5) Vierergruppen $D_2$: & \textquotedbl & $15^{\mathrm{ten}}$ & \textquotedbl \\
6) Cyklische Gruppen $C_3$: & \textquotedbl & $20^{\mathrm{sten}}$ & \textquotedbl \\
7) Cyklische Gruppen $C_2$: & \textquotedbl & $30^{\mathrm{sten}}$ & \textquotedbl \\
\end{tabular}

Da die Ikosaëdergruppe einfach ist, sind alle diese Gleichungen \gesperrt{Totalresolventen} von einander; wir beschränken uns hier auf die Betrachtung der wichtigsten unter ihnen.

Die Resolventen niedrigsten, nämlich $5^{\mathrm{ten}}$ Grades sind die zu der Tetraëdergruppe gehörigen.

Um sie zu finden, müssen wir eine Function von $\vartheta$ suchen, die durch die Substitutionen der Tetraëdergruppe ungeändert bleibt, während sie in der Ikosaëdergruppe fünfwerthig ist.

Wir nehmen die Ikosaëdergruppe in der in \S.\,58, (22) aufgestellten Normalform, und nehmen die darin enthaltene Tetraëdergruppe $Q$ (\S.\,59) heraus, die wir über der Vierergruppe
\weq{1}{1, \psi, \chi, \psi\chi}
construirt haben, worin, wenn $\varepsilon$ eine imaginäre fünfte Einheitswurzel ist, $\psi$ und $\chi$ die Bedeutung haben:
\weq{2}{\psi(x) = \dfrac{-1}{x}, \quad \chi(x) = \dfrac{\omega x+1}{x-\omega}, \quad \omega = \varepsilon + \varepsilon^{-1}.}

Setzen wir $\Theta(x) = \varepsilon x$, so wird die ganze Ikosaëdergruppe
\weq{3}{P = Q + Q\Theta + Q\Theta^2 + Q\Theta^3 + Q\Theta^4,}
und
\weq{4}{Q_h = \Theta^{-h} Q \Theta^h}
sind die zu $Q$ conjugirten Tetraëdergruppen.


% ===== scan leaf 441, printed 420 =====
% !PRINTED 420
\weberpage{420}
Nun haben wir im \S.\,56 zwei Formen $f$ und $H$ vom $6^{\mathrm{ten}}$ und $8^{\mathrm{ten}}$ Grade kennen gelernt, die durch die Tetraëdersubstitutionen, wenn sie homogen und mit der Determinante 1 genommen werden, absolut ungeändert bleiben. Diese Formen können wir freilich nicht geradezu in der Form anwenden, wie sie dort gegeben sind, weil dort die Tetraëdergruppe in anderer Gestalt vorausgesetzt war. Wir finden aber die erste dieser Formen sehr einfach, wenn wir uns erinnern, dass ihre Wurzeln die zweizähligen Pole der Tetraëdergruppe waren, also die Wurzeln der drei Gleichungen
\[
x = \psi(x), \quad x = \chi(x), \quad x = \psi\chi(x)
\]
oder
\[
x^2 + 1 = 0, \quad x^2 - 2\omega x - 1 = 0, \quad \omega x^2 + 2x - \omega = 0,
\]
und man erhält also mit Benutzung der Relation
\[
\omega^2 + \omega = 1
\]
die Gleichung für die zweizähligen Pole
\[
(x^2+1)(x^4+2x^3-6x^2-2x+1) = 0.
\]

Daher wird die gesuchte Tetraëderform $6^{\mathrm{ten}}$ Grades in den homogenen Variablen $x_1, x_2$
\weq{5}{t(x_1,x_2) = x_1^6 + 2x_1^5x_2 - 5x_1^4x_2^2 - 5x_1^2x_2^4 - 2x_1x_2^5 + x_2^6.}

Die zweite Tetraëderform, die wir suchen, ist die \gesperrt{Hesse}'sche Determinante hiervon.

Setzen wir
\[
\tilde\omega(x_1,x_2) = \dfrac{1}{5^2 \cdot 4^2} \left\{ t''(x_1,x_1)\, t''(x_2,x_2) - [t''(x_1,x_2)]^2 \right\},
\]
so ergiebt eine einfache Rechnung
\weq{6}{\tilde\omega(x_1,x_2) = -(x_1^8+x_2^8) + (x_1^7x_2 - x_1x_2^7) - 7(x_1^6x_2^2+x_1^2x_2^6) - 7(x_1^5x_2^3-x_1^3x_2^5).}

Wir führen noch die drei Ikosaëderformen $12^{\mathrm{ten}}$, $20^{\mathrm{sten}}$ und $30^{\mathrm{sten}}$ Grades an, die ja auch durch die Tetraëdersubstitutionen ungeändert bleiben (\S.\,60):
\weq{7}{f(x_1,x_2) = x_1x_2(x_1^{10}+11x_1^5x_2^5-x_2^{10})}
\weq{8}{H(x_1,x_2) = -(x_1^{20}+x_2^{20}) + 228(x_1^{15}x_2^5-x_1^5x_2^{15}) - 494x_1^{10}x_2^{10}}
\weq{9}{T(x_1,x_2) = (x_1^{30}+x_2^{30}) + 522(x_1^{25}x_2^5-x_1^5x_2^{25}) - 10005(x_1^{20}x_2^{10}+x_2^{20}x_1^{10}).}


% ===== scan leaf 442, printed 421 =====
% !PRINTED 421
\weberpage{421}

Bilden wir aus diesen invarianten Formen des Tetra\"eders Functionen von $x_1:x_2$ und setzen $\vartheta = x_1:x_2$, so ergeben sich Functionen, die durch die linearen gebrochenen Tetra\"edersubstitutionen unge\"andert bleiben, und die daher Wurzeln von Resolventen 5\textsuperscript{ten} Grades der Ikosa\"edergleichung sind. Die Ikosa\"edergleichung selbst erh\"alt man in zwei Formen, wenn man
\weq{10}{H^3 = zf^5,\quad T^2 = z_1f^5,\quad z + z_1 = 1728}
setzt, und die Resolventen h\"angen rational von $z$ ab.

Bezeichnen wir f\"ur den Augenblick mit $\varphi(x_1, x_2)$ irgend eine absolut invariante Form des Tetra\"eders, so geht diese Form durch irgend eine der Substitutionen $Q\Theta$ in
\weq{11}{\varphi_\varepsilon = \varphi(\varepsilon^{-2}x_1,\ \varepsilon^2x_2)}
\"uber, und diese Function bleibt unge\"andert durch die Substitutionen der Gruppe $\Theta^{-1}Q\Theta$.

\gesperrt{Daraus ergiebt sich, dass jede symmetrische Function der f\"unf Formen $\varphi_\varepsilon$ eine Invariante der Ikosa\"edergruppe ist, und daher nach dem Satze \S.\,61 durch die Grundformen $f, H, T$ ausgedr\"uckt werden kann.}

Dieser Satz gestattet eine verh\"altnissm\"assig leichte Berechnung der Resolventen.

Die einfachsten Functionen, die wir als Wurzeln von Resolventen verwenden k\"onnen, sind
\weq{12}{r = \frac{t^2}{f},\quad u = \frac{f^2t}{T},\quad v = \frac{f\tilde\omega}{H}.}

Bemerken wir zun\"achst, dass die symmetrischen Grundfunctionen der f\"unf Functionen $t_\varepsilon$ Ikosa\"ederinvarianten der Grade $6, 12, 18, 30$ sind, so folgt aus den S\"atzen in \S.\,61, da die Gradzahlen $6, 18$ unter diesen Invarianten nicht vorkommen, eine Gleichung von der Form
\weq{13}{t^5 + aft^3 + bf^2t + cT = 0,}
worin $a, b, c$ Zahlenco\"efficienten sind. Um sie zu bestimmen, bezeichnen wir mit $S(\varphi)$ die Summe der f\"unf Functionen $\varphi_\varepsilon$ in (11) und mit $\Pi(\varphi)$ ihr Product, und erhalten aus den \gesperrt{Newton}'schen Formeln (Bd.\,I, \S.\,42):
\weq{14}{af = -\tfrac{1}{2}S(t^2),\quad bf^2 = \tfrac{1}{2}a^2f^2 - \tfrac{1}{4}S(t^4),\quad cT = -\Pi(t).}

In den Summen $S(\varphi)$ haben nur solche Glieder $x_1^h x_2^k$ einen von Null verschiedenen Co\"efficienten, in denen $h - k$ durch 5 theilbar ist. Man erh\"alt also $a, b, c$ durch Gleichsetzen der
%CONT>>


% ===== scan leaf 443, printed 422 =====
% !PRINTED 422
\weberpage{422}
\pdfbookmark[2]{§ 112. Die Hauptresolvente fünften Grades}{bm128}

%<<CONT
Co\"efficienten von $x_1^{11}x_2,\ x_1^{22}x_2^2,\ x_1^{30}$ auf beiden Seiten der Gleichungen (14), und findet so die Relation
\weq{15}{t^5 - 10ft^3 + 45f^2t - T = 0,}
oder auch
\weq{16}{t^2(t^4 - 10ft^2 + 45f^2)^2 - T^2 = 0.}

Aus (16) und (15) ergeben sich nun nach (10) und (12) die Gleichungen 5\textsuperscript{ten} Grades f\"ur $r$ und f\"ur $u$
\weq{17}{r(r^2 - 10r + 45)^2 = z_1,}
\weq{18}{u^5 - \frac{10u^3}{z_1} + \frac{45u}{z_1^2} - \frac{1}{z_1^2} = 0.}

Die letzte dieser Gleichungen geht durch die Substitution
\weq{19}{y = -\frac{1}{3u},\quad z_1 = 27\gamma}
in die Form
\weq{20}{y^5 + 15y^4 - 10\gamma y^2 + 3\gamma^2 = 0}
\"uber, die wir bereits im \S.\,75 des ersten Bandes als eine Normalform der Gleichung 5\textsuperscript{ten} Grades kennen gelernt haben.

\begin{center}
\S.\,112.\\[4pt]
\gesperrt{Die Hauptresolvente f\"unften Grades.}
\end{center}

Wenn man die Functionen $u, v$ gleichzeitig benutzt, so kann man eine Schaar von Resolventen ableiten, die die Form der Hauptgleichung 5\textsuperscript{ten} Grades haben, und sich direct mit einer beliebig gegebenen Hauptgleichung 5\textsuperscript{ten} Grades in Uebereinstimmung bringen lassen\footnote{\gesperrt{Kiepert}, G\"ottinger Nachrichten 1878. Crelle's Journal, Bd.\,87. \gesperrt{Klein}, Ikosa\"eder, S.\,106.}.

Da es beim Ikosa\"eder keine Invarianten 8\textsuperscript{ten}, 14\textsuperscript{ten}, 16\textsuperscript{ten}, 22\textsuperscript{sten} oder 28\textsuperscript{sten} Grades giebt, so m\"ussen die Functionen $S(\tilde\omega)$, $S(t\tilde\omega)$, $S(\tilde\omega^2)$, $S(t\tilde\omega^2)$, $S(t^2\tilde\omega^2)$ verschwinden, und wenn wir also
\weq{1}{Y = \alpha\tilde\omega + \beta t\tilde\omega}
setzen, so werden auch die Functionen
\[ S(Y),\quad S(Y^2) \]
f\"ur beliebige Werthe der Parameter $\alpha, \beta$ verschwinden, und daraus ergiebt sich eine identische Gleichung von der Form
%CONT>>


% ===== scan leaf 444, printed 423 =====
% !PRINTED 423
\weberpage{423}

%<<CONT
\weq{2}{Y^5 + 5aY^2 + 5bY + c = 0,}
worin $a, b, c$ homogene Functionen 3\textsuperscript{ten}, 4\textsuperscript{ten}, 5\textsuperscript{ten} Grades von $\alpha, \beta$ bedeuten, deren Co\"efficienten Ikosa\"ederinvarianten sind.

Die Function $c$ k\"onnen wir leicht aus der Formel bestimmen
\[ c = -\Pi(\tilde\omega)\,\Pi(\alpha + \beta t). \]

Es ist n\"amlich nach (15) des vorigen Paragraphen f\"ur ein unbestimmtes $\lambda$
\[ \lambda^5 - 10f\lambda^3 + 45f^2\lambda - T = \Pi(\lambda - t), \]
also, wenn man $\lambda = -\alpha : \beta$ setzt
\[ \Pi(\alpha + \beta t) = \alpha^5 - 10f\alpha^3\beta^2 + 45f^2\alpha\beta^4 + T\beta^5, \]
und ferner ergiebt sich ohne Weiteres durch Vergleichung eines Gliedes [\S.\,111, (6), (8)]
\[ \Pi(\tilde\omega) = -H^2, \]
also
\weq{3}{c = H^2(\alpha^5 - 10f\alpha^3\beta^2 + 45f^2\alpha\beta^4 + T\beta^5).}

Die Co\"efficienten $a, b$ berechnet man wohl am einfachsten mit H\"ulfe der \gesperrt{Newton}'schen Formeln (Bd.\,I, \S.\,42):
\[
\begin{aligned}
-15a &= S(\alpha\tilde\omega + \beta t\tilde\omega)^3\\
-20b &= S(\alpha\tilde\omega + \beta t\tilde\omega)^4,
\end{aligned}
\]
indem man die Formeln anwendet, die sich nach kurzer Rechnung durch Vergleichung je eines Gliedes der rechten und linken Seite ergeben:
\[
\begin{array}{ll}
S(\tilde\omega^3) = -3.5.8f^2, & S(t\tilde\omega^3) = -5T,\\
S(t^2\tilde\omega^3) = -5.72f^3, & S(t^3\tilde\omega^3) = -3.5fT,\\
S(\tilde\omega^4) = 4.5fH, & S(t^2\tilde\omega^4) = -5.12f^2H,\\
S(t^3\tilde\omega^4) = -5HT, & S(t^4\tilde\omega^4) = -4.5.27f^3H.
\end{array}
\]

So findet sich
\weq{4}{a = 8f^2\alpha^3 + T\alpha^2\beta + 72f^3\alpha\beta^2 + fT\beta^3}
\weq{5}{b = -fH\alpha^4 + 18f^2H\alpha^2\beta^2 + HT\alpha\beta^3 + 27f^3H\beta^4.}

Um daraus die Resolvente der Ikosa\"edergleichung zu erhalten, m\"ussen wir statt $\tilde\omega$ und $t$ die Functionen $u, v$ [\S.\,111, (12)] einf\"uhren.

Setzen wir, indem wir mit $\lambda, \mu$ irgend zwei unbestimmte Gr\"ossen bezeichnen
\weq{6}{\alpha = \frac{\lambda f}{H},\quad \beta = \frac{\mu f^3}{HT},}
%CONT>>


% ===== scan leaf 445, printed 424 =====
% !PRINTED 424
\weberpage{424}

%<<CONT
so folgt nach (1) und \S.\,111, (12):
\weq{7}{Y = \lambda v + \mu vu,}
und wenn wir, wie in \S.\,111
\weq{8}{\frac{H^3}{f^5} = z,\quad \frac{T^2}{f^5} = z_1,\quad z + z_1 = 1728}
setzen, so gehen die Ausdr\"ucke (3), (4), (5) in folgende \"uber:
\weq{9}{az = 8\lambda^3 + \lambda^2\mu + \frac{72\lambda\mu^2 + \mu^3}{z_1}}
\[
bz = -\lambda^4 + \frac{18\lambda^2\mu^2 + \lambda\mu^3}{z_1} + 27\frac{\mu^4}{z_1^2}
\]
\[
cz = \lambda^5 - 10\frac{\lambda^3\mu^2}{z_1} + \frac{45\lambda\mu^4 + \mu^5}{z_1^2},
\]
wenn $Y$ die Wurzel der Gleichung
\weq{10}{Y^5 + 5aY^2 + 5bY + c = 0}
ist. Diese Gleichung hat die Form einer Hauptgleichung 5\textsuperscript{ten} Grades, und sie kann jede Hauptgleichung darstellen, wenn wir $a, b, c$ beliebig annehmen k\"onnen. Man hat also, wenn $a, b, c$ beliebig gegeben angenommen werden, aus (9) die drei Gr\"ossen $\lambda, \mu, z_1$ (und $z = 1728 - z_1$) zu bestimmen, und es ist eine sehr merkw\"urdige Eigenth\"umlichkeit dieses Gleichungssystemes, dass sich diese Unbekannten \gesperrt{rational} durch die gegebenen Gr\"ossen $a, b, c$ und die Discriminante der Gleichung (10) ausdr\"ucken lassen. Um dies nachzuweisen, leiten wir aus (9) zun\"achst einfachere Gleichungen ab.

Wenn wir die zweite von ihnen mit $\lambda$ multipliciren und zur dritten addiren, so folgt
\weq{11}{z_1(\lambda b + c) = \mu^2a,}
und wenn wir die dritte mit $\lambda$, die zweite mit $\mu^2 : z_1$ multipliciren und subtrahiren
\weq{12}{z\left(\lambda c - \frac{\mu^2}{z_1}b\right) = \left(\lambda^2 - 3\frac{\mu^2}{z_1}\right)^3.}

Ebenso ergiebt sich aus der ersten und zweiten
\[
z\,\frac{a\lambda + 8b}{\mu} = \lambda^3 + \frac{216\lambda^2\mu}{z_1} + 9\frac{\lambda\mu^2}{z_1} + \frac{216\mu^3}{z_1^2}.
\]


% ===== scan leaf 446, printed 425 =====
% !PRINTED 425
\weberpage{425}
Diese Gleichung ergiebt, wenn man beiderseits zum Quadrat erhebt:
\[
\begin{aligned}
z_1 z^2\left(\frac{a\lambda+8b}{\mu}\right)^2 &= \lambda^6 z_1 + 2.216\,\lambda^5\mu + \left(18+\frac{216^2}{z_1}\right)\lambda^4\mu^2 \\
&\quad + \frac{2.2160}{z_1}\,\lambda^3\mu^3 + \left(\frac{81}{z_1}+\frac{2.216^2}{z_1^2}\right)\lambda^2\mu^4 \\
&\quad + \frac{18.216}{z_1^2}\,\lambda\,\mu^5 + \frac{216^2\mu^6}{z_1^3},
\end{aligned}
\]
und wenn man dies abzieht von der aus der ersten Gleichung (9) abgeleiteten Formel:
\[
\begin{aligned}
27a^2z^2 &= 1728\,\lambda^6 + 2.216\,\lambda^5\mu + 27\left(1+\frac{16.72}{z_1}\right)\lambda^4\mu^2 \\
&\quad + \frac{2.2160}{z_1}\,\lambda^3\mu^3 + 27\left(\frac{2}{z_1}+\frac{72^2}{z_1^2}\right)\lambda^2\mu^4 \\
&\quad + \frac{18.216}{z_1^2}\,\lambda\,\mu^5 + \frac{27\mu^6}{z_1^2},
\end{aligned}
\]
so folgt mit R\"ucksicht auf (8):
\weq{13}{z\left[27a^2 - z_1\left(\frac{a\lambda+8b}{\mu}\right)^2\right] = \left(\lambda^2-\frac{3\mu^2}{z_1}\right)^3.}
Daraus ergiebt sich weiter durch Vergleichung mit (12)
\weq{14}{27a^2 - \frac{z_1}{\mu^2}(a\lambda+8b)^2 = \lambda c - \frac{\mu^2}{z_1}b,}
und wenn man hieraus durch (11) das Verh\"altniss $\mu^2:z_1$ eliminirt, so ergiebt sich die quadratische Gleichung f\"ur $\lambda$:
\weq{15}{\begin{split} \lambda^2(a^4+abc-b^3) - \lambda(11a^3b-ac^2+2b^2c) \\ - 27a^3c + 64a^2b^2 - bc^2 = 0. \end{split}}
Die Discriminante dieser quadratischen Gleichung ist
\[
\begin{aligned}
\Delta &= (11a^3b - ac^2 + 2b^2c)^2 \\
&\quad + 4(a^4+abc-b^3)(27a^3c - 64a^2b^2 + bc^2) \\
&= a^2(108a^5c - 135a^4b^2 + 90a^2bc^2 \\
&\quad - 320ab^3c + 256b^5 + c^4).
\end{aligned}
\]

Im \S.\,74 des ersten Bandes ist die Discriminante $D$ der Hauptgleichung $5^{\mathrm{ten}}$ Grades abgeleitet. Wenn man in der dort gefundenen Formel (3)
\[
a_0=1,\quad a_3=5a,\quad a_4=5b,\quad a_5=c
\]
%CONT>>


% ===== scan leaf 447, printed 426 =====
% !PRINTED 426
\weberpage{426}
\pdfbookmark[2]{§ 113. Resolventen sechsten Grades}{bm129}
%<<CONT
einsetzt, so findet sich
\[
D = 5^5(108a^5c-135a^4b^2+90a^2bc^2 - 320ab^3c+256b^5+c^4),
\]
so dass sich
\weq{16}{5^5\Delta = a^2D}
ergiebt, und man aus (15) f\"ur $\lambda$ den Ausdruck
\[
\lambda = \frac{11a^3b - ac^2 + 2b^2c + \frac{1}{25}a\sqrt{\frac{1}{5}D}}{2(a^4+abc-b^3)}
\]
erh\"alt, in dem die Quadratwurzel beide Vorzeichen haben kann.

Man sieht, dass ausser der Quadratwurzel aus der Discriminante, die rational durch die Wurzeln ausdr\"uckbar ist, noch $\sqrt5$ darin vorkommt.

Hat man $\lambda$ berechnet, so erh\"alt man aus (11) das Verh\"altniss $z_1:\mu^2$ und aus (12) die letzte Unbekannte $z$ \gesperrt{rational} durch $\lambda$ dargestellt.

\begin{center}
\S.\,113.

\gesperrt{Resolventen sechsten Grades.}
\end{center}

Um die Resolventen $6^{\mathrm{ten}}$ Grades der Ikosa\"edergleichung zu finden, gehen wir von einer der in der Ikosa\"edergruppe enthaltenen Di\"edergruppen $D_5$ aus (\S.\,111, 2.). Die Darstellung der Ikosa\"edergruppe [\S.\,58, (20)] giebt uns unmittelbar die Zerlegung in die Nebengruppen; ist
\weq{1}{Q = \Theta^r,\ \Theta^r\psi,\ r=0,1,2,3,4}
die Di\"edergruppe, von der wir ausgehen, so erhalten wir die volle Ikosa\"edergruppe
\weq{2}{P = Q + Q\chi + Q\chi\Theta + Q\chi\Theta^2 + Q\chi\Theta^3 + Q\chi\Theta^4.}
Ist dann $U_\infty$ eine Function, die durch die Substitutionen von $Q$ unge\"andert bleibt und durch
\[
\chi,\ \chi\Theta,\ \chi\Theta^2,\ \chi\Theta^3,\ \chi\Theta^4
\]
in
\[
U_0,\ U_1,\ U_2,\ U_3,\ U_4
\]
\"ubergeht, so sind die $U_\infty, U_0, U_1, U_2, U_3, U_4$ die Wurzeln einer Resolvente $6^{\mathrm{ten}}$ Grades. (Ueber die Bezeichnung $U_\infty$ vgl. \S.\,65.)


% ===== scan leaf 448, printed 427 =====
% !PRINTED 427
\weberpage{427}
Um die \gesperrt{Galois}'sche Gruppe dieser Gleichung $6^{\mathrm{ten}}$ Grades zu erhalten, haben wir den Einfluss zu untersuchen, den die Anwendung der Ikosa\"edersubstitutionen auf das System der Nebengruppen (2) hat. Es gen\"ugt dazu, die Substitutionen $\Theta$, $\psi$, $\chi$ zu betrachten. Es ist aber, wenn wir die Reihenfolge der Nebengruppen in (2) beachten, nach \S.\,58, (23), (25):
\[
\begin{aligned}
P\Theta &= Q + Q\chi\Theta + Q\chi\Theta^2 + Q\chi\Theta^3 + Q\chi\Theta^4 + Q\chi \\
P\psi &= Q + Q\chi + Q\chi\Theta^4 + Q\chi\Theta^3 + Q\chi\Theta^2 + Q\chi\Theta \\
P\chi &= Q\chi + Q + Q\chi\Theta + Q\chi\Theta^3 + Q\chi\Theta^2 + Q\chi\Theta^4,
\end{aligned}
\]
und daraus folgt, dass sich die Indices von $U$ folgendermaassen vertauschen:
\[
\begin{array}{c@{\quad}cccccc}
 & \infty, & 0, & 1, & 2, & 3, & 4 \\
\cline{2-7}
\Theta) & \infty & 1 & 2 & 3 & 4 & 0 \\
\psi) & \infty & 0 & 4 & 3 & 2 & 1 \\
\chi) & 0 & \infty & 1 & 3 & 2 & 4.
\end{array}
\]
Bezeichnen wir den Index allgemein mit $\xi$, und nehmen, wie im siebenten Abschnitte, $\xi$ nach dem Modul 5, so dass zwei nach dem Modul 5 congruente Zahlen als nicht verschieden gelten, so k\"onnen wir diese Vertauschungen so darstellen:
\[
\Theta = (\xi,\ \xi+1),\quad \psi = (\xi,\ -\xi),\quad \chi = \left(\xi,\ \frac{1}{\xi}\right)\!.
\]
\gesperrt{Daraus folgt aber, dass die Gruppe unserer Gleichung} $6^{\mathrm{ten}}$ \gesperrt{Grades mit der Congruenzgruppe} $L_5$ (\S.\,68) \gesperrt{\"ubereinstimmt.}

Um zur Bildung von Resolventen $6^{\mathrm{ten}}$ Grades zu gelangen, m\"ussen wir \"ahnlich wie bei den Resolventen $5^{\mathrm{ten}}$ Grades verfahren.

Wir haben schon im \S.\,55 die Grundformen der Di\"edergruppen kennen gelernt. Wir stellen sie jetzt in der Form dar:
\[
\varphi_1 = x_1x_2,\quad \varphi_2 = x_1^5+ix_2^5,\quad \varphi_3 = x_1^5-ix_2^5.
\]
Nehmen wir $\Theta$ und $\psi$ mit der Determinante 1, also
\[
\Theta = \begin{pmatrix} e^{\frac{\pi i}{5}}, & 0 \\ 0, & e^{-\frac{\pi i}{5}} \end{pmatrix},\quad \psi = \begin{pmatrix} 0, & 1 \\ -1, & 0 \end{pmatrix};
\]
so \"andern sich die Grundformen $\varphi_1, \varphi_2, \varphi_3$ in folgender Weise:
\[
\begin{array}{c@{\quad}ccc}
 & \varphi_1 & \varphi_2 & \varphi_3 \\
\cline{2-4}
\Theta) & \varphi_1 & -\varphi_2 & -\varphi_3 \\
\psi) & -\varphi_1 & -i\varphi_2 & i\varphi_3.
\end{array}
\]


% ===== scan leaf 449, printed 428 =====
% !PRINTED 428
\weberpage{428}
Es sind also $\varphi_1^2, \varphi_2^4, \varphi_3^4, \varphi_2\varphi_3$ absolut invariant, und wenn wir daher, wie im \S.\,111, unter $f, H, T$ die Ikosa\"ederinvarianten verstehen, so k\"onnen wir folgende Functionen als Wurzeln von Resolventen $6^{\mathrm{ten}}$ Grades einf\"uhren:
\[
\frac{\varphi_1^2H}{f^2},\quad \frac{\varphi_2^4}{H},\quad \frac{\varphi_3^4}{H},\quad \frac{\varphi_2\varphi_3H}{T}.
\]
Um die Ausdr\"ucke f\"ur die \"ubrigen Wurzeln zu erhalten, muss der Einfluss der mit der Determinante 1 genommenen Substitutionen $\Theta^r\chi$ auf die Functionen $\varphi_1, \varphi_2, \varphi_3$ untersucht werden. Wir wollen dies nur f\"ur die Function $\varphi_1$ durchf\"uhren, die zu der einfachsten Resolvente $6^{\mathrm{ten}}$ Grades f\"uhrt.

Wenn wir in $\varphi_1$ die Substitution $\chi$ in der Form \S.\,58, (26) anwenden, also $x_1, x_2$ durch
\[
\frac{x_1}{\varepsilon-\varepsilon^{-1}}+\frac{x_2}{\varepsilon^2-\varepsilon^{-2}},\quad \frac{x_1}{\varepsilon^2-\varepsilon^{-2}}-\frac{x_2}{\varepsilon-\varepsilon^{-1}}
\]
ersetzen, so geht $\varphi_1$ in
\[
\frac{-x_1^2+x_1x_2+x_2^2}{\sqrt5}
\]
\"uber. Wir setzen demnach
\weq{3}{w_\infty = 5\varphi_1^2 = 5x_1^2x_2^2,}
und erhalten durch Anwendung der Substitutionen $\chi\Theta^r$ daraus die f\"unf weiteren Formen
\weq{4}{w_r = (\varepsilon^rx_1^2+x_1x_2-\varepsilon^{-r}x_2^2)^2.}
Nun sind die symmetrischen Functionen der sechs Formen $w$ invariante Ikosa\"ederformen, und danach kann man leicht die Co\"efficienten der Gleichung bilden, deren Wurzeln die $w$ sind. Wir erhalten durch Vergleichung der Grade und je eines Co\"efficienten f\"ur die Potenzsummen der $w$
\[
S(w)=0,\ S(w^2)=0,\ S(w^3)=30f,\ S(w^4)=0,\ S(w^5)=-5H,
\]
und so das Product aller sechs $w$
\[
\Pi(w)=5f^2,
\]
und demnach ergiebt sich nach den \gesperrt{Newton}'schen Formeln (Bd.\,I, \S.\,42) f\"ur $w$ die Gleichung
\weq{5}{w^6-10fw^3+Hw+5f^2=0.}
Setzen wir also
\weq{6}{U=\frac{wH}{f^2},\ z=\frac{H^3}{f^2},}
% [sic] The source prints $z=H^3/f^2$; the homogeneous relation and the continuation indicate that $f^5$ was intended.
%CONT>>


% ===== scan leaf 450, printed 429 =====
% !PRINTED 429
\weberpage{429}
\pdfbookmark[2]{§ 114. Ueber die Lösung der Ikosaëdergleichung durch transcendente Functionen}{bm130}
%<<CONT
so folgt hieraus die Gleichung 6\textsuperscript{ten} Grades für $U$:
\weq{7}{U^6 - 10\,z\,U^3 + z^2\,U + 5\,z^2 = 0}
als \gesperrt{Resolvente} 6\textsuperscript{ten} \gesperrt{Grades} der Ikosaëdergleichung.

Setzen wir nach (3) und (4)
\weq{8}{\begin{aligned}
\sqrt{w_\infty} &= \sqrt{5}\,x_1\,x_2 \\
\sqrt{w_r} &= \varepsilon^r x_1^2 + x_1 x_2 - \varepsilon^{-r} x_2^2,
\end{aligned}}
so ergeben sich daraus die folgenden Relationen:
\weq{9}{\begin{aligned}
\sqrt{5\,w_\infty} &= \sqrt{w_0} + \sqrt{w_1} + \sqrt{w_2} + \sqrt{w_3} + \sqrt{w_4} \\
0 &= \sqrt{w_0} + \varepsilon^2\sqrt{w_1} + \varepsilon^4\sqrt{w_2} + \varepsilon\sqrt{w_3} + \varepsilon^3\sqrt{w_4} \\
0 &= \sqrt{w_0} + \varepsilon^3\sqrt{w_1} + \varepsilon\sqrt{w_2} + \varepsilon^4\sqrt{w_3} + \varepsilon^2\sqrt{w_4}
\end{aligned}}
\weq{10}{\frac{x_1\sqrt{5}}{x_2} = \frac{\sqrt{w_0} + \varepsilon^4\sqrt{w_1} + \varepsilon^3\sqrt{w_2} + \varepsilon^2\sqrt{w_3} + \varepsilon\sqrt{w_4}}{\sqrt{w_\infty}},}
und in den Gleichungen (9), (10) können, da sie homogen sind, an Stelle der $w$ auch die $U$ gesetzt werden.

Auf die Gleichungen 6\textsuperscript{ten} Grades, deren Wurzeln den Relationen (9) genügen, hat zuerst \gesperrt{Jacobi} hingewiesen; man nennt sie daher \gesperrt{Jacobi}'sche Gleichungen 6\textsuperscript{ten} Grades\footnote{Suite des notices sur les fonctions elliptiques. \gesperrt{Crelle's} Journal, Bd. III (1828); Werke, Bd. I, S. 261.}.

\begin{center}
\S.\,114.\\[4pt]
\gesperrt{Ueber die Lösung der Ikosaëdergleichung durch transcendente Functionen.}
\end{center}

Wir wollen hier, ohne auf Beweise oder nähere Ausführungen einzugehen, kurz auf die Beziehungen hinweisen, die zwischen den durchgeführten Untersuchungen und der Theorie gewisser transcendenter Functionen, besonders der Theorie der elliptischen Functionen und der Theorie der hypergeometrischen Reihen bestehen, woraus sich die Auflösung der Ikosaëdergleichung und damit aller ihrer Resolventen durch transcendente Functionen ergiebt. Dies führt aus dem Gebiete der Algebra hinaus, steht aber zu ihr etwa in derselben Beziehung, wie wenn man bei der Berechnung von Wurzelgrössen Logarithmen, oder bei den cubischen Gleichungen die Winkeltheilung anwendet.


% ===== scan leaf 451, printed 430 =====
% !PRINTED 430
\weberpage{430}
Wir müssen uns hier auf eine kurze Angabe der Resultate beschränken, und werden in der Folge keine Anwendungen davon machen, so dass diese Ausführungen, die nur einem mit der Theorie der transcendenten Functionen einigermaassen vertrauten Leser ganz verständlich sein können, ohne Störung des Zusammenhanges auch übergangen werden können.

Es möge $\omega$ eine Variable bedeuten, deren imaginärer Theil positiv ist, so dass
\weq{1}{q = e^{\pi i \omega}}
eine Grösse ist, deren absoluter Werth ein echter Bruch ist. Wir wollen die folgenden drei Functionen dieser Variablen $\omega$ betrachten:
\weq{2}{\eta(\omega) = q^{1/12}\prod_{1,\infty}^{\nu}(1 - q^{2\nu}) = q^{1/12}\sum_{-\infty,\infty}^{\nu}(-1)^\nu q^{3\nu^2+\nu} = \sum_{-\infty,\infty}^{\nu}(-1)^\nu q^{\frac{(6\nu+1)^2}{12}},}
\weq{3}{f(\omega) = q^{-\frac{1}{24}}\prod(1 + q^{2\nu-1}),}
\weq{4}{\gamma_2(\omega) = \frac{f(\omega)^{24} - 16}{f(\omega)^8}\footnote{Man vergl.: \gesperrt{Weber}, \glqq{}Elliptische Functionen und algebraische Zahlen\grqq{}. Braunschweig 1891. S. 63, 86, 149.}.}
In der Transformationstheorie der elliptischen Functionen\footnote{Ebend., S. 263, 316.} wird eine Gleichung 6\textsuperscript{ten} Grades abgeleitet, deren Wurzeln die sechs Grössen
\weq{5}{v_\infty = 5\left(\frac{\eta(5\omega)}{\eta(\omega)}\right)^2, \qquad v_\xi = \left(\frac{\eta\left(\dfrac{-24\xi+\omega}{5}\right)}{\eta(\omega)}\right)^2}
sind, wenn $\xi$ ein volles Restsystem nach dem Modul 5 durchläuft. Diese Gleichung lautet
\weq{6}{v^6 + 10v^3 - \gamma_2(\omega)v + 5 = 0,}
und hat, wie man sieht, grosse Aehnlichkeit mit der im vorigen Paragraphen betrachteten Resolvente 6\textsuperscript{ten} Grades
\weq{7}{U^6 - 10zU^3 + z^2U + 5z^2 = 0.}
Diese beiden Gleichungen gehen geradezu in einander über, wenn man
\weq{8}{z = \gamma_2(\omega)^3}
\weq{9}{U = -\gamma_2(\omega)v}
%CONT>>


% ===== scan leaf 452, printed 431 =====
% !PRINTED 431
\weberpage{431}
%<<CONT
setzt. Die Gleichung (6) hat dieselbe \gesperrt{Galois}'sche Gruppe, wie (7). Wir können nun eine beliebige unter den Wurzeln von (7) für $U_\infty$ setzen, und folglich können wir $U_\infty$ und $v_\infty$ in (9) einander entsprechen lassen. Dann können wir von den übrigen Wurzeln von (7) eine beliebige für $U_0$ nehmen, weil wir $x_1, x_2$ in \S.\,113, (4) durch $\varepsilon^r x_1$, $\varepsilon^{-r} x_2$ ersetzen können. Wir lassen also $U_0$ und $v_0$ einander entsprechen. Endlich können wir über die imaginäre fünfte Einheitswurzel $\varepsilon$ noch so verfügen, dass $U_1$ und $v_1$ sich entsprechen, und dann folgt aus der Uebereinstimmung der Gruppe, dass überhaupt $U_\xi$ und $v_\xi$ sich entsprechen, dass also
\weq{10}{U_\xi = -\gamma_2(\omega)v_\xi, \qquad \xi = \infty, 0, 1, 2, 3, 4}
ist. \gesperrt{Wenn wir also} $\omega$ \gesperrt{bei gegebenem} $z$ \gesperrt{aus der transcendenten Gleichung (8) bestimmen, so geben uns (10) und (5) die vollständige Lösung der Ikosaëderresolvente} 6\textsuperscript{ten} \gesperrt{Grades.}

Nach (2) und (5) ergiebt sich für $\xi = 0, 1, 2, 3, 4$
\weq{11}{\eta(\omega)\sqrt{v_\xi} = \sum_{-\infty,\infty}^{\nu}(-1)^\nu e^{\frac{\pi i}{60}(\omega - 24\xi)(6\nu+1)^2}.}
Der Exponent $(6\nu+1)^2$ kann nach dem Modul 5 nur den Rest $0, 1, -1$ geben. Lässt man also $\nu_1, \nu_2$ alle Zahlen durchlaufen, die den Bedingungen
\[
\nu_1 \equiv 1, 2, \qquad \nu_2 \equiv 0, 3 \pmod{5}
\]
genügen, so können wir den Ausdruck (11) so zerlegen:
\weq{12}{\eta(\omega)\sqrt{v_\xi} = \eta(5\omega) + \varepsilon^\xi \sum^{\nu_1}(-1)^{\nu_1} e^{\frac{\pi i}{60}\omega(6\nu_1+1)^2} + \varepsilon^{-\xi}\sum^{\nu_2}(-1)^{\nu_2} e^{\frac{\pi i}{60}\omega(6\nu_2+1)^2},}
wenn $\varepsilon = e^{\frac{2\pi i}{5}}$ gesetzt ist. Hieraus erkennt man, dass die Grössen $\sqrt{v_\xi}$ die Relationen (9) des vorigen Paragraphen befriedigen, und die Formel (10), \S.\,113 ergiebt für die Wurzel $x = x_1 : x_2$ der Ikosaëdergleichung den Ausdruck:
\weq{13}{\eta(5\omega)x = \sum^{\nu_1}(-1)^{\nu_1} e^{\frac{\pi i}{60}\omega(6\nu_1+1)^2}.}
Die übrigen Wurzeln erhält man, wenn man auf $x$ die Ikosaëdersubstitutionen anwendet. Man kann sie aber auch da-
%CONT>>


% ===== scan leaf 453, printed 432 =====
% !PRINTED 432
\weberpage{432}
%<<CONT
durch bilden, dass man die Variable $\omega$ durch einen äquivalenten Werth ersetzt (Bd. I, \S.\,120).

\begin{center}
\rule{0.25\linewidth}{0.4pt}
\end{center}

Wir wollen hier endlich auch noch die Lösung der Ikosaëdergleichung durch hypergeometrische Reihen kurz erwähnen. Nach \gesperrt{Gauss} verstehen wir unter
\[
F(\alpha, \beta, \gamma, \xi)
\]
die unendliche Reihe
\[
1 + \frac{\alpha\beta}{1.\gamma}\xi + \frac{\alpha(\alpha+1)\beta(\beta+1)}{1.2\,\gamma(\gamma+1)}\xi^2 + \frac{\alpha(\alpha+1)(\alpha+2)\beta(\beta+1)(\beta+2)}{1.2.3\,\gamma(\gamma+1)(\gamma+2)}\xi^2 + \cdots
\]
% [sic] The source clearly repeats $\xi^2$ in the third term; the hypergeometric series requires $\xi^3$.
Setzen wir darin
\[
z = 1728\,\xi,
\]
so erhält die Wurzel der Ikosaëdergleichung
\[
H^3 - zf^5 = 0
\]
den Ausdruck
\[
x = \xi^{\frac{1}{3}}\,\frac{F\left(\dfrac{11}{60}, \dfrac{19}{60}, \dfrac{4}{3}, \xi\right)}{F\left(\dfrac{11}{60}, \dfrac{-1}{60}, \dfrac{1}{3}, \xi\right)}.
\]
Diese Reihen convergiren nur, so lange der absolute Werth von $\xi$ ein echter Bruch ist, können aber für andere Werthe von $\xi$ durch ähnliche Reihen ersetzt werden.\footnote{\gesperrt{H. A. Schwarz}, \glqq{}Ueber diejenigen Fälle, in denen die \gesperrt{Gauss}'sche Reihe $F(\alpha,\beta,\gamma,x)$ eine algebraische Function ihres vierten Elementes ist\grqq{}. Crelle's Journal, Bd. 75, 1872 (gesammelte mathematische Werke, Bd. II). \gesperrt{F. Klein}, Vorlesungen über das Ikosaëder, Abschnitt I, Cap. III.}

\begin{center}
\rule{0.25\linewidth}{0.4pt}
\end{center}


% ===== scan leaf 454, printed 433 =====
\weberpage{433}% page number not printed on this page
\pdfbookmark[1]{Vierzehnter Abschnitt. Gruppen linearer ternärer Substitutionen.}{bm131}
\pdfbookmark[2]{§ 115. Ternäre lineare Substitutionsgruppe vom 168sten Grade}{bm132}
% !PRINTED none
% Chapter-opening page: no running head/page number is printed at the top of this
% page; instead the page carries the Abschnitt title block, then the §.115 heading
% and the start of the section text. (Weber suppresses the folio number on
% Abschnitt-opening pages.)

\begin{center}
\gesperrt{Vierzehnter Abschnitt.}

\gesperrt{Gruppen linearer ternärer Substitutionen.}

\rule{0.3\linewidth}{0.4pt}

\S.\,115.

\gesperrt{Ternäre lineare Substitutionsgruppe vom 168\textsuperscript{sten} Grade.}
\end{center}

Die Frage nach der Gesammtheit aller möglichen endlichen
Gruppen linearer Substitutionen von mehreren, insbesondere
derer von drei und vier Dimensionen, ist von \gesperrt{C. Jordan} behandelt,
der, ähnlich wie es für die binären Substitutionen geschehen,
eine endliche Anzahl von Typen solcher Gruppen aufgestellt hat\footnote{\gesperrt{C. Jordan}, Mémoire sur les équations différentielles linéaires à
intégrale algébrique. Crelle's Journal für Mathematik, Bd. 84 (1878). Sur
la détermination des groupes d'ordre fini etc. Atti della R. Accademia di
Napoli, Vol. VIII (1879). \gesperrt{Valentiner}, Kj\"ob. Skrift (6) V (1889).}. Diese allgemeinen Untersuchungen können wir
hier nicht verfolgen, wir begnügen uns, ein Beispiel einer ternären
Gruppe ausführlicher zu betrachten.

In dem Abschnitte über die Congruenzgruppen (\S.\,66) haben
wir eine Reihe \gesperrt{einfacher} Gruppen kennen gelernt, von denen
sich die erste, vom Grade 60, als isomorph mit der Ikosaëdergruppe
erwiesen hat. Wir wollen nun die nächste dieser einfachen
Gruppen vom Grade 168 betrachten, die wir mit $L_7$ bezeichnet
haben, und versuchen, eine mit dieser isomorphe, ternäre, lineare
Substitutionsgruppe zu construiren.

Dazu führt uns das Theorem \S.\,72, I., wenn wir vier Elemente
$\tau, \chi, \omega, \Theta$ aufsuchen, die bei der Composition den Bedingungen
dieses Theorems genügen.

Wir wollen zunächst eine Substitution $\tau$ in der Normalform
annehmen
\weq{1}{\tau = \begin{pmatrix} \varepsilon_1, & 0, & 0 \\ 0, & \varepsilon_2, & 0 \\ 0, & 0, & \varepsilon_3 \end{pmatrix},}
%CONT>>


% ===== scan leaf 455, printed 434 =====
% !PRINTED 434
\weberpage{434}
%<<CONT
und darin müssen, da $\tau$ vom 7\textsuperscript{ten} Grade sein soll, die $\varepsilon_1, \varepsilon_2, \varepsilon_3$
siebente Einheitswurzeln sein, die der Bedingung
\weq{2}{\varepsilon_1 \varepsilon_2 \varepsilon_3 = 1}
genügen.

Nun wollen wir die Substitution $\chi$ so zu bestimmen suchen,
dass die Bedingung $\chi \tau = \tau^4 \chi$ oder
\weq{3}{\tau \chi = \chi \tau^2}
erfüllt ist. Nehmen wir $\chi$ in der Form
\[
\chi = \begin{pmatrix} \alpha_1, & \alpha_2, & \alpha_3 \\ \beta_1, & \beta_2, & \beta_3 \\ \gamma_1, & \gamma_2, & \gamma_3 \end{pmatrix}
\]
an, so ergiebt die Bedingung (3)
\[
\begin{pmatrix} \alpha_1 \varepsilon_1, & \alpha_2 \varepsilon_1, & \alpha_3 \varepsilon_1 \\ \beta_1 \varepsilon_2, & \beta_2 \varepsilon_2, & \beta_3 \varepsilon_2 \\ \gamma_1 \varepsilon_3, & \gamma_2 \varepsilon_3, & \gamma_3 \varepsilon_3 \end{pmatrix}
=
\begin{pmatrix} \alpha_1 \varepsilon_1^2, & \alpha_2 \varepsilon_2^2, & \alpha_3 \varepsilon_3^2 \\ \beta_1 \varepsilon_1^2, & \beta_2 \varepsilon_2^2, & \beta_3 \varepsilon_3^2 \\ \gamma_1 \varepsilon_1^2, & \gamma_2 \varepsilon_2^2, & \gamma_3 \varepsilon_3^2 \end{pmatrix}.
\]

Da nun $\alpha_1, \beta_1, \gamma_1$ nicht alle drei verschwinden können, so
muss $\varepsilon_1^2$ mit einer der drei Grössen $\varepsilon_1, \varepsilon_2, \varepsilon_3$ übereinstimmen.
Wäre $\varepsilon_1 = \varepsilon_1^2$, also $\varepsilon_1 = 1$, und folglich $\varepsilon_3 = \varepsilon_2^{-1}$, so könnten
weder $\varepsilon_2$ noch $\varepsilon_3 = 1$ sein, weil sonst $\tau$ die identische Substitution
wäre. Es müsste also $\beta_1 = \gamma_1 = \alpha_2 = \alpha_3 = 0$ und $\varepsilon_2^2 = \varepsilon_3$,
$\varepsilon_2^3 = 1$ sein, was unmöglich ist. Es bleiben daher nach (2) noch
zwei mögliche Annahmen übrig:
\[
\varepsilon_1 = \varepsilon, \quad \varepsilon_2 = \varepsilon^2, \quad \varepsilon_3 = \varepsilon^4,
\]
\[
\varepsilon_1 = \varepsilon, \quad \varepsilon_2 = \varepsilon^4, \quad \varepsilon_3 = \varepsilon^2,
\]
worin $\varepsilon$ eine imaginäre siebente Einheitswurzel ist, und diese
beiden Annahmen führen zu zwei \gesperrt{verschiedenen} Gruppen, die
ganz gleichartig gebaut, übrigens auch isomorph sind. Wir verfolgen hier die erste Annahme weiter, setzen also
\weq{4}{\tau = \begin{pmatrix} \varepsilon, & 0, & 0 \\ 0, & \varepsilon^2, & 0 \\ 0, & 0, & \varepsilon^4 \end{pmatrix}, \quad \chi = \begin{pmatrix} 0, & 0, & 1 \\ 1, & 0, & 0 \\ 0, & 1, & 0 \end{pmatrix}, \quad \chi^2 = \begin{pmatrix} 0, & 1, & 0 \\ 0, & 0, & 1 \\ 1, & 0, & 0 \end{pmatrix}, \quad \chi^3 = 1.}

Dass wir die nicht verschwindenden Grössen $\alpha_3, \beta_1, \gamma_2$ gleich 1
gesetzt haben, ist keine Beschränkung, da wir jede andere Annahme durch eine Transformation der ganzen Gruppe darauf
zurückführen können.

Die Bedeutung von $\chi$ in dieser Form ist die einer cyklischen
Permutation der Variablen. Die Substitutionen $\tau, \chi$ erzeugen
zusammen eine Gruppe $\tau^\varrho \chi^\lambda$ vom 21\textsuperscript{sten} Grade.


% ===== scan leaf 456, printed 435 =====
% !PRINTED 435
\weberpage{435}
Um nun $\omega$ zu bestimmen, bedienen wir uns der Relation
\[
\omega \chi = \chi^2 \omega,
\]
und erhalten, wenn
\[
\omega = \begin{pmatrix} \alpha_1, & \alpha_2, & \alpha_3 \\ \beta_1, & \beta_2, & \beta_3 \\ \gamma_1, & \gamma_2, & \gamma_3 \end{pmatrix}
\]
gesetzt wird,
\[
\begin{pmatrix} \alpha_2, & \alpha_3, & \alpha_1 \\ \beta_2, & \beta_3, & \beta_1 \\ \gamma_2, & \gamma_3, & \gamma_1 \end{pmatrix}
=
\begin{pmatrix} \beta_1, & \beta_2, & \beta_3 \\ \gamma_1, & \gamma_2, & \gamma_3 \\ \alpha_1, & \alpha_2, & \alpha_3 \end{pmatrix},
\]
also
\[
\begin{aligned}
\alpha_1 &= \beta_3 = \gamma_2 = \alpha, \\
\alpha_2 &= \beta_1 = \gamma_3 = \beta, \\
\alpha_3 &= \beta_2 = \gamma_1 = \gamma,
\end{aligned}
\]
und $\omega$ erhält also die Form
\weq{5}{\omega = \begin{pmatrix} \alpha, & \beta, & \gamma \\ \beta, & \gamma, & \alpha \\ \gamma, & \alpha, & \beta \end{pmatrix}.}

Drücken wir die Bedingung aus, dass $\omega$ vom 2\textsuperscript{ten} Grade
sein soll, so ergeben sich die Relationen
\weq{6}{\begin{aligned} \alpha^2 + \beta^2 + \gamma^2 &= 1 \\ \beta\gamma + \gamma\alpha + \alpha\beta &= 0. \end{aligned}}

Aus (6) ergiebt sich $(\alpha + \beta + \gamma)^2 = 1$, und wenn man die
Determinante von $\omega$ gleich 1 setzt, so folgt mit Benutzung von (6)
$(\alpha + \beta + \gamma)^3 = -1$, folglich
\weq{7}{\alpha + \beta + \gamma = -1.}

Aus (6) und (7) folgern wir noch die Relationen:
\weq{8}{\alpha = \beta\gamma - \alpha^2, \quad \beta = \gamma\alpha - \beta^2, \quad \gamma = \alpha\beta - \gamma^2}
\weq{9}{\alpha^3 + \beta^3 + \gamma^3 - 3\alpha\beta\gamma = -1,}
von denen die letztere, die man aus
\[
(\alpha + \beta + \gamma)^3 - 3(\beta\gamma + \gamma\alpha + \alpha\beta)(\alpha + \beta + \gamma) = -1
\]
erhält, den negativen Werth der Determinante von $\omega$ darstellt.

Endlich bilden wir noch nach den Relationen [\S.\,72, (15)]
\[
\Theta = \chi \tau^3 \omega \tau^4, \quad \Theta^3 = \chi \tau^6 \omega \tau^2 \text{:}
\]
\weq{10}{\Theta = \begin{pmatrix} \gamma \varepsilon^2, & \alpha \varepsilon^6, & \beta \\ \alpha, & \beta \varepsilon^4, & \gamma \varepsilon^5 \\ \beta \varepsilon^3, & \gamma, & \alpha \varepsilon \end{pmatrix}, \quad \Theta^{-1} = \begin{pmatrix} \gamma \varepsilon^5, & \alpha, & \beta \varepsilon^4 \\ \alpha \varepsilon, & \beta \varepsilon^3, & \gamma \\ \beta, & \gamma \varepsilon^2, & \alpha \varepsilon^6 \end{pmatrix}.}


% ===== scan leaf 457, printed 436 =====
% !PRINTED 436
\weberpage{436}
Wenn wir nun die Bedingung aufsuchen, dass $\Theta^4 = 1$ wird,
so haben wir $\Theta^2 = \Theta^{-2}$ zu setzen. Bilden wir also $\Theta^2$ und $\Theta^{-2}$
nach (10), so kann man zuerst versuchen, die Uebereinstimmung
herzustellen, indem man eine der drei Zahlen $\alpha, \beta, \gamma$, etwa $\gamma = 0$,
setzt; dann muss aber nach (6) und (7) noch eine zweite, etwa
$\beta = 0$, und die dritte $\alpha = -1$ sein. Dann wird aber $\Theta^2$ nicht
mit $\Theta^{-2}$ übereinstimmend. Also sind alle drei Grössen $\alpha, \beta, \gamma$
von Null verschieden. Setzen wir nun das zweite und dritte
Glied der ersten Zeile in den aus (10) abgeleiteten Substitutionen
$\Theta^2$ und $\Theta^{-2}$ einander gleich, so ergiebt sich, wenn die Factoren
$\gamma$ und $\alpha$ abgeworfen werden:
\[
\alpha(\varepsilon - \varepsilon^{-2}) = \beta(\varepsilon^{-1} - 1), \quad \gamma(\varepsilon^4 - 1) = \beta(\varepsilon^3 - \varepsilon)
\]
oder
\[
\alpha(\varepsilon^2 - \varepsilon^{-2}) = \beta(\varepsilon^4 - \varepsilon^{-4}), \quad \gamma(\varepsilon^2 - \varepsilon^{-2}) = \beta(\varepsilon - \varepsilon^{-1}),
\]
und hieraus, wenn $h$ einen unbestimmten Factor bedeutet:
\weq{11}{\begin{aligned} \alpha &= h(\varepsilon^4 - \varepsilon^{-4}) \\ \beta &= h(\varepsilon^2 - \varepsilon^{-2}) \\ \gamma &= h(\varepsilon - \varepsilon^{-1}). \end{aligned}}

Der Factor $h$ wird nach (7) aus der Gleichung
\weq{12}{-1 = h(\varepsilon + \varepsilon^2 + \varepsilon^4 - \varepsilon^{-1} - \varepsilon^{-2} - \varepsilon^{-4})}
bestimmt, und es ergiebt sich nach Bd. I, \S.\,171
\weq{13}{h = \frac{i}{\sqrt{7}} = \frac{-1}{\sqrt{-7}}.}

Das Vorzeichen von $\sqrt{7}$ hängt von der Wahl von $\varepsilon$ ab und
ist positiv, wenn z.\,B.
\[
\varepsilon = e^{\frac{2\pi i}{7}}
\]
genommen wird. Dann erhält man für $\alpha, \beta, \gamma$
\[
\alpha = \frac{-2\sin\dfrac{8\pi}{7}}{\sqrt{7}}, \quad \beta = \frac{-2\sin\dfrac{4\pi}{7}}{\sqrt{7}}, \quad \gamma = \frac{-2\sin\dfrac{2\pi}{7}}{\sqrt{7}}.
\]

Aus (11) und (12) ergeben sich noch die Formeln
\weq{14}{\begin{aligned}
\alpha\varepsilon + \beta\varepsilon^4 + \gamma\varepsilon^2 &= 1 \\
\alpha\varepsilon^{-1} + \beta\varepsilon^{-4} + \gamma\varepsilon^{-2} &= 1, \\
\alpha + \beta\varepsilon^{-1} + \gamma\varepsilon^2 &= 0 \\
\alpha + \beta\varepsilon + \gamma\varepsilon^{-2} &= 0; \\
\alpha\beta\gamma &= \tfrac{1}{7},
\end{aligned}}
%CONT>>


% ===== scan leaf 458, printed 437 =====
% !PRINTED 437
\weberpage{437}
%<<CONT
und $\alpha$, $\beta$, $\gamma$ sind die Wurzeln der cubischen Gleichung
\[ \alpha^3 + \alpha^2 - \frac{1}{7} = 0. \]
Man kann nun nach \S.\,72, (15), (12), (13)
\weq{15}{\Theta^2 = \tau^4\,\omega\,\tau^5\,\chi}
setzen, und daraus erhält man leicht nach (4) und (5)
\weq{16}{\Theta^2 = \begin{pmatrix} \beta, & \gamma\varepsilon^3, & \alpha\varepsilon^2 \\ \gamma\varepsilon^{-3}, & \alpha, & \beta\varepsilon^{-1} \\ \alpha\varepsilon^{-2}, & \beta\varepsilon, & \gamma \end{pmatrix}.}

Vergleicht man dies mit dem Resultate, was man erhält, wenn man aus (10) direct $\Theta^2$ bildet, so ergiebt sich:
\weq{17}{\begin{aligned}
\alpha &= \alpha^2\varepsilon^{-1} + \beta^2\varepsilon + \gamma^2\varepsilon^{-2}, & \alpha &= \beta\gamma + \gamma\alpha\varepsilon^2 + \alpha\beta\varepsilon^{-1} \\
\beta &= \alpha^2\varepsilon^{-1} + \beta^2\varepsilon^3 + \gamma^2\varepsilon^{-3}, & \beta &= \beta\gamma\varepsilon^3 + \gamma\alpha + \alpha\beta\varepsilon \\
\gamma &= \alpha^2\varepsilon^2 + \beta^2\varepsilon^3 + \gamma^2\varepsilon^{-2}, & \gamma &= \beta\gamma\varepsilon^{-3} + \gamma\alpha\varepsilon^{-2} + \alpha\beta,
\end{aligned}}

und die Vergleichung mit dem aus (10) gebildeten $\Theta^{-2}$ ergiebt ein ganz ähnliches Formelsystem, das aus (17) hervorgeht, wenn $\varepsilon$ mit $\varepsilon^{-1}$ vertauscht wird, nämlich:
\[ \begin{aligned}
\alpha &= \alpha^2\varepsilon + \beta^2\varepsilon^{-1} + \gamma^2\varepsilon^2, & \alpha &= \beta\gamma + \gamma\alpha\varepsilon^{-2} + \alpha\beta\varepsilon \\
\beta &= \alpha^2\varepsilon + \beta^2\varepsilon^{-3} + \gamma^2\varepsilon^3, & \beta &= \beta\gamma\varepsilon^{-3} + \gamma\alpha + \alpha\beta\varepsilon^{-1} \\
\gamma &= \alpha^2\varepsilon^{-2} + \beta^2\varepsilon^{-3} + \gamma^2\varepsilon^2, & \gamma &= \beta\gamma\varepsilon^3 + \gamma\alpha\varepsilon^2 + \alpha\beta.
\end{aligned} \]

Damit aber hierin kein Cirkelschluss liege, ist zu zeigen, dass die Relationen (17) in Folge der Bestimmungen (11), (12) wirklich erfüllt sind, denn dann erst ist das Bestehen der Relation (15) und die Gleichheit von $\Theta^2$ mit $\Theta^{-2}$ nachgewiesen. Es genügt dazu, zwei dieser Relationen, etwa
\[ \begin{aligned}
\alpha &= \alpha^2\varepsilon^{-1} + \beta^2\varepsilon + \gamma^2\varepsilon^{-2} \\
\alpha &= \beta\gamma + \gamma\alpha\varepsilon^2 + \alpha\beta\varepsilon^{-1},
\end{aligned} \]
abzuleiten; denn hat man diese nachgewiesen, so kann man darin $\varepsilon$ mit $\varepsilon^{-1}$ vertauschen, wodurch $\alpha$, $\beta$, $\gamma$ nicht geändert werden, und sodann $\varepsilon$ mit $\varepsilon^2$ und $\varepsilon^4$, wodurch $\alpha$, $\beta$, $\gamma$ in $\gamma$, $\alpha$, $\beta$ und in $\beta$, $\gamma$, $\alpha$ übergehen; und daher erhält man das ganze Formelsystem (17).

Diese beiden Relationen ergeben sich aber durch eine einfache Rechnung nach (11), (12) in der Form:
\[ \begin{aligned}
&-(\varepsilon + \varepsilon^2 + \varepsilon^4 - \varepsilon^{-1} - \varepsilon^{-2} - \varepsilon^{-4})(\varepsilon^4 - \varepsilon^{-4}) \\
&= \varepsilon^{-1}(\varepsilon^4 - \varepsilon^{-4})^2 + \varepsilon(\varepsilon^2 - \varepsilon^{-2})^2 + \varepsilon^{-2}(\varepsilon - \varepsilon^{-1})^2 \\
&= (\varepsilon - \varepsilon^{-1})(\varepsilon^2 - \varepsilon^{-2}) + \varepsilon^2(\varepsilon - \varepsilon^{-1})(\varepsilon^4 - \varepsilon^{-4}) \\
&\phantom{=} + \varepsilon^{-1}(\varepsilon^2 - \varepsilon^{-2})(\varepsilon^4 - \varepsilon^{-4}).
\end{aligned} \]


% ===== scan leaf 459, printed 438 =====
% !PRINTED 438
\weberpage{438}
\pdfbookmark[2]{§ 116. Pole und Axen der ternären Gruppen}{bm133}
Aus diesen Formeln ergiebt sich leicht nach (5), (8), (14):
\weq{18}{\begin{aligned}
\omega\Theta &= \begin{pmatrix} \gamma, & \alpha\varepsilon^{-2}, & \beta\varepsilon \\ \alpha\varepsilon^2, & \beta, & \gamma\varepsilon^3 \\ \beta\varepsilon^{-1}, & \gamma\varepsilon^{-3}, & \alpha \end{pmatrix}, \quad
\omega\Theta^2 = \begin{pmatrix} \gamma, & \alpha\varepsilon, & \beta\varepsilon^3 \\ \alpha\varepsilon^{-1}, & \beta, & \gamma\varepsilon^2 \\ \beta\varepsilon^{-3}, & \gamma\varepsilon^{-2}, & \alpha \end{pmatrix}, \\
\Theta\omega &= \omega\Theta^3 = \begin{pmatrix} \alpha, & \beta\varepsilon^{-1}, & \gamma\varepsilon^{-3} \\ \beta\varepsilon, & \gamma, & \alpha\varepsilon^{-2} \\ \gamma\varepsilon^3, & \alpha\varepsilon^2, & \beta \end{pmatrix}.
\end{aligned}}

Hiernach ist es nun sehr einfach, die charakteristischen Relationen des Theorems I., \S.\,72 für unsere Gruppe durch wirkliche Ausrechnung zu bestätigen, und damit ist also die ternäre Gruppe 168\textsuperscript{sten} Grades hergestellt\footnote{Vergl. F. \gesperrt{Klein}, \glqq Ueber die Transformation siebenter Ordnung der elliptischen Functionen\grqq{}. \glqq Ueber die Auflösung gewisser Gleichungen vom 7\textsuperscript{ten} und 8\textsuperscript{ten} Grade\grqq{}. Mathem. Annalen, Bd. 14 u. 15. \gesperrt{Klein-Fricke}, Vorlesungen über Modul-Functionen, Bd. I, dritter Abschnitt, Capitel VII. \gesperrt{Gordan}, \glqq Ueber Gleichungen 7\textsuperscript{ten} Grades mit einer Gruppe von 168 Substitutionen\grqq{}. Mathem. Annalen, Bd. 20, 25.}).

\begin{center}
\S.\,116.\\[2pt]
\gesperrt{Pole und Axen der ternären Gruppen.}
\end{center}

Wir wollen uns jetzt der Kürze halber einer geometrischen Ausdrucksweise bedienen, indem wir die Variablen $x_1, x_2, x_3$ als Dreieckscoordinaten eines Punktes $x$ in einer Ebene betrachten, obwohl auch imaginäre Werthe von $x_1, x_2, x_3$ zulässig sein sollen.

Bedeutet $A$ eine ternäre lineare Substitution mit der Determinante 1, und ist
\weq{1}{(y) = A(x),}
so sind $y_1, y_2, y_3$ die Coordinaten eines zweiten Punktes $(y)$, bezogen auf dasselbe Coordinatensystem, der aus $(x)$ durch die Substitution $A$ abgeleitet ist.

Wenn $x$ eine gerade Linie durchläuft, so durchläuft auch $y$ eine gerade Linie, und wenn die Gleichungen dieser beiden Linien
\weq{2}{u_1x_1 + u_2x_2 + u_3x_3 = 0, \quad v_1y_1 + v_2y_2 + v_3y_3 = 0}
sind, so ergiebt sich durch Einsetzen von (1) in (2), dass $u_1, u_2, u_3$ mit $v_1, v_2, v_3$ durch die \gesperrt{transponirte Substitution} von $A$:
\[ (u) = A_1(v) \]
%CONT>>


% ===== scan leaf 460, printed 439 =====
% !PRINTED 439
\weberpage{439}
%<<CONT
zusammenhängen (\S.\,37, 9.). Es wird also durch $A$ nicht nur aus jedem Punkte ein Punkt, sondern auch aus jeder geraden Linie eine gerade Linie abgeleitet.

Die durch (1) ausgedrückte Beziehung der Punkte $y$ zu den Punkten $x$ kann als eine \gesperrt{Abbildung der Ebene in sich selbst} bezeichnet werden, wobei jedem Punkte ein bestimmter anderer Punkt und jeder geraden Linie eine gerade Linie entspricht. Eine solche Abbildung heisst in der Geometrie auch eine Collineation. Der Punkt $y$ heisst das \gesperrt{Bild} des Punktes $x$ und die gerade Linie $v$ das \gesperrt{Bild} der geraden Linie $u$. Hat man die Bilder zweier Punkte, so ist die Verbindungslinie dieser Bilder das Bild der Verbindungslinie der beiden Originalpunkte.

Bedeutet $S$ eine zweite ternäre lineare Substitution und
\[ A' = S^{-1}AS \]
die aus $A$ durch $S$ transformirte Substitution, so ist
\[ (y') = A'(x') \]
gleichbedeutend mit
\[ (y) = A(x), \]
wenn
\[ (y) = S(y'), \quad (x) = S(x') \]
gesetzt wird. Statt nun durch $S$ eine Abbildung zu definiren, kann man auch eine \gesperrt{Coordinatentransformation} darunter verstehen, indem man unter $x_1', x_2', x_3'$ die Coordinaten des Punktes $(x)$ in einem neuen, durch $S$ bestimmten Coordinatensysteme versteht. Dann wird der Zusammenhang zwischen den beiden Punkten $x$, $y$ ebenso wie durch $A$ auch durch die transformirte Substitution $A'$ ausgedrückt, und die Transformation der Substitutionen ist also gleichbedeutend mit der Transformation des Coordinatensystems. Die Composition der transformirten Substitutionen geschieht, wie wir schon früher gesehen haben, ebenso wie die der ursprünglichen, und es entsteht also durch Transformation aus jeder Gruppe eine isomorphe Gruppe.

Von grosser Wichtigkeit ist nun die Untersuchung der Punkte, die bei der Abbildung durch eine Substitution $A$ ungeändert bleiben.

Nehmen wir $A$ in der Form an:
\[ A = \begin{pmatrix} \alpha, & \beta, & \gamma \\ \alpha', & \beta', & \gamma' \\ \alpha'', & \beta'', & \gamma'' \end{pmatrix}, \]
%CONT>>


% ===== scan leaf 461, printed 440 =====
% !PRINTED 440
\weberpage{440}
%<<CONT
so erhalten wir als Bedingung dafür, dass der Punkt $y$ mit dem Punkte $x$ zusammenfällt, da ein Punkt durch die \gesperrt{Verhältnisse} seiner Coordinaten eindeutig bestimmt ist, die, dass für einen passend bestimmten Co\"efficienten $\lambda$:
\weq{3}{\begin{aligned}
\lambda x_1 &= \alpha\,x_1 + \beta\,x_2 + \gamma\,x_3 \\
\lambda x_2 &= \alpha'x_1 + \beta'x_2 + \gamma'x_3 \\
\lambda x_3 &= \alpha''x_1 + \beta''x_2 + \gamma''x_3,
\end{aligned}}
woraus man durch Elimination von $x_1, x_2, x_3$ die \gesperrt{cubische} Gleichung für $\lambda$ erhält:
\weq{4}{\Lambda = \begin{vmatrix} \alpha - \lambda, & \beta, & \gamma \\ \alpha', & \beta' - \lambda, & \gamma' \\ \alpha'', & \beta'', & \gamma'' - \lambda \end{vmatrix} = 0.}

Es giebt daher im Allgemeinen \gesperrt{drei Punkte}, die ihre eigenen Bilder sind, und diese wollen wir (wie bei den binären Substitutionen) die \gesperrt{Pole der Substitution $A$} nennen. Die drei Verbindungslinien dieser Pole sind dann gleichfalls ihre eigenen Bilder, jedoch so, dass auf jeder dieser drei Geraden nur zwei Punkte liegen, die auf sich selbst abgebildet sind. Die Bilder der übrigen Punkte sind in der Linie verschoben. Diese Linien wollen wir die \gesperrt{Axen der Substitution $A$} nennen.

Es können nun aber besondere Fälle eintreten, die hervorzuheben sind. Es können zwei oder selbst alle drei Pole in einen Punkt zusammenfallen, wie das Beispiel
\[ \begin{pmatrix} \alpha & 0 & 0 \\ \alpha' & \alpha & 0 \\ \alpha'' & \beta'' & \alpha^{-2} \end{pmatrix} \]
zeigt. Hier fallen, wenn $\alpha'$, $\alpha''$, $\beta''$ von Null verschieden sind, zwei, oder wenn $\alpha = 1$ ist, alle drei Pole in einen Punkt zusammen.

Wichtiger noch ist ein anderer besonderer Fall, nämlich der, dass unendlich viele Punkte ihre eigenen Bilder sind. Wenn wir von dem Falle absehen, dass alle Punkte ihre eigenen Bilder sind, der nur bei der Multiplication vorkommt, so müssen die Punkte, die ihre eigenen Bilder sind, wenn ihrer unendlich viele sind, auf einer geraden Linie liegen; denn es kann dieser Fall nur dann eintreten, wenn für eine Wurzel von (4) die drei Gleichungen (3) aus einer von ihnen folgen, oder, was dasselbe ist, wenn mit $\Lambda$ zugleich die sämmtlichen ersten Unterdeterminanten
%CONT>>


% ===== scan leaf 462, printed 441 =====
% !PRINTED 441
\weberpage{441}
%<<CONT
verschwinden. Eine solche gerade Linie wollen wir eine \gesperrt{Hauptaxe der Substitution} $A$ nennen. Aber nur gewisse besondere Substitutionen besitzen Hauptaxen. Eine Substitution mit einer Hauptaxe hat ausser dieser noch \gesperrt{einen Pol}, der in besonderen F\"allen auch in die Hauptaxe hineinfallen kann.

Legen wir, um diese Verh\"altnisse zu \"ubersehen, die Hauptaxe in die Linie $x_1=0$, so muss, wenn $x_1=0$ ist, $y_1=0$, $y_2:y_3=x_2:x_3$ sein. Daraus ergiebt sich
\[
\begin{aligned}
y_1 &= \alpha\, x_1\\
y_2 &= \alpha'\, x_1 + \beta\, x_2\\
y_3 &= \alpha''\, x_1 + \beta\, x_3.
\end{aligned}
\]
Den ausser der Hauptaxe existirenden Pol erhalten wir, wenn wir $y_1=\alpha\,x_1$, $y_2=\alpha\,x_2$, $y_3=\alpha\,x_3$ setzen; dann ergiebt sich
\[
(\beta-\alpha)\,x_2+\alpha'\,x_1=0,\quad (\beta-\alpha)\,x_3+\alpha''\,x_1=0,
\]
und dieser Punkt wird dann und nur dann in die Linie $x_1=0$ fallen, wenn $\beta=\alpha$ ist. Dieser Fall kann bei endlichen Gruppen nicht eintreten (ausser bei der Multiplication).

Wenn ein von der Hauptaxe verschiedener Pol existirt, so k\"onnen wir diesen in die Ecke $x_2=0$, $x_3=0$ legen, und die Substitution erh\"alt die Form
\[
y_1=\alpha\,x_1,\quad y_2=\beta\,x_2,\quad y_3=\beta\,x_3,
\]
wo $\alpha$ von $\beta$ verschieden ist.

Eine gerade Linie $u_1x_1+u_2x_2+u_3x_3=0$ ist dann und nur dann ihr eigenes Bild, wenn entweder $u_2=u_3=0$ oder $u_1=0$ ist.

\gesperrt{Wenn also eine Substitution einen Pol und eine Hauptaxe hat, so bleiben ausser dieser alle geraden Linien und nur die unge\"andert, die durch den Pol gehen.}

Es ist noch zu bemerken, dass, wenn drei getrennte Pole vorhanden sind, diese nicht in eine gerade Linie fallen k\"onnen, ohne dass die Linie zur Hauptaxe wird. Denn wenn eine gerade Linie ihr eigenes Bild ist, so k\"onnen, wenn sich die Linie nicht Punkt f\"ur Punkt selbst entspricht, nur zwei Punkte auf ihr liegen, die ihr eigenes Bild sind.

Fassen wir dies Alles zusammen, so finden wir folgende Arten von tern\"aren linearen Substitutionen, wobei zusammenfallende Pole nur einmal mitgez\"ahlt sind:
%CONT>>


% ===== scan leaf 463, printed 442 =====
% !PRINTED 442
\weberpage{442}
\pdfbookmark[2]{§ 117. Anwendung auf die Gruppe G168. Siebenzählige Pole}{bm134}
%<<CONT

\begin{tabular}{@{}l l l l@{}}
1. & Substitutionen mit drei Polen, & & \\
2. & \textquotedbl & \textquotedbl & zwei Polen,\\
3. & \textquotedbl & \textquotedbl & einem Pol,\\
4. & \textquotedbl & \textquotedbl & einer Hauptaxe und einem Pol,\\
5. & \textquotedbl & \textquotedbl & einer Hauptaxe.\\
\end{tabular}

Als letzter Fall w\"urde noch die Multiplication aufzuz\"ahlen sein, f\"ur den jeder beliebige Punkt sein eigenes Bild ist.

Die Art der Substitution bleibt erhalten bei jeder Transformation.

Wenn ein Punkt unge\"andert bleibt durch eine Substitution $A$, so bleibt er auch bei jeder Wiederholung von $A$, also bei $A^2, A^3,\dots$, unge\"andert. Die Pole von $A$ finden sich daher immer unter den Punkten, die bei der Abbildung durch $A^2$ ihre eigenen Bilder sind. Es kann aber wohl der Fall eintreten, dass z.\,B. $A$ drei Pole hat, w\"ahrend eine Potenz, $A^k$, eine Hauptaxe besitzt. Dann m\"ussen zwei der Pole von $A$ auf der Hauptaxe von $A^k$ liegen, und der dritte Pol von $A$ ist der einzelne Pol von $A^k$. Ist n\"amlich $A$ eine Substitution mit drei Polen, so k\"onnen wir sie, indem wir die Pole in die Ecken des Coordinatendreiecks legen, in die Normalform
\[
A=\begin{pmatrix}\alpha, & 0, & 0\\ 0, & \beta, & 0\\ 0, & 0, & \gamma\end{pmatrix}
\]
setzen, worin $\alpha,\beta,\gamma$ von einander verschieden sind. Ist nun $\beta^k=\gamma^k$, also $\beta$ von $\gamma$ um eine $k^{\text{te}}$ Einheitswurzel als Factor unterschieden, so hat $A^k$ die Linie $x_1=0$ zur Axe, und die gegen\"uberliegende Ecke des Coordinatendreiecks zum Pol.

Ist $A$ von endlichem Grade $\mu$, so sind $\alpha,\beta,\gamma$ $\mu^{\text{te}}$ Einheitswurzeln. Ist $A^k$ die niedrigste Potenz von $A$, die eine Hauptaxe hat, so muss $\beta:\gamma$ primitive $k^{\text{te}}$ und zugleich $\mu^{\text{te}}$ Einheitswurzel sein, und folglich muss $k$ ein echter Theiler von $\mu$ sein; wenn $k$ relativ prim zu $\mu$ ist, so hat $A^k$ dieselben drei Pole wie $A$.

\begin{center}
\S.\,117.

\gesperrt{Anwendung auf die Gruppe $G_{168}$. Siebenz\"ahlige Pole.}
\end{center}

Wir wollen nun die Pole und Axen der Substitutionen unserer Gruppe $G_{168}$ aufsuchen.


% ===== scan leaf 464, printed 443 =====
% !PRINTED 443
\weberpage{443}
Diese Arbeit wird wesentlich dadurch erleichtert, dass uns aus der Betrachtung der Congruenzgruppe $L_7$ die Grade und Cykeln der Gruppe schon bekannt sind (\S.\,68). Danach giebt es $p^2-1=48$ Elemente $7^{\text{ten}}$ Grades $A_7$, $\frac14 p(p-3)(p+1)=56$ Elemente $3^{\text{ten}}$ Grades $A_3$ und $\frac14 p(p-1)^2=63$ Elemente $4^{\text{ten}}$ oder $2^{\text{ten}}$ Grades $A_4, A_2$ in $G_{168}$.

Die Elemente $A_7$ ordnen sich (mit Zuziehung der identischen Substitution) in acht Cykeln von sieben Gliedern, die $A_3$ in 28 Cykeln von drei Elementen und die $A_4$ und $A_2$ in 21 Cykeln von vier Gliedern. Jeder dieser letzteren Cykeln enth\"alt ein Element $A_2$ und zwei Elemente $A_4$, und wenn wir also zusammenfassen, so erhalten wir:
\[
\begin{aligned}
&48 \text{ Elemente } 7^{\text{ten}} \text{ Grades}\\
&56 \text{ „} \quad 3^{\text{ten}} \text{ „}\\
&42 \text{ „} \quad 4^{\text{ten}} \text{ „}\\
&21 \text{ „} \quad 2^{\text{ten}} \text{ „}.
\end{aligned}
\]
Wenn ein Punkt durch $\nu$ Substitutionen der Gruppe (die identische eingeschlossen) unge\"andert bleibt, so wollen wir einen solchen Punkt einen \gesperrt{$\nu$-z\"ahligen Pol} nennen (\S.\,52).

F\"ur die Feststellung der Pole und Axen gen\"ugt die Betrachtung je eines Cyklus, da aus diesem die anderen durch Transformation abgeleitet sind (\S.\,69).

Nehmen wir die Substitution $7^{\text{ten}}$ Grades $\tau$ [\S.\,115, (4)], so finden wir die Ecken des Coordinatendreiecks als drei getrennte Pole. Alle Potenzen von $\tau$, deren Exponenten nicht durch 7 theilbar sind, haben dieselben drei Pole.

Durch Anwendung der Substitutionen $\chi,\chi^2$ auf die Coordinaten dieser drei Pole gehen dieselben cyklisch in einander \"uber; wendet man aber die Substitutionen $\Theta,\Theta^2,\Theta^3,\omega,\omega\Theta,\omega\Theta^2,\omega\Theta^3$ an, die in \S.\,115, (5), (10), (16), (18) vollst\"andig gebildet sind, so geht das ganze Coordinatendreieck in ein v\"ollig davon verschiedenes \"uber (z.\,B. durch $\omega$ in das Dreieck mit den Eckcoordinaten $\alpha,\beta,\gamma$; $\beta,\gamma,\alpha$; $\gamma,\alpha,\beta$), und daraus ist zu schliessen, dass diese Pole nicht mehr als siebenz\"ahlig sind. Es giebt acht Tripel solcher siebenz\"ahliger Pole, die alle von einander verschieden sind, und jeder von ihnen kann in jeden anderen durch eine Substitution der Gruppe $G_{168}$ transformirt werden.

Wir haben also den Satz:
%CONT>>


% ===== scan leaf 465, printed 444 =====
% !PRINTED 444
\weberpage{444}
\pdfbookmark[2]{§ 118. Die Hauptaxen}{bm135}
%<<CONT
\gesperrt{1. Es giebt 24 siebenz\"ahlige Pole, die sich, den acht siebengliedrigen Cykeln entsprechend, in acht Tripel theilen. Jeder dieser 24 Punkte kann in jeden anderen durch Substitutionen der Gruppe $G_{168}$ transformirt werden.}

\begin{center}
\S.\,118.

\gesperrt{Die Hauptaxen.}
\end{center}

Wir gehen zur Betrachtung der Substitutionen $4^{\text{ten}}$ und $2^{\text{ten}}$ Grades \"uber, als deren Repr\"asentanten wir $\Theta$ und $\Theta^2$ w\"ahlen. Nach \S.\,115, (10) und \S.\,116 haben wir f\"ur $\Theta$ die cubische Gleichung zu l\"osen:
\weq{1}{
\begin{vmatrix}
\gamma\varepsilon^2-\lambda, & \alpha\varepsilon^6, & \beta\\
\alpha, & \beta\varepsilon^4-\lambda, & \gamma\varepsilon^5\\
\beta\varepsilon^3, & \gamma, & \alpha\varepsilon-\lambda
\end{vmatrix} = 0,
}
die entwickelt so lautet:
\weq{2}{
\begin{aligned}
&\lambda^3-\lambda^2(\alpha\varepsilon+\beta\varepsilon^4+\gamma\varepsilon^2)\\
&\quad -\lambda\left[\varepsilon^6(\alpha^2-\beta\gamma)+\varepsilon^3(\beta^2-\gamma\alpha)+\varepsilon^5(\gamma^2-\alpha\beta)\right]\\
&\quad +\alpha^3+\beta^3+\gamma^3-3\alpha\beta\gamma=0.
\end{aligned}
}
Diese Gleichung reducirt sich aber mit Benutzung der Relationen (8), (9), (14) des \S.\,115 auf
\weq{3}{\lambda^3-\lambda^2+\lambda-1=0,}
und hat die Wurzeln
\[
\lambda=1,\quad \lambda=\pm i,
\]
und daraus lassen sich die Coordinaten der drei Pole berechnen, die sich als rationale Functionen von $\varepsilon$ und $i$ ergeben.

F\"ur die Pole von $\Theta^2$ erhalten wir nach \S.\,115, (16) zun\"achst die cubische Gleichung
\[
\begin{vmatrix}
\beta-\lambda, & \gamma\varepsilon^3, & \alpha\varepsilon^2\\
\gamma\varepsilon^{-3}, & \alpha-\lambda, & \beta\varepsilon^{-1}\\
\alpha\varepsilon^{-2}, & \beta\varepsilon, & \gamma-\lambda
\end{vmatrix} = 0,
\]
die nach den Relationen zwischen $\alpha,\beta,\gamma$ leicht in die Form
\weq{4}{(\lambda+1)^2(\lambda-1)=0}
gebracht wird und daher eine Doppelwurzel $\lambda=-1$ und eine einfache Wurzel $\lambda=1$ hat. Setzen wir den Werth $\lambda=-1$
%CONT>>


% ===== scan leaf 466, printed 445 =====
% !PRINTED 445
\weberpage{445}
%<<CONT
in die Gleichungen ein, durch die der unge\"anderte Punkt bestimmt wird [\S.\,116, (3)], so ergiebt sich:

\[
\begin{array}{l}
(\beta+1)\,x_1 + \gamma\,\varepsilon^3 x_2 + \alpha\,\varepsilon^2 x_3 = 0,\\
\gamma\,\varepsilon^{-3} x_1 + (\alpha+1)\,x_2 + \beta\,\varepsilon^{-1} x_3 = 0,\\
\alpha\,\varepsilon^{-2} x_1 + \beta\,\varepsilon\, x_2 + (\gamma+1)\,x_3 = 0.
\end{array}
\]
Diese drei Gleichungen reduciren sich alle auf die eine, die man z. B. aus der ersten durch Multiplication mit $\beta$ und Anwendung von \S.\,115, (8) ableitet:

\weq{5}{\alpha\,\gamma\,x_1 + \beta\,\gamma\,\varepsilon^3 x_2 + \alpha\,\beta\,\varepsilon^2 x_3 = 0,}
und die eine gerade Linie darstellt. Diese Linie ist also eine \gesperrt{Hauptaxe}.

Solcher Hauptaxen erhalten wir 21, da wir 21 Substitutionen 2\textsuperscript{ten} Grades haben.

Man kann die Functionen, die, gleich Null gesetzt, die 21 Hauptaxen darstellen, aus (5) leicht bilden, wenn man die Function

\weq{6}{A = \alpha\,\gamma\,x_1 + \beta\,\gamma\,\varepsilon^3 x_2 + \alpha\,\beta\,\varepsilon^2 x_3}
durch die Substitutionen $\tau^r$, $\tau^r\chi$, $\tau^r\chi^2$ transformirt ($r = 0, 1, 2, 3, 4, 5, 6$). Man erh\"alt so die 21 Functionen:

\weq{7}{\begin{array}{l}
A_{1,r} = \alpha\,\gamma\,\varepsilon^r x_1 + \gamma\,\beta\,\varepsilon^{2r+3} x_2 + \alpha\,\beta\,\varepsilon^{4r+2} x_3\\
A_{2,r} = \beta\,\gamma\,\varepsilon^{r+3} x_1 + \alpha\,\beta\,\varepsilon^{2r+2} x_2 + \alpha\,\gamma\,\varepsilon^{4r} x_3\\
A_{3,r} = \alpha\,\beta\,\varepsilon^{r+2} x_1 + \alpha\,\gamma\,\varepsilon^{2r} x_2 + \beta\,\gamma\,\varepsilon^{4r+3} x_3,
\end{array}}
aus denen leicht zu ersehen ist, dass die 21 Hauptaxen wirklich alle von einander verschieden sind.

Wir fassen dies so zusammen:

\noindent 2.\ \gesperrt{Es geh\"oren 21 verschiedene Hauptaxen zu der Gruppe $G_{168}$, von denen jede in jede andere durch Substitutionen der Gruppe transformirt werden kann.}

Da es nur 21 Hauptaxen giebt, so muss jede von ihnen durch acht Substitutionen unge\"andert bleiben, und diese acht Substitutionen bilden eine Gruppe. Um diese Gruppe zu ermitteln, wenden wir auf den in (6) dargestellten Ausdruck $A$ die Substitution $\omega$ [\S.\,115, (5)] an, und dadurch geht er \"uber in

\weq{8}{\begin{array}{l}
\gamma\,(\alpha^2 + \beta^2\varepsilon^3 + \alpha\,\beta\,\varepsilon^2)\,x_1 + \beta\,(\alpha\,\gamma + \gamma^2\varepsilon^3 + \alpha^2\varepsilon^2)\,x_2\\
{}+ \alpha\,(\gamma^2 + \beta\,\gamma\,\varepsilon^3 + \beta^2\varepsilon^2)\,x_3.
\end{array}}


% ===== scan leaf 467, printed 446 =====
% !PRINTED 446
\weberpage{446}
Nach \S.\,115, (14) und (8) ist aber

\[
\alpha^2 + \beta^2\varepsilon^3 + \alpha\,\beta\,\varepsilon^2 = \alpha^2 + \beta\,(\alpha\,\varepsilon^2 + \beta\,\varepsilon^3) = \alpha^2 - \beta\,\gamma = -\,\alpha,
\]
und wenn man die drei Co\"efficienten des Ausdrucks (8) in dieser Weise umformt, so ergiebt sich

\[
-\,\alpha\,\gamma\,x_1 - \beta\,\gamma\,\varepsilon^3 x_2 - \alpha\,\beta\,\varepsilon^2 x_3,
\]
d. h. $A$ \"andert durch Anwendung der Substitution $\omega$ sein Vorzeichen, und die Hauptaxe $A$ bleibt also durch die Substitution $\omega$ unge\"andert.

Da $A$ ausserdem durch die Substitution $\Theta$ unge\"andert bleibt, weil zwei Pole von $\Theta$ auf $A$ liegen, so haben wir die ganze Gruppe 8\textsuperscript{ten} Grades, durch die $A$ unge\"andert bleibt, die wir die \gesperrt{Gruppe von} $A$ nennen und mit $G_a$ bezeichnen, in der Form

\weq{9}{\omega^u\,\Theta^v.}
Unter diesen acht Substitutionen sind nur zwei, n\"amlich die identische und $\Theta^2$, die alle Punkte der Linie $A$ unge\"andert lassen. Bei den anderen bleiben nur je zwei Punkte in Ruhe. Denn berechnet man auf dem hier eingeschlagenen Wege aus \S.\,115, (5), (16) die Hauptaxen von $\omega$, $\omega\Theta$, $\omega\Theta^2$, $\omega\Theta^3$, so erh\"alt man, in der Bezeichnung (7),

\[
A_{2,1},\ A_{3,0},\ A_{3,3},\ A_{2,0},
\]
die alle von der Hauptaxe $A_{1,0}$ von $\Theta^2$ verschieden sind.

Die Substitution $\Theta^2$ hat ausser der Hauptaxe $A$ noch einen isolirten Pol, den wir erhalten, wenn wir die Wurzel $\lambda = 1$ der cubischen Gleichung (4) w\"ahlen. Dann ergeben sich f\"ur die Coordinaten dieses Poles die Gleichungen

\[
\begin{array}{l}
(\beta-1)\,x_1 + \gamma\,\varepsilon^3 x_2 + \alpha\,\varepsilon^2 x_3 = 0,\\
\gamma\,\varepsilon^{-3} x_2 + (\alpha-1)\,x_2 + \beta\,\varepsilon^{-1} x_3 = 0,\\ % [sic] printed x_2 twice, expected x_1 in first term
\alpha\,\varepsilon^{-2} x_1 + \beta\,\varepsilon\, x_2 + (\gamma-1)\,x_3 = 0,
\end{array}
\]
und daraus erh\"alt man

\weq{10}{x_1 : x_2 : x_3 = \alpha\,\gamma\,\varepsilon^4 : \beta\,\gamma\,\varepsilon : \alpha\,\beta\,\varepsilon^2.}
Solcher Punkte giebt es 21. Den Punkt (10) bezeichnen wir f\"ur den Augenblick mit $a$. Der Punkt $a$ bleibt nun nicht nur durch $\Theta$, sondern auch, wie eine Rechnung auf Grund der Formeln \S.\,115, (17), (14) zeigt, durch $\omega$, und folglich durch die ganze Gruppe $G_a$ unge\"andert, und ist also ein mindestens achtz\"ahliger Pol. Da er aber durch 21 Substitutionen in 21 verschiedene Lagen gebracht werden kann, so ist er auch nicht \gesperrt{mehr} als
%CONT>>


% ===== scan leaf 468, printed 447 =====
% !PRINTED 447
\weberpage{447}
%<<CONT
achtz\"ahlig. Da er durch die Substitutionen 2\textsuperscript{ten} Grades $\omega$, $\Theta\omega$, $\Theta^2\omega$, $\Theta^3\omega$ unge\"andert bleibt, und da die Pole dieser Substitutionen auf $A$ liegen (weil $A$ durch sie unge\"andert bleibt), w\"ahrend $a$ nicht auf $A$ liegt, so m\"ussen die Hauptaxen dieser vier Substitutionen $\omega$, $\Theta\omega$, $\Theta^2\omega$, $\Theta^3\omega$ durch den Punkt $a$ gehen.

Die Pole der vier Substitutionen $\omega$, $\Theta\omega$, $\Theta^2\omega$, $\Theta^3\omega$, die wir $a_1$, $a_2$, $a_3$, $a_4$ nennen wollen, sind als Bilder des Punktes $a$ gleichfalls achtz\"ahlige Pole, und m\"ussen, wie $a$, Pole von Substitutionen vierter Ordnung sein, die jedenfalls alle von $\Theta$ und $\Theta^3$ verschieden sind, weil sonst $\Theta^2$ einer jener vier Substitutionen gleich sein m\"usste, was nicht der Fall ist. Durch $\omega$ bleiben nur zwei Punkte der Axe $A$, n\"amlich die Pole von $\omega$ und von $\Theta^2\omega$, unge\"andert.

Betrachten wir nun die beiden anderen Pole $c_1$, $c_2$ der Substitution $\Theta$. Diese Punkte liegen gleichfalls auf der Hauptaxe $A$, sind aber von den Punkten $a_1$, $a_2$, $a_3$, $a_4$ verschieden, weil diese, wie schon bemerkt, nicht unter den Polen von $\Theta$ vorkommen. Nun ist $\omega\Theta\omega = \Theta$, also

\[
\omega\,\Theta(x) = \Theta\,\omega(x),
\]
und wenn nun $(x) = \Theta(x)$ ist, so ist auch

\[
\omega(x) = \Theta\,\omega(x);
\]
d. h. wenn $(x)$ die Coordinaten eines Poles von $\Theta$ sind, so haben die $\omega(x)$ die gleiche Bedeutung. Durch die Transformation $\omega$ gehen also die Pole von $\Theta$ nur in einander \"uber. Der Pol $a$ bleibt unge\"andert durch $\omega$; $c_1$ und $c_2$ dagegen k\"onnen nicht unge\"andert bleiben, weil sie unter den Polen von $\omega$ nicht vorkommen, und m\"ussen also in einander \"ubergehen.

Wir haben also folgende Uebersicht:

\begin{center}
\begin{tabular}{@{}lllllll@{}}
Unge\"andert durch & $\omega$ und $\Theta^2\omega$ & auf & $A$ & nur die Punkte & $a_1, a_3$ & \\
\textquotedbl\ \textquotedbl & $\Theta\omega$\ \textquotedbl\ $\Theta^3\omega$ & \textquotedbl & $A$ & \textquotedbl\ \textquotedbl\ \textquotedbl & $a_2, a_4$ & \\
\textquotedbl\ \textquotedbl & $\Theta$\ \textquotedbl\ $\Theta^3$ & \textquotedbl & $A$ & \textquotedbl\ \textquotedbl\ \textquotedbl & $c_1, c_2.$ &
\end{tabular}
\end{center}

Ein von diesen sechs verschiedener Punkt $x$ von $A$ geht nur durch die Substitutionen 1 und $\Theta^2$ in sich selbst \"uber und kann daher ein \gesperrt{zweiz\"ahliger Pol} genannt werden. Dann gilt der Satz:

\noindent 3.\ \gesperrt{Ein zweiz\"ahliger Pol geht durch die aus $\omega$ und $\Theta$ abgeleitete Gruppe 8\textsuperscript{ten} Grades in vier verschiedene Lagen \"uber.}


% ===== scan leaf 469, printed 448 =====
% !PRINTED 448
\weberpage{448}
\pdfbookmark[2]{§ 119. Die drei- und sechszähligen Pole}{bm136}
Ist n\"amlich $x$ ein zweiz\"ahliger Pol, und geht $x$ durch $\omega$ in $x'$, durch $\Theta$ in $x''$ \"uber, so sind nicht nur $x'$ und $x''$ von $x$, sondern sie sind auch unter einander verschieden, weil, wenn sie identisch w\"aren, $x$ durch $\omega\Theta^3$ unge\"andert bliebe. Durch $\omega\Theta$ geht dann $x$ in eine vierte Lage $x'''$ \"uber, und $x'''$ ist sowohl von $x'$ als von $x''$ verschieden, weil sonst $x$ durch $\omega\Theta\omega = \Theta$ oder durch $\omega\Theta\Theta^{-1} = \omega$ unge\"andert bliebe. Da $x$ durch $\Theta^2$ unge\"andert bleibt, so geht $x$ durch $\Theta^3$ in $x''$, durch $\Theta^2\omega = \omega\Theta^2$ in $x'$, durch $\omega\Theta^3$ in $x'''$ \"uber.

Da wir auf jeder Hauptaxe zwei Punkte $c_1$, $c_2$ haben, so giebt es im Ganzen 42, und jeder von ihnen kann in jeden anderen durch Substitutionen der Gruppe \"ubergef\"uhrt werden. Daraus folgt, dass diese Pole nicht mehr als vierz\"ahlig sind. Daher der Satz:

\noindent 4.\ \gesperrt{Es giebt 21 achtz\"ahlige und 42 vierz\"ahlige Pole. Auf jeder Hauptaxe liegen vier achtz\"ahlige und zwei vierz\"ahlige Pole. Durch jeden achtz\"ahligen Pol gehen vier Hauptaxen. Jeder achtz\"ahlige Pol kann in jeden anderen achtz\"ahligen und jeder vierz\"ahlige Pol in jeden anderen vierz\"ahligen Pol \"ubergef\"uhrt werden.}

\begin{center}
\S.\,119.

\gesperrt{Die drei- und sechsz\"ahligen Pole.}
\end{center}

Wir wenden uns nun zur Betrachtung der Substitutionen dritter Ordnung und ihrer Pole, und w\"ahlen als Repr\"asentanten einer solchen Substitution

\[
\chi = \begin{pmatrix} 0 & 0 & 1\\ 1 & 0 & 0\\ 0 & 1 & 0 \end{pmatrix},
\]
deren Pole man aus

\[
\lambda\,x_1 = x_3, \quad \lambda\,x_2 = x_1, \quad \lambda\,x_3 = x_2
\]
erh\"alt. Es folgt daraus $\lambda^3 = 1$; also ist $\lambda$ eine dritte Einheitswurzel, und wenn daher $\varrho$ eine imagin\"are dritte Einheitswurzel bedeutet, so ergeben sich die drei Pole:

\weq{1}{\begin{array}{l}
x_1 = x_2 = x_3\\
x_1 = \varrho\,x_2 = \varrho^2 x_3\\
x_1 = \varrho^2 x_2 = \varrho\,x_3.
\end{array}}


% ===== scan leaf 470, printed 449 =====
% !PRINTED 449
\weberpage{449}
Zwischen diesen drei Polen besteht aber ein wesentlicher Unterschied. Wenden wir nämlich die Substitution $\omega$ darauf an, so bleibt der erste von ihnen ungeändert. Er ist also mindestens sechszählig und gehört zu den Substitutionen
\[
1, \chi, \chi^2, \omega, \omega\chi, \omega\chi^2.
\]
Die beiden anderen Pole, nämlich
\weq{2}{x_1 = \varrho x_2 = \varrho^2 x_3, \quad x_1 = \varrho^2 x_2 = \varrho x_3,}
gehen durch die Substitution $\omega$ in einander über, und wir müssen nachweisen, dass sie durch jede andere Substitution der Gruppe, ausser $\chi, \chi^2$, verändert werden, dass sie also nur dreizählig sind. Bezeichnen wir für den Augenblick die beiden Punkte (2) mit $\pi$ und $\pi'$, so können wir zunächst sagen, dass alle Substitutionen aus $G_{168}$, welche $\pi$ ungeändert lassen, eine Gruppe $G_\pi$ bilden
müssen, und zwar einen Theiler von $G_{168}$; die Gruppe $G_\pi$ enthält ihrerseits die cyklische Gruppe 1, $\chi$, $\chi^2$.

Ist $\sigma$ irgend eine Substitution von $G_\pi$, so kommt also $\pi$ unter den Polen von $\sigma$ vor. Daraus folgt, dass $\sigma$ weder vom $7^{\mathrm{ten}}$ noch vom $4^{\mathrm{ten}}$ Grade sein kann. Denn die Coordinaten der Pole von diesen zwei Graden sind rationale Functionen von $\varepsilon, i$, während die Coordinaten von $\pi$ die davon unabhängige Irrationalität $\varrho$ enthalten (Bd. I, \S.\,134, vergl. den Nachtrag am Ende dieses Bandes). Dass aber $\sigma$ auch nicht vom $2^{\mathrm{ten}}$ Grade sein kann, ergiebt sich daraus, dass $\pi$ auf keiner der Hauptaxen liegt. Denn läge es auf einer Hauptaxe, so müsste, weil $\varrho$ nicht rational durch $\varepsilon$ ausgedrückt werden kann, wie aus (2) mittelst der Gleichung
\[
1 + \varrho + \varrho^2 = 0
\]
abzuleiten ist, nach \S.\,118, (7) eine der 21 Relationen bestehen
\[
\begin{aligned}
\alpha &= \beta\varepsilon^{r+3} &&= \gamma\varepsilon^{-2r+1}, \\
\alpha &= \beta\varepsilon^{-3r+3} &&= \gamma\varepsilon^{-r+1}, \\
\alpha &= \beta\varepsilon^{2r+3} &&= \gamma\varepsilon^{3r+1},
\end{aligned}
\]
von denen offenbar keine möglich ist. Die Gruppe $G_\pi$ enthält also ausser der identischen nur Substitutionen $3^{\mathrm{ten}}$ Grades, und der Grad von $G_\pi$ muss eine Potenz von 3 sein. Da aber 168 nicht durch 9 theilbar ist, so muss der Grad von $G_\pi$ gleich 3 sein. Hieraus folgt:

5. \gesperrt{Es giebt 56 dreizählige Pole, die zu je zweien zu derselben Substitution gehören, und sich danach in 28 Paare ordnen. Jeder dieser Pole}
%CONT>>


% ===== scan leaf 471, printed 450 =====
% !PRINTED 450
\weberpage{450}
%<<CONT
\gesperrt{kann durch Substitutionen der Gruppe} $G_{168}$ \gesperrt{in jeden anderen transformirt werden.}

Setzen wir
\weq{3}{T = x_1 + x_2 + x_3,}
so ist $T = 0$ die Gleichung der Verbindungslinie der Punkte $\pi, \pi'$. Die Function $T$ bleibt durch $\chi$ ungeändert und wechselt durch $\omega$ ihr Vorzeichen. Sie geht aber durch die 28 Substitutionen $\Theta^s\tau^r$ in 28 \gesperrt{verschiedene} Functionen $T_{s,r}$ über, die wir mit Hülfe der Formeln des \S.\,115, (11), (14) leicht so darstellen können:
\weq{4}{
\begin{aligned}
T_{0,r} &= \varepsilon^r x_1 + \varepsilon^{2r} x_2 + \varepsilon^{4r} x_3 \\
h\,T_{1,r} &= \beta^2\varepsilon^{1+r} x_1 + \alpha^2\varepsilon^{2+2r} x_2 + \gamma^2\varepsilon^{4+4r} x_3 \\
-h\,T_{2,r} &= \alpha^2\varepsilon^{2+r} x_1 + \gamma^2\varepsilon^{4+2r} x_2 + \beta^2\varepsilon^{1+4r} x_3 \\
h\,T_{3,r} &= \gamma^2\varepsilon^{4+r} x_1 + \beta^2\varepsilon^{1+2r} x_2 + \alpha^2\varepsilon^{2+4r} x_3,
\end{aligned}
}
und diese stellen, gleich Null gesetzt, 28 verschiedene gerade Linien dar, die alle aus einer von ihnen durch Substitutionen der Gruppe ableitbar sind.

Auf der Linie $T$ liegen die drei Punkte mit den Coordinaten
\[
\begin{aligned}
\alpha\gamma &: \beta\gamma : \beta\alpha \\
\beta\alpha &: \alpha\gamma : \beta\gamma \\
\beta\gamma &: \beta\alpha : \alpha\gamma,
\end{aligned}
\]
die aus (10), \S.\,118 durch die Substitutionen $\tau^{-4}, \tau^{-4}\chi, \tau^{-4}\chi^2$ hervorgehen und also zu den achtzähligen Polen gehören.

6. \gesperrt{Auf jeder Linie $T$ liegen drei achtzählige und zwei dreizählige Pole.}

Es bleibt noch der dritte Pol $\pi_0$ der Substitution $\chi$ mit den Coordinaten
\[
x_1 = x_2 = x_3
\]
zu untersuchen, von dem wir schon nachgewiesen haben, dass er durch die sechs Substitutionen
\weq{5}{1, \chi, \chi^2; \omega, \omega\chi, \omega\chi^2}
ungeändert bleibt. Es fragt sich nun, ob dieser Pol nicht mehr als sechszählig ist.

Wenn es ausser (5) noch andere Substitutionen giebt, die den Pol $\pi_0$ unverändert lassen, so müssen diese nothwendig vom $2^{\mathrm{ten}}$ Grade sein, da $\pi_0$ weder ein vierzähliger noch ein siebenzähliger Pol ist, und weil ferner der Grad der zu $\pi_0$ gehörigen Gruppe nicht durch 9 theilbar sein kann.


% ===== scan leaf 472, printed 451 =====
% !PRINTED 451
\weberpage{451}
\pdfbookmark[2]{§ 120. Die Configuration der Gruppe G168}{bm137}
Nun liegt der Punkt $\pi_0$ auf den drei Hauptaxen der Substitutionen $2^{\mathrm{ten}}$ Grades $\omega, \omega\chi, \omega\chi^2$. Giebt es noch eine weitere Substitution $2^{\mathrm{ten}}$ Grades, durch die $\pi_0$ ungeändert bleibt, so muss $\pi_0$ auf der Hauptaxe dieser Substitution liegen.

Setzen wir aber in den Gleichungen der Hauptaxen [\S.\,118, (7)] $x_1 = x_2 = x_3$, so ist jede dieser drei Gleichungen nur für einen Werth von $r$ befriedigt; also kann $\pi_0$ nicht auf mehr als dreien der Hauptaxen liegen, und $\pi_0$ ist sechszählig.

Daraus der Satz:

7. \gesperrt{Es giebt 42 sechszählige Pole, die alle aus einem von ihnen durch Substitutionen der Gruppe ableitbar sind. Durch jeden dieser Pole gehen drei Hauptaxen der Gruppe.}

Der Pol $\pi_0$ liegt auf keiner der Linien $T_{s,r}$. Dies lässt sich leicht zeigen, wenn man in den Ausdrücken (4) $x_1 = x_2 = x_3$, und dann für $\alpha^2, \beta^2, \gamma^2$ ihre Ausdrücke durch $\varepsilon$ setzt. Man sieht dann leicht (mit Anwendung der Irreducibilität der Kreistheilungsgleichung), dass von diesen Ausdrücken keiner verschwindet.

\S.\,120.

\gesperrt{Die Configuration der Gruppe $G_{168}$.}

Die Gesammtheit der geraden Linien und Punkte, die wir in den vorangegangenen Paragraphen betrachtet haben, wollen wir die \gesperrt{Configuration der Gruppe} $G_{168}$ nennen. Das Wort hat hier dieselbe Bedeutung, in der es in neuerer Zeit in der Geometrie gebraucht wird\footnote{Der Ausdruck ist von \gesperrt{Reye} eingeführt: \glqq{}Geometrie der Lage.\grqq{} Bd. I, 2. Aufl. (1876).}.

Wir beschreiben diese Configuration im Folgenden etwas näher, ohne etwas Neues hinzuzufügen, nur um die Sätze der letzten Paragraphen anschaulicher und übersichtlicher hervortreten zu lassen.

Wir haben in dieser Configuration:
\[
\begin{aligned}
&\text{21 Linien } A \text{ (die Hauptaxen)}, \\
&\text{21 Punkte } a \text{ (die achtzähligen Pole)}, \\
&\text{28 Linien } T, \\
&\text{28 Punkte } t \text{ (die sechszähligen Pole)}.
\end{aligned}
\]


% ===== scan leaf 473, printed 452 =====
% !PRINTED 452
\weberpage{452}
Das ganze System geht durch 168 Substitutionen der Gruppe in sich selbst über.

\[
\begin{aligned}
&\text{Auf jeder Linie } A \text{ liegen vier Punkte } a. \\
&\text{Durch jeden Punkt } a \text{ gehen vier Linien } A. \\
&\text{Auf jeder Linie } T \text{ liegen drei Punkte } a. \\
&\text{Durch jeden Punkt } t \text{ gehen drei Linien } A.
\end{aligned}
\]

\gesperrt{Die Punkte} $a$ \gesperrt{und} $t$ \gesperrt{bilden zusammen das vollständige System aller Schnittpunkte der Linien} $A$\gesperrt{.}

Denn 21 gerade Linien schneiden sich in 210 Punkten, und von diesen fallen drei oder sechs zusammen, wenn drei oder vier dieser Linien durch einen Punkt gehen. Es ist aber
\[
210 = 21.6 + 28.3.
\]

\gesperrt{Auf jeder Linie} $A$ \gesperrt{liegen vier Punkte} $t$\gesperrt{.}

Denn da durch jeden Punkt $t$ drei Linien $A$ gehen, und auf jeder Linie $A$ gleich viele Punkte $t$ liegen müssen (weil jede Linie $A$ in jede andere transformirbar ist), so ist, wenn $x$ die Anzahl dieser Punkte ist, $x . 21 : 3 = 28$, also $x = 4$. Ebenso schliesst man:

\gesperrt{Durch jeden Punkt} $a$ \gesperrt{gehen vier Linien} $T$\gesperrt{.}

Keiner der Punkte $t$ liegt auf einer Linie $T$.
Bezeichnen wir die $\nu$-zähligen Pole mit $P_\nu$, so haben wir also folgende Systeme von Polen:
\[
\begin{aligned}
&21 \text{ achtzählige Pole } P_8, \\
&28 \text{ sechszählige} \;\text{\textquotedbl}\; P_6, \\
&42 \text{ vierzählige} \;\text{\textquotedbl}\; P_4, \\
&56 \text{ dreizählige} \;\text{\textquotedbl}\; P_3, \\
&24 \text{ siebenzählige} \;\text{\textquotedbl}\; P_7, \\
&\text{unendlich viele zweizählige Pole } P_2.
\end{aligned}
\]

Die $P_2$ sind alle von $P_4, P_6, P_8$ verschiedene Punkte der Hauptaxen; von den Punkten $P_4$ liegen je zwei auf einer Linie $A$; von den Punkten $P_3$ liegen je zwei auf einer Linie $T$; die Punkte $P_7$ liegen nicht auf den Linien $A$ (dass sie auch nicht auf $T$ liegen, wird sich später ergeben).

Das System der Axen ist genau entsprechend dem System der Pole. Jede gerade Linie, die durch einen achtzähligen Pol geht, ist eine Axe (darunter die Hauptaxen, vier Linien $T$ und zwei Verbindungslinien des $P_8$ mit $P_4$). Demnach könnte
%CONT>>


% ===== scan leaf 474, printed 453 =====
% !PRINTED 453
\weberpage{453}
\pdfbookmark[2]{§ 121. Invariantencurven der Gruppe G168}{bm138}

%<<CONT
man die $P_8$ passend die \gesperrt{Hauptpole} nennen und die durch sie gehenden Axen mit $A_8, A_6, A_4, A_2$ bezeichnen.

Hervorzuheben sind ferner die 28 Paare von Verbindungslinien eines $P_6$ mit den beiden zugehörigen $P_3$, die wir mit $A_3$ bezeichnen können, und endlich die 24 Verbindungslinien je zweier zusammengehöriger $P_7$, als deren Repräsentanten die Seiten des Coordinatendreiecks zu betrachten sind, die mit $A_7$ bezeichnet werden können.

\begin{center}
\S.\,121.\\
\gesperrt{Invariantencurven der Gruppe $G_{168}$.}
\end{center}

Eine Form $\mu^{\mathrm{ten}}$ Grades $\varphi\,(x_1, x_2, x_3)$ wird durch die Substitutionen der Gruppe im Allgemeinen in 168 verschiedene Formen übergehen. Bei besonderen Formen von $\varphi$ kann diese Zahl sich aber verringern, und dann bilden die Substitutionen, durch die $\varphi$ ungeändert bleibt, eine in $G_{168}$ enthaltene Gruppe $G'$, deren Grad ein Theiler von 168 sein muss. Die Anzahl der Formen, in die $\varphi$ übergeht, ist der Index des Theilers $G'$, und jede dieser verschiedenen Formen von $\varphi$ bleibt durch eine mit $G'$ conjugirte Gruppe ungeändert. Die Form $\varphi$ ist eine absolute Invariante der Gruppe $G'$.

Die Gleichung $\varphi = 0$ stellt eine auf das Coordinatendreieck $x_1, x_2, x_3$ bezogene Curve dar, und diese Curve wird eben durch die Substitutionen von $G_{168}$ auf andere Curven abgebildet. Sind die 168 Bildcurven nicht alle von einander verschieden, so bleibt die Curve $\varphi$ durch die Substitutionen einer Gruppe $G'$ ungeändert oder ändert sich nur um einen constanten Factor, ist also relative Invariante von $G'$. Eine gerade Linie bleibt nur dann durch andere als die identische Substitution ungeändert, wenn sie zu den im vorigen Paragraphen beschriebenen Axen gehört.

Alle Eigenschaften und Beziehungen zwischen Punkten, Linien und Curven, die durch lineare Transformation unzerstörbar sind (die sogenannten projectiven Eigenschaften der Geometrie), bleiben in den Bildern erhalten. Wenn also z.\,B. eine gerade Linie Tangente oder Wendetangente oder Doppeltangente einer Curve ist, so stehen alle Bilder der geraden Linie in derselben Beziehung zu den Bildern der Curve. Ebenso wenn ein Punkt
%CONT>>


% ===== scan leaf 475, printed 454 =====
% !PRINTED 454
\weberpage{454}
\pdfbookmark[2]{§ 122. Die erste Invariante der Gruppe G168 und die Grundcurve}{bm139}

%<<CONT
Doppelpunkt, Wendepunkt, Rückkehrpunkt, Berührungspunkt einer Doppeltangente u.\,s.\,w. ist.

Unter den Formen $\varphi$ sind uns nun vor Allem die von Wichtigkeit, die bei der ganzen Gruppe ungeändert bleiben, die Invarianten, deren es hier, da die Gruppe einfach ist, nur absolute giebt (\S.\,40). Ist $\Phi\,(x_1, x_2, x_3)$ eine solche Form, so soll die durch die Gleichung $\Phi = 0$ dargestellte Curve eine invariante Curve der Gruppe heissen. Es giebt 168 Abbildungen der Ebene auf sich selbst, bei denen alle Punkte der Curve in Punkte derselben Curve übergehen. Liegt irgend ein Punkt auf dieser Curve, so liegen auch alle seine Bildpunkte darauf.

Im Allgemeinen ist die Zahl der so mit einander verbundenen Punkte der Curve 168, jedenfalls nicht grösser. Ist sie kleiner, so müssen die Punkte \gesperrt{Pole} sein, und die Anzahl der Bildpunkte ist ein Theiler von 168 (nämlich 24 für die $P_7$, 21 für die $P_8$, 28 für die $P_6$, 42 für die $P_4$, 56 für die $P_3$ und 84 für ein System zusammengehöriger $P_2$). Wenn ein Pol von einer dieser Arten auf der Curve liegt, so liegen alle Pole von derselben Art darauf.

\begin{center}
\S.\,122.\\
\gesperrt{Die erste Invariante der Gruppe $G_{168}$ und die Grundcurve.}
\end{center}

Um nun die Invarianten unserer Gruppe zu bilden, suchen wir Formen der drei Variablen $x_1, x_2, x_3$ auf, die durch Anwendung der drei Substitutionen $\chi, \tau, \omega$ (\S.\,115) ungeändert bleiben. Dies genügt, da die ganze Gruppe sich aus diesen drei Substitutionen zusammensetzen lässt (\S.\,72). Nun ist $\chi$ eine cyklische Vertauschung der drei Variablen $x_1, x_2, x_3$, und $\tau$ bedeutet die Substitution

\[
\begin{pmatrix} x_1, & x_2, & x_3 \\ \varepsilon x_1, & \varepsilon^2 x_2, & \varepsilon^4 x_3 \end{pmatrix}.
\]

Wenn also in einer Invariante ein Glied $x_1^{h_1} x_2^{h_2} x_3^{h_3}$ vorkommt, worin $h_1, h_2, h_3$ ganze nicht negative Zahlen sind, so verlangt die Unveränderlichkeit durch $\tau$, dass

\weq{1}{h_1 + 2h_2 + 4h_3 \equiv 0\ (\mathrm{mod}\ 7)}
%CONT>>


% ===== scan leaf 476, printed 455 =====
% !PRINTED 455
\weberpage{455}

%<<CONT
und die Unveränderlichkeit durch $\chi$, dass neben diesem einen Gliede noch die zwei entsprechenden

\[
x_2^{h_1} x_3^{h_2} x_1^{h_3}, \quad x_3^{h_1} x_1^{h_2} x_2^{h_3}
\]

in der Function vorkommen, wenn nicht $h_1 = h_2 = h_3$ ist. Dazu kommt noch die Bedingung der Unveränderlichkeit durch $\omega$. Wir wollen nun sehen, wie wir diesen Forderungen genügen können, und zwar zunächst so, dass der Grad $m$ der Invariante, d.\,h. die Summe $h_1 + h_2 + h_3$, möglichst klein wird.

Die Bedingung (1) kann offenbar nicht erfüllt sein, wenn $m < 3$ ist. Ist diese Summe $= 3$, so muss $h_1 = h_2 = h_3 = 1$ sein; aber das Product $x_1 x_2 x_3$ ist offenbar nicht unverändert durch $\omega$. Der kleinste Werth, der in Betracht kommt, ist also $m = 4$, und es sind also alle nicht negativen Lösungen von

\[
\begin{aligned}
h_1 + h_2 + h_3 &= 4 \\
h_1 + 2h_2 + 4h_3 &\equiv 0\ (\mathrm{mod}\ 7)
\end{aligned}
\]
%CONT>>

aufzusuchen. Eliminiren wir $h_1$, so folgt

\[
h_2 + 3h_3 \equiv 3\ (\mathrm{mod}\ 7),
\]

und daraus ergeben sich die einzig möglichen Lösungen:

\[
\begin{aligned}
h_3 &= 0, & h_2 &= 3, & h_1 &= 1 \\
h_3 &= 1, & h_2 &= 0, & h_1 &= 3 \\
h_3 &= 3, & h_2 &= 1, & h_1 &= 0,
\end{aligned}
\]

und folglich die \gesperrt{einzige}, durch $\chi$ und $\tau$ ungeänderte Form $4^{\mathrm{ten}}$ Grades

\weq{2}{f\,(x_1, x_2, x_3) = x_1^3 x_3 + x_2^3 x_1 + x_3^3 x_2.}

Um den Einfluss der Substitution $\omega$ auf die Function $f$ zu prüfen, setzen wir

\weq{3}{
\begin{aligned}
x_1 &= \alpha y_1 + \beta y_2 + \gamma y_3 \\
x_2 &= \beta y_1 + \gamma y_2 + \alpha y_3 \\
x_3 &= \gamma y_1 + \alpha y_2 + \beta y_3,
\end{aligned}
}

worin die Coëfficienten $\alpha, \beta, \gamma$ die in \S.\,115 festgesetzte Bedeutung haben, und nehmen an, dass durch (3) die Transformation

\weq{4}{F\,(y_1, y_2, y_3) = f\,(x_1, x_2, x_3)}

geleistet werde. Wir bilden die zweiten Ableitungen von $F$ nach $y_1, y_2, y_3$ mit Rücksicht auf (3) und auf die Formeln

\[
\begin{aligned}
&\tfrac{1}{6} f''(x_1, x_1) = x_1 x_3, \quad \tfrac{1}{6} f''(x_2, x_2) = x_2 x_1, \quad \tfrac{1}{6} f''(x_3, x_3) = x_3 x_2 \\
&\tfrac{1}{3} f''(x_2, x_3) = x_3^2, \quad \tfrac{1}{3} f''(x_3, x_1) = x_1^2, \quad \tfrac{1}{3} f''(x_1, x_2) = x_2^2.
\end{aligned}
\]


% ===== scan leaf 477, printed 456 =====
% !PRINTED 456
\weberpage{456}

%<<CONT
Daraus ergiebt sich:

\[
\begin{aligned}
\tfrac{1}{6} F''(y_1, y_1) &= x_1 x_3 \alpha^2 + x_2 x_1 \beta^2 + x_3 x_2 \gamma^2 \\
&\quad + x_3^2 \beta\gamma + x_1^2 \alpha\gamma + x_2^2 \alpha\beta, \\
\tfrac{1}{3} F'''(y_2, y_3) &= 2x_1 x_3 \beta\gamma + 2x_2 x_1 \alpha\gamma + 2x_3 x_2 \alpha\beta \\
&\quad + x_3^2 (\beta\gamma + \alpha^2) + x_1^2 (\alpha\gamma + \beta^2) + x_2^2 (\alpha\beta + \gamma^2).
\end{aligned}
\]
% [sic] mixed term printed with triple prime F''' here against F'' just below for the same pairing

Da nun die Auflösungen des Systems (3) von derselben Form sind ($y_1 = \alpha x_1 + \beta x_2 + \gamma x_3, \dots$), so erhält man hieraus:

\[
\begin{aligned}
\tfrac{1}{6} F''(y_1, y_1) - y_1 y_3 &= \\
x_1 x_3 (\alpha^2 - \alpha\beta - \gamma^2) + x_2 x_1 (\beta^2 - \beta\gamma - \alpha^2) &+ x_2 x_3 (\gamma^2 - \beta^2 - \alpha\gamma), \\
\tfrac{1}{3} F''(y_2, y_3) - y_3^2 &= \\
x_1^2 (\beta^2 + \alpha\gamma - \gamma^2) + x_2^2 (\gamma^2 + \alpha\beta - \alpha^2) &+ x_3^2 (\alpha^2 + \beta\gamma - \beta^2).
\end{aligned}
\]

Die Coëfficienten auf der rechten Seite dieser Gleichungen $(\alpha^2 - \alpha\beta - \gamma^2) \dots$ ergeben sich aber aus den Werthen von $\alpha, \beta, \gamma$ [\S.\,115, (11)] als verschwindend, und wir erhalten also, wenn wir noch eine cyklische Vertauschung der $x$, die eine cyklische Vertauschung der $y$ zur Folge hat, anwenden:

\[
\begin{aligned}
&\tfrac{1}{6} F'''(y_1, y_1) = y_1 y_3, \quad \tfrac{1}{6} F'''(y_2, y_2) = y_2 y_1, \quad \tfrac{1}{6} F'''(y_3, y_3) = y_3 y_2, \\
&\tfrac{1}{3} F''(y_2, y_3) = y_3^2, \quad \tfrac{1}{3} F''(y_3, y_1) = y_1^2, \quad \tfrac{1}{3} F''(y_1, y_2) = y_2^2.
\end{aligned}
\]
% [sic] prime-count on F wavers between the two lines above, reproduced as printed

Demnach ergiebt sich nach dem \gesperrt{Euler}'schen Satze [Bd.\ I, \S.\,17, (6)]:

\[
F\,(y_1, y_2, y_3) = y_1^3 y_3 + y_2^3 y_1 + y_3^3 y_2,
\]

und damit ist nachgewiesen, dass die Function $f\,(x_1, x_2, x_3)$ in der That eine Invariante unserer Gruppe $G_{168}$ ist. Es ist, abgesehen von einem willkürlich beizufügenden constanten Factor, die einzige Invariante $4^{\mathrm{ten}}$ Grades.

Die Gleichung

\weq{5}{f\,(x_1, x_2, x_3) = 0}

bedeutet, wenn $x_1, x_2, x_3$ Coordinaten in der Ebene sind, eine \gesperrt{Curve vierter Ordnung}, die wir die \gesperrt{Grundcurve} der Gruppe oder die Curve $f$ nennen wollen, und es giebt 168 Abbildungen der Ebene auf sich selbst, bei denen jedem Punkte dieser Curve ein Punkt der Curve entspricht. Wenn ein Punkt auf der Grundcurve liegt, so liegen auch alle seine Bildpunkte darauf. Die Anzahl der verschiedenen Bildpunkte eines Punktes kann nur dann auf eine kleinere Zahl als 168, und zwar immer auf einen Theiler von 168, heruntersinken, wenn der Punkt ein Pol ist.


% ===== scan leaf 478, printed 457 =====
% !PRINTED 457
\weberpage{457}
Die gerade Linie $x_1 = 0$ schneidet die Curve in drei zusammenfallenden
Punkten bei $x_3 = 0$ und in einem vierten davon getrennten Punkte
$x_2 = 0$. Die Eckpunkte des Coordinatendreiecks sind also Wendepunkte
der Curve und die Seiten sind die Wendetangenten. Bezeichnen wir die
Seiten des Coordinatendreiecks mit 1, 2, 3 und die gegenüberliegenden
Ecken durch dieselben Ziffern, so ist die Seite 1 Wendetangente im
Punkt 2, die Seite 2 Wendetangente im Punkt 3 und die Seite 3
Wendetangente im Punkt 1.

Wir untersuchen nun die Lage der Pole und Axen in Bezug auf die
Grundcurve. Da die Eckpunkte des Coordinatendreiecks zu den
siebenzähligen Polen gehören, so schliessen wir zunächst,
\gesperrt{dass alle siebenzähligen Pole $P_7$ auf der Grundcurve
liegen.}

Es sind, wie die Eckpunkte selbst, alles \gesperrt{Wendepunkte} der Curve, und
diese ordnen sich in acht Wendepunktsdreiecke.

\gesperrt{Die dreizähligen Pole} $P_3$ \gesperrt{liegen gleichfalls auf der Grundcurve;
denn setzt man}
\[
x_2 = \varrho\,x_1, \quad x_3 = \varrho^2 x_1
\]
oder
\[
x_2 = \varrho^2 x_1, \quad x_3 = \varrho\,x_1,
\]
worin $\varrho$ eine cubische Einheitswurzel ist, so reducirt sich
$f(x_1, x_2, x_3)$ auf
\[
1 + \varrho + \varrho^2 = 0.
\]
Um die Schnittpunkte ihrer Verbindungslinie
\[
T = x_1 + x_2 + x_3 = 0
\]
mit der Grundcurve zu finden, setzt man $x_3 = -x_1 - x_2$, wodurch
$f(x_1, x_2, x_3) = 0$ in
\[
x_1^3(x_1+x_2) - x_2^3 x_1 + (x_1+x_2)^3 x_2 = (x_1^2+x_1x_2+x_2^2)^2 = 0
\]
übergeht. Da die linke Seite ein Quadrat ist, so fallen die vier
Schnittpunkte zweimal zu zweien zusammen, und es folgt, dass die
Linie $T$ eine Doppeltangente der Grundcurve ist.

\gesperrt{Die 28 Linien $T$ sind Doppeltangenten der Grundcurve; ihre
Berührungspunkte sind die 56 Pole $P_3$.}

Hieraus folgt beiläufig, dass die Punkte $P_7$ nicht auf den Linien
$T$ liegen, da eine Doppeltangente ausser den Berührungspunkten
keinen weiteren Schnittpunkt mit der Curve haben kann.

\gesperrt{Die sechs- und achtzähligen Pole liegen nicht auf der
Grundcurve.}


% ===== scan leaf 479, printed 458 =====
% !PRINTED 458
\weberpage{458}
Für die Punkte $P_6$ ist dies unmittelbar einzusehen, wenn man
[nach \S.\,119, (1)] $x_1 = x_2 = x_3 = 1$ setzt, wodurch
$f(x_1, x_2, x_3) = 3$ wird, also nicht verschwindet. Um dasselbe für
die $P_8$ nachzuweisen, setzt man nach \S.\,118 (10):
\[
x_1 = \alpha\gamma\varepsilon^4, \quad x_2 = \beta\gamma\varepsilon,
\quad x_3 = \alpha\beta\varepsilon^2,
\]
wodurch, da $\alpha\beta\gamma = \frac{1}{7}$ ist [\S.\,115, (14)],
\[
f(x_1, x_2, x_3) = \frac{1}{7}(\alpha^3\gamma^2 + \beta^3\alpha^2 + \gamma^3\beta^2)
\]
wird. Setzt man darin [nach \S.\,115, (8)]:
\[
\alpha^3 = \alpha\beta\gamma - \alpha^2, \quad \beta^3 = \alpha\beta\gamma - \beta^2,
\quad \gamma^3 = \alpha\beta\gamma - \gamma^2,
\]
so erhält man
\weq{6}{f(x_1, x_2, x_3) = \tfrac{1}{7}\alpha\beta\gamma(\alpha^2+\beta^2+\gamma^2) - \tfrac{1}{7}(\alpha^2\gamma^2+\beta^2\alpha^2+\gamma^2\beta^2).}
Aus \S.\,115, (6), (7) aber folgt
\[
\alpha^2 + \beta^2 + \gamma^2 = 1,
\]
\[
\alpha^2\gamma^2 + \beta^2\alpha^2 + \gamma^2\beta^2 = -\alpha\beta\gamma(\alpha+\beta+\gamma) = \frac{1}{7},
\]
und daher ist
\[
f(x_1, x_2, x_3) = \frac{1}{42}
\]
von Null verschieden.

Die sämmtlichen Schnittpunkte der 21 Hauptaxen unter einander sind
die Pole $P_6$ und $P_8$, und folglich schneiden sich niemals zwei
Hauptaxen auf der Grundcurve. Die Schnittpunkte der Hauptaxen mit der
Grundcurve können also nur vierzählige oder zweizählige Pole sein. Um
darüber zu entscheiden, stellen wir folgende Erwägung an. Wenn unter
den Schnittpunkten der Hauptaxe $A$ mit der Grundcurve ein $P_2$
vorkommt, so sind alle vier Schnittpunkte von einander verschieden und
sind alle zweizählige Pole. Wenn aber ein $P_4$ auf $A$ und $f$ liegt,
so liegt auch der zweite auf $A$ liegende $P_4$ auf $f$, und es kann
keinen anderen Schnittpunkt von $A$ und $f$ geben, weil ein solcher
ein $P_2$ wäre und vier weitere $P_2$ zur Folge hätte. Die vier
Schnittpunkte von $A$ und $f$ müssen also paarweise zusammenfallen und
$A$ ist eine Doppeltangente. Nun haben wir aber gesehen, dass die 28
Linien $T$ Doppeltangenten sind, und eine Curve vierter Ordnung kann
nicht mehr als 28 Doppeltangenten haben. Daher ist die Annahme
unzulässig, und es folgt, dass die Schnittpunkte der Grundcurve mit
den Hauptaxen zweizählige Pole sind.

Bezeichnen wir also die besonderen zweizähligen Pole, die auf der
Curve $f$ liegen, mit $\bar P_2$, so haben wir folgende ausgezeichnete
Punktsysteme:
%CONT>>


% ===== scan leaf 480, printed 459 =====
% !PRINTED 459
\weberpage{459}
\pdfbookmark[2]{§ 123. Die höheren Invarianten}{bm140}
%<<CONT
\[
24\,P_7,\ 56\,P_3,\ 84\,\bar P_2 \ \text{auf der Curve } f.
\]
\[
21\,P_8,\ 42\,P_4,\ 28\,P_6 \ \text{nicht auf der Curve } f.
\]
Alle anderen Punkte der Curve $f$ gehen durch die Substitutionen der
Gruppe in 168 verschiedene Punkte über.

Daraus geht noch unmittelbar hervor, dass die Curve $f$ keinen
Doppelpunkt haben und also um so weniger in Curven niedrigeren Grades
zerfallen kann. Denn angenommen, sie hätte einen Doppelpunkt, so
müssten auch alle seine Bildpunkte Doppelpunkte sein, und weil eine
Curve vierter Ordnung, auch wenn sie zerfällt, nicht mehr als sechs
Doppelpunkte haben kann, wenn sie nicht unendlich viele Doppelpunkte,
d.\,h. doppelt gezählte Curventheile hat, so müsste $f$ das Quadrat
einer quadratischen Form sein. Die Wurzel aus $f$ müsste dann auch
eine Invariante sein, während es doch keine quadratischen Invarianten
giebt.

Man kann übrigens auch leicht direct zeigen, dass die Grundcurve
keinen Doppelpunkt hat; denn die Bedingungen für einen Doppelpunkt
sind:
\[
f'(x_1) = 3x_1^2x_3 + x_2^3 = 0,
\]
\[
f'(x_2) = 3x_2^2x_1 + x_3^3 = 0,
\]
\[
f'(x_3) = 3x_3^2x_2 + x_1^3 = 0,
\]
und diese Gleichungen können nicht anders erfüllt sein, als wenn
$x_1, x_2, x_3$ verschwinden (\S.\,86).

Der Kürze wegen wollen wir ein System von Punkten, deren jeder in
jeden anderen durch Substitutionen der Gruppe transformirbar ist, ein
\gesperrt{System verbundener Punkte} nennen.

\S.\,123.

\gesperrt{Die höheren Invarianten.}

Weitere Invarianten unserer Gruppe lassen sich nach dem Satze 4.,
\S.\,40 ableiten, indem wir Covarianten der Form $f$ bilden. Wir
fassen die in \S.\,87 allgemein besprochenen Covarianten ins Auge, und
bilden zunächst die \gesperrt{Hesse}'sche Covariante
\weq{1}{\Delta = \frac{1}{54}
\begin{vmatrix}
f''(x_1,x_1), & f''(x_1,x_2), & f''(x_1,x_3) \\
f''(x_2,x_1), & f''(x_2,x_2), & f''(x_2,x_3) \\
f''(x_3,x_1), & f''(x_3,x_2), & f''(x_3,x_3)
\end{vmatrix},}
%CONT>>


% ===== scan leaf 481, printed 460 =====
% !PRINTED 460
\weberpage{460}
%<<CONT
die entwickelt die Form erhält
\weq{2}{\Delta = 5x_1^2x_2^2x_3^2 - x_1x_3^5 - x_2x_1^5 - x_3x_2^5.}
Auf der Curve $\Delta$ liegen also die Ecken des Coordinatendreiecks,
und folglich liegen alle siebenzähligen Pole $P_7$ auf $\Delta$ und
bilden das vollständige System der Schnittpunkte von $f$ und $\Delta$.

Eine weitere Covariante, und zwar vom 14$^{\mathrm{ten}}$ Grade,
erhalten wir aus \S.\,87, (3), nämlich die Determinante:
\weq{3}{C = \frac{1}{9}
\begin{vmatrix}
f''(x_1,x_1), & f''(x_1,x_2), & f''(x_1,x_3), & \Delta'(x_1) \\
f''(x_2,x_1), & f''(x_2,x_2), & f''(x_2,x_3), & \Delta'(x_2) \\
f''(x_3,x_1), & f''(x_3,x_2), & f''(x_3,x_3), & \Delta'(x_3) \\
\Delta'(x_1), & \Delta'(x_2), & \Delta'(x_3), & 0
\end{vmatrix}.}
Die Determinante $C$ lässt sich aus (3), wenn auch etwas weitläufig,
berechnen, und erhält den Ausdruck:
\weq{4}{C = \Sigma x_1^{14} - 34x_1x_2x_3\,\Sigma x_1^{10}x_3 - 250x_1x_2x_3\,\Sigma x_1^8x_2^3 + 375x_1^2x_2^2x_3^2\,\Sigma x_1^6x_2^2 + 18\,\Sigma x_1^7x_2^7 - 126x_1^3x_2^3x_3^3\,\Sigma x_1^3x_2^2,}
worin das Zeichen $\Sigma$ bedeutet, dass die Summe der drei Glieder
genommen werden soll, die man aus dem ersten erhält, wenn man
$x_1, x_2, x_3$ cyklisch vertauscht. Es genügt schon die Berechnung
des ersten Gliedes $x_1^{14}$, um zu erkennen, dass die Curve $C = 0$
\gesperrt{nicht} durch die Ecken des Coordinatendreiecks geht.

Eine vierte Invariante, und zwar vom 21$^{\mathrm{sten}}$ Grade,
erhalten wir nach \S.\,87, (6), wenn wir die Functionaldeterminante
der drei Formen $f$, $\Delta$, $C$ bilden:
\weq{5}{K = \frac{1}{14}
\begin{vmatrix}
f'(x_1), & \Delta'(x_1), & C'(x_1) \\
f'(x_2), & \Delta'(x_2), & C'(x_2) \\
f'(x_3), & \Delta'(x_3), & C'(x_3)
\end{vmatrix},}
die gleichfalls hieraus berechnet werden kann. Wir führen hier nur die
drei ersten Glieder an, aus denen man sieht, dass auch diese Curve
nicht durch die Ecken des Coordinatendreiecks geht
\weq{6}{K = x_1^{21} + x_2^{21} + x_3^{21} + \cdots{}^{1)}.}

$^{1)}$ Die vollständig ausgerechneten Ausdrücke finden sich in der
Abhandlung von \gesperrt{Gordan}: \glqq Ueber die typische Darstellung
der ternären biquadratischen Form $f = x_1^3x_2 + x_2^3x_3 + x_3^3x_1$\grqq{}
[Mathem. Annalen, Bd. XVII, S. 366 (1880)]. Auch in \gesperrt{Klein-Fricke},
Modulfunctionen, Bd. I, S. 734.


% ===== scan leaf 482, printed 461 =====
% !PRINTED 461
\weberpage{461}
\pdfbookmark[2]{§ 124. Das volle Invariantensystem}{bm141}
\begin{center}
§.\,124.

\gesperrt{Das volle Invariantensystem.}
\end{center}

Wir können nun nachweisen, dass die Formen $f, \Delta, C, K$ ein
volles Invariantensystem der Gruppe sind, d.\,h. dass alle Invarianten
der Gruppe als ganze rationale Functionen von diesen vier dargestellt
werden können. Wir stützen uns dabei auf das Theorem von
\gesperrt{Bezout} (Bd.\,I, §.\,49), dass zwei Curven von der $m^{\text{ten}}$
und $n^{\text{ten}}$ Ordnung, die mehr als $mn$ Punkte gemeinsam haben,
einen gemeinsamen Curventheil haben müssen. Berührungspunkte sind
dabei als doppelt oder mehrfach zu zählen. Alle Schnittpunkte zweier
Invariantencurven bilden entweder ein verbundenes System, oder sie
zerfallen in Systeme verbundener Punkte, und alle Punkte eines
solchen Systemes sind gleich oft zu zählen.

Es sei $\Phi$ eine Invariante $m^{\text{ter}}$ Ordnung, die die Function
$f$ nicht als Factor enthält, und unter den Schnittpunkten der Curve
$\Phi$ und $f$ mögen $h_1$ mal die Pole $P_7$, $h_2$ mal die Pole $P_3$,
$h_3$ mal die Pole $\overline{P}_2$ vorkommen. Ausserdem sollen noch $h_4$
Systeme von je 168 verbundenen Punkten vorkommen, von denen auch
(bei Berührung) mehrere Systeme in ein mehrfach gezähltes zusammenfallen
können. Dann ist die Anzahl aller Schnittpunkte beider Curven
\[
4m = 24h_1 + 56h_2 + 84h_3 + 168h_4,
\]
und daher
\weq{1}{m = 6h_1 + 14h_2 + 21h_3 + 42h_4,}
worin $h_1, h_2, h_3, h_4$ nicht negative ganze Zahlen bedeuten, die
auch nicht alle verschwinden können. Der kleinste Werth, den $m$
haben kann, ist daher 6, und eine Invariantencurve $6^{\text{ter}}$
Ordnung muss durch die 24 Punkte $P_7$ gehen, wie wir es von $\Delta$
schon nachgewiesen haben. Ist $\Delta'$ eine zweite Invariante
$6^{\text{ter}}$ Ordnung, die also auch durch die Punkte $P_7$ gehen
muss, so können wir in $\Delta' - a\Delta$ die Constante $a$ so
bestimmen, dass die Invariantencurve $\Delta' - a\Delta = 0$ durch
irgend einen $25^{\text{sten}}$ Punkt von $f$ geht. Dann muss aber, wenn
$\Delta' - a\Delta$ nicht identisch verschwindet, nach dem
\gesperrt{Bezout}'schen Theorem $\Delta' - a\Delta$ durch $f$ theilbar
sein. Der Quotient wäre eine Invariante zweiter Ordnung, die nicht
existirt; folglich muss $\Delta' - a\Delta$ identisch Null sein.


% ===== scan leaf 483, printed 462 =====
% !PRINTED 462
\weberpage{462}
Ist $h_1 = 2$, $h_2 = h_3 = h_4 = 0$, so ist $m = 12$. Eine
Invariantencurve $12^{\text{ter}}$ Ordnung $\Phi$ muss die Curve $f$ in
den 24 Punkten $P_7$ berühren, und wenn wir in $\Phi - a\Delta^2$ die
Constante $a$ passend bestimmen, so ergiebt sich, dass diese Function
durch $f$ theilbar sein muss. Der Quotient kann, als Invariante
$8^{\text{ter}}$ Ordnung, nur von der Form $bf^2$ sein, und demnach ist
$\Phi$ von der Form $a\Delta^2 + bf^3$. Ebenso können wir schliessen,
dass eine Invariante $18^{\text{ter}}$ Ordnung die Form
$a\Delta^3 + b\Delta f^3$ haben muss.

\begin{center}
\gesperrt{Alle Invarianten $6^{\text{ter}}$, $12^{\text{ter}}$, $18^{\text{ter}}$
Ordnung sind also rationale Functionen von $\Delta$ und $f$.}
\end{center}

Der nächste Werth, den $m$ nach (1) haben kann, ist $m = 14$. In
diesem Falle ist $h_2 = 1$, während $h_1, h_3, h_4$ gleich Null sind.
Es gehen also alle Invarianten $14^{\text{ter}}$ Ordnung durch die 56
Punkte $P_3$, und diese bilden das vollständige Schnittpunktsystem
einer solchen Invariantencurve mit der Grundcurve. Dies gilt auch von
der Invariante $C$. Haben wir eine zweite Invariante $14^{\text{ter}}$
Ordnung $C'$, so können wir wieder, wie oben, die Constante $a$ so
bestimmen, dass $C' - aC$ durch $f$ theilbar ist. Der Quotient ist
eine Invariante $10^{\text{ter}}$ Ordnung, und daraus folgt, da es
ausser $f\Delta$ keine Invariante $10^{\text{ter}}$ Ordnung giebt:

\begin{center}
\gesperrt{Jede Invariante $14^{\text{ter}}$ Ordnung ist in der Form
darstellbar:}
\end{center}
\[
aC + bf^2\Delta,
\]
\begin{center}
\gesperrt{worin $a, b$ Constanten sind. Umgekehrt ist jeder Ausdruck von
dieser Form eine Invariante $14^{\text{ter}}$ Ordnung.}
\end{center}

Es kann sodann $m$ nach (1) den Werth 20 haben, nämlich für
$h_1 = h_2 = 1$. Eine Invariante $20^{\text{ster}}$ Ordnung muss also
durch die Punkte $P_7$ und $P_3$ gehen, d.\,h. durch die 80
Schnittpunkte von $f$ mit $\Delta C$. Daraus können wir ebenso wie
vorhin schliessen, dass eine Invariante $20^{\text{sten}}$ Grades in
der Form
\[
f(a\Delta^2 + bf^3) + cC\Delta.
\]
darstellbar ist, worin $a, b, c$ beliebige Constanten sind. Eine
unabhängige Invariante $20^{\text{ster}}$ Ordnung giebt es nicht.

Nehmen wir nun an, es seien $K'$ und $K$ zwei Invarianten
$21^{\text{ster}}$ Ordnung; beide müssen nach (1) durch die 84 Punkte
$\overline{P}_2$ gehen, und folglich kann man $a$ so bestimmen, dass
$K' - aK$ durch $f$ theilbar wird. Der Quotient wäre eine Invariante
%CONT>>


% ===== scan leaf 484, printed 463 =====
% !PRINTED 463
\weberpage{463}
%<<CONT
$17^{\text{ten}}$ Grades, die nicht existirt, und folglich ist $K'$ mit
$aK$ identisch.

\begin{center}
\gesperrt{Es giebt also, von einem constanten Factor abgesehen, nur
eine Invariante $21^{\text{sten}}$ Grades.}
\end{center}

Nun ist aber das System der 21 Hauptaxen auch eine Invariante
$21^{\text{sten}}$ Grades, und daraus ist zu schliessen:

\begin{center}
\gesperrt{Die Invariante $K$ zerfällt in 21 lineare Factoren, die,
gleich Null gesetzt, die Hauptaxen der Gruppe darstellen.}
\end{center}

Wir können sodann eine Invariante $42^{\text{sten}}$ Grades bilden,
nämlich
\weq{2}{\Delta^7 - kC^3 = Q,}
worin $k$ eine beliebige Constante ist, und diese Constante lässt
sich so bestimmen, dass die Curve $Q$ durch einen beliebig gegebenen
Punkt auf $f$ geht, und sie muss dann auch durch alle mit diesem
verbundenen Punkte hindurchgehen. Also können wir $k$ in $Q$ so
bestimmen, dass die Curve $Q$ aus der Curve $f$ ein beliebig
gegebenes System verbundener Punkte ausschneidet.

Ist nun $\Phi$ eine beliebige durch $f$ nicht theilbare Invariante,
die $h_1, h_2, h_3$ mal durch die Pole $P_7, P_3, \overline{P}_2$
geht, und ausserdem durch beliebige Systeme $S_1, S_2, \ldots$
verbundener Punkte auf $f$, die auch theilweise zusammenfallen
können, so bilden wir zunächst nach (2) die Formen
$Q_1, Q_2, \ldots$, die in den Systemen $S_1, S_2, \ldots$
verschwinden, und dann die Form
\[
\Psi = \Phi - a\Delta^{h_1} C^{h_2} K^{h_3} Q_1 Q_2 \ldots
\]

Die Curve $\Psi$ geht für jeden Werth der Constanten $a$ durch die
sämmtlichen Schnittpunkte von $\Phi$ mit $f$, und wenn wir also $a$ so
bestimmen, dass $\Psi$ durch irgend einen davon verschiedenen Punkt
von $f$ geht, so ist $\Psi$ durch $f$ theilbar, also
\[
\Phi = a\Delta^{h_1} C^{h_2} K^{h_3} Q_1 Q_2 \cdots + f\Phi_1.
\]
Darin ist nun $\Phi_1$ wieder eine Invariante, aber von niedrigerem
Grade als $\Phi$, und durch vollständige Induction ist hiermit der
Satz bewiesen:

\begin{center}
\gesperrt{Jede Invariante der Gruppe lässt sich als ganze rationale
Function der vier fundamentalen Invarianten $f, \Delta, C, K$
darstellen.}
\end{center}

Bestimmt man in (2) die Constante $k$ so, dass die Curve $Q$ durch
einen der Pole $\overline{P}_2$ geht, so kann sie durch keinen nicht
%CONT>>


% ===== scan leaf 485, printed 464 =====
% !PRINTED 464
\weberpage{464}
%<<CONT
mit $\overline{P}_2$ verbundenen Punkt der Curve $f$ gehen; denn sie
kann nicht durch die Punkte $P_7, P_3$ gehen, weil in diesen entweder
$\Delta$ oder $C$ verschwindet, sie kann aber auch nicht durch einen
Punkt der Grundcurve gehen, der kein Pol ist, weil sie sonst durch
die 168 verbundenen Punkte gehen müsste, und nicht durch den Pol
$\overline{P}_2$ gehen könnte. Dann kann man aber die Constante $h$
so bestimmen, dass $K^2 - hQ$ durch $f$ theilbar wird.

Der Quotient ist eine Invariante von niedrigerem als dem
$42^{\text{sten}}$, jedenfalls aber von geradem Grade, und wenn wir
ihn also durch die fundamentalen Invarianten darstellen, so kann
darin $K$ nicht vorkommen. Daraus folgt:

\begin{center}
\gesperrt{Die Invariante $K^2$ kann rational durch $f, \Delta, C$
ausgedrückt werden.}
\end{center}

Stellen wir nach diesem Satze $K^2$ in der Form dar:
\weq{3}{K^2 = \Sigma\, af^{\nu}\Delta^{\nu_1}C^{\nu_2},}
so können in dieser Summe, in der die $a$ numerische
Co\"efficienten sind, nur solche Glieder vorkommen, in denen
\[
2\nu + 3\nu_1 + 7\nu_2 = 21,
\]
und indem wir nun abzählen, welche Werthe von $\nu, \nu_1, \nu_2$
vorkommen können, erhalten wir das Resultat, dass zwischen den
Formen
\[
\begin{gathered}
K^2,\quad \Delta^7,\quad C^3,\quad f\Delta^4 C,\quad f^2\Delta C^2,\quad f^3\Delta^5,\\
f^4\Delta^2 C,\quad f^6\Delta^3,\quad f^7 C,\quad f^9\Delta
\end{gathered}
\]
eine lineare Relation mit numerischen Co\"efficienten besteht.
Stellen wir diese Relation in der Form dar:
\weq{4}{K^2 = \Phi\Delta + \Psi C,}
worin $\Phi, \Psi$ gleichfalls Invarianten sind, so können wir daraus
noch einen geometrischen Schluss ziehen.

Die Curven $\Delta$ und $C$ schneiden sich sicher nicht auf der Curve
$f$, weil die sämmtlichen Schnittpunkte von $\Delta$ mit $f$ die
$P_7$, die von $C$ mit $f$ die $P_3$ sind. Wenn aber $\Delta$ und
$C = 0$ sind, so ist auch $K = 0$, und folglich liegen
\gesperrt{alle Schnittpunkte von $\Delta$ und $C$ auf den 21 Hauptaxen der
Gruppe}\footnote{Die Relation (3) ist von \gesperrt{Gordan} durch die
Methoden der Invariantentheorie vollständig berechnet (Mathem.\
Annalen, Bd.\ XVII, S.\,371).

Wir wollen die dort gegebene Formel in
den von uns gebrauchten Zeichen hier angeben:
\[
K^2 = C^3 - 88f^2\Delta C^2 + 16.63f\Delta^4 C + 17.64f^4\Delta^2 C
- 256f^7 C + 27.64\Delta^7 - 128.469f^3\Delta^5 + 43.512f^6\Delta^3
- 2048f^9\Delta.
\]}.


% ===== scan leaf 486, printed 465 =====
% !PRINTED 465
\weberpage{465}

Die Relation (3) l\"asst sich benutzen, um aus einem Ausdrucke in den Invarianten alle Potenzen von $K$, mit Ausnahme der ersten, zu eliminiren, und daraus ergiebt sich noch:

\begin{quote}
\gesperrt{Eine Invariante geraden Grades l\"asst sich rational durch $f, \Delta, C$, eine Invariante ungeraden Grades als Product von $K$ mit einer Invariante geraden Grades darstellen.}
\end{quote}

\begin{center}
\rule{4cm}{0.4pt}
\end{center}


% ===== scan leaf 487, printed 466 =====
\weberpage{466}% page number not printed on this page
\pdfbookmark[1]{Fünfzehnter Abschnitt. Das Formenproblem der Gruppe G168 und die Theorie der Gleichungen siebenten Grades.}{bm142}
\pdfbookmark[2]{§ 125. Die Resolventen des Formenproblems}{bm143}
% !PRINTED none

\begin{center}
\gesperrt{Fünfzehnter Abschnitt.}

\textbf{Das Formenproblem der Gruppe $G_{168}$ und die Theorie der Gleichungen siebenten Grades.}

\rule{4cm}{0.4pt}

\S.\,125.

\gesperrt{Die Resolventen des Formenproblems.}
\end{center}

Das Formenproblem für die Gruppe $G_{168}$ (\S.\,43) liefert nach der allgemeinen Theorie zunächst eine Gleichung 168\textsuperscript{sten} Grades, deren Coëfficienten Invarianten sind. Jeder Theiler der Gruppe führt aber zu einer Resolvente niedrigeren Grades, und zwar vom Grade des Index des Theilers.

Wir wollen hier nur die beiden interessantesten Fälle solcher Theiler, nämlich die Octaëdergruppe

\weq{1}{G_{24} = \chi^\lambda \omega^\mu \Theta^\nu,}

und die Gruppe 21\textsuperscript{sten} Grades

\weq{2}{G_{21} = \tau^\varrho \chi^\lambda \ (\S.\,72)}

betrachten, die uns zu Resolventen 7\textsuperscript{ten} und 8\textsuperscript{ten} Grades führen.

Was zunächst die Resolventen 7\textsuperscript{ten} Grades anlangt, so können wir bei ihrer Bildung von den Hauptaxen der Gruppe ausgehen. Eine Hauptaxe nämlich bleibt durch die Gruppe $\omega^\mu \Theta^\nu$ ungeändert, und geht durch die ganze Gruppe $G_{24}$ in drei verschiedene Linien über. Es muss also sieben Tripel von Hauptaxen geben, die durch eine Gleichung 7\textsuperscript{ten} Grades bestimmt werden.

Setzen wir nach \S.\,118, (7):

\weq{3}{\begin{aligned}
A_1 &= \alpha\gamma x_1 + \gamma\beta\varepsilon^3 x_2 + \alpha\beta\varepsilon^2 x_3 \\
A_2 &= \beta\gamma\varepsilon^3 x_1 + \alpha\beta\varepsilon^2 x_2 + \alpha\gamma x_3 \\
A_3 &= \alpha\beta\varepsilon^2 x_1 + \alpha\gamma x_2 + \beta\gamma\varepsilon^3 x_3,
\end{aligned}}
%CONT>>


% ===== scan leaf 488, printed 467 =====
% !PRINTED 467
\weberpage{467}

%<<CONT
so sind $A_1=0$, $A_2=0$, $A_3=0$ die Gleichungen von dreien dieser Axen. Durch die Substitution $\chi$ gehen $A_1, A_2, A_3$ cyklisch in einander über. Durch die Substitution $\Theta$ erleiden die Functionen $A_1, A_2, A_3$ folgende Vertauschung:

\[
\begin{pmatrix} A_1, & A_2, & A_3 \\ A_1, & -A_3, & A_2 \end{pmatrix},
\]

wie eine einfache Rechnung zeigt, wenn man die Substitution $\Theta$ wirklich ausführt und die Formeln des \S.\,115 benutzt.

Um den Gang der Rechnung wenigstens anzudeuten, sei bemerkt, dass $A_1$ durch die Substitution $\Theta$ [\S.\,115, (10)] in

\[
\alpha(\gamma^2\varepsilon^2 + \beta\gamma\varepsilon^3 + \beta^2\varepsilon^5)x_1 + \gamma(\alpha^2\varepsilon^6 + \beta^2 + \alpha\beta\varepsilon^2)x_2 + \beta(\alpha\gamma + \gamma^2\varepsilon + \alpha^2\varepsilon^3)x_3
\]

übergeht. Es ist aber nach \S.\,115, (8), (17), (11):

\[
\begin{aligned}
\gamma^2\varepsilon^2 + \beta\gamma\varepsilon^3 + \beta^2\varepsilon^5 &= \gamma^2\varepsilon^2 + \alpha^2\varepsilon^3 + \beta^2\varepsilon^5 + \alpha\varepsilon^3 \\
&= \alpha(\varepsilon^4+\varepsilon^{-4}) = h(\varepsilon-\varepsilon^{-1}) = \gamma,
\end{aligned}
\]

und daher wird der Coëfficient von $x_1$ gleich $\alpha\gamma$, wie in $A_1$, und ebenso formt man die übrigen Ausdrücke um. Da sich nun $\omega$ aus $\Theta$ und $\chi$ zusammensetzen lässt (\S.\,72), so folgt, dass die Grössen $A_1^2, A_2^2, A_3^2$ durch die Gruppe $G_{24}$ nur unter einander permutirt werden, und dass demnach ihre symmetrischen Functionen Wurzeln von Resolventen 7\textsuperscript{ten} Grades sind.

Die einfachste symmetrische Function dieser Grössen ist die Summe

\weq{4}{A_1^2 + A_2^2 + A_3^2,}

und diese wollen wir, mit einem geeigneten numerischen Factor multiplicirt, als die Unbekannte der Resolvente einführen.

Ordnen wir die Summe (4) nach $x_1, x_2, x_3$, so erhalten wir dafür einen Ausdruck von der Form

\[
\lambda(x_1^2+x_2^2+x_3^2) + \mu(x_2x_3+x_3x_1+x_1x_2),
\]

worin nach (3):

\[
\begin{aligned}
\lambda &= (\alpha^2\gamma^2\varepsilon^{-1} + \beta^2\gamma^2\varepsilon^{-2} + \alpha^2\beta^2\varepsilon^3)\varepsilon \\
\mu &= 2\alpha\beta\gamma(\alpha\varepsilon+\beta\varepsilon^{-3}+\gamma\varepsilon^2)\varepsilon.
\end{aligned}
\]

Hierbei ist der Factor $\varepsilon$ aus der Klammer gezogen, damit der andere Factor durch die Vertauschung von $\varepsilon$ und $\varepsilon^2$ ungeändert bleibt, und sich daher rational durch $\sqrt{-7}$ ausdrücken l\"asst (Bd. I, \S.\,171).


% ===== scan leaf 489, printed 468 =====
% !PRINTED 468
\weberpage{468}

Nach den Formeln \S.\,115, (14) ergiebt sich zunächst sehr einfach:

\[
\mu = \frac{2\varepsilon}{7},
\]

und für $\lambda$ erhält man, wenn man die Werthe \S.\,115, (11) für $\alpha, \beta, \gamma$ einsetzt, und die Formeln \S.\,115, (12), (13) benutzt:

\weq{5}{\begin{aligned}
\lambda &= -\varepsilon\,\frac{1+\sqrt{-7}}{14} \\
\frac{\mu}{\lambda} &= -\frac{1-\sqrt{-7}}{2}.
\end{aligned}}

Setzt man daher

\[
A_1^2+A_2^2+A_3^2 = \lambda z,
\]

so ergiebt sich

\weq{6}{z = x_1^2+x_2^2+x_3^2 - \frac{1-\sqrt{-7}}{2}(x_2x_3+x_3x_1+x_1x_2),}

und diese Function wollen wir als Wurzel der Resolvente 7\textsuperscript{ten} Grades einführen. Die übrigen Wurzeln erhält man daraus durch die Substitutionen $\tau^r$, so dass sie alle in der gemeinschaftlichen Form

\weq{7}{\begin{aligned}
z_r &= \varepsilon^{2r}x_1^2 + \varepsilon^{4r}x_2^2 + \varepsilon^{r}x_3^2 \\
&\quad - \frac{1-\sqrt{-7}}{2}(\varepsilon^{-r}x_2x_3+\varepsilon^{-2r}x_3x_1+\varepsilon^{-4r}x_1x_2) \\
&\qquad r=0, 1, 2, 3, 4, 5, 6
\end{aligned}}

enthalten sind.

Die Coëfficienten der Gleichung 7\textsuperscript{ten} Grades, deren Wurzeln die sieben Grössen $z_r$ sind, sind Invarianten unserer Gruppe, deren Grad sich leicht angeben l\"asst.

Wenn nämlich $a_\nu$ der Coëfficient von $z^{7-\nu}$ in dieser Gleichung ist, nachdem der Coëfficient der siebenten Potenz auf 1 gebracht ist, so ist $a_\nu$ eine Invariante $2\nu^{\text{ten}}$ Grades. Es muss also zunächst $a_1=0$ sein, weil es keine quadratische Invariante giebt. Die übrigen Coëfficienten sind durch folgende Invariantenverbindungen linear und homogen ausgedrückt:

\weq{8}{\begin{aligned}
a_2 &\text{ durch } f, \\
a_3 &\text{ \textquotedbl} \ \varDelta, \\
a_4 &\text{ \textquotedbl} \ f^2, \\
a_5 &\text{ \textquotedbl} \ \varDelta f, \\
a_6 &\text{ \textquotedbl} \ \varDelta^2, f^3, \\
a_7 &\text{ \textquotedbl} \ C, \varDelta f^2
\end{aligned}}


% ===== scan leaf 490, printed 469 =====
% !PRINTED 469
\weberpage{469}

und es sind also acht numerische Co\"efficienten zu berechnen, die
sich durch Vergleichung einiger Glieder in den Ausdr\"ucken der
$a_\nu$ durch die Wurzeln einerseits, durch die Invarianten (\S.\,123)
andererseits finden lassen, und die ausser rationalen Zahlen nur
$\sqrt{-7}$ enthalten k\"onnen.

Wir wollen diese Co\"efficienten zun\"achst nur unter der
Voraussetzung berechnen, dass $f=0$ ist ${}^1$).

Man braucht dann nur die ersten Glieder (mit $x_1^5x_2$, $x_1^{10}x_2^2$)
in den Potenzsummen $\Sigma z_r^3$ und $\Sigma z_r^6$ zu berechnen und die
\gesperrt{Newton}'schen Formeln anzuwenden. Der letzte Co\"efficient ergiebt
sich aus dem einem Gliede $x_1^{14}$ in dem Producte der $z_r$. Man
findet zun\"achst
\[
\Sigma z_r^3 = -3\cdot7\cdot\frac{1-\sqrt{-7}}{2}\,x_1^5x_2\dots
\]
\[
\Sigma z_r^6 = \left[6\cdot7 + 15\cdot7\left(\frac{1-\sqrt{-7}}{2}\right)^2\right]x_1^{10}x_2^2\dots
\]
und daraus die gesuchte Resolvente
\weq{9}{z^7 - 7\cdot\frac{1-\sqrt{-7}}{2}\,\Delta z^4 - 7\cdot\frac{5+\sqrt{-7}}{2}\,\Delta^2 z - C = 0.}
Will man diese Gleichung auf eine andere zur\"uckf\"uhren, die
nur von den Verh\"altnissen der $x_1, x_2, x_3$ abh\"angt, so setzt man
\weq{10}{z = \frac{Cu}{\Delta^2}, \quad g = \frac{\Delta^7}{C^3},}
wodurch man aus (9) eine Gleichung erh\"alt, in der nur noch der
eine Parameter $g$ vorkommt, n\"amlich
\weq{11}{u^7 - 7\,\frac{1-\sqrt{-7}}{2}\,g\,u^4 - 7\,\frac{5+\sqrt{-7}}{2}\,g^2\,u - g^2 = 0.}
Macht man dieselbe Substitution in der allgemeinen Resol\-vente,
in der $f$ nicht $=0$ gesetzt ist, so hat man noch einen
weiteren Parameter
\weq{12}{h = \frac{f\Delta^4}{C^2}}
einzuf\"uhren, und die Co\"efficienten der Resolvente werden, von
%CONT>>

\par\noindent\rule{0.3\linewidth}{0.4pt}\par
${}^1$) Dieser Fall ist darum von besonderem Interesse, weil er, \"ahnlich
wie die Ikosa\"edergleichung, auf die Transformationsgleichungen aus der
Theorie der elliptischen Functionen f\"uhrt.


% ===== scan leaf 491, printed 470 =====
% !PRINTED 470
\weberpage{470}

%<<CONT
numerischen Factoren abgesehen, wie sich aus (8) leicht ergiebt,
der Reihe nach
\[
h,\ g,\ h^2,\ hg,\ (g^2,h^3),\ (g^2,h^2g),
\]
worin $(g^2,h^3)$, $(g^2,h^2g)$ lineare homogene Ausdr\"ucke mit numeri\-schen
Co\"efficienten bedeuten.

Die Rechnung kann ebenso ausgef\"uhrt werden, wie in dem
obigen speciellen Falle. Zur Vereinfachung kann man $x_3=0$
setzen und erh\"alt immer noch Gleichungen genug zur Be\-stimmung
aller Co\"efficienten. Man findet so die Co\"efficienten
der Reihe nach:
\[
7\,\frac{1-\sqrt{-7}}{2}\,h,
\]
\[
-7\,\frac{1-\sqrt{-7}}{2}\,g,
\]
\[
-7\,(4+\sqrt{-7})\,h^2,
\]
\[
14\,(2+\sqrt{-7})\,hg,
\]
\[
-7\,\frac{5+\sqrt{-7}}{2}\,g^2 - 7\cdot\frac{7+3\sqrt{-7}}{2}\,h^3,
\]
\[
-g^2 + \frac{167-7\sqrt{-7}}{2}\,gh^2.
\]
Die Functionen $g, h$ sind gebrochene Invarianten, die nur
von den Verh\"altnissen $x_1:x_2:x_3$ abh\"angen, und wir k\"onnen
jede andere Invariante von derselben Eigenschaft
\gesperrt{rational} durch $g, h$ ausdr\"ucken. Denn stellen wir eine
solche Invariante als Quotienten zweier Formen gleichen Grades
ohne gemeinsamen Theiler dar, so m\"ussen Z\"ahler und Nenner
ganze Invarianten gleichen Grades sein, weil n\"amlich zwei in
einfachster Form dargestellte gebrochene Functionen der
Variablen $x$ nur dann einander gleich sein k\"onnen, wenn Z\"ahler
und Nenner einzeln bis auf constante Factoren einander gleich
sind. Hier k\"onnen nun Z\"ahler und Nenner nicht von ungeradem
Grade sein, weil sie sonst den gemeinsamen Factor $K$ h\"atten
(\S.\,126). Also sind Z\"ahler und Nenner rational durch $f, C, \Delta$
darstellbar; und wenn ein im Z\"ahler oder im Nenner vorkommendes
Glied
\weq{13}{f^a\,\Delta^b\,C^c,}
und $m$ der Grad von Z\"ahler und Nenner ist, so ist
\weq{14}{4a + 6b + 14c = m.}


% ===== scan leaf 492, printed 471 =====
% !PRINTED 471
\weberpage{471}
\pdfbookmark[2]{§ 126. Reduction der allgemeinen Resolvente siebenten Grades auf die specielle}{bm144}
Setzen wir nun in (13)

\[
\varDelta = (g\,C^3)^{\frac{1}{7}}, \quad f = h\,C^2\,(g\,C^3)^{-\frac{4}{7}},
\]

so ergiebt sich für den Ausdruck (13) mit Benutzung von (14)

\[
g^{b+2c}\,h^a\,C^{\frac{m}{14}}\,g^{-\frac{m}{7}},
\]

und im Zähler und Nenner lässt sich der Factor $C^{\frac{m}{14}}
g^{-\frac{m}{7}}$ heben, so dass alles rational durch $g$ und $h$
ausgedrückt ist.

Ebenso wie $g, h$ hängt auch die Function $u$ nur von dem Verhältniss
der Variablen $x_1 : x_2 : x_3$ ab, und es folgt leicht, dass jede
andere Function von derselben Eigenschaft, die wie $u$ die
Substitutionen der Gruppe $G_{24}$ gestattet, rational durch $u, g,
h$ darstellbar ist. Denn eine solche Function lässt sich zunächst
nach den allgemeinen Sätzen des \S.\,43 als rationale Function von
$u$ und den Invarianten darstellen, und zwar nur auf eine Weise als
ganze Function von $u$, die den $6^{\text{ten}}$ Grad nicht
übersteigt (Bd.\,I, \S.\,142). Die Co\"efficienten in dieser
Darstellung sind Invarianten, und da $u$ nur von den Verhältnissen
abhängt, so können auch die Co\"efficienten nur von den
Verhältnissen abhängen, und sind daher rational durch $g, h$
darstellbar.

\begin{center}
\S.\ 126.

\gesperrt{Reduction der allgemeinen Resolvente siebenten Grades auf
die specielle.}
\end{center}

Wir haben im vorigen Paragraphen zwei Formen der Resolvente
$7^{\text{ten}}$ Grades des Formenproblems kennen gelernt, von denen
die eine, die specielle, für den Fall gilt, dass $f = 0$ ist, d.\,h.
für den Fall, dass der gesuchte Punkt auf der Grundcurve liegt, die
allgemeine für den Fall, dass er eine beliebige Lage hat.

Die Grössen $u, g, h$ wollen wir, wenn sie sich auf den Punkt $(x)$
beziehen, mit

\weq{1}{u_x = \dfrac{\varDelta^2 z}{C}, \quad g_x = \dfrac{\varDelta^7}{C^3}, \quad h_x = \dfrac{f\varDelta^4}{C^2}}

bezeichnen. Die allgemeine Resolvente soll dann mit

\weq{2}{R(u_x, g_x, h_x) = R_x = 0}

bezeichnet werden, und die specielle geht daraus hervor, wenn
%CONT>>


% ===== scan leaf 493, printed 472 =====
% !PRINTED 472
\weberpage{472}

%<<CONT
man $h_x = 0$ setzt. Die Gr\"ossen $u_x, g_x, h_x$ h\"angen nur von den
Variablen $x_1, x_2, x_3$ ab.

Nun l\"asst sich, wie \gesperrt{Klein} a.\,a.\,O. ${}^1$) bemerkt hat, die L\"osung
der allgemeinen Resolvente auf die der speciellen zur\"uckf\"uhren,
wenn man die Wurzel einer \gesperrt{biquadratischen Gleichung}
adjungirt.

Um dies nachzuweisen, f\"uhren wir neben dem Punkte $(x)$
einen zweiten Punkt $(y)$ ein und bilden die Polare von $f$:
\weq{3}{f_1(x,y) = y_1f'(x_1) + y_2f'(x_2) + y_3f'(x_3).}
Wenn wir dann $(x)$ und $(y)$ gleichzeitig derselben Substitu\-tion
der Gruppe $G_{168}$ unterwerfen, so bleibt die Function $f_1(x,y)$
unge\"andert (Bd.\ I, \S.\,60). Wir nehmen nun den Punkt $(x)$ beliebig
 an, verlangen aber von dem Punkte $(y)$, dass er gleichzeitig
 auf der Grundcurve und auf der Polaren des Punktes $(x)$
liegen soll, dass also gleichzeitig
\weq{4}{f(y_1,y_2,y_3) = 0, \quad f_1(x,y) = 0}
sein soll. Dann entsprechen jedem Punkte $x$ vier Punkte $y$, und
wenn wir f\"ur $(x)$ einen mit ihm verbundenen Punkt setzen, so
geht jeder dieser vier Punkte $(y)$ gleichfalls in einen verbundenen
Punkt \"uber. Die Function $h_y$ ist jetzt $=0$ und $g_y$ ist eine vierwerthige
 Function des Punktes $(x)$. Symmetrische Functionen
dieser vier Werthe bleiben unge\"andert durch die Substitutionen
von $G_{168}$, und folglich ist $g_y$ \gesperrt{Wurzel einer biquadratischen
Gleichung, deren Co\"efficienten rational von den $g_x, h_x$
abh\"angen}. Die Wurzel dieser biquadratischen Gleichung muss
adjungirt werden.

Die Function $u_y$ ist Wurzel der speciellen Resolvente
\weq{5}{R(u_y, g_y, 0) = R_y = 0.}
Wir bezeichnen nun die vier zu demselben $x$ geh\"origen Punkte
$y$ mit $y, y', y'', y'''$ und bilden die symmetrische Function dieser
vier Punkte
\weq{6}{(t-u_y)(t-u_{y'})(t-u_{y''})(t-u_{y'''}) = \Phi(t)}
f\"ur ein unbestimmtes $t$. Diese Function bleibt unge\"andert bei
allen Substitutionen der Gruppe $G_{24}$, und ist also rational durch
$u_x, g_x, h_x$ ausdr\"uckbar.

\par\noindent\rule{0.3\linewidth}{0.4pt}\par
${}^1$) Mathematische Annalen, Bd.\ XV, S.\ 280.


% ===== scan leaf 494, printed 473 =====
% !PRINTED 473
\weberpage{473}
\pdfbookmark[2]{§ 127. Permutationsgruppe von sieben Ziffern vom Grade 168}{bm145}
Da, wenn wir uns die Bestimmung von $x$ vorbehalten, die Gleichung (3) jede beliebige gerade Linie darstellen kann, so können wir $x$ so annehmen, dass unter den vier Punkten $y, y', y'', y'''$ keine zwei verbundenen Punkte vorkommen. Setzen wir dann $u_y$ für die unbestimmte Grösse $t$ in (6), so erhalten wir eine rationale Gleichung

\weq{7}{\Psi (u_y, u_x, g, h) = 0,}

und diese Gleichung ist nicht mehr befriedigt, wenn wir für $(x)$ und $(y)$ eine Substitution aus $G_{168}$ machen, durch die $u_x$ geändert und folglich $u_y$ in einen von $u_y, u_{y'}, u_{y''}, u_{y'''}$ verschiedenen Werth übergeführt wird.

Die Gleichung (7) hat also, als Gleichung für $u_x$ betrachtet, nur eine Wurzel mit der allgemeinen Resolvente $R_x$ gemein, und wenn man den grössten gemeinschaftlichen Theiler von $R$ und $\Psi$ aufsucht, so erhält man $u_x$ rational durch $u_y, g, h$ ausgedrückt. Damit ist die allgemeine Resolvente auf die specielle zurückgeführt.

\begin{center}\rule{0.3\linewidth}{0.4pt}\end{center}

Die Gruppe $G_{168}$ enthält noch einen Theiler 21\textsuperscript{sten} Grades $G_{21}$, der durch die Substitutionen $\chi, \tau$ erzeugt wird, und der zu einer Resolvente 8\textsuperscript{ten} Grades Anlass giebt. Als Wurzel dieser Resolvente kann man einfach das Product $x_1 x_2 x_3$ betrachten. Die Co\"efficienten dieser Resolvente sind Invarianten, die bis zum 24\textsuperscript{sten} Grade ansteigen. Wir wollen hier auf diese Resolvente nicht näher eingehen.

\begin{center}
\S.\,127.

\gesperrt{Permutationsgruppe von sieben Ziffern vom Grade 168.}
\end{center}

Da, wie wir gesehen haben, das Formenproblem der ternären Substitutionsgruppe 168\textsuperscript{sten} Grades eine Resolvente 7\textsuperscript{ten} Grades hat, und die Galois'sche Resolvente dieser Gleichung 7\textsuperscript{ten} Grades folglich vom Grade 168 ist, so ergiebt sich daraus, dass in der allgemeinen Permutationsgruppe von sieben Ziffern, deren Grad $1.2.3.4.5.6.7 = 5040 = 168.30$ ist, ein Theiler vom Grade 168 enthalten sein muss. Die Existenz dieses Theilers hat zuerst
%CONT>>


% ===== scan leaf 495, printed 474 =====
% !PRINTED 474
\weberpage{474}
%<<CONT
\gesperrt{Kronecker} erkannt\footnote{\gesperrt{Kronecker}, \glqq Ueber Gleichungen 7\textsuperscript{ten} Grades\grqq{}. Monatsberichte der Berliner Akademie 1858.}, und seine nähere Untersuchung ist für die allgemeine Theorie der Gleichung 7\textsuperscript{ten} Grades von grosser Wichtigkeit.

Wir können diese Permutationsgruppe dadurch erhalten, dass wir die Permutationen aufsuchen, die durch die Substitutionen der Gruppe $G_{168}$ unter den Grössen $z_0, z_1, z_2, z_3, z_4, z_5, z_6$ (\S.\,125) hervorgerufen werden. Hierbei ist dann in Bezug auf die Zusammensetzung zu beachten, dass, wenn die beiden linearen Substitutionen $\xi_1, \xi_2$ die Permutationen $\pi_1, \pi_2$ bewirken, die zusammengesetzte Permutation $\pi_1 \pi_2$ durch $\xi_2 \xi_1$ hervorgerufen wird. Denn nach der Definition erhält man $\xi_2 \xi_1 (x)$ dadurch, dass man auf die Variablen $(x)$ zunächst die Substitution $\xi_1$ und auf das Ergebniss die Substitution $\xi_2$ anwendet. Ebenso bedeutet $\pi_1 \pi_2$ die Permutation, die sich ergiebt, wenn man auf die sieben Ziffern zuerst $\pi_1$ und darauf $\pi_2$ anwendet. Wir erhalten so, der Gruppe $G_{168}$ entsprechend, eine Permutationsgruppe 168\textsuperscript{sten} Grades, die wir mit $P_{168}$ bezeichnen, und diese beiden Gruppen sind isomorph.

Da wir $\tau$ und $\omega$ als erzeugende Elemente der Gruppe $G_{168}$ erkannt haben (\S.\,72, I.), so genügt es, wenn wir die diesen beiden Substitutionen entsprechenden Permutationen bestimmen, um daraus die ganze Gruppe $P_{168}$ abzuleiten.

Nun geht aber aus der Substitution $\tau$ die cyklische Permutation der sieben Ziffern

\[
(\tau) \qquad (0, 1, 2, 3, 4, 5, 6)
\]

hervor, und der Einfluss von $\omega$ ergiebt sich aus der Bemerkung, dass $z_0$ durch die Substitutionen der Octa\"edergruppe $\chi^\lambda \omega^u \Theta^v$ ungeändert bleibt. Um also die Aenderung von $z_r$ durch $\omega$ zu erhalten, haben wir nur den Einfluss von $\omega \tau^r$ auf $z_0$ zu ermitteln. Dieser Einfluss aber ergiebt sich unmittelbar aus den Formeln \S.\,72, (15), wonach z. B. $\omega \tau = \tau^2 \chi^2 \omega \Theta^2$ ist, so dass also $z_1$ in $z_2$ übergeht u. s. f. Demnach entspricht der Substitution $\omega$ die Permutation

\[
(\omega) \qquad \begin{pmatrix} 0, & 1, & 2, & 3, & 4, & 5, & 6 \\ 0, & 2, & 1, & 5, & 4, & 3, & 6 \end{pmatrix} = (1, 2)\,(3, 5).
\]

Man kann die Gruppe $P_{168}$ durch die Congruenzgruppe ter-
%CONT>>


% ===== scan leaf 496, printed 475 =====
% !PRINTED 475
\weberpage{475}
%<<CONT
närer linearer Substitutionen für den Modul 2 darstellen, die wir im \S.\,79 untersucht haben\footnote{Nach einer mündlichen Mittheilung ist dies der Weg, auf dem sie \gesperrt{Kronecker} gebildet hat.}.

Wir haben zu diesem Zweck die sieben Grössen durch drei Indices $(x_1, x_2, x_3)$ zu bezeichnen, die nach dem Modul 2 genommen sind und wobei die Combination $(0, 0, 0)$ ausgeschlossen ist. Man erhält so die sieben Grössen

\weq{1}{(1,0,0),\ (0,1,0),\ (0,0,1);\ (0,1,1),\ (1,0,1),\ (1,1,0),\ (1,1,1),}

und wenn man die $x_1, x_2, x_3$ durch die linearen Verbindungen

\weq{2}{a x_1 + b x_2 + c x_3,\ a_1 x_1 + b_1 x_2 + c_1 x_3,\ a_2 x_1 + b_2 x_2 + c_2 x_3}

ersetzt, worin die Substitution

\weq{3}{\begin{pmatrix} a, & b, & c \\ a_1, & b_1, & c_1 \\ a_2, & b_2, & c_2 \end{pmatrix}}

nach dem Modul 2 zu nehmen ist, so erhält man die Permutationsgruppe $P_{168}$.

Wenn wir, wie im \S.\,79, für $\tau$ die Substitution der Congruenzgruppe

\weq{4}{\tau = \begin{pmatrix} 1, & 0, & 1 \\ 1, & 0, & 0 \\ 0, & 1, & 0 \end{pmatrix}}

wählen, so können wir die Grössen $z$ so bezeichnen, dass sie durch Anwendung von $\tau$ und seinen Wiederholungen cyklisch in einander übergehen, wobei wir eine beliebige der Grössen (1) für $z_0$ wählen können, etwa so:

\weq{5}{z_0 = (1,0,0),\ z_1 = (1,1,0),\ z_2 = (1,1,1),\ z_3 = (0,1,1),\ z_4 = (1,0,1),\ z_5 = (0,1,0),\ z_6 = (0,0,1).}

Es ist dann $\omega$ so zu wählen, dass $z_0, z_4, z_6$ dadurch ungeändert bleiben und $z_1$ mit $z_2$, $z_3$ mit $z_5$ vertauscht werden. Dies giebt

\weq{6}{\omega = \begin{pmatrix} 1, & 0, & 0 \\ 0, & 1, & 0 \\ 0, & 1, & 1 \end{pmatrix},}

und daraus erhält man nach \S.\,72, (17):

\weq{7}{\chi = \begin{pmatrix} 1, & 0, & 0 \\ 0, & 0, & 1 \\ 0, & 1, & 1 \end{pmatrix}, \quad \Theta = \begin{pmatrix} 1, & 1, & 1 \\ 0, & 1, & 0 \\ 0, & 1, & 1 \end{pmatrix}.}


% ===== scan leaf 497, printed 476 =====
% !PRINTED 476
\weberpage{476}
\pdfbookmark[2]{§ 128. Gleichungen siebenten Grades mit einer Gruppe 168sten Grades}{bm146}
Die beiden erzeugenden Permutationen $\tau, \omega$ gehören zur ersten Art (Bd. I, \S.\,153), und folglich ist $P_{168}$ ein Theiler der alternirenden \gesperrt{Permutationsgruppe} von sieben Ziffern.

\begin{center}
\S.\,128.

\gesperrt{Gleichungen siebenten Grades mit einer Gruppe 168\textsuperscript{sten} Grades.}
\end{center}

Wir wenden uns jetzt zu der allgemeinen Theorie der speciellen Art von Gleichungen 7\textsuperscript{ten} Grades, deren \gesperrt{Galois}'sche Gruppe sich auf $P_{168}$ reducirt.

Es seien zunächst $\lambda_0, \lambda_1, \lambda_2, \lambda_3, \lambda_4, \lambda_5, \lambda_6$ beliebige Grössen, und

\weq{1}{\tau = (\lambda_0, \lambda_1, \lambda_2, \lambda_3, \lambda_4, \lambda_5, \lambda_6)}

eine cyklische Permutation 7\textsuperscript{ten} Grades. Führen wir noch das Transpositionspaar

\weq{2}{\omega = (\lambda_1, \lambda_2)\,(\lambda_3, \lambda_5)}

ein, so erzeugen diese beiden Permutationen durch ihre Zusammensetzungen und Wiederholungen die ganze Gruppe $P_{168}$.

Setzt man nach \S.\,72, (17) $\chi$ und $\Theta$ daraus zusammen, so erhält man (da, wie schon oben bemerkt, die Zusammensetzung der Permutationen in umgekehrter Reihenfolge, wie bei den Substitutionen, geschehen muss) die folgenden Permutationen der sieben Indices:

\[
\begin{aligned}
\Theta^3 &= \tau^6 \omega \tau \omega = (1, 2, 5, 3)\,(4, 6),
\qquad \Theta = (1, 3, 5, 2)\,(4, 6) \\
\chi &= \omega \tau^3 \Theta \tau^2 = (1, 4, 2)\,(3, 5, 6),
\end{aligned}
\]

Ausdrücke, die sich auch sehr leicht aus \S.\,127, (7) ableiten lassen. Man sieht, dass $\lambda_0$ durch $\omega, \Theta, \chi$ und folglich durch die ganze in $P_{168}$ enthaltene Octa\"edergruppe $P_{24}$ ungeändert bleibt. Es ist leicht, eine Function der sieben Grössen $\lambda$ zu bilden, die zu der Gruppe $P_{168}$ gehört.

Man nehme z. B. das Product $\lambda_0 \lambda_4 \lambda_6$, das aus den durch $\omega$ unberührt bleibenden $\lambda$ besteht, und bilde die Summe der Producte, die sich daraus durch Anwendung der cyklischen Permutation $\tau$ und ihrer Wiederholungen ergeben:

\weq{3}{v = \lambda_0 \lambda_4 \lambda_6 + \lambda_1 \lambda_5 \lambda_0 + \lambda_2 \lambda_6 \lambda_1 + \lambda_3 \lambda_0 \lambda_2 + \lambda_4 \lambda_1 \lambda_3 + \lambda_5 \lambda_2 \lambda_4 + \lambda_6 \lambda_3 \lambda_5.}


% ===== scan leaf 498, printed 477 =====
% !PRINTED 477
\weberpage{477}
Wendet man darauf die Substitutionen $\omega$ und $\tau$ an, so sieht man, dass $v$ ungeändert bleibt und daher alle Permutationen der Gruppe $P_{168}$ gestattet.

Nun ist $P_{168}$ Theiler der symmetrischen Permutationsgruppe von sieben Elementen vom Index 30 und Theiler der alterniren\-den Gruppe vom Index 15. Folglich ist $v$ Wurzel einer Gleichung 30\,sten Grades, deren Co\"efficienten symmetrische Functionen der $\lambda$ sind, und Wurzel einer Gleichung 15\,ten Grades, deren Co\"efficienten noch das Differenzenproduct der $\lambda$ enthalten. Sind die $\lambda$ die Wurzeln einer Gleichung 7\,ten Grades ohne Affect, so ist $v$ die Wurzel einer Resolvente 30\,sten Grades, die durch Adjunction der Quadratwurzel aus der Discriminante in zwei Factoren 15\,ten Grades zerfällt.

Wenn ausser den symmetrischen Functionen der $\lambda$ die Grösse $v$ dem Rationalitätsbereiche angehört, sei es, dass sie von vornherein rational ist, oder dass der Rationalitätsbereich durch Adjunction von $v$ erweitert wird, so sind die $\lambda$ die Wurzeln einer speciellen Gleichung 7\,ten Grades, deren \gesperrt{Galois}'sche Gruppe vom 168\,sten Grade ist. Was wir nun noch zu beweisen haben, ist, dass sich diese specielle Art von Gleichungen 7\,ten Grades auf das Formenproblem der Gruppe $G_{168}$ zurückführen lässt.

Um dies zu erreichen, müssen wir drei rationale Functionen $X_1$, $X_2$, $X_3$ der Wurzeln $\lambda$ zu bilden suchen, die, wenn die Permutationen der Gruppe $P_{168}$ ausgeführt werden, die entsprechenden linearen Substitutionen der Gruppe $G_{168}$ erfahren, d.\,h. die, wenn $\pi$ eine Permutation aus $P_{168}$ und $A$ die entsprechende Substitution aus $G_{168}$ ist, durch Ausführung der Permutation $\pi$ in $A(X_1, X_2, X_3)$ übergehen.

Setzen wir diese Functionen $X_1, X_2, X_3$ für die Variablen in die Invarianten der Gruppe $G_{168}$ ein, so gehen diese Invarianten in Functionen der $\lambda_i$ über, die durch die Permutationen der Gruppe $P_{168}$ ungeändert bleiben, und die folglich dem Rationalitätsbereiche angehören. Die Berechnung der Werthe der Functionen $X_1, X_2, X_3$ aus diesen Werthen der Invarianten ist dann das Formenproblem für die $G_{168}$. Irgend eine durch die Substitutionen von $G_{168}$ veränderte Function der $X_1, X_2, X_3$ ist Wurzel einer Resolvente der gegebenen Gleichung 7\,ten Grades, und zwar, da die Gruppen $G_{168}$ und $P_{168}$ einfach sind, eine Totalresolvente.


% ===== scan leaf 499, printed 478 =====
% !PRINTED 478
\weberpage{478}
\pdfbookmark[2]{§ 129. Contragrediente Gruppen}{bm147}
Hat man irgend ein System nicht verschwindender Functionen $X_1, X_2, X_3$, so kann man daher solche Resolventen immer bilden. Durch das Formenproblem der Gruppe $G_{168}$ ist also dann die Gleichung 7\textsuperscript{ten} Grades mit der Gruppe $P_{168}$ zugleich gelöst.

Um die Lösung dieser Gleichungen 7\textsuperscript{ten} Grades auf das specielle Formenproblem, wie es bei den elliptischen Functionen auftritt (mit $f = 0$), zurückzuführen, ist dann noch eine accessorische Gleichung 4\textsuperscript{ten} Grades zu lösen, so wie wir, um die allgemeine Gleichung 5\textsuperscript{ten} Grades auf die Ikosaëdergleichung zurückzuführen, eine accessorische Quadratwurzel n\"othig fanden.

Alles kommt also jetzt noch darauf an, die Functionen $X_1, X_2, X_3$ der Wurzeln $\lambda$ diesen Forderungen gemäss zu bestimmen. Um dies zu ermöglichen, müssen wir einige einfache Sätze aus der allgemeinen Invariantentheorie benutzen, die wir im folgenden Paragraphen, soweit sie für unsere Aufgabe in Betracht kommen, ableiten wollen.

\begin{center}
§. 129.

\gesperrt{Contragrediente Gruppen.}
\end{center}

Wir haben schon im §.\,37 den Begriff der contragredienten Transformation erläutert. Sind nämlich

\weq{1}{A = \begin{pmatrix} a_1, & b_1, & c_1 \\ a_2, & b_2, & c_2 \\ a_3, & b_3, & c_3 \end{pmatrix}, \quad A_1 = \begin{pmatrix} a_1, & a_2, & a_3 \\ b_1, & b_2, & b_3 \\ c_1, & c_2, & c_3 \end{pmatrix}}

zwei zu einander transponirte Substitutionen, sind $x_1, x_2, x_3$ und $\xi_1, \xi_2, \xi_3$ zwei Reihen von Variablen, die durch die Substitutionen

\weq{2}{(y) = A(x), \quad \xi = A_1(\eta)}

in zwei neue Reihen von Variablen $y_1, y_2, y_3$ und $\eta_1, \eta_2, \eta_3$ transformirt werden, so haben wir diese beiden Reihen von Variablen und ebenso ihre Transformationen \gesperrt{contragredient} genannt.

Durch die Substitution $y = A(x)$ wird jede Function $\Phi(x_1, x_2, x_3)$ der $(x)$ in eine Function der $(y)$ transformirt, und die Bildung der Abgeleiteten ergiebt:

\weq{3}{\frac{\partial \Phi}{\partial x_1} = a_1 \frac{\partial \Phi}{\partial y_1} + a_2 \frac{\partial \Phi}{\partial y_2} + a_3 \frac{\partial \Phi}{\partial y_3}, \cdots}
%CONT>>


% ===== scan leaf 500, printed 479 =====
% !PRINTED 479
\weberpage{479}
%<<CONT
oder in unserer abgekürzten Schreibweise:

\[
\left(\frac{\partial \Phi}{\partial x_1}, \frac{\partial \Phi}{\partial x_2}, \frac{\partial \Phi}{\partial x_3}\right) = A_1 \left(\frac{\partial \Phi}{\partial y_1}, \frac{\partial \Phi}{\partial y_2}, \frac{\partial \Phi}{\partial y_3}\right).
\]

Dies beweist den Satz:

1. \gesperrt{Die Variablenreihen}

\[
(x_1, x_2, x_3) \quad \text{und} \quad \left(\frac{\partial \Phi}{\partial x_1}, \frac{\partial \Phi}{\partial x_2}, \frac{\partial \Phi}{\partial x_3}\right)
\]

\gesperrt{sind contragredient.}

Durch wiederholte Anwendung dieses Satzes lassen sich auch die höheren Differentialquotienten nach den $x$ durch die nach den $y$ bilden, wofür man folgende Regel erhält:

2. \gesperrt{Um die $m^{\text{ten}}$ Ableitungen}

\[
\frac{\partial^m \Phi}{\partial x_1^\alpha \partial x_2^\beta \partial x_3^\gamma}, \quad \alpha + \beta + \gamma = m
\]

\gesperrt{durch die Ableitungen nach $y$ auszudrücken, ersetze man in dem entwickelten Ausdrucke}

\weq{4}{\xi_1^\alpha \xi_2^\beta \xi_3^\gamma = \sum^{\varkappa, \lambda, \mu} C^{\alpha, \beta, \gamma}_{\varkappa, \lambda, \mu} \eta_1^\varkappa \eta_2^\lambda \eta_3^\mu, \quad \varkappa + \lambda + \mu = m}

\gesperrt{die Producte}

\[
\begin{aligned}
&\xi_1^\alpha \xi_2^\beta \xi_3^\gamma \quad \text{durch} \quad \frac{\partial^m \Phi}{\partial x_1^\alpha \partial x_2^\beta \partial x_3^\gamma} \\
&\eta_1^\varkappa \eta_2^\lambda \eta_3^\mu \quad \text{durch} \quad \frac{\partial^m \Phi}{\partial y_1^\varkappa \partial y_2^\lambda \partial y_3^\mu},
\end{aligned}
\]

\gesperrt{also}

\weq{5}{\frac{\partial^m \Phi}{\partial x_1^\alpha \partial x_2^\beta \partial x_3^\gamma} = \sum C^{\alpha, \beta, \gamma}_{\varkappa, \lambda, \mu} \frac{\partial^m \Phi}{\partial y_1^\varkappa \partial y_2^\lambda \partial y_3^\mu},}

\gesperrt{wo unter dem Summenzeichen $\varkappa, \lambda, \mu$ alle nicht negativen der Bedingung $\varkappa + \lambda + \mu = m$ genügenden Werthe durchlaufen.}

Die Co\"efficienten $C^{\alpha, \beta, \gamma}_{\varkappa, \lambda, \mu}$ sind ganze rationale Functionen der Substitutionscoëfficienten $a_1, a_2, \ldots$

Um diese Regel allgemein zu beweisen, braucht man nur die Formel (5) für $m - 1$ statt $m$ als bewiesen anzusehen, und mit Anwendung der Formel (3) die Ableitung nach einer der Variablen $x$ zu bilden, und dabei die aus der Definition (4) folgende Relation

\[
C^{\alpha, \beta, \gamma}_{\varkappa, \lambda, \mu} = C^{\alpha - 1, \beta, \gamma}_{\varkappa - 1, \lambda, \mu}\, a_1 + C^{\alpha - 1, \beta, \gamma}_{\varkappa, \lambda - 1, \mu}\, a_2 + C^{\alpha - 1, \beta, \gamma}_{\varkappa, \lambda, \mu - 1}\, a_3
\]

zu berücksichtigen.


% ===== scan leaf 501, printed 480 =====
% !PRINTED 480
\weberpage{480}
Diesen Satz können wir nun auch in folgender Weise verallgemeinern:

3. \gesperrt{Wenn durch die Substitution $y = A(x)$ irgend eine Form $\varphi(x)$ in $\Phi(y)$ übergeht, wenn irgend eine zweite Form $\psi(\xi)$ durch die transponirte Substitution $\xi = A_1(\eta)$ in $\Psi(\eta)$ übergeht, so erhält man eine neue Transformation durch $A$, wenn man in $\psi(\xi)$ und $\Psi(\eta)$ die Vertauschungen macht}

\[
\begin{aligned}
&\xi_1^\alpha \xi_2^\beta \xi_3^\gamma \quad \text{mit} \quad \frac{\partial^{\alpha + \beta + \gamma} \varphi}{\partial x_1^\alpha \partial x_2^\beta \partial x_3^\gamma} \\
&\eta_1^\alpha \eta_2^\beta \eta_3^\gamma \quad \text{mit} \quad \frac{\partial^{\alpha + \beta + \gamma} \Phi}{\partial y_1^\alpha \partial y_2^\beta \partial y_3^\gamma},
\end{aligned}
\]

\gesperrt{wenn man also, wie man sich auch ausdrücken kann, in $\psi$ und $\Psi$ die Potenzen und Producte der Variablen $\xi, \eta$ durch die entsprechenden Ableitungen von $\varphi, \Phi$ nach den Variablen $x, y$ setzt.}

Aus der Compositionsregel der linearen Substitutionen ergiebt sich nun sofort der folgende Satz:

4. \gesperrt{Durchläuft $A$ eine Gruppe $G$, so durchläuft die transponirte Substitution $A_1$ eine Gruppe $G_1$. Sind $A, B$ zwei Elemente aus $G$ und $A_1, B_1$ die entsprechenden Elemente aus $G_1$, so sind $AB$ und $B_1 A_1$ entsprechende Elemente. Die Gruppen $G$ und $G_1$ werden zu einander contragredient genannt.}

Die beiden Gruppen $G$ und $G_1$ sind aber nur dann isomorph auf einander bezogen, wenn man dem $A_1$ nicht das Element $A$, sondern das Element $A^{-1}$ entsprechen lässt; denn dann entspricht $A_1 B_1$ dem Elemente $A^{-1} B^{-1} = (BA)^{-1}$.

Die Invarianten der Gruppe $G_1$ heissen Contravarianten der Gruppe $G$. Demnach sind auch die Invarianten von $G$ die Contravarianten von $G_1$.

Aus (3) ergiebt sich dann der folgende Satz:

5. \gesperrt{Wenn man in einer Contravariante von $G$ die Potenzen und Producte der Variablen durch die entsprechenden Ableitungen einer Invariante ersetzt, so erhält man wieder eine Invariante von $G$}


% ===== scan leaf 502, printed 481 =====
% !PRINTED 481
\weberpage{481}
\pdfbookmark[2]{§ 130. Lösung der Gleichung siebenten Grades mit der Gruppe P168 durch das Formenproblem der Gruppe G168}{bm148}

Und ebenso:

6. \gesperrt{Wenn man in einer Invariante von $G$ die Potenzen und Producte der Variablen durch die entsprechenden Ableitungen einer Contravariante ersetzt, so ergiebt sich wieder eine Contravariante.}

\begin{center}
\S. 130.

\gesperrt{Lösung der Gleichung siebenten Grades mit der Gruppe $P_{168}$ durch das Formenproblem der Gruppe $G_{168}$.}
\end{center}

Die lineare Substitutionsgruppe $G_{168}$ hat die bemerkenswerthe Eigenschaft, dass sie mit sich selbst contragredient ist. Denn die erzeugenden Substitutionen $\tau$, $\omega$ von $G_{168}$ (\S.\,115) bleiben durch Transposition ungeändert, und wenn man also irgend eine Substitution der Gruppe transponirt, so erhält man eine Substitution, die gleichfalls in der Gruppe vorkommt.

Die Contravarianten von $G_{168}$ sind also (von der Bezeichnung der Variablen abgesehen) mit ihren Invarianten identisch.

Die in der Gruppe $G_{168}$ enthaltene Octaëdergruppe $G_{24}$, die aus den Substitutionen $\chi^\lambda\omega^\mu\Theta^\nu$ besteht, ist aber von ihrer contragredienten Gruppe verschieden; denn es ist z.\,B. nach \S.\,72, (17)
\[
\Theta^3 = \omega\tau\omega\tau^6,
\]
und die dazu transponirte Substitution
\[
\Theta_1^3 = \tau^6\omega\tau\omega
\]
lässt sich nach \S.\,72, (15) in die Form $\tau\chi^2\Theta^2$ bringen und ist also nicht in $G_{24}$ enthalten.

Es sei nun $\eta_0$ irgend eine zu der Gruppe $P_{24}$ gehörige Function der Grössen $\lambda_0, \lambda_1, \ldots, \lambda_6$ (\S.\,128), z.\,B. die Wurzel $\lambda_0$ selbst. Durch die cyklischen Permutationen $\tau^\nu$ gehe $\eta_0$ in $\eta_0, \eta_1, \ldots, \eta_6$ über. In den Resolventen
\weq{1}{\tilde\omega_r = \sum_{0,6}^{\nu} \varepsilon^{\nu r}\eta_\nu \qquad r = 1, 2, \ldots, 6}
ist dann ein System von Functionen gegeben, die sich durch die Permutationen von $P_{168}$ zwar linear, aber nicht ternär substituiren; da man ja die $\eta_\nu$ selbst linear durch die $\tilde\omega_r$ ausdrücken kann.


% ===== scan leaf 503, printed 482 =====
% !PRINTED 482
\weberpage{482}

Ein System von drei Functionen, die sich ternär substituiren, kann man auf folgende Weise bilden\footnote{F.\ \gesperrt{Klein}, \glqq Ueber die Auflösung gewisser Gleichungen vom 7\textsuperscript{ten} und 8\textsuperscript{ten} Grade\grqq. Mathem. Annalen, Bd. XV (1879).}.

Wir führen zunächst ein System von Hülfsvariablen $x_1, x_2, x_3$ ein, die wir den Substitutionen der Gruppe $G_{168}$ unterwerfen, und daraus bilden wir die zur Gruppe $G_{24}$ gehörige Function $z_0$ mit ihren conjugirten $z_r$ [\S.\,125, (7)]:
\weq{2}{z_r = \varepsilon^{2r}x_1^2 + \varepsilon^{4r}x_2^2 + \varepsilon^r x_3^2 - \frac{1-\sqrt{-7}}{2}(\varepsilon^{-r}x_2x_3 + \varepsilon^{-2r}x_3x_1 + \varepsilon^{-4r}x_1x_2).}

Hierzu nehmen wir nun eine Function $\eta_0$ der Wurzeln $\lambda_r$ unserer Gleichung 7\textsuperscript{ten} Grades, die zu der Gruppe $P_{24}$ gehört, z.\,B. eine rationale Function von $\lambda_0$, und die conjugirten Werthe $\eta_0, \eta_1, \ldots, \eta_6$, und bilden die Summe
\weq{3}{\psi = \eta_0 z_0 + \eta_1 z_1 + \eta_2 z_2 + \eta_3 z_3 + \eta_4 z_4 + \eta_5 z_5 + \eta_6 z_6,}
die eine quadratische Function der $x$ ist, deren Coëfficienten von den Wurzeln $\lambda_r$ abhängen.

\gesperrt{Diese Function $\psi$ ändert sich nicht, wenn die Variablen $(x)$ einer Substitution der Gruppe $G_{168}$ und die Wurzeln $\lambda_r$ gleichzeitig der entsprechenden Permutation aus $P_{168}$ unterworfen werden.} Denn durch diese gleichzeitige Operation werden in der Summe (2) nur die Summanden unter einander vertauscht, also die Summe selbst nicht geändert.

Wir können daher $\psi$ als \gesperrt{simultane Invariante} der Gruppen $G_{168}$ und $P_{168}$ bezeichnen.

Ordnen wir die Function $\psi$ nach den Variablen $x_1, x_2, x_3$, so ergiebt sich
\weq{4}{\psi = p_1x_1^2 + p_2x_2^2 + p_3x_3^2 + 2q_1x_2x_3 + 2q_2x_3x_1 + 2q_3x_1x_2,}
worin zur Abkürzung
\weq{5}{\begin{aligned}
p_1 &= \sum^r \varepsilon^{2r}\eta_r, & q_1 &= -\frac{1-\sqrt{-7}}{4}\sum^r \varepsilon^{-r}\eta_r\\
p_2 &= \sum^r \varepsilon^{4r}\eta_r, & q_2 &= -\frac{1-\sqrt{-7}}{4}\sum^r \varepsilon^{-2r}\eta_r\\
p_3 &= \sum^r \varepsilon^{r}\eta_r, & q_3 &= -\frac{1-\sqrt{-7}}{4}\sum^r \varepsilon^{-4r}\eta_r
\end{aligned}}


% ===== scan leaf 504, printed 483 =====
% !PRINTED 483
\weberpage{483}

gesetzt ist, so dass also die $p, q$ Functionen der Wurzeln $\lambda_r$ sind. Nun wählen wir drei verschiedene Functionen $\eta_0$, die den bisher ausgesprochenen Bedingungen genügen, und bezeichnen sie mit $\eta_0, \eta_0', \eta_0''$.

Die aus diesen drei Functionen abgeleiteten Formen $\psi$ seien $\psi, \psi', \psi''$, und deren Coëfficienten (5) $p_i, q_i$; $p_i', q_i'$; $p_i'', q_i''$. Dann ist nicht nur jede der drei Functionen $\psi, \psi', \psi''$ eine simultane Invariante der Gruppen $P_{168}, G_{168}$, sondern auch ihre Functionaldeterminante (Bd.\,I, \S.\,59):
\weq{6}{\Psi = \frac{1}{S}\begin{vmatrix} \dfrac{\partial\psi}{\partial x_1}, & \dfrac{\partial\psi}{\partial x_2}, & \dfrac{\partial\psi}{\partial x_3}\\[2mm] \dfrac{\partial\psi'}{\partial x_1}, & \dfrac{\partial\psi'}{\partial x_2}, & \dfrac{\partial\psi'}{\partial x_3}\\[2mm] \dfrac{\partial\psi''}{\partial x_1}, & \dfrac{\partial\psi''}{\partial x_2}, & \dfrac{\partial\psi''}{\partial x_3} \end{vmatrix}.}

Die Determinante $\Psi$ ist eine Form 3\textsuperscript{ten} Grades in den Variablen $x$, die nach der Bezeichnungsweise Bd.\,I, \S.\,15, (4) den Ausdruck haben mag:
\weq{7}{\Psi = \sum A_{h,i,k}\, x_h x_i x_k.}

Die Coëfficienten $A_{h,i,k}$ sind dann Functionen der Wurzeln $\lambda_r$, die linear und homogen von jedem der drei Systeme $\eta_r, \eta_r', \eta_r''$ abhängen. Solche Functionen nennt man \gesperrt{trilinear}.

Es bedeute nun $\xi_1, \xi_2, \xi_3$ ein System zu $(x)$ contragredienter Variablen, so dass die biquadratische Form
\weq{8}{f(\xi_1, \xi_2, \xi_3) = \xi_1^3\xi_3 + \xi_2^3\xi_1 + \xi_3^3\xi_2}
eine Contravariante von $G_{168}$ ist. Wenn wir nun in (7)
\[
x_h x_i x_k \quad \text{durch} \quad \frac{1}{6}\frac{\partial^3 f}{\partial\xi_h\partial\xi_i\partial\xi_k}
\]
ersetzen, so erhalten wir eine lineare Form
\weq{9}{L = \frac{1}{6}\sum A_{h,i,k}\,\frac{\partial^3 f}{\partial\xi_h\partial\xi_i\partial\xi_k} = X_1\xi_1 + X_2\xi_2 + X_3\xi_3,}
in der die $X_1, X_2, X_3$ Functionen von $\lambda_r$ sind, und $L$ bleibt nach dem Satze \S.\,129, 3. ungeändert, wenn die Wurzeln $\lambda_r$ durch irgend einer Permutation $\pi$ der Gruppe $P_{168}$ und gleichzeitig die Variablen $(\xi)$ mit den Variablen $(x)$ durch die entsprechende Substitution contragredient transformirt werden. % [sic] \textquotedbl{}irgend einer Permutation\textquotedbl{} as printed


% ===== scan leaf 505, printed 484 =====
% !PRINTED 484
\weberpage{484}
\pdfbookmark[2]{§ 131. Möglichkeit der Bestimmung der Functionen X1, X2, X3}{bm149}

Geht nämlich durch $\pi$ der Coëfficient $A_{h,i,k}$ in $A_{h,i,k}'$ über, so ist vermittelst der Transformation $(y) = A(x)$:
\[
\sum A_{h,i,k}' \, y_h y_i y_k = \sum A_{h,i,k} \, x_h x_i x_k,
\]
und vermittelst der Substitution $\xi = A_1(\eta)$ besteht die Identität $f(\xi_1, \xi_2, \xi_3) = f(\eta_1, \eta_2, \eta_3)$. Folglich ergiebt sich nach dem Satze \S.\,129, 3.:
\[
\sum A_{h,i,k}' \, \frac{\partial^3 f}{\partial\eta_h\partial\eta_i\partial\eta_k} = \sum A_{h,i,k} \, \frac{\partial^3 f}{\partial\xi_h\partial\xi_i\partial\xi_k}.
\]

Bedeutet also $\pi$ irgend eine Permutation der Gruppe $P_{168}$, durch die $X_1, X_2, X_3$ in $Y_1, Y_2, Y_3$ übergeht, und
\[
A = \begin{pmatrix} a_1, & b_1, & c_1\\ a_2, & b_2, & c_2\\ a_3, & b_3, & c_3 \end{pmatrix}
\]
die entsprechende Substitution aus der Gruppe $G_{168}$, so haben wir zu setzen:
\[
\begin{aligned}
\xi_1 &= a_1\eta_1 + a_2\eta_2 + a_3\eta_3\\
\xi_2 &= b_1\eta_1 + b_2\eta_2 + b_3\eta_3\\
\xi_3 &= c_1\eta_1 + c_2\eta_2 + c_3\eta_3,
\end{aligned}
\]
und die Invarianteneigenschaft von $L$ giebt die Relation
\weq{10}{Y_1\eta_1 + Y_2\eta_2 + Y_3\eta_3 = X_1\xi_1 + X_2\xi_2 + X_3\xi_3,}
oder entwickelt
\weq{11}{\begin{aligned} Y_1 &= a_1X_1 + b_1X_2 + c_1X_3\\ Y_2 &= a_2X_1 + b_2X_2 + c_2X_3\\ Y_3 &= a_3X_1 + b_3X_2 + c_3X_3, \end{aligned}}
d.\,h. die Permutation $\pi$, auf die Functionen $X_1, X_2, X_3$ angewandt, hat denselben Erfolg, wie die lineare Substitution $A(X_1, X_2, X_3)$, und demnach sind die $X_1, X_2, X_3$ solche Functionen, wie sie unser Problem verlangt (Schluss des \S.\,128).

\begin{center}
\S. 131.

\gesperrt{Möglichkeit der Bestimmung der Functionen $X_1, X_2, X_3$.}
\end{center}

Um die Zurückführung der Gleichung 7\textsuperscript{ten} Grades mit der Gruppe $P_{168}$ auf das Formenproblem der Gruppe $G_{168}$ vollständig sicher zu stellen, bleibt noch Eines übrig:
%CONT>>


% ===== scan leaf 506, printed 485 =====
% !PRINTED 485
\weberpage{485}
%<<CONT
Es handelt sich nämlich noch um den Nachweis, dass man über $\eta_0, \eta_0', \eta_0''$ (§.\,130) so verfügen kann, dass die Functionen $X_1, X_2, X_3$ nicht identisch verschwinden. Dazu müssen wir die Bildungsweise der Grössen $X$ etwas genauer betrachten.

Setzen wir in (9) zufolge (8) (§.\,130):
\[
\frac{\partial^3 f}{\partial \xi_1^3} = 6\,\xi_3, \quad \frac{\partial^3 f}{\partial \xi_1^2 \xi_2} = 0, \quad \frac{\partial^3 f}{\partial \xi_1^2 \partial \xi_3} = 6\,\xi_1, \quad \frac{\partial^3 f}{\partial \xi_1 \partial \xi_2 \partial \xi_3} = 0, \dots, % [sic] source omits the second \partial before \xi_2
\]
so ergiebt sich

\weq{1}{\begin{aligned}
X_1 &= A_{2,2,2} + 3\,A_{1,1,3}, \quad X_2 = A_{3,3,3} + 3\,A_{2,2,1}, \\
X_3 &= A_{1,1,1} + 3\,A_{3,3,2}.
\end{aligned}}

Nun ist aber ferner nach (4) und (6) (§.\,130):

\weq{2}{\Psi =
\begin{vmatrix}
p_1 x_1 + q_3 x_2 + q_2 x_3, & q_3 x_1 + p_2 x_2 + q_1 x_3, & q_2 x_1 + q_1 x_2 + p_3 x_3 \\
p_1' x_1 + q_3' x_2 + q_2' x_3, & q_3' x_1 + p_2' x_2 + q_1' x_3, & q_2' x_1 + q_1' x_2 + p_3' x_3 \\
p_1'' x_1 + q_3'' x_2 + q_2'' x_3, & q_3'' x_1 + p_2'' x_2 + q_1'' x_3, & q_2'' x_1 + q_1'' x_2 + p_3'' x_3
\end{vmatrix}.}

Dies lässt sich leicht nach Potenzen und Producten der $x$ ordnen, und wenn wir also die Bezeichnung gebrauchen
\[
(p_1, q_3, q_2) =
\begin{vmatrix}
p_1, & q_3, & q_2 \\
p_1', & q_3', & q_2' \\
p_1'', & q_3'', & q_2''
\end{vmatrix} \cdots,
\]
so erhält man
\[
\begin{aligned}
A_{1,1,1} &= (p_1, q_3, q_2), & 3\,A_{3,3,2} &= (q_3, q_1, p_3) + (q_2, p_2, p_3) \\
A_{2,2,2} &= (q_3, p_2, q_1), & 3\,A_{1,1,3} &= (p_1, q_1, q_2) + (p_1, q_3, p_3) \\
A_{3,3,3} &= (q_2, q_1, p_3), & 3\,A_{2,2,1} &= (q_3, p_2, q_2) + (p_1, p_2, q_1),
\end{aligned}
\]
und daraus nach (1)

\weq{3}{\begin{aligned}
X_1 &= (q_3, p_2, q_1) + (p_1, q_1, q_2) + (p_1, q_3, p_3) \\
X_2 &= (q_2, q_1, p_3) + (q_3, p_2, q_2) + (p_1, p_2, q_1) \\
X_3 &= (p_1, q_3, q_2) + (q_3, q_1, p_3) + (q_2, p_2, p_3).
\end{aligned}}

Diese Functionen $X_1, X_2, X_3$ sind, wie aus dem oben Bemerkten folgt [§.\,130, (5)], trilineare Formen der drei Variablenreihen $\eta_r, \eta_r', \eta_r''$, deren Co\"efficienten rational durch die siebente Einheitswurzel $\varepsilon$ ausgedrückt werden können. Die Co\"efficienten dieser Formen sind aber gewiss nicht alle gleich Null. Denn nach (3) kann man die Grössen $p, q$ so annehmen, dass $X_1, X_2, X_3$ nicht gleich Null werden, und die 21 Grössen $\eta_r, \eta_r', \eta_r''$ lassen sich aus §.\,130, (5) so bestimmen, dass die 18 Grössen $p, q$ (und ausserdem die drei Summen $\Sigma \eta_r$) beliebig vorgeschriebene Werthe
%CONT>>


% ===== scan leaf 507, printed 486 =====
% !PRINTED 486
\weberpage{486}
%<<CONT
bekommen. Machen wir dann für die sieben Variablen $\eta_r$ die Substitution

\weq{4}{\eta_r = a_0 + a_1 \lambda_r + a_2 \lambda_r^2 + a_3 \lambda_r^3 + a_4 \lambda_r^4 + a_5 \lambda_r^5 + a_6 \lambda_r^6,}

deren Determinante als das Product aller Differenzen $\lambda_i - \lambda_k$ von Null verschieden ist, und substituiren entsprechend

\weq{5}{\eta_r = \sum_{0,\,6}^{s} a_s \lambda_r^s, \quad \eta_r' = \sum_{0,\,6}^{s} a_s' \lambda_r^s, \quad \eta_r'' = \sum_{0,\,6}^{s} a_s'' \lambda_r^s,}

so geht dadurch $X_1$ in eine trilineare Form der drei Variablenreihen $a_s, a_s', a_s''$ über, die nicht identisch verschwinden kann, weil ja auch umgekehrt die $a_s, a_s', a_s''$ durch die $\eta_r, \eta_r', \eta_r''$ linear ausdrückbar sind. Nun kann man für die Variablen $a_s, a_s', a_s''$ solche rationale Zahlenwerthe annehmen (Bd. I, §.\,143), dass die Functionen $X_1, X_2, X_3$ von Null verschiedene Werthe annehmen, und dann stellt (5) eine geeignete Annahme für die Functionen $\eta_r, \eta_r', \eta_r''$ dar\footnote{Vergl. über die hiermit erledigte Frage: \gesperrt{Burckhardt}, \glqq Ueber einen fundamentalen Satz der Lehre von den endlichen Gruppen linearer Substitutionen\grqq{}. Mathem. Annalen, Bd.\ 42 (1892).}.


% ===== scan leaf 508, printed 487 =====
\weberpage{487}% page number not printed on this page
% !PRINTED none
% Part-title page (no running head, no page number).
\begin{center}
\gesperrt{VIERTES BUCH.}

\rule{0.3\textwidth}{0.4pt}

\gesperrt{ALGEBRAISCHE ZAHLEN.}

\rule{0.3\textwidth}{0.4pt}
\end{center}


% ===== scan leaf 509, printed 488 =====
\weberpage{488}% page number not printed on this page
% !PRINTED none
% Blank verso (no running head, page number, or body text).


% ===== scan leaf 510, printed 489 =====
\weberpage{489}% page number not printed on this page
\pdfbookmark[0]{Viertes Buch. Algebraische Zahlen.}{bm150}
\pdfbookmark[1]{Sechzehnter Abschnitt. Zahlen und Functionale eines algebraischen Körpers.}{bm151}
\pdfbookmark[2]{§ 132. Definition der algebraischen Zahlen}{bm152}
% !PRINTED none
\begin{center}
\gesperrt{Sechzehnter Abschnitt.}

\textbf{Zahlen und Functionale eines algebraischen K\"orpers.}

\rule{4cm}{0.4pt}
\end{center}

\begin{center}
\S.\,132.

\gesperrt{Definition der algebraischen Zahlen.}
\end{center}

Eine algebraische Gleichung
\[
F(x) = x^n + a_1 x^{n-1} + \cdots + a_{n-1} x + a_n = 0,
\]
deren Co\"efficienten $a_1, a_2, \ldots, a_n$ rationale Zahlen sind, nennen wir der K\"urze wegen eine rationale Gleichung. Sie hat, wie wir in fr\"uheren Abschnitten nachgewiesen haben, immer $n$ oder weniger, aber immer wenigstens eine Wurzel. Wie man jeder dieser Wurzeln durch rationale Zahlen, etwa durch Decimalbr\"uche oder durch Kettenbr\"uche, n\"othigenfalls mit Zuziehung der imagin\"aren Einheit $i = \sqrt{-1}$ bis auf jeden beliebigen Grad nahe kommen kann, d.\,h. wie man die Werthe der Wurzeln ann\"ahernd berechnen kann, ist im zweiten Buche des ersten Bandes gezeigt.

In den folgenden Betrachtungen soll es sich nun nicht um diese numerischen Werthe handeln, sondern um die arithmetischen Gesetze, denen diese Zahlen unterworfen sind, die sich aus der Definition selbst und nicht aus den numerischen Werthen ableiten lassen. Wir stellen also jetzt folgende Definition an die Spitze:

\gesperrt{Eine Zahl $\Theta$, die einer rationalen Gleichung}
\weq{1}{F(\Theta) = 0}
\gesperrt{gen\"ugt, heisst eine algebraische Zahl.}

Jede algebraische Gleichung mit rationalen Co\"efficienten liefert uns solche algebraische Zahlen, die sich also in beliebiger Menge angeben lassen. Die Frage, ob es auch nicht algebraische Zahlen giebt, wird uns sp\"ater besch\"aftigen.


% ===== scan leaf 511, printed 490 =====
% !PRINTED 490
\weberpage{490}
\pdfbookmark[2]{§ 133. Ganze algebraische Zahlen}{bm153}

Eine algebraische Zahl gen\"ugt nicht nur einer, sondern unendlich vielen rationalen Gleichungen; denn multiplicirt man zwei beliebige Functionen von der Form $F(x)$ mit einander, so erh\"alt man eine Function derselben Form, die f\"ur $x = \Theta$ verschwindet, wenn einer der Factoren diese Eigenschaft hat.

Unter allen rationalen Gleichungen, denen eine algebraische Zahl gen\"ugt, ist eine von m\"oglichst niedrigem Grade, $f(\Theta) = 0$, wenn $f(x)$ die Form hat:
\weq{2}{f(x) = x^n + a_1 x^{n-1} + a_2 x^{n-2} + \cdots + a_n,}
und es kann auch nur eine solche Gleichung geben, wenn wir, wie bisher immer, den Co\"efficienten der h\"ochsten Potenz von $x$ gleich 1 annehmen.

Denn sind $f(x), f_1(x)$ zwei Functionen von der Form (2) von gleichem Grade $n$, so ist $f(x) - f_1(x)$ von niedrigerem als dem $n^{\text{ten}}$ Grade, und wenn sowohl $f(\Theta)$ als $f_1(\Theta)$ verschwindet, so verschwindet auch $f(\Theta) - f_1(\Theta)$; wenn also diese Differenz nicht identisch verschwindet, so gen\"ugt $\Theta$ einer Gleichung von niedrigerem als dem $n^{\text{ten}}$ Grade, was gegen die Voraussetzung ist.

\gesperrt{Die Function $f(x)$ ist im K\"orper der rationalen Zahlen irreducibel.}

Denn zerf\"allt $f(x)$ in zwei rationale Factoren $f_1(x)$ und $f_2(x)$, von denen jeder von niedrigerem Grade ist als $f(x)$, so gen\"ugt $\Theta$ einer der beiden Gleichungen $f_1(\Theta) = 0$, $f_2(\Theta) = 0$, was unserer Voraussetzung widerspricht.

\gesperrt{Ist $n$ der Grad der rationalen Gleichung niedrigsten Grades, der die Zahl $\Theta$ gen\"ugt, so nennen wir $\Theta$ eine algebraische Zahl $n^{\text{ten}}$ Grades.}

\begin{center}
\S.\,133.

\gesperrt{Ganze algebraische Zahlen.}
\end{center}

\gesperrt{Eine algebraische Zahl $\Theta$ wird eine ganze algebraische Zahl genannt, wenn sie einer rationalen Gleichung}
\weq{1}{\Theta^m + A_1 \Theta^{m-1} + \cdots + A_{m-1} \Theta + A_m = 0}
%CONT>>


% ===== scan leaf 512, printed 491 =====
% !PRINTED 491
\weberpage{491}
%<<CONT
\gesperrt{gen\"ugt; deren Co\"efficienten $A_1, A_2, \ldots, A_m$ ganze Zahlen sind.}

Wir bemerken, dass es nach dieser Definition ausreicht, um eine algebraische Zahl $\Theta$ als ganz zu charakterisiren, wenn unter den unendlich vielen Gleichungen der Form (1), denen $\Theta$ gen\"ugt, \gesperrt{eine} ist, deren Co\"efficienten ganze Zahlen sind.

Die ganzen algebraischen Zahlen umfassen als speciellen Fall die gew\"ohnlichen ganzen Zahlen, die wir zur Unterscheidung \gesperrt{ganze rationale Zahlen} nennen. Die positiven ganzen rationalen Zahlen nennen wir auch, einem verbreiteten Sprachgebrauch folgend, \gesperrt{nat\"urliche Zahlen}.

Unter ganzen Zahlen schlechtweg verstehen wir dann ganze algebraische, rationale und irrationale Zahlen.

1. \gesperrt{Eine ganze algebraische Zahl, die zugleich rational ist, ist nothwendig eine ganze rationale Zahl.}

Nehmen wir n\"amlich an, es sei $\Theta = P:Q$ ein rationaler Bruch, und $P, Q$ ganze rationale Zahlen ohne gemeinsamen Theiler, etwa $Q$ positiv, so ergiebt sich aus (1)
\[
P^m + A_1 P^{m-1} Q + A_2 P^{m-2} Q^2 + \cdots + A_m Q^m = 0,
\]
und daraus ist zu ersehen, dass jeder Primtheiler von $Q$ in $P$ enthalten sein m\"usste. Es muss also $Q = 1$ sein, und $\Theta = P$ ist eine ganze rationale Zahl.

2. \gesperrt{Summe, Differenz und Product zweier ganzer Zahlen sind wieder ganze Zahlen.}

Um diesen Hauptsatz zu beweisen, nehmen wir an, es seien $\alpha, \beta$ zwei ganze Zahlen, die den Gleichungen
\weq{2}{\begin{array}{l}
\alpha^\mu + a_1 \alpha^{\mu-1} + \cdots + a_{\mu-1} \alpha + a_\mu = 0 \\
\beta^\nu + b_1 \beta^{\nu-1} + \cdots + b_{\nu-1} \beta + b_\nu = 0
\end{array}}
gen\"ugen, und machen eine der drei Annahmen
\[
\omega = \alpha + \beta, \ \alpha - \beta, \ \alpha \beta.
\]
Dann setzen wir $\mu \nu = m$ und bezeichnen die $m$ Gr\"ossen
\[
\alpha^r \beta^s \quad \begin{array}{l} r = 0, 1, \ldots, \mu - 1 \\ s = 0, 1, \ldots, \nu - 1 \end{array}
\]
in irgend einer Reihenfolge mit $\omega_1, \omega_2, \ldots, \omega_m$.


% ===== scan leaf 513, printed 492 =====
% !PRINTED 492
\weberpage{492}

Dann k\"onnen die Producte $\omega \omega_1, \omega \omega_2, \ldots, \omega \omega_m$ mit H\"ulfe der Gleichungen (2) in die Form gesetzt werden
\[
\omega \omega_r = c_{r,1} \omega_1 + c_{r,2} \omega_2 + \cdots + c_{r,m} \omega_m \quad r = 1, 2, \ldots, m,
\]
worin die Co\"efficienten $c_{s,r}$ ganze rationale Zahlen sind. Wenn man aus diesen Gleichungen aber die $\omega_1, \omega_2, \ldots, \omega_m$ eliminirt, so folgt
\[
\begin{vmatrix}
c_{1,1} - \omega, & c_{1,2}, & \ldots, & c_{1,m} \\
c_{2,1}, & c_{2,2} - \omega, & \ldots, & c_{2,m} \\
\cdots & \cdots & \cdots & \cdots \\
c_{m,1}, & c_{m,2}, & \ldots, & c_{m,m} - \omega
\end{vmatrix}
= 0,
\]
was entwickelt die Form enth\"alt
\[
\omega^m + C_1 \omega^{m-1} + \cdots + C_m = 0,
\]
worin die $C_1, C_2, \ldots, C_m$ gleichfalls ganze rationale Zahlen sind. Dies aber zeigt, dass $\omega$ eine ganze Zahl ist, wie bewiesen werden sollte.

3. \gesperrt{Ist $f(x)$ eine im K\"orper der rationalen Zahlen irreducible Function, und ist eine Wurzel $\Theta$ von $f(x) = 0$ eine ganze Zahl, so sind alle Wurzeln von $f(x)$ ganze Zahlen.}

Denn wenn eine rationale Function $F(x)$ f\"ur $x = \alpha$ verschwindet, so ist $F(x)$ durch $f(x)$ theilbar, und alle Wurzeln von $f(x)$ sind zugleich Wurzeln von $F(x)$ (Bd.\,I, \S.\,141). Wenn nun $\alpha$ eine ganze Zahl ist, so giebt es eine Function
\[
F(x) = x^m + A_1 x^{m-1} + \cdots + A_m
\]
mit ganzzahligen Co\"efficienten $A_1, \ldots, A_m$, die f\"ur $x = \alpha$ verschwindet, und $F(x)$ verschwindet also auch f\"ur alle anderen Wurzeln von $f(x)$, die sonach alle ganze Zahlen sind.

Ist
\[
f(x) = x^n + a_1 x^{n-1} + \cdots + a_n,
\]
so sind die $a_1, a_2, \ldots, a_n$ durch Multiplication und Addition aus den Wurzeln von $f(x)$ zusammengesetzt und sind also nach 2. ganze rationale Zahlen. Daraus folgt:

4. \gesperrt{Ist $\Theta$ eine ganze algebraische Zahl, so hat die Gleichung niedrigsten Grades $f(\Theta) = 0$ [\S.\,132, (2)] ganzzahlige Co\"efficienten.}


% ===== scan leaf 514, printed 493 =====
% !PRINTED 493
\weberpage{493}
\pdfbookmark[2]{§ 134. Algebraische Körper}{bm154}

Wir beweisen noch den Satz:

5. \gesperrt{Jede algebraische Zahl $\Theta$ l\"asst sich durch Multiplication mit
einer nat\"urlichen Zahl in eine ganze algebraische Zahl verwandeln.}

Denn ist
\[
\Theta^m + A_1 \Theta^{m-1} + \cdots + A_m = 0,
\]
und sind $A_1, \ldots, A_m$ rationale Zahlen mit dem gemeinsamen Nenner $a$, so
erh\"alt man durch Multiplication mit $a^m$:
\[
(a\Theta)^m + A_1 a (a\Theta)^{m-1} + A_2 a^2 (a\Theta)^{m-2} + \cdots + A_m a^m = 0,
\]
woraus hervorgeht, dass $a\Theta$ eine \gesperrt{ganze Zahl} ist.

\begin{center}
\S.\,134.

\gesperrt{Algebraische K\"orper.}
\end{center}

Im dreizehnten Abschnitte des ersten Bandes haben wir gesehen, wie man aus
jeder Wurzel $\Theta$ einer in irgend einem K\"orper $\Omega$ irreduciblen
Gleichung $n^{\mathrm{ten}}$ Grades $f(\Theta) = 0$ einen algebraischen
K\"orper $\Omega(\Theta)$ \"uber $\Omega$ ableitet. Die aus den $n$ Wurzeln
dieser Gleichung abgeleiteten $n$ K\"orper, die auch zum Theil oder alle
identisch sein k\"onnen, haben wir \gesperrt{conjugirte K\"orper} genannt.

Bezeichnen wir mit $R$ den K\"orper der rationalen Zahlen, so giebt also nach
unserer Definition jede algebraische Zahl $n^{\mathrm{ten}}$ Grades, $\Theta$,
Anlass zu einem algebraischen K\"orper $R(\Theta)$ \"uber $R$, den wir von
jetzt an kurz einen algebraischen Zahlk\"orper $n^{\mathrm{ten}}$
Grades nennen. Wir haben auch schon fr\"uher nachgewiesen (Bd.\,I,
\S.\,143), dass man immer einen \gesperrt{algebraischen Zahlk\"orper bestimmen kann,
der eine endliche Anzahl beliebig gegebener algebraischer Zahlen enth\"alt.}

Dieser Satz wird an dieser Stelle hervorgehoben, um darauf hinzuweisen, dass
die Allgemeinheit einer Betrachtung \"uber irgend eine endliche Anzahl
algebraischer Zahlen dadurch nicht beeintr\"achtigt wird, dass man diese
Zahlen alle in einem algebraischen Zahlk\"orper gelegen voraussetzt.

Jede Zahl $\omega$ eines solchen K\"orpers kann als ganze Function
$\varphi(\Theta)$ von $\Theta$ mit rationalen Co\"efficienten dargestellt
werden, und jeder Zahl $\omega$ entspricht in jedem der $n$ conjugirten
K\"orper eine bestimmte Zahl. Diese \gesperrt{conjugirten Zahlen} k\"onnen zum Theil
einander gleich sein, und wir haben danach primitive und
%CONT>>


% ===== scan leaf 515, printed 494 =====
% !PRINTED 494
\weberpage{494}

%<<CONT
imprimitive Zahlen des K\"orpers unterschieden. Eine symmetrische Function
der conjugirten Zahlen ist eine rationale Zahl. Unter diesen symmetrischen
Functionen sind zwei von besonderer Wichtigkeit, die Summe und das Product,
von denen die erste die \gesperrt{Spur}, die zweite die \gesperrt{Norm} von
$\omega$ genannt wird. Man bezeichnet diese beiden Zahlen durch $S(\omega)$
und $N(\omega)$.

Hierbei werden, wenn unter den conjugirten Zahlen dieselben Zahlen mehrfach
vorkommen, diese gleichen Zahlen so oft in die Summe oder das Product
aufgenommen, als der Grad ihrer H\"aufigkeit angiebt.

Da sich in jedem solchen Zahlk\"orper die vier fundamentalen
Rechenoperationen ebenso wie im K\"orper der rationalen Zahlen ausf\"uhren
lassen, so kann man auch die Frage aufwerfen, inwieweit sich die aus der
Theorie der rationalen Zahlen bekannten arithmetischen Grundgesetze in einem
beliebigen algebraischen Zahlk\"orper bew\"ahren. Es handelt sich hierbei in
erster Linie um die Zerlegung der ganzen Zahlen in ihre \gesperrt{Primfactoren}.

Da diese Zerlegung mit den Zahlen des algebraischen K\"orpers selbst im
Allgemeinen \gesperrt{nicht} gelingt, so ist eine Erweiterung des
Rechenmateriales n\"othig, um die einfachen Gesetze wieder herzustellen, und
eine solche Erweiterung ist in verschiedenem Sinne m\"oglich. Es m\"ussen
sich aber diese verschiedenen Erweiterungen auf einander zur\"uckf\"uhren,
oder, genauer gesagt, in eine eindeutige Beziehung zu einander setzen
lassen.

\gesperrt{Kummer} hat zuerst f\"ur die aus Einheitswurzeln gebildeten
algebraischen Zahlen (die Kreistheilungszahlen) das grosse Problem durch die
Sch\"opfung der \gesperrt{idealen Zahlen}\textsuperscript{1)} gel\"ost. Eine
andere ganz allgemeine, keiner Ausnahme unterworfene L\"osung hat die
Aufgabe durch \gesperrt{Dedekind} gefunden, der als die einfachsten Elemente
der Rechnung die von ihm so genannten \gesperrt{Ideale}\textsuperscript{2)}

\vspace{1em}\hrule\vspace{0.5em}

\noindent\textsuperscript{1)} \gesperrt{Kummer}, Theorie der idealen
Primfactoren der complexen Zahlen etc. \gesperrt{Crelle}'s Journal, Bd.\,35, 1846,
Bd.\,40, 1850. Abhandlungen der Berliner Akademie 1856. Sur la th\'eorie des
nombres complexes compos\'es de racines de l'unit\'e et de nombres entiers.
\gesperrt{Liouville}'s Journal, Bd.\,16, 1851.

\noindent\textsuperscript{2)} \gesperrt{Dedekind}, in dem letzten Supplement
der 2., 3. und 4. Auflage von \gesperrt{Dirichlet's} Vorlesungen \"uber
Zahlentheorie (Braunschweig 1871, 1879, 1894). Zu vergleichen ist auch: Sur
la th\'eorie des nombres entiers alg\'ebriques im Bulletin von
\gesperrt{Darboux} und \gesperrt{Ho\"uel} (1\textsuperscript{re} s\'er. XI,
1877). Ueber den Zusammenhang zwischen der Theorie der Ideale und der
h\"oheren Congruenzen (Abhandlungen der Ges. der Wissensch. in G\"ottingen,
Bd.\,23,
%CONT>>


% ===== scan leaf 516, printed 495 =====
% !PRINTED 495
\weberpage{495}
\pdfbookmark[2]{§ 135. Ganze Functionen in einem algebraischen Körper}{bm155}

%<<CONT
betrachtet. Einen davon verschiedenen Weg hat
\gesperrt{Kronecker}\textsuperscript{1)} eingeschlagen.

Die Theorie von \gesperrt{Dedekind} ist von ihrem Begr\"under umfassend und
in stets wachsender Einfachheit und Vollkommenheit dargestellt in dem
letzten Supplement der drei neuesten Auflagen von \gesperrt{Dirichlet's}
Vorlesungen \"uber Zahlentheorie.

Die Theorie von \gesperrt{Kronecker} ist erst im Jahre 1882 durch die
Festschrift zu \gesperrt{Kummer's} Jubil\"aum dem weiteren Kreise der
Mathematiker bekannt geworden, und ist auch jetzt noch nicht ganz leicht
zug\"anglich\textsuperscript{2)}.

Bez\"uglich der \gesperrt{Dedekind'schen} Theorie k\"onnen wir also den
Leser, der tiefer in diesen interessanten Gegenstand einzudringen
w\"unscht, auf die genannten Werke verweisen. Dagegen soll hier der Versuch
gemacht werden, den Zusammenhang zwischen den verschiedenen Theorieen
herzustellen und auf Grund der \gesperrt{Kronecker'schen} Voraussetzungen
und H\"ulfsmittel so viel daraus zu entwickeln, als f\"ur die folgenden
algebraischen Anwendungen nothwendig ist.

\begin{center}
\S.\,135.

\gesperrt{Ganze Functionen in einem algebraischen K\"orper.}
\end{center}

Schon im ersten Bande haben wir mehrfach Gelegenheit gehabt, ganze
Functionen einer beliebigen Anzahl von unabh\"angigen Ver\"anderlichen
einzuf\"uhren, und haben auch (in \S.\,141 f.)
%CONT>>

\vspace{1em}\hrule\vspace{0.5em}

\noindent 1878). Ueber die Discriminanten endlicher K\"orper (ebend.,
Bd.\,29, 1882). Ueber die Anzahl der Ideal-Classen in den verschiedenen
Ordnungen eines endlichen K\"orpers (Braunschweig 1877). Festschrift zur
S\"acularfeier des Geburtstages von \gesperrt{Gauss}. \glqq Zur Theorie der
Ideale\grqq{} und \glqq Ueber die Begr\"undung der Idealtheorie\grqq{}.
Nachrichten d. Ges. d. Wissensch. in G\"ottingen 1894, 1895. Hierher
geh\"oren auch die Abhandlungen von \gesperrt{Hilbert}, \glqq Ueber die
Zerlegung der Ideale etc.\grqq{}, Mathem. Annalen, Bd.\,44, 1893. \glqq
Grundz\"uge einer Theorie der \gesperrt{Galois}'schen Zahlk\"orper\grqq{}, G\"ottinger
Nachrichten 1894. \gesperrt{Hurwitz}, \glqq Zur Theorie der Ideale\grqq{}.
\glqq Ueber einen Fundamentalsatz etc.\grqq{}, G\"ottinger Nachrichten 1894,
1895.

\noindent\textsuperscript{1)} \gesperrt{Kronecker}, Grundz\"uge einer
arithmetischen Theorie der algebraischen Gr\"ossen. Festschrift zu
\gesperrt{Kummer's} 50j\"ahrigem Doctor-Jubil\"aum. Berlin 1882. (Auch in
Bd.\,92 von \gesperrt{Crelle}'s Journal.) Zu erw\"ahnen sind hier noch die Arbeiten von
\gesperrt{Hensel} in den B\"anden 101, 103, 105, 111, 113 des \gesperrt{Crelle}'schen
Journals.

\noindent\textsuperscript{2)} Vgl. die Vorrede zur 4. Auflage von
\gesperrt{Dirichlet-Dedekind's} Vorlesungen \"uber Zahlentheorie.


% ===== scan leaf 517, printed 496 =====
% !PRINTED 496
\weberpage{496}

%<<CONT
den Fall er\"ortert, dass die Co\"efficienten einem bestimmten K\"orper
$\Omega$ angeh\"oren. Wir machen jetzt die Annahme, dass dieser K\"orper
$\Omega$ ein \gesperrt{algebraischer Zahlk\"orper} sei, und betrachten also
Ausdr\"ucke $\varphi(x, y, z, \ldots)$, die als eine Summe von Gliedern der
Form
\[
\alpha x^r y^s z^t \ldots
\]
dargestellt sind, worin die Exponenten $r, s, t, \ldots$ positive oder
wenigstens nicht negative ganze Zahlen sind, w\"ahrend die Coëfficienten
$\alpha$ Zahlen in $\Omega$ bedeuten. Einen solchen Ausdruck
\weq{1}{\varphi(x, y, z, \ldots) = \sum \alpha x^r y^s z^t \ldots}
nennen wir eine \gesperrt{ganze Function in $\Omega$}. Wir nehmen den
Ausdruck immer so geordnet und zusammengefasst an, dass dieselbe
Combination der Exponenten $r, s, t, \ldots$ nicht zweimal darin vorkommt,
und nennen zwei solche Ausdr\"ucke nur dann einander gleich, wenn sie
dieselben Producte $x^r y^s z^t \ldots$, mit denselben Coëfficienten
behaftet, enthalten. \gesperrt{Eine ganze Function wird dann und nur dann gleich Null
gesetzt, wenn alle ihre Coëfficienten Null sind.}

Mehrere ganze Functionen geben durch Addition, Subtraction und
Multiplication immer wieder ganze Functionen. Nach Bd.\,I, \S.\,143, I. kann
man f\"ur die Variablen $x, y, z, \ldots$ solche rationale Zahlwerthe
setzen, dass eine oder eine beliebige Anzahl von gegebenen von Null
verschiedenen ganzen Functionen in $\Omega$ nicht verschwindende Zahlwerthe
(in $\Omega$) erhalten. Daraus ergiebt sich, \gesperrt{dass ein Product mehrerer
ganzer Functionen nur dann verschwindet, wenn einer seiner Factoren
verschwindet.}

Die Summe $r + s + t + \cdots$ der Exponenten in einem Gliede des Ausdrucks
(1) heisst der \gesperrt{Grad dieses Gliedes}, und der gr\"osste Werth, den
der Grad eines Gliedes mit nicht verschwindendem Coëfficienten in $\varphi$
annimmt, heisst der \gesperrt{Grad der Function $\varphi$}.

Der Grad eines Productes aus zweien oder mehreren ganzen Functionen ist
gleich der Summe der Grade der einzelnen Factoren.

Denn fasst man in jedem der Factoren die Summe der Glieder h\"ochsten
Grades zu einer homogenen Function zusammen, so erh\"alt man die Glieder
h\"ochsten Grades des Productes, wenn man alle diese homogenen Functionen
mit einander multiplicirt.


% ===== scan leaf 518, printed 497 =====
% !PRINTED 497
\weberpage{497}

Das Product dieser homogenen Functionen kann nach dem oben Bewiesenen nicht verschwinden, wenn keiner der Factoren verschwindet, und sein Grad ist gleich der Summe der Grade der einzelnen Factoren.

Die Zahlen des K\"orpers $\Omega$ sind unter den Functionen mit enthalten. Man erh\"alt sie, wenn man entweder den Grad oder die Anzahl der Variablen auf Null heruntersinken l\"asst.

Ueber die ganzen Functionen in $\Omega$ haben wir im ersten Bande mehrere wichtige S\"atze abgeleitet, von denen wir hier den Satz \S.\,141, IV hervorheben:

1. \gesperrt{Unter den ganzen Functionen in $\Omega$ m\"ussen reducible und irreducible unterschieden werden, von denen die ersten als Product aus mehreren ganzen Functionen in $\Omega$ darstellbar sind, die anderen nicht. Die reduciblen Functionen $\varphi$ lassen sich in eine endliche Anzahl von irreduciblen Factoren zerlegen, die selbst ganze Functionen in $\Omega$ sind, und die irreduciblen Factoren von $\varphi$ sind, von constanten Factoren, d. h. Zahlen in $\Omega$, abgesehen, eindeutig durch $\varphi$ selbst bestimmt.}

Als specielle F\"alle sind unter den ganzen Functionen in $\Omega$ auch die ganzen Functionen im K\"orper der rationalen Zahlen $R$ enthalten:
\weq{2}{\Phi(x,y,z,\ldots) = \sum a\,x^r y^s z^t \ldots,}
worin die Co\"efficienten $a$ rationale Zahlen sind.

Wenn diese Co\"efficienten \gesperrt{ganze Zahlen ohne gemeinsamen Theiler} sind, so heisst diese Function eine \gesperrt{urspr\"ungliche oder primitive}, von denen wir im Bd. I, \S.\,2 den Satz nachgewiesen haben:

2. \gesperrt{Das Product von zwei primitiven Functionen ist wieder eine primitive Function.}

Wenn die Co\"efficienten der Function (2), die wir f\"ur den Augenblick mit
\[ a_0, a_1, a_2, \ldots \]
bezeichnen wollen, ganze Zahlen mit dem gr\"ossten gemeinschaftlichen Theiler $m$ sind, so ist, wenn $m>1$ ist, $\Phi$ eine imprimitive \gesperrt{ganze ganzzahlige Function vom Theiler $m$, % [sic] \textquotedbl{}ganze ganzzahlige\textquotedbl{} so gedruckt
und der Theiler einer primitiven Function ist $= 1$.}


% ===== scan leaf 519, printed 498 =====
% !PRINTED 498
\weberpage{498}

Setzen wir
\weq{3}{a_0 = m\,e_0,\quad a_1 = m\,e_1,\quad a_2 = m\,e_2 \ldots,}
so sind die $e_0, e_1, e_2, \ldots$ ganze Zahlen ohne gemeinsamen Theiler;
\weq{4}{E(x,y,z,\ldots) = \sum e\,x^r y^s z^t \ldots}
ist eine primitive Function und es wird
\[ \Phi = m\,E. \]

Diese Functionen $\Phi$ haben wir im \S.\,2 des ersten Bandes betrachtet und haben dort von ihnen den Satz bewiesen:

3. \gesperrt{Der Theiler eines Productes von zwei oder mehr ganzen Functionen $\Phi$ ist gleich dem Product der Theiler der einzelnen Factoren.}

Die ganzen Functionen in $\Omega$ h\"angen ausser von den Variablen von einer algebraischen Zahl $\Theta$ ab, enthalten aber sonst nur rationale Zahlenco\"efficienten. Bezeichnen wir eine solche Function mit $\varphi(\Theta, x, y, z, \ldots)$, so erhalten wir die \gesperrt{conjugirten Functionen} $\varphi, \varphi_1, \varphi_2, \ldots$ oder
\weq{5}{\varphi(\Theta,x,y,z,\ldots),\quad \varphi(\Theta_1,x,y,z,\ldots),\quad \varphi(\Theta_2,x,y,z,\ldots),\ \ldots}
wenn wir f\"ur $\Theta$ die s\"ammtlichen Wurzeln der irreduciblen Gleichung $f(x)=0$ [\S.\,132, (2)] einsetzen. Diese conjugirten Functionen k\"onnen auch zum Theil einander gleich sein.

Sie sind alle einander gleich, wenn $\varphi$ eine Function in $R$ ist, und es ist umgekehrt $\varphi$ eine Function in $R$, wenn die conjugirten Functionen alle einander gleich sind; denn es ist dann, wenn wir unter $S(\varphi)$ die Summe der conjugirten Functionen (die Spur) verstehen,
\[ n\varphi = S(\varphi), \]
und $S(\varphi)$ ist eine ganze Function, deren Co\"efficienten symmetrische Functionen der $n$ Wurzeln $\Theta$, d. h. rationale Zahlen, sind.

Zu den ganzen Functionen in $R$ geh\"ort auch die \gesperrt{Norm von} $\varphi$, d. h. das Product
\weq{6}{N(\varphi) = \varphi\,\varphi_1\,\varphi_2 \ldots,}
denn alle Co\"efficienten dieser Function sind symmetrische Functionen der $\Theta$. Diese Function ist theilbar durch $\varphi$, und wenn wir
\[ N(\varphi) = \varphi\,\varphi' \]
setzen, so sind sowohl $\varphi$ als $\varphi'$ ganze Functionen in $\Omega$. Denn $\varphi'$ ist als Product $\varphi_1\varphi_2\ldots$ eine ganze Function, und die Co\"effi-
%CONT>>


% ===== scan leaf 520, printed 499 =====
% !PRINTED 499
\weberpage{499}
\pdfbookmark[2]{§ 136. Die Functionale eines algebraischen Körpers Ω und der erweiterte Körper Ω}{bm156}
%<<CONT
cienten von $\varphi'$ sind symmetrische Functionen der Wurzeln der Gleichung
\[ \frac{f(x)}{x - \Theta} = 0, \]
die ihrerseits in $\Omega$ enthalten sind.

Aus der Definition ergiebt sich, dass die Norm eines Productes gleich dem Producte der Normen der Factoren ist, dass also, wenn $\varphi, \psi$ Functionen in $\Omega$ sind,
\weq{7}{N(\varphi\psi) = N(\varphi)\,N(\psi)}
ist.

\begin{center}
\S.\,136.

\gesperrt{Die Functionale eines algebraischen K\"orpers $\Omega$ und der erweiterte K\"orper $\overline{\Omega}$.}
\end{center}

Die Variablen, die in der Theorie der algebraischen Zahlen verwendet werden, haben nicht die Bedeutung von Zeichen f\"ur ver\"anderliche Zahlenreihen, wie man es aus der Functionentheorie gew\"ohnt ist, sondern sie sind lediglich Rechnungssymbole ohne eine selbst\"andige Bedeutung. Bei den Functionen dieser Variablen kommt es eigentlich nur auf die Co\"efficientensysteme an, und die Variablen werden nur dazu benutzt, um die bekannten und gel\"aufigen Regeln der Buchstabenrechnung auf diese Co\"efficientensysteme anzuwenden. Es ist damit freilich nicht ausgeschlossen, dass gelegentlich auch die Zahlen betrachtet werden, die man erh\"alt, wenn man die Variablen durch gewisse Zahlen, z. B. durch rationale Zahlen, ersetzt.

\gesperrt{Demnach f\"uhren wir in unsere Betrachtungen sowohl ganze als gebrochene Functionen von beliebig vielen Ver\"anderlichen ein, deren Co\"efficienten Zahlen eines algebraischen K\"orpers $\Omega$ sind, und setzen fest, dass mit diesen Functionen so gerechnet wird, wie es die Buchstabenrechnung vorschreibt.}

Jede solche Function $\omega$ kann als Quotient zweier ganzer Functionen in $\Omega$,
\weq{1}{\omega = \frac{\varphi}{\psi},}
dargestellt werden, wobei $\psi$ immer von Null verschieden angenommen werden muss. Zwei solche Functionen sind nur dann einander gleich:
\[ \frac{\varphi}{\psi} = \frac{\varphi_1}{\psi_1}, \]
%CONT>>


% ===== scan leaf 521, printed 500 =====
% !PRINTED 500
\weberpage{500}
%<<CONT
wenn $\varphi\psi_1 = \varphi_1\psi$ ist. Haben die beiden Functionen $\varphi, \psi$ keinen gemeinsamen Theiler, so heisst $\varphi:\psi$ ein \gesperrt{irreducibler Bruch}. Unter den verschiedenen Darstellungen einer Function $\omega$ giebt es \gesperrt{eine} durch einen irreduciblen Bruch, und diese wollen wir die \gesperrt{einfachste Darstellung} nennen. In ihr sind Z\"ahler und Nenner bis auf einen gemeinsamen Factor, der eine beliebige Zahl in $\Omega$ sein kann, durch $\omega$ selbst v\"ollig bestimmt (Bd. I, \S.\,141).

Eine solche Function $\omega$ wollen wir ein \gesperrt{Functional des K\"orpers} $\Omega$ nennen. Als specielle F\"alle sind darunter die ganzen Functionen und die Zahlen selbst enthalten. Bei den Zahlen ist \gesperrt{jede} Darstellung als Quotient zweier Zahlen in $\Omega$ als einfachste Darstellung zu betrachten.

Auf die Functionale lassen sich die vier Grundrechnungsarten in demselben Umfange anwenden, wie auf die Zahlen, und der Inbegriff aller Functionale des K\"orpers $\Omega$ ist daher gleichfalls ein K\"orper, den wir mit $\overline{\Omega}$ bezeichnen und den \gesperrt{Functionalk\"orper} von $\Omega$ nennen wollen.\footnote{Es liegt nahe, die Functionale des K\"orpers $\Omega$, mit denen wie mit den Zahlen in $\Omega$ gerechnet wird, geradezu als \gesperrt{ideale Zahlen} des K\"orpers $\Omega$ zu bezeichnen. Diese Ausdrucksweise w\"urde sich einerseits an die Idealfactoren von \gesperrt{Kummer}, andererseits an die von \gesperrt{Dedekind} eingef\"uhrten Ideale anschliessen, die zu den Functionalen in einer sehr nahen, sp\"ater zu er\"orternden Beziehung stehen. So bestechend in mancher Hinsicht eine solche Terminologie w\"are, so erweist sie sich doch in anderer Hinsicht nicht als zweckm\"assig, weil dadurch den Functionalen, die doch immerhin nur Mittel zum Zweck, nicht selbst\"andig f\"ur sich Gegenstand unseres Interesses sind, anscheinend eine gr\"ossere Bedeutung beigelegt w\"urde, als ihnen in Wirklichkeit zukommt. Mit der Einf\"uhrung des Wortes \glqq Functional\grqq{} folgen wir einem Vorschlage von \gesperrt{Dedekind}.}

Der Functionalk\"orper $\overline{\Omega}$ enth\"alt den Zahlk\"orper $\Omega$ als Theiler.

Ein Functional, dessen Co\"efficienten rationale Zahlen sind, heisst ein \gesperrt{rationales Functional}. Der Inbegriff aller rationalen Functionale ist der Functionalk\"orper $\overline{R}$ des K\"orpers $R$ der rationalen Zahlen, und der K\"orper $\overline{R}$ ist in jedem algebraischen Functionalk\"orper $\overline{\Omega}$ enthalten.

Jedes rationale Functional $A$ kann als Quotient zweier ganzer Functionen in $R$ dargestellt werden, denen man auch ganzzahlige Co\"efficienten geben kann, indem man Z\"ahler und Nenner mit dem Hauptnenner aller Co\"efficienten multiplicirt.


% ===== scan leaf 522, printed 501 =====
% !PRINTED 501
\weberpage{501}

Sind in dieser Darstellung Zähler und Nenner imprimitiv, so
kann man den Theiler herausnehmen und erhält eine Darstellung
in der Form
\weq{2}{A = a\,\frac{E_1(x,y,z,\ldots)}{E_2(x,y,z,\ldots)} = a\,\frac{E_1}{E_2} = aE,}
in der $a$ eine \gesperrt{positive}, ganze oder gebrochene rationale Zahl
ist, während $E_1, E_2$ primitive Functionen in $R$ sind. Hieraus
folgt:

1. \gesperrt{Man kann jedes rationale Functional durch
Multiplication mit einer primitiven Function in $R$ in eine
ganze Function in $R$ verwandeln, und mehrere rationale
Functionale lassen sich als Brüche darstellen, deren
gemeinsamer Nenner eine primitive Function ist.}

Die positive Zahl $a$ und der Quotient $E_1:E_2$ sind durch $A$
vollständig bestimmt. Denn setzen wir $a$ in die Form eines
rationalen Bruches $q_1:q_2$ und nehmen an, es sei
\[
\frac{q_1}{q_2}\,\frac{E_1}{E_2} = \frac{q_1'}{q_2'}\,\frac{E_1'}{E_2'},
\]
so folgt
\[
q_1 q_2' E_1 E_2' = q_2 q_1' E_2 E_1'.
\]
% [sic: Weber prints ``Functionen $R$'' without ``in''.]
Hier haben wir also zwei ganze Functionen $R$ mit ganzzahligen
Co\"efficienten, die einander gleich sind, und deren Theiler, da
$E_1 E_2'$ und $E_2 E_1'$ primitive Functionen sind, $q_1 q_2'$
oder $q_2 q_1'$ ist. Folglich ist $q_1 q_2' = q_2 q_1'$, und daher
auch
\[
\frac{q_1}{q_2} = \frac{q_1'}{q_2'}, \quad \frac{E_1}{E_2} = \frac{E_1'}{E_2'}.
\]

2. \gesperrt{Wir nennen die positive rationale Zahl $a$ den
absoluten Werth des Functionales $A$.}

In dem Falle, dass $A$ selbst eine Zahl ist, ist $a$ der absolute
Werth von $A$ in dem gewöhnlichen Sinne dieses Wortes, und $E$
ist $= +1$ oder $= -1$ zu setzen, je nachdem $A$ positiv oder
negativ ist. Es ist also $E$ nichts weiter als das Vorzeichen von
$A$. In der weitgehenden Verallgemeinerung dieser elementaren
Begriffe liegt das Befremdende, was unsere Definition dem ersten
Blick bietet. Sie wird sich aber in der Folge als durchaus
sachgemäss und nützlich erweisen.

Aus \S.\,135, 3. ergiebt sich der Satz:
%CONT>>


% ===== scan leaf 523, printed 502 =====
% !PRINTED 502
\weberpage{502}
%<<CONT

3. \gesperrt{Der absolute Werth eines Productes zweier oder
mehrerer rationaler Functionale ist gleich dem Producte der
absoluten Werthe der Factoren.}

Ist $\omega$ ein Functional des Körpers $\Omega$, so ist $N(\omega)$ ein rationales
Functional, und die \gesperrt{Norm eines Productes oder eines Quotienten
zweier Functionale ist gleich dem Producte oder dem Quotienten
der Normen der Bestandtheile.}

4. \gesperrt{Wir nennen den absoluten Werth des rationalen
Functionales $N(\omega)$ die absolute Norm $N_a(\omega)$ von
$\omega$ und setzen}
\weq{3}{N(\omega) = N_a(\omega)\,E(\omega).}
$N_a(\omega)$ ist immer eine positive rationale Zahl, und $E(\omega)$
ist der Quotient zweier primitiver ganzer Functionen. Aus 3.
folgt, wenn $\alpha, \beta$ zwei Functionale in $\Omega$ sind, die Formel
\weq{4}{N_a(\alpha\beta) = N_a(\alpha)\,N_a(\beta),}
oder der Satz:

5. \gesperrt{Die absolute Norm eines Productes von Functionalen
in $\Omega$ ist gleich dem Producte der absoluten Normen der
Factoren.}

Wenn wir alle Co\"efficienten eines Functionales in $\Omega$ durch
die entsprechenden Zahlen eines zu $\Omega$ conjugirten Körpers $\Omega_1$
ersetzen, so erhalten wir ein \gesperrt{conjugirtes Functional}. Der
gesammte Functionalkörper $\overline{\Omega}$ geht dadurch in einen
\gesperrt{conjugirten Functionalkörper} $\overline{\Omega}_1$ über. Da jede Gleichung
zwischen Functionalen des Körpers $\Omega$ im Grunde nur die
Zusammenfassung einer Reihe von Gleichungen zwischen Zahlen des
Körpers $\Omega$ ist, so haben wir den Satz (Bd. I, \S.\,141):

6. \gesperrt{Jede Gleichung zwischen Functionalen des Körpers
$\overline{\Omega}$ bleibt richtig, wenn für alle Functionale die
entsprechenden Elemente eines conjugirten Körpers $\overline{\Omega}_1$
gesetzt werden.}

Bedeutet $t$ eine in $\omega$ nicht vorkommende Variable, so hat
die Function $N(t-\omega)$ die Form
\weq{5}{\Phi(t) = t^n + A_1 t^{n-1} + A_2 t^{n-2} + \cdots + A_n,}
worin die Co\"efficienten $A_i$ rationale Functionale sind, und diese
Function $\Phi(t)$ verschwindet, wenn $\omega$ für $t$ gesetzt wird. Es giebt
aber nicht bloss eine solche Function $\Phi(t)$, die für $t = \omega$ ver-
%CONT>>


% ===== scan leaf 524, printed 503 =====
% !PRINTED 503
\weberpage{503}
%<<CONT
schwindet, sondern beliebig viele, da man jedes $\Phi(t)$ mit einer
beliebigen Function der gleichen Form multipliciren kann. Auch
kann es vorkommen, dass schon ein Product von weniger als
$n$ Factoren von $N(t-\omega)$ rationale Co\"efficienten erhält, woraus
Functionen $\Phi(t)$ von niedrigerem als dem $n^{\text{ten}}$ Grade entspringen,
die für $t = \omega$ verschwinden. Daraus folgt:

7. \gesperrt{Es giebt für jedes Functional $\omega$ des Körpers $\Omega$
unendlich viele Functionen $\Phi(t)$, die für $t = \omega$
verschwinden, in denen die höchste Potenz von $t$
den Co\"efficienten 1 hat, und deren übrige Co\"effi\-
cienten in $\overline{R}$ enthalten sind.}

Wir sagen dann, $\omega$ ist eine Wurzel der Gleichung
\[
\Phi(t) = 0.
\]

Unter den verschiedenen Functionen $\Phi(t)$, deren Existenz
im Satze 7. ausgesprochen ist, giebt es eine und nur eine $F(t)$
von möglichst niedrigem Grade. Denn existiren zwei solche
Functionen $F(t)$ und $F_1(t)$ von gleichem Grade $m$, so ist die
Differenz $F(t) - F_1(t)$ in Bezug auf $t$ höchstens vom Grade
$m - 1$ und verschwindet für $t = \omega$. Dividirt man durch den
Co\"efficienten der höchsten Potenz von $t$, so erhält man eine
Function $\Phi(t)$ von niedrigerem Grade als $F(t)$, die nach der
Voraussetzung über $F(t)$ nicht existirt.

Die Function $F(t)$ kann im Körper $\overline{R}$ nicht in Factoren
zerlegt werden, die ganze Functionen von $t$ sind. Denn zerfiele
sie in mehrere Factoren derselben Form $F_1(t)$, $F_2(t)$, so müsste
einer dieser Factoren, die doch alle von niedrigerem Grade als
$F(t)$ selbst sind, für $t = \omega$ verschwinden, entgegen unserer
Voraussetzung.

Jede Function $\Phi(t)$ des Satzes 7. ist dann durch diese
irreducible Function $F(t)$ theilbar, so dass der Quotient
$\Phi(t):F(t)$ eine ganze Function von $t$ in $\overline{R}$ ist, in der die
höchste Potenz von $t$ den Co\"efficienten 1 hat (vgl. Bd. I,
\S.\,141, I.).

Unter den Functionen $\Phi(t)$ findet sich, wie schon bemerkt,
auch die Norm $N(t-\omega)$, und folglich ist $N(t-\omega)$ durch $F(t)$
theilbar. Sind beide Functionen von gleichem Grade, so ist
$N(t-\omega) = F(t)$. Ist aber der Grad von $F(t)$ niedriger als $n$,
so verschwindet der Quotient $N(t-\omega):F(t)$ wenigstens noch
für eines der mit $\omega$ conjugirten Functionale, und folglich für
alle; folglich auch für $t = \omega$, und daher ist dieser Quotient
%CONT>>


% ===== scan leaf 525, printed 504 =====
% !PRINTED 504
\weberpage{504}
\pdfbookmark[2]{§ 137. Ganze Functionale}{bm157}
%<<CONT
nochmals durch $F(t)$, d. h. $N(t-\omega)$ ist durch $F(t)^2$ theilbar.
Durch Fortsetzung dieses Schlussverfahrens erkennt man, dass
$N(t-\omega)$ eine Potenz von $F(t)$ sein muss, und dass folglich der
Grad von $F(t)$ ein Theiler von $n$ ist (Bd. I, \S.\,144).

Wir fassen dies noch in dem Satz zusammen:

8. \gesperrt{Jedes Functional in $\Omega$ ist die Wurzel einer und
nur einer irreduciblen Gleichung $F(t) = 0$ in $\overline{R}$
und $N(t-\omega)$ ist eine Potenz von $F(t)$.}

\begin{center}
\S.\ 137.

\gesperrt{Ganze Functionale.}
\end{center}

1. \textbf{Definition:} \gesperrt{Ein rationales Functional soll ganz
genannt werden, wenn sein absoluter Werth eine ganze Zahl ist.}

Es ergiebt sich aus dieser Definition zunächst, dass die
Summe, die Differenz und das Product von zwei und folglich
auch von beliebig vielen ganzen rationalen Functionalen wieder
ganz sind. Für das Product ist dies eine unmittelbare Folge
des Satzes \S.\,136, 3. Um aber den Beweis für die Summe und
die Differenz zu führen, stellen wir zwei ganze rationale Func\-
tionale $A_1, A_2$ so durch gebrochene Functionen dar, dass sie eine
primitive Form als gemeinschaftlichen Nenner erhalten:
\[
A_1 = a_1\,\frac{E_1}{E}, \quad A_2 = a_2\,\frac{E_2}{E}.
\]
Dann ist
\[
A_1 \pm A_2 = \frac{a_1 E_1 \pm a_2 E_2}{E},
\]
da nun $a_1, a_2$ ganze Zahlen sind, so ist der Theiler der ganzen
Function $a_1 E_1 \pm a_2 E_2$ der absolute Werth von $A_1 \pm A_2$. Dieser
ist eine ganze Zahl, also $A_1 \pm A_2$ ein ganzes Functional.

Für constante rationale Functionale, d. h. für rationale
Zahlen, giebt die Definition 1. die \gesperrt{ganzen rationalen Zahlen.}

Wir stellen ferner, ebenso wie im \S.\,133 für die Zahlen,
folgende Definition der ganzen Functionale des Körpers $\Omega$ auf.

2. \textbf{Definition:} \gesperrt{Ein Functional $\omega$ aus $\Omega$ heisst
ganz, wenn es die Wurzel einer Gleichung}
\[
\Phi(t) = t^m + A_1 t^{m-1} + A_2 t^{m-2} + \cdots + A_m = 0
\]
%CONT>>


% ===== scan leaf 526, printed 505 =====
% !PRINTED 505
\weberpage{505}
%<<CONT
\gesperrt{ist, in der die Coëfficienten $A_1, A_2, \ldots, A_m$ ganze rationale Functionale sind.}

Functionale, die nicht zu den ganzen gehören, werden wir gelegentlich auch der Kürze wegen als gebrochene Functionale bezeichnen.

Zur Rechtfertigung dieser Definition beweisen wir zunächst den Satz:

3. \gesperrt{Ist $\omega$ ein ganzes Functional nach der Definition 2. und zugleich rational, so ist es ein ganzes rationales Functional (nach der Definition 1.).}

Angenommen, es sei $\omega$ ein gebrochenes rationales Functional, und daher der absolute Werth von $\omega$ ein rationaler Bruch $p:q$, worin $p$ und $q$ ganze rationale Zahlen ohne gemeinsamen Theiler sind; dann ist $q\omega$ ein ganzes rationales Functional, dessen absoluter Werth $=p$, also relativ prim zu $q$ ist. Wenn aber andererseits $\omega$ zugleich ganz im Sinne der Definition 2. ist, so können wir
\[
\omega^m = -(A_1\omega^{m-1} + A_2\omega^{m-2} \ldots)
\]
setzen, woraus
\[
(q\omega)^m = -q[A_1(q\omega)^{m-1} + A_2\,q(q\omega)^{m-2} + \cdots) % [sic] opening square bracket is closed by a round parenthesis as printed
\]
folgt.

Hieraus aber ergiebt sich, dass der absolute Werth $p^m$ von $(q\omega)^m$ durch $q$ theilbar sein muss, was nur möglich ist, wenn $q=1$ ist. Die Definition 1. ist also in der Definition 2. als Specialfall enthalten.

Ist
\[
\Phi(t) = t^m + A_1t^{m-1} + A_2t^{m-2} + \cdots + A_m
\]
eine Function von $t$, in der die Coëfficienten $A_1, A_2, \ldots, A_m$ gebrochene rationale Functionale mit den Variablen $x, y, \ldots$, aber von der Variablen $t$ frei sind, so kann man eine Function $aE(x,y,\ldots)$ bestimmen, in der $a$ den Hauptnenner der absoluten Werthe von $A_1, A_2, \ldots$ und $E(x,y,\ldots)$ eine primitive ganze Function bedeutet, so dass
\weq{1}{aE(x\ y,\ldots)\,\Phi(t) = P(t,x,y,\ldots)} % [sic] printed without comma between x and y
selbst eine primitive ganze Function ist [§.\,136, 2.].

Zerfällt $\Phi(t)$ in zwei Factoren $\Phi_1(t)$, $\Phi_2(t)$ in $R$, ist also
\[
\Phi(t) = \Phi_1(t)\,\Phi_2(t),
\]
so bestimme man hiernach für die beiden Functionen $\Phi_1, \Phi_2$ die Factoren $a_1E_1, a_2E_2$, so dass
%CONT>>


% ===== scan leaf 527, printed 506 =====
% !PRINTED 506
\weberpage{506}
%<<CONT
\[
a_1E_1\Phi_1 = P_1, \quad a_2E_2\Phi_2 = P_2
\]
primitive Functionen werden, und dann ist auch
\weq{2}{a_1a_2E_1E_2\Phi = P_1P_2}
eine primitive Function. Aus (1) und (2) folgt
\weq{3}{a_1a_2E_1E_2P = aEP_1P_2}
und mithin
\weq{4}{a = a_1a_2.}
Wenn also $a=1$ ist, so müssen die natürlichen Zahlen $a_1, a_2$ auch $=1$ sein, und wir haben den Satz:

4. \gesperrt{Ist $\omega$ ein ganzes Functional in $\Omega$ und $\Phi(t)$ eine Function von $t$ mit ganzen Coëfficienten in $\overline{R}$, die für $t=\omega$ verschwindet, so hat auch jeder rationale Theiler von $\Phi(t)$ ganze Coëfficienten in $\overline{R}$; insbesondere hat die irreducible Function $F(t)$, von der $\omega$ nach §.\,136, 8. eine Wurzel ist, ganze rationale Functionale zu Coëfficienten.}

Dieser Satz ist deshalb von principieller Bedeutung, weil er lehrt, dass ein Functional $\omega$, das im Körper $\Omega$ zu den gebrochenen gehört, nicht etwa in einem höheren Körper als ganzes Functional erscheinen kann.

Wenn ein Functional $\omega$ nicht ganz ist, so genügt es einer Gleichung
\[
\Phi(\omega) = \omega^m + A_1\omega^{m-1} + A_2\omega^{m-2} + \cdots + A_m = 0,
\]
worin die Coëfficienten $A_1, A_2, \ldots, A_m$ zwar rationale, aber nicht ganze Functionale sind. Ist $a$ der Hauptnenner der absoluten Werthe von $A_1, A_2, \ldots, A_m$, so sind $aA_1, aA_2, \ldots, aA_m$ ganze Functionale. Nun ist
\[
a^m\Phi(\omega) = (a\omega)^m + aA_1(a\omega)^{m-1} + a^2A_2(a\omega)^{m-2} + \cdots
+ a^mA_m = 0,
\]
und $a\omega$ ist daher ein ganzes Functional. Daraus folgt:

5. \gesperrt{Jedes Functional $\omega$ des Körpers $\Omega$ lässt sich durch Multiplication mit einer ganzen rationalen Zahl in ein ganzes Functional verwandeln, und daher kann man jedes Functional $\omega$ als Quotienten zweier ganzer Functionale darstellen. Diese Darstellung ist auf unendlich viele verschiedene Arten möglich, unter anderem so, dass der Nenner eine natürliche Zahl ist.}


% ===== scan leaf 528, printed 507 =====
% !PRINTED 507
\weberpage{507}
6. \gesperrt{Ist $\omega$ ein Functional in $\Omega$ und giebt es $m$ Grössen $\omega_1, \omega_2, \ldots, \omega_m$, die nicht alle verschwinden, von der Art, dass die $m$ Producte $\omega\omega_i$ für $i=1,2,\ldots,m$ in die Form gesetzt werden können:}
\weq{5}{\omega\,\omega_i = A_{1,i}\omega_1 + A_{2,i}\omega_2 + \cdots + A_{m,i}\omega_m,}
\gesperrt{worin die $m^2$ Symbole $A_{k,i}$ ganze rationale Functionale sind, so ist $\omega$ ein ganzes Functional.}

Hierin ist $m$ irgend eine natürliche ganze Zahl. Die $\omega_i$ sind in der Anwendung immer Zahlen oder Functionale in $\Omega$, jedoch ist für die Gültigkeit des Satzes nur die Ausführbarkeit der in (5) angedeuteten Multiplication wesentlich.

Der Satz ist eine einfache Folge der Definition der ganzen Functionale; denn da die $\omega_i$ nicht alle verschwinden, so muss nach dem Determinantensatze (Bd. I, §.\,24, II.)
\weq{6}{\begin{vmatrix}
A_{1,1}-\omega, & A_{2,1}, & \ldots, & A_{m,1} \\
A_{1,2}, & A_{2,2}-\omega, & \ldots, & A_{m,2} \\
\hdotsfor{4} \\
A_{1,m}, & A_{2,m}, & \ldots, & A_{m,m}-\omega
\end{vmatrix} = 0}
sein. Durch Ordnen nach Potenzen von $\omega$ ergiebt sich hieraus eine Gleichung
\weq{7}{\omega^m + A_1\omega^{m-1} + \cdots + A_m = 0,}
worin die Coëfficienten $A_1, A_2, \ldots$ durch Addition und Multiplication aus den $A_{i,k}$ zusammengesetzt sind und daher nach 1. ganze rationale Functionale sind; demnach ist auch $\omega$ ein ganzes Functional.

Für den letzten Coëfficienten $A_m$ in der Gleichung (7) erhalten wir den Ausdruck
\weq{8}{(-1)^mA_m = \Sigma \pm A_{1,1}A_{2,2}\ldots A_{m,m},}
\gesperrt{der also gleich dem Producte der sämmtlichen Wurzeln der Gleichung (7) ist.}

Der Satz 4. ist wichtig als Kennzeichen für ganze Functionale; er dient uns hier zum Beweise des folgenden Satzes: % [sic: Weber prints Satz 4.; the determinant criterion just proved is Satz 6.]

7. \gesperrt{Ist $\Psi(x,y,\ldots)$ eine ganze Function, deren Coëfficienten ganze rationale Zahlen oder Functionale, aber frei von den Variablen $x,y,\ldots$ sind, sind ferner $\alpha,\beta,\ldots$ ganze Functionale in $\Omega$, so ist}
\weq{9}{\omega = \Psi(\alpha,\beta,\ldots)}
\gesperrt{auch ein ganzes Functional.}


% ===== scan leaf 529, printed 508 =====
% !PRINTED 508
\weberpage{508}
Es seien nämlich $\mu, \nu, \ldots$ die Grade der ganzzahligen Gleichungen, denen (nach 2.) die Functionale $\alpha, \beta, \ldots$ genügen, und $m=\mu\nu\ldots$ Wir verstehen unter $\omega_1, \omega_2, \ldots, \omega_m$ die $m$ Grössen
\[
\alpha^r\beta^s\ldots \qquad r=0,1,\ldots,\mu-1, \quad s=0,1,\ldots,\nu-1, \ldots
\]
Dann können mit Hülfe der Gleichungen $\mu^{\text{ten}}, \nu^{\text{ten}}, \ldots$ Grades, denen die Zahlen $\alpha, \beta, \ldots$ genügen, alle Producte $\alpha^r\beta^s\ldots$, in denen einer der Exponenten $r, s, \ldots$ grösser als $\mu-1, \nu-1, \ldots$ ist, linear in der Form (5) durch $\omega_1, \omega_2, \ldots, \omega_m$ ausgedrückt werden, und dasselbe gilt daher auch von jedem Functional $\omega$ und folglich auch von den Producten $\omega\omega_1, \omega\omega_2, \ldots, \omega\omega_m$. Daraus folgt nach 6., dass $\omega$ ganz ist, wie wir beweisen wollten.

Als speciellen Fall des Satzes 7., aus dem übrigens der allgemeine Satz leicht wieder gefolgert werden kann, heben wir hervor:

8. \gesperrt{Durch Addition, Subtraction und Multiplication ganzer Functionale entstehen immer wieder ganze Functionale.}

Wir haben schon früher (§.\,134) bemerkt, dass man immer einen algebraischen Körper bestimmen kann, der beliebig gegebene algebraische Zahlen enthält. Daraus folgt, dass in dem Satze 7. die $\alpha, \beta, \ldots$ beliebige ganze algebraische Zahlen oder Functionale sein können, und Entsprechendes gilt von dem Satze 8.

Aus der Definition 2. folgt noch nach dem Satze §.\,136, 8.:

9. \gesperrt{Ist $\omega$ ein ganzes Functional des Körpers $\Omega$, so sind auch alle mit $\omega$ conjugirte Functionale ganz.}

Als besondere Folgerung dieses Satzes sei noch erwähnt:

10. \gesperrt{Die absolute Norm eines ganzen Functionales ist eine natürliche ganze Zahl.}

Wir beweisen endlich noch den folgenden Satz:

11. \gesperrt{Wenn ein Functional $\omega$ einer Gleichung von der Form}
\weq{10}{\omega^m + \alpha_1\omega^{m-1} + \cdots + \alpha_m = 0}
\gesperrt{genügt, in der $\alpha_1, \alpha_2, \ldots, \alpha_m$ ganze Functionale in $\Omega$ sind, so ist auch $\omega$ ein ganzes Functional.}


% ===== scan leaf 530, printed 509 =====
% !PRINTED 509
\weberpage{509}
\pdfbookmark[2]{§ 138. Theilbarkeit. Associirte Functionale. Einheiten}{bm158}

Um ihn zu beweisen, bezeichnen wir mit $t$ eine in den $\alpha$ und in $\omega$ nicht vorkommende Variable und setzen

\weq{11}{\varphi(t) = t^m + \alpha_1 t^{m-1} + \cdots + \alpha_m.}

Dann ist, wenn $n$ der Grad des Körpers $\Omega$ ist,
\[
\Phi(t) = N[\varphi(t)]
\]
eine ganze Function $mn^{\text{ten}}$ Grades von $t$, deren Coëfficienten ganze rationale Functionale sind. Zugleich ist $\Phi(\omega) = 0$, und folglich $\omega$ eine ganze Zahl.

\gesperrt{Daraus folgt noch, dass der Satz 6. richtig bleibt, wenn die Coëfficienten $A_{k,i}$ nicht ganze Functionale in $R$, sondern in $\Omega$ sind.}

Ist ein ganzes Functional $\omega$ zugleich eine Zahl, so ist es eine ganze Zahl, weil in diesem Falle $N(t-\omega)$ ganze rationale Zahlen zu Coëfficienten hat. Die ganzen Zahlen, die wir im \S.\,133 betrachtet haben, sind demnach als specielle Fälle unter den ganzen Functionalen enthalten.

\S.\,138.

\gesperrt{Theilbarkeit. Associirte Functionale. Einheiten.}

Die ganzen Functionale unterliegen ähnlichen Gesetzen der Theilbarkeit, wie die ganzen rationalen Zahlen. Um sie zu erkennen, stellen wir folgende Definition an die Spitze:

1. \gesperrt{Ein ganzes Functional $\alpha$ heisst durch ein anderes, von Null verschiedenes ganzes Functional $\beta$ theilbar, wenn der Quotient $\alpha:\beta=\gamma$ ein ganzes Functional ist.}

Es ist dann $\alpha=\beta\gamma$ und $\beta$ und $\gamma$ heissen \gesperrt{Theiler} von $\alpha$, und man sagt auch, \gesperrt{$\beta$ und $\gamma$ gehen in $\alpha$ auf}. Die Zahl 0 ist durch \gesperrt{jedes ganze} Functional theilbar.

Aus dieser Definition ergeben sich ohne Schwierigkeit die folgenden fundamentalen Sätze über Theilbarkeit:

2. \gesperrt{Sind $\alpha$ und $\alpha_1$ theilbar durch $\beta$, so ist auch $\alpha\pm\alpha_1$ theilbar durch $\beta$.}

Denn ist $\alpha=\beta\gamma$, $\alpha_1=\beta\gamma_1$, so ist $\alpha\pm\alpha_1=\beta(\gamma\pm\gamma_1)$, nnd wenn $\gamma$ und $\gamma_1$ ganz sind, so ist auch $\gamma\pm\gamma_1$ ganz. % [sic] printed ``nnd''


% ===== scan leaf 531, printed 510 =====
% !PRINTED 510
\weberpage{510}

3. \gesperrt{Ist $\alpha$ theilbar durch $\beta$, und $\beta$ theilbar durch $\gamma$,
so ist auch $\alpha$ theilbar durch $\gamma$.}

Denn nach der Voraussetzung giebt es zwei ganze Functionale $\varkappa$, $\lambda$, die
den Bedingungen $\alpha = \varkappa\beta$, $\beta = \lambda\gamma$ genügen. Demnach ist
auch $\alpha = \varkappa\lambda\gamma$, und da $\varkappa\lambda$ ganz ist, so ist
$\alpha$ durch $\gamma$ theilbar.

Ein Product $\beta\gamma$ zweier ganzer Functionale ist sowohl durch $\beta$ als durch
$\gamma$ theilbar, und folglich ist, wenn $\beta$ durch $\alpha$ theilbar ist, auch
$\beta\gamma$ durch $\alpha$ theilbar. Aus diesen Sätzen ergiebt sich:

4. \gesperrt{Sind $\alpha$, $\beta$, \ldots\ durch $\delta$ theilbar, und sind $\xi$,
$\eta$, \ldots\ beliebige ganze Functionale, so ist $\xi\alpha + \eta\beta + \ldots$
durch $\delta$ theilbar.}

Selbstverständlich erstrecken sich diese Definitionen und Sätze auch auf den Fall, dass
Zahlen an die Stelle von Functionalen treten, und so erhalten wir die Theilbarkeit
ganzer Zahlen.

5. \gesperrt{Zwei ganze Functionale $\alpha$, $\beta$, die gegenseitig durch einander
theilbar sind, heissen associirt.}

6. \gesperrt{Ein mit der natürlichen Zahl 1 associirtes ganzes Functional $\varepsilon$,
d.\,h. jeder Theiler der Zahl 1, heisst eine Einheit.}

Je nachdem $\varepsilon$ ein Functional oder eine Zahl ist, ist $\varepsilon$ eine
\gesperrt{functionale} oder eine \gesperrt{numerische} Einheit.

Im Körper $R$ der rationalen Zahlen sind als functionale Einheiten die primitiven
ganzen Functionen und die Quotienten von zweien unter ihnen anzusehen. Numerische
Einheiten giebt es in $R$ nur zwei, nämlich $+\,1$ und $-\,1$.

Ueber die hierdurch eingeführten Begriffe, die in enger gegenseitiger Beziehung
stehen, leiten wir eine Reihe von Sätzen ab.

7. \gesperrt{Sind $\alpha$, $\beta$ associirte Functionale, so sind die Quotienten
$\beta : \alpha$ und $\alpha : \beta$ Einheiten.}

Denn setzen wir $\beta = \alpha\varepsilon$, so ist $\varepsilon$ ein ganzes
Functional und $1 : \varepsilon = \alpha : \beta$ gleichfalls ganz; also ist
$\varepsilon$ ein Theiler der Zahl 1, d.\,h. $\varepsilon$ ist eine Einheit.

8. \gesperrt{Ist $\alpha$ ein ganzes Functional und $\varepsilon$ eine Einheit, so
sind $\alpha$ und $\alpha\varepsilon$ associirt.}

Dies folgt unmittelbar aus der Definition; denn $\alpha : \alpha\varepsilon =
1 : \varepsilon$ und $\alpha\varepsilon : \alpha = \varepsilon$ sind beides ganze
Functionale.


% ===== scan leaf 532, printed 511 =====
% !PRINTED 511
\weberpage{511}

Eine Einheit ist ein ganzes Functional, dessen reciprokes gleichfalls ganz ist, und dieses reciproke Functional ist selbst eine Einheit. Durch eine Einheit ist jedes beliebige ganze Functional theilbar. Ueberhaupt gilt der Satz:

9. \gesperrt{Das Product und der Quotient zweier Einheiten sind wieder Einheiten.}

Denn sind $\varepsilon_1, \varepsilon_2$ zwei Einheiten, so sind $\varepsilon_1\varepsilon_2$ und $1:\varepsilon_1\varepsilon_2$ ganze Functionale, also $\varepsilon_1\varepsilon_2$ eine Einheit, und ebenso sind $\varepsilon_1:\varepsilon_2$ und $\varepsilon_2:\varepsilon_1$ ganz.

10. \gesperrt{Ist ein ganzes Functional $\mu$ theilbar durch ein anderes, $\alpha$, so ist jedes mit $\mu$ associirte Functional $\mu'$ auch durch jedes mit $\alpha$ associirte Functional $\alpha'$ theilbar.}

Denn wenn $\mu:\alpha$ ganz und $\varepsilon, \varepsilon_1$ Einheiten sind, so ist auch $\mu\varepsilon:\alpha\varepsilon_1$ ganz.

11. \gesperrt{Ist $\alpha$ associirt mit $\beta$ und mit $\gamma$, so sind auch $\beta$ und $\gamma$ unter einander associirt.}

Denn ist $\beta=\alpha\varepsilon$, $\gamma=\alpha\varepsilon_1$, so ist $\beta=\gamma\varepsilon:\varepsilon_1$, und $\varepsilon:\varepsilon_1$ ist nach 8. eine Einheit.

Ist $\alpha$ theilbar durch $\beta$, so ist die absolute Norm von $\alpha$ theilbar durch die absolute Norm von $\beta$, denn aus $\alpha=\beta\gamma$ folgt nach \S.\,136, 5.:

\weq{1}{N_a(\alpha) = N_a(\beta)\, N_a(\gamma).}

Daraus ergiebt sich weiter, dass die absolute Norm einer Einheit $=1$ sein muss. Es gilt aber auch das Umgekehrte:

12. \gesperrt{Ein ganzes Functional, dessen absolute Norm $=1$ ist, ist eine Einheit.}

Denn ist $\varepsilon$ ein ganzes Functional mit der absoluten Norm 1, so ist $N(\varepsilon)$ ein ganzes rationales Functional mit dem absoluten Werthe 1, und daher ist auch $1:N(\varepsilon)$ ein ganzes Functional. Setzen wir dann
\[
N(\varepsilon) = \varepsilon\varepsilon',
\]
so ist $\varepsilon'$ als Product von ganzen Functionalen (den Conjugirten zu $\varepsilon$) selbst ganz, und folglich ist auch
\[
\frac{1}{\varepsilon} = \frac{\varepsilon'}{N(\varepsilon)}
\]
ein ganzes Functional, also $\varepsilon$ eine Einheit.


% ===== scan leaf 533, printed 512 =====
% !PRINTED 512
\weberpage{512}
\pdfbookmark[2]{§ 139. Grösster gemeinschaftlicher Theiler}{bm159}

Daraus schliessen wir noch nach der Formel (1), dass \gesperrt{associirte Functionale dieselbe absolute Norm haben}.

Ein ganzes rationales Functional ist hiernach immer mit seinem absoluten Werthe associirt. Ist also $\alpha$ irgend ein ganzes Functional des Körpers $\overline{\Omega}$, so sind auch $N(\alpha)$ und $N_a(\alpha)$ associirt. Da nun in $N(\alpha)=\alpha\alpha'$ der Factor $\alpha'$ als Product von ganzen Functionalen selbst ganz ist, so ist $N(\alpha)$ und folglich auch die natürliche Zahl $N_a(\alpha)$ durch $\alpha$ theilbar. Wir formuliren also noch den Satz:

13. \gesperrt{Es giebt natürliche ganze Zahlen (in unendlicher Menge), darunter die absolute Norm von $\alpha$, die durch ein beliebiges ganzes Functional $\alpha$ theilbar sind. Ist $a$ die kleinste unter den durch $\alpha$ theilbaren natürlichen Zahlen, so ist jede durch $\alpha$ theilbare ganze rationale Zahl durch $a$ theilbar.}

Denn ist $m$ eine durch $\alpha$ theilbare ganze rationale Zahl, so ist auch der Rest der Division von $m$ durch $a$ eine durch $\alpha$ theilbare, ganze rationale Zahl. Da diese kleiner als $a$ ist, so muss sie $=0$ sein.

\S.\,139.

\gesperrt{Grösster gemeinschaftlicher Theiler.}

Es mögen $\alpha, \beta, \ldots$ von Null verschiedene ganze Functionale in $\Omega$ in beliebiger aber endlicher Anzahl bedeuten, und $x, y, \ldots$ Variable, die in $\alpha, \beta, \ldots$ nicht vorkommen. Dann ist

\weq{1}{\delta = \alpha x + \beta y + \cdots}

gleichfalls ein ganzes Functional in $\Omega$. Die Norm von $\delta$ ist eine ganze homogene Function $n^{\text{ten}}$ Grades der Variablen $x, y, \ldots$, deren Coëfficienten ganze rationale Functionale sind. Bezeichnen wir die absolute Norm von $\delta$ mit $D$, so ist

\weq{2}{N(\delta) = D\,E(x, y, \ldots) = D\,E,}

und darin ist $E(x, y, \ldots) = E$ eine ganze rationale Function der Variablen $x, y, \ldots$ und im Allgemeinen eine gebrochene Function der in $\alpha, \beta, \ldots$ vorkommenden Variablen, jedenfalls aber eine \gesperrt{rationale Einheit} [\S.\,136, (2)].


% ===== scan leaf 534, printed 513 =====
% !PRINTED 513
\weberpage{513}

Wir wollen jetzt unter $x_0, y_0, \dots$ irgend welche ganze rationale Zahlen verstehen. Dann ist auch
\[ \delta_0 = \alpha x_0 + \beta y_0 + \cdots \]
ein ganzes Functional, und wir wollen beweisen, \gesperrt{dass $\delta_0$ durch $\delta$ theilbar ist.}

Wir brauchen zu diesem Zwecke nur das ganze Functional
\[ \delta t - \delta_0 = \alpha (x t - x_0) + \beta (y t - y_0) \dots \]
zu betrachten, worin $t$ eine neue Variable bedeutet. Dann ist [nach (2)]:

\weq{3}{N (\delta t - \delta_0) = D E (x t - x_0, y t - y_0, \dots),}

oder, indem wir nach absteigenden Potenzen von $t$ ordnen:

\weq{4}{N (\delta t - \delta_0) = D (t^n E + t^{n-1} E_1 + t^{n-2} E_2 + \cdots),}

worin $E_1, E_2, \dots$ nach \S.\,137, 7., 8. ganze rationale Functionale sind. Setzen wir nun
\[ C_1 = \frac{E_1}{E}, \quad C_2 = \frac{E_2}{E}, \cdots, \]
so sind, da $E$ eine Einheit ist, auch $C_1, C_2, \dots$ ganze rationale Functionale, und es folgt, wenn wir noch
\[ \frac{\delta_0}{\delta} = \eta \]
setzen, aus (4)
\[ N (\delta) N (t - \eta) = D E (t^n + C_1 t^{n-1} + C_2 t^{n-2} + \cdots), \]
oder wegen (2)

\weq{5}{N (t - \eta) = t^n + C_1 t^{n-1} + C_2 t^{n-2} + \cdots}

Da diese Function nun verschwindet, wenn $t = \eta$ gesetzt wird, so folgt, dass $\eta$ ein ganzes Functional ist (\S.\,137, 2.), und damit ist bewiesen, dass $\delta_0$ durch $\delta$ theilbar ist.

Da $x_0, y_0, \dots$ beliebige ganze rationale Zahlen bedeuten können, so schliessen wir daraus, dass die Functionale $\alpha, \beta, \dots$ selbst durch $\delta$ theilbar sind, dass also $\delta$ ein gemeinsamer Theiler der Functionale $\alpha, \beta, \dots$ ist.

Andererseits ist aber auch (nach \S.\,138, 4.) $\delta$ durch jeden gemeinsamen Theiler von $\alpha, \beta, \dots$ theilbar, und $\delta$ hat also die \gesperrt{charakteristischen Eigenschaften des grössten gemeinschaftlichen Theilers von $\alpha, \beta, \dots$} Dieselben Eigenschaften kommen aber nach \S.\,138, 7. jedem mit $\delta$ associirten Functional zu, und ebenso sind auch zwei Functionale $\delta, \delta'$, die diese doppelte Eigen-
%CONT>>


% ===== scan leaf 535, printed 514 =====
% !PRINTED 514
\weberpage{514}
%<<CONT
schaft haben, durch einander theilbar, also associirt, und wir stellen also die Definition auf:

1. \gesperrt{Das Functional $\delta = \alpha x + \beta y + \cdots$ und jedes damit associirte Functional heisst grösster gemeinschaftlicher Theiler von $\alpha, \beta, \dots$}

Wenn $\delta$ eine Einheit ist, so sagen wir auch, $\alpha, \beta, \dots$ seien ohne gemeinsamen Theiler, denn dann giebt es ausser den Einheiten kein ganzes Functional, das in allen $\alpha, \beta, \dots$ aufgeht.

2. \gesperrt{Zwei Functionale $\alpha, \beta$, die keinen gemeinsamen Theiler haben, für die also $\alpha x + \beta y$ eine Einheit ist, heissen relative Primfunctionale oder theilerfremd.}

Aus diesen Definitionen ergeben sich sehr einfach folgende Sätze:

3. \gesperrt{Wenn ganze Functionale $\xi, \eta, \dots$ in $\Omega$ existiren, derart, dass
\[ \alpha \xi + \beta \eta + \cdots = \varepsilon \]
eine Einheit ist, so sind die ganzen Functionale $\alpha, \beta, \dots$ ohne gemeinsamen Theiler.}

Denn haben $\alpha, \beta, \dots$ einen gemeinsamen Theiler $\delta$, so ist (nach \S.\,138) $\delta$ auch Theiler von $\alpha \xi + \beta \eta + \cdots$, und $\delta$ muss also auch eine Einheit sein.

4. \gesperrt{Sind $\alpha, \beta, \gamma$ drei ganze Functionale und ist $\alpha$ relativ prim zu $\beta$ und zu $\gamma$, so ist $\alpha$ auch relativ prim zu $\beta \gamma$.}

Denn nach Voraussetzung sind, wenn $x, y, u, v$ vier Variable sind,
\[ \varepsilon = \alpha x + \beta y, \quad \varepsilon_1 = \alpha u + \gamma v \]
Einheiten. Demnach ist auch
\[ \alpha (\alpha u x + \gamma v x + \beta u y) + \beta \gamma v y = \varepsilon \varepsilon_1 \]
eine Einheit. Da aber $\alpha u x + \gamma v x + \beta u y$ und $v y$ ganz sind, so folgt nach dem Satze 3., dass $\alpha$ und $\beta \gamma$ relativ prim sind.

Hieran schliesst sich der Beweis des folgenden sehr wichtigen Satzes:

5. \gesperrt{Sind $\alpha, \beta, \mu$ ganze Functionale, $\alpha$ relativ prim zu $\beta$ und $\alpha \mu$ durch $\beta$ theilbar, so ist $\mu$ durch $\beta$ theilbar.}


% ===== scan leaf 536, printed 515 =====
% !PRINTED 515
\weberpage{515}
\pdfbookmark[2]{§ 140. Primfunctionale im Körper Ω}{bm160}

Denn nach der Voraussetzung über $\alpha, \beta$ ist
\[ \alpha x + \beta y = \varepsilon \]
eine Einheit. Durch Multiplication mit $\mu$ folgt daraus
\[ \alpha \mu x + \beta \mu y = \varepsilon \mu, \]
und da $\alpha \mu$ und $\beta \mu$ nach Voraussetzung durch $\beta$ theilbar sind, so ist auch $\varepsilon \mu$ und folglich auch das mit $\varepsilon \mu$ associirte $\mu$ durch $\beta$ theilbar.

\begin{center}
\S.\,140.

\gesperrt{Primfunctionale im Körper $\Omega$.}
\end{center}

Durch die Sätze des vorigen Paragraphen haben wir die Hülfsmittel gewonnen, um die Gesetze der Theilbarkeit der ganzen Functionale im Körper $\Omega$ genau auf demselben Wege abzuleiten, den man in den Elementen der Arithmetik auf die natürlichen ganzen Zahlen anwendet. Wir definiren folgendermaassen:

1. \gesperrt{Ein ganzes Functional $\pi$ des Körpers $\Omega$, welches keine Einheit ist, heisst ein Primfunctional, wenn es ausser durch die Einheiten nur noch durch die mit ihm selbst associirten Functionale theilbar ist. Jedes ganze Functional in $\Omega$, das ausser diesen noch andere Theiler hat, heisst zusammengesetzt.}

Der Begriff des Primfunctionales ist hiernach wesentlich von dem Körper $\Omega$ abhängig. Es können sehr wohl die Primfunctionale eines Körpers in einem anderen erweiterten Körper zusammengesetzt sein.

Wenn $\omega$ ein beliebiges ganzes und $\pi$ ein Primfunctional des Körpers $\Omega$ ist, so sind nur zwei Fälle möglich: entweder $\omega$ ist relativ prim zu $\pi$ oder $\omega$ ist durch $\pi$ theilbar; denn ein gemeinschaftlicher Theiler von $\omega$ und $\pi$ kann nur entweder eine Einheit oder mit $\pi$ associirt sein, und im letzteren Falle ist $\omega$ durch $\pi$ theilbar. Daraus ergiebt sich der Satz:

2. \gesperrt{Wenn das Product $\alpha \beta$ zweier ganzer Functionale $\alpha, \beta$ in $\Omega$ durch ein Primfunctional $\pi$ theilbar ist, so muss einer der beiden Factoren durch $\pi$ theilbar sein.}


% ===== scan leaf 537, printed 516 =====
% !PRINTED 516
\weberpage{516}
\pdfbookmark[2]{§ 141. Zerlegung der ganzen und gebrochenen Functionale in Primfactoren}{bm161}

Denn wenn $\alpha$ und $\beta$ beide nicht durch $\pi$ theilbar, also beide relativ prim zu $\pi$ sind, so ist nach \S.\,139, 4. auch $\alpha \beta$ relativ prim zu $\pi$.

Es ergiebt sich daraus durch wiederholte Anwendung, dass, wenn ein Product aus \gesperrt{mehreren} Factoren durch $\pi$ theilbar ist, mindestens einer der Factoren durch $\pi$ theilbar sein muss.

3. \gesperrt{Die kleinste natürliche ganze Zahl $p$, die durch ein Primfunctional $\pi$ in $\Omega$ theilbar ist, ist eine natürliche Primzahl, und die absolute Norm von $\pi$ ist eine Potenz von $p$. Jede durch $\pi$ theilbare ganze rationale Zahl ist auch durch $p$ theilbar.}

Nach \S.\,138, 13. giebt es natürliche Zahlen, die durch $\pi$ theilbar sind. Die kleinste unter ihnen, $p$, kann nicht in zwei natürliche Factoren, die grösser als 1 sind, zerlegbar sein; denn ist $p = p_1 p_2$, so muss entweder $p_1$ oder $p_2$ durch $\pi$ theilbar sein (nach 2.). Ist aber keiner der Factoren $p_1, p_2 = 1$, so sind sie beide kleiner als $p$, was der Voraussetzung über $p$ widerspricht. Folglich ist $p$ eine natürliche Primzahl. Dass jede durch $\pi$ theilbare natürliche Zahl $m$ durch $p$ theilbar ist, haben wir schon im \S.\,138, 13. bewiesen.

Setzen wir nun

\weq{1}{p = \pi \omega,}

so ist $\omega$ ein ganzes Functional, und wenn wir beiderseits die absoluten Normen nehmen, so folgt, da die absolute Norm einer natürlichen Zahl die $n^{\text{te}}$ Potenz dieser Zahl ist:

\weq{2}{p^n = N_a (\pi) N_a (\omega).}

Hieraus folgt der zweite Theil unseres Satzes, dass die natürliche ganze Zahl $N_a (\pi)$ eine Potenz von $p$ ist.

Setzen wir demnach

\weq{3}{N_a (\pi) = p^f,}

so ist $f$ eine Zahl, die nur einen der Werthe $1, 2, 3, \dots, n$ haben kann, wenn $n$ den Grad des Körpers $\Omega$ bedeutet.

Die Zahl $f$ heisst der \gesperrt{Grad} des Primfunctionals $\pi$.

\begin{center}
\S.\,141.

\gesperrt{Zerlegung der ganzen und gebrochenen Functionale in Primfactoren.}
\end{center}

1. \gesperrt{Jedes von Null verschiedene ganze Functional $\omega$ des Körpers $\Omega$, das keine Einheit ist, ist durch ein Primfunctional theilbar.}


% ===== scan leaf 538, printed 517 =====
% !PRINTED 517
\weberpage{517}
Wenn $\omega$ prim ist, so ist der Satz evident, weil $\omega$ durch sich selbst theilbar ist. Wenn aber $\omega$ nicht prim ist, so ist es durch ein ganzes Functional $\omega_1$ theilbar, das weder eine Einheit, noch mit $\omega$ associirt ist. Ist also
\weq{1}{\omega = \omega_1 \alpha,}
so ist weder $\omega_1$ noch $\alpha$ eine Einheit. Aus (1) folgt aber
\[
N_a(\omega) = N_a(\omega_1)\, N_a(\beta), % [sic] printed with $\beta$, though (1) introduces $\alpha$
\]
und da $N_a(\beta)$ grösser als 1 ist, so ist
\weq{2}{N_a(\omega_1) < N_a(\omega).}

Wenn $\omega_1$ noch nicht prim ist, so kann man denselben Schluss auf $\omega_1$ anwenden und findet einen Theiler $\omega_2$ von $\omega_1$ derart, dass
\weq{3}{N_a(\omega_2) < N_a(\omega_1) < N_a(\omega)}
ist. Da es aber nur eine endliche Anzahl natürlicher Zahlen giebt, die kleiner sind als $N_a(\omega)$, so muss die Reihe der Functionale $\omega, \omega_1, \omega_2, \ldots$ abbrechen, und dies ist nur möglich, wenn das letzte von ihnen ein Primfunctional ist, wodurch der Satz 1. bewiesen ist.

2. \gesperrt{Jedes ganze von Null und von den Einheiten verschiedene Functional $\omega$ im Körper $\Omega$ kann in eine endliche Anzahl von Primfactoren zerlegt werden.}

Ist nämlich $\pi_1$ ein Primfactor von $\omega$ und
\weq{4}{\omega = \pi_1 \omega_1,}
so ist, da $N_a(\pi_1) > 1$ ist,
\[
N_a(\omega) > N_a(\omega_1).
\]

Ist $\omega_1$ keine Einheit, so ist es durch ein Primfunctional $\pi_2$ theilbar, und aus
\[
\omega_1 = \pi_2 \omega_2
\]
folgt
\[
N_a(\omega_1) > N_a(\omega_2).
\]

Fährt man so fort, so erhält man eine Reihe ganzer Functionale $\omega_1, \omega_2, \ldots$, deren absolute Normen fortwährend abnehmen, und diese Reihe bricht also mit einer Einheit ab. Ist $\pi_\nu$ die letzte von ihnen, die keine Einheit ist, so ist $\pi_\nu$ selbst ein Primfunctional, und es folgt
\weq{5}{\omega = \pi_1 \pi_2 \ldots \pi_\nu,}
w. z. b. w.


% ===== scan leaf 539, printed 518 =====
% !PRINTED 518
\weberpage{518}
3. \gesperrt{Ein ganzes Functional $\omega$ im Körper $\Omega$ ist nur auf eine Weise in Primfactoren zerlegbar, wenn associirte Primfactoren als nicht verschieden betrachtet werden. Associirte Functionale enthalten dieselben Primfactoren.}

Nehmen wir nämlich an, es seien die beiden Producte von Primfactoren
\weq{6}{\pi_1 \pi_2 \ldots \pi_\nu, \qquad \varkappa_1 \varkappa_2 \ldots \varkappa_\mu}
mit einander associirt, so ist das Product $\pi_1 \pi_2 \ldots \pi_\nu$ durch den Primfactor $\varkappa_1$ theilbar, und es muss daher, nach \S.\,140, 2., einer der Factoren, etwa $\pi_1$, durch $\varkappa_1$ theilbar und folglich mit $\varkappa_1$ associirt sein. Dann sind auch die Producte
\[
\pi_2 \ldots \pi_\nu, \qquad \varkappa_2 \ldots \varkappa_\mu
\]
associirt, und folglich ist einer der Factoren des ersten Productes, etwa $\pi_2$, durch $\varkappa_2$ theilbar und daher mit $\varkappa_2$ associirt. So kann man weiter schliessen, und es ergiebt sich, dass nicht nur die Anzahl der $\varkappa$ mit der Anzahl der $\pi$ übereinstimmen muss, sondern dass auch die $\varkappa$ einzeln den $\pi$ associirt sind.

Unter den Primfactoren eines ganzen Functionales $\omega$ kann derselbe mehrmals vorkommen, und diese einander gleichen (oder associirten) Factoren können zu einer Potenz zusammengefasst werden. Ist $\omega$ durch $\pi^h$ theilbar, so sagen wir, \gesperrt{das Primfunctional $\pi$ geht $h$ mal in $\omega$ auf.}

Hiernach hat es einen ganz bestimmten Sinn, wenn von \gesperrt{den Primfactoren} eines ganzen Functionales gesprochen wird. Wir folgern noch aus dem Bewiesenen:

4. \gesperrt{Ein ganzes Functional $\alpha$ ist dann und nur dann durch ein anderes $\beta$ theilbar, wenn alle Primfactoren von $\beta$ unter den Primfactoren von $\alpha$ vorkommen, und jeder von ihnen mindestens so oft in $\alpha$ aufgeht, als in $\beta$.}

Denn ist $\alpha$ durch $\beta$ und $\beta$ durch $\pi^h$ theilbar, so ist auch $\alpha$ durch $\pi^h$ theilbar.

Sind $\alpha, \beta, \ldots$ ganze Functionale, $\pi_1, \pi_2, \ldots$ verschiedene Primfunctionale, $\varepsilon_1, \varepsilon_2, \ldots$ Einheiten, so können wir Exponenten $a_1, a_2, \ldots, b_1, b_2, \ldots$ so bestimmen, dass
\weq{7}{\begin{aligned}
\varepsilon_1 \alpha &= \pi_1^{a_1} \pi_2^{a_2} \ldots \\
\varepsilon_2 \beta &= \pi_1^{b_1} \pi_2^{b_2} \ldots \\
&\ldots\ldots\ldots\ldots
\end{aligned}}
%CONT>>


% ===== scan leaf 540, printed 519 =====
% !PRINTED 519
\weberpage{519}
%<<CONT
wird, falls wir den Exponenten gleich Null setzen, wenn einer der Primfactoren in dem betreffenden Functionale nicht aufgeht. Ein gemeinsamer Theiler der Zahlen $\alpha, \beta, \ldots$ kann keine anderen Primfactoren als $\pi_1, \pi_2, \ldots$ enthalten, und jeder gemeinsame Theiler von $\alpha, \beta, \ldots$ hat die Form
\weq{8}{\varepsilon \delta = \pi_1^a \pi_2^b \ldots,}
worin $a$ nicht grösser als die kleinste der Zahlen $a_1, b_1, \ldots$ sein darf, $b$ nicht grösser als die kleinste der Zahlen $a_2, b_2, \ldots$ u. s. f.

Ist $a$ die kleinste unter den Zahlen $a_1, b_1, \ldots$, $b$ die kleinste unter den Zahlen $a_2, b_2, \ldots$, so ist die in (8) dargestellte Zahl $\delta$ der grösste gemeinschaftliche Theiler der Zahlen $\alpha, \beta, \ldots$ In Worten ausgedrückt:

Man erhält den grössten gemeinschaftlichen Theiler mehrerer ganzer Functionale $\alpha, \beta, \ldots$, wenn man ein Product aus Primfactoren bildet, in das man jeden Primfactor so oft aufnimmt, als er in jeder der Zahlen $\alpha, \beta, \ldots$ aufgeht.

Zwei Zahlen sind relativ prim, wenn sie keinen gemeinschaftlichen Primfactor enthalten.

Dem entsprechend definiren wir als das \gesperrt{kleinste gemeinschaftliche Multiplum} $\mu$ der Functionale $\alpha, \beta, \ldots$ ein Product aus Primfactoren, in das wir einen Factor $\pi$ so oft aufnehmen, dass er in keinem der Functionale $\alpha, \beta, \ldots$ öfter als in $\mu$ aufgeht. Dieses Functional $\mu$ hat dann, ebenso wie jedes mit $\mu$ associirte Functional, die doppelte, und, wie wir hinzufügen können, charakteristische Eigenschaft, dass es durch jedes der Functionale $\alpha, \beta, \ldots$ theilbar ist, und dass jedes andere Functional, das durch $\alpha, \beta, \ldots$ theilbar ist, auch durch $\mu$ theilbar ist.

Das kleinste gemeinschaftliche Multiplum zweier relativer Primfunctionale ist ihr Product.

Man kann diese Sätze anwenden, um gebrochene Functionale \gesperrt{in der einfachsten Gestalt oder als reducirte Brüche dar}zustellen, indem man Zähler und Nenner in ihre Primfactoren zerlegt und den grössten gemeinschaftlichen Theiler weghebt. Zähler und Nenner eines reducirten Bruches sind durch den Bruch selbst völlig bestimmt, abgesehen von einem gemeinschaftlichen Einheitsfactor, der unbestimmt bleibt.

Ebenso kann man eine beliebige Zahl gegebener Brüche auf gemeinsamen Nenner, den Hauptnenner, bringen, indem man ein gemeinsames Multiplum aller gegebenen Nenner als gemeinschaftlichen Nenner wählt.


% ===== scan leaf 541, printed 520 =====
% !PRINTED 520
\weberpage{520}
\pdfbookmark[2]{§ 142. Ganze Functionen im Körper Ω}{bm162}
Sind die gegebenen Functionale wirkliche Brüche, d. h. reduciren sie sich nicht alle auf ganze Functionale, so muss der gemeinsame Nenner wenigstens einen Primfactor enthalten, der nicht in allen Zählern vorkommt.

Alles das ist in vollkommener Uebereinstimmung mit den Regeln der elementaren Arithmetik, und auch die Beweismethoden, die wir hier angewandt haben, sind wesentlich dieselben, die dort gebraucht werden. Der Kernpunkt der Deduction ist einerseits die weitgehende Verallgemeinerung des Begriffes der Einheit, andererseits die darauf gegründete Definition des grössten gemeinschaftlichen Theilers im \S.\,139.

\begin{center}
\S.\,142.\\
\gesperrt{Ganze Functionen im Körper $\overline{\Omega}$.}
\end{center}

Die Hülfsmittel, über die wir jetzt verfügen, reichen aus, um die Schlussweise, deren wir uns im \S.\,2 des ersten Bandes zum Beweis des \gesperrt{Gauss}'schen Satzes bedient haben, auf die Functionale anzuwenden.

In der That ist ja der Satz, auf den sich jener Beweis wesentlich stützt, dass ein Product zweier ganzer Zahlen nur dann durch eine Primzahl theilbar ist, wenn diese Primzahl in einem der Factoren aufgeht, auch für die jetzt eingeführten Primfunctionale als gültig erwiesen.

Wir können demnach, indem wir dem erwähnten Beweise Schritt für Schritt folgen, den Satz als erwiesen ansehen:

1. \gesperrt{Wenn zwei ganze Functionen einer Variablen $t$ der Grade $h$ und $k$:}
\[
\varphi = \alpha_0 t^h + \alpha_1 t^{h-1} + \cdots + \alpha_h
\]
\[
\psi = \beta_0 t^k + \beta_1 t^{k-1} + \cdots + \beta_k,
\]
\gesperrt{deren Coëfficienten ganze Functionale sind, die die Variable $t$ nicht enthalten, ein Product}
\[
\chi = \gamma_0 t^{h+k} + \gamma_1 t^{h+k-1} + \cdots + \gamma_{h+k}
\]
\gesperrt{haben, in dem die Coëfficienten $\gamma_0, \gamma_1, \ldots, \gamma_{h+k}$ einen gemeinschaftlichen Primtheiler $\pi$ haben, so muss $\pi$ entweder in allen Coëfficienten von $\varphi$ oder in allen Coëfficienten von $\psi$ aufgehen.}


% ===== scan leaf 542, printed 521 =====
% !PRINTED 521
\weberpage{521}

Aus diesem Satze ziehen wir hier eine wichtige Folgerung: Die Functionen

\weq{1}{\varphi = \varphi_0 t^h + \varphi_1 t^{h-1} + \cdots + \varphi_h,}

in denen die Co\"efficienten $\varphi_0, \varphi_1, \ldots, \varphi_h$ ganze oder gebrochene Functionale in $\Omega$ sind, aber von den Variablen $t$ frei angenommen werden, geh\"oren selbst zu den Functionalen im K\"orper $\Omega$. Von ihnen gilt der Satz:

2. \gesperrt{Ein Functional $\varphi$ ist nur dann ganz, wenn die Co\"efficienten $\varphi_0, \varphi_1, \ldots, \varphi_h$ ganze Functionale sind.}

Um ihn zu beweisen, nehmen wir an, es seien $\varphi_0, \varphi_1, \ldots, \varphi_h$ nicht alle zugleich ganz. Bestimmen wir ihren Hauptnenner $\mu$ und setzen

\[
\mu\varphi_0 = \alpha_0, \quad \mu\varphi_1 = \alpha_1, \quad \ldots, \quad \mu\varphi_h = \alpha_h,
\]

so sind $\alpha_0, \alpha_1, \ldots, \alpha_h$ ganze Functionale, und $\mu$ enth\"alt wenigstens einen Primfactor $\pi$, der nicht zugleich in allen Z\"ahlern $\alpha_0,, \alpha_1, \ldots, \alpha_h$ aufgeht % [sic] doubled comma in print
(vgl. den Schluss des vor. Paragraphen).

Dann ist die Function

\weq{2}{\chi = \mu\varphi = \alpha_0 t^h + \alpha_1 t^{h-1} + \cdots + \alpha_h,}

deren Co\"efficienten nun ganz sind, gewiss ein ganzes Functional.

Nehmen wir nun an, es sei $\varphi$ selbst ganz, so ist $\mu\varphi$ durch $\pi$ theilbar, w\"ahrend doch nicht s\"ammtliche Co\"efficienten $\alpha_0, \alpha_1, \ldots, \alpha_h$ durch $\pi$ theilbar sind. Wenn nun $\varphi$ ein ganzes Functional ist, so gen\"ugt es einer Gleichung von der Form

\weq{3}{\varphi^m = C_1 \varphi^{m-1} + C_2 \varphi^{m-2} + \cdots + C_m,}

in der die Co\"efficienten $C_1, C_2, \ldots, C_m$ ganze rationale Functionale sind.

Diese Functionale setzen wir nach \S.\,136, (2) in die Form

\[
C_1 = \frac{a_1 E_1}{E}, \quad C_2 = \frac{a_2 E_2}{E}, \quad \ldots, \quad C_m = \frac{a_m E_m}{E},
\]

worin $a_1, a_2, \ldots, a_m$ die absoluten Werthe von $C_1, C_2, \ldots, C_m$, also nat\"urliche ganze Zahlen (oder Null), und $E, E_1, \ldots, E_m$ primitive ganze Functionen, also Einheiten, sind.

Dann ergiebt die Gleichung (3):

\[
E\varphi^m = a_1 E_1 \varphi^{m-1} + a_2 E_2 \varphi^{m-2} + \cdots + a_m E_m,
\]

und durch Multiplication mit $\mu^m$:

\weq{4}{E\chi^m = \mu\left(a_1 E_1 \chi^{m-1} + a_2 E_2 \mu \chi^{m-2} + \cdots + a_m E_m \mu^{m-1}\right).}


% ===== scan leaf 543, printed 522 =====
% !PRINTED 522
\weberpage{522}

Hier stehen nun rechter und linker Hand ganze Functionen von $t$, deren Co\"efficienten ganze Functionale sind. Auf der rechten Seite haben alle diese Co\"efficienten den Factor $\mu$, also auch den Factor $\pi$, w\"ahrend nach Voraussetzung nicht alle Co\"efficienten von $\chi$ diesen Factor haben. Da $E$ eine Einheit ist, so enthalten auch die Co\"efficienten von $E$ den Factor $\pi$ nicht, und folglich k\"onnen nach dem Satze 1. auch die Co\"efficienten von $E\chi^m$ nicht alle durch $\pi$ theilbar sein, was doch die Gleichung (4) verlangen w\"urde. Daraus ergiebt sich, dass unsere Annahme, $\varphi$ sei ganz, $\varphi_0, \varphi_1, \ldots, \varphi_h$ dagegen nicht alle ganz, unstatthaft ist, und der Satz 2. ist somit bewiesen.

Nehmen wir nun an, in $\varphi$ seien die Co\"efficienten $\varphi_0, \varphi_1, \ldots$ selbst wieder ganze Functionen einer Variablen, und wenden den Satz 2. wiederholt darauf an, so gelangen wir zu dem Schlusse:

3. \gesperrt{Eine ganze rationale Function beliebig vieler Ver\"anderlicher, deren Co\"efficienten Zahlen oder Functionale mit anderen Variablen sind, ist nur dann ein ganzes Functional, wenn die Co\"efficienten ganz sind.}

Wir k\"onnen jedes Functional $\omega$ als Quotienten zweier ganzer Functionen in der Weise darstellen, dass der Nenner eine primitive Function im K\"orper $R$ wird.

Denn sind $\varphi, \psi$ ganze Functionen in $\Omega$, und ist

\[
\omega = \frac{\varphi}{\psi}, \quad N(\psi) = \psi\psi',
\]

so k\"onnen wir den Bruch $\omega$ durch $\psi'$ erweitern und erhalten im Nenner eine ganze Function mit rationalen Co\"efficienten. Den Theiler dieser Function k\"onnen wir dann zum Z\"ahler von $\omega$ rechnen, und erhalten, wenn $E$ eine primitive Function, $\chi$ eine ganze Function in $\Omega$ bedeutet:

\[
E\omega = \chi,
\]

d. h. man kann jedes Functional $\omega$ in $\Omega$ durch Multiplication mit einer primitiven Function $E$ in eine ganze Function der Variablen verwandeln, deren Co\"efficienten Zahlen in $\Omega$ sind.

Mit Zuziehung des Satzes 3. ergiebt sich hieraus:

4. \gesperrt{Jedes ganze Functional $\omega$ im K\"orper $\Omega$ ist associirt mit einer ganzen Function $\chi$, deren Co\"effi-}
%CONT>>


% ===== scan leaf 544, printed 523 =====
% !PRINTED 523
\weberpage{523}
\pdfbookmark[2]{§ 143. Die Primfactoren der Zahlen des Körpers Ω}{bm163}

%<<CONT
\gesperrt{cienten ganze Zahlen in $\Omega$ sind, und geht durch Multiplication mit einem rationalen Functional, welches eine Einheit ist, in $\chi$ \"uber.}

\begin{center}
\S.\,143.

\gesperrt{Die Primfactoren der Zahlen des K\"orpers $\Omega$.}
\end{center}

Da unter den Functionalen des K\"orpers $\Omega$ auch die Zahlen enthalten sind, so ergiebt sich aus den Resultaten des \S.\,141, dass sich auch die ganzen Zahlen des K\"orpers $\Omega$ in Primfactoren zerlegen lassen, aber in Primfactoren, die im Allgemeinen nicht Zahlen, sondern Functionale sind. Es ist aber nicht ausgeschlossen, dass in besonderen F\"allen unter den Primfunctionalen auch Zahlen auftreten k\"onnen, die dann \gesperrt{Primzahlen des K\"orpers $\Omega$} heissen.

Alle Gleichungen zwischen Functionalen sind in letzter Instanz identische Gleichungen zwischen ganzen rationalen Functionen und l\"osen sich in eine Reihe von Gleichungen zwischen Zahlen auf. Sie bleiben also richtig, wenn f\"ur die Variablen andere Zeichen gesetzt werden.

Es folgt daraus, dass ein ganzes Functional, eine Einheit, ein Primfunctional nicht aufh\"oren, ganze Functionale, Einheiten, Primfunctionale zu sein, wenn f\"ur die Variablen andere Symbole gesetzt werden.

Zerlegen wir also eine Zahl in ihre Primfactoren, so k\"onnen, wenn unter diesen Primfactoren Functionale vorkommen, in der diese Zerlegung darstellenden Gleichung die Variablen durch beliebige andere Variable ersetzt werden, und daraus folgt die Verallgemeinerung:

1. \gesperrt{Man kann die Primfactoren eines ganzen Functionals $\omega$ so darstellen, dass sie die Variablen, von denen $\omega$ abh\"angt, nicht enthalten.}

Denn jedes ganze Functional $\omega$ ist Theiler von Zahlen, sogar von rationalen Zahlen, z. B. von der absoluten Norm von $\omega$. Die Primfactoren von $\omega$ sind daher unter den Primfactoren einer dieser Zahlen zu suchen und k\"onnen also durch Variable dargestellt werden, die von den in $\omega$ vorkommenden verschieden sind.


% ===== scan leaf 545, printed 524 =====
% !PRINTED 524
\weberpage{524}

Nach \S.\,142, 4. k\"onnen wir, wenn $\omega$ ein gegebenes ganzes Functional, $E$ eine Einheit, $\varphi$ eine ganze Function in $\Omega$ ist, setzen:

\weq{1}{E\omega = \varphi(x, y, \ldots).}

Wenn wir andererseits $\omega$ in seine Primfactoren zerlegen, so k\"onnen wir nach 1. diese Primfactoren von den Variablen $x, y, \ldots$ frei annehmen, und wenn wir wieder das Product dieser Primfactoren bilden, so erhalten wir eine mit $\omega$ associirte Zahl $\omega_1$, die von den Variablen $x, y, \mathinner{.\,.}$ frei ist.

Ordnen wir den Quotienten $\varphi : \omega_1$, der eine ganze Zahl (sogar eine Einheit) ist, nach den Variablen $x, y, \ldots$, so schliessen wir aus \S.\,142, 3., dass alle Co\"efficienten der Function $\varphi$ durch $\omega_1$ und folglich auch durch $\omega$ und $\varphi$ theilbar sind.

Wenn andererseits alle Co\"efficienten von $\varphi$ durch irgend einen Factor $\delta$ in $\Omega$ theilbar sind, so ist auch $\varphi$ durch $\delta$ theilbar, und daraus ergiebt sich:

2. \gesperrt{Eine ganze Function $\varphi$ mit ganzen Zahlen als Co\"efficienten ist der gr\"osste gemeinschaftliche Theiler aller ihrer Co\"efficienten.}

Die ganze Zahl $\varphi$ geht also nach \S.\,139 in eine associirte Zahl \"uber, wenn die einzelnen Potenzen und Producte der Variablen durch je eine Variable ersetzt werden. Daraus ergiebt sich auch als Corollar, dass jedes Functional $\omega$ in ein associirtes \"ubergeht, wenn die Variablen irgendwie anders bezeichnet werden.

Es ist nun der folgende wichtige Satz zu beweisen:

3. \gesperrt{Ist $\omega$ ein beliebiges ganzes Functional, so kann man eine durch $\omega$ theilbare Zahl $\alpha$ so w\"ahlen, dass der Quotient $\alpha : \omega$ zu einem beliebig gegebenen Functionale $\mu$ relativ prim ist.}

Wir beweisen zun\"achst, dass es eine ganze Zahl $\alpha$ giebt, die durch $\omega$, aber nicht durch $\omega\pi$ theilbar ist, wenn $\pi$ ein beliebiges Primfunctional ist.

Bilden wir n\"amlich nach \S.\,142, 4. eine mit $\omega$ associirte ganze Function $\varphi$, so sind die Co\"efficienten dieser Function zwar alle durch $\omega$, aber nicht alle durch $\omega\pi$ theilbar, weil sonst auch $\varphi$ und mithin $\omega$ selbst durch $\omega\pi$ theilbar w\"are, was nicht m\"oglich ist. Es giebt also unter den Co\"efficienten von $\varphi$ wenigstens einen, der die verlangte Eigenschaft hat.


% ===== scan leaf 546, printed 525 =====
% !PRINTED 525
\weberpage{525}
Es sei jetzt $\pi_1, \pi_2, \pi_3, \ldots$ eine beliebige Anzahl von einander verschiedener gegebener Primfunctionale. Wir setzen

\[
\omega_1 = \omega\pi_2\pi_3\ldots, \quad \omega_2 = \omega\pi_1\pi_3\ldots, \quad \omega_3 = \omega\pi_1\pi_2\ldots,
\]
und bestimmen nach dem, was soeben bewiesen ist, die ganzen Zahlen $\alpha_1, \alpha_2, \alpha_3, \ldots$ in $\Omega$, so dass
\[
\begin{array}{l}
\alpha_1 \text{ theilbar wird durch } \omega_1 \text{, aber nicht durch } \omega_1\pi_1\\
\alpha_2 \quad \text{\textquotedbl} \quad \text{\textquotedbl} \quad \text{\textquotedbl} \quad \omega_2 \quad \text{\textquotedbl} \quad \text{\textquotedbl} \quad \text{\textquotedbl} \quad \omega_2\pi_2\\
\alpha_3 \quad \text{\textquotedbl} \quad \text{\textquotedbl} \quad \text{\textquotedbl} \quad \omega_3 \quad \text{\textquotedbl} \quad \text{\textquotedbl} \quad \text{\textquotedbl} \quad \omega_3\pi_3\\
\cdots\cdots\cdots\cdots\cdots\cdots\cdots\cdots\cdots\cdots\cdots\cdots\cdots\cdots\cdots\cdots,
\end{array}
\]
und leiten daraus die ganze Zahl
\[
\alpha = \alpha_1 + \alpha_2 + \alpha_3 + \cdots
\]
ab. Diese Zahl ist offenbar theilbar durch $\omega$, da alle Summanden $\alpha_1, \alpha_2, \alpha_3, \ldots$ durch $\omega$ theilbar sind. Sie ist aber nicht theilbar durch $\omega\pi_1$, weil zwar $\alpha_2, \alpha_3, \ldots$, nicht aber $\alpha_1$ durch $\omega\pi_1$ theilbar ist; und ebenso ist sie nicht durch $\omega\pi_2, \omega\pi_3, \ldots$ theilbar. Wenn wir also
\[
\alpha = \omega\eta
\]
setzen, so ist $\eta$ ein ganzes Functional, das nicht durch $\pi_1, \pi_2, \pi_3, \ldots$ theilbar ist, und das daher, wenn $\pi_1, \pi_2, \pi_3, \ldots$ die von einander verschiedenen Primfactoren von $\mu$ sind, relativ prim zu $\mu$ ist.

Nehmen wir nun beliebig eine durch $\omega$ theilbare ganze Zahl $\beta$ in $\Omega$ an und setzen $\beta = \omega\mu$, dann können wir $\alpha = \omega\eta$ so bestimmen, dass $\eta$ relativ prim zu $\mu$ wird, und dann ist $\omega$ der grösste gemeinschaftliche Theiler von $\alpha$ und $\beta$. Daraus folgt:

4. \gesperrt{Jedes ganze Functional $\omega$ des Körpers $\Omega$ ist der grösste gemeinschaftliche Theiler zweier ganzer Zahlen, und folglich ist $\omega$ associirt mit einer binären Linearform $\alpha x + \beta y$, in der $\alpha$ und $\beta$ ganze Zahlen sind.}

\begin{center}
\rule{4cm}{0.4pt}
\end{center}

Es mag hier noch eine allgemeine, auf ein anderes Gebiet hinübergreifende Bemerkung ihren Platz finden.

Es ist das Hauptergebniss dieses Abschnittes, dass sich die ganzen algebraischen Zahlen in einem bestimmten Körper in eindeutiger Weise in Primfactoren zerlegen lassen, genau in derselben Weise, wie dies bei den ganzen rationalen Zahlen bekannt ist; freilich aber nur dadurch, dass der Inhalt des Körpers
%CONT>>


% ===== scan leaf 547, printed 526 =====
% !PRINTED 526
\weberpage{526}
%<<CONT
(durch Adjunction von Variablen) vergrössert wird. Es entsteht so ein erweiterter Körper, in dem die Gesetze der Zerlegbarkeit rein gelten.

Den Ausgangspunkt der Definitionen bildete der Körper $R$ der rationalen Zahlen, und wenn wir uns Rechenschaft darüber geben wollen, auf welchen Eigenschaften des Körpers $R$ die Möglichkeit dieser Erweiterung beruht, so finden wir, dass es einerseits die Existenz der ganzen Zahlen in $R$, andererseits die eindeutige Zerlegbarkeit dieser ganzen Zahlen in Primfactoren ist, die allein bei der ganzen Deduction benutzt wurden. Wenn wir also an Stelle des Körpers $R$ irgend einen anderen Körper treten lassen, dem diese beiden Eigenschaften zukommen, so werden wir dieselben Folgerungen ziehen können. Nehmen wir für $R$ einen anderen algebraischen Zahlkörper, so bekommen wir freilich nichts Neues, wohl aber, wenn wir z.\,B. an Stelle des Körpers $R$ den Körper der rationalen Functionen einer Variablen setzen, dem ja die beiden fundamentalen Eigenschaften auch zukommen. So gewinnen wir einen Ausgangspunkt für die Theorie der algebraischen Functionen einer Variablen\footnote{Vergl. \gesperrt{Dedekind-Weber}, Theorie der algebraischen Functionen einer Veränderlichen. \gesperrt{Crelle's} Journal, Bd. 92.}.

\begin{center}
\rule{4cm}{0.4pt}
\end{center}


% ===== scan leaf 548, printed 527 =====
\weberpage{527}% page number not printed on this page
\pdfbookmark[1]{Siebzehnter Abschnitt. Theorie der algebraischen Körper.}{bm164}
\pdfbookmark[2]{§ 144. Basis eines algebraischen Zahlkörpers. Discriminanten}{bm165}
% !PRINTED none

\begin{center}
\gesperrt{Siebzehnter Abschnitt.}

\gesperrt{Theorie der algebraischen K\"orper.}

\rule{4cm}{0.4pt}

\S.\,144.

\gesperrt{Basis eines algebraischen Zahlk\"orpers.}

\gesperrt{Discriminanten.}
\end{center}

Es sei
\weq{1}{f(\Theta) = \Theta^n + a_1\Theta^{n-1} + \cdots + a_n = 0}
die irreducible Gleichung $n^{\text{ten}}$ Grades mit rationalen Co\"efficienten $a_1, a_2, \ldots, a_n$, die uns einen algebraischen K\"orper $\Omega = R(\Theta)$ definirt. Der K\"orper $\Omega$ ist dann der Inbegriff aller Zahlen der Form
\weq{2}{\omega = h_1 + h_2\Theta + h_3\Theta^2 + \cdots + h_n\Theta^{n-1},}
worin $h_1, h_2, \ldots, h_n$ \gesperrt{rationale Zahlen} sind. Setzen wir aber f\"ur $h_1, h_2, \ldots, h_n$ rationale Functionale, so erhalten wir aus (2) alle \gesperrt{Functionale} des K\"orpers $\Omega$.

Betrachten wir ein beliebiges System von $n$ Zahlen
\weq{3}{\omega_r = h_{1,r} + h_{2,r}\Theta + \ldots + h_{n,r}\Theta^{n-1}, \quad r = 1,2,\ldots,n}
unter der Voraussetzung, dass die Determinante
\weq{4}{H = \Sigma \pm h_{1,1}h_{2,2}\ldots h_{n,n}}
von Null verschieden ist, so kann, wegen der Irreducibilit\"at von $f$, eine Gleichung der Form
\weq{5}{k_1\omega_1 + k_2\omega_2 + \cdots + k_n\omega_n = 0,}
in der $k_1, k_2, \ldots, k_n$ rationale Zahlen sind, nur dann bestehen, wenn diese Co\"efficienten alle Null sind, und dies gilt auch dann noch, wenn in der Gleichung (5) f\"ur die Co\"efficienten \gesperrt{rationale Functionale} zugelassen werden.


% ===== scan leaf 549, printed 528 =====
% !PRINTED 528
\weberpage{528}

Eliminiren wir aber aus den Gleichungen (2) und (3) die Potenzen von $\Theta$, so ergiebt sich eine Relation von der Form
\weq{6}{\omega = k_1\omega_1 + k_2\omega_2 + \cdots + k_n\omega_n.}

Man kann also jede Zahl und jedes Functional des Körpers $\Omega$ in der Form (6) darstellen, wenn man für $k_1, k_2, \ldots, k_n$ rationale Zahlen oder Functionale setzt. Diese Darstellung ist für ein gegebenes $\omega$ [wegen (5)] nur auf eine Art möglich, und $\omega$ ist eine Zahl, wenn die $k_1, k_2, \ldots, k_n$ Zahlen sind, dagegen ein Functional, wenn unter den $k$ auch Functionale vorkommen.

Ein solches System von Zahlen, wie
\[
\omega_1, \omega_2, \ldots, \omega_n,
\]
nennen wir eine \gesperrt{Basis des Körpers $\Omega$}. Eine solche Basis bilden auch die Potenzen von $\Theta$:
\[
1, \Theta, \Theta^2, \ldots, \Theta^{n-1}.
\]

Bezeichnen wir mit $\omega_{r,1}, \omega_{r,2}, \ldots, \omega_{r,n}$ die mit $\omega_r$ conjugirten Zahlen, so ist das Determinantenquadrat
\weq{7}{\Delta(\omega_1, \omega_2, \ldots, \omega_n) = (\Sigma \pm \omega_{1,1}\omega_{2,2}\ldots\omega_{n,n})^2}
eine symmetrische Function der conjugirten Werthe $\Theta_1, \Theta_2, \ldots, \Theta_n$ und ist folglich eine rationale Zahl.

Diese Zahl heisst die \gesperrt{Discriminante des Systemes} $\omega_1, \omega_2, \ldots, \omega_n$. Insbesondere ist
\weq{8}{\Delta(1, \Theta, \Theta^2, \ldots, \Theta^{n-1}) = \begin{vmatrix} 1, & \Theta_1, & \ldots, & \Theta_1^{\,n-1}\\ 1, & \Theta_2, & \ldots, & \Theta_2^{\,n-1}\\ \multicolumn{4}{c}{\cdots\cdots\cdots\cdots\cdots\cdots}\\ 1, & \Theta_n, & \ldots, & \Theta_n^{\,n-1} \end{vmatrix}^2}
die Discriminante der Gleichung (1) (Bd.\ I, \S.\,46), für die man auch, wenn $N$ das Zeichen für die Norm ist,
\weq{9}{(-1)^{\frac{n(n-1)}{2}} N f'(\Theta)}
setzen kann. Diese Zahl ist also sicher von Null verschieden. Nach dem Multiplicationssatze der Determinanten (Bd.\ I, \S.\,27) ergiebt sich aus (3) und (7):
\weq{10}{\Delta(\omega_1, \omega_2, \ldots, \omega_n) = H^2\Delta(1, \Theta, \Theta^2, \ldots, \Theta^{n-1}),}
woraus man schliesst, dass die Discriminante einer Basis von $\Omega$ immer von Null verschieden ist. Da $H$ eine rationale Zahl ist, so folgt, dass das Verhältniss der Discriminanten verschiedener Basen das Quadrat einer rationalen Zahl ist, und dass die
%CONT>>


% ===== scan leaf 550, printed 529 =====
% !PRINTED 529
\weberpage{529}
\pdfbookmark[2]{§ 145. Die Minimalbasis und die Körperdiscriminante}{bm166}
%<<CONT
Discriminanten aller Basen von $\Omega$ dasselbe Vorzeichen haben. Die Formel (10) zeigt auch, dass irgend ein System von $n$ Zahlen $\omega_r$ des Körpers $\Omega$ immer dann eine Basis von $\Omega$ ist, wenn das Determinantenquadrat (7) nicht verschwindet.

Ist $\omega_1, \omega_2, \ldots, \omega_n$ eine Basis von $\Omega$, und bedeuten $c_{r,s}$ rationale Zahlen, so bilden auch die $n$ Zahlen
\[
\omega_r' = c_{r,1}\omega_1 + c_{r,2}\omega_2 + \cdots + c_{r,n}\omega_n \qquad r = 1, 2, \ldots, n
\]
eine Basis von $\Omega$, wenn die Determinante
\[
C = \Sigma \pm c_{1,1}c_{2,2} \ldots c_{n,n}
\]
nicht verschwindet, und es ist
\weq{11}{\varDelta(\omega_1', \omega_2', \ldots, \omega_n') = C^2 \varDelta(\omega_1, \omega_2, \ldots, \omega_n).}

Wenn wir z. B. die Elemente $\omega_1, \omega_2, \ldots, \omega_n$ einer Basis mit rationalen Coëfficienten $c_1, c_2, \ldots, c_n$ multipliciren, deren keiner verschwindet, so erhalten wir eine neue Basis
\[
c_1\omega_1, c_2\omega_2, \ldots, c_n\omega_n.
\]

\begin{center}
\S.\,145.

\gesperrt{Die Minimalbasis und die Körperdiscriminante.}
\end{center}

Nach der zuletzt gemachten Bemerkung verliert eine Basis von $\Omega$ die Eigenschaft, eine Basis zu sein, nicht, wenn man jede ihrer Zahlen mit einer von Null verschiedenen rationalen Zahl multiplicirt. Nun kann man nach \S.\,133,5. jede Zahl durch Multiplication mit einer ganzen rationalen Zahl in eine ganze Zahl verwandeln, und daraus folgt, \gesperrt{dass es Basen von $\Omega$ giebt, deren Elemente lauter ganze Zahlen sind. Die Discriminante einer solchen Basis ist eine ganze rationale Zahl.} Diese ganze rationale Zahl ist von Null verschieden. Sie ändert sich, wenn eine andere ganzzahlige Basis gewählt wird, behält aber für einen bestimmten Körper ein unverändertes Vorzeichen.

Unter all diesen ganzen Zahlen, die als Discriminanten einer ganzzahligen Basis auftreten können, und die alle in quadratischem Verhältniss zu einander stehen, muss nun eine dem absoluten Werthe nach die kleinste sein. Diese kleinste Discriminante bezeichnen wir mit $\varDelta$ und nennen sie die \gesperrt{Grundzahl} oder auch die \gesperrt{Discriminante des Körpers} $\Omega$.


% ===== scan leaf 551, printed 530 =====
% !PRINTED 530
\weberpage{530}
Dies $\varDelta$ ist eine durch $\Omega$ völlig bestimmte positive oder negative, aber niemals verschwindende ganze rationale Zahl, und es giebt immer eine aus ganzen Zahlen $\omega_1, \omega_2, \ldots, \omega_n$ bestehende Basis von $\Omega$, deren Discriminante gleich $\varDelta$ ist.

Eine solche Basis wollen wir kurz eine \gesperrt{Minimalbasis von} $\Omega$ nennen.

Verstehen wir unter $k_1, k_2, \ldots, k_n$ irgend welche ganze rationale Zahlen, so ist jede Zahl von der Gestalt
\weq{1}{\omega = k_1\omega_1 + k_2\omega_2 + \cdots + k_n\omega_n}
eine ganze algebraische Zahl, und wir beweisen jetzt den fundamentalen Satz:

1. \gesperrt{Wenn $\omega_1, \omega_2, \ldots, \omega_n$ eine Minimalbasis ist, so sind in der Form (1) alle ganzen Zahlen des Körpers $\Omega$ enthalten.}

Da $\omega_1, \omega_2, \ldots, \omega_n$ eine Basis ist, so kann zunächst jede Zahl in $\Omega$ in der Form (1) dargestellt werden, wenn für $k_1, k_2, \ldots, k_n$ rationale Brüche zugelassen werden. Nehmen wir also an, es sei eine ganze Zahl $\omega$ in der Form
\weq{2}{\omega = \frac{k_1\omega_1 + k_2\omega_2 + \cdots + k_n\omega_n}{k}}
darstellbar, worin $k_1, k_2, \ldots, k_n, k$ ganze rationale Zahlen sind, so dass nicht alle $k_1, k_2, \ldots, k_n$ mit $k$ einen gemeinschaftlichen Theiler haben. Ist $p$ irgend eine in $k$ aufgehende natürliche Primzahl und $k = pk'$, so muss wenigstens einer der Coëfficienten $k_1, k_2, \ldots, k_n$ durch $p$ untheilbar sein. Es sei etwa $k_1$ durch $p$ nicht theilbar; dann lässt sich die ganze rationale Zahl $l$ so bestimmen, dass $l k_1 \equiv 1 \; (\mathrm{mod}\ p)$, oder $(l k_1 - 1)$ durch $p$ theilbar wird. Es folgt dann aus (2):
\weq{3}{l k' \omega - \frac{l k_1 - 1}{p}\omega_1 = \frac{\omega_1 + l k_2\omega_2 + \cdots + l k_n\omega_n}{p} = \omega_1',}
und $\omega_1'$ ist gleichfalls eine ganze algebraische Zahl. Setzen wir
\weq{4}{\omega_2' = \omega_2, \ldots, \omega_n' = \omega_n,}
so bilden die Zahlen $\omega_1', \omega_2', \ldots, \omega_n'$ eine \gesperrt{ganzzahlige Basis} von $\Omega$, weil sich die Zahlen $\omega_1, \omega_2, \ldots, \omega_n$ und folglich alle Zahlen $\omega$ linear durch $\omega_1', \omega_2', \ldots, \omega_n'$ ausdrücken lassen, und die Formel (11) des vorigen Paragraphen ergiebt
\weq{5}{\varDelta(\omega_1', \omega_2', \ldots, \omega_n') = \frac{1}{p^2}\varDelta(\omega_1, \omega_2, \ldots, \omega_n).}


% ===== scan leaf 552, printed 531 =====
% !PRINTED 531
\weberpage{531}
Die Discriminante $\varDelta(\omega_1', \omega_2', \ldots, \omega_n')$ ist also kleiner als $\varDelta(\omega_1, \omega_2, \ldots, \omega_n)$, und dies widerspricht der Annahme, dass $\omega_1, \omega_2, \ldots, \omega_n$ eine Minimalbasis sei. Damit ist unser Satz erwiesen.

Bezeichnet man die Gesammtheit aller \gesperrt{ganzen Zahlen} des Körpers $\Omega$ mit $\mathfrak{o}$, so erhält man alle Zahlen von $\mathfrak{o}$, und jede nur einmal, wenn man in
\weq{6}{k_1\omega_1 + k_2\omega_2 + \cdots + k_n\omega_n}
die Coefficienten $k_1, k_2, \ldots, k_n$ die sämmtlichen ganzen rationalen Zahlen durchlaufen lässt.

Aus diesem Grunde wird eine Minimalbasis von $\Omega$ auch eine \gesperrt{Basis von $\mathfrak{o}$} genannt.

\gesperrt{Die Discriminante einer Basis von $\mathfrak{o}$ ist gleich der Grundzahl des Körpers $\Omega$.}

Nach der Formel (11) des vorigen Paragraphen können wir aus einer Basis von $\mathfrak{o}$ beliebig viele andere durch lineare Substitution ableiten:
\weq{7}{(\omega_1', \omega_2', \ldots, \omega_n') = C(\omega_1, \omega_2, \ldots, \omega_n),}
wenn $C$ (nach den Bezeichnungen des \S.\,37) eine lineare Substitution mit rationalen ganzzahligen Coëfficienten und der Determinante $\pm 1$ bedeutet.

Da nach der Bedeutung der Basis von $\mathfrak{o}$ alle ganzen Zahlen in $\Omega$, also auch die Elemente $\omega_1', \omega_2', \ldots, \omega_n'$ einer zweiten Basis von $\mathfrak{o}$ linear mit ganzen rationalen Coëfficienten durch $\omega_1, \omega_2, \ldots, \omega_n$ darstellbar sind, so folgt auch umgekehrt, dass man durch solche lineare Substitutionen mit der Determinante $\pm 1$ aus einer Basis von $\mathfrak{o}$ alle anderen ableiten kann.

Denn ist
\[
(\omega_1, \omega_2, \ldots, \omega_n) = C'(\omega_1', \omega_2', \ldots, \omega_n'),
\]
so muss die zusammengesetzte Substitution $CC'$ die identische sein, und das Product beider Determinanten ist also 1, also jede von ihnen $= \pm 1$.

Da unter den Zahlen von $\mathfrak{o}$ immer die Zahl 1 enthalten ist, so kann man, wenn $\omega_1, \omega_2, \ldots, \omega_n$ eine Basis von $\mathfrak{o}$ ist, die \gesperrt{ganzen rationalen Zahlen} $c_1, c_2, \ldots, c_n$ so bestimmen, dass die Relation
\weq{8}{c_1\omega_1 + c_2\omega_2 + \cdots + c_n\omega_n = 1}
befriedigt ist.


% ===== scan leaf 553, printed 532 =====
% !PRINTED 532
\weberpage{532}
\pdfbookmark[2]{§ 146. Die Basen der Functionale}{bm167}
Es gilt aber auch in Bezug auf die \gesperrt{Functionale} der Satz:

2. \gesperrt{Wenn $\omega_1, \omega_2, \ldots, \omega_n$ eine Basis von $\mathfrak{o}$ ist, so sind in der Form}
\weq{9}{\omega = u_1\omega_1 + u_2\omega_2 + \cdots + u_n\omega_n,}
\gesperrt{in der $u_1, u_2, \ldots, u_n$ ganze rationale Functionale sind, alle ganzen Functionale in $\Omega$ enthalten.}

Denn wir haben schon oben gezeigt, dass alle Functionale überhaupt in der Form (9) enthalten sind, wenn die Coëfficienten $u_1, u_2, \ldots, u_n$ ganze oder gebrochene rationale Functionale sind. Wir können aber immer eine ganze primitive Function $e$ so bestimmen, dass
\[
e u_1 = y_1, \; e u_2 = y_2, \; \ldots, \; e u_n = y_n
\]
ganze Functionen der Variablen sind, und dann wird
\weq{10}{e\omega = y_1\omega_1 + y_2\omega_2 + \cdots + y_n\omega_n,}
und da $e$ eine Einheit ist, so ist $e\omega$ zugleich mit $\omega$ ganz. Da nun die Coëfficienten der Potenzen und Producte der Variablen in der Function (10) nach \S.\,142, 3. ganze Zahlen sein müssen, so folgt nach 1., dass die Coëfficienten in den Functionen $y_1, y_2, \ldots, y_n$ ganze rationale Zahlen sein müssen, und dass folglich $u_1, u_2, \ldots, u_n$ ganze rationale Functionale sind.

\begin{center}
\S.\,146.

\gesperrt{Die Basen der Functionale.}
\end{center}

\gesperrt{Ist $\mu$ ein ganzes Functional des Körpers $\Omega$, so verstehen wir unter einer Basis von $\mu$ ein System von $n$ ganzen Zahlen in $\Omega$:}
\weq{1}{\alpha_1, \alpha_2, \ldots, \alpha_n,}
\gesperrt{das eine Basis des Körpers $\Omega$ ist, und dem die Eigenschaft zukommt, dass in der Form}
\weq{2}{\alpha = x_1\alpha_1 + x_2\alpha_2 + \cdots + x_n\alpha_n}
\gesperrt{alle durch $\mu$ theilbaren ganzen Zahlen des Körpers $\Omega$ und keine anderen enthalten sind, wenn für $x_1, x_2, \ldots, x_n$ ganze rationale Zahlen gesetzt werden.}

Es soll jetzt bewiesen werden, dass jedes ganze Functional eine Basis hat.


% ===== scan leaf 554, printed 533 =====
% !PRINTED 533
\weberpage{533}

Zunächst ist klar, dass eine Basis eines Functionals $\mu$ zugleich Basis
aller mit $\mu$ associirten Functionale ist, ferner dass jedes Element
$\alpha_1, \alpha_2, \ldots$ einer solchen Basis durch $\mu$ theilbar sein
muss, endlich dass aus einer Basis von $\mu$ alle anderen abgeleitet werden
können, wenn man eine lineare Substitution mit ganzen rationalen
Co\"efficienten und der Determinante $\pm 1$ anwendet.

Um nun eine Basis von $\mu$ zu bilden, gehen wir von einer Minimalbasis
(Basis von $\mathfrak{o}$) aus, die wir, wie oben, mit
\[
\omega_1, \omega_2, \ldots, \omega_n
\]
bezeichnen.

Da es positive ganze rationale Zahlen giebt, die durch $\mu$ theilbar sind,
so giebt es auch ganze rationale Zahlen $a$, für die das Product
$a\,\omega_1$ durch $\mu$ theilbar ist. Die \gesperrt{kleinste} positive unter diesen
Zahlen wollen wir mit $a_{1,1}$ bezeichnen und
\[
\alpha_1 = a_{1,1}\,\omega_1
\]
setzen. Sodann bezeichnen wir mit $a_{2,2}$ die \gesperrt{kleinste positive} ganze
rationale Zahl, für die sich ein ganzes rationales $a_{1,2}$ so bestimmen
lässt, dass
\[
\alpha_2 = a_{1,2}\,\omega_1 + a_{2,2}\,\omega_2
\]
durch $\mu$ theilbar wird, und fahren so fort. Wir bilden auf diese Weise
das System
\weq{3}{
\begin{aligned}
\alpha_1 &= a_{1,1}\,\omega_1 \\
\alpha_2 &= a_{1,2}\,\omega_1 + a_{2,2}\,\omega_2 \\
&\cdots\cdots\cdots\cdots\cdots\cdots\cdots\cdots\cdots\cdots \\
\alpha_n &= a_{1,n}\,\omega_1 + a_{2,n}\,\omega_2 + \cdots + a_{n,n}\,\omega_n,
\end{aligned}
}
\gesperrt{worin $a_{1,1}, a_{2,2}, \ldots, a_{n,n}$ die kleinsten positiven
ganzen rationalen Zahlen sind, die eine Bestimmung der ganzen rationalen
Zahlen $a_{1,2}, \ldots, a_{n-1,n}$ so ermöglichen, dass die ganzen Zahlen
$\alpha_1, \alpha_2, \ldots, \alpha_n$ durch $\mu$ theilbar werden.}

Jede Zahl $a_{r,r}$ ist hiernach unabhängig von den übrigen
dadurch definirt, dass sie die kleinste natürliche Zahl ist, für die sich
die zugehörigen ganzen rationalen Zahlen $a_{1,r}, a_{2,r}, \ldots,
a_{r-1,r}$ so bestimmen lassen, dass
\weq{4}{
\alpha_r = a_{1,r}\,\omega_1 + a_{2,r}\,\omega_2 + \cdots +
a_{r-1,r}\,\omega_{r-1} + a_{r,r}\,\omega_r
}
durch $\mu$ theilbar wird. In dem besonderen Falle, wo
$\omega_r$ selbst durch $\mu$ theilbar ist, hat man demnach $a_{r,r} = 1$
zu setzen, und
%CONT>>


% ===== scan leaf 555, printed 534 =====
% !PRINTED 534
\weberpage{534}

%<<CONT
die $a_{1,r}, a_{2,r}, \ldots, a_{r-1,r}$ können alle gleich Null
angenommen werden.

Bezeichnen wir mit $\Delta$ die Grundzahl des Körpers $\Omega$, so
ergiebt sich die Discriminante des durch (3) bestimmten Systemes
$\alpha_1, \alpha_2, \ldots, \alpha_n$ nach der Formel \S.\,144, (11):
\weq{5}{
\Delta(\alpha_1, \alpha_2, \ldots, \alpha_n) = a_{1,1}^2\, a_{2,2}^2
\ldots a_{n,n}^2\, \Delta.
}
Dies ist eine von Null verschiedene ganze rationale Zahl, und folglich
sind die Grössen $\alpha$ eine Basis von $\Omega$.

Um also zu zeigen, dass das so bestimmte System $\alpha_1, \alpha_2,
\ldots, \alpha_n$ eine Basis von $\mu$ ist, bleibt noch nachzuweisen,
dass jede durch $\mu$ theilbare ganze Zahl $\alpha$ in der Form (2)
dargestellt werden kann. Nehmen wir, um diesen Beweis zu führen, irgend
einen Index $r \leqq n$ an, und suchen die Bedingung dafür, dass eine
ganze Zahl von der Form
\weq{6}{
\gamma_r = h_1\,\omega_1 + h_2\,\omega_2 + \cdots + h_r\,\omega_r
}
durch $\mu$ theilbar ist, wenn $h_1, h_2, \ldots, h_r$ ganze rationale
Zahlen sind.

Zunächst folgt, dass $h_r$ durch $a_{r,r}$ theilbar sein muss. Denn
bezeichnen wir mit $q_r$ den Rest der Division von $h_r$ durch $a_{r,r}$,
und setzen
\[
h_r = l_r\,a_{r,r} + q_r, \qquad 0 \leqq q_r < a_{r,r},
\]
so ergiebt sich nach (6)
\[
\gamma_r - l_r\,\alpha_r = (h_1 - l_r\,a_{1,r})\,\omega_1 +
(h_2 - l_r\,a_{2,r})\,\omega_2 + \cdots + q_r\,\omega_r,
\]
und diese Zahl müsste auch durch $\mu$ theilbar sein. Dies ist aber nach
der Definition von $a_{r,r}$ nur möglich, wenn $q_r = 0$ ist. Dann aber
erhält $\gamma_r - l_r\,\alpha_r$ den Ausdruck:
\weq{7}{
\gamma_r - l_r\,\alpha_r = h_1'\,\omega_1 + h_1'\,\omega_2 + \cdots +
h_{r-1}'\,\omega_{r-1}, % [sic]
}
wird also von derselben Form, wie (6), nur dass $r - 1$ an Stelle von $r$
tritt, und in (7) ist dieselbe Schlussweise zu wiederholen. Demnach
ergiebt sich durch vollständige Induction der Satz:

\gesperrt{Eine in der Form $h_1\,\omega_1 + h_2\,\omega_2 + \cdots +
h_r\,\omega_r$ darstellbare ganze Zahl des Körpers $\Omega$ ist immer
dann und nur dann durch $\mu$ theilbar, wenn sie in der Form}
\[
l_1\,\alpha_1 + l_2\,\alpha_2 + \cdots + l_r\,\alpha_r
\]
\gesperrt{darstellbar ist, in der die Co\"efficienten $l_1, l_2, \ldots,
l_r$ ganze rationale Zahlen sind.}

Setzt man in diesem Satze $r = n$, so hat man den Beweis dafür, dass das
Zahlensystem (3) eine Basis von $\mu$ ist.


% ===== scan leaf 556, printed 535 =====
% !PRINTED 535
\weberpage{535}

Wie man aus der einen Basis von $\mu$ alle anderen ableiten kann, haben
wir schon oben gesehen.

Wir verstehen jetzt unter $\alpha_1, \alpha_2, \ldots, \alpha_n$ eine
beliebige Basis von $\mu$ und stellen das Functional $\mu$ nach
\S.\,142, 4. als Quotienten zweier ganzer Functionen dar, dessen Nenner
eine rationale Einheit ist. Die Co\"efficienten des Zählers sind dann
ganze durch $\mu$ theilbare Zahlen (\S.\,143, 2.) und können daher in
der Form (2) dargestellt werden.

Fassen wir diese Darstellung gehörig zusammen, so ergiebt sich also für
$\mu$ ein Ausdruck
\weq{8}{
\mu = \alpha_1 u_1 + \alpha_2 u_2 + \cdots + \alpha_n u_n,
}
worin die $u_1, u_2, \ldots, u_n$ \gesperrt{ganze rationale Functionale}
sind.

Hiermit wollen wir die Linearform
\weq{9}{
\lambda = \alpha_1 t_1 + \alpha_2 t_2 + \cdots + \alpha_n t_n
}
vergleichen, in der $t_1, t_2, \ldots, t_n$ Variable sind. Diese
Linearform ist (nach \S.\,139) der grösste gemeinschaftliche Theiler der
Zahlen $\alpha_1, \alpha_2, \ldots, \alpha_n$ und ist durch $\mu$
theilbar, weil die Zahlen $\alpha_1, \alpha_2, \ldots, \alpha_n$ durch
$\mu$ theilbar sind. Andererseits ist aber auch, wie die Darstellung (8)
zeigt, $\mu$ durch $\lambda$ theilbar, und folglich sind die beiden
Functionale $\mu$ und $\lambda$ mit einander associirt.

\gesperrt{Das Functional $\lambda$ wollen wir eine Basisform des
Functionals $\mu$ nennen; es ist dann $\lambda$ zugleich Basisform von
allen mit $\mu$ associirten Functionalen.}

Eine Basisform $\lambda$ von $\mu$ hat die Eigenschaft, dass
man aus ihr alle durch $\mu$ (oder durch $\lambda$) theilbaren ganzen
Zahlen in $\Omega$ erhält, wenn man für die Variablen ganze rationale
Zahlen setzt.

Die Basis $\alpha_1, \alpha_2, \ldots, \alpha_n$ steht zu dem
Functionale $\mu$ in einer ähnlichen Beziehung, wie die Basis $\omega_1,
\omega_2, \ldots, \omega_n$ des Systemes $\mathfrak{o}$ aller ganzen Zahlen in
$\Omega$ zu den Einheiten. In der That ist die Linearform
\weq{10}{
\tau = \omega_1 t_1 + \omega_2 t_2 + \cdots + \omega_n t_n
}
eine Einheit; denn sie ist der grösste gemeinschaftliche Theiler von
$\omega_1, \omega_2, \ldots, \omega_n$ und muss also, wie die Formel
\S.\,145, (8) zeigt, ein Theiler von 1, also eine Einheit sein.

\gesperrt{Demnach wollen wir das Functional $\tau$ eine Basisform von $\mathfrak{o}$
nennen.}


% ===== scan leaf 557, printed 536 =====
% !PRINTED 536
\weberpage{536}
\pdfbookmark[2]{§ 147. Die absoluten Normen der Functionale}{bm168}

Aus einer solchen Basisform erhält man alle ganzen Zahlen des Körpers
$\Omega$, wenn man für die Variablen ganze rationale Zahlen setzt.

Die Grundform $\tau$ ist die Wurzel einer irreduciblen Gleichung
$n^{\text{ten}}$ Grades
\[
F(t) = N(t - \tau) = 0,
\]
in der die Co\"efficienten der Potenzen von $t$ ganze rationale (und
homogene) Functionen der Variablen $t_1, t_2, \ldots, t_n$ sind.

\begin{center}
\S. 147.

\gesperrt{Die absoluten Normen der Functionale.}
\end{center}

Die Elemente $\alpha_1, \alpha_2, \ldots, \alpha_n$ einer Basis des
Functionals $\mu$ können als ganze Zahlen in $\Omega$ linear und
ganzzahlig ausgedrückt werden durch eine Basis $\omega_1, \omega_2,
\ldots, \omega_n$ von $\rho$ in der Form
\weq{1}{
(\alpha_1, \alpha_2, \ldots, \alpha_n) = A(\omega_1, \omega_2, \ldots,
\omega_n),
}
worin $A$ eine lineare Substitution mit ganzen rationalen
Co\"efficienten bedeutet. Einen Specialfall hiervon bieten die Formeln
(3) des vorigen Paragraphen.

Eine Basisform $\lambda$ von $\mu$ wollen wir so darstellen:
\weq{2}{
\lambda = \sum^{\nu} \alpha_\nu t_\nu,
}
und die Substitution (1) schreiben wir ausführlicher:
\weq{3}{
\alpha_s = \sum^{\nu} a_{s,\nu}\,\omega_\nu,
}
worin $t_\nu$ Variable, $a_{s,\nu}$ die Substitutionsco\"efficienten
sind, und der Summationsbuchstabe $\nu$ von 1 bis $n$ läuft.

Da die Producte $\alpha_s\,\omega_r$ alle durch $\mu$ theilbar sind, so
können sie nach der Bedeutung der Basis in der Form dargestellt werden
\weq{4}{
\alpha_\nu\,\omega_r = \sum^{s} g^{(s)}_{\nu,r}\,\alpha_s,
}
worin die $g^{(s)}_{\nu,r}$ ganze rationale Zahlen sind. Hieraus erhält
man dann nach (2)
\weq{5}{
\lambda\,\omega_r = \sum^{s} \alpha_s\,t_{s,r},
}
wenn
\weq{6}{
t_{s,r} = \sum^{\nu} g^{(s)}_{\nu,r}\,t_\nu
}
ganze rationale Linearformen sind.

Substituirt man in (5) wieder die Ausdrücke (3), so folgt
\weq{7}{
\lambda\,\omega_r = \sum^{\nu} \omega_\nu \sum^{s} a_{s,\nu}\,t_{s,r}.
}


% ===== scan leaf 558, printed 537 =====
% !PRINTED 537
\weberpage{537}

Eliminiren wir aus diesen linearen Gleichungen die $\omega_\nu$, so
können wir die Gleichung $n^{\text{ten}}$ Grades für $\lambda$ in Determinanten\-
form darstellen, und das Product der Wurzeln dieser Gleichung,
also die Norm von $\lambda$, erhalten wir als die Determinante aus
den $n^2$ Grössen
\[
\sum^{s} a_{s,\nu}\,t_{s,r}
\]
(\S.\,137, 8.). Diese Determinante lässt sich aber nach dem Multi\-
plicationssatze der Determinanten (s. Bd. I, \S.\,27) zerlegen, und
giebt, wenn
\[
A = \Sigma \pm a_{1,1}a_{2,2}\ldots a_{n,n}, \quad T = \Sigma \pm t_{1,1}t_{2,2}\ldots t_{n,n}
\]
gesetzt wird,
\weq{8}{N(\lambda) = A\,T.}

Hierin ist $T$ eine ganze Function der Variablen $t_\nu$ mit ganzen
rationalen Zahlenco\"efficienten, von der wir nun noch nachweisen
werden, dass sie primitiv ist.

Nehmen wir an, im Theiler von $T$ gehe irgend eine natür\-
liche Primzahl $p$ auf. Dann können wir $n$ ganze rationale
Formen $y_1, y_2, \ldots, y_n$ der Variablen $t$ so bestimmen, dass die
$n$ Summen
\weq{9}{u_s = t_{s,1}y_1 + t_{s,2}y_2 + \cdots + t_{s,n}y_n}
alle durch $p$ theilbar sind, ohne dass alle $y_1, y_2, \ldots, y_n$ durch $p$
theilbar sind.

Die Richtigkeit dieser Behauptung folgt aus elementaren
Determinantensätzen (Bd. I, zweiter Abschnitt).

Wenn nämlich die $t_{s,r}$ alle durch $p$ theilbar sind, so können
wir die $y_r$ ganz beliebig z.\,B. gleich 1 annehmen.

Andernfalls nehmen wir an, dass ausser der Determinante $T$
auch alle $m$-reihigen Unterdeterminanten durch $p$ theilbar seien
$(m \leqq n)$,
und dass unter den $(m-1)$-reihigen Unterdeter\-
minanten wenigstens eine nicht durch $p$ theilbar sei. Ist dann
etwa unter den $(m-1)$-reihigen Determinanten der Matrix
\[
\begin{array}{cccc}
t_{1,1}, & t_{1,2}, & \ldots, & t_{1,m} \\
\multicolumn{4}{c}{\cdots\cdots\cdots\cdots\cdots\cdots\cdots\cdots\cdots\cdots} \\
t_{m-1,1}, & t_{m-1,2}, & \ldots, & t_{m-1,m}
\end{array}
\]
eine durch $p$ nicht theilbar, so setzen wir für $y_1, y_2, \ldots, y_m$ eben
diese $(m-1)$-reihigen Determinanten und nehmen $y_{m+1}, \ldots, y_n$
$= 0$ an. Diese $y$ genügen dann der gestellten Forderung.


% ===== scan leaf 559, printed 538 =====
% !PRINTED 538
\weberpage{538}

Sind die $y$ dann so bestimmt, so setzen wir
\weq{10}{\omega = \omega_1 y_1 + \omega_2 y_2 + \cdots + \omega_n y_n,}
und leiten aus (5) die Gleichung ab
\weq{11}{\lambda\,\omega = \sum^{s} \alpha_s u_s.}

Da nun die $\alpha_s$ durch $\lambda$ und die $u_s$ durch $p$ theilbar sind,
so folgt, dass $\omega$ durch $p$ theilbar ist, und dass mithin, nach
\S.\,145, 2., $y_1, y_2, \ldots, y_n$ durch $p$ theilbar sein müssen, was
unserer Voraussetzung entgegen ist.

Damit ist bewiesen, dass $T$ eine primitive Function ist, \gesperrt{und
dass also der absolute Werth der Determinante $A$
gleich der absoluten Norm von $\lambda$ und folglich auch
von $\mu$ ist} (\S.\,138).

Wenden wir das Ergebniss auf die specielle Basis (3), \S.\,146
an, so ergiebt sich die Formel:
\weq{12}{N_a(\mu) = a_{1,1}\,a_{2,2}\ldots a_{n,n}.}

Hieran knüpfen wir noch folgende Bemerkungen: Wenn
man in einer Basisform eines Functionales $\mu$
\[
\lambda = \alpha_1 t_1 + \alpha_2 t_2 + \cdots + \alpha_n t_n
\]
auf die Variablen $t_1, t_2, \ldots, t_n$ eine ganzzahlige lineare Sub\-
stitution mit der Determinante $\pm 1$ anwendet, so entsteht eine
neue Basisform von $\mu$. Denn allen ganzzahligen rationalen Werthen
der Variablen $t_1, t_2, \ldots, t_n$ entsprechen ganzzahlige rationale
Werthe der neuen Variablen und umgekehrt.

Die Anwendung einer linearen Substitution auf die Variablen $t$
ist aber gleichbedeutend mit der Anwendung der transponirten
Substitution auf die Co\"efficienten $\alpha_1, \alpha_2, \ldots, \alpha_n$ (\S.\,37), d.\,h. mit
dem Uebergange zu einer neuen Basis von $\mu$:
\weq{13}{(\beta_1, \beta_2, \ldots, \beta_n) = (B)\,(\alpha_1, \alpha_2, \ldots, \alpha_n),}
die dann durch Zusammensetzung mit (1) ergiebt:
\weq{14}{(\beta_1, \beta_2, \ldots, \beta_n) = (B)\,(A)\,(\omega_1, \omega_2, \ldots, \omega_n).}

Wenn die Discriminante (\S.\,144) der durch eine Substitu\-
tion (13) bestimmten Zahlen $\beta$ mit der Discriminante der $\alpha$
übereinstimmt, so muss die Substitutionsdeterminante $B = \pm 1$
sein; und wir kommen also zu den Sätzen:
\begin{enumerate}
\item[1.] \gesperrt{Die Discriminante einer Basis von $\mu$ ist gleich
dem Quadrate der absoluten Norm von $\mu$, multi\-
plicirt mit der Grundzahl des Körpers;}
\end{enumerate}
%CONT>>


% ===== scan leaf 560, printed 539 =====
% !PRINTED 539
\weberpage{539}
\pdfbookmark[2]{§ 148. Volles Restsystem nach einem Modul}{bm169}
%<<CONT
\begin{enumerate}
\item[] und umgekehrt:
\item[2.] \gesperrt{Ist $\beta_1, \beta_2, \ldots, \beta_n$ ein System ganzer durch $\mu$
theilbarer Zahlen, dessen Discriminante gleich
ist dem Quadrate der absoluten Norm von $\mu$,
multiplicirt mit der Grundzahl des Körpers, so
ist $\beta_1, \beta_2, \ldots, \beta_n$ eine Basis von $\mu$.}
\end{enumerate}

\S.\,148.

\gesperrt{Volles Restsystem nach einem Modul.}

\begin{enumerate}
\item[1.] \textbf{Definition:} \gesperrt{Zwei ganze algebraische Zahlen
$\xi, \eta$, deren Differenz $\xi - \eta$ durch ein ganzes
von Null verschiedenes Functional $\mu$ theilbar
ist, heissen mit einander congruent nach dem
Modul $\mu$.}
\end{enumerate}

Der Begriff der Congruenz lässt sich auch auf Functionale
ausdehnen, was wir aber fürs Erste noch nicht thun. Dagegen
ist es wesentlich, als Moduln der Congruenz nicht bloss die
Zahlen, sondern auch die Functionale zu berücksichtigen. Jede
Congruenz bleibt bestehen, wenn der Modul durch ein associirtes
Functional ersetzt wird. Beim Modul können beliebige Einheits\-
factoren hinzugefügt oder weggelassen werden.

Wir gebrauchen für die Congruenz das \gesperrt{Gauss}'sche Zeichen
(Bd. I, \S.\,115)
\weq{1}{\xi \equiv \eta \ (\mathrm{mod}\ \mu).}

Eine solche Congruenz ist gleichbedeutend mit der Gleichung
\weq{2}{\xi = \eta + \omega\mu,}
worin $\omega$ irgend ein ganzes Functional sein kann.

Aus dieser Darstellung erkennt man dann sofort, dass, ebenso
wie in Congruenzen zwischen rationalen Zahlen, wenn man in
einem durch \gesperrt{Addition}, \gesperrt{Subtraction} und \gesperrt{Multiplication} von
ganzen Zahlen zusammengesetzten Ausdruck jede Zahl durch
eine nach dem Modul $\mu$ congruente Zahl ersetzt, eine nach dem\-
selben Modul congruente Zahl das Resultat ist; in Zeichen:

Sind $\xi_1, \eta_1, \xi_2, \eta_2, \ldots$ ganze Zahlen, die den Congruenzen
\[
\xi_1 \equiv \eta_1, \ \xi_2 \equiv \eta_2, \ldots \ (\mathrm{mod}\ \mu)
\]
%CONT>>


% ===== scan leaf 561, printed 540 =====
% !PRINTED 540
\weberpage{540}
%<<CONT
genügen, und ist $\psi(x_1, x_2, \ldots)$ eine ganze Function mit ganz\-
zahligen rationalen Co\"efficienten, so ist auch
\weq{3}{\psi(\xi_1, \xi_2, \ldots) \equiv \psi(\eta_1, \eta_2, \ldots) \ (\mathrm{mod}\ \mu).}

\begin{enumerate}
\item[2.] \gesperrt{Ist $m$ die kleinste positive ganze rationale Zahl,
die durch $\mu$ theilbar ist, so sind zwei nach dem
Modul $\mu$ congruente ganze rationale Zahlen auch
nach dem Modul $m$ congruent} (\S.\,138, 13.).
\end{enumerate}

Ein Hauptsatz über die Congruenzen, zu dessen Beweis wir
jetzt schreiten, ist der:

\begin{enumerate}
\item[3.] \gesperrt{Die Anzahl der nach einem Modul $\mu$ incon\-
gruenten ganzen Zahlen eines Körpers $\Omega$ ist
endlich, und zwar gleich der absoluten Norm
von $\mu$.}
\end{enumerate}

Um ihn zu beweisen, nehmen wir eine Basis des Func\-
tionals $\mu$ an, und zwar wählen wir gerade die im \S.\,146, (3)
bestimmte specielle Basis:
\weq{4}{
\begin{array}{l}
\alpha_1 = a_{1,1}\,\omega_1 \\
\alpha_2 = a_{1,2}\,\omega_1 + a_{2,2}\,\omega_2 \\
\cdots\cdots\cdots\cdots\cdots\cdots\cdots\cdots\cdots\cdots \\
\alpha_n = a_{1,n}\,\omega_1 + a_{2,n}\,\omega_2 + \cdots + a_{n,n}\,\omega_n,
\end{array}
}
worin $\omega_1, \omega_2, \ldots, \omega_n$ eine Basis von $\mathfrak{o}$ ist.

Jede ganze Zahl $\eta$ in $\Omega$ lässt sich dann in der Form dar\-
stellen:
\weq{5}{\eta = y_1\omega_1 + y_2\omega_2 + \cdots + y_n\omega_n,}
wenn die $y_1, y_2, \ldots, y_n$ ganze rationale Zahlen sind. Damit ver\-
binden wir das Zahlensystem
\weq{6}{\xi = x_1\omega_1 + x_2\omega_2 + \cdots + x_n\omega_n,}
in dem wir $x_1$ ein volles Restsystem nach dem Modul $a_{1,1}$,
$x_2$ nach dem Modul $a_{2,2} \ldots$, $x_n$ nach dem Modul $a_{n,n}$ durch\-
laufen lassen. Die Anzahl der so gebildeten Zahlen $\xi$ beträgt
$a_{1,1}a_{2,2}\ldots a_{n,n}$, eine Zahl, die nach \S.\,147, (12) gleich der
absoluten Norm von $\mu$ ist.

Wir wollen nachweisen, dass jede Zahl $\eta$ einer und nur
einer dieser Zahlen $\xi$ nach dem Modul $\mu$ congruent ist.

Wenn $\eta - \xi$ durch $\mu$ theilbar sein soll, so muss nach der
Bedeutung der Basis diese Differenz in der Form
\weq{7}{\eta - \xi = h_1\alpha_1 + h_2\alpha_2 + \cdots + h_n\alpha_n}
%CONT>>


% ===== scan leaf 562, printed 541 =====
% !PRINTED 541
\weberpage{541}
\pdfbookmark[2]{§ 149. Congruenzen}{bm170}
%<<CONT
darstellbar sein, worin die $h_1, h_2, \ldots, h_n$ ganze rationale Zahlen
sind, und wenn diese Bedingung erfüllt ist, so ist auch umgekehrt $\eta$
mit $\xi$ nach dem Modul $\mu$ congruent. Ordnen wir die rechte Seite
mittelst (4) nach $\omega_1, \omega_2, \ldots, \omega_n$, so ergiebt sich
aus (5) und (6):
\weq{8}{\begin{array}{rclllll}
y_1 &=& x_1 &+ h_1a_{1,1} &+ h_2a_{1,2} &+ \cdots &+ h_na_{1,n} \\
y_2 &=& x_2 & &+ h_2a_{2,2} &+ \cdots &+ h_na_{2,n} \\
\multicolumn{7}{c}{.\quad.\quad.\quad.\quad.\quad.\quad.\quad.\quad.\quad.\quad.\quad.\quad.\quad.\quad.} \\
y_n &=& x_n & & & &+ h_na_{n,n}.
\end{array}}

Darin aber liegt der Beweis unserer Behauptung. Denn sind $y_1, y_2,
\ldots, y_n$ irgend welche gegebene ganze rationale Zahlen, so ist
zunächst aus
\[
y_n = x_n + h_na_{n,n}
\]
$x_n$ und $h_n$ vollständig bestimmt. Darauf erhalten wir aus der
vorletzten Gleichung
\[
y_{n-1} - h_na_{n-1,n} = x_{n-1} + h_{n-1}a_{n-1,n-1},
\]
wodurch $x_{n-1}$ und $h_{n-1}$ völlig bestimmt sind. Die nächst
vorhergehende Gleichung ergiebt
\[
y_{n-2} - h_na_{n-2,n} - h_{n-1}a_{n-2,n-1} = x_{n-2} + h_{n-2}a_{n-2,n-2},
\]
wodurch $x_{n-2}$ und $h_{n-2}$ völlig bestimmt sind, und so fahren wir
fort, bis alle $x_i$ und $h_i$ bestimmt sind.

Hiernach können wir zweckmässig die Gesammtheit der Zahlen $\xi$ oder
irgend ein System mit diesen congruenter Zahlen ein \gesperrt{volles
Restsystem nach dem Modul $\mu$ nennen.}

Die Anzahl der Individuen eines vollen Restsystems für den Modul $\mu$
ist gleich der absoluten Norm von $\mu$.

\begin{center}
\S.\,149.

\gesperrt{Congruenzen.}
\end{center}

Aus der Existenz eines vollen Restsystems nach einem Modul $\mu$, die
wir im vorigen Paragraphen nachgewiesen haben, ergiebt sich eine Reihe
von Sätzen, die mit bekannten elementaren Sätzen der rationalen
Zahlentheorie analog sind.

Bedeutet $\alpha$ irgend eine ganze Zahl in $\Omega$ und $\mu$ ein als
Modul dienendes ganzes Functional, und ist $\alpha$ relativ prim zu
$\mu$, so sind zwei Zahlen $\alpha\xi$ und $\alpha\xi'$ nur dann
congruent nach dem
%CONT>>


% ===== scan leaf 563, printed 542 =====
% !PRINTED 542
\weberpage{542}
%<<CONT
Modul $\mu$, wenn die ganzen Zahlen $\xi, \xi'$ congruent sind. Hieraus
ergiebt sich, dass, wenn $\xi$ ein volles Restsystem nach dem Modul
$\mu$ durchläuft, dasselbe auch von dem Producte $\alpha\xi$ gilt,
wodurch der Satz bewiesen ist:

1. \gesperrt{Ist $\alpha$ eine ganze Zahl in $\Omega$, relativ prim zu
dem Modul $\mu$, ferner $\gamma$ eine beliebige ganze Zahl in $\Omega$,
so ist die Congruenz:}
\weq{1}{\alpha\xi \equiv \gamma \ (\mathrm{mod}\ \mu)}
\gesperrt{immer durch eine ganze Zahl $\xi$ lösbar, und auch nur durch
eine, wenn für $\xi$ ein volles Restsystem nach dem Modul $\mu$
vorgeschrieben ist.}

Wenn hierin $\mu$ selbst eine Zahl ist, die wir mit $\beta$ bezeichnen,
so ist auch der Quotient
\[
\frac{\alpha\xi - \gamma}{\beta} = -\eta
\]
eine ganze Zahl, und dann nimmt der vorstehende Satz die Form an:

2. \gesperrt{Sind $\alpha, \beta, \gamma$ drei ganze Zahlen in $\Omega$
und $\alpha, \beta$ relativ prim, so kann man zwei andere ganze Zahlen
$\xi, \eta$ in $\Omega$ so bestimmen, dass}
\weq{2}{\alpha\xi + \beta\eta = \gamma}
\gesperrt{wird. Insbesondere kann man also auch für zwei beliebige
relative Primzahlen $\alpha, \beta$ die Gleichung}
\weq{3}{\alpha\xi + \beta\eta = 1}
\gesperrt{durch ganze Zahlen $\xi, \eta$ befriedigen.}

Dieser Satz lässt sich auf ein System von mehreren Zahlen übertragen.

Wenn die Zahlen $\alpha, \beta, \gamma, \ldots$ in $\mathfrak{o}$ keinen
gemeinsamen Theiler haben, so können wir zunächst Zahlen $\eta_1, \zeta_1,
\ldots$ in $\mathfrak{o}$ so bestimmen, dass
\[
\beta_1 = \beta\eta_1 + \gamma\zeta_1 + \cdots
\]
relativ prim zu $\alpha$ wird. Wir haben nämlich, wenn $\pi_1, \pi_2,
\ldots$ die verschiedenen Primfactoren von $\alpha$ sind, deren keiner
in allen $\beta, \gamma, \ldots$ aufgehen kann, und wenn etwa $\beta$
durch $\pi_1$ nicht theilbar ist, $\eta_1$ durch $\pi_1$ untheilbar, die
übrigen Zahlen $\zeta_1, \ldots$ durch $\pi_1$ theilbar u.\,s.\,f.
anzunehmen, und erhalten für jede der Zahlen $\eta_1, \zeta_1, \ldots$ die
Bedingung, dass sie durch einige der Primfactoren
%CONT>>


% ===== scan leaf 564, printed 543 =====
% !PRINTED 543
\weberpage{543}
%<<CONT
$\pi_i$ theilbar, durch andere nicht theilbar sein soll, und dieser
Forderung kann nach \S.\,143, 3. immer genügt werden. Dann können wir
nach 2. die ganze Zahl $\tau$ so bestimmen, dass
\[
\alpha\xi + \beta_1\tau = 1
\]
wird, und wenn wir $\tau\eta_1 = \eta$, $\tau\zeta_1 = \zeta, \ldots$
setzen, so erhalten wir den Satz:

3. \gesperrt{Sind $\alpha, \beta, \gamma, \ldots$ Zahlen in
$\mathfrak{o}$ ohne gemeinschaftlichen Theiler, so lassen sich andere
Zahlen $\xi, \eta, \zeta, \ldots$ in $\mathfrak{o}$ so bestimmen, dass}
\weq{4}{\alpha\xi + \beta\eta + \gamma\zeta + \cdots = 1}
\gesperrt{wird.}

Daraus lässt sich weiter auf folgenden Satz schliessen:

4. \gesperrt{Sind $\alpha, \beta, \gamma, \ldots$ beliebige Zahlen in
$\mathfrak{o}$, $\mu$ eine durch den grössten gemeinschaftlichen Theiler
aller dieser Zahlen theilbare Zahl in $\mathfrak{o}$, so kann man die
Zahlen $\xi, \eta, \zeta, \ldots$ in $\mathfrak{o}$ so bestimmen, dass}
\weq{5}{\mu = \alpha\xi + \beta\eta + \gamma\zeta + \cdots}
\gesperrt{wird.}

Wenn wir nämlich den grössten gemeinschaftlichen Theiler der Zahlen
$\alpha, \beta, \gamma, \ldots$ nach \S.\,139 in der Form
\[
\delta = \alpha u + \beta v + \gamma w \ldots
\]
annehmen, worin $u, v, w, \ldots$ Variable sind, so giebt es nach
Voraussetzung ein ganzes Functional $\omega$, so dass $\mu = \delta
\omega$ ist. Das Functional $\omega$ stellen wir als Quotienten zweier
ganzer Functionen $\varphi : \varepsilon$ dar, von denen $\varepsilon$
eine Einheit, und folglich $\varphi$ eine Function mit ganzzahligen
Co\"efficienten ist, und erhalten so
\weq{6}{\varepsilon\mu = (\alpha u + \beta v + \gamma w + \cdots)\varphi.}
Ordnet man beide Seiten dieser Gleichung nach den darin vorkommenden
Variablen, so erhält man, wenn $\varepsilon_1, \varepsilon_2, \ldots$
die Co\"efficienten in $\varepsilon$ sind, ein System von Gleichungen
von folgender Form:
\weq{7}{\begin{aligned}
\varepsilon_1\mu &= \alpha\xi_1 + \beta\eta_1 + \gamma\zeta_1 + \cdots \\
\varepsilon_2\mu &= \alpha\xi_2 + \beta\eta_2 + \gamma\zeta_2 + \cdots \\
& \text{.\,.\,.\,.\,.\,.\,.\,.\,.\,.\,.\,.\,.\,.\,.\,.\,.\,.\,.\,.\,.\,.\,.,}
\end{aligned}}
worin die $\xi_i, \eta_i, \zeta_i, \ldots$ Zahlen in $\mathfrak{o}$ sind.
Nun haben aber die Zahlen $\varepsilon_1, \varepsilon_2, \ldots$ als
Co\"efficienten einer Einheit, keinen gemein
%CONT>>


% ===== scan leaf 565, printed 544 =====
% !PRINTED 544
\weberpage{544}
\pdfbookmark[2]{§ 150. Der Fermat'sche Satz}{bm171}
%<<CONT
samen Theiler, und folglich kann man nach 3. die Zahlen $\tau_1,
\tau_2, \ldots$ in $\mathfrak{o}$ so bestimmen, dass
\weq{8}{\varepsilon_1\tau_1 + \varepsilon_2\tau_2 + \cdots = 1}
wird. Aus (7) folgt aber, wenn man mit $\tau_1, \tau_2, \ldots$
multiplicirt und addirt, und dann
\[
\xi = \tau_1\xi_1 + \tau_2\xi_2 + \cdots, \quad
\eta = \tau_1\eta_1 + \tau_2\eta_2 + \cdots
\]
setzt, mittelst (8) die zu beweisende Gleichung (5).

\begin{center}
\S.\,150.

\gesperrt{Der Fermat'sche Satz.}
\end{center}

Aus der Theorie der Congruenzen lassen sich Folgerungen ziehen, die den
aus dem \gesperrt{Fermat'schen Lehrsatze} abgeleiteten Sätzen der
rationalen Zahlentheorie genau entsprechen, von denen hier die
wichtigsten, späterer Anwendung wegen, besprochen werden müssen.

Wir wollen unter $\pi$ ein Primfunctional $f^{\mathrm{ten}}$ Grades
verstehen, also, wenn $p$ die durch $\pi$ theilbare natürliche
Primzahl ist,
\weq{1}{N_a(\pi) = p^f}
setzen (\S.\,140). Ist dann $\alpha$ irgend eine durch $\pi$ nicht
theilbare Zahl in $\mathfrak{o}$, so wird, wie wir schon im vorigen
Paragraphen gesehen haben, das Product $\alpha\xi$ zugleich mit der
Zahl $\xi$ ein volles Restsystem nach dem Modul $\pi$ durchlaufen.
Lassen wir die durch $\pi$ theilbare Zahl weg, so bleiben $N_a(\pi) - 1$
Zahlen übrig, und wenn wir das Product bilden, so folgt
\weq{2}{\alpha^{p^f-1}\Pi(\xi) \equiv \Pi(\xi) \ (\mathrm{mod}\ \pi),}
wenn $\Pi(\xi)$ das Product aller Zahlen eines vollen Restsystems (mit
Ausschluss der Null) bedeutet, und daher durch $\pi$ nicht theilbar
ist. Demnach folgt aus (2)
\weq{3}{\alpha^{p^f-1} \equiv 1 \ (\mathrm{mod}\ \pi),}
oder, wenn man mit $\alpha$ multiplicirt,
\weq{4}{\alpha^{p^f} \equiv \alpha \ (\mathrm{mod}\ \pi),}
und in der letzten Form gilt der Satz auch noch, wenn $\alpha$ durch
$\pi$ theilbar ist.

Die Formeln (3), (4), die genau dem \gesperrt{Fermat'schen Lehrsatze}
entsprechen (Bd.\,I, \S.\,136), sollen auch hier als \gesperrt{Fermat'scher
Lehrsatz} bezeichnet werden.


% ===== scan leaf 566, printed 545 =====
% !PRINTED 545
\weberpage{545}
Wir sprechen ihn so aus:

1. \gesperrt{Ist $\pi$ ein Primtheiler der natürlichen Primzahl $p$, so ist für jede ganze Zahl $\omega$ im Körper $\Omega$}
\[ \omega^{N_a(\pi)} \equiv \omega \pmod{\pi}. \]

Hieran knüpfen sich nun wichtige Folgerungen:

2. \gesperrt{Bezeichnet $f(t)$ eine ganze Function $m^{\text{ten}}$ Grades, deren Coëfficienten ganze Zahlen in $\Omega$ sind, und $\pi$ ein Primfunctional, so hat die Congruenz}
\weq{5}{f(t) \equiv 0 \pmod{\pi}}
\gesperrt{höchstens $m$ Wurzeln, d. h. es giebt höchstens $m$ incongruente ganze Zahlen in $\Omega$, die, für $t$ gesetzt, die Congruenz befriedigen.}

Bedeutet nämlich $\alpha$ irgend eine Zahl in $\mathfrak{o}$, so können wir
\weq{6}{f(t) = (t-\alpha)\,f_1(t) + f(\alpha)}
setzen, worin $f_1(t)$ eine ebensolche Function wie $f(t)$ ist, aber nur vom $(m-1)^{\text{ten}}$ Grade. Ist aber $f(\alpha) \equiv 0 \pmod{\pi}$, so muss jede Wurzel von (5) der Congruenz
\[ (t-\alpha)\,f_1(t) \equiv 0 \pmod{\pi} \]
genügen. Sie muss also entweder mit $\alpha$ congruent oder eine Wurzel der Congruenz $(m-1)^{\text{ten}}$ Grades
\[ f_1(t) \equiv 0 \pmod{\pi} \]
sein. Setzen wir unseren Satz als bewiesen voraus für Congruenzen $(m-1)^{\text{ten}}$ Grades, so gilt er demnach auch für Congruenzen $m^{\text{ten}}$ Grades; und da er für Congruenzen ersten Grades gilt, so ist er allgemein richtig.

Jede durch $\pi$ nicht theilbare Zahl in $\mathfrak{o}$ genügt, wie wir gesehen haben, der Congruenz
\weq{7}{\omega^{p^f-1} \equiv 1 \pmod{\pi}.}
Ist nun $a$ die kleinste natürliche Zahl, für die die Congruenz
\weq{8}{\omega^a \equiv 1 \pmod{\pi}}
befriedigt ist, so lässt sich durch das schon oft angewandte Schlussverfahren zeigen, dass jeder andere Exponent $l$, für den $\omega^l \equiv 1$ ist, ein Vielfaches von $a$ sein muss. Denn wäre $l$ nicht durch $a$ theilbar, so wäre auch, wenn $a'$ der Rest der Division von $a$ durch $l$ ist, $\omega^{a'} \equiv 1$, was nach der Voraussetzung über $a$ nicht möglich ist. % [sic] the print has ``von $a$ durch $l$''; the minimality argument requires the remainder of $l$ by $a$
Also ist $a$ ein Theiler von $p^f-1$, und wir nennen $\omega$ eine \gesperrt{zum Exponenten $a$ gehörige Zahl}.


% ===== scan leaf 567, printed 546 =====
% !PRINTED 546
\weberpage{546}
Gehört $\omega$ zum Exponenten $a$, so sind die Potenzen
\weq{9}{1, \omega, \omega^2, \ldots, \omega^{a-1}}
alle incongruent und bilden also, da sie alle der Congruenz (8) genügen (nach 2.), die Gesammtheit der Wurzeln dieser Congruenz. Unter den Zahlen (9) müssen daher alle anderen zum Exponenten $a$ gehörigen Zahlen $\omega$ gesucht werden. Es wird aber $\omega^l$ nur dann zum Exponenten $a$ gehören, wenn $l$ relativ prim zu $a$ ist, und es folgt:

Wenn es überhaupt Zahlen $\omega$ giebt, die zum Exponenten $a$ gehören, so ist ihre Anzahl so gross, wie die Anzahl der relativen Primzahlen zu $a$ in der Reihe der Zahlen $0, 1, 2, \ldots, a-1$. Diese Zahl bezeichnen wir, wie schon früher (Bd. I, \S.\,132), mit $\varphi(a)$. Dass aber zu jedem Theiler $a$ von $p^f-1$ immer wenigstens eine Zahl $\omega$ und folglich $\varphi(a)$ Zahlen gehören, kann man ganz so beweisen, wie der entsprechende Satz der rationalen Zahlentheorie im \S.\,136 des ersten Bandes bewiesen ist.

Die Zahlen $\omega$, die zu dem Exponenten $p^f-1$ gehören, deren es hiernach immer $\varphi(p^f-1)$ giebt, heissen \gesperrt{primitive Wurzeln} von $\pi$. Ist $\gamma$ eine solche primitive Wurzel, so bilden die Potenzen
\[ 1, \gamma, \gamma^2, \ldots, \gamma^{p^f-2} \]
ein volles Restsystem nach dem Modul $\pi$ mit Ausschluss der durch $\pi$ theilbaren Zahl.

Nach dem \gesperrt{Fermat}'schen Lehrsatze für rationale Zahlen ist $t^{p-1}-1$ für $t = 1, 2, \ldots, p-1$ durch $p$, und folglich auch durch $\pi$ theilbar. Die Congruenz
\weq{10}{t^{p-1} \equiv 1 \pmod{\pi}}
hat daher die Wurzeln
\[ 1, 2, \ldots, p-1, \]
und diese sind, da nach \S.\,140,3. eine rationale Zahl nur dann durch $\pi$ theilbar ist, wenn sie durch $p$ theilbar ist, unter einander incongruent. Nach dem Satze 2. hat also die Congruenz (10) keine anderen Wurzeln als diese.

Multipliciren wir die Congruenz (10) noch mit $t$, so folgt, dass die Congruenz $p^{\text{ten}}$ Grades
\[ t^p - t \equiv 0 \pmod{\pi} \]
die $p$ Wurzeln
\[ 0, 1, 2, \ldots, p-1, \]
%CONT>>


% ===== scan leaf 568, printed 547 =====
% !PRINTED 547
\weberpage{547}
\pdfbookmark[2]{§ 151. Die Dedekind'schen Ideale}{bm172}
%<<CONT
und keine anderen hat. Darin liegt der Beweis des folgenden Satzes:

3. \gesperrt{Eine Zahl $\omega$ in $\mathfrak{o}$ ist dann und nur dann nach dem Modul $\pi$ mit einer rationalen Zahl congruent, wenn sie der Bedingung}
\[ \omega^p \equiv \omega \pmod{\pi} \]
\gesperrt{genügt.}

Beachtet man noch, dass die Polynomialcoëfficienten in der $p^{\text{ten}}$ Potenz eines Polynoms alle durch $p$ theilbar sind, mit Ausnahme derer, die zu den $p^{\text{ten}}$ Potenzen der einzelnen Glieder des Polynoms gehören, so ergiebt sich noch folgender Satz, der sich auf ganze Functionen in $\Omega$ von beliebigen Veränderlichen $x, y, \ldots$ bezieht:

4. \gesperrt{Ist $\psi(x,y,\ldots)$ eine ganze Function der Variablen $x, y, \ldots$ mit ganzzahligen Coëfficienten aus $\Omega$, so ist das Bestehen der Congruenz}
\[ [\psi(x,y,\ldots)]^p \equiv \psi(x^p, y^p, \ldots) \pmod{\pi} \]
\gesperrt{die nothwendige und hinreichende Bedingung dafür, dass alle Coëfficienten von $\psi$ mit ganzen rationalen Zahlen nach dem Modul $\pi$ congruent sind.}

\begin{center}
\S.\,151.\\
\gesperrt{Die Dedekind'schen Ideale.}
\end{center}

\gesperrt{Dedekind} gründet in den schon oben erwähnten Arbeiten die Theorie der algebraischen Zahlen auf den Begriff des \gesperrt{Ideals}.

Wir wollen jetzt nachweisen, dass die Theorie der Ideale im Wesen übereinstimmt mit der Theorie der Functionale, indem wir zeigen, wie der Uebergang von der einen zur anderen bewirkt werden kann.

Das System aller ganzen Zahlen eines algebraischen Zahlkörpers $\Omega$ soll, wie oben, mit $\mathfrak{o}$ bezeichnet werden\footnote{Vgl. \gesperrt{Dirichlet-Dedekind}, Vorlesungen über Zahlentheorie im \S.\,167 der dritten, \S.\,177 der vierten Auflage.}. Ein in $\mathfrak{o}$ enthaltenes Zahlensystem $\mathfrak{a}$ wird ein \gesperrt{Ideal} genannt, wenn es den beiden Forderungen genügt:
%CONT>>


% ===== scan leaf 569, printed 548 =====
% !PRINTED 548
\weberpage{548}
%<<CONT
I. \gesperrt{Summe und Differenz irgend zweier Zahlen in $\mathfrak{a}$ geben immer wieder Zahlen in $\mathfrak{a}$.}

II. \gesperrt{Das Product irgend einer Zahl in $\mathfrak{a}$ und einer Zahl in $\mathfrak{o}$ gehört dem System $\mathfrak{a}$ an.}

Dieser Forderung würde das aus der einzigen Zahl Null bestehende System genügen, was aber der Einfachheit halber nicht als ein Ideal bezeichnet wird.

Das System $\mathfrak{o}$ dagegen ist ein eigentliches Ideal. Ebenso ist das System aller durch eine bestimmte Zahl $\mu$ in $\mathfrak{o}$ theilbaren Zahlen $\mathfrak{o}\mu$ ein Ideal, und ein solches wird ein \gesperrt{Hauptideal} genannt.

Unter dem Producte $\mathfrak{a}\mathfrak{b}$ zweier Ideale $\mathfrak{a}$ und $\mathfrak{b}$ versteht man den Inbegriff aller Zahlen, die man erhält, wenn man irgend eine Zahl $\alpha$ aus $\mathfrak{a}$ mit einer Zahl $\beta$ aus $\mathfrak{b}$ multiplicirt und eine beliebige Anzahl solcher Zahlenproducte addirt, also den Inbegriff aller Zahlen von der Form $\Sigma \alpha\beta$. Dass dieses Product $\mathfrak{a}\mathfrak{b}$ wieder ein Ideal ist, leuchtet unmittelbar ein. Nach dieser Definition ist z. B. $\mathfrak{o}\mathfrak{a} = \mathfrak{a}$, und das Ideal $\mathfrak{o}$ spielt bei dieser Multiplication die Rolle der Einheit.

Man kann nun die Ideale und Functionale in der Weise auf einander beziehen, dass dabei folgende Gesetze obwalten:

1. \gesperrt{Jedem ganzen Functional entspricht ein bestimmtes Ideal, und associirten Functionalen entspricht dasselbe Ideal.}

2. \gesperrt{Jedem Ideal entsprechen unendlich viele, aber nur associirte ganze Functionale.}

3. \gesperrt{Dem Producte zweier oder mehrerer ganzer Functionale entsprechen die Producte der den Factoren entsprechenden Ideale.}

4. \gesperrt{Einer ganzen Zahl entspricht ein Hauptideal.}

5. \gesperrt{Den Einheiten entspricht das Ideal $\mathfrak{o}$.}

Um dieses Entsprechen zu definiren, ordnen wir zunächst dem Systeme aller Einheiten das Ideal $\mathfrak{o}$ zu. Ist dann ferner $\varphi$ irgend ein ganzes Functional, was keine Einheit ist, so genügt der Inbegriff aller durch $\varphi$ theilbaren ganzen Zahlen $\alpha$ des Körpers $\Omega$ nach \S.\,138 den Forderungen I, II, und ist also ein Ideal, das wir mit $\mathfrak{a}$ bezeichnen\footnote{Zur Bezeichnung der Ideale gebrauchen wir mit \gesperrt{Dedekind} die kleinen \gesperrt{deutschen Buchstaben}.} und dem Functional $\varphi$ zu-
%CONT>>


% ===== scan leaf 570, printed 549 =====
% !PRINTED 549
\weberpage{549}
% running head: \S.\,151. Ideale. 549
%<<CONT
ordnen. Dasselbe Ideal $\mathfrak{a}$ ist dann auch sämmtlichen mit $\varphi$
associirten Functionalen zugeordnet. Diese Zuordnung hat die
Eigenschaften 1., 4., 5.

Sind $\varphi$ und $\varphi_1$ zwei nicht associirte Functionale, so ist
gewiss eines von ihnen, etwa $\varphi$, nicht durch das andere $\varphi_1$
theilbar, und folglich giebt es (nach \S.\,143, 3.) ganze Zahlen, die
durch $\varphi$, aber nicht durch $\varphi_1$ theilbar sind. Folglich sind nicht
associirten Functionalen immer \gesperrt{verschiedene} Ideale $\mathfrak{a}$, $\mathfrak{a}_1$ zugeordnet.

Es ist aber nun auch zu zeigen, dass auf diese Weise \gesperrt{alle}
Ideale des Körpers $\Omega$ erhalten werden können, mit anderen
Worten, dass jedes von $\mathfrak{o}$ verschiedene Ideal $\mathfrak{a}$ aus der Gesammtheit
der durch einen gewissen Functionalfactor theilbaren Zahlen
besteht.

Wir gehen also jetzt von irgend einem Ideal $\mathfrak{a}$ aus und
wählen eine beliebige endliche Menge von Zahlen daraus,
$\alpha_1, \alpha_2, \ldots, \alpha_r$, deren grösster gemeinschaftlicher Theiler $\delta_r$ sein
mag. Dieses Functional $\delta_r$ hat eine endliche Anzahl von Primfactoren.

Giebt es nun eine Zahl $\alpha_{r+1}$ in $\mathfrak{a}$, die nicht durch $\delta_r$ theilbar
 ist, so hat der grösste gemeinschaftliche Theiler $\delta_{r+1}$ von
$\delta_r$ und $\alpha_{r+1}$ weniger Primfactoren als $\delta_r$. Wenn wir mit dieser
Schlussweise fortfahren, so kommen wir zu dem Ergebniss, dass
sich aus $\mathfrak{a}$ eine endliche Zahl von Zahlen $\alpha_1, \alpha_2, \ldots, \alpha_m$ so auswählen
 lässt, dass der grösste gemeinschaftliche Theiler $\delta$ dieser
Zahlen in \gesperrt{allen Zahlen von $\mathfrak{a}$ aufgeht.}

Andererseits gehört jede durch $\delta$ theilbare Zahl in $\mathfrak{o}$ zu $\mathfrak{a}$.
Denn nach \S.\,149, 4. kann jede durch $\delta$ theilbare Zahl $\alpha$ in die
Form gesetzt werden

\[
\alpha = \alpha_1 \xi_1 + \alpha_2 \xi_2 + \cdots + \alpha_m \xi_m,
\]

worin $\xi_1, \xi_2, \ldots, \xi_m$ Zahlen in $\mathfrak{o}$ sind, und folglich gehört nach
I. und II. $\alpha$ zum Ideal $\mathfrak{a}$.

Es ergiebt sich hieraus, dass, wenn $\delta$ eine Einheit ist, das
Ideal $\mathfrak{a}$ mit $\mathfrak{o}$ identisch ist. Jedes Ideal $\mathfrak{a}$ ist also dadurch
charakterisirt, dass alle seine Zahlen einen gewissen grössten
gemeinschaftlichen Theiler haben.

Es bleibt noch zu zeigen, dass, wenn die Functionale $\varphi, \psi$
den beiden Idealen $\mathfrak{a}, \mathfrak{b}$ entsprechen, das Product $\varphi\psi$ dem
Ideal $\mathfrak{a}\mathfrak{b}$ entspricht.


% ===== scan leaf 571, printed 550 =====
% !PRINTED 550
\weberpage{550}
% running head: 550 Siebzehnter Abschnitt. \S.\,151.
Da alle Zahlen aus $\mathfrak{a}\mathfrak{b}$ von der Form $\Sigma\alpha\beta$ sind, so ist zunächst
 klar, dass alle diese Zahlen durch $\varphi\psi$ theilbar sind.

Wenn wir aber nach \S.\,143, 4. das Functional $\varphi$ als grössten
gemeinschaftlichen Theiler zweier Zahlen $\alpha_1, \alpha_2$ darstellen, so
gehören diese Zahlen, als durch $\varphi$ theilbar, dem Ideal $\mathfrak{a}$ an,
und wir können, da es auf einen Einheitsfactor bei $\varphi$ nicht
ankommt,

\[
\varphi = \alpha_1 x_1 + \alpha_2 x_2
\]

setzen, wenn $x_1, x_2$ Variable sind. Ebenso können wir, wenn
$\beta_1, \beta_2$ zwei Zahlen aus $\mathfrak{b}$ und $y_1, y_2$ Variable bedeuten,

\[
\psi = \beta_1 y_1 + \beta_2 y_2
\]

setzen, und daraus ergiebt sich

\[
\varphi\psi = \alpha_1 \beta_1 x_1 y_1 + \alpha_1 \beta_2 x_1 y_2 + \alpha_2 \beta_1 x_2 y_1 + \alpha_2 \beta_2 x_2 y_2.
\]

Es ist also $\varphi\psi$ nach \S.\,143, 2. der grösste gemeinschaftliche
Theiler der vier Zahlen $\alpha_1\beta_1, \alpha_1\beta_2, \alpha_2\beta_1, \alpha_2\beta_2$, die dem Ideal
$\mathfrak{a}\mathfrak{b}$ angehören, und folglich ist $\varphi\psi$ der \gesperrt{grösste gemeinschaftliche
Theiler} aller Zahlen des Ideals $\mathfrak{a}\mathfrak{b}$.

Damit ist die gegenseitige Zuordnung der Functionale und
Ideale den Forderungen 1. bis 5. gemäss bewerkstelligt.

Den \gesperrt{Primfunctionalen} entsprechen bei dieser Zuordnung
\gesperrt{Primideale}, und die Zerlegung der Ideale in Primfactoren und
überhaupt die Gesetze der Theilbarkeit der Ideale ergeben sich
in völliger Uebereinstimmung mit den entsprechenden Sätzen
aus der Theorie der Functionale.

Bei dieser völligen Uebereinstimmung kann es zu keiner
Unzuträglichkeit führen, wenn wir das ganze System aller unter
einander associirten Functionale zu einem Gemeinbegriffe zusammenfassen
 und dafür den Namen \gesperrt{Ideal} brauchen.

Das System aller mit einer ganzen Zahl associirten Functionale
 ist dann ein \gesperrt{Hauptideal}.

Wenn irgend eine Zahl oder ein Functional durch die
Functionale eines Ideals theilbar ist, so nennen wir es \gesperrt{durch
das Ideal theilbar}, und wenn ein Functional $\varphi$ in Factoren
zerlegt ist, denen die Ideale $\mathfrak{a}, \mathfrak{b}, \ldots$ entsprechen, so setzen wir
auch, indem die Einheitsfactoren in der Bezeichnung weggelassen
 werden:

\weq{1}{\varphi = \mathfrak{a}\mathfrak{b} \ldots}

Eine Basis des Functionals ist zugleich eine Basis des Ideals,
und die absolute Norm des repräsentirenden Functionals stimmt
%CONT>>


% ===== scan leaf 572, printed 551 =====
% !PRINTED 551
\weberpage{551}
% running head: \S.\,151. Ideale. 551
%<<CONT
mit der Zahl überein, die bei \gesperrt{Dedekind} die \gesperrt{Norm des Ideals}
heisst. Sie soll also auch hier so genannt werden, und wir setzen
demnach, wenn das Functional $\varphi$ zu dem Ideal $\mathfrak{a}$ gehört,

\weq{2}{N_a(\varphi) = N(\mathfrak{a}).}

Die Norm eines Ideals ist also immer eine natürliche Zahl.

Ist $\mathfrak{p}$ ein Primideal und $p$ die durch $\mathfrak{p}$ theilbare natürliche
Primzahl, so ist

\weq{3}{N(\mathfrak{p}) = p^f,}

und $f$ heisst der \gesperrt{Grad} des Primideals $\mathfrak{p}$.

In den Congruenzen (\S.\,149, 150) können nach diesen Festsetzungen
 statt der Functionale auch die entsprechenden Ideale
als Moduln angesehen werden.

Es sei schliesslich noch ein Wort über die \gesperrt{Kummer'sche}
Schöpfung der idealen Zahlen beigefügt\footnote{Vgl. \gesperrt{Dirichlet-Dedekind}, Zahlentheorie, \S.\,176.}.

Sind $\alpha, \beta$ irgend zwei Zahlen aus $\mathfrak{o}$, die den grössten
gemeinschaftlichen Theiler $\chi$ haben, und ist $\beta = \chi\varphi$, so wird
die Congruenz

\weq{4}{\alpha\omega \equiv 0 \pmod{\beta}}

nur für solche Zahlen $\omega$ aus $\mathfrak{o}$ befriedigt sein, die durch $\varphi$
theilbar sind. Das Functional $\varphi$ wird aber im Allgemeinen
keiner Zahl associirt sein. Statt nun die Functionale zu benutzen,
kann man sich ein neues Begriffssystem schaffen, in dem
man die sämmtlichen Zahlen von $\mathfrak{o}$ zu Paaren zusammenfasst
$(\alpha, \beta)$, und kann diese Paare \glqq \gesperrt{ideale Zahlen}\grqq{} in $\mathfrak{o}$ nennen;
genau in derselben Weise, wie man durch Paarung der ganzen
rationalen Zahlen die Brüche, oder der reellen Zahlen die
complexen Zahlen gebildet hat.

Man nennt dann die Zahl $\omega$ durch die ideale Zahl $(\alpha, \beta)$
theilbar, wenn sie der Congruenz (1) genügt. % [sic] printed reference reads \textquotedbl{}(1)\textquotedbl{}; the congruence introduced on this page is numbered (4)
Dies ist der Standpunkt
 der \gesperrt{Kummer'schen} Theorie. Damit aber diese Definition
der Theilbarkeit fruchtbar sei, muss noch festgestellt werden,
unter welchen Voraussetzungen man zwei ideale Zahlen mit verschiedenen
 Elementen, $(\alpha, \beta), (\alpha', \beta')$, als gleich zu betrachten
hat, wie sich die wirklichen Zahlen in $\mathfrak{o}$ unter die idealen Zahlen
einordnen, endlich was man unter dem Product zweier idealen
Zahlen zu verstehen hat, damit eine durch eine ideale Zahl
%CONT>>


% ===== scan leaf 573, printed 552 =====
% !PRINTED 552
\weberpage{552}
\pdfbookmark[2]{§ 152. Aequivalenz}{bm173}
% running head: 552 Siebzehnter Abschnitt. \S.\,152.
%<<CONT
theilbare Zahl sich als Product von idealen Zahlen darstellen
lasse. Alle diese Fragen werden beantwortet durch die Vermittelung
 der Functionaltheorie (oder der Idealtheorie), wenn man
der idealen Zahl $(\alpha, \beta)$ das Functional $\varphi$ zuordnet, und können
wohl kaum auf einfachere Weise entschieden werden. Die Einführung
 des Systemes der idealen Zahlen bietet dann auch keine
wesentlichen Vortheile mehr, und man kann ebenso gut die
Functionale selbst, oder die Ideale als Elemente der neu einzuführenden
 Mannigfaltigkeit betrachten.

\begin{center}
\S.\,152.

\gesperrt{Aequivalenz.}
\end{center}

Wir nennen jetzt zwei Functionale, die sich nur durch einen
Einheitsfactor unterscheiden, auch wenn sie gebrochen sind,
\gesperrt{associirt}, und stellen folgende Definition auf:

1. \gesperrt{Zwei ganze oder gebrochene Functionale $\varphi, \psi$
im Körper $\Omega$ heissen äquivalent, wenn ihr
Quotient $\varphi : \psi$ mit einer Zahl associirt ist.}

Es heissen also die beiden Functionale $\varphi$ und $\psi$ äquivalent,
wenn eine Einheit $\varepsilon$ und eine Zahl $\alpha$ in $\Omega$ existiren, so dass

\weq{1}{\frac{\varphi}{\psi} = \alpha\varepsilon}

ist. Ist $\varphi$ äquivalent mit $\psi$ und mit $\psi_1$, so folgt aus (1)

\[
\frac{\varphi}{\psi} = \alpha\varepsilon, \qquad \frac{\varphi}{\psi_1} = \alpha_1\varepsilon_1,
\]

folglich

\[
\frac{\psi}{\psi_1} = \frac{\alpha_1}{\alpha} \frac{\varepsilon_1}{\varepsilon},
\]

und da $\alpha_1 : \alpha$ eine Zahl, $\varepsilon_1 : \varepsilon$ eine Einheit ist, so folgt der
erste Satz:

2. \gesperrt{Zwei Functionale, die mit einem dritten äquivalent
 sind, sind auch unter einander äquivalent
.}

Theilt man hiernach alle Functionale des Körpers $\Omega$ in
\gesperrt{Classen} ein, indem man zwei Functionale in dieselbe oder in
verschiedene Classen wirft, je nachdem sie äquivalent sind oder
nicht, so ergiebt sich, dass zwei dieser Classen, die ein einziges
%CONT>>


% ===== scan leaf 574, printed 553 =====
% !PRINTED 553
\weberpage{553}
%<<CONT
gemeinsames Element enthalten, vollständig identisch sein müssen, und die Classeneintheilung ist also durchaus eindeutig. Jede Classe ist durch ein beliebiges in ihr enthaltenes Functional, einen \gesperrt{Repräsentanten}, völlig bestimmt.

3. \gesperrt{Zwei mit einander associirte Functionale sind auch äquivalent und kommen daher in derselben Classe vor.}

Denn wenn $\varphi$ und $\psi$ associirt sind, so ist ihr Quotient $\varphi:\psi$ eine Einheit, und $\varphi$ und $\psi$ sind also auch äquivalent.

Eine Classe enthält aber nicht bloss die einzelnen Functionale $\varphi$, sondern alle durch diese Functionale bestimmten \gesperrt{Ideale}, und wir nennen diese Classen daher, wenn eine genauere Bezeichnung nöthig ist, \gesperrt{Functionalclassen} oder, häufiger noch, dem üblichen Sprachgebrauche gemäss, \gesperrt{Idealclassen}.

4. \gesperrt{Die Gesammtheit aller ganzen und gebrochenen Zahlen des Körpers $\Omega$, verbunden mit den sämmtlichen Einheiten und den Producten von Zahlen mit Einheiten, bilden unter sich eine Classe, die die Hauptclasse genannt wird.}

Als Repräsentanten der Hauptclasse kann man z. B. die Zahl 1 betrachten. Die Hauptclasse, als Idealclasse aufgefasst, enthält das Ideal $\mathfrak{o}$ und wird daher in der Folge durch den Buchstaben $O$ bezeichnet.

5. \gesperrt{In jeder Idealclasse giebt es ganze Functionale.}

Denn nach der Definition ist, wenn $\varphi$ irgend ein Functional und $\alpha$ eine Zahl ist, $\varphi$ mit $\alpha\varphi$ äquivalent. Wir können aber nach \S.\,137, 5. die Zahl $\alpha$, sogar rational, so bestimmen, dass $\alpha\varphi$ ein ganzes Functional wird.

6. \gesperrt{Aus jeder Idealclasse $C$ können wir einen Repräsentanten $\varphi$ auswählen, der nicht nur selbst ein ganzes Functional ist, sondern auch zu einem beliebig gegebenen ganzen Functional $\omega$ relativ prim ist.}

Nehmen wir, um diesen Satz zu beweisen, zunächst nach 5. einen beliebigen ganzen Repräsentanten $\varphi$ der Classe $C$ und eine durch $\varphi$ theilbare ganze Zahl $\alpha$, so ist

\weq{2}{\varphi\chi = \alpha,}
%CONT>>


% ===== scan leaf 575, printed 554 =====
% !PRINTED 554
\weberpage{554}
\pdfbookmark[2]{§ 153. Die Classenzahl des Körpers Ω}{bm174}
%<<CONT
und $\chi$ ein ganzes Functional. Nun wählen wir (nach \S.\,143, 3.) eine durch $\chi$ theilbare Zahl $\beta$ so, dass $\beta:\chi$ relativ prim zu $\omega$ wird, und setzen

\weq{3}{\psi\chi = \beta.}

Da jetzt $\varphi:\psi$ eine Zahl ist, so ist $\psi$ mit $\varphi$ äquivalent, und $\psi$ ist ein zu $\omega$ theilerfremder Repräsentant der Classe $C^1$).

\S.\,153.

\gesperrt{Die Classenzahl des Körpers $\Omega$.}

Wir kommen nun zum Beweise des wichtigen Satzes:

1. \gesperrt{Die Anzahl der Idealclassen eines Körpers $\Omega$ ist endlich.}

Dieser Satz ist gleichbedeutend mit dem folgenden:

2. \gesperrt{In jeder Classe giebt es ganze Functionale, deren absolute Norm eine bestimmte endliche, nur von der Natur des Körpers $\Omega$ abhängige Zahl nicht übersteigt.}

Denn weil jedes ganze Functional ein Factor seiner absoluten Norm ist, und jede ganze Zahl nur eine endliche Anzahl von nichtassociirten Factoren hat, so giebt es, von den associirten abgesehen, nur eine endliche Zahl von ganzen Functionalen, die eine gegebene Zahl zur absoluten Norm haben. Wenn nun bewiesen werden kann, dass in jeder Classe ein ganzes Functional vorkommt, dessen absolute Norm unter einer endlichen Zahl

\begin{footnotesize}
$^1$) Wir wollen hier im Vorübergehen auf eine Analogie der Aequivalenz der Ideale hinweisen. Die ganze Theorie der algebraischen Zahlen lässt sich mit den nothwendigen Modificationen übertragen auf die Theorie der algebraischen Functionen, die wieder ihren geometrischen Ausdruck in der Theorie der algebraischen Curven findet. Den Idealen entsprechen dann Punktsysteme auf einer festen Grundcurve und den Hauptidealen volle Schnittpunktsysteme der Grundcurve mit einer anderen algebraischen Curve. Wenn sich zwei Punktsysteme zu einem vollen Schnittpunktsystem ergänzen, was in der obigen Formel $\varphi\chi = \alpha$ seinen Ausdruck finden würde, so wird von den beiden Punktsystemen $\varphi$, $\chi$ jedes der \gesperrt{Rest} des anderen genannt. Zwei Punktsysteme, die, wie $\varphi$, $\psi$, denselben Rest haben, werden in der Geometrie corresidual genannt. Dieser Begriff entspricht also der Aequivalenz. (Vergl. \gesperrt{Brill} u. \gesperrt{Nöther}, Mathem. Annalen, Bd. 7. \gesperrt{Salmon}, \glqq higher plane Curves\grqq{}, deutsch von \gesperrt{Fiedler}.)
\end{footnotesize}
%CONT>>


% ===== scan leaf 576, printed 555 =====
% !PRINTED 555
\weberpage{555}
%<<CONT
liegt, so ist die Endlichkeit der Anzahl der Classen nachgewiesen.

Lassen wir, wie bisher, $\omega_1, \omega_2, \ldots, \omega_n$ eine Minimalbasis von $\Omega$ bedeuten, so werden alle ganzen Zahlen des Körpers aus

\weq{1}{\omega = \omega_1 x_1 + \omega_2 x_2 + \cdots + \omega_n x_n}

erhalten, wenn wir den $x_1, x_2, \ldots, x_n$ ganze rationale Zahlwerthe ertheilen. Wenn wir eine positive ganze Zahl $k$ annehmen und festsetzen, dass keine der Zahlen $x_i$ aus dem Intervall $\pm k$ heraustreten soll, so wird der \gesperrt{absolute Werth} von $\omega$ nicht über der Grenze

\[
(|\omega_1| + |\omega_2| + \cdots + |\omega_n|) k = rk
\]

liegen, wenn unter $|\omega_1|, |\omega_2|, \ldots$ die absoluten Werthe (Einleitung, Bd. I, S. 18) der (reellen oder complexen) Grössen $\omega_i$ verstanden sind, und $r$ die Summe $|\omega_1| + |\omega_2| + \cdots$ bedeutet. Bilden wir den Ausdruck (1) für die $n$ conjugirten Körper, und nehmen das Product, so erhalten wir, wenn wir mit $R$ eine positive reelle Zahl bezeichnen, die über dem Product der $n$ Werthe $r$ liegt,

\weq{2}{N_a(\omega) < R k^n.}

Diese Zahl $R$ ist nur von der Natur des Körpers $\Omega$, nicht aber von $k$ abhängig.

Jetzt sei $\mu$ irgend ein ganzes Functional in $\Omega$, und $N_a(\mu)$ seine absolute Norm. Wenn wir die ganze Zahl $k$ so bestimmen, dass

\weq{3}{k^n \leqq N_a(\mu) < (k+1)^n,}

und wenn wir ferner in (1) den Zahlen $x_i$ die Werthe $0, 1, 2, \ldots, k$ ertheilen, so ist die Anzahl der verschiedenen Werthe, die aus (1) hervorgehen, $(k+1)^n$, also grösser als $N_a(\mu)$. Nach \S.\,148, 3. ist aber die Zahl der nach dem Modul $\mu$ incongruenten Zahlen gleich $N_a(\mu)$, und folglich müssen unter den so bestimmten Zahlen $\omega$ mindestens zwei verschiedene nach dem Modul $\mu$ congruente Zahlen vorkommen. Ist also $\omega' \equiv \omega'' (\mathrm{mod}\ \mu)$, so wird die Differenz

\weq{4}{\alpha = \omega' - \omega'' = (x_1' - x_1'') \omega_1 + \cdots + (x_n' - x_n'') \omega_n}

durch $\mu$ theilbar sein, und zugleich sind die ganzen Zahlen

\[
x_1' - x_1'' = a_1, \ldots, x_n' - x_n'' = a_n
\]
%CONT>>


% ===== scan leaf 577, printed 556 =====
% !PRINTED 556
\weberpage{556}
%<<CONT
absolut genommen nicht grösser als $k$. Es giebt eine durch $\mu$ theilbare von Null verschiedene Zahl:

\weq{5}{\alpha = a_1 \omega_1 + a_2 \omega_2 + \cdots + a_n \omega_n,}

in der die ganzzahligen Coëfficienten $a_1, a_2 \ldots a_n$ die Grenzen $\pm k$ nicht überschreiten, und folglich ist nach (2) und (3)

\weq{6}{N_a(\alpha) < R k^n < R N_a(\mu).}

Da nun $\alpha$ durch $\mu$ theilbar ist, so setzen wir

\weq{7}{\alpha = \mu\varphi,\ N_a(\alpha) = N_a(\mu) N_a(\varphi),}

und erhalten aus (6)

\weq{8}{N_a(\varphi) < R.}

Wenn nun $\psi$ ein Repräsentant einer beliebig gegebenen Classe $C$ ist, so wählen wir $\mu$ so, dass

\[
\beta = \mu\psi
\]

eine Zahl ist, und wenn dann nach (7)

\[
\alpha = \mu\varphi
\]

ist, so ist

\[
\frac{\varphi}{\psi} = \frac{\alpha}{\beta},
\]

also $\varphi$ und $\psi$ äquivalent. $\varphi$ ist also gleichfalls ein Repräsentant der Classe $C$, und dieser genügt der Bedingung (8). Es kommt also, wie bewiesen werden sollte, in jeder Classe ein Functional vor, dessen absolute Norm unter $R$ liegt.

Die Anzahl der Idealclassen, die wir mit $h$ bezeichnen wollen, ist hiernach eine dem Körper $\Omega$ eigenthümliche natürliche Zahl, die die \gesperrt{Classenzahl} genannt wird.

In dem einfachsten Falle, wo die Classenzahl gleich 1 ist, ist jedes Functional mit einer Zahl associirt, d. h. es kann jedes (ganze oder gebrochene) Functional durch Absonderung eines Zahlenfactors in eine Einheit verwandelt werden. In diesem Falle lässt sich jede ganze Zahl des Körpers $\Omega$ in \gesperrt{Primzahlfactoren} zerlegen, und diese Körper haben eine Theorie, die im Wesentlichen mit der rationalen Zahlentheorie übereinstimmt. Für solche Körper ist die Einführung der Functionale und Ideale nicht nothwendig.

Hierher gehören neben dem Körper der rationalen Zahlen unter anderen der Körper der \gesperrt{Gauss}'schen imaginären Zahlen und der aus dritten Einheitswurzeln gebildeten Zahlen (Bd. I, \S.\,173, 174).


% ===== scan leaf 578, printed 557 =====
% !PRINTED 557
\weberpage{557}
\pdfbookmark[2]{§ 154. Die Gruppe der Idealclassen}{bm175}

\S.\,154.

\begin{center}
\gesperrt{Die Gruppe der Idealclassen.}
\end{center}

Die Idealclassen können auf Grund des folgenden Satzes componirt werden:

1.~\gesperrt{Sind $\varphi$, $\psi$, $\varphi_1$, $\psi_1$ Functionale und $\varphi$ äquivalent mit $\varphi_1$, $\psi$ äquivalent mit $\psi_1$, so ist auch $\varphi\psi$ äquivalent mit $\varphi_1 \psi_1$.}

Die Richtigkeit hiervon ergiebt sich unmittelbar aus der Definition. Denn wenn $\varphi:\varphi_1$ und $\psi:\psi_1$ mit Zahlen associirt sind, so ist auch $\varphi\psi:\varphi_1\psi_1$ mit einer Zahl associirt.

Ebenso ergiebt sich auch der umgekehrte Satz:

2.~\gesperrt{Ist $\varphi$ äquivalent mit $\varphi_1$ und $\varphi\psi$ äquivalent mit $\varphi_1\psi_1$, so ist auch $\psi$ äquivalent mit $\psi_1$.}

Betrachten wir also zwei Idealclassen $A$, $B$, die auch identisch sein können, und bilden das Product $\varphi\psi$ irgend eines Functionals $\varphi$ aus $A$ und eines Functionals $\psi$ aus $B$, so ist die Classe $C$, in der das Product $\varphi\psi$ vorkommt, unabhängig von der Wahl von $\varphi$ und $\psi$, und die Classe $C$ ist durch die beiden Classen $A$, $B$ völlig bestimmt. \gesperrt{Wir nennen $C$ aus $A$ und $B$ componirt} und schreiben symbolisch

\weq{1}{C = AB = BA.}

Die Classe $C$ enthält alle Producte eines Elementes von $A$ mit einem Element von $B$, kann aber auch noch andere Functionale enthalten.

Da diese Composition aus der wahren Multiplication abgeleitet ist, so gelten auch die Gesetze der Multiplication für diese Composition, nämlich das commutative und das associative Gesetz.

Es folgt ferner aus dem Satze 2., dass, wenn $AB = AB_1$ ist, auch $B = B_1$ sein muss, und folglich erzeugen die Idealclassen bei dieser Composition eine \gesperrt{endliche Abel'sche Gruppe vom Grade $h$}, auf die wir alle Sätze anwenden können, die wir im zweiten Abschnitt dieses Bandes über solche Gruppen kennen gelernt haben.

Die Einheit dieser Gruppe ist die schon im \S.\,152 definirte Hauptclasse $O$, die ja, wie wir gesehen haben, den Repräsen-
%CONT>>


% ===== scan leaf 579, printed 558 =====
% !PRINTED 558
\weberpage{558}
%<<CONT
tanten 1 hat. Um entgegengesetzte Classen zu definiren, nehmen wir einen Repräsentanten $\varphi$ einer Classe $A$ und eine durch $\varphi$ theilbare Zahl $\alpha$. Ist dann $\alpha = \varphi\chi$, so ist $\chi$ ein Repräsentant der Classe $A^{-1}$, und es ist $AA^{-1} = O$.

Ist $A$ eine beliebige Classe, und $h$ die Classenzahl, so ist immer

\[ A^h = O, \]

und wenn $k$ die kleinste positive Zahl ist, die der Bedingung

\[ A^k = O \]

genügt, so ist $k$ ein Theiler von $h$. Daraus ergiebt sich der Satz:

3.~\gesperrt{Jedes Functional $\varphi$ in $\Omega$ gehört zu einem bestimmten Exponenten $k$, der ein Theiler der Classenzahl $h$ ist, so dass $\varphi^k$ associirt ist mit einer Zahl in $\Omega$.}

\begin{center}
\rule{4cm}{0.4pt}
\end{center}


% ===== scan leaf 580, printed 559 =====
% !PRINTED 559
\weberpage{559}
\pdfbookmark[1]{Achtzehnter Abschnitt. Discriminanten.}{bm176}
\pdfbookmark[2]{§ 155. Minimum einer quadratischen Form}{bm177}
\begin{center}
\gesperrt{Achtzehnter Abschnitt.}

\gesperrt{Discriminanten.}

\rule{0.3\textwidth}{0.4pt}

\S.\,155.

\gesperrt{Minimum einer quadratischen Form.}
\end{center}

Auf einem ganz neuen Weg hat \gesperrt{Minkowski} die Endlichkeit der Anzahl der Classen bewiesen\footnote{Ueber die positiven quadratischen Formen etc. \gesperrt{Crelle}, Bd.\ 107, 1891. Eingehender noch in dem Buche von \gesperrt{Minkowski}, \glqq{}Geometrie der Zahlen\grqq{}, Leipzig 1896.}, der von eigenthümlichen, zahlreicher Anwendungen fähigen Betrachtungen ausgeht und dabei zu dem lange vergeblich gesuchten Beweis eines Satzes über die Körperdiscriminante führt. Wir wollen hier in der Kürze die Hauptmomente dieser Untersuchungen mittheilen.

\gesperrt{Minkowski} geht aus von der Betrachtung einer quadratischen Form von $n$ Variablen

\weq{1}{f(x_1, x_2, \ldots, x_n) = \Sigma^{i,k} a_{i,k}\, x_i x_k}

mit reellen Co\"efficienten $a_{i,k}$, von der vorausgesetzt wird, dass sie sich in $n$ und nicht weniger positive Quadrate linearer Functionen zerlegen lässt (einer positiven Form, Bd.\ I, \S.\,81 bis 83). Es möge eine dieser Darstellungen sein

\weq{2}{f(x_1, x_2, \ldots, x_n) = \xi_1^2 + \xi_2^2 + \cdots + \xi_n^2,}
worin
\weq{3}{\xi_i = \alpha_{1,i} x_1 + \alpha_{2,i} x_2 + \cdots + \alpha_{n,i} x_n}

ein System linearer Functionen der $x_i$ mit nicht verschwindender Determinante ist. Setzen wir

\weq{4}{A = \Sigma \pm \alpha_{1,1}\alpha_{2,2}\ldots\alpha_{n,n},}
%CONT>>


% ===== scan leaf 581, printed 560 =====
% !PRINTED 560
\weberpage{560}
%<<CONT
so ist die Determinante $D$ der quadratischen Form (1) (Bd.\,I, \S.\,59)
\weq{5}{D = \Sigma \pm a_{1,1}a_{2,2}\ldots a_{n,n} = A^2.}

Die $\alpha_{i,k}$ sind dann gleichfalls reelle Zahlen, die auf unendlich viele Arten bestimmt werden k\"onnen. Wir nehmen eine beliebige, aber weiterhin unver\"anderliche Transformation (2), (3) an.

Das Gleichungssystem (3) m\"oge, nach den $x$ aufgel\"ost, ergeben:
\weq{6}{x_i = \beta_{1,i}\xi_1 + \cdots + \beta_{n,i}\xi_n.}

Wir wollen nun den Variablen $x_i$ nur ganze rationale Zahlwerthe beilegen. Dann ist klar, dass die Function $f$ nur verschwinden kann, wenn alle $x_i$ verschwinden, und dass ihr Werth, wenn wir diesen Fall ausschliessen, nicht unter einen gewissen endlichen Grenzwerth $M$ heruntersinken kann, den wir das \gesperrt{Minimum der Form} $f$ nennen. Denn sinkt der Werth von $f$ unter einen Werth $\alpha^2$ herunter, so m\"ussen nach (2) alle $\xi_i$ dem absoluten Werthe nach unter $\alpha$ herunter sinken, und nach (6) w\"urden f\"ur ein hinl\"anglich kleines $\alpha$ s\"ammtliche $x_i$ unter 1 herabgedr\"uckt werden und m\"ussten dann, da es ganze Zahlen sind, alle gleich Null sein.

Es kommt vor Allem darauf an, nicht sowohl das Minimum genau zu bestimmen, als eine m\"oglichst kleine obere Grenze daf\"ur zu finden, \"uber die es nicht hinausgehen kann.

Wir bedienen uns jetzt mit \gesperrt{Minkowski} einer Ausdrucksweise, die einer Geometrie im Raume von $n$ Dimensionen entnommen ist, die ausserordentlich einfach zu dem gew\"unschten Resultat f\"uhrt. F\"ur die F\"alle $n=2$, $n=3$ sind die dabei gebrauchten Ausdr\"ucke vollst\"andig anschaulich und allereinfachster Art. Wir empfehlen dem Leser, sich die nun auszuf\"uhrenden Schritte an diesen beiden speciellen F\"allen anschaulich klar zu machen. Bei gr\"osseren Werthen von $n$ geht freilich die eigentliche Anschauung verloren. Man hat dann diese Ausdr\"ucke als kurze Bezeichnung gewisser analytischer Verh\"altnisse anzusehen, wobei noch die Analogie mit dem gew\"ohnlichen Raume die Auffassung erleichtert. Es hat nicht die geringste Schwierigkeit, alle diese Ausdr\"ucke durch rein analytische Formeln zu ersetzen; es w\"urde darunter aber sehr die K\"urze des Ausdrucks und die Leichtigkeit der Auffassung leiden.

Wir bezeichnen also jetzt die Variablen $\xi_1, \ldots, \xi_n$ als Coor-
%CONT>>


% ===== scan leaf 582, printed 561 =====
% !PRINTED 561
\weberpage{561}
%<<CONT
dinaten eines Punktes $\xi$ in einem Raume von $n$ Dimensionen
$R_\xi$ und nennen den Ausdruck

\weq{7}{\varrho = \sqrt{(\xi_1' - \xi_1'')^2 + (\xi_2' - \xi_2'')^2 + \cdots + (\xi_n' - \xi_n'')^2},}

die \gesperrt{Entfernung} der beiden Punkte $\xi'$, $\xi''$.

Wenn wir in (3) die Variablen $x_i$ alle ganzzahligen Werthe
durchlaufen lassen, so geben uns die zugehörigen Werthe der $\xi_i$
ein Punktsystem im Raume $R_\xi$. Der Punkt, der den Werthen $x_i =
0$ entspricht, ist der Coordinaten-Anfangspunkt oder der \gesperrt{Nullpunkt}.
Die Gesammtheit dieser Punkte bleibt ungeändert, wenn wir $x_i$
durch $x_i + x_i^{(0)}$ ersetzen, wenn $x_1^{(0)}, x_2^{(0)}, \ldots,
x_n^{(0)}$ ein festes System ganzer Zahlen bedeutet. Wir können diese
Thatsache auch so ausdrücken, dass das Punktsystem der $\xi_i$ mit
sich selbst zur Deckung kommt, wenn es eine Parallelverschiebung
erfährt, bei dem der Nullpunkt an den Ort eines beliebig gegebenen
anderen Punktes des Systems gelangt; oder auch: Die Punkte des
Systems sind zu jedem seiner Punkte in derselben Weise gelagert, wie
zu jedem anderen.

Wir wollen dies System ein \gesperrt{Punktgitter} nennen und seine
einzelnen Punkte die \gesperrt{Gitterpunkte}.

Nach (2) und (7) ist der Werth von $f$ das Quadrat der Entfernung des
Punktes $\xi$ vom Nullpunkte und folglich ist $\sqrt{M}$ die
\gesperrt{kleinste Entfernung zweier Gitterpunkte}, die in dem ganzen System
vorkommt.

Ausser den Entfernungen in dem Raume $R_n$ haben wir auch noch
Volumina zu betrachten, d.\,h. die Werthe des $n$-fachen Integrals

\weq{8}{V_\xi = \int\!\!\int \ldots \int d\xi_1\, d\xi_2 \ldots d\xi_n,}

worin die Variablen $\xi_1, \xi_2, \ldots, \xi_n$ irgend ein
endliches Gebiet $G_\xi$ erfüllen, was etwa durch eine oder mehrere
Ungleichungen bestimmt ist. So ist das durch die Bedingung

\weq{9}{\xi_1^2 + \xi_2^2 + \cdots + \xi_n^2 \leqq 1 \text{ oder auch } \Sigma(\xi_i - \xi_i^0)^2 \leqq 1}

bei beliebigen Werthen $\xi_i^0$ begrenzte Gebiet eine
\gesperrt{$n$-dimensionale Kugel} vom Radius 1, deren Volumen wir mit $K_n$
bezeichnen wollen. Aus jedem Gebiete $G_\xi$ können wir ein
\gesperrt{ähnliches Gebiet} $G_\xi^{(t)}$ ableiten, indem wir die Coordinaten
alle im Verhältniss $t^{\frac{1}{n}} : 1$ vergrössern. Ist $V^{(t)}$
das Volumen des neuen Gebietes, so ist nach (8)

\weq{10}{V^{(t)} = t\,V.}


% ===== scan leaf 583, printed 562 =====
% !PRINTED 562
\weberpage{562}
Betrachten wir ausser dem Raume $R_\xi$ noch einen zweiten Raum $R_x$,
in dem die durch (3) oder (6) bestimmten Variablen $x$ die Coordinaten
sind, so entspricht jedem Punkte des einen Raumes ein und nur ein
Punkt des anderen Raumes. Einem endliche % [sic]
Gebiete $G_\xi$ entspricht ein endliches Gebiet $G_x$.

Um die Beziehung zwischen dem Volumen $V_\xi$ und $V_x$ zu ermitteln,
müssen wir das Integral (8) nach den Regeln über die Umformung
mehrfacher Integrale behandeln, wofür die allgemeine Formel gilt:

\weq{11}{\int\!\int\!\ldots\!\int d\xi_1\, d\xi_2 \ldots d\xi_n
= \int\!\int\!\cdots\!\int
\sqrt{\left(\Sigma\pm\frac{d\xi_1}{dx_1}\frac{d\xi_2}{dx_2}\cdots\frac{d\xi_n}{dx_n}\right)^{2}}\,
dx_1\, dx_2 \cdots dx_n\ \footnotemark.}

\footnotetext{Die allgemeine Formulirung dieser Umformung rührt von
\gesperrt{Jacobi} her (De determinantibus functionalibus; \gesperrt{Crelle},
Bd. 22, 1891, % [sic: Weber prints 1891.]
gesammelte Werke, Bd. 3). Die Ableitung findet sich in
den meisten ausführlicheren Lehrbüchern der Integralrechnung, z. B.
\gesperrt{Serret-Harnack}, Bd. 2; \gesperrt{Lipschitz}, Lehrbuch der
Analysis Bd. 2; auch \gesperrt{Baltzer}, Determinanten.}

Setzen wir darin für die Functionaldeterminante den Werth (4) oder
(5) ein, so ergiebt sich

\weq{12}{V_\xi = \sqrt{D}\ V_x,}

worin nun das Integral

\weq{13}{V_x = \int\!\int\!\cdots\!\int d x_1\, d x_2 \cdots d x_n}

sich auf das Gebiet $G_x$ bezieht.

Das Integral $V_x$ können wir als Grenzwerth einer Summe von
Elementen $\Delta x_1 \Delta x_2 \ldots \Delta x_n$ auffassen, und
wenn wir den Incrementen $\Delta x_i$ allen die Grösse $1 : t^{\frac1n}$
geben, so wird jedes dieser Elemente den Werth $1 : t$ haben.
Bedeutet $Z_t$ die Anzahl dieser Elemente, so ist $V_x$ der
Grenzwerth, dem sich das Verhältniss $Z_t : t$ bei unendlich
wachsendem $t$ nähert, also

\weq{14}{V_x = \mathrm{Lim}_{t\,=\,\infty} \frac{Z_t}{t}.}

Nun ist die Anzahl der Elemente $Z_t$ ebenso gross, wie die Anzahl
der Punkte im Inneren von $G_x$, deren Coordinaten den Ausdruck

\[
x_1 = h_1 t^{-\frac1n},\ x_2 = h_2 t^{-\frac1n},\ \ldots,\ x_n = h_n t^{-\frac1n}
\]

mit ganzzahligen $h_1, h_2, \ldots, h_n$ haben, oder auch gleich der An-
%CONT>>


% ===== scan leaf 584, printed 563 =====
% !PRINTED 563
\weberpage{563}
%<<CONT
zahl der Punkte mit ganzzahligen Coordinaten, die in einem im
Verhältniss $t^{\frac1n} : 1$ vergrösserten Gebiet $G_x^{(t)}$
liegen, und diese Zahl ist ebenso gross wie die Anzahl der
Gitterpunkte, die in dem Gebiete $G_\xi^{(t)}$ liegen. Wir haben
also folgenden Satz:

\gesperrt{Ist $G_\xi$ irgend ein endliches Gebiet im Raume $R_\xi$
und $G_\xi^{(t)}$ das aus $G_\xi$ durch Vergrösserung im
Linearverhältniss $t^{\frac1n} : 1$ abgeleitete Gebiet; ist $Z_t$
die Anzahl der im Inneren von $G_\xi^{(t)}$ liegenden Gitterpunkte,
so ist das durch $G_\xi$ bestimmte Volumen}

\weq{15}{V_\xi = \sqrt{D}\ \mathrm{Lim}_{t\,=\,\infty} \frac{Z_t}{t}.}

Wir legen jetzt um jeden Gitterpunkt als Mittelpunkt eine
$n$-dimensionale Kugel vom Radius $\frac12\sqrt{M}$. Da $\sqrt{M}$
die kleinste Entfernung von irgend zwei Gitterpunkten ist, so können
je zwei dieser Kugeln kein Gebiet gemein haben; sie können sich
höchstens in einem Punkte berühren, und keine dieser Kugeln enthält
ausser ihrem Mittelpunkte einen weiteren Gitterpunkt. Eine solche
Kugel hat nach (10) das Volumen

\[
K_n \left(\frac{\sqrt{M}}{2}\right)^{n},
\]

und das Gesammtvolumen aller der Kugeln, deren Mittelpunkt im
Inneren von $G_\xi^{(t)}$ liegt, ist

\[
V_1^{(t)} = Z_t\, K_n \left(\frac{\sqrt{M}}{2}\right)^{n}.
\]

Dies Volumen $V_1^{(t)}$ ist aber offenbar kleiner als $V^{(t)}$,
wenigstens wenn wir von den Kugeln, die zum Theil aus dem Gebiete
$G^{(t)}$ herausragen, absehen, deren Gesammtvolumen, durch $t$
dividirt, mit unendlich wachsendem $t$ gegen Null convergirt, weil
sein Gebiet nur einen Theil von $G^{(t)}$ erfüllt, und folglich ist
auch $V$ grösser als der Grenzwerth von $V_1^{(t)} : t$, also

\[
V > \mathrm{Lim}\, \frac{Z_t\, K_n}{t} \left(\frac{\sqrt{M}}{2}\right)^{n},
\]

und daraus nach (15)

\weq{16}{\sqrt{D} > K_n \left(\frac{\sqrt{M}}{2}\right)^{n}.}

Da es nun leicht ist, das Volumen $K_n$ zu bestimmen, so ist
hierdurch die Aufgabe gelöst, und eine obere Grenze für $M$
gefunden.


% ===== scan leaf 585, printed 564 =====
% !PRINTED 564
\weberpage{564}
Wir wollen hier nicht den genauen Werth für $K_n$ einsetzen, sondern
einen einfacher auszudrückenden kleineren Werth, wodurch die
Ungleichung (11) um so mehr richtig wird.

Nach der Definition ist

\weq{17}{K_n = \int\!\int\!\cdots\!\int d\xi_1\, d\xi_2 \ldots d\xi_n,}

mit der Begrenzung

\weq{18}{\xi_1^2 + \xi_2^2 + \cdots + \xi_n^2 \leqq 1.}

Die Grenzgleichung (18) ist sicher erfüllt, wenn wir die $\xi_i$
auf die Intervalle

\weq{19}{-\frac{1}{\sqrt{n}} < \xi_i < \frac{1}{\sqrt{n}}}

beschränken, was geometrisch die Bedeutung hat, dass wir die Kugel
durch den eingeschriebenen Würfel ersetzen. Demnach ist

\[
K_n > \int\limits_{-\frac{1}{\sqrt n}}^{\frac{1}{\sqrt n}}
\int\limits_{-\frac{1}{\sqrt n}}^{\frac{1}{\sqrt n}}
\cdots
\int\limits_{-\frac{1}{\sqrt n}}^{\frac{1}{\sqrt n}}
d\xi_1\, d\xi_2 \cdots d\xi_n,
\]

und wenn diese Integration ausgeführt wird,

\weq{20}{K_n > \left(\frac{2}{\sqrt{n}}\right)^{n}.}

Die hierdurch bestimmte obere Grenze wird von $K_n$ niemals
erreicht, wenn wir von dem Falle $n = 1$ absehen.

Die Vergleichung von (20) mit (16) ergiebt einen sehr einfachen
Ausdruck für den oberen Grenzwerth von $M$, nämlich

\weq{21}{M < n \sqrt[n]{D},}

und $M$ ist (von dem Falle $n = 1$ abgesehen) wirklich kleiner, und
nicht gleich dem Werthe $n \sqrt[n]{D}$.

Dies zeigt sich klar, wenn man eine noch etwas engere Grenze für
$M$ aufsucht.

Die Bedingung (18) ist sicher dann erfüllt, wenn wir

\weq{22}{-\alpha < \xi_n < \alpha,\qquad \xi_1^2 + \xi_2^2 + \cdots + \xi_{n-1}^2 < 1 - \alpha^2}

setzen, wo $\alpha$ ein beliebiger positiver echter Bruch sein kann.
(Die geometrische Bedeutung hiervon ist die, dass man die Kugel
durch einen eingeschriebenen Cylinder von der Höhe $2\alpha$
ersetzt). Es ist also

\[
\begin{gathered}
K_n > 2\alpha \int\!\int\!\cdots\!\int d\xi_1\, d\xi_2 \ldots d\xi_{n-1}\\
\xi_1^2 + \xi_2^2 + \cdots + \xi_{n-1}^2 < 1 - \alpha^2.
\end{gathered}
\]


% ===== scan leaf 586, printed 565 =====
% !PRINTED 565
\weberpage{565}
Wenn man in dem $(n-1)$fachen Integral $\xi_i$ durch
$\sqrt{1-\alpha^2}\,\xi_i$ ersetzt, so findet man dafür den Werth
$\sqrt{(1-\alpha^2)}^{\,n-1} K_{n-1}$, und es findet sich also die
Ungleichung

\[
K_n > 2\alpha \sqrt{(1-\alpha^2)}^{\,n-1} K_{n-1}.
\]

Nun wollen wir $\alpha$ so wählen, dass $2\alpha
\sqrt{(1-\alpha^2)}^{\,n-1}$ so gross als möglich wird, und dafür
findet man nach der gewöhnlichen Regel durch Nullsetzen des
Differentialquotienten

\[
\alpha = \frac{1}{\sqrt{n}}.
\]

Daraus

\[
K_n > 2\, \frac{\sqrt{(n-1)}^{\,n-1}}{\sqrt{n}^{\,n}}\, K_{n-1}
\]

\[
\left(\frac{\sqrt{n}}{2}\right)^{n} K_n > \left(\frac{\sqrt{n-1}}{2}\right)^{n-1} K_{n-1}.
\]

Setzt man also

\[
K_n = \left(\frac{2}{\sqrt{\Theta_n\, n}}\right)^{n},
\]

so ist

\weq{23}{\Theta_n^n < \Theta_{n-1}^{n-1},}

also $\Theta_n^n$ eine mit wachsenden $n$ abnehmende Zahl. Zugleich
ergiebt sich aus (16)

\weq{24}{M < \Theta_n\, n\, \sqrt[n]{D}.}

Für $n = 2$ ist $K_n$ gleich dem Flächeninhalt des Kreises mit dem
Radius 1, d. h. gleich der \gesperrt{Ludolfischen} Zahl $\pi$, und
daraus ergiebt sich

\[
\Theta_2 = \frac{2}{\pi} < 0{,}636,\qquad \Theta_2^2 < 0{,}404,
\]

% [sic: Weber prints references (22) and (23); the bound uses (24) and (23).]
und es gilt also nach (22) und (23) die für alle Werthe von $n$,
die grösser als 1 sind, gültige Grenzbestimmung:

\weq{25}{M < n \sqrt[n]{0{,}404\, D}\ \footnotemark.}

\footnotetext{Für specielle Anwendungen ist es nützlich, die Grenze
so genau als möglich zu kennen. Dazu dient eine genaue Berechnung
des Integrals $K_n$ nach seiner Definition (17), (18), die sich
leicht nach verschiedenen Methoden der Integralrechnung, z. B.
mittelst des discontinuirlichen Factors von \gesperrt{Dirichlet},
ausführen lässt:

\[
K_n = \frac{\Gamma\left(\frac12\right)^{n}}{\Gamma\left(1+\frac{n}{2}\right)},
\]
}


% ===== scan leaf 587, printed 566 =====
% !PRINTED 566
\weberpage{566}
\pdfbookmark[2]{§ 156. Anwendung auf algebraische Körper}{bm178}

\begin{center}
\S.\,156.

\gesperrt{Anwendung auf algebraische K\"orper.}
\end{center}

Aus dem Minimum f\"ur eine positive quadratische Form kann
man ein Minimum von Linearformen ableiten, indem man eine
positive quadratische Form als Summe von Quadraten von Linearformen
 auffasst. Darauf beruhen die folgenden Anwendungen
auf die Theorie der algebraischen K\"orper, denen wir einen H\"ulfssatz
 vorausschicken:

1. \gesperrt{Sind $a_1, a_2, \ldots, a_n$ reelle positive Zahlwerthe, so ist}
\weq{1}{a_1 + a_2 + \cdots + a_n \geqq n \sqrt[n]{a_1 a_2 \ldots a_n}.}

Dieser Satz ist offenbar richtig f\"ur $n=2$, denn es ist
$a_1 + a_2 - 2\sqrt{a_1 a_2} = (\sqrt{a_1} - \sqrt{a_2})^2$, also nie negativ. Wir
wenden also die vollst\"andige Induction an, indem wir die Formel
(1) f\"ur $n-1$ Glieder als schon erwiesen annehmen. Nehmen
wir dann an, dass $a_n$ von keinem der \"ubrigen $a$ an Gr\"osse \"ubertroffen
 wird, so ist
\[ a_n \geqq \sqrt[n-1]{a_1 a_2 \ldots a_{n-1}} = b, \]
und
\[ a_1 + a_2 + \cdots + a_{n-1} \geqq (n-1)\,b. \]
Also ist
\[ a_1 + a_2 + \cdots + a_n \geqq a_n + (n-1)\,b, \]
und folglich
\[ \frac{a_1 + a_2 + \cdots + a_n}{n} \geqq b + \frac{a_n - b}{n}, \]

\begin{footnotesize}
\noindent\rule{4cm}{0.4pt}

worin $\Gamma(x)$ die Gammafunction bedeutet, f\"ur die die Gleichungen
\[ \Gamma(x) = (x-1)\,\Gamma(x-1), \quad \Gamma\left(\frac{1}{2}\right) = \sqrt{\pi} \]
bestehen, so dass $K_n$ durch rationale Zahlen und durch $\sqrt{\pi}$ ausgedr\"uckt ist.

(Vergl. z.\,B. Vorlesungen \"uber die Theorie der bestimmten Integrale
nach \gesperrt{Lejeune-Dirichlet} von G.\,F. \gesperrt{Meyer} herausgegeben, \S.\,175.
Leipzig 1871).

Daraus ergiebt sich
\[ \Theta_n = \frac{4}{n\pi}\,\Gamma\left(1 + \frac{n}{2}\right)^{\frac{2}{n}}, \]
was nach den N\"aherungsformeln f\"ur die $\Gamma$-Function f\"ur grosse
Werthe von $n$ nahe mit
\[ \frac{2}{\pi e} = 0{,}234\ldots \]
\"ubereinstimmt. (\gesperrt{Minkowski}, Crelle, Bd. 107, S. 292).
\end{footnotesize}
%CONT>>


% ===== scan leaf 588, printed 567 =====
% !PRINTED 567
\weberpage{567}
%<<CONT

und nach dem binomischen Satze, da $a_n - b$ nicht negativ ist,
\[ \left(\frac{a_1 + a_2 + \cdots + a_n}{n}\right)^n \geqq b^n + b^{n-1}(a_n - b) = a_n b^{n-1}, \]
dies aber f\"allt mit der zu beweisenden Ungleichung (1) zusammen.

Es sei nun $\omega_1, \omega_2, \ldots, \omega_n$ eine Minimalbasis eines algebraischen
 K\"orpers $\Omega$. Es sei ferner $\varphi$ ein beliebiges ganzes Functional
.

Wir nehmen eine Basis von $\varphi$ an (\S.\,146):
\weq{2}{(\beta_1, \beta_2, \ldots, \beta_n) = \begin{pmatrix} b_{1,1} & \ldots & b_{1,n} \\ \cdot & \cdots & \cdot \\ b_{n,1} & \ldots & b_{n,n} \end{pmatrix} (\omega_1, \omega_2, \ldots, \omega_n),}
so dass
\weq{3}{N_a(\varphi) = \pm \Sigma \pm b_{1,1} \ldots b_{n,n}}
wird. Aus der mit $\varphi$ associirten Linearform
\weq{4}{\eta = \beta_1 x_1 + \beta_2 x_2 + \cdots + \beta_n x_n}
erhalten wir dann alle durch $\varphi$ theilbaren ganzen Zahlen des
K\"orpers $\Omega$, wenn f\"ur die $x_1, x_2, \ldots, x_n$ alle m\"oglichen Combinationen
 ganzer Zahlen gesetzt werden.

Aus (4) ergeben sich $n$ verschiedene Formen $\eta_1, \eta_2, \ldots, \eta_n$,
wenn wir zu den $n$ conjugirten K\"orpern, $\Omega_1, \Omega_2, \ldots, \Omega_n$ \"ubergehen
. Das Quadrat der Determinante dieses Systems von
$n$ Formen ist nach \S.\,147 gleich der Discriminante
\weq{5}{\Delta(\beta_1, \beta_2, \ldots, \beta_n) = \pm N_a(\varphi)^2 \Delta,}
wenn $\Delta$ die Grundzahl des K\"orpers bedeutet.

Die conjugirten K\"orper $\Omega_1, \Omega_2 \ldots$ k\"onnen theils reell, theils
paarweise conjugirt imagin\"ar sein. Wir bezeichnen, wie fr\"uher
schon mehrfach, mit $|\omega|$ den \gesperrt{absoluten Werth} einer reellen
oder complexen Gr\"osse und betrachten, indem wir unter
$x_1, x_2, \ldots, x_n$ reelle Variable verstehen, die quadratische Form
\weq{6}{f(x_1, x_2, \ldots, x_n) = |\eta_1|^2 + |\eta_2|^2 + \cdots + |\eta_n|^2 = S\,|\eta|^2.}
Ist $\Omega$ reell, so ist
\weq{7}{|\eta|^2 = \eta^2,}
also das Quadrat einer reellen linearen Form. Wenn aber $\Omega, \Omega'$
zwei conjugirt imagin\"are unter den conjugirten K\"orpern sind, so
sind auch die entsprechenden Linearformen $\eta, \eta'$ conjugirt imagin
\"ar, und es ist
\weq{8}{|\eta|^2 = |\eta'|^2 = \eta \eta'}
%CONT>>


% ===== scan leaf 589, printed 568 =====
% !PRINTED 568
\weberpage{568}
%<<CONT
die Summe von zwei Quadraten reeller Linearformen. Das Product
 $\eta \eta'$ tritt dann mit dem Factor 2 in der Function $f$ auf.
Daher ist $f$ unter allen Umst\"anden eine positive Form.

Nach (7), (8) l\"asst sich aber $f$ auch als quadratische Form
der Variablen $\eta_1, \eta_2, \ldots, \eta_n$
\weq{9}{f = F(\eta_1, \eta_2, \ldots, \eta_n) = \sum^{i,k} a_{i,k}\,\eta_i \eta_k}
auffassen, und hierin ist, wenn $\eta_i$ reell ist, $a_{i,i} = 1$, und wenn
$\eta_i, \eta_k$ ein conjugirt imagin\"ares Paar ist, $a_{i,k} = 1$. Alle \"ubrigen
Co\"efficienten sind $= 0$. Die Determinante von $F$ ist also $= \pm 1$,
weil in jeder Zeile und in jeder Colonne nur ein von Null verschiedenes
 Element mit dem Werthe 1 vorkommt. Nach Bd. I,
\S.\,56 ist die Determinante der Form $f$ gleich dem Product der
Determinante von $F$ und dem Quadrate der Determinante der
linearen Formen $\eta$, also nach (5)
\weq{10}{= \pm [N_a(\varphi)]^2 \Delta,}
wo das Zeichen so zu w\"ahlen ist, dass der Werth positiv ausf\"allt.

Nach \S.\,155, (24) kann man nun f\"ur die $x_i$ solche ganze
rationale Zahlen setzen, dass
\weq{11}{f(x_1, x_2, \ldots, x_n) < \Theta_n\, n \sqrt[n]{\pm \Delta N_a(\varphi)^2}}
wird, worin $\Theta_n$, wenn $n$ gr\"osser als 1 ist, ein echter Bruch und
immer kleiner als $\sqrt[n]{0{,}404}$ ist. Setzt man diese Werthe von $x_i$ in
(4) ein, so geht $\eta$ in eine ganze durch $\varphi$ theilbare Zahl $\beta$ \"uber,
und wir haben also damit den Satz:

2. \gesperrt{Ist $\varphi$ irgend ein ganzes Functional in $\Omega$, so giebt}
\gesperrt{es eine durch $\varphi$ theilbare ganze Zahl $\beta$ in $\Omega$ von}
\gesperrt{der Eigenschaft, dass}
\weq{12}{S\,|\eta|^2 < \Theta_n\, n \sqrt[n]{\pm \Delta [N_a(\varphi)]^2}}
\gesperrt{wird, worin $\Delta$ die Grundzahl des K\"orpers und}
\gesperrt{$\pm \Delta$ positiv ist.}

Das Product der conjugirten Werthe $\eta$ ist gleich der Norm
von $\eta$, und daher ist nach dem H\"ulfssatze 1. [mit R\"ucksicht auf
(7) und (8)]
\weq{13}{n \sqrt[n]{[N(\eta)]^2} \leqq S\,|\eta|^2,}
und also folgt nach (12) die auf die absoluten Werthe bez\"ugliche
 Ungleichung:
\weq{14}{N(\eta)^2 < \Theta_n^n\, [N_a(\varphi)]^2\, \Delta.}
Da nun $\eta$ durch $\varphi$ theilbar ist, so k\"onnen wir
\[ \varepsilon \eta = \varphi \psi, \quad \pm N(\eta) = N_a(\varphi)\, N_a(\psi) \]
%CONT>>


% ===== scan leaf 590, printed 569 =====
% !PRINTED 569
\weberpage{569}
\pdfbookmark[2]{§ 157. Primfactoren der natürlichen Primzahlen}{bm179}
%<<CONT
setzen, worin $\psi$ ein ganzes Functional, $\varepsilon$ eine Einheit ist, und
erhalten aus (14)
\weq{15}{N_a(\psi) < \sqrt{\pm 0{,}404\, \Delta}.}

Da wir f\"ur $\varphi$ jedes ganze Functional setzen k\"onnen, so kann
$\psi$ ein Repr\"asentant eines Ideals einer beliebigen Classe sein, und
wir haben damit den wichtigen Satz bewiesen:

3. \gesperrt{In jeder Idealclasse k\"onnen wir einen Repr\"asen-}
\gesperrt{tanten finden, dessen Norm kleiner als die}
\gesperrt{Quadratwurzel aus dem absoluten Werthe der}
\gesperrt{Grundzahl des K\"orpers ist.}

Die im \S.\,153 zum Beweise der Endlichkeit der Classenzahl
benutzte obere Grenze f\"ur die Norm eines Ideals einer Classe
wird hierdurch in viel sch\"arferer Weise bestimmt.

Aber noch einen anderen sehr bemerkenswerthen Schluss
gestattet die Ungleichung (15). Da n\"amlich $N_a(\psi)$ eine positive
ganze Zahl, also mindestens gleich 1 ist, so geht aus (15) hervor
, dass, abgesehen von dem Falle $n=1$, $\Delta$ absolut gr\"osser als
1 sein muss, \gesperrt{dass es also ausser dem K\"orper der rationalen
Zahlen keinen K\"orper giebt, dessen Grundzahl $\pm 1$ w\"are.} Dies
ist der von \gesperrt{Minkowski} zuerst erbrachte, fr\"uher lange vergeblich
 gesuchte Beweis dieses wichtigen Satzes\footnote{In einem Briefe an \gesperrt{Hermite} (Comptes rendus d. Pariser Akademie,
26. Januar 1891) und auch in dem oben erw\"ahnten Buche findet \gesperrt{Min-}
\gesperrt{kowski} eine weit engere Grenze, indem er nicht nur das Minimum
der quadratischen Formen, sondern noch anderer Functionen betrachtet,
wie $\eta_1^p + \eta_2^p + \cdots + \eta_n^p$, worin $p$ eine beliebige positive Zahl ist, die
nicht einmal ganz zu sein braucht. F\"ur $p=1$ ergiebt sich, wenn $n-\nu$
die Anzahl der conjugirt imagin\"aren Paare ist:
\[ N_a(\psi) < \left(\frac{4}{\pi}\right)^{n-\nu} \frac{1.2.3\ldots n}{n^n}\, \sqrt{\pm \Delta}. \]
Daraus lassen sich noch weitere Schl\"usse ziehen, wie die: Die Grundzahl
 eines quadratischen K\"orpers muss gr\"osser als 4 oder keiner als % [sic: Weber prints ``keiner'' for ``kleiner''.]
$-2$ sein; die Grundzahl eines cubischen K\"orpers muss gr\"osser als 20
oder kleiner als $-12$ sein etc.}.

\begin{center}
\S.\,157.

\gesperrt{Primfactoren der nat\"urlichen Primzahlen.}
\end{center}

Eine genauere Untersuchung der im \S.\,146 erkl\"arten Basisform
 $\tau$ von $\mathfrak{o}$ des K\"orpers $\Omega$ soll uns nun das Mittel geben,
%CONT>>


% ===== scan leaf 591, printed 570 =====
% !PRINTED 570
\weberpage{570}
%<<CONT
jede beliebig gegebene natürliche Primzahl $p$ in ihre Primfactoren
wirklich zu zerlegen, und damit alle Primideale des Körpers $\Omega$,
oder wenigstens Repräsentanten aller Primideale wirklich
darzustellen\footnote{Wir folgen hier einer Arbeit von \gesperrt{Hensel}: \glqq Untersuchung der
Fundamentalgleichung einer Gattung für eine reelle Primzahl als Modul
und Bestimmung der Theiler der Discriminante\grqq{} (\gesperrt{Crelle}'s Journ., Bd. 113).
Zu erwähnen sind auch die anderen Arbeiten von \gesperrt{Hensel}, Ebend.,
Bd. 101, 104, 105, 113.}.

Es sei $\mathfrak p$ ein beliebiges Primideal vom Grade $f$, so dass
\weq{1}{N(\mathfrak p) = p^f}
ist, worin $p$ eine natürliche Primzahl bedeutet, die durch den
Primfactor $\mathfrak p$ theilbar ist. $f$ ist ein positiver Exponent
$\leqq n$. Die kleinste ganze rationale Zahl, die durch $\mathfrak p$
theilbar ist, ist $p$, und jede andere ganze rationale Zahl, die durch
$\mathfrak p$ theilbar ist, ist daher auch durch $p$ theilbar
(\S.\,138).

Wir müssen nun auch Congruenzen mit dem Modul $\mathfrak p$ zwischen
ganzen Functionen beliebiger Variablen mit Coëfficienten in
$\mathfrak o$ betrachten und bemerken, dass zwei solche Functionen
nach \S.\,143, 2. \gesperrt{dann und nur dann congruent sind, wenn die
entsprechenden Coëfficienten congruent sind.}

Den Körper der rationalen Zahlen bezeichnen wir mit $R$ und nennen
demnach eine Function mit rationalen Coëfficienten auch eine \gesperrt{Function
in} $R$.

Die Basisform von $\mathfrak o$
\weq{2}{\tau = \omega_1 t_1 + \omega_2 t_2 + \cdots + \omega_n t_n}
genügt, wie wir im \S.\,146 gesehen haben, einer Gleichung
$n^{\text{ten}}$ Grades $F(\tau) = 0$, deren Coëfficienten ganze
rationale Functionen von $t_1, t_2, \ldots, t_n$ sind. Es ist also
$F(\tau)$ jedenfalls durch $\mathfrak p$ theilbar, und daraus folgt,
dass es ganze Functionen $\Phi(t)$ in $R$ giebt, die ausser $t$
irgend welche Variable enthalten können, die durch die Substitution
$t = \tau$ in durch $\mathfrak p$ theilbare Functionale in
$\mathfrak o$ übergehen, die also der Congruenz
\weq{3}{\Phi(\tau) \equiv 0\ (\mathrm{mod}\ \mathfrak p)}
genügen, und wir werden also sagen können, $\tau$ \gesperrt{ist eine Wurzel der
Congruenz}
\weq{4}{\Phi(t) \equiv 0\ (\mathrm{mod}\ \mathfrak p).}

Die Function $\Phi$ wird gewiss die Variablen $t_1, t_2, \ldots, t_n$
enthalten müssen; sie kann aber auch noch andere Variable ent-
%CONT>>


% ===== scan leaf 592, printed 571 =====
% !PRINTED 571
\weberpage{571}
%<<CONT
halten, und wenn wir also die Variablen von $\Phi$ mit $t, u_1, u_2,
\ldots$ bezeichnen, werden wir auch setzen
\weq{5}{\Phi(t) = \Phi(t, u_1, u_2, \ldots),}
oder kürzer $\Phi(t, u)$, und diese Function enthält ganze rationale
Zahlencoëfficienten. Wenn wir die Function $\Phi(t)$ in die
$p^{\text{te}}$ Potenz erheben, und die Formel \S.\,150, 4. anwenden,
so folgt aus der Congruenz (3) eine neue Congruenz
\weq{6}{\Phi(\tau^p, u_1^p, u_2^p, \ldots) \equiv 0\ (\mathrm{mod}\ \mathfrak p).}

Ebenso ist aber auch
\weq{7}{\tau^p \equiv \omega_1^p t_1^p + \omega_2^p t_2^p + \cdots + \omega_n^p t_n^p.}

Wenn wir dies in die Congruenz (6) substituiren, so entsteht eine
durch $\mathfrak p$ theilbare Function, in der die Variablen
$u_1, u_2, \ldots$, unter denen ja die $t_1, t_2, \ldots, t_n$ mit
enthalten sind, nur in der $p^{\text{ten}}$ Potenz vorkommen. Die
Congruenz muss also richtig bleiben, wenn wir die $u_1^p, u_2^p,
\ldots$ und also auch die $t_1^p, t_2^p, \ldots, t_n^p$ durch
unabhängige Variable $u_1, u_2, \ldots t_1, t_2, \ldots, t_n$ ersetzen
(\S.\,143).

Setzen wir demnach
\weq{8}{\tau_1 = \omega_1^p t_1 + \omega_2^p t_2 + \cdots + \omega_n^p t_n,}
so ergiebt sich aus (6)
\[
\Phi(\tau_1, u_1, u_2, \ldots) \equiv 0\ (\mathrm{mod}\ \mathfrak p),
\]
und folglich ist $\tau_1$ auch eine Wurzel der Congruenz (4).

Dieses nämliche Verfahren lässt sich wiederholt anwenden, und wir
finden, dass auch
\[
\tau_2 = \omega_1^{p^2} t_1 + \omega_2^{p^2} t_2 + \cdots + \omega_p^{p^2} t_n % [sic: Weber prints \omega_p here; (9) has \omega_n.]
\]
Wurzel der Congruenz (4) ist, u.\,s.\,f.

Wenn wir also ein System von
Formen $\tau_r$ definiren durch
\weq{9}{\tau_r = \omega_1^{p^r} t_1 + \omega_2^{p^r} t_2 + \cdots + \omega_n^{p^r} t_n}
für beliebige positive Exponenten $r$, so sind alle diese Grössen
$\tau_r$ zugleich Wurzeln der Congruenz (4). Es ist noch die Frage zu
beantworten, wie viele von diesen Formen $\tau_r$ von einander
verschieden sind.

Nach \S.\,150, (4) ist sicher
\[
\tau_r \equiv \tau_{r'}\ (\mathrm{mod}\ \mathfrak p),
\]
wenn
\[
r \equiv r'\ (\mathrm{mod}\ f)
\]
ist, und folglich giebt es unter den $\tau_r$ gewiss nicht mehr als
$f$ nach dem Modul $\mathfrak p$ verschiedene
\weq{10}{\tau, \tau_1, \tau_2, \ldots, \tau_{f-1}.}


% ===== scan leaf 593, printed 572 =====
% !PRINTED 572
\weberpage{572}
Dass diese Formen aber wirklich von einander verschieden sind,
ergiebt sich daraus, dass wir (nach \S.\,145, 150) für die Variablen
$t_1, t_2, \ldots, t_n$ solche ganze rationale Zahlen setzen können,
dass $\tau$ in eine primitive Wurzel $\gamma$ des Primideals
$\mathfrak p$ übergeht. Durch dieselbe Substitution werden die
Grössen (10)
\weq{11}{\equiv \gamma, \gamma^p, \gamma^{p^2}, \ldots, \gamma^{p^{f-1}}\ (\mathrm{mod}\ \mathfrak p),}
die nach dem Modul $\mathfrak p$ alle von einander verschieden sind.
Es können also auch nicht zwei der Formen (10) nach dem Modul
$\mathfrak p$ congruent sein, weil sonst auch die beiden
entsprechenden Zahlen (11) congruent ausfallen würden.

Dies fassen wir als Satz so zusammen:

1. \gesperrt{Jede Congruenz (4), deren eine Wurzel $t = \tau$ ist,
hat die $f$ verschiedenen Wurzeln}
\[
\tau, \tau_1, \tau_2, \ldots, \tau_{f-1}.
\]

Hiernach können wir, indem wir $\Phi(t)$ durch $t - \tau$ algebraisch
dividiren,
\[
\Phi(t) \equiv (t - \tau)\, \Phi_1(t)\ (\mathrm{mod}\ \mathfrak p)
\]
setzen, und $\tau_1, \tau_2, \ldots, \tau_{f-1}$ sind Wurzeln von
$\Phi_1(t) \equiv 0$, worin aber $\Phi_1(t)$ noch nicht rationale
Coëfficienten, sondern Coëfficienten in $\mathfrak o$ hat. Dividiren
wir $\Phi_1(t)$ wieder durch $t - \tau_1$, und fahren so fort, so
folgt endlich, wenn wir also nun die Function $f^{\text{ten}}$ Grades
\weq{12}{\Pi(t) = (t - \tau)(t - \tau_1) \ldots (t - \tau_{f-1})}
setzen,
\weq{13}{\Phi(t) \equiv \Pi(t)\, \Phi_0(t)\ (\mathrm{mod}\ \mathfrak p),}
worin $\Phi_0(t)$ eine ganze Function mit Coëfficienten in
$\mathfrak o$ ist.

Die Function $\Pi(t)$ hängt von den Variablen $t, t_1, \ldots, t_n$
ab, und um dies auszudrücken, setzen wir
\[
\Pi(t) = \Pi(t, t_1, \ldots, t_n).
\]

Die Coëfficienten dieser Form sind Zahlen in $\mathfrak o$, und es
lässt sich noch nachweisen, dass sie mit ganzen rationalen Zahlen
nach dem Modul $\mathfrak p$ congruent sind. Dieser Beweis ergiebt
sich durch Erheben in die $p^{\text{te}}$ Potenz:
\[
[\Pi(t)]^p \equiv (t^p - \tau^p)(t^p - \tau_1^p) \ldots (t^p - \tau_{f-1}^p).
\]

Nun ist aber nach (9)
\[
\tau_r^p \equiv \omega_1^{p^{r+1}} t_1^p + \omega_2^{p^{r+1}} t_2^p + \cdots + \omega_n^{p^{r+1}} t_n^p\ (\mathrm{mod}\ \mathfrak p),
\]
und wir erhalten also $\tau_r^p$ aus $\tau_{r+1}$, wenn wir $t_i$
durch $t_i^p$ ersetzen,
%CONT>>


% ===== scan leaf 594, printed 573 =====
% !PRINTED 573
\weberpage{573}
%<<CONT
und $\tau_f$ ist congruent mit $\tau$. Demnach erhalten wir die
Congruenz:
\[
[\Pi(t, t_1, t_2, \ldots, t_n)]^p \equiv \Pi(t^p, t_1^p, t_2^p, \ldots, t_n^p)\ (\mathrm{mod}\ \mathfrak p).
\]
Damit ist nach \S.\,150, 4. der Satz bewiesen:

2. \gesperrt{Die Form}
\[
\Pi(t) = (t - \tau)(t - \tau_1) \ldots (t - \tau_{f-1})
\]
\gesperrt{ist nach dem Modul $\mathfrak p$ mit einer ganzen und
homogenen Form $f^{\text{ten}}$ Grades in $R$ der Variablen
$t, t_1, t_2, \ldots, t_n$ congruent.}

Diese ganze rationale Form, deren Coëfficienten bis auf Vielfache der
Primzahl $p$ völlig bestimmt sind, wollen wir mit $P(t)$ bezeichnen.

Dann folgt aus (13)
\weq{14}{\Phi(t) \equiv P(t)\, \Phi_0(t)\ (\mathrm{mod}\ \mathfrak p),}
und da nun $\Phi(t)$ eine ganze Function in $R$ ist, so ergiebt sich
durch Erheben zur Potenz $p$:
\[
\begin{aligned}
\Phi(t^p, u^p) &\equiv P(t^p, u^p)\ \Phi_0(t^p, u^p) \\
[\Phi(t, u)]^p &\equiv [P(t, u)]^p\ [\Phi_0(t, u)]^p
\end{aligned}
\qquad (\mathrm{mod}\ \mathfrak p).
\]
Weil aber die linken Seiten nach dem Modul $\mathfrak p$ congruent
sind, so folgt
\[
[P(t)]^p\ [\Phi_0(t, u)^p - \Phi_0(t^p, u^p)] \equiv 0\ (\mathrm{mod}\ \mathfrak p),
\]
und daraus, da $P(t)$ nicht durch $\mathfrak p$ theilbar ist (weil
der Coëfficient von $t^f$ den Werth 1 hat):
\weq{15}{[\Phi_0(t, u)]^p \equiv \Phi_0(t^p, u^p),}
woraus nach \S.\,150, 4. hervorgeht, dass auch $\Phi_0(t)$ mit einer
ganzen rationalen Form nach dem Modul $\mathfrak p$ congruent ist.

Demnach können wir in (14) auch $\Phi_0(t)$ als ganze Form in $R$
annehmen, und dann muss nach \S.\,140, 3. die Congruenz (14) nicht
nur für den Modul $\mathfrak p$, sondern für den Modul $p$ bestehen.
Daraus erhalten wir den Satz:

3. \gesperrt{Unter den Formen $\Phi(t)$, die für $t = \tau$ in ein
durch $\mathfrak p$ theilbares Functional übergehen, ist $P(t)$ vom
Grade $f$ in Bezug auf $t$ die Form niedrigsten Grades, und wenn
$\Phi(t)$ eine beliebige unter ihnen ist, so lässt sich eine ganze
Function $\Phi_0(t)$ in $R$ so bestimmen, dass}
\weq{15}{\Phi(t) \equiv P(t)\, \Phi_0(t)\ (\mathrm{mod}\ p)} % [sic] number (15) reused for a second equation on this page, as printed
\gesperrt{wird.}


% ===== scan leaf 595, printed 574 =====
% !PRINTED 574
\weberpage{574}
Lassen wir in $\Phi(t)$ und $\Phi_0(t)$ Glieder weg, deren Coëfficienten durch $p$ theilbar sind, so ist der Grad von $\Phi(t)$ in Bezug auf $t$ um $f$ grösser als der Grad von $\Phi_0(t)$.

$P(t)$ ist eine ganze Function in $R$ der Variablen $t, t_1, t_2, \ldots, t_n$, die durch die Substitution $t=\tau$ in ein durch $\mathfrak p$ theilbares Functional übergeht, die natürlich nicht mehr in $R$ enthalten ist.

Die Bedeutung dieser Form $P(t)$ tritt nun noch deutlicher hervor, wenn wir den Satz beweisen:

4. \gesperrt{Der Primfactor $\mathfrak p$ ist der grösste gemeinschaftliche Theiler von $p$ und $P(\tau)$.}

Dazu haben wir nur nachzuweisen, dass, wenn $p$ durch $\mathfrak p \mathfrak p_1$ theilbar ist, wo $\mathfrak p, \mathfrak p_1$ zwei gleiche oder verschiedene Primfactoren sind, $P(\tau)$ zwar durch $\mathfrak p$, nicht aber durch $\mathfrak p \mathfrak p_1$ theilbar ist.

Nach \S.\,143, 3. existirt immer eine \gesperrt{Zahl} $\xi$ in $\mathfrak{o}$, die zwar durch $\mathfrak p$, aber nicht durch $\mathfrak p \mathfrak p_1$ theilbar ist. Diese Zahl wird die Form haben
\weq{16}{\xi = a_1 \omega_1 + a_2 \omega_2 + \cdots + a_n \omega_n,}
worin die $a_1, a_2, \ldots, a_n$ ganze rationale Zahlen sind, und geht also aus der Form $\tau$ hervor durch die Substitution
\weq{17}{(t_1, t_2, \ldots, t_n) = (a_1, a_2, \ldots, a_n).}
Wenn wir dieselbe Substitution in den in (9) definirten Formen $\tau_r$ machen, so geht $\tau_r$ in eine Zahl $\xi_r$ über, die nach dem \gesperrt{Fermat}'schen Satze der Congruenz
\[ \xi^{p^r} \equiv \xi_r \pmod{\mathfrak p} \]
genügt und also sicher auch durch $\mathfrak p$ theilbar ist. Wenn wir daher in der Form
\[ \Pi(t) = (t-\tau)(t-\tau_1) \cdots (t - \tau_{f-1}) \]
$\tau_r + \xi_r$ an Stelle von $\tau_r$ setzen, so bleibt diese Form mit sich selbst nach dem Modul $\mathfrak p$ congruent. Diese Substitution kommt aber darauf hinaus, dass wir $t_1+a_1, t_2+a_2, \ldots, t_n+a_n$ an Stelle von $t_1, t_2, \ldots, t_n$ substituiren, und wir erhalten demnach
\[ \Pi(t, t_1+a_1, \ldots, t_n+a_n) \equiv \Pi(t, t_1, \ldots, t_n) \pmod{\mathfrak p}, \]
% [sic: Weber prints ``nnd'' for ``und''.]
nnd wenn wir $\Pi$ durch die congruente Form $P$ ersetzen
\[ P(t, t_1+a_1, \ldots, t_n+a_n) \equiv P(t, t_1, \ldots, t_n) \pmod{\mathfrak p}. \]


% ===== scan leaf 596, printed 575 =====
% !PRINTED 575
\weberpage{575}
Da aber die Form $P$ lauter \gesperrt{rationale} Coëfficienten hat, so muss die letztere Congruenz auch nach dem Modul $p$ stattfinden, also
\weq{18}{P(t, t_1+a_1, \ldots, t_n+a_n) \equiv P(t, t_1, \ldots, t_n) \pmod{p}.}
Diese Congruenz besteht für variable $t, t_1, \ldots, t_n$.

Nehmen wir nun an, dass, entgegen dem zu beweisenden Satze, $P(\tau)$ durch $\mathfrak p \mathfrak p_1$ theilbar sei, so besteht die Congruenz
\weq{19}{P(\tau, t_1, \ldots, t_n) \equiv 0 \pmod{\mathfrak p \mathfrak p_1},}
und diese bleibt richtig, wenn für die unabhängigen Variablen $t_1, t_2, \ldots, t_n$ die Substitution $t_1+a_1, t_2+a_2, \ldots, t_n+a_n$ gemacht wird, und da hierdurch $\tau$ in $\tau+\xi$ übergeht, so folgt
\weq{20}{P(\tau+\xi, t_1+a_1, \ldots, t_n+a_n) \equiv 0 \pmod{\mathfrak p \mathfrak p_1}.}
Machen wir andererseits in der Congruenz (18) die Substitution $t=\tau+\xi$, so folgt
\weq{21}{P(\tau+\xi, t_1, \ldots, t_n) = P(\tau+\xi) \equiv 0 \pmod{\mathfrak p \mathfrak p_1}.}
Nehmen wir nun zunächst an, $\mathfrak p_1$ sei von $\mathfrak p$ verschieden, dann können wir so schliessen.

Wir setzen in (21) $t_1=t_2=\cdots=t_n=0$, also auch $\tau=0$. Dadurch aber geht $P(t)$ in $t^f$ über (abgesehen von Vielfachen von $p$), und es ergiebt sich aus (21)
\[ \xi^f \equiv 0 \pmod{\mathfrak p \mathfrak p_1}, \]
was aber der Annahme widerspricht, dass $\xi$ nicht durch $\mathfrak p_1$ theilbar sein soll.

Ist aber $\mathfrak p_1 = \mathfrak p$, so ordnen wir (21) nach Potenzen von $\xi$ und erhalten
\weq{22}{P(\tau+\xi) = P(\tau) + \xi P'(\tau) + \frac{\xi^2}{2} P''(\tau) \ldots,}
wenn $P'(t), P''(t), \ldots$ die Derivirten von $P(t)$ sind, wobei zu beachten ist, dass die Formen
% [sic: Weber prints $P'''(\tau)$ here although the sentence refers to derivatives of $P(t)$.]
\[ \frac{1}{2} P''(t), \ \frac{1}{2.3} P'''(\tau) \ldots \]
trotz der scheinbaren Nenner ganze Formen in $R$ sind. Da nun nach (19) und (21) $P(\tau+\xi)$ und $P(\tau)$ durch $\mathfrak p^2$ theilbar sind, und ebenso nach Voraussetzung $\xi^2, \xi^3, \ldots$, während $\xi$ nicht durch $\mathfrak p^2$ theilbar ist, so folgt aus (22)
\[ P'(\tau) \equiv 0 \pmod{\mathfrak p}, \]
woraus wegen
\[ P(t) \equiv \Pi(t) \pmod{\mathfrak p} \]
%CONT>>


% ===== scan leaf 597, printed 576 =====
% !PRINTED 576
\weberpage{576}
%<<CONT
nach der Bedeutung (12) von $\Pi(t)$ folgt:
\[ \Pi'(\tau) \equiv (\tau-\tau_1)(\tau-\tau_2) \cdots (\tau-\tau_{f-1}) \equiv 0 \pmod{\mathfrak p}. \]
Dies ist aber unmöglich, weil die $\tau, \tau_1, \ldots \tau_{f-1}$, wie wir gesehen haben, nach dem Modul $\mathfrak p$ incongruent sind. Damit ist also unser Satz 4. vollständig bewiesen. Wir können hiernach, wenn $x, y$ zwei neue Variable bedeuten, ein Functional $\pi$ des Ideals $\mathfrak p$ bilden:
\weq{23}{\pi = x p + y P(\tau).}
Wir betrachten nun ganze Formen $\Phi(t)$ in $R$, die durch die Substitution $t=\tau$ nicht nur durch $\mathfrak p$, sondern durch die natürliche Primzahl $p$ theilbar werden, also der Congruenz
\weq{24}{\Phi(\tau) \equiv 0 \pmod p}
genügen. Die Primzahl $p$ möge folgendermaassen in ihre Primfactoren in $\Omega$ zerlegt sein:
\weq{25}{p = \mathfrak p \mathfrak p_1 \mathfrak p_2 \cdots,}
worin
\[ \mathfrak p, \mathfrak p_1, \mathfrak p_2, \ldots \]
gleiche oder verschiedene Primideale der Grade
\[ f, f_1, f_2, \ldots \]
sind. Diesen Primfactoren entspricht (nach dem Satze 3.) eine Reihe ganzer rationaler Formen
\[ P(t), P_1(t), P_2(t), \ldots \]
der Grade $f, f_1, f_2, \ldots$, und wenn etwa $\mathfrak p$ mit $\mathfrak p_1$ identisch ist, so ist auch $P(t)$ mit $P_1(t)$ identisch. Wenn man in (25) rechts und links die Norm nimmt, und die Formel $N(\mathfrak p) = p^f$ berücksichtigt [\S.\,140, (3), \S.\,151, (3)], so folgt
\weq{26}{n = f + f_1 + f_2 + \cdots}
Wenn nun $\Phi(t)$ eine der Bedingung (24) genügende ganze rationale Form ist, so folgt aus dem Satze 3.
\weq{27}{\Phi(t) \equiv P(t) \Phi_1(t) \pmod p,}
worin $\Phi_1(t)$ eine ganze Function in $R$ ist, die der Bedingung
\[ P(\tau) \Phi_1(\tau) \equiv 0 \pmod{\mathfrak p \mathfrak p_1 \mathfrak p_2 \ldots} \]
genügt. Nach dem Satze 4. folgt hieraus
\weq{28}{\Phi_1(\tau) \equiv 0 \pmod{\mathfrak p_1 \mathfrak p_2 \ldots},}
%CONT>>


% ===== scan leaf 598, printed 577 =====
% !PRINTED 577
\weberpage{577}
%<<CONT
und daraus schliesst man wieder nach Satz 3. (auf $\mathfrak p_1$ angewandt)
\[ \Phi_1(t) \equiv P_1(t) \Phi_2(t) \pmod p. \]
Hierin lässt sich dieselbe Betrachtung wiederholen, die zu der Congruenz
\[ \Phi_2(t) \equiv 0 \pmod{\mathfrak p_2 \ldots} \]
führt, woraus wieder nach 3.
\[ \Phi_2(t) \equiv P_2(t) \Phi_3(t) \pmod p \]
zu schliessen ist. Fährt man damit fort, bis alle Primfactoren von $p$ berücksichtigt sind, so ergiebt sich der folgende Satz:

5. \gesperrt{Ist $\Phi(t)$ eine ganze Function in $R$, die durch die Substitution $t=\tau$ durch $p$ theilbar wird, so lässt sich eine andere ganze Function $\Phi_0(t)$ in $R$ so bestimmen, dass}
\weq{29}{\Phi(t) \equiv \Phi_0(t) P(t) P_1(t) P_2(t) \ldots \pmod p}
\gesperrt{wird.}

\gesperrt{Das Product $P(t) P_1(t) P_2(t) \ldots$ ist vom Grade $n = f+f_1+f_2+\cdots$ und ist die ganze rationale Form niedrigsten Grades, die durch die Substitution $t=\tau$ durch $p$ theilbar wird.}

Zu den im Satze 5. vorkommenden Functionen $\Phi(t)$ gehört auch die ganze Function $n^{\mathrm{ten}}$ Grades
\[ F(t) = N(t-\tau), \]
die für $t=\tau$ verschwindet. Für diese Function wird $\Phi_0(t)$ von $t$ unabhängig, und da sowohl in $F(t)$ als in $P(t), P_1(t), P_2(t) \ldots$ die höchste Potenz von $t$ den Coëfficienten 1 hat, so wird $\Phi_0(t) = 1$. Es gilt also die Congruenz
\weq{30}{N(t-\tau) \equiv P(t) P_1(t) P_2(t) \ldots \pmod p.}
Fassen wir in der Zerlegung (25) der Primzahl $p$ die gleichen Primfactoren zu Potenzen zusammen, so können wir, wenn $\mathfrak p_1, \mathfrak p_2, \ldots$ verschiedene Primideale der Grade $f_1, f_2, \ldots$ sind, die positiven Exponenten $e_1, e_2, \ldots$ so annehmen, dass
\weq{31}{p = \mathfrak p_1^{e_1} \mathfrak p_2^{e_2} \ldots,}
und
\weq{32}{n = e_1 f_1 + e_2 f_2 + \cdots}
wird. Die Formel (30) nimmt dann die Gestalt an
\weq{33}{N(t-\tau) \equiv [P_1(t)]^{e_1} [P_2(t)]^{e_2} \ldots \pmod p.}


% ===== scan leaf 599, printed 578 =====
% !PRINTED 578
\weberpage{578}
\pdfbookmark[2]{§ 158. Dedekind's Satz über die Körperdiscriminante}{bm180}

\begin{center}
\S.\,158.

\gesperrt{Dedekind's Satz über die Körperdiscriminante.}
\end{center}

Die in der Discriminante $\Delta$ des K\"orpers $\Omega$ aufgehenden natürlichen Primzahlen haben in Bezug auf ihre Zerlegung in Primfactoren einen besonderen Charakter, über den ein Satz von \gesperrt{Dedekind} die bündigste Auskunft giebt, zu dessen Ableitung wir jetzt schreiten\footnote{\gesperrt{Dedekind}, \glqq Ueber die Discriminanten endlicher K\"orper\grqq{} im XXIX. Bande der Abhandlungen der Gesellschaft der Wissenschaften in G\"ottingen (1882).}.

Wir bilden nach \S.\,145, 2. die Potenzen der Basisform
\[
\tau = t_1\omega_1 + t_2\omega_2 + \cdots + t_n\omega_n,
\]
und erhalten die Ausdrücke:
\weq{1}{\tau^k = u_{1,k}\,\omega_1 + u_{2,k}\,\omega_2 + \cdots + u_{n,k}\,\omega_n,}
worin die $u_{i,k}$ ganze rationale Functionen der Variablen $t_1, t_2, \ldots, t_n$ sind. Wir setzen wie oben
\weq{2}{F(t) = N(t-\tau),}
so dass $F(\tau)=0$ ist. Nun bilden wir nach \S.\,144 die Discriminante
\weq{3}{\Delta(1,\tau,\tau^2,\ldots,\tau^{n-1}) = (-1)^{\frac{n(n-1)}{2}}\, N F'(\tau),}
wofür sich nach (1) der Werth $\Delta U^2$ ergiebt, wenn $\Delta$ die Körperdiscriminante und
\weq{4}{U = \left|\begin{array}{cccc} u_{1,0}, & u_{2,0}, & \cdots, & u_{n,0} \\ u_{1,1}, & u_{2,1}, & \cdots, & u_{n,1} \\ \cdots & \cdots & \cdots & \cdots \\ u_{1,n-1}, & u_{2,n-1}, & \cdots, & u_{n,n-1} \end{array}\right|,}
also eine ganze Function in $R$ ist.

Wir müssen nachweisen, dass die Form $U$ eine Einheit ist. Angenommen, es gehe in $U$ irgend eine Primzahl $p$ auf, so können wir, wie schon im \S.\,147 bewiesen ist, ein System ganzer Functionen $y_0, y_1, \ldots, y_{n-1}$ in $R$, \gesperrt{die nicht alle durch $p$ theilbar sind}, so bestimmen, dass für $i=1,2,\ldots,n$
\[
y_0\, u_{i,0} + y_1\, u_{i,1} + \cdots + y_{n-1}\, u_{i,n-1} \equiv 0\ (\mathrm{mod}\ p)
\]
wird. Dann aber ergiebt sich aus (1)
\[
y_0 + y_1\tau + \cdots + y_{n-1}\tau^{n-1} \equiv 0\ (\mathrm{mod}\ p).
\]
%CONT>>


% ===== scan leaf 600, printed 579 =====
% !PRINTED 579
\weberpage{579}

%<<CONT
Dies widerspricht aber dem Satze 5. des vorigen Paragraphen, nach dem $\tau$ nach einem Primzahlmodul $p$ keiner Congruenz von niedrigerem als $n^{\mathrm{tem}}$ Grade genügen kann.

Demnach ergiebt sich aus (3), \gesperrt{dass die Grundzahl $\Delta$ des Körpers, vom Vorzeichen abgesehen, die absolute Norm der Function $F'(\tau)$ ist}, dass also
\weq{5}{\pm \Delta = N_a[F'(\tau)]}
zu setzen ist.

Wenn nun statt der $\omega_i$ eine andere Basis $\omega_i'$ von $\mathfrak{o}$ zu Grunde gelegt wird, so tritt an Stelle von $\tau$ eine andere Form
\weq{6}{\tau' = \omega_1' t_1 + \omega_2' t_2 + \cdots + \omega_n' t_n,}
wofür wir auch setzen können
\weq{7}{\tau' = \omega_1 t_1' + \omega_2 t_2' + \cdots + \omega_n t_n',}
und darin sind die $\omega_i'$ mit den $\omega_i$ durch eine lineare Substitution mit der Determinante $\pm 1$ verbunden (\S.\,145):
\weq{8}{(\omega_1', \omega_2', \ldots, \omega_n') = C(\omega_1, \omega_2, \ldots, \omega_n).}
Führt man diese Substitution in (6) aus, und ordnet nach $\omega_1, \omega_2, \ldots, \omega_n$, so ergiebt sich
\weq{9}{(t_1', t_2', \ldots, t_n') = C_1(t_1, t_2, \ldots, t_n),}
wenn $C_1$ die transponirte Substitution von $C$ ist (\S.\,37). Bildet man nun die Function $F(t)$ für die Function $\tau'$, so mag sich $F_1(t)$ ergeben; die Ableitung für $t=\tau'$ sei $F_1'(\tau')$. Setzen wir nun
\[
F'(\tau) = \Psi(t_1, t_2, \ldots, t_n),
\]
so ist $\Psi$ eine ganze Function in $R$, und es ergiebt sich wegen (7)
\[
F_1'(\tau') = \Psi(t_1', t_2', \ldots, t_n').
\]

% [sic: The source prints $F'(\tau')$ as the first functional; the preceding definitions compare $F'(\tau)=\Psi(t_1,\ldots,t_n)$ with $F_1'(\tau')=\Psi(t_1',\ldots,t_n')$. The primed argument is retained as printed.]
Da die Substitution (9) umkehrbar ist, so folgt hieraus, dass die beiden Functionale $F'(\tau')$ und $F_1'(\tau')$ gegenseitig durch einander theilbar sind (\S.\,143), und dass sie mithin associirt sind.

Die Function $F'(\tau)$, deren absolute Norm gleich dem absoluten Werthe der Grundzahl ist, nennen wir daher das \gesperrt{Grundfunctional}, und das durch $F'(\tau)$ repräsentirte Ideal das \gesperrt{Grundideal} des Körpers $\Omega$. Nun sei $\mathfrak{p}^e$ die höchste Potenz des Primideals $\mathfrak{p}$, die in $p$ aufgeht, und $e \geqq 1$, dann können wir nach 5. des vorigen Paragraphen
\weq{10}{F(t) \equiv P(t)^e\, \Phi(t)\ (\mathrm{mod}\ p)}
setzen, worin $\Phi(t)$ eine ganze Function in $R$ von der Beschaffenheit ist, dass $\Phi(\tau)$ nicht durch $\mathfrak{p}$ theilbar ist. Aus (10) aber
%CONT>>


% ===== scan leaf 601, printed 580 =====
% !PRINTED 580
\weberpage{580}

%<<CONT
folgt, indem wir die Ableitung nach $t$ bilden, was offenbar gestattet ist:
\[
F'(t) \equiv e\,[P(t)]^{e-1}\,P'(t)\,\Phi(t) + [P(t)]^e\,\Phi'(t)\ (\mathrm{mod}\ p).
\]
Setzen wir hierin $t=\tau$, so geht $P'(t)$ in eine durch $\mathfrak{p}$ untheilbare Form $P'(\tau)$ über (was wir schon im Beweis von \S.\,157, 4. gezeigt und benutzt haben) und wir erhalten
% [sic: Weber prints $F_1(\tau)$ here; the preceding factorization and following sentence use $\Phi(\tau)$.]
\weq{11}{F'(\tau) \equiv e\, P(\tau)^{e-1}\, P'(\tau)\, F_1(\tau)\ (\mathrm{mod}\ \mathfrak{p}^e).}
Nun ist (nach \S.\,157, 4.) $P(\tau)$ durch $\mathfrak{p}$, aber nicht durch $\mathfrak{p}^2$ theilbar, $P'(\tau)$ und $\Phi(\tau)$ sind durch $\mathfrak{p}$ nicht theilbar, und so giebt uns also die Formel (11) den Beweis des folgenden Satzes:

1. \gesperrt{Ist $\mathfrak{p}$ ein beliebiges Primideal, $p$ die durch $\mathfrak{p}$ theilbare natürliche Primzahl, und $\mathfrak{p}^e$ die höchste in $p$ aufgehende Potenz von $\mathfrak{p}$, so ist die Grundform $F'(\tau)$ allemal theilbar durch $\mathfrak{p}^{e-1}$; ist ferner der Exponent $e$ nicht theilbar durch $p$, so ist $F'(\tau)$ nicht theilbar durch $\mathfrak{p}^e$; ist aber $e$ theilbar durch $p$, so ist $F'(\tau)$ theilbar durch $\mathfrak{p}^e$ und vielleicht durch noch höhere Potenzen von $\mathfrak{p}$.}

Wenn $e$ grösser als 1 ist, so ist hiernach $F'(\tau)$ durch $\mathfrak{p}$ theilbar, und folglich ist $N_a[F'(\tau)]$ und also auch die Grundzahl $\Delta$ durch $p$ theilbar.

Wenn aber alle Primideale $\mathfrak{p}$ nur in erster Potenz in $p$ aufgehen, so ist $F'(\tau)$ relativ prim zu $p$. Zerlegt man also $F'(\tau)$ in seine Primfactoren, so kommt darunter keiner vor, dessen Norm eine Potenz von $p$ ist, folglich ist auch $N_a[F'(\tau)]$ und $\Delta$ durch $p$ nicht theilbar. Daraus folgt dann der Satz:

2. \gesperrt{Eine natürliche Primzahl $p$ ist dann und nur dann im Körper $\Omega$ durch das Quadrat eines Primfactors theilbar, wenn $p$ in der Grundzahl von $\Omega$ aufgeht.}

Aus diesen Betrachtungen können wir noch andere wichtige Schlüsse ziehen.

Wenn wir das System der linearen Gleichungen (1) in Bezug auf $\omega_1, \omega_2, \ldots, \omega_n$ auflösen, so erhalten wir Ausdrücke mit dem Nenner $U$, der, wie wir gesehen haben, eine Einheit ist. Substituirt man diese Ausdrücke in
\weq{12}{\omega = x_1\omega_1 + x_2\omega_2 + \cdots + x_n\omega_n,}
%CONT>>


% ===== scan leaf 602, printed 581 =====
% !PRINTED 581
\weberpage{581}

%<<CONT
worin die $x_1, x_2, \ldots, x_n$ ganze rationale Zahlen oder ganze rationale Functionale sind, so erhält man einen Ausdruck von der Form
\weq{13}{\omega = A_0 + A_1\tau + A_2\tau^2 + \cdots + A_{n-1}\tau^{n-1},}
worin die $A_0, A_1, \ldots, A_{n-1}$ gleichfalls \gesperrt{ganze rationale Functionale} bedeuten. Nach \S.\,145, 2. wird aber in dieser Form jede ganze Zahl und jedes ganze Functional des Körpers $\Omega$ dargestellt, und wir können daher den Satz aussprechen:

3. \gesperrt{Die Potenzen}
\[
1, \tau, \tau^2, \ldots, \tau^{n-1}
\]
\gesperrt{bilden eine Basis der ganzen Functionale des Körpers $\Omega$.}

Statt der Potenzen von $\tau$ können auch die Functionen
\[
F_0, F_1, F_2, \ldots, F_{n-1}
\]
als Elemente der Basis eingeführt werden, die durch die Gleichung
\weq{14}{\frac{F(t)}{t-\tau} = t^{n-1}F_0(\tau) + t^{n-2}F_1(\tau) + \cdots + F_{n-1}(\tau)}
definirt sind, die, wie schon im \S.\,4 des ersten Bandes gezeigt ist, wenn
\[
F(t) = t^n + a_1t^{n-1} + \cdots + a_n
\]
ist, den Ausdruck haben:
\weq{15}{\begin{aligned}
F_0(t) &= 1 \\
F_1(t) &= t + a_1 \\
F_2(t) &= t^2 + a_1 t + a_2 \\
&\cdots\cdots\cdots\cdots\cdots\cdots\cdots\cdots\cdots\cdots \\
F_{n-1}(t) &= t^{n-1} + a_1 t^{n-2} + a_2 t^{n-3} + \cdots + a_{n-1},
\end{aligned}}
so dass die Potenzen $1, t, t^2, \ldots, t^{n-1}$ ohne Nenner durch $F_0, F_1, \ldots, F_{n-1}$ dargestellt werden können, und dass jede ganze Zahl oder jedes ganze Functional $\omega$ auch den Ausdruck erhält:
\weq{16}{\omega = B_0 F_0(\tau) + B_1 F_1(\tau) + \cdots + B_{n-1} F_{n-1}(\tau),}
dessen Coëfficienten $B_0, B_1, \ldots, B_{n-1}$ \gesperrt{ganze rationale Functionale} sind.

Betrachten wir jetzt die ganze Function von $t$:
\[
S\,\frac{F(t)}{(t-\tau)\,F'(\tau)},
\]
% [sic: A lowered comma is printed after "so" at the line break. It is retained although the sentence does not require this comma.]
worin $S$ das Zeichen für die in \S.\,134 erklärte Spur ist, so, ergiebt sich, dass sie den Werth 1 erhält für $t=\tau_1,\tau_2,\ldots,\tau_n$
%CONT>>


% ===== scan leaf 603, printed 582 =====
% !PRINTED 582
\weberpage{582}
%<<CONT
und dass sie also, da sie nur vom $(n-1)^{\text{ten}}$ Grade ist, identisch $=1$ sein
muss. Daraus folgt nach (14):
\[
1 = t^{n-1} S\,\frac{F_0(\tau)}{F'(\tau)} + t^{n-2} S\,\frac{F_1(\tau)}{F'(\tau)} + \cdots + S\,\frac{F_{n-1}(\tau)}{F'(\tau)}
\]
oder
\weq{17}{\begin{aligned} S\,\frac{F_\nu(\tau)}{F'(\tau)} &= 0, \quad S\,\frac{F_{n-1}(\tau)}{F'(\tau)} = 1,\\ \nu &= 0, 1, \ldots, n-2. \end{aligned}}
Hieraus ergiebt sich mit Benutzung von (16):
\weq{18}{S\,\frac{\omega}{F'(\tau)} = B_{n-1},}
oder der Satz:

4. \gesperrt{Ist $\omega$ eine ganze Zahl oder ein ganzes Functional des K\"orpers $\Omega$, so ist die Spur von $\dfrac{\omega}{F'(\tau)}$ ein ganzes rationales Functional.}

\begin{center}
\rule{0.2\linewidth}{0.4pt}
\end{center}


% ===== scan leaf 604, printed 583 =====
\weberpage{583}% page number not printed on this page
\pdfbookmark[1]{Neunzehnter Abschnitt. Beziehung eines Körpers auf seine Theiler.}{bm181}
\pdfbookmark[2]{§ 159. Relativnormen}{bm182}
% !PRINTED none
\begin{center}
\gesperrt{Neunzehnter Abschnitt.}

\gesperrt{Beziehungen eines K\"orpers auf seine Theiler.}

\rule{0.2\linewidth}{0.4pt}

\S.\,159.

\gesperrt{Relativnormen.}
\end{center}

Wir wollen in diesem Abschnitte die Modificationen betrachten, die in den Definitionen und S\"atzen der Theorie der algebraischen Zahlen eintreten, wenn an Stelle des \gesperrt{absoluten Rationalit\"atsbereiches}, d. h. des K\"orpers der rationalen Zahlen, ein beliebiger algebraischer Zahlk\"orper gesetzt wird. Es werden sich dabei einige wichtige Erg\"anzungen auch f\"ur die allgemeine Theorie der Primfactoren ergeben, die in einer Reihe von \gesperrt{Dedekind} und \gesperrt{Hilbert} aufgestellter S\"atze ihren Ausdruck finden${}^{1}$).

Es sei also $R$ ein algebraischer Zahlk\"orper $m^{\text{ten}}$ Grades und
\weq{1}{f(\Theta) = \Theta^n + \alpha_1\Theta^{n-1} + \cdots + \alpha_n = 0}
eine in $R$ rationale und irreducible Gleichung $n^{\text{ten}}$ Grades. Der K\"orper $\Omega = R(\Theta)$ ist ein algebraischer K\"orper \gesperrt{\"uber} $R$, der in Bezug auf $R$ vom $n^{\text{ten}}$ Grade ist (Bd. I, \S.\,142). Im absoluten Rationalit\"atsbereiche ist $\Omega$ vom Grade $mn$, und wenn $\varrho$ eine Primitivzahl des K\"orpers $R$ ist, so ist
\weq{2}{\Theta^s \varrho^t, \qquad \begin{aligned} s &= 0, 1, \ldots, n-1\\ t &= 0, 1, \ldots, m-1 \end{aligned}}
eine Basis des K\"orpers $\Omega$; denn wegen der Irreducibilit\"at der
%CONT>>

${}^{1}$) \gesperrt{Dedekind}, \glqq Ueber die Discriminanten endlicher K\"orper\grqq{}, Abhandlungen der G\"ottinger Gesellschaft der Wissenschaften (1882). \glqq Zur Theorie der Ideale\grqq{}, Nachrichten der Gesellschaft der Wissenschaften zu G\"ottingen, 1894, Nr. 4. \gesperrt{Hilbert}, \glqq Grundz\"uge einer Theorie des \gesperrt{Galois'schen} Zahlk\"orpers\grqq{}. Ebend. 1894, Nr. 3.


% ===== scan leaf 605, printed 584 =====
% !PRINTED 584
\weberpage{584}
%<<CONT
Gleichung (1) kann zwischen den $mn$ Zahlen (2) keine lineare Relation mit rationalen Zahlencoëfficienten bestehen.

Jede Zahl $\omega$ des K\"orpers $\Omega$ kann als ganze Function $(n-1)^{\text{ten}}$ Grades von $\Theta$ mit Coëfficienten in $R$ dargestellt werden, und jede solche Zahl $\omega$ gen\"ugt einer in $R$ irreduciblen Gleichung h\"ochstens vom $n^{\text{ten}}$ Grade; $\omega$ ist dann und nur dann eine ganze Zahl, wenn diese Gleichung, nachdem der Coëfficient der h\"ochsten Potenz auf 1 gebracht ist, ganze Zahlen in $R$ zu Coëfficienten hat.

Den Inbegriff aller ganzen Zahlen in $\Omega$ bezeichnen wir, wie fr\"uher, mit $\mathfrak{o}$.

Wenn wir, ohne die Zahlen des K\"orpers $R$ zu ver\"andern, f\"ur $\Theta$ die s\"ammtlichen Wurzeln $\Theta_1, \Theta_2, \ldots, \Theta_n$ der Gleichung (1) nehmen, so erhalten wir $n$ K\"orper
\weq{3}{\Omega_1 = R(\Theta_1), \ \Omega_2 = R(\Theta_2), \ \ldots, \ \Omega_n = R(\Theta_n),}
\gesperrt{die in Bezug auf den K\"orper $R$ conjugirt sind.}

Sind $\omega_1, \omega_2, \ldots, \omega_n$ entsprechende Zahlen dieser K\"orper, so heisst das Product
\weq{4}{\mathfrak{N}_R(\omega) = \omega_1 \omega_2 \ldots \omega_n}
\gesperrt{die in Bezug auf $R$ genommene Partialnorm (Relativnorm)} der Zahl $\omega$.

Eine Zahl $\eta$ des K\"orpers $R$ hat eine in diesem K\"orper genommene gew\"ohnliche Norm $N_R(\eta)$, die eine rationale Zahl ist. Die Relativnorm $\mathfrak{N}_R(\omega)$ ist nun eine solche Zahl $\eta$, und wenn wir ihre Norm im K\"orper $R$ nehmen, so erhalten wir die \gesperrt{Totalnorm} der Zahl $\omega$ im K\"orper der rationalen Zahlen:
\weq{5}{N_\Omega(\omega) = N_R\,\mathfrak{N}_R(\omega).}

Dies ergiebt sich, wenn man $\omega$ durch die Potenzen (2) von $\Theta$ und $\varrho$ ausdr\"uckt, sodann f\"ur $\varrho$ die $m$ conjugirten Werthe, und zu jedem dieser $\varrho$ die nach $R$ conjugirten Werthe $\Theta$ setzt.

In demselben Sinne haben nun auch die \gesperrt{Functionale} des K\"orpers $\Omega$ und die durch diese repr\"asentirten \gesperrt{Ideale ihre Partialnormen}. Die Partialnorm eines Functionals in $\Omega$ ist ein Functional in $R$, und demnach ist die Partialnorm eines Ideals in $\Omega$ ein Ideal in $R$.

Von den conjugirten Idealen k\"onnen auch zwei oder mehrere einander gleich sein, und zwar kann dies auch dann eintreten, wenn die die Ideale repr\"asentirenden conjugirten Functionale nicht identisch sind, wenn sie nur associirt sind. Wenn man
%CONT>>


% ===== scan leaf 606, printed 585 =====
% !PRINTED 585
\weberpage{585}
%<<CONT
daher nur das Product aller von einander verschiedenen unter den conjugirten Idealen nimmt, so braucht dies keineswegs ein Ideal in $R$ zu sein.

Was hier von den Normen ausgef\"uhrt ist, gilt auch von den anderen symmetrischen Functionen, insbesondere von den Spuren. Wir nennen die Summe
\weq{6}{\mathfrak{S}_R(\omega) = \omega_1 + \omega_2 + \cdots + \omega_n}
die \gesperrt{Partialspur (oder Relativspur) von $\omega$ in Bezug auf} $R$. Sie ist eine Zahl oder ein Functional in $R$, und wir erhalten f\"ur die Totalspur
\weq{7}{S_\Omega(\omega) = S_R\,\mathfrak{S}_R(\omega),}
worin die Bedeutung der Zeichen unmittelbar verst\"andlich ist.

Eine besonders wichtige Beziehung ergiebt sich aber durch die Betrachtung der Grundideale.

Es sei $\tau$ eine Basisform von $\Omega$ (also eine Linearform von $mn$ Variablen, \S.\,146) und
\[
F(t) = N_\Omega(t - \tau)
\]
die irreducible Function $mn^{\text{ten}}$ Grades (im absoluten Rationalit\"atsbereiche), deren Wurzel $\tau$ ist. Dann haben wir im \S.\,158 das durch das Functional $F'(\tau)$ definirte Ideal als das \gesperrt{Grundideal} des K\"orpers $\Omega$ bezeichnet.

Ist nun $\sigma$ eine Basisform von $R$ und
\[
\psi(t) = N_R(t - \sigma),
\]
so ist durch $\psi'(\sigma)$ das Grundideal des K\"orpers $R$ definirt.

Aus der Function $F(t)$ spaltet sich aber im K\"orper $R$ ein irreducibler Factor $n^{\text{ten}}$ Grades ab,
\[
f(t) = \mathfrak{N}_R(t - \tau),
\]
der f\"ur $t = \tau$ verschwindet. Wir definiren durch $f'(\tau)$ das \gesperrt{Partial-Grundideal von $\Omega$ in Bezug auf} $R$, und beweisen nun den Satz:
\weq{8}{F'(\tau) = \varepsilon\, f'(\tau)\, \psi'(\sigma),}
worin $\varepsilon$ eine Einheit ist.

Um ihn zu beweisen, setzen wir nach der im \S.\,158 angewandten Bezeichnung
\weq{9}{\frac{F(t)}{t - \tau} = t^{mn-1}F_0(\tau) + t^{mn-2}F_1(\tau) + \cdots + F_{mn-1}(\tau),}
\weq{10}{\frac{f(t)}{t - \tau} = t^{n-1}f_0(\tau) + t^{n-2}f_1(\tau) + \cdots + f_{n-1}(\tau),}
\weq{11}{\frac{\psi(t)}{t - \sigma} = t^{m-1}\psi_0(\sigma) + t^{m-2}\psi_1(\sigma) + \cdots + \psi_{m-1}(\sigma).}


% ===== scan leaf 607, printed 586 =====
% !PRINTED 586
\weberpage{586}
Hierin sind $F_\nu$, $\psi_\nu$ ganze Functionale des absoluten Rationalitätsbereiches,
während $f_\nu$ ganze Functionale in $R$ sind. Die Ausdrücke für $F_\nu, f_\nu, \psi_\nu$
ergeben sich aus den Formeln \S.\,158, (15). Auf diese Functionale kann man dieselben
Schlüsse anwenden, die zu dem Theorem \S.\,158, 4. geführt haben, und es ergiebt sich:

\weq{12}{
\begin{aligned}
\mathfrak{S}_R \frac{f_\nu(\tau)}{f'(\tau)} &= 0, & \nu &= 0, 1, \ldots, n-2.\\
\mathfrak{S}_R \frac{f_{n-1}(\tau)}{f'(\tau)} &= 1.
\end{aligned}
}
Wenn man nun die Functionen $F_\nu(t)$, wenn sie von höherem als $(n-1)^{\mathrm{ten}}$
Grade sind, durch $f(t)$ dividirt und den Rest der Division nimmt; dann in diesem Reste
die Potenzen von $t$ [nach \S.\,158, (15), auf $f_\nu$ angewandt], durch die Functionen
$f_\nu(t)$ ausdrückt, so ergiebt sich die Darstellung

\weq{13}{
F_\nu(\tau) = \alpha_0 f_0 + \alpha_1 f_1 + \cdots + \alpha_{n-1} f_{n-1},
}
worin die $\alpha_0, \alpha_1, \ldots, \alpha_{n-1}$ \gesperrt{ganze Functionale} in $R$ sind. Diese
Darstellung gilt für jede Wurzel $\tau$ von $f(t) = 0$, d.\,h. für $\tau = \tau_1, \tau_2,
\ldots, \tau_n$. Hieraus ergiebt sich aber nach (12):
\[
\mathfrak{S}_R \left( \frac{F_\nu(\tau)}{f'(\tau)} \right) = \alpha_{n-1}.
\]
Dividirt man hier durch $\psi'(\sigma)$ und bildet die Spur $S_R$, so erhält man aus (7)
mit Anwendung des Satzes 4., \S.\,158:
\weq{14}{
S_\Omega \left( \frac{F_\nu(\tau)}{f'(\tau)\,\psi'(\sigma)} \right) = a_\nu,
}
worin $a_\nu$ ein ganzes rationales Functional bedeutet.

Nach (9) ist nun für jede Wurzel $\tau_i$ von $F(t) = 0$ mit Ausnahme von $\tau_i = \tau$
\[
\tau^{n\,m-1} F_0(\tau_i) + \tau^{n\,m-2} F_1(\tau_i) + \cdots + F_{n\,m-1}(\tau_i) = 0,
\]
und
\[
\tau^{n\,m-1} F_0(\tau) + \tau^{n\,m-2} F_1(\tau) + \cdots + F_{n\,m-1}(\tau) = F'(\tau),
\]
so dass sich, wenn man (14) mit $\tau^{n\,m-\nu-1}$ multiplicirt und die Summe in Bezug
auf $\nu$ nimmt,
\weq{15}{
\frac{F'(\tau)}{f'(\tau)\,\psi'(\sigma)} = a_0 \tau^{n\,m-1} + a_1 \tau^{n\,m-2} + \cdots
}
ergiebt. \gesperrt{Es ist daher $F'(\tau)$ durch $f'(\tau)\,\psi'(\sigma)$ theilbar.}


% ===== scan leaf 608, printed 587 =====
% !PRINTED 587
\weberpage{587}
Andererseits ist $f_\nu(\tau)\,\psi_\mu(\sigma)$ ein ganzes Functional des Körpers
$\Omega$, und folglich nach dem vorhin schon angewandten Satze \S.\,158, 4.:

\weq{16}{
\begin{aligned}
S_\Omega \left( \frac{f_\nu(\tau)\,\psi_\mu(\sigma)}{F'(\tau)} \right) &= b_{\mu,\nu} &
\mu &= 0, 1, \ldots, m-1\\
& & \nu &= 0, 1, \ldots, n-1
\end{aligned}
}
ein ganzes rationales Functional. Multiplicirt man diese Gleichung mit
$\tau^{n-\nu-1}\sigma^{m-\mu-1}$ und nimmt die Summe über alle $\mu$ und $\nu$, so folgt,
wie vorhin, da für zwei von $\sigma, \tau$ verschiedene Wurzeln $\sigma_k, \tau_i$ von
$\psi=0, f=0$
\[
\begin{aligned}
\sigma^{m-1}\psi_0(\sigma_k) + \sigma^{m-2}\psi_1(\sigma_k) + \cdots + \psi_{m-1}(\sigma_k) &= 0,\\
\tau^{n-1} f_0(\tau_i) + \tau^{n-2} f_1(\tau_i) + \cdots + f_{n-1}(\tau_i) &= 0,
\end{aligned}
\]
und
\[
\begin{aligned}
\sigma^{m-1}\psi_0(\sigma) + \sigma^{m-2}\psi_1(\sigma) + \cdots + \psi_{m-1}(\sigma) &= \psi'(\sigma),\\
\tau^{n-1} f_0(\tau) + \tau^{n-2} f_1(\tau) + \cdots + f_{n-1}(\tau) &= f'(\tau)
\end{aligned}
\]
ist, die Gleichung
\weq{17}{
\frac{f'(\tau)\,\psi'(\sigma)}{F'(\tau)} = \sum^{\mu,\nu} b_{\mu,\nu}\, \tau^{n-\nu-1}\sigma^{m-\mu-1}.
}

Also ist auch $f'(\tau)\,\psi'(\sigma)$ durch $F'(\tau)$ theilbar und unser Theorem (8)
daher bewiesen.

Bezeichnen wir mit $G_\Omega, G_R$ die Grundideale der Körper $\Omega, R$, mit
$\mathfrak{G}_R$ das relative Grundideal von $\Omega$ in Bezug auf $R$, so kann man dem
bewiesenen Theorem auch den Ausdruck geben:
\weq{18}{
G_\Omega = \mathfrak{G}_R\, G_R.
}
Es ist klar, wie sich dieses Theorem verallgemeinern lässt: Es sei
\[
\Omega_1, \Omega_2, \Omega_3, \ldots
\]
eine Reihe von Körpern, deren jeder die folgenden als Theiler enthält. Es sei ferner
$\mathfrak{G}_{i,k}$ das Grundideal von $\Omega_i$ in Bezug auf $\Omega_k$ ($k > i$),
dann ist
\weq{19}{
\mathfrak{G}_{i,k} = \mathfrak{G}_{i,i+1}\,\mathfrak{G}_{i+1,i+2} \cdots \mathfrak{G}_{k-1,k}.
}
Die tiefere Bedeutung dieser Grundideale wird sich in einem der nächsten Paragraphen noch
deutlicher herausstellen.

Wir wollen aber schon hier auf eine merkwürdige Erscheinung aufmerksam machen, dass es
nämlich vorkommen kann, dass $f'(\tau)$ eine Einheit ist, was nach dem
\gesperrt{Minkowski'schen} Satze nicht möglich ist, wenn $R$ der Körper der rationalen
Zahlen ist. In diesen Fällen, von denen die Theorie der com-
%CONT>>


% ===== scan leaf 609, printed 588 =====
% !PRINTED 588
\weberpage{588}
\pdfbookmark[2]{§ 160. Primideale im relativ normalen Körper}{bm183}
%<<CONT
plexen Multiplication der elliptischen Functionen Beispiele giebt, ist das
Partial-Grundideal $\mathfrak{G}_R$ gleich 1 zu setzen.

Wenn die $n$ Körper (3) mit einander identisch sind, so heisst $\Omega$ ein
\gesperrt{Normalkörper in Bezug auf} $R$ (relativ normal, wenn $R$ nicht der absolute
Rationalitätsbereich ist). Die Gesammtheit der $n$ Substitutionen $(\Theta, \Theta_n)$,
die den Uebergang der Zahlen des Körpers $\Omega$ zu den conjugirten Zahlen vermitteln,
bilden eine Gruppe vom Grade $n$, wie wir schon im \S.\,147 des ersten Bandes gesehen
haben, die nichts anderes ist, als die \gesperrt{Galois'sche} Gruppe der Gleichung (1) im
Rationalitätsbereiche $R$, die wir hier die \gesperrt{Gruppe des Körpers $\Omega$ (in Bezug auf $R$)}
nennen.

Wenn insbesondere diese Gruppe commutativ ist, so ist $\Omega$ ein
\gesperrt{Abel'scher Körper in Bezug auf} $R$ (ein relativ \gesperrt{Abel'scher Körper}). In diesem
Falle ist (1) eine \gesperrt{Abel'sche} Gleichung im Rationalitätsbereiche $R$.

\begin{center}
\S.\,160.\\[6pt]
\gesperrt{Primideale im relativ normalen Körper.}
\end{center}

Wir nehmen jetzt an, dass der Körper $\Omega$ normal sei in Bezug auf den beliebigen
Körper $R$, und untersuchen unter dieser Voraussetzung die Primideale des Körpers
$\Omega$.

Die Gruppe von $\Omega$ in Bezug auf $R$ sei $\Phi$, und die Substitutionen von $\Phi$
mögen mit $\varphi, \varphi_1, \ldots$ bezeichnet sein. Wir bezeichnen ferner mit
\weq{1}{
\omega \mid \varphi
}
\gesperrt{die Zahl, die durch die Substitution $\varphi$ aus $\omega$ hervorgeht}, so dass $\omega$
und $\omega \mid \varphi$ in $\Omega$ enthalten sind.

Ebenso wie die Zahlen $\omega$ gehen auch alle Functionale und damit zugleich alle
Ideale $\mathfrak{a}$ des Körpers $\Omega$ durch eine Substitution $\varphi$ in bestimmte
Functionale und Ideale $\mathfrak{a} \mid \varphi$ über, und wenn $\omega$ eine durch
$\mathfrak{a}$ theilbare ganze Zahl ist, so ist $\omega \mid \varphi$ durch
$\mathfrak{a} \mid \varphi$ theilbar.

Wenn $\mathfrak{a}$ irgend ein Ideal in $\Omega$ ist, so giebt es gewisse Substitutionen
$\psi$ in $\Phi$, die der Bedingung
\weq{2}{
\mathfrak{a} \mid \psi = \mathfrak{a}
}
genügen, darunter immer die identische Substitution, und diese Substitutionen $\psi$
bilden eine Gruppe $\Psi$, die ein Theiler von $\Phi$
%CONT>>


% ===== scan leaf 610, printed 589 =====
% !PRINTED 589
\weberpage{589}
%<<CONT
ist. \gesperrt{Wir nennen $\Psi$ die zum Ideal $\mathfrak{a}$ gehörige Gruppe, und wir sagen auch,
$\mathfrak{a}$ gehört zu der Gruppe $\Psi$.}

Da man in jeder Gleichung zwischen Zahlen und Functionalen in $\Omega$ alle
Substitutionen der Gruppe $\Phi$ ausführen darf, wobei die Grössen des Körpers $R$
ungeändert bleiben, ohne dass die Gleichung zu bestehen aufhört, so folgt, dass
Einheiten, ganze Functionale, associirte Functionale, durch einander theilbare
Functionale diese Eigenschaften nicht verlieren, wenn irgend eine der Substitutionen von
$\Phi$ ausgeführt wird. Wir heben den Satz hervor:

1. \gesperrt{Ist $\mathfrak{p}$ ein Primideal, so sind auch alle mit $\mathfrak{p}$ in
Bezug auf $R$ conjugirten Ideale $\mathfrak{p} \mid \varphi$ Primideale. Ist $\mathfrak{a}$
durch irgend eine Potenz von $\mathfrak{p}$ theilbar so ist $\mathfrak{a} \mid \varphi$
durch die gleiche Potenz von $\mathfrak{p} \mid \varphi$, theilbar.} % [sic] source omits the comma after "theilbar" and prints a stray comma after $\mathfrak p\mid\varphi$

Ist $\mathfrak{p}$ ein Primideal in $\Omega$, so giebt es eine und nur eine durch
$\mathfrak{p}$ theilbare natürliche Primzahl $p$.

Zerlegt man diese in ihre Primfactoren in $R$, so muss einer dieser Primfactoren, den
wir mit $\mathfrak{P}$ bezeichnen, durch $\mathfrak{p}$ theilbar sein.

Ist dann $\mathfrak{A}$ irgend ein Ideal in $R$ und zugleich durch $\mathfrak{p}$
theilbar, so ist der grösste gemeinschaftliche Theiler von $\mathfrak{P}$ und
$\mathfrak{A}$, der nach \S.\,139 auch in $R$ enthalten ist, durch $\mathfrak{p}$
theilbar, und ist also gewiss keine Einheit. $\mathfrak{P}$ und $\mathfrak{A}$ sind also
nicht relativ prim, und $\mathfrak{A}$ ist folglich durch das Primideal $\mathfrak{P}$
theilbar. Ist $\mathfrak{A}$ auch ein Primideal, so müssen $\mathfrak{A}$ und
$\mathfrak{P}$ identisch sein. Damit ist, mit Rücksicht auf 1., bewiesen:

2. \gesperrt{Ist $\mathfrak{p}$ ein Primideal in $\Omega$, so giebt es ein und nur ein
Primideal $\mathfrak{P}$ in $R$, das durch $\mathfrak{p}$ theilbar ist, und jedes durch
$\mathfrak{p}$ theilbare Ideal in $R$ ist zugleich durch $\mathfrak{P}$ theilbar. Das
Primideal $\mathfrak{P}$ ist zugleich durch die sämmtlichen zu $\mathfrak{p}$
conjugirten Primideale $\mathfrak{p} \mid \varphi$ theilbar.}

Die Partialnorm $\mathfrak{N}_R(\mathfrak{p})$ ist durch $\mathfrak{p}$, und folglich als
Ideal in $R$ durch $\mathfrak{P}$ theilbar. Sie kann aber auch nach dem Satze 1. durch
kein von $\mathfrak{P}$ verschiedenes Ideal in $R$ theilbar sein. Denn ist $\mathfrak{P}'$
irgend ein Primtheiler von $\mathfrak{N}_R(\mathfrak{p})$, so ist $\mathfrak{P}'$
wenigstens durch einen der Primfactoren $\mathfrak{p} \mid \varphi$ theilbar und mithin
gleich $\mathfrak{P}$. Folglich
%CONT>>


% ===== scan leaf 611, printed 590 =====
% !PRINTED 590
\weberpage{590}
%<<CONT
ist die Partialnorm von $\mathfrak{p}$ eine Potenz von $\mathfrak{P}$. Wir setzen
\weq{3}{\mathfrak{N}_R(\mathfrak{p}) = \mathfrak{P}^f,}
und nennen $f$ \gesperrt{den Grad von $\mathfrak{p}$ in Bezug auf den Körper $R$.}

Die Norm von $\mathfrak{P}$ im Körper $R$ ist gleichfalls eine Potenz von $p$, die wir mit $P$ bezeichnen wollen, also
\weq{4}{N_R(\mathfrak{P}) = P.}

Es ergiebt sich dann für die Totalnorm von $\mathfrak{P}$ im Körper $\Omega$ nach \S.\,159, (5) % [sic] source names $\mathfrak P$ here, while (5) and the norm identity require $\mathfrak p$
\weq{5}{N_\Omega(\mathfrak{p}) = P^f.}

Ist $\mathfrak{p}^g$ die höchste in $\mathfrak{P}$ aufgehende Potenz von $\mathfrak{p}$, so ist $\mathfrak{P}$ nach 1. auch durch die $g^{\text{te}}$ Potenz aller mit $\mathfrak{p}$ conjugirten Primfactoren und durch keine höhere Potenz theilbar. Wenn nun $\mathfrak{p}_1, \mathfrak{p}_2, \ldots, \mathfrak{p}_e$ \gesperrt{die von einander verschiedenen} unter den mit $\mathfrak{p}$ conjugirten Primidealen sind, so ist
\weq{6}{\mathfrak{P} = (\mathfrak{p}_1\mathfrak{p}_2\ldots\mathfrak{p}_e)^g,}
und wenn man die Partialnorm auf beiden Seiten nimmt, und wie früher mit $n$ den Grad des Körpers $\Omega$ in Bezug auf $R$ bezeichnet, so dass die Partialnorm von $\mathfrak{P}$ gleich $\mathfrak{P}^n$ wird, so folgt aus (3):
\weq{7}{n = efg.}
Es sind also die natürlichen Zahlen $e, f, g$ Theiler von $n$.

Wenn $\mathfrak{p}$ durch $\varphi_1$ in $\mathfrak{p}_1$ übergeht, wenn also
\[
\mathfrak{p}_1 = \mathfrak{p} \,|\, \varphi_1,
\]
so ist $\mathfrak{p} = \mathfrak{p}_1 \,|\, \varphi_1^{-1}$, und wenn $\Psi$ die zu $\mathfrak{p}$ gehörige Gruppe ist, so ist auch für jede in $\Psi$ enthaltene Substitution $\psi$
\[
\mathfrak{p}_1 = \mathfrak{p} \,|\, \psi \varphi_1, \quad \mathfrak{p}_1 = \mathfrak{p}_1 \,|\, \varphi_1^{-1} \psi \varphi_1.
\]

Ist umgekehrt $\mathfrak{p}_1 \,|\, \varphi = \mathfrak{p}_1$, so folgt $\mathfrak{p} \,|\, \varphi_1 \varphi \varphi_1^{-1} = \mathfrak{p}$, d. h. $\varphi_1 \varphi \varphi_1^{-1} = \psi$ oder $\varphi = \varphi_1^{-1} \psi \varphi_1$.

Es gehört daher $\mathfrak{p}_1$ zur Gruppe $\varphi_1^{-1} \Psi \varphi_1$, und wenn
\[
\mathfrak{p}_1 = \mathfrak{p} \,|\, \varphi_1, \quad \mathfrak{p}_2 = \mathfrak{p} \,|\, \varphi_2, \ldots, \mathfrak{p}_e = \mathfrak{p} \,|\, \varphi_e
\]
ist, so ist
\weq{8}{\Phi = \Psi\varphi_1 + \Psi\varphi_2 + \cdots + \Psi\varphi_e.}
$\Psi$ ist also ein Theiler von $\Phi$ vom Index $e$ und vom Grade $gf$.


% ===== scan leaf 612, printed 591 =====
% !PRINTED 591
\weberpage{591}
\pdfbookmark[2]{§ 161. Primitivwurzeln der Primideale}{bm184}
\begin{center}
\S.\,161.\\[4pt]
\gesperrt{Primitivwurzeln der Primideale.}
\end{center}

Aus den Formeln (4), (5) des vorigen Paragraphen ergiebt sich, mit Rücksicht auf \S.\,150, für jede ganze Zahl $\eta$ des Körpers $R$
\[
\eta^P \equiv \eta \pmod{\mathfrak{P}},
\]
also auch
\weq{1}{\eta^P \equiv \eta \pmod{\mathfrak{p}},}
und für jede Zahl $\omega$ in $\mathfrak{o}$
\weq{2}{\omega^{P^f} \equiv \omega \pmod{\mathfrak{p}}.}
Die Anzahl der nach $\mathfrak{P}$ incongruenten ganzen Zahlen in $R$ ist $P$, und die Anzahl der nach $\mathfrak{p}$ incongruenten Zahlen in $\Omega$ ist $P^f$.

Wir haben schon früher (\S.\,150) gezeigt, dass es zu jedem Primideal $\mathfrak{p}$ \gesperrt{Primitivwurzeln} $\gamma$ giebt, d. h. Zahlen in $\Omega$, die der Congruenz
\weq{3}{\gamma^{P^f-1} \equiv 1 \pmod{\mathfrak{p}}}
genügen, und zugleich die Eigenschaft haben, dass keine niedrigere Potenz mit positiven Exponenten der Einheit congruent wird. Dann ist jede durch $\mathfrak{p}$ nicht theilbare ganze Zahl in $\Omega$ mit einer und nur mit einer der Potenzen von $\gamma$
\weq{4}{1, \gamma, \gamma^2, \ldots, \gamma^{P^f-2}}
nach dem Modul $\mathfrak{p}$ congruent.

Jede ganze Zahl $\Theta$ in $\Omega$ genügt einer Gleichung höchstens vom $n^{\text{ten}}$ Grade, deren Coëfficienten ganze Zahlen in $R$ sind.

Diese Gleichung ist zugleich eine Congruenz nach dem Modul $\mathfrak{p}$, und unter allen solchen Congruenzen wird es eine von \gesperrt{möglichst niedrigem Grade}
\[
F(\Theta) \equiv 0 \pmod{\mathfrak{p}}
\]
geben, in der wir überdies den Coëfficienten der höchsten Potenz von $\Theta$ gleich 1 annehmen können. Denn ist der höchste (durch $\mathfrak{p}$ und folglich durch $\mathfrak{P}$ untheilbare) Coëfficient $\eta_0$, so können wir immer der Congruenz
\[
\eta_0 \eta \equiv 1 \pmod{\mathfrak{P}}
\]
durch ein ganzzahliges $\eta$ in $R$ genügen, und haben dann nur die Function $F$ mit diesem $\eta$ zu multipliciren. Wir können
%CONT>>


% ===== scan leaf 613, printed 592 =====
% !PRINTED 592
\weberpage{592}
%<<CONT
also, wenn $t$ eine Variable ist, die Function $F$ in der Form annehmen:
\[
F(t) = t^\nu + \alpha_1 t^{\nu-1} + \alpha_2 t^{\nu-2} + \cdots + \alpha_\nu,
\]
deren Coëfficienten $\alpha_1, \alpha_2, \ldots, \alpha_\nu$ ganze Zahlen in $R$ sind. Durch Division können wir dann für jede andere ganze Function $F_1(t)$ den Quotienten $Q(t)$ und den Rest $\varphi(t)$ so bestimmen, dass
\weq{5}{F_1(t) = Q(t) F(t) + \varphi(t),}
worin
\weq{6}{\varphi(t) = \varrho_0 + \varrho_1 t + \cdots + \varrho_{\nu-1} t^{\nu-1}}
eine durch $F_1(t)$ eindeutig bestimmte Function $(\nu-1)^{\text{ten}}$ Grades ist, mit ganzzahligen Coëfficienten $\varrho$.

Setzen wir für $\Theta$ die Primitivwurzel $\gamma$ und für $F_1(t)$ eine Potenz von $t$, so ergiebt sich aus (4), (5) und (6), dass jede durch $\mathfrak{p}$ nicht theilbare Zahl $\omega$ einer Congruenz
\weq{7}{\omega \equiv \varrho_0 + \varrho_1 \gamma + \cdots + \varrho_{\nu-1} \gamma^{\nu-1} \pmod{\mathfrak{p}}}
genügt, und da die $\varrho$ auch $= 0$ sein können, so gilt dies auch noch für ein durch $\mathfrak{p}$ theilbares $\omega$.

Die Coëfficienten $\varrho$ in dem Ausdrucke (7) sind durch $\omega$ selbst nach dem Modul $\mathfrak{P}$ völlig bestimmt, da nach unserer Voraussetzung $\gamma$ keiner Congruenz in $R$ von niedrigerem als dem $\nu^{\text{ten}}$ Grade genügen soll, deren Coëfficienten nicht alle durch $\mathfrak{P}$ theilbar sind. Folglich giebt es, da jeder der Coëfficienten in (7) nur $P$ incongruente Werthe haben kann, $P^\nu$ und nicht mehr incongruente Zahlen $\omega$, und daraus folgt, dass $\nu = f$ sein muss.

Aus einer Congruenz $F(\gamma) \equiv 0 \pmod{\mathfrak{p}}$ folgt nun aber durch wiederholtes Potenziren mit Rücksicht auf (1):
\[
F(\gamma^P) \equiv 0, \ F(\gamma^{P^2}) \equiv 0, \ldots \pmod{\mathfrak{p}},
\]
und wir haben also den Satz bewiesen (\S.\,150, 2.):

\gesperrt{Eine Primitivwurzel $\gamma$ von $\mathfrak{p}$ genügt nach dem Modul $\mathfrak{p}$ einer Congruenz in $R$ vom $f^{\text{ten}}$ Grade, deren sämmtliche Wurzeln}
\[
\gamma, \gamma^P, \gamma^{P^2}, \ldots, \gamma^{P^f-1}
\]
\gesperrt{sind.}
% [sic: The source prints $\gamma^{P^f-1}$ as the last root, with -1 on the P-exponent level rather than inside the exponent of P. The P-power orbit requires $\gamma^{P^{f-1}}$; the printed expression is retained.]

Diesem Satze kann man auch den Ausdruck geben:

\gesperrt{Das Product}
% [sic: The last factor repeats the source's $\gamma^{P^f-1}$; the P-power-orbit factor would be $\gamma^{P^{f-1}}$.]
\weq{8}{(t-\gamma)(t-\gamma^P) \ldots (t-\gamma^{P^f-1})}
\gesperrt{ist nach dem Modul $\mathfrak{p}$ mit einer ganzen Function $F(t)$ in $R$ congruent.}


% ===== scan leaf 614, printed 593 =====
% !PRINTED 593
\weberpage{593}
\pdfbookmark[2]{§ 162. Das Partial-Grundideal}{bm185}
\begin{center}
\S.\,162.\\[4pt]
\gesperrt{Das Partial-Grundideal.}
\end{center}

Es möge jetzt
\weq{1}{\tau = t_1\omega_1 + t_2\omega_2 + \cdots + t_{mn}\omega_{mn}}
eine Basisform von $\mathfrak{o}$ sein. Dann wird es gewisse Substitutionen $\chi$ in $\Phi$ geben, und darunter immer die identische, die der Congruenz
\weq{2}{\tau \,|\, \chi \equiv \tau \pmod{\mathfrak{p}}}
genügen. Alle diese Substitutionen bilden eine in $\Phi$ enthaltene Gruppe $X$, die auch ein Theiler der Gruppe $\Psi$ ist, zu der $\mathfrak{p}$ gehört. Denn aus $\tau$ erhält man (nach \S.\,146) eine Basisform von $\mathfrak{p}$, wenn man für die Variablen $t$ gewisse lineare Functionen neuer Variablen mit ganzen rationalen Coëfficienten setzt; wenn also $\pi$ eine solche Basisform von $\mathfrak{p}$ ist, so folgt aus (2), dass $\pi \,|\, \chi$ durch $\pi$ theilbar und daher auch $\mathfrak{p} \,|\, \chi = \mathfrak{p}$ ist.

Nun genügt $\tau$ einer Gleichung $n^{\text{ten}}$ Grades $f(t) = 0$, deren Coëfficienten ganze Functionen in $R$ mit den Variablen $t_1, t_2 \ldots$ sind, und es ist, wenn $t$ eine neue Variable bedeutet, und $\tau_1, \tau_2, \ldots, \tau_n$ die in Bezug auf $R$ zu $\tau$ conjugirten Formen sind,
\weq{3}{f(t, t_1, t_2, \ldots) = f(t) = (t-\tau_1)(t-\tau_2) \ldots (t-\tau_n).}
Hieraus ergiebt sich nun, wenn wir wiederholt in die Potenz $p$ erheben, und auf die Zahlencoëfficienten von $f$ die Formel \S.\,161, (1) anwenden:
\weq{4}{f(t^P, t_1^P, t_2^P, \ldots) \equiv [f(t)]^P \equiv (t^P-\tau_1')(t^P-\tau_2') \ldots (t^P-\tau_n') \pmod{\mathfrak{p}},}
wenn $\tau_1', \tau_2', \ldots, \tau_n'$ dadurch aus $\tau_1, \tau_2, \ldots, \tau_n$ abgeleitet sind, dass die Variablen $t_1, t_2, \ldots$ durch $t_1^P, t_2^P, \ldots$ ersetzt sind.

Setzen wir nun $t = \tau$ in (4), so verschwindet $f(t)$ und es folgt, dass einer der Factoren des letzten Productes durch $\mathfrak{p}$ theilbar sein muss. Dies aber kann so ausgedrückt werden, dass es in $\Phi$ eine Substitution $\psi_0$ giebt, die der Bedingung
\weq{5}{\tau^P \equiv \tau' \,|\, \psi_0 \pmod{\mathfrak{p}}}
genügt.

Aus $\tau$ entstehen alle ganzen Zahlen in $\Omega$, wenn man für die Variablen ganze rationale Zahlen setzt. Wendet man dann
%CONT>>


% ===== scan leaf 615, printed 594 =====
% !PRINTED 594
\weberpage{594}
%<<CONT
noch den \gesperrt{Fermat'schen} Lehrsatz für rationale Zahlen an, so erkennt man, dass die Congruenzen (2) und (5) \gesperrt{gleichbedeutend} sind mit den für \gesperrt{jede Zahl} $\omega$ in $\mathfrak{o}$ gültigen Formeln

\weq{6}{\omega\mkern1mu|\mkern1mu\chi \equiv \omega,\quad \omega^P \equiv \omega\mkern1mu|\mkern1mu\psi_0\ (\mathrm{mod}\ \mathfrak{p}).}
Aus der zweiten Congruenz (6) folgt, dass, wenn $\omega$ durch $\mathfrak{p}$ theilbar ist, immer auch $\omega\mkern1mu|\mkern1mu\psi_0$ durch $\mathfrak{p}$ theilbar sein muss. Folglich ist $\mathfrak{p}$ durch $\mathfrak{p}\mkern1mu|\mkern1mu\psi_0$ theilbar, und als Primideal $= \mathfrak{p}\mkern1mu|\mkern1mu\psi_0$. Daraus folgt, dass $\psi_0$ in der Gruppe $\Psi$ enthalten ist.

Wir verstehen jetzt unter $\gamma$ eine Primitivwurzel von $\mathfrak{p}$ und wenden die am Schlusse des \S.\,161 bewiesene Formel an:
% [sic: The source repeats $\gamma^{P^f-1}$ in the last factor, as on printed page 592; the P-power orbit instead requires $\gamma^{P^{f-1}}$. The printed expression is retained.]
\weq{7}{(t-\gamma)(t-\gamma^P)\ldots(t-\gamma^{P^f-1}) \equiv F(t)\ (\mathrm{mod}\ \mathfrak{p}),}
in der $F(t)$ eine ganze Function von $t$ in $R$ ist.

Daraus folgt
\[ F(\gamma) \equiv 0\ (\mathrm{mod}\ \mathfrak{p}), \]
und wenn nun $\psi$ irgend eine Substitution aus $\Psi$ ist:
\[ F(\gamma\mkern1mu|\mkern1mu\psi) \equiv 0\ (\mathrm{mod}\ \mathfrak{p}). \]

Daraus ergiebt sich aber, dass $\gamma\mkern1mu|\mkern1mu\psi$ mit einer der in (7) vorkommenden Potenzen von $\gamma$ nach $\mathfrak{p}$ congruent sein muss, also etwa
\weq{8}{\gamma^{P^\nu} \equiv \gamma\mkern1mu|\mkern1mu\psi\ (\mathrm{mod}\ \mathfrak{p}).}
Nun ist jede durch $\mathfrak{p}$ nicht theilbare Zahl $\omega$ in $\mathfrak{o}$ mit einer Potenz von $\gamma$ congruent, und wenn man also (8) zu dieser Potenz erhebt, so folgt
\weq{9}{\omega^{P^\nu} \equiv \omega\mkern1mu|\mkern1mu\psi\ (\mathrm{mod}\ \mathfrak{p}),}
und diese Congruenz gilt offenbar auch noch, wenn $\omega$ durch $\mathfrak{p}$ theilbar ist, also für \gesperrt{alle Zahlen} in $\mathfrak{o}$.

Wenn wir nun die zweite der Congruenzen (6) $\nu$ mal nach einander anwenden, so folgt
\weq{10}{\omega^{P^\nu} \equiv \omega\mkern1mu|\mkern1mu\psi_0^\nu\ (\mathrm{mod}\ \mathfrak{p}),}
also
\weq{11}{\omega\mkern1mu|\mkern1mu\psi \equiv \omega\mkern1mu|\mkern1mu\psi_0^\nu\ (\mathrm{mod}\ \mathfrak{p}),}
und wenn wir $\omega$ durch $\omega\mkern1mu|\mkern1mu\psi^{-1}$ ersetzen oder auf (11) die Substitution $\psi^{-1}$ anwenden:
\weq{12}{\omega \equiv \omega\mkern1mu|\mkern1mu\psi^{-1}\psi_0^\nu \equiv \omega\mkern1mu|\mkern1mu\psi_0^\nu\psi^{-1}\ (\mathrm{mod}\ \mathfrak{p}).}
Es sind also sowohl $\psi^{-1}\psi_0^\nu$ als auch $\psi_0^\nu\psi^{-1}$ in $X$ enthalten,
%CONT>>


% ===== scan leaf 616, printed 595 =====
% !PRINTED 595
\weberpage{595}
%<<CONT
woraus sich ergiebt, dass \gesperrt{jede Substitution $\psi$ aus $\Psi$ in einem der Systeme} (Nebengruppen)
\[ X,\ X\psi_0,\ X\psi_0^2,\ \ldots \]
\gesperrt{enthalten ist.}

Wenn $\mathfrak{p}$ vom Grade $f$ (in Bezug auf $R$) ist, so ist
\[ \omega^{P^f} \equiv \omega\ (\mathrm{mod}\ \mathfrak{p}), \]
und $P^f$ ist die niedrigste Potenz von $P$ mit positiven Exponenten, die dieser Congruenz für alle $\omega$ genügt.

Setzt man also $\nu = f$ in (10), so folgt, dass $\psi_0^f$ in $X$ enthalten ist, dass aber keine Potenz von $\psi_0$ mit niedrigerem positiven Exponenten diese Eigenschaft hat. Die Nebengruppen
\[ X,\ X\psi_0,\ X\psi_0^2,\ \ldots,\ X\psi_0^{f-1} \]
sind also alle von einander verschieden, und es ergiebt sich:
\weq{13}{\Psi = X + X\psi_0 + X\psi_0^2 + \cdots + X\psi_0^{f-1}.}

Nach (12) ist aber die Gesammtheit der Substitutionen $\psi_0^\nu X$ mit $X\psi_0^\nu$ identisch (für jedes $\nu$), also
\weq{14}{\psi_0^\nu X = X\psi_0^\nu,}
woraus folgt, dass \gesperrt{$X$ ein Normaltheiler von $\Psi$ ist.}

\gesperrt{Da, wie wir oben gesehen haben, $\Psi$ vom Grade $gf$ ist, so ist $X$ vom Grade $g$.}

In der Gruppe $X$ giebt es nun solche Substitutionen $\chi_1$, darunter immer die identische, für die die Congruenz
\weq{15}{\tau\mkern1mu|\mkern1mu\chi_1 \equiv \tau\ (\mathrm{mod}\ \mathfrak{p}^2)}
erfüllt ist, und diese bilden eine Gruppe $X^{(1)}$, deren Grad mit $g_1$ bezeichnet sein mag.

\gesperrt{Es ist zu beweisen, dass $g_1$ eine Potenz von $p$ ist.} Bezeichnen wir nämlich mit $\pi$ ein durch $\mathfrak{p}$ theilbares Primfunctional, mit $\alpha$ ein ganzes Functional, so folgt aus (15):
\weq{16}{\tau\mkern1mu|\mkern1mu\chi_1 = \tau + \pi^2\alpha,}
und da ebenso nach (15)
\[ \alpha\mkern1mu|\mkern1mu\chi_1 \equiv \alpha\ (\mathrm{mod}\ \mathfrak{p}^2),\quad (\pi\mkern1mu|\mkern1mu\chi_1)^2 \equiv \pi^2\ (\mathrm{mod}\ \mathfrak{p}^3) \]
ist, so ergiebt sich aus (16)
\[ \tau\mkern1mu|\mkern1mu\chi_1^2 \equiv \tau + 2\pi^2\alpha\ (\mathrm{mod}\ \mathfrak{p}^3), \]
und allgemein für jeden positiven Exponenten $h$
\[ \tau\mkern1mu|\mkern1mu\chi_1^h \equiv \tau + h\pi^2\alpha\ (\mathrm{mod}\ \mathfrak{p}^3). \]


% ===== scan leaf 617, printed 596 =====
% !PRINTED 596
\weberpage{596}
Setzt man $h = p$, so ergiebt sich hieraus
\[ \tau\mkern1mu|\mkern1mu\chi_1^p \equiv \tau\ (\mathrm{mod}\ \mathfrak{p}^3), \]
und auf die gleiche Weise schliesst man hieraus für jeden Exponenten $\nu$
\weq{17}{\tau\mkern1mu|\mkern1mu\chi_1^{p^\nu} \equiv \tau\ (\mathrm{mod}\ \mathfrak{p}^{\nu+2}).}
Um diese Formel allgemein zu beweisen, wendet man die vollständige Induction an. Man nimmt (17) als bewiesen an und setzt
\[ \tau\mkern1mu|\mkern1mu\chi_1^{p^\nu} = \tau + \pi^{\nu+2}\alpha; \]
daraus durch wiederholte Anwendung
\[
\begin{aligned}
\tau\mkern1mu|\mkern1mu\chi_1^{2p^\nu} &\equiv \tau + 2\pi^{\nu+2}\alpha \\
\tau\mkern1mu|\mkern1mu\chi_1^{3p^\nu} &\equiv \tau + 3\pi^{\nu+2}\alpha\ (\mathrm{mod}\ \mathfrak{p}^{\nu+3}) \\
&\cdots\cdots\cdots\cdots \\
\tau\mkern1mu|\mkern1mu\chi_1^{p^{\nu+1}} &\equiv \tau + p\pi^{\nu+2}\alpha,
\end{aligned}
\]
also die Formel (17) für $\nu + 1$.

Nun kann man $\nu$ immer so gross annehmen, dass von den Differenzen
\[ \tau - \tau_1,\ \tau - \tau_2,\ \tau - \tau_3,\ \ldots, \]
wenn $\tau$ von $\tau_1, \tau_2, \tau_3, \ldots$ verschieden ist, keine durch $\mathfrak{p}^{\nu-2}$ theilbar ist, und dann folgt aus (17), dass $\chi_1^{p^\nu}$ die identische Substitution sein muss. \gesperrt{Es ist hiernach der Grad eines jeden Elementes $\chi_1$ eine Potenz von $p$, und folglich kann nach dem Cauchy-Sylow'schen Satze (\S.\,29, I.) im Grade von $X^{(1)}$ keine von $p$ verschiedene Primzahl aufgehen, wie bewiesen werden sollte.}

Wenn also $g$ nicht durch $p$ theilbar ist, dann ist sicher $X^{(1)}$ die Einheitsgruppe.

Man kann auf diese Weise fortfahren und eine Gruppe $X^{(2)}$ vom Grade $g_2$ bilden aus der Congruenz
\[ \tau\mkern1mu|\mkern1mu\chi_2 \equiv \tau\ (\mathrm{mod}\ \mathfrak{p}^3), \]
$X^{(2)}$ ist ein Theiler von $X^{(1)}$, also $g_2$ gleichfalls eine Potenz von $p$, und so muss man schliesslich zur Einheitsgruppe gelangen. Unter den auf einander folgenden Gruppen $X^{(1)}, X^{(2)}, \ldots$ können unter Umständen auch mehrere einander gleich sein\footnote{\gesperrt{Hilbert} nennt $\Psi$ die Zerlegungsgruppe, $X$ die Trägheitsgruppe, $X^{(1)}$ die Verzweigungsgruppe des Ideals $\mathfrak{p}$. In der oben erwähnten Abhandlung ist noch bewiesen, dass $g_1$ die höchste in $g$ aufgehende Potenz von $p$ ist, dass $X^{(1)}$ ein Normaltheiler von $X$ ist, und dass die ganze Gruppe $\Psi$ metacyklisch ist.}.


% ===== scan leaf 618, printed 597 =====
% !PRINTED 597
\weberpage{597}
\pdfbookmark[2]{§ 163. Die Ideale in den Theilern des Körpers Ω}{bm186}
Hiernach lässt sich die höchste Potenz des Primideals $\mathfrak{p}$ angeben, die in dem Partialgrundideal $\mathfrak{G}_R$ aufgeht. Nach \S.\,159 nämlich ist dies Ideal definirt durch das Product
\[ f'(\tau) = (\tau - \tau_1)(\tau - \tau_2)\ldots(\tau - \tau_{n-1}), \]
worin $\tau_1, \tau_2, \ldots, \tau_{n-1}$ die von $\tau$ verschiedenen unter den mit $\tau$ conjugirten Functionalen sind. Einer der Factoren dieses Productes, $\tau - \tau\mkern1mu|\mkern1mu\varphi$, ist dann und nur dann durch $\mathfrak{p}$ theilbar, wenn $\varphi$ zu $X$ gehört. Da nun die identische Substitution ausgeschlossen ist, so folgt, dass $\mathfrak{G}$ den Factor $\mathfrak{p}$ mindestens $g-1$ mal enthält.

Es ist aber $\tau - \tau\mkern1mu|\mkern1mu\chi$ dann und nur dann durch $\mathfrak{p}^2$ theilbar, wenn $\chi$ zu $X^{(1)}$ gehört, und folglich ist $\mathfrak{G}$ durch $\mathfrak{p}^{g-1+g_1-1}$ theilbar. Wir können so fortfahren und finden die genaue Potenz von $\mathfrak{p}$, die in $\mathfrak{G}$ enthalten ist. Sie hat den Exponenten
\[ (g-1) + (g_1-1) + (g_2-1) + \cdots \]
Jede der Zahlen $g, g_1, g_2, \ldots$ ist Theiler der vorangehenden, und $g_1, g_2, \ldots$ sind Potenzen von $p$, und die Reihe setzt sich so lange fort, bis eine der Zahlen $g = 1$ geworden ist. Die Zusammensetzung des Grundideals $\mathfrak{G}$ ist also vollkommen durch die Zerlegung der Gruppe $X$ bestimmt.

Die Zahlen $g, g_1, g_2, \ldots$ sind überdies für alle conjugirten Primideale $\mathfrak{p}_1, \mathfrak{p}_2, \ldots, \mathfrak{p}_e$ dieselben, und es folgt also
\[ \mathfrak{G} = \Pi(\mathfrak{p}_1\mathfrak{p}_2\ldots\mathfrak{p}_e)^{g-1+g_1-1+g_2-2\cdots}, \] % [sic] exponent pattern breaks from g_2-1 to g_2-2 as printed
wo sich das Productzeichen auf die Primfactoren aller Primzahlen $p$ des Körpers $\Omega$ bezieht. Es brauchen dabei aber selbstverständlich nur die in endlicher Anzahl vorhandenen Primzahlen $p$ berücksichtigt zu werden, für die $g > 1$ ist.

Die Partialnorm des Ideals $\mathfrak{G}$
\[ \mathfrak{D} = \mathfrak{N}_R(\mathfrak{G}) = \Pi\, \mathfrak{P}^{ef(g-1+g_1-1+g_2-1+\cdots)} \]
heisst die Partial-Discriminante des Körpers $\Omega$ in Bezug auf den Körper $R$. Sie ist ein im Körper $R$ gelegenes Ideal, die für den Fall, dass $R$ der Körper der rationalen Zahlen ist, von einem Einheitsfactor abgesehen, in die Grundzahl des Körpers $\Omega$ übergeht (\S.\,158).

\begin{center}
\S.\,163.

\gesperrt{Die Ideale in den Theilern des Körpers $\Omega$.}
\end{center}

Wenn die Gruppe $\Phi$ des Körpers $\Omega$ einen Theiler $\Phi'$ hat, so gehört zu $\Phi'$ ein Körper $\Omega'$, der ein Theiler von $\Omega$ ist und
%CONT>>


% ===== scan leaf 619, printed 598 =====
% !PRINTED 598
\weberpage{598}
%<<CONT
seinerseits $R$ als Theiler enthält. Die Zahlen von $\Omega'$ bleiben
durch die Substitutionen $\Phi'$ ungeändert, und $\Phi'$ ist die Gruppe des
Körpers $\Omega$ in Bezug auf $\Omega'$ (vergl. Bd. I, \S.\,155). Wir bezeichnen
den Grad von $\Phi'$, der ein Theiler von $n$ ist, mit $n'$ und setzen

\weq{1}{n = m'\,n'.}

Es ist dann $n'$ der Grad von $\Omega$ in Bezug auf $\Omega'$, und $m'$
der Grad von $\Omega'$ in Bezug auf $R$.

Es ist dabei zu beachten, dass zwar $\Omega$ ein Normalkörper
auch in Bezug auf $\Omega'$ ist, dass aber $\Omega'$ im Allgemeinen kein
Normalkörper in Bezug auf $R$ sein wird.

Die Primideale im Körper $\Omega'$ lassen sich nun vollständig aus
denen von $\Omega$ ableiten.

Wenn $\mathfrak{p}$ irgend ein Primideal in $\Omega$ ist, so giebt es \gesperrt{ein und
nur ein durch $\mathfrak{p}$ theilbares Ideal $\mathfrak{p}'$ in $\Omega'$}, und $\mathfrak{p}'$ kann als
Theiler von $\mathfrak{P}$ durch keine anderen als die mit $\mathfrak{p}$ conjugirten Ideale
theilbar sein. Durchläuft $\varphi'$ die Substitutionen von $\Phi'$, so ist $\mathfrak{p}'$,
da es durch $\varphi'$ ungeändert bleibt, auch durch $\mathfrak{p}\,|\,\varphi'$ theilbar.

Wenn also $\mathfrak{p}_1'$ durch $\mathfrak{p}\,|\,\varphi_1$ theilbar ist, so ist es auch durch
$\mathfrak{p}\,|\,\psi\,\varphi_1\,\varphi'$ theilbar, wenn $\psi$ ein beliebiges Element der Gruppe
$\boldsymbol{\Psi}$ (\S.\,162) bedeutet.

Wenn daher $\varphi$ irgend ein Element des Systemes

\[
\Phi_1 = \boldsymbol{\Psi}\,\varphi_1\,\Phi'
\]

ist, so ist $\mathfrak{p}_1'$ durch $\mathfrak{p}\,|\,\varphi$ theilbar. Es kommt demnach jetzt die
Gruppenzerlegung in Betracht, die wir im \S.\,28 dieses Bandes
auseinandergesetzt haben.

Wenn $\varphi_2$ ein nicht in $\Phi_1$ enthaltenes Element aus $\Phi$ ist, so
hat das System

\[
\Phi_2 = \boldsymbol{\Psi}\,\varphi_2\,\Phi'
\]

mit $\Phi_1$ gar kein Element gemein. Giebt es noch ein Element $\varphi_3$,
das weder in $\Phi_1$ noch in $\Phi_2$ vorkommt, so bildet man ebenso

\[
\Phi_3 = \boldsymbol{\Psi}\,\varphi_3\,\Phi',
\]

und fährt so fort, bis die ganze Gruppe $\Phi$ erschöpft ist. Man
erhält so die Zerlegung

\weq{2}{\Phi = \Phi_1 + \Phi_2 + \Phi_3 + \cdots + \Phi_{e'}.}

Den Grad eines dieser Systeme $\Phi_r$ können wir nach \S.\,28 bestimmen.
Er ist gleich dem Producte des Grades von $\Phi$ multiplicirt mit % [sic] source prints $\Phi$; the equivalent expression below requires $\boldsymbol\Psi$
dem Index des Durchschnittes $\Psi_r'$ von $\Phi'$ mit $\Psi_r = \varphi_r^{-1}\boldsymbol{\Psi}\varphi_r$ in
Bezug auf $\Phi'$, oder was dasselbe ist, gleich dem Producte der
Grade von $\boldsymbol{\Psi}$ und $\Phi'$, getheilt durch den
%CONT>>


% ===== scan leaf 620, printed 599 =====
% !PRINTED 599
\weberpage{599}
%<<CONT
Grad $h_r$ des Durchschnittes von $\Phi'$ und $\Psi_r$. Der Grad von $\Phi_r$
ist also gleich $f g n' : h_r$.

Wenn wir jetzt mit $\varphi_r$ alle Elemente von $\Phi_r$ bezeichnen, so
führen alle Primideale $\mathfrak{p}\,|\,\varphi_r$ in $\Omega$ zu demselben Primideale
$\mathfrak{p}_r'$ in $\Omega'$. Es ist aber noch zu zeigen, dass zwei verschiedene
dieser Systeme $\Phi_1, \Phi_2$ zu \gesperrt{verschiedenen} Primidealen $\mathfrak{p}_1', \mathfrak{p}_2'$
führen.

Zunächst ist ersichtlich, dass die sämmtlichen $\mathfrak{p}\,|\,\varphi_2$ von den
$\mathfrak{p}\,|\,\varphi_1$ verschieden sind. Denn wenn $\mathfrak{p}\,|\,\varphi_1 = \mathfrak{p}\,|\,\varphi_2$ ist, so ist
auch $\mathfrak{p} = \mathfrak{p}\,|\,\varphi_2\,\varphi_1^{-1}$, also $\varphi_2\,\varphi_1^{-1}$ in $\boldsymbol{\Psi}$ und folglich $\varphi_2$ in
$\boldsymbol{\Psi}\varphi_1$, d.\,h. in $\Phi_1$ enthalten, gegen die Voraussetzung. Wir
können also eine ganze Zahl $\alpha$ in $\Omega$ annehmen, die durch
$\mathfrak{p}\,|\,\varphi_1$, aber durch keines der Ideale $\mathfrak{p}\,|\,\varphi_2$ theilbar ist. Dann
ist auch keine der Zahlen $\alpha\,|\,\varphi'$ durch $\mathfrak{p}\,|\,\varphi_2$ theilbar, weil sonst
$\mathfrak{p}\,|\,\varphi_1\,\varphi' = \mathfrak{p}\,|\,\varphi_2$, also gegen die Voraussetzung $\varphi_2$ in $\Phi_1$
enthalten wäre. Das Product $\alpha'$ aller von einander verschiedener
$\alpha\,|\,\varphi'$, welches eine in $\Omega'$ enthaltene Zahl ist, ist zwar durch
$\mathfrak{p}_1'$, nicht aber durch $\mathfrak{p}_2'$ theilbar. Folglich ist $\mathfrak{p}_1'$ von $\mathfrak{p}_2'$
verschieden, und es ergiebt sich also, dass $e'$ und
nicht mehr \gesperrt{verschiedene} Primideale $\mathfrak{p}'$ in $\mathfrak{P}$ aufgehen.
Wir setzen daher

\weq{3}{\mathfrak{P} = \mathfrak{p}_1'^{\,a_1}\,\mathfrak{p}_2'^{\,a_2}\cdots\mathfrak{p}_{e'}'^{\,a_{e'}},}

worin die Exponenten $a_1, a_2, \ldots, a_{e'}$ noch näher zu bestimmende
natürliche Zahlen sind.

Um nun die Primideale $\mathfrak{p}_r'$ im Körper $\Omega$ zu zerlegen, können
wir die Resultate der \S.\,160, 162 benutzen, indem wir einfach
$\Omega'$ an Stelle von $R$ treten lassen. Bedeutet dann $X_r'$ den Durchschnitt
 von $\Phi'$ mit $X_r = \varphi_r^{-1}\,X\,\varphi_r$, so ist $X_r$ der Inbegriff aller % [sic] source omits the prime; the next sentence correctly refers to $X_r'$
% [sic: The source says "Körpers" although $\Phi'$ denotes the group, not the field.]
Substitutionen des Körpers $\Phi'$, die der Bedingung

\[
\tau\,|\,\chi_r' \equiv \tau \ (\mathrm{mod}\ \mathfrak{p}_r)
\]

genügen. Bezeichnen wir daher den Grad von $X_r'$ mit $g_r$, so ist
$\mathfrak{p}_r'$ durch $\mathfrak{p}_r^{g_r}$, aber durch keine höhere Potenz von $\mathfrak{p}_r$ theilbar.

Der Durchschnitt $\Psi_r'$ von $\Psi_r$ mit $\Phi'$, dessen Grad wir schon
oben mit $h_r$ bezeichnet haben, ist der Inbegriff aller Substitutionen
 $\psi_r'$ in $\Phi'$, die der Bedingung

\[
\mathfrak{p}_r\,|\,\psi_r' = \mathfrak{p}_r
\]

genügen, und wenn wir daher $\Phi'$ in die Nebengruppen

\weq{4}{\Phi' = \Psi_r'\varphi_{r,1}' + \Psi_r'\varphi_{r,2}' + \cdots + \Psi_r'\varphi_{r,e_r}} % [sic] source omits the prime on the final $\varphi$, unlike the two preceding terms

zerlegen, so ist

\weq{5}{n' = e_r\,h_r.}


% ===== scan leaf 621, printed 600 =====
% !PRINTED 600
\weberpage{600}
Wenn nun

\[
\mathfrak{p}_{r,s} = \mathfrak{p}_r\,|\,\varphi_{r,s}
\]

gesetzt wird, so ist nach \S.\,160, (6)

\weq{6}{\mathfrak{p}_r' = (\mathfrak{p}_{r,1}\,\mathfrak{p}_{r,2}\cdots\mathfrak{p}_{r,e_r})^{g_r},}

und die Substitution von (6) in (3) ergiebt durch Vergleichung
mit \S.\,160, (6)

\weq{7}{g = a_r\,g_r, \quad e_1 + e_2 + \cdots + e_{e'} = e,}

wodurch die Exponenten $a_r$ definirt sind.

Die in Bezug auf den Körper $\Omega'$ genommene Partialnorm
von $\mathfrak{p}_{r,s}$ ist nach \S.\,160, (3), (7):

\weq{8}{\mathfrak{N}_{\Omega'}(\mathfrak{p}_{r,s}) = \mathfrak{p}_r'^{\,f_r},}

wenn $f_r$ durch die Gleichung

\weq{9}{n' = e_r\,f_r\,g_r, \quad (h_r = f_r\,g_r)}

definirt wird.

Wir wollen nun die Gruppe $\Phi$ nach $\Phi'$ in Nebengruppen zerlegen,
und setzen, wenn $\vartheta_1, \vartheta_2, \ldots, \vartheta_{m'}$ passend ausgewählte
Substitutionen aus $\Phi$ sind:

\weq{10}{\Phi = \Phi'\vartheta_1 + \Phi'\vartheta_2 + \cdots + \Phi'\vartheta_{m'},}

so dass durch die Substitutionen $\vartheta_1, \vartheta_2, \ldots, \vartheta_{m'}$ der Körper
$\Omega'$ in $m'$ conjugirte (gleiche oder verschiedene) Körper

\[
\Omega_1', \ \Omega_2', \ \ldots, \ \Omega_{m'}'
\]

übergeht. Der Körper $\Omega_t'$ gehört dann zu der Gruppe $\vartheta_t^{-1}\Phi'\vartheta_t$.

Wenn nun durch $\vartheta_t$ die Ideale $\mathfrak{p}_r'$ und $\mathfrak{p}_{r,s}$ in $\mathfrak{p}_{r,t}'$ und
$\mathfrak{p}_{r,s,t}$ übergehen, so ist nach \S.\,160, (6)

\weq{11}{\mathfrak{p}_{r,t}' = (\mathfrak{p}_{r,1,t}\,\mathfrak{p}_{r,2,t}\cdots\mathfrak{p}_{r,e_r,t})^{g_r},}

und das Product aller dieser Ideale für $t = 1, 2, \ldots, m'$ ist die
im Körper $\Omega'$ genommene Partialnorm von $\mathfrak{p}_r'$ in Bezug auf den
Körper $R$, die wir mit $\mathfrak{N}_R'(\mathfrak{p}_r')$ bezeichnen. Sie muss eine Potenz
von $\mathfrak{P}$ sein, und wir setzen

\weq{12}{\mathfrak{N}_R'(\mathfrak{p}_r') = \mathfrak{P}^{\,f_r'}.}

Um $f_r'$ zu finden, brauchen wir nur $\mathfrak{P}$ in (12) nach \S.\,160, (6)
in seine Primfactoren $\mathfrak{p}$ zu zerlegen, wodurch sich $e\,g\,f_r'$ Primfactoren
 $\mathfrak{p}$ in $\mathfrak{P}^{\,f_r'}$ ergeben. Bilden wir andererseits das Product
der $m'$ Factoren (11), so erhalten wir $g_r\,e_r\,m'$ solcher Primfactoren.
Es ist daher

\[
e\,g\,f_r' = m'\,g_r\,e_r,
\]

also nach (1), (9) und \S.\,160, (7):

\weq{13}{f = f_r\,f_r'.}

\begin{center}
\rule{0.3\linewidth}{0.4pt}
\end{center}


% ===== scan leaf 622, printed 601 =====
\weberpage{601}% page number not printed on this page
\pdfbookmark[1]{Zwanzigster Abschnitt. Quadratische Körper.}{bm187}
\pdfbookmark[2]{§ 164. Basis eines quadratischen Körpers}{bm188}
% !PRINTED none
\begin{center}
\gesperrt{Zwanzigster Abschnitt.}

\gesperrt{Quadratische Körper.}
\end{center}

\begin{center}
\rule{0.2\linewidth}{0.4pt}
\end{center}

\begin{center}
\S.\,164.

\gesperrt{Basis eines quadratischen Körpers.}
\end{center}

Wir wollen die Resultate der allgemeinen Theorie der algebraischen Zahlen auf den besonderen Fall der Körper vom
$2^{\text{ten}}$ Grade, die auch quadratische Körper heissen, anwenden.
Ein quadratischer Körper $\Omega$ entspringt aus einer ganzzahligen
quadratischen Gleichung, und wird also aus dem Körper der
rationalen Zahlen durch Adjunction einer Quadratwurzel $\sqrt{d}$ abgeleitet.
Hierin kann $d$ als positive oder negative ganze Zahl
ohne quadratischen Theiler angenommen werden, und der einzige
 Werth $d = 1$ ist auszunehmen. Je nachdem $d$ positiv oder
negativ ist, erhalten wir einen reellen oder imaginären quadratischen
 Körper.

Den conjugirten Körper erhalten wir, wenn wir in allen
Zahlen in $\Omega$ $\sqrt{d}$ in $-\sqrt{d}$ verwandeln. Dadurch bekommen wir
aber keinen neuen Körper. Der quadratische Körper ist also
ein (absoluter) Normalkörper (\S.\,160).

Die Zahlen des Körpers $\Omega$ sind keine anderen als die in
\S.\,122 des ersten Bandes definirten quadratischen Irrationalzahlen,
die alle in der Form $x + y\sqrt{d}$ enthalten sind, worin
$x$ und $y$ ganze oder gebrochene rationale Zahlen bedeuten. Jede
irrationale Zahl $\omega$ dieses Körpers genügt einer quadratischen
Gleichung

\weq{1}{c\,\omega^2 = a + b\,\omega,}

in der $a, b, c$ ganze rationale Zahlen ohne gemeinsamen Theiler
sind, von denen $c$ als positiv angenommen werden kann, und
die Verbindung

\weq{2}{D = b^2 + 4\,a\,c}
%CONT>>


% ===== scan leaf 623, printed 602 =====
% !PRINTED 602
\weberpage{602}
%<<CONT
haben wir die Discriminante der Zahl $\omega$ genannt. Das Verh\"altniss $D:d$ ist eine ganze rationale Quadratzahl. Die beiden Wurzeln von (1) erhalten den Ausdruck

\weq{3}{\omega = \frac{b+\sqrt{D}}{2c}, \quad \omega' = \frac{b-\sqrt{D}}{2c}.}
Setzen wir
\weq{4}{D = e^2 d,}
so wird
\weq{5}{\omega = \frac{b+e\sqrt{d}}{2c}, \quad \omega' = \frac{b-e\sqrt{d}}{2c},}
und wegen (1) sind dies dann und nur dann ganze Zahlen, wenn $c=1$ ist.

Nun ist $D\equiv 0$ oder $\equiv 1\ (\mathrm{mod}\ 4)$, und wenn also $D$ gerade ist, so ist es durch 4 theilbar; $b$ ist in diesem Falle gerade, und weil $d$ nicht durch 4 theilbar sein kann, so muss auch $e$ gerade sein.

Ist aber $D$ ungerade, so sind auch $b$ und $e$ ungerade. Dieser Fall kann aber, wie (4) zeigt, nur dann eintreten, wenn $d\equiv 1\ (\mathrm{mod}\ 4)$ ist.

Daraus erhalten wir folgendes Resultat:

1. \gesperrt{Ist $d\equiv 2$ oder $\equiv 3\ (\mathrm{mod}\ 4)$, so ist die Zahl}
\[ \omega = x + y\sqrt{d} \]
\gesperrt{dann und nur dann eine ganze Zahl, wenn die rationalen Zahlen $x, y$ ganze Zahlen sind.} (Quadratische Irrationalzahlen erster Art, nach \gesperrt{Dirichlet}.)

2. \gesperrt{Ist $d\equiv 1\ (\mathrm{mod}\ 4)$, so ist}
\[ \omega = \frac{x+y\sqrt{d}}{2} \]
\gesperrt{dann und nur dann eine ganze Zahl, wenn die rationalen Zahlen $x,y$ entweder zwei gerade oder zwei ungerade ganze Zahlen sind.} (Quadratische Irrationalzahlen erster oder zweiter Art, je nachdem $e$ gerade oder ungerade ist.)

Letzteres k\"onnen wir auch so ausdr\"ucken:

3. \gesperrt{Ist $d\equiv 1\ (\mathrm{mod}\ 4)$, so ist}
\[ \omega = x + y\,\frac{1+\sqrt{d}}{2} \]
\gesperrt{dann und nur dann eine ganze Zahl, wenn die rationalen Zahlen $x$ und $y$ ganze Zahlen sind.}


% ===== scan leaf 624, printed 603 =====
% !PRINTED 603
\weberpage{603}
\pdfbookmark[2]{§ 165. Die Primideale in Ω}{bm189}
Daraus ergiebt sich eine Basis von $\mathfrak{o}$:

a) \quad $1, \sqrt{d},$ \hfill $d\equiv 2$ oder $\equiv 3\ (\mathrm{mod}\ 4),$

b) \quad $1, \dfrac{1+\sqrt{d}}{2},$ \hfill $d\equiv 1\ (\mathrm{mod}\ 4).$

Wenn wir also

\weq{6}{\begin{aligned}
\text{a)}\quad \Theta &= \sqrt{d}, & d &\equiv 2, 3\ (\mathrm{mod}\ 4)\\
\text{b)}\quad \Theta &= \frac{1+\sqrt{d}}{2}, & d &\equiv 1\ (\mathrm{mod}\ 4)
\end{aligned}}
setzen, so ist $1, \Theta$ eine Basis von $\mathfrak{o}$.

F\"ur die Grundzahl erhalten wir demnach, wenn $\Theta'$ mit $\Theta$ conjugirt ist,
\[ \Delta = \left|\begin{array}{cc} 1, & \Theta \\ 1, & \Theta' \end{array}\right|^2 = (\Theta'-\Theta)^2, \]
also

\weq{7}{\begin{aligned}
\text{a)}\quad \Delta &= 4d, & d &\equiv 2, 3\ (\mathrm{mod}\ 4)\\
\text{b)}\quad \Delta &= d, & d &\equiv 1\ (\mathrm{mod}\ 4),
\end{aligned}}
\weq{8}{\Theta - \Theta' = \sqrt{\Delta}.}

Die Grundzahl ist also $\equiv 0$ oder $\equiv 1\ (\mathrm{mod}\ 4)$, und muss, wenn sie durch 4 theilbar ist, $\equiv 8, 12\ (\mathrm{mod}\ 16)$ sein\footnote{Eine ganze rationale Zahl $\Delta$, die nach dem Modul 4 mit 0 oder mit 1 congruent ist, hat \gesperrt{Kronecker} eine Zahl von Discriminantenform genannt. Hat $\Delta$ die Eigenschaft, dass sich kein quadratischer Factor so daraus absondern l\"asst, dass eine Zahl von Discriminantenform \"ubrig bleibt, so heisst $\Delta$ eine \gesperrt{Stammdiscriminante}. Die Grundzahl eines quadratischen K\"orpers ist also immer eine \gesperrt{Stammdiscriminante}, und umgekehrt ist auch jede Stammdiscriminante Grundzahl eines quadratischen K\"orpers. Vergl. \gesperrt{H. Weber}, Zahlentheoretische Untersuchungen aus dem Gebiete der elliptischen Functionen. I, II, III. Nachrichten der Gesellschaft der Wissenschaften zu G\"ottingen. 1893.}. Sie kann daher (in Uebereinstimmung mit dem Satze von \gesperrt{Minkowski}) weder $=\pm 1$ noch $=\pm 2$ sein. Die Grundzahl vom absolut kleinsten Werthe ist $\Delta = -3.$

\begin{center}
\S.\,165.\\
\gesperrt{Die Primideale in $\Omega$.}
\end{center}

Es bedeute nun $\varphi$ ein beliebiges ganzes Functional in $\Omega$. Wir wollen nach \S.\,146 eine Basis $\alpha_1, \alpha_2$ f\"ur $\varphi$, und damit die
%CONT>>


% ===== scan leaf 625, printed 604 =====
% !PRINTED 604
\weberpage{604}
%<<CONT
entsprechende Basisform $\alpha_1 t_1 + \alpha_2 t_2$, die mit $\varphi$ associirt ist, bestimmen. Wir haben nach \S.\,146, (3), wenn $\omega_1, \omega_2$ eine Basis von $\mathfrak{o}$ ist, die ganzen Zahlen $a_{1,1}, a_{1,2}, a_{2,2}$ so zu bestimmen, dass $a_{1,1}, a_{2,2}$ positiv und m\"oglichst klein sind, und dass
\[ \alpha_1 = a_{1,1}\omega_1, \quad \alpha_2 = a_{1,2}\omega_1 + a_{2,2}\omega_2 \]
durch $\varphi$ theilbar werden. Setzen wir also $(\omega_1, \omega_2) = (1, \Theta)$, so wird

\weq{1}{\alpha_1 = a_{1,1}, \quad \alpha_2 = a_{1,2} + a_{2,2}\Theta.}

Es ist also $a_{1,1}$ die kleinste positive durch $\varphi$ theilbare rationale Zahl, und $a_{2,2}$ ist die kleinste positive Zahl, f\"ur die $a_{2,2}\Theta$ nach dem Modul $\varphi$ mit einer rationalen Zahl congruent wird.

Wir wollen dies auf den Fall anwenden, dass $\varphi$ ein Primfunctional $\pi$ ist. Ein solches kann hier nur vom ersten oder vom zweiten Grade sein, je nachdem seine absolute Norm gleich $p$ oder $=p^2$ ist. Wir unterscheiden also die beiden F\"alle:

1) \quad $N_a(\pi) = p$

oder

2) \quad $N_a(\pi) = p^2.$

Da hier $N(p) = p^2$ ist, so ist im ersten Falle $p$ in zwei Primfactoren zerlegbar, im zweiten Falle enth\"alt $p$ nur einen Primfactor, d. h. $p$ ist im K\"orper $\Omega$ selbst noch als Primzahl zu betrachten.

Bestimmen wir nun f\"ur $\varphi = \pi$ die Basis (1), so ergiebt sich zun\"achst $a_{1,1} = p$, und $a_{2,2}$ ist die kleinste positive ganze rationale Zahl, f\"ur die das Product $a_{2,2}\Theta$ modulo $\pi$ mit einer rationalen Zahl congruent wird. Daf\"ur ist aber nach \S.\,150, 3 die nothwendige und hinreichende Bedingung
\[ (a_{2,2}\Theta)^p \equiv a_{2,2}\Theta \pmod{\pi}, \]
oder, da nach dem \gesperrt{Fermat}'schen Satze $a_{2,2}^p \equiv a_{2,2}$ ist,

\weq{2}{a_{2,2}(\Theta^p - \Theta) \equiv 0 \pmod{\pi}.}

Es ist also entweder $\Theta^p - \Theta \equiv 0$, und dann ist $a_{2,2} = 1$, oder $\Theta^p - \Theta$ nicht $\equiv 0$, und dann ist $a_{2,2} = p$.

Im ersten Falle ist $\Theta$ und damit jede Zahl $\omega$ mit einer rationalen Zahl congruent. Es giebt also nicht mehr als $p$ nach dem Modul $\pi$ incongruente Zahlen, und folglich ist (\S.\,148, 3.) $p^f = p$, $f = 1$, d. h. $p$ in zwei Primfactoren ersten Grades zerlegbar.


% ===== scan leaf 626, printed 605 =====
% !PRINTED 605
\weberpage{605}
Umgekehrt bilden in diesem Falle 1) die $p$ rationalen Zahlen $0, 1, 2, \ldots, p-1$ ein volles Restsystem nach dem Modul $\pi$, und es ist folglich jede Zahl in $\mathfrak{o}$ und mithin auch $\Theta$ mit einer rationalen Zahl congruent. Also haben wir

1) \quad $N_a(\pi) = p$ \quad $a_{2,2} = 1$

2) \quad $N_a(\pi) = p^2$ \quad $a_{2,2} = p.$

Im ersten Falle ist $a_{1,2}$ aus der Bedingung

\weq{3}{a_{1,2} + \Theta \equiv 0 \pmod{\pi}}
zu bestimmen, im zweiten kann $a_{1,2} = 0$ genommen werden. Die Basisform von $\pi$ ergiebt sich also, wenn $u, v$ die Variablen sind,

1) \quad $\pi = pu + (a_{1,2}+\Theta)v$

2) \quad $\pi = p(u+\Theta v),$

woraus man sieht, dass $\pi$ im Falle 2), wie zu erwarten war, ein \gesperrt{Hauptfunctional} ist.

Um im Falle 1) die Zahl $a_{1,2}$ zu bestimmen, bilden wir die Norm von $a_{1,2}+\Theta$, die durch $p$ theilbar sein muss. Wenn $\Theta$ durch Aenderung des Vorzeichens von $\sqrt{d}$ in $\Theta'$ \"ubergeht, so ist nach \S.\,164, (8) $\Theta - \Theta' = \sqrt{\Delta}$, und f\"ur die Norm von $a_{1,2}+\Theta$ ergiebt sich nach \S.\,164 (6), a), b):

\weq{4}{N(a_{1,2}+\Theta) = (a_{1,2}+\Theta)(a_{1,2}+\Theta')}
\[ {} = a_{1,2}^2 - d, \quad d\equiv 2, 3\ (\mathrm{mod}\ 4) \]
\[ {} = a_{1,2}^2 + a_{1,2} + \frac{1-d}{4}, \quad d\equiv 1\ (\mathrm{mod}\ 4). \]

Wenn diese Zahl durch $p$ theilbar ist, so ist einer der beiden Factoren $a_{1,2}+\Theta$, $a_{1,2}+\Theta'$ durch $\pi$ theilbar. Da aber $\Theta' = -\Theta$ oder $=1-\Theta$ ist, also

a) \quad $a_{1,2} + \Theta' = -(-a_{1,2}+\Theta)$

b) \quad $a_{1,2} + \Theta' = -(-a_{1,2}-1+\Theta)$

ist, so kann man $a_{1,2}$ immer so annehmen, dass gerade $a_{1,2}+\Theta$ durch $\pi$ theilbar wird, und $\Theta$ ist also wirklich mit einer rationalen Zahl congruent. Die nothwendige und hinreichende Bedingung daf\"ur, dass $p$ zwei Primfactoren enth\"alt, ist folglich die, dass $a_{1,2}^2 - d$ oder $a_{1,2}^2 + a_{1,2} + \frac14(1-d)$ durch geeignete Annahme \"uber $a_{1,2}$ durch $p$ theilbar wird. Durch Multiplication mit 4 gehen diese beiden Bedingungen \"uber in
\[ 4a_{1,2}^2 - 4d \equiv 0, \quad (2a_{1,2}+1)^2 - d \equiv 0 \pmod{4p}, \]
%CONT>>


% ===== scan leaf 627, printed 606 =====
% !PRINTED 606
\weberpage{606}
%<<CONT
und wenn man bedenkt, dass die Grundzahl $\Delta$ im Falle a) gleich $4d$, im Falle b) gleich $d$ ist, so erhalten wir, wenn wir
\weq{5}{x = 2a_{1,2} \text{ oder } {}= 2a_{1,2}+1}
setzen, folgendes Resultat:

1. \gesperrt{Die nothwendige und hinreichende Bedingung daf\"ur, dass $p$ zwei Primfactoren enth\"alt, ist die, dass die Congruenz}
\weq{6}{x^2 \equiv \Delta\ (\mathrm{mod}\ 4p)}
\gesperrt{durch eine ganze rationale Zahl $x$ befriedigt werden kann, oder dass die Grundzahl $\Delta$ quadratischer Rest von $4p$ ist.}

Bezeichnen wir mit $\pi, \pi'$ die beiden conjugirten Formen
\weq{7}{\pi = pu + (a_{1,2}+\Theta)v,\quad \pi' = pu + (a_{1,2}+\Theta')v,}
so ergiebt sich durch Multiplication
\weq{8}{\pi\pi' = p\left(pu^2 + xuv + \frac{x^2-\Delta}{4p}v^2\right).}

Hierin ist nun
\[
pu^2 + xuv + \frac{x^2-\Delta}{4p}v^2
\]
eine Einheit, weil, wenn $p$ in $\Delta$ nicht aufgeht, $x$ durch $p$ nicht theilbar ist, und wenn $p$ in $\Delta$ aufgeht, zwar $x$, aber nicht $(x^2-\Delta:4p)$ durch $p$ theilbar ist. Denn ist $p$ ungerade, so kann $p^2$ nicht in $\Delta$ aufgehen, und ist $p=2$ und $\Delta$ durch 2 theilbar, so ist $\Delta = 4d$, $d \equiv 2, 3\ (\mathrm{mod}\ 4)$. W\"are aber
\[
x^2 - \Delta \equiv 0\ (\mathrm{mod}\ 16),
\]
so w\"urde folgen
\[
d \equiv \left(\frac{x}{2}\right)^2 (\mathrm{mod}\ 4),
\]
w\"ahrend doch $\left(\dfrac{x}{2}\right)^2$ als Quadrat nicht mit 2 oder 3 congruent sein kann. Demnach ist $p$ in die beiden Primfactoren $\pi$ und $\pi'$ zerlegt.

Die beiden Factoren $\pi$ und $\pi'$ von $p$ k\"onnen auch mit einander associirt sein. Dies tritt nur dann ein, wenn $\pi'$ und folglich auch $\pi - \pi'$ durch $\pi$ theilbar ist. Nun ist aber nach (7)
\[
\pi - \pi' = (\Theta - \Theta')v,
\]
und folglich muss
\[
(\Theta - \Theta')^2 = \Delta
\]
%CONT>>


% ===== scan leaf 628, printed 607 =====
% !PRINTED 607
\weberpage{607}
\pdfbookmark[2]{§ 166. Functionale im quadratischen Körper und quadratische Irrationalzahlen}{bm190}
%<<CONT
durch $p$ theilbar sein, und dies ist auch die hinreichende Bedingung.

2. \gesperrt{Es ist also $p$ nur dann mit dem Quadrat eines Primfunctionals associirt, wenn $p$ in der Grundzahl $\Delta$ aufgeht, in Uebereinstimmung mit dem allgemeinen Satze (\S.\,158).}

\begin{center}
\S.\,166.

\gesperrt{Functionale im quadratischen K\"orper und quadratische Irrationalzahlen.}
\end{center}

Wir betrachten nun in unserem quadratischen K\"orper ein beliebiges ganzes Functional und bilden seine Basisform
\[
\varphi = \alpha_1 t_1 + \alpha_2 t_2,
\]
indem wir $\alpha_1, \alpha_2$ nach \S.\,165 bestimmen:
\weq{1}{\alpha_1 = a_{1,1}, \quad \alpha_2 = a_{1,2} + a_{2,2}\Theta.}
Nach \S.\,146 ist dann
\weq{2}{N_a(\varphi) = a_{1,1}a_{2,2}.}
Geben wir der Form $\varphi$ den Ausdruck
\weq{3}{\varphi = a_{1,1}t_1 + (a_{1,2}+a_{2,2}\Theta)t_2,}
und bilden die Norm, so ergiebt sich
\weq{4}{N(\varphi) = a_{1,1}^2 t_1^2 + [2a_{1,1}a_{1,2} + a_{1,1}a_{2,2}(\Theta+\Theta')]t_1 t_2 + t_2^2 N(\alpha_2)}
und der Theiler dieser Form ist also nach (2) gleich $a_{1,1}a_{2,2}$.
Wenn wir also
\weq{5}{\begin{aligned}
a_{1,1}^2 &= c\,a_{1,1}a_{2,2}\\
2a_{1,1}a_{1,2}+a_{1,1}a_{2,2}(\Theta+\Theta') &= b\,a_{1,1}a_{2,2}\\
N(\alpha_2) &= -a\,a_{1,1}a_{2,2}
\end{aligned}}
setzen, so sind $a, b, c$ ganze rationale Zahlen ohne gemeinschaftlichen Theiler, und der Quotient
\weq{6}{\frac{\alpha_2}{\alpha_1} = \omega}
ist eine Wurzel der quadratischen Gleichung
\weq{7}{c\omega^2 = a + b\omega.}
% [sic: The source prints the incomplete subscript $a_{1,}$ here; equations (1) and (5) identify the intended quantity as $a_{1,2}$.]
Wenn wir in den Ausdr\"ucken (1) die Werthe von $a_{1,1}$ und $a_{1,}$ aus (5) substituiren, so folgt
\weq{8}{\begin{aligned}
2\alpha_2 &= a_{2,2}(b+\Theta-\Theta')\\
\alpha_1 &= c\,a_{2,2}.
\end{aligned}}


% ===== scan leaf 629, printed 608 =====
% !PRINTED 608
\weberpage{608}
Nun ist nach \S.\,164, (8) die Differenz $\Theta - \Theta'$ gleich der Quadratwurzel aus der Grundzahl von $\Omega$, so dass wir aus (6) und (8) erhalten:
\weq{9}{\omega = \frac{b+\sqrt{\Delta}}{2c}, \quad \Delta = b^2 + 4ac,}
worin das Vorzeichen von $\sqrt{\Delta}$ durch die in \S.\,164, (6) gemachte Annahme \"uber das Vorzeichen von $\sqrt{d}$ bestimmt ist.

Es ist also $\omega$ eine zur Discriminante $\Delta$ geh\"orige \gesperrt{quadratische Irrationalzahl}, wie wir sie im \S.\,122 des ersten Bandes betrachtet haben, und die Form $\varphi$ kann, wenn wir $a_{2,2}=e$ setzen, so ausgedr\"uckt werden:
\weq{10}{\varphi = ec(t_1+\omega t_2).}

Hierin ist $c$ die kleinste positive ganze rationale Zahl, f\"ur die das Product $c\omega$ eine ganze Zahl wird, und jede andere ganze rationale Zahl $k$, durch die das Product $k\omega$ eine ganze Zahl wird, muss durch $c$ theilbar sein.

Man erkennt dies am einfachsten, wenn man $\omega$ nach \S.\,164, je nachdem $b$ gerade oder ungerade ist, in eine der beiden Formen setzt:
\weq{11}{\omega = \frac{1}{c}\left(\frac{b}{2}+\Theta\right), \quad \frac{1}{c}\left(\frac{b-1}{2}+\Theta\right),}
woraus, da $1, \Theta$ eine Basis von $\mathfrak{o}$ ist, das Gesagte hervorgeht.

Darauf gr\"undet sich nun der Beweis, dass, wenn $\omega$ irgend eine zur Discriminante $\Delta$ geh\"orige quadratische Irrationalzahl bedeutet, $c$ die kleinste positive ganze rationale Zahl, die $c\omega$ zu einer ganzen Zahl macht, und $e$ eine beliebige positive ganze rationale Zahl, auch umgekehrt durch
\weq{12}{\varphi = ec(t_1+\omega t_2)}
immer eine Basisform dargestellt ist.

Dazu haben wir nur zu zeigen, dass
\[
\alpha_1 = ec, \quad \alpha_2 = ec\omega
\]
eine Basis des durch (12) definirten Functionals $\varphi$ ist.

Bilden wir zu diesem Zwecke zun\"achst die Norm von $\varphi$, so erhalten wir nach (7)
\weq{13}{N(\varphi) = e^2 c(ct_1^2 + bt_1t_2 - at_2^2),}
und folglich die absolute Norm
\weq{14}{N_a(\varphi) = e^2 c.}


% ===== scan leaf 630, printed 609 =====
% !PRINTED 609
\weberpage{609}
Ist nun
\weq{15}{\alpha_1 = a_{1,1}, \quad \alpha_2 = a_{1,2}+a_{2,2}\Theta}
eine Basis von $\varphi$, so ist andererseits (nach \S.\,146)
\weq{16}{N_a(\varphi) = a_{1,1}a_{2,2},}
und $a_{1,1}$ ist die kleinste positive ganze rationale Zahl, die durch $\varphi$ theilbar ist.

Nun k\"onnen wir andererseits zeigen, dass diese Zahl $= ec$ ist, woraus $ec = a_{1,1}$, und nach (14) und (16) $e = a_{2,2}$ folgt.

Wenn n\"amlich eine ganze rationale Zahl $m$ durch $\varphi$ theilbar sein soll, so muss, wenn $\omega'$ zu $\omega$ conjugirt ist, da
\[
ct_1^2 + bt_1t_2 - at_2^2
\]
ein Einheitsfunctional ist,
\[
\frac{m(ct_1^2+bt_1t_2-at_2^2)}{\varphi} = \frac{m}{e}(t_1+\omega' t_2)
\]
ein ganzes Functional sein. Folglich m\"ussen
\[
\frac{m}{e} \quad \text{und} \quad \frac{m}{e}\omega'
\]
ganze Zahlen sein, woraus folgt, dass $m : e$ eine durch $c$ theilbare, also $m$ eine durch $ec$ theilbare ganze rationale Zahl ist. Da andererseits $ec$ durch $\varphi$ theilbar ist, so folgt, dass $ec$ die kleinste diesen Forderungen gen\"ugende Zahl ist. Ausserdem ergiebt sich noch, dass die ganze Zahl $ec\omega$ durch $\varphi$ theilbar ist. Denn $ec\omega$ ist in der Darstellung (12) ein Co\"efficient der ganzen Function $\varphi$ und daher durch $\varphi$ theilbar (\S.\,142).

Dr\"ucken wir $\Theta$ in (15) durch $\omega$ aus, so ergiebt sich nach (11)
\[
\begin{aligned}
\alpha_2 &= a_{1,2} - \frac{eb}{2} + ec\omega, \quad b \text{ gerade}\\
&= a_{1,2} - \frac{e(b-1)}{2} + ec\omega, \quad b \text{ ungerade,}
\end{aligned}
\]
und es ist also
\[
a_{1,2} \equiv \frac{eb}{2} \quad \text{oder} \quad \equiv \frac{e(b-1)}{2}
\]
nach dem Modul $\varphi$, und folglich, da beide Seiten rational sind, auch nach dem Modul $a_{1,1}=ec$. Demnach kann $a_{1,2}$ dem Werthe $\frac{1}{2}eb$ oder $\frac{1}{2}e(b-1)$ gleich gesetzt werden, und wir erhalten $\alpha_2 = ec\omega$, was zu beweisen war.

Um das hierdurch Bewiesene \"ubersichtlich zusammenzufassen, k\"onnen wir folgende S\"atze aussprechen:
%CONT>>


% ===== scan leaf 631, printed 610 =====
% !PRINTED 610
\weberpage{610}
\pdfbookmark[2]{§ 167. Aequivalente Functionale und äquivalente Zahlen}{bm191}
%<<CONT

1. \gesperrt{Ist $\Omega$ ein quadratischer Körper mit der Grundzahl $\varDelta$, und $\varphi$ ein ganzes Functional in diesem Körper, so lässt sich eine zur Discriminante $\varDelta$ gehörige quadratische Irrationalzahl}

\weq{17}{\omega = \frac{b + \sqrt{\varDelta}}{2c}}

\gesperrt{und eine positive ganze rationale Zahl $e$ so bestimmen, dass}

\weq{18}{ec(t_1 + \omega t_2)}

\gesperrt{eine Basisform von $\varphi$ wird. Umgekehrt ist, wenn $\omega$ eine beliebige zur Discriminante $\varDelta$ gehörige quadratische Irrationalzahl ist, die Form (18) immer eine Basisform eines Functionals $\varphi$.}

\begin{center}
§. 167.

\gesperrt{Aequivalente Functionale und äquivalente Zahlen.}
\end{center}

Wenn nun $\omega, \omega_1$ irgend zwei zu der Discriminante $\varDelta$ gehörige quadratische Irrationalzahlen sind, so erhalten wir zwei Basisformen

\weq{1}{\begin{aligned} \varphi &= ec(t_1+\omega t_2)\\ \varphi_1 &= e_1c_1(t_1+\omega_1 t_2). \end{aligned}}

Wenn nun die beiden Zahlen $\omega, \omega_1$ in dem Sinne äquivalent sind, wie wir diesen Begriff im \S.\,120 des ersten Bandes festgesetzt haben, so ist

\weq{2}{\omega_1 = \frac{p\omega+q}{r\omega+s},}

worin $p,q,r,s$ ganze rationale Zahlen sind, die der Bedingung

\weq{3}{ps - qr = \pm 1}

genügen. Machen wir dann die Substitution (2) in $\varphi_1$, so folgt

\[ (r\omega+s)\varphi_1 = e_1c_1[st_1+qt_2+\omega(rt_1+pt_2)], \]

oder wenn wir

\weq{4}{u_1 = st_1+qt_2,\quad u_2 = rt_1+pt_2}

setzen,

\[ (r\omega+s)\varphi_1 = e_1c_1(u_1+\omega u_2). \]

Nun ist $ec(u_1+\omega u_2)$ mit $\varphi$ associirt, weil (4) eine ganz-
%CONT>>


% ===== scan leaf 632, printed 611 =====
% !PRINTED 611
\weberpage{611}

%<<CONT
zahlige lineare Substitution mit der Determinante $\pm 1$ ist (\S.\,147), und folglich sind die beiden Formen

\[ (r\omega+s)\,ec\,\varphi_1,\quad e_1c_1\varphi \]

mit einander associirt, und dies hat zur Folge, dass die beiden Formen $\varphi,\varphi_1$ in dem in \S.\,152 definirten Sinne äquivalent sind. Damit haben wir bewiesen:

2. \gesperrt{Wenn die quadratischen Irrationalzahlen $\omega,\omega_1$ äquivalent sind, so sind auch die entsprechenden Functionale $\varphi,\varphi_1$ äquivalent.}

Es lässt sich aber auch der umgekehrte Satz beweisen:

3. \gesperrt{Wenn die Functionale $\varphi,\varphi_1$ äquivalent sind, so sind auch die quadratischen Irrationalzahlen $\omega$ und $\omega_1$ äquivalent.}

Um diesen Beweis zu führen, bezeichnen wir mit $\psi$ ein Functional der Classe, die zu der Classe von $\varphi$ reciprok ist, so dass $\varphi\psi$ mit einer Zahl äquivalent ist. Dann können wir setzen:

\weq{5}{\varphi\psi = \varepsilon\xi,}

worin $\varepsilon$ eine Einheit, $\xi$ eine \gesperrt{Zahl} in $\mathfrak{o}$ ist. Hierin können wir $\varphi$ und $\psi$ beide als Basisformen annehmen und demnach

\weq{6}{\varphi = ec(t_1+t_2\omega),\quad \psi = e'c'(t_1+t_2\eta)}

setzen, worin $e', c', \eta$ dieselbe Bedeutung für $\psi$ haben, wie $e,c,\omega$ für $\varphi$. Wir setzen zur Abkürzung [\S.\,166, (13)]:

\weq{7}{N_a(\varphi) = e^2c = P,\quad N_a(\psi) = e'^2c' = Q.}

Wenn wir mit $\psi'$ die zu $\psi$ conjugirte Form bezeichnen, und für den Augenblick die Einheitsform

\[ c't_1^2 + b't_1t_2 - a't_2^2 = \varepsilon_1 \]

setzen, so folgt aus (5) durch Multiplication mit $\psi'$:

\[ Q\varepsilon_1\varphi = \varepsilon\xi\psi', \]

und folglich

\weq{8}{\frac{Q\varphi}{\xi} = \frac{\varepsilon}{\varepsilon_1}\psi' = \chi,}

worin $\chi$ und $\psi'$ associirt sind. Setzen wir

\weq{9}{\frac{Qec}{\xi} = \beta_1,\quad \frac{Qec\omega}{\xi} = \beta_2,}
%CONT>>


% ===== scan leaf 633, printed 612 =====
% !PRINTED 612
\weberpage{612}

%<<CONT
so wird

\[ \chi = \frac{Q\varphi}{\xi} = \beta_1t_1+\beta_2t_2, \]

und folglich sind $\beta_1,\beta_2$ ganze durch $\psi'$ theilbare Zahlen.

Bilden wir aus (9) die Discriminante $\varDelta(\beta_1,\beta_2)$, so erhalten wir

\weq{10}{\varDelta(\beta_1,\beta_2) = \begin{vmatrix}\beta_1, & \beta_2\\ \beta_1', & \beta_2'\end{vmatrix}^2 = \frac{e^4c^4Q^4}{(\xi\xi')^2}(\omega-\omega')^2.}

Aus (5) und (7) folgt aber

\[ \xi\xi' = N(\xi) = PQ = e^2cQ, \]

und wenn wir noch $\omega-\omega' = \dfrac{\sqrt{\varDelta}}{c}$ setzen, so findet sich

\weq{11}{\varDelta(\beta_1,\beta_2) = Q^2\varDelta.}

Hieraus folgt aber nach dem Satze \S.\,147, 2., dass $\beta_1,\beta_2$ eine Basis von $\psi'$ ist.

Andererseits ist nach (6) auch $e'c', e'c'\eta'$ eine Basis von $\psi'$, und folglich lassen sich die ganzen rationalen Zahlen $p,q,r,s$ so bestimmen, dass

\weq{12}{\begin{aligned}
\beta_1 &= \frac{Qec}{\xi} = e'c'(r\eta'+s)\\
\beta_2 &= \frac{Qec}{\xi}\,\omega = e'c'(p\eta'+q),
\end{aligned}}

und zugleich

\weq{13}{ps-qr = \pm 1.}

Aus (12) folgt aber durch Division:

\weq{14}{\omega = \frac{p\eta'+q}{r\eta'+s},}

d. h. $\omega$ ist mit $\eta'$ äquivalent.

Ersetzen wir nun $\varphi$ durch die äquivalente Form $\varphi_1$, so besteht eine Gleichung wie (5):

\[ \varphi_1\psi = \varepsilon_1\xi_1, \]

worin $\psi$ ungeändert geblieben ist. Folglich ist auch $\omega_1$ mit $\eta$ äquivalent, und mithin sind $\omega$ und $\omega_1$ unter einander äquivalent. % [sic] source prints $\eta$; repeating the derivation with unchanged $\psi$ requires $\eta'$
Das ist aber die in unserem Satze 3. ausgesprochene Behauptung, die also hiermit erwiesen ist.

Hiernach entspricht also jeder Idealclasse des Körpers $\Omega$ eine Classe äquivalenter Zahlen $\omega$ und umgekehrt, und wir können daher noch den Satz hinzufügen:
%CONT>>


% ===== scan leaf 634, printed 613 =====
% !PRINTED 613
\weberpage{613}
\pdfbookmark[2]{§ 168. Composition der quadratischen Irrationalzahlen}{bm192}
%<<CONT

4. \gesperrt{Die Anzahl der nicht äquivalenten Irrationalzahlen der Discriminante $\varDelta$ ist gleich der Anzahl der Formenclassen des Körpers $\Omega$}\footnote{Nach der am Schlusse des \S.\,128 des ersten Bandes gemachten Bemerkung über die \gesperrt{Gauss}'sche Theorie der quadratischen Formen stimmt die Classenzahl der quadratischen Formen nach der \gesperrt{Gauss}'schen Definition nicht immer überein mit der Anzahl der Idealclassen des Körpers $\Omega$. Um die Uebereinstimmung herzustellen, muss man in manchen Fällen die Idealclassen noch weiter in je zwei Classen theilen. Bei unserer Theorie der Aequivalenz der Zahlen $\omega$, wo die eigentliche nicht von der uneigentlichen Aequivalenz unterschieden wird, ist diese weitere Eintheilung nicht erforderlich. Vergl. \gesperrt{Dirichlet-Dedekind}, Zahlentheorie. 4. Auflage, S. 578, 584.}\gesperrt{.}

\begin{center}
§. 168.

\gesperrt{Composition der quadratischen Irrationalzahlen.}
\end{center}

Wir haben im \S.\,166 gesehen, dass zu jedem ganzen Functional $\varphi$ im Körper $\Omega$ eine gewisse quadratische Irrationalzahl $\omega$ von der Discriminante $\varDelta$ gefunden werden kann.

Wir müssen aber jetzt noch untersuchen, inwiefern diese Zahl $\omega$ durch den Körper $\Omega$ und durch das Functional $\varphi$ bestimmt ist.

Nach \S.\,166, (1) ist die Basis von $\varphi$

\weq{1}{\alpha_1 = a_{1,1},\quad \alpha_2 = a_{1,2}+a_{2,2}\Theta}

dadurch definirt, dass $a_{1,1}$ die kleinste positive ganze rationale Zahl ist, die durch $\varphi$ theilbar ist, und $a_{2,2}$ ist dann durch

\weq{2}{N_a(\varphi) = a_{1,1}a_{2,2}}

bestimmt, also sind die beiden Zahlen $a_{1,1},a_{2,2}$ durch $\varphi$ eindeutig bestimmt. Es fragt sich noch, inwiefern auch $a_{1,2}$ bestimmt ist. Die Zahl $a_{1,2}$ ist aber nach dem Modul $\varphi$ durch die Congruenz $a_{1,2}+a_{2,2}\Theta\equiv 0$ gegeben, und da $a_{1,2}$ auch eine ganze rationale Zahl sein muss, so ist $a_{1,2}$ bis auf ein Vielfaches von $a_{1,1}$ bestimmt. Dieses Vielfache von $a_{1,1}$ bleibt aber in der That willkürlich, da, wenn $\alpha_1,\alpha_2$ eine Basis von $\varphi$ ist, dasselbe für jedes ganze rationale $k$ auch von $\alpha_1,\alpha_2+k\alpha_1$ gilt.

Da nun $\alpha_2:\alpha_1=\omega$ die gesuchte quadratische Irrationalzahl ist, so erhalten wir den Satz:
%CONT>>


% ===== scan leaf 635, printed 614 =====
% !PRINTED 614
\weberpage{614}
%<<CONT

5. \gesperrt{Zu jedem ganzen Functional des K\"orpers $\Omega$ geh\"ort eine und nur eine Reihe von quadratischen Irrationalzahlen der Discriminante $\Delta$, und die Zahlen dieser Reihe unterscheiden sich um beliebige ganze rationale Zahlen und sind folglich unter einander \"aquivalent. Associirte Formen f\"uhren auf dieselbe Reihe.}

Wenn $\varphi_1$, $\varphi_2$ zwei ganze Functionale in $\Omega$ bedeuten, und
\[
\varphi_1 \varphi_2 = \varphi_3
\]
gesetzt wird, so k\"onnen wir zu jeder dieser drei Formen eine quadratische Irrationalzahl $\omega_1$, $\omega_2$, $\omega_3$ bestimmen, und $\omega_3$ ist durch $\omega_1$, $\omega_2$ bis auf eine additive ganze rationale Zahl bestimmt. Man nennt dann $\omega_3$ aus $\omega_1$ und $\omega_2$ \gesperrt{componirt} oder zusammengesetzt, und setzt in symbolischer Bezeichnung
\weq{3}{\omega_1 \omega_2 = \omega_3.}

Ist $\psi_1$ \"aquivalent mit $\varphi_1$, $\psi_2$ mit $\varphi_2$, so ist nach \S.\,152 auch $\psi_1 \psi_2 = \psi_3$ mit $\varphi_3$ \"aquivalent.

Wenn nun $\eta_1$, $\eta_2$, $\eta_3$ die durch den Satz 5. zu den Formen $\psi_1$, $\psi_2$, $\psi_3$ geh\"origen Irrationalzahlen sind, so bilden nach 3., \S.\,167 die Zahlen $\eta_1$, $\omega_1$; $\eta_2$, $\omega_2$; $\eta_3$, $\omega_3$ drei Paare \"aquivalenter Zahlen und wir k\"onnen daher folgenden Satz aussprechen:

6. \gesperrt{Ersetzt man in einer Composition von quadratischen Irrationalzahlen jedes Element durch eine \"aquivalente Zahl, so geht auch die componirte Zahl in eine \"aquivalente \"uber.}

Dadurch sind die Classen der quadratischen Irrationalzahlen einer Discriminante $\Delta$ zu einer mit der Gruppe der Idealclassen isomorphen \gesperrt{Abel'schen} Gruppe verbunden.

Um die aus zwei gegebenen Zahlen componirte Zahl wirklich zu finden, hat man so zu verfahren.

Wir wollen die nach (1) gebildeten Basen der Formen $\varphi_1$, $\varphi_2$, $\varphi_3$ so bezeichnen:
\[
\begin{array}{lll}
\varphi_1) & \alpha_1 = a_{1,1} & \alpha_2 = a_{1,2} + a_{2,2}\,\Theta \\
\varphi_2) & \beta_1 = b_{1,1} & \beta_2 = b_{1,2} + b_{2,2}\,\Theta \\
\varphi_3) & \gamma_1 = c_{1,1} & \gamma_2 = c_{1,2} + c_{2,2}\,\Theta,
\end{array}
\]
so dass
\[
\omega_1 = \frac{\alpha_2}{\alpha_1}, \quad \omega_2 = \frac{\beta_2}{\beta_1}, \quad \omega_3 = \frac{\gamma_2}{\gamma_1}
\]
%CONT>>


% ===== scan leaf 636, printed 615 =====
% !PRINTED 615
\weberpage{615}

%<<CONT
wird. Es ist dann $c_{1,1}$ die kleinste positive ganze rationale Zahl, die durch $\varphi_1 \varphi_2$ theilbar ist, und diese muss jedenfalls ein Theiler von $a_{1,1} b_{1,1}$ sein. Hat man
\[
a_{1,1} b_{1,1} = k\, c_{1,1},
\]
wo $k$ eine positive ganze rationale Zahl ist, so ist, weil $N_a(\varphi_1)\, N_a(\varphi_2) = N_a(\varphi_3)$ sein muss,
\[
k\, a_{2,2} b_{2,2} = c_{2,2}.
\]
Die letzte Zahl $c_{1,2}$ ist dann aus der Congruenz
\weq{4}{c_{1,2} + c_{2,2}\,\Theta \equiv 0 \ (\mathrm{mod}\ \varphi_1 \varphi_2)}
zu bestimmen.

Wenn es nur darauf ankommt, einen Repr\"asentanten der zusammengesetzten Classe zu finden, so verf\"ahrt man am einfachsten so.

Ist $\varphi_1$ beliebig angenommen, so kann man (nach \S.\,152, 6.) $\varphi_2$ in seiner Classe so w\"ahlen, dass es relativ prim zu $a_{1,1}$ wird. Dann ist auch $\varphi_1$ relativ prim zu $\varphi_2$ und $b_{1,1}$ relativ prim zu $a_{1,1}$.

Denn haben $a_{1,1}$ und $b_{1,1}$ einen gr\"ossten gemeinschaftlichen Theiler $d$, so ist dieser auch relativ prim zu $\varphi_2$, und also ist $b_{1,1} : d$ durch $\varphi_2$ theilbar; dies aber widerspricht, wenn $d > 1$ ist, der Definition von $b_{1,1}$. Ebenso folgt auch umgekehrt, dass, wenn $a_{1,1}$ und $b_{1,1}$ relativ prim sind, $\varphi_2$ zu $a_{1,1}$ relativ prim ist. Denn wenn keiner der rationalen Primfactoren von $b_{1,1}$ in $a_{1,1}$ aufgeht, so kann auch keiner der in $\varphi_2$ aufgehenden Primfactoren in $a_{1,1}$ aufgehen.

In diesem Falle ist nun:
\weq{5}{a_{1,1} b_{1,1} = c_{1,1}, \quad a_{2,2} b_{2,2} = c_{2,2}.}
Die Congruenz (4) wird dann
\weq{6}{c_{1,2} + a_{2,2} b_{2,2}\,\Theta \equiv 0 \ (\mathrm{mod}\ \varphi_1 \varphi_2),}
und da
\[
a_{1,2} + a_{2,2}\,\Theta \equiv 0 \ (\mathrm{mod}\ \varphi_1), \quad b_{1,2} + b_{2,2}\,\Theta \equiv 0 \ (\mathrm{mod}\ \varphi_2)
\]
ist, so zerf\"allt (6) in die beiden Congruenzen:
\[
\begin{aligned}
c_{1,2} &\equiv b_{2,2} a_{1,2} \ (\mathrm{mod}\ \varphi_1) \\
&\equiv a_{2,2} b_{1,2} \ (\mathrm{mod}\ \varphi_2).
\end{aligned}
\]
Diese Congruenzen sind aber nach \S.\,138, 13. gleichwerthig mit den gew\"ohnlichen Zahlencongruenzen
\weq{7}{
\begin{aligned}
c_{1,2} &\equiv b_{2,2} a_{1,2} \ (\mathrm{mod}\ a_{1,1}) \\
&\equiv a_{2,2} b_{1,2} \ (\mathrm{mod}\ b_{1,1}),
\end{aligned}}
wodurch $c_{1,2}$ nach dem Modul $c_{1,1}$ bestimmt ist.


% ===== scan leaf 637, printed 616 =====
% !PRINTED 616
\weberpage{616}

Es bleibt nur noch \"ubrig, zu zeigen, wie man $\varphi_2$ der hier gestellten Forderung gem\"ass bestimmen kann. Da zwei Formen, die sich nur durch einen Zahlenfactor von einander unterscheiden, \"aquivalent sind, so ist $\varphi_2$ \"aquivalent mit
\[
c_2\, (t_1 + \omega_2 t_2),
\]
[\S.\,166, (10)] und nehmen wir diese Form f\"ur $\varphi_2$, so ist $b_{1,1} = c_2$. Die Zahl $\omega_2$ kann man aber unter den mit ihr \"aquivalenten immer so w\"ahlen, dass $c_2$ zu $a_{1,1}$ relativ prim wird, und dann ist auch $a_{1,1}$ zu $\varphi_2$ relativ prim.

Ersetzt man n\"amlich $\omega_2$ durch
\[
\frac{p\, \omega_2 + q}{r\, \omega_2 + s}, \quad p s - q r = \pm 1,
\]
so geht $c_2$ nach Bd.\ I, \S.\,122, (13) in
\[
c_2' = -a_2 r^2 + b_2 r s + c_2 s^2
\]
\"uber, und nun kann man \"uber die ganzen rationalen Zahlen $r, s$ so verf\"ugen, dass sie unter einander theilerfremd sind und dass $c_2'$ durch irgend welche gegebene Primzahlen nicht theilbar ist.

Denn ist $n$ eine solche Primzahl, so nehme man, wenn $n$ in $a_2$ und $c_2$, also nicht in $b_2$ aufgeht, $r$ und $s$ durch $n$ untheilbar, wenn $n$ in $a_2$ nicht aufgeht, $s$ durch $n$ theilbar und $r$ durch $n$ untheilbar, und wenn $n$ in $c_2$ nicht aufgeht, $r$ durch $n$ theilbar, $s$ durch $n$ untheilbar. (Sind $a_2$ und $c_2$ beide nicht durch $n$ theilbar, so hat man zwischen den beiden letzteren Annahmen die Wahl.) Diese Forderungen sind f\"ur verschiedene beliebig gegebene Primzahlen $n$ mit einander vertr\"aglich, und dann l\"asst sich $p, q$ aus der Bedingung $p s - q r = \pm 1$ bestimmen.

Diese Gesetze der Composition der quadratischen Irrationalzahlen haben wir hier nur unter der Voraussetzung abgeleitet, dass die Discriminante zugleich die Grundzahl eines quadratischen K\"orpers ist, d.\,h. dass $\Delta$ Stammdiscriminante ist. Dieselben Gesetze gelten aber auch, mit geringen Modificationen, allgemein $^{1)}$.

\noindent\rule{4cm}{0.4pt}

$^{1)}$ \gesperrt{Gauss}, Disquisitiones arithm., Art.\,234 ff. \gesperrt{Dirichlet-Dedekind}, Zahlentheorie. 4. Auflage, \S.\,187.

\begin{center}
\rule{0.2\linewidth}{0.4pt}
\end{center}


% ===== scan leaf 638, printed 617 =====
\weberpage{617}% page number not printed on this page
\pdfbookmark[1]{Einundzwanzigster Abschnitt. Kreistheilungskörper.}{bm193}
\pdfbookmark[2]{§ 169. Zerlegung der Primzahl q im Kreistheilungskörper Ωq}{bm194}
% !PRINTED none
% Abschnitt-title page (part divider): no running head/page number visible in
% the scan for this page; §. 169 opens on it.

\begin{center}
\gesperrt{Einundzwanzigster Abschnitt.}

\gesperrt{Kreistheilungsk\"orper.}
\end{center}

\begin{center}
\rule{4cm}{0.4pt}
\end{center}

\begin{center}
\S.\,169.

\gesperrt{Zerlegung der Primzahl $q$ in Primfactoren im Kreistheilungsk\"orper $\Omega_{q^{\varkappa}}$.}
\end{center}

Wir machen eine zweite Anwendung der allgemeinen Theorie der algebraischen Zahlen auf die Kreistheilungsk\"orper, und gewinnen dadurch das Mittel, die Theorie der \gesperrt{Abel'schen} Zahlk\"orper \"uberhaupt zu einem sch\"onen Abschlusse zu bringen.

Die Betrachtungen, die wir zun\"achst anzustellen haben, lassen sich mit kleinen Modificationen auf jeden vollen Kreistheilungsk\"orper (\S.\,18) anwenden. Der Einfachheit halber beschr\"anken wir uns hier aber auf den Fall, dass der Grad der Einheitswurzel eine Primzahlpotenz ist, der f\"ur die beabsichtigte Anwendung ausreicht.

Es sei $q$ eine nat\"urliche Primzahl (mit Einschluss von 2) und
\weq{1}{m = q^{\varkappa}}
eine Potenz von $q$, deren positiver Exponent $\varkappa$ f\"ur $q=2$ gr\"osser als 1 vorausgesetzt wird.

Es sei ferner $r$ eine primitive $m^{\text{te}}$ Einheitswurzel und $\Omega_m$ der volle Kreistheilungsk\"orper f\"ur den Exponenten $m$, dessen Grad
\weq{2}{\mu = \varphi(m) = q^{\varkappa-1}(q-1)}
ist (Bd. I, \S.\,134).

Wir bezeichnen durchweg mit $n$ die $\mu$ Zahlen eines vollen Restsystems f\"ur den Modul $m$ mit Ausschluss der durch $q$ theilbaren Zahlen, und es sind dann die $\mu$ Zahlen $r^n$ die Wurzeln der irreduciblen Gleichung $\mu^{\text{ten}}$ Grades $f(x)=0$, worin
\weq{3}{f(x) = \dfrac{x^m-1}{x^{\frac{m}{q}}-1} = x^{q^{\varkappa-1}(q-1)} + x^{q^{\varkappa-1}(q-2)} + \cdots + 1}
%CONT>>


% ===== scan leaf 639, printed 618 =====
% !PRINTED 618
\weberpage{618}
%<<CONT
zu setzen ist. $r$ ist daher eine \gesperrt{ganze Zahl} in $\Omega_m$, und ihre Norm ist $=+1$; folglich ist $r$ eine Einheit.

$\Omega_m$ ist ein Normalk\"orper, der durch die $\mu$ Substitutionen
\[ s_n = (r, r^n) \]
in sich selber \"ubergeht.

Diese $\mu$ Substitutionen $s_n$ bilden eine \gesperrt{Abel'sche} Gruppe $\mathfrak{N}$, welche die \gesperrt{Galois'sche} Gruppe des K\"orpers $\Omega_m$ ist. Sie ist isomorph mit der gleichfalls durch $\mathfrak{N}$ zu bezeichnenden Gruppe der nach dem Modul $m$ genommenen Zahlclassen $n$ (\S.\,16).

Die Function $f(x)$ l\"asst sich in die $\mu$ Factoren $x - r^n$ zerlegen, und wir setzen daher
\weq{4}{f(x) = \prod\limits^{n} (x - r^n).}

Setzen wir darin $x=1$, und beachten, dass [nach (3)] $f(1) = q$ ist, so folgt
\weq{5}{q = \prod\limits^{n} (1 - r^n).}

Die Zahlen
\weq{6}{\sigma_n = 1 - r^n, \quad \sigma = 1 - r}
sind aber ganze Zahlen des K\"orpers $\Omega_m$, und die Formel (5) zeigt zun\"achst, \gesperrt{dass $q$ im K\"orper $\Omega_m$ in $\mu$ Factoren $\sigma_n$ zerlegbar ist.}

Wir beweisen zun\"achst, dass die $\mu$ Zahlen $\sigma_n$ mit einander associirt sind. Bedeuten $n, n'$ zwei durch $q$ nicht theilbare Zahlen, so k\"onnen wir eine nat\"urliche Zahl $a$ so bestimmen, dass
\[ n' \equiv a n \ (\mathrm{mod}\ m) \]
wird. Dann ist aber
\[ \frac{\sigma_{n'}}{\sigma_n} = \frac{1 - r^{an}}{1 - r^n} = 1 + r^n + r^{2n} + \cdots + r^{(a-1)n} \]
eine ganze Zahl, also $\sigma_{n'}$ durch $\sigma_n$ theilbar. Da hierin $n$ mit $n'$ vertauscht werden kann, so ist auch $\sigma_n$ durch $\sigma_{n'}$ theilbar, also sind beide Zahlen associirt (\S.\,138). Wenn daher $\varepsilon$ eine Einheit bedeutet, so ist nach (5):
\weq{7}{q = \varepsilon \sigma^{\mu}, \quad N(\sigma) = q.}

Es ist also die nat\"urliche Primzahl $q$ associirt mit der $\mu^{\text{ten}}$ Potenz einer ganzen Zahl $\sigma$ im K\"orper $\Omega_m$. Diese Zahl $\sigma$ ist aber auch im K\"orper $\Omega_m$ noch Primzahl, und zwar vom ersten Grade. Denn hat $\sigma$ irgend einen Theiler $\sigma'$, so muss die Norm von $\sigma'$ ein Theiler der Norm von $\sigma$ sein, also, wenn $\sigma'$ keine
%CONT>>


% ===== scan leaf 640, printed 619 =====
% !PRINTED 619
\weberpage{619}
%<<CONT
Einheit ist, so ist $N(\sigma') = q$. Folglich ist die Norm von $\sigma : \sigma'$ gleich 1, d.\,h. $\sigma : \sigma'$ ist eine Einheit und $\sigma$ mit $\sigma'$ associirt. Wir haben also den ersten Satz:

\begin{quote}
I. \gesperrt{Die nat\"urliche Primzahl $q$ ist in dem vollen Kreistheilungsk\"orper $\Omega_m$ mit der $\mu^{\text{ten}}$ Potenz eines Primfactors ersten Grades associirt.}
\end{quote}

Wenn wir mit $m_1$ einen Theiler von $m$ bezeichnen, kleiner als $m$ und gr\"osser als 1, und
\weq{8}{m = m_1 m_2}
setzen, so dass $m_1$ und $m_2$ selbst Potenzen von $q$ sind, so ergiebt sich die Zerlegung:
\[
% [sic: The printed product needs the factor (-1)^{m_1-1} on its right-hand side. For even m_1 (the permitted q=2 case), the product as printed equals 1-x^{m_1}, not x^{m_1}-1.]
x^{m_1} - 1 = (x-1)(xr^{m_2}-1)(xr^{2m_2}-1)\ldots(xr^{(m_1-1)m_2}-1),
\]
und wenn wir darin $x = r$ setzen, nach (6):
% [sic: The prefactor in (9) should be -1 for every m_1. The printed (-1)^{m_1} repeats the preceding sign error when q=2; for m=4,m_1=2,r=i, the two printed sides are -2 and +2. The following norm identity (10) is unaffected because mu is even.]
\weq{9}{r^{m_1} - 1 = (-1)^{m_1}\prod_{0,\,m_1-1}^{s}\sigma_{1+sm_2},}
und daraus, da $\mu$ immer gerade ist,
\weq{10}{N(r^{m_1}-1) = q^{m_1}.}

Hiernach l\"asst sich leicht die Discriminante $\Delta$ der Gleichung (3) bilden, die nach Bd.\,I, \S.\,47 den Ausdruck hat:
\[
\Delta = (-1)^{\frac{\mu(\mu-1)}{2}} N f'(r).
\]

Es ist n\"amlich nach (3):
\[
f'(r) = \dfrac{m\,r^{m-1}}{r^{\frac{m}{q}} - 1},
\]
also, da die Norm von $r$ gleich 1 ist:
\[
N f'(r) = \dfrac{m^{\mu}}{q^{\frac{m}{q}}},
\]
und folglich (da $\mu$ immer gerade ist):
\weq{11}{\Delta = (-1)^{\frac{\mu}{2}} q^{q^{\varkappa-1}[\varkappa(q-1)-1]}.}

Die Discriminante $\Delta$ ist also, vom Vorzeichen abgesehen, eine Potenz von $q$.


% ===== scan leaf 641, printed 620 =====
% !PRINTED 620
\weberpage{620}
\pdfbookmark[2]{§ 170. Minimalbasis des Körpers Ωm}{bm195}

\begin{center}
\S. 170.

\gesperrt{Minimalbasis des K\"orpers $\Omega_m$.}
\end{center}

Die S\"atze, die im vorigen Paragraphen bewiesen sind, f\"uhren uns zu dem folgenden Theorem:

II. \gesperrt{Die Zahlen}
\weq{1}{1, r, r^2, \ldots, r^{\mu-1}}
\gesperrt{bilden eine Basis des Systemes $\mathfrak{o}$ der ganzen Zahlen in $\Omega_m$.}

Um dies zu zeigen, gen\"ugt nach \S.\,145 der Nachweis, dass die ganze Zahl in $\Omega_m$
\weq{2}{\omega = x_0 + x_1 r + x_2 r^2 + \cdots + x_{\mu-1} r^{\mu-1},}
worin $x_0, x_1, \ldots, x_{\mu-1}$ ganze rationale Zahlen sind, nur dann durch eine rationale Primzahl $p$ theilbar sein kann, wenn die Zahlen $x_0, x_1, \ldots, x_{\mu-1}$ alle durch $p$ theilbar sind.

Nehmen wir also an, es sei $\omega$ durch $p$ theilbar; dann sind auch alle Zahlen $\omega_n$, die aus $\omega$ durch eine der $\mu$ Substitutionen $(r, r^n)$ entstehen, durch $p$ theilbar, und wir erhalten also aus (2) ein System von $\mu$ Gleichungen, die in Bezug auf die $x_i$ linear sind:
\weq{3}{\omega_n = x_0 + x_1 r^n + x_2 r^{2n} + \cdots + x_{\mu-1} r^{(\mu-1)n}.}

Die Determinante dieses Systemes, d.\,h. die aus den $\mu$ Reihen
\[ 1, r^n, r^{2n}, \ldots, r^{(\mu-1)n} \]
gebildete Determinante ist aber nach Bd. I, \S.\,46 gleich der Quadratwurzel $\sqrt{\varDelta}$, und wenn wir also das System der Gleichungen (3) aufl\"osen, so folgt, dass die s\"ammtlichen rationalen Zahlen $\varDelta x_i$ durch $p$ theilbar sein m\"ussen.

Ist nun $p$ nicht $=q$, so ist $\varDelta$ nicht durch $p$ theilbar, und s\"ammtliche $x_i$ m\"ussen durch $p$ theilbar sein.

Ist aber $p = q$, so setzen wir
\[ f(t) = x_0 + x_1 t + x_2 t^2 + \cdots + x_{\mu-1} t^{\mu-1}, \]
so dass $\omega_n = f(r^n)$ wird, und setzen darin $r = 1 - \sigma$ [\S.\,169, (6)]. Dann ist nach der \gesperrt{Taylor'schen} Entwickelung:
\weq{4}{\begin{aligned}
\omega &= f(1-\sigma) = f(1) - \sigma f'(1) + \dfrac{\sigma^2}{1.2} f''(1) - \cdots\\
&\qquad - \dfrac{\sigma^{\mu-1}}{\Pi(\mu-1)} f^{(\mu-1)}(1);
\end{aligned}}
%CONT>>


% ===== scan leaf 642, printed 621 =====
% !PRINTED 621
\weberpage{621}
%<<CONT
und hierin sind die Co\"efficienten
\[
f(1), f'(1), \tfrac12 f''(1), \cdots, \dfrac{1}{\Pi(\mu-1)} f^{(\mu-1)}(1)
\]
ganze rationale Zahlen, n\"amlich
\begin{align*}
\dfrac{1}{\Pi(\mu-1)} f^{(\mu-1)}(1) &= x_{\mu-1} \\[6pt]
\dfrac{1}{\Pi(\mu-2)} f^{(\mu-2)}(1) &= x_{\mu-2}+(\mu-1)x_{\mu-1} \\[6pt]
\dfrac{1}{\Pi(\mu-3)} f^{(\mu-3)}(1) &= x_{\mu-3}+(\mu-2)x_{\mu-2}+\dfrac{(\mu-1)(\mu-2)}{1.2}x_{\mu-1} \tag{5}\\[6pt]
\cdots\cdots\cdots &\cdots\cdots\cdots\cdots\cdots\cdots\cdots \\[6pt]
f(1) &= x_0+x_1+\cdots\cdots\cdots\cdots+x_{\mu-1}.
\end{align*}
Da nun $\omega$ durch $q$ und folglich durch $\sigma^\mu$ theilbar ist, so
muss $f(1)$ nach (4) durch $q$ theilbar sein. Da aber $\mu>1$ und
folglich $q:\sigma$ noch durch $\sigma$ theilbar ist, so folgt, dass auch
\[
f'(1) - \dfrac{\sigma}{1.2} f''(1) \ldots
\]
durch $\sigma$, folglich $f'(1)$ durch $q$ theilbar ist, und wenn man so
weiter schliesst, ergiebt sich, dass die s\"ammtlichen rationalen
Zahlen (5) durch $q$ theilbar sind, und folglich sind auch die
$x_{\mu-1}, x_{\mu-2}, \ldots, x_0$ durch $q$ theilbar. Damit ist der Satz II. vollst\"andig bewiesen.

Anstatt der Reihe (1) kann man auch eine beliebige Reihe von
$\mu$ auf einander folgenden Potenzen von $r$ als Basis von $\mathfrak{o}$ nehmen:
\weq{6}{r^a, r^{a+1}, r^{a+2}, \ldots, r^{a+\mu-1},}
die sich ergiebt, wenn man alle Glieder der Reihe (1) mit $r^a$
multiplicirt.

Die Richtigkeit dieser Bemerkung ergiebt sich daraus, dass
die $r^{\pm a}$ ganze Zahlen sind, und dass also $\omega$ und $r^a\omega$ stets gleichzeitig ganze oder gebrochene Zahlen sind.

Beispielsweise bilden die Zahlen
\weq{7}{r^{-\frac12\mu+1}, r^{-\frac12\mu+2}, \ldots, r^{-1}, 1, r, r^2, \ldots, r^{\frac12\mu}}
eine Basis von $\mathfrak{o}$.

Damit ist auch die Grundzahl $\Delta$ des K\"orpers $\Omega_m$ bestimmt.
Sie ist nach der allgemeinen Definition (\S.\,145) gleich der Discirminante % [sic] "Discirminante" as printed; "Discriminante" is meant
\[
\Delta(1, r, r^2, \ldots, r^{\mu-1}),
\]
d.\,h. gleich der in \S.\,169 bestimmten Discriminante der Kreistheilungsgleichung.


% ===== scan leaf 643, printed 622 =====
% !PRINTED 622
\weberpage{622}
\pdfbookmark[2]{§ 171. Die Primideale im Körper Ωm}{bm196}
\begin{center}
\S.\,171.

\gesperrt{Die Primideale im K\"orper $\Omega_m$.}
\end{center}

% [sic: The printed reference 164 should be 169, theorem I and (7), where the prime-power factorisation is established.]
% [Source shorthand: q is associated with sigma^mu, not generally equal to it as an element; the corresponding principal ideals are equal.]
Wir haben im \S.\,164 gesehen, dass die nat\"urliche Primzahl $q$
die $\mu^{\text{te}}$ Potenz einer Primzahl $\sigma$ im K\"orper $\Omega_m$ ist.

Um alle Primideale, die im K\"orper $\Omega_m$ existiren, zu ermitteln, sind also noch s\"ammtliche von $q$ verschiedene nat\"urliche
Primzahlen $p$ in Primfactoren zu zerlegen, und es sind darunter
die nicht associirten auszusuchen.

Die Grundlage f\"ur diese Untersuchung bilden die S\"atze des
\S.\,162, wenn der dort mit $R$ bezeichnete K\"orper durch den
K\"orper der rationalen Zahlen ersetzt wird.

Jede Zahl $\omega$ des K\"orpers $\Omega_m$ geht durch eine der $\mu$ Substitutionen
\[
s_n = (r, r^n)
\]
in eine bestimmte andere Zahl $\omega_n$ \"uber, die gleichfalls in $\Omega_m$
enthalten ist. Ist
\weq{1}{\omega = \omega_1 = a_0 + a_1 r + a_2 r^2 + \cdots + a_{\mu-1} r^{\mu-1},}
so ist
\weq{2}{\omega_n = a_0 + a_1 r^n + a_2 r^{2n} + \cdots + a_{\mu-1} r^{(\mu-1)n}.}
Wenn eine dieser Zahlen $\omega_n$ eine ganze Zahl ist, wie wir
jetzt annehmen wollen, so sind es auch alle anderen, und dies
tritt immer dann und nur dann ein, wenn die rationalen Zahlen
$a_0, a_1, \ldots, a_{\mu-1}$ ganz sind.

Sind $h, k$ irgend zwei Exponenten, so ist
\[
r^h - r^k = r^k(r^{h-k} - 1),
\]
und dies ist nach \S.\,169, (9) mit einer Potenz von $\sigma$ associirt,
also, wenn $\mathfrak{p}$ ein in $p$ aufgehendes Primideal bedeutet, durch $\mathfrak{p}$
nicht theilbar, ausser wenn $h \equiv k \pmod{m}$ ist. Daraus ergiebt
% [sic: The following assertion needs congruence for every integral omega, not for one fixed omega (omega=1 is a counterexample). Choosing omega=r proves the required assertion about the inertia group X, so the unramifiedness conclusion remains valid.]
sich, dass zwei Zahlen $\omega_h, \omega_k$, nur dann nach dem Modul $\mathfrak{p}$ congruent sein k\"onnen, wenn $h \equiv k \pmod{m}$ ist, und dass also die
Gruppe $X$ des \S.\,162, die der Bedingung
\[
\omega\,|\,\chi \equiv \omega \pmod{\mathfrak{p}}
\]
gen\"ugt, nur die \gesperrt{identische Substitution} enth\"alt. Daraus folgt
aber, dass $g = 1$, d.\,h. dass $p$ \gesperrt{nicht durch das Quadrat
eines Primideales theilbar ist.}


% ===== scan leaf 644, printed 623 =====
% !PRINTED 623
\weberpage{623}
\pdfbookmark[2]{§ 172. Die conjugirten Primideale}{bm197}
Bilden wir aus (2) die $p^{\text{te}}$ Potenz von $\omega$, und beachten den
\gesperrt{Fermat'schen} Lehrsatz f\"ur rationale Zahlen $[a_0^p \equiv a_0 \pmod{p}]$,
so ergiebt sich
\[
\omega^p \equiv a_0 + a_1 r^p + a_2 r^{2p} + \cdots + a_{\mu-1} r^{(\mu-1)p} \pmod{p},
\]
oder nach (2)
\weq{3}{\omega^p \equiv \omega_p \pmod{p}.}

Dadurch ist die Gruppe $\Psi$ bestimmt, zu der das Ideal $\mathfrak{p}$
geh\"ort. Es ist n\"amlich nach (3) auch
\[
\omega^p \equiv \omega_p \pmod{\mathfrak{p}},
\]
daher ist $\psi_0 = (r, r^p)$ und die Gruppe $\Psi$ besteht nach \S.\,162, (13)
aus den Potenzen dieser Substitution, soweit sie von einander
verschieden sind.

Ist daher $f$ der kleinste positive Exponent, f\"ur den
\weq{4}{p^f \equiv 1 \pmod{m}}
ist, d.\,h. geh\"ort $p$ zu dem Exponenten $f$ f\"ur den Modul $m$, so ist
\weq{5}{\Psi = (r,r), (r,r^p), (r,r^{p^2}), \ldots (r,r^{p^{f-1}})}
und ist vom $f^{\text{ten}}$ Grade. Demnach ist auch $\mathfrak{p}$ ein Primideal
$f^{\text{ten}}$ Grades und
\weq{6}{N(\mathfrak{p}) = p^f.}

Die Anzahl $e$ der von einander verschiedenen conjugirten
Primfactoren von $f$ ist $\mu : f$, und wir haben also den Satz: % [sic] printed $f$ after "von"; mathematically $p$ is required

III. \gesperrt{Ist $p$ eine von $q$ verschiedene Primzahl, die f\"ur den Modul $m$ zum Exponenten $f$ geh\"ort, ist ferner $\mu = \varphi(m) = ef$, so zerf\"allt $p$ in $\Omega_m$ in $e$ von einander verschiedene conjugirte Primfactoren $f^{\text{ten}}$ Grades.}

\begin{center}
\S.\,172.

\gesperrt{Die conjugirten Primideale.}
\end{center}

Wir m\"ussen noch etwas genauer auf die Bildung der conjugirten Primfactoren $\mathfrak{p}_n$ einer von $q$ verschiedenen Primzahl $p$
eingehen, die nach dem zuletzt bewiesenen Satze in $\Omega_m$ in $e$
von einander verschiedene Primfactoren zerf\"allt, wenn
\weq{1}{\varphi(m) = ef}
ist, und $f$ der kleinste positive Exponent, f\"ur den
\weq{2}{p^f \equiv 1 \pmod{m}.}


% ===== scan leaf 645, printed 624 =====
% !PRINTED 624
\weberpage{624}
Die Gruppe $\Psi$ ist isomorph mit einer in $\mathfrak{N}$ (\S.\,169) enthaltenen Gruppe nach dem Modul $m$ genommener ganzer
rationaler Zahlen, die wir mit $\mathfrak{A}$ bezeichnen, die aus allen, einer
Congruenz
\weq{3}{a \equiv p^h \pmod{m}}
gen\"ugenden Zahlen $a$ besteht, und die wir daher symbolisch durch
\weq{4}{\mathfrak{A} \equiv p^h \pmod{m}}
darstellen k\"onnen, wenn $h$ einen beliebigen ganzzahligen Exponenten bedeutet (vergl. \S.\,18).

Wir k\"onnen also, wenn wir die Zahlen $\xi_1, \xi_2, \ldots, \xi_e$ aus $\mathfrak{N}$
passend ausw\"ahlen:
\weq{5}{\mathfrak{N} = \mathfrak{A}\xi_1 + \mathfrak{A}\xi_2 + \mathfrak{A}\xi_3 + \cdots + \mathfrak{A}\xi}
% [sic] source prints the final \xi without the expected subscript e; (6)--(7) supply the context
setzen, und jeder dieser Nebengruppen entspricht eines der conjugirten Functionale $\mathfrak{p}_{\xi_1}, \mathfrak{p}_{\xi_2}, \ldots, \mathfrak{p}_{\xi_e}$, die alle von einander verschieden sind und deren Product mit $p$ associirt ist.
% [sic] "Functionale" as printed; context elsewhere on the page uses "Primfactoren"/"Primideale" for this notion

Wir setzen also
\weq{6}{p = \mathfrak{p}_{\xi_1} \mathfrak{p}_{\xi_2} \ldots \mathfrak{p}_{\xi_e}.}
Um aber das Zahlensystem
\weq{7}{\xi_1, \xi_2, \ldots, \xi_e,}
auf das es haupts\"achlich ankommt, genauer zu charakterisiren,
bemerken wir, dass wir dazu jedes System von $e$ Zahlen der
Gruppe $\mathfrak{N}$ nehmen k\"onnen, wenn nur keine zwei unter ihnen,
$\xi, \xi'$, eine Congruenz
\[
\xi'\xi^{-1} \equiv p^h \pmod{m}
\]
erf\"ullen.

Um ein solches Zahlensystem zu finden, ist es gut, den Fall
eines ungeraden $m$ von dem anderen zu trennen, in dem $m$ eine
Potenz von 2 ist.

1. Ist zun\"achst $m = q^{\varkappa}$ ungerade, so nehmen wir eine
primitive Wurzel $g$ von $m$ (\S.\,14) und setzen
\weq{8}{p \equiv g^{\gamma} \pmod{m}.}
Denn $f$ ist [nach (1), (2)] die kleinste positive Zahl, die der
Bedingung
\[
\gamma f \equiv 0 \pmod{\mu}
\]
gen\"ugt, und folglich ist $e$ der gr\"osste gemeinschaftliche Theiler
von $\mu$ und $\gamma$, und irgend eine Zahl $n \equiv g^{\lambda}$ ist nur dann mit
%CONT>>


% ===== scan leaf 646, printed 625 =====
% !PRINTED 625
\weberpage{625}
%<<CONT
einer Potenz von $p$ congruent, wenn $\lambda$ durch $e$ theilbar ist. Denn nur dann kann die Congruenz
\[ \lambda \equiv \gamma x \, (\mathrm{mod}\ \mu) \]
befriedigt werden. Setzen wir
\[ \gamma = e \nu, \]
so ist $\nu$ relativ prim zu $f$. Daraus folgt, dass die Reihe der Zahlen
\weq{9}{1, g, g^2, \ldots, g^{e-1}}
für die $\xi$ genommen werden kann. Denn sind $h, h'$ zwei verschiedene Zahlen der Reihe $0, 1, 2, \ldots, e-1$, so ist $h'-h$ niemals durch $e$ theilbar, und folglich
\[ \xi' \xi^{-1} = g^{h'-h} \]
niemals mit einer Potenz von $p$ congruent.

Wir wollen ferner noch die beiden Fälle unterscheiden, dass $p-1$ durch $q$ theilbar ist oder nicht.

a) Ist $p \equiv 1 \, (\mathrm{mod}\ q)$, so ist $\gamma \equiv 0 \, (\mathrm{mod}\ q-1)$, und folglich $e$ theilbar durch $q-1$. Wir können daher
\[ e = \varphi(q^{\varkappa_1}) = q^{\varkappa_1 - 1}(q-1) \]
setzen, worin $0 < \varkappa_1 \leqq \varkappa$ ist.

Wenn $\varkappa_1 = \varkappa$, also $e = \varphi(m)$, $f=1$, und daher $p-1$ durch $m$ theilbar ist, so durchläuft $\xi$ die ganze Gruppe $\mathfrak{N}$, und alle Primfunctionale $\mathfrak{p}_n$ sind von einander verschieden.

Ist $\varkappa_1 < \varkappa$, so ist $\nu$ durch $q$ nicht theilbar, weil sonst $e$ nicht der grösste gemeinschaftliche Theiler von $\mu$ und $\gamma$ wäre, und
\[ p - 1 \equiv g^{\nu\varphi(q^{\varkappa_1})} - 1 \, (\mathrm{mod}\ m) \]
ist durch $q^{\varkappa_1}$ theilbar, aber durch keine höhere Potenz von $q$. Setzen wir also $q^{\varkappa_1} = m_1$ und
\weq{10}{m = m_1 m_2,}
so ist $m_1$ der \gesperrt{grösste gemeinschaftliche Theiler von $p-1$ und $m$}. Die Primzahl $p$ zerfällt im Körper $\Omega_m$ in $\varphi(m_1)$ verschiedene Primfactoren. In ebenso viele Primfactoren zerfällt aber $p$ auch im Körper $\Omega_{m_1}$, der ein Theiler des Körpers $\Omega_m$ ist und aus allen rationalen Functionen der $m_1{}^{\mathrm{ten}}$ Einheitswurzel $r_1 = r^{m_2}$ besteht. Daraus folgt, dass die Primfactoren von $p$ in $\Omega_{m_1}$ auch im Körper $\Omega_m$ nicht weiter zerlegbar sind, und dass wir daher die Primfactoren $p$ in $\Omega_m$ als Functionale des Körpers $\Omega_{m_1}$ darstellen können.
% [sic] Source prints 'Primfactoren p' without 'von'; retain the printed wording.

Es durchläuft $\xi$ ein volles System durch $q$ nicht theilbarer Reste nach dem Modul $m_1$, und die $\mathfrak{p}_\xi$ sind Ideale in $\Omega_{m_1}$.


% ===== scan leaf 647, printed 626 =====
% !PRINTED 626
\weberpage{626}
b) Ist umgekehrt $\gamma$ durch $q-1$ theilbar, so folgt aus (8), dass $p-1$ durch $q$ theilbar ist. Nehmen wir also $p-1$ nicht durch $q$ theilbar an, so ist auch $e$ nicht durch $q-1$ theilbar, und $\xi$ durchläuft die Reihe der Zahlen (9).

Bilden wir die Summe dieser Zahlen $\xi$, so folgt:
\[ (g-1) \Sigma \xi \equiv g^e - 1 \, (\mathrm{mod}\ m), \]
und diese Zahl ist durch $q$ theilbar oder nicht theilbar, je nachdem $e$ durch $q-1$ theilbar oder nicht theilbar ist. Da ausserdem $g-1$ nicht durch $q$ theilbar ist, so ist $\Sigma\xi$ durch $q$ theilbar oder nicht theilbar, je nachdem $p-1$ durch $q$ theilbar oder nicht theilbar ist.

2. Wir wenden uns zu dem Falle, dass $m = 2^\varkappa$ eine Potenz von 2 ist, und hier haben wir wieder zu unterscheiden, ob $p-1$ durch 4 oder nur durch 2 theilbar ist.

a) Ist $p-1$ durch 4 theilbar, so ist (\S.\,15):
\[ p \equiv 5^\beta \, (\mathrm{mod}\ m), \]
und $f$ ist die kleinste positive Zahl, die der Bedingung
\[ \beta f \equiv 0 \, [\mathrm{mod}\ \tfrac{1}{2}\varphi(m)] \]
genügt. Ist $f=1$, also $\beta$ durch $\tfrac{1}{2}\varphi(m) = 2^{\varkappa-2}$ theilbar, so ist $e = \varphi(m)$, und $\xi$ durchläuft die ganze Gruppe $\mathfrak{N}$, d.\,h. alle ungeraden Zahlen zwischen 0 und $m$. Dies ist der Fall, wo $p-1$ durch $m$ theilbar ist.

Ist $2^{\varkappa_1-2}$ die höchste Potenz von 2, die in $\beta$ aufgeht, und $2 \leqq \varkappa_1 < \varkappa$, so ist, wenn wir $m_1 = 2^{\varkappa_1}$ und $m = m_1 m_2$ setzen,
\[ f = 2^{\varkappa - \varkappa_1}, \quad e = 2^{\varkappa_1 - 1} = \varphi(m_1), \]
und es ist, wenn $s$ eine ungerade Zahl ist,
\[ p - 1 \equiv 5^{s\tfrac{1}{2}\varphi(m_1)} - 1 \ldots \, (\mathrm{mod}\ m). \]
% [sic] Source prints three low dots after -1 in this congruence; no continuation term is specified.
Nach \S.\,15 ist $5^{\tfrac{1}{2}\varphi(m_1)} - 1$ durch $m_1$, aber durch keine höhere Potenz von 2 theilbar. Da man ausserdem nach denselben Sätzen
\[ 5^{s\tfrac{1}{2}\varphi(m_1)} \equiv 5^{\tfrac{1}{2}\varphi(m_1)} \, (\mathrm{mod}\ 2m_1) \]
hat, so folgt, dass $m_1$ der grösste gemeinschaftliche Theiler von $p-1$ und $m$ ist.

Wenn wir also, wie oben, den Körper $\Omega_{m_1}$ zuziehen, so ergiebt sich, dass, wenn $m_1 \geqq 4$ der grösste gemeinschaftliche Theiler von $p-1$ und $m$ ist, $\xi$ ein volles System ungerader Reste von $m_1$ durchläuft, und dass $\mathfrak{p}_\xi$ zugleich die Primfactoren von $p$ im Körper $\Omega_{m_1}$ sind.


% ===== scan leaf 648, printed 627 =====
% !PRINTED 627
\weberpage{627}
b) Ist $p-1$ nicht durch 4 theilbar, also
\weq{11}{p \equiv -5^\beta \, (\mathrm{mod}\ m),}
so ist $f$ die kleinste positive Zahl, die den Congruenzen
\[ f \equiv 0 \, (\mathrm{mod}\ 2), \quad 2\beta f \equiv 0 \, [\mathrm{mod}\ \varphi(m)] \]
genügt, und folglich ist $e$ der grösste gemeinschaftliche Theiler von $\varphi(m)$ und $2\beta$, oder, wenn $2\beta$ durch $\varphi(m)$ theilbar ist, gleich $\tfrac{1}{2}\varphi(m)$. (Eine Ausnahme bildet hier der Fall $m=4$, in dem $e=1$ ist. Diesen Fall, der ja schon in früheren Abschnitten vollständig erledigt ist, lassen wir daher jetzt bei Seite.)

Es ergiebt sich dann aus (11), dass eine Zahl von der Form $5^\lambda$ nur dann einer Potenz von $p$ congruent sein kann, wenn $\lambda$ ein Vielfaches von $2\beta$, also ein Vielfaches von $e$ ist.

Hier kann nun in jeder der Nebengruppen $\mathfrak{A}\xi$ eine Zahl gefunden werden, die $\equiv 1 \, (\mathrm{mod}\ 4)$ ist, da $p \equiv -1 \, (\mathrm{mod}\ 4)$ eine Zahl in $\mathfrak{A}$ ist, und wir erhalten demnach ein vollständiges System der $\xi$ in der Form:
\[ 1, 5, 5^2, \ldots, 5^{e-1}, \]
und darin ist $e$ höchstens $= \tfrac{1}{2}\varphi(m) = 2^{\varkappa-2}$.

Wir fassen das Bewiesene im folgenden Theorem zusammen:

III. \gesperrt{Zerlegt man die von $q$ verschiedene Primzahl $p$ im Körper $\Omega_m$ in ihre Primfactoren:}
\[ p = \mathfrak{p}_{\xi_1} \mathfrak{p}_{\xi_2} \cdots \mathfrak{p}_{\xi_e}, \]
\gesperrt{so kann man, wenn $p-1$ durch $q$, oder im Falle eines geraden $m$ durch 4 theilbar ist, die $\xi$ ein volles System durch $q$ nicht theilbarer Reste nach dem Modul $m_1$ durchlaufen lassen, und die $\mathfrak{p}_\xi$ als Primfunctionale im Körper $\Omega_{m_1}$ annehmen, wenn $m_1$ der grösste gemeinschaftliche Theiler von $p-1$ und $m$ ist.}

\gesperrt{Ist aber $p-1$ nicht durch $q$ oder für ein gerades $m$ nicht durch 4 theilbar, und bedeutet $g$ im ersten Falle eine Primitivwurzel von $m$, im zweiten Falle die Zahl 5, so durchläuft $\xi$ die Reihe der Zahlen}
\[ 1, g, g^2, \ldots, g^{e-1}, \]
\gesperrt{und $e$ ist bei einem ungeraden $m$ nicht durch $q-1$ theilbar, und bei geradem $m$ mindestens $=2$ und höchstens $=\tfrac{1}{2}\varphi(m)$. Ausgeschlossen ist hierbei der Fall $m=4$.}


% ===== scan leaf 649, printed 628 =====
% !PRINTED 628
\weberpage{628}
\pdfbookmark[2]{§ 173. Darstellung der Primfactoren von p}{bm198}
\begin{center}
\S.\,173.

\gesperrt{Darstellung der Primfactoren von $p$.}
\end{center}

Zur Darstellung der Primfunctionale des Körpers $\Omega_m$ können wir ein Verfahren anwenden, was wir im \S.\,157 kennen gelernt haben, welches sich hier in Folge des Umstandes, dass
\[ 1, r, r^2, \ldots, r^{\mu-1} \]
eine Basis von $\mathfrak{o}$ ist, wesentlich vereinfacht.

Die natürliche Primzahl $q$ ist, wie wir schon gesehen haben, die $m^{\mathrm{te}}$ Potenz einer im Körper $\Omega_m$ existirenden Primzahl $\sigma$; mit dieser brauchen wir uns nicht weiter zu beschäftigen, und betrachten daher hier nur die von $q$ verschiedenen Primzahlen $p$.
% [sic] Source prints the m-th power of sigma here; the ramification index of the prime above q is mu = phi(m). Preserve the printed m.
Es sei $\mathfrak{p}$ ein Primfactor von $p$, und $\mathfrak{A}$ habe dieselbe Bedeutung wie oben. Wenn $p$ zum Exponenten $f$ gehört und
\[ ef = \varphi(m) = \mu \]
ist, so besteht die Gruppe $\mathfrak{A}$ aus den Zahlen
\[ 1, p, p^2, \ldots, p^{f-1}. \]
Wenn nun
\[ f(t) = N(t-r) = \prod^n (t - r^n) \]
das Polynom von $t$ vom Grade $\mu$ bedeutet, dessen Wurzeln die $\mu$ Grössen $r^n$ sind, so können wir dies so in Factoren zerlegen, dass jeder Factor einer der Nebengruppen von $\mathfrak{A}$ entspricht. Setzen wir nämlich
\weq{1}{F_1(t) = \prod_{0,\,f-1}^{h} (t - r^{p^h}),}
und allgemein für jedes durch $q$ nicht theilbare $n$
\weq{2}{F_n(t) = \prod_{0,\,f-1}^{h} (t - r^{np^h}),}
so ist $F_n(t)$ mit $F_{n'}(t)$ identisch, wenn $n$ und $n'$ in dieselbe Nebengruppe $\mathfrak{A}\xi_i$ gehören. Wenn wir also mit $\xi_1, \xi_2, \ldots, \xi_e$ das im \S.\,172, (5) angewandte Zahlensystem verstehen, so ist
\weq{3}{f(t) = F_{\xi_1}(t)\, F_{\xi_2}(t) \ldots F_{\xi_e}(t).}

Nun ist aber nach dem binomischen Lehrsatze:
\[ (t - r^{p^h})^p \equiv (t^p - r^{p^{h+1}}) \, (\mathrm{mod}\ p), \]
%CONT>>


% ===== scan leaf 650, printed 629 =====
% !PRINTED 629
\weberpage{629}
%<<CONT
und wenn wir dies auf jeden Factor von $F_n(t)$ anwenden, und
beachten, dass $p^{h+1}$ nach dem Modul $m$ dieselbe Zahlenreihe
durchläuft, wie $p^h$, so folgt:
\[
[F_n(t)]^p \equiv F_n(t^p) \ (\mathrm{mod}\ p),
\]
und diese Congruenz gilt dann natürlich auch für den Modul $\mathfrak p$.

Dies ist aber das Kennzeichen dafür, dass $F_n(t)$ mit einem
Polynom $P_n(t)$ mit ganzen rationalen Zahlenco\"efficienten nach
dem Modul $\mathfrak p$ congruent ist (\S.\,150, 4.), also
\weq{4}{F_n(t) \equiv P_n(t) \ (\mathrm{mod}\ \mathfrak p).}

Setzen wir dies in (3) ein, so ergiebt sich zunächst eine
Congruenz nach dem Modul $\mathfrak p$, die aber, da es eine Congruenz
zwischen \gesperrt{rationalen} Functionen ist, auch für den Modul $p$
bestehen muss, also
\weq{5}{f(t) \equiv P_{\xi_1}(t)\ P_{\xi_2}(t)\ \ldots\ P_{\xi_e}(t) \ (\mathrm{mod}\ p).}

Nach der Definition (1), (4) ist
\weq{6}{P_1(r) \equiv 0 \ (\mathrm{mod}\ \mathfrak p),}
und die sämmtlichen Wurzeln dieser Congruenz sind
\weq{7}{r,\ r^p,\ \ldots,\ r^{p^{f-1}}.}

Ebenso hat die Congruenz
\[
P_n(t) \equiv 0 \ (\mathrm{mod}\ \mathfrak p)
\]
die Wurzeln
\[
r^n,\ r^{np},\ \ldots,\ r^{np^{f-1}},
\]
und diese Wurzeln sind, wenn man $n$ das System der Zahlen
$\xi_1, \xi_2, \ldots, \xi_e$ durchlaufen lässt, alle incongruent nach dem
Modul $\mathfrak p$.

Macht man in der Congruenz
\[
F_1(t) \equiv P_1(t) \ (\mathrm{mod}\ \mathfrak p)
\]
die Substitution $(r, r^n)$, so geht $F_1(t)$ in $F_n(t)$, $\mathfrak p$ in das
conjugirte Ideal $\mathfrak p_n$ über, und $P_1(t)$ bleibt als rationale Form
ungeändert. Wir haben also
\weq{8}{F_n(t) \equiv P_1(t) \ (\mathrm{mod}\ \mathfrak p_n),}
und wenn wir hierin $t=r$ setzen, so folgt:
\weq{9}{F_n(r) \equiv P_1(r) \ (\mathrm{mod}\ \mathfrak p_n),}
$F_n(r)$ ist aber, wenn $n$ nicht in $\mathfrak A$ enthalten ist, relativ prim
zu $p$, also nicht durch $\mathfrak p_n$ theilbar.

Es ist also $P_1(r)$ nach (9) durch $\mathfrak p_1$, aber durch keinen der
%CONT>>


% ===== scan leaf 651, printed 630 =====
% !PRINTED 630
\weberpage{630}
%<<CONT
mit $\mathfrak p_1$ conjugirten Primfactoren theilbar, und es folgt, da $p$ nicht
durch $\mathfrak p^2$ theilbar ist:

1. \gesperrt{Der Primfactor $\mathfrak p_1$ ist der grösste gemeinschaftliche
Theiler von $p$ und $P_1(r)$; in gleicher Weise ergiebt sich, dass
$\mathfrak p_n$ der grösste gemeinschaftliche Theiler von $p$ und $P_1(r^n)$
ist, oder auch der grösste gemeinschaftliche Theiler von $p$ und
$P_{n'}(r)$, wenn $n n' \equiv 1 \ (\mathrm{mod}\ m)$ ist.}

Das letztere ergiebt sich aus der aus (4) durch die Substitution
 $(r, r^{n'})$ folgenden Congruenz
\[
F_1(t) \equiv P_n(t) \ (\mathrm{mod}\ \mathfrak p_{n'}).
\]

Man erhält also die sämmtlichen Primfactoren von $p$, wenn
man die grössten gemeinschaftlichen Theiler von $p$ mit jeder
der rationalen Functionen
\weq{10}{P_{\xi_1}(r),\ P_{\xi_2}(r),\ \ldots,\ P_{\xi_e}(r)}
aufsucht, und zwar ist $\mathfrak p_{\xi_i}$ der grösste gemeinschaftliche Theiler
von $p$ und $P_{\xi_i^{-1}}$.
% [sic] Source omits the argument (r) after P_{xi_i^{-1}} here; (10) and theorem 1 give the intended evaluation.

Wir wollen den Fall noch etwas näher betrachten, wo $f=1$
ist, also
\weq{11}{p \equiv 1 \ (\mathrm{mod}\ m).}

In diesem Falle ist $e = \varphi(m)$; die Gruppe $\mathfrak A$ reducirt sich
auf die Einheit, und die sämmtlichen $\mu$ conjugirten Factoren $\mathfrak p$
von $p$ sind von einander verschieden. Die Formen
\[
P_{\xi_1}(t),\ P_{\xi_2}(t),\ \ldots,\ P_{\xi_e}(t)
\]
werden alle linear, d.\,h. $r$ ist einer rationalen Zahl nach jedem
der Moduln $\mathfrak p_i$ congruent.

Dies ergiebt sich auch daraus, dass hier die Anzahl der
nach dem Modul $\mathfrak p$ incongruenten Zahlen gleich $N(\mathfrak p) = p$ ist, und
dass also $0, 1, 2, \ldots, p-1$ ein volles Restsystem ist.

Ist hiernach etwa $r \equiv c \ (\mathrm{mod}\ \mathfrak p)$, so muss, da $r^m = 1$ ist,
$c^m \equiv 1 \ (\mathrm{mod}\ \mathfrak p)$ also auch $(\mathrm{mod}\ p)$ sein, und wenn also $g$ eine
primitive Wurzel der Primzahl $p$ ist, so muss
\weq{12}{c \equiv g^{-n\frac{p-1}{m}} \ (\mathrm{mod}\ p)}
sein, worin $n$ relativ prim zu $m$ ist. Lassen wir $\mathfrak p$ das System
der conjugirten Ideale durchlaufen, so muss $n$ in (12) die
Gruppe $\mathfrak N$ durchlaufen, und es ist also nur Sache der Bezeichnung,
wenn wir festsetzen:
\weq{13}{r \equiv g^{-\frac{p-1}{m}} \ (\mathrm{mod}\ \mathfrak p).}


% ===== scan leaf 652, printed 631 =====
% !PRINTED 631
\weberpage{631}
Machen wir darin die Substitution $(r, r^n)$, so wird
\weq{14}{r^n \equiv g^{-\frac{p-1}{m}} \ (\mathrm{mod}\ \mathfrak p_n),}
oder wenn wir $n'$ durch die Congruenz
\weq{15}{n n' \equiv 1 \ (\mathrm{mod}\ m)}
definiren:
\weq{16}{r \equiv g^{-n'\frac{p-1}{m}} \ (\mathrm{mod}\ \mathfrak p_n).}

Es ist also in diesem Falle, übereinstimmend mit der allgemeinen
 Regel, $\mathfrak p_n$ der grösste gemeinschaftliche Theiler von
\[
p \ \text{und} \ \left(r - g^{-n'\frac{p-1}{m}}\right).
\]

Durch diese Bestimmung sind die einzelnen Primformen $\mathfrak p_n$
genau charakterisirt. Wir erhalten also den Satz:

2. \gesperrt{Eine Primzahl $p$, die nach dem Modul $m$ mit 1
congruent ist, zerfällt im Körper $\Omega_m$ in $\varphi(m)$ von
einander verschiedene Primfactoren $\mathfrak p_n$ vom
ersten Grade, die man als grösste gemeinschaftliche
 Theiler von $p$ mit den verschiedenen
Zahlen $r - g^{-n'\frac{p-1}{m}}$ erhält, wenn $g$ eine primitive
Wurzel von $p$ ist und $n'$ durch die Congruenz
$n n' \equiv 1 \ (\mathrm{mod}\ m)$ bestimmt ist.}

Nach \S.\,141 können wir jede ganze oder gebrochene Zahl $\omega$
und jedes Functional des Körpers $\Omega_m$ in der Weise in Primfactoren
 zerlegen, dass Zähler und Nenner keinen gemeinschaftlichen
 Primfactor enthalten, und diese Zerlegung ist, von
Einheitsfactoren abgesehen, völlig bestimmt, so dass, wenn
$\mathfrak p$ ein Primideal ist, ein ganz bestimmter positiver oder negativer
Exponent $k$ existirt, so dass das Functional $\omega \mathfrak p^{-k}$ den Primfactor
 $k$ weder im Zähler noch im Nenner enthält. % [sic] printed $k$; context requires $\mathfrak p$
Wir sagen
dann der Kürze halber, es sei $\mathfrak p^k$ \gesperrt{die in $\omega$ enthaltene Potenz
von $\mathfrak p$.}

Ist $k=0$, so sagen wir, $\mathfrak p$ ist in $\omega$ \gesperrt{nicht enthalten}.

Zwei Zahlen (oder auch Functionale) $\omega, \omega'$ eines Körpers, in
denen ein und dasselbe Primfunctional $\mathfrak p$ enthalten ist, \gesperrt{wollen
wir theilerverwandt oder kurz verwandt nennen}; wenn dagegen
kein Primfunctional zugleich in $\omega$ und in $\omega'$ enthalten ist, so
heissen $\omega$ und $\omega'$ \gesperrt{theilerfremd oder fremd}. Besonders nützlich
wird uns dieser Ausdruck in der Folge sein, wenn an
%CONT>>


% ===== scan leaf 653, printed 632 =====
% !PRINTED 632
\weberpage{632}
\pdfbookmark[2]{§ 174. Das Kummer'sche Theorem}{bm199}
%<<CONT
Stelle von $\omega'$ die natürliche Primzahl $p$ tritt, die durch $\mathfrak p$ theilbar
 ist, und wir reden danach von den \gesperrt{mit einer Zahl $\omega$
verwandten oder zu ihr fremden Primzahlen}.

\gesperrt{Eine Zahl $\omega$, die gar keine verwandten Primzahlen hat, ist
eine Einheit.}

Im Körper $\Omega_m$ (wie überhaupt in jedem Normalkörper) sind
alle conjugirten Zahlen $\omega_n$ mit denselben Primzahlen verwandt.

\begin{center}
\S.\ 174.

\gesperrt{Das Kummer'sche Theorem.}
\end{center}

Wir haben schon in der Kreistheilungstheorie (Bd.\,I, \S.\,169)
gewisse Zerlegungen der natürlichen Primzahlen in Factoren, die
aus Einheitswurzeln zusammengesetzt waren, kennen gelernt,
und es ist nun von Wichtigkeit, zu untersuchen, wie sich diese
Factoren zu den Primfactoren im Körper $\Omega_m$ verhalten.

Wir betrachten eine Primzahl $p$, die zum Exponenten 1
gehört, so dass also $m$ ein Theiler von $p-1$ ist, und setzen
\[
p - 1 = m m'.
\]

Aus den $p^{\mathrm{ten}}$ Einheitswurzeln $\varrho$ lassen sich nach Bd.\,I, \S.\,167
Perioden von je $m'$ Gliedern bilden, die, wenn $g$ eine primitive
Wurzel von $p$ (nicht, wie oben, von $m$) ist, den Ausdruck haben:
\weq{1}{
\begin{aligned}
\eta &= \varrho \phantom{^g} + \varrho^{g^m} + \varrho^{g^{2m}} + \cdots + \varrho^{g^{(m'-1)m}} \\
\eta_1 &= \varrho^g + \varrho^{g^{m+1}} + \varrho^{g^{2m+1}} + \cdots + \varrho^{g^{(m'-1)m+1}} \\
&\ \, \cdots\cdots\cdots\cdots\cdots\cdots\cdots\cdots\cdots\cdots \\
\eta_{m-1} &= \varrho^{g^{m-1}} + \varrho^{g^{2m-1}} + \varrho^{g^{3m-1}} + \cdots + \varrho^{g^{p-2}}.
\end{aligned}
}

Bezeichnen wir, wie im \S.\,169 des ersten Bandes, mit $\varepsilon$ eine
primitive $(p-1)^{\mathrm{te}}$ Einheitswurzel, so dass
\weq{2}{\varepsilon^{m'} = r,}
wie bisher, eine primitive $m^{\mathrm{te}}$ Einheitswurzel ist, so sind die
Resolventen
\begin{equation*}\tag*{3)}
(\varepsilon^\lambda, \varrho) = \varrho + \varepsilon^\lambda \varrho^g + \varepsilon^{2\lambda}\varrho^{g^2} + \cdots + \varepsilon^{(p-2)\lambda}\varrho^{g^{p-2}},
\end{equation*}
% [sic] printed "3)" without an opening parenthesis, unlike the other display numbers on this page
worin $\lambda$ jede ganze rationale Zahl sein kann, und speciell für
$\lambda = m'$:
\weq{4}{(r, \eta) = \eta + r\eta_1 + r^2\eta_2 + \cdots + r^{m-1}\eta_{m-1}.}


% ===== scan leaf 654, printed 633 =====
% !PRINTED 633
\weberpage{633}
Diese Resolventen sind ganze algebraische Zahlen in gewissen
Kreistheilungskörpern, und wir haben, wenn $\varepsilon^\lambda$ nicht $=1$ ist,
\[
(\varepsilon^\lambda,\varrho)\,(\varepsilon^{-\lambda},\varrho) = (-1)^\lambda p,
\]
und für $\lambda = m'n$
\weq{5}{(r^n,\eta)\,(r^{-n},\eta) = \pm\, p.}

Es sind also alle diese Resolventen Theiler der natürlichen
Primzahl $p$.

Ausser diesen Resolventen haben wir an der angeführten
Stelle noch andere ganze Zahlen der Kreistheilungskörper abgeleitet, nämlich
\weq{6}{\psi_{\lambda,\mu}(\varepsilon) = \sum\limits_{1,\,p-2}^{t} \varepsilon^{\mu\,\mathrm{ind}\,t - (\lambda+\mu)\,\mathrm{ind}\,(t+1)},}
(worin sich die Indices auf die Basis $g$ beziehen), die im Allgemeinen dem Körper $\Omega_{p-1}$ angehören, in dem besonderen Falle, wo $\lambda$ und $\mu$ durch $m'$ theilbar sind, dem Körper $\Omega_m$. Von besonderer Wichtigkeit sind darunter die Zahlen des Körpers $\Omega_m$:
\weq{7}{\psi_{sm',\,m'}(\varepsilon) = \psi_s(r) = \sum\limits_{1,\,p-2}^{t} r^{\mathrm{ind}\,t - (s+1)\,\mathrm{ind}\,(t+1)}.}

Auch diese Zahlen $\psi_{\lambda,\mu}(\varepsilon)$ sind, wenn keine der Zahlen $\lambda$, $\mu$, $\lambda+\mu$ durch $p-1$ theilbar ist, Theiler von $p$ nach der Formel
\[
\psi_{\lambda,\mu}(\varepsilon)\,\psi_{\lambda,\mu}(\varepsilon^{-1}) = p,
\]
aus der sich wieder nach (7) ergiebt:
\weq{8}{\psi_s(r)\,\psi_s(r^{-1}) = p,}
vorausgesetzt, dass weder $s$ noch $s+1$ durch $m$ theilbar ist. Gewisse Verbindungen der Resolventen (4), darunter die $m^{\mathrm{te}}$ Potenz, haben wir durch die Functionen $\psi$ ausgedrückt [Bd.\ I, \S.\,169, (14), (15)]:
\weq{9}{(r,\eta)^n\,(r^n,\eta)^{-1} = \psi_1(r)\,\psi_2(r)\cdots\psi_{n-1}(r),}
so lange $n<m$ ist, und
\weq{10}{(r,\eta)\,(r^{-1},\eta) = \pm\,p}
\weq{11}{(r,\eta)^m = (-1)^{\frac{p-1}{m}}\, p\, \psi_1(r)\,\psi_2(r)\cdots\psi_{m-2}(r),}
woraus hervorgeht, dass diese $m^{\mathrm{ten}}$ Potenzen dem Körper $\Omega_m$ angehören.

Wir bezeichnen nun, wie früher, mit $\mathfrak{p}_n$ die verschiedenen Primtheiler von $p$ im Körper $\Omega_m$. Die Zahlen $\psi_s(r)$, die gleichfalls diesem Körper angehören und Factoren von $p$ sind, müssen
%CONT>>


% ===== scan leaf 655, printed 634 =====
% !PRINTED 634
\weberpage{634}
%<<CONT
daher durch einige dieser Primfactoren theilbar sein, und sie können, so wenig wie $p$ selbst, durch das Quadrat eines $\mathfrak{p}_n$ theilbar sein. Wir müssen feststellen, durch welche Primtheiler $\mathfrak{p}_n$ jede dieser Zahlen $\psi_s(r)$ theilbar ist.

Dazu aber führen einerseits die Congruenzen des vorigen Paragraphen:
\weq{12}{r \equiv g^{-n\frac{p-1}{m}} \ (\mathrm{mod}\ \mathfrak{p}_{n'}), \qquad n n' \equiv 1\ (\mathrm{mod}\ m),}
andererseits die im \S.\,170 des ersten Bandes abgeleiteten Congruenzen (17):
\weq{13}{\begin{aligned}
&\psi_{\lambda,\mu}(g) \equiv 0, && \lambda+\mu<p-1\\[4pt]
&\psi_{\lambda,\mu}(g) \equiv -\dfrac{\Pi(2p-\lambda-\mu-2)}{\Pi(p-\lambda-1)\,\Pi(p-\mu-1)}, && (\mathrm{mod}\ p)\\
&&& p-1<\lambda+\mu<2p-2,
\end{aligned}}
worin $\lambda$ und $\mu$ zwischen $0$ und $p-1$ genommen sind, so dass der Binominalcoëfficient % [sic] "Binominalcoëfficient" as printed; "Binomialcoëfficient" is meant
\[
\frac{\Pi(2p-\lambda-\mu-2)}{\Pi(p-\lambda-1)\,\Pi(p-\mu-1)}
\]
eine durch $p$ nicht theilbare ganze Zahl ist.

Nach der Congruenz (12) ist
\[
\psi_s(r) \equiv \psi_s(g^{-nm'}) \equiv \sum\limits^{t} g^{-nm'\,[\mathrm{ind}\,t - (s+1)\,\mathrm{ind}\,(t+1)]},
\]
oder nach (6):
\[
\psi_s(r) \equiv \psi_{\lambda,\mu}(g) \ (\mathrm{mod}\ \mathfrak{p}_{n'}),
\]
wenn
\weq{14}{\begin{aligned}
\lambda &\ \text{den kleinsten positiven Rest von } {-}nm's\\
\mu &\ \text{\textquotedbl\ \textquotedbl\ \textquotedbl\ \textquotedbl\ \textquotedbl}\ {-}nm'
\end{aligned}\quad (\mathrm{mod}\ p-1)}
bedeutet, und die Congruenzen (13) zeigen dann, dass $\psi_s(r)$ durch $\mathfrak{p}_{n'}$ theilbar ist, wenn $\lambda+\mu<p-1$, und nicht theilbar, wenn $\lambda+\mu>p-1$ ist.

Da die zu $m$ theilerfremde Zahl $n$ nur nach dem Modul $m$ bestimmt ist, so nehmen wir sie jetzt \gesperrt{positiv und kleiner als $m$ an}. Dann ist
\[
\mu = m'(m-n).
\]

Die aus (14) bestimmte Zahl $\lambda$ ist ein Vielfaches von $m'$ und wenn wir $\lambda = m'\varkappa$ setzen, so ist
\weq{15}{\varkappa \ \text{der kleinste positive Rest von } {-}ns \ (\mathrm{mod}\ m).}

Hiernach wird $\psi_s(r)$ durch $\mathfrak{p}_{\varkappa'}$ theilbar sein, wenn
% [sic] Source prints the index varkappa' in this sentence; the preceding derivation and following factor count indicate mathfrak p_{n'}. Preserve the printed index.
\[
m'(\varkappa+m-n) < m m'
\]
oder $\varkappa < n$, und nicht theilbar, wenn $\varkappa > n$ ist.


% ===== scan leaf 656, printed 635 =====
% !PRINTED 635
\weberpage{635}
Da wir jetzt im Stande sind, für jedes $s$ die Primfactoren von $\psi_s(r)$ zu ermitteln, gehen wir dazu über, die Zahl $(r,\eta)^m$ nach der Formel (11) in ihre Primfactoren zu zerlegen.

Wir bilden zu diesem Zwecke nach (15) die kleinsten positiven Reste $\varkappa$ der Zahlen, $-n, -2n, \cdots -(m-2)n$ nach dem Modul $m$, und finden, von der Ordnung abgesehen, die Zahlen
\[
\varkappa = 1, 2, \ldots, n-1, n+1, \ldots, m-1;
\]
die Reste $0$ und $n$ kommen darunter nicht vor, weil $s$ keinen Werth erhält, für den $s$ oder $s+1$ durch $m$ theilbar wird. Von diesen Zahlen sind aber gerade $n-1$ kleiner als $n$, und so oft kommt also der Factor $\mathfrak{p}_{n'}$ in dem Producte $\psi_1\psi_2\ldots\psi_{m-2}$ vor. Dann enthält $p$ diesen Factor noch einmal und folglich kommt er nach (11) genau $n$ mal in $(r,\eta)^m$ vor. Da keine anderen als die Primfactoren $\mathfrak{p}_{n'}$ in $(r,\eta)^m$ aufgehen können, so ist hierdurch die Zerlegung vollständig ausgeführt:
\weq{16}{(r,\eta)^m = \prod^{n} \mathfrak{p}_{n'}^{\,n};}
hierin durchläuft $n$ die Reihe der positiven Zahlen, die kleiner als $m$ und relativ prim zu $m$ sind, und $n'$ ist jedesmal aus der Congruenz
\weq{17}{n n' \equiv 1\ (\mathrm{mod}\ m)}
zu bestimmen\footnote{Dieser Satz rührt von \gesperrt{Kummer} her. Vgl. Theorie der idealen Primfactoren der complexen Zahlen etc. Abhandlungen der Berliner Akademie, 1856.}.

Die Formel (16) gilt für ein gerades wie für ein ungerades $m$ und ist auch noch für den Fall $m=2$ gültig, für den sie ein schon bekanntes Resultat giebt, da dann $(r,\eta) = (-1,\eta)$ mit einer \gesperrt{Gauss}'schen Summe übereinstimmt (Bd.\ I, \S.\,171).

Wir wollen aber nun, indem wir den Fall $m=2$ jetzt ausschliessen, der Formel (16) eine etwas allgemeinere Gestalt geben.

Es sei $m=q^{\varkappa}$ eine beliebige Primzahlpotenz, nur grösser als $2$, und $p$ eine von $q$ verschiedene Primzahl von der Eigenschaft, dass $p-1$ durch $q$ oder, wenn $q=2$ ist, durch $4$ theilbar ist.

Es sei ferner $m_1$ der grösste gemeinschaftliche Theiler von $p-1$ und $m$, und
\[
m = m_1 m_2,
\]
%CONT>>


% ===== scan leaf 657, printed 636 =====
% !PRINTED 636
\weberpage{636}
%<<CONT
dann ist nach \S.\,172, III. die Zerlegung von $p$ in Primfactoren im Körper $\Omega_m$ dieselbe, wie im Körper $\Omega_{m_1}$, und es ist
\[
\mathfrak{p}_n = \mathfrak{p}_{n+sm_1},
\]
wenn $s$ eine beliebige ganze Zahl ist. Wir setzen jetzt
\weq{18}{t = n + s m_1}
und lassen $n$ die Reihe der durch $q$ nicht theilbaren positiven ganzen Zahlen $<m_1$, und $s$ die Zahlenreihe $0,1,2,\ldots,m_2-1$ durchlaufen, dann durchläuft $t$ ein volles System durch $q$ nicht theilbarer Reste von $m$, und zwar das ganz bestimmte System, dessen Elemente alle $<m$ sind. Es ist dann nach \S.\,172
\weq{19}{p = \prod^{n} \mathfrak{p}_n.}

Ist ferner $t'$ eine Zahl, die der Congruenz
\weq{20}{t\,t' \equiv 1\ (\mathrm{mod}\ m)}
genügt, so ist, wenn $n'$ durch die Bedingung $n n' \equiv 1\ (\mathrm{mod}\ m_1)$ bestimmt ist,
\weq{21}{t' \equiv n'\ (\mathrm{mod}\ m_1).}

Danach lässt sich das Product
\[
\prod^{t} \mathfrak{p}_{t'}^{\,t}
\]
nach der Formel (16) bestimmen, die, auf den Körper $\Omega_{m_1}$ angewandt, der aus allen rationalen Functionen der $m_1^{\mathrm{ten}}$ Einheitswurzel $r_1 = r^{m_2}$ besteht, die Formel ergiebt:
\weq{22}{(r_1,\eta)^{m_1} = \prod^{n} \mathfrak{p}_{n'}^{\,n}.}

Es ist aber nach (18) und (21)
\[
\prod^{t} \mathfrak{p}_{t'}^{\,t} = \prod^{n} \mathfrak{p}_{n'}^{\,n m_2} \prod^{n'}\prod^{s} \mathfrak{p}_{n'}^{\,s m_1},
\]
und da $n'$ nach dem Modul $m_1$ dieselbe Zahlenreihe durchläuft wie $n$, so ist nach (19):
\[
\prod^{s}\prod^{n'} \mathfrak{p}_{n'}^{\,s m_1} = p^{m_1 \Sigma s} = p^{\frac{1}{2}m(m_2-1)},
\]
und es ergiebt sich:
\[
\prod^{t} \mathfrak{p}_{t'}^{\,t} = \left[p^{\frac{1}{2}(m_2-1)}\,(r^{m_2},\eta)\right]^m.
\]

Setzen wir
\weq{23}{\varrho = p^{\frac{1}{2}(m_2-1)}\,(r_1,\eta), \qquad \varrho_n = p^{\frac{1}{2}(m_2-1)}\,(r_1^{\,n},\eta),}
so ist $\varrho$ eine \gesperrt{Kreistheilungszahl}, weil ja auch $\sqrt{p}$ als Werth
%CONT>>


% ===== scan leaf 658, printed 637 =====
% !PRINTED 637
\weberpage{637}
%<<CONT
einer \gesperrt{Gauss}'schen Summe zu den Kreistheilungszahlen gehört,
und es ist
\weq{24}{\varrho^m = \Pi^{t}\, \mathfrak{p}^t_{t'}.}

Nach (9) und (11) ist aber, wenn darin $m$ durch $m_1$ ersetzt wird,
\weq{25}{\varrho_n^{m_1} = \pm\, p^{\frac{1}{2}\,m_1(m_2-1)+1}\, \psi_1(r_1^n)\,\dots\,\psi_{m_1-2}(r_1^n),}
\weq{26}{\varrho^n \varrho_n^{-1} = p^{\frac{1}{2}(n-1)(m_2-1)}\, \psi_1(r_1)\,\psi_2(r_1)\,\dots\,\psi_{n-1}(r_1)}
und es sind also $\varrho^n \varrho_n^{-1}$, $\varrho_n^{m_1}$ und folglich auch $\varrho_n^m$ Zahlen des
Körpers $\Omega_m$, und nach (25) bleibt die in (26) vorkommende
Verbindung $\varrho^n \varrho_n^{-1}$ in $\Omega_m$ enthalten, wenn $n$ um ein Vielfaches
von $m_1$ vermehrt wird. Es ist demnach $\varrho^n \varrho_n^{-1}$ auch dann noch
eine Zahl in $\Omega_m$, wenn $n$ nicht gerade die beschränkte Bedeutung
hat wie bisher, sondern eine beliebige durch $q$ nicht theilbare
Zahl bedeutet.

Diese Betrachtung hat den Zweck, das \gesperrt{Kummer}'sche
Theorem (16) gleichzeitig auf mehrere verschiedene Primzahlen
anwendbar zu machen, die zu verschiedenen Werthen von $m_1$
gehören, und das Resultat spricht sich dann in folgendem
Theorem aus:

3. \gesperrt{Es sei $\varphi_n$ ein System conjugirter Functionale des
Körpers $\Omega_m$, das mit keinen anderen natürlichen
Primzahlen verwandt ist, als solchen, die nach
dem Modul $q$ (oder bei $q = 2$ nach dem Modul 4)
mit 1 congruent sind; es durchlaufe ferner $t$
die durch $q$ nicht theilbaren Zahlen der Reihe
$1, 2, \dots, m - 1$, und $t'$ sei aus der Congruenz
$t\,t' \equiv 1 \pmod{m}$ bestimmt.}

\gesperrt{Dann ist das Functional des Körpers $\Omega_m$:}
\[ \text{a)}\qquad \Phi_n = \Pi^{t}\, \varphi^t_{nt'} = \varepsilon\,\vartheta_n^m \]
\gesperrt{mit der $m^{\text{ten}}$ Potenz einer Kreistheilungszahl $\vartheta_n$
associirt, die zwar selbst in einem höheren
Körper enthalten ist, deren $m^{\text{te}}$ Potenz aber dem
Körper $\Omega_m$ angehört, und es ist}
\[ \text{b)}\qquad \vartheta_n^{-1}\vartheta_1^n = \alpha \]
\gesperrt{eine Zahl des Körpers $\Omega_m$.}

Der Beweis ist in den vorangegangenen Ausführungen enthalten
und ergiebt sich aus den Formeln (25), (26), wenn man
%CONT>>


% ===== scan leaf 659, printed 638 =====
% !PRINTED 638
\weberpage{638}
\pdfbookmark[2]{§ 175. Die Einheitswurzeln im Körper Ωm}{bm200}
%<<CONT
unter $\vartheta$ ein Product aus Potenzen der Zahlen $\varrho$ versteht, die
nach den Formeln (24) aus den in $\varphi$ enthaltenen Primfactoren $\mathfrak{p}$
abgeleitet sind.

\begin{center}
§. 175.

\gesperrt{Die Einheitswurzeln im Körper $\Omega_m$.}
\end{center}

Für die weiter zu machenden Anwendungen ist eine genauere
Kenntniss der Einheiten des Körpers $\Omega_m$ erforderlich, die ja
auch an sich von grossem Interesse ist. Wir beschäftigen uns
hier nicht mit den functionalen, sondern nur mit den numerischen
Einheiten, worunter, wie wir uns erinnern, \gesperrt{ganze Zahlen des
Körpers $\Omega_m$ zu verstehen sind, deren Norm $\pm 1$ ist,
oder eine ganze Zahl, deren reciproker Werth gleichfalls ganz ist.}

Zu den Einheiten gehören sicher alle in $\Omega_m$ enthaltenen
Einheitswurzeln; wir können aber leicht beweisen, dass diese
Einheitswurzeln nur Potenzen von $r$ sein können, mit positivem
und negativem Vorzeichen.

Es sei nämlich $\varrho$ eine in $\Omega_m$ enthaltene Einheitswurzel, und
\weq{1}{\varrho = \varphi(r).}

Ist $q^h$ die höchste im Grade von $\varrho$ aufgehende Potenz von $q$,
so ist $\varrho^{q^h}$ eine Einheitswurzel, deren Grad nicht durch $q$ theilbar
ist, und
\weq{2}{\varrho^{q^h} = [\varphi(r)]^{q^h}.}

Wegen der Irreducibilität der Kreistheilungsgleichung im
weiteren Sinne\footnote{Vgl. Bd.~I, §.\,134 und den Nachtrag zu diesem Bande.} können daher in (2) die sämmtlichen Substitutionen
 $(r, r^n)$ gemacht werden, ohne dass die linke Seite sich
ändert, und folglich ist $\varrho^{q^h}$ eine rationale Zahl, und da es eine
Einheit ist, muss es $= \pm 1$ sein. Mithin ist $\varrho$, vom Vorzeichen
abgesehen, einer Einheitswurzel vom Grade $q^h$ oder, wenn $q = 2$
ist, vom Grade $q^{h+1}$ gleich. Es ist nachzuweisen, dass $q^h$ oder
$q^{h+1}$ nicht grösser als $m$ sein kann. Dies ergiebt sich aber
aus (1). Denn wäre $\varrho$ eine Einheitswurzel höheren als $m^{\text{ten}}$
Grades, so wäre $r$ eine Potenz von $\varrho$, und $\varrho$ müsste einer irreduciblen
 Gleichung höheren Grades als $r$, also auch höheren
Grades als $\varphi(r)$ genügen, was nach (1) ein Widerspruch ist. Also
haben wir den Satz:
%CONT>>


% ===== scan leaf 660, printed 639 =====
% !PRINTED 639
\weberpage{639}
%<<CONT
1. \gesperrt{Die einzigen in $\Omega_m$ enthaltenen Einheitswurzeln
sind die mit positivem und negativem Zeichen
genommenen Potenzen von $r$.}

Hieran schliesst sich ein ganz allgemeiner, d.\,h. in jedem
algebraischen Zahlkörper gültiger Satz über die Einheitswurzeln,
den wir jetzt beweisen wollen.

Der Satz lautet:

2. \gesperrt{Ist $\alpha$ eine ganze Zahl eines algebraischen Zahlkörpers
 $n^{\text{ten}}$ Grades, und haben alle mit $\alpha$ conjugirten
 Zahlen $\alpha_1, \alpha_2, \dots, \alpha_n$ den absoluten
Werth 1, so ist $\alpha$ eine Einheitswurzel}\footnote{\gesperrt{Kronecker}, \glqq Zwei Sätze über Gleichungen mit ganzzahligen
Co\"efficienten\grqq. Crelle's Journal, Bd. 53 (1857). \gesperrt{Minkowski}, \glqq Geometrie
der Zahlen\grqq, Art. 43.}.

Zum Beweis ist zunächst zu bemerken, dass der absolute
Werth der Norm einer ganzen Zahl $\omega$ dem Product der absoluten
Werthe der conjugirten Zahlen $\omega_1, \omega_2, \dots, \omega_n$ gleich, und als
ganze rationale Zahl jedenfalls grösser oder gleich 1 ist. Es ist
daher nicht möglich, dass alle diese absoluten Werthe kleiner
als 1 sind. Sind sie aber alle gleich 1, wie bei der in unserem
Satze angenommenen Zahl $\alpha$, so muss die Norm $= \pm 1$ sein,
und $\alpha$ ist eine Einheit. Ist $\beta$ eine zweite Einheit, die mit ihren
conjugirten zugleich den absoluten Werth 1 hat, so hat der
Quotient $\alpha : \beta$ dieselbe Eigenschaft. Wenn daher unter den zu
diesem Quotienten conjugirten Zahlen auch nur eine reell ist,
so muss
\[ \frac{\alpha}{\beta} = \pm 1 \]
oder $\alpha = \pm \beta$ sein, und dies gilt dann auch noch, wenn $\alpha, \beta$
durch die entsprechenden Zahlen eines conjugirten Körpers
ersetzt werden.

Wenn also $\alpha$ und $\beta$ weder gleich noch entgegengesetzt sind,
so ist ihr Verhältniss nicht reell, und es besteht daher zwischen
den absoluten Werthen die Ungleichung:
\[ |\alpha \pm \beta| < |\alpha| + |\beta|. \]
(Vergl. die Einleitung zum ersten Bande, S. 19.)

Es ist also
\[ |\alpha \pm \beta| < 2, \]
%CONT>>


% ===== scan leaf 661, printed 640 =====
% !PRINTED 640
\weberpage{640}
%<<CONT
und folglich kann $\frac{1}{2}(\alpha \pm \beta)$ keine ganze Zahl sein, weil der
absolute Werth der Norm dieser Zahl kleiner als 1 ist. Wenn
nun $\omega_1, \omega_2, \dots, \omega_n$ eine Basis von $\mathfrak{o}$ ist, und
\[
\begin{aligned}
\alpha &= a_1\omega_1 + a_2\omega_2 + \dots + a_n\omega_n \\
\beta &= b_1\omega_1 + b_2\omega_2 + \dots + b_n\omega_n,
\end{aligned}
\]
so können die ganzen rationalen Zahlen $a, b$ nicht den Congruenzen
\[ a_1 \equiv b_1,\ a_2 \equiv b_2,\ \dots,\ a_n \equiv b_n \pmod{2} \]
genügen, weil sonst $\frac{1}{2}(\alpha - \beta)$ eine ganze Zahl wäre.

Wenn wir also die sämmtlichen Zahlen $\alpha$ in $2^n$ Fächer vertheilen,
indem wir alle Zahlen in ein Fach werfen, in denen
$a_1, a_2, \dots, a_n$ dieselben Reste (0 oder 1) nach dem Modul 2
lassen, so können in jedem dieser Fächer höchstens 2 Zahlen,
nämlich $\alpha$ und $-\alpha$, vorkommen, und wenn wir $\alpha = 0$ noch
ausschliessen, so giebt es sogar in einem dieser Fächer, in dem
die $a_1, a_2, \dots, a_n$ alle gerade sind, gar keine Zahl $\alpha$. Die Anzahl
 aller möglichen Zahlen $\alpha$ ist also endlich und höchstens
$= 2^{n+1} - 2$.

Wenn nun der absolute Werth von $\alpha = 1$ ist, so gilt dasselbe
von allen Potenzen von $\alpha$. Folglich muss in der unbegrenzten
Reihe
\[ 1,\ \alpha,\ \alpha^2,\ \alpha^3,\ \dots \]
nothwendig dieselbe Zahl zum zweiten Male wiederkehren, also
$\alpha^h = \alpha^k$ und $k > h$ sein. Dann ist aber
\[ \alpha^{k-h} = 1, \]
d.\,h. $\alpha$ ist eine Einheitswurzel, wie bewiesen werden sollte.

Wir fügen noch die Bemerkung bei, dass, wenn eine Zahl $\alpha$
eines Körpers $\Omega$ eine Einheitswurzel $m^{\text{ten}}$ Grades ist, auch alle
mit $\alpha$ conjugirten Zahlen Einheitswurzeln desselben Grades sind.
Denn in der rationalen Gleichung $\alpha^m - 1 = 0$ kann $\alpha$ durch
jeden der conjugirten Werthe ersetzt werden.

Die Anzahl der Einheitswurzeln, die in einem Körper $\Omega$
enthalten sind, kann immer nur eine endliche sein, weil jede
Zahl in $\Omega$ einer rationalen Gleichung genügt, deren Grad dem
Körpergrad höchstens gleich ist. Der Grad der Einheitswurzeln
in $\Omega$ kann daher einen endlichen Werth nicht übersteigen.


% ===== scan leaf 662, printed 641 =====
% !PRINTED 641
\weberpage{641}
\pdfbookmark[2]{§ 176. Der in Ωm enthaltene reelle Körper Hm}{bm201}
\begin{center}
\S.\,176.\\[4pt]
\gesperrt{Der in $\Omega_m$ enthaltene reelle K\"orper $H_m$.}
\end{center}

Jede Zahl des K\"orpers $\Omega_m$ geht durch die Substitution $(r, r^{-1})$ in die conjugirt imagin\"are Zahl \"uber, die auch durch die Vertauschung $(i, -i)$ erhalten wird. Das Product zweier solcher conjugirt imagin\"arer Zahlen ist das \gesperrt{Quadrat des absoluten Werthes} einer jeden von ihnen.

Eine Zahl in $\Omega_m$ ist reell, wenn sie durch die Vertauschung $(r, r^{-1})$ unge\"andert bleibt, und nur unter dieser Voraussetzung. Jede solche Zahl l\"asst sich rational durch die \gesperrt{zweigliedrige Periode} $r + r^{-1}$ ausdr\"ucken, und geh\"ort also einem reellen K\"orper vom Grade $\frac12\varphi(m)$ an, der ein Theiler von $\Omega_m$ ist. Diesen wollen wir den in $\Omega_m$ \gesperrt{enthaltenen reellen K\"orper} nennen und mit $H_m$ bezeichnen (obwohl auch noch andere reelle K\"orper, n\"amlich alle Theiler von $H_m$, in $\Omega_m$ enthalten sind).

Die $\frac12\varphi(m) = \frac12\mu$ Zahlen
\weq{1}{1,\ r+r^{-1},\ r^2+r^{-2},\ \dots,\ r^{\frac12\mu-1}+r^{-\frac12\mu+1}}
bilden eine Basis der ganzen Zahlen des K\"orpers $H_m$. Denn nach \S.\,170, (7) ist
\weq{2}{\omega = x_0+x_1r+x_2r^2+\cdots+x_{\frac12\mu-1}r^{\frac12\mu-1}+x_{\frac12\mu}r^{\frac12\mu} +x_{-1}r^{-1}+x_{-2}r^{-2}+\cdots+x_{-\frac12\mu+1}r^{-\frac12\mu+1}}
dann und nur dann eine ganze Zahl des K\"orpers $\Omega_m$, wenn die Co\"efficienten $x_0, x_1, \dots$ ganze rationale Zahlen sind.

Die Bedingung daf\"ur, dass diese Zahl dem K\"orper $H_m$ angeh\"ort, erh\"alt man, wenn man die Substitution $(r, r^{-1})$ ausf\"uhrt und die Differenz $=0$ setzt. Dividirt man noch durch $r-r^{-1}$, so ergiebt sich so:
\[
(x_1-x_{-1})+(x_2-x_{-2})(r+r^{-1})+\cdots
+(x_{\frac12\mu-1}-x_{-\frac12\mu+1})\frac{r^{\frac12\mu-1}-r^{-\frac12\mu+1}}{r-r^{-1}}
+x_{\frac12\mu}\frac{r^{\frac12\mu}-r^{-\frac12\mu}}{r-r^{-1}}=0.
\]
Die Divisionen lassen sich hier ausf\"uhren, und wenn man dann mit $r^{\frac12\mu-1}$ multiplicirt, so ergiebt sich f\"ur $r$ eine rationale Gleichung von niedrigerem als $\mu^{\text{ten}}$ Grade, die nur dann be-
%CONT>>


% ===== scan leaf 663, printed 642 =====
% !PRINTED 642
\weberpage{642}
%<<CONT
friedigt sein kann, wenn alle ihre Co\"efficienten verschwinden. Dies f\"uhrt zu den Gleichungen
\[
x_{\frac12\mu}=0,\ x_{\frac12\mu-1}=x_{-\frac12\mu+1},\ \dots,\ x_1=x_{-1},
\]
und es ist also jede ganze Zahl $\omega$ des K\"orpers $H_m$ in der Form enthalten:
\[
\omega = x_0+x_1(r+r^{-1})+x_2(r^2+r^{-2})+\cdots
+x_{\frac12\mu-1}(r^{\frac12\mu-1}+r^{-\frac12\mu+1}).
\]

Statt der Basis (1) kann man auch, wenn man
\[
\varrho = r+r^{-1}
\]
setzt, die Potenzen von $\varrho$:
\weq{3}{1,\ \varrho,\ \varrho^2,\ \dots,\ \varrho^{\frac12\mu-1}}
als Basis der ganzen Zahlen von $H_m$ w\"ahlen. Denn die Gr\"ossen (3) lassen sich ganzzahlig durch die (1) ausdr\"ucken und umgekehrt.

Die Zahl $\varrho$ ist die Wurzel einer irreduciblen Gleichung vom Grade $\frac12\mu$, die wir mit $\psi(\varrho)=0$ bezeichnen wollen. Ist $f(r)=0$ die Gleichung f\"ur $r$ (\S.\,169), so besteht die identische Relation
\weq{4}{f(x) = x^{\frac12\mu}\,\psi\Bigl(x+\frac1x\Bigr),}
woraus durch Bildung der Ableitung
\weq{5}{f'(r) = r^{\frac12\mu}\,\psi'(\varrho)\,(1-r^{-2}).}

Hieraus k\"onnen wir die Grundzahl $\Delta_1$ des K\"orpers $H_m$ bilden, die, da $H_m$ nur reelle Zahlen enth\"alt, positiv sein muss. Diese Grundzahl ist, wenn wir die Norm in Bezug auf $H_m$ mit $N_1$ bezeichnen:
\[
\Delta_1 = \pm\,N_1\,\psi'(\varrho).
\]

Ist aber $N$ die in Bezug auf $\Omega_m$ gebildete Norm, so ist
\[
N\,\psi'(\varrho) = [N_1\,\psi'(\varrho)]^2 = \Delta_1^2,
\]
und demnach ergiebt die Formel (5) mit R\"ucksicht auf \S.\,169:
\[
\Delta = N(1-r^{-2})\,\Delta_1^2.
\]

Nach \S.\,169, (7) und (10) ist:
\weq{6}{
\begin{array}{rll}
N(1-r^{-2}) &{}=q & \text{bei ungeradem } q\\
&{}=4 & \text{f\"ur } q=2,\\[4pt]
\pm\,\Delta &{}=q\,\Delta_1^2 & \text{bei ungeradem } q\\
&{}=4\,\Delta_1^2 & \text{f\"ur } q=2.
\end{array}
}


% ===== scan leaf 664, printed 643 =====
% !PRINTED 643
\weberpage{643}
\pdfbookmark[2]{§ 177. Die Primideale des Körpers Hm}{bm202}
Setzt man darin
\[
\Delta = \pm\,q^{\,q^{\varkappa-1}[\varkappa(q-1)-1]}
\]
so folgt
\weq{7}{
\begin{array}{rl}
\Delta_1 = q^{\varkappa\, q^{\varkappa-1}\frac{q-1}{2} - \frac{q^{\varkappa-1}+1}{2}} & \text{bei ungeradem } q\\
= 2^{(\varkappa-1)2^{\varkappa-2}-1} & \text{f\"ur } q=2,
\end{array}
}
so dass $\Delta_1$ immer eine Potenz von $q$ ist.

\begin{center}
\S.\,177.\\[4pt]
\gesperrt{Die Primideale im K\"orper $H_m$.}
\end{center}

Es ist jetzt die Frage zu untersuchen, in welcher Beziehung die Primideale des K\"orpers $H_m$ zu denen des K\"orpers $\Omega_m$ stehen.

Wenn $\mathfrak{P}$ irgend ein Primideal in $H_m$ vom Grade $f_1$ bedeutet, so muss $\mathfrak{P}$ Theiler einer nat\"urlichen Primzahl $p$ sein, und $\mathfrak{P}$ kann also (nach \S.\,171) nur dann durch die zweite oder eine h\"ohere Potenz eines Primfactors $\mathfrak{p}$ in $\Omega_m$ theilbar sein, wenn $p=q$ ist.

Ist aber $\mathfrak{P}$ ein Theiler von $q$, so muss es eine Potenz von $\sigma=1-r$ sein. Wir setzen etwa
\[
\mathfrak{P}=\sigma^\lambda.
\]
Nehmen wir hiervon die Norm in Bezug auf $\Omega_m$, die das Quadrat der Norm in Bezug auf $H_m$ ist, so folgt (nach \S.\,169)
\[
q^{2f_1}=q^\lambda,
\]
also $f_1=\frac12\lambda$. Da $f_1$ eine ganze Zahl ist, so muss $\lambda$ mindestens $=2$ sein. Es kann aber $\lambda$ auch nicht gr\"osser als $2$ sein, da $\sigma^2$ mit der in $H_m$ enthaltenen ganzen Zahl $(1-r)(1-r^{-1})$ associirt ist, also $\sigma^2$ durch $\mathfrak{P}$ theilbar sein muss. Demnach ist $f_1=1$, und wir erhalten den Satz:

1. \gesperrt{Die Primzahl $q$ ist die $\frac12\mu^{\text{te}}$ Potenz einer in $H_m$ existirenden Primzahl ersten Grades.}

Es sei nun $p$ von $q$ verschieden, und $\mathfrak{p}$ ein Primfactor von $\mathfrak{P}$ im K\"orper $\Omega_m$ vom Grade $f$, der durch die Substitution $(r,r^{-1})$ in $\mathfrak{p}'$ \"ubergeht. Dann ist $\mathfrak{P}$ sowohl durch $\mathfrak{p}$ als durch $\mathfrak{p}'$ theilbar, und andererseits ist $\mathfrak{p}\mathfrak{p}'$ in $H_m$ enthalten und also durch $\mathfrak{P}$ theilbar, und wir haben zu unterscheiden, ob $\mathfrak{p}, \mathfrak{p}'$ verschieden sind oder nicht.

Ist $\mathfrak{p}$ von $\mathfrak{p}'$ verschieden, so ist $\mathfrak{P}$ durch $\mathfrak{p}\mathfrak{p}'$ theilbar, und es folgt $\mathfrak{P}=\mathfrak{p}\mathfrak{p}'$, und durch Normbildung $f_1=f$.


% ===== scan leaf 665, printed 644 =====
% !PRINTED 644
\weberpage{644}
Ist aber $\mathfrak{p}=\mathfrak{p}'$, so ist $\mathfrak{P}=\mathfrak{p}$, und die Normbildung ergiebt $2f_1=f$.

Welcher dieser beiden F\"alle eintritt, das h\"angt also davon ab, ob das Primideal $\mathfrak{p}$ die Substitution $(r,r^{-1})$ gestattet oder nicht, d.\,h. ob $(r,r^{-1})$ in der Gruppe, zu der das Primideal geh\"ort, enthalten ist oder nicht. Nach \S.\,172 ist dieser Unterschied dadurch bedingt, ob es irgend einen Exponenten $h$ giebt, f\"ur den die Congruenz
\[
p^h \equiv -1 \ (\mathrm{mod}\ m)
\]
erf\"ullt ist, oder ob es keinen solchen Exponenten giebt.

Wir fassen das Ergebniss folgendermaassen zusammen:

2. \gesperrt{Die von $q$ verschiedenen nat\"urlichen Primzahlen $p$ zerfallen in zwei Arten $p_1, p_2$.}

\gesperrt{Die erste Art $p_1$ ist dadurch charakterisirt, dass f\"ur irgend einen Exponenten $h$}
\[
p_1^h \equiv -1 \ (\mathrm{mod}\ m),
\]
\gesperrt{und diese Primzahlen zerfallen, wenn sie f\"ur den Modul $m$ zum Exponenten $f$ geh\"oren und $\mu=ef$ gesetzt wird, im K\"orper $H_m$ in $e$ Primfactoren $\frac12 f^{\text{ten}}$ Grades.}

\gesperrt{Ihre Primfactoren sind auch in $\Omega_m$ nicht weiter zerlegbar.}

\gesperrt{Die Primzahlen der zweiten Art $p_2$ sind dadurch bestimmt, dass keine ihrer Potenzen nach dem Modul $m$ mit $-1$ congruent wird, und sie zerfallen im K\"orper $H_m$ in $\frac12 e$ Primfactoren $f^{\text{ten}}$ Grades. Ihre Primfactoren sind in $\Omega_m$ noch in je zwei Factoren zerlegbar.}

Bei den Primzahlen der ersten Art muss $f$ sicher eine gerade Zahl sein. Wenn $m$ ungerade ist, so gen\"ugt es auch, damit $p$ eine Primzahl von der ersten Art sei, dass sie zu einem geraden Exponenten geh\"ore; denn dann ist
\[
(p^{\frac12 f}-1)(p^{\frac12 f}+1)\equiv 0\ (\mathrm{mod}\ m).
\]
Der erste dieser Factoren $p^{\frac12 f}-1$ ist aber nicht durch $m$ theilbar, weil sonst $p$ zum Exponenten $\frac12 f$ geh\"oren w\"urde. Folglich muss $p^{\frac12 f}+1$ durch $q$ theilbar sein. Hier k\"onnen aber nicht beide Factoren $p^{\frac12 f}-1$ und $p^{\frac12 f}+1$ durch $q$ theilbar sein, weil sonst auch ihre Differenz, die $=2$ ist, durch $q$ theilbar w\"are, und folglich ist $p^{\frac12 f}+1$ durch $m$ theilbar.


% ===== scan leaf 666, printed 645 =====
% !PRINTED 645
\weberpage{645}
\pdfbookmark[2]{§ 178. Die Einheiten des Körpers Hm}{bm203}

Ist aber $m$ eine Potenz von 2, so ist die Congruenz
\[
p^h \equiv -1 \ (\mathrm{mod}\ m)
\]
nur dann möglich, wenn $p \equiv -1\ (\mathrm{mod}\ m)$ ist. Denn jede gerade Potenz einer ungeraden Zahl ist $\equiv + 1\ (\mathrm{mod}\ 8)$ und kann also nicht $\equiv -1\ (\mathrm{mod}\ m)$ sein. Ist aber $h$ ungerade, so ist
\[
\frac{p^h + 1}{p + 1} = p^{h-1} - p^{h-2} + \cdots + 1
\]
selbst ungerade, und folglich ist $p^h + 1$ durch keine höhere Potenz von 2 theilbar, als $p + 1$. Demnach können wir dem Satze 2. noch folgende nähere Bestimmung hinzufügen:

3. \gesperrt{Wenn $m$ ungerade ist, so gehören die Primzahlen $p_1$ zu einem geraden, die Primzahlen $p_2$ zu einem ungeraden Exponenten.}

\gesperrt{Ist $m$ gerade, so gehören die Primzahlen der Form $km - 1$ zur ersten und alle anderen ungeraden Primzahlen zur zweiten Art.}

\begin{center}
\S.\,178.\\[4pt]
\gesperrt{Die Einheiten des Körpers $H_m$.}
\end{center}

Die dem Körper $H_m$ angehörigen Einheiten sind von besonderer Wichtigkeit für die Anwendungen. Auf sie lassen sich auch die Einheiten des Körpers $\Omega_m$ zurückführen nach folgendem Satze:

1. \gesperrt{Jede Einheit des Körpers $\Omega_m$ ist das Product einer Einheitswurzel und einer reellen Einheit des Körpers $\Omega_m$.}

Bezeichnen wir irgend eine Einheit des Körpers $\Omega_m$ mit $\mathfrak{E}(r)$, so ist $\mathfrak{E}(r^{-1})$ der conjugirt imaginäre Werth, der gleichfalls eine Einheit darstellt, und der Quotient $\mathfrak{E}(r) : \mathfrak{E}(r^{-1})$ ist wieder eine Einheit, deren absoluter Werth
\weq{1}{\sqrt{\dfrac{\mathfrak{E}(r)}{\mathfrak{E}(r^{-1})}\ \dfrac{\mathfrak{E}(r^{-1})}{\mathfrak{E}(r)}}}
gleich 1 ist, und folglich ist dieser Quotient eine Einheitswurzel und mithin $= \pm r^h$, worin $h$ irgend ein ganzzahliger Exponent ist (§. 175, Satz 1., 2.). Wir haben also
\weq{2}{\mathfrak{E}(r) = \pm r^h \mathfrak{E}(r^{-1}).}


% ===== scan leaf 667, printed 646 =====
% !PRINTED 646
\weberpage{646}

Es ist nun im weiteren Beweise ein kleiner Unterschied zu machen, je nachdem $m$ ungerade oder gerade ist.

Ist $m$ zunächst \gesperrt{ungerade}, so können wir in (2) den Exponenten $h$ gerade annehmen, da wir ihn nöthigenfalls durch $h + m$ ersetzen können. Ferner ist in der Formel (2) nur das obere Zeichen zulässig. Denn wenn das untere Zeichen stände, so würde folgen:
\[
\mathfrak{E}(r) \equiv \mathfrak{E}(1) \equiv -\mathfrak{E}(1), \quad 2\,\mathfrak{E}(1) \equiv 0\ [\mathrm{mod}\ (1 - r)].
\]

Nach §. 169 ist aber $1 - r$ ein Theiler von $q$, also relativ prim zu 2, und daher ist $\mathfrak{E}(1)$ und folglich auch $\mathfrak{E}(r)$ durch $(1 - r)$ theilbar. Dies aber widerspricht der Voraussetzung, dass $\mathfrak{E}(r)$ eine Einheit sein soll.

Demnach ergiebt sich aus (2):
\[
r^{-1/2\,h}\, \mathfrak{E}(r) = r^{1/2\,h}\, \mathfrak{E}(r^{-1}),
\]
woraus folgt, dass $r^{-1/2\,h}\, \mathfrak{E}(r)$ eine reelle Einheit ist. Bezeichnen wir sie mit $e(r)$, so ist also
\weq{3}{\mathfrak{E}(r) = r^{1/2\,h}\, e(r),}
woraus für diesen Fall das Theorem 1. bewiesen ist.

Wenn aber zweitens $m$ \gesperrt{eine Potenz von 2 ist}, so ist $\mu = \tfrac{1}{2} m$ und $r^\mu = -1$, und zugleich mit $r$ ist $-r$ eine $m^{\text{te}}$ Einheitswurzel. Wenn wir nöthigenfalls $h$ durch $h + \tfrac{1}{2} m$ ersetzen, so können wir in (2) das obere Zeichen annehmen und setzen:
\weq{4}{\mathfrak{E}(r) = r^h\, \mathfrak{E}(r^{-1}).}

Es kommt jetzt darauf an, nachzuweisen, dass $h$ gerade sein muss.

Wenn wir $\mathfrak{E}(r)$ durch die Basis $1, r, r^2, \ldots r^{\mu-1}$ darstellen, so ergiebt sich
\weq{5}{\mathfrak{E}(r) = \sum_{0,\,\mu-1}^{s} x_s\, r^s,}
worin die $x_s$ ganze rationale Zahlen sind. Nach (4) ist aber dann
\[
\sum_{0,\,\mu-1}^{s} x_s\, r^s = \sum_{1,\,\mu-1}^{s} x_s\, r^{h-s}, % [sic] source prints lower limit 1 on the right; substitution in (4) requires 0
\]
und daraus folgt, dass $x_s = \pm x_{s'}$ ist, wenn $s + s' \equiv h\ (\mathrm{mod}\ \mu)$. Ist nun $h$ \gesperrt{ungerade}, so ist aus zwei so verbundenen Zahlen die eine gerade, die andere ungerade (da $\mu$ eine Potenz von 2 ist). Der Ausdruck (5) ergiebt also für $\mathfrak{E}(r)$ ein Aggregat von Gliedern der Form
\[
x_s\,(r^s \pm r^{s'}).
\]
%CONT>>


% ===== scan leaf 668, printed 647 =====
% !PRINTED 647
\weberpage{647}
%<<CONT
Es ist aber $r^s \pm r^{s'} = r^s (1 \pm r^{h-2s}) = r^s (1 \mp r^{\mu+h-2s})$ immer durch $1 - r$ theilbar, und daraus würde folgen, dass $\mathfrak{E}(r)$ durch $1 - r$, was keine Einheit ist, theilbar wäre, was dem Begriffe der Einheit widerspricht. Also ist die Annahme unzulässig, dass in der Formel (4) der Exponent $h$ ungerade sei, und wir können wie oben
\[
\mathfrak{E}(r) = r^{\tfrac12 h}\, e(r)
\]
setzen, worin $e(r)$ eine reelle Einheit ist. Damit ist der Satz 1. bewiesen.

Wir wollen ein System von reellen Einheiten hervorheben: Nach §. 169 sind die $\mu$ Grössen $1 - r^n$ alle unter einander associirt. Demnach ist der Quotient
\[
\frac{1 - r^n}{1 - r^{n'}},
\]
worin $n, n'$ zwei relative Primzahlen zu $m$ bedeuten, eine Einheit. Daraus leitet man aber, wenn
\[
r = e^{\frac{2\pi i}{m}}
\]
gesetzt wird, die reelle Einheit
\weq{6}{\dfrac{\sin \dfrac{n\pi}{m}}{\sin \dfrac{n'\pi}{m}} = \dfrac{r^{\frac{n}{2}} - r^{-\frac{n}{2}}}{r^{\frac{n'}{2}} - r^{-\frac{n'}{2}}} = r^{\frac{n'-n}{2}}\, \dfrac{1 - r^n}{1 - r^{n'}}}
her. Ist $m$ gerade, so kann man $n' = n + \dfrac{m}{2}$ setzen, und erhält für diesen Fall die reellen Einheiten
\weq{7}{\tau_n = \mathrm{tang}\,\dfrac{n\pi}{m},}
worin $n$ jede ungerade Zahl bedeuten kann.

Ersetzt man $n$ durch $\tfrac{1}{2} m - n$, so geht $\tau_n$ in den reciproken Werth über. Lässt man $n$ die Reihe der Zahlen
\[
1,\ 3,\ 5,\ \cdots,\ \frac{m}{4} - 1
\]
durchlaufen, so erhält $\tau_n$ lauter positive echt gebrochene Werthe.

\begin{center}
\rule{0.2\linewidth}{0.4pt}
\end{center}


% ===== scan leaf 669, printed 648 =====
\weberpage{648}% page number not printed on this page
\pdfbookmark[1]{Zweiundzwanzigster Abschnitt. Abel'sche Körper und Kreistheilungskörper.}{bm204}
\pdfbookmark[2]{§ 179. Einfache Abel'sche Körper}{bm205}
% !PRINTED none

\begin{center}
\gesperrt{Zweiundzwanzigster Abschnitt.}

\gesperrt{Abel'sche Körper und Kreistheilungskörper.}

\rule{3cm}{0.4pt}

\S.\,179.

\gesperrt{Einfache Abel'sche Körper.}
\end{center}

Wir wenden uns nun zu einer der interessantesten Anwendungen der Theorie der algebraischen Zahlen.

\begingroup
\renewcommand{\thefootnote}{\arabic{footnote})}
Es handelt sich um den Beweis des allgemeinen Satzes\footnote{\gesperrt{Kronecker} hat diesen Satz zuerst ausgesprochen, aber keinen vollständigen Beweis dafür veröffentlicht. Den ersten Beweis hat der Verfasser des vorliegenden Werkes in Bd. 8 der Acta Mathematica bekannt gemacht (1886). In neuester Zeit ist ein zweiter Beweis von \gesperrt{Hilbert} veröffentlicht, der sich auf die im neunzehnten Abschnitte entwickelten Sätze stützt (Nachrichten d. Gesellschaft d. Wissenschaften zu Göttingen, 1896, Heft 1).}:
\endgroup

I. \gesperrt{Alle im absoluten Rationalitätsbereich Abel'schen Zahlkörper sind Kreistheilungskörper.}

Haben wir diesen Satz bewiesen, so gewinnen die Untersuchungen des dritten Abschnittes über Kreistheilungskörper ein erhöhtes Interesse, weil damit gezeigt ist, dass auf dem dort angegebenen Wege nicht nur alle Kreistheilungskörper, sondern alle \gesperrt{Abel'schen} Körper überhaupt gefunden werden.

Es sei $\Omega = R(x)$ ein \gesperrt{Abel'scher} Körper $m^{\text{ten}}$ Grades, d. h. ein Körper, der aus allen rationalen Functionen einer Wurzel $x$ einer irreduciblen \gesperrt{Abel'schen} Gleichung $m^{\text{ten}}$ Grades besteht, wenn als Rationalitätsbereich der Körper $R$ der rationalen Zahlen betrachtet wird. Die \gesperrt{Galois'sche} Gruppe $\mathfrak{G}$ dieser Gleichung, die also eine \gesperrt{Abel'sche} Gruppe ist und denselben Grad $m$ hat, wie der Körper $\Omega$, ist auch die Gruppe des Körpers $\Omega$.

Ist $x_1$ eine nicht primitive Zahl aus $\Omega$, so ist $\Omega_1 = R(x_1)$ gleichfalls ein \gesperrt{Abel'scher} Körper, den wir einen echten Theiler von $\Omega$ nennen und dessen Grad ein echter Theiler des Grades von $\Omega$ ist. Denn ist $\mathfrak{A}$ die Gruppe, zu der die Zahl $x_1$ gehört,
%CONT>>


% ===== scan leaf 670, printed 649 =====
% !PRINTED 649
\weberpage{649}
%<<CONT
so ist $\mathfrak{G}\,|\,\mathfrak{A}$ (\S.\,4 dieses Bandes) die Gruppe des Körpers $\Omega_1$ und
dies ist zugleich mit $\mathfrak{G}$ eine \gesperrt{Abel'sche} Gruppe. Ein Theiler
von $\Omega$, der mit $\Omega$ von gleichem Grade ist, ist mit $\Omega$ identisch
(vgl. Bd. I, \S.\,144, 149, 156).

Giebt es nun zwei nicht primitive Zahlen $x_1, x_2$ in $\Omega$ von
der Art, dass
\[
\Omega = R(x_1, x_2)
\]
gesetzt werden kann, so sind die Körper $\Omega_1 = R(x_1)$, $\Omega_2 = R(x_2)$
echte Theiler von $\Omega$, und $\Omega$ heisst aus den beiden Körpern
$\Omega_1, \Omega_2$ \gesperrt{zusammengesetzt}.

Gestattet der Körper $\Omega$ keine solche Darstellung, so heisst
er \gesperrt{einfach}.

Wenn also als bewiesen vorausgesetzt wird, dass $\Omega_1, \Omega_2$
Kreistheilungskörper sind, d.\,h. dass alle ihre Zahlen rational
durch Einheitswurzeln darstellbar sind, so folgt dasselbe für $\Omega$.

Wenn einer der Körper $\Omega_1, \Omega_2$ zerlegbar ist, so können
wir die Zerlegung wiederholen, und müssen, da die Grade abnehmende
 ganze positive Zahlen sind, endlich zu einfachen Körpern
 gelangen.

\gesperrt{Das Theorem I. braucht also nur für einfache Abel'sche
Körper bewiesen zu werden.}

Es ist aber daran zu erinnern, dass die Zerlegbarkeit eines
Körpers $\Omega$ keineswegs schon aus der Existenz eines Theilers
folgt, und dass bei einem zerlegbaren Körper die Zerlegung in
einfache Körper auf ganz verschiedene Arten geschehen kann, so
dass bei den Körpern nicht die Gesetze der Theilbarkeit gelten,
wie im Gebiete der Zahlen.

Ein Kennzeichen für die Einfachheit eines Körpers können
wir aus seiner Gruppe ableiten.

Wenn die Gruppe $\mathfrak{G}$ zwei echte Theiler $\mathfrak{A}$ und $\mathfrak{B}$ von der
Art hat, dass jedes Element ein- und nur einmal aus einem Element
 von $\mathfrak{A}$ und einem Element von $\mathfrak{B}$ zusammengesetzt werden
kann, so nennen wir die Gruppe $\mathfrak{G}$ zerlegbar in die beiden
Componenten $\mathfrak{A}, \mathfrak{B}$; $\mathfrak{G}$ ist also zerlegbar, wenn in dem nach der
Composition der Theile (\S.\,4) gebildeten Producte
\weq{1}{\mathfrak{G} = \mathfrak{A}\mathfrak{B}}
jedes Element von $\mathfrak{G}$ ein- und nur einmal erscheint. Der Grad
von $\mathfrak{G}$ ist dann gleich dem Producte der Grade von $\mathfrak{A}$ und $\mathfrak{B}$.

Giebt es solche Theiler nicht, so heisst die Gruppe $\mathfrak{G}$
unzerlegbar.


% ===== scan leaf 671, printed 650 =====
% !PRINTED 650
\weberpage{650}

Dann können wir den Satz beweisen:

1. \gesperrt{Ist die Gruppe eines Abel'schen Körpers zerlegbar,
so ist auch der Körper zusammengesetzt.}

Nehmen wir, um dies zu beweisen, an, die Gruppe $\mathfrak{G}$ des
Körpers $\Omega$ sei nach der Formel (1) zerlegbar und $m, a, b$
seien die Grade von $\mathfrak{G}, \mathfrak{A}, \mathfrak{B}$. Wir nehmen eine zur Gruppe $\mathfrak{B}$
\gesperrt{gehörige} Zahl $\xi$ des Körpers $\Omega$, die also durch die Substitutionen
 von $\mathfrak{G}$ in $a$ verschiedene conjugirte Werthe übergeht und
ebenso eine zu $\mathfrak{A}$ gehörige $b$-werthige Zahl $\eta$. Wir können
dann eine Zahl
\[
x = \alpha \xi + \beta \eta
\]
mit \gesperrt{rationalen} Zahlencoëfficienten $\alpha, \beta$ bilden, die $m$ verschiedene
 conjugirte Werthe hat, und dann ist $\Omega = R(x)$. Also
ist $\Omega$ aus $R(\xi)$ und $R(\eta)$ zusammengesetzt (Bd. I, \S.\,143).

Erinnern wir uns nun an die in \S.\,9, 10 dieses Bandes abgeleitete
 Darstellung einer Abel'schen Gruppe durch eine Basis, so
ergiebt sich durch wiederholte Anwendung des eben Bewiesenen:

2. \gesperrt{Jeder Abel'sche Körper $\Omega$ lässt sich aus solchen
zusammensetzen, deren Grad eine Primzahlpotenz
 und deren Gruppe cyklisch ist. Die Grade
dieser zusammensetzenden Körper sind die Invarianten
 der Gruppe von $\Omega$.}

Diese Sätze werden durch den folgenden ergänzt:

3. \gesperrt{Ein Abel'scher Körper von Primzahlpotenzgrad
mit cyklischer Gruppe ist einfach.}

Wenn die Gruppe $\mathfrak{G}$ cyklisch ist, so lassen sich ihre Elemente
 durch die Potenzen einer Basis, $G^\gamma$, darstellen, wenn $\gamma$
ein volles Restsystem nach dem Modul $m$ durchläuft. Wir erhalten
 alle Theiler $\mathfrak{A}$ von $\mathfrak{G}$ dargestellt durch
\[
\mathfrak{A} = G^{b\gamma},
\]
wenn $b$ ein Theiler von $m$ ist und $\gamma$ ein volles Restsystem nach
dem Modul $a = m : b$ durchläuft.

Ist
\[
\mathfrak{A}' = G^{b'\gamma}
\]
ein zweiter Theiler von $\mathfrak{G}$ vom Grade $a'$ und ist
\[
b' \leqq b, \quad a' \geqq a,
\]
so ist, wenn wir jetzt annehmen, dass $m$ und mithin auch
$a, b, a', b'$ Potenzen einer Primzahl $q$ sind, $b'$ ein Theiler von $b$
und folglich $\mathfrak{A}$ ein Theiler von $\mathfrak{A}'$.


% ===== scan leaf 672, printed 651 =====
% !PRINTED 651
\weberpage{651}
\pdfbookmark[2]{§ 180. Die Resolventen}{bm206}

Wenn also $\xi$ eine Zahl des Körpers $\Omega$ ist, die zu der
Gruppe $\mathfrak{A}$ gehört, so ist $\xi$ eine primitive Zahl des Körpers $b^{\mathrm{ten}}$
Grades $R(\xi)$, und in diesem Körper ist auch jede Zahl $\xi'$ enthalten,
die die Substitutionen der Gruppe $\mathfrak{A}'$ gestattet. Sind
also $\mathfrak{A}, \mathfrak{B}$ und $\mathfrak{A}', \mathfrak{B}'$ echte Theiler von $\mathfrak{G}$, so ist $R(\xi)$ von $\Omega$
verschieden, und $R(\xi)$ wird durch Adjunction von $\xi'$ nicht erweitert.
Es kann also auch nicht $\Omega$ aus $R(\xi)$ und $R(\xi')$ zusammengesetzt
werden, wodurch 3. bewiesen ist.

\begin{center}
\S.\,180.

\gesperrt{Die Resolventen.}
\end{center}

Nach dem, was jetzt bewiesen ist, können wir uns in den
weiteren Betrachtungen, die den Beweis des Satzes I. zum Ziele
haben, auf solche \gesperrt{Abel'sche} Körper beschränken, deren Grad
eine Primzahlpotenz und deren Gruppe cyklisch ist. Aus diesem
Grunde genügte es auch für die beabsichtigte Anwendung, dass
wir im vorigen Abschnitte uns auf die Betrachtung der Kreistheilungskörper
 $\Omega_m$ beschränkt haben, in denen $m$ eine Primzahlpotenz
 war. Wir behalten so viel als möglich die dort
gebrauchte Bezeichnung bei und setzen
\weq{1}{m = q^{\varkappa},}
worin $q$ eine Primzahl, $\varkappa$ ein positiver Exponent ist, und in dem
Falle, dass $q = 2$ ist, nehmen wir $\varkappa > 1$ an. Mit $n$ bezeichnen
wir jede durch $q$ nicht theilbare Zahl. Es ist $r$ eine primitive
$m^{\mathrm{te}}$ Einheitswurzel und $\Omega_m$ der Körper der rationalen Functionen
 von $r$.

Es sei nun $\Omega = R(x)$ ein einfacher \gesperrt{Abel'scher} Körper $m^{\mathrm{ten}}$
Grades, so dass $x$ Wurzel einer rationalen irreduciblen cyklischen
Gleichung $m^{\mathrm{ten}}$ Grades ist. Bezeichnen wir die Wurzeln dieser
Gleichung in geeigneter Reihenfolge mit
\weq{2}{x_0, x_1, x_2, \ldots, x_{m-1},}
und nehmen $x_m = x_0$, $x_{m+1} = x_1, \ldots$ an, so ist
\[
x_1 = \Phi(x_0), \quad x_2 = \Phi(x_1), \quad \ldots, \quad x_0 = \Phi(x_{m-1}),
\]
worin $\Phi$ eine ganze Function mit rationalen Coëfficienten bedeutet,
und die Gruppe des Körpers $\Omega$ besteht aus den Potenzen
der Substitution $(x_0, x_1)$ oder aus den cyklischen Permutationen
%CONT>>


% ===== scan leaf 673, printed 652 =====
% !PRINTED 652
\weberpage{652}
%<<CONT
der Wurzeln (2). Jede cyklische Function dieser Grössen ist
eine rationale Zahl. Was wir zu beweisen haben, ist, dass sich
die $x$ rational durch Einheitswurzeln ausdrücken lassen. Adjungiren
wir dem Körper $\Omega$ noch die $m^{\mathrm{te}}$ Einheitswurzel $r$, so entsteht
ein Körper $R(x, r)$, der die beiden Körper $\Omega = R(x)$ und
$R(r) = \Omega_m$ als Theiler enthält.

\gesperrt{Wenn eine Zahl} $\omega = \Phi(x, r)$ des Körpers $R(x, r)$ die
sämmtlichen Substitutionen $(x_0, x_1), (x_0, x_2), \ldots, (x_0, x_{m-1})$ gestattet,
so ist sie im Körper $\Omega_m$ enthalten; denn dann ist
\[
m\,\omega = \Phi(x_0, r) + \Phi(x_1, r) + \cdots + \Phi(x_{m-1}, r),
\]
und darin sind die Coëfficienten der Potenzen von $r$ nicht nur
cyklische, sondern sogar symmetrische Functionen der Wurzeln
(2), also rationale Zahlen. Dies gilt selbst noch in dem
Falle, dass die Gleichung für $x$ durch Adjunction von $r$ reducirt
wird.

Dies wenden wir an auf die \gesperrt{Resolventen}
\weq{3}{(r, x_0) = x_0 + r x_1 + r^2 x_2 + \cdots + r^{m-1} x_{m-1},}
und allgemeiner, wenn $s$ einen beliebigen ganzzahligen Exponenten
 bedeutet:
\weq{4}{(r^s, x_0) = x_0 + r^s x_1 + r^{2s} x_2 + \cdots + r^{(m-1)s} x_{m-1}.}

Durch diese Resolventen lässt sich $x_0$ selbst wieder rational
ausdrücken, nämlich:
\weq{5}{m x_0 = \sum\limits_{0,\,m-1}^{s} (r^s, x_0),}
und wenn wir also nachweisen können, dass alle diese Resolventen
 $(r^s, x_0)$ Kreistheilungszahlen sind, so haben wir unser
Ziel erreicht. Wenn wir
\weq{6}{(r^s, x_{\varkappa}) = x_{\varkappa} + r^s x_{\varkappa+1} + r^{2s} x_{\varkappa+2} + \cdots + r^{(m-1)s} x_{\varkappa+m-1}}
setzen, so geht $(r^s, x_{\varkappa})$ aus $(r^s, x_0)$ durch die Substitution
$(x_0, x_{\varkappa})$ hervor, und es besteht die fundamentale Relation
\weq{7}{r^{\varkappa s} (r^s, x_{\varkappa}) = (r^s, x_0).}

Alle diese Resolventen sind Zahlen des Körpers $R(x, r)$;
wir betrachten vorzugsweise die beiden folgenden Verbindungen:
\weq{8}{(r^s, x_0)^m = \omega_s,}
\weq{9}{(r^s, x_0)^{-1} (r, x_0)^s = \alpha.}


% ===== scan leaf 674, printed 653 =====
% !PRINTED 653
\weberpage{653}

Diese beiden Zahlen gestatten, wie aus (7) unmittelbar hervorgeht, die Substitutionen $(x_0, x_\varkappa)$, und sind also Zahlen des Körpers $\Omega_m$.

Ist $n$ durch $q$ nicht theilbar, so kann wegen der Irreducibilität der Kreistheilungsgleichung keine von den Resolventen $(r^n, x_0)$ verschwinden, wenn sie nicht alle verschwinden; denn die Gleichung
\[
(r, x_0)^m = 0
\]
wäre eine Gleichung für $r$ mit rationalen Co\"efficienten, die alle Substitutionen $(r, r^n)$ gestattet.

Ebenso kann keine der Resolventen $(r^s, x_0)$ verschwinden, wenn nicht die ganze Schaar verschwindet, in der $s$ einen bestimmten grössten gemeinschaftlichen Theiler mit $m$ hat.

Wenn alle $(r, x_0)$ verschwinden, so ist die Formel (9) nicht anwendbar. Dann können wir aber die Resolventen $(r^q, x_0)$ direct betrachten, die zu einem Körper vom Grade $m q^{-1}$ gehören, der aus
\[
y_0 = x_0 + x_q + x_{2q} + \cdots + x_{m-q}
\]
hervorgeht, und wenn auch diese Grössen verschwinden, so gehen wir zu den Resolventen $(r^{q^2}, x_0)$ über u.\,s.\,f.

Wenn aber $(r, x_0)$ nicht verschwindet, so lassen sich durch die Formel (9) alle $(r^s, x_0)$ rational durch $(r, x_0)$ und durch Kreistheilungszahlen darstellen.

Demnach ist es für unseren Zweck ausreichend, wenn wir beweisen können, dass die \gesperrt{nicht verschwindenden Resolventen} $(r, x_0)$ \gesperrt{Kreistheilungszahlen sind}.

Lassen wir $n$ ein volles System incongruenter, durch $q$ nicht theilbarer Zahlen durchlaufen, so ist
\[
\sum n \equiv 0 \ (\mathrm{mod}\ m),
\]
denn die Zahlen $n$ zerfallen in Paare einander zu $m$ ergänzender Zahlen. Daraus folgt nach (7), dass das Product
\weq{10}{\prod^n (r^n, x_0) = a}
eine Zahl in $\Omega_m$ ist. Da aber diese Zahl sich überdies nicht ändert, wenn irgend eine der Substitutionen $(r, r^n)$ darin ausgeführt wird, so ist $a$ eine \gesperrt{rationale Zahl}.


% ===== scan leaf 675, printed 654 =====
% !PRINTED 654
\weberpage{654}
\pdfbookmark[2]{§ 181. Vorbereitung zum Beweis}{bm207}

\begin{center}
\S.\,181.\\
\gesperrt{Vorbereitung zum Beweis.}
\end{center}

Nach den Ausführungen des vorigen Paragraphen ist zum Beweis des grossen Theorems I nur noch der Nachweis erforderlich, dass die in $\Omega_m$ enthaltene Zahl
\[
\omega = (r, x_0)^m
\]
die $m^{\text{te}}$ Potenz einer Kreistheilungszahl ist.

1. \gesperrt{Die wesentliche Eigenschaft der zu $\omega$ conjugirten Zahlen $\omega_n$, auf die sich der Beweis stützt, ist die, dass}
\weq{1}{\omega_n^{-1} \omega_1^n = \alpha^m}
\gesperrt{die $m^{\text{te}}$ Potenz einer Zahl in $\Omega_m$ ist.}

Hierin kann $n$ jede durch $q$ nicht theilbare Zahl sein. Bilden wir von der rechten und linken Seite von (1) die Normen im Körper $\Omega_m$, so ergiebt sich, wenn wir die Norm der Zahl $\alpha$, die eine rationale Zahl ist, mit $a$ bezeichnen, da jedes $\omega_n$ dieselbe Norm wie $\omega$ hat,
\[
N(\omega)^{n-1} = a^m.
\]
Wenn nun $m$ ungerade ist, so können wir hierin $n = 2$ annehmen, und erhalten als \gesperrt{Folgerung aus 1.}:

2. \gesperrt{Die Norm von $\omega$ ist die $m^{\text{te}}$ Potenz einer rationalen Zahl}
\weq{2}{N(\omega) = a^m.}

% [sic] The German source prints the anomalous punctuation ``Potenz von 2.,''; retained verbatim.
Ist $m$ eine Potenz von 2., so ist diese Folgerung nicht mehr zulässig. Es ergiebt sich dann aus 1. nur, dass $N(\omega)^2$ die $m^{\text{te}}$ Potenz einer rationalen Zahl ist. Gleichwohl müssen die Zahlen $\omega$, wie aus (10) des vorigen Paragraphen zu sehen ist, auch in diesem Falle noch der Bedingung 2. genügen, die dann aber nicht mehr von selbst schon in 1. enthalten ist, sondern als neue Eigenschaft hinzukommt.

\gesperrt{Was wir zu beweisen haben, ist, dass aus 1. und 2. folgt, dass auch $\omega$ selbst die $m^{\text{te}}$ Potenz einer Kreistheilungszahl, wenn auch in einem höheren Körper, ist.}

Dazu ist aber erforderlich, die Zerlegung der Zahlen $\omega$ in ihre Primfactoren im Körper $\Omega_m$ zu untersuchen.

Die Primzahl $q$ ist, wie im \S.\,169 nachgewiesen, mit der $\varphi(m)^{\text{ten}}$ Potenz einer in $\Omega_m$ existirenden Primzahl $\sigma = 1 - r$
%CONT>>


% ===== scan leaf 676, printed 655 =====
% !PRINTED 655
\weberpage{655}
%<<CONT
associirt. Ist also $\sigma^k$ die in $\omega$ enthaltene Potenz von $\sigma$, so setzen wir
\[
\omega = \sigma^k \omega',
\]
worin $\omega'$ mit der Primzahl $q$ theilerfremd ist. Die Bildung der Norm ergiebt nach 2.
\[
a^m = q^k b,
\]
worin $b$ eine rationale Zahl ist, die in einfachster Gestalt die Primzahl $q$ weder im Zähler noch im Nenner enthält. Daraus ergiebt sich aber, dass $k$ durch $m$ theilbar sein muss, und daraus das erste Resultat:

3. \gesperrt{Der Exponent der in $\omega$ enthaltenen Primzahl $\sigma$ ist immer durch $m$ theilbar.}

Es sei nun ferner $p$ eine von $q$ verschiedene Primzahl und $\mathfrak{p}_\xi$ ein Primfactor von $p$ im Körper $\Omega_m$. Wir lassen $\xi$ die im \S.\,172 in den verschiedenen Fällen näher bestimmte Zahlenreihe $\xi_1, \xi_2, \ldots, \xi_e$ durchlaufen, so dass
\weq{3}{p = \mathfrak{p}_{\xi_1}\mathfrak{p}_{\xi_2}\ldots\mathfrak{p}_{\xi_e}}
wird, und nehmen an, es sei $\mathfrak{p}_\xi^{k_\xi}$ die in $\omega$ enthaltene Potenz von $\mathfrak{p}_\xi$. Die in $\omega_n$ enthaltene Potenz von $\mathfrak{p}_{n\xi}$ ist dann $\mathfrak{p}_{n\xi}^{k_\xi}$, und wenn wir $n'$ aus der Congruenz
\weq{4}{n n' \equiv 1 \ (\mathrm{mod}\ m)}
bestimmen, so ist $\mathfrak{p}_\xi^{k_{n'\xi}}$ die in $\omega_n$ enthaltene Potenz von $\mathfrak{p}_\xi$. Demnach ist in
\weq{5}{\omega_n^{-1} \omega_1^n \text{ die Potenz } \mathfrak{p}_\xi^{n k_\xi - k_{n'\xi}}}
enthalten, und wegen 1. muss der Exponent dieser Potenz durch $m$ theilbar sein.

Wir erhalten also für jedes durch $q$ nicht theilbare $n$ die Congruenz
\weq{6}{n k_\xi \equiv k_{n'\xi} \ (\mathrm{mod}\ m),}
in der wir nun auch $\xi$ durch $n\xi$ ersetzen können, wodurch sich
\weq{7}{n k_{n\xi} \equiv k_\xi \ (\mathrm{mod}\ m)}
ergiebt. Nehmen wir hierin $\xi = 1$ und setzen dann $\xi, \xi'$ für $n, n'$ und $k$ für $k_1$, so folgt aus (7)
\weq{8}{k_\xi \equiv \xi' k \ (\mathrm{mod}\ m),}
worin nun $k$ für alle $\xi$ denselben Werth hat und $\xi, \xi'$ durch die Congruenz
\[
\xi \xi' \equiv 1 \ (\mathrm{mod}\ m)
\]
zusammenhängen.


% ===== scan leaf 677, printed 656 =====
% !PRINTED 656
\weberpage{656}

Nehmen wir $n = p$ an, so wird $\mathfrak{p}_{p\xi} = \mathfrak{p}_\xi$ (\S.\,172), also ist auch $k_{p\xi} = k_\xi$, und aus (7) und (8) folgt:
\weq{9}{k(p-1) \equiv 0 \ (\mathrm{mod}\ m).}

Hieraus ergeben sich nun wichtige Folgerungen: Wenn zunächst $p - 1$ nicht durch $q$ theilbar ist, so folgt, dass $k$ durch $m$ theilbar sein muss, und wir haben also den Satz:

4. \gesperrt{Die Primfactoren einer Primzahl $p$, die nach dem Modul $q$ nicht mit 1 congruent ist, sind in $\omega$ nur in solchen Potenzen enthalten, deren Exponenten durch $m$ theilbar sind.}

Wenn in dem Falle, wo $m$ eine Potenz von 2 ist, $p - 1$ durch 2, aber nicht durch 4 theilbar ist, wenn also $p \equiv -1 \ (\mathrm{mod}\ 4)$ ist, so folgt aus (9) nur, dass $k$ durch $\tfrac{1}{2} m$, nicht aber, dass $k$ durch $m$ theilbar ist. Der Exponent der in $\omega$ enthaltenen Potenz von $\mathfrak{p}_\xi$ ist, von Vielfachen von $m$ abgesehen, gleich $k$, und wenn $k$ durch $m$ theilbar ist, so ist alles wie im Falle 4. Ist aber $k$ nur durch $\tfrac{1}{2} m$, nicht durch $m$ theilbar, also
\[
k \equiv \tfrac{1}{2} m \ (\mathrm{mod}\ m),
\]
so ist $k - \tfrac{1}{2} m$ durch $m$ theilbar, und wir können mit Rücksicht auf die Zerlegung (3) den Satz so formuliren:

5. \gesperrt{Ist $m$ eine Potenz von 2 und $p$ eine Primzahl congruent mit $-1$ nach dem Modul 4, so lässt sich der Exponent $h$ so bestimmen, dass alle Primfactoren von $p$ in}
\[
\omega \sqrt{p}^{-hm}
\]
\gesperrt{nur zu solchen Potenzen enthalten sind, deren Exponent durch $m$ theilbar ist.}

Es sei endlich $p - 1$ durch $q$, oder, falls $m$ gerade ist, durch 4 theilbar. Ist $m_1$ der grösste gemeinschaftliche Theiler von $m$ und $p - 1$, und ist
\[
m = m_1 m_2,
\]
so muss $k$ nach (9) durch $m_2$ theilbar sein, und wenn wir
\[
k = h m_2
\]
setzen, so ist in $\omega$ nach (8) das Product enthalten:
\weq{10}{\left(\prod^\xi \mathfrak{p}_\xi^{\xi'}\right)^{h m_2}.}
Ausserdem kommen in $\omega$ die Primfactoren $\mathfrak{p}_\xi$ nur noch in solchen Potenzen vor, deren Exponenten durch $m$ theilbar sind.


% ===== scan leaf 678, printed 657 =====
% !PRINTED 657
\weberpage{657}

In diesem Falle durchl\"auft nun $\xi$ nach \S.\,172, III ein volles
System durch $q$ nicht theilbarer Reste nach dem Modul $m_1$, und
wir k\"onnen in (10) unter $\xi'$ die kleinste positive L\"osung der
Congruenz
\[
\xi\xi' \equiv 1\ (\mathrm{mod}\ m_1)
\]
verstehen, weil. wenn wir die Exponenten $\xi'$ in (10) nach dem % [source sic: full stop after ``weil'' as printed; a comma is grammatically expected]
Modul $m_1$ ver\"andern, nur $m^{\text{te}}$ Potenzen der
Primfactoren $\mathfrak{p}$ hinzutreten.

Dann k\"onnen wir auf das Product (10) das \gesperrt{Kummer'sche}
Theorem [\S.\,174, (16)] anwenden, indem wir dort $m$ durch $m_1$
ersetzen, und erhalten, wenn wir unter $r_1 = r^{m_2}$ eine
$m_1^{\text{te}}$ Einheitswurzel, unter den $\eta$ gewisse Perioden
der $p^{\text{ten}}$ Einheitswurzeln verstehen:
\[
\prod^{\xi} \mathfrak{p}_\xi^{\xi'} = (r_1, \eta)^{m_1},
\]
und folglich
\weq{11}{(\prod^{\xi} \mathfrak{p}_\xi^{\xi'})^{hm_2} = (r_1, \eta)^{hm}.}
Demnach haben wir das Theorem:

6. \gesperrt{Ist $p$ eine Primzahl und $p-1$ durch $q$, oder,
wenn $q=2$ ist, durch 4 theilbar, so l\"asst sich der positive
Exponent $h$ so bestimmen, dass die in $\Omega_m$ enthaltene Zahl}
\[
\omega (r_1, \eta)^{-hm}
\]
\gesperrt{die Primfactoren von $p$ nur zu solchen Potenzen
enth\"alt, deren Exponent durch $m$ theilbar ist.}

Es ist jetzt auch nicht mehr n\"othig, den Satz 5. von dem
Satze 6. zu unterscheiden; denn f\"ur $m_1 = 2$ geht 6. in 5.
\"uber, weil dann $r_1 = -1$ wird, und nach Bd. I, \S.\,171
\[
(r_1, \eta)^m = (-1, \eta)^m = \sqrt{p}^{\,m}
\]
ist $^1$).

Fassen wir das Ergebniss der S\"atze 3. bis 6. in eine Formel
zusammen, so ergiebt sich, wenn $\varepsilon$ ein
Einheitsfunctional, $\varphi$ irgend ein Functional des K\"orpers
$\Omega_m$ bedeutet, und $p$ eine Reihe von Primzahlen durchl\"auft,
die congruent mit 1 nach dem Modul $q$ sind, worunter dieselbe
Primzahl $p$ auch mehrmals vorkommen kann,
\weq{12}{\omega = (r, x_0)^m = \varepsilon \varphi^m \prod^{p} (r^{m_2}, \eta)^m.}

\noindent
$^1$) In des Verfassers Abhandlung in Bd. 8 der Acta mathematica ist
es vers\"aumt, den Fall des Satzes 5. besonders hervorzuheben.


% ===== scan leaf 679, printed 658 =====
% !PRINTED 658
\weberpage{658}

Die in diesen Formeln vorkommenden Resolventen der Kreistheilung
$(r^{m_2}, \eta)$ sind Kreistheilungszahlen in einem h\"oheren
K\"orper. Ihre $m^{\text{ten}}$ Potenzen sind aber im K\"orper
$\Omega_m$ enthaltene Zahlen, die durch die Substitution $(r, r^n)$
in
\[
(r^{n m_2}, \eta)^m
\]
\"ubergehen. Ebenso sind die Verbindungen
\weq{13}{(r^{n m_2}, \eta)^{-1} (r^{m_2}, \eta)^n}
im K\"orper $\Omega_m$ enthalten [\S.\,174, (9), (11)]. Diese
Resolventen sind specielle F\"alle der allgemeinen Resolvente
$(r, x_0)$. Aus (12) erhalten wir, wenn wir mit $\varphi_n,
\varepsilon_n$ die conjugirten Functionale zu $\varphi$ und
$\varepsilon$ bezeichnen
\weq{14}{\omega_n = \varepsilon_n \varphi_n^m \prod^{p} (r^{m_2 n}, \eta)^m.}
Diese Formel lehrt uns die wichtigsten Eigenschaften der
Functionale $\varphi_n$ und $\varepsilon_n$ kennen, zun\"achst:

7. \gesperrt{Die Functionale $\varphi_n^m$ sind mit Zahlen des
K\"orpers $\Omega_m$ associirt, geh\"oren also der Hauptclasse an
(\S.\,152).}

Nach 1. ist $\omega_n^{-1} \omega_1^n = \alpha^m$ die % [sic] source omits "in" before $\Omega_m$
$m^{\text{te}}$ Potenz einer Zahl $\Omega_m$. Da auch die
Verbindungen (13) Zahlen in $\Omega_m$ sind, so k\"onnen wir
\[
\prod^{p} (r^{m_2 n}, \eta)^{-1} (r^{m_2}, \eta)^n = \frac{\alpha}{\beta}
\]
setzen, und erhalten aus (14)
\[
\varepsilon_n^{-1} \varepsilon_1^n (\varphi_n^{-1} \varphi_1^n)^m = \beta^m,
\]
worin $\beta$ eine Zahl in $\Omega_m$ ist.

Daraus geht hervor, dass das Product
\[
\beta \varphi_n \varphi_1^{-n} = \varepsilon
\]
ein Einheitsfunctional des K\"orpers $\Omega_m$ ist, und daraus
ergeben sich f\"ur die Functionale $\varphi_n$ und die
Einheitsfunctionale $\varepsilon_n$ die folgenden beiden S\"atze:

8. \gesperrt{Die conjugirten Functionale $\varphi_n$ haben die
Eigenschaft, dass alle Producte}
\[
\varphi_n \varphi_1^{-n}
\]
\gesperrt{der Hauptclasse angeh\"oren.}

9. \gesperrt{Die Einheitsfunctionale $\varepsilon_n$ haben die
Eigenschaft, dass die Producte}
\[
\varepsilon_n \varepsilon_1^{-n} = \varepsilon^m
\]
\gesperrt{$m^{\text{te}}$ Potenzen von Einheitsfunctionalen sind.}


% ===== scan leaf 680, printed 659 =====
% !PRINTED 659
\weberpage{659}
\pdfbookmark[2]{§ 182. Beweis des ersten Hülfssatzes für ein ungerades m}{bm208}

Es kommt nun f\"ur den Beweis des Haupttheorems I. alles auf den
Beweis der beiden S\"atze an, die aus den Eigenschaften der
Functionale $\varphi$ und $\varepsilon$ zu folgern sind:

A. \gesperrt{Die Functionale $\varphi_n$ geh\"oren der Hauptclasse an.}

K\"onnen wir dieses Lemma voraussetzen, so ist $\varphi_n$ mit
einer Zahl $\alpha_n$ associirt, und wir k\"onnen die Formel (14)
mit etwas ver\"anderter Bedeutung von $\varepsilon_n$ so
darstellen:
\weq{15}{\omega_n = \varepsilon_n \alpha_n^m \prod^{p} (r^{m_2 n}, \eta)^m.}
In dieser Formel ist aber $\varepsilon_n$ eine \gesperrt{numerische Einheit}
des K\"orpers $\Omega_m$, die dem Satze 9. gen\"ugt.

K\"onnen wir also aus dem Satze 9. noch das zweite Lemma
beweisen:

B. \gesperrt{Eine numerische Einheit $\varepsilon_n$ des K\"orpers
$\Omega_m$, die dem Satze 9. gen\"ugt, ist die $m^{\text{te}}$
Potenz einer Einheit in $\Omega_m$, multiplicirt mit einer Potenz
von $r$;}

so k\"onnen wir in (15) die $m^{\text{te}}$ Wurzel ziehen, und
erhalten, wenn mit $\varrho$ eine Einheitswurzel von
m\"oglicherweise h\"oherem Grade als $m$ und mit $\alpha$ eine
Zahl des K\"orpers $\Omega_m$ bezeichnet wird:
\weq{16}{(r, x_0) = \varrho \alpha \prod^{p} (r^{m_2}, \eta),}
wo nun rechts lauter Kreistheilungszahlen stehen und wodurch
daher das Theorem I. erwiesen ist.

\begin{center}
\S.\,182.

\gesperrt{Beweis des ersten H\"ulfssatzes f\"ur ein ungerades $m$.}
\end{center}

Wir haben im vorigen Paragraphen den Beweis des Theorems I. von
zwei H\"ulfss\"atzen abh\"angig gemacht, deren Beweis nun ferner
zu suchen ist.

Wir formuliren den ersten dieser H\"ulfss\"atze so:

a) \gesperrt{Ist $\varphi_n$ ein System von Null verschiedener
conjugirter Functionale des K\"orpers $\Omega_m$, von denen
bekannt ist, dass}
\[
\varphi_n^m \text{ \gesperrt{und} } \varphi_n^{-1} \varphi_1^n
\]
\gesperrt{der Hauptclasse angeh\"oren, so geh\"oren die
Functionale $\varphi$ selbst der Hauptclasse an.}


% ===== scan leaf 681, printed 660 =====
% !PRINTED 660
\weberpage{660}

Der Satz wird bewiesen sein, wenn von irgend einer Potenz
$\varphi^k$ von $\varphi$, deren Exponent nicht durch $q$ theilbar
ist, bewiesen werden kann, dass sie der Hauptclasse angeh\"ort.
Denn sind $\alpha, \beta$ Zahlen in $\Omega_m$ und ist
\[
\varphi^m = \varepsilon' \alpha, \quad \varphi^k = \varepsilon'' \beta,
\]
so k\"onnen wir die ganzen rationalen Zahlen $x, y$ so bestimmen,
dass
\[
mx + ky = 1
\]
wird, und dann wird
\[
\varphi = \varphi^{mx} \varphi^{ky} = \varepsilon''' \alpha^x \beta^y,
\]
also ist auch, wenn $\varepsilon', \varepsilon'', \varepsilon'''$
Einheiten sind, $\varphi$ selbst mit einer Zahl associirt.

Nach dem heutigen Stande der Sache ist der Beweis f\"ur ein
gerades $m$ auf einem ganz anderen Wege zu f\"uhren, als f\"ur
ein ungerades, und wir besch\"aftigen uns also zun\"achst mit dem
leichteren Falle des ungeraden $m$ $^1$).

Die in a) \"uber $\varphi_n$ gemachten Voraussetzungen k\"onnen
wir, wenn $\alpha, \beta$ wieder irgend welche Zahlen des
K\"orpers $\Omega_m$ und $\varepsilon$ Einheitsfunctionale
bedeuten, durch die beiden Formeln ausdr\"ucken:
\weq{1}{\varphi_n^m = \varepsilon \beta,}
\weq{2}{\varphi_n = \varepsilon \alpha \varphi_1^n.}
Wir werden nun die Aufgabe schrittweise l\"osen.

Ist zun\"achst, wie fr\"uher, $\sigma$ der in $q$ enthaltene
Primfactor, der eine in $\Omega_m$ existirende Zahl ist, und
$\sigma^k$ die in $\varphi$ enthaltene Potenz von $\sigma$, so
setzen wir
\weq{3}{\varphi = \sigma^k \psi,}
worin $\psi$ ein zu $q$ theilerfremdes Functional bedeutet, das
denselben Bedingungen (1), (2) wie $\varphi$ gen\"ugt, und wenn
eines von beiden in die Hauptclasse geh\"ort, so gilt das auch
vom anderen. Der Satz a) braucht also nur noch unter der
Voraussetzung bewiesen zu werden, dass $\varphi$ theilerfremd zu
$q$ ist.

Es sei nun $p$ eine von $q$ verschiedene Primzahl, die in
Primfactoren zerlegt den Ausdruck hat:
\weq{4}{p = \prod^{\xi} \mathfrak{p}_\xi,}
%CONT>>

\noindent
$^1$) Vgl. \"ubrigens den neuen Beweis von \gesperrt{Hilbert}, bei
dem der Unterschied zwischen geradem und ungeradem $m$ weit mehr
zur\"ucktritt (S. 648).


% ===== scan leaf 682, printed 661 =====
% !PRINTED 661
\weberpage{661}
%<<CONT
worin $\xi$ die in \S.\,172 näher bestimmte Zahlenreihe durchläuft.
Es möge $\mathfrak{p}_\xi^{k_\xi}$ die in $\varphi$ enthaltene Potenz von
$\mathfrak{p}_\xi$ sein, so dass, wenn wir
\[
\varphi = \chi \prod^{\xi} \mathfrak{p}_\xi^{k_\xi}
\]
setzen, $\chi$ theilerfremd zu $p$ und mit keinen anderen Primzahlen verwandt
ist (s. den Schluss von \S.\,173), als $\varphi$ selbst. Daraus ergiebt sich
\weq{5}{\varphi_n = \chi_n \prod^{\xi} \mathfrak{p}_{n\xi}^{k_\xi},}
und hierin lassen wir $n$ abermals die Reihe der Zahlen $\xi$ durchlaufen.
Setzen wir noch
\[
\psi = \prod^{\xi} \chi_\xi,
\]
so ist auch $\psi$ theilerfremd zu $p$ und mit keinen anderen Primzahlen
 verwandt als $\varphi$, und wir erhalten aus (5) mit Benutzung von (4):
\weq{6}{\prod^{\xi} \varphi_\xi = \psi\, p^{\Sigma k_\xi},}
und daraus ergiebt sich, dass das Functional $\psi$ denselben beiden
Bedingungen (1), (2) genügt, wie $\varphi$.

Es zeigt sich ferner, dass, wenn die $\varphi_\xi$ in der Hauptclasse
liegen, auch $\psi$ in der Hauptclasse enthalten ist.

Wenn aber $p-1$ nicht durch $q$ theilbar ist, so können wir
auch das Umgekehrte schliessen.

Denn nach (2) ist $\varphi_\xi$ äquivalent mit $\varphi^\xi$ und folglich ist:
\[
\begin{gathered}
\prod^{\xi} \varphi_\xi \ \text{äquivalent}\ \varphi^{\Sigma\xi}, \\
\text{äquivalent}\ \psi,
\end{gathered}
\]
und wenn $\psi$ mit einer Zahl associirt ist, so gilt dasselbe
von $\varphi^{\Sigma\xi}$. Nach dem Satze \S.\,172 ist aber, wenn $p-1$ nicht
durch $q$ theilbar ist, $\Sigma \xi$ auch nicht durch $q$ theilbar, und folglich
 ist auch $\varphi$ selbst in der Hauptclasse enthalten.

Diese Betrachtung können wir nun, wenn $\psi$ noch mit
weiteren Primzahlen, die nach dem Modul $q$ nicht mit 1 congruent
 sind, verwandt ist, wiederholen, und kommen so zu dem
Satze:

\gesperrt{Der Satz a) ist allgemein bewiesen, wenn er
unter der Voraussetzung bewiesen werden kann,
dass $\varphi$ nur mit solchen Primzahlen verwandt ist,
die nach dem Modul $q$ mit 1 congruent sind.}


% ===== scan leaf 683, printed 662 =====
% !PRINTED 662
\weberpage{662}
Machen wir diese Voraussetzung über die Functionale $\varphi_n$, so
können wir das \gesperrt{Kummer'sche} Theorem in der Fassung des
\S.\,174, 3. anwenden.

Wir bezeichnen in der Folge mit $\varepsilon$ Einheitsfunctionale, auf
deren genauere Kenntniss nichts ankommt, und setzen
\weq{7}{\Phi_n = \prod^{t} \varphi_{nt'}^{t} = \varepsilon\, \vartheta_n^m,}
wenn $t$ die durch $q$ nicht theilbaren Reste von $m$, die kleiner
als $m$ sind, durchläuft, und $t'$ durch $t\,t' \equiv 1\ (\mathrm{mod}\ m)$ bestimmt
ist. Dann ist $\vartheta_n$ eine Kreistheilungszahl in einem höheren
Körper, deren $m^{\text{te}}$ Potenz in $\Omega_m$ enthalten ist und die der zweiten
Bedingung genügt, dass
\weq{8}{\vartheta_n^{-1} \vartheta_1^n = \alpha}
eine Zahl des Körpers $\Omega_m$ ist.

Wir führen nun ein vielfach gebrauchtes Zeichen ein, indem
wir unter $E(x)$ die grösste ganze Zahl verstehen, die in der
reellen Zahl $x$ enthalten ist, und unter $(x)$ den Rest, der stets
ein echter Bruch ist, und, wenn $x$ selbst eine ganze Zahl sein
sollte, gleich Null zu setzen ist. Dann ist
\weq{9}{x = E(x) + (x).}

Ist $x$ eine ganze Zahl, so ist nach dieser Bezeichnungsweise
\[
m \left(\frac{x}{m}\right)
\]
auch eine ganze Zahl, und zwar der kleinste positive Rest der
Theilung von $x$ durch $m$.

In der Formel (7) ist der Index der Functionale $\varphi$ nur
nach dem Modul $m$ bestimmt. Der Exponent aber ist auf seinen
kleinsten Rest nach dem Modul $m$ reducirt. Ersetzen wir also
in (7) den Index $t'$ durch $n' t'$, so muss $t$ durch den Rest der
Division von $n t$ durch $m$, d.\,h. [nach (9)] durch
\[
m \left(\frac{n t}{m}\right) = n t - m E\left(\frac{n t}{m}\right)
\]
ersetzt werden, und wir können setzen
\[
\Phi_n = \prod^{t} \varphi_{t'}^{\,n t - m E\left(\frac{n t}{m}\right)} = \varepsilon\, \vartheta_n^m,
\]
andererseits ist
\[
\Phi_1 = \prod^{t} \varphi_{t'}^{t} = \varepsilon\, \vartheta_1^m,
\]
%CONT>>


% ===== scan leaf 684, printed 663 =====
% !PRINTED 663
\weberpage{663}
%<<CONT
also nach (8)
\weq{10}{\Phi_1^n \Phi_n^{-1} = \prod^{t} \varphi_{t'}^{\,m E\left(\frac{n t}{m}\right)} = \varepsilon\, \alpha^m.}

Hiernach ist der Quotient
\[
\left(\frac{\prod\limits^{t} \varphi_{t'}^{\,E\left(\frac{n t}{m}\right)}}{\alpha}\right)^m
\]
ein Einheitsfunctional, und da er zugleich die $m^{\text{te}}$ Potenz eines
Functionals in $\Omega_m$ ist, so ist die Basis dieser Potenz auch eine
Einheit, d.\,h. es ist
\[
\prod^{t} \varphi_{t'}^{\,E\left(\frac{n t}{m}\right)} = \varepsilon \alpha,
\]
und es gehört dies Product der Hauptclasse an. Nun ist aber
nach unserer Voraussetzung
\[
\varphi_{t'}\ \text{äquivalent mit}\ \varphi^{t'},
\]
und daher ist
\[
\prod^{t} \varphi_{t'}^{\,E\left(\frac{n t}{m}\right)}\ \text{äquivalent mit}\ \varphi^{\sum^{t} t' E\left(\frac{n t}{m}\right)},
\]
und dies gilt für jedes beliebige durch $q$ nicht theilbare $n$.

Können wir also nachweisen, dass sich $n$ so bestimmen
lässt, dass auch die Summe
\[
S_n = \sum^{t} t' E\left(\frac{n t}{m}\right)
\]
durch $q$ nicht theilbar ist, so haben wir das Theorem a) bewiesen.

Wir wollen zunächst eine Umformung der Summe $S_n$ vornehmen,
durch die die Frage auf den weit einfacheren Fall
zurückgeführt wird, in dem $m$ eine Primzahl ist.

Wir setzen, wenn $m$ eine höhere als die erste Potenz von
$q$ ist,
\[
m = q m', \quad t = t_1 + q t_2, \quad
\begin{aligned}
t_1 &= 1, 2, \dots, q-1 \\
t_2 &= 0, 1, \dots, m'-1.
\end{aligned}
\]

Darin ist $m'$ noch eine Potenz von $q$, der Summationsbuchstabe
 $t_1$ durchläuft die Zahlenreihe $1, 2, \dots, q-1$, und $t_2$
die Zahlenreihe $0, 1, 2, \dots, m'-1$. Wir bestimmen $t_1'$ aus der
Congruenz
\[
t_1 t_1' \equiv 1\ (\mathrm{mod}\ q)
\]
und erhalten
\[
t' \equiv t_1'\ (\mathrm{mod}\ q).
\]


% ===== scan leaf 685, printed 664 =====
% !PRINTED 664
\weberpage{664}
Es ist dann
\weq{11}{S_n = \sum^{t} t' E\left(\frac{t n}{m}\right) \equiv \sum^{t_1} t_1' \sum^{t_2} E\left(\frac{n t_2}{m'} + \frac{n t_1}{m}\right)\ (\mathrm{mod}\ q).}

Wir nehmen jetzt $n < q$ an, also auch $n t_1 < m$, so dass
$n t_1 : m$ ein echter Bruch ist. Dann ist
\[
\begin{aligned}
\text{a)} \quad E\left(\frac{n t_2}{m'} + \frac{n t_1}{m}\right)
  &\quad \text{entweder}\quad = E\left(\frac{n t_2}{m'}\right) \\
\text{b)} \quad
  &\quad \text{oder}\quad = E\left(\frac{n t_2}{m'}\right) + 1,
\end{aligned}
\]
und zwar tritt der Fall b) nur dann ein, wenn zwischen
\[
\frac{n t_2}{m'}\ \text{und}\ \frac{n t_2}{m'} + \frac{n t_1}{m}
\]
eine ganze Zahl liegt. Dies ist aber nur dann der Fall, wenn
der Unterschied zwischen $n t_2 : m'$ und der nächst grösseren
ganzen Zahl kleiner als $n t_1 : m$ ist, also wenn
\weq{12}{E\left(\frac{n t_2}{m'}\right) + 1 - \frac{n t_2}{m'} < \frac{n t_1}{m}.}

Lassen wir in dem Ausdrucke auf der linken Seite von (12)
\weq{13}{E\left(\frac{n t_2}{m'}\right) + 1 - \frac{n t_2}{m'}}
$t_2$ alle seine Werthe durchlaufen, so erhalten wir lauter positive,
die Einheit nicht übersteigende rationale Brüche mit dem
Nenner $m'$, von denen keine zwei einander gleich sind. Denn
wären zwei der Werthe, die aus (13) durch Substitution von
$t_2', t_2''$ für $t_2$ entstehen, einander gleich, so müsste $\dfrac{n(t_2' - t_2'')}{m'}$ eine
ganze Zahl sein, was, da $n$ relativ prim zu $m'$ und $t_2' - t_2''$ kleiner
als $m'$ ist, nicht sein kann. Die Ausdrücke (13) durchlaufen
daher in irgend einer Reihenfolge die Zahlenwerthe
\weq{14}{\frac{1}{m'}, \frac{2}{m'}, \dots, \frac{m'-1}{m'}, 1,}
und wenn $a$ einen beliebigen positiven Bruch bedeutet, so ist
$E(m' a)$ die Anzahl der Zahlen der Reihe (14), die nicht grösser
als $a$ sind. Nach (12) ist daher die Anzahl der Werthe von $t_2$,
für die der Fall b) eintritt, gleich $E\left(\dfrac{n t_1}{q}\right)$, und es ergiebt sich
\[
\sum^{t_2} E\left(\frac{n t_2}{m'} + \frac{n t_1}{m}\right) = \sum^{t_2} E\left(\frac{n t_2}{m}\right) + E\left(\frac{n t_1}{q}\right). % [source sic: the first denominator on the right is printed m, not the mathematically required m']
\]


% ===== scan leaf 686, printed 665 =====
% !PRINTED 665
\weberpage{665}

%% [running head, not body text -- commented out 2026-08-07] Page furniture (running
%% head, folio, signature) is excluded from this edition by standing policy.
%% \S.\,182. \hfill \gesperrt{Beweis der H\"ulfss\"atze.} \hfill 665

Um $S_n$ zu bilden, multipliciren wir diese Gleichung mit $t_1'$ und
summiren. Dadurch ergiebt sich

\[
S_n \equiv \sum^{t_1'} t_1' \sum^{t_2} E\left(\frac{n t_2}{m}\right)
+ \sum^{t_1} t_1' E\left(\frac{n t_1}{q}\right) \pmod q.
% [source sic: the denominator m repeats the source's m-for-m' error from the preceding leaf]
\]

Darin ist nun $\sum t_1' = \sum t_1 \equiv 0 \pmod q$, weil die $t_1$
paarweise die Summe $q$ ergeben, und es folgt:

\weq{15}{S_n \equiv \sum^{t_1} t_1' E\left(\frac{n t_1}{q}\right) \pmod q.}

Die Summe, die hier auf der rechten Seite steht, ist ebenso gebildet,
wie $S_n$ selbst in dem Falle, wo $m$ eine Primzahl ist. F\"ur diesen
Fall ergiebt sich aber

\[
\frac{n t_1}{q} = E\left(\frac{n t_1}{q}\right) + \left(\frac{n t_1}{q}\right)
\]
\[
\frac{(q-n) t_1}{q} = E\left(\frac{(q-n) t_1}{q}\right) + \left(\frac{(q-n) t_1}{q}\right),
\]

und daraus durch Addition

\[
t_1 = E\left(\frac{n t_1}{q}\right) + E\left(\frac{(q-n) t_1}{q}\right)
+ \left(\frac{n t_1}{q}\right) + \left(\frac{(q-n) t_1}{q}\right),
\]

und dabei kann die Summe der beiden positiven echten Br\"uche, die
doch eine ganze Zahl sein muss, nur den Werth 1 haben. Es folgt also

\[
E\left(\frac{n t_1}{q}\right) + E\left(\frac{(q-n) t_1}{q}\right) = t_1 - 1,
\]

und daraus durch Multiplication mit $t_1'$ und Summation \"uber alle
Werthe von $t_1$ von 1 bis $q - 1$

\[
\sum^{t_1} t_1' E\left(\frac{n t_1}{q}\right)
+ \sum^{t_1} t_1' E\left(\frac{(q-n) t_1}{q}\right) \equiv -1 \pmod q.
\]

Folglich k\"onnen sicher nicht beide Summen auf der linken Seite durch
$q$ theilbar sein. Dann zeigt der Ausdruck (15), dass von den beiden
Summen $S_n$ und $S_{q-n}$ gewiss eine durch $q$ untheilbar ist.

Dies aber war zu beweisen.


% ===== scan leaf 687, printed 666 =====
% !PRINTED 666
\weberpage{666}
\pdfbookmark[2]{§ 183. Beweis des zweiten Hülfssatzes für eiu ungerades m}{bm209}

%% [running head, not body text -- commented out 2026-08-07] Page furniture (running
%% head, folio, signature) is excluded from this edition by standing policy.
%% 666 \hfill \gesperrt{Zweiundzwanzigster Abschnitt.} \hfill \S.\,183.

\begin{center}
\S.\,183.

\gesperrt{Beweis des zweiten H\"ulfssatzes f\"ur ein ungerades $m$.}
\end{center}

Das zweite Lemma, was nach \S.\,181 noch zu beweisen ist, ist
folgendes:

2. \gesperrt{Ist $\varepsilon_n$ ein System conjugirter numerischer
Einheiten des K\"orpers $\Omega_m$, das der Bedingung gen\"ugt, dass}

\weq{1}{\varepsilon_n \varepsilon_1^{-n} = \mathfrak{E}^m}

\gesperrt{die $m^{\text{te}}$ Potenz einer Einheit in $\Omega_m$ ist,
so sind die $\varepsilon_n$ selbst $m^{\text{te}}$ Potenzen von
Einheiten in $\Omega_m$, multiplicirt mit einer Einheitswurzel der
Form $r^{nk}$.}

Auch dieser Beweis ist f\"ur den Fall eines ungeraden $m$ ganz
elementar. Wir betrachten daher hier zun\"achst diesen Fall.

Nach \S.\,178, 1. k\"onnen wir jede Einheit des K\"orpers $\Omega_m$
in der Form darstellen:

\weq{2}{\varepsilon = \pm r^k \mathfrak{E}(r), \quad \varepsilon_n = \pm r^{nk} \mathfrak{E}(r^n),}

worin $\mathfrak{E}(r)$ eine reelle Einheit des K\"orpers $\Omega_m$
ist. Es gen\"ugt also $\mathfrak{E}(r)$ der Bedingung

\weq{3}{\mathfrak{E}(r) = \mathfrak{E}(r^{-1}),}

woraus sich nach (2) ergiebt:

\weq{4}{\varepsilon_1 \varepsilon_{-1} = \mathfrak{E}(r)^2;}

aus (1) aber findet sich, wenn man $n = -1$ setzt,

\[
\varepsilon_1 \varepsilon_{-1} = \mathfrak{E}^m,
\]

d.\,h. die linke Seite von (4) ist die $m^{\text{te}}$ Potenz einer
Einheit in $\Omega_m$, und es ist also auch

\weq{5}{\mathfrak{E}(r)^2 = \mathfrak{E}^m}

die $m^{\text{te}}$ Potenz einer solchen Einheit.

Da nun $m$ ungerade ist, so l\"asst sich die ganze rationale Zahl $x$
so bestimmen, dass

\weq{6}{2 x + m = 1}

wird, und daraus folgt

\[
\mathfrak{E}(r) = \mathfrak{E}(r)^{2x} \mathfrak{E}(r)^m.
\]
%CONT>>


% ===== scan leaf 688, printed 667 =====
% !PRINTED 667
%% [running head, not body text -- commented out 2026-08-07] Page furniture (running
%% head, folio, signature) is excluded from this edition by standing policy.
%% \S.\,184. \hfill \gesperrt{Fall eines geraden $m$.} \hfill 667
\weberpage{667}
\pdfbookmark[2]{§ 184. Vorläufiges über den Fall eines geraden m}{bm210}
%<<CONT
Wenn wir also nach (5)

\[
\mathfrak{E}^x \mathfrak{E}(r) = e(r)
\]

setzen, so folgt

\weq{7}{\mathfrak{E}(r) = e(r)^m.}

Durch (7) ist aber der Satz 2. bewiesen und damit zugleich f\"ur ein
\gesperrt{ungerades $m$ das ganze Theorem I.}

Auch dieser Schluss versagt f\"ur ein gerades $m$, weil dann die
Gleichung (6) nicht mehr l\"osbar ist.

\begin{center}
\S.\,184.

\gesperrt{Vorl\"aufiges \"uber den Fall eines geraden $m$.}
\end{center}

F\"ur den Fall, dass $m$ eine Potenz von 2 ist, schlagen wir einen
ganz anderen Weg ein, \"uber den hier zun\"achst einige vorl\"aufige
Bemerkungen Platz finden m\"ogen.

Wir haben schon im \S.\,182 gesehen, dass das erste Lemma bewiesen
ist, wenn wir nachweisen k\"onnen, dass irgend eine Potenz von
$\varphi$, deren Exponent nicht durch $q$ theilbar ist, zur
Hauptclasse geh\"ort. Dabei ist von den beiden Eigenschaften des
Functionals $\varphi$ nur die erste benutzt, dass die $m^{\text{te}}$
Potenz von $\varphi$ in der Hauptclasse enthalten ist.

Nun kennen wir nach den allgemeinen S\"atzen in \S.\,154 in jedem
K\"orper einen Exponenten $h$, n\"amlich die Anzahl der
Idealclassen, f\"ur den $\varphi^h$ zur Hauptclasse geh\"ort, wenn
$\varphi$ ein beliebiges Functional in diesem K\"orper ist.

Wenn wir also beweisen k\"onnen, dass die \gesperrt{Classenzahl $h$ im K\"orper
$\Omega_m$, wenn $m$ eine Potenz von 2 ist, eine ungerade Zahl ist,}
so ist damit ohne alles Weitere das Lemma 1. bewiesen.

Dies soll das Ziel unserer Betrachtungen in den n\"achsten
Abschnitten sein.

In Bezug auf das zweite Lemma liegt die Sache \"ahnlich. Hier kann
man aus den Voraussetzungen in \S.\,183, 2. ganz wie oben die Formel
\S.\,183, (5) herleiten, aus der aber nur zu schliessen ist, dass
$\mathfrak{E}(r)$ die $\frac{1}{2} m^{\text{te}}$ Potenz einer
reellen Einheit $e(r)$ ist.

Nehmen wir also das erste Lemma als bewiesen an, so ergiebt die
Formel \S.\,181, (14):

\weq{1}{(r^n, x_0) = \varrho_n \sqrt{e(r^n)}\ \alpha_n\ \Pi(r^{m_2 n}, \eta),}
%CONT>>


% ===== scan leaf 689, printed 668 =====
% !PRINTED 668
%% [running head, not body text -- commented out 2026-08-07] Page furniture (running
%% head, folio, signature) is excluded from this edition by standing policy.
%% 668 \hfill \gesperrt{Zweiundzwanzigster Abschnitt.} \hfill \S.\,184.
\weberpage{668}
%<<CONT
worin $\varrho_n$ eine Einheitswurzel vom Grade $m^2$, und $\alpha_n$
eine Zahl in $\Omega_m$ ist, und es w\"are noch zu beweisen, dass
$e(r^n)$ das \gesperrt{Quadrat einer Einheit} in $\Omega_m$ ist.

Da $e(r)$ eine reelle Einheit ist, so ist

\weq{2}{e(r^n) = e(r^{-n}),}

und wir wollen noch festsetzen, dass das Vorzeichen der Wurzel so
bestimmt sei (was wir bei passender Annahme \"uber $\varrho_n$ immer
annehmen k\"onnen), dass

\weq{3}{\sqrt{e(r^n)} = \sqrt{e(r^{-n})}.}

Das Product $(r^n, x_0)\,(r^{-n}, x_0)$ ist nach der Voraussetzung
eine Zahl des K\"orpers $\Omega_m$, die sich durch die Substitution
$(r, r^{-1})$ nicht \"andert, d.\,h. eine reelle Zahl. Ferner ist
nach \S.\,174, (5)

\[
(r^{m_2 n}, \eta)\,(r^{-m_2 n}, \eta) = \pm p,
\]

also gleichfalls reell. Ebenso ist $\alpha_n \alpha_{-n}$ reell, und
daraus ergiebt sich nach (1), dass

\[
\varrho_n \varrho_{-n} = \frac{(r^n, x_0)\,(r^{-n}, x_0)}
{e(r^n)\ \alpha_n \alpha_{-n}\ \Pi(r^{m_2 n}, \eta)\,(r^{-m_2 n}, \eta)}
\]

eine reelle Zahl ist, die, weil sie eine Einheitswurzel ist, nur $=
\pm 1$ sein kann.

Nun k\"onnen wir in der hieraus f\"ur $n = 1$ sich ergebenden
Gleichung

\[
(r, x_0)\,(r^{-1}, x_0) = \pm\ e(r)\ \alpha_1 \alpha_{-1}\ \Pi(\pm p)
\]

jede Substitution $(r, r^n)$ ausf\"uhren, woraus hervorgeht, dass das
Zeichen von $\varrho_n \varrho_{-n}$ von $n$ unabh\"angig ist, dass
also

\weq{4}{\varrho_n \varrho_{-n} = \varrho_1 \varrho_{-1} = \pm 1}

ist. Nun leiten wir aus (1) weiter her:

\weq{5}{\varrho_n \varrho_1^{-1} \sqrt{e(r^n)}\ \sqrt{e(r)}^{-n}
= \frac{(r^n, x_0)\,(r, x_0)^{-n}}
{\alpha_n \alpha_1^{-1}\ \Pi(r^{m_2 n}, \eta)\,(r^{m_2}, \eta)^{-n}},}
% [source sic: (5) prints both \varrho_1^{-1} and \alpha_1^{-1}; formation from (1) and the displayed -n powers requires \varrho_1^{-n} and \alpha_1^{-n}]

und die rechte Seite ist eine Zahl des K\"orpers $\Omega_m$. Die
linke Seite zeigt aber, dass es eine Einheit ist, und wir k\"onnen
demnach, wenn $\mathfrak{E}(r)$ eine reelle Einheit bedeutet,

\weq{6}{\varrho_n \varrho_1^{-1} \sqrt{e(r^n)}\ \sqrt{e(r)}^{-n} = r^k \mathfrak{E}(r)}
% [source sic: (6) repeats the inherited \varrho_1^{-1}; the derivation requires exponent -n]

setzen. Substituiren wir diesen Werth in die Formel (5) und machen
die Substitution $(r, r^{-1})$, so ergiebt sich

\weq{7}{\varrho_{-n} \varrho_{-1}^{-1} \sqrt{e(r^n)}\ \sqrt{e(r)}^{-n} = r^{-k} \mathfrak{E}(r),}
% [source sic: (7) prints \varrho_{-1}^{-1}; the substituted form of the preceding relation requires exponent -n]
%CONT>>


% ===== scan leaf 690, printed 669 =====
% !PRINTED 669
\weberpage{669}
%<<CONT
und durch Multiplication von (6) und (7) mit Rücksicht auf (4)

\weq{8}{e(r^n)\, e(r)^{-n} = \mathfrak{E}(r)^2.}

In unseren Betrachtungen über die Classenzahl wird sich noch der Satz ergeben:

\gesperrt{Eine Einheit $e(r)$ des reellen Körpers $H_m$, die mit allen ihren Conjugirten positiv ist, ist das Quadrat einer Einheit im Körper $H_m$.}

Die Formel (8) zeigt aber, dass der Einheit $\pm e(r)$ diese Eigenschaft zukommt, und dass daher $e(r)$ (da auch $-1 = i^2$ das Quadrat einer Einheit ist) das Quadrat einer Einheit des Körpers $\Omega_m$ ist.

Damit ist die Frage auf den Beweis der beiden erwähnten Sätze zurückgeführt, der sich als Schlussstein einer langen, aber an interessanten Beziehungen äusserst reichen Kette von Betrachtungen ergiebt, denen die nächsten Abschnitte gewidmet sein sollen.

\begin{center}
\rule{6em}{0.4pt}
\end{center}


% ===== scan leaf 691, printed 670 =====
\weberpage{670}% page number not printed on this page
\pdfbookmark[1]{Dreiundzwanzigster Abschnitt. Classenzahl.}{bm211}
\pdfbookmark[2]{§ 185. Der Dirichlet'sche Satz über die Einheiten}{bm212}
% !PRINTED none
\weberpage{none}
% [source note: this Abschnitt opener has no printed folio; sequence position corresponds to inferred page 670.]

\begin{center}
\gesperrt{Dreiundzwanzigster Abschnitt.}

\gesperrt{Classenzahl.}

\rule{4em}{0.4pt}

\S.\,185.

\gesperrt{Der Dirichlet'sche Satz über die Einheiten.}
\end{center}

Die Untersuchungen über die Anzahl der Idealclassen in einem algebraischen Körper, die uns nun die nothwendigen Ergänzungen der Beweise des vorigen Abschnittes geben sollen, bedienen sich der Methoden, die \gesperrt{Dirichlet} in die Zahlentheorie eingeführt hat\begingroup\renewcommand{\thefootnote}{\arabic{footnote})}\footnote{\gesperrt{Dirichlet's}, Recherches sur diverses applications de l'analyse infinitésimale à la théorie des nombres. \gesperrt{Crelle's} Journal, Bd. 19, 21. \gesperrt{Dirichlet's} Werke, Bd. I (1839, 1840). (Hierher gehören übrigens auch die nachgelassenen Abhandlungen von \gesperrt{Gauss}, \glqq De nexu inter multitudinem classium etc.\grqq{} \gesperrt{Gauss'} Werke, Bd. II.) Die Anwendung dieser Principien auf die Kreistheilungszahlen hat zuerst \gesperrt{Kummer} gemacht [\gesperrt{Crelle's} Journal, Bd. 35, 40 (1847, 1850); \gesperrt{Liouville's} Journal, Bd. 16 (1851)]. Die allgemeine Theorie der Einheiten ist gleichfalls von \gesperrt{Dirichlet} begründet [\gesperrt{Dirichlet's} Werke, Bd. I, S. 622, 633, 639 (1840, 1842, 1846)]. Die allgemeine Theorie der Einheiten und die Bestimmung der Classenzahlen hat \gesperrt{Dedekind} gegeben: \gesperrt{Dirichlet-Dedekind}, Vorlesungen über Zahlentheorie, 4. Auflage, \S.\,183 f. \gesperrt{Minkowski}, Geometrie der Zahlen.}\endgroup.

Sie sind nicht mehr rein algebraisch, insofern dabei auch transcendente Functionen und Zahlen, wie der natürliche Logarithmus und die Zahl $\pi$ vorkommen. Auch wird von Grenzübergängen, wie in der Integralrechnung, Gebrauch gemacht, aus denen wir schon in \S.\,156 ein schönes Resultat gezogen haben.

Die Untersuchungen, von denen wir zunächst zu handeln haben, sind ganz allgemein, d.\,h. sie gelten für jeden algebraischen Zahlkörper, und es bietet keinen Vortheil der Einfachheit, sie auf specielle Körper, wie etwa die Kreistheilungskörper, zu beschränken.


% ===== scan leaf 692, printed 671 =====
% !PRINTED 671
\weberpage{671}
Es sei also $\Omega$ ein Körper $n^{\text{ten}}$ Grades, und

\weq{1}{\Omega_1,\ \Omega_2,\ \ldots,\ \Omega_n}

die conjugirten Körper. Die irreducible Gleichung $n^{\text{ten}}$ Grades, deren Wurzeln uns diese conjugirten Körper bestimmen, wird im Allgemeinen sowohl reelle als imaginäre Wurzeln haben, von denen aber die letzteren immer paarweise vorkommen, so dass zu jedem imaginären Körper $\Omega^{(i)}$ ein anderer gehört, dessen Zahlen mit denen von $\Omega^{(i)}$ conjugirt imaginär sind, der also aus $\Omega^{(i)}$ durch die Substitution $(i, -i)$ hervorgeht.

Solche imaginäre Paare fassen wir zu einer Einheit zusammen und bezeichnen \gesperrt{die Anzahl der reellen Körper und der Paare imaginärer Körper der Reihe (1)} mit $\nu$. Sind alle conjugirten Körper reell, so ist $n = \nu$; sind sie alle imaginär, so ist $n = 2\nu$.

Ausser im Falle des rationalen Körpers ist $\nu$ nur noch in dem Falle eines imaginären quadratischen Körpers $= 1$. Sonst ist $\nu$ immer grösser als 1.

Die Anzahl der imaginären Paare ist $n - \nu$, und die Anzahl der reellen Körper $2\nu - n$.

In der Determinante, durch die in \S.\,145 die Quadratwurzel aus der Grundzahl, $\sqrt{\Delta}$, definirt ist, vertauschen sich durch die Substitution $(i, -i)$ je zwei conjugirt imaginäre Reihen, und $\sqrt{\Delta}$ wechselt also sein Zeichen dann, wenn diese Zahl ungerade ist, sonst nicht. Im ersten Falle ist $\sqrt{\Delta}$ imaginär, im zweiten reell. Daraus ergiebt sich also, dass die \gesperrt{Grundzahl $\Delta$ des Körpers $\Omega$ positiv oder negativ ist, je nachdem $n - \nu$ gerade oder ungerade ist.}

Behalten wir von zwei conjugirt imaginären Körpern nur den einen bei, so ergiebt sich die Reihe der conjugirten Körper

\[ \Omega_1,\ \Omega_2,\ \ldots,\ \Omega_\nu, \]

denen wir ein Zeichensystem

\[ \delta_1,\ \delta_2,\ \ldots,\ \delta_\nu \]

mit der Bestimmung zuordnen, dass $\delta_s = 1$ sein soll, wenn $\Omega_s$ reell, und $= 2$, wenn $\Omega_s$ imaginär ist, also $\Sigma \delta_s = n$.

Es möge $\eta$ eine von Null verschiedene Zahl des Körpers $\Omega$ bedeuten und

\[ \eta_1,\ \eta_2,\ \ldots,\ \eta_n \]

die conjugirten Zahlen. Diesem Zahlensysteme ordnen wir ein anderes Zahlensystem zu

\[ \lambda_1,\ \lambda_2,\ \ldots,\ \lambda_n, \]
%CONT>>


% ===== scan leaf 693, printed 672 =====
% !PRINTED 672
\weberpage{672}
%<<CONT
das wir durch die Gleichung

\weq{2}{\lambda_s = \delta_s \log|\eta_s|}

definiren, worin $|\eta_s|$ den absoluten Werth von $\eta_s$ und $\log|\eta_s|$ den reellen natürlichen Logarithmus der positiven Zahl $|\eta_s|$ bedeutet. Ist $\Omega_s$ reell und das Zeichen $\pm$ so gewählt, dass $\pm \eta_s$ positiv ist, so ist

\[ \lambda_s = \log(\pm \eta_s). \]

Ist aber $\eta_s, \eta_s'$ ein imaginäres Paar, so ist

\[ \lambda_s = \log \eta_s \eta_s', \]

und zu $\eta_s$ und $\eta_s'$ gehört dasselbe $\lambda_s$, so dass die Anzahl der verschiedenen $\lambda_s$ gleich $\nu$ ist. Aus dieser Bestimmung ergiebt sich noch

\weq{3}{\lambda_1 + \lambda_2 + \cdots + \lambda_\nu = \log N_a(\eta),}

wenn, wie früher, $N_a$ die absolute Norm bedeutet.

Die Zahlen $\lambda_1, \lambda_2, \ldots, \lambda_\nu$ wollen wir die \gesperrt{conjugirten Logarithmen} der Zahl $\eta$ nennen.

Wenn $\eta$ eine \gesperrt{ganze Zahl} des Körpers $\Omega$ ist, so ist die absolute Norm eine \gesperrt{natürliche Zahl}, die nur dann $= 1$ ist, wenn $\eta$ eine Einheit ist, und daraus folgt nach (3):

1. \gesperrt{Die Summe der conjugirten Logarithmen einer ganzen Zahl $\eta$ ist positiv, und nur dann gleich Null, wenn $\eta$ eine Einheit ist.}

Bedeutet $\omega_1, \omega_2, \ldots, \omega_n$ eine Basis der ganzen Zahlen in $\Omega$ (eine Minimalbasis von $\Omega$, \S.\,145), und $\omega_{1,s}, \omega_{2,s}, \ldots, \omega_{n,s}$ für $s = 1, 2, \ldots, n$ die conjugirten Zahlen, so lassen sich alle ganzen Zahlen des Körpers $\Omega_s$ in der Form darstellen:

\weq{4}{\eta_s = x_1 \omega_{1,s} + x_2 \omega_{2,s} + \cdots + x_n \omega_{n,s},}

mit ganzen rationalen Co\"efficienten $x_1, x_2, \ldots, x_n$. Die Determinante dieses Systemes

\[ \Sigma \pm \omega_{1,1} \omega_{2,2} \ldots \omega_{n,n} = \sqrt{\Delta} \]

ist die Quadratwurzel aus der Discriminante $\Delta$ des Körpers, und demnach immer von Null verschieden. Demnach lässt sich das System (4) nach den Unbekannten $x_1, x_2, \ldots, x_n$ auflösen, und wenn wir nun festsetzen, dass die absoluten Werthe der Zahlen $\eta_s$ nicht über eine gewisse Grenze hinausgehen sollen, so ergeben sich aus diesen Auflösungen Grenzen für die ganzen rationalen Zahlen $x_i$.


% ===== scan leaf 694, printed 673 =====
% !PRINTED 673
\weberpage{673}
Diese einfache Bemerkung liefert uns den folgenden wichtigen Satz:

2. \gesperrt{Es giebt in einem algebraischen K\"orper $\Omega$ nur eine endliche Anzahl ganzer Zahlen $\eta$ von der Beschaffenheit, dass die absoluten Werthe der conjugirten Zahlen $\eta_s$ unter einer gegebenen Grenze liegen.}

Wir machen jetzt von dem allgemeinen Satze \S.\,155, (25) Gebrauch, nach dem, wenn $f(x_1,x_2,\dots,x_n)$ irgend eine \gesperrt{positive quadratische Form} von $n$ Variablen mit der Determinante $D$ ist, die Variablen $x_i$ sich als ganze rationale Zahlen, nicht alle verschwindend, so bestimmen lassen, dass
\[
f(x_1,x_2,\dots,x_n) < n\sqrt[n]{0{,}404\,D},
\]
d.\,h. kleiner als eine von der Determinante $D$ allein abh\"angige Zahl wird.

Diesen Satz wenden wir, wie im \S.\,156, auf die Form an:
\[
f(x_1,x_2,\dots,x_n) = \frac{|\eta_1|^2}{c_1^2} + \frac{|\eta_2|^2}{c_2^2} + \cdots + \frac{|\eta_n|^2}{c_n^2},
\]
worin die $\eta_s$ die Linearformen (4) sind, und im Falle eines reellen $\eta_s$
\[
|\eta_s|^2 = (\eta_s)^2,
\]
im Falle eines imagin\"aren Paares $\eta_s, \eta_s'$
\[
|\eta_s|^2 = |\eta_s'|^2 = \eta_s\eta_s'.
\]
Im letzteren Falle ist $\eta_s\eta_s'$ die Summe zweier reeller Quadrate, und demnach ist, wenn die $c_s$ reell angenommen werden, $f$ eine positive Form. Wir wollen \"uber die bis jetzt willk\"urlichen reellen Gr\"ossen $c_s$ noch festsetzen, dass, wenn $\eta_s, \eta_s'$ ein conjugirt imagin\"ares Paar bilden,
\weq{5}{c_s = c_s'}
sein soll; ausserdem wollen wir die $c_s$ noch der Bedingung unterwerfen:
\weq{6}{c_1 c_2 \dots c_n = 1.}
Die Determinante der Function $f$ bilden wir, wie im \S.\,156, dadurch, dass wir sie zun\"achst als Form der Variablen $\eta_1, \eta_2, \dots, \eta_n$ auffassen, und die Determinante dieser Form ist wegen (6) gleich $\pm 1$. Die Determinante von $f$ in Bezug auf die Variablen $x_i$ ist also (Bd.\ I, \S.\,56), vom Vorzeichen abgesehen, gleich dem
%CONT>>


% ===== scan leaf 695, printed 674 =====
% !PRINTED 674
\weberpage{674}
%<<CONT
Quadrate der Determinante der Linearformen $\eta_s$, d.\,h. der Grundzahl des K\"orpers $\Omega$.

Daraus folgt, dass man eine von den $c_s$ unabh\"angige, nur durch den K\"orper $\Omega$ bestimmte Zahl $A$ finden kann, von der Art, dass, was auch die $c_s$ sein m\"ogen, $f(x_1,x_2,\dots,x_n)$ f\"ur ganzzahlige nicht verschwindende $x$ unter $A$ heruntersinkt. Da $f$ die Summe von positiven Summanden ist, so gilt dasselbe von jedem dieser Summanden, und damit ist der folgende Satz bewiesen:

3. \gesperrt{Ist $c_s$ ein beliebiges, den Bedingungen (5), (6) gen\"ugendes System reeller positiver Zahlen, so giebt es eine durch den K\"orper $\Omega$ allein bestimmte reelle Zahl $A$ von der Art, dass man immer eine ganze Zahl $\eta$ des K\"orpers $\Omega$ bestimmen kann, die mit ihren conjugirten Zahlen $\eta_s$ den Bedingungen gen\"ugt}
\weq{7}{|\eta_s| < c_s A.}
Beil\"aufig folgt hieraus, dass man in jedem algebraischen K\"orper $\Omega$, ausgenommen dem rationalen und dem imagin\"aren quadratischen, von Null verschiedene ganze Zahlen finden kann, deren absoluter Werth unter jede gegebene Grenze heruntersinkt.

Wir wollen den Ausdruck des Satzes 3. durch Einf\"uhrung der conjugirten Logarithmen von $\eta$ etwas umformen. Wir setzen $\log A = k$, so dass $k$ eine durch $\Omega$ bestimmte reelle Zahl ist, ferner
\weq{8}{c_s = e^{\frac{\gamma_s u}{\delta_s}},}
worin $u$ eine beliebige reelle Gr\"osse sein kann, und die Zahlen $\gamma_s$ wegen (5) und (6) der Bedingung gen\"ugen:
\weq{9}{\gamma_1 + \gamma_2 + \cdots + \gamma_\nu = \sum_{1,\nu}^{i} \gamma_i = 0.}
Da \"uberdies
\[
|\eta_s| = e^{\frac{\lambda_s}{\delta_s}}
\]
ist, so folgt aus (7) und (8):
\weq{10}{\lambda_s - \gamma'_s u < k\delta_s.}
% [sic] The printed prime has no referent; the unprimed \gamma_s is required by (8), (9), and the following sum.
Die Summe der $\nu$ Ausdr\"ucke auf der linken Seite dieser Ungleichung ist wegen (9) gleich $\Sigma\lambda_s$, also nach dem Satze 1. nicht negativ, und wir finden
\[
0 \leqq \sum^{s} (\lambda_s - \gamma_s u) < k n.
\]


% ===== scan leaf 696, printed 675 =====
% !PRINTED 675
\weberpage{675}
Daraus ergiebt sich wegen (10), dass ein Glied dieser Summe, $\lambda_s - \gamma_s u$, nicht kleiner sein kann, als $-(n - \delta_s) k$, weil sonst die Gesammtsumme negativ w\"are, und wenn wir der Einfachheit halber die untere Grenze noch etwas kleiner nehmen
\weq{11}{-n k \delta_s < \lambda_s - \gamma_s u < k \delta_s.}
Damit ist bewiesen, dass es f\"ur jedes gegebene System der Zahlen $u, \gamma_s$, wenn nur die Bedingung (9) erf\"ullt ist, eine ganze Zahl $\eta$ des K\"orpers $\Omega$ giebt, die mit ihren conjugirten Zahlen den Grenzbedingungen (11) gen\"ugt.

Es sei nun ferner $g_s$ ein System reeller Gr\"ossen, von dem wir nur verlangen, dass
\[
\gamma_1 g_1 + \gamma_2 g_2 + \cdots + \gamma_\nu g_\nu = \alpha
\]
von Null verschieden sei. Damit ist [nach (9)] ausgeschlossen, dass alle $g_s$ einander gleich sind; wenn sie aber das nicht sind, so werden sich zu jedem gegebenen Systeme der $g_s$ unendlich viele Systeme $\gamma_s$ bestimmen lassen, die der Forderung (9) gen\"ugen.

Wird nun $k \Sigma \delta_s g_s = g$ gesetzt, so ergiebt sich aus (11) durch Multiplication mit $g_s$ und Summation
\weq{12}{-n g + \alpha u < \Sigma g_s \lambda_s < g + \alpha u.}
% [source note] Multiplying the inequalities (11) in this direction requires nonnegative g_s, not arbitrary real g_s as stated. One may first add a sufficiently large common constant to all g_s: this leaves alpha unchanged since sum gamma_s=0 and leaves the final linear form on unit logarithms unchanged since sum l_s=0. The printed proof omits this reduction.
Nach dieser Grenzbedingung bestimmen wir nun, bei festgehaltenen $g_s$ und $\alpha$, eine Reihe von ganzen Zahlen
\weq{13}{\eta, \eta', \eta'', \eta''', \dots,}
indem wir f\"ur $u$ eine Reihe von Annahmen machen:
\[
u, u', u'', u''', \dots,
\]
die folgendermaassen n\"aher bestimmt sind:

Wir w\"ahlen $u$ zun\"achst so, dass
\[
-n g + \alpha u = a, \quad g + \alpha u = a'
\]
positiv werden; dann setzen wir
\[
\delta = \frac{a' - a}{\alpha} = \frac{(n+1)g}{\alpha},
\]
und setzen
\[
\begin{gathered}
u' = u + \delta, \quad u'' = u' + \delta, \quad u''' = u'' + \delta, \dots,\\
a + \alpha\delta = a', \quad a' + \alpha\delta = a'', \quad a'' + \alpha\delta = a''', \dots,
\end{gathered}
\]
%CONT>>


% ===== scan leaf 697, printed 676 =====
% !PRINTED 676
\weberpage{676}
%<<CONT
und erhalten so aus (12), wenn $\lambda_s', \lambda_s'', \dots$ die conjugirten Logarithmen von $\eta', \eta'', \dots$ sind:
\weq{14}{
\begin{aligned}
a &< \Sigma g_s \lambda_s < a'\\
a' &< \Sigma g_s \lambda_s' < a''\\
a'' &< \Sigma g_s \lambda_s'' < a'''\\
&\cdots\cdots\cdots\cdots,
\end{aligned}
}
und die Reihe dieser Ungleichungen l\"asst sich unbegrenzt fortsetzen.

Alle diese Zahlen $\eta, \eta' \eta'', \dots$ haben \"uberdies nach (6) und (7) die Eigenschaft, dass ihre absoluten Normen $N, N', N'', \dots$ unter einer bestimmten endlichen Grenze $e^{nk}$ bleiben.

% [sic] The source omits the comma between \eta' and \eta'' in the preceding list.

Die Anzahl der m\"oglichen Werthe von $N, N', N'', \dots$ ist also endlich, n\"amlich gewiss nicht gr\"osser, als die gr\"osste in $e^{nk}$ enthaltene ganze Zahl. Ebenso ist die Anzahl aller m\"oglichen Reste der Zahlen $x_1, x_2, \dots, x_n$ [in (4)] nach allen Moduln $N, N', N'', \dots$ nur eine endliche, und daraus folgt, dass, wenn wir die Reihe der Zahlen (13) nur weit genug fortsetzen, in der Reihe \gesperrt{zwei verschiedene Zahlen $\eta^{(h)}, \eta^{(k)}$ auftreten m\"ussen, in denen $N^{(h)} = N^{(k)}$ und die Reste von $x_1, x_2, \dots, x_n$ nach dem Modul $N^{(h)}$ genau dieselben sind.}

Aus (14) aber folgt, wenn $h > k$ vorausgesetzt wird:
\weq{15}{\Sigma g_s (\lambda_s^{(h)} - \lambda_s^{(k)}) > 0.}
Ist $N$ die absolute Norm dieser beiden Zahlen $\eta^{(h)}, \eta^{(k)}$, so ist wegen der Uebereinstimmung der Reste der $x$ die Differenz $\eta^{(h)} - \eta^{(k)}$ durch $N$ theilbar. Andererseits ist nach \S.\,138\,(13) die Zahl $N$ durch $\eta^{(h)}$ theilbar, und folglich ist auch $\eta^{(h)} - \eta^{(k)}$ durch $\eta^{(h)}$ theilbar. Daraus ergiebt sich, dass auch $\eta^{(k)}$ durch $\eta^{(h)}$ und ebenso $\eta^{(h)}$ durch $\eta^{(k)}$ theilbar ist. Beide Zahlen sind also associirt, und wenn wir
\[
\frac{\eta^{(h)}}{\eta^{(k)}} = \varepsilon
\]
setzen, so ist $\varepsilon$ eine \gesperrt{Einheit} des K\"orpers $\Omega$.

Setzen wir ausserdem, indem wir mit $|\varepsilon_s|$ den absoluten Werth von $\varepsilon_s$ bezeichnen,
\weq{16}{\delta_s \log|\varepsilon_s| = l_s,}
so ergiebt sich aus (15):
\weq{17}{g_1 l_1 + g_2 l_2 + \cdots + g_\nu l_\nu > 0.}


% ===== scan leaf 698, printed 677 =====
% !PRINTED 677
\weberpage{677}
\pdfbookmark[2]{§ 186. Systeme unabhängiger Einheiten und Exponentensysteme der Einheiten}{bm213}
Dies beweist nun der Satz:
% [source note] The source prints der Satz; den Satz is apparently intended, since the preceding argument establishes the theorem which follows. The printed wording reverses that logical direction if read literally.

4. \gesperrt{Ist $g_1, g_2, \ldots, g_\nu$ ein beliebiges System reeller Zahlen, die nicht alle einander gleich sind, so l\"asst sich in jedem algebraischen K\"orper $\Omega$ eine Einheit $\varepsilon$ so bestimmen, dass die Summe}
\[
g_1 l_1 + g_2 l_2 + \cdots + g_\nu l_\nu
\]
\gesperrt{von Null verschieden ist, wenn $l_s$ das System der conjugirten Logarithmen von $\varepsilon$ ist.}

Dieser wichtige Satz, der das Fundament f\"ur das Folgende ist, r\"uhrt, wenn auch in etwas ver\"anderter Fassung, von \gesperrt{Dirichlet} her. Diese Formulirung ist von \gesperrt{Minkowski} gegeben.

\begin{center}
\S.\,186.\\
\gesperrt{Systeme unabh\"angiger Einheiten und Exponentensysteme der Einheiten.}
\end{center}

Wenn wir in dem zuletzt bewiesenen Satze $g_1 = 1, g_2 = 0, \ldots, g_\nu = 0$ setzen, was, wenn wir von jetzt an $\nu > 1$ voraussetzen, gestattet ist, so folgt, dass es in $\Omega$ eine Einheit $\varepsilon$ giebt, f\"ur die einer der conjugirten Logarithmen, $l_1$, von Null verschieden ist. Da jede Potenz von $\varepsilon$ dieselbe Eigenschaft hat, so giebt es auch unendlich viele solche Einheiten.

Dieser Satz gestattet nun eine sehr wichtige Verallgemeinerung:

5. \gesperrt{Ist $s \leqq \nu - 1$, so l\"asst sich im K\"orper $\Omega$ ein System von Einheiten mit den zugeh\"origen conjugirten Logarithmen}
\weq{1}{
\begin{array}{l}
\varepsilon_1 : l_{1,1}, l_{1,2}, \ldots, l_{1,\nu} \\
\varepsilon_2 : l_{2,1}, l_{2,2}, \ldots, l_{2,\nu} \\
\cdots\cdots\cdots\cdots\cdots\cdots \\
\varepsilon_s : l_{s,1}, l_{s,2}, \ldots, l_{s,\nu}
\end{array}
}
\gesperrt{derart bestimmen, dass eine beliebige der $s$-reihigen Determinanten der Matrix der $l_{h,k}$, etwa}
\[
L_s = \Sigma \pm l_{1,1} l_{2,2} \ldots l_{s,s}
\]
\gesperrt{von Null verschieden ist.}

F\"ur $s = 1$ ist der ausgesprochene Satz nach der am Eingange gemachten Bemerkung richtig. Wir nehmen also an, er
%CONT>>


% ===== scan leaf 699, printed 678 =====
% !PRINTED 678
\weberpage{678}
%<<CONT
sei bewiesen f\"ur $s - 1$ und leiten ihn durch vollst\"andige Induction allgemein her. Dazu ordnen wir die Determinante $L_s$ nach den Elementen der letzten Zeile, und schreiben sie so:
\weq{2}{
L_s = g_1 l_{s,1} + g_2 l_{s,2} + \cdots + g_s l_{s,s},
}
worin $g_s = L_{s-1}$ ist, und daher nach der gemachten Voraussetzung von Null verschieden angenommen werden kann. Substituiren wir diese Werthe von $g$ in die Formel des Satzes 4., \S.\,185, indem wir $g_{s+1} = 0, \ldots, g_\nu = 0$ annehmen, w\"ahrend $g_s$ nicht $= 0$ ist, so sind, so lange $s < \nu$ ist, gewiss nicht alle $g_1, g_2, \ldots, g_\nu$ einander gleich, und es ergiebt sich, dass sich $\varepsilon_s$ so annehmen l\"asst, dass $L_s$ von Null verschieden ist, wie bewiesen werden sollte.

Auf $s = \nu$ ist diese Schlussweise nicht auszudehnen, und die Determinante $L_\nu$ muss auch in der That immer Null sein, weil f\"ur jede Einheit $\sum\limits_{1,\nu}^{s} l_s$ verschwindet.

Wir wollen jetzt f\"ur den Fall, dass $s = \nu - 1$ ist, die Determinante $L_s$ so bezeichnen:
\weq{3}{
L_{\nu-1} = L(\varepsilon_1, \varepsilon_2, \ldots, \varepsilon_{\nu-1}),
}
und ein System von Einheiten $\varepsilon_1, \varepsilon_2, \ldots, \varepsilon_{\nu-1}$, f\"ur welches diese Determinante von Null verschieden ist, ein \gesperrt{System unabh\"angiger Einheiten} nennen.

Der absolute Werth $L$ der Determinante $L(\varepsilon_1, \varepsilon_2, \ldots, \varepsilon_{\nu-1})$ heisst der \gesperrt{Regulator} des Systems $\varepsilon_1, \varepsilon_2, \ldots, \varepsilon_{\nu-1}$\footnote{\gesperrt{Dirichlet-Dedekind}, Vorlesungen \"uber Zahlentheorie 4. Auflage, \S.\,183.}.

Wenn wir in der Determinante
\weq{4}{
\begin{vmatrix}
l_{1,1}, & l_{1,2}, & \ldots, & l_{1,\nu} \\
\cdots & \cdots & \cdots & \cdots \\
l_{\nu-1,1}, & l_{\nu-1,2}, & \ldots, & l_{\nu-1,\nu} \\
x_1, & x_2, & \ldots, & x_\nu
\end{vmatrix},
}
in der die $x_1, x_2, \ldots, x_\nu$ willk\"urliche Gr\"ossen sind, alle Colonnen zu der letzten addiren, so erhalten wir mit R\"ucksicht auf die Relationen $\sum\limits^{i} l_{s,i} = 0$
ihren Werth gleich
\[
\pm (x_1 + x_2 + \cdots + x_\nu) L.
\]
Wenn wir also alle $x$ mit Ausnahme von einem beliebigen
%CONT>>


% ===== scan leaf 700, printed 679 =====
% !PRINTED 679
\weberpage{679}
%<<CONT
$= 0$, das eine $= 1$ annehmen, so ergiebt sich aus (3) irgend eine der $(\nu-1)$-reihigen Determinanten der Matrix
% [sic] The source cites (3), but the determinant specialized by the x_i choices is (4) on printed p. 678.
\[
\begin{matrix}
l_{1,1}, & l_{1,2}, & \ldots, & l_{1,\nu} \\
\cdots & \cdots & \cdots & \cdots \\
l_{\nu-1,1}, & l_{\nu-1,2}, & \ldots, & l_{\nu-1,\nu}
\end{matrix}
\]
und alle diese Determinanten haben also (vom Vorzeichen abgesehen) denselben Werth. Es ist daher gleichg\"ultig, welcher unter den $\nu$ conjugirten Werthen bei der Bildung des Regulators $L$ weggelassen wird.

Ist $\varepsilon_1, \varepsilon_2, \ldots, \varepsilon_{\nu-1}$ ein unabh\"angiges System von Einheiten, und sind $\lambda_1, \lambda_2, \ldots, \lambda_\nu$ die conjugirten Logarithmen irgend einer Einheit $\varepsilon$ in $\Omega$, so kann man die Zahlen $\xi_1, \xi_2, \ldots, \xi_{\nu-1}$ immer und nur auf eine Weise so bestimmen, dass
\weq{5}{
\begin{array}{l}
\xi_1 l_{1,1} + \xi_2 l_{2,1} + \cdots + \xi_{\nu-1} l_{\nu-1,1} = \lambda_1, \\
\xi_1 l_{1,2} + \xi_2 l_{2,2} + \cdots + \xi_{\nu-1} l_{\nu-1,2} = \lambda_2, \\
\cdots\cdots\cdots\cdots\cdots\cdots\cdots \\
\xi_1 l_{1,\nu} + \xi_2 l_{2,\nu} + \cdots + \xi_{\nu-1} l_{\nu-1,\nu} = \lambda_\nu
\end{array}
}
wird. Denn von diesen Gleichungen sind die ersten $\nu - 1$ ein System linearer Gleichungen mit nicht verschwindender Determinante, und die $\nu^{\mathrm{te}}$ Gleichung folgt aus den \"ubrigen, weil die Summe der conjugirten Logarithmen einer jeden Einheit verschwindet.

Das System der Zahlen $\xi_1, \xi_2, \ldots, \xi_{\nu-1}$ nennen wir das \gesperrt{Exponentensystem} der Einheit $\varepsilon$ in Bezug auf das System $\varepsilon_1, \varepsilon_2, \ldots, \varepsilon_{\nu-1}$, und es ist nun zu beweisen:

6. \gesperrt{Dass die Exponenten jeder Einheit $\varepsilon$ rationale Zahlen sind, deren Nenner eine gewisse endliche Grenze nicht \"ubersteigt, und die daher alle mit einem gemeinschaftlichen, nur von dem Systeme $\varepsilon_i$ abh\"angigen Nenner behaftet angenommen werden k\"onnen.}

Um dies einzusehen, hat man Folgendes zu erw\"agen:

1) Wenn wir die Gr\"ossen $\xi_1, \xi_2, \ldots, \xi_{\nu-1}$ wie Variable betrachten, und jede von ihnen, von den anderen unabh\"angig, von 0 bis 1 gehen lassen, so bleiben die linken Seiten von (5) in endlichen, nur von den Einheiten $\varepsilon_1, \varepsilon_2, \ldots, \varepsilon_{\nu-1}$ abh\"angigen Grenzen eingeschlossen. In diesen Grenzen m\"ussen also auch die conjugirten Logarithmen $\lambda$ der Einheit $\varepsilon$ bleiben, so lange
%CONT>>


% ===== scan leaf 701, printed 680 =====
% !PRINTED 680
\weberpage{680}
%<<CONT
die Exponenten $\xi_i$ auf nicht negative echte gebrochene Werthe beschr\"ankt bleiben, und daraus ergiebt sich nach der Definition der conjugirten Logarithmen eine obere Grenze f\"ur $\varepsilon$ und seine Conjugirten.

Nach dem Satze 2., \S.\,185 giebt es aber in $\Omega$ nur eine endliche Anzahl ganzer Zahlen, also um so mehr nur eine endliche Anzahl von Einheiten, die dem absoluten Werthe nach mit allen ihren Conjugirten unter einer endlichen Grenze liegen.

Wenn wir nun eine Einheit, deren Exponenten in den Grenzen 0 und 1 liegen, mit Einschluss der unteren und mit Ausschluss der oberen Grenze eine \gesperrt{in Bezug auf das System $\varepsilon_1, \varepsilon_2, \ldots, \varepsilon_{\nu-1}$ reducirte Einheit nennen}, so haben wir den Satz:

\gesperrt{Es giebt nur eine endliche Anzahl von Einheiten in $\Omega$, die in Bezug auf ein unabh\"angiges System $\varepsilon_1, \varepsilon_2, \ldots, \varepsilon_{\nu-1}$ reducirt sind.}

2) Die Einheiten reproduciren sich durch Multiplication und Division. Sind
\[
\begin{array}{l}
\xi_1', \xi_2', \ldots, \xi_{\nu-1}' \\
\xi_1'', \xi_2'', \ldots, \xi_{\nu-1}''
\end{array}
\]
die Exponenten von zwei Einheiten $\varepsilon', \varepsilon''$, so sind die Summen und die Differenzen
\[
\xi_1' \pm \xi_1'', \xi_2' \pm \xi_2'', \ldots, \xi_{\nu-1}' \pm \xi_{\nu-1}''
\]
die Exponenten der Einheiten $\varepsilon' \varepsilon''$ und $\varepsilon' : \varepsilon''$.

Dies ergiebt sich unmittelbar aus dem Satze, dass der Logarithmus eines Productes oder eines Quotienten gleich der Summe oder der Differenz der Logarithmen der beiden Bestandtheile ist.

3) Die Einheiten $\varepsilon_1, \varepsilon_2, \ldots$ selbst haben die Exponenten $1, 0, \ldots, 0; 0, 1, \ldots, 0; \ldots$, und daraus folgt nach 2), dass jedes System von ganzen Zahlen, f\"ur $\xi_1, \xi_2, \ldots, \xi_{\nu-1}$ gesetzt, das Exponentensystem einer Einheit giebt, die sich durch Multiplication von Potenzen der $\varepsilon_1, \varepsilon_2, \ldots$ bilden l\"asst.

4) Aus 2) und 3) ergiebt sich, dass ein Exponentensystem $\xi_1, \xi_2, \ldots, \xi_{\nu-1}$ einer Einheit ein Exponentensystem bleibt, wenn alle seine Elemente mit einer und derselben ganzen Zahl multiplicirt werden, und dass es auch dann noch diesen Charakter beh\"alt, wenn seine Elemente $\xi_i$ um beliebige ganze Zahlen vermehrt oder vermindert werden. Gebrauchen wir also das Zeichen $(x)$ in dem Sinne, wie im \S.\,182, dass es den Ueber-
%CONT>>


% ===== scan leaf 702, printed 681 =====
% !PRINTED 681
\weberpage{681}
\pdfbookmark[2]{§ 187. Fundamentalsysteme von Einheiten}{bm214}
%<<CONT
schuss der Zahl $x$ über die grösste in $x$ enthaltene ganze Zahl
bedeutet, so ergiebt sich Folgendes:

Ist $\xi_1, \xi_2, \ldots, \xi_{\nu-1}$ das Exponentensystem einer Einheit, und
$m$ eine ganze Zahl, so ist auch

\weq{6}{(m\,\xi_1), (m\,\xi_2), \ldots, (m\,\xi_{\nu-1})}

das Exponentensystem einer Einheit.

Nun sind die Grössen $(m\,\xi_i)$, wenn sie nicht Null sind,
positive echte Brüche, und nach 1) muss es sich also, wenn die
Reihe der ganzen Zahlen hinlänglich weit fortgesetzt wird, ereignen,
dass für zwei verschiedene Werthe $m'$, $m''$ von $m$ die
Zahlenreihe (6) übereinstimmt. Da hiernach $m'\xi_i$, $m''\xi_i$ denselben
Ueberschuss über eine ganze Zahl haben, so ist $(m' - m'')\,\xi_i$
selbst eine ganze Zahl, und es giebt also eine ganze rationale
Zahl $m$, die höchstens gleich der Anzahl der Exponentensysteme
reducirter Einheiten ist, von der Eigenschaft, dass

\[
m\,\xi_1, m\,\xi_2, \ldots, m\,\xi_{\nu-1}
\]

ganze rationale Zahlen werden.

Diese Zahl $m$ kann sich ändern, wenn die Einheit $\varepsilon$ geändert
wird. Da aber alle möglichen Werthe dieser Zahl unter einer
endlichen Grenze liegen, so giebt es eine endliche Zahl, in der
alle diese Werthe von $m$ enthalten sind und die selbst für $m$
genommen werden kann. Es giebt daher eine gewisse positive
ganze Zahl $m$, die als Nenner aller der rationalen Zahlen genommen
 werden kann, die als Exponenten von Einheiten auftreten
 können. Damit ist der Satz 6. bewiesen.

\begin{center}
\S.\,187.\\[4pt]
\gesperrt{Fundamentalsysteme von Einheiten.}
\end{center}

Wir wählen jetzt an Stelle des Systemes unabhängiger Einheiten
 $\varepsilon_i$ ein anderes, dessen conjugirte Logarithmen

\weq{1}{\lambda_{i,1}, \lambda_{i,2}, \ldots, \lambda_{i,\nu} \qquad i = 1, 2, \ldots, \nu-1}

sein mögen. Nach \S.\,186, (5) ist

\weq{2}{\begin{aligned}
&\lambda_{i,k} = \xi_{1,i}\, l_{1,k} + \xi_{2,i}\, l_{2,k} + \cdots + \xi_{\nu-1,i}\, l_{\nu-1,k} \\
&i = 1, 2, \ldots, \nu-1; \quad k = 1, 2, \ldots, \nu,
\end{aligned}}

wenn $\xi_{1,i}, \xi_{2,i}, \ldots, \xi_{\nu-1,i}$ die Exponenten einer der neuen Einheiten
 in Bezug auf das System $\varepsilon_1, \varepsilon_2, \ldots, \varepsilon_{\nu-1}$ bedeuten, die
wir als rationale Zahlen mit dem gemeinsamen Nenner $m$ annehmen
 können.


% ===== scan leaf 703, printed 682 =====
% !PRINTED 682
\weberpage{682}
Es ist aber nach dem Multiplicationssatze der Determinanten

\weq{3}{\begin{vmatrix} \lambda_{1,1}, & \cdots, & \lambda_{\nu-1,1} \\ \cdot & \cdot & \cdot \\ \lambda_{1,\nu-1}, & \cdots, & \lambda_{\nu-1,\nu-1} \end{vmatrix}
=
\begin{vmatrix} \xi_{1,1}, & \cdots, & \xi_{1,\nu-1} \\ \cdot & \cdot & \cdot \\ \xi_{\nu-1,1}, & \cdots, & \xi_{\nu-1,\nu-1} \end{vmatrix}
\begin{vmatrix} l_{1,1}, & \cdots, & l_{\nu-1,1} \\ \cdot & \cdot & \cdot \\ l_{1,\nu-1}, & \cdots, & l_{\nu-1,\nu-1} \end{vmatrix},}

oder, wenn wir

\[
\Lambda_{\nu-1} = \Sigma \pm \lambda_{1,1}\lambda_{2,2} \cdots \lambda_{\nu-1,\nu-1}
\]

setzen, und beachten, dass die Determinante

\[
\Sigma \pm \xi_{1,1}\xi_{2,2} \cdots \xi_{\nu-1,\nu-1}
\]

eine rationale Zahl mit dem Nenner $m^{\nu-1}$ ist:

\weq{4}{\Lambda_{\nu-1} = \frac{a}{m^{\nu-1}} \, L(\varepsilon_1, \varepsilon_2, \ldots, \varepsilon_{\nu-1}),}

worin $a$ eine ganze rationale Zahl ist.

Daraus ergiebt sich, dass die neu eingeführten Einheiten
immer und nur dann ein unabhängiges System bilden, wenn die
Determinante der $\xi_{i,k}$, d. h. die Zahl $a$, von Null verschieden ist.

Nun giebt es unter einer endlichen oder unendlichen Anzahl
nicht verschwindender rationaler Brüche mit demselben Nenner
immer einen dem absoluten Werthe nach kleinsten, und folglich
können wir nach (4) das neue System unabhängiger Einheiten
so wählen, dass sein Regulator $\pm \Lambda_{\nu-1}$ \gesperrt{so klein als möglich}
wird. Ein solches System nennen wir ein \gesperrt{Fundamentalsystem
von Einheiten}, und den Minimalwerth des Regulators selbst,
also den Regulator eines Fundamentalsystemes, nennen wir den
\gesperrt{Regulator des Körpers}. Wir nehmen jetzt an, dass unser
System $\varepsilon_1, \varepsilon_2, \ldots, \varepsilon_{\nu-1}$ selbst ein Fundamentalsystem sei.

Unter dieser Voraussetzung lässt sich beweisen, dass
alle Exponenten von Einheiten \gesperrt{ganze Zahlen} sein müssen.

Wenn wir nämlich annehmen, es existire eine Einheit $\varepsilon$,
deren nach den Formeln \S.\,186, (5) bestimmte Exponenten nicht
alle ganze Zahlen sind, so giebt es nach \S.\,186, 4) auch eine
Einheit $\varepsilon_0$, deren Exponenten echte Brüche sind, die nicht alle
verschwinden.

Verstehen wir unter $\lambda_1, \lambda_2, \ldots, \lambda_{\nu-1}$ in den Formeln (5),
\S.\,186 das System der conjugirten Logarithmen von $\varepsilon_0$, so ergiebt
 sich

\[
L(\varepsilon_0, \varepsilon_2, \ldots, \varepsilon_{\nu-1}) = \xi_1\, L(\varepsilon_1, \varepsilon_2, \ldots, \varepsilon_{\nu-1}),
\]
%CONT>>


% ===== scan leaf 704, printed 683 =====
% !PRINTED 683
\weberpage{683}
%<<CONT
und wenn wir also annehmen, was offenbar keine Beschränkung
der Allgemeinheit ist, dass $\xi_1$ ein nicht verschwindender echter
Bruch sei, so widerspricht diese Formel der Annahme, dass

\[
\varepsilon_1, \varepsilon_2, \ldots, \varepsilon_{\nu-1}
\]

ein Fundamentalsystem sei, weil die Determinante $L_{\nu-1}$ verkleinert
 würde, wenn $\varepsilon_1$ durch $\varepsilon_0$ ersetzt wird.

Haben wir ein Fundamentalsystem, so können wir ein beliebiges
 System ganzer rationaler Zahlen $\xi_1, \xi_2, \ldots, \xi_{\nu-1}$ als
Exponentensystem annehmen und eine Einheit $\varepsilon$ mit diesen Exponenten
 bilden:

\[
\varepsilon = \varepsilon_1^{\xi_1} \varepsilon_2^{\xi_2} \ldots \varepsilon_{\nu-1}^{\xi_{\nu-1}}.
\]

Es bleibt also noch die Frage zu beantworten, inwieweit
eine Einheit durch das Exponentensystem bestimmt ist.

Nehmen wir an, es seien $\varepsilon', \varepsilon''$ zwei Einheiten mit demselben
Exponentensysteme, dann besteht das Exponentensystem der
Einheit $\varepsilon' : \varepsilon'' = \varrho$ aus lauter Nullen, und folglich sind die conjugirten
 Logarithmen von $\varrho$ alle gleich Null; $\varrho$ ist daher eine
ganze Zahl von der Eigenschaft, dass die absoluten Werthe aller
mit $\varrho$ conjugirten Zahlen gleich 1 sind, und daher, wie im
\S.\,175, 2. bewiesen ist, eine \gesperrt{Einheitswurzel}. Damit ist das
folgende Theorem bewiesen:

I. \gesperrt{Es giebt im Körper $\Omega$ ein System von $\nu - 1$ fundamentalen
 Einheiten}

\weq{5}{\varepsilon_1, \varepsilon_2, \ldots, \varepsilon_{\nu-1},}

\gesperrt{welches die Eigenschaft hat, dass in der Form}

\weq{6}{\varepsilon = \varrho\, \varepsilon_1^{\xi_1} \varepsilon_2^{\xi_2} \ldots \varepsilon_{\nu-1}^{\xi_{\nu-1}}}

\gesperrt{alle Einheiten des Körpers, jede nur einmal, enthalten sind,
wenn $\xi_1, \xi_2, \ldots, \xi_{\nu-1}$ alle ganzen rationalen Zahlen und $\varrho$ alle
in $\Omega$ vorhandenen Einheitswurzeln durchläuft.}

Es ist nun leicht, aus einem Fundamentalsysteme alle anderen
abzuleiten, indem man in (6) für die Exponenten $\nu - 1$ verschiedene
 Systeme ganzer Zahlen $\xi_{i,k}$ setzt, deren Determinante
$= \pm 1$ ist. Denn setzt man

\[
\varepsilon_i' = \varrho_i\, \varepsilon_1^{\xi_{1,i}} \varepsilon_2^{\xi_{2,i}} \ldots \varepsilon_{\nu-1}^{\xi_{\nu-1,i}},
\]

so folgt aus (3):

\[
L(\varepsilon_1', \varepsilon_2', \ldots, \varepsilon_{\nu-1}') = \pm L(\varepsilon_1, \varepsilon_2, \ldots, \varepsilon_{\nu-1}).
\]


% ===== scan leaf 705, printed 684 =====
% !PRINTED 684
\weberpage{684}
\pdfbookmark[2]{§ 188. Reducirte Zahlen}{bm215}
% [source note] The same-parity restriction on x,y below applies only to odd m. For a fundamental discriminant -m the general condition is x congruent to m*y modulo 2; for m=4, x=0,y=1 represents i. The printed blanket restriction is retained.
Beide Determinanten haben also den Minimalwerth, und
beide Systeme sind Fundamentalsysteme.

Der Fall $\nu = 1$, den wir oben ausgeschlossen haben, tritt
ausser im Falle des rationalen Körpers, in dem nur die beiden
Einheiten $\pm 1$ existiren, im Falle des \gesperrt{imaginären quadratischen
 Körpers} ein. In diesem Falle haben wir im zwanzigsten
 Abschnitte die ganzen Zahlen des Körpers in der Form

\[
\frac{x + y\sqrt{-m}}{2}
\]

dargestellt, worin $-m$ eine Stammdiscriminante bedeutet, so
dass $m \equiv 0$ oder $\equiv 3\ (\mathrm{mod}\ 4)$ ist, und $x, y$ beide gerade oder
beide ungerade sind. Die Einheiten erhalten wir, wenn wir die
Gleichung

\[
x^2 + m y^2 = 4
\]

auf alle mögliche Arten in ganzen rationalen Zahlen lösen. Diese
Gleichung hat aber, wenn $m > 4$ ist, nur die zwei Lösungen
$x = \pm 2$, $y = 0$, und es giebt also in diesen Fällen, wie im
Körper der rationalen Zahlen, nur die zwei Einheiten $\pm 1$. Ist
$m = 4$, so findet man die vier Einheiten $\pm 1, \pm i$, und ist endlich
 $m = 3$, die sechs Einheiten

\[
\pm 1, \; \pm \frac{-1 + i\sqrt{3}}{2}, \; \pm \frac{-1 - i\sqrt{3}}{2}.
\]

In dem nächst einfachen Falle, $n = 2$, $\nu = 2$, d. h. im
reellen quadratischen Körper, fällt die Theorie der Einheiten
zusammen mit der im \S.\,127 des ersten Bandes behandelten
Theorie der Pell'schen Gleichung.

\begin{center}
\S.\,188.\\[4pt]
\gesperrt{Reducirte Zahlen.}
\end{center}

Ist $\alpha$ irgend eine von Null verschiedene ganze oder gebrochene
 Zahl des Körpers $\Omega$, so betrachten wir, wie bei den
Einheiten, das \gesperrt{System der conjugirten Logarithmen} von $\alpha$,
worunter wir, wie im \S.\,185, die \gesperrt{reellen Zahlen}

\weq{1}{\lambda_1 = \delta_1 \log|\alpha_1|, \; \lambda_2 = \delta_2 \log|\alpha_2|, \; \ldots, \; \lambda_\nu = \delta_\nu \log|\alpha_\nu|}

verstehen, die wir auch mit $\lambda_s(\alpha)$ bezeichnen. Ziehen wir das
Fundamentalsystem $\varepsilon_1, \varepsilon_2, \ldots, \varepsilon_{\nu-1}$ von Einheiten zu Rathe, so
%CONT>>


% ===== scan leaf 706, printed 685 =====
% !PRINTED 685
\weberpage{685}
%<<CONT
können wir die linearen Gleichungen, analog wie im \S.\,186, (5)
bilden:
\weq{2}{
\begin{aligned}
\xi_1 l_{1,1} + \xi_2 l_{2,1} + \cdots + \xi_{\nu-1} l_{\nu-1,1} + \delta_1 \xi_\nu &= \lambda_1, \\
\xi_1 l_{1,2} + \xi_2 l_{2,2} + \cdots + \xi_{\nu-1} l_{\nu-1,2} + \delta_2 \xi_\nu &= \lambda_2, \\
\cdot\ \cdot\ \cdot\ \cdot\ \cdot\ \cdot\ \cdot\ \cdot\ \cdot\ \cdot\ \cdot\ \cdot & \\
\xi_1 l_{1,\nu} + \xi_2 l_{2,\nu} + \cdots + \xi_{\nu-1} l_{\nu-1,\nu} + \delta_\nu \xi_\nu &= \lambda_\nu,
\end{aligned}
}
dessen Determinante [\S.\,186, (4)] den von Null verschiedenen
Werth hat:
\weq{3}{
\left|
\begin{array}{ccccc}
l_{1,1}, & l_{2,1}, & \ldots, & l_{\nu-1,1}, & \delta_1\\
l_{1,2}, & l_{2,2}, & \ldots, & l_{\nu-1,2}, & \delta_2\\
\cdot & \cdot & \cdot & \cdot & \cdot\\
l_{1,\nu}, & l_{2,\nu}, & \ldots, & l_{\nu-1,\nu}, & \delta_\nu
\end{array}
\right| = n\,L\,(\varepsilon_1, \varepsilon_2, \ldots, \varepsilon_{\nu-1})
}
Demnach sind die Unbekannten $\xi_1, \xi_2, \ldots, \xi_{\nu-1}, \xi_\nu$ aus (2)
\gesperrt{eindeutig} bestimmt.

Durch Addition der Gleichungen (2) findet man zunächst
sehr einfach
\weq{4}{
n\,\xi_\nu = \log N_a(\alpha),
}
und $\xi_\nu$ bleibt also dasselbe für alle \gesperrt{Zahlen $\alpha$, die unter einander
 associirt sind.}

Die Zahlen $\xi_1, \xi_2, \ldots, \xi_{\nu-1}$ heissen die Exponenten der
Zahl $\alpha$ [in Bezug auf das Fundamentalsystem $(\varepsilon_1, \varepsilon_2, \ldots, \varepsilon_{\nu-1})$].

Aus dieser Definition ergiebt sich ohne Weiteres, dass die
Exponenten eines Productes oder eines Quotienten gleich der
Summe oder der Differenz der entsprechenden Exponenten der
einzelnen Bestandtheile sind. Die Exponenten einer Einheit sind
ganze rationale Zahlen, und jedes System ganzer rationaler Zahlen
$\xi_1, \xi_2, \ldots, \xi_{\nu-1}$ ist auch das Exponentensystem einer Einheit.

Die Exponenten associirter Zahlen unterscheiden sich also
um ganze Zahlen von einander, und man kann zu jeder Zahl $\alpha$
eine associirte Zahl $\alpha_0$ finden, deren Exponenten zwischen 0
und 1 liegen. Solche Zahlen heissen \gesperrt{reducirte Zahlen} (in
Bezug auf das Fundamentalsystem $\varepsilon_1, \ldots, \varepsilon_{\nu-1}$). Reducirte Einheiten
sind alle und nur die in $\Omega$ vorhandenen \gesperrt{Einheitswurzeln,
deren Zahl wir mit $w$ bezeichnen wollen.} Da
die beiden Einheitswurzeln $\pm 1$ in jedem Körper enthalten sind,
so ist $w$ mindestens $= 2$ und immer eine gerade Zahl.

Wenn zwei associirte Zahlen dasselbe Exponentensystem
haben, so unterscheiden sie sich nur durch einen Factor, der
%CONT>>


% ===== scan leaf 707, printed 686 =====
% !PRINTED 686
\weberpage{686}
\pdfbookmark[2]{§ 189. Grenzen der Anzahl der durch ein Ideal theilbaren ganzen Zahlen des Körpers Ω}{bm216}
%<<CONT
eine Einheitswurzel ist. Das Exponentensystem der aus $\alpha$ abgeleiteten
 reducirten Zahl $\alpha_0$ ist durch $\alpha$ selbst völlig bestimmt.
Hierdurch ist der Satz bewiesen:

II. \gesperrt{Zu jeder (von Null verschiedenen) Zahl $\alpha$ des
Körpers $\Omega$ giebt es $w$ und nicht mehr reducirte
Zahlen $\varrho\,\alpha_0$.}

Wenn statt des Fundamentalsystemes $(\varepsilon_1, \varepsilon_2, \ldots, \varepsilon_{\nu-1})$ ein
anderes angewendet wird, so erleiden die Exponenten $\xi_1, \xi_2, \ldots, \xi_{\nu-1}$
eine ganzzahlige lineare Substitution von der Determinante $\pm 1$.
Eine reducirte Zahl kann dann aufhören, für das neue Fundamentalsystem
 reducirt zu sein. Der Begriff der reducirten Zahl
ist also von der Wahl des Fundamentalsystemes abhängig. Die
Anzahl der aus einer gegebenen Zahl abgeleiteten reducirten
Zahlen ist aber immer dieselbe.

Durch diesen Satz haben wir den Zweck erreicht, aus
dem ganzen unendlichen Systeme der unter einander associirten
Zahlen eine bestimmte endliche Anzahl von Repräsentanten
herausgehoben zu haben.

\begin{center}
\S.\,189.\\[4pt]
\gesperrt{Grenzen der Anzahl der durch ein Ideal theilbaren
ganzen Zahlen des Körpers $\Omega$.}
\end{center}
% [source note] The following finiteness assertion and definition of T omit counting modulo association. Infinitely many units already have absolute norm 1. The intended finite count is of non-associated integers (equivalently the corresponding principal ideals), explicitly specified on printed p. 690. The source also changes the endpoint from norm at most t here to natural numbers less than t below; this can alter a boundary term but not the limiting quotient. The wording is preserved.

Unsere früheren Betrachtungen haben ergeben, dass jede
ganze Zahl eines Körpers $\Omega$ nur eine endliche Anzahl von Idealen
zu Theilern hat. Da jedes ganze Ideal ein Theiler seiner Norm
ist, so giebt es also auch nur eine endliche Anzahl von ganzen
Zahlen in $\Omega$, deren absolute Norm eine gegebene Grenze nicht
übersteigt.

Aus diesen Zahlen greifen wir jetzt wieder einen Theil
heraus, und fragen nach der Anzahl $T$ aller \gesperrt{ganzen Zahlen}
des Körpers $\Omega$, deren absolute Norm eine positive Grösse $t$
nicht überschreitet, und die durch ein gegebenes Ideal $\mathfrak{a}$ theilbar
 sind.

Nehmen wir als vorläufiges Beispiel den rationalen Körper,
so ist dort $T$ die Anzahl der natürlichen Zahlen, die kleiner als
$t$ und durch eine gegebene ganze Zahl $m$ theilbar sind, also, in
der früher (\S.\,182) gebrauchten Bezeichnung $T = E\left(\frac{t}{m}\right)$.


% ===== scan leaf 708, printed 687 =====
% !PRINTED 687
\weberpage{687}
Die Zahl $T$ ist eine im Allgemeinen nicht genauer zu bestimmende
 Function von $t$, die mit $t$ zugleich ins Unendliche
wächst, die sich aber, da sie nur ganzzahlige Werthe hat, nur
stufenweise, also unstetig ändert. Worauf es hier ankommt, ist
\gesperrt{der Grenzwerth, dem sich das Verhältniss $T:t$ mit
unendlich wachsendem $t$ nähert} (der in dem oben angeführten
 Beispiele $1:m$ ist).

Dieser Grenzwerth lässt sich durch Betrachtungen, wie wir
sie schon im \S.\,155 benutzt haben, finden.

Wir bilden uns zunächst eine Basisform des gegebenen Ideals $\mathfrak{a}$:
\weq{1}{
\alpha = \alpha_1 x_1 + \alpha_2 x_2 + \cdots + \alpha_n x_n,
}
d.\,h. eine lineare Function der $n$ Variablen $x_i$, von der Eigenschaft,
dass durch Substitution aller ganzen rationalen Zahlen für
$x_i$ aus $\alpha$ alle durch $\mathfrak{a}$ theilbare ganze Zahlen des Körpers $\Omega$ entstehen
 (\S.\,146). Mit $\alpha_{i,1}, \alpha_{i,2}, \ldots, \alpha_{i,n}$ bezeichnen wir die mit $\alpha_i$
conjugirten Werthe und bilden die lineare Substitution für $x_i$:
\weq{2}{
\begin{aligned}
y_1 &= \alpha_{1,1} x_1 + \alpha_{2,1} x_2 + \cdots + \alpha_{n,1} x_n\\
y_2 &= \alpha_{1,2} x_1 + \alpha_{2,2} x_2 + \cdots + \alpha_{n,2} x_n\\
&\ \cdot\ \cdot\ \cdot\ \cdot\ \cdot\ \cdot\ \cdot\ \cdot\ \cdot\ \cdot\ \cdot\ \cdot\\
y_n &= \alpha_{1,n} x_1 + \alpha_{2,n} x_2 + \cdots + \alpha_{n,n} x_n.
\end{aligned}
}
Ist $A$ die Determinante dieses Gleichungssystems
\[
A = \Sigma \pm \alpha_{1,1} \alpha_{2,2} \ldots \alpha_{n,n},
\]
so ist $A^2$ die Discriminante des Functionensystems $(\alpha_1, \alpha_2, \ldots, \alpha_n)$,
und daher nach \S.\,147 gleich dem Producte aus dem Quadrate
der absoluten Norm von $\alpha$ und der Körperdiscriminante $\Delta$.
Da die absolute Norm von $\alpha$ zugleich die Norm des Ideals $\mathfrak{a}$
ist (\S.\,151), so setzen wir
\weq{3}{
A^2 = N(\mathfrak{a})^2 \, \Delta.
}
Nun können wir, wie im \S.\,155, die $x_1, x_2, \ldots, x_n$ als
Coordinaten eines Punktes $(x)$ in einem Raume von $n$ Dimensionen
 deuten, und wir lassen jedem Punkt $(x)$ einen Punkt $(x')$
entsprechen, dessen Coordinaten durch
\weq{4}{
x_1' = t^{\frac{1}{n}} x_1, \quad x_2' = t^{\frac{1}{n}} x_2, \quad \ldots, \quad x_n' = t^{\frac{1}{n}} x_n
}
bestimmt sind, worin $t$ eine beliebige positive Grösse bedeutet.
Jedem Gebiete $S$ von Punkten $(x)$ entspricht ein Gebiet $S'$ von
Punkten $(x')$. Alle Figuren in $S$ sind den entsprechenden Figuren
in $S'$ ähnlich, und die Lineardimensionen in $S$ verhalten sich
zu denen in $S'$ wie $1 : t^{\frac{1}{n}}$.


% ===== scan leaf 709, printed 688 =====
% !PRINTED 688
\weberpage{688}
Denken wir uns das Gebiet $S$ gegeben, so wird das entsprechende
 Gebiet $S'$ von $t$ abhängig sein und sich mit $t$ vergrössern.
Die Punkte mit \gesperrt{ganzzahligen} Coordinaten $x_i$ nennen
wir, wie früher, die \gesperrt{Gitterpunkte}. Die Anzahl der in einem
endlichen Gebiete $S'$ liegenden Gitterpunkte ist immer endlich,
wächst aber mit $t$ und soll mit $Z_t$ bezeichnet werden. Dann ist
nach den Sätzen \S.\,155, (14) das Volumen des Gebietes $S$, d.\,h. das
über alle Punkte von $S$ erstreckte $n$-fache Integral
\[
V = \int\!\!\int \ldots \int dx_1\, dx_2 \ldots dx_n
\]
gleich dem Grenzwerthe des Verhältnisses $Z_t : t$ für unendlich
wachsendes $t$:
\weq{5}{
V = \mathop{\mathrm{Lim}}\limits_{t=\infty} \frac{Z_t}{t}.
}
Hierin liegt zugleich die genauere Präcisirung, die wir über
das Gebiet $S$ machen müssen, wenn unsere Betrachtung anwendbar
 sein soll. Es muss $S$ so beschaffen sein, dass das
$n$-fache Integral $V$ eine bestimmte Bedeutung hat. Jeder Gitterpunkt
 ist nun durch (1) das Bild einer ganzen durch $\mathfrak{a}$ theilbaren
 Zahl des Körpers $\Omega$, und wir können also auch sagen,
dass $Z_t$ die Anzahl der durch $\mathfrak{a}$ theilbaren ganzen Zahlen in $\Omega$
ist, deren Bilder in dem Gebiete $S'$ liegen.

Um das Gebiet $S'$ in geeigneter Weise abzugrenzen, wählen
wir ein System fundamentaler Einheiten des Körpers $\Omega$
\[
\varepsilon_1, \varepsilon_2, \ldots, \varepsilon_{\nu-1}
\]
mit dem System der conjugirten Logarithmen
\[
l_{i,1}, l_{i,2}, \ldots, l_{i,\nu}
\]
und dem Regulator
\weq{6}{
\pm L(\varepsilon_1, \varepsilon_2, \ldots, \varepsilon_{\nu-1}) = \pm \Sigma \pm l_{1,1} l_{2,2} \ldots l_{\nu-1,\nu-1},
}
der mit $L$ bezeichnet werden mag.

Wir definiren jetzt ein System von Functionen $\xi_1, \xi_2, \ldots, \xi_\nu$
der Variablen $x_i$ durch folgende Bedingungen:

Es sei, wenn die $y_i$ die Bedeutung (2) haben,
\weq{7}{
z_i = \delta_i \log |y_i|,
}
so dass, wenn für die $x_i$ ein System ganzer rationaler Zahlen
gesetzt wird, $z_i$ nach \S.\,185 in das System der conjugirten
Logarithmen der Zahl $\alpha$ übergeht.

Es ist dann, wenn $y_k$ zu einem reellen Körper gehört,
\weq{8}{
z_k = \tfrac{1}{2} \log y_k^2,
}
%CONT>>


% ===== scan leaf 710, printed 689 =====
% !PRINTED 689
\weberpage{689}
%<<CONT
und wenn $y_k, y_k'$ einem imaginären Paare angehören,
\weq{9}{z_k = \log y_k y_k',}
wodurch die Definition völlig eindeutig ist.

Nun bestimmen wir die Variablen $\xi_1, \xi_2, \ldots, \xi_{\nu-1}, \xi_\nu$ aus den linearen Gleichungen
\weq{10}{\begin{array}{l}
\xi_1 l_{1,1} + \cdots + \xi_{\nu-1} l_{\nu-1,1} + \delta_1 \xi_\nu = z_1\\
\xi_1 l_{1,2} + \cdots + \xi_{\nu-1} l_{\nu-1,2} + \delta_2 \xi_\nu = z_2\\
. \; . \; . \; . \; . \; . \; . \; . \; . \; . \; . \; . \; . \; .\\
\xi_1 l_{1,\nu} + \cdots + \xi_{\nu-1} l_{\nu-1,\nu} + \delta_\nu \xi_\nu = z_\nu,
\end{array}}
deren Determinante nicht verschwindet.

Giebt man den Variablen $x_i$ ganzzahlige Werthe, so gehen nach \S.\,188, (2), (3) die $\xi_1, \xi_2, \ldots, \xi_{\nu-1}$ in das Exponentensystem der Zahl $\alpha$ über, und es ergiebt sich durch Addition der Gleichungen (10) [\S.\,185, (3)]:
\weq{11}{\xi_\nu = \frac{1}{n} \log N_a(\alpha).}

Wir gehen jetzt von dem Punkte $(x)$ zu dem Punkte $(x')$ über, indem wir $x_i' = t^{\frac{1}{n}} x_i$ setzen. Dadurch geht $y_i$ in
\[ y_i' = t^{\frac{1}{n}} y_i \]
über, und $z_i$ in
\[ z_i' = z_i + \frac{1}{n} \delta_i \log t. \]

Wenn wir also $z_i'$ an Stelle von $z_i$ setzen, und die dann aus (10) sich ergebenden Werthe der $\xi_1, \xi_2, \ldots, \xi_\nu$ mit $\xi_1', \xi_2', \ldots, \xi_\nu'$ bezeichnen, so ist
\weq{12}{\xi_1' = \xi_1, \xi_2' = \xi_2, \ldots, \xi_{\nu-1}' = \xi_{\nu-1}, \xi_\nu' = \xi_\nu + \frac{1}{n} \log t.}

Nunmehr wollen wir das Gebiet $S'$ dadurch abgrenzen, dass
\weq{13}{\begin{array}{c}
0 \leqq \xi_1' < 1, \ldots, 0 \leqq \xi_{\nu-1}' < 1, \\
\xi_\nu' < \frac{1}{n} \log t
\end{array}}
sein soll. Dadurch erreichen wir, dass die in dem Gebiete $S'$ liegenden Gitterpunkte $(x')$ nur reducirte ganze Zahlen $\alpha$ darstellen, und alle und nur solche, deren absolute Norm kleiner als $t$ ist.

Unter den reducirten Zahlen sind aber immer je $w$ mit einander associirt, wenn $w$ die Anzahl der in $\Omega$ enthaltenen
%CONT>>


% ===== scan leaf 713, printed 690 =====
% !PRINTED 690
\weberpage{690}
\pdfbookmark[2]{§ 190. Bestimmung des Volumens}{bm217}
%<<CONT
Einheitswurzeln bedeutet, und wenn wir also diese $w$ associirten Zahlen zu einem Complex zusammenfassen, so ist die Anzahl dieser Complexe gleich der Anzahl $T$ aller nicht associirten, durch $\mathfrak{a}$ theilbaren ganzen Zahlen in $\Omega$, deren absolute Norm kleiner als $t$ ist. Demnach ist die Anzahl der in $S'$ liegenden Gitterpunkte
\weq{14}{Z_t = w T.}

Für die Begrenzung des Gebietes $S$ erhält man aus (12) und (13):
\weq{15}{0 \leqq \xi_1 < 1, \ldots, 0 \leqq \xi_{\nu-1} < 1, \xi_\nu < 0,}
und nach (5) erhalten wir
\weq{16}{\mathrm{Lim}\, \frac{T}{t} = \frac{1}{w} \int\!\int \ldots \int dx_1\, dx_2 \ldots dx_n = \frac{1}{w} V,}
worin das $n$-fache Integral über das durch (15) bestimmte Gebiet $S$ auszudehnen ist.

\begin{center}
\S.\,190.

\gesperrt{Bestimmung des Volumens.}
\end{center}

Das Volumen $V$ lässt sich bestimmen nach der schon im \S.\,155 angeführten Transformationsformel für mehrfache Integrale, nach der, wenn $x_1, x_2, \ldots, x_n$ Functionen der Variablen $u_1, u_2, \ldots, u_n$ sind,
\weq{1}{\begin{aligned}
V &= \int\!\int \ldots \int dx_1\, dx_2 \ldots dx_n\\
&= \int\!\int \ldots \int \left( \Sigma \pm \frac{\partial x_1}{\partial u_1} \frac{\partial x_2}{\partial u_2} \cdots \frac{\partial x_n}{\partial u_n} \right) du_1\, du_2 \ldots du_n
\end{aligned}}
ist, wobei die Functionaldeterminante in dem nach den Variablen $u$ genommenen Integrale mit ihrem absoluten Werthe zu nehmen ist.

Diese Functionaldeterminante ist auch gleich dem reciproken Werthe der Determinante
\[ \Sigma \pm \frac{\partial u_1}{\partial x_1} \frac{\partial u_2}{\partial x_2} \cdots \frac{\partial u_n}{\partial x_n}. \]

Wir führen zunächst die Variablen $u$ durch eine lineare Substitution ein, und bemerken, dass unter den mit $\Omega$ conjugirten Körpern nach unserer Annahme $n - \nu$ conjugirt imaginäre Paare
%CONT>>


% ===== scan leaf 714, printed 691 =====
% !PRINTED 691
\weberpage{691}
%<<CONT
und $2\nu - n$ reelle vorkommen. Demnach sind von den $n$ Functionen $y$ [§. 189, (2)]

\weq{2}{y_k = \alpha_{1,k}x_1 + \alpha_{2,k}x_2 + \cdots + \alpha_{n,k}x_n.}

$2\nu - n$ reell, $2n-2\nu$ paarweise conjugirt. Wir wollen unter den $n$ Variablen $u$ reelle lineare Verbindungen der $x$ verstehen, nämlich die $2\nu-n$ reellen $y$ und die $2n-2\nu$ Bestandtheile der conjugirt imaginären Paare der $y$, so dass, wenn $y, y'$ ein imaginäres Paar ist,
\[
2u_1 = y+y', \quad 2iu_2 = y-y'
\]
oder

\weq{3}{y = u_1+iu_2, \quad y' = u_1-iu_2}

gesetzt wird. Wenn man in der Functionaldeterminante der $y$ auf jedes Paar conjugirt imaginärer Zeilen zweimal den elementaren Determinantensatz von der Addition der Zeilen anwendet (Bd. I, §. 22), so ergiebt sich

\[
\begin{aligned}
\Sigma \pm \frac{\partial u_1}{\partial x_1}\frac{\partial u_2}{\partial x_2}\cdots\frac{\partial u_n}{\partial x_n}
&= \frac{1}{(-2i)^{n-\nu}}\Sigma\pm\frac{\partial y_1}{\partial x_1}\frac{\partial y_2}{\partial x_2}\cdots\frac{\partial y_n}{\partial x_n} \\
&= \frac{A}{(-2i)^{n-\nu}},
\end{aligned}
\]

worin $A$ die in §. 189, (3) festgesetzte Bedeutung hat.

Die Körperdiscriminante ist positiv oder negativ, je nachdem die Anzahl der conjugirt imaginären Paare, also $n-\nu$, gerade oder ungerade ist, und demnach erhält man, wenn mit $\pm\varDelta$ der absolute Werth der Körperdiscriminante bezeichnet wird, nach §. 189, (4), (16):

\weq{4}{\mathrm{Lim}\,\frac{T}{t} = \frac{2^{n-\nu}}{w\,N(\mathfrak{a})\sqrt{\pm\varDelta}}\int\int\cdots\int du_1\,du_2\ldots du_n.}

Es ist also weiter noch das Integral

\[
U = \int\int\cdots\int du_1\,du_2\ldots du_n
\]

zu behandeln.

Aus den Gleichungen §. 189, (10) erhält man zu jedem den Bedingungen (15) genügenden Werthsysteme der Variablen $\xi_1, \xi_2, \ldots, \xi_\nu$ ein bestimmtes endliches System der Werthe für $z_1, z_2, \ldots, z_\nu$, und daraus nach §. 189, (7) für jedes $y$ einen von Null verschiedenen absoluten Werth. Setzt man im Falle eines imaginären Paares (3)

\weq{5}{u_1 = e^{1/2\,z}\cos\vartheta, \quad u_2 = e^{1/2\,z}\sin\vartheta,}

\weq{6}{y_1y_1' = u_1^2+u_2^2 = e^z,}
%CONT>>


% ===== scan leaf 715, printed 692 =====
% !PRINTED 692
\weberpage{692}
%<<CONT
so ergiebt sich, wenn $\vartheta$ in den Grenzen
\[
0 \leqq \vartheta < 2\pi
\]
genommen wird, bei feststehendem $z$ zu jedem Werthe von $\vartheta$ ein zugehöriges Werthepaar $u_1, u_2$. Ein solcher Winkel $\vartheta$ gehört zu jedem imaginären Paare.

Es gehört dann zu jedem Punkte des Gebietes $S$ ein und nur ein Werthsystem der Variablen $\xi_1, \xi_2, \ldots, \xi_\nu, \vartheta \ldots$ in der Begrenzung

\weq{7}{\begin{aligned}
&0 \leqq \xi_1 < 1, \ldots, 0 \leqq \xi_{\nu-1} < 1 \\
&\xi_\nu < 0, \qquad 0 \leqq \vartheta < 2\pi, \ldots
\end{aligned}}
% [source notation] Both first-row less-than-or-equal signs place the equality bars above the angle; other occurrences on this page place them below. Both variants express the same relation and are represented by \leqq. No damaged-type or reversed-inequality claim is made.

Wenn umgekehrt irgend ein Werthsystem der Variablen $\xi, \vartheta$ in den Grenzen (7) gegeben ist, so erhält man hieraus die entsprechenden Bestandtheile $u_1, u_2$ eines imaginären $y$ eindeutig, von den reellen $y$ aber nur die absoluten Werthe, und es kann also jedem der reellen $y$ noch jedes der beiden Vorzeichen gegeben werden. Die Anzahl dieser möglichen Bestimmungen ist $2^{2\nu-n}$. Ist eine dieser Bestimmungen herausgegriffen, so sind die Variablen $x$ eindeutig bestimmt und sind für jedes endliche Werthsystem der $\xi, \vartheta$ endlich. Für ein gegen $-\infty$ abnehmendes $\xi_\nu$ werden die Werthe der $y$ zum Theil gegen Null convergiren.

Das Gebiet $S$ zerfällt demnach in $2^{2\nu-n}$ Theilgebiete $S_1, S_2, \ldots$, deren jedes durch eine bestimmte Vorzeichen-Combination der reellen $y_1, y_2, \ldots$ charakterisirt ist, und demnach zerfällt das Integral $U$ in ebenso viele Theilintegrale $U_1, U_2, \ldots$

\[
U = U_1+U_2+\cdots
\]

Bei der Ausführung der Integration in einem dieser Bestandtheile, die sich alle als von gleichem Werthe ergeben, fangen wir mit einem imaginären Paare an, wenn ein solches vorhanden ist.

Betrachten wir $u_1, u_2$ als rechtwinkelige Coordinaten in einer Ebene, so können wir $e^{1/2\,z}, \vartheta$ nach (5) als Polarcoordinaten auffassen, und erhalten

\[
\int\int du_1\,du_2 = \frac12\int\int e^z\,dz\,d\vartheta,
\]

worin die Integration in Bezug auf $\vartheta$, die sich von $0$ bis $2\pi$ erstreckt, ausgeführt werden kann. Man findet so

\[
\int\int du_1\,du_2 = \pi\int e^z\,dz.
\]


% ===== scan leaf 716, printed 693 =====
% !PRINTED 693
\weberpage{693}
Nun ändert man die Reihenfolge der Integration, und stellt ein etwa noch vorhandenes zweites imaginäres Paar voran, das man wieder in derselben Weise behandelt.

Die verschiedenen Bestandtheile $U_1, U_2, \ldots$ unterscheiden sich durch die Vorzeichen der den reellen $y$ entsprechenden $u$. Nehmen wir das Zeichen so, dass $\pm u$ positiv wird, und setzen

\weq{8}{z = \log(\pm u), \quad \pm du = e^z\,dz,}

so erhalten wir für die verschiedenen $U_1, U_2, \ldots$ denselben Ausdruck durch ein $\nu$-faches Integral

\weq{9}{U_1 = \pi^{n-\nu}\int\int\cdots\int e^{z_1+z_2+\cdots+z_\nu}\,dz_1\,dz_2\ldots dz_\nu,}

wenn wir $dz_1, dz_2, \ldots, dz_\nu$ positiv annehmen.

Die Variablen $z_k$ sind aber nach (6) und (8) hier keine anderen, als in §. 189, (8), (9), (10), und danach ist

\[
z_1+z_2+\cdots+z_\nu = n\xi_\nu.
\]

Wenn wir nun die Transformationsformel (1) auf das $\nu$-fache Integral (9) anwenden, um die Variablen $\xi_1, \xi_2, \ldots, \xi_\nu$ einzuführen, so haben wir die Determinante

\[
\Sigma\pm\frac{\partial z_1}{\partial\xi_1}\frac{\partial z_2}{\partial\xi_2}\cdots\frac{\partial z_\nu}{\partial\xi_\nu}
\]

zu bilden, die nach §. 188, (3) dem absoluten Werthe nach mit $nL$, d. h. mit dem $n$fachen Werthe des \gesperrt{Körperregulators} $L$ übereinstimmt, und danach ist

\[
U_1 = \pi^{n-\nu}nL\int\limits_0^1 d\xi_1\ldots\int\limits_0^1 d\xi_{\nu-1}\int\limits_{-\infty}^0 e^{n\xi_\nu}\,d\xi_\nu = \pi^{n-\nu}L
\]

und

\[
U = 2^{2\nu-n}\pi^{n-\nu}L.
\]

Daraus ergiebt sich also endlich nach (4) der Satz:

1. \gesperrt{Bedeutet $T$ die Anzahl der durch $\mathfrak{a}$ theilbaren nicht associirten ganzen Zahlen, deren absolute Norm kleiner als $t$ ist, so ist}

\weq{10}{\mathop{\mathrm{Lim}}\limits_{t=\infty}\frac{T}{t} = \frac{2^\nu\pi^{n-\nu}L}{w\,N(\mathfrak{a})\sqrt{\pm\varDelta}} = \frac{g}{N(\mathfrak{a})},}

\gesperrt{worin $g$ eine durch die Natur des Körpers $\Omega$ völlig bestimmte positive Zahl ist, nämlich:}

\weq{11}{g = \frac{2^\nu\pi^{n-\nu}L}{w\sqrt{\pm\varDelta}}.}


% ===== scan leaf 717, printed 694 =====
% !PRINTED 694
\weberpage{694}
\pdfbookmark[2]{§ 191. Sätze aus der Reihenlehre}{bm218}
\begin{center}
§. 191.

\gesperrt{Sätze aus der Reihenlehre.}
\end{center}

Bei den weiteren Anwendungen der bisherigen Resultate sind einige Sätze aus der Lehre von den unendlichen Reihen erforderlich, die zunächst hier abgeleitet werden sollen.

1. \gesperrt{Ist $a_1, a_2, a_3, \ldots, a_n\ldots$ ein unbegrenztes System reeller (positiver, negativer oder auch verschwindender) Grössen, von dem wir voraussetzen, dass die Summe}

\weq{1}{\sigma_n = a_1+a_2+\cdots+a_n}

\gesperrt{für jedes beliebige $n$ dem absoluten Werthe nach unter einer endlichen Grenze $C$ bleibt, so ist die Reihe}

\weq{2}{S = \frac{a_1}{1^s}+\frac{a_2}{2^s}+\frac{a_3}{3^s}+\cdots}

\gesperrt{für jedes positive $s$ nicht nur convergent, sondern auch eine stetige Function von $s$.}

Um diesen Satz zu beweisen, zerlegen wir $S$ in der Weise:
\[
S = S_m+R_m,
\]
worin
\[
S_m = \frac{a_1}{1^s}+\frac{a_2}{2^s}+\cdots+\frac{a_{m-1}}{(m-1)^s},
\]
\[
R_m = \frac{a_m}{m^s}+\frac{a_{m+1}}{(m+1)^s}+\frac{a_{m+2}}{(m+2)^s}+\cdots
\]
Nun ist nach (1):
\[
a_m = \sigma_m-\sigma_{m-1}, \quad a_{m+1}=\sigma_{m+1}-\sigma_m, \quad a_{m+2}=\sigma_{m+2}-\sigma_{m+1}, \ldots
\]
und daher
\[
\begin{aligned}
R_m+\frac{\sigma_{m-1}}{m^s}
&= \sigma_m\left(\frac{1}{m^s}-\frac{1}{(m+1)^s}\right) \\
&\quad +\sigma_{m+1}\left(\frac{1}{(m+1)^s}-\frac{1}{(m+2)^s}\right)+\cdots
\end{aligned}
\]
Da nun nach der Voraussetzung $\sigma_{m-1}, \sigma_m, \sigma_{m+1}, \ldots$ zwischen endlichen Grenzen $\pm C$ eingeschlossen sind, so ist hiernach $R_m+\dfrac{\sigma_{m-1}}{m^s}$ zwischen den beiden Grenzen
\[
\pm C\left(\frac{1}{m^s}-\frac{1}{(m+1)^s}+\frac{1}{(m+1)^s}-\frac{1}{(m+2)^s}+\cdots\right) = \pm\frac{C}{m^s}
\]
%CONT>>


% ===== scan leaf 718, printed 695 =====
% !PRINTED 695
\weberpage{695}
%<<CONT
eingeschlossen, und es ist dem absoluten Werthe nach
\[
R_m < \frac{2\,C}{m^s},
\]
oder, wenn $s > c$ ist,
\[
R_m < \frac{2\,C}{m^c}.
\]

Diese obere Grenze für $R_m$, die von $s$ unabhängig ist, kann aber, wenn $c$ positiv ist, dadurch, dass man $m$ hinlänglich gross annimmt, unter jeden noch so kleinen Werth herabgedrückt werden.

Daraus folgt aber nicht nur die Convergenz, sondern auch die Stetigkeit von $S$. Denn nimmt man $m$ hinlänglich gross, so wird nicht nur $R_m$, sondern es werden auch die Schwankung von $R_m$ bei veränderlichem $s$ unendlich klein, und $S_m$ ist für ein feststehendes $m$ eine stetige Function von $s$. Also ist auch $S$ stetig.
% [sic] The source has a singular subject (die Schwankung) with the plural verb werden.

Es ist dabei noch zu bemerken, dass die Voraussetzung über $\sigma_m$ keineswegs die Convergenz der unendlichen Reihe $\Sigma a_k$ voraussetzt. Es ist nicht einmal erforderlich, dass die $a_k$ reell seien; denn wenn sie imaginär sind, so braucht man nur den Satz auf den reellen und den imaginären Bestandtheil anzuwenden. Endlich ist es auch nicht nothwendig, die $a_1, a_2, a_3, \dots$ als Constanten vorauszusetzen. Alles bleibt gültig, wenn es \gesperrt{stetige Functionen} von $s$ sind.

2. \gesperrt{Bedeutet $\mu_1, \mu_2, \dots, \mu_n, \dots$ eine unendliche Menge positiver Zahlen von der Beschaffenheit, dass zwei endliche Zahlen $\alpha, \beta$ so angegeben werden können, dass für jedes noch so grosse $n$}
\weq{3}{\alpha < \dfrac{n}{\mu_n} < \beta,}
\gesperrt{so ist die unendliche Reihe}
\[
S = \frac{1}{\mu_1^s} + \frac{1}{\mu_2^s} + \frac{1}{\mu_3^s} + \cdots
\]
\gesperrt{convergent, so lange der Exponent $s$ grösser als 1 ist, und das Product $(s-1)\,S$ bleibt bei hinlänglicher Annäherung von $s$ an den Werth 1 gleichfalls zwischen den Grenzen $\alpha$ und $\beta$ oder nähert sich wenigstens einem dieser Werthe unbegrenzt an.}


% ===== scan leaf 719, printed 696 =====
% !PRINTED 696
\weberpage{696}
Der Beweis ist durch Zurückführung auf ein Integral sehr einfach zu führen. Offenbar ist nämlich
\[
\frac{1}{(n+1)^s} < \int\limits_n^{n+1} \frac{dx}{x^s} < \frac{1}{n^s},
\]
weil sich der erste oder der zweite Werth, $(n+1)^{-s}$ oder $n^{-s}$, ergiebt, wenn unter dem Integralzeichen für $x^{-s}$ der zu kleine oder der zu grosse Werth $(n+1)^{-s}$ oder $n^{-s}$ gesetzt wird.

Demnach ergiebt sich aus (3) für jedes positive $s$
\[
\alpha^s \int\limits_n^{n+1} \frac{dx}{x^s} < \frac{1}{\mu_n^s} < \beta^s \int\limits_{n-1}^n \frac{dx}{x^s}.
\]
Setzen wir
\[
R_m = \frac{1}{\mu_{m+1}^s} + \frac{1}{\mu_{m+2}^s} + \frac{1}{\mu_{m+3}^s} + \cdots,
\]
so ergiebt sich hieraus
\[
\alpha^s \int\limits_{m+1}^{\infty} \frac{dx}{x^s} < R_m < \beta^s \int\limits_m^{\infty} \frac{dx}{x^s}
\]
oder
\weq{4}{\dfrac{\alpha^s}{(s-1)(m+1)^{s-1}} < R_m < \dfrac{\beta^s}{(s-1)\,m^{s-1}}.}

Es bleibt also $R_m$, wie viele Glieder auch summirt werden mögen, wenn $s > 1$ ist, endlich, und dies ist bei Reihen mit positiven Gliedern eine ausreichende Bedingung für die Convergenz.

Setzt man ferner (4) in die Form
\[
\frac{\alpha^s}{(m+1)^{s-1}} < (s-1)\,R_m < \frac{\beta^s}{m^{s-1}},
\]
so sieht man, dass die beiden Grenzen beliebig nahe an $\alpha, \beta$ gebracht werden können, wenn man $s-1$ klein genug annimmt, und daraus ergiebt sich, dass $(s-1)\,R_m$ bei hinlänglich kleinem $s-1$ nicht mehr ausserhalb der Grenzen $\alpha, \beta$ liegen kann, wenn es sich nicht mit abnehmendem $s-1$ einem dieser Werthe unbegrenzt annähert. Da sich nun $R_m$ von $S$ nur durch einen Bestandtheil unterscheidet, der für jedes positive $s$ endlich ist, so ist der Satz hiermit bewiesen.

Es gilt der Satz aber auch dann noch, wenn die Ungleichheit (3) nicht für alle $n$ gilt, sondern wenn sie nur besteht, sobald $n$ eine beliebig gegebene Grenze $m$ überschritten hat. Auf
%CONT>>


% ===== scan leaf 720, printed 697 =====
% !PRINTED 697
\weberpage{697}
%<<CONT
Grund dieser Bemerkung kann man, wenn sich $n : \mu_n$ einem endlichen Grenzwerthe $\gamma$ nähert, $\alpha$ und $\beta$ diesem Grenzwerthe beliebig nahe annehmen, und erhält so die etwas präcisere Fassung des Theorems:

3. \gesperrt{Sind die positiven Zahlen $\mu_1, \mu_2, \mu_3, \dots$ so beschaffen, dass}
\[
\mathop{\mathrm{Lim}}\limits_{n=\infty} \frac{n}{\mu_n} = \gamma
\]
\gesperrt{ein endlicher Grenzwerth ist, so ist die unendliche Reihe}
\[
S = \frac{1}{\mu_1^s} + \frac{1}{\mu_2^s} + \frac{1}{\mu_3^s} + \cdots
\]
\gesperrt{für jedes $s$, was grösser als 1 ist, convergent, und es ist}
\[
\mathop{\mathrm{Lim}}\limits_{s=1} (s-1)\,S = \gamma.
\]

Es sei jetzt irgend ein System $M$ von \gesperrt{unendlich vielen} positiven Zahlen $\mu_n$ gegeben, die wir so ordnen
\weq{5}{\mu_1 \leqq \mu_2 \leqq \mu_3 \leqq \dots \qquad (M),}
so dass in der Reihe (5) niemals ein kleineres Glied auf ein grösseres folgt, und es bedeute $T$ die Anzahl von Gliedern in $M$, die nicht grösser als eine beliebig gegebene positive Grösse $t$ sind. Es gilt dann folgender Satz:

4. \gesperrt{Wenn einer der beiden Grenzwerthe}
\weq{6}{\mathop{\mathrm{Lim}}\limits_{n=\infty} \frac{n}{\mu_n}, \quad \mathop{\mathrm{Lim}}\limits_{t=\infty} \frac{T}{t}}
\gesperrt{endlich ist, so hat der andere denselben endlichen Werth.}

Es sei zunächst
\weq{7}{\mathop{\mathrm{Lim}}\limits_{n=\infty} \frac{n}{\mu_n} = \gamma}
ein endlicher Grenzwerth. Dann werden die $\mu_n$ mit unendlich wachsendem $n$ nothwendig ins Unendliche wachsen müssen, und es giebt für jedes positive $t$ einen Werth $m$, so dass
\weq{8}{\mu_m \leqq t < \mu_{m+1};}
$m$ wächst zugleich mit $t$ ins Unendliche und es ist $T = m$ [wegen (5)]. Daher nach (8)
\weq{9}{\frac{m}{\mu_m} \geqq \frac{T}{t} > \frac{m}{\mu_{m+1}}.}


% ===== scan leaf 721, printed 698 =====
% !PRINTED 698
\weberpage{698}
Wenn nun $\gamma = 0$ ist, so folgt hieraus, indem man $t$ und $m$ ins Unendliche wachsen lässt, dass auch $T : t$ den Grenzwerth 0 hat. Ist aber $\gamma$ von 0 verschieden, so ist nach (7)
\[
\operatorname{Lim} \frac{\mu_m}{\mu_{m+1}} = \operatorname{Lim} \frac{m}{m+1} = 1,
\]
und daher nähern sich die beiden Grenzwerthe in (9) dem Werthe $\gamma$, und es folgt also:
\weq{10}{\mathop{\mathrm{Lim}}\limits_{t=\infty} \frac{T}{t} = \gamma.}

Setzen wir zweitens umgekehrt die Grenzgleichung (10) voraus, so wird auch jetzt $\mu_n$ mit $n$ ins Unendliche wachsen müssen, denn sonst würde, da die Gesammtzahl aller $\mu$ unendlich ist, $T$ schon für ein endliches $t$ unendlich, und $\gamma$ könnte nicht endlich sein.

Nehmen wir nun einen Werth $\mu_n$, der in der Reihe der unter einander gleichen
\[
\mu_{m+1}, \; \mu_{m+2}, \; \dots, \; \mu_{m+l}
\]
vorkommt, so dass
\[
m + 1 \leqq n \leqq m + l,
\]
so wird $T$, wenn $t$ durch den Werth $\mu_n$ geht, plötzlich um $l$ Einheiten wachsen, und das Verhältniss $T : t$ wächst um $l : \mu_n$. Da aber $T : t$ einen endlichen Grenzwerth haben soll, so muss
\weq{11}{\operatorname{Lim} \frac{l}{\mu_n} = 0}
sein. Wenn nun $t = \mu_n$ ist, so ist $T = m + l$, und folglich ist
\weq{12}{\frac{T}{t} = \frac{m+l}{\mu_n}, \quad \operatorname{Lim} \frac{m+l}{\mu_n} = \gamma.}

Ferner
\[
\frac{m}{\mu_n} < \frac{n}{\mu_n} \leqq \frac{m+l}{\mu_n},
\]
und folglich ist wegen (11) und (12)
\[
\operatorname{Lim} \frac{n}{\mu_n} = \operatorname{Lim} \frac{m}{\mu_n} = \gamma.
\]
Also ist aus der Gleichung (10) die Gleichung (7) gefolgert \textsuperscript{1}).

Mit Benutzung des Satzes 4. können wir nun dem Satze 3. auch die Form geben:

\begingroup
\footnotesize
\textsuperscript{1}) Dieser Satz ist im Wesentlichen auf dieselbe Weise bewiesen bei \gesperrt{Dirichlet-Dedekind}, Supplement II. Vgl. auch \gesperrt{Riemann's} mathematische Werke, 2. Auflage, Nr. XXX.
\endgroup


% ===== scan leaf 722, printed 699 =====
% !PRINTED 699
\weberpage{699}
5. \gesperrt{Ist $T$ die Anzahl der Grössen $\mu_n$, die nicht grösser als $t$ sind, so ist}
\weq{13}{\mathop{\mathrm{Lim}}\limits_{s=1}\sum_{1,\infty}^{n}\frac{s-1}{\mu_n^{s-1}} = \mathop{\mathrm{Lim}}\limits_{t=\infty}\frac{T}{t},}
% [sic] The source prints exponent s-1 in the denominator. Theorem 3 and the application in section 192,(1) require exponent s; for mu_n=n the printed series diverges for 1<s<=2. Printed bytes are retained.
\gesperrt{wenn der Grenzwerth von $T:t$ endlich ist.}

Nehmen wir an, dass die Grössen $\mu_1, \mu_2, \mu_3, \ldots$ nicht nur den Bedingungen des Satzes 3., sondern noch den weiteren genügen, dass, wenn
\weq{14}{\mu_n = \frac{n+c_n}{\gamma}}
gesetzt wird, die $c_n$ mit unendlich wachsendem $n$ nicht unendlich werden (d. h. für jedes noch so grosse $n$ zwischen endlichen Grenzen eingeschlossen sind), so können wir den Satz 3. durch den noch schärferen ersetzen:

6. \gesperrt{Haben die Zahlen $\mu_n$ die Eigenschaft, dass sich ein endliches $\gamma$ so bestimmen lässt, dass}
\weq{15}{\gamma \mu_n - n = c_n}
\gesperrt{mit unendlich wachsendem $n$ nicht unendlich wird, so hat die für jedes $s>1$ definirte Function von $s$}
\weq{16}{S = \frac{1}{\mu_1^s} + \frac{1}{\mu_2^s} + \cdots}
\gesperrt{die Eigenschaft, dass die Differenz}
\weq{17}{S - \frac{\gamma}{s-1} = C_s}
\gesperrt{sich mit unendlich abnehmendem $s-1$ einer endlichen Grenze nähert.}

Um dies zu beweisen, setzen wir nach (14) und (16):
\weq{18}{\sum^{n} \frac{\gamma^s}{n^s} - S = \gamma^s \sum^{n} \frac{1}{n^s}\left[1-\left(1+\frac{c_n}{n}\right)^{-s}\right].}

Wenn nun $\varepsilon$ ein positiver echter Bruch ist, so ist, wenn wir
\[
a_n = \frac{1}{n^\varepsilon}\left[1-\left(1+\frac{c_n}{n}\right)^{-s}\right]
\]
setzen, $a_n n^{1+\varepsilon}$ für ein unendlich wachsendes $n$ nicht unendlich, und folglich ist nach einem bekannten elementaren Satze der Reihenlehre $a_1 + a_2 + \cdots$ eine unbedingt convergente Reihe.


% ===== scan leaf 723, printed 700 =====
% !PRINTED 700
\weberpage{700}
Setzen wir nun
\[
s = s_1 + \varepsilon,
\]
so ergiebt sich aus (18):
\[
\sum \frac{\gamma^s}{n^s} - S = \gamma^s \sum \frac{a_s}{n^{s_1}},
\]
% [sic] The source prints a_s; a_n is required by the definition on printed p. 699 and by the summation index.
und dies ist nach 1. eine für positive $s_1$ endliche und stetige Function von $s_1$. Für $s_1 = 1-\varepsilon$ wird aber $s=1$, und folglich ist
\weq{19}{\sum \frac{\gamma^s}{n^s} - S = D_s}
für $s=1$ endlich, nämlich gleich
\[
D_1 = \sum \frac{\gamma c_n}{n(n+c_n)}.
\]
Mit Anwendung einfacher Sätze aus der Theorie der $\Gamma$-Functionen lässt sich aber die Summe $\sum \dfrac{1}{n^s}$ durch ein bestimmtes Integral ausdrücken. Es ist
\[
\frac{1}{n^s} = \frac{1}{\Gamma(s)}\int_0^{\infty} e^{-nx} x^{s-1}\, dx,
\]
und folglich
\[
\sum \frac{1}{n^s} = \frac{1}{\Gamma(s)}\int_0^{\infty} \frac{x^{s-1}e^{-x}\,dx}{1-e^{-x}}.
\]
Ferner ist [nach dem Satze $(s-1)\Gamma(s-1) = \Gamma(s)$]:
\[
\frac{1}{s-1} = \frac{1}{\Gamma(s)}\int_0^{\infty} e^{-x} x^{s-2}\, dx,
\]
und daraus
\[
\sum \frac{1}{n^s} - \frac{1}{s-1} = \frac{1}{\Gamma(s)}\int_0^{\infty} x^{s-1}e^{-x}\left(\frac{1}{1-e^{-x}} - \frac{1}{x}\right)dx.
\]
Hierin ist die rechte Seite eine für alle positiven Werthe von $s$ stetige Function, die für $s=1$ in die \gesperrt{Euler}'sche Constante
\[
\int_0^{\infty} e^{-x}\left(\frac{1}{1-e^{-x}} - \frac{1}{x}\right)dx = -\Gamma'(1) = 0{,}57721566\ldots
\]
übergeht. Nun ist nach (19):
\[
D_s = \gamma^s\left(\sum \frac{1}{n^s} - \frac{1}{s-1}\right) - \left(S - \frac{\gamma_s}{s-1}\right),
\]
% [sic] The source has a subscript in the final gamma_s, unlike the leading gamma^s. Equation (19) requires gamma^s in both places; the printed subscript is preserved.
%CONT>>


% ===== scan leaf 724, printed 701 =====
% !PRINTED 701
\weberpage{701}
\pdfbookmark[2]{§ 192. Anwendung auf die Bestimmung der Classenzahl}{bm219}
%<<CONT
und daher
\weq{20}{\mathop{\mathrm{Lim}}\limits_{s=1}\left(S - \frac{\gamma_s}{s-1}\right) = -\gamma\,\Gamma'(1) + D_1 = C_1,}
% [sic] The source prints gamma_s, +D_1, and =C_1. With the intended gamma^s, (19) gives the limit -gamma*Gamma'(1)-D_1. Definition (17), which subtracts gamma/(s-1), instead gives C_1=-gamma*Gamma'(1)+gamma*log(gamma)-D_1. The printed subscript, sign, and identification are retained and distinguished here.
wie zu beweisen war, endlich\footnote{Die hier gebrauchten Sätze über $\Gamma$-Functionen finden sich in den ausführlicheren Lehrbüchern der Integralrechnung, z.\,B. \gesperrt{Serret-Harnack}, Lehrbuch der Differential- und Integralrechnung, Bd.\ II, S.\,170\,f. (Leipzig 1885).}.

\begin{center}
\S.\,192.\\[4pt]
\gesperrt{Anwendung auf die Bestimmung der Classenzahl.}
\end{center}

Wir machen von diesen Sätzen jetzt die Anwendung auf die Theorie des Körpers $\Omega$. Wir verstehen unter den Zahlen $\mu_n$ des Theorems 5. die Normen der sämmtlichen Ideale des Körpers $\Omega$, also unter $T$ die Anzahl der Ideale in $\Omega$, deren Norm nicht grösser als eine gegebene positive Grösse $t$ ist, und erhalten
\weq{1}{\mathop{\mathrm{Lim}}\limits_{s=1}\sum \frac{s-1}{N(\mathfrak{a})^s} = \mathop{\mathrm{Lim}}\limits_{t=\infty}\frac{T}{t},}
und darin erstreckt sich die Summe der linken Seite auf \gesperrt{alle} Ideale $\mathfrak{a}$ des Körpers. Die Ideale zerfallen nun nach \S.\,153 in eine endliche Anzahl von Classen
\[
A_1, A_2, \ldots, A_h,
\]
und das grosse Ziel ist die Bestimmung dieser Zahl $h$, der \gesperrt{Classenzahl}.

Die Ideale $\mathfrak{a}_1$ einer dieser Classen $A_1$ sind dadurch charakterisirt, dass ihre Producte mit einem und demselben ganzen Functionale $\varphi_1$, das der Classe $A_1^{-1}$ angehört, \gesperrt{Hauptideale} sind, dass also
\[
\varphi_1 \mathfrak{a}_1 = \alpha
\]
eine ganze Zahl des Körpers $\Omega$ ist.

Ist die Norm von $\mathfrak{a}_1$ nicht grösser als $t$, so ist
\[
N_a(\alpha) \leqq N_a(\varphi_1)\, t = t_1.
\]
Ist $T_1$ die Anzahl der Ideale der Classe $A_1$, deren Normen nicht grösser als $t$ sind, so ist $T_1$ zugleich die Anzahl der nicht
%CONT>>


% ===== scan leaf 725, printed 702 =====
% !PRINTED 702
\weberpage{702}
%<<CONT
associirten durch $\varphi_1$ theilbaren ganzen Zahlen, deren Normen nicht grösser als $t_1$ sind, und nach dem Satze \S.\,190,\,1. ist
\weq{2}{\mathop{\mathrm{Lim}}\limits_{t_1=\infty}\frac{T_1}{t_1} = \frac{g}{N_a(\varphi_1)},}
wenn $g$ die an der erwähnten Stelle angegebene Bedeutung hat, also eine \gesperrt{von Null verschiedene}, durch die Natur des Körpers $\Omega$ bestimmte Zahl ist. Wenn jetzt $T_2, t_2, \ldots, T_h, t_h$ die entsprechende Bedeutung für die Classen $A_2, \ldots, A_h$ haben, wie $T_1, t_1$ für $A_1$, so ist
\[
T = T_1 + T_2 + \cdots + T_h,
\]
\[
\frac{T}{t} = \frac{T_1}{t_1}\, N_a(\varphi_1) + \frac{T_2}{t_2}\, N_a(\varphi_2) + \cdots + \frac{T_h}{t_h}\, N_a(\varphi_h),
\]
und aus (1) und (2) ergiebt sich die fundamentale Formel:
\weq{3}{\mathop{\mathrm{Lim}}\limits_{s=1}\sum \frac{s-1}{N(\mathfrak{a})^s} = gh.}
% [sic] The following source citation is section 185,(11); the formula defining g is section 190,(11) on printed p. 693.
Die Zahl $g$ können wir nach \S.\,185,\,(11) als bekannt betrachten, wenn auch ihre wirkliche Berechnung noch auf der Voraussetzung beruht, dass ein Fundamentalsystem von Einheiten bekannt sei, und wenn auch die Ermittelung eines solchen Systemes in höheren Körpern immer als eine der grössten Schwierigkeiten betrachtet worden ist. Dann hängt die Berechnung der Classenzahl noch von der Bestimmung des Grenzwerthes auf der linken Seite von (3) ab, die natürlich auch nur in besonderen Fällen gelingt, doch aber oft wichtige Schlüsse über die Natur der Classenzahl gestattet. Wir wollen noch eine die Berechnung vorbereitende Umformung der Summe
\weq{4}{\sum \frac{1}{N(\mathfrak{a})^s} = \Phi(s)}
in ein unendliches Product entwickeln.

Es sei $\mathfrak{p}$ irgend ein Primideal $f^{\mathrm{ten}}$ Grades im Körper $\Omega$, also
\[
N(\mathfrak{p}) = p^f.
\]
Dann ist die unendliche geometrische Reihe
\[
1 + \frac{1}{N(\mathfrak{p})^s} + \frac{1}{N(\mathfrak{p})^{2s}} + \cdots = \frac{1}{1-p^{-sf}};
\]
wenn man diese Reihen für alle verschiedenen Primideale $\mathfrak{p}$ mit einander multiplicirt, so erhält man nach bekannten Sätzen aus
%CONT>>


% ===== scan leaf 726, printed 703 =====
% !PRINTED 703
\weberpage{703}
%<<CONT
der Lehre von den unendlichen Reihen eine Summe von Gliedern
der Form
\[
\left(\frac{1}{N(\mathfrak{p})^k}\,\frac{1}{N(\mathfrak{p}')^{k'}}\,\frac{1}{N(\mathfrak{p}'')^{k''}}\cdots\right)^s = \left(\frac{1}{N(\mathfrak{p}^k\mathfrak{p}'^{k'}\mathfrak{p}''^{k''}\dots)}\right)^s,
\]
die alle in der Form $N(\mathfrak{a})^{-s}$ enthalten sind, und jedes solche
Glied ergiebt sich ein- und nur einmal. Danach ist also
\weq{5}{\varPhi(s) = \prod \frac{1}{1-N(\mathfrak{p})^{-s}}.}

Sind daher $\mathfrak{p}_1, \mathfrak{p}_2, \dots, \mathfrak{p}_e$ die von einander verschiedenen
Primfactoren von $p$, und $f_1, f_2, \dots, f_e$ ihre Grade (\S.\,143), so
ergiebt sich
\weq{6}{\varPhi(s) = \overset{p}{\prod} \frac{1}{(1-p^{-sf_1})(1-p^{-sf_2})\dots(1-p^{-sf_e})},}
worin das unendliche Product $\prod$ über alle natürlichen Primzahlen
 $p$ auszudehnen ist, und nach (3), (4)
\weq{7}{gh = \mathop{\mathrm{Lim}}\limits_{s=1}\,(s-1)\varPhi(s).}

Hieraus lässt sich ein für mannigfache Anwendungen wichtiger
Schluss ziehen:

Wenn wir aus der Entwickelung
\[
1 + \frac{1}{p^s} + \frac{1}{p^{2s}} + \cdots = \frac{1}{1-p^{-s}}
\]
das Product für alle Primzahlen $p$ bilden, so ergiebt sich:
\weq{8}{\overset{p}{\prod}\frac{1}{1-p^{-s}} = \overset{n}{\sum}\frac{1}{n^s},}
worin sich die Summe auf der rechten Seite auf alle natürlichen
Zahlen $n$ erstreckt, und es hat also dies Product für alle Werthe
von $s$, die grösser als 1 sind, einen endlichen Werth. Der reciproke
Werth, nämlich das Product
\weq{9}{P = \overset{p}{\prod}(1-p^{-s}),}
dessen Factoren sämmtlich kleiner als 1 sind, hat daher, so
lange $s>1$ ist, einen von Null verschiedenen Werth. Diese
Eigenschaft bleibt nun erhalten, wenn wir bei der Bildung des
Productes $P$ nur einen Theil aller Primzahlen $p$ berücksichtigen,
weil das Product durch Weglassen beliebiger Factoren nur
vergrössert wird, ohne die Einheit je zu übersteigen.


% ===== scan leaf 727, printed 704 =====
% !PRINTED 704
\weberpage{704}
Daraus ergiebt sich, dass in dem Producte (6) die Theilproducte
\[
\prod \frac{1}{1-p^{-sf}},
\]
die sich über alle Primzahlen $p$ erstrecken, für die $f>1$ ist, auch
für $s=1$ noch endlich bleiben. Da andererseits $g, h$ bestimmte
von Null verschiedene Werthe haben, so ergiebt sich aus (7) der
wichtige Satz:

I. \gesperrt{Durchläuft $\mathfrak{p}$ die Gesammtheit der Primideale
ersten Grades irgend eines algebraischen Körpers $\Omega$, so hat das Product}
\weq{10}{(s-1)\prod\frac{1}{1-N(\mathfrak{p})^{-s}}}
\gesperrt{für $s=1$ einen endlichen von Null verschiedenen
Grenzwerth.}

Diese Eigenschaft bleibt auch dann noch erhalten, wenn bei
der Bildung des Productes eine beliebige \gesperrt{endliche} Anzahl von
Primidealen ersten Grades ausgelassen wird.

Die Normen $N(\mathfrak{p})$ sind in der Formel (10) gleich natürlichen
Primzahlen $p$. Eine Primzahl $p$ kommt aber darin so oft
vor, als sie Primfactoren ersten Grades in $\Omega$ enthält, also höchstens
$n$ mal. Ist $\Omega$ ein \gesperrt{Normalkörper}, so kommt auch wirklich
jede Primzahl $p$, die in Primfactoren ersten Grades zerlegbar
ist, genau $n$ mal darin vor, und wir können für das Product (10)
auch setzen:
\[
(s-1)\prod\frac{1}{(1-p^{-s})^n},
\]
worin sich das Product $\prod$ auf alle in \gesperrt{Primfactoren ersten
Grades zerlegbaren} Primzahlen $p$ erstreckt (\S.\,160).
% [source note] The assertion of exactly n distinct degree-one prime factors requires an unramified prime. For Q(i), 2 ramifies as a unit times (1+i)^2 and has only one distinct degree-one prime. Omitting the finitely many ramified Euler factors preserves the finite nonzero limiting-value conclusion, but the displayed products are not literally identical without that qualification.

Eine unmittelbare Folgerung des Satzes I. ist die:

\gesperrt{In jedem algebraischen Körper giebt es unendlich
viele Primideale ersten Grades.}

\begin{center}
\rule{6em}{0.4pt}
\end{center}


% ===== scan leaf 728, printed 705 =====
\weberpage{705}% page number not printed on this page
\pdfbookmark[1]{Vierundzwanzigster Abschnitt. Classenzahl der Kreistheilungskörper.}{bm220}
\pdfbookmark[2]{§ 193. Classenzahldarstellung im Kreistheilungskörper Ωm}{bm221}
% !PRINTED none
\weberpage{none}
% [source note: this Abschnitt opener has no printed folio; sequence position corresponds to inferred page 705.]

\begin{center}
\gesperrt{Vierundzwanzigster Abschnitt.}

\gesperrt{Classenzahl der Kreistheilungskörper.}

\rule{4em}{0.4pt}
\end{center}

\begin{center}
\S.\,193.

\gesperrt{Classenzahldarstellung im Kreistheilungskörper $\Omega_m$.}
\end{center}

Die allgemeinen Resultate über die Classenzahl $h$ sollen
nun angewandt werden auf den vollen Kreistheilungskörper $\Omega_m$
unter der bisherigen Voraussetzung, dass $m$ eine Potenz einer
Primzahl $q$ sei, und dass also
\weq{1}{m = q^{\varkappa}, \quad \varphi(m) = q^{\varkappa-1}(q-1) = \mu}
gesetzt sei.

Im einundzwanzigsten Abschnitte (\S.\,169, 171) haben wir
gesehen, dass in $q$ nur ein Primideal $1^{\text{sten}}$ Grades aufgeht, und
dass eine zum Exponenten $f$ \gesperrt{gehörige}, von $q$ verschiedene
Primzahl $p$ in $e$ verschiedene Primideale $f^{\text{ten}}$ Grades zerfällt.
Darin bedeutet $f$ den kleinsten positiven Exponenten, für den
\weq{2}{p^f \equiv 1\ (\mathrm{mod}\ m)}
ist, und $e$ ist durch die Gleichung
\weq{3}{\mu = ef}
bestimmt.

Die am Ende des vorigen Paragraphen durch die Formel (6)
bestimmte Function $\varPhi(s)$ erhält daher hier den Ausdruck:
\weq{4}{\varPhi(s) = \frac{1}{1-q^{-s}}\prod\frac{1}{(1-p^{-sf})^e},}
und darin erstreckt sich das Productzeichen $\varPi$ auf alle Primzahlen
 $p$, die von $q$ verschieden sind. Darin ist $s$ eine Variable,
die immer grösser als 1 ist, die sich schliesslich der Grenze 1
nähern soll.


% ===== scan leaf 729, printed 706 =====
% !PRINTED 706
\weberpage{706}
Um nun diesen Ausdruck für die Function $\varPhi$ weiter umzuformen
und zur Berechnung vorzubereiten, hat man die Sätze
über die \gesperrt{Abel}'schen Gruppen und Gruppencharaktere zu benutzen,
die wir in den Paragraphen 14 bis 16 dieses Bandes
kennen gelernt haben.

Die Gesammtheit der durch $q$ nicht theilbaren Zahlen $n$
bildet, nach dem Modul $m$ genommen, eine \gesperrt{Abel}'sche Gruppe $\mathfrak{N}$
vom Grade $\mu$, deren Charaktere $\chi(n)$ im \S.\,16 bestimmt sind. Es
war dabei ein kleiner Unterschied, je nachdem $q$ ungerade oder
$q=2$ ist.

1) Ist $q$ ungerade, so giebt es eine primitive Wurzel $c$ von $m$,
und für jede Zahl $n$ giebt es einen Index $\gamma$, so dass
\weq{5}{n \equiv c^{\gamma}\ (\mathrm{mod}\ m).}
Ist dann $\Theta$ eine $\mu^{\text{te}}$ Einheitswurzel, so ist
\weq{6}{\chi(n) = \Theta^{\gamma},}
und die sämmtlichen $\mu$ Charaktere erhält man, wenn man für $\Theta$
die verschiedenen $\mu^{\text{ten}}$ Einheitswurzeln setzt.

Gehört $n$ zum Exponenten $f$, so ist $e$ der grösste gemeinschaftliche
Theiler von $\gamma$ und $\mu$, und $\chi(n)$ ist $f^{\text{te}}$ Einheitswurzel.

2) Für $q=2$, $m>4$ (den Fall $m=4$ berücksichtigen wir
hier nicht) hat jede ungerade Zahl $n$ zwei Indices $\alpha, \beta$, die nach
den Moduln $2, \frac{1}{2}\mu$ bestimmt sind, für die
\weq{7}{n \equiv (-1)^{\alpha}\,5^{\beta}\ (\mathrm{mod}\ m)}
ist. Als Charaktere erhält man, wenn $\varepsilon = \pm 1$ und $\Theta$ gleich
einer $\frac{1}{2}\mu^{\text{ten}}$ Einheitswurzel gesetzt wird:
\weq{8}{\chi(n) = \varepsilon^{\alpha}\Theta^{\beta};}
$f$ ist die kleinste positive Lösung der beiden Congruenzen:
\[
\alpha f \equiv 0\ (\mathrm{mod}\ 2), \quad \beta f \equiv 0\ (\mathrm{mod}\ \tfrac{1}{2}\mu),
\]
und $\chi(n)$ ist auch hier $f^{\text{te}}$ Einheitswurzel.

Für beide Fälle gilt aber die Bemerkung:

3) Gehört $n$ zum Exponenten $f$, so erhält man, wenn man $\chi$
die gesammte Gruppe der $\mu$ Charaktere durchlaufen lässt, aus
$\chi(n)$ jede $f^{\text{te}}$ Einheitswurzel $e$ mal.

Denn da die $\chi$ eine \gesperrt{Abel}'sche Gruppe bilden, so erhält
man zunächst jede $f^{\text{te}}$ Einheitswurzel, die überhaupt vorkommt,
gleich oft, nämlich so oft als den Werth 1, weil, wenn $\chi_1(n)=\chi_2(n)$
ist, $\chi_1\chi_2^{-1}(n) = \chi_0(n) = 1$ ist. Andererseits erhält man aber
%CONT>>


% ===== scan leaf 730, printed 707 =====
% !PRINTED 707
\weberpage{707}
%<<CONT
jede $f^{\text{te}}$ Einheitswurzel; denn wenn man im Falle 1) für $\Theta$ eine
primitive $\mu^{\text{te}}$, im Falle 2) eine primitive $\frac{1}{2}\mu^{\text{te}}$
Einheitswurzel, und im letzten Falle zugleich $\varepsilon = -1$ setzt, so ist
$\chi(n)$ eine primitive $f^{\text{te}}$ Einheitswurzel, und aus deren Potenzen
kann man alle $f^{\text{ten}}$ Einheitswurzeln herleiten.

Bedeutet daher jetzt $x$ eine Variable, so ist, wenn $n$ zum Exponenten $f$
gehört, das über sämmtliche Charaktere $\chi$ erstreckte Product
\weq{9}{\overset{\chi}{\Pi} [1 - \chi(n)\,x] = (1-x^f)^e,}
und wenn man daher $x = p^{-s}$ setzt, so ergiebt sich aus (4):
\weq{10}{\Phi(s) = \frac{1}{1-q^{-s}} \overset{\chi}{\Pi} \overset{p}{\Pi} \frac{1}{1-\chi(p)\,p^{-s}},}
wenn sich das erste Productzeichen auf alle Charaktere $\chi$, das zweite auf
alle Primzahlen $p$ erstreckt. Es ist also $\Phi(s)$ ein Product aus einer
endlichen Zahl, $\mu$, von Factoren, deren jeder ein unendliches Product ist.

Jetzt wollen wir jeden Factor eines dieser unendlichen Producte nach
steigenden Potenzen von $p^{-s}$ entwickeln. Dadurch ergiebt sich
\[
\frac{1}{1-\chi(p)\,p^{-s}} = 1 + \chi(p)\,p^{-s} + \chi(p^2)\,p^{-2s} + \chi(p^3)\,p^{-3s} + \cdots
\]
Das Product aus allen Factoren, die man hieraus erhält, wenn $\chi$
festgehalten wird, während $p$ alle von $q$ verschiedenen Primzahlen
durchläuft, ist ein Aggregat von Gliedern der Form
\[
\chi(p^k)\,\chi(p'^{k'})\,\chi(p''^{k''}) \cdots p^{-sk}\,p'^{-sk'}\,p''^{-sk''} \cdots = \chi(n)\,n^{-s},
\]
und jedes Glied dieser Form kommt in dem Producte ein- und nur einmal vor
[vgl. \S.\,192, (8)]. Danach ergiebt sich die Umformung von $\Phi(s)$ in ein
endliches Product von unendlichen Reihen:
\weq{11}{\Phi(s) = \frac{1}{1-q^{-s}} \overset{\chi}{\Pi} \overset{n}{\Sigma} \frac{\chi(n)}{n^s},}
worin sich die Summe auf alle durch $q$ nicht theilbaren Zahlen $n$ erstreckt
und das Product auf alle $\mu$ Charaktere $\chi$.

Dieser Ausdruck ist geeignet, um den Grenzwerth von $(s-1)\,\Phi(s)$ für
$s=1$ zu bestimmen.

Betrachten wir zunächst die dem Hauptcharakter $\chi = 1$ entsprechende
Summe
\[
\overset{n}{\Sigma} \frac{1}{n^s}.
\]


% ===== scan leaf 731, printed 708 =====
% !PRINTED 708
\weberpage{708}
\pdfbookmark[2]{§ 194. Bestimmung der Summen X}{bm222}
Man erhält die Zahlen $n$, wenn man von der Gesammtheit aller natürlichen
Zahlen $k$ die Vielfachen von $q$, d.\,h. die Gesammtheit der Zahlen $qk$
wegnimmt. Hiernach ist
\[
\Sigma \frac{1}{n^s} = \Sigma \frac{1}{k^s} - \Sigma \frac{1}{q^s k^s} = (1-q^{-s}) \Sigma \frac{1}{k^s}.
\]
Wenn man aber in dem Satze 3., \S.\,191, für die $\mu_n$ die Reihe der
natürlichen Zahlen $k$ setzt, so folgt:
\[
\underset{s=1}{\mathrm{Lim}} (s-1) \Sigma \frac{1}{k^s} = 1,
\]
und wir erhalten daher
\weq{12}{\underset{s=1}{\mathrm{Lim}} \frac{s-1}{1-q^{-s}} \Sigma \frac{1}{n^s} = 1.}
Die anderen Factoren des Productes $\Phi(s)$ aber sind, wenn die
unendlichen Reihen nach steigenden Werthen von $n$ geordnet werden, nach
dem Satze 1., \S.\,191, stetige Functionen von $s$.

Denn die nach steigenden Werthen von $n$ geordnete Summe
\weq{13}{\Sigma \chi(n)}
ist zwar nicht convergent, kann aber doch dem absoluten Werthe nach nicht
über eine endliche Grenze hinausgehen. Denn so oft $n$ ein volles
Restsystem nach dem Modul $m$ durchläuft, kommt zu der Summe (13) ein
Beitrag hinzu, der nach \S.\,11, 6. den Werth 0 hat. Demnach ist, wenn $\chi$
nicht der Hauptcharakter ist,
\weq{14}{\underset{s=1}{\mathrm{Lim}} \overset{n}{\Sigma} \frac{\chi(n)}{n^s} = \overset{n}{\Sigma} \frac{\chi(n)}{n},}
und demnach erhalten wir für die Classenzahl $h$ im Kreistheilungskörper
$\Omega_m$ jetzt den Ausdruck
\weq{15}{g\,h = \overset{\chi}{\Pi} \overset{n}{\Sigma} \frac{\chi(n)}{n},}
in dem sich das Productzeichen $\Pi$ nur noch auf die \gesperrt{vom Hauptcharakter
verschiedenen} Charaktere $\chi$ bezieht.

\begin{center}
\S.\,194.

\gesperrt{Bestimmung der Summen X.}
\end{center}

Die unendlichen Reihen
\weq{1}{X = \overset{n}{\Sigma} \frac{\chi(n)}{n},}
%CONT>>


% ===== scan leaf 732, printed 709 =====
% !PRINTED 709
\weberpage{709}
%<<CONT
von denen hiernach die Bestimmung der Classenzahl noch abhängt, lassen
sich durch endliche Ausdrücke darstellen. Es ist nämlich
\[
\frac{1}{n} = \int_0^1 x^{n-1}\,dx,
\]
und demnach
\[
X = \int_0^1 \Sigma \chi(n)\,x^{n-1}\,dx.
\]

Nun bleibt $\chi(n)$ ungeändert, wenn $n$ um ein Vielfaches von $m$ wächst.
Versteht man aber unter $t$ den kleinsten positiven Rest von $n$ und setzt
\[
n = t + ml,
\]
so ist $\chi(n) = \chi(t)$, und es folgt:
\[
X = \int_0^1 \overset{t}{\Sigma} \chi(t)\,x^{t-1} \underset{0,\,\infty}{\overset{l}{\Sigma}} x^{ml}\,dx.
\]

Setzt man jetzt
\weq{2}{f(x) = \overset{t}{\Sigma} \chi(t)\,x^t,}
so ist $f(x)$ eine ganze Function von $x$ vom Grade $m-1$, und zugleich ist
wegen der Relation $\Sigma \chi(t) = 0$:
\weq{3}{f(0) = 0,\quad f(1) = 0,}
also $f(x)$ durch $x(1-x)$ theilbar. Für $X$ ergiebt sich dann nach der
Summenformel für die geometrische Reihe $\Sigma x^{ml}$:
\weq{4}{X = \int_0^1 \frac{f(x)\,dx}{x(1-x^m)}.}

Um ein solches Integral zu finden, schreibt die Integralrechnung vor, den
rationalen Bruch unter dem Integralzeichen in Partialbrüche zu zerlegen.
Diese Partialbruch-Zerlegung ergiebt aber, wegen (3), wenn
\[
r = e^{\frac{2\pi i}{m}}
\]
gesetzt wird:
\[
\frac{f(x)}{x(1-x^m)} = -\frac{1}{m} \sum\limits_{1,\,m-1}^{s} \frac{f(r^s)}{x-r^s}
\]
[Bd. I, \S.\,29, (11)], und nun haben wir weiter nach bekannten
%CONT>>


% ===== scan leaf 733, printed 710 =====
% !PRINTED 710
\weberpage{710}
%<<CONT
elementaren Sätzen der Integralrechnung, \gesperrt{wenn $\alpha$ irgend einen Winkel
zwischen 0 und $2\pi$ bedeutet,}
\[
\int_0^1 \frac{dx}{x-e^{i\alpha}} = \tfrac{1}{2}\log\left(4\sin^2\frac{\alpha}{2}\right) + i\,\frac{\pi-\alpha}{2}.
\]
Hieraus ergiebt sich nach (4):
\weq{5}{X = -\frac{1}{m}\underset{1,\,m-1}{\overset{s}{\Sigma}} f(r^s)\left[\tfrac{1}{2}\log 4\left(\sin\frac{s\pi}{m}\right)^2 - \frac{i\pi(m-2s)}{2m}\right].}
Hierin ist nach (2)
\weq{6}{f(r^s) = \overset{t}{\Sigma}\chi(t)\,r^{ts},}
und daraus ergiebt sich [wegen (3)]:
\weq{7}{\underset{1,\,m-1}{\overset{s}{\Sigma}} f(r^s) = \underset{0,\,m-1}{\overset{s}{\Sigma}} f(r^s) = 0,}
da
\[
\underset{0,\,m-1}{\overset{s}{\Sigma}} r^{st} = 0
\]
ist. Hiernach vereinfacht sich die Formel (5) noch etwas und ergiebt
\weq{8}{X = -\frac{1}{m}\underset{1,\,m-1}{\overset{s}{\Sigma}} f(r^s)\left[\tfrac{1}{2}\log\left(\sin\frac{s\pi}{m}\right)^2 + \frac{i\pi s}{m}\right],}
wofür man auch, wenn man $s$ durch $m-s$ ersetzt und wieder (7) benutzt,
\weq{9}{X = -\frac{1}{m}\underset{1,\,m-1}{\overset{s}{\Sigma}} f(r^{-s})\left[\tfrac{1}{2}\log\left(\sin\frac{s\pi}{m}\right)^2 - \frac{i\pi s}{m}\right]}
setzen kann.

Aus (6) folgt aber
\[
f(r^{-s}) = \overset{t}{\Sigma}\chi(t)\,r^{-ts},
\]
und die nach $t$ genommene Summe ändert sich nicht, wenn $t$ durch $m-t$
ersetzt wird. Daraus folgt:
\[
f(r^{-s}) = \chi(-1)\,f(r^s).
\]
Der Factor $\chi(-1)$ hat den Werth $+1$ oder $-1$, da sein Quadrat $=+1$
ist. Wir unterscheiden daher jetzt zwei Arten von Charakteren $\chi_1(n)$,
$\chi_2(n)$, so dass
\weq{10}{\chi_1(-1) = 1,\quad \chi_2(-1) = -1}
%CONT>>


% ===== scan leaf 734, printed 711 =====
% !PRINTED 711
\weberpage{711}
\pdfbookmark[2]{§ 195. Ueber die Classenzahl in dem in Ωm enthaltenen reellen Körper}{bm223}
%<<CONT
ist, und danach sind auch die Functionen $f$ in $f_1$ und $f_2$ zu
unterscheiden, so dass

\weq{11}{f_1(r^{-s}) = f_1(r^s), \quad f_2(r^{-s}) = -\,f_2(r^s)}
wird. Ebenso unterscheiden wir die Ausdrücke (8), (9) als $X_1$
und $X_2$, und es ergiebt sich, wenn man beide Ausdrücke addirt,
nach (11):
\weq{12}{\begin{aligned}
X_1 &= -\frac{1}{2m}\sum_{1,\,m-1}^{s} f_1(r^s)\,\log\left(\sin\frac{s\pi}{m}\right)^{2}, \\
X_2 &= -\frac{i\pi}{m^2}\sum_{1,\,m-1}^{s} s\,f_2(r^s).
\end{aligned}}

\begin{center}
\S.\,195.\\[4pt]
\gesperrt{Ueber die Classenzahl in dem in $\Omega_m$ enthaltenen
reellen Körper.}
\end{center}

Unsere allgemeinen Principien können auch angewandt werden, um die
Classenzahlen in den Theilern des Körpers $\Omega_m$ zu bestimmen.

Insbesondere sind die Formeln (6), (7), \S.\,192 anwendbar, wenn
wir die Grade der Primideale in einem solchen Theiler kennen. Die
Frage nach dem Verhältniss der Classenzahlen in den Theilern eines
Körpers $\Omega_m$ zu der Classenzahl in $\Omega_m$ selbst ist von
grossem Interesse und soll hier für den Fall behandelt werden, dass
es sich um den in $\Omega_m$ enthaltenen reellen Körper $H_m$
handelt, den wir ja schon im \S.\,176 f. untersucht haben\footnote{Vgl.
\gesperrt{Kummer}, \glqq Bestimmung der Anzahl u.\,s.\,f.\grqq{}.
\gesperrt{Crelle's Journ.}, Bd.\,40 (1849).}. Es hat sich dort ergeben, dass
$q$ im Körper $H_m$ die $\tfrac12\mu^{\text{te}}$ Potenz einer
Primzahl ersten Grades ist, dass die anderen Primzahlen $p$ in $e_1$
Primfactoren $f_1^{\text{ten}}$ Grades zerfallen, so dass

\weq{1}{f_1 e_1 = \tfrac{1}{2}\mu,}
und dass zwei Arten von Primzahlen $p_1$, $p_2$ zu unterscheiden
sind, so dass bei den Primzahlen erster Art $f_1 = \tfrac12 f$,
$e_1 = e$, bei denen der zweiten Art $f_1 = f$, $e_1 = \tfrac12 e$
ist.

Für den Körper $H_m$ erhalten wir demnach aus \S.\,192, (6), (7)
die Classenzahl $h_1$ nach der Formel:
%CONT>>


% ===== scan leaf 735, printed 712 =====
% !PRINTED 712
\weberpage{712}
%<<CONT
\weq{2}{g_1 h_1 = \mathop{\mathrm{Lim}}\limits_{s=1} (s-1)\,\Phi_1(s),}
\weq{3}{\Phi_1(s) = \frac{1}{1-q^{-s}}\;\prod^{p_1}\frac{1}{(1-p_1^{-1/2\,sf})^{e}}\;\prod^{p_2}\frac{1}{(1-p_2^{-sf})^{1/2\,e}},}
worin sich die beiden Productzeichen auf alle Primzahlen $p_1$ und
$p_2$ beziehen. Mit Benutzung der oben erklärten Zeichen $f_1$,
$e_1$ kann man dafür auch setzen:
\weq{4}{\Phi_1(s) = \frac{1}{1-q^{-s}}\;\prod^{p}\frac{1}{(1-p^{-sf_1})^{e_1}},}
und das Productzeichen auf \gesperrt{alle} Primzahlen $p$ erstrecken.

Die Bedeutung von $g_1$ ergiebt sich aus \S.\,190, (11).

Nun wollen wir aus der Gruppe $C$ der Charaktere $\chi$ einen
Theiler $C_1$ aussondern, der aus allen den Charakteren $\chi_1$
bestehen soll, für die
\weq{5}{\chi_1(-1) = +1}
ist. Zunächst ist klar, dass diese Charaktere $\chi_1$ eine Gruppe
$C_1$ bilden. Diese Gruppe ist aber nicht mit $C$ identisch, weil
es gewiss einen Charakter $\chi_0$ giebt, für den $\chi_0(-1) = -1$
ist. Wir brauchen, um einen solchen zu erhalten, nur in \S.\,193,
(6) für $\Theta$ eine primitive $\mu^{\text{te}}$ Einheitswurzel,
oder in \S.\,193, (8) $\varepsilon = -1$ zu setzen. Dann bildet die
Gesammtheit der Charaktere $\chi_0\chi_1$, die alle nicht zu $C_1$
gehören, eine Nebengruppe $C_2$, deren Elemente mit $\chi_2$
bezeichnet werden mögen. Jeder Charakter $\chi$ ist also dann
entweder in $C_1$ oder in $C_2$ enthalten, denn wenn $\chi$ nicht
in $C_1$ vorkommt, so kommt $\chi_0^{-1}\chi$ darin vor, und
folglich $\chi$ in $C_2$. Es ist daher $C_1$ ein Theiler von $C$
vom Index 2.

Die Charaktere $\chi_1$ kann man auf folgende Art bilden:

1) Ist $m$ ungerade, so setzen wir, wie in \S.\,193,
\[ n \equiv c^{\gamma} \pmod{m}, \]
und lassen in
\weq{6}{\chi_1(n) = \Theta^{\gamma}}
$\Theta$ die Gesammtheit der $\tfrac12\mu^{\text{ten}}$
Einheitswurzeln durchlaufen.

2) Ist $m$ gerade, so setzen wir
\[ n \equiv (-1)^{\alpha}5^{\beta}, \]
und erhalten
\weq{7}{\chi_1(n) = \Theta^{\beta},}
worin $\Theta$ wieder die sämmtlichen $\tfrac12\mu^{\text{ten}}$
Einheitswurzeln durchläuft.


% ===== scan leaf 736, printed 713 =====
% !PRINTED 713
\weberpage{713}
Beides ergiebt sich leicht dadurch, dass für $n = -1$ im ersten
Falle $\gamma = \tfrac12\mu$, im zweiten $\alpha=1$, $\beta=0$ ist,
also beide Male $\chi_1(-1)=1$ ist, und dass jede der Formeln (6),
(7) genau $\tfrac12\mu$ verschiedene Charaktere darstellt.

Worauf es nun wesentlich ankommt, ist Folgendes:

\begin{quote}
\gesperrt{Ist $p$ irgend eine von $q$ verschiedene Primzahl, so
erhält man in der Form $\chi_1(p)$, wenn $\chi_1$ die Gruppe $C_1$
durchläuft, jede $f_1^{\text{te}}$ Einheitswurzel, und jede gleich
oft, also $e_1$ mal.}
\end{quote}

Dass alle $\chi_1(p)$ Einheitswurzeln vom Grade $f_1$ sind, ist
klar, denn es ist immer (vgl. \S.\,177) $p^{f_1}\equiv\pm1\pmod{m}$,
und daher
\[ [\chi_1(p)]^{f_1} = \chi_1(\pm1) = 1. \]
Wenn aber unter den $\chi_1(p)$ eine primitive $f_1^{\text{te}}$
Einheitswurzel vorkommt, so kommen alle anderen $f_1^{\text{ten}}$
Einheitswurzeln auch darunter vor, und gleich oft, nämlich so oft
wie die Einheitswurzel 1.

Dies folgt unmittelbar aus der Gruppennatur der $\chi_1$. Es ist
also nun noch zu zeigen, dass unter den $\chi_1(p)$ \gesperrt{eine primitive}
$f_1^{\text{te}}$ Einheitswurzel vorkommt.

Diese Möglichkeit ergiebt sich aber aus der Bestimmung von $\chi_1$
sehr einfach. Ist zunächst $\Theta$ in (6) bei \gesperrt{ungeradem $m$ eine
primitive Einheitswurzel vom Grade} $\tfrac12\mu$, und ist $\gamma$
der Index von $p$, so ist eine Potenz von $\Theta^{\gamma}$,
$\Theta^{\gamma\lambda}$ dann und nur dann $=1$, wenn
\[ 2\gamma\lambda \equiv 0 \pmod{\mu}. \]
Nun ist $\mu=ef$ und $e$ der grösste gemeinschaftliche Theiler von
$\gamma$ und $\mu$; und daher wird diese Bedingung
$2\lambda\equiv0\pmod{f}$. Es ist aber nach \S.\,177, 3. $f_1=f$
oder $=\tfrac12 f$, je nachdem $f$ ungerade oder gerade ist, und
folglich muss $\lambda\equiv0\pmod{f_1}$ sein. Also ist
$\Theta^{\gamma}$ primitive $f_1^{\text{te}}$ Einheitswurzel.

Ist aber $m$ \gesperrt{gerade}, so nehmen wir in (7) gleichfalls für $\Theta$
eine primitive $\tfrac12\mu^{\text{te}}$ Einheitswurzel und setzen
\[ p \equiv (-1)^{\alpha}5^{\beta} \pmod{m}, \]
so dass $f$ die kleinste positive Zahl ist, die den Congruenzen
\[ \alpha f \equiv 0 \pmod{2}, \quad \beta f \equiv 0 \pmod{\tfrac12\mu} \]
genügt. Es ist nun $\Theta^{\beta\lambda}$ nur dann $=1$, wenn
$\beta\lambda\equiv0\pmod{\tfrac12\mu}$ ist. Wenn nicht $\lambda=1$,
also $\beta=0$ ist, so ist der kleinste Werth, den $\lambda$ haben
kann, gerade, und daher ist dieser kleinste Werth von $\lambda=f$.


% ===== scan leaf 737, printed 714 =====
% !PRINTED 714
\weberpage{714}
Wenn aber $\lambda=1$ ist, so muss $\beta=0$ sein, und es sind
noch zwei Fälle möglich:
\[
\begin{aligned}
&\text{entweder}\quad f=1,\ \alpha=0,\ p\equiv\phantom{-}1\pmod{m},\ \lambda=f\\
&\text{oder}\qquad\quad\, f=2,\ \alpha=1,\ p\equiv -1\pmod{m},\ \lambda=\tfrac12 f.
\end{aligned}
\]
Nur in diesem letzten Falle gehört $p$ zur ersten Art, und es
folgt also, dass auch hier allgemein $\Theta^{\beta}$ eine
primitive $f_1^{\text{te}}$ Einheitswurzel ist. Hiermit ist dann
der Satz bewiesen.

Dieser Satz gestattet uns die Anwendung der Formel \S.\,193, (9),
die hier ergiebt:
\weq{8}{(1-p^{-sf_1})^{e_1} = \prod^{\chi_1}\bigl[1-\chi_1(p)p^{-s}\bigr],}
worin sich das Product auf alle Charaktere $\chi_1$ der Gruppe
$C_1$ erstreckt, und nun lässt sich die Function $\Phi_1(s)$
ebenso wie $\Phi(s)$ im \S.\,193 weiter entwickeln. Man erhält auf
demselben Wege, wie dort, die Formel (11):
\weq{9}{\Phi_1(s) = \frac{1}{1-q^{-s}}\;\prod^{\chi_1}\sum^{n}\frac{\chi_1(n)}{n^s},}
und wenn man zur Grenze $s=1$ übergeht,
\weq{10}{g_1 h_1 = \prod^{\chi_1}\sum^{n}\frac{\chi_1(n)}{n},}
worin aber jetzt das Product auf alle Charaktere der Gruppe $C_1$
\gesperrt{mit Ausnahme des Hauptcharakters} zu erstrecken ist. Dieses Product
ist also ein Theil des Productes, durch welches im \S.\,193 die
Zahl $gh$ ausgedrückt ist.

Die Zahlen $g$, $g_1$ ergeben sich aus dem im \S.\,190, (11)
angegebenen allgemeinen Ausdrucke:
\weq{11}{\frac{2^{\nu}\pi^{n-\nu}L}{w\sqrt{\pm\Delta}}.}
Da, wie im \S.\,178 nachgewiesen ist, die Einheiten in $\Omega_m$
und $H_m$, von den Einheitswurzeln abgesehen, dieselben sind, so
ist auch ein Fundamentalsystem von Einheiten in $H_m$ zugleich ein
Fundamentalsystem in $\Omega_m$.

Bedeutet $\varepsilon_1,\varepsilon_2,\dots,\varepsilon_{\nu-1}$
daher ein Fundamentalsystem reeller Einheiten, und sind
$\varepsilon_{i,1},\varepsilon_{i,2},\dots,\varepsilon_{i,\nu}$ die
conjugirten Einheiten zu $\varepsilon_i$, so sind die conjugirten
Logarithmen im Körper $H_m$:
\[ \log|\varepsilon_{i,1}|,\ \log|\varepsilon_{i,2}|,\ \dots,\ \log|\varepsilon_{i,\nu}|, \]
dagegen im Körper $\Omega_m$, in dem nur conjugirt imaginäre Paare
vorkommen (\S.\,185),
\[ 2\log|\varepsilon_{i,1}|,\ 2\log|\varepsilon_{i,2}|,\ \dots,\ 2\log|\varepsilon_{i,\nu}|. \]


% ===== scan leaf 738, printed 715 =====
% !PRINTED 715
\weberpage{715}

Bezeichnen wir daher mit
\[
(12)\quad E = \pm
\begin{vmatrix}
\log|\varepsilon_{1,1}|, & \log|\varepsilon_{1,2}|, & \ldots, & \log|\varepsilon_{1,\nu-1}| \\
\log|\varepsilon_{2,1}|, & \log|\varepsilon_{2,2}|, & \ldots, & \log|\varepsilon_{2,\nu-1}| \\
\cdots & \cdots & \cdots & \cdots \\
\log|\varepsilon_{\nu-1,1}|, & \log|\varepsilon_{\nu-1,2}|, & \ldots, & \log|\varepsilon_{\nu-1,\nu-1}|
\end{vmatrix}
\]
den Regulator des Körpers $H_m$ (\S.\,187), so ist bei der Anwendung des Ausdruckes (11) im Körper $\Omega_m$
\[
L = 2^{\nu-1} E,
\]
und im Körper $H_m$
\[
L = E
\]
zu setzen. Es ist ferner
\[
\begin{aligned}
\text{in } \Omega_m: & \quad n = \mu, \quad \nu = \tfrac{1}{2}\mu, \\
\text{in } H_m: & \quad n = \tfrac{1}{2}\mu, \quad \nu = \tfrac{1}{2}\mu.
\end{aligned}
\]
Im Körper $H_m$ sind nur die zwei Einheitswurzeln $\pm 1$ enthalten. In $\Omega_m$ sind die mit positivem und negativem Zeichen genommenen Potenzen von $r$, sonst aber keine Einheitswurzeln enthalten (\S.\,175), und bei der Zählung ist noch zu beachten, dass bei geradem $m$ unter den Potenzen von $r$ auch $-1$ enthalten ist, bei ungeradem $m$ nicht. Daher haben wir
\[
\begin{aligned}
\text{in } \Omega_m: & \quad w = 2m, \quad m \text{ ungerade} \\
& \quad = m, \quad m \text{ gerade,} \\
\text{in } H_m: & \quad w = 2.
\end{aligned}
\]
Endlich haben wir noch, wenn $\varDelta$, $\varDelta_1$ die Grundzahlen der Körper $\Omega_m$, $H_m$ sind (\S.\,176):
\[
\begin{aligned}
\pm \varDelta & = q\varDelta_1^2 \quad m \text{ ungerade} \\
& = 4\varDelta_1^2 \quad m \text{ gerade,}
\end{aligned}
\]
und es folgt, wenn wir der Einfachheit wegen $\nu$ für $\tfrac{1}{2}\mu$ setzen:
\[
\begin{aligned}
g & = \frac{2^{2(\nu-1)}\pi^\nu E}{m\sqrt{q}\,\varDelta_1} \quad m \text{ ungerade} \\
g & = \frac{2^{2(\nu-1)}\pi^\nu E}{m\varDelta_1} \quad m \text{ gerade,}
\end{aligned}
\]
\weq{13}{g_1 = \frac{2^{\nu-1}E}{\sqrt{\varDelta_1}},}
woraus
\weq{14}{\begin{aligned}
g & = \frac{2^{\nu-1}\pi^\nu g_1}{m\sqrt{q\varDelta_1}} \quad m \text{ ungerade} \\
& = \frac{2^{\nu-1}\pi^\nu g_1}{m\sqrt{\varDelta_1}} \quad m \text{ gerade.}
\end{aligned}}


% ===== scan leaf 739, printed 716 =====
% !PRINTED 716
\weberpage{716}

Wenn wir die Formel (10) mit der Formel \S.\,193, (15) verbinden und die im \S.\,194 eingeführte Bezeichnung
\[
X_1 = \sum{}^n \frac{\chi_1(n)}{n}, \quad X_2 = \sum{}^n \frac{\chi_2(n)}{n}
\]
gebrauchen, so ergiebt sich
\[
g\,h = g_1 h_1 \prod X_2.
\]
Ersetzt man $g$ durch seinen Werth (14), so folgt für die Classenzahl
\weq{15}{h = k h_1,}
wenn
\weq{16}{
\begin{aligned}
k & = \frac{m\sqrt{q\varDelta_1}}{2^{\nu-1}\pi^\nu} \prod X_2 \quad \text{bei ungeradem } m \\
& = \frac{m\sqrt{\varDelta_1}}{2^{\nu-1}\pi^\nu} \prod X_2 \quad \text{bei geradem } m,
\end{aligned}
}
und worin [nach (10), (13)]:
\weq{17}{h_1 = \frac{\sqrt{\varDelta_1}}{2^{\nu-1}E} \prod X_1}
die Classenzahl des Körpers $H_m$ ist. Im letzteren Ausdrucke erstreckt sich das Product auf alle Charaktere $\chi_1$ der Gruppe $C_1$ \gesperrt{mit Ausnahme des Hauptcharakters}.

Dass $k$ eine ganze Zahl ist, und daher $h$ ein Vielfaches von $h_1$, wird sich bald zeigen. Die Zahlen $k$ und $h_1$ heissen der \gesperrt{erste und zweite Factor der Classenzahl}.

In dem Falle, dass
\[
m = 2^\varkappa
\]
eine Potenz von 2 ist, haben wir nach \S.\,176, (7):
\[
\varDelta_1 = 2^{(\varkappa-1)2^{\varkappa-2}-1},
\]
und wenn wir diesen Werth einsetzen, so können wir für die vorstehenden Formeln in diesem Falle schreiben:
\weq{18}{
\begin{aligned}
k & = \frac{\sqrt{2}\,2^{2^{\varkappa-3}(\varkappa-3)}\,2^\varkappa}{\pi^\nu} \prod X_2 \\
h_1 & = \frac{\sqrt{2}\,2^{2^{\varkappa-3}(\varkappa-3)}}{E} \prod X_1.
\end{aligned}
}


% ===== scan leaf 740, printed 717 =====
% !PRINTED 717
\weberpage{717}
\pdfbookmark[2]{§ 196. Classenzahl im Körper der achten Einheitswurzeln}{bm224}

\begin{center}
\S.\,196.

\gesperrt{Classenzahl im Körper der achten Einheitswurzeln.}
\end{center}

Wir wollen, einerseits zur Veranschaulichung der bisherigen Resultate, andererseits um für die später anzuwendende vollständige Induction eine Basis zu gewinnen, den Fall $m = 8$, $\mu = 4$, $\nu = 2$ zunächst behandeln.

Hier haben wir ausser dem Hauptcharakter nur drei Charaktere, nämlich, wenn
\[
n \equiv (-1)^\alpha 5^\beta \pmod{8}
\]
ist,
\[
\begin{aligned}
\chi_1(n) & = (-1)^\beta \\
\chi_2(n) & = (-1)^\alpha \\
\chi_3(n) & = (-1)^{\alpha+\beta}.
\end{aligned}
\]
Für $n = -1$ ist $\alpha = 1$, $\beta = 0$, und folglich gehört $\chi_1$ zur ersten, $\chi_2$, $\chi_3$ zur zweiten Art [\S.\,195, (5)], und es ergeben sich für die vier Zahlen $n = 1, 3, 5, 7$ folgende Werthe dieser Charaktere:
\[
\begin{aligned}
\chi_1: & \quad +1, \; -1, \; -1, \; +1 \\
\chi_2: & \quad +1, \; -1, \; +1, \; -1 \\
\chi_3: & \quad +1, \; +1, \; -1, \; -1.
\end{aligned}
\]
Daraus erhält man die drei Functionen $f_1, f_2, f_3$ [\S.\,194, (2)]:
\[
\begin{aligned}
f_1(x) & = x - x^3 - x^5 + x^7 = x(1-x^2)(1-x^4) \\
f_2(x) & = x - x^3 + x^5 - x^7 = x(1-x^2)(1+x^4) \\
f_3(x) & = x + x^3 - x^5 - x^7 = x(1+x^2)(1-x^4).
\end{aligned}
\]
Für $r$ können wir setzen:
\[
r = \frac{1+i}{\sqrt{2}}, \quad r^2 = i, \quad r^3 = -\frac{1-i}{\sqrt{2}}, \quad r^4 = -1,
\]
und danach ergiebt sich:
\[
\begin{aligned}
f_1(r) & = 2\sqrt{2}, & f_2(r) & = 0, & f_3(r) & = 2\sqrt{2}\,i, \\
f_1(r^2) & = 0, & f_2(r^2) & = 4i, & f_3(r^2) & = 0, \\
f_1(r^3) & = -2\sqrt{2}, & f_2(r^3) & = 0, & f_3(r^3) & = 2\sqrt{2}\,i.
\end{aligned}
\]
Für $x = r^4$ verschwinden alle drei Functionen $f(x)$, und für $x = r^5, r^6, r^7$ erhält $f_1(x)$ dieselben, $f_2(x)$ und $f_3(x)$ die entgegengesetzten Werthe, wie für $x = r^3, r^2, r$.


% ===== scan leaf 741, printed 718 =====
% !PRINTED 718
\weberpage{718}

Demnach ergiebt sich aus \S.\,194, (12) mit Rücksicht auf die trigonometrische Formel
\[
\sin\frac{3\pi}{8} = \cos\frac{\pi}{8}
\]
\[
X_1 = -\frac{1}{\sqrt{2}} \log \operatorname{tg}\frac{\pi}{8}, \quad X_2 = \frac{-\pi}{4}, \quad X_3 = \frac{-\pi}{2\sqrt{2}}.
\]
Nun ist
\[
\tau = \frac{r-1}{r^2(r+1)} = \operatorname{tg}\frac{\pi}{8} = \sqrt{2} - 1
\]
eine Einheit des Körpers $H_8$, der hier nichts Anderes ist, als der aus $\sqrt{2}$ entspringende quadratische Körper, und die Einheiten darin erhält man also aus den Lösungen der \gesperrt{Pell}'schen Gleichung
\[
u^2 - 2v^2 = \pm 1
\]
in der Form $u + v\sqrt{2}$.

Die kleinste positive Lösung dieser \gesperrt{Pell}'schen Gleichung ist aber $u = 1$, $v = 1$, und folglich findet man nach Bd.\,I, \S.\,128 alle Einheiten in $H_8$ in der Form $\pm(1+\sqrt{2})^a$, worin $a$ ein positiver oder negativer ganzzahliger Exponent ist; $\tau$ ergiebt sich hieraus für $a = -1$, und daher sind alle Einheiten auch in der Form
\[
\pm \tau^a
\]
enthalten. Es ist also $\tau$ eine \gesperrt{fundamentale Einheit}, und die in den Formeln des vorigen Paragraphen vorkommende Determinante $E$ [\S.\,195, (12)] wird hier (da $\tau$ ein echter Bruch ist):
\[
E = -\log \tau.
\]
Man erhält also aus \S.\,195, (15), (18):
\[
h_1 = 1, \quad k = 1, \quad h = 1.
\]
Der Körper $\Omega_8$ ist also ein \gesperrt{einclassiger}, d.\,h. ein solcher, in dem die Zerlegung der ganzen Zahlen in Primfactoren nach denselben Gesetzen, wie im Körper der rationalen Zahlen, möglich ist.

Es hat sich dabei zugleich ergeben, dass alle Einheiten im Körper $H_8$ in der Form $\pm \tau^a$ darstellbar sind. Die conjugirten Werthe der Einheit $\tau$ sind aber
\[
\tau_1 = \operatorname{tg}\frac{\pi}{8}, \quad \tau_3 = -\operatorname{tg}\frac{3\pi}{8} = -\tau_1^{-1},
\]
und haben also verschiedene Zeichen. Wir schliessen daraus für diesen Fall auf den Satz:
%CONT>>


% ===== scan leaf 742, printed 719 =====
% !PRINTED 719
\weberpage{719}
\pdfbookmark[2]{§ 197. Recurrente Berechnung der Classenzahl im Körpar Ωm, wenn m eine Potenz von 2 ist}{bm225}
%<<CONT

\gesperrt{Eine Einheit in $H_8$, die mit ihren Conjugirten von einerlei Zeichen ist,
ist, vom Zeichen abgesehen, das Quadrat einer Einheit.}

\begin{center}
\S.\,197.

\gesperrt{Recurrente Berechnung der Classenzahl im K\"orper $\Omega_m$,
wenn $m$ eine Potenz von 2 ist.}
\end{center}

Zur allgemeinen Bestimmung der Classenzahl im K\"orper $\Omega_m$ unter
der Voraussetzung, dass $m$ irgend eine Potenz von 2 und gr\"osser als 8
ist, wenden wir ein recurrirendes Verfahren an. Es sei also
\weq{1}{m = 2^\varkappa, \quad \mu = 2^{\varkappa-1}, \quad
\nu = 2^{\varkappa-2}, \quad \nu' = 2^{\varkappa-3}.}
Neben dem K\"orper $\Omega_m$ betrachten wir den K\"orper $\Omega_\mu$,
der ein Theiler von $\Omega_m$ ist, und den wir hier auch mit $\Omega_m'$
bezeichnen wollen. Es sei $h'$ die Classenzahl im K\"orper $\Omega_m'$,
d.\,h. die Zahl, die sich aus $h$ ergiebt, wenn $\varkappa$ in
$\varkappa - 1$ verwandelt wird, und wir setzen
\weq{2}{h = h' H:}
Setzen wir $h'$ als schon bekannt voraus, so kommt es nur noch auf die
Berechnung von $H$ an, was, wie sich zeigen wird, eine ganze Zahl ist.
Indem wir $h$ und $h'$ wie im \S.\,195 zerlegen, setzen wir
\[
\begin{aligned}
(3)\quad h &= k\,h_1, & \qquad h' &= k'\,h_1', \\
(4)\quad k &= k'A, & \qquad h_1 &= h_1'B,
\end{aligned}
\qquad H = AB,
\]
worin $k'$ aus $k$ und $h_1'$ aus $h_1$ durch die Vertauschung von
$\varkappa$ mit $\varkappa - 1$ hervorgeht. Wir wollen \"uberhaupt jetzt
durch Accente an den Buchstaben andeuten, dass $\varkappa - 1$ statt
$\varkappa$ gesetzt sein soll.

Nach \S.\,176, (7) ist dann
\[
\begin{aligned}
\Delta_1 &= 2^{(\varkappa-1)2^{\varkappa-2}-1}, \\
\Delta_1' &= 2^{(\varkappa-2)2^{\varkappa-3}-1},
\end{aligned}
\]
und folglich
\weq{5}{\Delta_1 = 2^{\varkappa\,2^{\varkappa-3}} \Delta_1'.}
Hat $E$ die Bedeutung \S.\,195, (12), und geht $E'$ aus $E$ hervor, wenn
$\varkappa$ durch $\varkappa - 1$ ersetzt wird, so setzen wir noch
\weq{6}{E = E'D,}
%CONT>>


% ===== scan leaf 743, printed 720 =====
% !PRINTED 720
\weberpage{720}

%<<CONT
wodurch $D$ als eine von Null verschiedene positive Zahl definirt ist.
Was die Summen
\[
X = \overset{n}{\Sigma}\, \frac{\chi(n)}{n}
\]
betrifft, so ist daran zu erinnern, dass $\chi(n)$ [nach \S.\,193, (8)]
definirt ist durch
\weq{7}{n \equiv (-1)^\alpha 5^\beta \pmod{m}, \qquad
\chi(n) = (\pm 1)^\alpha \Theta^\beta,}
worin $\Theta$ eine $\nu^{\mathrm{te}}$ Einheitswurzel bedeutet.

Unter den $\nu^{\mathrm{ten}}$ Einheitswurzeln sind aber auch die
$\nu'^{\mathrm{ten}}$ Einheitswurzeln enthalten, und unter den
Charakteren $\chi$ auch die Charaktere f\"ur den K\"orper $\Omega_m'$.
Wir unterscheiden demnach \gesperrt{primitive und imprimitive Charaktere des
K\"orpers $\Omega_m$, indem wir die dem K\"orper $\Omega_m$
eigenth\"umlichen Charaktere $\chi$, d.\,h. die, in denen $\Theta$ eine
primitive $\nu^{\mathrm{te}}$ Einheitswurzel} ist, als primitiv
bezeichnen.

Unter den Summen $X$ kommen auch alle zu dem K\"orper $\Omega_m'$
geh\"origen Summen vor:
\[
X' = \overset{n}{\Sigma}\, \frac{\chi'(n)}{n}.
\]

Wenn man nun in (4) f\"ur $k, k'$ und $h_1, h_1'$ die im \S.\,195, (16),
(17) gebildeten Ausdr\"ucke einsetzt, und dann die Summen $X'$
weghebt, so ergiebt sich:
\weq{8}{\begin{aligned}
\pi^{\nu'} A &= 2 \cdot 2^{2^{\varkappa-4}(\varkappa-2)} \prod X_2 \\
D B &= 2^{2^{\varkappa-4}(\varkappa-2)} \prod X_1,
\end{aligned}}
\gesperrt{worin sich jetzt aber die Producte $\prod$ nur noch auf die mit den
primitiven Charakteren $\chi_1, \chi_2$ gebildeten Summen $X_1, X_2$
erstrecken.}

In den im \S.\,194, (12) gegebenen Ausdr\"ucken f\"ur die $X$:
\weq{9}{\begin{aligned}
X_1 &= -\frac{1}{2m} \sum_{1,\,m-1}^{s} f_1(r^s)
\log\left(\sin \frac{s\pi}{m}\right)^2 \\
X_2 &= -\frac{i\pi}{m^2} \sum_{1,\,m-1}^{s} s\, f_2(r^s)
\end{aligned}}
treten nun, unter der Voraussetzung, dass $m$ eine Potenz von 2 und
dass $\chi_1, \chi_2$ primitive Charaktere sind, bedeutende
Vereinfachungen ein.


% ===== scan leaf 744, printed 721 =====
% !PRINTED 721
\weberpage{721}

Die Zahl $1 + \mu$ hat die Indices $\alpha = 0$, $\beta = \nu'$, und
folglich ist nach (7)
\[
\chi(1+\mu) = \Theta^{\nu'} = -1,
\]
wenn $\chi$ ein primitiver Charakter ist.

Ferner ist f\"ur ein ungerades $n$:
\[
n(1+\mu) \equiv n + \mu \pmod{m},
\]
und folglich
\[
\chi(n+\mu) = -\chi(n).
\]
Wenn wir nun die Function $f(r^s)$ betrachten:
\weq{10}{f(r^s) = \overset{n}{\Sigma} \chi(n)\, r^{sn},}
worin $n$ die ungeraden Zahlen eines vollen Restsystems f\"ur den Modul
$m$ durchl\"auft, so erhalten wir, da wir in (10) unter dem
Summenzeichen $n$ durch $n+\mu$ ersetzen d\"urfen:
\[
f(r^s) = -r^{\mu s} \sum \chi(n)\, r^{sn},
\]
und da $r^\mu = -1$ ist:
\[
f(r^s) = -(-1)^s f(r^s).
\]
Es ist also
\weq{11}{f(r^s) = 0, \quad \text{\gesperrt{wenn $s$ gerade ist,}}}
\gesperrt{und demnach k\"onnen wir in den Formeln (9) den Summationsbuchstaben $s$
auf die ungeraden Zahlen beschr\"anken, die kleiner als $m$ sind.}

Es ist aber nach (10):
\[
f(r^s) = \chi(s)^{-1} \overset{n}{\Sigma} \chi(ns)\, r^{sn},
\]
und wenn man $sn$ durch $n$ ersetzt, was bei ungeradem $s$ gestattet
ist:
\weq{12}{f(r^s) = \chi(s)^{-1} f(r), \quad \text{\gesperrt{wenn $s$ ungerade ist.}}}
Die Function $f(r)$ hat, wenn $\alpha, \beta$ die Indices von $n$ sind,
den Ausdruck
\weq{13}{f(r) = \overset{n}{\Sigma} (\pm 1)^\alpha \Theta^\beta r^n,}
und ist also mit der im \S.\,17 dieses Bandes betrachteten
Kreistheilungsresolvente
\[
(\pm 1, \Theta, r)
\]
gleichbedeutend, f\"ur die im \S.\,17, (17) die Relation aufgestellt
war:
\weq{14}{(\pm 1, \Theta, r)(\pm 1, \Theta^{-1}, r) = \pm m.}
Hier gelten durchweg die oberen Zeichen f\"ur die Summen $X_1$, die
unteren f\"ur die Summen $X_2$ [\S.\,195, (7)].

Wir theilen jetzt die Reihe der in (9) vorkommenden Zahlen $s$ in vier
Theile. Es soll $t$ ein Zeichen sein, welches die
%CONT>>


% ===== scan leaf 745, printed 722 =====
% !PRINTED 722
\weberpage{722}
\pdfbookmark[2]{§ 198. Der Classenzahlfactor A}{bm226}

%<<CONT
Reihe der positiven ungeraden Zahlen von 1 bis $\nu - 1$ durchl\"auft;
dann durchl\"auft das ungerade $s$ die vier Zahlenreihen:
\[
t, \quad t+\mu, \quad m-t, \quad m-t-\mu,
\]
und es ist:
\[
\begin{aligned}
\chi(t+\mu) &= -\chi(t), \\
\chi_1(m-t) &= \chi_1(t), \qquad \chi_1(m-\mu-t) = -\chi_1(t), \\
\chi_2(m-t) &= -\chi_2(t), \qquad \chi_2(m-\mu-t) = \chi_2(t).
\end{aligned}
\]
Nach (12) ergiebt sich hieraus:
\[
\begin{aligned}
f(r^{t+\mu}) &= -f(r^t), \\
f_1(r^{m-t}) &= f_1(r^t), \qquad f_1(r^{m-\mu-t}) = -f_1(r^t), \\
f_2(r^{m-t}) &= -f_2(r^t), \qquad f_2(r^{m-\mu-t}) = f_2(r^t).
\end{aligned}
\]
Theilt man demnach die Summen (9) in je vier Theile und benutzt noch
die trigonometrische Gleichung
\[
\sin \frac{(t+\mu)\pi}{m} = \cos \frac{t\pi}{m},
\]
so ergiebt sich
\[
\begin{aligned}
X_1 &= \frac{-1}{m} \sum^{t} f_1(r^t)
\log\left(\operatorname{tg} \frac{t\pi}{m}\right)^2, \\
X_2 &= \frac{i\pi}{m} \sum^{t} f_2(r^t).
\end{aligned}
\]
In dieser Summe ist
\[
\frac{t\pi}{m} < \frac{\pi}{4}, \qquad
\operatorname{tg} \frac{t\pi}{m} > 0,
\]
und folglich erhalten wir mit R\"ucksicht auf (12), (13), wenn wir
$\Theta^{-1}$ f\"ur $\Theta$ setzen, wodurch $\chi(t)^{-1}$ in $\chi(t)$
\"ubergeht:
\weq{15}{X_1 = -\frac{2}{m} (+1, \Theta^{-1}, r) \sum^{t} \chi_1(t)
\log\left(\operatorname{tg} \frac{t\pi}{m}\right),}
\weq{16}{X_2 = \frac{i\pi}{m} (-1, \Theta^{-1}, r) \sum^{t} \chi_2(t),
\qquad 1 \leqq t \leqq \nu - 1.}

\begin{center}
\S.\,198.

\gesperrt{Der Classenzahlfactor $A$.}
\end{center}

Um nun zun\"achst $A$ zu berechnen, haben wir den Ausdruck (16) f\"ur
$X_2$ in die erste Formel (8) des vorigen Paragraphen einzusetzen.
Diese Formel k\"onnen wir auch so darstellen:
\weq{1}{\pi^{\nu'} A = 2\, m^{1/4\,\nu} 2^{-\nu'} \prod X_2.}


% ===== scan leaf 746, printed 723 =====
% !PRINTED 723
\weberpage{723}

Bezeichnen wir mit $\alpha$, $\beta$ die Indices von $t$, so ist $\chi_2(t) = (-1)^\alpha\Theta^\beta;$

\weq{2}{\overset{t}{\Sigma}\,\chi_2(t) = \Sigma\,(-1)^\alpha\Theta^\beta = \varphi(\Theta)}

ist eine \gesperrt{ganze Zahl des Körpers} $\Omega_\nu$, und es wird:
\[
X_2 = \frac{i\pi}{m}\,(-1,\Theta^{-1},r)\,\varphi(\Theta).
\]

Nun hat man für $\Theta$ alle Wurzeln der Gleichung
\weq{3}{\Theta^{\nu'} + 1 = 0}
zu setzen, und erhält für $A$, mit Rücksicht auf \S.\,197, (14):
\weq{4}{A = 2^{1-\nu'}\overset{\Theta}{\prod}\,\varphi(\Theta).}

Die Summe, durch die in (2) die Zahl $\varphi(\Theta)$ definirt ist,
enthält $\nu'$ Glieder von der Form $(-1)^\alpha\Theta^\beta$, worin $\alpha$, $\beta$ so zu
bestimmen sind, dass
\weq{5}{(-1)^\alpha\,5^\beta \equiv t \pmod{m}}
und $t$ zwischen 0 und $\nu$ liegt. Nun ist zunächst ersichtlich, dass
unter den Exponenten $\beta$ nicht zwei nach dem Modul $\nu'$ congruente vorkommen können. Denn ist $\beta \equiv \beta' \pmod{\nu'}$, so ist
$5^\beta \equiv 5^{\beta'} \pmod{\mu}$, und aus
\[
(-1)^\alpha\,5^\beta \equiv t, \quad (-1)^{\alpha'}\,5^{\beta'} \equiv t' \pmod{m}
\]
folgt:
\[
(-1)^\alpha\,t - (-1)^{\alpha'}\,t' \equiv 0 \pmod{\mu}.
\]

Da aber $t$ und $t'$ absolut kleiner als $\nu = \tfrac{1}{2}\mu$ sind, so ist
dies nur möglich, wenn $t = t'$, $\alpha = \alpha'$ ist. Demnach durchläuft
$\beta$ in (2) ein volles Restsystem nach dem Modul $\nu'$.

Da aber $\beta$ nach dem Modul $\nu$ zu nehmen ist, so werden die
nach (5) bestimmten $\beta$ theils kleiner, theils grösser als $\nu'$ ausfallen, und wenn wir also $\beta_1$ die Reihe der Zahlen $0, 1, \dots, \nu' - 1$
durchlaufen lassen, so erhalten wir in (2) zweierlei Arten von
Gliedern:
\weq{6}{\begin{aligned}
1.\ &(-1)^\alpha\Theta^\beta = (-1)^\alpha\Theta^{\beta_1}, & \beta = \beta_1 < \nu', \\
2.\ &(-1)^\nu\Theta^\beta = -(-1)^\alpha\Theta^{\beta_1}, & \beta = \beta_1 + \nu' \geqq \nu'.
\end{aligned}}
% [sic] in row 2 the page prints the exponent \nu where \alpha is required; reproduced as printed

Nun ist nach (5) der absolut kleinste Rest von $(-1)^\alpha\,5^\beta$
nach dem Modul $m$ positiv, und wenn wir daher eine Zahl
$c_\beta = \pm 1$ so bestimmen, dass der absolut kleinste Rest von $c_\beta\,5^\beta$
positiv wird, so ist in dem Falle (6), 1.
\[
c_{\beta_1} = (-1)^\alpha.
\]


% ===== scan leaf 747, printed 724 =====
% !PRINTED 724
\weberpage{724}

Ferner ist $5^{\nu'} \equiv 1 \pmod{\mu}$, und da $5^{\nu'}$ nicht auch nach dem
Modul $m$ mit 1 congruent ist, so folgt:
\weq{7}{5^{\nu'} \equiv 1 + \mu \pmod{m}.}
Demnach ist im Falle (6), 2.
\[
(-1)^\alpha\,5^\beta = (-1)^\alpha\,5^{\beta_1}\,5^{\nu'} \equiv (-1)^\alpha\,5^{\beta_1} + \mu \pmod{m},
\]
und folglich ist hier der absolut kleinste Rest von $(-1)^\alpha\,5^{\beta_1}$
nach dem Modul $m$ negativ. Es ist also im Falle (6), 2.
\[
c_{\beta_1} = -(-1)^\alpha
\]
zu setzen.

Daraus ergiebt sich nach (2) (wenn wir wieder $\beta$ für $\beta_1$
schreiben):
\weq{8}{\varphi(\Theta) = \sum_{0,\,\nu'-1}^{\beta} c_\beta\,\Theta^\beta = c_0 + c_1\Theta + c_2\Theta^2 + \cdots + c_{\nu'-1}\Theta^{\nu'-1},}
und hierin ist
\weq{9}{c_\beta = +1 \ \text{oder} \ = -1,}
je nachdem der absolut kleinste Rest von $5^\beta$ positiv oder
negativ ist. Hiernach ist also $\varphi(\Theta)$ eine \gesperrt{ganze Zahl des
Körpers} $\Omega_\nu$.

Es kommt noch darauf an, das Verhalten dieser Zahl zu
der Zahl 2, oder vielmehr zu dem in $\Omega_\nu$ enthaltenen Primfactor
$(1 - \Theta)$ von 2 zu ermitteln. Bilden wir zu diesem Zwecke
\weq{10}{(1 - \Theta)\,\varphi(\Theta) = 2\,\psi(\Theta),}
so ergiebt sich
\weq{11}{\psi(\Theta) = \frac{c_0 + c_{\nu'-1}}{2} + \frac{c_1 - c_0}{2}\,\Theta + \frac{c_2 - c_1}{2}\,\Theta^2 + \cdots + \frac{c_{\nu'-1} - c_{\nu'-2}}{2}\,\Theta^{\nu'-1}.}

Die Co\"efficienten in diesem Ausdrucke sind alle $= 0$ oder
$= 1$, und daher ist auch $\psi(\Theta)$ eine ganze Zahl des Körpers $\Omega_\nu$.
Für diese Zahl ergiebt sich aber, wenn man $\Theta \equiv 1$ setzt:
\weq{12}{\psi(\Theta) \equiv c_{\nu'-1} \equiv \pm 1\ [\mathrm{mod}\ (1-\Theta)],}
und folglich ist $\psi(\Theta)$ relativ prim zu 2.

Das in (4) vorkommende Product ist nichts Anderes, als die
in Bezug auf den Körper $\Omega_\nu$ genommene Norm der Zahl $\varphi(\Theta)$,
und wenn wir diese Norm mit $N$ bezeichnen, so ist nach \S.\,169
\[
N(1-\Theta) = 2,
\]
%CONT>>


% ===== scan leaf 748, printed 725 =====
% !PRINTED 725
\weberpage{725}
%<<CONT
und folglich nach (10)
\[
N\varphi(\Theta) = 2^{\nu'-1}\,N\psi(\Theta).
\]
Daraus ergiebt sich
\weq{13}{A = N\psi(\Theta),}
und hieraus folgt mit Rücksicht auf (12), dass $A$ eine \gesperrt{ungerade
ganze Zahl ist.}

Bezeichnen wir die zu $m = 2^\varkappa$ gehörige Zahl $A$ mit $A_\varkappa$, so
ergiebt sich aus \S.\,197, (4) durch vollständige Induction der
Ausdruck für den ersten Factor der Classenzahl:
\weq{14}{k = A_4\,A_5\,\dots\,A_\varkappa,}
und daraus der Satz:

1. \gesperrt{Der erste Factor der Classenzahl ist eine ungerade ganze Zahl.}

Die Zahl $A_\varkappa$ ist nach der Formel (10) für die ersten Werthe
von $\varkappa$ nicht schwer zu berechnen.

Für $m = 16$ findet man $\nu' = 2$, $\beta = 0, 1$, $c_0 = c_1 = 1$
folglich
\[
\psi(\Theta) = 1, \quad A_4 = 1.
\]

Für $m = 32$, $\nu' = 4$, $\beta = 0, 1, 2, 3$.
\[
c_0 = 1, \quad c_1 = 1, \quad c_2 = -1, \quad c_3 = -1,
\]
\[
\psi(\Theta) = -\Theta^2,
\]
also auch $A_5 = 1$.

Für $m = 64$, $\nu' = 8$, $\beta = 0, 1, 2, 3, 4, 5, 6, 7$ sind die
absolut kleinsten Reste von $5^\beta$
\[
1,\ 5,\ 25,\ -3,\ -15,\ -11,\ 9,\ -19;
\]
also sind die $c_\beta$:
\[
+1, +1, +1, -1, -1, -1, +1, -1
\]
und
\[
\psi(\Theta) = -\Theta^3 + \Theta^6 - \Theta^7, \quad \psi(\Theta)\,\psi(-\Theta) = \Theta^6(\Theta^6 - 2i),
\]
woraus sich leicht ergiebt:
\[
A_6 = 17.
\]
Eine etwas längere Rechnung ergiebt noch
\[
A_7 = 21121.
\]


% ===== scan leaf 749, printed 726 =====
% !PRINTED 726
\weberpage{726}
\pdfbookmark[2]{§ 199. Der Classenzahlfactor B}{bm227}

\begin{center}
\S.\,199.

\gesperrt{Der Classenzahlfactor} $B$.
\end{center}

Auf ganz andere Weise muss der Classenzahlfactor $B$ berechnet werden. Hier ist, wenn $\Theta$ wieder eine Primitivwurzel
der Gleichung \S.\,198, (3) ist ($\Theta^{\nu'} + 1 = 0$), und
\weq{1}{t \equiv (-1)^\alpha\,5^\beta \pmod{m},}
\weq{2}{\chi_1(t) = \Theta^\beta}
zu setzen, und die Formel \S.\,197, (15) ergiebt:
\weq{3}{X_1 = -\frac{2}{m}\,(1, \Theta^{-1}, r)\,\overset{t}{\Sigma}\,\Theta^\beta\,\log\,\mathrm{tg}\,\frac{t\pi}{m}.}

Nun ist für jedes ungerade $t$, wenn [mit Rücksicht auf (1)]
\[
r = e^{\frac{2\pi i}{m}}, \quad r^\nu = i, \quad r^{\nu t} = (-1)^\alpha i
\]
gesetzt wird,
\weq{4}{\mathrm{tg}\,\frac{t\pi}{m} = (-1)^\alpha\,r^{t\nu}\,\frac{1 - r^t}{1 + r^t},}
und dies ist, wie wir schon im \S.\,178 gesehen haben, eine Einheit des reellen Körpers $H_m$. Man erhält die conjugirten Werthe
dieser Einheit in Bezug auf diesen Körper, wenn man $t \equiv 5^\beta \pmod{m}$ setzt, und $\beta$ ein volles Restsystem nach dem Modul $\mu$
durchlaufen lässt. Wir setzen demnach:
\weq{5}{\tau_\beta = \mathrm{tg}\,\frac{5^\beta\pi}{m}.}

Nach \S.\,198, (7) ist $5^{\beta+\nu'} \equiv 5^\beta + \mu \pmod{m}$, und daraus
ergiebt sich nach einer bekannten trigonometrischen Formel:
\weq{6}{\tau_{\beta+\nu'} = \frac{-1}{\tau_\beta}.}

Setzen wir also nun
\weq{7}{\lambda_\beta = \tfrac{1}{2}\,\log\,\tau_\beta^2,}
so ist nach (6)
\weq{8}{\lambda_{\beta+\nu'} = -\lambda_\beta,}
und demnach
\weq{9}{\Theta^{\beta+\nu'}\,\lambda_{\beta+\nu'} = \Theta^\beta\,\lambda_\beta;}
unter den in der Summe (3) vorkommenden Werthen von $\beta$
sind, wie wir schon im vorigen Paragraphen gesehen haben, nicht
zwei nach dem Modul $\nu'$ congruent, und wenn wir daher mittelst
%CONT>>


% ===== scan leaf 750, printed 727 =====
% !PRINTED 727
\weberpage{727}

%<<CONT
der Formel (9) alle Werthe von $\beta$ auf das Intervall von $0$ bis
$\nu' - 1$ reduciren, so bekommen wir alle Zahlen dieses Intervalles.

Demnach lässt sich die Formel (3) so darstellen:
\weq{10}{
X_1 = -\frac{2}{m}\,(1,\,\Theta^{-1},\,r)\sum_{0,\,\nu'-1}^{\beta}\Theta^\beta\lambda_\beta.
}
Dies haben wir in dem im \S.\,197, (8) angegebenen Ausdrucke
\[
D B = 2^{2^{\varkappa-4}(\varkappa-2)}\prod X_1 = m^{\nu''}2^{-\nu'}\prod X_1
\]
einzusetzen, und dann ergiebt sich, wenn wir
\weq{11}{
U_\Theta = \sum_{0,\,\nu'-1}^{\beta}\Theta^\beta\lambda_\beta
}
setzen, mit Rücksicht auf \S.\,197, (14):
\weq{12}{
D B = \prod^{\Theta} U_\Theta.
}
Dies Product lässt sich nun auf folgende Weise durch eine
Determinante darstellen, die nur die conjugirten Logarithmen
der Einheiten $\tau$ enthält.

Lassen wir den Index $\Theta$ bei $U$ weg, so können wir das
Gleichungssystem ansetzen:
\[
\begin{array}{l}
U \;=\; \lambda_0 \;+\; \lambda_1\Theta \;+\; \lambda_2\Theta^2 \;+\; \cdots \;+\; \lambda_{\nu'-1}\Theta^{\nu'-1} \\[4pt]
\Theta U \;=\; -\lambda_{\nu'-1} \;+\; \lambda_0\Theta \;+\; \lambda_1\Theta^2 \;+\; \cdots \;+\; \lambda_{\nu'-2}\Theta^{\nu'-1} \\[4pt]
\text{\normalfont .\;.\;.\;.\;.\;.\;.\;.\;.\;.\;.\;.\;.\;.\;.\;.\;.\;.\;.\;.\;.\;.\;.\;.\;.\;.\;.\;.\;.\;.} \\[4pt]
\Theta^{\nu'-1} U \;=\; -\lambda_1 \;-\; \lambda_2\Theta \;-\; \lambda_3\Theta^2 \;+\; \cdots \;+\; \lambda_0\Theta^{\nu'-1},
\end{array}
\]
woraus durch Elimination der Potenzen von $\Theta$ folgt:
\weq{13}{
\begin{vmatrix}
\lambda_0 - U, & \lambda_1, & \lambda_2, & \ldots, & \lambda_{\nu'-1} \\
-\lambda_{\nu'-1}, & \lambda_0 - U, & \lambda_1, & \ldots, & \lambda_{\nu'-2} \\
\multicolumn{5}{c}{\text{\normalfont .\;.\;.\;.\;.\;.\;.\;.\;.\;.\;.\;.\;.\;.\;.\;.}} \\
-\lambda_1, & -\lambda_2, & -\lambda_3, & \ldots, & \lambda_0 - U
\end{vmatrix} = 0.
}
% [source note: the first row prints a semicolon-shaped mark after \lambda_0-U; the mathematically required list punctuation is normalized to a comma.]
Dies ist für $U$ eine Gleichung $\nu'^{\text{ten}}$ Grades, deren Wurzeln
die Grössen $U_\Theta$ sind. Ordnet man nach Potenzen von $U$, so
erhält, da $\nu'$ gerade ist, die höchste Potenz $U^{\nu'}$ von $U$ den
Co\"efficienten $1$, und man findet also das in (12) vorkommende
Product der Wurzeln, wenn man $U = 0$ setzt, aus der Determinante (13).

Die so gewonnene Determinante wollen wir etwas anders anordnen,
indem wir die erste Zeile an ihrer Stelle lassen, die
%CONT>>


% ===== scan leaf 751, printed 728 =====
% !PRINTED 728
\weberpage{728}
\pdfbookmark[2]{§ 200. Normaleinheiten in Hm}{bm228}
%<<CONT
übrigen aber in umgekehrter Ordnung und mit verändertem Vorzeichen schreiben.
Dabei sind $\frac{1}{2}\nu'$ Vorzeichenänderungen nöthig,
und wir erhalten daher für das Product (12) die durch folgende
Gleichung definirte Determinante, deren absoluten Werth wir
mit $T$ bezeichnen:
\weq{14}{(-1)^{\frac{1}{2}\nu'}
\begin{vmatrix}
\lambda_0, & \lambda_1, & \ldots, & \lambda_{\nu'-2}, & \lambda_{\nu'-1}\\
\lambda_1, & \lambda_2, & \ldots, & \lambda_{\nu'-1}, & -\lambda_0\\
\lambda_2, & \lambda_3, & \ldots, & -\lambda_0, & -\lambda_1\\
\cdots & \cdots & \cdots & \cdots & \cdots\\
\lambda_{\nu'-1}, & -\lambda_0, & \ldots, & -\lambda_{\nu'-3}, & -\lambda_{\nu'-2}
\end{vmatrix} = \pm T,}
worin rechts von der von links unten nach rechts oben laufenden
Diagonalreihe die negativen Zeichen stehen.

Die Determinante $T$ ist hier so geordnet, dass jede Zeile
aus der vorangehenden durch die Substitution $(r, r^{5^\beta})$ entsteht.
Jede Colonne enthält also ein System conjugirter Logarithmen.

Nach (12) ist jetzt
\weq{15}{B = \frac{T}{D}.}
Um über die Natur dieses Resultates Klarheit zu bekommen,
ist es nöthig, die Zahl $D$ etwas genauer zu betrachten, die durch
\S.\,195, (12), \S.\,197, (6) definirt war; und dies erfordert ein Eingehen
 auf die Theorie der Einheiten des Körpers $H_m$.

\begin{center}
\S.\,200.\\
\gesperrt{Normaleinheiten in $H_m$.}
\end{center}

Wir führen jetzt eine vereinfachende Bezeichnung ein, indem wir
\weq{1}{r^{5^\beta} = r_\beta}
setzen, und $\beta$ ein volles Restsystem nach dem Modul $\nu$ durchlaufen
 lassen. Es folgt hieraus
\weq{2}{r_{\beta+\nu'} = -r_\beta,}
und in der Form $(r, r_\beta)$ sind dann zwar nicht alle Substitutionen
der Gruppe des Körpers $\Omega_m$ enthalten, sondern es müssen noch
die Substitutionen $(r, r_\beta^{-1})$ hinzukommen; wohl aber liefert uns
$(r, r_\beta)$ alle Substitutionen der Gruppe von $H_m$.

Wir bezeichnen jetzt mit $\mathfrak{E}(r)$ eine Einheit des Körpers $H_m$.
Unter diesen sind auch die Einheiten des Körpers $H_\mu$ (als
%CONT>>


% ===== scan leaf 752, printed 729 =====
% !PRINTED 729
\weberpage{729}
%<<CONT
imprimitive Zahlen) enthalten und durch die Gleichung $\mathfrak{E}(r)
= \mathfrak{E}(-r)$ charakterisirt.

Wir wollen nun unter einer \gesperrt{Normaleinheit} des Körpers $H_m$
eine Einheit verstehen, die nicht $= \pm 1$ ist, aber der Gleichung
\weq{3}{\mathfrak{E}(r)\mathfrak{E}(-r) = \pm 1}
genügt.

Es ist klar, dass eine Einheit des Körpers $H_\mu$ dieser Bedingung
 jedenfalls nicht genügen kann. Aber nicht jede Einheit
in $H_m$, die dem Körper $H_\mu$ nicht angehört, wird Normaleinheit
sein $^{1)}$.

Die Einheiten $\tau_\beta$ [\S.\,199, (5)] gehören aber zu den Normaleinheiten
. Es ist nämlich
\weq{4}{\tau_\beta = \mathrm{tg}\,\frac{5^\beta\pi}{m} = \frac{1-r_\beta}{1+r_\beta}r_\beta^{\nu},}
und wenn darin die Substitution $(r,-r) = (r_\beta, r_{\beta+\nu'})$ gemacht
wird, so geht $\tau_\beta$ in
\weq{5}{\tau_{\beta+\nu'} = -\frac{1}{\tau_\beta}}
über, was der Relation (3) entspricht.

Ein System von $\nu'$ Normaleinheiten
\[
\mathfrak{E}_0(r), \mathfrak{E}_1(r), \ldots, \mathfrak{E}_{\nu'-1}(r)
\]
soll ein \gesperrt{unabhängiges System} heissen, wenn die Determinante
\weq{6}{\pm\Sigma\pm\log|\mathfrak{E}_0(r_0)|\log|\mathfrak{E}_1(r_1)|\ldots\log|\mathfrak{E}_{\nu'-1}(r_{\nu'-1})|
= L(\mathfrak{E}_0, \mathfrak{E}_1, \ldots, \mathfrak{E}_{\nu'-1})}
einen von Null verschiedenen Werth hat.

Den absoluten Werth dieser Determinante, $L$, nennen wir
auch hier den \gesperrt{Regulator} des Systems $\mathfrak{E}_0, \mathfrak{E}_1, \ldots, \mathfrak{E}_{\nu'-1}$.

Die Einheiten $\tau_0, \tau_1, \ldots, \tau_{\nu'-1}$ bilden ein unabhängiges System
normaler Einheiten, denn ihr Regulator
\[
L(\tau_0, \tau_1, \ldots, \tau_{\nu'-1})
\]
ist eben die im vorigen Paragraphen definirte Determinante $T$
[\S.\,199, (14), (15)]. Und dass diese nicht Null sein kann, ergiebt
sich unmittelbar daraus, dass sie in dem Ausdrucke für die
%CONT>>

\par\noindent{\footnotesize $^{1)}$ Der Verfasser hat in der oben erwähnten Arbeit (Acta
mathematica, Bd.\,8) die Normaleinheiten als \glqq primitive
Einheiten\grqq{} bezeichnet. Dieser Name ist hier verworfen, weil er
zu dem Irrthum Veranlassung geben könnte, als ob jede Einheit
des Körpers $H_m$, die nicht zugleich dem Körper $H_\mu$ angehört,
dazu zu rechnen sei.}


% ===== scan leaf 753, printed 730 =====
% !PRINTED 730
\weberpage{730}
%<<CONT
Classenzahl als Factor vorkommt. Die Classenzahl ist aber, da
gewiss immer wenigstens eine Classe, die Hauptclasse, existirt,
von Null verschieden, und also kann auch $T$ nicht gleich Null
sein. Wir haben daher den Satz:

1. \gesperrt{Die Zahlen $\tau_0, \tau_1, \ldots, \tau_{\nu'-1}$ sind ein unabhängiges
System von Normaleinheiten des Körpers $H_m$.}

Ist $\mathfrak{E}(r)$ eine Normaleinheit, so ist, wenn wir
\weq{7}{L_\beta = \log|\mathfrak{E}(r_\beta)|}
setzen, wegen (3)
\weq{8}{L_\beta = -L_{\beta+\nu'}.}
Ist nun $\mathfrak{E}_0, \mathfrak{E}_1, \ldots, \mathfrak{E}_{\nu'-1}$ irgend ein System unabhängiger
% [source note: the separator after \mathfrak{E}_0 is printed as a period-shaped or damaged comma; the mathematically required list punctuation is normalized to a comma.]
Normaleinheiten, so setzen wir:
\weq{9}{L_{i,k} = \log|\mathfrak{E}_i(r_k)|, \quad L_{i,k+\nu'} = -L_{i,k},}
so dass der Regulator des Systems der $\mathfrak{E}_i$ den Ausdruck erhält:
\weq{10}{\pm\Sigma\pm L_{0,0}L_{1,1}\cdots L_{\nu'-1,\nu'-1} = L(\mathfrak{E}_0, \mathfrak{E}_1, \ldots, \mathfrak{E}_{\nu'-1}),}
und eine von Null verschiedene Zahl ist, die wir mit $T_0$ bezeichnen.

Bedeutet $\mathfrak{E}(r)$ irgend eine andere Normaleinheit, so können
wir die Unbekannten $\xi_0, \xi_1, \ldots, \xi_{\nu'-1}$ aus den $\nu'$ linearen Gleichungen
 bestimmen, die wir erhalten, wenn wir in
\weq{11}{L_\beta = \xi_0 L_{0,\beta} + \xi_1 L_{1,\beta} + \cdots + \xi_{\nu'-1}L_{\nu'-1,\beta}}
den Index $\beta$ von 0 bis $\nu'-1$ gehen lassen, und wegen (8) und (9)
ist diese Gleichung auch für die übrigen Werthe von $\beta$ richtig,
so dass man daraus die für alle $r$ gültige Formel ableiten kann:
\weq{12}{\mathfrak{E}(r) = \pm\mathfrak{E}_0(r)^{\xi_0}\mathfrak{E}_1(r)^{\xi_1}\cdots\mathfrak{E}_{\nu'-1}(r)^{\xi_{\nu'-1}}.}
Bedeuten andererseits $m_0, m_1, \ldots, m_{\nu'-1}$ irgend welche ganze
rationale Zahlen, so sind alle in der Form
\[
\pm\mathfrak{E}_0(r)^{m_0}\mathfrak{E}_1(r)^{m_1}\cdots\mathfrak{E}_{\nu'-1}(r)^{m_{\nu'-1}}
\]
enthaltenen Zahlen Normaleinheiten des Körpers $H_m$. Demnach
lassen sich auf das Gleichungssystem (10) die Betrachtungen
anwenden, die wir im \S.\,186 ausführlich durchgeführt haben,
woraus sich ergiebt:

2. \gesperrt{Die aus (10) oder (11) bestimmten Exponenten
$\xi_0, \xi_1, \ldots, \xi_{\nu'-1}$ sind rationale Zahlen.}


% ===== scan leaf 754, printed 731 =====
% !PRINTED 731
\weberpage{731}

Hieraus kann, wie im \S.\,187, gefolgert werden, dass es Systeme von $\nu'$ unabhängigen Normaleinheiten giebt, deren Regulator $T_0$ \gesperrt{so klein als möglich} wird, und solche Systeme heissen \gesperrt{Fundamentalsysteme von Normaleinheiten}.

Ist $\mathcal{E}_0, \mathcal{E}_1, \ldots, \mathcal{E}_{\nu'-1}$ ein solches Fundamentalsystem, so ergeben sich die Exponenten $\xi_0, \xi_1, \ldots, \xi_{\nu'-1}$ in (10) und (11) als \gesperrt{ganze rationale Zahlen}. Auch dies kann ebenso bewiesen werden, wie im \S.\,187.

Stellen wir unter dieser Voraussetzung über die $\mathcal{E}_i$ die Gleichungen (11) für irgend ein anderes System unabhängiger Einheiten $\mathcal{E}_0^{(1)}, \mathcal{E}_1^{(1)}, \ldots, \mathcal{E}_{\nu'-1}^{(1)}$ auf, so erhalten wir Ausdrücke von der Form:

\[
L_{i,k}^{(1)} = \xi_{i,0} L_{0,k} + \xi_{i,1} L_{1,k} + \cdots + \xi_{i,\nu'-1} L_{\nu'-1,k}.
\]

Ist also $T_1$ der Regulator des Systems der $\mathcal{E}_i^{(1)}$, so ergiebt sich nach dem Multiplicationssatze der Determinanten, wenn mit $C$ der absolute Werth der Determinante der ganzzahligen Exponenten $\xi_{i,k}$, also eine \gesperrt{ganze rationale positive Zahl} bezeichnet wird:

\weq{13}{T_1 = C\,T_0.}

Es kann nicht bewiesen werden und trifft für die höheren Werthe von $m$ wahrscheinlich auch nicht zu, dass $\tau_0, \tau_1, \ldots, \tau_{\nu'-1}$ ein Fundamentalsystem bilden\footnote{Es ist hier, wie bei allen Fragen, die die Einheiten betreffen, sehr schwierig, zu bestimmten numerischen Resultaten zu gelangen. Nach den Resultaten, die in dem reichhaltigen Werke von \gesperrt{Reuschle}, \glqq Tafeln complexer Primzahlen, welche aus Wurzeln der Einheiten gebildet sind\grqq{} (Berlin 1875), enthalten sind, bilden für $m=16$ die $\tau$ noch ein Fundamentalsystem. Denn hier ist die Grundzahl des Körpers $H_m$ gleich $2^{11}$, der Grad des Körpers gleich 4, und daher giebt es nach der Anmerkung zu \S.\,156 in jeder Idealclasse ein Ideal, dessen Norm kleiner als 6 ist. Nach den Tafeln von \gesperrt{Reuschle} kann aber jede natürliche Primzahl unter 100 im Körper $H_m$ in wirklich existirende Primfactoren zerlegt werden. Folglich gehen gewiss in allen Zahlen unter 6 nur Hauptideale auf, und folglich giebt es hier nur die eine Classe, die Hauptclasse. Da, wie wir gesehen haben, auch der erste Classenzahlfactor $=1$ ist, so ist der Kreistheilungskörper $\Omega_{16}$ einclassig. Es gelten in ihm dieselben Gesetze der Primzahlen, wie im Körper der rationalen Zahlen. Zweifelhaft ist dies für den Körper $\Omega_{32}$ und sicher nicht mehr zutreffend im Körper $\Omega_{64}$, in dem die Classenzahl ein Vielfaches von 17 ist.}.

Für unseren Zweck ist es aber von Wichtigkeit, dass sich diese Einheiten in Bezug auf den Nenner 2 in den Exponenten
%CONT>>


% ===== scan leaf 9732, printed 732 =====
% !PRINTED 732
\weberpage{732}
%<<CONT
gewissermaassen wie ein Fundamentalsystem verhalten, was durch
folgenden Satz bestimmter ausgedrückt wird:

2. \gesperrt{Wenn eine Normaleinheit des Körpers $H_m$
von der Form}
\[
\mathfrak{E}(r) = \tau_0^{e_0}\tau_1^{e_1}\ldots\tau_{\nu'-1}^{e_{\nu'-1}},
\]
\gesperrt{in der die Exponenten $e_0, e_1, \ldots, e_{\nu'-1}$ ganze rationale
 Zahlen sind, mit allen ihren conjugirten
Werthen einerlei Vorzeichen hat, so müssen die
sämmtlichen $e_0, e_1, \ldots, e_{\nu'-1}$ gerade Zahlen sein,
und $\mathfrak{E}(r)$ ist das Quadrat einer Normaleinheit.}

Den Beweis können wir wieder durch vollständige Induction
führen, wollen aber zunächst den Ausdruck für $\tau_\beta$ etwas umformen
.

Die conjugirten Werthe zu der Einheit $\mathfrak{E}(r)$ erhalten wir
in der Form
\[
\mathfrak{E}(r_\beta) = \tau_\beta^{e_0}\tau_{\beta+1}^{e_1}\ldots\tau_{\beta+\nu'-1}^{e_{\nu'-1}},
\]
und diese Ausdrücke sollen also für alle Werthe von $\beta$ dasselbe
Vorzeichen haben.

Es ist nach §.\,199, (5)
\[
\tau_\beta = \mathrm{tg}\,\frac{5^\beta\pi}{m} = \frac{1}{2}\frac{\sin\left(\dfrac{2\pi}{m}5^\beta\right)}{\left[\cos\left(\dfrac{5^\beta\pi}{m}\right)\right]^2}.
\]

Hierin setzen wir zur Abkürzung
\weq{14}{\sigma_\beta = \sin\left(\frac{2\pi}{m}5^\beta\right),}
so dass $\sigma_\beta$ immer dasselbe Vorzeichen hat, wie $\tau_\beta$.

Daraus bilden wir das Product
\weq{15}{S_\beta = \sigma_\beta^{e_0}\sigma_{\beta+1}^{e_1}\ldots\sigma_{\beta+\nu'-1}^{e_{\nu'-1}},}
und unsere Voraussetzung ist also die, dass $S_\beta$ für alle Werthe
von $\beta$ dasselbe Vorzeichen hat.

Es soll nachgewiesen werden, dass die ganzen Zahlen
$e_0, e_1, \ldots, e_{\nu'-1}$ alle gerade sein müssen. Der Satz ist
richtig für $m=8$; denn in diesem Falle ist
\[
\sigma_0 = \sin\frac{2\pi}{8}, \sigma_1 = \sin\frac{10\pi}{8} = -\sin\frac{6\pi}{8},
\]
%CONT>>


% ===== scan leaf 9733, printed 733 =====
% !PRINTED 733
\weberpage{733}
%<<CONT

also $\sigma_0$ positiv, $\sigma_1$ negativ. Hier ist nun $S_\beta = \sigma_\beta^{e_0}$,
und es kann also $S_0$ und $S_1$ nur dann gleiches Vorzeichen haben, wenn $e_0$
gerade ist.

Hiernach wenden wir die vollständige Induction an. Wir
setzen zur Vereinfachung $\nu' = 2\nu''$ und erinnern an die Bezeichnung

\weq{16}{m = 2^{\varkappa}, \quad \mu = \tfrac{1}{2}m, \quad \nu = \tfrac{1}{2}\mu, \quad \nu' = \tfrac{1}{2}\nu, \quad \nu'' = \tfrac{1}{2}\nu',}
und setzen voraus, dass unser Satz als richtig erwiesen sei, wenn
$\mu$ an die Stelle von $m$ tritt.

Es ist dann
\[
\begin{gathered}
5^{\nu'} \equiv 1 \pmod{m}, \quad 5^{\nu'} \equiv 1+\mu \pmod{m}, \quad 5^{\nu''} \equiv 1+\nu \pmod{\mu} \\
5^{\nu''} \equiv 1+\nu+\mu \pmod{m}
\end{gathered}
\]
und daraus ergiebt sich

\[
\sigma_{\beta+\nu'} = -\sigma_\beta,
\]
\[
\sigma_{\beta+\nu''} = -\cos\left(\frac{2\pi}{m}5^\beta\right),
\]
\weq{17}{\sigma_\beta \sigma_{\beta+\nu''} = -\tfrac{1}{2}\sin\left(\frac{2\pi}{\mu}5^\beta\right) = -\tfrac{1}{2}\sigma'_\beta,}
\weq{18}{\sigma'_{\beta+\nu''} = -\sigma'_\beta,}
wenn die $\sigma'_\beta$ aus den $\sigma_\beta$ durch den Uebergang von $\varkappa$ zu
$\varkappa - 1$ hervorgehen. Nun bilden wir nach (17), (18) das Product
$S_\beta S_{\beta+\nu''}$, und erhalten, wenn noch zur Abkürzung
\[
e' = e_0 + e_1 + \cdots + e_{\nu'-1}, \quad e'' = e_0 + e_1 + \cdots + e_{\nu''-1}
\]
gesetzt wird,
\[
S_\beta S_{\beta+\nu''} = (-1)^{e''}\left(\tfrac{1}{2}\right)^{e'}\sigma_\beta'^{\,e_0+e_{\nu''}}\sigma_{\beta+1}'^{\,e_1+e_{\nu''+1}}\cdots\sigma_{\beta+\nu''-1}'^{\,e_{\nu''-1}+e_{\nu'-1}}.
\]

Dies ist aber ein Product von derselben Form wie $S_\beta$, in
dem $\mu$ an Stelle von $m$ getreten ist, was gleichfalls mit allen
seinen Conjugirten einerlei Vorzeichen hat. Es ist daher nach
der gemachten Voraussetzung
\weq{19}{e_0 \equiv e_{\nu''+1}, \; e_1 \equiv e_{\nu''+1}, \; \ldots, \; e_{\nu''-1} \equiv e_{\nu'-1} \pmod{2}.} % [sic: the source prints the same right-hand subscript in the first two congruences; the preceding product indicates that the first should be e_{\nu''}.]
Setzen wir also
\[
S'_\beta = \sigma_\beta'^{\,e_0}\sigma_{\beta+1}'^{\,e_1}\cdots\sigma_{\beta+\nu''-1}'^{\,e_{\nu''-1}},
\]
\[
R_\beta = \sigma_{\beta+\nu''}^{\frac{1}{2}(e_0-e_{\nu''})}\sigma_{\beta+\nu''+1}^{\frac{1}{2}(e_1-e_{\nu''+1})}\cdots\sigma_{\beta+\nu'-1}^{\frac{1}{2}(e_{\nu''-1}-e_{\nu'-1})},
\]
so folgt aus (15) und (17)
\[
S_\beta R_\beta^2 = \left(-\tfrac{1}{2}\right)^{e''}S'_\beta,
\]
%CONT>>


% ===== scan leaf 755, printed 734 =====
% !PRINTED 734
\weberpage{734}

%<<CONT
und hieraus ergiebt sich, dass $S'_\beta$ ebenso wie $S_\beta$ mit allen seinen conjugirten Werthen dasselbe Zeichen hat. $S'_\beta$ entsteht aber aus $S_\beta$ dadurch, dass $\mu$ an Stelle von $m$ gesetzt wird, und nach unserer Voraussetzung ist daher

\[
e_0 \equiv 0,\ e_1 \equiv 0,\ \ldots,\ e_{\nu''-1} \equiv 0 \ (\mathrm{mod}\ 2),
\]

und mit Hülfe von (19)

\[
e_{\nu''} \equiv 0,\ e_{\nu''+1} \equiv 0,\ \ldots,\ e_{\nu'-1} \equiv 0 \ (\mathrm{mod}\ 2),
\]

wodurch der Beweis des Satzes 2. geführt ist.

Wenn wir nun in der Formel (12) für $\mathcal{E}_0, \mathcal{E}_1, \ldots, \mathcal{E}_{\nu'-1}$ das System $\tau_0, \tau_1, \ldots, \tau_{\nu'-1}$ wählen und die Exponenten $\xi$ in die Form setzen:

\[
\xi_0 = \frac{e_0}{e},\ \xi_1 = \frac{e_1}{e},\ \cdots,\ \xi_{\nu'-1} = \frac{e_{\nu'-1}}{e},
\]

worin $e, e_0, e_1, \ldots, e_{\nu'-1}$ ganze rationale Zahlen ohne gemeinsamen Theiler sind, so erhalten wir

\[
\mathcal{E}(r)^e = \pm\, \tau_0^{e_0}\tau_1^{e_1} \cdots \tau_{\nu'-1}^{e_{\nu'-1}},
\]

und es folgt aus unserem Satze, dass $e$ \gesperrt{ungerade} sein muss; denn wäre es gerade, so wäre $\mathcal{E}(r)^e$ ein Quadrat einer reellen Grösse, und folglich hätte $\pm\, \mathcal{E}(r)^e$ mit seinen Conjugirten dasselbe Vorzeichen. Dann aber müssten nach unserem Satze alle $e_0, e_1, \ldots, e_{\nu'-1}$ durch 2 theilbar sein, was wieder gegen die Voraussetzung ist, dass die Exponenten $e, e_0, \ldots, e_{\nu'-1}$ keinen gemeinsamen Theiler haben sollen. Wir haben daher den Satz:

3. \gesperrt{Ist $\mathcal{E}(r)$ irgend eine Normaleinheit des Körpers $H_m$, so lassen sich die ungerade Zahl $e$ und die ganzen Zahlen $e_0, e_1, \ldots, e_{\nu'-1}$ so bestimmen, dass}

\weq{20}{\mathcal{E}(r_\beta)^e = \pm\, \tau_\beta^{e_0}\tau_{\beta+1}^{e_1} \cdots \tau_{\beta+\nu'-1}^{e_{\nu'-1}}}

\gesperrt{wird.}

Dies Ergebniss führt zu einigen wichtigen Consequenzen:

4. \gesperrt{Ist $T$ der Regulator des Systems der $\tau_\beta$, und $T_0$ der Regulator eines Fundamentalsystems normaler Einheiten in $H_m$, d.\,h. der Minimalwerth, den der Regulator irgend eines Systems unabhängiger Normaleinheiten annehmen kann, so ist}

\weq{21}{T = C\,T_0, \quad C \equiv 1\ (\mathrm{mod}\ 2),}

\gesperrt{worin $C$ eine ungerade natürliche Zahl bedeutet.}


% ===== scan leaf 756, printed 735 =====
% !PRINTED 735
\weberpage{735}
\pdfbookmark[2]{§ 201. Fundamentalsystem von Einheiten des Körpers Hm}{bm229}

Dass nämlich $T:T_0 = C$ eine ganze Zahl ist, folgt aus der Formel (13).

Wenn wir aber den Satz 3. auf alle Elemente des Fundamentalsystems $\mathcal{E}_i(r)$ anwenden, so ergiebt sich ein System ganzer Zahlen $e_{i,k}$ und eine ungerade Zahl $e$, so dass

\[
\mathcal{E}_i(r)^e = \pm\, \tau_0^{e_{i,0}}\tau_1^{e_{i,1}} \cdots \tau_{\nu'-1}^{e_{i,\nu'-1}}.
\]

Nehmen wir hiervon die Logarithmen, und bilden dann den Regulator $T_0$, so ergiebt sich

\[
e^{\nu'} T_0 = \pm\, \Sigma \pm e_{0,0}\, e_{1,1} \cdots e_{\nu'-1,\nu'-1}\, T,
\]

folglich

\[
e^{\nu'} = \pm\, C\, \Sigma \pm e_{0,0}\, e_{1,1} \cdots e_{\nu'-1,\nu'-1},
\]

und da hierin die linke Seite ungerade ist, so kann $C$ nicht gerade sein, w. z. b. w.

Wenn alle conjugirten Werthe der Normaleinheit $\mathcal{E}(r_\beta)$ einerlei Zeichen haben, so gilt dasselbe auch von $\mathcal{E}(r_\beta)^e$, und folglich müssen in diesem Falle in der Formel (20) die Exponenten $e_0, e_1, \ldots, e_{\nu'-1}$ nach 2. gerade Zahlen sein. Schreibt man dann die Formel (20) so:

\[
\mathcal{E}(r) = \pm\, \tau_0^{e_0}\tau_1^{e_1} \cdots \tau_{\nu'-1}^{e_{\nu'-1}}\, \mathcal{E}(r)^{1-e},
\]

so erhält man daraus den Satz:

5. \gesperrt{Eine Normaleinheit des Körpers $H_m$, die mit ihren Conjugirten gleiches Vorzeichen hat, ist, vom Vorzeichen abgesehen, ein Quadrat einer Normaleinheit.}

\begin{center}
\S.\,201.

\gesperrt{Fundamentalsystem von Einheiten des Körpers $H_m$.}
\end{center}

Es bleibt noch übrig, den Nenner $D$ in der Formel für $B$ [\S.\,197, (6)] zu bestimmen, der durch die Gleichung

\[
E = E'\, D
\]

definirt war, worin $E$ und $E'$ die Regulatoren der Körper $H_m, H_\mu$ bedeuteten.

Hier handelt es sich also nicht mehr um die Normaleinheiten, sondern um die Einheiten des Körpers $H_m$ überhaupt.


% ===== scan leaf 757, printed 736 =====
% !PRINTED 736
\weberpage{736}

Wir nehmen ein Fundamentalsystem von Einheiten im Körper $H_m$ an:

\weq{1}{\varepsilon_1(r_i),\ \varepsilon_2(r_i),\ \ldots,\ \varepsilon_{\nu-1}(r_i),}

und die conjugirten Logarithmen dieses Systems seien:

\weq{2}{l_{1,k},\ l_{2,k},\ \ldots,\ l_{\nu-1,k} \quad k=0,1,2,\ldots,\nu-1.}

Daneben betrachten wir ein Fundamentalsystem im Körper $H_\mu$

\weq{3}{\varepsilon_1',\ \varepsilon_2',\ \ldots,\ \varepsilon_{\nu'-1}'}

mit den conjugirten Logarithmen

\weq{4}{l_{1,k}',\ l_{2,k}',\ \ldots,\ l_{\nu'-1,k}' \quad k=0,1,2,\ldots,\nu'-1.}

Bei der Bildung der Regulatoren wird einer der conjugirten Werthe, etwa $k=0$, weggelassen (vgl. \S.\,186), und wir erhalten daher

\weq{5}{\begin{aligned} E &= \pm\, \Sigma \pm l_{1,1} l_{2,2} \ldots l_{\nu-1,\nu-1} \\ E' &= \pm\, \Sigma \pm l_{1,1}' l_{2,2}' \ldots l_{\nu'-1,\nu'-1}'. \end{aligned}}

Nun fügen wir zu dem Systeme (3) ein Fundamentalsystem von Normaleinheiten des Körpers $H_m$

\weq{6}{\mathcal{E}_0,\ \mathcal{E}_1,\ \ldots,\ \mathcal{E}_{\nu'-1},}

und wir behaupten, dass das System

\weq{7}{\varepsilon_1',\ \varepsilon_2',\ \ldots,\ \varepsilon_{\nu'-1}',\ \mathcal{E}_0,\ \mathcal{E}_1,\ \ldots,\ \mathcal{E}_{\nu'-1}}

ein unabhängiges System von Einheiten des Körpers $H_m$ ist (wenn auch nicht ein Fundamentalsystem).

Bezeichnen wir mit

\weq{8}{L_{0,k},\ L_{1,k},\ \ldots,\ L_{\nu'-1,k} \quad k=0,1,\ldots,\nu'-1}

die conjugirten Logarithmen des Systems (6), so ist nach dem vorigen Paragraphen

\weq{9}{\pm\, \Sigma \pm L_{0,0} L_{1,1} \ldots L_{\nu'-1,\nu'-1} = T_0}

der Regulator dieses Systems, und wir haben mit Rücksicht auf die Definition der Normaleinheiten [\S.\,200, (9)]:

\weq{10}{L_{i,k+\nu'} = -L_{i,k},}

und weil die $\varepsilon_i'$ als Zahlen des Körpers $H_\mu$ durch die Substitution $(r,-r)$ ungeändert bleiben:

\weq{11}{l_{i,k+\nu'}' = l_{i,k}'.}


% ===== scan leaf 758, printed 737 =====
% !PRINTED 737
\weberpage{737}

Hiernach ergiebt sich der Regulator $R$ des Systems (7) aus
der Determinante:

\[
\left|
\begin{array}{cccccccc}
l'_{1,1}, & l'_{2,1}, & \ldots, & l'_{\nu'-1,1}, & L_{0,1}, & L_{1,1}, & \ldots, & L_{\nu'-1,1} \\
\cdots & \cdots & \cdots & \cdots & \cdots & \cdots & \cdots & \cdots \\
l'_{1,\nu'-1}, & l'_{2,\nu'-1}, & \ldots, & l'_{\nu'-1,\nu'-1}, & L_{0,\nu'-1}, & L_{1,\nu'-1}, & \ldots, & L_{\nu'-1,\nu'-1} \\
l'_{1,0}, & l'_{2,0}, & \ldots, & l'_{\nu'-1,0}, & -L_{0,0}, & -L_{1,0}, & \ldots, & -L_{\nu'-1,0} \\
l'_{1,1}, & l'_{2,1}, & \ldots, & l'_{\nu'-1,1}, & -L_{0,1}, & -L_{1,1}, & \ldots, & -L_{\nu'-1,1} \\
\cdots & \cdots & \cdots & \cdots & \cdots & \cdots & \cdots & \cdots \\
l'_{1,\nu'-1}, & l'_{2,\nu'-1}, & \ldots, & l'_{\nu'-1,\nu'-1}, & -L_{0,\nu'-1}, & -L_{1,\nu'-1}, & \ldots, & -L_{\nu'-1,\nu'-1}
\end{array}
\right|.
\]

Diese Determinante hat $\nu-1$ Zeilen; wir addiren die erste
zur $(\nu'+1)^{\mathrm{ten}}$, die zweite zur $(\nu'+2)^{\mathrm{ten}}$ u.\,s.\,f., die $(\nu'-1)^{\mathrm{te}}$ zur
letzten. Dadurch heben sich in den letzten $(\nu'-1)$ Zeilen die
$L$ heraus und die $l$ bekommen den Factor 2 [nach (10), (11)],
und hieraus ergiebt sich nach einem Determinantensatze (Bd.\,I,
\S.\,23, V.):
\weq{12}{R = 2^{\nu'-1} E' T_0.}

Da hiernach $R$ nicht verschwindet und daher in (7) ein
System von $\nu-1$ unabhängigen Einheiten des Körpers $H_m$ vorliegt,
so können wir [nach \S.\,186, (5)] die rationalen (ganzen
oder gebrochenen) Zahlen $m_{i,k}$, $M_{i,k}$ aus den Gleichungen
\weq{13}{
\begin{aligned}
2\, l_{i,k} & = m_{1,i}\, l'_{1,k} + \cdots + m_{\nu'-1,i}\, l'_{\nu'-1,k} \\
& \quad + M_{0,i}\, L_{0,k} + M_{1,i}\, L_{1,k} + \cdots + M_{\nu'-1,i}\, L_{\nu'-1,k}
\end{aligned}
}

bestimmen. Es lässt sich aber leicht durch die Relationen (10), (11)
nachweisen, dass die Zahlen $m_{i,k}$, $M_{i,k}$ hier \gesperrt{ganze} sein müssen.
Denn es ist
\[
\begin{aligned}
l_{i,k} + l_{i,k+\nu'} & = \log \lvert \varepsilon_i(r_k)\, \varepsilon_i(-r_k) \rvert = \\
& \quad m_{1,i}\, l'_{1,k} + \cdots + m_{\nu'-1,i}\, l'_{\nu'-1,k} \\
l_{i,k} - l_{i,k+\nu'} & = \log \lvert \varepsilon_i(r_k)\, \varepsilon_i(-r_k)^{-1} \rvert = \\
& \quad M_{0,i}\, L_{0,k} + M_{1,i}\, L_{1,k} + \cdots + M_{\nu'-1,i}\, L_{\nu'-1,k}.
\end{aligned}
\]

Die Einheiten $\varepsilon_i(r)\, \varepsilon_i(-r)$ gehören aber dem Körper $H_\mu$
an, und da die $\varepsilon'_i$ nach Voraussetzung ein Fundamentalsystem
dieses Körpers sind, so müssen die $m_{i,k}$ nach \S.\,187 ganze rationale
 Zahlen sein.

Die Einheiten $\varepsilon_i(r)\, \varepsilon_i(-r)^{-1}$ sind Normaleinheiten des Körpers
 $H_m$, und weil die $\mathcal{E}_i$ ein Fundamentalsystem von Normaleinheiten
 sein sollten, so sind nach \S.\,200 die $M_{i,k}$ ganze Zahlen.


% ===== scan leaf 759, printed 738 =====
% !PRINTED 738
\weberpage{738}

Zur besseren Uebersicht setzen wir die Gleichungen (13)
noch einmal ohne den Index $k$, der die conjugirten Werthe von
einander unterscheidet, hierher:
\weq{14}{
\begin{aligned}
2\, l_i & = m_{1,i}\, l'_1 + m_{2,i}\, l'_2 + \cdots + m_{\nu'-1,i}\, l'_{\nu'-1} \\
& \quad + M_{0,i}\, L_0 + M_{1,i}\, L_1 + M_{2,i}\, L_2 + \cdots + M_{\nu'-1,i}\, L_{\nu'-1}.
\end{aligned}
}

Bilden wir nun nach (13) den Regulator $E$ und bezeichnen
mit $M$ den absoluten Werth der Determinante der $\nu-1$ Reihen
der Zahlen $m_{k,i}$, $M_{k,i}$:
\[
M = \pm \Sigma \pm m_{1,1} \ldots m_{\nu'-1,\nu'-1}\, M_{0,0}\, M_{1,1} \ldots M_{\nu'-1,\nu'-1},
\]
so folgt:
\weq{15}{2^{\nu-1} E = M R,}
und daraus mit Rücksicht auf (12):
\weq{16}{E = 2^{-\nu'} M E' T_0.}

Es lässt sich nun weiter nachweisen, dass $M$ eine Potenz
von 2 und durch $2^{\nu'}$ theilbar ist.

Da nämlich das System (1) als Fundamentalsystem in $H_m$
vorausgesetzt war, so giebt es \gesperrt{ganze rationale Zahlen}
$n_{k,i}$, $N_{k,i}$, die den Gleichungen genügen:
\weq{17}{
\begin{aligned}
l'_i & = n_{1,i}\, l_1 + n_{2,i}\, l_2 + \cdots + n_{\nu-1,i}\, l_{\nu-1}, \quad i = 1, 2, \ldots, \nu'-1 \\
L_i & = N_{1,i}\, l_1 + N_{2,i}\, l_2 + \cdots + N_{\nu-1,i}\, l_{\nu-1}, \quad i = 0, 1, 2, \ldots, \nu'-1.
\end{aligned}
}

Bilden wir daraus die Determinante $R$ und bezeichnen mit $N$
den absoluten Werth der Determinante der $n_{k,i}$, $N_{k,i}$, so folgt:
\[
R = N E,
\]
und daraus nach (15):
\[
M N = 2^{\nu-1};
\]
daraus folgt, dass jede der beiden ganzen Zahlen $M$, $N$ eine
Potenz von 2 sein muss.

Um zu entscheiden, durch welche Potenz von 2 die Determinante
 $M$ theilbar ist, bestimmen wir zunächst ein System
ganzer rationaler Zahlen $a_i$ ohne gemeinschaftlichen Theiler, das
den linearen Gleichungen
\weq{18}{\sum\limits_{1,\, \nu-1}^{i} a_i\, m_{k,i} = 0 \quad (k = 1, 2, \ldots, \nu'-1)}
genügt. Setzen wir dann
\[
\sum\limits_{1,\, \nu-1}^{i} a_i\, M_{k,i} = \xi_k,
\]
%CONT>>


% ===== scan leaf 760, printed 739 =====
% !PRINTED 739
\weberpage{739}

%<<CONT
so erhalten wir aus den Gleichungen (13)
\[
2 \sum\limits_{1,\, \nu-1}^{i} a_i\, l_{i,k} = \xi_0 L_{0,k} + \xi_1 L_{1,k} + \cdots + \xi_{\nu'-1} L_{\nu'-1,k}.
\]
Daraus ergiebt sich nach (10)
\[
\sum^{i} a_i\, l_{i,k+\nu'} = -\sum^{i} a_i\, l_{i,k},
\]
und dies bedeutet, dass
\[
\delta = \varepsilon_1^{a_1}\, \varepsilon_2^{a_2} \ldots \varepsilon_{\nu-1}^{a_{\nu-1}}
\]
eine Normaleinheit des Körpers $H_m$ ist. Daraus folgt dann weiter,
weil $\mathcal{E}_0, \mathcal{E}_1, \ldots, \mathcal{E}_{\nu'-1}$ ein Fundamentalsystem von Normaleinheiten
 ist, dass die Zahlen $\xi_0, \xi_1, \ldots, \xi_{\nu'-1}$ gerade sein
müssen. Daraus ergiebt sich:

1. \gesperrt{Das System der Gleichungen (18) hat die Congruenzen
}
\weq{19}{\begin{gathered}
\sum\limits_{1,\, \nu-1}^{i} a_i\, M_{k,i} \equiv 0 \pmod{2} \\
k = 0, 1, \ldots, \nu'-1
\end{gathered}}
\gesperrt{zur Folge.}

In den Gleichungen (18), deren Anzahl nur $\nu'-1$ beträgt,
sind aber $\nu-1$ Unbekannte $a_i$ enthalten. Daher giebt es mehrere
Systeme von einander unabhängiger Lösungen, was wir etwas
genauer dahin präcisiren können:

2. \gesperrt{Es giebt $\nu'$ ganzzahlige Lösungen des Systems
(18):}
\weq{20}{
\begin{aligned}
& a_{1,1}, a_{2,1}, \ldots, a_{\nu-1,1} \\
& a_{1,2}, a_{2,2}, \ldots, a_{\nu-1,2} \\
& \ldots\ldots\ldots\ldots \\
& a_{1,\nu'}, a_{2,\nu'}, \ldots, a_{\nu-1,\nu'}
\end{aligned}
}
\gesperrt{von der Art, dass aus der Matrix (20) eine
Determinante von $\nu'$ Reihen gebildet werden
kann, die eine ungerade Zahl ist.}

Denn es giebt zunächst eine Lösung von (18):
\[
a_{1,1}, a_{2,1}, a_{3,1}, \ldots, a_{\nu-1,1},
\]
in der eine ungerade Zahl vorkommt. Nehmen wir an, was
gestattet ist, da hier auf die Anordnung der Indices nichts an-
%CONT>>


% ===== scan leaf 761, printed 740 =====
% !PRINTED 740
\weberpage{740}

%<<CONT
kommt, es sei $a_{1,1}$ diese ungerade Zahl, dann bestimmen wir
eine zweite Lösung:
\[
0, a_{2,2}, a_{3,2}, \ldots, a_{\nu-1,2},
\]
in der wieder eine ungerade Zahl, etwa $a_{2,2}$, vorkommt. Dann
bestimmen wir eine dritte Lösung:
\[
0, 0, a_{3,3}, \ldots, a_{\nu-1,3},
\]
und fahren mit Nullsetzen so lange fort, als die Anzahl der
übrig bleibenden Unbekannten noch grösser ist, als die Anzahl
der Gleichungen.

Im letzten Systeme werden also $\nu'-1$ Nullen vorkommen,
damit $\nu'$ Unbekannte übrigbleiben.

Diese Zahlenreihen $a_{\nu,k}$ können wir für das System (20)
setzen, denn darin reducirt sich die Determinante
\weq{21}{\Sigma \pm a_{1,1}\, a_{2,2} \ldots a_{\nu',\nu'}}
auf das Diagonalglied $a_{1,1}\, a_{2,2} \ldots a_{\nu',\nu'}$, dessen Factoren alle ungerade
 sind.

Ist nun die Forderung des Satzes 2. erfüllt, so lässt sich
die Matrix (20) durch Hinzufügung von $\nu'-1$ Zeilen zu einer
Determinante von $\nu-1$ Reihen ergänzen, die gleichfalls einen
ungeraden ganzzahligen Werth hat:
\weq{22}{a = \Sigma \pm a_{1,1}\, a_{2,2} \ldots a_{\nu',\nu'}\, a_{\nu'+1,\nu'+1} \ldots a_{\nu-1,\nu-1}.}

Man braucht nur, wenn die Determinante (21) als ungerade
vorausgesetzt wird, in den hinzugefügten Zeilen für die Elemente
lauter Nullen zu setzen, ausgenommen die in der Diagonalreihe
 stehenden, $a_{\nu'+1,\nu'+1} \ldots a_{\nu-1,\nu-1}$, die gleich 1 genommen
werden können. Dann erhält $a$ denselben Werth wie die Determinante
 (21).

Wenn wir aber jetzt nach dem Multiplicationssatze der
Determinanten das Product $a M$ bilden, so ergiebt sich eine
Determinante, in der die $\nu-1$ Summen
\[
\begin{aligned}
& \sum\limits_{1,\, \nu-1}^{i} a_{i,s}\, m_{k,i}, \qquad \sum\limits_{1,\, \nu-1}^{i} a_{i,s}\, M_{k,i} \\
& k = 1, 2, \ldots \nu'-1, \qquad k = 0, 1, 2 \ldots \nu'-1
\end{aligned}
\]
für jeden Werth $s = 1, 2, \ldots, \nu'$ eine Reihe bilden, und in der
also $\nu'$ Reihen vorkommen, deren Elemente alle durch 2 theilbar
 sind.


% ===== scan leaf 762, printed 741 =====
% !PRINTED 741
\weberpage{741}
Demnach ist $aM$ und folglich auch $M$ durch $2^{\nu'}$ theilbar,
und wenn wir
\weq{23}{M = 2^{\nu'+\sigma}, \quad N = 2^{\nu'-\sigma-1}}
setzen, so ist $\sigma$ eine nicht negative ganze rationale Zahl.

Setzen wir dies in (16) ein, so folgt:
\[ E = 2^\sigma E' T_0. \]
Nun war in \S.\,197, (6)
\[ E = E' D \]
gesetzt, und es ergiebt sich also:
\[ D = 2^\sigma T_0. \]
Hieraus folgt für den Classenzahlfactor $B$ nach \S.\,199, (15)
und \S.\,200, (21):
\weq{24}{B = 2^{-\sigma} \frac{T}{T_0} = 2^{-\sigma} C,}
worin $C$ eine ungerade ganze Zahl ist.

Demnach ergiebt sich nach \S.\,197, (4) für den zweiten Factor
der Classenzahl $h_1$, der, wie wir gesehen haben, eine \gesperrt{ganze
Zahl} ist:
\weq{25}{h_1 = 2^{-\sigma} h_1' C.}
Nehmen wir nun als bewiesen an, dass $h_1'$ ungerade ist, so
folgt, da auch $C$ ungerade ist, dass $\sigma = 0$ und $h_1$
ungerade ist, und beides ist also damit durch vollständige
Induction bewiesen.

Damit werden die abgeleiteten Resultate folgendermaassen
vereinfacht:

3. \gesperrt{Es ist}
\weq{26}{M = 2^{\nu'}, \quad D = T_0, \quad B = \frac{T}{T_0};}
\gesperrt{der Classenzahlfactor $B$ ist gleich dem Verhältniss
des Regulators $T$ der Einheiten $\tau$ zum Regulator $T_0$
eines Fundamentalsystems von Normaleinheiten, und ist eine
ungerade ganze Zahl.}

Bezeichnen wir die zu $m = 2^\varkappa$ gehörige Zahl $B$ mit
$B_\varkappa$, so ist die Classenzahl $h_1$, wie durch
vollständige Induction folgt,
\weq{27}{h_1 = B_4 B_5 \ldots B_\varkappa,}
und ist also auch eine ungerade Zahl. Damit ist dann, mit
Rücksicht auf \S.\,198, 1., das Haupttheorem bewiesen:

A. \gesperrt{Die Classenzahl im Körper $\Omega_m$ ist, wenn
$m$ eine Potenz von 2 ist, ungerade.}


% ===== scan leaf 763, printed 742 =====
% !PRINTED 742
\weberpage{742}
\pdfbookmark[2]{§ 202. Positive Einheiten}{bm230}

\S.\,202.

\gesperrt{Positive Einheiten.}

Die zuletzt gefundenen Resultate zeigen, dass das System
unabhängiger Einheiten \S.\,201, (7)
\[ \varepsilon'_1, \varepsilon'_2, \ldots, \varepsilon'_{\nu'-1}, \; \mathcal{E}_0, \mathcal{E}_1, \ldots, \mathcal{E}_{\nu'-1} \]
im Körper $H_m$ kein \gesperrt{Fundamentalsystem} bildet. Denn sonst
müssten in den Formeln \S.\,201, (13) die Zahlen
% [sic] -- the source omits a separator between the coefficient families $m_{k,i}$ and $M_{k,i}$; the printed form is preserved.
$m_{k,i} M_{k,i}$ sämmtlich durch 2 theilbar sein, und es wäre
die Determinante $M$, deren Werth wir $= 2^{\nu'}$ gefunden
haben, durch $2^{\nu-1}$ theilbar.

Wir betrachten nun an Stelle des Systems der Gleichungen (18),
\S.\,201 die folgenden Congruenzen:
\weq{1}{
\begin{aligned}
\sum_{1,\,\nu-1}^{i} b_i m_{1,i} &\equiv 1 \\
\sum_{1,\,\nu-1}^{i} b_i m_{2,i} &\equiv 0 \pmod{2}, \\
& \cdots\cdots\cdots\cdots \\
\sum_{1,\,\nu-1}^{i} b_i m_{\nu'-1,i} &\equiv 0
\end{aligned}
}
und beweisen, dass es möglich ist, ihnen durch ganzzahlige
Werthe der $b_i$ zu genügen.

Wäre es nicht möglich, so müsste sich aus den $\nu'-2$
Congruenzen
\[ \sum_{1,\,\nu-1}^{i} b_i m_{k,i} \equiv 0 \pmod{2} \quad k = 2, 3, \ldots, \nu'-1 \]
die letzte
\[ \sum_{1,\,\nu-1}^{i} b_i m_{1,i} \equiv 0 \pmod{2} \]
als \gesperrt{Folgerung} ergeben, und dann würde, wie im vorigen
Paragraphen, weiter folgen:
\[ \sum_{1,\,\nu-1}^{i} b_i M_{k,i} \equiv 0 \pmod{2} \quad k = 0, 1, \ldots, \nu'-1. \]


% ===== scan leaf 764, printed 743 =====
% !PRINTED 743
\weberpage{743}
Man könnte dann an Stelle der Matrix \S.\,201, (20) eine Matrix
der $b_{k,i}$ von $\nu'+1$ Reihen setzen, und es würde sich auf
dem im vorigen Paragraphen eingeschlagenen Wege schliessen
lassen, dass $M$ durch $2^{\nu'+1}$ theilbar sein müsste, was
nicht der Fall ist.

Wenden wir aber das System der Multiplicatoren $b_i$, das den
Bedingungen (1) genügt, auf die Gleichungen (14) des vorigen
Paragraphen an, indem wir mit $b_i$ multipliciren und dann
summiren, so ergiebt sich, wenn wir Vielfache von 2 weglassen
und dies auch hier (wo irrationale Zahlen, nämlich die
Logarithmen, auftreten) durch das Zeichen der Congruenz
ausdrücken:
\weq{2}{
\begin{aligned}
l'_1 \equiv{} & L_0 \sum M_{0,i} b_i + L_1 \sum M_{1,i} b_i + \cdots \\
& + L_{\nu'-1} \sum M_{\nu'-1,i} b_i \pmod{2}.
\end{aligned}
}
Wollen wir die Congruenz in eine Gleichung verwandeln, so kommt
eine Summe von Grössen
\[ l_1, l_2, \ldots, l_{\nu-1}, \quad l'_1, l'_2, \ldots, l'_{\nu'-1} \]
mit \gesperrt{geraden ganzen Zahlen} als Co\"efficienten hinzu.

Dieselbe Betrachtung lässt sich aber auf $l'_2, l'_3, \ldots,
l'_{\nu'-1}$ anwenden, und danach kann man, wenn man von den
Logarithmen wieder zu den Zahlen übergeht, die Formel (2) so
in Worte fassen:

1. \gesperrt{Jede Einheit des Körpers $H_\mu$ ist als Product
einer Normaleinheit des Körpers $H_m$ und eines Quadrates
einer Einheit in $H_m$ darstellbar.}

Wenn wir darauf den Satz 2., \S.\,200 anwenden, so ergiebt
sich, dass jede Einheit $\mathcal{E}(r)$ des Körpers $H_\mu$,
die mit allen ihren Conjugirten dasselbe Zeichen hat, vom
Vorzeichen abgesehen, das Quadrat einer Einheit in $H_m$ sein
muss. Ist aber
\[ \pm\, \mathcal{E}(r) = \Delta(r)^2, \]
worin $\mathcal{E}(r)$ eine Einheit in $H_\mu$ und $\Delta(r)$
eine Einheit in $H_m$ ist, so muss auch $\Delta(r)$ in $H_\mu$
enthalten sein.

Denn setzen wir
\[ \Delta(r) = \Delta_1(r^2) + r \Delta_2(r^2), \]
so kann das Quadrat davon nur dann in $H_\mu$ enthalten sein,
also durch die Vorzeichenänderung von $r$ ungeändert bleiben,
wenn entweder $\Delta_1$ oder $\Delta_2 = 0$ ist. Es kann aber
nicht $\Delta(r) = r \Delta_2(r^2)$ sein, denn sonst wäre
$\Delta_2(r^2)$ eine Einheit in $\Omega_\mu$ und müsste
%CONT>>


% ===== scan leaf 765, printed 744 =====
% !PRINTED 744
\weberpage{744}
%<<CONT
nach \S.\,178, 1. das Product einer reellen Einheit und einer
Einheitswurzel $r^{2\lambda}$ sein, und da auch $\Delta(r)$
reell ist, so müsste $r^{2\lambda+1}$ reell sein, was nicht
möglich ist. Demnach ist $\Delta(r)$ in $H_\mu$ enthalten.

Da nun $\mu$ ebenso wie $m$ jede beliebige Potenz von 2 sein
kann, so ergiebt sich hieraus das zweite Haupttheorem:

B. \gesperrt{Eine Einheit des Körpers $H_m$, die mit allen
ihren Conjugirten positiv ist, ist das Quadrat einer Einheit
desselben Körpers.}

Von den beiden Theoremen A. und B. haben wir aber im \S.\,184
den Beweis des Hauptsatzes von den \gesperrt{Abel}'schen Zahlkörpern
(\S.\,179, I.) abhängig gemacht.

\begin{center}
\rule{4cm}{0.4pt}
\end{center}


% ===== scan leaf 766, printed 745 =====
\weberpage{745}% page number not printed on this page
\pdfbookmark[1]{Fünfundzwanzigster Abschnitt. Transcendente Zahlen.}{bm231}
\pdfbookmark[2]{§ 203. Abzählbare Mengen}{bm232}
% !PRINTED none
\weberpage{none}

\begin{center}
\gesperrt{Fünfundzwanzigster Abschnitt.}

\gesperrt{Transcendente Zahlen.}
\end{center}

\begin{center}
\rule{4cm}{0.4pt}
\end{center}

\begin{center}
\S.\,203.

\gesperrt{Abzählbare Mengen.}
\end{center}

In der Einleitung zu unserem Werke ist der allgemeine Begriff der Zahlen definirt worden, und mit diesem allgemeinen Zahlbegriffe haben wir zunächst, z.\,B. bei dem Beweise der Wurzelexistenz, operirt. Im weiteren Verlaufe unserer Betrachtungen haben wir uns dann nur mit \gesperrt{algebraischen Zahlen} beschäftigt, ohne uns darüber Rechenschaft zu geben, ob damit der Umfang des Zahlenbereiches erschöpft sei, oder ob es auch nichtalgebraische Zahlen giebt. Die Existenz von nichtalgebraischen Zahlen, die man auch transcendente Zahlen nennt, ist zuerst von \gesperrt{Liouville} bewiesen. Andere Beweise für diesen Satz rühren von \gesperrt{G. Cantor} her\footnote{\gesperrt{G. Cantor}, Crelle's Journal, Bd. 77 (1873). Ueber eine Eigenschaft des Inbegriffes aller reellen algebraischen Zahlen.}.

Wir knüpfen hier an den schon in der Einleitung aus einander gesetzten Begriff einer \gesperrt{Menge} oder \gesperrt{Mannigfaltigkeit} an, worunter ein System von Elementen irgend welcher Art zu verstehen ist, was so genau abgegrenzt ist, dass von jedem beliebigen Objecte völlig entschieden ist, ob es zu dem Systeme gehört oder nicht.

Wir unterscheiden endliche und unendliche Mengen, und führen als erstes und wichtigstes Beispiel einer unendlichen Menge die Gesammtheit der \gesperrt{natürlichen Zahlen} $1, 2, 3, \ldots$ an. Dann gilt folgende Definition:

1. \textbf{Definition.} \gesperrt{Eine Menge heisst abzählbar, wenn ihre Elemente mit der natürlichen Zahlenreihe}
%CONT>>


% ===== scan leaf 767, printed 746 =====
% !PRINTED 746
\weberpage{746}
%<<CONT
\gesperrt{oder einem Theil derselben in eine gegenseitig eindeutige Beziehung gesetzt werden kann}\footnote{\gesperrt{Cantor}, l.\,c. Der Begriff der abzählbaren Mengen deckt sich mit dem von \gesperrt{Dedekind} ohne die \gesperrt{Voraussetzung des Zahlensystems} definirten Begriffe der \gesperrt{einfach unendlichen Systeme}. (\gesperrt{Dedekind}, Was sind und was sollen die Zahlen? \S.\,6. Braunschweig 1887. Zweite unveränderte Auflage 1893.)}.

Jede endliche Menge ist hiernach abzählbar, und bei der Abzählung wird die natürliche Zahlenreihe nur bis zu einer gewissen höchsten Zahl verwendet. Im Weiteren betrachten wir vorzugsweise unendliche Mengen.

Wir können der Definition für eine unendliche abzählbare Menge auch den Ausdruck geben, dass es eine solche Menge ist, bei der jedem Elemente eine bestimmte Zahl der natürlichen Zahlenreihe als Name beigelegt werden kann, so dass auch jede Zahl der Zahlenreihe dabei verwandt wird, also eine Menge, in der es ein erstes, zweites, drittes \ldots hundertstes \ldots Element giebt.

Endlich können wir auch sagen, dass eine unendliche abzählbare Menge eine solche ist, die sich derart in eine Reihe ordnen lässt, dass ein \gesperrt{erstes} Element vorhanden ist, und dass auf jedes Element ein \gesperrt{bestimmtes anderes der Menge} folgt, und jedem Elemente, mit Ausnahme des ersten, ein \gesperrt{bestimmtes anderes Element vorangeht}. Eine solche Reihe können wir eine zählbare Anordnung nennen.

Es ist klar, dass eine abzählbare Menge nicht nur auf eine, sondern auf unendlich viele verschiedene Arten abzählbar ist.

Ausser der natürlichen Zahlenreihe selbst, die sicher abzählbar ist, können wir als zweites Beispiel das System der rationalen positiven echten Brüche anführen, die unter Anderem in folgender Weise abgezählt werden können:
\[
\frac{1}{2}, \frac{1}{3}, \frac{2}{3}, \frac{1}{4}, \frac{3}{4}, \frac{1}{5}, \frac{2}{5}, \frac{3}{5}, \frac{4}{5}, \frac{1}{6}, \frac{5}{6}, \ldots
\]
d.\,h. so, dass jeder grössere Nenner dem kleineren Nenner folgt, und dass die Brüche von gleichem Nenner nach der Grösse des Zählers geordnet sind. Wollte man aber die Brüche nach der Grösse ihres numerischen Werthes ordnen, so würde man keine zählbare Anordnung bekommen.

2. \gesperrt{Die Gesammtheit aller algebraischen Zahlen ist eine abzählbare Menge.}


% ===== scan leaf 768, printed 747 =====
% !PRINTED 747
\weberpage{747}
Um diesen wichtigen Satz nachzuweisen, erinnern wir uns daran, dass jede algebraische Zahl $\Theta$ die Wurzel einer und nur einer irreduciblen Gleichung
\weq{1}{f(\Theta) = a_0 \Theta^n + a_1 \Theta^{n-1} + \cdots + a_n = 0}
ist, in der $a_0, a_1, \ldots, a_n$ ganze rationale Zahlen ohne gemeinsamen Theiler sind, und $a_0$ von Null verschieden und positiv ist. Der Grad $n$ der Gleichung (1) ist eine positive ganze Zahl, also mindestens $=1$. Die Zahlen $a_1, \ldots, a_{n-1}$ können zum Theil Null sein.

Wir wollen nun die Vorzeichen $\pm$ so bestimmen, dass $\pm a_1, \pm a_2, \pm a_3, \ldots, \pm a_n$ nicht negativ sind, und nennen die Summe
\weq{2}{N = (n-1) + a_0 \pm a_1 \pm a_2 \cdots \pm a_n}
die \gesperrt{Höhe} der algebraischen Zahl $\Theta$. Die Höhe ist dann immer eine \gesperrt{positive} ganze Zahl.

Nun ist leicht zu sehen, dass es für einen gegebenen Werth der Höhe $N$ immer nur eine endliche Anzahl von algebraischen Zahlen geben kann. Denn zunächst kann $n$ nach (2) niemals grösser als $N$ sein, und zu jedem gegebenen $N$ und $n$ können die Zahlen $a_0, a_1, \ldots, a_n$ nur auf eine endliche Anzahl von Arten bestimmt werden. Man hat dann unter den so bestimmten Functionen $f(\Theta)$ nur die irreduciblen beizubehalten. Wenn man nun die algebraischen Zahlen in der Weise ordnet, dass man die Zahlen von geringerer Höhe denen von grösserer Höhe voranstellt, dass man unter Zahlen gleicher Höhe die voranstellt, deren reeller Theil kleiner ist, und unter Zahlen von gleicher Höhe und gleichem reellen Theile die von kleinerem imaginären Theile vorangehen lässt, so haben wir eine zählbare Anordnung der algebraischen Zahlen, und es ist erwiesen, dass die Gesammtheit aller algebraischen Zahlen eine abzählbare Menge ist.

Wir erhalten beispielsweise:
\[
\begin{array}{llll}
N=1, & n=1, & a_0=1, & a_1=0\\
N=2, & n=1, & a_0=1, & a_1=\pm 1\\
N=3, & n=1, & a_0=1, & a_1=\pm 2\\
     &      & a_0=2, & a_1=\pm 1\\
     & n=2, & a_0=1, & a_1=0, \quad a_2=1,
\end{array}
\]
und der Anfang der geordneten Reihe der algebraischen Zahlen wird also:
\[
0, -1, +1, -2, -\tfrac{1}{2}, -\sqrt{-1}, \sqrt{-1}, \tfrac{1}{2}, 2, \ldots
\]


% ===== scan leaf 769, printed 748 =====
% !PRINTED 748
\weberpage{748}
\pdfbookmark[2]{§ 204. Unzählbare Mengen}{bm233}
Jeder Theil der natürlichen Zahlenreihe ist eine abzählbare Menge; denn man braucht ja die Zahlen eines solchen Theiles nur nach ihrer Grösse zu ordnen, um eine zählbare Anordnung zu erhalten.

Daraus aber ergiebt sich, dass jeder Theil einer abzählbaren Menge selbst eine abzählbare Menge ist. Es folgt daraus unter Anderem, dass auch die Menge der \gesperrt{reellen algebraischen} Zahlen abzählbar ist.

\begin{center}
\S.\,204.

\gesperrt{Unzählbare Mengen.}
\end{center}

Wir kommen nun zweitens zu dem Beweise des Satzes, dass es unter den Zahlenmengen auch nicht abzählbare giebt, und wir werden speciell nachweisen:

\begin{center}
\gesperrt{Dass die Gesammtheit aller reellen Zahlen, selbst wenn wir uns auf ein endliches Intervall beschränken, nicht abzählbar ist.}
\end{center}

Wir betrachten zu diesem Zwecke irgend eine abzählbare Menge reeller von einander verschiedener Zahlen, die wir, in eine zählbare Anordnung $\Omega$ gesetzt, so bezeichnen wollen:
\weq{$\Omega$}{\omega_1, \omega_2, \omega_3, \omega_4, \ldots}
(wobei daran zu erinnern ist, dass dies nicht etwa eine Aufeinanderfolge nach der Grösse bedeuten soll).

Der Kürze wegen wollen wir von zwei Elementen der Reihe $\Omega$ das mit kleinerem Index das \gesperrt{frühere}, das mit grösserem das \gesperrt{spätere} nennen.

Wir nehmen nun irgend zwei reelle Zahlen $\alpha, \beta$ an, so dass $\alpha < \beta$ ist, und zeigen, dass es in dem Intervalle $\delta = (\alpha, \beta)$ mindestens eine Zahl giebt, die nicht in der Reihe $\Omega$ vorkommt. Haben wir eine solche Zahl für jedes Intervall nachgewiesen, so giebt es auch deren unendlich viele, da man ja dieselbe Schlussweise auf jeden Theil des Intervalles $(\alpha, \beta)$ anwenden kann.

Zunächst ist klar, dass unsere Behauptung richtig ist, wenn in irgend einem endlichen Theile des Intervalles $\delta$ nur eine endliche Anzahl von Zahlen der Reihe $\Omega$ enthalten ist, und wir können also ohne Weiteres zu dem Falle übergehen, dass in jedem noch so kleinen Theile des Intervalles $\delta$ eine unendliche Menge von Zahlen der Reihe $\Omega$ liegt.


% ===== scan leaf 770, printed 749 =====
% !PRINTED 749
\weberpage{749}

% [sic] -- the source uses $\delta_1$ and $\delta_2$ both for scalar lengths and for intervals; the printed symbols are preserved.
Wir bezeichnen mit $\alpha_1$, $\beta_1$ die beiden \gesperrt{fr\"uhesten} Zahlen der Reihe $\Omega$, die in dem Intervalle $\delta$ liegen, nehmen $\alpha_1 < \beta_1$ an, und setzen $\delta_1 = \beta_1 - \alpha_1$, so dass $\delta_1 = (\alpha_1, \beta_1)$ ein Theil des Intervalles $\delta$ ist.

Nun bezeichnen wir ebenso mit $\alpha_2$, $\beta_2$ die beiden fr\"uhesten Zahlen von $\Omega$, die \gesperrt{im Inneren} des Intervalles $\delta_1$ (mit Ausschluss der Grenzen) liegen, setzen $\alpha_2 < \beta_2$ voraus, und setzen $\delta_2 = \beta_2 - \alpha_2$.

Auf diese Weise k\"onnen wir fortfahren und erhalten eine unbegrenzte Reihe von Intervallen:
\[
\delta,\ \delta_1,\ \delta_2,\ \delta_3,\ \dots,
\]
deren jedes alle folgenden einschliesst, und zwei Reihen von Zahlen:
\[
\begin{gathered}
\alpha,\ \alpha_1,\ \alpha_2,\ \dots \\
\beta,\ \beta_1,\ \beta_2,\ \dots
\end{gathered}
\]
die, mit etwaiger Ausnahme der ersten, $\alpha$ und $\beta$, alle der Reihe $\Omega$ angeh\"oren. Die $\alpha, \alpha_1, \alpha_2, \dots$ bilden eine wachsende, die $\beta, \beta_1, \beta_2, \dots$ eine fallende Zahlenreihe, und zugleich ist jedes $\alpha$ kleiner als jedes $\beta$.

1. \gesperrt{Daraus ergiebt sich, dass die Zahlen $\alpha_\nu$ eine obere Grenze $a$, die Zahlen $\beta_\nu$ eine untere Grenze $b$ haben, und dass $a$ jedenfalls nicht gr\"osser als $b$ ist.} (Vgl. Bd. I, \S.\,38, 1.)

Es kann aber m\"oglicherweise $a = b$ sein.

Aus der Bildungsweise der Intervalle $\delta, \delta_1, \delta_2, \dots$ geht noch Folgendes hervor:

2. \gesperrt{Wenn irgend eine Zahl $\omega$ der Reihe $\Omega$ in dem Intervalle $\delta_\nu$, mit Ausschluss seiner Grenzen $\alpha_\nu, \beta_\nu$, liegt, so ist $\omega$ in der Reihe $\Omega$ sp\"ater als das derselben Reihe angeh\"orige Zahlenpaar $\alpha_\nu, \beta_\nu$.}

Denn $\alpha_\nu, \beta_\nu$ waren ja die beiden \gesperrt{fr\"uhesten} im Intervalle $\delta_{\nu-1}$ gelegenen Zahlen von $\Omega$. Daraus folgt:

3. \gesperrt{Die der Reihe $\Omega$ angeh\"origen Zahlenpaare $\alpha_\nu, \beta_\nu$ sind um so sp\"atere Glieder der Reihe $\Omega$, je gr\"osser der Index $\nu$ ist, und da wir angenommen haben, die Reihe der Intervalle $\delta_\nu$ breche nicht ab, so k\"onnen wir $\alpha_\nu, \beta_\nu$ f\"ur ein hinl\"anglich grosses $\nu$ beliebig weit in der Reihe $\Omega$ hinausr\"ucken lassen.}


% ===== scan leaf 771, printed 750 =====
% !PRINTED 750
\weberpage{750}

Nun ergiebt sich sehr einfach, dass keine Zahl $g$, die mit einer der Zahlen $a, b$ zusammenf\"allt, oder auch, wenn $a$ und $b$ verschieden sind, zwischen ihnen liegt, zu der Reihe $\Omega$ geh\"oren kann.

% [sic] -- the source prints singular ``kommt'' with the pair $\alpha_\nu,\beta_\nu$; the printed German is preserved.
Denn die Zahl $g$ liegt im Inneren eines jeden der Intervalle $\delta_\nu$. Nehmen wir an, es komme $g$ in $\Omega$ vor, und gehen nach 3. mit $\nu$ so weit, dass $\alpha_\nu, \beta_\nu$ in $\Omega$ sp\"ater als $g$ kommt, so kann $g$ nach 2. nicht im Intervall $\delta_\nu$ liegen, und damit ist die Unm\"oglichkeit unserer Annahme erwiesen.

\gesperrt{Daraus folgt also, dass die Gesammtheit der Zahlen eines Intervalles $\alpha, \beta$ keine abz\"ahlbare Menge bildet.}

\begin{center}
\rule{4cm}{0.4pt}
\end{center}

Diese Thatsache l\"asst sich noch auf einem anderen Wege beweisen, der in gewisser Beziehung noch einfacher ist und den wir mit wenig Worten darlegen wollen. Wir beschr\"anken die Allgemeinheit nicht, wenn wir das Intervall von 0 bis 1 zu Grunde legen. Alle Zahlen dieses Intervalles denken wir uns durch unendliche Decimalbr\"uche dargestellt. Darunter sind auch die endlichen Decimalbr\"uche enthalten, wenn wir alle Ziffern, von einer gewissen an, gleich Null setzen. Um die Darstellung durch Decimalbr\"uche zu einer eindeutigen zu machen, mag noch festgesetzt sein, dass f\"ur einen endlichen Decimalbruch \gesperrt{immer} diese Darstellung gew\"ahlt, also z.\,B. nicht 0,4999\dots\ f\"ur 0,5000\dots\ gesetzt werden soll.

Wir wollen nun annehmen, diese Decimalbr\"uche bilden eine abz\"ahlbare Menge; sie lassen sich also in eine z\"ahlbare Reihe anordnen, die wir so darstellen:
\weq{$\Omega$}{
\begin{aligned}
\omega_1 &= 0,\ \alpha_1^{(1)}\, \alpha_2^{(1)}\, \alpha_3^{(1)}\, \dots\\
\omega_2 &= 0,\ \alpha_1^{(2)}\, \alpha_2^{(2)}\, \alpha_3^{(2)}\, \dots\\
\omega_3 &= 0,\ \alpha_1^{(3)}\, \alpha_2^{(3)}\, \alpha_3^{(3)}\, \dots\\
&\quad\cdots\cdots\cdots\cdots\cdots
\end{aligned}
}
worin die $\alpha_\mu^{(\nu)}$ Ziffern des dekadischen Systems bedeuten.

Es ist nun aber sehr leicht, einen Decimalbruch (oder auch beliebig viele) nachzuweisen, die in der Reihe $\Omega$ nicht enthalten sind. Wir brauchen nur
\[
\eta = 0,\ \beta_1\ \beta_2\ \beta_3\ \dots
\]
% [sic] -- the printed condition permits a recurring-9 expansion equivalent to an earlier terminating decimal; one repair is to choose $\beta_\nu\in\{1,2\}$ unequal to $\alpha_\nu^{(\nu)}$.
zu bilden, wobei die $\beta_\nu$ Ziffern des dekadischen Systems sind, die der einen Bedingung gen\"ugen, dass $\beta_\nu$ f\"ur jedes $\nu$ von $\alpha_\nu^{(\nu)}$ ver-
%CONT>>


% ===== scan leaf 772, printed 751 =====
% !PRINTED 751
\weberpage{751}
\pdfbookmark[2]{§ 205. Transcendenz der Zahl e}{bm234}

%<<CONT
schieden ist. Diese Zahl $\eta$, die doch auch dem Intervalle $(0, 1)$ angeh\"ort, kann mit keiner Zahl der Reihe $\Omega$ \"ubereinstimmen.

Man kann die Bildung von $\eta$ noch dadurch verallgemeinern, dass man die ersten $\beta$ bis zu einem beliebig weit entfernten willk\"urlich annimmt, und erst von da an das Gesetz: $\beta_\nu$ verschieden von $\alpha_\nu^{(\nu)}$, gelten l\"asst.

Da nun bewiesen ist, dass die reellen algebraischen Zahlen eine abz\"ahlbare Menge bilden, und dass es in jedem Intervalle Zahlen giebt, die einer gegebenen abz\"ahlbaren Menge nicht angeh\"oren, so folgt hieraus ganz unmittelbar:

\begin{center}
\gesperrt{Es giebt in jedem reellen Intervalle transcendente Zahlen.}
\end{center}

\begin{center}
\S.\,205.

\gesperrt{Transcendenz der Zahl $e$.}
\end{center}

Eine weit schwierigere Aufgabe ist es nun, von einer bestimmt vorgelegten Zahl zu entscheiden, ob sie algebraisch oder transcendent ist. Hier hat sich das Interesse haupts\"achlich auf die beiden in der Analysis so h\"aufig vorkommenden Zahlen $e$, d.\,h. die Basis des nat\"urlichen Logarithmensystems, und die \gesperrt{Ludolph}'sche Zahl $\pi$, das Verh\"altniss des Kreisumfanges zum Durchmesser, concentrirt.

F\"ur die Zahl $e$ ist die Frage von \gesperrt{Hermite} in einer ber\"uhmten Abhandlung entschieden, die f\"ur die sp\"ateren Untersuchungen \"uber die Zahl $\pi$ die Grundlage geworden ist$^{1)}$. Die Entscheidung f\"ur die Zahl $\pi$, die wegen ihrer Beziehung zu dem altber\"uhmten Problem der Quadratur des Kreises ganz besonders interessant war, bot aber noch lange Zeit un\"uberwindliche Schwierigkeiten. Endlich ist von \gesperrt{Lindemann} der Nachweis gef\"uhrt, dass auch $\pi$ zu den transcendenten Zahlen geh\"ort. Der von \gesperrt{Lindemann} gegebene Beweis bot aber dem Verst\"andniss zun\"achst noch grosse Schwierigkeiten, die durch sp\"atere Untersuchungen von \gesperrt{Weierstrass}, \gesperrt{Hilbert}, \gesperrt{Hurwitz} und \gesperrt{Gordan}$^{2)}$

\noindent\rule{4cm}{0.4pt}

$^{1)}$ \gesperrt{Hermite}, Sur la fonction exponentielle. Comptes rendus, T.\,LXXVII, 1873.

$^{2)}$ \gesperrt{Lindemann}, Ueber die Zahl $\pi$. Mathem. Annalen, Bd. 20, 1882. \gesperrt{Weierstrass}, Zu \gesperrt{Lindemann}'s Abhandlung \glqq Ueber die \gesperrt{Ludolph}'sche
%CONT>>


% ===== scan leaf 773, printed 752 =====
% !PRINTED 752
\weberpage{752}

%<<CONT
allm\"ahlich so v\"ollig beseitigt sind, dass sich der Beweis jetzt mit ganz elementaren Mitteln und auf die einfachste Weise f\"uhren l\"asst.

Wenn wir, wie schon fr\"uher, mit $\Pi(n)$ das Product
\[
\Pi(n) = 1.2.3\dots n
\]
bezeichnen, so wird, wenn $x$ eine beliebige Gr\"osse ist,
\weq{1}{\frac{x^n}{\Pi(n)}}
mit unendlich wachsendem $n$ sich der Grenze Null n\"ahern, und zwar st\"arker, als die Glieder einer fallenden geometrischen Reihe.

Denn ist $k$ eine ganze Zahl, die gr\"osser ist als der absolute Werth von $x$, so ist der absolute Werth von $x : h$ immer dann ein echter Bruch, wenn $h$ gleich oder gr\"osser als $k$ ist. Folglich ist
\[
\left|\frac{x^n}{\Pi(n)}\right| = \left|\frac{x^k}{\Pi(k)}\, \frac{x}{k+1}\, \frac{x}{k+2} \cdots \frac{x}{n}\right| < \left|\frac{x^k}{\Pi(k)}\right| \left|\frac{x}{k}\right|^{n-k}.
\]
% [sic] -- the strict inequality tacitly assumes $n>k$ and $x\ne0$; $x=0$ is the trivial equality and convergence case.
Hieraus folgt, dass die unendliche Reihe
\weq{2}{e^x = 1 + x + \frac{x^2}{\Pi(2)} + \frac{x^3}{\Pi(3)} + \cdots}
f\"ur alle Werthe von $x$ convergirt, und durch sie definiren wir die Exponentialfunction $e^x$, deren einfachste Eigenschaften wir hier aus der Analysis voraussetzen. Insbesondere gilt f\"ur zwei beliebige Werthe $x, y$ die Relation
\[
e^{x+y} = e^x e^y,
\]
aus der dann folgt, dass $e^n$ f\"ur ein ganzes positives $n$ die $n^{\text{te}}$ Potenz der Zahl
\weq{3}{e = 2 + \frac{1}{\Pi(2)} + \frac{1}{\Pi(3)} + \frac{1}{\Pi(4)} + \cdots = 2{,}718281828459\cdots}
ist.

Wenn nun $r$ irgend eine ganze positive Zahl bedeutet, so k\"onnen wir die Formel (2) so darstellen:
\weq{4}{\Pi(r)\, e^x = \Pi(r) + \frac{\Pi(r)}{\Pi(1)} x + \frac{\Pi(r)}{\Pi(2)} x^2 + \cdots + x^r + x^r\, U_r,}

\noindent\rule{4cm}{0.4pt}

Zahl\grqq{}. Sitzungsbericht der Berliner Akademie, 3. December 1885. Die Arbeiten von \gesperrt{Hilbert}, \gesperrt{Hurwitz} und \gesperrt{Gordan} finden sich alle drei in Band 43 der Mathematischen Annalen (1893), die beiden ersten auch in den G\"ottinger Nachrichten von 1893.
%CONT>>


% ===== scan leaf 774, printed 753 =====
% !PRINTED 753
\weberpage{753}
%<<CONT
worin

\weq{5}{U_r = \frac{x}{r+1}+\frac{x^2}{(r+1)(r+2)}+\frac{x^3}{(r+1)(r+2)(r+3)}+\cdots,}

und diese Formel gilt auch für $r=0$, wenn man $\Pi(0)=1$ annimmt.

Bezeichnen wir den absoluten Werth von $x$ mit $\xi$, so ist nach dem in der Einleitung bewiesenen Satze, dass der absolute Werth einer Summe niemals grösser ist als die Summe der absoluten Werthe ihrer Theile, der absolute Werth von $U_r$ gewiss nicht grösser als
\[
\frac{\xi}{r+1}+\frac{\xi^2}{(r+1)(r+2)}+\frac{\xi^3}{(r+1)(r+2)(r+3)}+\cdots,
\]
also kleiner als
\[
1+\frac{\xi}{1}+\frac{\xi^2}{1.2}+\frac{\xi^3}{1.2.3}+\cdots=e^{\xi}.
\]
Wenn wir daher
\[
U_r = q_r\,e^{\xi}
\]
setzen, so ist $q_r$ eine Function von $x$, von der wir nur so viel feststellen, dass ihr absoluter Werth kleiner als 1 ist.

Hiernach stellen wir die Formel (4) für $r=0,1,\ldots,n$ auf, und schreiben die Glieder, in umgekehrter Reihenfolge geordnet, übersichtlich so:
\[
\begin{aligned}
\Pi(n)\,e^x &= q_n\,x^n e^{\xi} + x^n + n\,x^{n-1} + n(n-1)x^{n-2}+\cdots\\
&\qquad{}+\Pi(n)\\[4pt]
(6)\qquad \Pi(n-1)\,e^x &= q_{n-1}\,x^{n-1}e^{\xi}+x^{n-1}+(n-1)x^{n-2}\\
&\qquad{}+(n-1)(n-2)x^{n-3}+\cdots\\[4pt]
&\hphantom{\Pi(n)\,e^x{}}\;\cdots\cdots\cdots\cdots\cdots\cdots\cdots\cdots\cdots\cdots\cdots\cdots\cdots\\[4pt]
e^x &= q_0\,e^{\xi}+1.
\end{aligned}
\]

Nun nehmen wir eine beliebige ganze Function $n^{\text{ten}}$ Grades von $x$ an:

\weq{7}{f(x) = c_n x^n + c_{n-1}x^{n-1} + c_{n-2}x^{n-2}+\cdots+c_0,}

und ihre Derivirten
\[
\begin{aligned}
f'(x) &= n\,c_n x^{n-1} + (n-1)c_{n-1}x^{n-2}+\cdots\\[4pt]
(8)\qquad f''(x) &= n(n-1)c_n x^{n-2}+(n-1)(n-2)c_{n-1}x^{n-3}+\cdots\\[4pt]
&\hphantom{f''(x){}}\;\cdots\cdots\cdots\cdots\cdots\cdots\cdots\cdots\cdots\cdots\cdots\cdots\cdots\\[4pt]
f^{(n)}(x) &= \Pi(n)\,c_n,
\end{aligned}
\]

und setzen noch zur Abkürzung

\weq{9}{F(x) = f(x)+f'(x)+f''(x)+\cdots+f^{(n)}(x).}
% signature "Weber, Algebra. II." and gathering-mark "48" at foot -- not text, omitted


% ===== scan leaf 775, printed 754 =====
% !PRINTED 754
\weberpage{754}
Es ist dann $F(x)$ eine ganze Function von $x$ vom Grade $n$. Setzen wir endlich noch

\weq{10}{Q(x) = c_n q_n x^n + c_{n-1}q_{n-1}x^{n-1}+\cdots+c_0 q_0,}

\weq{11}{P = c_n\Pi(n)+c_{n-1}\Pi(n-1)+\cdots+c_0,}

so ist $Q(x)$ von $x$ abhängig, wenn auch nicht rational durch $x$ ausdrückbar; $P$ aber ist von $x$ unabhängig.

Wenn wir nun die Gleichungen (6) der Reihe nach mit $c_n, c_{n-1},\ldots, c_0$ multipliciren und addiren, so folgt

\weq{12}{e^x P = F(x) + e^{\xi}Q(x).}

Diese Formel ist nun der Ausgangspunkt der weiteren Schlüsse.

Nehmen wir an, es sei $e$ eine algebraische Zahl, so muss eine Gleichung, deren Grad $m$ sei, bestehen:

\weq{13}{C_0 + C_1 e + C_2 e^2 + \cdots + C_m e^m = 0,}

deren Coëfficienten $C_0, C_1, \ldots, C_m$ \gesperrt{ganze rationale Zahlen} sind, von denen $C_0$ und $C_m$ von Null verschieden sind. Es handelt sich darum, die Unmöglichkeit dieser Annahme darzuthun.

Zu diesem Zwecke setzen wir in der Gleichung (12) für $x$ der Reihe nach die ganzen Zahlen $0,1,2,\ldots,m$, so dass $\xi$ mit $x$ identisch wird, multipliciren mit $C_0, C_1, \ldots, C_m$ und addiren. Dann ergiebt sich nach (13), da $P$ von $x$ unabhängig ist:

\weq{14}{0 = C_0 F(0) + C_1 F(1) + \cdots + C_m F(m) + C_0 Q(0) + C_1 e\,Q(1) + \cdots + C_m e^m Q(m),}

und nun soll aus einer passenden Annahme über die noch willkürliche Function $f(x)$ nachgewiesen werden, dass die Gleichung (14) unmöglich ist.

Wir wählen eine Primzahl $p$, die grösser ist als $m$\,\footnote{Dass es immer eine Primzahl $p$ giebt, die grösser ist, als eine beliebige Zahl $\mu$, ist schon bei \gesperrt{Euklid} bewiesen (Elemente, Buch IX, Nr. XX, Bd. 2 der \gesperrt{Heiberg}'schen Ausgabe). Der Beweis ist einfach der, dass die ganze Zahl $\Pi(\mu)+1$, die offenbar grösser als $\mu$ ist, durch keine Primzahl theilbar ist, die nicht grösser als $\mu$ ist, weil diese Zahl bei der Theilung durch jede der Zahlen $2,3,\ldots,\mu$ den Rest 1 ergiebt. Es ist also unmöglich, dass keine Primzahl über $\mu$ liegt.}, und setzen

\weq{15}{f(x) = \frac{x^{p-1}(1-x)^p(2-x)^p\ldots(m-x)^p}{\Pi(p-1)},}
%CONT>>


% ===== scan leaf 776, printed 755 =====
% !PRINTED 755
\weberpage{755}
%<<CONT
so dass der Grad $n$ von $f(x)$ gleich $(m+1)p-1$ ist, und wir beweisen nun zweierlei:

1) $C_0 F(0) + C_1 F(1) + \cdots + C_m F(m)$ ist eine von Null verschiedene ganze Zahl, also, vom Zeichen abgesehen, mindestens gleich 1,

2) $C_0 Q(0) + C_1 e\,Q(1) + \cdots + C_m e^m Q(m)$ ist kleiner als 1,

beides unter der Voraussetzung, dass über $p$ passend verfügt wird. Ist dies beides bewiesen, so erkennt man die Unmöglichkeit der Gleichung (14) und also die der Gleichung (13), aus der (14) gefolgert war.

Ordnen wir den Zähler von $f(x)$ nach Potenzen von $x$, so ergiebt sich ein Ausdruck der Form:

\weq{16}{f(x) = \frac{A_{p-1}x^{p-1}+A_p x^p + A_{p+1}x^{p+1}+\cdots}{\Pi(p-1)},}

worin $A_{p-1}, A_p, A_{p+1}, \ldots$ ganze Zahlen sind, und $A_{p-1}=[\Pi(m)]^p$, also gewiss nicht durch $p$ theilbar. Es ist daher, wenn man (16) mit der \gesperrt{Taylor}'schen Entwickelung
\[
f(x) = f(0) + xf'(0) + \frac{x^2}{\Pi(2)}f''(0)+\cdots
\]
vergleicht,
\[
f(0)=f'(0)=f''(0)=\cdots=f^{(p-2)}(0)=0,
\]
\[
f^{(p-1)}(0)=A_{p-1},\quad f^{(p)}(0)=pA_p,\quad f^{(p+1)}(0)=p(p+1)A_{p+1}\cdots,
\]
also
\[
F(0) = A_{p-1}+pA_{p+1}+p(p+1)A_{p+1}+\cdots, % [sic] -- printed with A_{p+1} repeated where the pattern would expect A_p, A_{p+1}
\]
eine \gesperrt{durch $p$ nicht theilbare ganze Zahl}. Ordnen wir aber $f(x)$ nach Potenzen von $x-1$, so folgt:
\[
f(x) = \frac{B_p(x-1)^p+B_{p+1}(x-1)^{p+1}+\cdots}{\Pi(p-1)},
\]
worin die $B_p, B_{p+1}, \ldots$ wieder ganze Zahlen sind. Daraus folgt, wie oben, durch Vergleichung mit
\[
f(x) = f(1) + (x-1)f'(1) + \frac{(x-1)^2}{\Pi(2)}f''(1)+\cdots
\]
\[
F(1) = pB_p + p(p+1)B_{p+1}+\cdots,
\]
folglich ist $F(1)$ eine \gesperrt{durch $p$ theilbare ganze Zahl}, und genau auf demselben Wege ergiebt sich, dass $F(2), F(3), \ldots, F(m)$ durch $p$ theilbare ganze Zahlen sind. Da man nun auch $p$ so gross wählen kann, dass $C_0$ nicht durch $p$ theilbar ist, so folgt
%CONT>>
% signature "48*" at foot -- not text, omitted


% ===== scan leaf 777, printed 756 =====
% !PRINTED 756
\weberpage{756}
\pdfbookmark[2]{§ 206. Transcendenz der Zahl π}{bm235}
%<<CONT
dass $C_0 F(0) + C_1 F(1) + \cdots + C_m F(m)$ eine nicht durch $p$ theilbare, also auch nicht verschwindende ganze Zahl ist, und damit ist 1) bewiesen.

Wenn wir nun noch nachweisen können, dass $Q(x)$ für jedes positive $x$ durch Vergrösserung von $p$ beliebig klein gemacht werden kann, so folgt, dass bei hinlänglich grossem $p$ der Ausdruck 2) gewiss kleiner als 1 ist, und unser Beweis ist vollendet.

Gehen wir aber zu dem Ausdrucke (10) für $Q(x)$ zurück und bezeichnen mit $\gamma_n, \gamma_{n-1}, \ldots, \gamma_0$ die absoluten Werthe der Coëfficienten $c_n, c_{n-1}, \ldots, c_0$ von $f(x)$, so sehen wir, da die $q_n, q_{n-1}, \ldots$ dem absoluten Werthe nach kleiner als 1 sind, dass für jedes positive $x$ der absolute Werth von $Q(x)$ kleiner ist als

\weq{17}{\psi(x) = \gamma_n x^n + \gamma_{n-1}x^{n-1}+\cdots+\gamma_0.}

Die Coëfficienten $c_n, c_{n-1}, \ldots, c_0$ der Function $f(x)$ in (15) unterscheiden sich aber von den $\gamma_n, \gamma_{n-1}, \ldots, \gamma_0$ nur durch das Vorzeichen. Ersetzen wir daher $x$ durch $-x$, bilden also die Function
\[
f(-x) = \frac{x^{p-1}(1+x)^p(2+x)^p\ldots(m+x)^p}{\Pi(p-1)},
\]
so hat diese Function dieselben Coëfficienten, wie $f(x)$, aber alle mit positiven Vorzeichen. Sie ist also keine andere, als die Function $\psi(x)$.

Setzen wir also noch
\[
X = x(1+x)(2+x)\ldots(m+x),
\]
so wird

\weq{18}{\psi(x) = \frac{X}{x}\cdot\frac{X^{p-1}}{\Pi(p-1)},}

und dies nähert sich, wie am Anfange dieses Paragraphen nachgewiesen ist, mit unendlich wachsendem $p$ der Grenze \gesperrt{Null}.

\gesperrt{Es ist also $e$ eine transcendente Zahl.}

\begin{center}
§. 206.

\gesperrt{Transcendenz der Zahl $\pi$.}
\end{center}

Mit denselben Hülfsmitteln lässt sich nun auch die Transcendenz der Zahl $\pi$ beweisen. Als Definition dieser Zahl dient uns dabei, dass es die kleinste positive Zahl ist, die der Gleichung

\weq{1}{e^{i\pi} = -1}

genügt, wenn $e^x$ durch die Formel \S.\,205, (2) definirt wird.


% ===== scan leaf 778, printed 757 =====
% !PRINTED 757
\weberpage{757}

Nehmen wir also an, es sei $\pi$ und folglich auch $i\pi$ eine algebraische Zahl, so ist $i\pi$ eine der Wurzeln einer irreduciblen rationalen Gleichung $\chi(x)=0$, deren Coëfficienten ganze rationale Zahlen sind.

Bezeichnen wir die sämmtlichen Wurzeln dieser algebraischen Gleichung mit $\beta_1,\beta_2,\ldots,\beta_\nu$, und bezeichnen den Coëfficienten von $x^\nu$ in $\chi$ mit $a$, so ist
\weq{2}{\chi(x) = a(x-\beta_1)(x-\beta_2)\ldots(x-\beta_\nu),}
und die Producte $a\beta_1,a\beta_2,\ldots,a\beta_\nu$ sind \gesperrt{ganze} algebraische Zahlen, ihre ganzen symmetrischen Functionen also \gesperrt{ganze rationale Zahlen} (\S.\,133).

Da nun die Zahl $i\pi$ unter den $\beta$ vorkommen soll, so ist nach (1):
\[
(1+e^{\beta_1})(1+e^{\beta_2})\ldots(1+e^{\beta_\nu}) = 0,
\]
und wenn wir die Multiplication ausführen, so ergiebt sich eine Gleichung:
\[
1 + \sum e^{\beta_i} + \sum e^{\beta_i+\beta_h} + \sum e^{\beta_i+\beta_h+\beta_k} + \cdots = 0.
\]
Unter den Exponenten in dieser Summe kann mehrmals die Zahl Null vorkommen; wir wollen annehmen $(C-1)$mal, so dass $C$ eine positive ganze Zahl, mindestens $=1$, ist. Die übrigen Exponenten $\beta_i,\beta_i+\beta_h,\beta_i+\beta_h+\beta_k,\ldots$, die zum Theil auch unter einander gleich sein können, wollen wir mit $\alpha_1,\alpha_2,\ldots,\alpha_\mu$ bezeichnen, so dass die Gleichung besteht:
\weq{3}{C + e^{\alpha_1} + e^{\alpha_2} + \cdots + e^{\alpha_\mu} = 0.}
Hierin sind nun die $\alpha_1,\alpha_2,\ldots,\alpha_\mu$ algebraische Zahlen, die, mit der ganzen rationalen Zahl $a$ multiplicirt, in ganze algebraische Zahlen übergehen.

Die symmetrischen Functionen der sämmtlichen Summen $\beta_i,\beta_i+\beta_h,\beta_i+\beta_h+\beta_k,\ldots$ sind aber zugleich symmetrische Functionen der $\beta_1,\ldots,\beta_\nu$, und folglich rationale Zahlen. Diese Summen sind folglich auch die Wurzeln einer rationalen Gleichung, und da man die Wurzel $0$, so oft sie vorhanden ist, absondern kann, so sind auch die $\alpha_1,\alpha_2,\ldots,\alpha_\mu$ die Wurzeln einer rationalen Gleichung; die symmetrischen Grundfunctionen der $a\alpha_1,a\alpha_2,\ldots,a\alpha_\mu$ sind \gesperrt{ganze rationale Zahlen}.

Die absoluten Werthe der Zahlen
\[
\alpha_1,\alpha_2,\ldots,\alpha_\mu
\]
sollen jetzt mit
\[
a_1,a_2,\ldots,a_\mu
\]
bezeichnet werden.


% ===== scan leaf 779, printed 758 =====
% !PRINTED 758
\weberpage{758}

Wenn wir nun in der Gleichung (12) des vorigen Paragraphen $x=0,\alpha_1,\alpha_2,\ldots,\alpha_\mu$ setzen und die Gleichung (3) anwenden, so ergiebt sich:
\weq{4}{0 = CF(0) + F(\alpha_1) + F(\alpha_2) + \cdots + F(\alpha_\mu) + CQ(0) + e^{a_1}Q(\alpha_1) + e^{a_2}Q(\alpha_2) + \cdots + e^{a_\mu}Q(\alpha_\mu),}
und wir haben die Unmöglichkeit einer solchen Gleichung bei einer geeigneten Wahl von $f(x)$ darzuthun.

Es sei wieder $p$ eine zu unbegrenztem Wachsen bestimmte Primzahl, und
\weq{5}{f(x) = \dfrac{a^{\mu p+p-1}x^{p-1}(x-\alpha_1)^p(x-\alpha_2)^p\ldots(x-\alpha_\mu)^p}{\Pi(p-1)},}
so dass $f(x)$ eine Function mit rationalen Coëfficienten ist.

Wir bilden nun, wie im vorigen Paragraphen, durch Ordnen nach Potenzen von $x$:
\[
f(x) = \dfrac{A_{p-1}x^{p-1} + A_px^p + A_{p+1}x^{p+1}\ldots}{\Pi(p-1)},
\]
worin die $A_{p-1},A_p,\ldots$ ganze rationale Zahlen sind, und
\[
A_{p-1} = (-1)^\mu a^{\mu p+p-1}\alpha_1^p\alpha_2^p\ldots\alpha_\mu^p.
\]
Wenn wir also $p$ grösser als jede der ganzen Zahlen
\[
a, a^\mu\alpha_1\alpha_2\ldots\alpha_\mu
\]
annehmen, so ist $A_{p-1}$ durch $p$ nicht theilbar, und es wird
\[
F(0) = A_{p-1} + pA_p + p(p+1)A_{p+1} + \cdots,
\]
d.\,h. eine durch $p$ nicht theilbare ganze Zahl. Ebenso ist, wenn $p$ gross genug genommen wird, $C$ durch $p$ nicht theilbar.

Andererseits ordnen wir den Zähler von $f(x)$ nach Potenzen von $a(x-\alpha_1)$ und erhalten
\[
f(x) = \dfrac{B_pa^p(x-\alpha_1)^p + B_{p+1}a^{p+1}(x-\alpha_1)^{p+1} + \cdots}{\Pi(p-1)},
\]
worin die $B_p,B_{p+1},\ldots$ zwar nicht mehr rationale, wohl aber \gesperrt{ganze algebraische Zahlen} sind, da der Zähler von $f(x)$ eine ganzzahlige Function von $ax,a\alpha_1,\ldots,a\alpha_\mu$ ist.

Hieraus folgt nun, wie im vorigen Paragraphen:
\[
F(\alpha_1) = pB_pa^p + p(p+1)B_{p+1}a^{p+1} + \cdots
\]
Bildet man auf die gleiche Weise $F(\alpha_2),\ldots,F(\alpha_\mu)$, und beachtet, dass die Summe $B_p(\alpha_1)+B_p(\alpha_2)+\cdots+B_p(\alpha_\mu)$ eine ganze rationale Zahl ist, so folgt, dass
\[
F(\alpha_1) + F(\alpha_2) + \cdots + F(\alpha_\mu)
\]
%CONT>>


% ===== scan leaf 780, printed 759 =====
% !PRINTED 759
\weberpage{759}
%<<CONT
eine durch $p$ theilbare ganze rationale Zahl ist. Daraus ergiebt sich endlich, dass
\[
CF(0) + F(\alpha_1) + F(\alpha_2) + \cdots + F(\alpha_\mu)
\]
eine nicht durch $p$ theilbare, also auch nicht verschwindende, ganze rationale Zahl ist, die daher, vom Vorzeichen abgesehen, mindestens $=1$ ist.

Können wir nun endlich noch beweisen, dass auch bei der jetzigen Annahme über $f$ die Function $Q(x)$ für jedes endliche $x$ bei hinlänglicher Vergrösserung von $p$ beliebig klein gemacht werden kann, so ergiebt sich, genau wie oben, die Unmöglichkeit der Gleichung (4).

Dies lässt sich aber folgendermaassen einsehen: Wir betrachten statt der Function
\weq{6}{f(x) = c_nx^n + c_{n-1}x^{n-1} + \cdots + c_0}
die Function
\[
\begin{aligned}
\psi(x) &= \dfrac{a^{\mu p+p-1}x^{p-1}(x+a_1)^p(x+a_2)^p\ldots(x+a_\mu)^p}{\Pi(p-1)} \\
&= \gamma_nx^n + \gamma_{n-1}x^{n-1} + \cdots + \gamma_0,
\end{aligned}
\]
die nun lauter positive (aber nicht nothwendig rationale) Coëfficienten hat. Die Coëfficienten $c_n,c_{n-1}\ldots$ sind durch Multiplication und Addition aus den Zahlen $a,-\alpha_1,-\alpha_2,\ldots,-\alpha_\mu$ gebildet, und die entsprechenden Coëfficienten $\gamma_n,\gamma_{n-1}\ldots$ erhält man daraus, wenn man die $-\alpha_1,-\alpha_2,\ldots,-\alpha_\mu$ durch ihre absoluten Werthe $a_1,a_2,\ldots,a_\mu$ ersetzt, woraus nach dem schon oben erwähnten Satze der Einleitung folgt, dass die Coëfficienten $\gamma_n,\gamma_{n-1},\ldots$ gewiss nicht kleiner sind, als die absoluten Werthe der entsprechenden $c_n,c_{n-1}\ldots$

Nun ist für jedes endliche $x$, dessen absoluter Werth $\xi$ ist, der absolute Werth von $Q(x)$ nach \S.\,205, (10) kleiner, oder wenigstens nicht grösser als
\[
\gamma_n\xi^n + \gamma_{n-1}\xi^{n-1} + \cdots + \gamma_0 = \psi(\xi),
\]
und dass $\psi(\xi)$ unter jede Grenze heruntersinkt, wenn $p$ gross genug wird, schliesst man aus der Darstellung:
\[
\psi(x) = \dfrac{X}{ax}\dfrac{X^{p-1}}{\Pi(p-1)},
\]
wenn
\[
% [sic] -- the source prints Latin $u+1$; the derivation from the exponent $\mu p+p-1$ requires $\mu+1$.
X = a^{u+1}x(x+a_1)(x+a_2)\ldots(x+a_\mu)
\]
gesetzt ist.


% ===== scan leaf 781, printed 760 =====
% !PRINTED 760
\weberpage{760}
\pdfbookmark[2]{§ 207. Der allgemeine Satz von Lindemann über die Exponentialfunction}{bm236}

Damit ist bewiesen:

\gesperrt{Die Zahl $\pi$ ist eine transcendente Zahl. Die \glqq Quadratur des Kreises\grqq{} kann nicht durch geometrische Construction, bei der nur algebraische Curven und Flächen angewandt werden, gelöst werden.}

\begin{center}
\S.\,207.

\gesperrt{Der allgemeine Satz von Lindemann über die Exponentialfunction.}
\end{center}

Die Transcendenz der Zahlen $e$ und $\pi$, die hierdurch bewiesen ist, ist als specieller Fall in einem sehr allgemeinen Theoreme über die Exponentialfunction enthalten, das \gesperrt{Lindemann} in der oben erwähnten Abhandlung angekündigt hat, von dem ein ausgeführter Beweis in der citirten Abhandlung von \gesperrt{Weierstrass} enthalten ist. Wir wollen diesen Satz hier zum Schluss noch beweisen mit Anwendung derselben Hülfsmittel, die wir für die beiden speciellen Fälle benutzt haben. Der Satz lautet so:

I. \gesperrt{Es besteht keine Gleichung von der Form} \weq{1}{C_1e^{z_1} + C_2e^{z_2} + \cdots + C_me^{z_m} = 0,} \gesperrt{in der die Coëfficienten $C_1,C_2,\ldots,C_m$ algebraische Zahlen und die Exponenten $z_1,z_2,\ldots,z_m$ von einander verschiedene algebraische Zahlen sind, es sei denn, dass alle Coëfficienten $C_1,C_2,\ldots,C_m$ gleich Null sind.}

Um ihn zu beweisen, leiten wir zunächst einen Hülfssatz ab:

Es seien
\[
\alpha_1,\alpha_2,\ldots,\alpha_r
\]
beliebige reelle oder imaginäre, jedoch \gesperrt{von einander verschiedene} Grössen, und
\[
A_1,A_2,\ldots,A_r
\]
ebenfalls beliebige Grössen, die \gesperrt{nicht alle verschwinden}.

Dasselbe soll von den zwei Grössenreihen
\[
\begin{gathered}
\beta_1,\beta_2,\ldots,\beta_s \\
B_1,B_2,\ldots,B_s
\end{gathered}
\]
gelten. Wir bezeichnen mit
\[
\gamma_1,\gamma_2,\ldots,\gamma_t
\]
%CONT>>


% ===== scan leaf 782, printed 761 =====
% !PRINTED 761
\weberpage{761}
%<<CONT
\gesperrt{die von einander verschiedenen} unter den $rs$ Summen
$\alpha_i+\beta_k$, und setzen

\weq{2}{\begin{aligned} A &= A_1e^{\alpha_1}+A_2e^{\alpha_2}+\cdots+A_re^{\alpha_r}\\ B &= B_1e^{\beta_1}+B_2e^{\beta_2}+\cdots+B_se^{\beta_s}, \end{aligned}}

\weq{3}{AB = C_1e^{\gamma_1}+C_2e^{\gamma_2}+\cdots+C_te^{\gamma_r}.} % [sic] printed exponent index $\gamma_r$, inconsistent with $C_t$

Der zu beweisende Hülfssatz lautet dann:

1. \gesperrt{Die Co\"efficienten $C_1, C_2, \ldots, C_t$ können nicht alle verschwinden.}

Beim Beweise dieses Satzes können wir offenbar annehmen, dass von den Coëfficienten $A_1, A_2, \ldots, A_r, B_1, B_2, \ldots, B_s$ keiner verschwindet (da wir die etwa verschwindenden einfach weglassen können).

Des kürzeren Ausdrucks wegen nennen wir für den Augenblick, wenn $a, b$ zwei verschiedene complexe Zahlen sind, $a$ kleiner als $b$, $(a<b)$, wenn der reelle Theil von $a$ kleiner ist als der reelle Theil von $b$, oder wenn die reellen Theile gleich und der imaginäre Theil von $a$ kleiner ist, als der imaginäre Theil von $b$.

Ist dann in diesem Sinne $a<b$, $b<c$, so ist auch $a<c$, und ist $a<b$, $c<d$, so ist $a+c<b+d$.

Unter jeder endlichen Reihe von einander verschiedener complexer Zahlen giebt es dann eine bestimmte kleinste, und wenn also $\alpha_1$ die kleinste unter den Zahlen $\alpha$, $\beta_1$ die kleinste unter den Zahlen $\beta$ ist, so ist $\alpha_1+\beta_1$ die kleinste unter den Zahlen $\gamma$, und diese Summe ist keiner der anderen Summen $\alpha_i+\beta_k$ gleich. Es ist also $C_1=A_1B_1$ und $C_1$ von Null verschieden, w. z. b. w.

Dieser Satz lässt sich nun durch vollständige Induction sofort verallgemeinern.

2. \gesperrt{Sind}

\[
\begin{aligned}
A' &= A_1'e^{\alpha_1'}+A_2'e^{\alpha_2'}+\cdots\\
A'' &= A_1''e^{\alpha_1''}+A_2''e^{\alpha_2''}+\cdots\\
A''' &= A_1'''e^{\alpha_1'''}+A_2'''e^{\alpha_2'''}+\cdots\\
&\ \cdots\cdots\cdots\cdots\cdots\cdots
\end{aligned}
\]

\gesperrt{Summen von der Form (2) in beliebiger Anzahl, und $\gamma_1,\gamma_2,\ldots$ die von einander verschiedenen unter den Summen $\alpha_h'+\alpha_i''+\alpha_k'''+\cdots$, so sind in dem Producte}

\weq{4}{A'A''A'''\cdots = C_1e^{\gamma_1}+C_2e^{\gamma_2}+\cdots}

\gesperrt{nicht alle Coëfficienten $C_1, C_2, \ldots$ gleich Null.}


% ===== scan leaf 783, printed 762 =====
% !PRINTED 762
\weberpage{762}
Dieser Hülfssatz gestattet uns zunächst, beim Beweise des Theorems (1) die vereinfachende Voraussetzung zu machen, die Coëfficienten $C_1, C_2, \ldots, C_m$ seien ganze rationale Zahlen.

Angenommen nämlich, es bestehe eine Gleichung

\weq{5}{X_1e^{x_1}+X_2e^{x_2}+\cdots = 0}

mit algebraischen Coëfficienten $X_1, X_2, \ldots$ und von einander verschiedenen Exponenten $x_1, x_2, \ldots$, so können wir immer annehmen, die $X_1, X_2, \ldots$ gehören einem und demselben algebraischen Körper $\Omega$ an.

Wir bezeichnen mit $u_1, u_2, \ldots$ Variable und bilden die Norm der Linearform $X_1u_1+X_2u_2+\cdots$

\weq{6}{N(X_1u_1+X_2u_2+\cdots) = \Phi(u_1,u_2,\ldots),}

die eine ganze homogene Function der Variablen $u$ mit \gesperrt{rationalen Coëfficienten ist}, deren Grad gleich dem Grade des Körpers $\Omega$ ist. Setzen wir in (6)

\[
u_1 = e^{x_1},\ u_2 = e^{x_2},\ \ldots,
\]

so geht jedes Product von Variablen $u$ in eine Grösse von der Form $e^z$ über, worin $z$ eine Summe von Zahlen $x$ ist, und $\Phi(u_1,u_2,\ldots)$ wird ein Ausdruck von der Form $C_1e^{z_1}+C_2e^{z_2}+\cdots$, dessen Coëfficienten rationale Zahlen sind, die, wenn die $X_1, X_2, \ldots$ nicht alle verschwinden, nach 2. auch dann nicht alle Null sein können, wenn die unter einander gleichen unter den Gliedern der Function $\Phi(e^{x_1},e^{x_2},\ldots)$ in ein einziges Glied $Ce^z$ zusammengefasst werden. Wenn nun aber die Gleichung (5) besteht, so ist auch $\Phi(e^{x_1},e^{x_2},\ldots)=0$, und folglich

\[
C_1e^{z_1}+C_2e^{z_2}+\cdots = 0.
\]

Sind unter den $C_1, C_2, \ldots$ gebrochene Zahlen, so können wir sie durch Multiplication der ganzen Gleichung mit dem Hauptnenner in ganze Zahlen verwandeln.

3. \gesperrt{Es braucht also jetzt nur noch bewiesen zu werden, dass die Gleichung (1) für kein System ganzer rationaler Zahlen $C_1, C_2, \ldots, C_m$, die nicht alle verschwinden, bestehen kann.}

Demnach nehmen wir jetzt an, es bestehe eine Gleichung

\weq{7}{A_1e^{x_1}+A_2e^{x_2}+\cdots = 0,}

deren Coëfficienten $A_1, A_2, \ldots$ ganze rationale Zahlen sind,
%CONT>>


% ===== scan leaf 784, printed 763 =====
% !PRINTED 763
\weberpage{763}
%<<CONT
die nicht alle $=0$ sind, worin die Exponenten $x_1, x_2, \ldots$ von einander verschiedene algebraische Zahlen sind. Wir bezeichnen mit $\Omega$ einen \gesperrt{Normalkörper}, dem alle diese Zahlen $x_1, x_2, \ldots$ angehören, und mit $\Theta$ eine primitive Zahl dieses Körpers, ferner mit

\[
\sigma' = (\Theta,\Theta'),\ \sigma'' = (\Theta,\Theta''),\ \ldots
\]

die Substitutionen des Körpers $\Omega$. Wenn durch eine dieser Substitutionen, $\sigma'$, die Zahlen $x_1, x_2, \ldots$ in $x_1', x_2', \ldots$ übergehen, so müssen auch diese von einander verschieden sein.

Es sei jetzt $u$ eine Variable und

\weq{8}{U = A_1e^{ux_1}+A_2e^{ux_2}+\cdots}

Hierin machen wir die sämmtlichen Substitutionen $\sigma', \sigma'', \ldots$ und bezeichnen die so entstehenden Functionen mit $U', U'', \ldots$ Das Product aller dieser Functionen können wir, obwohl es sich nicht bloss um algebraische Zahlen handelt, die \gesperrt{Norm} von $U$ nennen. Dieses Product hat die Form

\weq{9}{N(U) = C_1e^{uz_1}+C_2e^{uz_2}+\cdots,}

worin die $C_1, C_2, \ldots$ gleichfalls ganze rationale Zahlen sind. Die $z_1, z_2, \ldots$ sind Zahlen des Körpers $\Omega$, die wir, wenn wir die Glieder mit gleichen Exponenten in ein einziges Glied zusammenfassen, als \gesperrt{von einander verschieden} voraussetzen dürfen. Zugleich ist wegen (7)

\weq{10}{C_1e^{z_1}+C_2e^{z_2}+\cdots = 0.}

Der Ausdruck auf der rechten Seite von (9) hat aber noch, wie seine Entstehung als Norm zeigt, die Eigenschaft, ungeändert zu bleiben, wenn irgend eine der Substitutionen $\sigma', \sigma'', \ldots$ gemacht wird. Entwickelt man aber (9) nach Potenzen von $u$, so muss diese Unveränderlichkeit von jedem Gliede dieser Reihenentwickelung gelten, also wenn $h$ ein beliebiger ganzer positiver Exponent ist, für

\[
C_1z_1^h+C_2z_2^h+\cdots;
\]

diese Summe ist aber eine Zahl des Körpers $\Omega$, und daher eine \gesperrt{rationale Zahl}. Dies können wir dahin zusammenfassen:

Bedeutet $g(z)$ irgend eine ganze Function von $z$ mit rationalen Coëfficienten, so ist

\[
C_1g(z_1)+C_2g(z_2)+\cdots
\]

eine \gesperrt{rationale Zahl}.


% ===== scan leaf 785, printed 764 =====
% !PRINTED 764
\weberpage{764}
4. \gesperrt{Der Beweis des Theorems I. ist hierdurch auf den Beweis des folgenden speciellen Falles zurückgeführt: Besteht eine Gleichung}

\weq{11}{C_1e^{z_1}+C_2e^{z_2}+\cdots+C_me^{z_m} = 0,}

\gesperrt{in der die $z_1, z_2, \ldots, z_m$ von einander verschiedene algebraische Zahlen sind, in der $C_1, C_2, \ldots, C_m$ und die Verbindungen}

\weq{12}{C_1g(z_1)+C_2g(z_2)+\cdots+C_mg(z_m)}

\gesperrt{rationale Zahlen sind, wenn $g(z)$ irgend eine ganze Function mit rationalen Coëfficienten ist, so müssen die $C_1, C_2, \ldots, C_m$ alle verschwinden.}

Um diesen Satz zu beweisen, bestimmen wir zunächst eine positive ganze rationale Zahl $c$ so, dass

\weq{13}{x_1 = cz_1,\ x_2 = cz_2,\ \ldots,\ x_m = cz_m}

ganze algebraische Zahlen werden (§.\,133, 5.). Wir nehmen die Coëfficienten $C_1, C_2, \ldots, C_m$ als \gesperrt{ganze} rationale Zahlen an, und bilden eine ganze Function

\weq{14}{\varphi(x) = ax^r+a_1x^{r-1}+\cdots+a_r}

mit ganzen rationalen Zahlencoëfficienten, die folgende Eigenschaften hat:

1) Die Zahlen $x_1, x_2, \ldots, x_m$ kommen unter den Wurzeln von $\varphi(x)=0$ vor.

2) Die Summe

\[
C_1\varphi'(x_1)+C_2\varphi'(x_2)+\cdots+C_m\varphi'(x_m) = k,
\]

die nach den Voraussetzungen (12), (13) eine ganze rationale Zahl ist, ist von Null verschieden.

Um einzusehen, dass eine solche Function $\varphi(x)$ immer existirt, nehme man zunächst eine Function $\chi(x)$, die nur der Bedingung 1) genügt, und der zweiten, dass keine der Zahlen $x_1, x_2, \ldots, x_m$ eine Doppelwurzel von $\chi(x)$ ist, eine solche Function existirt offenbar [man kann z. B., um eine Function $\chi$ zu bilden, die Norm des Productes $(x-x_1)(x-x_2)\cdots(x-x_m)$ von gemeinsamen Theilern mit ihren Derivirten befreien]. Setzt man dann

\[
\varphi(x) = x^h\chi(x),\ \varphi'(x_1) = x_1^h\chi'(x_1),\ \varphi'(x_2) = x_2^h\chi'(x_2)\ldots,
\]
%CONT>>


% ===== scan leaf 786, printed 765 =====
% !PRINTED 765
\weberpage{765}
%<<CONT
so kann die Summe

\[
C_1\,\varphi'(x_1)+\cdots+C_m\,\varphi'(x_m) = C_1\,\chi'(x_1)\,x_1^h+\cdots+C_m\,\chi'(x_m)\,x_m^h
\]

gewiss nicht für jeden Exponenten $h$ verschwinden, da sonst gegen die Voraussetzung $C_1\,\chi'(x_1)\ldots C_m\,\chi'(x_m)$ alle gleich Null sein müssten.

Nachdem so die Existenz einer Function $\varphi(x)$, wie sie verlangt war, festgestellt ist, wenden wir die Formel \S.\,205, (12) an:

% [sic] -- the source prints $P(x)$, although $P$ in \S.\,205, (12) is constant and the deduction below requires it to be independent of $x$.
\[
e^x\,P(x) = F(x) + e^{\xi}\,Q(x).
\]

Wir setzen darin $x=z_1, z_2, \ldots, z_m$, und bezeichnen die absoluten Werthe von $z_1, z_2, \ldots, z_m$ mit $\xi_1, \xi_2, \ldots, \xi_m$.

Dann ergiebt sich nach (11):

\weq{15}{0 = C_1\,F(z_1)+C_2\,F(z_2)+\cdots+C_m\,F(z_m)+C_1\,e^{\xi_1}\,Q(z_1)+C_2\,e^{\xi_2}\,Q(z_2)+\cdots+C_m\,e^{\xi_m}\,Q(z_m),}

und es ist, wenn $f(x)$ eine beliebige ganze Function $n^{\mathrm{ten}}$ Grades ist

\weq{16}{F(z) = f(z)+f'(z)+\cdots+f^{(n)}(z).}

Wir setzen, indem wir mit $p$ eine natürliche Primzahl bezeichnen, die beliebig gross genommen werden kann, $x=cz$ [nach (13)] und

\weq{17}{f(z) = \dfrac{\varphi(x)^{p-1}\,\varphi'(x)}{\Pi(p-1)},}

wenn $\varphi(x)$ die den Bedingungen 1), 2) genügende Function ist.

Durch Ordnen nach Potenzen von $(x-x_1)$ mag sich ergeben:

\[
\varphi(x)^{p-1}\,\varphi'(x) = \varphi'(x_1)^p\,(x-x_1)^{p-1}+A_p(x_1)\,(x-x_1)^p+A_{p+1}(x_1)\,(x-x_1)^{p+1}+\cdots,
\]

und darin sind die $A_p(x_1), A_{p+1}(x_1), \ldots$ ganze algebraische Zahlen, und zwar rationale Functionen von $x_1$. Andererseits ist nach dem \gesperrt{Taylor}'schen Lehrsatze:

\[
f(z) = f(z_1)+(z-z_1)\,f'(z_1)+\dfrac{(z-z_1)^2}{1\,.\,2}\,f''(z_1)+\cdots
\]

und $c(z-z_1)=x-x_1$. Die Vergleichung ergiebt alsdann

\[
f(z_1)=0,\ f'(z_1)=0,\ \ldots,\ f^{(p-2)}(z_1)=0,\ f^{(p-1)}(z_1)=c^{p-1}\,\varphi'(x_1)^p,
\]
\[
f^{(p)}(z_1)=p\,c^p\,A_p(x_1),\ f^{(p+1)}(z_1)=p(p+1)\,c^{p+1}\,A_{p+1}(x_1),\ \ldots,
\]

und folglich

\[
F(z_1) = c^{p-1}\,\varphi'(x_1)^p+p\,c^p\,A_p(x_1)+p(p+1)\,c^{p+1}\,A_{p+1}(x_1)+\cdots;
\]

hierin kann $x_1, z_1$ durch $x_2, z_2$ oder durch $x_3, z_3$ u.\,s.\,f. ersetzt werden.


% ===== scan leaf 787, printed 766 =====
% !PRINTED 766
\weberpage{766}
Danach ist die ganze rationale Zahl

\[
C_1\,F(z_1)+C_2\,F(z_2)+\cdots+C_m\,F(z_m)
\]

nach dem Modul $p$ mit

\[
c^{p-1}\left[C_1\,\varphi'(x_1)^p+C_2\,\varphi'(x_2)^p+\cdots+C_m\,\varphi'(x_m)^p\right]
\]

congruent. Da nun $C_1, C_2, \ldots, C_p$ ganze rationale Zahlen sind, % [sic] -- printed $C_p$; the surrounding system has $m$ constants
so ist $C_1^p\equiv C_1, C_2^p\equiv C_2, \ldots \pmod{p}$, und es ergiebt sich durch Anwendung des polynomischen Lehrsatzes:

\[
\begin{aligned}
C_1\,\varphi'(x_1)^p+C_2\,\varphi'(x_2)^p+\cdots+C_m\,\varphi'(x_m)^p &\equiv \\
\left[C_1\,\varphi'(x_1)+C_2\,\varphi'(x_2)+\cdots+C_m\,\varphi'(x_m)\right]^p &\equiv k^p \pmod{p}.
\end{aligned}
\]

Danach erhalten wir also die Congruenz

\weq{18}{C_1\,F(z_1)+C_2\,F(z_2)+\cdots+C_m\,F(z_m) \equiv c^{p-1}\,k^p \pmod{p}.}

Nun sind die Zahlen $c, k$ von $p$ unabhängig, und wir können daher $p$ so gross annehmen, dass es nicht in $c$ und in $k$ aufgeht.

5. \gesperrt{Dann ist die Summe $C_1\,F(z_1)+C_2\,F(z_2)+\cdots+C_m\,F(z_m)$ eine von Null verschiedene ganze rationale Zahl, also dem absoluten Werthe nach mindestens gleich 1.}

Wir bezeichnen nun mit $\varphi_1(x)$ die ganze Function, die sich aus $\varphi(x)$ ergiebt, wenn die negativen unter den Coëfficienten $a, a_1, a_2, \ldots$ durch ihre positiven Werthe ersetzt werden. Die Function

\[
f_1(z) = \dfrac{\varphi_1(x)^{p-1}\,\varphi_1'(x)}{\Pi(p-1)},
\]

die sich nach (17) ergiebt, hat dann nur positive Coëfficienten, und diese Coëfficienten sind dem absoluten Werthe nach gewiss nicht kleiner, als die entsprechenden Coëfficienten von $f(z)$, weil die Coëfficienten von $f_1(z)$ aus denselben Zahlen durch Addition entstehen, die bei der Bildung der Coëfficienten von $f(z)$ theils addirt, theils subtrahirt werden.

Hieraus aber ergiebt sich nach der Formel \S.\,205, (10), wenn wir mit $\xi$ den absoluten Werth von $x$ bezeichnen, für den absoluten Werth von $Q(z)$:

\[
|Q(z)| \leqq \dfrac{\varphi_1(\xi)^{p-1}\,\varphi_1'(\xi)}{\Pi(p-1)},
\]

6. \gesperrt{und es kann also $Q(z)$ für jedes endliche $z$ durch hinlängliche Vergrösserung von $p$ beliebig klein gemacht werden.}


% ===== scan leaf 788, printed 767 =====
% !PRINTED 767
\weberpage{767}
Durch 5. und 6. ist aber nachgewiesen, dass, wenn $p$ gross genug ist, die Gleichung (15) nicht bestehen kann. Dadurch ist 4. und damit das ganze Theorem I. bewiesen.

Dieser Satz gestattet nun mannigfaltige Anwendungen. Er giebt uns zunächst die Transcendenz von $e$, wenn wir $C_1, C_2, \ldots, z_1, z_2, \ldots$ als ganze rationale Zahlen annehmen. Er giebt ferner die Transcendenz von $\pi$. Denn aus der Gleichung $1+e^{i\pi}=0$ folgt nach I., dass $i\pi$ und folglich auch $\pi$ nicht algebraisch sein kann.

Es folgt ferner daraus:

\gesperrt{Für jede algebraische Zahl $x$, mit Ausnahme von $x=0$, ist $X=e^x$ eine transcendente Zahl.}

\gesperrt{Für jedes algebraische $X$, mit Ausnahme von $X=1$, ist jeder natürliche Logarithmus $x=\log X$ eine transcendente Zahl.}

\gesperrt{Für jeden Bogen, der zum Radius in einem algebraisch ausdrückbaren Verhältnisse $x$ steht, mit Ausnahme von $x=0$, ist $X=\sin x$ eine transcendente Zahl.}

Dies folgt nach I. aus $2i\,X=e^{ix}-e^{-ix}$.

% [sic] -- the source prints the chord as $\frac12\sin(x/2)$; for unit radius the chord is $2\sin(x/2)$.
% [sic] -- the source's claim for $\mathrm{tg}\,x=\alpha x$ omits the algebraic root $x=0$; the intended claim concerns nonzero roots.
Dasselbe gilt für die anderen trigonometrischen Linien, $\cos x$, $\mathrm{tg}\,x$ und für die Sehne $\dfrac{1}{2}\sin\dfrac{x}{2}$. Um noch eins anzuführen: Die transcendente Gleichung $\mathrm{tg}\,x=\alpha x$ hat für ein algebraisches $\alpha$ ausser 0 nur transcendente Zahlen zu Wurzeln.

\begin{center}
\rule{4cm}{0.4pt}
\end{center}


% ===== scan leaf 789, printed 768 =====
\weberpage{768}% page number not printed on this page
% !PRINTED none
% part-title page (no running head, no printed page number): centred heading
% "Nachträge." with a horizontal rule beneath it, then subsection "I." and
% "Irreducibilität der Kreistheilungsgleichung.", followed by body text.
\begin{center}
{\Large \gesperrt{Nachträge.}}

\rule{4em}{0.4pt}

I.

\textbf{\gesperrt{Irreducibilität der Kreistheilungsgleichung.}}
\end{center}

Der \S.\,134 des ersten Bandes bedarf in seinem letzten Theile, der von der Irreducibilität der Kreistheilungsgleichung in dem allgemeinen Falle handelt, in dem der Grad der Einheitswurzel eine aus mehreren verschiedenen Primzahlen zusammengesetzte Zahl ist, einer Berichtigung, die hier nachgetragen werden soll, und wir bitten den Leser, die genannte Stelle nach dem, was hier ausgeführt wird, zu ergänzen.

Wenn $n=ab$ gesetzt wird, $a$ und $b$ relativ prim sind, und $\varrho, \alpha, \beta$ primitive $n^{\mathrm{te}}, a^{\mathrm{te}}, b^{\mathrm{te}}$ Einheitswurzeln sind, so ist $\varrho=\alpha\beta$, und es ist an der angeführten Stelle bewiesen, dass, wenn $\varrho$ Wurzel einer rationalen Gleichung $\Phi(x)=0$ ist, unter den Wurzeln dieser Gleichung $\varrho=\alpha\beta$ alle primitiven $a^{\mathrm{ten}}$ Einheitswurzeln $\alpha$ und alle primitiven $b^{\mathrm{ten}}$ Einheitswurzeln $\beta$ vorkommen müssen. Daraus folgt aber nicht, wie dort irrthümlich angenommen ist, dass darunter \gesperrt{alle} primitiven $n^{\mathrm{ten}}$ Einheitswurzeln vorkommen, weil nicht gezeigt ist, dass jedes $\alpha$ mit jedem $\beta$ combinirt vorkommt.

Zur Ausfüllung dieser Lücke möge hier zunächst ein ganz elementarer Beweis Platz finden.\footnote{\gesperrt{Arndt}, \gesperrt{Crelle}'s Journ., Bd. 56, S. 178 (1859). \gesperrt{Lebesgue}, \gesperrt{Liouville}'s Journ., 2. Ser., Bd. 4, S. 105 (1859).}

Ist $p$ eine Primzahl, so sind alle Binomialcoëfficienten für die $p^{\mathrm{te}}$ Potenz

\[
B_h^{(p)} = \dfrac{p(p-1)\ldots(p-h+1)}{1\,.\,2\ldots h},
\]

mit Ausnahme von $B_0^{(p)}$ und $B_p^{(p)}$, durch $p$ theilbare ganze Zahlen,
%CONT>>


% ===== scan leaf 790, printed 769 =====
% !PRINTED 769
\weberpage{769}
\pdfbookmark[0]{ Nachträge.}{bm237}
\pdfbookmark[1]{I. Irreducibilität der Kreistheilungsgleichung}{bm238}
%<<CONT

und folglich ist, wenn $u$, $v$ unbestimmte Grössen sind und die
\gesperrt{Congruenz nur auf die Coëfficienten bezogen wird},
\[
(u+v)^p \equiv u^p+v^p \ (\mathrm{mod}\ p).
\]

Ersetzen wir darin $v$ wieder durch eine Summe $v+w+\cdots$,
so erhalten wir durch vollständige Induction (oder auch aus
dem polynomischen Lehrsatze):
\[
(u+v+w+\cdots)^p \equiv u^p+v^p+w^p+\cdots \ (\mathrm{mod}\ p),
\]
und wenn wir diese Formel wiederholt anwenden, und unter $k$
einen beliebigen positiven Exponenten verstehen,
\weq{1}{(u+v+w+\cdots)^{p^k} \equiv u^{p^k}+v^{p^k}+w^{p^k}+\cdots \ (\mathrm{mod}\ p).}

Es sei jetzt $x$ eine Variable, $a_1, a_2, \dots, a_\nu$ ganze rationale
Zahlen,
\[
f(x) = x^\nu + a_1 x^{\nu-1} + a_2 x^{\nu-2} + \cdots + a_\nu
\]
eine ganze Function von $x$, deren Wurzeln $\alpha, \beta, \gamma, \dots$ sein
mögen, so dass
\[
-a_1 = \Sigma\alpha, \quad a_2 = \Sigma\alpha\beta, \quad -a_3 = \Sigma\alpha\beta\gamma \dots
\]
die symmetrischen Grundfunctionen der $\alpha, \beta, \gamma, \dots$ sind.

Wir bilden die Function
\[
F(x) = x^\nu + A_1 x^{\nu-1} + A_2 x^{\nu-2} + \cdots + A_\nu,
\]
deren Wurzeln die $p^{k\,\mathrm{ten}}$ Potenzen der Wurzeln von $f$, also
\[
\alpha^{p^k}, \ \beta^{p^k}, \ \gamma^{p^k}, \dots
\]
sind. Die Coëfficienten von $F$:
\[
A_1 = -\Sigma\alpha^{p^k}, \quad A_2 = \Sigma\alpha^{p^k}\beta^{p^k}, \quad A_3 = -\Sigma\alpha^{p^k}\beta^{p^k}\gamma^{p^k} \dots
\]
sind als ganzzahlige symmetrische Functionen der $\alpha, \beta, \gamma, \dots$
rational und ganzzahlig durch die $a_1, a_2, \dots$ darstellbar, und sind
daher ganze Zahlen. Die Quotienten
\[
\frac{A_1 - a_1^{p^k}}{p}, \quad \frac{A_2 - a_2^{p^k}}{p}, \quad \frac{A_3 - a_3^{p^k}}{p}, \dots
\]
sind aber nach der Formel (1) gleichfalls ganze Zahlen, und es
ist also nach dem \gesperrt{Fermat'schen Satze}
\weq{2}{F(x) \equiv f(x) \ (\mathrm{mod}\ p).}

Diese Sätze sollen nun auf die Function $f_n(x)$ angewandt
werden, deren Wurzeln die sämmtlichen \gesperrt{primitiven} $n^{\mathrm{ten}}$ Einheitswurzeln $\alpha, \beta, \gamma, \dots$ sind, und die, wie im \S.\,133 des ersten
Bandes bewiesen ist, ganzzahlige Coëfficienten $a_1, a_2, \dots$ hat.


% ===== scan leaf 791, printed 770 =====
% !PRINTED 770
\weberpage{770}

Ist $p$ eine in $n$ aufgehende Primzahl, so ist
\weq{3}{\frac{x^n-1}{x^{\frac{n}{p}}-1} = x^{\frac{n}{p}(p-1)} + x^{\frac{n}{p}(p-2)} + \cdots + x^{\frac{n}{p}} + 1,}
und diese Function verschwindet für $x=\alpha, \beta, \gamma, \dots$ und ist
daher durch $f_n(x)$ theilbar. Wir wollen sie $= f_n(x)\, g(x)$ setzen,
worin dann $g(x)$ (nach Bd. I, \S.\,2) eine ganze Function von $x$
mit ganzzahligen Coëfficienten ist.

Ist nun
\[
n = p^k n',
\]
$n'$ nicht mehr durch $p$ theilbar, und $\alpha'$ eine primitive $n'^{\mathrm{te}}$ Einheitswurzel, so ist
\[
\alpha'^{\frac{n}{p}} = 1,
\]
und folglich, wenn man $x = \alpha'$ setzt [nach (3)]:
\weq{4}{f_n(\alpha')\, g(\alpha') = p.}

Wir leiten hieraus zunächst einen Beweis für die Irreducibilität von $f_n(x)$ unter der Voraussetzung ab, dass $n$ eine
\gesperrt{Primzahlpotenz} ist, wenn auch für diesen Fall der in Bd. I,
\S.\,134 gegebene Beweis einwandsfrei ist.

Angenommen, es zerfalle $f_n(x)$ in zwei rationale Factoren:
\weq{5}{f_n(x) = \varphi(x)\, \psi(x).}

Dann sind alle Wurzeln, sowohl von $\varphi(x)$ als von $\psi(x)$,
$n^{\mathrm{te}}$ Einheitswurzeln, und ihre $n^{\mathrm{ten}}$ Potenzen also gleich 1. Wenn
nun $n = p^k$ ist, so können wir aus $\varphi(x)$ und $\psi(x)$ zwei ganze
Functionen $\Phi(x)$, $\Psi(x)$ ableiten, deren Wurzeln die $n^{\mathrm{ten}}$ Potenzen der Wurzeln von $\varphi(x)$ und $\psi(x)$ sind und die daher alle
gleich 1 sind, und aus (2) ergiebt sich:
\[
\varphi(x) \equiv \Phi(x), \quad \psi(x) \equiv \Psi(x) \ (\mathrm{mod}\ p),
\]
woraus für $x = 1$ folgt:
\[
\varphi(1) \equiv 0, \quad \psi(1) \equiv 0 \ (\mathrm{mod}\ p).
\]
Hiernach ergiebt sich aus (5):
\weq{6}{f_n(1) \equiv 0 \ (\mathrm{mod}\ p^2).}

Nun ist aber in diesem Falle $f_n(1) = p$ (Bd. I, \S.\,133, V.),
was der Formel (6) widerspricht. Unsere Annahme war also
unzulässig und $f_n(x)$ ist irreducibel.


% ===== scan leaf 792, printed 771 =====
% !PRINTED 771
\weberpage{771}

Um die Irreducibilität für ein beliebig zusammengesetztes $n$
nachzuweisen, wenden wir die vollständige Induction an. Wir
setzen
\[
n = p^k n',
\]
und verstehen unter $p$ eine Primzahl, die in $n$, aber nicht in $n'$
aufgeht, und setzen die \gesperrt{Irreducibilität von} $f_{n'}(x)$ \gesperrt{schon als
bewiesen voraus}. Ist dann $f_n(x)$ in zwei rationale Factoren
$\varphi(x)$, $\psi(x)$ zerlegbar, deren keiner constant ist:
\weq{7}{f_n(x) = \varphi(x)\, \psi(x),}
so leiten wir aus $\varphi(x)$ die Function $\Phi(x)$ ab, deren Wurzeln
die $p^{k\,\mathrm{ten}}$ Potenzen der Wurzeln von $\varphi(x)$ sind, die der Bedingung
\[
\varphi(x) \equiv \Phi(x) \ (\mathrm{mod}\ p)
\]
genügt, oder auch, indem wir mit $\chi(x)$ eine zweite ganzzahlige
Function von $x$ bezeichnen:
\[
\varphi(x) = \Phi(x) + p\, \chi(x).
\]

Die Wurzeln von $\Phi(x)$ sind aber primitive Einheitswurzeln
vom Grade $n'$, und daher muss unter diesen wenigstens eine
sein, $\alpha'$, für die $\Phi(\alpha')$ verschwindet, also
\[
\varphi(\alpha') = p\, \chi(\alpha')
\]
ist. Wegen der vorausgesetzten Irreducibilität von $f_{n'}(x)$ muss
aber diese Gleichung für \gesperrt{alle} Einheitswurzeln $\alpha'$ bestehen, und
aus dem gleichen Grunde ist, wenn $\vartheta(x)$ eine ganzzahlige
Function ist, für alle $\alpha'$:
\[
\psi(\alpha') = p\, \vartheta(\alpha').
\]

Hiernach folgt aus (7):
\[
f_n(\alpha') = p^2\, \chi(\alpha')\, \vartheta(\alpha'),
\]
und nach (4):
\[
1 = p\, \chi(\alpha')\, \vartheta(\alpha')\, g(\alpha').
\]

Das Product $\chi(x)\, \vartheta(x)\, g(x)$ ist aber eine ganze ganzzahlige
Function von $x$, deren Rest bei der Division mit $f_{n'}(x)$ gleichfalls ganzzahlige Coëfficienten hat und mit $\Theta(x)$ bezeichnet sein
mag. Dieser Rest ist von niedrigerem Grade, als $f_{n'}(x)$, und die
Function
\[
p\, \Theta(x) - 1
\]
verschwindet für \gesperrt{alle Wurzeln} von $f_{n'}(x)$. Folglich muss die
Gleichung
\[
p\, \Theta(x) = 1
\]
identisch sein, was aber offenbar unmöglich ist, da $\Theta(x)$ für ein
%CONT>>


% ===== scan leaf 793, printed 772 =====
% !PRINTED 772
\weberpage{772}

%<<CONT
ganzzahliges $x$ in eine ganze Zahl übergeht, die, mit $p$ multiplicirt, nicht $= 1$ sein kann.

Hiermit ist die Irreducibilität von $f_n(x)$ allgemein bewiesen.

\begin{center}
\rule{4cm}{0.4pt}
\end{center}

Nachdem so die Irreducibilität der Kreistheilungsgleichung
$f_n(x) = 0$ allgemein nachgewiesen ist, folgt, wie im \S.\,134 des
ersten Bandes, dass sie auch irreducibel bleibt, wenn eine beliebige Einheitswurzel adjungirt wird, deren Grad zu $n$ relativ
prim ist.

\gesperrt{Kronecker} geht gewissermaassen den umgekehrten Weg\footnote{Mém. sur les facteurs irréd. de l'expression $(x^n-1)$. \gesperrt{Liouville}'s
Journal, Bd. 19 (1854). Einen auf der Theorie der höheren Congruenzen
beruhenden einfachen Beweis der Irreducibilität der allgemeinen Kreistheilungsgleichung hat \gesperrt{Dedekind} gegeben [\gesperrt{Crelle}'s Journal für Mathematik, Bd. 58 (1859)].}.

Er beweist zunächst unter der Voraussetzung, dass $n$ eine
Potenz einer Primzahl $p$ ist, die Irreducibilität von $f_n(x)$ nach
Adjunction einer Einheitswurzel, deren Grad nicht durch $p$
theilbar ist, und leitet daraus den allgemeinen Irreducibilitätsbeweis ab, wie wir jetzt noch kurz zeigen wollen.

\begin{center}
\rule{4cm}{0.4pt}
\end{center}

Es sei jetzt
\[
n = p^k
\]
eine Potenz einer Primzahl $p$ und
\weq{8}{f_n(x) = x^{p^{k-1}(p-1)} + x^{p^{k-1}(p-2)} + \cdots + x^{p^{k-1}} + 1}
die Function, deren Wurzeln die primitiven $n^{\mathrm{ten}}$ Einheitswurzeln
sind. Es sei ferner $\alpha$ irgend eine Einheitswurzel, deren Grad $m$
durch $p$ nicht theilbar ist, und $\nu$ der Grad des Körpers $R(\alpha)$.
\gesperrt{Es soll bewiesen werden, dass $f_n(x)$ im Körper $R(\alpha)$ irreducibel ist.}

Jede Zahl des Körpers $R(\alpha)$ lässt sich auf eine und nur
auf eine Weise in der Form darstellen:
\weq{9}{\varphi(\alpha) = m_0 + m_1 \alpha + \cdots + m_{\nu-1} \alpha^{\nu-1},}
worin $m_0, m_1, \dots, m_{\nu-1}$ rationale Zahlencoëfficienten sind.

Ist $\chi(x) = 0$ die rationale Gleichung niedrigsten Grades,
deren Wurzel $\alpha$ ist, und
\weq{10}{\chi(x) = x^\nu + c_1 x^{\nu-1} + \cdots + c_{\nu-1} x + c_\nu,}
%CONT>>


% ===== scan leaf 794, printed 773 =====
% !PRINTED 773
\weberpage{773}
%<<CONT
so sind die $c_1, c_2, \ldots$ \gesperrt{ganze Zahlen}, weil $\chi(x)$ ein Theiler von $x^m - 1$ sein muss (nach dem \gesperrt{Gauss}'schen Satze, Bd.~I, \S.\,2).

Daraus ergiebt sich, dass, wenn $\varphi(x)$ eine ganze Function von $x$ mit ganzzahligen rationalen Coëfficienten bedeutet, auch wenn der Grad von $\varphi(x)$ ursprünglich höher als $\nu$ ist, die Coëfficienten $m_0, m_1, \ldots, m_{\nu-1}$ in (9) ganze Zahlen werden; denn diese Coëfficienten erhält man, wenn man den Rest der Division von $\varphi(x)$ durch $\chi(x)$ aufsucht.

Wir nehmen nun an, $f_n(x)$ sei im Körper $R(\alpha)$ zerlegbar, und es sei demnach:

\weq{11}{f_n(x) = \varphi(x,\alpha)\ \psi(x,\alpha),}

worin $\varphi(x,\alpha)$ und $\psi(x,\alpha)$ zwei ganze Functionen von $x$ sind, deren keine constant ist, deren Coëfficienten in $R(\alpha)$ enthalten sind. Aus (8) und (11) ergiebt sich, wenn wir $x=1$ setzen,

\weq{12}{p = \varphi(1,\alpha)\ \psi(1,\alpha),}

worin $\varphi(1,\alpha)$, $\psi(1,\alpha)$ als Zahlen des Körpers $R(\alpha)$ in der Form (9) darstellbar sind. Bezeichnen wir die kleinsten gemeinschaftlichen Nenner dieser beiden Darstellungen mit $a$ und $b$, so erhalten wir:

\weq{13}{\begin{aligned}
a\,\varphi(1,\alpha) &= a_0 + a_1\alpha + a_2\alpha^2 + \cdots + a_{\nu-1}\alpha^{\nu-1} = A(\alpha)\\
b\,\psi(1,\alpha) &= b_0 + b_1\alpha + b_2\alpha^2 + \cdots + b_{\nu-1}\alpha^{\nu-1} = B(\alpha),
\end{aligned}}

worin $a, a_0, a_1, \ldots, a_{\nu-1}$ ganze Zahlen ohne gemeinsamen Theiler und $A$ ein Functionszeichen ist. Gleiches gilt für $b, b_0, b_1, \ldots, b_{\nu-1}$ und $B$.

Die Wurzeln der Gleichung $\varphi(x,\alpha) = 0$ finden sich unter den primitiven $n^{\text{ten}}$ Einheitswurzeln, und diese sind sämmtlich Potenzen von einer unter ihnen, $r$. Demnach giebt es gewisse Potenzen $r^h, r^{h'}, r^{h''}, \ldots$, so dass

\[
\varphi(x,\alpha) = (x - r^h)(x - r^{h'})(x - r^{h''}) \ldots
\]

wird, und folglich nach (13):

\[
A(\alpha) = a(1 - r^h)(1 - r^{h'})(1 - r^{h''}) \ldots
\]

Nehmen wir hiervon die $n^{\text{te}}$ Potenz, so ergiebt sich nach (1), mit Rücksicht auf die Gleichung $r^n = 1$:

\[
[A(\alpha)]^n = p\,\chi(r),
\]

worin $\chi(x)$ eine ganze Function der Variablen $x$ mit ganzzahligen Coëfficienten bedeutet.


% ===== scan leaf 795, printed 774 =====
% !PRINTED 774
\weberpage{774}
Nun setzen wir, indem wir unter $t$ eine Variable verstehen:

\[
\begin{aligned}
{}[t - \chi(1)][t - \chi(r)] \ldots [t - \chi(r^{n-1})] &= \\
t^n + A_1t^{n-1} + A_2t^{n-2} + \cdots + A_n, &
\end{aligned}
\]

worin die Coëfficienten $A_1, A_2, \ldots, A_n$ als ganze symmetrische Functionen der $n$ Grössen $\chi(1), \chi(r), \ldots, \chi(r^{n-1})$ ganze rationale Zahlen sind. Wenn wir aber darin

\[
t = \chi(r) = \dfrac{[A(\alpha)]^n}{p}
\]

setzen, so verschwindet die linke Seite, und es ergiebt sich durch Multiplication mit $p^n$:

\weq{14}{[A(\alpha)]^{n^2} = p\,C(\alpha),}

wenn $C(\alpha)$ eine Zahl von der Form (9) mit \gesperrt{ganzzahligen Coëfficienten} bedeutet.

Wenn man die Formel (14) wiederholt in die $p^{\text{te}}$ Potenz erhebt und den Satz (1) anwendet, so ergiebt sich daraus für jede Potenz, die nicht kleiner als $n^2$ ist:

\weq{15}{A(\alpha^{p^\lambda}) = p\,\xi(\alpha),}

worin die ganze Function $\xi(\alpha)$ gleichfalls ganzzahlige Coëfficienten hat.

Hierin können wir $\lambda$ durch $\varphi(m)$ theilbar annehmen, und dann ist, da $m$ durch $p$ nicht theilbar ist, nach dem verallgemeinerten \gesperrt{Fermat}'schen Satze (Bd.~II, \S.\,14):

\[
p^\lambda \equiv 1 \ (\mathrm{mod}\ m), \quad \alpha^{p^\lambda} = \alpha,
\]

so dass aus (14) folgt: % [sic] im Druck; die Gleichung (16) folgt sachlich aus (15), nicht aus (14)

\weq{16}{A(\alpha) = p\,\xi(\alpha).}

Es sind daher in (13) die Coëfficienten $a_0, a_1, \ldots, a_{\nu-1}$ alle durch $p$ theilbar und $a$ ist folglich nicht durch $p$ theilbar. Ebenso wird gezeigt, dass $b_0, b_1, \ldots, b_{\nu-1}$ durch $p$ theilbar sind, während $b$ durch $p$ untheilbar ist.

Demnach ergiebt sich aus (12), wenn $A(\alpha)B(\alpha) = p^2\Theta(\alpha)$ gesetzt wird:

\weq{17}{ab = p\,\Theta(\alpha),}

worin $\Theta(\alpha)$ eine Function von $\alpha$ von der Form (9) mit ganzzahligen Coëfficienten bedeutet. Die Gleichung (17) müsste aber, da sie höchstens vom $(\nu-1)^{\text{ten}}$ Grade ist, in Bezug auf $\alpha$ identisch sein, und es müsste $ab$ durch $p$ theilbar sein, was nicht möglich ist.


% ===== scan leaf 796, printed 775 =====
% !PRINTED 775
\weberpage{775}
Hieraus folgt die Unzulässigkeit unserer Annahme, dass $f_n(x)$ nach der Formel (11) im Körper $R(\alpha)$ in zwei Factoren zerlegbar ist.

\begin{center}
\rule{4cm}{0.4pt}
\end{center}

Damit lässt sich nun aufs Neue allgemein beweisen, dass die Gleichung, deren Wurzeln die sämmtlichen primitiven $m^{\text{ten}}$ Einheitswurzeln sind, deren Grad $\varphi(m)$ ist, auch wenn $m$ eine beliebige zusammengesetzte Zahl ist, irreducibel ist. Dies ist mit folgendem Satze gleichbedeutend:

\gesperrt{Wenn eine ganze Function $\varphi(x)$ mit rationalen Coëfficienten verschwindet, wenn für $x$ eine primitive $m^{\text{te}}$ Einheitswurzel gesetzt wird, so verschwindet $\varphi(x)$ auch, wenn für $x$ irgend eine andere primitive $m^{\text{te}}$ Einheitswurzel gesetzt wird.}

Wenn $m$ eine Primzahlpotenz ist, so ist der Satz in dem Vorhergehenden schon bewiesen. Wir wenden daher die vollständige Induction an, nehmen ihn für $m^{\text{te}}$ Einheitswurzeln als erwiesen an und leiten ihn für $mn^{\text{te}}$ Einheitswurzeln her, wenn $n = p^k$ und $p$ eine in $m$ nicht aufgehende Primzahl ist.

Jede $mn^{\text{te}}$ Einheitswurzel $\beta$ kann in der Form

\[
\beta = \alpha r
\]

dargestellt werden, worin $\alpha$ eine $m^{\text{te}}$, $r$ eine $n^{\text{te}}$ Einheitswurzel ist. Was zu beweisen ist, lässt sich so aussprechen, dass, wenn $\beta$ die Wurzel einer rationalen Gleichung $\Phi(x) = 1$ ist, % [sic]
auch alle anderen primitiven Einheitswurzeln vom Grade $mn$ Wurzeln dieser Gleichung sein müssen.

Ist, wie oben, $f_n(x) = 0$ die Gleichung, deren Wurzeln die Grössen $r$ sind, so ist, wie wir bewiesen haben, $f_n(x)$ nicht in rationale Factoren zerlegbar, auch nicht in solche, die rational von $\alpha$ abhängen. Die Function $f_n(\alpha^{-1}x)$, die für $x = \beta$ verschwindet, kann also auch nicht in rationale Factoren zerlegt werden, weil aus einer Zerlegung der Form

\[
f_n(\alpha^{-1}x) = \varphi(x,\alpha)\ \psi(x,\alpha)
\]

folgen würde:

\[
f_n(x) = \varphi(\alpha x,\alpha)\ \psi(\alpha x,\alpha),
\]

was nicht bestehen kann. Da hiernach $f_n(\alpha^{-1}x)$ eine Wurzel
%CONT>>


% ===== scan leaf 797, printed 776 =====
% !PRINTED 776
\weberpage{776}
%<<CONT
mit $\Phi(x)$ gemein hat, so muss $\Phi(x)$ durch $f_n(\alpha^{-1}x)$ theilbar sein, und wir können demnach für ein unbestimmtes $x$

\weq{18}{\Phi(x) = f_n(\alpha^{-1}x)\ \Psi(x,\alpha)}

setzen, worin $\Psi(x,\alpha)$ eine ganze rationale Function von $x$ ist.

Sind $\alpha, \alpha_1, \alpha_2, \ldots$ die sämmtlichen primitiven $m^{\text{ten}}$ Einheitswurzeln, die, wie wir angenommen haben, Wurzeln einer \gesperrt{irreduciblen} rationalen Gleichung

\[
F_m(x) = 0
\]

sind, so kann $\alpha$ nach dieser Voraussetzung in der Identität (18) durch jedes andere $\alpha_1, \alpha_2, \ldots$ ersetzt werden, und es folgt, dass $\Phi(x)$ durch jeden der Factoren

\[
f_n(\alpha^{-1}x),\ f_n(\alpha_1^{-1}x),\ f_n(\alpha_2^{-1}x) \ldots
\]

theilbar ist. Zwei dieser Factoren können aber keinen gemeinschaftlichen Theiler haben, denn wenn etwa

\[
f_n(\alpha^{-1}x) = 0,\ f_n(\alpha_1^{-1}x) = 0
\]

eine gemeinsame Wurzel $x$ hätten, so müssten zwei primitive $n^{\text{te}}$ Einheitswurzeln $r, r_1$ existiren, so dass

\[
x = \alpha r = \alpha_1 r_1
\]

wäre; es müsste also $\alpha^n = \alpha_1^n$, und folglich, da $n$ relativ prim zu $m$ ist, $\alpha = \alpha_1$ sein.

Demnach ist $\Phi(x)$ theilbar durch das Product

\[
F_{mn}(x) = f_n(\alpha^{-1}x)\ f_n(\alpha_1^{-1}x)\ f_n(\alpha_2^{-1}x) \ldots,
\]

dessen Wurzeln die sämmtlichen $\beta$ sind, und diese Function ist daher auch unzerlegbar, wie bewiesen werden sollte.

\begin{center}
\rule{4cm}{0.4pt}
\end{center}


% ===== scan leaf 798, printed 777 =====
% !PRINTED 777
\weberpage{777}
\pdfbookmark[1]{II. Die Irreducibilität der Kreistheilungsgleichung und der Satz über die in einer Linearform enthaltenen Primzahlen}{bm239}

\begin{center}
II.

\gesperrt{Die Irreducibilität der Kreistheilungsgleichung
und der Satz über die in einer Linearform enthaltenen
Primzahlen.}
\end{center}

Während im ersten Nachtrage, der als Ergänzung zum ersten Bande zu betrachten
ist, absichtlich kein Gebrauch von der Theorie der ganzen algebraischen Zahlen
gemacht ist, geben wir zum Schluss noch einen auf anderen Grundlagen, nämlich
auf den Sätzen über die Classenzahl (\S.\,192) beruhenden Beweis für die
Irreducibilität der Kreistheilungsgleichung, der wegen der Verallgemeinerungen,
die er zulässt, merkwürdig ist.

Es sei $m$ eine beliebige natürliche Zahl, und $n$ sei das System der zu $m$
theilerfremden positiven Zahlen; unter diesen giebt es
\weq{1}{\varphi(m) = \mu,}
die kleiner als $m$ sind, die wir mit $a$ bezeichnen.

Wenn wir alle nach dem Modul $m$ mit einem $a$ congruenten Zahlen $n$ in eine
Classe $A$ vereinigen, so erhalten wir $\mu$ Zahlclassen:
\weq{2}{A_1, A_2, \ldots, A_\mu,}
die bei der Composition durch Multiplication eine \gesperrt{Abel}'sche Gruppe
$\mu^{\text{ten}}$ Grades, $\mathfrak{N}$, bilden. Diese Gruppe ist schon im
\S.\,16 dieses Bandes betrachtet.

Die $\mu$ Charaktere dieser Gruppe bezeichnen wir mit
\weq{3}{\chi_1, \chi_2, \ldots, \chi_\mu,}
und setzen, wenn $\chi$ einer dieser Charaktere ist, und $n$ eine Zahl aus der
Classe $A$ bedeutet,
\weq{4}{\chi(A) = \chi(n).}
Diese Charaktere, die im \S.\,16 näher bestimmt sind, sind sämmtlich
$\mu^{\text{te}}$ Einheitswurzeln.


% ===== scan leaf 799, printed 778 =====
% !PRINTED 778
\weberpage{778}
% running head "II. Nachtrag." at the top of this page is not transcribed as body text.

Wir haben aber hier nicht nöthig, von den speciellen Ausdrücken dieser
Charaktere Gebrauch zu machen. Dahingegen stützen wir uns auf folgenden
Lehrsatz, der für alle \gesperrt{Abel}'schen Gruppen gilt:

1. \gesperrt{Ist} $A$ \gesperrt{ein Element} $f^{\text{ten}}$ \gesperrt{Grades
einer} \gesperrt{Abel}'\gesperrt{schen Gruppe} $\mu^{\text{ten}}$
\gesperrt{Grades, und ist} $\mu = ef$, \gesperrt{so sind alle} $\chi_i(A)$
$f^{\text{te}}$ \gesperrt{Einheitswurzeln, und darunter kommt jede}
$f^{\text{te}}$ \gesperrt{Einheitswurzel genau} $e$ \gesperrt{mal vor.}

Der Satz ist implicite in dem Satze 7., \S.\,12 dieses Bandes enthalten. Denn
wenden wir diesen Satz auf die Gruppe
\[
T = 1, A, A^2, \ldots, A^{f-1}
\]
an, deren Index in Bezug auf $\mathfrak{N}$ gleich $e$ ist, so folgt, dass es
eine Gruppe $\Xi$ von $e$ Charakteren $\xi$ giebt, die den
Bedingungen
\[
\xi(A) = 1
\]
genügen.

Zerlegt man also die Gruppe $X$ der Charaktere $\chi$ in die Nebengruppen
\[
\Xi_1, \Xi_2, \ldots, \Xi_f,
\]
so sind $\chi_1(A)$ und $\chi_2(A)$ einander gleich oder von einander
verschieden, je nachdem $\chi_1$ und $\chi_2$ in derselben oder in
verschiedenen dieser Nebengruppen vorkommen. Da ausserdem wegen
$\chi(A)^f = \chi(A^f) = 1$ alle $\chi(A)$ Einheitswurzeln vom Grade $f$
sind, und es nur $f$ verschiedene solche giebt, so ist damit unser Theorem
bewiesen.

Ist $n$ in der Classe $A$ enthalten, so ist der Grad $f$ von $A$ der
Exponent, zu dem $n$ nach dem Modul $m$ \gesperrt{gehört}, d.\,h. der kleinste positive
Exponent, für den
\weq{5}{n^f \equiv 1 \pmod{m}}
ist. Wenn also $n$ zum Exponenten $f$ gehört, so sind alle $\chi_i(n)$
$f^{\text{te}}$ Einheitswurzeln, und jede $f^{\text{te}}$ Einheitswurzel
kommt darunter $e$ mal vor.

Bedeutet $a$ irgend einen der $\varphi(m)$ Reste der Zahlen $n$, so ist in der
Form
\weq{6}{\mu_x = mx + a - m}
jede Zahl aus der durch $a$ repräsentirten Zahlclasse (2) enthalten;
verstehen wir unter $a$ den kleinsten positiven Rest, so erhalten wir die
Gesammtheit der Zahlen dieser Classe, wenn wir $x$ alle positiven
ganzzahligen Werthe durchlaufen lassen.


% ===== scan leaf 800, printed 779 =====
% !PRINTED 779
\weberpage{779}
% running head "Primzahlen in Linearformen." / "779" at the top of this page is not
% transcribed as body text.

Nun ist
\[
\frac{\mu_x}{m} - x = \frac{a-m}{m},
\]
und folglich können wir auf die Grössen $\mu_x$ den Satz \S.\,191, 6.
anwenden, in dem wir $n = x$, $\gamma = 1:m$, $c_n = (a-m):m$ zu setzen
haben. Es ist also
\weq{7}{A(s) = \overset{x}{\underset{0,\,\infty}{\Sigma}} \frac{1}{(mx+a)^s}
= \overset{x}{\underset{1,\,\infty}{\Sigma}} \frac{1}{(mx+a-m)^s},}
so lange $s > 1$ ist, eine stetige Function von $s$, und es ist
\weq{8}{A(s) = \frac{1}{m(s-1)} + C(s),}
worin $C(s)$ für $s=1$ endlich bleibt\footnote{Durch Benutzung bekannter
Sätze über $\Gamma$-Functionen findet man nach \S.\,191, (20):
\[
C(1) = \frac{1}{m} \frac{\Gamma'\left(\frac{a}{m}\right)}{\Gamma\left(\frac{a}{m}\right)}.
\]}.

Wir lassen nun $A$ die sämmtlichen Classen (2) durchlaufen, bezeichnen mit
$\chi$ irgend einen der Charaktere der Gruppe $\mathfrak{N}$ und setzen
\weq{9}{Q(s) = \overset{A}{\Sigma} \chi(A)\, A(s),}
oder, was dasselbe ist,
\weq{10}{Q(s) = \overset{n}{\Sigma} \frac{\chi(n)}{n^s},}
worin sich die Summe auf alle Zahlen $n$ erstreckt. Die Anzahl dieser Summen
$Q$ ist so gross, als die Anzahl der Charaktere, d.\,h. $= \varphi(m)$. Wir
bezeichnen sie, wenn eine Unterscheidung nöthig ist, mit
\[
Q_1, Q_2, \ldots, Q_\mu.
\]

Eine von ihnen, $Q_1$, entspricht dem Hauptcharakter und hat den Ausdruck
\[
Q_1(s) = \overset{n}{\Sigma} \frac{1}{n^s}.
\]

Nun ist [\S.\,11, (21)]:
\[
\overset{A}{\Sigma} \chi(A) = \mu \ \text{oder} = 0,
\]
je nachdem $\chi$ der Hauptcharakter ist oder nicht. Ferner ergiebt sich aus
(8)
%CONT>>


% ===== scan leaf 801, printed 780 =====
% !PRINTED 780
\weberpage{780}
% running head "II. Nachtrag." at the top of this page is not transcribed as body text.
%<<CONT
\[
Q(s) = \frac{1}{m(s-1)} \overset{A}{\Sigma} \chi(A) + \overset{A}{\Sigma} \chi(A)\, C(s),
\]
und daraus der Satz:

2. \gesperrt{Die Summen $Q_2, \ldots, Q_\mu$ erhalten für $s=1$ endliche
Werthe,} $Q_1$ \gesperrt{wird unendlich, und zwar wird}
\[
\mathrm{Lim}_{s=1} (s-1)\, Q_1(s) = \frac{\mu}{m}.
\]

Um die Summen $Q(s)$ umzuformen, bezeichnen wir mit $p$ jede in $m$ nicht
aufgehende Primzahl und multipliciren alle die unendlichen Reihen von der
Form
\[
1 + \frac{\chi(p)}{p^s} + \frac{\chi(p^2)}{p^{2s}} + \cdots
= \frac{1}{1 - \chi(p)\, p^{-s}}
\]
mit einander. Dadurch erhalten wir [vgl. \S.\,192, (8)]:
\weq{11}{Q(s) = \overset{p}{\Pi} \frac{1}{1 - \chi(p)\, p^{-s}}.}
Wenn nun $p$ zum Exponenten $f$ gehört, wenn also $f$ die kleinste positive,
der Congruenz
\weq{12}{p^f \equiv 1 \pmod{m}}
genügende Zahl, und $\mu = ef$ ist, so kommt unter den $\chi(p)$ jede
$f^{\text{te}}$ Einheitswurzel genau $e$ mal vor, und es ist also (für ein
variables $x$):
\[
\overset{\chi}{\Pi} [1 - \chi(p)\, x] = (1 - x^f)^e,
\]
demnach ergiebt sich aus (11):
\weq{13}{Q_1 Q_2 \ldots Q_\mu = \overset{p}{\Pi} \frac{1}{(1-p^{-fs})^e}.}

Wir führen nun, ohne die Irreducibilität der Kreistheilungsgleichung
vorauszusetzen, den Kreistheilungskörper $\Omega_m$ ein, dessen noch
unbekannten Grad wir mit $\nu$ bezeichnen wollen. Da aber die primitive
$m^{\text{te}}$ Einheitswurzel $r$ einer rationalen Gleichung
$\mu^{\text{ten}}$ Grades genügt, so ist
\weq{14}{\nu \leqq \mu.}

Die Zahlen $\omega$ des Körpers $\Omega_m$ sind alle in der Form
\weq{15}{a\omega = a_0 + a_1 r + \cdots + a_{\nu-1} r^{\nu-1}}
mit ganzen rationalen Coëfficienten $a, a_0, a_1, \ldots, u_{\nu-1}$ % [sic] -- printed $u_{\nu-1}$; display (15) has $a_{\nu-1}$
darstellbar, und wenn $\omega$ eine \gesperrt{ganze Zahl} ist, so können wir
für $a$ eine feste ganze Zahl setzen, deren nähere Kenntniss nicht erforder-
%CONT>>


% ===== scan leaf 802, printed 781 =====
% !PRINTED 781
\weberpage{781}
%<<CONT
lich ist. Ihr Werth ist gleich der Quadratwurzel aus dem
Quotienten der Discriminanten
\[
\Delta (1, r, r^2, \ldots, r^{\nu-1}) : \Delta (\omega_1, \omega_2, \ldots, \omega_\nu),
\]
worin $\omega_1, \omega_2, \ldots, \omega_\nu$ eine Minimalbasis ist (\S.\,144, 145).

Wir schliessen alle in der Discriminante der Kreistheilungs\-gleichung und alle in $m$ und in $a$ aufgehenden Primzahlen, deren Anzahl jedenfalls endlich ist, von dem Systeme der $p$ aus, und bezeichnen mit $\mathfrak{p}$ ein in $p$ aufgehendes Primideal. Dann ist die Congruenz
\[
r^p \equiv r \ (\mathrm{mod}\ \mathfrak{p})
\]
immer dann, aber auch nur dann, erfüllt, wenn

\weq{16}{p \equiv 1 \ (\mathrm{mod}\ m)}

ist, d.\,h. wenn $p$ zum Exponenten 1 gehört, wenn also $r^p = r$ ist, und unter derselben Bedingung ist daher auch nach (15) für jede ganze Zahl $\omega$ in $\Omega_m$:
\[
\omega^p \equiv \omega \ (\mathrm{mod}\ \mathfrak{p}).
\]
Dies ist aber nach \S.\,162 die nothwendige und hinreichende Bedingung dafür, dass die in der Körperdiscriminante von $\Omega_m$ nicht aufgehende Primzahl $p$ in lauter Primideale ersten Grades zerfällt.

Bezeichnen wir also mit
\[
P_f(s) = \prod (1 - p^{-sf})
\]
das über alle zum Exponenten $f$ gehörigen Primzahlen $p$ erstreckte Product, so ergiebt sich aus \S.\,192, dass, wenn $f>1$ ist, $P_f(1)$ endlich und von Null verschieden ist, und dass

\weq{17}{\dfrac{s-1}{P_1(s)^{\nu}}}

\gesperrt{für} $s=1$ \gesperrt{endlich und von Null verschieden ist.}

Nun ist nach (13), wenn mit $S$ das über die \gesperrt{ausgeschlossenen} Primzahlen $p$ erstreckte \gesperrt{endliche} Product $\Pi(1-p^{-sf})^{-e}$ bezeichnet wird,
\[
Q_1 Q_2 \ldots Q_\mu = S \overset{f}{\Pi} \frac{1}{P_f(s)^{e}},
\]
und da $e = \mu$ für $f = 1$, so ergiebt sich aus (17):

\weq{18}{(s-1)^{\frac{\mu}{\nu}} Q_1 Q_2 \ldots Q_\mu}

\gesperrt{für} $s = 1$ \gesperrt{endlich und von Null verschieden.}


% ===== scan leaf 803, printed 782 =====
% !PRINTED 782
\weberpage{782}
Andererseits ist aber nach 2.:

\weq{19}{(s-1) Q_1 Q_2 \ldots Q_\mu}

für $s = 1$ endlich, und wenn wir also (19) durch (18) dividiren, so folgt:
\[
(s-1)^{1-\frac{\mu}{\nu}} \text{ für } s = 1 \text{ endlich,}
\]
d.\,h.

\weq{20}{\nu \geqq \mu.}

Dies mit (14) zusammen ergiebt aber $\nu = \mu$, wodurch der folgende Satz bewiesen ist:

3. \gesperrt{Der Grad des Kreistheilungskörpers $\Omega_m$ ist gleich $\varphi(m)$, also die Kreistheilungsgleichung $\varphi(m)^{\mathrm{ten}}$ Grades, deren Wurzeln die primitiven $m^{\mathrm{ten}}$ Einheitswurzeln sind, irreducibel.}

Diese Betrachtung leistet uns aber noch einen anderen wichtigen Dienst, indem sie uns den Beweis des berühmten Satzes liefert, \gesperrt{dass in jeder arithmetischen Progression, deren Differenz und Anfangsglied natürliche Zahlen ohne gemeinsamen Theiler sind, unendlich viele Primzahlen enthalten sind}\textsuperscript{1)}.

Es folgt nämlich jetzt aus (18), dass das Product
\[
(s-1) Q_1 Q_2 \ldots Q_\mu
\]
für $s = 1$ einen endlichen und von Null verschiedenen Werth hat, und da nach 2. keiner der Factoren unendlich ist, so folgt:

4. \gesperrt{Die Summen}
\[
(s-1) Q_1(s), Q_2(s), \ldots, Q_\mu(s)
\]
\gesperrt{haben für $s = 1$ endliche und von Null verschiedene Werthe.}

Wenden wir auf alle Factoren von (11) die Formel an:
\[
\log \left[1 - \chi(p) p^{-s}\right] = \frac{\chi(p)}{p^s} + \frac{1}{2} \frac{\chi(p^2)}{p^{2s}} + \frac{1}{3} \frac{\chi(p^3)}{p^{3s}} + \cdots % [sic] im Druck fehlt das mathematisch erforderliche Minuszeichen vor \log
\]

\noindent\rule{4cm}{0.4pt}

\textsuperscript{1)} Der erste vollständige und allgemeine Beweis dieses Satzes ist von \gesperrt{Dirichlet} gegeben (Abhandlungen der Berl. Akademie vom Jahre 1837. Gesammelte Werke Nr.\ XXI). Auf die wesentliche Vereinfachung dieses Beweises durch Benutzung der \gesperrt{Kummer}'schen Formeln für die Classenzahl in den Kreistheilungskörpern hat \gesperrt{Dedekind} aufmerksam gemacht (Vorlesungen über Zahlentheorie, 3. Auflage, S. 596).
%CONT>>


% ===== scan leaf 804, printed 783 =====
% !PRINTED 783
\weberpage{783}
%<<CONT
(worin der imaginäre Theil des Logarithmus zwischen $\pm \pi i$ zu nehmen ist), so folgt, wenn der imaginäre Theil von $\log Q(s)$ passend bestimmt wird,

\weq{21}{\log Q(s) = \Sigma \frac{\chi(p)}{p^s} + \frac{1}{2}\Sigma \frac{\chi(p^2)}{p^{2s}} + \frac{1}{3}\Sigma \frac{\chi(p^3)}{p^{3s}} + \cdots}

Wenn nun $A$ irgend eine der Classen (2) und $a$ eine Zahl aus $A$ bedeutet, so ist nach \S.\,11, (6) dieses Bandes

\weq{22}{\overset{\chi}{\Sigma} \chi(A) \chi(n) = \Sigma \chi(an) = \mu \text{ oder } = 0,}

je nachdem $an \equiv 1$ oder nicht $\equiv 1$ nach dem Modul $m$ ist, d.\,h. je nachdem $n$ in der Classe $A^{-1}$ enthalten ist oder nicht enthalten ist. Wenn wir also die Formel (21) mit $\chi(A)$ multipliciren und in Bezug auf $\chi$ summiren, so folgt

\weq{23}{\frac{1}{\mu} \overset{\chi}{\Sigma} \chi(A) \log Q(s) = \Sigma \frac{1}{p^s} + \frac{1}{2}\Sigma \frac{1}{p^{2s}} + \frac{1}{3}\Sigma \frac{1}{p^{3s}} + \cdots,}

wo sich rechts die erste, zweite, dritte, \ldots\ Summe auf alle Primzahlen $p$ bezieht, deren erste, zweite, dritte, \ldots\ Potenz in $A^{-1}$ enthalten ist.

Der zweite Theil dieser Summe
\[
R(s) = \frac{1}{2}\Sigma \frac{1}{p^{2s}} + \frac{1}{3}\Sigma \frac{1}{p^{3s}} + \cdots
\]
wird vergrössert, wenn man jede seiner Theilsummen auf \gesperrt{alle} Primzahlen $p$ erstreckt, und demnach ist
\[
R(s) < \frac{1}{2}\Sigma \left(\frac{1}{p^{2s}} + \frac{1}{p^{3s}} + \cdots\right) = \frac{1}{2}\Sigma \frac{1}{p^s (p^s - 1)},
\]
also um so mehr:
\[
R(s) < \frac{1}{2} \sum_{2,\infty}^{n} \frac{1}{n^{2s} (1 - n^{-s})},
\]
oder, da $1 - n^{-s} > \frac{1}{2}$ ist:
\[
R(s) < \sum_{2,\infty}^{n} \frac{1}{n^{2s}},
\]
und hieraus schliesst man, dass $R$ für $s = 1$ endlich bleibt (\S.\,191, 2.).

Wenn wir nun in der Formel (23),
\[
\frac{1}{\mu} \overset{\chi}{\Sigma} \chi(A) \log Q(s) = \Sigma \frac{1}{p^s} + R(s),
\]
%CONT>>


% ===== scan leaf 805, printed 784 =====
% !PRINTED 784
\weberpage{784}
%<<CONT
$s$ in 1 übergehen lassen, so wird $\log Q_1(s)$ nach 4. unendlich, während die übrigen Glieder der linken Seite, $\log Q_2(s), \ldots \log Q_\mu(s)$, endlich bleiben. Da $R(s)$ gleichfalls endlich bleibt, so muss die Summe $\Sigma p^{-s}$ unendlich werden, und dies ist sicher nur dann möglich, wenn die Summe aus unendlich vielen Gliedern besteht. Hiermit ist bewiesen:

5. \gesperrt{In jeder der Zahlclassen $A$ (nach dem Modul $m$) sind unendlich viele Primzahlen enthalten.}

Man kann dies auch so ausdrücken, dass die Linearform $mx + a$, in der $m$ und $a$ ganze Zahlen ohne gemeinsamen Theiler sind, für unendlich viele ganzzahlige Werthe von $x$ eine Primzahl darstellt.

\begin{center}
\rule{4cm}{0.4pt}
\end{center}

% [The remainder of this leaf and the following leaf are blank in the source.]


% ===== scan leaf 806, printed 785 =====
\weberpage{785}% page number not printed on this page
\pdfbookmark[1]{Register zu Band I und II}{bm240}
% !PRINTED none
% Register (index) opening page: no folio number appears in the running head on
% this page image (title page of the alphabetical index). Footer at the foot of
% the page reads "Weber, Algebra. II." followed by the gathering signature "50".
\begin{center}
\gesperrt{REGISTER.}
\end{center}

\begin{center}
(Die r\"omischen Ziffern bezeichnen den Band, die arabischen Ziffern die Seite.)
\end{center}

\begin{center}
\rule{3cm}{0.5pt}
\end{center}

\noindent
\gesperrt{\textbf{A}bel}'sche Gleichungen I, 533.\\
--- --- cubische II, 107.\\
--- --- biquadratische II, 110.\\
--- Gruppen I, 476, 536; II, 6, 32.\\
--- K\"orper II, 588, 648.\\
Abgeleitete Functionen I, 52.\\
Absolute Norm II, 502, 536.\\
Absoluter Werth einer imagin\"aren Zahl I, 18.\\
--- --- eines Functionals II, 501.\\
Abz\"ahlbare Mengen II, 745.\\
Addition I, 14.\\
Adjunction I, 451.\\
Aehnliche Substitutionen II, 154.\\
Aequivalente Functionale II, 552.\\
--- Zahlen I, 371.\\
--- --- und Functionale im quadratischen K\"orper II, 610.\\
Affect I, 480.\\
Algebraische Functionen II, 526.\\
--- K\"orper I, 455; II, 527.\\
--- Zahlen II, 489.\\
Algorithmus des gr\"ossten gemeinschaftlichen Theilers bei Zahlen I, 2.\\
--- bei ganzen Functionen I, 34.\\
Alternirende Gruppe I, 496, 600.\\
--- --- von f\"unf Ziffern II, 138, 230.\\
--- Function I, 496.\\
Arithmetische Reihen I, 47.\\
\gesperrt{Aronhold}'sche Systeme von Doppeltangenten II, 367.\\
Associirte Functionale II, 510.\\

Associirte Zahlen I, 586.\\
Aufl\"osung von Gleichungen durch Wurzelzeichen I, 595.\\
Axen einer tern\"aren Substitutionsgruppe II, 438.\\

\medskip
\noindent
\textbf{B}asen der Functionale II, 532, 581.\\
--- der Ideale II, 550.\\
Basis einer \gesperrt{Abel}'schen Gruppe II, 33.\\
--- eines algebraischen K\"orpers II, 528.\\
--- des nat\"url. Logarithmensystems II, 752.\\
Basisform eines Functionals II, 535.\\
--- --- --- von $\mathfrak{o}$ II, 535.\\
--- --- --- im quadratischen K\"orper II, 604, 608.\\
Befreiung einer Gleichung vom zweiten Gliede I, 114.\\
\gesperrt{Bernoulli}'sche N\"aherungsmethode der Aufl\"osung von Gleichungen I, 341.\\
\gesperrt{Bertrand}'scher Satz \"uber die Grenzen des Index von Permutationsgruppen II, 143.\\
\gesperrt{Bezout}'sches Theorem I, 161.\\
Bezoutiante I, 225, 265.\\
--- und Wurzelrealit\"at I, 252.\\
Binomischer Lehrsatz I, 42.\\
Binomial-Co\"efficienten I, 42.\\
Bin\"are Formen I, 189.\\
--- lineare Substitution II, 191.\\
Biquadratische \gesperrt{Abel}'sche Gleichungen II, 110.\\

\bigskip
\noindent
%% [page furniture, not body text -- commented out 2026-08-10] Page furniture (running
%% head, folio, signature) is excluded from this edition by standing policy.
%% Weber, Algebra. II. \hfill 50


% ===== scan leaf 807, printed 786 =====
% !PRINTED 786
\weberpage{786}
\noindent
Biquadratische Formen I, 199.\\
--- Gleichungen I, 119, 201, 528; II, 319.\\
--- Kreistheilungsk\"orper II, 101.\\
\gesperrt{Bring-Jerrard}'sche Form der Gleichung f\"unften Grades I, 176, 233.\\
\gesperrt{Brioschi}'sche Normalform der Gleichung f\"unften Grades I, 236.\\
Br\"uche I, 10.\\
Buchstabenrechnung I, 20.\\
\gesperrt{Budan-Fourier}'sches Theorem I, 299.\\

\medskip
\noindent
\textbf{C}anonische Form der tern\"aren cubischen Form II, 332.\\
Cardanische Formel I, 116.\\
Cartesischer (\gesperrt{Harriot}'scher) Lehrsatz \"uber die Zahl der reellen Wurzeln I, 308.\\
Casus irreducibilis der cubischen Gleichungen I, 608.\\
\gesperrt{Cayley}'scher Ausdruck der Cardanischen Formel I, 118.\\
Charaktere einer \gesperrt{Abel}'schen Gruppe II, 44.\\
Charakteristik eines Functionensystems I, 285.\\
Classen von Functionalen II, 553.\\
--- --- Idealen II, 553.\\
--- --- quadratischen Formen I, 400.\\
--- --- quadratischen Irrationalzahlen I, 383.\\
Classenzahl eines algebraischen K\"orpers II, 554.\\
--- im K\"orper der achten Einheitswurzeln II, 717.\\
Classenzahlformel II, 702.\\
--- im Kreistheilungsk\"orper II, 705.\\
--- im reellen Kreistheilungsk\"orper II, 711.\\
--- im K\"orper der $2^{\varkappa\text{ten}}$ Einheitswurzeln II, 719.\\
Classenzahlfactoren II, 716.\\
Classenzahlfactor $A$ II, 722.\\
--- $B$ II, 726.\\
Commutative Gruppen I, 476; II, 6.\\
Commutative Normaltheiler einer metacyklischen Gruppe II, 27.\\
Complexe von Doppeltangenten II, 357.\\

Complexe Wurzeln (ihre Eingrenzung) I, 292.\\
--- Zahlen I, 18.\\
--- --- von \gesperrt{Gauss} I, 585.\\
Complexpaare, Complextripel II, 363.\\
Composition in einer Gruppe II, 3.\\
--- der Theile II, 12.\\
--- quadratischer Irrationalzahlen II, 613.\\
--- von Permutationen I, 474.\\
--- von Substitutionen I, 470.\\
--- linearer Substitutionen II, 153.\\
Compositionsreihe einer Gruppe II, 17.\\
Configuration der tern\"aren Substitutionsgruppe 168\textsuperscript{sten} Grades II, 451.\\
Congruenz der Zahlen I, 358, 368; II, 539.\\
Congruenzen ersten Grades I, 368.\\
--- --- --- in algebraischen K\"orpern II, 542.\\
Congruenzgruppe II, 250.\\
Congruenzk\"orper II, 245.\\
--- zweiten Grades II, 259.\\
Congruenzwurzeln I, 426.\\
Conjugirte Functionen I, 505.\\
--- Gruppen I, 506; II. 10.\\ % [sic] gedruckter Punkt nach der Bandziffer
--- K\"orper I, 460; II, 493.\\
--- Logarithmen II, 672.\\
--- Pole einer linearen Substitutionsgruppe II, 201.\\
Constanz der Indexreihe II, 18.\\
Contragradiente Gruppen II, 480.\\
--- Transformationen II, 157, 478.\\
Contravarianten II, 480.\\
Covarianten I, 187.\\
--- der tern\"aren Formen II, 327.\\
--- der tern\"aren cubischen Formen II, 333.\\
Cubische \gesperrt{Abel}'sche Gleichungen I, 107.\\
--- Formen, bin\"are I, 192.\\
--- --- tern\"are II, 331.\\
--- Gleichungen I, 116, 522.\\
--- --- trigonometrische L\"osung I, 349.\\
--- Kreistheilungsk\"orper II, 93.\\
Curven, algebraische II, 324.\\
Cyklische Functionen I, 531, 538.\\
--- Gleichungen I, 538.\\
--- Gruppen I, 476.\\
--- --- linearer Substitutionen II, 209.\\


% ===== scan leaf 808, printed 787 =====
% !PRINTED 787
\weberpage{787}
\noindent
Cyklische Gruppen im Congruenzk\"orper II, 269, 275.\\
--- Permutationen I, 476, 493.\\

\medskip
\noindent
\gesperrt{\textbf{D}edekind}'sche Ideale II, 547.\\
\gesperrt{Dedekind}'scher Satz \"uber die K\"orperdiscriminante II, 578.\\
Derivirte Functionen I, 51.\\
Derivirte eines Productes I, 54.\\
--- von Functionen mehrerer Ver\"anderlichen I, 59.\\
Determinanten I, 68.\\
--- ger\"anderte I, 81.\\
--- aus Unterdeterminanten I, 97.\\
Determinante eines Systems linearer Gleichungen I, 84.\\
--- einer quadratischen Form I, 180.\\
Dichte Mengen I, 4.\\
Di\"edergruppe II, 210.\\
--- im Congruenzk\"orper II, 276.\\
Differentialquotienten I, 54, 60.\\
Differenzen I, 46.\\
Differenzenproduct der $m^{\text{ten}}$ Einheitswurzeln I, 423.\\
\gesperrt{Dirichlet}'scher Satz \"uber die Einheiten II, 670.\\
Discrete Mengen I, 4.\\
Discriminante I, 150.\\
--- der cubischen Form I, 37, 153, 193.\\
--- der biquadratischen Form I, 155, 200.\\
--- der quadratischen Irrationalzahlen I, 378; II, 602.\\
--- der Kreistheilungsgleichung f\"ur einen Primzahlgrad I, 421.\\
--- einer algebraischen Curve II, 325.\\
--- eines algebraischen K\"orpers II, 529.\\
--- des Kreistheilungsk\"orpers II, 619. 643.\\ % [sic] gedruckter Punkt statt Komma nach 619
Discriminanten in einem algebraischen K\"orper II, 528.\\
Discriminantenfl\"ache I, 249.\\
Division von Zahlen I, 1.\\
--- ganzer Functionen I, 28.\\
Divisoren einer \gesperrt{Abel}'schen Gruppe II, 48.\\
--- einer Gruppe I, 476, 501; II, 7.\\
Drehungen II, 186.\\
Drehungsgruppe II, 189.\\

Dreieckscoordinaten II, 322.\\
% [sic] Im Druck folgen Doppelpunkte und Doppeltangenten erst nach Drehungsgruppe und Dreieckscoordinaten.
Doppelpunkte einer Curve II, 326.\\
Doppeltangenten II, 327.\\
--- einer Curve vierter Ordnung II, 351.\\
Durchschnitt von Gruppen I, 510; II, 9.\\

\medskip
\noindent
\textbf{E}igentliche und uneigentliche Aequivalenz I, 371.\\
Eigentliche und uneigentliche orthogonale Substitutionen II, 195.\\
Einfache Gruppen I, 511; II, 136.\\
--- --- 60\textsuperscript{sten} Grades II, 138.\\
--- \gesperrt{Abel}'sche K\"orper II, 649.\\
Einfachheit der alternirenden Gruppe I, 600.\\
--- --- Congruenzgruppen II, 254.\\
Einheit einer Gruppe II, 5.\\
Einheiten I, 399.\\
--- im K\"orper $R(i)$ I, 586.\\
--- in algebraischen K\"orpern II, 510.\\
--- im reellen Kreistheilungsk\"orper II, 645.\\
--- im K\"orper der achten Einheitswurzeln II, 718.\\
--- unabh\"angige II, 677.\\
--- Fundamentalsystem II, 682.\\
Einheitswurzeln I, 114, 408.\\
--- dritte I, 118.\\
--- primitive I, 411.\\
--- in den Kreistheilungsk\"orpern II, 638.\\
--- im Congruenzk\"orper II, 260.\\
Elimination aus linearen Gleichungen I, 87.\\
--- aus h\"oheren Gleichungen I, 157, 159.\\
--- aus mehreren Gleichungen I, 161.\\
Endgleichung I, 157.\\
Endliche Gruppen II, 4.\\
--- --- linearer Substitutionen II, 195.\\
Endlichkeit des Invariantensystems einer linearen Gruppe II, 169.\\
Entgegengesetzte Elemente einer Gruppe II, 5.\\
\gesperrt{Euler}'s Theorem \"uber homogene Functionen I, 62.\\
Exponentensystem von Einheiten II, 679.\\

% Druckersignatur „50*“ als Seitenelement ausgeschlossen.


% ===== scan leaf 809, printed 788 =====
% !PRINTED 788
\weberpage{788}
\noindent
Exponentensystem von Zahlen II, 685.\\
Exponentialfunctionen II, 752.\\

\medskip
\noindent
\textbf{F}actoren ganzer Functionen I, 39, 103.\\
--- der nat\"urlichen Primzahlen in algebraischen K\"orpern II, 576.\\
\gesperrt{Fermat}'scher Lehrsatz I, 427; II, 55.\\
--- --- im Congruenzk\"orper II, 247.\\
--- --- in algebraischen K\"orpern II, 544.\\
Formen I, 58.\\
Formenproblem der linearen Substitutionsgruppe II, 175.\\
Formensystem der cubischen Form I, 196.\\
--- der biquadratischen Form I, 203, 206.\\
Functionen, ganze I, 23.\\
--- --- Zerlegung in lineare Factoren I, 103.\\
--- --- Zerlegung in Primfunctionen I, 164, 452; II, 495.\\
--- homogene I, 56.\\
--- in einem K\"orper I, 452; II, 520.\\
Functionencongruenzen II, 243.\\
Functionaldeterminanten I, 185.\\
Functionale II, 499.\\
Functionalclassen II, 552.\\
\gesperrt{Frobenius}'sche S\"atze \"uber Gruppen II, 129, 134.\\
Fundamentalsatz der Algebra I, 121, 127, 296.\\
Fundamentalsysteme von Einheiten II, 681.\\
--- von Normaleinheiten II, 731.\\

\medskip
\noindent
\gesperrt{\textbf{G}alois}'sche Gruppe I, 477, 511; II, 588.\\
--- Resolventen I, 467.\\
--- K\"orper I, 465; II, 588.\\
Ganze Functionen von einer Ver\"anderlichen I, 23.\\
--- --- von mehreren Ver\"anderlichen I, 24.\\
--- --- in einem K\"orper II, 520.\\
--- algebraische Zahlen, II, 491.\\ % [sic] Komma vor der Bandangabe wie gedruckt.
--- Functionale II, 504.\\
\gesperrt{Gauss}'sche Summen I, 575.\\
Gebrochene Functionen I, 32.\\
Geordnete Mengen I, 4.\\

Gewicht einer Invariante I, 187.\\
Gitter II, 561.\\
Gleichungen I, 101.\\
--- lineare homogene I, 81.\\
--- --- unhomogene I, 89.\\
--- dritten Grades I, 116, 118, 522.\\
--- vierten Grades I, 119, 201, 528.\\
--- f\"unften Grades I, 233, 621; II, 404.\\
--- siebenten Grades II, 476, 481.\\
Grad einer Gruppe I, 472, 475; II, 4.\\
--- eines Elementes einer Gruppe II, 10.\\
--- einer Permutation I, 500.\\
--- eines Primideals und Primfunctionals II, 516.\\
\gesperrt{Gr\"affe}'sche Methode der gen\"aherten Aufl\"osung einer Gleichung I, 344.\\
Grenzen I, 121.\\
Gr\"osster gemeinschaftlicher Theiler von Zahlen I, 2.\\
--- --- --- von ganzen Functionen I, 34.\\
--- --- --- von Gruppen I, 510.\\
--- --- --- von Functionalen II, 512, 519.\\
Gr\"osster Normaltheiler einer Gruppe II, 17.\\
Grundcurve der Gruppe $G_{168}$ II, 456.\\
Grundformen der Poly\"edergruppen II, 205.\\
--- --- cyklischen Gruppen II, 210.\\
--- --- Di\"edergruppen II, 210.\\
--- --- Tetra\"edergruppe II, 214.\\
--- --- Octa\"edergruppe II, 216, 218.\\
--- --- Ikosa\"edergruppe II, 221, 230.\\
Grundideal II, 579, 593.\\
Grundzahl eines K\"orpers II, 529.\\
Gruppe (allgemein) II, 3.\\
--- der Doppeltangentengleichung II, 380.\\
--- der Tripelgleichungen II, 344.\\
--- des Tetra\"eders II, 212.\\
--- des Octa\"eders II, 216.\\
--- des Ikosa\"eders II, 220.\\
--- einer Gleichung I, 477.\\
--- eines Normalk\"orpers II, 588.\\
--- der Kreistheilungsk\"orper II, 68.\\
--- eines Ideals im Normalk\"orper II, 589.\\
--- der Idealclassen II, 557.\\
--- 168\textsuperscript{sten} Grades II, 282, 433.\\


% ===== scan leaf 810, printed 789 =====
% !PRINTED 789
\weberpage{789}

\noindent
Gruppen, \gesperrt{Abel}'sche I, 476, 536; II, 6.\\
--- endliche II, 4.\\
--- von Permutationen I, 475.\\
--- von Substitutionen I, 472.\\
--- vom Grade $p^{\alpha}$ II, 127.\\
--- vom Grade $p\,q$ II, 141.\\
--- linearer Substitutionen II, 152.\\
--- orthogon. Substitutionen II, 184.\\
--- linearer gebrochener Substitutionen II, 189, 195.\\
Gruppencharaktere II, 44.\\
Gruppentafel II, 115.

\medskip
\noindent
\textbf{H}albmetacyklische Gruppen I, 616.\\
Hauptaxen einer ternären linearen Substitution II, 441.\\
--- der Gruppe $G_{168}$ II, 444.\\
Hauptgleichung fünften Grades I, 175, 230.\\
Hauptreihe II, 24.\\
Hauptunterdeterminanten I, 256.\\
\gesperrt{Hermite}'s Lösung des \gesperrt{Sturm}'schen Problems I, 280.\\
\gesperrt{Hesse}'sche Curve II, 330.\\
--- Determinante I, 186.\\
\gesperrt{Hesse-Cayley}'sche Bezeichnung der Doppeltangenten II, 369.\\
\gesperrt{Hilbert}'scher Satz II, 165.\\
Homogene Functionen I, 56.

\medskip
\noindent
\textbf{I}deale II, 547.\\
Ideale Zahlen nach \gesperrt{Kummer} II, 551.\\
Identische Gruppe I, 496.\\
--- Substitution II, 152.\\
Ikosa\"edergleichung II, 232, 418.\\
Ikosa\"edergruppe II, 220, 279.\\
Ikosa\"ederinvarianten II, 232.\\
Ikosa\"ederresolventen II, 422, 426.\\
Imaginäre Zahlen I, 17.\\
Imaginäre Form der linearen Congruenzgruppe II, 266.\\
Imaginärer Congruenzkörper zweiten Grades II, 259.\\
Imaginäre quadratische Irrationalzahlen I, 379.\\
Imprimitive Gruppen I, 483, 516.\\
--- Körper I, 460, 483.\\
Index einer Invariante II, 163.\\
--- eines Theilers einer Gruppe I, 502; II, 8.

Indexmoduln II, 61.\\
Indexreihe einer Gruppe II, 17.\\
Indices I, 431.\\
--- der Elemente einer Gruppe II, 44.\\
--- nach einer ungeraden Primzahlpotenz II, 54.\\
--- nach einer Potenz von 2 II, 58.\\
--- nach einem zusammengesetzten Modul II, 60.\\
Inflexionspunkte II, 326.\\
--- einer Curve dritter Ordnung II, 331.\\
Interpolation I, 44, 98, 331.\\
Intransitive Gruppen I, 482.\\
--- Normaltheiler I, 520.\\
Invarianten I, 186.\\
--- der binären cubischen Form I, 193.\\
--- --- --- biquadratischen Form I, 200, 206.\\
--- einer \gesperrt{Abel}'schen Gruppe II, 42.\\
--- endlicher Gruppen linearer Substitutionen II, 161.\\
--- absolute und relative II, 162, 177.\\
--- im weiteren Sinne II, 179.\\
--- der ternären Formen II, 328.\\
--- der Curven dritter Ordnung II, 337.\\
--- der Gruppe $G_{168}$ II, 453, 461.\\
Invariantencurven II, 454.\\
Inverse Substitution II, 155.\\
Irrationale Verhältnisse I, 10.\\
--- Zahlen I, 12, 370.\\
Irreducibilität I, 453.\\
--- reiner Gleichungen I, 607.\\
--- der Kreistheilungsgleichung I, 417; II, 768.\\
Isomorphe Gruppen I, 477; II, 6.\\
Isomorphismus (mehrstufiger) II, 15.

\medskip
\noindent
\gesperrt{\textbf{J}acobi}'s Abschätzung der Zahl der Wurzeln zwischen zwei Grenzen I, 311.\\
\gesperrt{Jordan}'scher Satz über zusammengesetzte Gruppen II, 18.

\medskip
\noindent
\textbf{K}ettenbrüche I, 358.\\
--- für äquivalente Zahlen I, 374.\\
Ketten von Hauptunterdeterminanten I, 257.\\
--- \gesperrt{Sturm}'sche I, 272.


% ===== scan leaf 811, printed 790 =====
% !PRINTED 790
\weberpage{790}

\noindent
Kleinstes gemeinschaftliches Vielfaches von Zahlen I, 3.\\
--- --- --- von Functionalen II, 519.\\
--- --- --- von Gruppen II, 29.\\
Körper I, 449; II, 493.\\
Kugelfunctionen I, 273.\\
\gesperrt{Kummer}'sches Theorem über die Resolventen II, 632.\\
Kreistheilungsgleichung I, 422, 554.\\
Kreistheilungskörper II, 67.\\
--- von gegebener Gruppe II, 79.\\
--- vom Primzahlpotenzgrad II, 617.\\
Kreistheilungstheorie I, 408, 554.\\
Krystallographische Gruppen II, 241.

\medskip
\noindent
\gesperrt{\textbf{L}agrange}'scher Satz über Functionen, die die Permutationen einer Gruppe gestatten I, 508.\\
\gesperrt{Laguerre}'sche Sätze über Gleichungen mit nur reellen Wurzeln I, 322.\\
\gesperrt{Legendre}'sches Symbol I, 441.\\
Lineare Congruenzgruppen I, 611; II, 250, 302.\\
--- --- homogene II, 303.\\
--- --- reelle II, 261.\\
--- --- imaginäre Form II, 266.\\
--- --- vom Grade 168 II, 282.\\
--- --- binäre für den Modul 3 II, 305.\\
--- --- ternäre für den Modul 2 II, 306.\\
--- gebrochene Substitutionen II, 189.\\
--- Functionen I, 30.\\
--- Gleichungen I, 81.\\
--- Transformation I, 178.\\
--- Substitutionen I, 178, 371; II, 151.\\
\gesperrt{Lüroth}'scher Satz II, 404.

\medskip
\noindent
\textbf{M}annigfaltigkeit I, 4.\\
Matrix I, 82.\\
Menge I, 4.\\
Messbare Menge I, 8.\\
Metacyklische Gleichungen I, 597; II, 292.\\
--- --- vom Primzahlgrad I, 609.\\
--- --- fünften Grades I, 621, 648.\\
--- --- sechsten Grades II, 296.\\
--- --- achten Grades II, 314.\\
--- --- neunten Grades II, 305.

Metacyklische Gruppen I, 598; II, 27.\\
--- Functionen I, 617, 635.\\
Minimalbasis eines Körpers II, 529.\\
--- des Kreistheilungskörpers II, 620.\\
--- eines quadratischen Körpers II, 603.\\
Minimum I, 122.\\
--- einer quadratischen Form II, 559.\\
--- der Grundzahl eines Körpers II, 569.\\
Modul I, 358.\\
Multiplicative Substitutionen II, 152.\\
Multiplication von Zahlen I, 14.\\
--- von Determinanten I, 92.\\
--- trigonometrischer Functionen I, 433.

\medskip
\noindent
\textbf{N}äherungsbrüche eines Kettenbruches I, 362.\\
Näherungsmethode zur Berechnung von Gleichungswurzeln durch die regula falsi I, 331.\\
--- zur Berechnung von Gleichungswurzeln von \gesperrt{Daniell Bernoulli} I, 341.\\ % [sic] gedruckt \glqq{}Daniell\grqq{} statt \glqq{}Daniel\grqq{}
--- \gesperrt{Gräffe} I, 344.\\
--- \gesperrt{Newton} I, 335.\\
--- durch Kettenbrüche I, 401.\\
Näherungsmethode zur Auflösung trinomischer Gleichungen I, 352.\\
Natürliche Irrationalitäten I, 516.\\
Nebengruppen I, 502; II, 8.\\
Negative Zahlen I, 16.\\
Neunte Einheitswurzeln I, 582.\\
\gesperrt{Newton}'sche Formeln für die Potenzsummen I, 142.\\
\gesperrt{Newton}'sche Regel für die Wurzelabschätzung I, 305.\\
Nichtmetacyklische Gleichungen I, 603.\\
Norm I, 461; II, 494, 498.\\
--- absolute II, 502, 538.\\
--- der \gesperrt{Gauss}'schen complexen Zahlen I, 586.\\
--- eines Ideals II, 551.\\
--- eines Körpers I, 465.\\
--- einer quadratischen Irrationalzahl I, 377.\\
Normaleinheiten im Kreistheilungskörper II, 729.


% ===== scan leaf 812, printed 791 =====
% !PRINTED 791
\weberpage{791}

\noindent
Normalformen der linearen Substitutionsgruppen II, 181.\\
Normalgleichung I, 465.\\
Normalkörper I, 464; II, 588.\\
Normaltheiler einer Gruppe I, 511; II, 11.\\
Null I, 16.

\medskip
\noindent
\textbf{O}bere Grenze für die Wurzeln I, 316.\\
Octa\"edergruppe II. 216.\\ % [sic] gedruckter Punkt nach der Bandziffer
--- im Congruenzkörper II, 277.\\
Orthogonale Substitution II, 184.

\medskip
\noindent
\textbf{P}artialdiscriminante II, 597.\\
Partialgrundideal II, 585, 597.\\
Partialnorm II, 584.\\
Partialspur II, 585.\\
Partialresolventen I, 512.\\
\gesperrt{Pell}'sche Gleichung I, 395.\\
Periode einer Permutation I, 500.\\
Perioden der Kreistheilung I, 557, 579; II, 74.\\
--- der Wurzeln einer cyklischen Gleichung I, 545.\\
--- der reducirten Zahlen I, 388.\\
Periodische Kettenbrüche I, 387.\\
Permutationen I, 64, 473, 500.\\
--- ihre analytische Darstellung I, 613; II, 299.\\
--- erster und zweiter Art I, 65, 496.\\
Permutationsgruppen I, 475, 489; II, 117.\\
--- $168^{\mathrm{sten}}$ Grades von sieben Ziffern II, 473.\\
Polaren I, 188.\\
Pole linearer gebrochener Substitutionen II, 196.\\
--- einer Gruppe II, 198.\\
--- --- --- im Congruenzkörper II, 274.\\
--- einer ternären Gruppe II, 438.\\
Poly\"edergruppen II, 202.\\
--- der zweiten Art II, 234.\\
--- im Congruenzkörper II, 273.\\
Polynomialco\"efficienten I, 51.\\
Polynomischer Lehrsatz I, 50.\\
Positive Einheiten im Kreistheilungskörper II, 742.\\
Potenzreste I, 430.\\
Potenzsummen I, 140.

Primäre Theiler eines Kreistheilungskörpers II, 71.\\
--- --- einer \gesperrt{Abel}'schen Gruppe II, 72.\\
Primfactoren der Zahlen eines algebraischen Körpers II, 523.\\
--- der natürlichen Primzahlen II, 576.\\
--- der Kreistheilungsresolventen II, 635.\\
Primfunctionale II, 515, 576.\\
--- relative II, 514.\\
Primideale II, 550.\\
--- im quadratischen Körper II, 604.\\
--- im relativ normalen Körper II, 588.\\
--- in den Theilern eines Normalkörpers II, 599.\\
--- in den Kreistheilungskörpern II, 619, 622.\\
--- im reellen Kreistheilungskörper II, 643.\\
Primitive Congruenzwurzeln I, 428; II, 546.\\
--- Einheitswurzeln I, 410.\\
--- Functionen I, 25; II, 497.\\
--- Wurzeln von Primzahlen I, 429.\\
--- --- von Primzahlquadraten II, 57.\\
--- --- von Primfunctionalen II, 546.\\
--- --- von Primidealen II, 591.\\
--- --- eines Congruenzkörpers II, 248.\\
--- und imprimitive Charaktere in Kreistheilungskörpern II, 720.\\
--- --- --- Gruppen I, 484.\\
--- --- --- Körper I, 463.\\
--- --- --- Formenprobleme II, 180.\\
Primzahlen I, 1.\\
--- relative I, 359.\\
--- im Körper $R(i)$ I, 587.\\
--- im Körper der dritten Einheitswurzeln I, 593.

\medskip
\noindent
\textbf{Q}uadrate im Congruenzkörper II, 248.\\
Quadratische Formen I, 179.\\
--- Gleichungen I, 116.\\
--- Irrationalzahlen I, 377.\\
--- Reste I, 444.\\
--- Körper II, 601.\\
Quadratur des Kreises II, 760.

\medskip
\noindent
\textbf{R}ationale Wurzeln einer Gleichung I, 403.


% ===== scan leaf 813, printed 792 =====
% !PRINTED 792
\weberpage{792}

\noindent
Rationale Zahlen I, 12.\\
Rationales Verhältniss I, 10.\\
Rationalitätsbereich I, 452.\\
Raum von $n$ Dimensionen II, 560.\\
Realität der Wurzeln I, 241.\\
--- --- --- bei quadratischen und cubischen Gleichungen I, 243.\\
--- --- --- bei biquadratischen Gleichungen I, 246.\\
--- --- --- bei cyklischen Gleichungen I, 551.\\
--- --- --- bei metacykl. Gleichungen I, 647.\\
--- --- --- der Doppeltangenten einer Curve vierter Ordnung II, 391.\\
Realitätsbedingungen der orthogonalen Gruppe II, 193.\\
Realitätsverhältnisse der Wendepunkte einer Curve dritter Ordnung II, 340.\\
--- bei Tripelgleichungen II, 348.\\
Rechenoperationen I, 1.\\
Rechnen mit ganzen Functionen I, 24.\\
--- mit Zahlen I, 14.\\
Reciprocitätsgesetz der quadratischen Reste I, 443.\\
Reducible und irreducible Functionen I, 453.\\
--- --- --- nach einem Primzahlmodul II, 243.\\
--- --- --- Gleichungen I, 404, 482.\\
Reducirte Zahlen I, 380, 384.\\
--- Einheiten II, 680.\\
Reelle Kreistheilungskörper II, 641.\\
--- Radicale I, 606.\\
Regula falsi I, 334.\\
Reguläre Körper II, 188.\\
Regulator eines Systemes von Einheiten II, 678, 729.\\
--- des Körpers II, 682.\\
Reine Gleichungen I, 109.\\
Reihen II, 694.\\
Relative Primzahlen I, 2.\\
--- Primfunctionale II, 514.\\
Relativnorm II, 584.\\
Relativspur II, 585.\\
Resolventen I, 511.\\
--- der biquadratischen Gleichung I, 120.

Resolventen der Compositionsreihe II, 289.\\
--- der Gleichungen siebenten Grades II, 476.\\
--- der Gleichungen achten Grades II, 308.\\
--- der Gruppe $G_{168}$ II, 466, 473.\\
--- der Ikosa\"edergleichung II, 419.\\
--- in der Kreistheilung I, 564; II, 63, 651.\\
--- von \gesperrt{Lagrange} I, 542.\\
--- --- bei metacyklisch. Gleichungen I, 631.\\
--- mit einem Parameter II, 404.\\
Rest der Division von ganzen Functionen I, 28.\\
Restsystem I, 359; II, 541.\\
Resultanten I, 156.\\
Resultante zweier quadratischer Functionen I, 36, 158.\\
\gesperrt{Rolle}'s Theorem über die Anzahl der reellen Wurzeln I, 319.

\medskip
\noindent
\textbf{S}äculargleichung I, 276.\\
Schlusszahlen eines Kettenbruches I, 361.\\
Schnitt I, 5.\\
Siebener-Systeme von Doppeltangenten II, 367.\\
Siebenzehneck I, 568.\\
Singuläre Punkte einer Curve II, 325.\\
Spur I, 461; II, 494.\\
\gesperrt{Steiner}'sche Complexe von Doppeltangenten II, 357.\\
Stetigkeit I, 5.\\
--- ganzer Functionen I, 105.\\
--- der Wurzeln I, 132.\\
\gesperrt{Sturm}'sche Functionen I, 279.\\
--- Ketten I, 271.\\
\gesperrt{Sturm}'sches Problem I, 270.\\
Substitution, lineare I, 92, 372; II, 151.\\
--- der Verhältnisse II, 158.\\
Substitutionen eines Normalkörpers I, 468; II, 588.\\
Substitutionsdeterminante I, 92, 178.\\
\gesperrt{Sylow}'sche Sätze II, 121, 125.\\
Symmetrische Determinanten I, 69.\\
--- Functionen I, 138, 144, 147.\\
--- Grundfunctionen I, 139.\\
--- Gruppe I, 492.


% ===== scan leaf 814, printed 793 =====
% !PRINTED 793
\weberpage{793}
\begin{flushleft}
Symmetrische Gruppe, ihre Normaltheiler I, 602. \\
Syzygetische und azygetische Systeme von Doppeltangenten II, 362. \\
— — — Complexe von Doppeltangenten II, 366. \\[6pt]
\textbf{T}angenten einer Curve II, 326. \\
\gesperrt{Taylor}'sche Entwickelung I, 59. \\
Ternäre Formen II, 322. \\
— lineare Gruppen vom Grade 168 II, 433. \\
— — Congruenzgruppe v. Grade 168 II, 475. \\
Tetra\"edergruppe II, 212. \\
— im Congruenzk\"orper II, 276. \\
Theilbarkeit ganzer Functionen I, 32. \\
— ganzer Functionale II, 509. \\
— ganzer Zahlen II, 510. \\
Theiler von ganzen Zahlen I, 1. \\
— von ganzen Zahlen, gr\"osster gemeinschaftlicher I, 359. \\
— ganzer Functionen I, 27; II, 497. \\
— ganzer Functionen, gr\"osster gemeinschaftlicher I, 34. \\
— einer Gruppe I, 476, 501; II, 7. \\
— der Ikosa\"edergruppe II, 227. \\
— eines K\"orpers I, 450. \\
— der linearen Congruenzgruppe II, 271. \\
Theilerfremde Zahlen I, 2, 359. \\
— Functionen I, 35. \\
— Functionale II, 514. \\
Theilnenner eines Kettenbruches I, 361. \\
Theilung des Winkels I, 435. \\
Totalresolventen I, 512. \\
Tr\"agheitsgesetz quadratischer Formen I, 183, 255. \\
Transcendente Functionen zur L\"osung der Ikosa\"edergleichung II, 429. \\
— Zahlen II, 745. \\
Transcendenz der Zahl $e$ II, 751. \\
— der Zahl $\pi$ II, 756. \\
Transformation einer Gruppe I, 506. \\
— der quadratischen Form in eine Summe von Quadraten I, 181. \\
— von Formen $n^{\text{ten}}$ Grades I, 185. \\
— der cubischen Gleichung I, 219. \\
Transformation der Gleichung f\"unften Grades I, 230. \\
Transformirte Substitutionen II, 156. \\
Transitive und intransitive Gruppen I, 482, 506. \\
— Permutationsgruppe vom Index 6 von sechs Ziffern I, 628. \\
Transponirte Substitutionen II, 156. \\
Transpositionen I, 65, 492. \\
Trigonometrische L\"osung reiner Gleichungen I, 111. \\
— — cubischer Gleichungen I, 349. \\
Trinomische Gleichungen I, 352. \\
Tripelgleichungen II, 342. \\
Tripelsysteme II, 310. \\
\gesperrt{Tschirnhausen}-Transformation I, 170, 210. \\
— — der cubischen Gleichung I, 219. \\[6pt]
\textbf{U}mkehrbare Perioden I, 391. \\
Unbekannte I, 20. \\
Unbestimmte Gleichungen I, 365. \\
Unendliche Wurzeln einer Gleichung I, 136. \\
Unendlichkeit ganzer rationaler Functionen I, 104. \\
Unicursalcurven II, 416. \\
Unterdeterminanten I, 72. \\
— h\"ohere I, 77. \\
— complement\"are I, 79. \\
Unz\"ahlbare Mengen II, 748. \\
Urspr\"ungliche Functionen I, 25. \\[6pt]
\textbf{V}ariable in der K\"orpertheorie II, 499. \\
Verh\"altnisse I, 10. \\
Verwandte Zahlen und Functionale II, 631. \\
Volles Restsystem I, 359; II, 541. \\
Vollst\"andige Systeme von Doppeltangenten II, 367. \\
Volumen II, 561, 690. \\
Vorzeichenwechsel ganzer Functionen I, 107. \\[6pt]
\textbf{W}endepunkte II, 326. \\
Wendetangenten II, 326. \\
— einer Curve dritter Ordnung II, 331. \\
\gesperrt{Wilson}'scher Lehrsatz I, 428. \\
Winkeltheilung I, 551, 609.
\end{flushleft}


% ===== scan leaf 815, printed 794 =====
% !PRINTED 794
\weberpage{794}
\begin{flushleft}
W\"urfelverdoppelung I, 609. \\
Wurzeln von Gleichungen I, 101. \\
— von Gleichungen ungeraden Grades I, 109. \\
— von reinen Gleichnngen I, 109. % [sic] printed \textquotedbl{}Gleichnngen\textquotedbl{} for \textquotedbl{}Gleichungen\textquotedbl{}
\\
— von metacyklischen Gleichungen I, 638. \\
— rationale I, 403. \\[6pt]
\textbf{Z}ahl I, 1, 12. \\
— der Glieder einer homogenen Function I, 58. \\
— der Permutationen I, 64. \\
Zahlclassen nach einem zusammengesetzten Modul II, 60.
\\
Zahlenreihen I, 15. \\
Zahlk\"orper I, 449; II, 527.
\end{flushleft}

\begin{flushleft}
Zeichenwechsel und Zeichenfolge I, 262. \\
Zerlegbare und unzerlegbare Functionen mehrerer Variablen I, 164. \\
Zerlegung ganzer Functioaen in lineare Factoren I, 102. % [sic] printed \textquotedbl{}Functioaen\textquotedbl{} for \textquotedbl{}Functionen\textquotedbl{}
\\
— von Gruppen in Nebengruppen I, 502; II, 8. \\
— — — nach zwei Theilern II, 120. \\
Zusammengesetzte Abel'sche K\"orper II, 649. \\
Zusammensetzung linearer Substitutionen I, 372; II, 153. \\
— von Substitutionen eines Normalk\"orpers I, 470. \\
— von Permutationen I, 473. \\
Zweiseitige Zahlen I, 390.
\end{flushleft}


% ===== scan leaf 816, printed 795 =====
\weberpage{795}% page number not printed on this page
\pdfbookmark[1]{Berichtigungen}{bm241}
% !PRINTED none
\weberpage{none}
\begin{center}
\textbf{\gesperrt{Berichtigungen zum ersten Bande.}}
\end{center}

\begin{center}
\rule{4cm}{0.4pt}
\end{center}

\begin{flushleft}
Seite 93, Zeile 6 v.\,u. lies $B_h^{(i)}$ statt $B^{(i)}$. \\
„ 117, Formel (10) lies $\sqrt[3]{\dfrac{-b}{2}+\sqrt{R}}$ statt $\sqrt{\dfrac{-b}{2}+\sqrt{R}}$. \\
„ 158, Formeln (14) rechts $m$ und $n$ zu vertauschen. \\
„ 172, Formel (14) lies $n-1$ statt $n+1$. \\
„ 185, Zeile 9 v.\,o. „ $\Phi'$ statt $\Phi$. \\
„ 190, „ 11 v.\,u. „ $a_1'$ „ $(n-1)a_1'$. \\
„ 194, „ 7 v.\,u. „ $+9$ „ $-9$. \\
„ 200, „ 7 v.\,u. „ \S.\,61 „ \S.\,67. \\
„ 202, „ 3 v.\,u. „ $u$ „ $v_1$ (in beiden Formeln). \\
„ 213, „ 6 v.\,u. „ $\Phi(\tau,\xi)$ statt $\Phi(\tau,k)$. \\
„ 221, „ 12 v.\,u. „ $\dfrac{2}{9}$ statt $\dfrac{2}{3}$. \\
„ 222, „ 4,5 v.\,u. „ $+q_0, -q_0$ statt $-q_0, +q_0$. \\
„ 233, „ 4,7 v.\,o. „ $\tfrac{1}{5}D$ statt $5D$. \\
„ 236, „ 12 v.\,o. „ $5c^2$ „ $-5c^2$. \\
„ 251, „ 10 v.\,o. „ $-\psi_3$ „ $\psi_3$. \\
„ 268, „ 5 v.\,u. „ $\tfrac{1}{5}D$ „ $D$. \\
„ 282, Formel (2) „ $f_s$ „ $f^s$. \\
„ 318, Zeile 2 v.\,u. „ $F(z)$ statt $F(x)$. \\
„ 332, Formel (5) „ $f(\beta)-$ statt $f(\beta)+$. \\
„ 334, Zeile 6 v.\,u. „ 0,00164 „ 0,0164. \\
„ 344, „ 6 v.\,u., Formel, lies $\alpha_n\varphi(\alpha_n)$ und $\varphi(\alpha_n)$ statt $\alpha_m\varphi(\alpha_m)$
in $\varphi(\alpha_m)$. \\
„ 357, „ 3 v.\,u. lies 6 statt 5. \\
„ 378, 379, Formel (13), (14) ist $a_1, b_1, c_1$ durch $\pm a_1, \pm b_1, \pm c_1$ zu
ersetzen, und das Zeichen so zu bestimmen, dass $c_1$ positiv wird. \\
„ 402, Zeile 6,7 v.\,u. ist $a_4=1$ und $a_3=2$ zu vertauschen. \\
„ — „ 4 v.\,u. lies 9083 statt 9183. \\
„ 419, Hierzu vergleiche man den I. Nachtrag zum zweiten Bande. \\
„ 423, Zeile 15 v.\,o. lies $\dfrac{(n-1)(n-2)}{2}$ statt $\dfrac{n\cdot n-1}{2}$. \\
„ — „ 24 v.\,o. ist zu lesen
$\dfrac{1}{2}\left(\dfrac{(n-1)(n-2)}{2}-\dfrac{n-1}{2}\right)
=\dfrac{(n-1)(n-3)}{4}\equiv 0\ (\mathrm{mod}\ 2)$.
\end{flushleft}


% ===== scan leaf 817, printed 796 =====
% !PRINTED 796
\weberpage{796}
\begin{flushleft}
Seite 423, Zeile 3 v.\,u. lies $(-1)^{\frac{(n-1)(n-3)}{4}}=+1$ statt $(-1)^{\frac{n-1}{2}}$. \\
„ 424, Formel (14) „ $P=i^{\frac{n-1}{2}}n^{\frac{n-3}{2}}\sqrt{n}$. \\
„ 436, Zeile 3, 22 v.\,o. lies $\dfrac{\pi}{n}$ statt $\pi$. \\
„ 482, „ 8 v.\,u. ist transitiv und intransitiv zu vertauschen. \\
„ 502, „ 18 v.\,o. lies $\pi_1$ statt $\pi$. \\
„ 504, „ 9 v.\,u. „ $k_{q-1}$ statt $k_{\nu-1}$. \\
„ 514, „ 5 v.\,u. „ $\Theta_1$ statt $\Theta_i$. \\
„ 531, „ 11 v.\,o. „ $y_1^2$ „ $y^2$. \\
„ 540, Formel (10) „ $\gamma_{e-1}$ statt $\gamma_{c-1}$. \\
„ 544, Zeile 15 v.\,o. „ $mB_k=\overset{\varepsilon}{\Sigma}\varepsilon^{-k}(\varepsilon,\alpha)^\nu
(\varepsilon^{\lambda_1},\alpha)^{\nu_1}(\varepsilon^{\lambda_2},\alpha)^{\nu_2}\ldots$ \\
„ — „ 23 v.\,o. „ $\varepsilon^{\lambda_2}$ statt $\varepsilon^2$. \\
„ 551, „ 15 v.\,o. „ $\alpha_k$ „ $\alpha^k$. \\
„ 592, „ 8 v.\,u. „ die nur erf\"ullt ist f\"ur $b=0$, $a=\pm 1$ oder
$a=0$, $b=\pm 1$ oder $a=b=\pm 1$. \\
„ 612, „ 12 v.\,o. „ $1-a$ statt $a-b$. \\
„ 618, „ 1 v.\,u. „ $c_1^h$ statt $c_1$. \\
„ 623, „ 6 v.\,u. „ $a_6$ „ $\alpha_6$.
\\
„ 636, „ 18, 21 v.\,o. lies $-g^{n-1}+1$ statt $g^{n-1}-1$. \\
„ 639, „ 10, 27 v.\,o. „ $n-2$, $n-3$ statt $n-1$, $n-2$. \\
„ 641, „ 1 v.\,o. lies (6) statt (5). \\
„ 646, „ 11 v.\,o. „ rational in $\mathfrak{K}(f_0)$ statt rational. \\
„ — „ 14, 19, 22, 26 v.\,o. lies $\mathfrak{K}(f_0)$ statt $\mathfrak{K}$. \\
„ — „ 17, 18 v.\,o. von \glqq{}Alle Gr\"ossen\grqq{} an zu streichen. \\
„ — „ 1, 2 v.\,u. „ \glqq{}die anderen\grqq{} an zu streichen. \\
„ 653, „ 4 v.\,o. lies $-A_4 r\varrho$ statt $+A_4 r\varrho$. \\
„ — „ 5 v.\,o. „ $+A_2\varrho'$ „ $-A_2\varrho'$.
\end{flushleft}

\begin{center}
\rule{6em}{0.4pt}
\end{center}

\begin{center}
\textbf{\gesperrt{Berichtigungen zum zweiten Bande.}}
\end{center}

\begin{center}
\rule{6em}{0.4pt}
\end{center}

\begin{flushleft}
Seite 15, Zeile 19 v.\,o. lies $N$ statt $P$. \\
„ 16, „ 12, 13 v.\,o. sind $P$ und $Q$ zu vertauschen. \\
„ 21, „ 19 v.\,o. lies (17) statt (13). \\
„ 26, „ 17 v.\,o. „ \S.\,6, 4 statt \S.\,6, 2. \\
„ 569, „ 5 v.\,u. im Exponenten lies $n-\nu$ statt $\mu-\nu$. \\
„ 592, „ 6, 3 v.\,u. lies $\gamma^{P^{f-1}}$ statt $\gamma^{P^f-1}$.
\end{flushleft}

\begin{center}
\rule{6em}{0.4pt}
\end{center}


% ===== scan leaf 818, printed None =====
% !PRINTED none
% Verlagsanzeige von Friedrich Vieweg & Sohn in Braunschweig; ohne Seitenzahl. Auf den
% Verlagskopf und einen kurzen Mittelstrich folgen fünf vollständige, durch kurze
% Mittelstriche getrennte Anzeigen:
% 1. „Elliptische Functionen und algebraische Zahlen. Akademische Vorlesungen“ von H. Weber,
% Professor der Mathematik an der Universität Straßburg; gr. 8.; geh.; Preis 13 Mark.
% 2. „Elektrodynamik mit Berücksichtigung der Thermoelektricität, der Elektrolyse und der
% Thermochemie“ von Dr. Heinrich Weber, Professor an der Herzoglichen Technischen Hochschule
% zu Braunschweig; mit Holzstichen; gr. 8.; geh.; Preis 6 Mark.
% 3. „Fünf populäre wissenschaftliche Vorträge“, gehalten in der Aula der Herzoglichen
% Technischen Hochschule zu Braunschweig, von Dr. Heinrich Weber, Professor der Physik an der
% Herzogl. Technischen Hochschule; mit 84 Illustrationen; 8.; geh.; Preis 2 Mark 50 Pfennig.
% Inhalt: Übersicht über die Entstehung und die Wirkungen des galvanischen Stromes;
% Telegraphie und Telephonie; elektromagnetische, magnetelektrische und Dynamo-Maschinen;
% Galvanoplastik und elektrisches Licht; über das Perpetuum mobile.
% 4. „Stetigkeit und irrationale Zahlen“ von Richard Dedekind, Professor der Mathematik an
% der Technischen Hochschule zu Braunschweig; zweite unveränderte Auflage; gr. 8.; geh.;
% Preis 1 Mark. Die Editionszeile ist als selbständige Versalzeile gesetzt.
% 5. „Lehrbuch der Differential-Gleichungen“ von Dr. Andrew Russell Forsyth, Professor am
% Trinity College zu Cambridge; mit einem Anhange: „Die Resultate der im Lehrbuche
% angeführten Übungsaufgaben enthaltend“; herausgegeben von H. Maser; autorisirte
% Uebersetzung; gr. 8.; geh.; Preis 14 Mark.


% ===== scan leaf 819, printed None =====
% !PRINTED none
% Verlagsanzeige von Friedrich Vieweg & Sohn in Braunschweig; ohne Seitenzahl. Auf den
% Verlagskopf und einen kurzen Mittelstrich folgen sechs vollständige, durch kurze
% Mittelstriche getrennte Anzeigen:
% 1. „Was sind und was sollen die Zahlen?“ von Richard Dedekind, Professor der Mathematik an
% der technischen Hochschule zu Braunschweig; zweite unveränderte Auflage; 8.; geh.;
% Preis 1 Mark 60 Pfennig.
% 2. „Vorlesungen über Zahlentheorie“ von P. G. Lejeune-Dirichlet; herausgegeben und mit
% Zusätzen versehen von R. Dedekind, Professor an der technischen Hochschule
% Carolo-Wilhelmina zu Braunschweig; vierte umgearbeitete und vermehrte Auflage; gr. 8.;
% geh.; Preis 14 Mark.
% 3. „Lehrbuch der politischen Arithmetik“, für höhere Handelsschulen (Handelsakademien)
% und zum Selbstunterricht bearbeitet von F. S. Holzinger, Professor an der öffentlichen
% Handelsakademie in Linz; gr. 8.; geh.; Preis 3 Mark 50 Pfennig.
% 4. „Anleitung zur Durchmusterung des Himmels. Astronomische Objecte für gewöhnliche
% Teleskope. Ein Hand- und Hülfsbuch für alle Freunde der Himmelskunde, besonders für die
% Besitzer von Fernrohren“ von Dr. Hermann J. Klein; zweite verbesserte Auflage; mit 75
% Holzstichen, 5 Tafeln, zum Theil in Farbendruck, 4 Sternkarten und 1 Titelbilde; 8.; geh.;
% Preis 24 Mark.
% 5. „Sammlung von Formeln der reinen und angewandten Mathematik“ von Dr. W. Láska; mit
% drei Tafeln; gr. 8.; geh. 26 Mark; in Halbfranz geb. 28 Mark.
% 6. „Die Beobachtung der Sterne sonst und jetzt“ von J. Norman Lockyer, Mitglied der Royal
% Society, corr. Mitglied des Instituts von Frankreich; autorisirte deutsche Ausgabe;
% Uebersetzt von G. Siebert; mit 217 Holzstichen; 8.; geh.; Preis 18 Mark.


% ===== scan leaf 820, printed None =====
% !PRINTED none
% Verlagsanzeige von Friedrich Vieweg & Sohn in Braunschweig; ohne Seitenzahl. Auf den
% Verlagskopf und einen kurzen Mittelstrich folgen fünf vollständige, durch kurze
% Mittelstriche getrennte Anzeigen:
% 1. „Praktische Anleitung zur Anstellung astronomischer Beobachtungen mit besonderer
% Rücksicht auf die Astrophysik. Nebst einer modernen Instrumentenkunde“ von Nicolaus von
% Konkoly, Dr. phil., Mitglied der Akademie der Wissenschaften in Budapest, der Royal
% Astronomical Society in London etc.; mit 345 Holzstichen; gr. 8.; geh.; Preis 24 Mark.
% 2. „Vorträge und Reden“ von Hermann von Helmholtz; vierte Auflage; mit dem Bildniss des
% Verfassers und zahlreichen eingedruckten Holzstichen; zwei Bände; gr. 8.; pro Band geh.
% 8 Mark, geb. 9 Mark 50 Pfennig.
% 3. „Der Mond und die Beschaffenheit und Gestaltung seiner Oberfläche“ von Edmund Neison,
% Mitglied der Königlichen astronomischen Gesellschaft zu London etc.; autorisirte deutsche
% Original-Ausgabe; zweite Auflage, vermehrt mit einem Anhange, enthaltend die Untersuchungen
% des Verfassers über die Neubildung Hyginus N auf dem Monde; nebst einem Atlas in 4° von
% 26 Karten und 5 Tafeln in Farbendruck; gr. 8.; geh.; mit Atlas zusammen 18 Mark.
% 4. „Photogrammetrie und internationale Wolkenmessung“ von Dr. Carl Koppe, Professor an
% der Herzogl. technischen Hochschule zu Braunschweig; mit Abbildungen und fünf Tafeln;
% gr. 8.; geh.; Preis 7 Mark.
% 5. „Fünfstellige logarithmische und trigonometrische Tafeln“, herausgegeben von
% Dr. O. Schlömilch, K. S. Geheimerath a. D., Mitglied der Königl. Sächsischen Gesellschaft
% der Wissenschaften zu Leipzig, der Königl. Schwedischen Akademie zu Stockholm, der
% Kaiserl. Leopoldinischen Akademie etc.; galvanoplastische Stereotypie; wohlfeile
% Schulausgabe; zwölfte Auflage; 8.; geh.; Preis 1 Mark. Die Anzeige endet auf dieser Seite.


% ===== scan leaf 821, printed None =====
% !PRINTED none
% Verlagsanzeige von Friedrich Vieweg & Sohn in Braunschweig; ohne Seitenzahl. Auf den
% Verlagskopf und einen kurzen Mittelstrich folgen fünf vollständige Anzeigen:
% 1. „Partielle Differentialgleichungen und deren Anwendung auf physikalische Fragen.
% Vorlesungen“ von Bernhard Riemann; für den Druck bearbeitet und herausgegeben von Karl
% Hattendorff; dritte Auflage; mit Holzstichen; gr. 8.; geh.; Preis 8 Mark.
% 2. „Die polydimensionalen Grössen und die vollkommenen Primzahlen“ von Dr. Hermann
% Scheffler, Mitglied der Kaiserlichen Societät der Naturforscher zu Moskau,
% Korrespondirendem Mitglied der Königlichen Akademie der Wissenschaften zu Padua; mit
% 42 Holzstichen; gr. 8.; geh.; Preis 5 Mark 60 Pfennig.
% 3. „Compendium der höheren Analysis“ von Dr. Oskar Schlömilch, K. S. Geheimrath a. D.,
% Mitglied der Königl. Sächsischen Gesellschaft der Wissenschaften zu Leipzig, der
% Königlich Schwedischen Akademie zu Stockholm, der Kaiserlich Leopoldinischen Akademie
% etc.; in zwei Bänden; mit Holzstichen; gr. 8.; geh.; erster Band, fünfte verbesserte
% Auflage, 9 Mark; zweiter Band, vierte Auflage, 9 Mark.
% 4. „Das Buch der Natur, die Lehren der Physik, Astronomie, Chemie, Mineralogie, Geologie,
% Botanik, Zoologie und Physiologie umfassend“, allen Freunden der Naturwissenschaft,
% insbesondere den Gymnasien, Realschulen und höheren Bürgerschulen gewidmet, von
% Dr. Friedrich Schoedler, Director der Grossherzogl. Hessischen Realschule I. O. in Mainz.
% Erster Theil: Physik, Astronomie und Chemie; zweiundzwanzigste verbesserte Auflage mit dem
% Portrait des Verfassers; 404 Holzstiche, Spectraltafel in Farbendruck, Sternkarten und
% Mondkarte; gr. 8.; geh.; 4 Mark 80 Pfennig. Zweiter Theil: Mineralogie, Geologie, Botanik,
% Zoologie und Physiologie; zweiundzwanzigste verbesserte Auflage; 683 Holzstiche und eine
% geognostische Tafel in Farbendruck; gr. 8.; geh.; 4 Mark 80 Pfennig. Der Trennstrich zur
% folgenden Anzeige steht in derselben Zeile wie der Preis des zweiten Theils.
% 5. „Die Fundamentaltheorieen der neueren Geometrie und die Elemente der Lehre von den
% Kegelschnitten für den Schulunterricht bearbeitet“ von H. Seeger; mit 60 Holzstichen;
% gr. 8.; geh.; Preis 2 Mark 80 Pfennig. Die Anzeige endet mit anschließendem Leerraum.


% ===== scan leaf 822, printed None =====
% !PRINTED none
% Scan shows only faint speckling/noise on all three bands; no printed text, no running head, no legible content. Appears to be a blank or near-blank leaf.


% ===== scan leaf 823, printed None =====
% !PRINTED none
% A library "NON-CIRCULATING BOOK" ownership stamp lies across the upper part of this otherwise blank leaf. Faint verso show-through is visible but not legible. A pencil call-number annotation, best read as "QH 11-80", appears near the lower-left corner. No printed page number or Weber body text is present.


% ===== scan leaf 824, printed None =====
% !PRINTED none
% The leaf is blank apart from faint speckling, tonal variation, and the dark scanned gutter at the left edge. No handwriting, printed text, running head, page number, or Weber body content is present.


% ===== scan leaf 825, printed None =====
% !PRINTED none
% This scan is the library's rear pastedown and due-date slip, set against a marbled endpaper. The heading reads "RETURN TO -> Astronomy/Mathematics/Statistics/Computer Science Library, 100 Evans Hall, 642-3381". In the six-cell loan-period table, the first cell's "7 DAYS" is struck out and the second is stamped "1 MONTH". Below "ALL BOOKS MAY BE RECALLED AFTER 7 DAYS" and "DUE AS STAMPED BELOW" are a cancelled "MAR - 5 1988", heavily overprinted and only partly legible 1989 stamps, a cancelled "OCT 14 1998", and "OCT 05 2002". The footer reads "FORM NO. DD3, 1/83" and "UNIVERSITY OF CALIFORNIA, BERKELEY / BERKELEY, CA 94720", with a small circled-P device. No Weber text is present.

\end{document}
